% Auto-generated from data/segments.json + layout.json (+ transforms) — do not edit by hand.
% volume: 24Khu07
% mode: reading
% segments: 2820
% transform_rules: 0
% toc_marks: 36

% Volume layout defaults (layout.json ``layout``).
\csromanlayoutapply{1}{1.5}{21.6pt}{5pt}{6.3pt}{65pt}{2.5em}%

% Reading physical-page layout (layout.json ``page_layout_reading_mode``).
\csromanreadingpagelayoutvolume{1}{1.5}{21.6pt}{5pt}{6.3pt}{65pt}{2.5em}%
\csromanreadingpagelayoutenable

\csromanheader{2}{\textbf{ขุทฺทกนิกาย}}
\csromanmatikaanchor{mtk.0}
\csromantocmarknopage{chapter}{มหานิทฺเทสปาฬิ}{mtk.0}
\bookmark[dest={mtk.0},level=0]{มหานิทฺเทสปาฬิ}
\gambhira{\textbf{มหานิทฺเทสปาฬิ}}
\namakkaram{นโม ตสฺส ภควโต อรหโต สมฺมาสมฺพุทฺธสฺส.}
\csromanmatikaanchor{mtk.1}
\csromantocmarknopage{section}{อฏฺฐกวคฺค}{mtk.1}
\bookmark[dest={mtk.1},level=1]{อฏฺฐกวคฺค}
\csromanheader{1}{1. อฏฺฐกวคฺค}
\csromanmatikaanchor{mtk.2}
\csromantocmarkref{subsection}{1. กามสุตฺตนิทฺเทส}{mtk.2}
\bookmark[dest={mtk.2},level=2]{1. กามสุตฺตนิทฺเทส}
\proseitem{1}{\csromanfolio{1}\csromansymbolfootnote{*}{ขุ 1. 399; ขุ 10. 6, 59 ปิฏฺเฐสุปิ.}กามํ \textbf{กามยมานสฺส, ตสฺส เจ ตํ สมิชฺฌติ.}}

\csromangathameasure{อทฺธา ปีติมโน โหติ, ลทฺธา มจฺโจ ยทิจฺฉติ.}
\csromangathagroup{\csromangathastanza{อทฺธา ปีติมโน โหติ, ลทฺธา มจฺโจ ยทิจฺฉติ.}}

\prose{\textbf{กามํ กามยมานสฺสา}ติ \textbf{กามา}ติ อุทฺทานโต เทฺว \textbf{กามา} วตฺถุ\textbf{กามา} จ กิเลส\textbf{กามา} จ.\csromanspacer{} กตเม \textbf{วตฺถุกามา\csromandash{}}มนาปิกา รูปา มนาปิกา สทฺทา มนาปิกา คนฺธา มนาปิกา รสา มนาปิกา โผฏฺฐพฺพา, อตฺถรณา ปาวุรณา\footnote{ปาปุรณา (สี, สฺยา) ขุ 8. 36 ปิฏฺเฐปิ ปสฺสิตพฺพํ.} ทาสิทาสา อเชฬกา กุกฺกุฏสูกรา หตฺถิควาสฺสวฬวา เขตฺตํ วตฺถุ หิรญฺญํ สุวณฺณํ คามนิคมราชธานิโย รฏฺฐญฺจ ชนปโท จ โกโส จ โกฏฺฐาคารญฺจ, ยํ กิญฺจิ รชนียํ วตฺถุ วตฺถุ\textbf{กามา}.}

\prose{อปิ จ อตีตา \textbf{กามา} อนาคตา \textbf{กามา} ปจฺจุปฺปนฺนา \textbf{กามา}, อชฺฌตฺตา \textbf{กามา} พหิทฺธา \textbf{กามา} อชฺฌตฺตพหิทฺธา \textbf{กามา}, หีนา \textbf{กามา} มชฺฌิมา \textbf{กามา} ปณีตา \textbf{กามา}, อาปายิกา \textbf{กามา} มานุสิกา \textbf{กามา} ทิพฺพา \textbf{กามา}, ปจฺจุปฏฺฐิตา \textbf{กามา} นิมฺมิตา \textbf{กามา} อนิมฺมิตา \textbf{กามา} ปรนิมฺมิตา \textbf{กามา} ปริคฺคหิตา \textbf{กามา} อปริคฺคหิตา \textbf{กามา} มมายิตา \textbf{กามา} อมมายิตา \textbf{กามา}, สพฺเพปิ \textbf{กามา}วจรา ธมฺมา สพฺเพปิ รูปาวจรา ธมฺมา สพฺเพปิ \csromanfolio{2}อรูปาวจรา ธมฺมา ตณฺหาวตฺถุกา ตณฺหารมฺมณา กามนียฏฺเฐน รชนียฏฺเฐน มทนียฏฺเฐน กามา, อิเม วุจฺจนฺติ วตฺถุกามา.}

\prose{กตเม \textbf{กิเลสกามา\csromandash{}}ฉนฺโท กาโม ราโค กาโม ฉนฺทราโค กาโม, สงฺกปฺโป กาโม ราโค กาโม สงฺกปฺปราโค กาโม, โย กาเมสุ กามจฺฉนฺโท กามราโค กามนนฺที กามตณฺหา กามสฺเนโห กามปริฬาโห กามมุจฺฉา กามชฺโฌสานํ กาโมโฆ กามโยโค กามุปาทานํ กามจฺฉนฺทนีวรณํ.}

\csromangathameasure{อทฺทสํ กาม เต มูลํ, สงฺกปฺปา กาม ชายสิ.}
\csromangathagroup{\csromangathastanza{อทฺทสํ กาม เต มูลํ, สงฺกปฺปา กาม ชายสิ.}}

\vspace{\tipitakaparskip}
\csromancenter{น ตํ สงฺกปฺปยิสฺสามิ, เอวํ กาม น เหหิสีติ\footnote{น โหหิสีติ (ก) ขุ 8. 37 ปิฏฺเฐปิ.}\csromandash{}}
\prose{อิเม วุจฺจนฺติ กิเลสกามา.\csromanspacer{}\csromanspacer{}\textbf{กามยมานสฺสา}ติ กามยมานสฺส อิจฺฉมานสฺส สาทิยมานสฺส ปตฺถยมานสฺส ปิหยมานสฺส อภิชปฺปมานสฺสาติ กามํ กามยมานสฺส.}

\prose{\textbf{ตสฺส เจ ตํ สมิชฺฌตี}ติ \textbf{ตสฺส เจ}ติ ตสฺส ขตฺติยสฺส วา พฺราหฺมณสฺส วา เวสฺสสฺส วา สุทฺทสฺส วา คหฏฺฐสฺส วา ปพฺพชิตสฺส วา เทวสฺส วา มนุสฺสสฺส วา.\csromanspacer{}\csromanspacer{}\textbf{ตนฺ}ติ วตฺถุกามา วุจฺจนฺติ มนาปิกา รูปา มนาปิกา สทฺทา มนาปิกา คนฺธา มนาปิกา รสา มนาปิกา โผฏฺฐพฺพา.\csromanspacer{}\csromanspacer{}\textbf{สมิชฺฌตี}ติ อิชฺฌติ สมิชฺฌติ ลภติ ปฏิลภติ อธิคจฺฉติ วินฺทตีติ \textbf{ตสฺส เจ} ตํ สมิชฺฌติ.}

\prose{\textbf{อทฺธา ปีติมโน โหตี}ติ \textbf{อทฺธา}ติ เอกํสวจนํ นิสฺสํสยวจนํ นิกฺกงฺขาวจนํ อเทฺวชฺฌวจนํ อเทฺวฬฺหกวจนํ นิโยควจนํ อปณฺณกวจนํ อวตฺถาปนวจนเมตํ อทฺทาติ.\csromanspacer{}\csromanspacer{}\textbf{ปีตี}ติ ยา ปญฺจกามคุณปฏิสญฺญุตฺตา ปีติ ปามุชฺชํ อาโมทนา ปโมทนา หาโส ปหาโส วิตฺติ ตุฏฺฐิ โอทคฺยํ อตฺตมนตา อภิผรณตา จิตฺตสฺส.\csromanspacer{}\csromanspacer{}\textbf{มโน}ติ ยํ จิตฺตํ มโน มานสํ หทยํ ปณฺฑรํ มโน มนายตนํ มนินฺทฺริยํ วิญฺญาณํ วิญฺญาณกฺขนฺโธ ตชฺชา มโนวิญฺญาณธาตุ, อยํ วุจฺจติ มโน.\csromanspacer{} อยํ มโน อิมาย ปีติยา สหคโต โหติ สหชาโต สํสฏฺโฐ สมฺปยุตฺโต เอกุปฺปาโท เอกนิโรโธ เอกวตฺถุโก เอการมฺมโณ.\csromanspacer{}\csromanspacer{}\textbf{ปีติมโน โหตี}ติ ปีติมโน โหติ ตุฏฺฐมโน หฏฺฐมโน \csromanfolio{3}ปหฏฺฐมโน อตฺตมโน อุทคฺคมโน มุทิตมโน ปโมทิตมโน โหตีติ อทฺธา ปีติมโน โหติ.}

\prose{\textbf{ลทฺธา มจฺโจ ยทิจฺฉตี}ติ \textbf{ลทฺธา}ติ ลภิตฺวา ปฏิลภิตฺวา อธิคนฺตฺวา วินฺทิตฺวา.\csromanspacer{}\csromanspacer{}\textbf{มจฺโจ}ติ สตฺโต นโร มานโว โปโส ปุคฺคโล ชีโว ชาคุ\footnote{ชาตุ (สฺยา), ชคุ (ก) 2. หินฺทคู (สี, สฺยา) * ขุ 1. 399; ขุ 10. 6, 59 ปิฏฺเฐสุปิ.} ชนฺตุ อินฺทคุ2 มนุโช.\csromanspacer{}\csromanspacer{}\textbf{ยทิจฺฉตี}ติ ยํ อิจฺฉติ ยํ สาทิยติ ยํ ปตฺเถติ ยํ ปิเหติ ยํ อภิชปฺปติ รูปํ วา สทฺทํ วา คนฺธํ วา รสํ วา โผฏฺฐพฺพํ วาติ \textbf{ลทฺธา} มจฺโจ ยทิจฺฉติ.\csromanspacer{} เตนาห ภควา\csromandash{}}

\csromangathameasure{กามํ กามยมานสฺส, ตสฺส เจ ตํ สมิชฺฌติ. \\ อทฺธา ปีติมโน โหติ, ลทฺธา มจฺโจ ยทิจฺฉตีติ.}
\csromangathagroup{\csromangathastanza{กามํ กามยมานสฺส, ตสฺส เจ ตํ สมิชฺฌติ. \\ อทฺธา ปีติมโน โหติ, ลทฺธา มจฺโจ ยทิจฺฉตีติ.}}

\proseitem{2}{\csromansymbolmark{*}ตสฺส เจ กามยานสฺส, ฉนฺทชาตสฺส ชนฺตุโน.}

\csromangathameasure{เต กามาปริหายนฺติ, สลฺลวิทฺโธว รุปฺปติ.}
\csromangathagroup{\csromangathastanza{เต กามาปริหายนฺติ, สลฺลวิทฺโธว รุปฺปติ.}}

\prose{\textbf{ตสฺส เจ กามยานสฺสา}ติ \textbf{ตสฺส เจ}ติ ตสฺส ขตฺติยสฺส วา พฺราหฺมณสฺส วา เวสฺสสฺส วา สุทฺทสฺส วา คหฏฺฐสฺส วา ปพฺพชิตสฺส วา เทวสฺส วา มนุสฺสสฺส วา.\csromanspacer{}\csromanspacer{}\textbf{กามยานสฺสา}ติ กาเม อิจฺฉมานสฺส สาทิยมานสฺส ปตฺถยมานสฺส ปิหยมานสฺส อภิชปฺปมานสฺส.\csromanspacer{}\csromanspacer{}อถ วา กามตณฺหาย ยายติ นิยฺยติ วุยฺหติ สํหรียติ, ยถา หตฺถิยาเนน วา อสฺสยาเนน วา โคยาเนน วา อชยาเนน วา เมณฺฑยาเนน วา โอฏฺฐยาเนน วา ขรยาเนน วา ยายติ นิยฺยติ วุยฺหติ สํหรียติ.\csromanspacer{} เอวเมวํ กามตณฺหาย ยายติ นิยฺยติ วุยฺหติ สํหรียตีติ \textbf{ตสฺส เจ} กามยานสฺส.}

\prose{\textbf{ฉนฺทชาตสฺส ชนฺตุโน}ติ \textbf{ฉนฺโท}ติ โย กาเมสุ กามจฺ\textbf{ฉนฺโท} กามราโค กามนนฺที กามตณฺหา กามสฺเนโห กามปริฬาโห กามมุจฺฉา กามชฺโฌสานํ กาโมโฆ กามโยโค กามุปาทานํ กามจฺฉนฺทนีวรณํ, ตสฺส โส กามจฺ\textbf{ฉนฺโท} ชาโต โหติ สญฺชาโต นิพฺพตฺโต อภินิพฺพตฺโต ปาตุภูโต.\csromanspacer{}\csromanspacer{}\textbf{ชนฺตุโน}ติ สตฺตสฺส นรสฺส มานวสฺส โปสสฺส ปุคฺคลสฺส ชีวสฺส ชาคุสฺส ชนฺตุสฺส อินฺทคุสฺส มนุชสฺสาติ ฉนฺทชาตสฺส ชนฺตุโน.}

\prose{\textbf{เต กามา ปริหายนฺตี}ติ เต วา กามา ปริหายนฺติ, โส วา กาเมหิ ปริหายติ.\csromanspacer{} กถํ เต กามา ปริหายนฺติ\csromandash{}ตสฺส \csromanfolio{4}ติฏฺฐนฺตสฺเสว เต โภเค ราชาโน วา หรนฺติ, โจรา วา หรนฺติ, อคฺคิ วา ทหติ, อุทกํ วา วหติ, อปฺปิยา วา ทายาทา หรนฺติ, นิหิตํ วา นาธิคจฺฉติ, ทุปฺปยุตฺตา วา กมฺมนฺตา ภิชฺชนฺติ, กุเล วา กุลงฺคาโร อุปฺปชฺชติ, \textbf{โย} เต โภเค วิกิรติ วิธมติ\footnote{วิธเมติ (สฺยา)} วิทฺธํเสติ อนิจฺจตาเยว อฏฺฐมี.\csromanspacer{} เอวํ เต \textbf{กามา} หายนฺติ ปริหายนฺติ ปริธํเสนฺติ ปริปตนฺติ อนฺตรธายนฺติ วิปฺปลุชฺชนฺติ.\csromanspacer{} กถํ โส กาเมหิ ปริหายติ\csromandash{}ติฏฺฐนฺเตว เต โภเค โส จวติ มรติ วิปฺปลุชฺชติ.\csromanspacer{} เอวํ โส กาเมหิ หายติ ปริหายติ ปริธํสติ ปริปตติ อนฺตรธายติ วิปฺปลุชฺชติ.}

\csromangathameasure{โจรา หรนฺติ ราชาโน, อคฺคิ ทหติ นสฺสติ. \\ อถ อนฺเตน ชหติ, สรีรํ สปริคฺคหํ. \\ เอตทญฺญาย เมธาวี, ภุญฺเชถ จ ทเทถ จ. \\ ทตฺวา จ ภุตฺวา จ ยถานุภาวํ, อนนฺทิโต สคฺคมุเปติ ฐานนฺติ เต กามา ปริหายนฺติ.}
\csromangathagroup{\csromangathastanza{โจรา หรนฺติ ราชาโน, อคฺคิ ทหติ นสฺสติ. \\ อถ อนฺเตน ชหติ\footnote{อโถ อนฺเตน ชหติ (สฺยา), อสหนฺเตน ทหติ (ก)}, สรีรํ สปริคฺคหํ.}\csromangathastanza{เอตทญฺญาย เมธาวี, ภุญฺเชถ จ ทเทถ จ. \\ ทตฺวา จ ภุตฺวา จ ยถานุภาวํ, อนนฺทิโต สคฺคมุเปติ ฐานนฺติ เต กามา ปริหายนฺติ.}}

\prose{\textbf{สลฺลวิทฺโธว รุปฺปตี}ติ ยถา อ\textbf{โย}มเยน วา สลฺเลน วิทฺโธ, อฏฺฐิมเยน วา สลฺเลน ทนฺตมเยน วา สลฺเลน วิสาณมเยน วา สลฺเลน กฏฺฐมเยน วา สลฺเลน วิทฺโธ รุปฺปติ กุปฺปติ ฆฏฺฏียติ ปีฬียติ, พฺยาธิโต โทมนสฺสิโต โหติ.\csromanspacer{} เอวเมวํ วตฺถุกามานํ วิปริณามญฺญถาภาวา อุปฺปชฺชนฺติ โสกปริเทวทุกฺขโทมนสฺสุปายาสา.\csromanspacer{} โส กามสลฺเลน จ โสกสลฺเลน จ วิทฺโธ รุปฺปติ กุปฺปติ ฆฏฺฏียติ ปีฬียติ, พฺยาธิโต โทมนสฺสิโต โหตีติ สลฺลวิทฺโธว รุปฺปติ.\csromanspacer{} เตนาห ภควา\csromandash{}}

\csromangathameasure{ตสฺส เจ กามยานสฺส, ฉนฺทชาตสฺส ชนฺตุโน. \\ เต กามา ปริหายนฺติ, สลฺลวิทฺโธว รุปฺปตีติ.}
\csromangathagroup{\csromangathastanza{ตสฺส เจ กามยานสฺส, ฉนฺทชาตสฺส ชนฺตุโน. \\ เต กามา ปริหายนฺติ, สลฺลวิทฺโธว รุปฺปตีติ.}}

\proseitem{3}{\csromansymbolfootnote{*}{ขุ 1. 399; ขุ 10. 6, 59 ปิฏฺเฐสุปิ.}โย กาเม ปริวชฺเชติ, สปฺปสฺเสว ปทา สิโร.}

\csromangathameasure{โส’มํ วิสตฺติกํ โลเก, สโต สมติวตฺตติ.}
\csromangathagroup{\csromangathastanza{โส’มํ วิสตฺติกํ โลเก, สโต สมติวตฺตติ.}}

\prose{\textbf{โย กาเม ปริวชฺเชตี}ติ \textbf{โย}ติ \textbf{โย} ยาทิโส ยถายุตฺโต ยถาวิหิโต ยถาปกาโร ยํฐานปฺปตฺโต ยํธมฺมสมนฺนาคโต ขตฺติ\textbf{โย} วา พฺราหฺมโณ วา เวสฺโส วา สุทฺโท วา คหฏฺโฐ วา \csromanfolio{5}ปพฺพชิโต วา เทโว วา มนุสฺ\textbf{โส} วา.\csromanspacer{}\csromanspacer{}\textbf{กาเม ปริวชฺเชตี}ติ \textbf{กามา}ติ อุทฺทานโต เทฺว \textbf{กามา} วตฺถุ\textbf{กามา} จ กิเลส\textbf{กามา} จ ฯเปฯ อิเม วุจฺจนฺติ วตฺถุ\textbf{กามา} ฯเปฯ อิเม วุจฺจนฺติ กิเลส\textbf{กามา}.\csromanspacer{}\csromanspacer{}\textbf{กาเม ปริวชฺเชตี}ติ ทฺวีหิ การเณหิ กาเม ปริวชฺเชติ วิกฺขมฺภนโต วา สมุจฺเฉทโต วา.\csromanspacer{} กถํ วิกฺขมฺภนโต กาเม ปริวชฺเชติ\csromandash{}“อฏฺฐิกงฺกลูปมา \textbf{กามา} อปฺปสฺสาทฏฺเฐนา”ติ * ปสฺสนฺโต วิกฺขมฺภนโต กาเม ปริวชฺเชติ, “มํสเปสูปมา \textbf{กามา} พหุสาธารณฏฺเฐนา”ติ ปสฺสนฺโต วิกฺขมฺภนโต กาเม ปริวชฺเชติ, “ติณุกฺกูปมา \textbf{กามา} อนุทหนฏฺเฐนา”ติ ปสฺสนฺโต วิกฺขมฺภนโต กาเม ปริวชฺเชติ, “องฺคารกาสูปมา \textbf{กามา} มหาปริฬาหฏฺเฐนา”ติ ปสฺสนฺโต วิกฺขมฺภนโต กาเม ปริวชฺเชติ, “สุปินกูปมา \textbf{กามา} อิตฺตรปจฺจุปฏฺฐานฏฺเฐนา”ติ ปสฺสนฺโต วิกฺขมฺภนโต กาเม ปริวชฺเชติ, “ยาจิตกูปมา \textbf{กามา} ตาวกาลิกฏฺเฐนา”ติ ปสฺสนฺโต วิกฺขมฺภนโต กาเม ปริวชฺเชติ, “รุกฺขผลูปมา \textbf{กามา} สมฺภญฺชนปริภญฺชนฏฺเฐนา”ติ ปสฺสนฺโต วิกฺขมฺภนโต กาเม ปริวชฺเชติ, “อสิสูนูปมา \textbf{กามา} อธิกุฏฺฏนฏฺเฐนา”ติ\footnote{อธิกนฺตนฏฺเฐนาติ (สฺยา)} ปสฺสนฺโต วิกฺขมฺภนโต กาเม ปริวชฺเชติ, “สตฺติสูลูปมา \textbf{กามา} วินิวิชฺฌนฏฺเฐนา”ติ ปสฺสนฺโต วิกฺขมฺภนโต กาเม ปริวชฺเชติ, “สปฺปสิรูปมา \textbf{กามา} สปฺปฏิภยฏฺเฐนา”ติ ปสฺสนฺโต วิกฺขมฺภนโต กาเม ปริวชฺเชติ, “อคฺคิกฺขนฺธูปมา \textbf{กามา} มหาภิตาปนฏฺเฐนา”ติ ปสฺสนฺโต วิกฺขมฺภนโต กาเม ปริวชฺเชติ.}

\prose{พุทฺธานุสฺสติํ ภาเวนฺโตปิ วิกฺขมฺภนโต กาเม ปริวชฺเชติ, ธมฺมานุสฺสติํ ภาเวนฺโตปิ ฯเปฯ สํฆานุสฺสติํ ภาเวนฺโตปิ.\csromanspacer{} สีลานุสฺสติํ ภาเวนฺโตปิ.\csromanspacer{} จาคานุสฺสติํ ภาเวนฺโตปิ.\csromanspacer{} เทวตานุสฺสติํ ภาเวนฺโตปิ.\csromanspacer{} อานาปานสฺสติํ\footnote{อานาปานสติํ (สี)} ภาเวนฺโตปิ.\csromanspacer{} มรณสฺสติํ ภาเวนฺโตปิ.\csromanspacer{} กายคตาสติํ ภาเวนฺโตปิ.\csromanspacer{} อุปสมานุสฺสติํ ภาเวนฺโตปิ วิกฺขมฺภนโต กาเม ปริวชฺเชติ.}

\prose{ปฐมํ ฌานํ ภาเวนฺโตปิ วิกฺขมฺภนโต กาเม ปริวชฺเชติ.\csromanspacer{} ทุติยํ ฌานํ ภาเวนฺโตปิ ฯเปฯ.\csromanspacer{} ตติยํ ฌานํ ภาเวนฺโตปิ.\csromanspacer{} จตุตฺถํ ฌานํ ภาเวนฺโตปิ.\csromanspacer{} อากาสานญฺจายตนสมาปตฺติํ ภาเวนฺโตปิ.\csromanspacer{} วิญฺญาณญฺจายตนสมาปตฺติํ ภาเวนฺโตปิ.\csromanspacer{} อากิญฺจญฺญายตนสมาปตฺติํ \csromanfolio{6}ภาเวนฺโตปิ.\csromanspacer{} เนวสญฺญานาสญฺญายตนสมาปตฺติํ ภาเวนฺโตปิ วิกฺขมฺภนโต กาเม ปริวชฺเชติ.\csromanspacer{} เอวํ วิกฺขมฺภนโต กาเม ปริวชฺเชติ.}

\prose{กถํ สมุจฺเฉทโต กาเม ปริวชฺเชติ\csromandash{}\textbf{โส}ตาปตฺติมคฺคํ ภาเวนฺโตปิ อปายคมนีเย กาเม สมุจฺเฉทโต ปริวชฺเชติ, สกทาคามิมคฺคํ ภาเวนฺโตปิ โอฬาริเก กาเม สมุจฺเฉทโต ปริวชฺเชติ, อนาคามิมคฺคํ ภาเวนฺโตปิ อณุสหคเต กาเม สมุจฺเฉทโต ปริวชฺเชติ, อรหตฺตมคฺคํ ภาเวนฺโตปิ สพฺเพน สพฺพํ สพฺพถา สพฺพํ อเสสํ นิสฺเสสํ สมุจฺเฉทโต กาเม ปริวชฺเชติ.\csromanspacer{} เอวํ สมุจฺเฉทโต กาเม ปริวชฺเชตีติ โย กาเม ปริวชฺเชติ.}

\prose{\textbf{สปฺปสฺเสว ปทา สิโร}ติ \textbf{สปฺโป} วุจฺจติ อหิ.\csromanspacer{} เกนฏฺเฐน \textbf{สปฺโป}.\csromanspacer{} สํสปฺปนฺโต คจฺฉตีติ \textbf{สปฺโป}, ภุชนฺโต คจฺฉตีติ ภุชโค, อุเรน คจฺฉตีติ อุรโค, ปนฺนสิโร คจฺฉตีติ ปนฺนโค, สิเรน สุปตีติ\footnote{สปฺปตีติ (ก)} สรีสโป\footnote{สิริํสโป (สี)}, พิเล สยตีติ พิลาสโย, คุหายํ สยตีติ คุหาสโย, ทาฐา ตสฺส อาวุโธติ ทาฐาวุโธ, วิสํ ตสฺส โฆรนฺติ โฆรวิ\textbf{โส}, ชิวฺหา ตสฺส ทุวิธาติ ทฺวิชิโวฺห, ทฺวีหิ ชิวฺหาหิ รสํ สายตีติ ทฺวิรสญฺญู.\csromanspacer{} ยถา ปุริ\textbf{โส} ชีวิตุกาโม อมริตุกาโม สุขกาโม ทุกฺขปฏิกฺกูโล ปาเทน สปฺปสิรํ วชฺเชยฺย วิวชฺเชยฺย ปริวชฺเชยฺย อภินิวชฺเชยฺย.\csromanspacer{} เอวเมวํ สุขกาโม ทุกฺขปฏิกฺกูโล กาเม วชฺเชยฺย วิวชฺเชยฺย ปริวชฺเชยฺย อภินิวชฺเชยฺยาติ สปฺปสฺเสว ปทา สิโร.}

\prose{\textbf{โส’มํ วิสตฺติกํ โลเก, สโต สมติวตฺตตี}ติ \textbf{โส}ติ โย กาเม ปริวชฺเชติ.\csromanspacer{} วิสตฺติกา วุจฺจติ ตณฺหา. * โย ราโค สาราโค อนุนโย อนุโรโธ นนฺที นนฺทิราโค จิตฺตสฺส สาราโค อิจฺฉา มุจฺฉา อชฺโฌสานํ เคโธ ปลิเคโธ\footnote{ปฬิเคโธ (สี)} สงฺโค ปงฺโก เอชา มายา ชนิกา สญฺชนนี สิพฺพินี ชาลินี สริตา วิสตฺติกา สุตฺตํ วิสตา อายูหินี\footnote{อายูหนี (สี, สฺยา)} ทุติยา ปณิธิ ภวเนตฺติ วนํ วนโถ สนฺธโว สฺเนโห อเปกฺขา ปฏิพนฺธุ อาสา อาสีสนา อาสีสิตตฺตํ รูปาสา สทฺทาสา คนฺธาสา รสาสา โผฏฺฐพฺพาสา ลาภาสา ธนาสา ปุตฺตาสา ชีวิตาสา ชปฺปา ปชปฺปา อภิชปฺปา ชปฺปนา ชปฺปิตตฺตํ โลลุปฺปํ โลลุปฺปายนา โลลุปฺปายิตตฺตํ ปุจฺฉญฺชิกตา สาธุกมฺยตา อธมฺมราโค วิสมโลโภ \csromanfolio{7}นิกนฺติ นิกามนา ปตฺถนา ปิหนา สมฺปตฺถนา, กามตณฺหา ภวตณฺหา วิภวตณฺหา รูปตณฺหา อรูปตณฺหา นิโรธตณฺหา, รูปตณฺหา สทฺทตณฺหา คนฺธตณฺห, รสตณฺหา โผฏฺฐพฺพตณฺหา ธมฺมตณฺหา โอโฆ โยโค คนฺโถ อุปาทานํ อาวรณํ นีวรณํ ฉทนํ พนฺธนํ อุปฺปกฺกิเลโส อนุสโย ปริยุฏฺฐานํ ลตา เววิจฺฉํ ทุกฺขมูลํ ทุกฺขนิทานํ ทุกฺขปฺปภโว มารปาโส มารพฬิสํ มารวิสโย ตณฺหานที ตณฺหาชาลํ ตณฺหาคทฺทูลํ ตณฺหาสมุทฺโท อภิชฺฌา โลโภ อกุสลมูลํ.}

\prose{\textbf{วิสตฺติกา}ติ เกนฏฺเฐน วิสตฺติกา.\csromanspacer{} วิสตาติ วิสตฺติกา, วิสาลาติ วิสตฺติกา, วิสฏาติ วิสตฺติกา, วิสกฺกตีติ วิสตฺติกา, วิสํหรตีติ วิสตฺติกา, วิสํวาทิกาติ วิสตฺติกา, วิสมูลาติ วิสตฺติกา, วิสผลาติ วิสตฺติกา, วิสปริโภโคติ วิสตฺติกา, วิสาลา วา ปน สา ตณฺหา รูเป สทฺเท คนฺเธ รเส โผฏฺฐพฺเพ, กุเล คเณ อาวาเส ลาเภ ยเส ปสํสาย สุเข จีวเร ปิณฺฑปาเต เสนาสเน คิลานปจฺจยเภสชฺชปริกฺขาเร, กามธาตุยา รูปธาตุยา อรูปธาตุยา, กามภเว รูปภเว อรูปภเว, สญฺญาภเว อสญฺญาภเว เนวสญฺญานาสญฺญาภเว, เอกโวการภเว จตุโวการภเว ปญฺจโวการภเว, อตีเต อนาคเต ปจฺจุปฺปนฺเน, ทิฏฺฐสุตมุตวิญฺญาตพฺเพสุ ธมฺเมสุ วิสฏา วิตฺถตาติ วิสตฺติกา.}

\prose{\textbf{โลเก}ติ อปายโลเก มนุสฺสโลเก เทวโลเก, ขนฺธโลเก ธาตุโลเก อายตนโลเก.\csromanspacer{}\csromanspacer{}* \textbf{สโต}ติ จตูหิ การเณหิ สโต, กาเย กายานุปสฺสนาสติปฏฺฐานํ ภาเวนฺโต สโต, เวทนาสุ.\csromanspacer{} จิตฺเต.\csromanspacer{} ธมฺเมสุ ธมฺมานุปสฺสนาสติปฏฺฐานํ ภาเวนฺโต สโต.}

\prose{อปเรหิปิ จตูหิ การเณหิ สโต อสติปริวชฺชนาย สโต, สติกรณียานํ ธมฺมานํ กตตฺตา สโต, สติปริพนฺธานํ ธมฺมานํ หตตฺตา สโต, สตินิมิตฺตานํ ธมฺมานํ อสมฺมุฏฺฐตฺตา สโต.}

\prose{อปเรหิปิ จตูหิ การเณปิ สโต สติยา สมนฺนาคตตฺตา สโต, สติยา วสิตตฺตา สโต, สติยา ปาคุญฺญตาย สโต, สติยา อปจฺโจโรหณตาย\footnote{อปจฺโจโรปนตาย (สี)} สโต.}

\prose{\csromanfolio{8}อปเรหิปิ จตูหิ การเณหิ สโต, สตฺตตฺตา สโต, สนฺตตฺตา สโต, สมิตตฺตา สโต, สนฺตธมฺมสมนฺนาคตตฺตา สโต.\csromanspacer{} พุทฺธานุสฺสติยา สโต, ธมฺมานุสฺสติยา สโต, สํฆานุสฺสติยา สโต, สีลานุสฺสติยา สโต, จาคานุสฺสติยา สโต, เทวตานุสฺสติยา สโต, อานาปานสฺสติยา สโต, มรณสฺสติยา สโต, กายคตาสติยา สโต, อุปสมานุสฺสติยา สโต.\csromanspacer{} ยา สติ อนุสฺสติ ปฏิสฺสติ สติ สรณตา ธารณตา อปิลาปนตา อสมฺมุสฺสนตา สติ สตินฺทฺริยํ สติพลํ สมฺมาสติ สติสมฺโพชฺฌงฺโค เอกายนมคฺโค, อยํ วุจฺจติ สติ.\csromanspacer{} อิมาย สติยา อุเปโต โหติ สมุเปโต อุปคโต สมุปคโต อุปปนฺโน สมุปปนฺโน สมนฺนาคโต, โส วุจฺจติ สโต.}

\prose{\textbf{โส’มํ วิสตฺติกํ โลเก, สโต สมติวตฺตตี}ติ โลเก วา สา วิสตฺติกา, โลเก วา ตํ วิสตฺติกํ สโต ตรติ อุตฺตรติ ปตรติ สมติกฺกมติ วีติวตฺตตีติ โส’มํ วิสตฺติกํ โลเก, สโต สมติวตฺตติ.\csromanspacer{} เตนาห ภควา\csromandash{}}

\csromangathameasure{โย กาเม ปริวชฺเชติ, สปฺปสฺเสว ปทา สิโร.}
\csromangathagroup{\csromangathastanza{โย กาเม ปริวชฺเชติ, สปฺปสฺเสว ปทา สิโร.}}

\prose{\textbf{โส’มํ วิสตฺติกํ โลเก, สโต สมติวตฺตตี}ติ.}

\proseitem{4}{\csromansymbolfootnote{*}{ขุ 1. 399; ขุ 10. 6 ปิฏฺเฐสุปิ.}\textbf{เขตฺตํ วตฺถุํ หิรญฺญํ วา}, ควาสฺสํ ทาสโปริสํ.}

\csromangathameasure{ถิโย พนฺธู ปุถุ กาเม, โย นโร อนุคิชฺฌติ.}
\csromangathagroup{\csromangathastanza{ถิโย พนฺธู ปุถุ กาเม, โย นโร อนุคิชฺฌติ.}}

\prose{\textbf{เขตฺตํ วตฺถุํ หิรญฺญํ วา}ติ \textbf{เขตฺตนฺ}ติ สาลิกฺเขตฺตํ วีหิกฺเขตฺตํ มุคฺคกฺเขตฺตํ มาสกฺเขตฺตํ ยวกฺเขตฺตํ โคธุมกฺเขตฺตํ ติลกฺเขตฺตํ.\csromanspacer{}\csromanspacer{}\textbf{วตฺถุนฺ}ติ ฆรวตฺถุํ โกฏฺฐกวตฺถุํ ปุเรวตฺถุํ ปจฺฉาวตฺถุํ อารามวตฺถุํ วิหารวตฺถุํ.\csromanspacer{}\csromanspacer{}\textbf{หิรญฺญนฺ}ติ หิรญฺญํ วุจฺจติ กหาปโณติ เขตฺตํ วตฺถุํ หิรญฺญํ วา.}

\prose{\textbf{ควาสฺสํ ทาสโปริสนฺ}ติ \textbf{ควนฺ}ติ ควา\footnote{คาโว (ก)} วุจฺจนฺติ.\csromanspacer{}\csromanspacer{}\textbf{อสฺสา}ติ ปสุกาทโย วุจฺจนฺติ.\csromanspacer{}\csromanspacer{}\textbf{ทาสา}ติ จตฺตาโร ทาสา อนฺโตชาตโก ทาโส ธนกฺกีตโก ทาโส สามํ วา ทาสพฺยํ อุเปติ อกามโก วา ทาสวิสยํ อุเปติ.}

\csromangathameasure{อามาย ทาสาปิ ภวนฺติ เหเก, \\ ธเนน กีตาปิ ภวนฺติ ทาสา. \\ สามญฺจ เอเก อุปยนฺติ ทาสฺยํ, \\ ภยาปนุณฺณาปิ ภวนฺติ ทาสาติ.}
\csromangathagroup{\csromangathastanza{\csromanfolio{9}อามาย ทาสาปิ ภวนฺติ เหเก, \\ ธเนน กีตาปิ ภวนฺติ ทาสา. \\ สามญฺจ เอเก อุปยนฺติ ทาสฺยํ, \\ ภยาปนุณฺณาปิ ภวนฺติ ทาสาติ.}}

\prose{\textbf{ปุริสา}ติ ต\textbf{โย} ปุริสา ภตกา กมฺมกรา อุปชีวิโนติ ควาสฺสํ ทาสโปริสํ.}

\prose{\textbf{ถิโย พนฺธู ปุถุ กาเม}ติ \textbf{ถิโย}ติ อิตฺถิปริคฺคโห วุจฺจติ.\csromanspacer{} \textbf{พนฺธู}ติ จตฺตาโร พนฺธู ญาติพนฺธวาปิ พนฺธุ, โคตฺตพนฺธวาปิ พนฺธุ, มนฺตพนฺธวาปิ พนฺธุ, สิปฺปพนฺธวาปิ พนฺธุ.\csromanspacer{}\csromanspacer{}\textbf{ปุถุ กาเม}ติ พหู กาเม.\csromanspacer{} เอเต ปุถุ กามา มนาปิกา รูปา ฯเปฯ มนาปิกา โผฏฺฐพฺพาติ \textbf{ถิโย} พนฺธู ปุถุ กาเม.}

\prose{\textbf{โย นโร อนุคิชฺฌตี}ติ \textbf{โย}ติ \textbf{โย} ยาทิโส ยถายุตฺโต ยถาวิหิโต ยถาปกาโร ยํฐานปฺปตฺโต ยํธมฺมสมนฺนาคโต ขตฺติ\textbf{โย} วา พฺราหฺมโณ วา เวสฺโส วา สุทฺโท วา คหฏฺโฐ วา ปพฺพชิโต วา เทโว วา มนุสฺโส วา.\csromanspacer{} \textbf{นโร}ติ สตฺโต นโร มานโว โปโส ปุคฺคโล ชีโว ชาคุ ชนฺตุ อินฺทคุ มนุโช.\csromanspacer{}\csromanspacer{}\textbf{อนุคิชฺฌตี}ติ กิเลสกาเมน วตฺถุกาเมสุ คิชฺฌติ อนุคิชฺฌติ ปลิคิชฺฌติ ปลิพชฺฌตีติ \textbf{โย} นโร อนุคิชฺฌติ.\csromanspacer{} เตนาห ภควา\csromandash{}}

\csromangathameasure{เขตฺตํ วตฺถุํ หิรญฺญํ วา, ควาสฺสํ ทาสโปริสํ. \\ \textbf{ถิโย พนฺธู ปุถุ กาเม}, \textbf{โย} นโร อนุคิชฺฌตีติ.}
\csromangathagroup{\csromangathastanza{เขตฺตํ วตฺถุํ หิรญฺญํ วา, ควาสฺสํ ทาสโปริสํ. \\ \textbf{ถิโย พนฺธู ปุถุ กาเม}, \textbf{โย} นโร อนุคิชฺฌตีติ.}}

\proseitem{5}{\csromansymbolfootnote{*}{ขุ 1. 399; ขุ 10. 6 ปิฏฺเฐสุปิ.}อพลา นํ พลียนฺติ, มทฺทนฺเต นํ ปริสฺสยา.}

\csromangathameasure{ตโต นํ ทุกฺขมเนฺวติ, นาวํ ภินฺนมิโวทกํ.}
\csromangathagroup{\csromangathastanza{ตโต นํ ทุกฺขมเนฺวติ, นาวํ ภินฺนมิโวทกํ.}}

\prose{\textbf{อพลานํ พลียนฺตี}ติ \textbf{อพลา}ติ \textbf{อพลา} กิเลสา ทุพฺพลา อปฺปฺ\textbf{อพลา} อปฺปถามกา หีนา นิหีนา ( )\footnote{(ปริหีนา) (สี, สฺยา)} โอมกา ลามกา ฉตุกฺกา ปริตฺตา.\csromanspacer{} เต กิเลสา ตํ ปุคฺคลํ สหนฺติ ปริสหนฺติ อภิภวนฺติ อชฺโฌตฺถรนฺติ ปริยาทิยนฺติ มทฺทนฺตีติ เอวมฺปิ \textbf{อพลา} นํ พลียนฺติ.\csromanspacer{} อถ วา อพลํ ปุคฺคลํ ทุพฺพลํ อปฺปพลํ อปฺปถามกํ หีนํ นิหีนํ โอมกํ ลามกํ ฉตุกฺกํ ปริตฺตํ.\csromanspacer{} ยสฺส นตฺถิ สทฺธาพลํ วีริยพลํ สติพลํ สมาธิพลํ ปญฺญาพลํ หิริพลํ \csromanfolio{10}โอตฺตปฺปพลํ.\csromanspacer{} เต กิเลสา ตํ ปุคฺคลํ สหนฺติ ปริสหนฺติ อภิภวนฺติ อชฺโฌตฺถรนฺติ ปริยาทิยนฺติ มทฺทนฺตีติ เอวมฺปิ อพลา นํ พลียนฺตีติ.}

\prose{\textbf{มทฺทนฺเต นํ ปริสฺสยา}ติ * เทฺว ปริสฺสยา \textbf{ปากฏปริสฺสยา} จ ปฏิจฺฉนฺนปริสฺสยา จ.\csromanspacer{} กตเม \textbf{ปากฏปริสฺสยา}\csromandash{}สีหา พฺยคฺฆา ทีปี อจฺฉา ตรจฺฉา โกกา มหิํสา\footnote{มหิสา (สี, สฺยา)} หตฺถี อหี วิจฺฉิกา สตปที, โจรา วา อสฺสุ มานวา วา กตกมฺมา วา อกตกมฺมา วา, จกฺขุโรโค โสตโรโค ฆานโรโค ชิวฺหาโรโค กายโรโค สีสโรโค กณฺณโรโค มุขโรโค ทนฺตโรโค กาโส สาโส ปินาโส ฑาโห ชโร กุจฺฉิโรโค มุจฺฉา ปกฺขนฺทิกา สูลา วิสูจิกา กุฏฺฐํ คณฺโฑ กิลาโส โสโส อปมาโร ททฺทุ กณฺฑุ กจฺฉุ รขสา\footnote{รกฺขสา (ก)} วิตจฺฉิกา โลหิตปิตฺตํ มธุเมโห อํสา ปิฬกา ภคนฺทลา ปิตฺตสมุฏฺฐานา อาพาธา เสมฺหสมุฏฺฐานา อาพาธา วาตสมุฏฺฐานา อาพาธา สนฺนิปาติกา อาพาธา อุตุปริณามชา อาพาธา วิสมปริหารชา อาพาธา โอปกฺกมิกา อาพาธา กมฺมวิปากชา อาพาธา สีตํ อุณฺหํ ชิฆจฺฉา ปิปาสา อุจฺจาโร ปสฺสาโว ฑํส มกส วาตาตป สรีสป สมฺผสฺสา อิติ วา, อิเม วุจฺจนฺติ \textbf{ปากฏปริสฺสยา}.}

\prose{กตเม \textbf{ปฏิจฺฉนฺนปริสฺสยา\csromandash{}}กายทุจฺจริตํ วจีทุจฺจริตํ มโนทุจฺจริตํ, กามจฺฉนฺทนีวรณํ พฺยาปาทนีวรณํ ถินมิทฺธนีวรณํ อุทฺธจฺจกุกฺกุจฺจนีวรณํ วิจิกิจฺฉานีวรณํ, ราโค โทโส โมโห, โกโธ อุปนาโห มกฺโข ปฬาโส อิสฺสา มจฺฉริยํ มายา สาเฐยฺยํ ถมฺโภ สารมฺโภ มาโน อติมาโน มโท ปมาโท สพฺเพ กิเลสา สพฺเพ ทุจฺจริตา สพฺเพ ทรถา สพฺเพ ปริฬาหา สพฺเพ สนฺตาปา สพฺพากุสลาภิสงฺขารา, อิเม วุจฺจนฺติ ปฏิจฺฉนฺนปริสฺสยา.}

\prose{\textbf{ปริสฺสยา}ติ เกนฏฺเฐน ปริสฺสยา\csromandash{}ปริสหนฺตีติ ปริสฺสยา, ปริหานาย สํวตฺตนฺตีติ ปริสฺสยา, ตตฺราสยาติ ปริสฺสยา.\csromanspacer{} กถํ ปริสหนฺตีติ ปริสฺสยา\csromandash{}เต ปริสฺสยา ตํ ปุคฺคลํ สหนฺติ ปริสหนฺติ อภิภวนฺติ อชฺโฌตฺถรนฺติ ปริยาทิยนฺติ มทฺทนฺติ.\csromanspacer{} เอวํ ปริสหนฺตีติ ปริสฺสยา.\csromanspacer{}\csromanspacer{}กถํ ปริหานาย สํวตฺตนฺตีติ ปริสฺสยา\csromandash{}เต ปริสฺสยา กุสลานํ ธมฺมานํ อนฺตรายาย ปริหานาย สํวตฺตนฺติ. \csromanfolio{11}กตเมสํ กุสลานํ ธมฺมานํ สมฺมาปฏิปทาย อนุโลมปฏิปทาย อปจฺจนีกปฏิปทาย อวิรุทฺธปฏิปทาย อนฺวตฺถปฏิปทาย ธมฺมานุธมฺมปฏิปทาย สีเลสุ ปริปูริการิตาย อินฺทฺริเยสุ คุตฺตทฺวารตาย โภชเน มตฺตญฺญุตาย ชาคริยานุโยคสฺส สติสมฺปชญฺญสฺส จตุนฺนํ สติปฏฺฐานานํ ภาวนานุโยคสฺส จตุนฺนํ สมฺมปฺปธานานํ ภาวนานุโยคสฺส จตุนฺนํ อิทฺธิปาทานํ ภาวนานุโยคสฺส ปญฺจนฺนํ อินฺทฺริยานํ ภาวนานุโยคสฺส ปญฺจนฺนํ พลานํ ภาวนานุโยคสฺส สตฺตนฺนํ โพชฺฌงฺคานํ ภาวนานุโยคสฺส อริยสฺส อฏฺฐงฺคิกสฺส มคฺคสฺส ภาวนานุโยคสฺส อิเมสํ กุสลานํ ธมฺมานํ อนฺตรายาย ปริหานาย สํวตฺตนฺติ.\csromanspacer{} เอวํ ปริหานาย สํวตฺตนฺตีติ ปริสฺสยา.}

\prose{กถํ ตตฺราสยาติ ปริสฺสยา\csromandash{}ตตฺเถเต ปาปกา อกุสลา ธมฺมา อุปฺปชฺชนฺติ อตฺตภาวสนฺนิสฺสยา, ยถา พิเล พิลาสยา ปาณา สยนฺติ, ทเก ทกาสยา ปาณา สยนฺติ, วเน วนาสยา ปาณา สยนฺติ, รุกฺเข รุกฺขาสยา ปาณา สยนฺติ.\csromanspacer{} เอวเมวํ ตตฺเถเต ปาปกา อกุสลา ธมฺมา อุปฺปชฺชนฺติ อตฺตภาวสนฺนิสฺสยา.\csromanspacer{} เอวมฺปิ ตตฺราสยาติ ปริสฺสยา.}

\prose{วุตฺตํ เหตํ ภควตา\csromandash{} * สานฺเตวาสิโก ภิกฺขเว ภิกฺขุ สาจริยโก ทุกฺขํ น ผาสุ วิหรติ.\csromanspacer{} กถญฺจ ภิกฺขเว ภิกฺขุ สานฺเตวาสิโก สาจริยโก ทุกฺขํ น ผาสุ วิหรติ\csromandash{}อิธ ภิกฺขเว ภิกฺขุโน จกฺขุนา รูปํ ทิสฺวา อุปฺปชฺชนฺติ เย ปาปกา อกุสลา ธมฺมา สรสงฺกปฺปา สญฺโญชนิยา, ตฺยสฺส อนฺโต วสนฺติ อนฺวาสวนฺติ ปาปกา อกุสลา ธมฺมาติ ตสฺมา “สานฺเตวาสิโก”ติ วุจฺจติ.\csromanspacer{} เต นํ สมุทาจรนฺติ, สมุทาจรนฺติ นํ ปาปกา อกุสลา ธมฺมาติ ตสฺมา “สาจริยโก”ติ วุจฺจติ.}

\prose{ปุน จปรํ ภิกฺขเว ภิกฺขุโน โสเตน สทฺทํ สุตฺวา, ฆาเนน คนฺธํ ฆายิตฺวา, ชิวฺหาย รสํ สายิตฺวา, กาเยน โผฏฺฐพฺพํ ผุสิตฺวา, มนสา ธมฺมํ วิญฺญาย อุปฺปชฺชนฺติ เย ปาปกา อกุสลา ธมฺมา สรสงฺกปฺปา สญฺโญชนิยา, ตฺยสฺส อนฺโต วสนฺติ อนฺวาสวนฺติ ปาปกา อกุสลา ธมฺมาติ ตสฺมา “สานฺเตวาสิโก”ติ วุจฺจติ.\csromanspacer{} เต นํ สมุทาจรนฺติ, สมุทาจรนฺติ นํ ปาปกา อกุสลา ธมฺมาติ ตสฺมา “สาจริยโก”ติ วุจฺจติ.\csromanspacer{} เอวํ โข ภิกฺขเว ภิกฺขุ สานฺเตวาสิโก สาจริยโก ทุกฺขํ น ผาสุ วิหรตีติ.\csromanspacer{} เอวมฺปิ ตตฺราสยาติ ปริสฺสยา.}

\prose{\csromanfolio{12}วุตฺตํ เหตํ ภควตา\csromandash{} * ตโยเม ภิกฺขเว อนฺตรามลา อนฺตรา-อมิตฺตา อนฺตราสปตฺตา อนฺตราวธกา อนฺตราปจฺจตฺถิกา.\csromanspacer{} กตเม ตโย\csromandash{} โลโภ ภิกฺขเว อนฺตรามโล\footnote{อนฺตรามลํ (สี, ก)} อนฺตรา-อมิตฺโต อนฺตราสปตฺโต อนฺตราวธโก อนฺตราปจฺจตฺถิโก.\csromanspacer{} โทโส ฯเปฯ.\csromanspacer{} โมโห ภิกฺขเว อนฺตรามโล อนฺตรา-อมิตฺโต อนฺตราสปตฺโต อนฺตราวธโก อนฺตราปจฺจตฺถิโก.\csromanspacer{} อิเม โข ภิกฺขเว ตโย อนฺตรามลา อนฺตรา อมิตฺตา อนฺตราสปตฺตา อนฺตราวธกา อนฺตราปจฺจตฺถิกา.}

\csromangathameasure{*อนตฺถชนโน โลโภ, โลโภ จิตฺตปฺปโกปโน. \\ ภยมนฺตรโต ชาตํ, ตํ ชโน นาวพุชฺฌติ. \\ ลุทฺโธ อตฺถํ น ชานาติ, ลุทฺโธ ธมฺมํ น ปสฺสติ. \\ อนฺธนฺตมํ ตทา โหติ, ยํ โลโภ สหเต นรํ. \\ อนตฺถชนโน โทโส, โทโส จิตฺตปฺปโกปโน. \\ ภยมนฺตรโต ชาตํ, ตํ ชโน นาวพุชฺฌติ. \\ กุทฺโธ อตฺถํ น ชานาติ, กุทฺโธ ธมฺมํ น ปสฺสติ. \\ อนฺธนฺตมํ ตทา โหติ, ยํ โทโส สหเต นรํ. \\ อนตฺถชนโน โมโห, โมโห จิตฺตปฺปโกปโน. \\ ภยมนฺตรโต ชาตํ, ตํ ชโน นาวพุชฺฌติ. \\ มูโฬฺห อตฺถํ น ชนาติ, มูโฬฺห ธมฺมํ น ปสฺสติ.}
\csromangathagroup{\csromangathastanza{\csromansymbolfootnote{*}{อํ 2. 471; ขุ 1. 252; ขุ 8. 251.}อนตฺถชนโน โลโภ, โลโภ จิตฺตปฺปโกปโน. \\ ภยมนฺตรโต ชาตํ, ตํ ชโน นาวพุชฺฌติ.}\csromangathastanza{ลุทฺโธ อตฺถํ น ชานาติ, ลุทฺโธ ธมฺมํ น ปสฺสติ. \\ อนฺธนฺตมํ\footnote{อนฺธตมํ (สฺยา, ก)} ตทา โหติ, ยํ โลโภ สหเต นรํ.}\csromangathastanza{อนตฺถชนโน โทโส, โทโส จิตฺตปฺปโกปโน. \\ ภยมนฺตรโต ชาตํ, ตํ ชโน นาวพุชฺฌติ.}\csromangathastanza{กุทฺโธ อตฺถํ น ชานาติ, กุทฺโธ ธมฺมํ น ปสฺสติ. \\ อนฺธนฺตมํ ตทา โหติ, ยํ โทโส สหเต นรํ.}\csromangathastanza{อนตฺถชนโน โมโห, โมโห จิตฺตปฺปโกปโน. \\ ภยมนฺตรโต ชาตํ, ตํ ชโน นาวพุชฺฌติ. \\ มูโฬฺห อตฺถํ น ชนาติ, มูโฬฺห ธมฺมํ น ปสฺสติ.}}

\prose{อนฺธนฺตมํ ตทา โหติ, ยํ โมโห สหเต นรนฺติ\csromandash{}}

\prose{เอวมฺปิ ตตฺราสยาติ ปริสฺสยา.}

\prose{วุตฺตมฺปิ เหตํ ภควตา\csromandash{} + ตโย โข มหาราช ปุริสสฺส ธมฺมา อชฺฌตฺตํ อุปฺปชฺชมานา อุปฺปชฺชนฺติ อหิตาย ทุกฺขาย อผาสุวิหาราย.\csromanspacer{} กตเม ตโย\csromandash{}โลโภ โข มหาราช ปุริสสฺส ธมฺโม อชฺฌตฺตํ อุปฺปชฺชมาโน อุปฺปชฺชติ อหิตาย ทุกฺขาย อผาสุวิหาราย.\csromanspacer{} โทโส โข มหาราช ฯเปฯ.\csromanspacer{} โมโห โข มหาราช ปุริสสฺส ธมฺโม อชฺฌตฺตํ อุปฺปชฺชมาโน อุปฺปชฺชติ อหิตาย ทุกฺขาย อผาสุวิหาราย.\csromanspacer{} อิเม โข มหาราช ตโย ปุริสสฺส ธมฺมา อชฺฌตฺตํ อุปฺปชฺชมานา อุปฺปชฺชนฺติ อหิตาย ทุกฺขาย อผาสุวิหาราย.}

\csromangathameasure{+โลโภ โทโส จ โมโห จ, ปุริสํ ปาปเจตสํ. \\ หิํสนฺติ อตฺตสมฺภูตา, ตจสารํว สมฺผลนฺติ.}
\csromangathagroup{\csromangathastanza{\csromanfolio{13}\csromansymbolfootnote{+}{ขุ 1. 226; ขุ 8. 252 ปิฏฺเฐสุปิ.}โลโภ โทโส จ โมโห จ, ปุริสํ ปาปเจตสํ. \\ หิํสนฺติ อตฺตสมฺภูตา, ตจสารํว สมฺผลนฺติ.}}

\prose{เอวมฺปิ ตตฺราสยาติ ปริสฺสยา.}

\csromanheader{2}{วุตฺตมฺปิ เจตํ ภควตา\csromandash{}}
\csromangathameasure{*ราโค จ โทโส จ อิโตนิทานา, \\ อรติ รติ โลมหํโส อิโตชา. \\ อิโต สมุฏฺฐาย มโนวิตกฺกา,}
\csromangathagroup{\csromangathastanza{\csromansymbolfootnote{*}{สํ 1. 209; ขุ 1. 320; ขุ 8. 252; ขุ 10. 126 ปิฏฺเฐสุปิ.}ราโค จ โทโส จ อิโตนิทานา, \\ อรติ รติ โลมหํโส อิโตชา. \\ อิโต สมุฏฺฐาย มโนวิตกฺกา,}}

\csromanheader{2}{กุมารกา ธงฺกมิโวสฺสชนฺตีติ\footnote{ธงฺกมิโวสฺสชฺชนฺติ (สฺยา)}\csromandash{}}
\prose{เอวมฺปิ ตตฺราสยาติ ปริสฺสยา.\csromanspacer{}\csromanspacer{}\textbf{มทฺทนฺเต นํ ปริสฺสยา}ติ เต ปริสฺสยา ตํ ปุคฺคลํ สหนฺติ ปริสหนฺติ อภิภวนฺติ อชฺโฌตฺถรนฺติ ปริยาทิยนฺติ มทฺทนฺตีติ มทฺทนฺเต นํ ปริสฺสยา.}

\prose{\textbf{ตโต นํ ทุกฺขมเนฺวตี}ติ \textbf{ตโต}ติ \textbf{ตโต} \textbf{ตโต} ปริสฺสยโต ตํ ปุคฺคลํ ทุกฺขํ อเนฺวติ อนุคจฺฉติ อนฺวายิกํ โหติ, ชาติทุกฺขํ อเนฺวติ อนุคจฺฉติ อนฺวายิกํ โหติ, ชราทุกฺขํ อเนฺวติ อนุคจฺฉติ อนฺวายิกํ โหติ, พฺยาธิทุกฺขํ อเนฺวติ อนุคจฺฉติ อนฺวายิกํ โหติ, มรณทุกฺขํ อเนฺวติ อนุคจฺฉติ อนฺวายิกํ โหติ, โสกปริเทวทุกฺขโทมนสฺสุปายาสทุกฺขํ อเนฺวติ อนุคจฺฉติ อนฺวายิกํ โหติ, เนรยิกํ ทุกฺขํ ติรจฺฉานโยนิกํ ทุกฺขํ เปตฺติวิสยิกํ ทุกฺขํ อเนฺวติ อนุคจฺฉติ อนฺวายิกํ โหติ.\csromanspacer{} มานุสิกํ ทุกฺขํ.\csromanspacer{} คมฺโภกฺกนฺติมูลกํ ทุกฺขํ.\csromanspacer{} คมฺเภ ฐิติมูลกํ ทุกฺขํ.\csromanspacer{} คมฺภา วุฏฺฐานมูลกํ ทุกฺขํ.\csromanspacer{} ชาตสฺสูปนิพนฺธกํ ทุกฺขํ.\csromanspacer{} ชาตสฺส ปราเธยฺยกํ ทุกฺขํ.\csromanspacer{} อตฺตูปกฺกมํ ทุกฺขํ.\csromanspacer{} ปรูปกฺกมํ ทุกฺขํ อเนฺวติ อนุคจฺฉติ อนฺวายิกํ โหติ, ทุกฺขทุกฺขํ อเนฺวติ อนุคจฺฉติ อนฺวายิกํ โหติ, สงฺขารทุกฺขํ วิปริณามทุกฺขํ จกฺขุโรโค โสตโรโค ฆานโรโค ชิวฺหาโรโค กายโรโค สีสโรโค กณฺณโรโค มุขโรโค ทนฺตโรโค กาโส สาโส ปินาโส ฑาโห ชโร กุจฺฉิโรโค มุจฺฉา ปกฺขนฺทิกา สูลา วิสูจิกา กุฏฺฐํ คณฺโฑ กิลาโส โสโส อปมาโร ททฺทุ กณฺฑุ กจฺฉุ รขสา วิตจฺฉิกา โลหิตปิตฺตํ มธุเมโห อํสา ปิฬกา ภคนฺทลา ปิตฺตสมุฏฺฐานา อาพาธา}

\prose{\csromanfolio{14}เสมฺหสมุฏฺฐานา อาพาธา วาตสมุฏฺฐานา อาพาธา สนฺนิปาติกา อาพาธา อุตุปริณามชา อาพาธา วิสมปริหารชา อาพาธา โอปกฺกมิกา อาพาธา กมฺมวิปากชา อาพาธา สีตํ อุณฺหํ ชิฆจฺฉา ปิปาสา อุจฺจาโร ปสฺสาโว ฑํสมกสวาตาตปสรีสปสมฺผสฺสทุกฺขํ มาตุมรณํ ทุกฺขํ ปิตุมรณํ ทุกฺขํ ภาตุมรณํ ทุกฺขํ ภคินิมรณํ ทุกฺขํ ปุตฺตมรณํ ทุกฺขํ ธีตุมรณํ ทุกฺขํ ญาติพฺยสนํ ทุกฺขํ โภคพฺยสนํ ทุกฺขํ โรคพฺยสนํ ทุกฺขํ สีลพฺยสนํ ทุกฺขํ ทิฏฺฐิพฺยสนํ ทุกฺขํ อเนฺวติ อนุคจฺฉติ อนฺวายิกํ โหตีติ ตโต นํ ทุกฺขมเนฺวติ.}

\prose{\textbf{นาวํ ภินฺนมิโวทกนฺ}ติ ยถา ภินฺนํ นาวํ ทกเมสิํ\footnote{อุทกทายิโต (สี), อุทกํ อนฺวายิกํ (สฺยา)} ตโต ตโต อุทกํ อเนฺวติ อนุคจฺฉติ อนฺวายิกํ โหติ, ปุรโตปิ อุทกํ อเนฺวติ อนุคจฺฉติ อนฺวายิกํ โหติ, ปจฺฉโตปิ เหฏฺฐโตปิ ปสฺสโตปิ อุทกํ อเนฺวติ อนุคจฺฉติ อนฺวายิกํ โหติ.\csromanspacer{} เอวเมวํ ตโต ตโต ปริสฺสยโต ตํ ปุคฺคลํ ทุกฺขํ อเนฺวติ อนุคจฺฉติ อนฺวายิกํ โหติ, ชาติทุกฺขํ อเนฺวติ อนุคจฺฉติ อนฺวายิกํ โหติ ฯเปฯ ทิฏฺฐิพฺยสนํ ทุกฺขํ อเนฺวติ อนุคจฺฉติ อนฺวายิกํ โหตีติ นาวํ ภินฺนมิโวทกํ.\csromanspacer{} เตนาห ภควา\csromandash{}}

\csromangathameasure{อพลา นํ พลียนฺติ, มทฺทนฺเต นํ ปริสฺสยา. \\ ตโต นํ ทุกฺขมเนฺวติ, นาวํ ภินฺนมิโวทกนฺติ.}
\csromangathagroup{\csromangathastanza{อพลา นํ พลียนฺติ, มทฺทนฺเต นํ ปริสฺสยา. \\ ตโต นํ ทุกฺขมเนฺวติ, นาวํ ภินฺนมิโวทกนฺติ.}}

\proseitem{6}{\csromansymbolfootnote{*}{ขุ 1. 399; ขุ 10. 6 ปิฏฺเฐสุปิ. 2. ผุสิตํ (สี, สฺยา) ขุ 8. 43 ปิฏฺเฐปิ.}\textbf{ตสฺมา ชนฺตุ สทา สโต}, กามานิ ปริวชฺชเย.}

\csromangathameasure{เต ปหาย ตเร โอฆํ, นาวํ สิตฺวาว ปารคู.}
\csromangathagroup{\csromangathastanza{เต ปหาย ตเร โอฆํ, นาวํ สิตฺวาว ปารคู.}}

\prose{\textbf{ตสฺมา ชนฺตุ สทา สโต}ติ \textbf{ตสฺมา}ติ \textbf{ตสฺมา} ตํการณา ตํเหตุ ตปฺปจฺจยา ตํนิทานา เอตํ อาทีนวํ สมฺปสฺสมาโน กาเมสูติ \textbf{ตสฺมา}.\csromanspacer{}\csromanspacer{}\textbf{ชนฺตู}ติ สตฺโต นโร มานโว โปโส ปุคฺคโล ชีโว ชาคุ ชนฺตุ อินฺทคุ มนุโช.\csromanspacer{}\csromanspacer{}\textbf{สทา}ติ สทา สพฺพทา สพฺพกาลํ นิจฺจกาลํ ธุวกาลํ สตตํ สมิตํ อพฺโพกิณฺณํ โปงฺขานุโปงฺขํ อุทกูมิกชาตํ อวีจิ สนฺตติ สหิตํ ผสฺสิตํ2 ปุเรภตฺตํ ปจฺฉาภตฺตํ ปุริมยามํ มชฺฌิมยามํ ปจฺฉิมยามํ กาเฬ ชุเณฺห วสฺเส เหมนฺเต คิเมฺห ปุริเม วโยขนฺเธ มชฺฌิเม วโยขนฺเธ ปจฺฉิเม วโยขนฺเธ.\csromanspacer{}\csromanspacer{}\textbf{สโต}ติ จตูหิ การเณหิ สโต, กาเย กายานุปสฺสนาสติปฏฺฐานํ ภาเวนฺโต สโต, เวทนาสุ. \csromanfolio{15}จิตฺ\textbf{เต}.\csromanspacer{} ธมฺเมสุ ธมฺมานุปสฺสนาสติปฏฺฐานํ ภาเวนฺโต สโต.\csromanspacer{} อปเรหิ จตูหิ การเณหิ สโต ฯเปฯ โส วุจฺจติ สโตติ ตสฺมา ชนฺตุ สทา สโต.}

\prose{\textbf{กามานิ ปริวชฺชเย}ติ \textbf{กามา}ติ อุทฺทานโต เทฺว \textbf{กามา} วตฺถุ\textbf{กามา} จ กิเลส\textbf{กามา} จ ฯเปฯ อิเม วุจฺจนฺติ วตฺถุ\textbf{กามา} ฯเปฯ อิเม วุจฺจนฺติ กิเลส\textbf{กามา}.\csromanspacer{} \textbf{กามานิ ปริวชฺชเย}ติ ทฺวีหิ การเณหิ กาเม ปริวชฺเชยฺย วิกฺขมฺภนโต วา สมุจฺเฉทโต วา.\csromanspacer{} กถํ วิกฺขมฺภนโต กาเม ปริวชฺเชยฺย\csromandash{} * “อฏฺฐิกงฺกลูปมา \textbf{กามา} อปฺปสฺสาทฏฺเฐนา”ติ ปสฺสนฺโต วิกฺขมฺภนโต กาเม ปริวชฺเชยฺย, “มํสเปสูปมา \textbf{กามา} พหุสาธารณฏฺเฐนา”ติ ปสฺสนฺโต วิกฺขมฺภนโต กาเม ปริวชฺเชยฺย, “ติณุกฺกูปมา \textbf{กามา} อนุทหนฏฺเฐนา”ติ ปสฺสนฺโต วิกฺขมฺภนโต กาเม ปริวชฺเชยฺย ฯเปฯ เนวสญฺญานาสญฺญายตนสมาปตฺติํ ภาเวนฺโต วิกฺขมฺภนโต กาเม ปริวชฺเชยฺย.\csromanspacer{} เอวํ วิกฺขมฺภนโต กาเม ปริวชฺเชยฺย ฯเปฯ.\csromanspacer{} เอวํ สมุจฺเฉทโต กาเม ปริวชฺเชยฺยาติ \textbf{กามา}นิ ปริวชฺชเย.}

\prose{\textbf{เต ปหาย ตเร โอฆนฺ}ติ \textbf{เต}ติ วตฺถุกาเม ปริชานิตฺวา กิเลสกาเม ปหาย ปชหิตฺวา วิโนเทตฺวา พฺยนฺติํ กริตฺวา อนภาวํ คมิตฺวา, กามจฺฉนฺทนีวรณํ ปหาย ปชหิตฺวา วิโนเทตฺวา พฺยนฺติํ กริตฺวา อนภาวํ คมิตฺวา, พฺยาปาทนีวรณํ ฯเปฯ ถินมิทฺธนีวรณํ.\csromanspacer{} อุทฺธจฺจกุกฺกุจฺจนีวรณํ.\csromanspacer{} วิจิกิจฺฉานีวรณํ ปหาย ปชหิตฺวา วิโนเทตฺวา พฺยนฺติํ กริตฺวา อนภาวํ คมิตฺวา กาโมฆํ ภโวฆํ ทิฏฺโฐฆํ อวิชฺโชฆํ ตเรยฺย อุตฺตเรยฺย ปตเรยฺย สมติกฺกเมยฺย วีติวตฺ\textbf{เต}ยฺยาติ \textbf{เต} ปหาย ตเร โอฆํ.}

\prose{\textbf{นาวํ สิตฺวาว ปารคู}ติ ยถา ครุกํ นาวํ ภาริกํ อุทกํ สิตฺวา\footnote{สิญฺจิตฺวา (สี, สฺยา)} โอสิญฺจิตฺวา ฉฑฺเฑตฺวา ลหุกาย นาวาย ขิปฺปํ ลหุํ อปฺปกสิเรเนว ปารํ คจฺเฉยฺย.\csromanspacer{} เอวเมวํ วตฺถุกาเม ปริชานิตฺวา กิเลสกาเม ปหาย ปชหิตฺวา วิโนเทตฺวา พฺยนฺติํ กริตฺวา อนภาวํ คมิตฺวา กามจฺฉนฺทนีวรณํ.\csromanspacer{} พฺยาปาทนีวรณํ.\csromanspacer{} ถินมิทฺธนีวรณํ.\csromanspacer{} อุทฺธจฺจกุกฺกุจฺจนีวรณํ.\csromanspacer{} วิจิกิจฺฉานีวรณํ ปหาย ปชหิตฺวา วิโนเทตฺวา พฺยนฺติํ กริตฺวา อนภาวํ คมิตฺวา ขิปฺปํ ลหุํ อปฺปกสิเรเนว ปารํ คจฺเฉยฺย.\csromanspacer{} ปารํ วุจฺจติ อมตํ นิพฺพานํ.\csromanspacer{} โย โส สพฺพสงฺขารสมโถ สพฺพูปธิปฏินิสฺสคฺโค ตณฺหกฺขโย วิราโค นิโรโธ นิพฺพานํ.\csromanspacer{} \textbf{ปารํ คจฺเฉยฺยา}ติ ปารํ อธิคจฺเฉยฺย, ปารํ ผุเสยฺย, ปารํ \csromanfolio{16}สจฺฉิกเรยฺย.\csromanspacer{}\csromanspacer{}\textbf{ปารคู}ติ โยปิ ปารํ คนฺตุกาโม, โสปิ ปารคู, โยปิ ปารํ คจฺฉติ, โสปิ ปารคู, โยปิ ปารํ คโต, โสปิ ปารคู.\csromanspacer{} วุตฺตมฺปิ เหตํ ภควตา * \csromandash{}}

\prose{“ติณฺโณ ปารงฺคโต ถเล ติฏฺฐติ พฺราหฺมโณ”ติ โข ภิกฺขเว อรหโต เอวํ อธิวจนํ.\csromanspacer{} โส อภิญฺญาปารคู ปริญฺญาปารคู ปหานปารคู ภาวนาปารคู สจฺฉิกิริยาปารคู สมาปตฺติปารคู.\csromanspacer{} อภิญฺญาปารคู สพฺพธมฺมานํ, ปริญฺญาปารคู สพฺพทุกฺขานํ, ปหานปารคู สพฺพกิเลสานํ, ภาวนาปารคู จตุนฺนํ อริยมคฺคานํ, สจฺฉิกิริยาปารคู นิโรธสฺส, สมาปตฺติปารคู สพฺพสมาปตฺตีนํ.\csromanspacer{} โส วสิปฺปตฺโต ปารมิปฺปตฺโต อริยสฺมิํ สีลสฺมิํ, วสิปฺปตฺโต ปารมิปฺปตฺโต อริยสฺมิํ สมาธิสฺมิํ, วสิปฺปตฺโต ปารมิปฺปตฺโต อริยาย ปญฺญาย, วสิปฺปตฺโต ปารมิปฺปตฺโต อริยาย วิมุตฺติยา.\csromanspacer{} โส ปารํ คโต ปารปฺปตฺโต อนฺตคโต อนฺตปฺปตฺโต โกฏิคโต โกฏิปฺปตฺโต ปริยนฺตคโต ปริยนฺตปฺปตฺโต โวสานคโต โวสานปฺปตฺโต ตาณคโต ตาณปฺปตฺโต เลณคโต เลณปฺปตฺโต สรณคโต สรณปฺปตฺโต อภยคโต อภยปฺปตฺโต อจฺจุตคโต อจฺจุตปฺปตฺโต อมตคโต อมตปฺปตฺโต นิพฺพานคโต นิพฺพานปฺปตฺโต, โส วุฏฺฐวาโส จิณฺณจรโณ คตทฺโธ คตทิโส คตโกฏิโก ปาลิตพฺรหฺมจริโย อุตฺตมทิฏฺฐิปฺปตฺโต ภาวิตมคฺโค ปหีนกิเลโส ปฏิวิทฺธากุปฺโป สจฺฉิกตนิโรโธ, ทุกฺขํ ตสฺส ปริญฺญาตํ, สมุทโย ปหีโน, มคฺโค ภาวิโต, นิโรโธ สจฺฉิกโต, อภิญฺเญยฺยํ อภิญฺญาตํ, ปริญฺเญยฺยํ ปริญฺญาตํ, ปหาตพฺพํ ปหีนํ, ภาเวตพฺพํ ภาวิตํ, สจฺฉิกาตพฺพํ สจฺฉิกตํ.}

\prose{โส อุกฺขิตฺตปลิโฆ สํกิณฺณปริกฺโข อพฺพุเฬฺหสิโก นิรคฺคโฬ อริโย ปนฺนทฺธโช ปนฺนภาโร วิสญฺญุตฺโต + ปญฺจงฺควิปฺปหีโน ฉฬงฺคสมนฺนา คโต เอการกฺโข จตุราปสฺเสโน ปณุนฺนปจฺเจกสจฺโจ สมวยสฏฺเฐสโน อนาวิลสงฺกปฺโป ปสฺสทฺธกายสงฺขาโร สุวิมุตฺตจิตฺโต สุวิมุตฺตปญฺโญ เกวลี วุสิตวา อุตฺตมปุริโส ปรมปุริโส ปรมปตฺติปฺปตฺโต, โส เนวาจินติ\footnote{เนว อาจินาติ (สี, สฺยา)} นาปจินติ อปจินิตฺวา ฐิโต, เนว ปชหติ น อุปาทิยติ ปชหิตฺวา ฐิโต, เนว สํสิพฺพติ\footnote{เนว สิเนติ (สี), เนว วิสีเนติ (สฺยา)} น อุสฺสิเนติ}

\vspace{\tipitakaparskip}
\csromancenter{คุหฏฺฐกสุตฺตนิทฺเทส วิสินิตฺวา ฐิโต, เนว วิธูเปติ น สนฺธูเปติ วิธูเปตฺวา ฐิโต, อเสกฺเขน สีลกฺขนฺเธน สมนฺนาคตตฺตา ฐิโต, อเสกฺเขน สมาธิกฺขนฺเธน.\csromanspacer{} อเสกฺเขน ปญฺญกฺขนฺเธน.\csromanspacer{} อเสกฺเขน วิมุตฺติกฺขนฺเธน.\csromanspacer{} อเสกฺเขน วิมุตฺติญาณทสฺสนกฺขนฺเธน สมนฺนาคตตฺตา ฐิโต, สจฺจํ สมฺปฏิปาทิยิตฺวา ฐิโต, เอชํ สมติกฺกมิตฺวา ฐิโต, กิเลสคฺคิํ ปริยาทิยิตฺวา ฐิโต, อปริคมนตาย ฐิโต, กฏํ สมาทาย ฐิโต, มุตฺติปฏิเสวนตาย ฐิโต, เมตฺตาย ปาริสุทฺธิยา ฐิโต, กรุณาย.\csromanspacer{} มุทิตาย.\csromanspacer{} อุเปกฺขาย ปาริสุทฺธิยา ฐิโต, อจฺจนฺตปาริสุทฺธิยา ฐิโต, อตมฺมยตาย\footnote{อกมฺมยตาย (ก), อกมฺมญฺญตาย (สฺยา) อํ 1. 148 ปิฏฺเฐปิ ปสฺสิตพฺพํ.} ปาริสุทฺธิยา ฐิโต, วิมุตฺตตฺตา ฐิโต, สนฺตุสฺสิตตฺตา ฐิโต, ขนฺธปริยนฺเต ฐิโต, ธาตุปริยนฺเต ฐิโต, อายตนปริยนฺเต ฐิโต, คติปริยนฺเต ฐิโต, อุปปตฺติปริยนฺเต ฐิโต, ปฏิสนฺธิปริยนฺเต ฐิโต, (ภวปริยนฺเต ฐิโต, สํสารปริยนฺเต ฐิโต, วฏฺฏปริยนฺเต ฐิโต, อนฺติเม ภเว ฐิโต,)\footnote{( ) นตฺถิ สีหฬ-โปตฺถเก.} อนฺติเม สมุสฺสเย ฐิโต อนฺติมเทหธโร อรหา.}
\vspace{0.5\baselineskip}
\csromangathameasure{ตสฺสายํ ปจฺฉิมโก ภโว, จริโมยํ สมุสฺสโย. \\ ชาติมรณสํสาโร, นตฺถิ ตสฺส ปุนพฺภโวติ.}
\csromangathagroup{\csromangathastanza{\csromanfolio{17}ตสฺสายํ ปจฺฉิมโก ภโว, จริโมยํ สมุสฺสโย. \\ ชาติมรณสํสาโร, นตฺถิ ตสฺส ปุนพฺภโวติ.}}

\csromanheader{2}{นาวํ สิตฺวาว ปารคูติ. เตนาห ภควา\csromandash{}}
\csromangathameasure{ตสฺมา ชนฺตุ สทา สโต, กามานิ ปริวชฺชเย. \\ เต ปหาย ตเร โอฆํ, นาวํ สิตฺวาว ปารคูติ. กามสุตฺตนิทฺเทโส ปฐโม.}
\csromangathagroup{\csromangathastanza{ตสฺมา ชนฺตุ สทา สโต, กามานิ ปริวชฺชเย.}\csromangathastanza{เต ปหาย ตเร โอฆํ, นาวํ สิตฺวาว ปารคูติ.\csromanspacer{} กามสุตฺตนิทฺเทโส ปฐโม.}}

\csromanheader{2}{2. คุหฏฺฐกสุตฺตนิทฺเทส}
\csromanheader{2}{อถ คุหฏฺฐกสุตฺตนิทฺเทสํ วกฺขติ\csromandash{}}
\proseitem{7}{\csromansymbolfootnote{*}{ขุ 1. 399 ปิฏฺเฐปิ.}\textbf{สตฺโต คุหายํ พหุนาภิฉนฺโน},}

\prose{\textbf{ติฏฺฐํ นโร โมหนสฺมิํ ปคาโฬฺห.}}

\csromangathameasure{\textbf{ทูเร วิเวกา หิ ตถาวิโธ โส,} \\ \textbf{กามา หิ โลเก น หิ สุปฺปหายา.}}
\csromangathagroup{\csromangathastanza{\textbf{ทูเร วิเวกา หิ ตถาวิโธ โส,} \\ \textbf{กามา หิ โลเก น หิ สุปฺปหายา.}}}

\csromanmatikaanchor{mtk.3}
\csromantocmarkref{subsection}{2. คุหฏฺฐกสุตฺตนิทฺเทส}{mtk.3}
\bookmark[dest={mtk.3},level=2]{2. คุหฏฺฐกสุตฺตนิทฺเทส}
\prose{\textbf{สตฺโต คุหายํ พหุนาภิฉนฺโน}ติ สตฺโตติ หิ โข วุตฺตํ, อปิ จ คุหา ตาว วตฺตพฺพา.\csromanspacer{} \textbf{คุหา} วุจฺจติ กาโย, กาโยติ วา คุหาติ วา \csromanfolio{18}เทโหติ วา สนฺเทโหติ วา นาวาติ วา รโถติ วา ธโชติ วา วมฺมิโกติ วา นครนฺติ วา นีฬนฺติ\footnote{นิฑฺฑนฺติ (สี, สฺยา)} วา กุฏีติ วา คณฺโฑติ วา กุมฺโภติ วา นาโคติ วา กายสฺเสตํ อธิวจนํ.\csromanspacer{} \textbf{สตฺโต คุหายนฺ}ติ คุหายํ สตฺโต วิสตฺโต อาสตฺโต ลคฺโค ลคฺคิโต ปลิพุทฺโธ, ยถา ภิตฺติขิเล วา นาคทนฺเต วา คณฺฑํ สตฺตํ วิสตฺตํ อาสตฺตํ ลคฺคํ ลคฺคิตํ ปลิพุทฺธํ.\csromanspacer{} เอวเมวํ คุหายํ สตฺโต วิสตฺโต อาสตฺโต ลคฺโค ลคฺคิโต ปลิพุทฺโธ.\csromanspacer{} วุตฺตํ เหตํ ภควตา\csromandash{}}

\prose{\csromansymbolfootnote{*}{สํ 2. 155 ปิฏฺเฐ.}รูเป โข ราธ โย ฉนฺโท โย ราโค ยา นนฺที ยา ตณฺหา เย อุปยูปาทานา เจตโส อธิฏฺฐานาภินิเวสานุสยา, ตตฺร สตฺโต ตตฺร วิสตฺโต, ตสฺมา “สตฺโต”ติ วุจฺจติ.\csromanspacer{}\csromanspacer{}เวทนาย โข ราธ ฯเปฯ.\csromanspacer{} สญฺญาย โข ราธ.\csromanspacer{} สงฺขาเรสุ โข ราธ.\csromanspacer{} วิญฺญาเณ โข ราธ โย ฉนฺโท โย ราโค ยา นนฺที ยา ตณฺหา เย อุปยูปาทานา เจตโส อธิฏฺฐานาภินิเวสานุสยา.\csromanspacer{} ตตฺร สตฺโต ตตฺร วิสตฺโต, ตสฺมา “สตฺโต”ติ วุจฺจติ.\csromanspacer{} สตฺโตติ ลคฺคนาธิวจนนฺติ สตฺโต คุหายํ.\csromanspacer{}\csromanspacer{}\textbf{พหูนาภิฉนฺโน}ติ พหุเกหิ กิเลเสหิ ฉนฺโน ราเคน ฉนฺโน โทเสน ฉนฺโน โมเหน ฉนฺโน โกเธน ฉนฺโน อุปนาเหน ฉนฺโน มกฺเขน ฉนฺโน ปฬาเสน ฉนฺโน อิสฺสาย ฉนฺโน มจฺฉริเยน ฉนฺโน มายาย ฉนฺโน สาเฐยฺเยน ฉนฺโน ถมฺเภน ฉนฺโน สารมฺเภน ฉนฺโน มาเนน ฉนฺโน อติมาเนน ฉนฺโน มเทน ฉนฺโน ปมาเทน ฉนฺโน สพฺพกิเลเสหิ.\csromanspacer{} สพฺพทุจฺจริเตหิ.\csromanspacer{} สพฺพทรเถหิ.\csromanspacer{} สพฺพปริฬาเหหิ.\csromanspacer{} สพฺพสนฺตาเปหิ.\csromanspacer{} สพฺพากุสลาภิสงฺขาเรหิ ฉนฺโน วิฉนฺโน อุจฺฉนฺโน อาวุโต นิวิโต โอวุโต\footnote{โอผุโต (สฺยา)} ปิหิโต ปฏิจฺฉนฺโน ปฏิกุชฺชิโตติ สตฺโต คุหายํ พหุนาภิฉนฺโน.}

\prose{\textbf{ติฏฺฐํ นโร โมหนสฺมิํ ปคาโฬฺห}ติ ติฏฺฐนฺโต นโร.\csromanspacer{} รตฺโต ราควเสน ติฏฺฐติ, ทุฏฺโฐ โทสวเสน ติฏฺฐติ, มูโฬฺห โมหวเสน ติฏฺฐติ, วินิพทฺโธ มานวเสน ติฏฺฐติ, ปรามฏฺโฐ ทิฏฺฐิวเสน ติฏฺฐติ, วิกฺเขปคโต อุทฺธจฺจวเสน ติฏฺฐติ, อนิฏฺฐงฺคโต วิจิกิจฺฉาวเสน ติฏฺฐติ, ถามคโต อนุสยวเสน ติฏฺฐติ.\csromanspacer{} เอวมฺปิ ติฏฺฐํ นโร.}

\prose{วุตฺตํ เหตํ ภควตา\csromandash{} + สนฺติ ภิกฺขเว จกฺขุวิญฺเญยฺยา รูปา อิฏฺฐา กนฺตา มนาปา ปิยรูปา กามูปสํหิตา รชนียา, ตญฺเจ ภิกฺขุ อภินนฺทติ}

\vspace{\tipitakaparskip}
\csromancenter{คุหฏฺฐกสุตฺตนิทฺเทส อภิวทติ อชฺโฌสาย ติฏฺฐติ.\csromanspacer{}\csromanspacer{}สนฺติ ภิกฺขเว โสตวิญฺเญยฺยา สทฺทา.\csromanspacer{} ฆานวิญฺเญยฺยา คนฺธา.\csromanspacer{} ชิวฺหาวิญฺเญยฺยา รสา.\csromanspacer{} กายวิญฺเญยฺยา โผฏฺฐพฺพา.\csromanspacer{} มโนวิญฺเญยฺยา ธมฺมา อิฏฺฐา กนฺตา มนาปา ปิยรูปา กามูปสํหิตา รชนียา, ตญฺเจ ภิกฺขุ อภินนฺทติ อภิวทติ อชฺโฌสาย ติฏฺฐตีติ.\csromanspacer{} เอวมฺปิ ติฏฺฐํ นโร.}
\prose{\csromanfolio{19}วุตฺตํ เหตํ ภควตา\csromandash{} * รูปูปยํ วา ภิกฺขเว วิญฺญาณํ ติฏฺฐมานํ ติฏฺฐติ, รูปารมฺมณํ รูปปติฏฺฐํ นนฺทูปเสจนํ วุทฺธิํ วิรูฬฺหิํ เวปุลฺลํ อาปชฺชติ.\csromanspacer{}\csromanspacer{}เวทนูปยํ วา ภิกฺขเว.\csromanspacer{} สญฺญูปยํ.\csromanspacer{} สงฺขารูปยํ วา ภิกฺขเว วิญฺญาณํ ติฏฺฐมานํ ติฏฺฐติ, สงฺขารารมฺมณํ สงฺขารปติฏฺฐํ นนฺทูปเสจนํ วุทฺธิํ วิรูฬฺหิํ เวปุลฺลํ อาปชฺชตีติ.\csromanspacer{} เอวมฺปิ ติฏฺฐํ นโร.}

\prose{วุตฺตมฺปิ เหตํ ภควตา\csromandash{} + กพฬีกาเร เจ ภิกฺขเว อาหาเร อตฺถิ ราโค อตฺถิ นนฺที อตฺถิ ตณฺหา, ปติฏฺฐิตํ ตตฺถ วิญฺญาณํ วิรูฬฺหํ.\csromanspacer{} ยตฺถ ปติฏฺฐิตํ วิญฺญาณํ วิรูฬฺหํ, อตฺถิ ตตฺถ นามรูปสฺสาวกฺกนฺติ.\csromanspacer{} ยตฺถ อตฺถิ นามรูปสฺสาวกฺกนฺติ, อตฺถิ ตตฺถ สงฺขารานํ วุทฺธิ.\csromanspacer{} ยตฺถ อตฺถิ สงฺขารานํ วุทฺธิ, อตฺถิ ตตฺถ อายติํ ปุนพฺภวาภินิพฺพตฺติ.\csromanspacer{} ยตฺถ อตฺติ อายติํ ปุนพฺภวาภินิพฺพตฺติ, อตฺถิ ตตฺถ อายติํ ชาติชรามรณํ.\csromanspacer{} ยตฺถ อตฺถิ อายติํ ชาติชรามรณํ, สโสกํ ตํ ภิกฺขเว สรชํ ส-อุปายาสนฺติ วทามีติ.\csromanspacer{} เอวมฺปิ ติฏฺฐํ นโร.}

\prose{ผสฺเส เจ ภิกฺขเว อาหาเร ฯเปฯ.\csromanspacer{} มโนสญฺเจตนาย เจ ภิกฺขเว อาหาเร.\csromanspacer{} วิญฺญาเณ เจ ภิกฺขเว อาหาเร อตฺถิ ราโค อตฺถิ นนฺที อตฺถิ ตณฺหา, ปติฏฺฐิตํ ตตฺถ วิญฺญาณํ วิรูฬฺหํ.\csromanspacer{} ยตฺถ ปติฏฺฐิตํ วิญฺญาณํ วิรูฬฺหํ, อตฺถิ ตตฺถ นามรูปสฺสาวกฺกนฺติ.\csromanspacer{} ยตฺถ อตฺถิ นามรูปสฺสาวกฺกนฺติ อตฺถิ ตตฺถ สงฺขารานํ วุทฺธิ.\csromanspacer{} ยตฺถ อตฺถิ สงฺขารานํ วุทฺธิ, อตฺถิ ตตฺถ อายติํ ปุนพฺภวาภินิพฺพตฺติ.\csromanspacer{} ยตฺถ อตฺถิ อายติํ ปุนพฺภวาภินิพฺพตฺติ, อตฺถิ ตตฺถ อายติํ ชาติชรามรณํ.\csromanspacer{} ยตฺถ อตฺถิ อายติํ ชาติชรามรณํ, สโสกํ ตํ ภิกฺขเว สรชํ ส-อุปายาสนฺติ วทามีติ.\csromanspacer{} เอวมฺปิ ติฏฺฐํ นโร.}

\prose{\textbf{โมหนสฺมิํ ปคาโฬฺห}ติ โมหนา วุจฺจนฺติ ปญฺจ กามคุณา.\csromanspacer{} จกฺขุวิญฺเญยฺยา รูปา อิฏฺฐา กนฺตา มนาปา ปิยรูปา กามูปสํหิตา รชนียา.\csromanspacer{} โสตวิญฺเญยฺยา สทฺทา.\csromanspacer{} ฆานวิญฺเญยฺยา คนฺธา.\csromanspacer{} ชิวฺหาวิญฺเญยฺยา รสา.\csromanspacer{} กายวิญฺเญยฺยา โผฏฺฐพฺพา อิฏฺฐา กนฺตา มนาปา ปิยรูปา กามูปสํหิตา รชนียา.\csromanspacer{} กิํ การณา โมหนา วุจฺจนฺติ ปญฺจ กามคุณา.\csromanspacer{} เยภุยฺเยน \csromanfolio{20}เทวมนุสฺสา ปญฺจสุ กามคุเณสุ มุยฺหนฺติ สมฺมุยฺหนฺติ สมฺปมุยฺหนฺติ, มูฬฺหา สมฺมูฬฺหา สมฺปมูฬฺหา อวิชฺชาย อนฺธีกตา อาวุตา นิวุตา โอวุตา ปิหิตา ปฏิจฺฉนฺนา ปฏิกุชฺชิตา, ตํการณา โมหนา วุจฺจนฺติ ปญฺจ กามคุณา.\csromanspacer{} \textbf{โมหนสฺมิํ ปคาโฬฺห}ติ โมหนสฺมิํ ปคาโฬฺห โอคาโฬฺห อชฺโฌคาโฬฺห นิมุคฺโคติ ติฏฺฐํ นโร โมหนสฺมิํ ปคาโฬฺห.}

\prose{\textbf{ทูเร วิเวกา หิ ตถาวิโธ โส}ติ \textbf{วิเวกา}ติ ตโย \textbf{วิเวกา} กายวิเวโก จิตฺตวิเวโก อุปธิวิเวโก.\csromanspacer{} กตโม กายวิเวโก\csromandash{}อิธ ภิกฺขุ วิวิตฺตํ เสนาสนํ ภชติ อรญฺญํ รุกฺขมูลํ ปพฺพตํ กนฺทรํ คิริคุหํ สุสานํ วนปตฺถํ อพฺโภกาสํ ปลาลปุญฺชํ.\csromanspacer{} กาเยน วิวิตฺโต วิหรติ.\csromanspacer{} โส เอโก คจฺฉติ, เอโก ติฏฺฐติ, เอโก นิสีทติ, เอโก เสยฺยํ กปฺเปติ, เอโก คามํ ปิณฺฑาย ปวิสติ, เอโก ปฏิกฺกมติ, เอโก รโห นิสีทติ, เอโก จงฺกมํ อธิฏฺฐาติ, เอโก จรติ วิหรติ อิริยติ วตฺตติ ปาเลติ ยเปติ ยาเปติ.\csromanspacer{} อยํ กายวิเวโก.}

\prose{กตโม จิตฺตวิเวโก\csromandash{}ปฐมํ ฌานํ สมาปนฺนสฺส นีวรเณหิ จิตฺตํ วิวิตฺตํ โหติ.\csromanspacer{} ทุติยํ ฌานํ สมาปนฺนสฺส วิตกฺกวิจาเรหิ จิตฺตํ วิวิตฺตํ โหติ.\csromanspacer{} ตติยํ ฌานํ สมาปนฺนสฺส ปีติยา จิตฺตํ วิวิตฺตํ โหติ.\csromanspacer{} จตุตฺถํ ฌานํ สมาปนฺนสฺส สุขทุกฺเขหิ จิตฺตํ วิวิตฺตํ โหติ.\csromanspacer{} อากาสานญฺจายตนํ สมาปนฺนสฺส รูปสญฺญาย ปฏิฆสญฺญาย นานตฺตสญฺญาย จิตฺตํ วิวิตฺตํ โหติ.\csromanspacer{} วิญฺญาณญฺจายตนํ สมาปนฺนสฺส อากาสานญฺจายตนสญฺญาย จิตฺตํ วิวิตฺตํ โหติ.\csromanspacer{} อากิญฺจญฺญายตนํ สมาปนฺนสฺส วิญฺญาณญฺจายตนสญฺญาย จิตฺตํ วิวิตฺตํ โหติ.\csromanspacer{} เนวสญฺญานาสญฺญายตนํ สมาปนฺนสฺส อากิญฺจญฺญายตนสญฺญาย จิตฺตํ วิวิตฺตํ โหติ.\csromanspacer{} โสตาปนฺนสฺส สกฺกายทิฏฺฐิยา วิจิกิจฺฉาย สีลพฺพตปรามาสา ทิฏฺฐานุสยา วิจิกิจฺฉานุสยา, ตเทกฏฺเฐหิ จ กิเลเสหิ จิตฺตํ วิวิตฺตํ โหติ.\csromanspacer{} สกทาคามิสฺส โอฬาริกา กามราคสญฺโญชนา ปฏิฆสญฺโญชนา โอฬาริกา กามราคานุสยา ปฏิฆานุสยา, ตเทกฏฺเฐหิ จ กิเลเสหิ จิตฺตํ วิวิตฺตํ โหติ.\csromanspacer{} อนาคามิสฺส อณุสหคตา กามราคสญฺโญชนา ปฏิฆสญฺโญชนา อณุสหคตา กามราคานุสยา ปฏิฆานุสยา, ตเทกฏฺเฐหิ จ กิเลเสหิ จิตฺตํ วิวิตฺตํ โหติ.\csromanspacer{} อรหโต รูปารูปราคา มานา อุทฺธจฺจา อวิชฺชาย มานานุสยา ภวราคานุสยา}

\vspace{\tipitakaparskip}
\csromancenter{คุหฏฺฐกสุตฺตนิทฺเทส อวิชฺชานุสยา, ตเทกฏฺเฐหิ จ กิเลเสหิ พหิทฺธา จ สพฺพนิมิตฺเตหิ จิตฺตํ วิวิตฺตํ โหติ.\csromanspacer{} อยํ จิตฺตวิเวโก.}
\prose{\csromanfolio{21}กตโม \textbf{อุปธิ}วิเวโก\csromandash{}\textbf{อุปธิ} วุจฺจนฺติ กิเลสา จ ขนฺธา จ อภิสงฺขารา จ.\csromanspacer{} อุปธิวิเวโก วุจฺจติ อมตํ นิพฺพานํ.\csromanspacer{} โย สพฺพสงฺขารสมโถ สพฺพูปธิปฏินิสฺสคฺโค ตณฺหกฺขโย วิราโค นิโรโธ นิพฺพานํ, อยํ \textbf{อุปธิ}วิเวโก.\csromanspacer{} กายวิเวโก จ วิเวกฏฺฐกายานํ\footnote{วูปกฏฺฐกายานํ (สฺยา) อุปริ 109 ปิฏฺเฐปิ.} เนกฺขมฺมาภิรตานํ, จิตฺตวิเวโก จ ปริสุทฺธจิตฺตานํ ปรมโวทานปฺปตฺตานํ, \textbf{อุปธิ}วิเวโก จ นิรูปธีนํ ปุคฺคลานํ วิสงฺขารคตานํ.}

\prose{\textbf{ทูเร วิเวกาหี}ติ โย โส เอวํ คุหายํ สตฺโต เอวํ พหุเกหิ กิเลเสหิ ฉนฺโน เอวํ โมหนสฺมิํ ปคาโฬฺห, โส กายวิเวกาปิ ทูเร, จิตฺตวิเวกาปิ ทูเร, \textbf{อุปธิ}วิเวกาปิ ทูเร วิทูเร สุวิทูเร น สนฺติเก น สามนฺตา อนาสนฺเน วิเวกฏฺเฐ\footnote{ว ว กฏฺเฐ (สี), อนุปกฏฺเฐ (สฺยา)}.\csromanspacer{} \textbf{ตถาวิโธ}ติ ตาทิโส ตสฺสณฺฐิโต ตปฺปกาโร ตปฺปฏิภาโค โย โส โมหนสฺมิํ ปคาโฬฺหติ ทูเร วิเวกาหิ ตถาวิโธ โส.}

\prose{\textbf{กามา หิ โลเก น หิ สุปฺปหายา}ติ * \textbf{กามา}ติ อุทฺทานโต เทฺว \textbf{กามา} วตฺถุ\textbf{กามา} จ กิเลส\textbf{กามา} จ.\csromanspacer{} กตเม \textbf{วตฺถุกามา\csromandash{}}มนาปิกา รูปา มนาปิกา สทฺทา มนาปิกา คนฺธา มนาปิกา รสา มนาปิกา โผฏฺฐพฺพา อตฺถรณา ปาวุรณา ทาสิทาสา อเชฬกา กุกฺกุฏสูกรา หตฺถิควาสฺสวฬวา เขตฺตํ วตฺถุํ หิรญฺญํ สุวณฺณํ คามนิคมราชธานิโย รฏฺฐญฺจ ชนปโท จ โกโส จ โกฏฺฐาคารญฺจ, ยํ กิญฺจิ รชนียํ วตฺถุ วตฺถุ\textbf{กามา}.\csromanspacer{} อปิ จ อตีตา \textbf{กามา} อนาคตา \textbf{กามา} ปจฺจุปฺปนฺนา \textbf{กามา}, อชฺฌตฺตา \textbf{กามา} พหิทฺธา \textbf{กามา} อชฺฌตฺตพหิทฺธา \textbf{กามา}, หีนา \textbf{กามา} มชฺฌิมา \textbf{กามา} ปณีตา \textbf{กามา}, อาปายิกา \textbf{กามา} มานุสิกา \textbf{กามา} ทิพฺพา \textbf{กามา}, ปจฺจุปฏฺฐิตา \textbf{กามา} นิมฺมิตา \textbf{กามา} อนิมฺมิตา \textbf{กามา} ปรนิมฺมิตา \textbf{กามา} ปริคฺคหิตา \textbf{กามา} อปริคฺคหิตา \textbf{กามา} มมายิตา \textbf{กามา} อมมายิตา \textbf{กามา}, สพฺเพปิ \textbf{กามา}วจรา ธมฺมา สพฺเพปิ รูปาวจรา ธมฺมา สพฺเพปิ อรูปาวจรา ธมฺมา ตณฺหาวตฺถุกา ตณฺหารมฺมณา กามนียฏฺเฐน รชนียฏฺเฐน มทนียฏฺเฐน \textbf{กามา}.\csromanspacer{} อิเม วุจฺจนฺติ วตฺถุ\textbf{กามา}.}

\prose{\csromanfolio{22}กตเม กิเลสกามา\csromandash{}ฉนฺโท กาโม ราโค กาโม ฉนฺท ราโค กาโม, สงฺกปฺโป กาโม ราโค กาโม สงฺกปฺปราโค กาโม, โย กาเมสุ กามจฺฉนฺโท กามราโค กามนนฺที กามตณฺหา กามสฺเนโห กามปริฬาโห กามมุจฺฉา กามชฺโฌสานํ กาโมโฆ กามโยโค กามุปาทานํ กามจฺฉนฺทนีวรณํ.}

\csromangathameasure{อทฺทสํ กาม เต มูลํ, สงฺกปฺปา กาม ชายสิ.}
\csromangathagroup{\csromangathastanza{อทฺทสํ กาม เต มูลํ, สงฺกปฺปา กาม ชายสิ.}}

\prose{น ตํ สงฺกปฺปยิสฺสามิ, เอวํ กาม น เหหิสีติ\csromandash{}}

\prose{อิเม วุจฺจนฺติ กิเลสกามา.\csromanspacer{}\csromanspacer{}\textbf{โลเก}ติ อปายโลเก มนุสฺสโลเก เทวโลเก, ขนฺธโลเก ธาตุโลเก อายตนโลเก.\csromanspacer{}\csromanspacer{}\textbf{กามา หิ โลเก น หิ สุปฺปหายา}ติ กามา หิ โลเก ทุปฺปหายา ทุจฺจชฺชา ทุปฺปริจฺจชฺชา ทุนฺนิมฺมทยา ทุนฺนิเวฐยา ทุพฺพินิเวฐยา ทุตฺตรา ทุปฺปตรา ทุสฺสมติกฺกมา ทุพฺพินิวตฺตาติ กามา หิ โลเก น หิ สุปฺปหายา.\csromanspacer{} เตนาห ภควา\csromandash{}}

\csromangathameasure{สตฺโต คุหายํ พหุนาภิฉนฺโน, \\ ติฏฺฐํ นโร โมหนสฺมิํ ปคาโฬฺห. \\ ทูเร วิเวกา หิ ตถาวิโธ โส, \\ \textbf{กามา หิ โลเก น หิ สุปฺปหายา}ติ.}
\csromangathagroup{\csromangathastanza{สตฺโต คุหายํ พหุนาภิฉนฺโน, \\ ติฏฺฐํ นโร โมหนสฺมิํ ปคาโฬฺห. \\ ทูเร วิเวกา หิ ตถาวิโธ โส, \\ \textbf{กามา หิ โลเก น หิ สุปฺปหายา}ติ.}}

\proseitem{8}{\csromansymbolfootnote{*}{ขุ 1. 400 ปิฏฺเฐปิ.}อิจฺฉานิทานา ภวสาตพทฺทา,}

\prose{\textbf{เต ทุปฺปมุญฺจา น หิ อญฺญโมกฺขา.}}

\csromangathameasure{\textbf{ปจฺฉา ปุเร วาปิ อเปกฺขมานา,} \\ \textbf{อิเม ว กาเม ปุริเม ว ชปฺปํ.}}
\csromangathagroup{\csromangathastanza{\textbf{ปจฺฉา ปุเร วาปิ อเปกฺขมานา,} \\ \textbf{อิเม ว กาเม ปุริเม ว ชปฺปํ.}}}

\prose{\textbf{อิจฺฉานิทานา ภวสาตพทฺธา}ติ อิจฺฉา วุจฺจติ ตณฺหา. + โย ราโค สาราโค อนุนโย อนุโรโธ นนฺที นนฺทิราโค จิตฺตสฺส สาราโค อิจฺฉา มุจฺฉา อชฺโฌสานํ เคโธ ปลิเคโธ สงฺโค ปงฺโก เอชา มายา ชนิกา สญฺชนนี สิพฺพินี ชาลินี สริตา วิสตฺติกา สุตฺตํ วิสฏา อายูหินี ทุติยา ปณิธี ภวเนตฺติ วนํ วนโถ สนฺธโว สฺเนโห อเปกฺขา ปฏิพนฺธุ อาสา อาสีสนา อาสีสิตตฺตํ รูปาสา สทฺทาสา คนฺธาสา รสาสา โผฏฺฐพฺพาสา ลาภาสา ธนาสา ปุตฺตาสา}

\vspace{\tipitakaparskip}
\csromancenter{คุหฏฺฐกสุตฺตนิทฺเทส ชีวิตาสา ชปฺปา ปชปฺปา อภิชปฺปา ชปฺปนา ชปฺปิตตฺตํ โลลุปฺปํ โลลุปฺปายนา โลลุปฺปายิตตฺตํ ปุจฺฉญฺฉิกตา สาธุกมฺยตา อธมฺมราโค วิสมโลโภ นิกนฺติ นิกามนา ปตฺถนา ปิหนา สมฺปตฺถนา, กามตณฺหา ภวตณฺหา วิภวตณฺหา, รูปตณฺหา อรูปตณฺหา นิโรธตณฺหา, รูปตณฺหา สทฺทตณฺหา คนฺธตณฺหา รสตณฺหา โผฏฺฐพฺพตณฺหา ธมฺมตณฺหา,โอโฆ โยโค คนฺโถ อุปาทานํ อาวรณํ นีวรณํ ฉทนํ พนฺธนํ อุปกฺกิเลโส อนุสโย ปริยุฏฺฐานํ ลตา เววิจฺฉํ ทุกฺขมูลํ ทุกฺขนิทานํ ทุกฺขปฺปภโว มารปาโส มารพฬิสํ มารวิสโย ตณฺหานที ตณฺหาชาลํ ตณฺหาคทฺทุลํ ตณฺหาสมุทฺโท อภิชฺฌา โลโภ อกุสลมูลํ.\csromanspacer{}\csromanspacer{}\textbf{อิจฺฉานิทานา}ติ อิจฺฉานิทานกา อิจฺฉาเหตุกา อิจฺฉาปจฺจยา อิจฺฉาการณา อิจฺฉาปภวาติ อิจฺฉานิทานา.}
\prose{\csromanfolio{23}\textbf{ภวสาตพทฺธา}ติ เอกํ ภวสาตํ\csromandash{}สุขา เวทนา.\csromanspacer{} เทฺว ภวสาตานิ\csromandash{}สุขา จ เวทนา อิฏฺฐญฺจ วตฺถุ.\csromanspacer{} ตีณิ ภวสาตานิ\csromandash{} โยพฺพญฺญํ อาโรคฺยํ ชีวิตํ.\csromanspacer{} จตฺตาริ ภวสาตานิ\csromandash{}ลาโภ ยโส ปสํสา สุขํ.\csromanspacer{} ปญฺจ ภวสาตานิ\csromandash{}มนาปิกา รูปา มนาปิกา สทฺทา มนาปิกา คนฺธา มนาปิกา รสา มนาปิกา โผฏฺฐพฺพา.\csromanspacer{} ฉ ภวสาตานิ\csromandash{} จกฺขุสมฺปทา โสตสมฺปทา ฆานสมฺปทา ชิวฺหาสมฺปทา กายสมฺปทา มโนสมฺปทา.\csromanspacer{} \textbf{ภวสาตพทฺธา}, สุขาย เวทนาย สาตพทฺธา, อิฏฺฐสฺมิํ วตฺถุสฺมิํ พทฺธา, โยพฺพญฺเญพทฺธา, อาโรเคฺย พทฺธา, ชีวิเต พทฺธา, ลาเภ พทฺธา, ยเส พทฺธา, ปสํสายํ พทฺธา, สุเข พทฺธา, มนาปิเกสุ รูเปสุ พทฺธา, สทฺเทสุ.\csromanspacer{} คนฺเธสุ.\csromanspacer{} รเสสุ.\csromanspacer{} มนาปิเกสุ โผฏฺฐพฺเพสุ พทฺธา, จกฺขุสมฺปทาย พทฺธา, โสต.\csromanspacer{} ฆาน.\csromanspacer{} ชิวฺหา.\csromanspacer{} กาย.\csromanspacer{} มโนสมฺปทาย พทฺธา, วิพทฺธา อาพทฺธา ลคฺคา ลคฺคิตา ปลิพทฺธาติ อิจฺฉานิทานา ภวสาตพทฺธา.}

\prose{\textbf{เต ทุปฺปมุญฺจา น หิ อญฺญโมกฺขา}ติ เต วา ภวสาตวตฺถู ทุปฺปมุญฺจา, สตฺตา วา เอตฺโต ทุมฺโมจยา.\csromanspacer{} กถํ เต ภวสาตวตฺถู ทุปฺปมุญฺจา\csromandash{}สุขา เวทนา ทุปฺปมุญฺจา, อิฏฺฐํ วตฺถุ ทุปฺปมุญฺจํ, โยพฺพญฺญํ ทุปฺปมุญฺจํ, อาโรคฺยํ ทุปฺปมุญฺจํ, ชีวิตํ ทุปฺปมุญฺจํ, ลาโภ ทุปฺปมุญฺโจ, ยโส ทุปฺปมุญฺโจ, ปสํสา ทุปฺปมุญฺจา, สุขํ ทุปฺปมุญฺจํ, มนาปิกา รูปา ทุปฺปมุญฺจา, มนาปิกา สทฺทา.\csromanspacer{} คนฺธา.\csromanspacer{} รสา.\csromanspacer{} โผฏฺฐพฺพา ทุปฺปมุญฺจา, จกฺขุสมฺปทา ทุปฺปมุญฺจา, โสต.\csromanspacer{} ฆาน.\csromanspacer{} ชิวฺหา.\csromanspacer{} กาย.\csromanspacer{} มโนสมฺปทา ทุปฺปมุญฺจา ทุมฺโมจยา ทุปฺปโมจยา ทุนฺนิเวฐยา \csromanfolio{24}ทุพฺพินิเวฐยา ทุตฺตรา ทุปฺปตรา ทุสฺสมติกฺกมา ทุพฺพินิวตฺตา.\csromanspacer{} เอวํ เต ภวสาตวตฺถู ทุปฺปมุญฺจา.}

\prose{กถํ สตฺตา เอตฺโต ทุมฺโมจยา\csromandash{}สุขาย เวทนาย สตฺตา ทุมฺโมจยา, อิฏฺฐสฺมา วตฺถุสฺมา ทุมฺโมจยา, โยพฺพญฺญา ทุมฺโมจยา, อาโรคฺยา ทุมฺโมจยา, ชีวิตา ทุมฺโมจยา, ลาภา ทุมฺโมจยา, ยสา ทุมฺโมจยา, ปสํสาย ทุมฺโมจยา, สุขา ทุมฺโมจยา, มนาปิเกหิ รูเปหิ ทุมฺโมจยา, มนาปิเกหิ สทฺเทหิ.\csromanspacer{} คนฺเธหิ.\csromanspacer{} รเสหิ.\csromanspacer{} โผฏฺฐพฺเพหิ ทุมฺโมจยา, จกฺขุสมฺปทาย ทุมฺโมจยา, โสต.\csromanspacer{} ฆาน.\csromanspacer{} ชิวฺหา.\csromanspacer{} กาย.\csromanspacer{} มโนสมฺปทาย ทุมฺโมจยา.\csromanspacer{} ทุรุทฺธรา\footnote{ทุทฺธรา (ก)} ทุสฺสมุทฺธรา ทุพฺพุฏฺฐาปยา ทุสฺสมุฏฺฐาปยา ทุนฺนิเวฐยา ทุพฺพินิเวฐยา ทุตฺตรา ทุปฺปตรา ทุสฺสมติกฺกมา ทุพฺพินิวตฺตา, เอวํ สตฺตา เอตฺโต ทุมฺโมจยาติ เต ทุปฺปมุญฺจา.}

\prose{\textbf{น หิ อญฺญโมกฺขา}ติ เต อตฺตนา ปลิปปลิปนฺนา น สกฺโกนฺติ ปรํ ปลิปปลิปนฺนํ อุทฺธริตุํ.\csromanspacer{} วุตฺตํ เหตํ ภควตา\csromandash{} * โส วต จุนฺท อตฺตนา ปลิปปลิปนฺโน ปรํ ปลิปปลิปนฺนํ อุทฺธริสฺสตีติ เนตํ ฐานํ วิชฺชติ.\csromanspacer{} โส วต จุนฺท อตฺตนา อทนฺโต อวินีโต อปรินิพฺพุโต ปรํ ทเมสฺสติ วิเนสฺสติ ปรินิพฺพาเปสฺสตีติ เนตํ ฐานํ วิชฺชตีติ.\csromanspacer{} เอวมฺปิ น หิ อญฺญโมกฺขา.}

\prose{อถ วา นตฺถญฺโญ โกจิ โมเจตา, เต ยทิ มุญฺเจยฺยุํ, สเกน ถาเมน สเกน พเลน สเกน วีริเยน สเกน ปรกฺกเมน สเกน ปุริสถาเมน สเกน ปุริสพเลน สเกน ปุริสวีริเยน สเกน ปุริสปรกฺกเมน อตฺตนา สมฺมาปฏิปทํ อนุโลมปฏิปทํ อปจฺจนีกปฏิปทํ อนฺวตฺถปฏิปทํ ธมฺมานุธมฺมปฏิปทํ ปฏิปชฺชมานา มุญฺเจยฺยุนฺติ.\csromanspacer{} เอวมฺปิ น หิ อญฺญโมกฺขา.}

\csromanheader{2}{วุตฺตมฺปิ เหตํ ภควตา\csromandash{}}
\csromangathameasure{+“นาหํ สหิสฺสามิ ปโมจนาย, \\ กถํกถิํ โธตก กิญฺจิ โลเก. \\ ธมฺมญฺจ เสฏฺฐํ อภิชานมาโน,}
\csromangathagroup{\csromangathastanza{\csromansymbolfootnote{+}{ขุ 1. 439; ขุ 8. 11 ปิฏฺฐาทีสุปิ.}“นาหํ สหิสฺสามิ ปโมจนาย, \\ กถํกถิํ โธตก กิญฺจิ โลเก. \\ ธมฺมญฺจ เสฏฺฐํ อภิชานมาโน,}}

\csromanheader{2}{เอวํ ตุวํ โอฆมิมํ ตเรสี”ติ\csromandash{}}
\prose{เอวมฺปิ น หิ อญฺญโมกฺขา.}

\csromanheader{2}{2. คุหฏฺฐกสุตฺตนิทฺเทส}
\csromanheader{2}{วุตฺตมฺปิ เหตํ ภควตา\csromandash{}}
\csromangathameasure{*“อตฺตนา หิ กตํ ปาปํ, อตฺตนา สํกิลิสฺสติ. \\ อตฺตนา อกตํ ปาปํ, อตฺตนาว วิสุชฺฌติ. \\ สุทฺธี อสุทฺธิ ปจฺจตฺตํ, นาญฺโญ อญฺญํ วิโสธเย”ติ.}
\csromangathagroup{\csromangathastanza{\csromanfolio{25}\csromansymbolfootnote{*}{ขุ 1. 38; ขุ 8. 95; อภิ 4. 381 ปิฏฺเฐสุปิ.}“อตฺตนา หิ กตํ ปาปํ, อตฺตนา สํกิลิสฺสติ. \\ อตฺตนา อกตํ ปาปํ, อตฺตนาว วิสุชฺฌติ. \\ สุทฺธี อสุทฺธิ ปจฺจตฺตํ, นาญฺโญ อญฺญํ วิโสธเย”ติ.}}

\prose{เอวมฺปิ น หิ อญฺญโมกฺขา.}

\prose{วุตฺตมฺปิ เหตํ ภควตา\csromandash{} + “เอวเมวํ โข พฺราหฺมณ ติฏฺฐเตว นิพฺพานํ, ติฏฺฐติ นิพฺพานคามี มคฺโค, ติฏฺฐามหํ สมาทเปตา.\csromanspacer{} อถ จ ปน มม สาวกา มยา เอวํ โอวทิยมานา เอวํ อนุสาสิยมานา อปฺเปกจฺเจ อจฺจนฺตนิฏฺฐํ นิพฺพานํ อาราเธนฺติ, เอกจฺเจ นาราเธนฺติ, เอตฺถ กฺยาหํ พฺราหฺมณ กโรมิ, มคฺคกฺขายี พฺราหฺมณ ตถาคโต, มคฺคํ พุทฺโธ อาจิกฺขติ, อตฺตนา ปฏิปชฺชมานา มุจฺเจยฺยุนฺ”ติ.\csromanspacer{} เอวมฺปิ น หิ อญฺญโมกฺขาติ เต ทุปฺปมุญฺจา น หิ อญฺญโมกฺขา.}

\prose{\textbf{ปจฺฉา ปุเร วาปิ อเปกฺขมานา}ติ \textbf{ปจฺฉา} วุจฺจติ อนาคตํ, \textbf{ปุเร} วุจฺจติ อตีตํ.\csromanspacer{} อปิ จ อตีตํ อุปาทาย อนาคตญฺจ ปจฺจุปฺปนฺนญฺจ \textbf{ปจฺฉา}, อนาคตํ อุปาทาย อตีตญฺจ ปจฺจุปฺปนฺนญฺจ \textbf{ปุเร}.\csromanspacer{}\csromanspacer{}กถํ \textbf{ปุเร} อเปกฺขํ กโรติ\csromandash{}“เอวํรูโป อโหสิํ อตีตมทฺธานนฺ”ติ ตตฺถ นนฺทิํ สมนฺนาเนติ. “เอวํเวทโน อโหสิํ.\csromanspacer{} เอวํสญฺโญ อโหสิํ.\csromanspacer{} เอวํสงฺขาโร อโหสิํ.\csromanspacer{} เอวํวิญฺญาโณ อโหสิํ อตีตมทฺธานนฺ”ติ ตตฺถ นนฺทิํ สมนฺนาเนติ.\csromanspacer{} เอวมฺปิ \textbf{ปุเร} อเปกฺขํ กโรติ.}

\prose{อถ วา “อิติ เม จกฺขุ อโหสิ อตีตมทฺธานํ อิติ รูปา”ติ ตตฺถ ฉนฺทราคปฏิพทฺธํ โหติ วิญฺญาณํ, ฉนฺทราคปฏิพทฺธตฺตา วิญฺญาณสฺส ตทภินนฺทติ, ตทภินนฺทนฺโต เอวมฺปิ \textbf{ปุเร} อเปกฺขํ กโรติ.\csromanspacer{}\csromanspacer{}อิติ เม โสตํ อโหสิ อตีตมทฺธานํ อิติ สทฺทาติ.\csromanspacer{} อิติ เม ฆานํ อโหสิ อตีตมทฺธานํ อิติ คนฺธาติ.\csromanspacer{} อิติ เม ชิวฺหา อโหสิ อตีตมทฺธานํ อิติ รสาติ.\csromanspacer{} อิติ เม กาโย อโหสิ อตีตมทฺธานํ อิติ โผฏฺฐพฺพาติ. “อิติ เม มโน อโหสิ อตีตมทฺธานํ อิติ ธมฺมา”ติ ตตฺถ ฉนฺทราคปฏิพทฺธํ โหติ วิญฺญาณํ, ฉนฺทราคปฏิพทฺธตฺตา วิญฺญาณสฺส ตทภินนฺทติ, ตทภินนฺทนฺโต เอวมฺปิ \textbf{ปุเร} อเปกฺขํ กโรติ.}

\prose{\csromanfolio{26}อถ วา ยานิสฺส ตานิ ปุพฺเพ มาตุคาเมน สทฺธิํ หสิตลปิตกีฬิตานิ ตทสฺสาเทติ, ตํ นิกาเมติ, เตน จ วิตฺติํ อาปชฺชติ.\csromanspacer{} เอวมฺปิ ปุเร อเปกฺขํ กโรติ.}

\prose{กถํ ปจฺฉา อเปกฺขํ กโรติ\csromandash{}“เอวํรูโป สิยํ อนาคตมทฺธานนฺ”ติ ตตฺถ นนฺทิํ สมนฺนาเนติ. “เอวํเวทโน สิยํ.\csromanspacer{} เอวํสญฺโญ สิยํ.\csromanspacer{} เอวํสงฺขาโร สิยํ.\csromanspacer{} เอวํวิญฺญาโณ สิยํ อนาคตมทฺธานนฺ”ติ ตตฺถ นนฺทิํ สมนฺนาเนติ.\csromanspacer{} เอวมฺปิ ปจฺฉา อเปกฺขํ กโรติ.}

\prose{อถ วา “อิติ เม จกฺขุ สิยา อนาคตมทฺธานํ อิติ รูปา”ติ อปฺปฏิลทฺธสฺส ปฏิลาภาย จิตฺตํ ปณิทหติ, เจตโส ปณิธานปจฺจยา ตทภินนฺทติ, ตทภินนฺทนฺโต เอวมฺปิ ปจฺฉา อเปกฺขํ กโรติ.\csromanspacer{}\csromanspacer{}อิติ เม โสตํ สิยา อนาคตมทฺธานํ อิติ สทฺทาติ.\csromanspacer{} อิติ เม ฆานํ สิยา อนาคตมทฺธานํ อิติ คนฺธาติ.\csromanspacer{} อิติ เม ชิวฺหา สิยา อนาคตมทฺธานํ อิติ รสาติ.\csromanspacer{} อิติ เม กาโย สิยา อนาคตมทฺธานํ อิติ โผฏฺฐพฺพาติ. “อิติ เม มโน สิยา อนาคตมทฺธานํ อิติ ธมฺมา”ติ อปฺปฏิลทฺธสฺส ปฏิลาภาย จิตฺตํ ปณิทหติ, เจตโส ปณิธานปจฺจยา ตทภินนฺทติ, ตทภินนฺทนฺโต เอวมฺปิ ปจฺฉา อเปกฺขํ กโรติ.}

\prose{อถ วา “อิมินาหํ สีเลน วา วเตน วา ตเปน วา พฺรหฺมจริเยน วา เทโว วา ภวิสฺสามิ เทวญฺญตโร วา”ติ อปฺปฏิลทฺธสฺส ปฏิลาภาย จิตฺตํ ปณิทหติ, เจตโส ปณิธานปจฺจยา ตทภินนฺทติ, ตทภินนฺทนฺโต เอวมฺปิ ปจฺฉา อเปกฺขํ กโรตีติ ปจฺฉา ปุเร วาปิ อเปกฺขมานา.}

\prose{\textbf{อิเม ว กาเม ปุริเม ว ชปฺปนฺ}ติ \textbf{อิเม ว กาเม}ติ ปจฺจุปฺปนฺเน ปญฺจ กามคุเณ อิจฺฉนฺตา สาทิยนฺตา ปตฺถยนฺตา ปิหยนฺตา อภิชปฺปนฺตา.\csromanspacer{} ปุริเม ว ชปฺปนฺติ อตีเต ปญฺจ กามคุเณ ชปฺปนฺตา ปชปฺปนฺตา อภิชปฺปนฺตาติ \textbf{อิเม ว กาเม} ปุริเม ว ชปฺปํ.\csromanspacer{} เตนาห ภควา\csromandash{}}

\csromangathameasure{“อิจฺฉานิทานา ภวสาตพทฺธา, \\ เต ทุปฺปมุญฺจา น หิ อญฺญโมกฺขา. \\ ปจฺฉา ปุเร วาปิ อเปกฺขมานา,}
\csromangathagroup{\csromangathastanza{“อิจฺฉานิทานา ภวสาตพทฺธา, \\ เต ทุปฺปมุญฺจา น หิ อญฺญโมกฺขา. \\ ปจฺฉา ปุเร วาปิ อเปกฺขมานา,}}

\prose{\textbf{อิเม ว กาเม ปุริเม ว ชปฺปนฺ}”ติ.}

\csromanheader{2}{2. คุหฏฺฐกสุตฺตนิทฺเทส}
\proseitem{9}{\csromanfolio{27}\csromansymbolfootnote{*}{ขุ 1. 400; ขุ 10. 172 ปิฏฺเฐ.}\textbf{กาเมสุ คิทฺธา ปสุตา ปมูฬฺหา},}

\prose{\textbf{อวทานิยา เต วิสเม นิวิฏฺฐา.}}

\csromangathameasure{\textbf{ทุกฺขูปนีตา ปริเทวยนฺติ,} \\ \textbf{กิํสู ภวิสฺสาม อิโต จุตาเส.}}
\csromangathagroup{\csromangathastanza{\textbf{ทุกฺขูปนีตา ปริเทวยนฺติ,} \\ \textbf{กิํสู ภวิสฺสาม อิโต จุตาเส.}}}

\prose{\textbf{กาเมสุ คิทฺธา ปสุตา ปมูฬฺหา}ติ \textbf{กามา}ติ อุทฺทานโต เทฺว \textbf{กามา} วตฺถุ\textbf{กามา} จ กิเลส\textbf{กามา} จ ฯเปฯ อิเม วุจฺจนฺติ วตฺถุ\textbf{กามา} ฯเปฯ อิเม วุจฺจนฺติ กิเลส\textbf{กามา}.\csromanspacer{}\csromanspacer{}\textbf{เคโธ} วุจฺจติ ตณฺหา.\csromanspacer{} โย ราโค สาราโค ฯเปฯ อภิชฺฌา โลโภ อกุสลมูลํ.\csromanspacer{} กิเลสกาเมน วตฺถุกาเมสุ รตฺตา คิทฺธา คธิตา มุจฺฉิตา อชฺโฌสนฺนา\footnote{อชฺโฌปนฺนา (สี, สฺยา)} ลคฺคา ลคฺคิตา ปลิพุทฺธาติ กาเมสุ คิทฺธา.}

\prose{\textbf{ปสุตา}ติ เยปิ เม เอสนฺติ คเวสนฺติ ปริเยสนฺติ ตจฺจริตา ตพฺพหุลา ตคฺครุกา ตนฺนินฺนา ตปฺโปณา ตปฺปพฺภารา ตทธิมุตฺตา ตทธิปเตยฺยา, เตปิ กามปสุตา.\csromanspacer{}\csromanspacer{}เยปิ ตณฺหาวเสน รูเป เอสนฺติ คเวสนฺติ ปริเยสนฺติ.\csromanspacer{} สทฺเท.\csromanspacer{} คนฺเธ.\csromanspacer{} รเส.\csromanspacer{} โผฏฺฐพฺเพ ฯเปฯ ปริเยสนฺติ ตจฺจริตา ตพฺพหุลา ตคฺครุกา ตนฺนินฺนา ตปฺโปณา ตปฺปพฺภารา ตทธิมุตฺตา ตทธิปเตยฺยา, เตปิ กามปสุตา.\csromanspacer{}\csromanspacer{}เยปิ ตณฺหาวเสน รูเป ปฏิลภนฺติ.\csromanspacer{} สทฺเท.\csromanspacer{} คนฺเธ.\csromanspacer{} รเส.\csromanspacer{} โผฏฺฐพฺเพ ปฏิลภนฺติ ตจฺจริตา ตพฺพหุลา ตคฺครุกา ตนฺนินฺนา ตปฺโปณา ตปฺปพฺภารา ตทธิมุตฺตา ตทธิปเตยฺยา, เตปิ กามปสุตา.\csromanspacer{}\csromanspacer{}เยปิ ตณฺหาวเสน รูเป ปริภุญฺชนฺติ.\csromanspacer{} สทฺเท.\csromanspacer{} คนฺเธ.\csromanspacer{} รเส.\csromanspacer{} โผฏฺฐพฺเพ ปริภุญฺชนฺติ ตจฺจริตา ตพฺพหุลา ตคฺครุกา ตนฺนินฺนา ตปฺโปณา ตปฺปพฺภารา ตทธิมุตฺตา ตทธิปตเตยฺยา, เตปิ กามปสุตา.\csromanspacer{} ยถา กลหการโก กลหปสุโต, กมฺมการโก กมฺมปสุโต, โคจเร จรนฺโต โคจรปสุโต ฌายี ฌานปสุโต.\csromanspacer{} เอวเมวํ เยปิ กาเม เอสนฺติ คเวสนฺติ ปริเยสนฺติ, ตจฺจริตา ตพฺพหุลา ตคฺครุกา ตนฺนินฺนา ตปฺโปณา ตปฺปพฺภารา ตทธิมุตฺตา ตทธิปเตยฺยา, เตปิ กามปสุตา.\csromanspacer{}\csromanspacer{}เยปิ ตณฺหาวเสน รูเป เอสนฺติ คเวสนฺติ ปริเยสนฺติ.\csromanspacer{} สทฺเท.\csromanspacer{} คนฺเธ.\csromanspacer{} รเส.\csromanspacer{} โผฏฺฐพฺเพ ฯเปฯ ปริเยสนฺติ ตจฺจริตา ตพฺพหุลา ตคฺครุกา ตนฺนินฺนา ตปฺโปณา ตปฺปพฺภารา ตทธิมุตฺตา ตทธิปเตยฺยา, เตปิ กามปสุตา.\csromanspacer{}\csromanspacer{}เยปิ ตณฺหาวเสน รูเป ปฏิลภนฺติ.\csromanspacer{} สทฺเท.\csromanspacer{} คนฺเธ.\csromanspacer{} รเส.\csromanspacer{} โผฏฺฐพฺเพ ปฏิลภนฺติ ตจฺจริตา ตพฺพหุลา ตคฺครุกา \csromanfolio{28}ตนฺนินฺนา ตปฺโปณา ตปฺปพฺภารา ตทธิมุตฺตา ตทธิปเตยฺยา, เตปิ กามปสุตา.\csromanspacer{}\csromanspacer{}เยปิ ตณฺหาวเสน รูเป ปริภุญฺชนฺติ.\csromanspacer{} สทฺเท.\csromanspacer{} คนฺเธ.\csromanspacer{} รเส.\csromanspacer{} โผฏฺฐพฺเพ ปริภุญฺชนฺติ ตจฺจริตา ตพฺพหุลา ตคฺครุกา ตนฺนินฺนา ตปฺโปณา ตปฺปพฺภารา ตทธิมุตฺตา ตทธิปเตยฺยา, เตปิ กามปสุตา.}

\prose{\textbf{ปมูฬฺหา}ติ เยภุยฺเยน เทวมนุสฺสา ปญฺจสุ กามคุเณสุ มุยฺหนฺติ สมฺมุยฺหนฺติ สมฺปมุยฺหนฺติ มูฬฺหา สมฺมูฬฺหา สมฺปมูฬฺหา, อวิชฺชาย อนฺธีกตา อาวุตา นิวุตา โอวุตา ปิหิตา ปฏิจฺฉนฺนา ปฏิกุชฺชิตาติ กาเมสุ คิทฺธา ปสุตา ปมูฬฺหา.}

\prose{\textbf{อวทานิยา เต วิสเม นิวิฏฺฐา}ติ \textbf{อวทานิยา}ติ อวคจฺฉนฺตีติปิ \textbf{อวทานิยา}, มจฺฉริโนปิ วุจฺจนฺติ \textbf{อวทานิยา}, พุทฺธานํ สาวกานํ วจนํ พฺยปฺปถํ เทสนํ อนุสิฏฺฐิํ นาทิยนฺตีติ \textbf{อวทานิยา}.\csromanspacer{}\csromanspacer{}กถํ อวคจฺฉนฺตีติ \textbf{อวทานิยา}\csromandash{}นิรยํ คจฺฉนฺติ, ติรจฺฉานโยนิํ คจฺฉนฺติ, เปตฺติวิสยํ คจฺฉนฺตีติ, เอวํ อวคจฺฉนฺตีติ \textbf{อวทานิยา}.\csromanspacer{}\csromanspacer{}กถํ มจฺฉริโน วุจฺจนฺติ \textbf{อวทานิยา}\csromandash{}}

\prose{\csromansymbolfootnote{*}{อภิ 1. 226 ปิฏฺเฐปิ.}ปญฺจ มจฺฉริยานิ อาวาสมจฺฉริยํ กุลมจฺฉริยํ ลาภมจฺฉริยํ วณฺณมจฺฉริยํ ธมฺมมจฺฉริยํ.\csromanspacer{} ยํ เอวรูปํ มจฺฉริยํ มจฺฉรายนา มจฺฉรายิตตฺตํ เววิจฺฉํ กทริยํ กฏุกญฺจุกตา อคฺคหิตตฺตํ จิตฺตสฺส.\csromanspacer{} อิทํ วุจฺจติ มจฺฉริยํ.\csromanspacer{}\csromanspacer{}อปิ จ ขนฺธมจฺฉริยมฺปิ มจฺฉริยํ, ธาตุมจฺฉริยมฺปิ มจฺฉริยํ, อายตนมจฺฉริยมฺปิ มจฺฉริยํ, คาโห.\csromanspacer{} อิทํ วุจฺจติ มจฺฉริยํ.\csromanspacer{} อิมินา มจฺฉริเยน อวทญฺญุตาย สมนฺนาคตา ชนา ปมตฺตา.\csromanspacer{} เอวํ มจฺฉริโน วุจฺจนฺติ \textbf{อวทานิยา}.\csromanspacer{}\csromanspacer{}กถํ พุทฺธานํ สาวกานํ วจนํ พฺยปฺปถํ เทสนํ อนุสิฏฺฐิํ นาทิยนฺตีติ \textbf{อวทานิยา}\csromandash{} พุทฺธานํ สาวกานํ วจนํ พฺยปฺปถํ เทสนํ อนุสิฏฺฐิํ น อาทิยนฺติ น สุสฺสุสนฺติ, น โสตํ โอทหนฺติ, น อญฺญา จิตฺตํ อุปฏฺฐเปนฺติ อนสฺสวา อวจนกรา ปฏิโลมวุตฺติโน, อญฺเญเนว มุขํ กโรนฺติ.\csromanspacer{} เอวํ พุทฺธานํ สาวกานํ\footnote{พุทฺธานํ พุทฺธสาวกานํ (สี, สฺยา)} วจนํ พฺยปฺปถํ เทสนํ อนุสิฏฺฐิํ นาทิยนฺตีติ \textbf{อวทานิยา}ติ \textbf{อวทานิยา}.}

\prose{\textbf{เต วิสเม นิวิฏฺฐา}ติ วิสเม กายกมฺเม นิวิฏฺฐา, วิสเม วจีกมฺเม นิวิฏฺฐา, วิสเม มโนกมฺเม นิวิฏฺฐา, วิสเม ปาณาติปาเต นิวิฏฺฐา, วิสเม อทินฺนาทาเน นิวิฏฺฐา, วิสเม กาเมสุมิจฺฉาจาเร นิวิฏฺฐา, วิสเม มุสาวาเท นิวิฏฺฐา, วิสมาย ปิสุณาย วาจาย.\csromanspacer{} วิสมาย ผรุสาย วาจาย.\csromanspacer{} วิสเม สมฺผปฺปลาเป.\csromanspacer{} วิสมาย อภิชฺฌาย นิวิฏฺฐา, วิสเม พฺยาปาเท.\csromanspacer{} วิสมาย มิจฺฉาทิฏฺฐิยา นิวิฏฺฐา, วิสเมสุ สงฺขาเรสุ นิวิฏฺฐา, วิสเมสุ ปญฺจสุ กามคุเณสุ นิวิฏฺฐา.\csromanspacer{} วิสเมสุ ปญฺจสุ นีวรเณสุ นิวิฏฺฐา}

\vspace{\tipitakaparskip}
\csromancenter{คุหฏฺฐกสุตฺตนิทฺเทส วินิวิฏฺฐา ปติฏฺฐิตา อลฺลีนา อุปคตา อชฺโฌสิตา อธิมุตฺตา ลคฺคา ลคฺคิตา ปลิพุทฺธาติ อวทานิยา เต วิสเม นิวิฏฺฐา.}
\prose{\csromanfolio{29}\textbf{ทุกฺขูปนีตา ปริเทวยนฺตี}ติ \textbf{ทุกฺขูปนีตา}ติ ทุกฺขปฺปตฺตา ทุกฺขสมฺปตฺตา ทุกฺขูปคตา, มารปฺปตฺตา มารสมฺปตฺตา มารูปคตา, มรณปฺปตฺตา มรณสมฺปตฺตา มรณูปคตา.\csromanspacer{}\csromanspacer{}\textbf{ปริเทวยนฺตี}ติ ลปนฺติ ลาลปนฺติ\footnote{สลฺลปนฺติ (สี)} โสจนฺติ กิลมนฺติ ปริเทวนฺติ อุรตฺตาฬิํ กนฺทนฺติ สมฺโมหํ \textbf{อาปชฺชนฺตีติ ทุกฺขูปนีตา} ปริเทวยนฺติ.}

\prose{\textbf{กิํสู ภวิสฺสาม อิโต จุตาเส}ติ อิโต จุตา กิํ ภวิสฺสาม, เนรยิกา ภวิสฺสาม, ติรจฺฉานโยนิกา ภวิสฺสาม, เปตฺติวิสยิกา ภวิสฺสาม, มนุสฺสา ภวิสฺสาม, เทวา ภวิสฺสาม, รูปี ภวิสฺสาม, อรูปี ภวิสฺสาม, สญฺญี ภวิสฺสาม, อสญฺญี ภวิสฺสาม, เนวสญฺญีนาสญฺญี ภวิสฺสาม “\textbf{ภวิสฺสาม นุ โข มยํ อนาคตมทฺธานํ, น นุ โข} ภวิสฺสาม อนาคตมทฺธานํ, กิํ นุ โข ภวิสฺสาม อนาคตมทฺธานํ, กถํ นุโข ภวิสฺสาม อนาคตมทฺธานํ, กิํ หุตฺวา กิํ ภวิสฺสาม นุ โข มยํ อนาคตมทฺธานนฺ”ติ สํสยปกฺขนฺทา วิมติปกฺขนฺทา เทฺวฬฺหกชาตา ลปนฺติ ลาลปนฺติ โสจนฺติ กิลมนฺติ ปริเทวนฺติ อุรตฺตาฬิํ กนฺทนฺติ สมฺโมหํ อาปชฺชนฺตีติ กิํสู ภวิสฺสาม อิโต จุตาเส.\csromanspacer{} เตนาห ภควา\csromandash{}}

\csromangathameasure{“กาเมสุ คิทฺธา ปสุตา ปมูฬฺหา, \\ อวทานิยา เต วิสเม นิวิฏฺฐา. \\ ทุกฺขูปนีตา ปริเทวยนฺติ, \\ \textbf{กิํสู ภวิสฺสาม อิโต จุตาเส}”ติ.}
\csromangathagroup{\csromangathastanza{“กาเมสุ คิทฺธา ปสุตา ปมูฬฺหา, \\ อวทานิยา เต วิสเม นิวิฏฺฐา. \\ ทุกฺขูปนีตา ปริเทวยนฺติ, \\ \textbf{กิํสู ภวิสฺสาม อิโต จุตาเส}”ติ.}}

\csromanhangingbegin
\hangingheaditem{10}{ตสฺมา หิ สิกฺเขถ อิเธว ชนฺตุ,}
\hangingbody{\textbf{ยํ กิญฺจิ ชญฺญา วิสมนฺติ โลเก.}}
\hangingbody{\textbf{น ตสฺส เหตู วิสมํ จเรยฺย,}}
\hangingbody{\textbf{อปฺปญฺหิํ’ทํ ชีวิตมาหุ ธีรา.}}
\csromanhangingend

\prose{\textbf{ตสฺมา หิ สิกฺเขถ อิเธว ชนฺตู}ติ \textbf{ตสฺมา}ติ ตํการณา ตํเหตุ ตปฺปจฺจยา ตนฺนิทานา, เอตมาทีนวํ สมฺปสฺสมาโน กาเมสูติ \textbf{ตสฺมา}.\csromanspacer{}\csromanspacer{}\textbf{สิกฺเขถา}ติ ติสฺโส สิกฺขา อธิสีลสิกฺขา, อธิจิตฺตสิกฺขา, อธิปญฺญาสิกฺขา.\csromanspacer{} กตมา อธิสีลสิกฺขา\csromandash{}อิธ ภิกฺขุ สีลวา โหติ, \csromanfolio{30}ปาติโมกฺขสํวรสํวุโต วิหรติ อาจารโคจรสมฺปนฺโน, อณุมตฺเตสุ วชฺเชสุ ภยทสฺสาวี สมาทาย สิกฺขติ สิกฺขาปเทสุ.\csromanspacer{} ขุทฺทโก สีลกฺขนฺโธ มหนฺโต สีลกฺขนฺโธ สีลํ ปติฏฺฐา อาทิ จรณํ สํยโม สํวโร โมกฺขํ ปาโมกฺขํ กุสลานํ ธมฺมานํ สมาปตฺติยา.\csromanspacer{} อ\textbf{ยํ} อธิสีลสิกฺขา.}

\prose{กตมา อธิจิตฺตสิกฺขา\csromandash{}อิธ ภิกฺขุ วิวิจฺเจว กาเมหิ วิวิจฺจ อกุสเลหิ ธมฺเมหิ สวิตกฺกํ สวิจารํ วิเวกชํ ปีติสุขํ ปฐมํ ฌานํ อุปสมฺปชฺช วิหรติ.\csromanspacer{} วิตกฺกวิจารานํ วูปสมา อชฺฌตฺตํ สมฺปสาทนํ เจตโส เอโกทิภาวํ อวิตกฺกํ อวิจารํ สมาธิชํ ปีติสุขํ ทุติยํ ฌานํ อุปสมฺปชฺช วิหรติ.\csromanspacer{} ปีติยา จ วิราคา อุเปกฺขโก จ วิหรติ, สโต จ สมฺปชาโน, สุขญฺจ กาเยน ปฏิสํเวเทติ.\csromanspacer{} ยํ ตํ อริยา อาจิกฺขนฺติ “อุเปกฺขโก สติมา สุขวิหารี”ติ ตติยํ ฌานํ อุปสมฺปชฺช วิหรติ.\csromanspacer{} สุขสฺส จ ปหานา ทุกฺขสฺส จ ปหานา ปุพฺเพว โสมนสฺสโทมนสฺสานํ อตฺถงฺคมา อทุกฺขมสุขํ อุเปกฺขาสติปาริสุทฺธิํ จตุตฺถํ ฌานํ อุปสมฺปชฺช วิหรติ.\csromanspacer{} อยํ อธิจิตฺตสิกฺขา.}

\prose{กตมา อธิปญฺญาสิกฺขา\csromandash{}อิธ ภิกฺขุ ปญฺญวา โหติ อุทยตฺถคามินิยา ปญฺญาย สมนฺนาคโต อริยาย นิพฺเพธิกาย สมฺมาทุกฺขกฺขยคามินิยา, โส “อิทํ ทุกฺขนฺ”ติ ยถาภูตํ ปชานาติ. “อยํ ทุกฺขสมุทโย”ติ ยถาภูตํ ปชานาติ. “อยํ ทุกฺขนิโรโธ”ติ ยถาภูตํ ปชานาติ. “อยํ ทุกฺขนิโรธคามินี ปฏิปทา”ติ ยถาภูตํ ปชานาติ. “อิเม อาสวา”ติ ยถาภูตํ ปชานาติ. “อยํ อาสวสมุทโย”ติ ยถาภูตํ ปชานาติ. “อยํ อาสวนิโรโธ”ติ ยถาภูตํ ปชานาติ. “อยํ อาสวนิโรธคามินี ปฏิปทา”ติ ยถาภูตํ ปชานาติ.\csromanspacer{} อยํ อธิปญฺญาสิกฺขา.\csromanspacer{} อิมา ติสฺโส สิกฺขาโย อาวชฺชนฺโต สิกฺเขยฺย, ชานนฺโต สิกฺเขยฺย, ปสฺสนฺโต สิกฺเขยฺย, ปจฺจเวกฺขนฺโต สิกฺเขยฺย, จิตฺตํ อธิฏฺฐหนฺโต สิกฺเขยฺย, สทฺธาย อธิมุจฺจนฺโต สิกฺเขยฺย, วีริยํ ปคฺคณฺหนฺโต สิกฺเขยฺย, สติํ อุปฏฺฐเปนฺโต สิกฺเขยฺย, จิตฺตํ สมาทหนฺโต สิกฺเขยฺย, ปญฺญาย ปชานนฺโต สิกฺเขยฺย, อภิญฺเญยฺยํ อภิชานนฺโต สิกฺเขยฺย, ปริญฺเญยฺยํ ปริชานนฺโต สิกฺเขยฺย, ปหาตพฺพํ ปชหนฺโต สิกฺเขยฺย, ภาเวตพฺพํ ภาเวนฺโต}

\vspace{\tipitakaparskip}
\csromancenter{คุหฏฺฐกสุตฺตนิทฺเทส สิกฺเขยฺย, สจฺฉิกาตพฺพํ สจฺฉิกโรนฺโต สิกฺเขยฺย อาจเรยฺย สมาจเรยฺย สมาทาย วตฺเตยฺย.}
\prose{\csromanfolio{31}\textbf{อิธา}ติ อิมิสฺสา ทิฏฺฐิยา อิมิสฺสา ขนฺติยา อิมิสฺสา รุจิยา อิมสฺมิํ อาทาเย อิมสฺมิํ ธมฺเม อิมสฺมิํ วินเย อิมสฺมิํ ธมฺมวินเย อิมสฺมิํ ปาวจเน อิมสฺมิํ พฺรหฺมจริเย อิมสฺมิํ สตฺถุสาสเน อิมสฺมิํ อตฺตภาเว อิมสฺมิํ มนุสฺสโลเก, เตน วุจฺจติ อิธาติ.\csromanspacer{}\csromanspacer{}\textbf{ชนฺตู}ติ สตฺโต นโร ฯเปฯ มนุโชติ ตสฺมา หิ สิกฺเขถ อิเธว ชนฺตุ.}

\prose{\textbf{ยํ} \textbf{กิญฺจิ ชญฺญา วิสมนฺติ โลเก}ติ \textbf{ยํ} \textbf{กิญฺจี}ติ สพฺเพน สพฺพํ สพฺพถา สพฺพํ อเสสํ นิสฺเสสํ ปริยาทิยนวจนเมตํ\footnote{ปริทายวจนเมตํ (สฺยา)} \textbf{ยํ} \textbf{กิญฺจี}ติ.\csromanspacer{}\csromanspacer{}\textbf{วิสมนฺติ ชญฺญา}ติ วิสมํ กายกมฺมํ วิสมนฺติ ชาเนยฺย, วิสมํ วจีกมฺมํ วิสมนฺติ ชาเนยฺย, วิสมํ มโนกมฺมํ วิสมนฺติ ชาเนยฺย, วิสมํ ปาณาติปาตํ วิสโมติ ชาเนยฺย, วิสมํ อทินฺนาทานํ วิสมนฺติ ชาเนยฺย, วิสมํ กาเมสุมิจฺฉาจารํ วิสโมติ ชาเนยฺย,วิสมํ มุสาวาทํ วิสโมติ ชาเนยฺย, วิสมํ ปิสุณํ วาจํ วิสมาติ ชาเนยฺย, วิสมํ ผรุสํ วาจํ วิสมาติ ชาเนยฺย, วิสมํ สมฺผปฺปลาปํ วิสโมติ ชาเนยฺย, วิสมํ อภิชฺฌํ วิสมาติ ชาเนยฺย, วิสมํ พฺยาปาทํ วิสโมติ ชาเนยฺย, วิสมํ มิจฺฉาทิฏฺฐิํ วิสมาติ ชาเนยฺย, วิสเม สงฺขาเร วิสมาติ ชาเนยฺย, วิสเม ปญฺจ กามคุเณ วิสมาติ ชาเนยฺย,วิสเม ปญฺจ นีวรเณ วิสมาติ ชาเนยฺย อาชาเนยฺย วิชาเนยฺย ปฏิวิชาเนยฺย ปฏิวิชฺเฌยฺย.\csromanspacer{}\csromanspacer{}\textbf{โลเก}ติ อปายโลเก ฯเปฯ อายตนโลเกติ \textbf{ยํ} \textbf{กิญฺจิ ชญฺญา วิสมนฺติ โลเก}.}

\prose{\textbf{น ตสฺส เหตู วิสมํ จเรยฺยา}ติ วิสมสฺส กายกมฺมสฺส เหตุ วิสมํ น จเรยฺย, วิสมสฺส วจีกมฺมสฺส เหตุ วิสมํ น จเรยฺย, วิสมสฺส มโนกมฺมสฺส เหตุ วิสมํ น จเรยฺย, วิสมสฺส ปาณาติปาตสฺส เหตุ วิสมํ น จเรยฺย, วิสมสฺส อทินฺนาทานสฺส เหตุ วิสมํ น จเรยฺย, วิสมสฺส กาเมสุมิจฺฉาจารสฺส เหตุ วิสมํ น จเรยฺย, วิสมสฺส มุสาวาทสฺส เหตุ วิสมํ น จเรยฺย, วิสมาย ปิสุณาย วาจาย เหตุ วิสมํ น จเรยฺย, วิสมาย ผรุสาย วาจาย เหตุ วิสมํ น จเรยฺย, วิสมสฺส สมฺผปฺปลาปสฺส เหตุ วิสมํ น จเรยฺย, วิสมาย \csromanfolio{32}อภิชฺฌาย เหตุ วิสมํ น จเรยฺย, วิสมสฺส พฺยาปาทสฺส เหตุ วิสมํ น จเรยฺย, วิสมาย มิจฺฉาทิฏฺฐิยา เหตุ วิสมํ น จเรยฺย, วิสมานํ สงฺขารานํ เหตุ วิสมํ น จเรยฺย, วิสมานํ ปญฺจนฺนํ กามคุณานํ เหตุ วิสมํ น จเรยฺย, วิสมานํ ปญฺจนฺนํ นีวรณานํ เหตุ วิสมํ น จเรยฺย, วิสมาย เจตนาย เหตุ วิสมํ น จเรยฺย, วิสมาย ปตฺถนาย เหตุ วิสมํ น \textbf{จเรยฺย, วิสมาย ปณิธิยา เหตุ วิสมํ น จเรยฺย, น อาจเรยฺย น} สมาจเรยฺย น สมาทาย วตฺเตยฺยาติ น ตสฺส เหตู วิสมํ จเรยฺย.}

\prose{\textbf{อปฺปญฺหิ’ทํ ชีวิตมาหุ ธีรา}ติ \textbf{ชีวิตนฺ}ติ * อายุ ฐิติ ยปนา ยาปนา อิริยนา วตฺตนา ปาลนา ชีวิตํ ชีวิตินฺทฺริยํ.\csromanspacer{} อปิ จ ทฺวีหิ การเณหิ อปฺปกํ ชีวิตํ ฐิติปริตฺตตาย วา อปฺปกํ ชีวิตํ, สรสปริตฺตตาย วา อปฺปกํ ชีวิตํ.\csromanspacer{} กถํ ฐิติปริตฺตตายา อปฺปกํ ชีวิตํ\csromandash{}อตีเต จิตฺตกฺขเณ ชีวิตฺถ น ชีวติ น ชีวิสฺสติ, อนาคเต จิตฺตกฺขเณ ชีวิสฺสติ, น ชีวติ น ชีวิตฺถ, ปจฺจุปฺปนฺเน จิตฺตกฺขเณ ชีวิติ, น ชีวิตฺถ น ชีวิสฺสติ.}

\csromangathameasure{ชีวิตํ อตฺตภาโว จ, สุขทุกฺขา จ เกวลา. \\ เอกจิตฺตสมายุตฺตา, ลหุโส วตฺตเต ขโณ. \\ จุลฺลาสีติสหสฺสานิ, กปฺปา ติฏฺฐนฺติ เย มรู. \\ นเตฺวว เตปิ ชีวนฺติ, ทฺวีหิ จิตฺเตหิ สํยุตา. \\ เย นิรุทฺธา มรนฺตสฺส, ติฏฺฐมานสฺส วา อิธ. \\ สพฺเพปิ สทิสา ขนฺธา, คตา อปฺปฏิสนฺธิกา. \\ อนนฺตรา จ เย ภคฺคา, เย จ ภคฺคา1 อนาคตา. \\ ตทนฺตเร นิรุทฺธานํ, เวสมํ นตฺถิ ลกฺขเณ. \\ อนิพฺพตฺเตน น ชาโต, ปจฺจุปฺปนฺเนน ชีวติ. \\ จิตฺตภคฺคา มโต โลโก, ปญฺญตฺติ ปรมตฺถิยา. \\ ยถา นินฺนา ปวตฺตนฺติ, ฉนฺเทน ปริณามิตา. \\ อจฺฉินฺนธารา วตฺตนฺติ, สฬายตนปจฺจยา. \\ อนิธานคตา ภคฺคา, ปุญฺโช นตฺถิ อนาคเต. \\ นิพฺพตฺตา เย จ ติฏฺฐนฺติ, อารคฺเค สาสปูปมา.}
\csromangathagroup{\csromangathastanza{ชีวิตํ อตฺตภาโว จ, สุขทุกฺขา จ เกวลา. \\ เอกจิตฺตสมายุตฺตา, ลหุโส วตฺตเต ขโณ.}\csromangathastanza{จุลฺลาสีติสหสฺสานิ, กปฺปา ติฏฺฐนฺติ เย มรู. \\ นเตฺวว เตปิ ชีวนฺติ, ทฺวีหิ จิตฺเตหิ สํยุตา.}\csromangathastanza{เย นิรุทฺธา มรนฺตสฺส, ติฏฺฐมานสฺส วา อิธ. \\ สพฺเพปิ สทิสา ขนฺธา, คตา อปฺปฏิสนฺธิกา.}\csromangathastanza{อนนฺตรา จ เย ภคฺคา\footnote{ภงฺคา (สี, สฺยา)}, เย จ ภคฺคา1 อนาคตา. \\ ตทนฺตเร นิรุทฺธานํ, เวสมํ นตฺถิ ลกฺขเณ.}\csromangathastanza{อนิพฺพตฺเตน น ชาโต, ปจฺจุปฺปนฺเนน ชีวติ. \\ จิตฺตภคฺคา มโต โลโก, ปญฺญตฺติ ปรมตฺถิยา.}\csromangathastanza{ยถา นินฺนา ปวตฺตนฺติ, ฉนฺเทน ปริณามิตา. \\ อจฺฉินฺนธารา วตฺตนฺติ, สฬายตนปจฺจยา.}\csromangathastanza{อนิธานคตา ภคฺคา, ปุญฺโช นตฺถิ อนาคเต. \\ นิพฺพตฺตา เย จ\footnote{นิพฺพตฺตาเยว (สพฺพตฺถ)} ติฏฺฐนฺติ, อารคฺเค สาสปูปมา.}}

\csromanheader{2}{2. คุหฏฺฐกสุตฺตนิทฺเทส}
\csromangathameasure{นิพฺพตฺตานญฺจ ธมฺมานํ, ภงฺโค เนสํ ปุรกฺขโต. \\ ปโลกธมฺมา ติฏฺฐนฺติ, ปุราเณหิ อมิสฺสิตา. \\ อทสฺสนโต อายนฺติ, ภงฺคา คจฺฉนฺตฺยทสฺสนํ.}
\csromangathagroup{\csromangathastanza{\csromanfolio{33}นิพฺพตฺตานญฺจ ธมฺมานํ, ภงฺโค เนสํ ปุรกฺขโต. \\ ปโลกธมฺมา ติฏฺฐนฺติ, ปุราเณหิ อมิสฺสิตา. \\ อทสฺสนโต อายนฺติ, ภงฺคา คจฺฉนฺตฺยทสฺสนํ.}}

\vspace{\tipitakaparskip}
\csromancenter{วิชฺชุปฺปาโทว อากาเส, อุปฺปชฺชนฺติ วยนฺติ จาติ\csromandash{}}
\prose{เอวํ ฐิติปริตฺตตาย อปฺปกํ ชีวิตํ.}

\prose{กถํ สรสปริตฺตตาย อปฺปกํ ชีวิตํ\csromandash{}อสฺสาสูปนิพนฺธํ ชีวิตํ, ปสฺสาสูปนิพนฺธํ ชีวิตํ, อสฺสาสปสฺสาสูปนิพนฺธํ ชีวิตํ, มหาภูตูปนิพนฺธํ ชีวิตํ, กพฬีการาหารูปนิพนฺธํ ชีวิตํ, อุสฺมูปนิพนฺธํ ชีวิตํ, วิญฺญาณูปนิพนฺธํ ชีวิตํ, มูลมฺปิ อิเมสํ ทุพฺพลํ, ปุพฺพเหตูปิ อิเมสํ ทุพฺพลา, เย ปจฺจยา เตปิ ทุพฺพลา, เยปิ ปภาวิกา เตปิ ทุพฺพลา, สหภูมิ อิเมสํ ทุพฺพลา สมฺปโยคาปิ อิเมสํ ทุพฺพลา สหชาปิ อิเมสํ ทุพฺพลา, ยาปิ ปโยชิกา, สาปิ ทุพฺพลา, อญฺญมญฺญํ อิเม นิจฺจทุพฺพลา, อญฺญมญฺญํ อนวฏฺฐิตา อิเม, อญฺญมญฺญํ ปริปาตยนฺติ อิเม, อญฺญมญฺญสฺส หิ นตฺถิ ตายิตา, น จาปิ ฐเปนฺติ อญฺญมญฺญํ อิเม, โยปิ นิพฺพตฺตโก โส น วิชฺชติ.}

\csromangathameasure{น จ เกนจิ โกจิ หายติ, \\ คนฺธพฺพา จ อิเม หิ สพฺพโส. \\ ปุริเมหิ ปภาวิกา อิเม, \\ เยปิ ปภาวิกา เต ปุเร มตา. \\ ปุริมาปิ จ ปจฺฉิมาปิ จ, \\ อญฺญมญฺญํ น กทาจิ มทฺทสํสูติ.}
\csromangathagroup{\csromangathastanza{น จ เกนจิ โกจิ หายติ, \\ คนฺธพฺพา จ อิเม หิ สพฺพโส. \\ ปุริเมหิ ปภาวิกา อิเม, \\ เยปิ ปภาวิกา เต ปุเร มตา. \\ ปุริมาปิ จ ปจฺฉิมาปิ จ, \\ อญฺญมญฺญํ น กทาจิ มทฺทสํสูติ.}}

\prose{เอวํ สรสปริตฺตตาย อปฺปกํ ชีวิตํ.}

\prose{อปิ จ จาตุมหาราชิกานํ เทวานํ ชีวิตํ อุปาทาย มนุสฺสานํ อปฺปกํ ชีวิตํ ปริตฺตกํ ชีวิตํ โถกํ\footnote{โถกกํ (ก) อุปริ 91 ปิฏฺเฐปิ.} ชีวิตํ ขณิกํ ชีวิตํ ลหุกํ ชีวิตํ อิตฺตรํ ชีวิตํ อนทฺธนียํ ชีวิตํ นจิรฏฺฐิติกํ ชีวิตํ, ตาวติํสานํ เทวานํ.\csromanspacer{} ยามานํ เทวานํ.\csromanspacer{} ตุสิตานํ เทวานํ.\csromanspacer{} นิมฺมานรตีนํ เทวานํ.\csromanspacer{} ปรนิมฺมิตวสวตฺตีนํ เทวานํ.\csromanspacer{} พฺรหฺมกายิกานํ เทวานํ ชีวิตํ อุปาทาย มนุสฺสานํ อปฺปกํ ชีวิตํ ปริตฺตกํ ชีวิตํ โถกํ ชีวิตํ ขณิกํ ชีวิตํ ลหุกํ ชีวิตํ อิตฺตรํ ชีวิตํ อนทฺธนียํ ชีวิตํ นจิรฏฺฐิติกํ ชีวิตํ.\csromanspacer{} วุตฺตํ เหตํ ภควตา\csromandash{}}

\prose{\csromanfolio{34}\csromansymbolfootnote{*}{สํ 1. 109, 110 ปิฏฺเฐสุ.}อปฺปมิทํ ภิกฺขเว มนุสฺสานํ อายุ, คมนิโย สมฺปราโย, มนฺตาย โพทฺธพฺพํ, กตฺตพฺพํ กุสลํ, จริตพฺพํ พฺรหฺมจริยํ, นตฺถิ ชาตสฺส อมรณํ.\csromanspacer{} โย ภิกฺขเว จิรํ ชีวติ, โส วสฺสสตํ, อปฺปํ วา ภิยฺโย.}

\csromangathameasure{*อปฺปมายุ มนุสฺสานํ, หีเฬยฺย นํ สุโปริโส. \\ จเรยฺยาทิตฺตสีโสว, นตฺถิ มจฺจุสฺสนาคโม. \\ อจฺจยนฺติ อโหรตฺตา, ชีวิตํ อุปรุชฺฌติ. \\ อายุ ขิยฺยติ มจฺจานํ, กุนฺนทีนํว โอทกนฺติ.}
\csromangathagroup{\csromangathastanza{\csromansymbolmark{*}อปฺปมายุ มนุสฺสานํ, หีเฬยฺย นํ สุโปริโส. \\ จเรยฺยาทิตฺตสีโสว, นตฺถิ มจฺจุสฺสนาคโม.}\csromangathastanza{อจฺจยนฺติ อโหรตฺตา, ชีวิตํ อุปรุชฺฌติ. \\ อายุ ขิยฺยติ มจฺจานํ, กุนฺนทีนํว โอทกนฺติ.}}

\prose{\textbf{อปฺปญฺหิ ทํ ชีวิตมาหุ ธีรา}ติ \textbf{ธี}ราติ \textbf{ธี}รา, ธิติมาติ \textbf{ธี}รา, ธิติสมฺปํนฺนาติ \textbf{ธี}รา, \textbf{ธี}กตปาปาติ \textbf{ธี}รา, \textbf{ธี} วุจฺจติ ปญฺญา, + ยา ปญฺญา ปชานนา วิจโย ปวิจโย ธมฺมวิจโย สลฺลกฺขณา อุปลกฺขณา ปจฺจุปลกฺขณา ปณฺฑิจฺจํ โกสลฺลํ เนปุญฺญํ เวภพฺยา จินฺตา อุปปริกฺขา ภูริ เมธา ปริณายิกา วิปสฺสนา สมฺปชญฺญํ ปโตโท ปญฺญา ปญฺญินฺทฺริยํ ปญฺญาพลํ ปญฺญาสตฺถํ ปญฺญาปาสาโท ปญฺญา-อาโลโก ปญฺญา-โอภาโส ปญฺญาปชฺโชโต ปญฺญารตนํ อโมโห ธมฺมวิจโย สมฺมาทิฏฺฐิ, ตาย ปญฺญาย สมนฺนาคตตฺตา \textbf{ธี}รา.\csromanspacer{}\csromanspacer{}อปิ จ ขนฺธ\textbf{ธี}รา ธาตุ\textbf{ธี}รา อายตน\textbf{ธี}รา ปฏิจฺจสมุปฺปาท\textbf{ธี}รา สติปฏฺฐาน\textbf{ธี}รา สมฺมปฺปธาน\textbf{ธี}รา อิทฺธิปาท\textbf{ธี}รา อินฺทฺริย\textbf{ธี}รา พล\textbf{ธี}รา โพชฺฌงฺค\textbf{ธี}รา มคฺค\textbf{ธี}รา ผล\textbf{ธี}รา นิพฺพาน\textbf{ธี}รา, เต \textbf{ธี}รา เอวมาหํสุ “มนุสฺสานํ อปฺปกํ ชีวิตํ, ปริตฺตกํ ชีวิตํ, โถกํ ชีวิตํ, ขณิกํ ชีวิตํ, ลหุกํ ชีวิตํ, อิตฺตรํ ชีวิตํ, อนทฺธนียํ ชีวิตํ, นจิรฏฺฐิติกํ ชีวิตนฺ”ติ เอวมาหํสุ เอวํ กเถนฺติ เอวํ ภณนฺติ เอวํ ทีปยนฺติ เอวํ โวหรนฺตีติ อปฺปญฺหิทํ ชีวิตมาหุ \textbf{ธี}รา.\csromanspacer{} เตนาห ภควา\csromandash{}}

\csromangathameasure{ตสฺมา หิ สิกฺเขถ อิเธว ชนฺตุ, \\ ยํ กิญฺจิ ชญฺญา วิสมนฺติ โลเก. \\ น ตสฺส เหตู วิสมํ จเรยฺย, \\ อปฺปญฺหิ’ทํ ชีวิตมาหุ \textbf{ธี}ราติ.}
\csromangathagroup{\csromangathastanza{ตสฺมา หิ สิกฺเขถ อิเธว ชนฺตุ, \\ ยํ กิญฺจิ ชญฺญา วิสมนฺติ โลเก. \\ น ตสฺส เหตู วิสมํ จเรยฺย, \\ อปฺปญฺหิ’ทํ ชีวิตมาหุ \textbf{ธี}ราติ.}}

\proseitem{11}{\csromansymbolfootnote{+}{อภิ 1. 20 ปิฏฺฐาทีสุปิ. ++ ขุ 1. 400 ปิฏฺเฐปิ.}+ ปสฺสามิ โลเก ปริผนฺทมานํ,}

\prose{\textbf{ปชํ อิมํ ตณฺหคตํ ภเวสุ.}}

\csromangathameasure{\textbf{หีนา นรา มจฺจุมุเข ลปนฺติ,} \\ \textbf{อวีตตณฺหาเส ภวาภเวสุ.}}
\csromangathagroup{\csromangathastanza{\textbf{หีนา นรา มจฺจุมุเข ลปนฺติ,} \\ \textbf{อวีตตณฺหาเส ภวาภเวสุ.}}}

\csromanheader{2}{2. คุหฏฺฐกสุตฺตนิทฺเทส}
\prose{\csromanfolio{35}\textbf{ปสฺสามิ โลเก ปริผนฺทมานนฺ}ติ \textbf{ปสฺสามี}ติ มํสจกฺขุนาปิ ปสฺสามิ, ทิพฺพจกฺขุนาปิ ปสฺสามิ, ปญฺญาจกฺขุนาปิ ปสฺสามิ, พุทฺธจกฺขุนาปิ ปสฺสามิ, สมนฺตจกฺขุนาปิ ปสฺสามิ ทกฺขามิ โอโลเกมิ นิชฺฌายามิ อุปปริกฺขามิ.\csromanspacer{}\csromanspacer{}\textbf{โลเก}ติ อปายโลเก มนุสฺสโลเก เทวโลเก, ขนฺธโลเก ธาตุโลเก อายตนโลเก.}

\prose{\textbf{ปริผนฺทมานนฺ}ติ ตณฺหาผนฺทนาย ผนฺทมานํ ทิฏฺฐิผนฺทนาย ผนฺทมานํ กิเลสผนฺทนาย ผนฺทมานํ ปโยคผนฺทนาย ผนฺทมานํ วิปากผนฺทนาย ผนฺทมานํ ทุจฺจริตผนฺทนาย ผนฺทมานํ, รตฺตํ ราเคน ผนฺทมานํ ทุฏฺฐํ โทเสน ผนฺทมานํ มูฬฺหํ โมเหน ผนฺทมานํ, วินิพทฺธํ มาเนน ผนฺทมานํ ปรามฏฺฐํ ทิฏฺฐิยา ผนฺทมานํ วิกฺเขปคตํ อุทฺธจฺเจน ผนฺทมานํ อนิฏฺฐงฺคตํ วิจิกิจฺฉาย ผนฺทมานํ ถามคตํ อนุสเยหิ ผนฺทมานํ, ลาเภน ผนฺทมานํ อลาเภน ผนฺทมานํ ยเสน ผนฺทมานํ อยเสน ผนฺทมานํ ปสํสาย ผนฺทมานํ นินฺทาย ผนฺทมานํ สุเขน ผนฺทมานํ ทุกฺเขน ผนฺทมานํ, ชาติยา ผนฺทมานํ ชราย ผนฺทมานํ พฺยาธินา ผนฺทมานํ มรเณน ผนฺทมานํ โสกปริเทวทุกฺขโทมนสฺสุปายาเสหิ ผนฺทมานํ, เนรยิเกน ทุกฺเขน ผนฺทมานํ ติรจฺฉานโยนิเกน ทุกฺเขน ผนฺทมานํ เปตฺติวิสยิเกน ทุกฺเขน ผนฺทมานํ มานุสิเกน ทุกฺเขน ผนฺทมานํ, คพฺโภกฺกนฺติมูลเกน ทุกฺเขน ผนฺทมานํ คพฺเภ ฐิติมูลเกน ทุกฺเขน ผนฺทมานํ คพฺภา วุฏฺฐานมูลเกน ทุกฺเขน ผนฺทมานํ ชาตสฺสูปนิพนฺธเกน ทุกฺเขน ผนฺทมานํ ชาตสฺส ปราเธยฺยเกน ทุกฺเขน ผนฺทมานํ อตฺตูปกฺกเมน ทุกฺเขน ผนฺทมานํ ปรูปกฺกเมน ทุกฺเขน ผนฺทมานํ, ทุกฺขทุกฺเขน ผนฺทมานํ สงฺขารทุกฺเขน ผนฺทมานํ วิปริณามทุกฺเขน ผนฺทมานํ, จกฺขุโรเคน ทุกฺเขน ผนฺทมานํ โสตโรเคน ทุกฺเขน ผนฺทมานํ ฆานโรเคน ทุกฺเขน.\csromanspacer{} ชิวฺหาโรเคน.\csromanspacer{} กายโรเคน.\csromanspacer{} สีสโรเคน.\csromanspacer{} กณฺณโรเคน.\csromanspacer{} มุขโรเคน.\csromanspacer{} ทนฺตโรเคน.\csromanspacer{} กาเสน.\csromanspacer{} สาเสน.\csromanspacer{} ปินาเสน.\csromanspacer{} ทาเหน.\csromanspacer{} ชเรน.\csromanspacer{} กุจฺฉิโรเคน.\csromanspacer{} มุจฺฉาย.\csromanspacer{} ปกฺขนฺทิกาย.\csromanspacer{} สูลาย.\csromanspacer{} วิสุจิกาย.\csromanspacer{} กุฏฺเฐน.\csromanspacer{} คณฺเฑน.\csromanspacer{} กิลาเสน.\csromanspacer{} โสเสน.\csromanspacer{} อปมาเรน.\csromanspacer{} ททฺทุยา.\csromanspacer{} กณฺฑุยา.\csromanspacer{} กจฺฉุยา.\csromanspacer{} รขสาย.\csromanspacer{} วิตจฺฉิกาย.\csromanspacer{} โลหิเตน.\csromanspacer{} ปิตฺเตน.\csromanspacer{} มธุเมเหน.\csromanspacer{} อํสาย.\csromanspacer{} ปิฬกาย.\csromanspacer{} ภคนฺทลาย\footnote{ภคนฺทเลน (สี, ก)}.\csromanspacer{} ปิตฺตสมุฏฺฐาเนน อาพาเธน.\csromanspacer{} เสมฺหสมุฏฺฐาเนน อาพาเธน.\csromanspacer{} วาตสมุฏฺฐาเนน อาพาเธน.\csromanspacer{} สนฺนิปาติเกน อาพาเธน.\csromanspacer{} อุตุปริณามเชน อาพาเธน.\csromanspacer{} วิสมปริหารเชน อาพาเธน.\csromanspacer{} โอปกฺกมิเกน}

\prose{\csromanfolio{36}อาพาเธน.\csromanspacer{} กมฺมวิปากเชน อาพาเธน.\csromanspacer{} สีเตน.\csromanspacer{} อุเณฺหน.\csromanspacer{} ชิฆจฺฉาย.\csromanspacer{} ปิปาสาย.\csromanspacer{} อุจฺจาเรน.\csromanspacer{} ปสฺสาเวน.\csromanspacer{} qอํสมกสวาตาตปสรีสปสมฺผสฺเสน ทุกฺเขน.\csromanspacer{} มาตุมรเณน ทุกฺเขน.\csromanspacer{} ปิตุมรเณน ทุกฺเขน.\csromanspacer{} ภาตุมรเณน ทุกฺเขน.\csromanspacer{} ภคินิมรเณน ทุกฺเขน.\csromanspacer{} ปุตฺตมรเณน ทุกฺเขน.\csromanspacer{} มีตุมรเณน ทุกฺเขน.\csromanspacer{} ญาติพฺยสเนน.\csromanspacer{} โภคพฺยสเนน.\csromanspacer{} โรคพฺยสเนน.\csromanspacer{} สีลพฺยสเนน.\csromanspacer{} ทิฏฺฐิพฺยสเนน ทุกฺเขน ผนฺทมานํ สมฺผนฺทมานํ วิปฺผนฺทมานํ เวธมานํ ปเวธมานํ สมฺปเวธมานํ ปสฺสามิ ทกฺขามิ โอโลเกมิ นิชฺฌายามิ อุปปริกฺขามีติ ปสฺสามิ โลเก ปริผนฺทมานํ.}

\prose{\textbf{ปชํ อิมํ ตณฺหคตํ ภเวสู}ติ \textbf{ปชา}ติ สตฺตาธิวจนํ.\csromanspacer{}\csromanspacer{}\textbf{ตณฺหา}ติ รูป\textbf{ตณฺหา} สทฺท\textbf{ตณฺหา} คนฺธ\textbf{ตณฺหา} รส\textbf{ตณฺหา} โผฏฺฐพฺพ\textbf{ตณฺหา} ธมฺม\textbf{ตณฺหา}.\csromanspacer{}\csromanspacer{}\textbf{ตณฺหคตนฺ}ติ \textbf{ตณฺหา}คตํ \textbf{ตณฺหา}นุคตํ \textbf{ตณฺหา}ยานุสฏํ \textbf{ตณฺหา}ยาสนฺนํ \textbf{ตณฺหา}ย ปาติตํ อภิภูตํ ปริยาทินฺนจิตฺตํ.\csromanspacer{}\csromanspacer{}\textbf{ภเวสู}ติ กามภเว รูปภเว อรูปภเวติ ปชํ อิมํ ตณฺหคตํ ภเวสุ.}

\prose{\textbf{หีนา นรา มจฺจุมุเข ลปนฺตี}ติ \textbf{หีนา นรา}ติ \textbf{หีนา นรา}.\csromanspacer{} หีเนน กายกมฺเมน สมนฺนาคตาติ \textbf{หีนา นรา}, หีเนน วจีกมฺเมน สมนฺนาคตาติ \textbf{หีนา นรา}, หีเนน มโนกมฺเมน สมนฺนาคตาติ \textbf{หีนา นรา}.\csromanspacer{}\csromanspacer{}หีเนน ปาณาติปาเตน สมนฺนาคตาติ \textbf{หีนา นรา}, หีเนน อทินฺนาทาเนน.\csromanspacer{} หีเนน กาเมสุมิจฺฉาจาเรน.\csromanspacer{} หีเนน มุสาวาเทน.\csromanspacer{} หีนาย ปิสุณาย วาจาย.\csromanspacer{} หีนาย ผรุสาย วาจาย.\csromanspacer{} หีเนน สมฺผปฺปลาเปน.\csromanspacer{} หีนาย อภิชฺฌาย.\csromanspacer{} หีเนน พฺยาปาเทน.\csromanspacer{} หีนาย มิจฺฉาทิฏฺฐิยา.\csromanspacer{} หีเนหิ สงฺขาเรหิ.\csromanspacer{} หีเนหิ ปญฺจหิ กามคุเณหิ นีวรเณหิ.\csromanspacer{} หีนาย เจตนาย.\csromanspacer{} หีนาย ปตฺถนาย.\csromanspacer{} หีนาย ปณิธิยา สมนฺนาคตาติ \textbf{หีนา นรา}.\csromanspacer{}\csromanspacer{}หีนา นิหีนา โอหีนา โอมกา ลามกา ฉตุกฺกา ปริตฺตาติ \textbf{หีนา นรา}.\csromanspacer{}\csromanspacer{}\textbf{มจฺจุมุเข ลปนฺตี}ติ \textbf{มจฺจุมุเข}ติ มารมุเข มรณมุเข, มจฺจุปฺปตฺตา มจฺจุสมฺปตฺตา มจฺจูปาคตา, มารปฺปตฺตา มารสมฺปตฺตา มารูปาคตา, มรณปฺปตฺตา มรณสมฺปตฺตา มรณูปาคตา ลปนฺติ ลาลปนฺติ โสจนฺติ กิลมนฺติ ปริเทวนฺติ อุรตฺตาฬิํ กนฺทนฺติ สมฺโมหํ อาปชฺชนฺตีติ \textbf{หีนา นรา} \textbf{มจฺจุมุเข} ลปนฺติ.}

\prose{\textbf{อวีตตณฺหาเส ภวาภเวสู}ติ \textbf{ตณฺหา}ติ รูป\textbf{ตณฺหา} ฯเปฯ ธมฺม\textbf{ตณฺหา}.\csromanspacer{} \textbf{ภวาภเวสู}ติ ภวาภเว กมฺมภเว ปุนพฺภเว กามภเว, กมฺมภเว กามภเว ปุนพฺภเว รูปภเว, กมฺมภเว รูปภเว ปุนพฺภเว อรูปภเว,}

\vspace{\tipitakaparskip}
\csromancenter{คุหฏฺฐกสุตฺตนิทฺเทส กมฺมภเว อรูปภเว ปุนพฺภเว ปุนปฺปุนพฺภเว, ปุนปฺปุนคติยา ปุนปฺปุน-อุปปตฺติยา ปุนปฺปุนปฏิสนฺธิยา ปุนปฺปุน- อตฺตภาวาภินิพฺพตฺติยา, อวีตตณฺหา อวิคตตณฺหา อจตฺตตณฺหา อวนฺตตณฺหา อมุตฺตตณฺหา อปฺปหีนตณฺหา อปฺปฏินิสฺสฏฺฐตณฺหาติ อวีตตณฺหาเส ภวาภเวสุ.\csromanspacer{} เตนาห ภควา\csromandash{}}
\vspace{0.5\baselineskip}
\csromangathameasure{ปสฺสามิ โลเก ปริผนฺทมานํ, \\ ปชํ อิมํ ตณฺหคตํ ภเวสุ. \\ หีนา นรา มจฺจุมุเข ลปนฺติ, \\ อวีตตณฺหาเส ภวาภเวสูติ.}
\csromangathagroup{\csromangathastanza{\csromanfolio{37}ปสฺสามิ โลเก ปริผนฺทมานํ, \\ ปชํ อิมํ ตณฺหคตํ ภเวสุ. \\ หีนา นรา มจฺจุมุเข ลปนฺติ, \\ อวีตตณฺหาเส ภวาภเวสูติ.}}

\proseitem{12}{\csromansymbolfootnote{*}{ขุ 1. 400 ปิฏฺเฐ.}\textbf{มมายิเต ปสฺสถ ผนฺทมาเน},}

\prose{\textbf{มจฺเฉว}\footnote{มจฺโฉว (สี)} \textbf{อปฺโปทเก ขิณโสเต.}}

\csromangathameasure{\textbf{เอตมฺปิ ทิสฺวา อมโม จเรยฺย,} \\ \textbf{ภเวสุ อาสตฺติมกุพฺพมาโน.}}
\csromangathagroup{\csromangathastanza{\textbf{เอตมฺปิ ทิสฺวา อมโม จเรยฺย,} \\ \textbf{ภเวสุ อาสตฺติมกุพฺพมาโน.}}}

\prose{\textbf{มมายิเต ปสฺสถ ผนฺทมาเน}ติ \textbf{มมตฺตา}ติ เทฺว \textbf{มมตฺตา} ตณฺหามมตฺตญฺจ ทิฏฺฐิมมตฺตญฺจ.\csromanspacer{} กตมํ ตณฺหามมตฺตํ\csromandash{}ยาวตา ตณฺหาสงฺขาเตน สีมกตํ มริยาทิกตํ โอธิกตํ ปริยนฺตกตํ ปริคฺคหิตํ มมายิตํ, อิทํ มมํ, เอตํ มมํ, เอตฺตกํ มมํ, เอตฺตาวตา มมํ, มม รูปา สทฺทา คนฺธา รสา โผฏฺฐพฺพา, อตฺถรณา ปาวุรณา ทาสิทาสา อเชฬกา กุกฺกุฏสูกรา หตฺถิควาสฺสวฬวา เขตฺตํ วตฺถุ หิรญฺญํ สุวณฺณํ คามนิคมราชธานิโย รฏฺฐญฺจ ชนปโท จ โกโส จ โกฏฺฐาคารญฺจ, เกวลมฺปิ มหาปถวิํ ตณฺหาวเสน มมายติ.\csromanspacer{} ยาวตา อฏฺฐสตํ ตณฺหาวิจริตํ, อิทํ ตณฺหามมตฺตํ.}

\prose{กตมํ ทิฏฺฐิมมตฺตํ\csromandash{}วีสติวตฺถุกา สกฺกายทิฏฺฐิ, ทสวตฺถุกา มิจฺฉาทิฏฺฐิ, ทสวตฺถุกา อนฺตคฺคาหิกา ทิฏฺฐิ, + ยา เอวรูปา ทิฏฺฐิ ทิฏฺฐิคตํ ทิฏฺฐิคหนํ ทิฏฺฐิกนฺตโร ทิฏฺฐิวิสูกายิกํ ทิฏฺฐิวิปฺผนฺทิกํ ทิฏฺฐิสญฺโญชนํ, คาโห ปฏิคฺคาโห อภินิเวโส ปรามาโส กุมฺมคฺโค มิจฺฉาปโถ มิจฺฉตฺตํ ติตฺถายตนํ วิปริเยสคฺคาโห วิปรีตคฺคาโห วิปลฺลาสคฺคาโห มิจฺฉาคาโห “อยาถาวกสฺมิํ ยาถาวกนฺ”ติ คาโห, ยาวตา ทฺวาสฏฺฐิทิฏฺฐิคตานิ, อิทํ ทิฏฺฐิมมตฺตํ.\csromanspacer{}\csromanspacer{}\textbf{มมายิเตปสฺสถ ผนฺทมาเน}ติ มมายิตํ วตฺถุํ อจฺเฉทสํกิโนปิ ผนฺทนฺติ, อจฺฉินฺทนฺเตปิ ผนฺทนฺติ, อจฺฉินฺเนปิ ผนฺทนฺติ, มมายิตํ วตฺถุํ วิปริณามสํกิโนปิ ผนฺทนฺติ, วิปริณามนฺเตปิ ผนฺทนฺติ, วิปริณเตปิ ผนฺทนฺติ ปผนฺทนฺติ สมฺผนฺทนฺติ วิปฺผนฺทนฺติ เวธนฺติ\footnote{เวเธนฺติ (สฺยา)} ปเวธนฺติ สมฺปเวธนฺติ.\csromanspacer{} เอวํ ผนฺทมาเน ปผนฺทมาเน สมฺผนฺทมาเน วิปฺผนฺทมาเน เวธมาเน ปเวธมาเน สมฺปเวธมาเน \csromanfolio{38}ปสฺสถ ทกฺขถ โอโลเกถ นิชฺฌายถ อุปปริกฺขถาติ มมายิเต ปสฺสถ ผนฺทมาเน.}

\prose{\textbf{มจฺเฉว อปฺโปทเก ขีณโสเต}ติ ยถา มจฺฉา อปฺโปทเก ปริตฺโตทเก อุทกปริยาทาเน กาเกหิ วา กุลเลหิ วา พลากาหิ วา ปริปาติยมานา อุกฺขิปิยมานา ขชฺชมานา ผนฺทนฺติ ปผนฺทนฺติ สมฺผนฺทนฺติ วิปฺผนฺทนฺติ เวธนฺติ ปเวธนฺติ สมฺปเวธนฺติ.\csromanspacer{} เอวเมวํ ปชา มมายิตํ วตฺถุํ อจฺเฉทสํกิโนปิ ผนฺทนฺติ, อจฺฉินฺทนฺเตปิ ผนฺทนฺติ, อจฺฉินฺเนปิ ผนฺทนฺติ, มมายิตํ วตฺถุํ วิปริณามสํกิโนปิ ผนฺทนฺติ, วิปริณามนฺเตปิ ผนฺทนฺติ, วิปริณเตปิ ผนฺทนฺติ ปผนฺทนฺติ สมฺผนฺทนฺติ วิปฺผนฺทนฺติ เวธนฺติ ปเวธนฺติ สมฺปเวธนฺตีติ มจฺเฉว อปฺโปทเก ขีณโสเต.}

\prose{\textbf{เอตมฺปิ ทิสฺวา อมโม จเรยฺยา}ติ เอตํ อาทีนวํ ทิสฺวา ปสฺสิตฺวา ตุลยิตฺวา ตีรยิตฺวา\footnote{ติรยิตฺวา (ก)} วิภาวยิตฺวา วิภูตํ กตฺวา มมตฺเตสูติ เอตมฺปิ ทิสฺวา.\csromanspacer{} \textbf{อมโม จเรยฺยา}ติ \textbf{มมตฺตา}ติ เทฺว \textbf{มมตฺตา} ตณฺหามมตฺตญฺจ ทิฏฺฐิมมตฺตญฺจ ฯเปฯ อิทํ ตณฺหามมตฺตํ ฯเปฯ อิทํ ทิฏฺฐิมมตฺตํ.\csromanspacer{} ตณฺหามมตฺตํ ปหาย ทิฏฺฐิมมตฺตํ ปฏินิสฺสชฺชิตฺวา จกฺขุํ อมมายนฺโต โสตํ อมมายนฺโต ฆานํ อมมายนฺโต ชิวฺหํ อมมายนฺโต กายํ อมมายนฺโต มนํ อมมายนฺโต, รูเป.\csromanspacer{} สทฺเท.\csromanspacer{} คนฺเธ.\csromanspacer{} รเส.\csromanspacer{} โผฏฺฐพฺเพ.\csromanspacer{} ธมฺเม.\csromanspacer{} กุลํ.\csromanspacer{} คณํ.\csromanspacer{} อาวาสํ.\csromanspacer{} ลาภํ.\csromanspacer{} ยสํ.\csromanspacer{} ปสํสํ.\csromanspacer{} สุขํ.\csromanspacer{} จีวรํ.\csromanspacer{} ปิณฺฑปาตํ.\csromanspacer{} เสนาสนํ.\csromanspacer{} คิลานปจฺจยเภสชฺชปริกฺขารํ.\csromanspacer{} กามธาตุํ.\csromanspacer{} รูปธาตุํ.\csromanspacer{} อรูปธาตุํ.\csromanspacer{} กามภวํ.\csromanspacer{} รูปภวํ.\csromanspacer{} อรูปภวํ.\csromanspacer{} สญฺญาภวํ.\csromanspacer{} อสญฺญาภวํ.\csromanspacer{} เนวสญฺญานาสญฺญาภวํ.\csromanspacer{} เอกโวการภวํ.\csromanspacer{} จตุโวการภวํ.\csromanspacer{} ปญฺจโวการภวํ.\csromanspacer{} อตีตํ.\csromanspacer{} อนาคตํ.\csromanspacer{} ปจฺจุปฺปนฺนํ.\csromanspacer{} ทิฏฺฐสุตมุตวิญฺญาตพฺเพ ธมฺเม อมมายนฺโต อคณฺหนฺโต อปรามสนฺโต อนภินิวิสนฺโต จเรยฺย วิหเรยฺย อิริเยยฺย วตฺเตยฺย ปาเลยฺย ยเปยฺย ยาเปยฺยาติ เอตมฺปิ ทิสฺวา อมโม จเรยฺย.}

\prose{\textbf{ภเวสุ อาสตฺติมกุพฺพมาโน}ติ \textbf{ภเวสู}ติ กามภเว รูปภเว อรูปภเว.\csromanspacer{} \textbf{อาสตฺติ} วุจฺจติ ตณฺหา.\csromanspacer{} โย ราโค สาราโค ฯเปฯ อภิชฺฌา โลโภ อกุสลมูลํ.\csromanspacer{} \textbf{ภเวสุ อาสตฺติมกุพฺพมาโน}ติ ภเวสุ อาสตฺติํ อกุพฺพมาโน, ฉนฺทํ เปมํ ราคํ ขนฺติํ อกุพฺพามาโน อชนยมาโน อสญฺชนยมาโน อนิพฺพตฺตยมาโน อนภินิพฺพตฺตยมาโนติ ภเวสุ อาสตฺติมกุพฺพมาโน.\csromanspacer{} เตนาห ภควา\csromandash{}}

\csromanheader{2}{2. คุหฏฺฐกสุตฺตนิทฺเทส}
\csromangathameasure{มมายิเต ปสฺสถ ผนฺทมาเน, \\ มจฺเฉว อปฺโปทเก ขีณโสเต. \\ เอตมฺปิ ทิสฺวา อมโม จเรยฺย, \\ ภเวสุ อาสตฺติมกุพฺพมาโนติ.}
\csromangathagroup{\csromangathastanza{\csromanfolio{39}มมายิเต ปสฺสถ ผนฺทมาเน, \\ มจฺเฉว อปฺโปทเก ขีณโสเต. \\ เอตมฺปิ ทิสฺวา อมโม จเรยฺย, \\ ภเวสุ อาสตฺติมกุพฺพมาโนติ.}}

\proseitem{13}{\csromansymbolfootnote{*}{ขุ 1. 400 ปิฏฺเฐ.}อุโภสุ อนฺเตสุ วิเนยฺย ฉนฺทํ,}

\prose{\textbf{ผสฺสํ ปริญฺญาย อนานุคิทฺโธ.}}

\csromangathameasure{\textbf{ยทตฺตครหี ตทกุพฺพมาโน,} \\ \textbf{น ลิมฺปตี} \textbf{ทิฏฺฐสุเตสุ ธีโร.}}
\csromangathagroup{\csromangathastanza{\textbf{ยทตฺตครหี ตทกุพฺพมาโน,} \\ \textbf{น ลิมฺปตี}\footnote{น ลิปฺปติ (สี)} \textbf{ทิฏฺฐสุเตสุ ธีโร.}}}

\prose{\textbf{อุโภสุ อนฺเตสุ วิเนยฺย ฉนฺทนฺ}ติ \textbf{อนฺตา}ติ \textbf{ผสฺโส} เอโก อนฺโต ผสฺสสมุทโย ทุติโย อนฺโต, อตีโต เอโก อนฺโต อนาคโต ทุติโย อนฺโต, สุขา เวทนา เอโก อนฺโต ทุกฺขา เวทนา ทุติโย อนฺโต, นามํ เอโก อนฺโต รูปํ ทุติโย อนฺโต, ฉ อชฺฌตฺติกานิ อายตนานิ เอโก อนฺโต ฉ พาหิรานิ อายตนานิ ทุติโย อนฺโต, สกฺกาโย เอโก อนฺโต สกฺกายสมุทโย ทุติโย อนฺโต.\csromanspacer{}\csromanspacer{}\textbf{ฉนฺโท}ติ โย กาเมสุ กามจฺฉนฺโท กามราโค กามนนฺที กามตณฺหา กามสฺเนโห กามปริฬาโห กามมุจฺฉา กามชฺโฌสานํ กาโมโฆ กามโยโค กามุปาทานํ กามจฺฉนฺทนีวรณํ.\csromanspacer{}\csromanspacer{}\textbf{อุโภสุ อนฺเตสุ วิเนยฺย ฉนฺทนฺ}ติ อุโภสุ อนฺเตสุ ฉนฺทํ วิเนยฺย ปฏิวิเนยฺย ปชเหยฺย วิโนเทยฺย พฺยนฺติํ กเรยฺย อนภาวํ คเมยฺยาติ อุโภสุ อนฺเตสุ วิเนยฺย ฉนฺทํ.}

\prose{\textbf{ผสฺสํ ปริญฺญาย อนานุคิทฺโธ}ติ \textbf{ผสฺโส}ติ จกฺขุสมฺ\textbf{ผสฺโส} โสตสมฺ\textbf{ผสฺโส} ฆานสมฺ\textbf{ผสฺโส} ชิวฺหาสมฺ\textbf{ผสฺโส} กายสมฺ\textbf{ผสฺโส} มโนสมฺ\textbf{ผสฺโส}, อธิวจนสมฺ\textbf{ผสฺโส} ปฏิฆสมฺ\textbf{ผสฺโส} สุขเวทนีโย สมฺ\textbf{ผสฺโส} ทุกฺขเวทนีโย สมฺ\textbf{ผสฺโส} อทุกฺขมสุขเวทนีโย สมฺ\textbf{ผสฺโส}, กุสโล \textbf{ผสฺโส} อกุสโล \textbf{ผสฺโส} อพฺยากโต \textbf{ผสฺโส}, กามาวจโร \textbf{ผสฺโส} รูปาวจโร \textbf{ผสฺโส} อรูปาวจโร \textbf{ผสฺโส}, สุญฺญโต \textbf{ผสฺโส} อนิมิตฺโต \textbf{ผสฺโส} อปฺปณิหิโต \textbf{ผสฺโส}, โลกิโย \textbf{ผสฺโส} โลกุตฺตโร \textbf{ผสฺโส}, อตีโต \textbf{ผสฺโส} อนาคโต \textbf{ผสฺโส} \csromanfolio{40}ปจฺจุปฺปนฺโน ผสฺโส, โย เอวรูโป ผสฺโส ผุสนา สมฺผุสนา สมฺผุสิตตฺตํ, อยํ วุจฺจติ ผสฺโส.}

\prose{\textbf{ผสฺสํ ปริญฺญายา}ติ ผสฺสํ ตีหิ ปริญฺญาหิ ปริชานิตฺวา ญาตปริญฺญาย ตีรณปริญฺญาย\footnote{ติรณปริญฺญาย (สฺยา)} ปหานปริญฺญาย. * กตมา ญาตปริญฺญา\csromandash{}ผสฺสํ ชานาติ “อยํ จกฺขุสมฺผสฺโส อยํ โสตสมฺผสฺโส อยํ ฆานสมฺผสฺโส อยํ ชิวฺหาสมฺผสฺโส อยํ กายสมฺผสฺโส อยํ มโนสมฺผสฺโส, อยํ อธิวจนสมฺผสฺโส อยํ ปฏิฆสมฺผสฺโส, อยํ สุขเวทนีโย ผสฺโส อยํ ทุกฺขเวทนีโย ผสฺโส อยํ อทุกฺขมสุขเวทนีโย ผสฺโส, อยํ กุสโล ผสฺโส อยํ อกุสโล ผสฺโส อยํ อพฺยากโต ผสฺโส, อยํ กามาวจโร ผสฺโส อยํ รูปาวจโร ผสฺโส อยํ อรูปาวจโร ผสฺโส, อยํ สุญฺญโต ผสฺโส อยํ อนิมิตฺโต ผสฺโส อยํ อปฺปณิหิโต ผสฺโส, อยํ โลกิโย ผสฺโส อยํ โลกุตฺตโร ผสฺโส, อยํ อตีโต ผสฺโส อยํ อนาคโต ผสฺโส อยํ ปจฺจุปฺปนฺโน ผสฺโส”ติ ชานาติ ปสฺสติ.\csromanspacer{} อยํ ญาตปริญฺญา.}

\prose{\csromansymbolfootnote{*}{ขุ 8. 38, 121 ปิฏฺเฐสุปิ.}กตมา ตีรณปริญฺญา\csromandash{}เอวํ ญาตํ กตฺวา ผสฺสํ ตีเรติ อนิจฺจโต ทุกฺขโต โรคโต คณฺฑโต สลฺลโต อฆโต อาพาธโต ปรโต ปโลกโต อีติโต อุปทฺทวโต ภยโต อุปสคฺคโต จลโต ปภงฺคุโต อธุวโต อตาณโต อเลณโต อสรณโต ริตฺตโต ตุจฺฉโต สุญฺญโต อนตฺตโต อาทีนวโต วิปริณามธมฺมโต อสารกโต อฆมูลโต วธกโต วิภวโต สาสวโต สงฺขตโต มารามิสโต ชาติชราพฺยาธิมรณธมฺมโต โสกปริเทวทุกฺขโทมนสฺสุปายาสธมฺมโต สํกิเลสธมฺมโต สมุทยโต อตฺถงฺคมโต อสฺสาทโต อาทีนวโต นิสฺสรณโต ตีเรติ.\csromanspacer{} อยํ ตีรณปริญฺญา.}

\prose{กตมา ปหานปริญฺญา\csromandash{}เอวํ ตีรยิตฺวา ผสฺเส ฉนฺทราคํ ปชหติ วิโนเทติ พฺยนฺติํ กโรติ อนภาวํ คเมติ.\csromanspacer{} วุตฺตํ เหตํ ภควตา\csromandash{} + โย ภิกฺขเว ผสฺเสสุ ฉนฺทราโค, ตํ ปชหถ.\csromanspacer{} เอวํ โส ผสฺโส ปหีโน ภวิสฺสติ อุจฺฉินฺนมูโล ตาลาวตฺถุกโต อนภาวํ กโต}

\vspace{\tipitakaparskip}
\csromancenter{คุหฏฺฐกสุตฺตนิทฺเทส อายติํ อนุปฺปาทธมฺโมติ.\csromanspacer{} อยํ ปหานปริญฺญา.\csromanspacer{}\csromanspacer{}\textbf{ผสฺสํ ปริญฺญายา}ติ ผสฺสํ อิมาหิ ตีหิ ปริญฺญาหิ ปริชานิตฺวา.\csromanspacer{} \textbf{อนานุคิทฺโธ}ติ \textbf{เคโธ} วุจฺจติ ตณฺหา.\csromanspacer{} โย ราโค สาราโค ฯเปฯ อภิชฺฌา โลโภ อกุสลมูลํ ยสฺเสโส \textbf{เคโธ} ปหีโน สมุจฺฉินฺโน วูปสนฺโต ปฏิปสฺสทฺโธ อภพฺพุปฺปตฺติโก ญาณคฺคินา ทฑฺโฒ, โส วุจฺจติ อคิทฺโธ.\csromanspacer{} โส รูเป อคิทฺโธ, สทฺเท อคิทฺโธ, คนฺเธ อคิทฺโธ, รเส อคิทฺโธ, โผฏฺฐพฺเพ อคิทฺโธ, กุเล.\csromanspacer{} คเณ.\csromanspacer{} อาวาเส.\csromanspacer{} ลาเภ.\csromanspacer{} ยเส.\csromanspacer{} ปสํสาย.\csromanspacer{} สุเข.\csromanspacer{} จีวเร.\csromanspacer{} ปิณฺฑปาเต.\csromanspacer{} เสนาสเน.\csromanspacer{} คิลานปจฺจยเภสชฺชปริกฺขาเร.\csromanspacer{} กามธาตุยา.\csromanspacer{} รูปธาตุยา.\csromanspacer{} อรูปธาตุยา.\csromanspacer{} กามภเว.\csromanspacer{} รูปภเว.\csromanspacer{} อรูปภเว.\csromanspacer{} สญฺญาภเว.\csromanspacer{} อสญฺญาภเว.\csromanspacer{} เนวสญฺญานาสญฺญาภเว.\csromanspacer{} เอกโวการภเว.\csromanspacer{} จตุโวการภเว.\csromanspacer{} ปญฺจโวการภเว.\csromanspacer{} อตีเต.\csromanspacer{} อนาคเต.\csromanspacer{} ปจฺจุปฺปนฺเน.\csromanspacer{} ทิฏฺฐสุตมุตวิญฺญาตพฺเพสุ ธมฺเมสุ อคิทฺโธ อคธิโต อมุจฺฉิโต อนชฺโฌสนฺโน, วีต\textbf{เคโธ} วิคต\textbf{เคโธ} จตฺต\textbf{เคโธ} วนฺต\textbf{เคโธ} มุตฺต\textbf{เคโธ}, ปหีน\textbf{เคโธ} ปฏินิสฺสฏฺฐ\textbf{เคโธ}, วีตราโค วิคตราโค จตฺตราโค วนฺตราโค มุตฺตราโค ปหีนราโค ปฏินิสฺสฏฺฐราโค, นิจฺฉาโต นิพฺพุโต สีติภูโต สุขปฏิสํเวที พฺรหฺมภูเตน อตฺตนา วิหรตีติ ผสฺสํ ปริญฺญาย อนานุคิทฺโธ.}
\prose{\csromanfolio{41}\textbf{ยทตฺตครหี ตทกุพฺพมาโน}ติ \textbf{ยทิ}ติ ยํ.\csromanspacer{} \textbf{อตฺตครหี}ติ ทฺวีหิ การเณหิ อตฺตานํ ครหติ กตตฺตา จ อกตตฺตา จ.\csromanspacer{} กถํ กตตฺตา จ อกตตฺตา จ อตฺตานํ ครหติ\csromandash{}กตํ เม กายทุจฺจริตํ, อกตํ เม กายสุจริตนฺติ อตฺตานํ ครหติ.\csromanspacer{} กตํ เม วจีทุจฺจริตํ, อกตํ เม วจีสุจริตนฺติ อตฺตานํ ครหติ.\csromanspacer{} กตํ เม มโนทุจฺจริตํ, อกตํ เม มโนสุจริตนฺติ อตฺตานํ ครหติ.\csromanspacer{} กโต เม ปาณาติปาโต, อกตา เม ปาณาติปาตา เวรมณีติ อตฺตานํ ครหติ.\csromanspacer{} กตํ เม อทินฺนาทานํ, อกตา เม อทินฺนาทานา เวรมณีติ อตฺตานํ ครหติ.\csromanspacer{} กโต เม กาเมสุมิจฺฉาจาโร, อกตา เม กาเมสุมิจฺฉาจารา เวรมณีติ อตฺตานํ ครหติ.\csromanspacer{} กโต เม มุสาวาโท, อกตา เม มุสาวาทา เวรมณีติ อตฺตานํ ครหติ.\csromanspacer{} กตา เม ปิสุณา วาจา, อกตา เม ปิสุณาย วาจาย เวรมณีติ อตฺตานํ ครหติ.\csromanspacer{} กตา เม ผรุสา วาจา, อกตา เม ผรุสาย วาจาย เวรมณีติ อตฺตานํ ครหติ.\csromanspacer{} กโต เม สมฺผปฺปลาโป, อกตา เม สมฺผปฺปลาปา เวรมณีติ \csromanfolio{42}อตฺตานํ ครหติ.\csromanspacer{} กตา เม อภิชฺฌา, อกตา เม อนภิชฺฌาติ อตฺตานํ ครหติ.\csromanspacer{} กโต เม พฺยาปาโท, อกโต เม อพฺยาปาโทติ อตฺตานํ ครหติ.\csromanspacer{} กตา เม มิจฺฉาทิฏฺฐิ, อกตา เม สมฺมาทิฏฺฐีติ อตฺตานํ ครหติ.\csromanspacer{} เอวํ กตตฺตา จ อกตตฺตา จ อตฺตานํ ครหติ.\csromanspacer{} อถ วา สีเลสุ’มฺหิ น ปริปูรการีติ อตฺตานํ ครหติ.\csromanspacer{} อินฺทฺริเยสุ’มฺหิ อคุตฺตทฺวาโรติ อตฺตานํ ครหติ, โภชเน’มฺหิ\footnote{โภชเน (สฺยา)} อมตฺตญฺญูติ อตฺตานํ ครหติ, ชาคริยํ อนนุยุตฺโตติ อตฺตานํ ครหติ, สติสมฺปชญฺเญน อสมนฺนาคโตติ อตฺตานํ ครหติ, อภาวิตา เม จตฺตาโร สติปฏฺฐานาติ อตฺตานํ ครหติ.\csromanspacer{} อภาวิตา เม จตฺตาโร สมฺมปฺปธานาติ อตฺตานํ ครหติ, อภาวิตา เม จตฺตาโร อิทฺธิปาทาติ อตฺตานํ ครหติ, อภาวิตานิ เม ปญฺจินฺทฺริยานีติ อตฺตานํ ครหติ, อภาวิตานิ เม ปญฺจ พลานีติ อตฺตานํ ครหติ, อภาวิตา เม สตฺต โพชฺฌงฺคาติ อตฺตานํ ครหติ, อภาวิโต เม อริโย อฏฺฐงฺคิโก มคฺโคติ อตฺตานํ ครหติ.\csromanspacer{} ทุกฺขํ เม อปริญฺญาตนฺติ อตฺตานํ ครหติ, สมุทโย เม อปฺปหีโนติ อตฺตานํ ครหติ, มคฺโค เม อภาวิโตติ อตฺตานํ ครหติ, นิโรโธ เม อสจฺฉิกโตติ อตฺตานํ ครหติ.\csromanspacer{} เอวํ กตตฺตา จ อกตตฺตา จ อตฺตานํ ครหติ.\csromanspacer{}\csromanspacer{}เอวํ อตฺตครหิตํ กมฺมํ อกุพฺพมาโน อชนยมาโน อสญฺชนยมาโน อนิพฺพตฺตยมาโน อนภินิพฺพตฺตยมาโนติ ยทตฺตครหี ตทกุพฺพมาโน.\csromanspacer{} \textbf{น ลิมฺปตี ทิฏฺฐสุเตสุ ธีโร}ติ \textbf{เลโป}ติ เทฺว เลปา ตณฺหา\textbf{เลโป} จ ทิฏฺฐิ\textbf{เลโป} จ ฯเปฯ อยํ ตณฺหา\textbf{เลโป} ฯเปฯ อยํ ทิฏฺฐิ\textbf{เลโป}.\csromanspacer{} \textbf{ธีโร}ติ ปณฺฑิโต ปญฺญวา พุทฺธิมา ญาณี วิภาวี เมธาวี, ธีโร ตณฺหาเลปํ ปหาย ทิฏฺฐิเลปํ ปฏินิสฺสชฺชิตฺวา ทิฏฺเฐ น ลิมฺปติ สุเต น ลิมฺปติ มุเต น ลิมฺปติ วิญฺญาเต น ลิมฺปติ น ปลิมฺปติ\footnote{น สํลิมฺปติ (สฺยา)} น อุปลิมฺปติ, อลิตฺโต อปลิตฺโต\footnote{อสํลิตฺโต (สฺยา)} อนุปลิตฺโต นิกฺขนฺโต นิสฺสโฏ วิปฺปมุตฺโต วิสญฺญุตฺโต วิมริยาทิกเตน เจตสา วิหรตีติ น ลิมฺปตี ทิฏฺฐสุเตสุ ธีโรติ.\csromanspacer{} เตนาห ภควา\csromandash{}}

\csromangathameasure{อุโภสุ อนฺเตสุ วิเนยฺย ฉนฺทํ, \\ ผสฺสํ ปริญฺญาย อนานุคิทฺโธ. \\ ยทตฺตครหี ตทกุพฺพมาโน, \\ \textbf{น ลิมฺปตี ทิฏฺฐสุเตสุ ธีโร}ติ.}
\csromangathagroup{\csromangathastanza{อุโภสุ อนฺเตสุ วิเนยฺย ฉนฺทํ, \\ ผสฺสํ ปริญฺญาย อนานุคิทฺโธ. \\ ยทตฺตครหี ตทกุพฺพมาโน, \\ \textbf{น ลิมฺปตี ทิฏฺฐสุเตสุ ธีโร}ติ.}}

\csromanheader{2}{2. คุหฏฺฐกสุตฺตนิทฺเทส}
\proseitem{14}{\csromanfolio{43}\csromansymbolfootnote{*}{ขุ 1. 400 ปิฏฺเฐปิ.}\textbf{สญฺญํ ปริญฺญา วิตเรยฺย โอฆ}ํ,}

\prose{\textbf{ปริคฺคเหสุ มุนิโนปลิตฺโต}.}

\csromangathameasure{\textbf{อพฺพูฬฺหสลฺโล จร’มปฺปมตฺโต,} \\ \textbf{นาสีสตี โลกมิมํ ปรญฺจ.}}
\csromangathagroup{\csromangathastanza{\textbf{อพฺพูฬฺหสลฺโล จร’มปฺปมตฺโต,} \\ \textbf{นาสีสตี โลกมิมํ ปรญฺจ.}}}

\prose{\textbf{สญฺญํ ปริญฺญา วิตเรยฺย โอฆ}นฺติ \textbf{สญฺญา}ติ กาม\textbf{สญฺญา} พฺยาปาท\textbf{สญฺญา} วิหิํสา\textbf{สญฺญา}, เนกฺขมฺม\textbf{สญฺญา} อพฺยาปาท\textbf{สญฺญา} อวิหิํสา\textbf{สญฺญา}, รูป\textbf{สญฺญา} สทฺท\textbf{สญฺญา} คนฺธ\textbf{สญฺญา} รส\textbf{สญฺญา} โผฏฺฐพฺพ\textbf{สญฺญา} ธมฺม\textbf{สญฺญา}, ยา เอวรูปา \textbf{สญฺญา} สญฺจนนา สญฺจานิตตฺตํ, อยํ วุจฺจติ \textbf{สญฺญา}.\csromanspacer{}\csromanspacer{}\textbf{สญฺญํ ปริญฺญา}ติ สญฺญํ ตีหิ ปริญฺญาหิ ปริชานิตฺวา ญาตปริญฺญาย ตีรณปริญฺญาย ปหานปริญฺญาย.\csromanspacer{}\csromanspacer{}กตมา ญาตปริญฺญา\csromandash{}สญฺญํ ชานาติ “อยํ กาม\textbf{สญฺญา} อยํ พฺยาปาท\textbf{สญฺญา} อยํ วิหิํสา\textbf{สญฺญา}, อยํ เนกฺขมฺม\textbf{สญฺญา} อยํ อพฺยาปาท\textbf{สญฺญา} อยํ อวิหิํสา\textbf{สญฺญา}, อยํ รูป\textbf{สญฺญา} อยํ สทฺท\textbf{สญฺญา} อยํ คนฺธ\textbf{สญฺญา} อยํ รส\textbf{สญฺญา} อยํ โผฏฺฐพฺพ\textbf{สญฺญา} อยํ ธมฺม\textbf{สญฺญา}”ติ ชานาติ ปสฺสติ, อยํ ญาตปริญฺญา.}

\prose{กตมา ตีรณปริญฺญา\csromandash{}เอวํ ญาตํ กตฺวา สญฺญํ ตีเรติ อนิจฺจโต ทุกฺขโต โรคโต คณฺฑโต สลฺลโต อฆโต อาพาธโต ปรโต ปโลกโต อีติโต อุปทฺทวโต ภยโต อุปสคฺคโต จลโต ปภงฺคุโต ฯเปฯ สมุทยโต อตฺถงฺคมโต อสฺสาทโต อาทีนวโต นิสฺสรณโต ตีเรติ, อยํ ตีรณปริญฺญา.}

\prose{กตมา ปหานปริญฺญา\csromandash{}เอวํ ตีรยิตฺวา \textbf{สญฺญา}ย ฉนฺทราคํ ปชหติ วิโนเทติ อนภาวํ คเมติ.\csromanspacer{} วุตฺตมฺปิ เหตํ ภควตา\csromandash{} + โย ภิกฺขเว \textbf{สญฺญา}ย ฉนฺทราโค, ตํ ปชหถ.\csromanspacer{} เอวํ สา \textbf{สญฺญา} ปหีนา ภวิสฺสติ อุจฺฉินฺนมูลา ตาลาวตฺถุกตา อนภาวํ กตา อายติํ อนุปฺปาทธมฺมาติ, อยํ ปหานปริญฺญา.\csromanspacer{}\csromanspacer{}\textbf{สญฺญํ ปริญฺญา}ติ สญฺญํ อิมาหิ ตีหิ ปริญฺญาหิ ปริชานิตฺวา.\csromanspacer{}\csromanspacer{}\textbf{วิตเรยฺย โอฆนฺ}ติ กาโมฆํ ภโวฆํ ทิฏฺโฐฆํ อวิชฺโชฆํ ตเรยฺย อุตฺตเรยฺย ปตเรยฺย สมติกฺกเมยฺย วีติวตฺเตยฺยาติ สญฺญํ ปริญฺญา วิตเรยฺย โอฆํ.}

\prose{\textbf{ปริคฺคเหสุ มุนิ โนปลิตฺโต}ติ \textbf{ปริคฺคหา}ติ เทฺว \textbf{ปริคฺคหา} ตณฺหาปริคฺคโห จ ทิฏฺฐิปริคฺคโห จ ฯเปฯ อยํ ตณฺหาปริคฺคโห ฯเปฯ อยํ ทิฏฺฐิปริคฺคโห.\csromanspacer{}\csromanspacer{}\textbf{มุนี}ติ โมนํ วุจฺจติ ญาณํ.\csromanspacer{} ยา ปญฺญา ปชานนา อโมโห \csromanfolio{44}ธมฺมวิจโย สมฺมาทิฏฺฐิ, เตน ญาเณน สมนฺนาคโต มุนิ โมนปฺปตฺโตติ.\csromanspacer{} ตีณิ โมเนยฺยานิ กายโมเนยฺยํ วจีโมเนยฺยํ มโนโมเนยฺยํ.}

\prose{กตมํ กายโมเนยฺยํ\csromandash{}ติวิธานํ กายทุจฺจริตานํ ปหานํ กายโมเนยฺยํ, ติวิธํ กายสุจริตํ กายโมเนยฺยํ, กายารมฺมเณ ญาณํ กายโมเนยฺยํ, กายปริญฺญา กายโมเนยฺยํ, ปริญฺญาสหคโต มคฺโค กายโมเนยฺยํ, กาเย ฉนฺทราคสฺส ปหานํ กายโมเนยฺยํ, กายสงฺขารนิโรโธ จตุตฺถชฺฌานสมาปตฺติ กายโมเนยฺยํ, อิทํ กายโมเนยฺยํ.}

\prose{กตมํ วจีโมเนยฺยํ\csromandash{}จตุพฺพิธานํ วจีทุจฺจริตานํ ปหานํ วจีโมเนยฺยํ, จตุพฺพิธํ วจีสุจริตํ วจีโมเนยฺยํ, วาจารมฺมเณ ญาณํ วจีโมเนยฺยํ, วาจาปริญฺญา วจีโมเนยฺยํ, ปริญฺญาสหคโต มคฺโค วจีโมเนยฺยํ, วาจาย ฉนฺทราคสฺส ปหานํ วจีโมเนยฺยํ, วจีสงฺขารนิโรโธ ทุติยชฺฌานสมาปตฺติ วจีโมเนยฺยํ, อิทํ วจีโมเนยฺยํ.}

\prose{กตมํ มโนโมเนยฺยํ\csromandash{}ติวิธานํ มโนทุจฺจริตานํ ปหานํ มโนโมเนยฺยํ, ติวิธํ มโนสุจริยํ มโนโมเนยฺยํ, จิตฺตารมฺมเณ ญาณํ มโนโมเนยฺยํ, จิตฺตปริญฺญา มโนโมเนยฺยํ, ปริญฺญาสหคโต มคฺโค มโนโมเนยฺยํ, จิตฺเต ฉนฺทราคสฺส ปหานํ มโนโมเนยฺยํ, จิตฺตสงฺขารนิโรโธ สญฺญาเวทยิตนิโรธํ มโนโมเนยฺยํ, อิทํ มโนโมเนยฺยํ.}

\csromangathameasure{กายมุนิํ วาจามุนิํ, มโนมุนิมนาสวํ. \\ มุนิํ โมเนยฺยสมฺปนฺนํ, อาหุ สพฺพปฺปหายินํ. \\ กายมุนิํ วาจามุนิํ, มโนมุนิมนาสวํ. \\ มุนิํ โมเนยฺยสมฺปนฺนํ, อาหุ นินฺหาตปาปกนฺติ.}
\csromangathagroup{\csromangathastanza{กายมุนิํ วาจามุนิํ, มโนมุนิมนาสวํ. \\ มุนิํ โมเนยฺยสมฺปนฺนํ, อาหุ สพฺพปฺปหายินํ.}\csromangathastanza{กายมุนิํ วาจามุนิํ, มโนมุนิมนาสวํ. \\ มุนิํ โมเนยฺยสมฺปนฺนํ, อาหุ นินฺหาตปาปกนฺติ\footnote{นิํนฺหาตปาปกนฺติ (สฺยา) อํ 1. 277; ขุ 1. 234; ขุ 8. 73 ปิฏฺเฐสุปิ.}.}}

\prose{อิเมหิ ตีหิ โมเนยฺเยหิ ธมฺเมหิ สมนฺนาคตา ฉ มุนิโน\footnote{ฉ มุนโย (สฺยา)} อคารมุนิโน อนคารมุนิโน เสขมุนิโน อเสขมุนิโน ปจฺเจกมุนิโน มุนิมุนิโนติ.\csromanspacer{} กตเม อคารมุนิโน\csromandash{}เย เต อคาริกา ทิฏฺฐปทา วิญฺญาตสาสนา, อิเม อคารมุนิโน.\csromanspacer{}\csromanspacer{}กตเม อนคารมุนิโน\csromandash{}เย เต ปพฺพชิตา ทิฏฺฐปทา วิญฺญาตสาสนา, อิเม อนคารมุนิโน.\csromanspacer{}\csromanspacer{}สตฺต เสขา เสขมุนิโน.\csromanspacer{}\csromanspacer{}อรหนฺโต อเสขมุนิโน.\csromanspacer{}\csromanspacer{}ปจฺเจกพุทฺธา ปจฺเจกมุนิโน.\csromanspacer{} มุนิมุนิโน วุจฺจนฺติ ตถาคตา อรหนฺโต สมฺมาสมฺพุทฺธา.}

\csromanheader{2}{2. คุหฏฺฐกสุตฺตนิทฺเทส}
\csromangathameasure{*น โมเนน มุนี โหติ, มูฬฺหรูโป อวิทฺทสุ. \\ โย จ ตุลํว ปคฺคยฺห, วรมาทาย ปณฺฑิโต. \\ ปาปานิ ปริวชฺเชติ, ส มุนี เตน โส มุนิ. \\ โย มุนาติ อุโภ โลเก, “มุนิ” เตน ปวุจฺจติ. \\ อสตญฺจ สตญฺจ ญตฺวา ธมฺมํ, \\ อชฺฌตฺตํ พหิทฺธา จ สพฺพโลเก. \\ เทวมนุสฺเสหิ ปูชิโต โย, \\ สงฺคชาลมติจฺจ โส มุนีติ.}
\csromangathagroup{\csromangathastanza{\csromanfolio{45}\csromansymbolfootnote{*}{ขุ 1. 51; ขุ 8. 74 ปิฏฺเฐสุปิ.}น โมเนน มุนี โหติ, มูฬฺหรูโป อวิทฺทสุ. \\ โย จ ตุลํว ปคฺคยฺห, วรมาทาย ปณฺฑิโต.}\csromangathastanza{ปาปานิ ปริวชฺเชติ, ส มุนี เตน โส มุนิ. \\ โย มุนาติ อุโภ โลเก, “มุนิ” เตน ปวุจฺจติ.}\csromangathastanza{อสตญฺจ สตญฺจ ญตฺวา ธมฺมํ, \\ อชฺฌตฺตํ พหิทฺธา จ สพฺพโลเก. \\ เทวมนุสฺเสหิ ปูชิโต โย, \\ สงฺคชาลมติจฺจ โส มุนีติ.}}

\prose{\textbf{เลปา}ติ เทฺว เลปา ตณฺหาเลโป จ ทิฏฺฐิเลโป จ ฯเปฯ อยํ ตณฺหาเลโป ฯเปฯ อยํ ทิฏฺฐิเลโป.\csromanspacer{} มุนิ ตณฺหาเลปํ ปหาย ทิฏฺฐิเลปํ ปฏินิสฺสชฺชิตฺวา ปริคฺคเหสุ น ลิมฺปติ น ปลิมฺปติ น อุปลิมฺปติ, อลิตฺโต อปลิตฺโต อนุปลิตฺโต นิกฺขนฺโต นิสฺสโฏ วิปฺปมุตฺโต วิสญฺญุตฺโต, วิมริยาทิกเตน เจตสา วิหรตีติ ปริคฺคเหสุ มุนิ โนปลิตฺโต.}

\prose{\textbf{อพฺพูฬฺหสลฺโล จร’มปฺปมตฺโต}ติ \textbf{สลฺลนฺ}ติ สตฺต สลฺลานิ ราคสลฺลํ โทสสลฺลํ โมหสลฺลํ มานสลฺลํ ทิฏฺฐิสลฺลํ โสกสลฺลํ กตํกถาสลฺลํ\footnote{ทุจฺจริตสลฺลํ (สี)}, ยสฺเสเต สลฺลา ปหีนา สมุจฺฉินฺนา วูปสนฺตา ปฏิปสฺสทฺธา อภพฺพุปฺปตฺติกา ญาณคฺคินา ทฑฺฒา, โส วุจฺจติ อพฺพูฬฺหสลฺโล อพฺพหิตสลฺโล อุทฺธตสลฺโล สมุทฺธตสลฺโล อุปฺปาฏิตสลฺโล สมุปฺปาฏิตสลฺโล จตฺตสลฺโล วนฺตสลฺโล มุตฺตสลฺโล ปหีนสลฺโล ปฏินิสฺสฏฺฐสลฺโล นิจฺฉาโต นิพฺพุโต สีติภูโต สุขปฏิสํเวที พฺรหฺมภูเตน อตฺตนา วิหรตีติ อพฺพูฬฺหสลฺโล.}

\prose{\textbf{จรนฺ}ติ จรนฺโต วิหรนฺโต อิริยนฺโต วตฺตนฺโต ปาเลนฺโต ยเปนฺโต ยาเปนฺโต.\csromanspacer{}\csromanspacer{}\textbf{อปฺปมตฺโต}ติ + สกฺกจฺจการี สาตจฺจการี อฏฺฐิตการี อโนลีนวุตฺติโก อนิกฺขิตฺตจฺฉนฺโท อนิกฺขิตฺตธุโร กุสเลสุ ธมฺเมสุ “กถาหํ อปริปูรํ วา สีลกฺขนฺธํ ปริปูเรยฺยํ, ปริปูรํ วา สีลกฺขนฺธํ ตตฺถ ตตฺถ ปญฺญาย อนุคฺคเณฺหยฺยนฺ”ติ โย ตตฺถ ฉนฺโท จ วายาโม จ อุสฺสาโห จ อุสฺเสฬฺหี จ อปฺปฏิวานี จ สติ จ สมฺปชญฺญญฺจ อาตปฺปํ ปธานํ อธิฏฺฐานํ อนุโยโค อปฺปมาโท กุสเลสุ ธมฺเมสุ.\csromanspacer{}\csromanspacer{}กถาหํ อปริปูรํ วา สมาธิกฺขนฺธํ ปริปูเรยฺยํ, ปริปูรํ วา สมาธิกฺขนฺธํ ตตฺถ ตตฺถ}

\prose{\csromanfolio{46}ปญฺญาย อนุคฺคเณยฺยํ ฯเปฯ กุสเลสุ ธมฺเมสุ.\csromanspacer{} กถาหํ อปริปูรํ วา ปญฺญากฺขนฺธํ ปริปูเรยฺยํ, วิมุตฺติกฺขนฺธํ.\csromanspacer{} วิมุตฺติญาณทสฺสนกฺขนฺธํ ปริปูเรยฺยํ.\csromanspacer{} ปริปูรํ วา วิมุตฺติญาณทสฺสนกฺขนฺธํ ตตฺถ ตตฺถ ปญฺญาย อนุคฺคเณฺหยฺยนฺติ โย ตตฺถ ฉนฺโท จ วายาโม จ อุสฺสาโห จ อุสฺเสฬฺหี จ อปฺปฏิวานี จ สติ จ สมฺปชญฺญญฺจ อาตปฺปํ ปธานํ อธิฏฺฐานํ อนุโยโค อปฺปมาโท กุสเลสุ ธมฺเมสุ.\csromanspacer{}\csromanspacer{}“กถาหํ อปริญฺญาตํ วา ทุกฺขํ ปริชาเนยฺยํ, อปฺปหีเน วา กิเลเส ปชเหยฺยํ, อภาวิตํ วา มคฺคํ ภาเวยฺยํ.\csromanspacer{} อสจฺฉิกตํ วา นิโรธํ สจฺฉิกเรยฺยนฺ”ติ โย ตตฺถ ฉนฺโท จ วายาโม จ อุสฺสาโห จ อุสฺโสฬฺหี จ อปฺปฏิวานี จ สติ จ สมฺปชญฺญญฺจ อาตปฺปํ ปธานํ อธิฏฺฐานํ อนุโยโค อปฺปมาโท กุสเลสุ ธมฺเมสูติ อพฺพูฬฺหสลฺโล จร’มปฺปมตฺโต.}

\prose{\textbf{นาสีสตี โลกมิมํ ปรญฺจา}ติ อิมํ โลกํ นาสีสติ สกตฺตภาวํ, ปรโลกํ นาสีสติ, ปรตฺตภาวํ.\csromanspacer{} อิมํ โลกํ นาสีสติ สกรูปเวทนาสญฺญาสงฺขารวิญฺญาณํ, ปรํ โลกํ นาสีสติ ปรรูปเวทนาสญฺญาสงฺขารวิญฺญาณํ.\csromanspacer{} อิมํ โลกํ นาสีสติ ฉ อชฺฌตฺติกานิ อายตนานิ, ปรํ โลกํ นาสีสติ ฉ พาหิรานิ อายตนานิ.\csromanspacer{} อิมํ โลกํ นาสีสติ มนุสฺสโลกํ, ปรํ โลกํ นาสีสติ เทวโลกํ.\csromanspacer{} อิมํ โลกํ นาสีสติ กามธาตุํ, ปรํ โลกํ นาสีสติ รูปธาตุํ อรูปธาตุํ.\csromanspacer{} อิมํ โลกํ นาสีสติ กามธาตุํ รูปธาตุํ, ปรํ โลกํ นาสีสติ อรูปธาตุํ.\csromanspacer{} ปุน คติํ วา อุปปตฺติํ วา ปฏิสนฺธิํ วา ภวํ วา สํสารํ วา วฏฺฏํ วา นาสีสติ น อิจฺฉติ น สาทิยติ น ปตฺเถติ น ปิเหติ นาติชปฺปตีติ นาสีสตี โลกมิมํ ปรญฺจาติ.\csromanspacer{} เตนาห ภควา\csromandash{}}

\csromangathameasure{สญฺญํ ปริญฺญา วิตเรยฺย โอฆํ, \\ ปริคฺคเหสุ มุนิ โนปลิตฺโต. \\ อพฺพูฬฺหสลฺโล จร’มปฺปมตฺโต,}
\csromangathagroup{\csromangathastanza{สญฺญํ ปริญฺญา วิตเรยฺย โอฆํ, \\ ปริคฺคเหสุ มุนิ โนปลิตฺโต. \\ อพฺพูฬฺหสลฺโล จร’มปฺปมตฺโต,}}

\prose{\textbf{นาสีสตี โลกมิมํ ปรญฺจา}ติ.}

\vspace{\tipitakaparskip}
\csromancenter{คุหฏฺฐกสุตฺตนิทฺเทโส ทุติโย.}
\csromanheader{2}{3. ทุฏฺฐฏฺฐกสุตฺตนิทฺเทส \textbf{3. ทุฏฺฐฏฺฐกสุตฺตนิทฺเทส}}
\csromanheader{2}{อถ ทุฏฺฐฏฺฐกสุตฺตนิทฺเทสํ วกฺขติ\csromandash{}}
\proseitem{15}{\csromanfolio{47}\csromansymbolfootnote{*}{ขุ 1. 401 ปิฏฺเฐปิ.}\textbf{วทนฺติ เว ทุฏฺฐมนาปิ เอเก},}

\prose{\textbf{อโถปิ}\footnote{อญฺเญปิ เต (สี), อญฺเญปิ (สฺยา)} \textbf{เว สจฺจมนา วทนฺติ.}}

\csromangathameasure{\textbf{วาทญฺจ ชาตํ มุนิ โน อุเปติ,} \\ \textbf{ตสฺมา มุนี นตฺถิ ขิโล กุหิญฺจิ.}}
\csromangathagroup{\csromangathastanza{\textbf{วาทญฺจ ชาตํ มุนิ โน อุเปติ,} \\ \textbf{ตสฺมา มุนี นตฺถิ ขิโล กุหิญฺจิ.}}}

\prose{\textbf{วทนฺติ เว ทุฏฺฐมนาปิ เอเก}ติ เต ติตฺถิยา ทุฏฺฐมนา ปทุฏฺฐมนา วิรุทฺธมนา ปฏิวิรุทฺธมนา อาหตมนา ปจฺจาหตมนา อาฆาติตมนา ปจฺจาฆาติตมนา วทนฺติ อุปวทนฺติ ภควนฺตญฺจ ภิกฺขุสํฆญฺจ อภูเตนาติ วทนฺติ เว ทุฏฺฐมนาปิ เอเก.}

\prose{\textbf{อโถปิ เว สจฺจมนา วทนฺตี}ติ เย เตสํ ติตฺถิยานํ สทฺทหนฺตา โอกปฺเปนฺตา อธิมุจฺจนฺตา สจฺจมนา สจฺจสญฺญิโน ภูตมนา ภูตสญฺญิโน ตถมนา ตถสญฺญิโน ยาถาวมนา ยาถาวสญฺญิโน, อวิปรีตมนา อวิปรีตสญฺญิโน วทนฺติ อุปวทนฺติ ภควนฺตญฺจ ภิกฺขุสํฆญฺจ อภูเตนาติ อโถปิ \textbf{เว สจฺจมนา วทนฺติ.}}

\prose{\textbf{วาทญฺจ ชาตํ มุนิ โน อุเปตี}ติ โส วาโท ชาโต โหติ สญฺชาโต นิพฺพตฺโต อภินิพฺพตฺโต ปาตุภูโต ปรโตโฆโส อกฺโกโส อุปวาโท ภควโต จ ภิกฺขุสํฆสฺส จ อภูเตนาติ วาทญฺจ ชาตํ.\csromanspacer{}\csromanspacer{}\textbf{มุนิ โน อุเปตี}ติ \textbf{มุนี}ติ โมนํ วุจฺจติ ญาณํ.\csromanspacer{} ยา ปญฺญา ปชานนา อโมโห ธมฺมวิจโย สมฺมาทิฏฺฐิ เตน ญาเณน สมนฺนาคโต มุนิ โมนปฺปตฺโต ฯเปฯ สงฺคชาลมติจฺจ โส มุนิ.\csromanspacer{}\csromanspacer{}โย วาทํ อุเปติ, โส ทฺวีหิ การเณหิ วาทํ อุเปติ การโก การกตาย วาทํ อุเปติ.\csromanspacer{}\csromanspacer{}อถ วา วุจฺจมาโน อุปวทิยมาโน กุปฺปติ พฺยาปชฺชติ ปติฏฺฐิยติ, โกปญฺจ โทสญฺจ อปฺปจฺจยญฺจ ปาตุกโรติ อการโกมฺหีติ โย วาทํ อุเปติ, โส อิเมหิ ทฺวีหิ การเณหิ วาทํ อุเปติ.\csromanspacer{} มุนิ ทฺวีหิ การเณหิ วาทํ น อุเปติ.\csromanspacer{} อการโก มุนิ อการกตาย วาทํ น อุเปติ.\csromanspacer{}\csromanspacer{}อถ วา วุจฺจมาโน อุปวทิยมาโน น กุปฺปติ น พฺยาปชฺชติ น ปติตฺถียติ\footnote{ปติฏฺฐิยติ (สี, สฺยา) อํ 1. 122 ปิฏฺเฐปิ.}, น โกปญฺจ โทสญฺจ อปฺปจฺจยญฺจ ปาตุกโรติ อการโกมฺหีติ มุนิ อิเมหิ ทฺวีหิ การเณหิ วาทํ น อุเปติ น \csromanfolio{48}อุปคจฺฉติ น คณฺหาติ น ปรามสติ น อภินิวิสตีติ วาทญฺจ ชาตํ มุนิ โน อุเปติ.}

\prose{\textbf{ตสฺมา มุนี นตฺถิ ขิโล กุหิญฺจี}ติ \textbf{ตสฺมา}ติ \textbf{ตสฺมา} ตํการณา ตํเหตุ ตปฺปจฺจยา ตํนิทานํ มุนิโน อาหตจิตฺตตา ขิลชาตตาปิ นตฺถิ.\csromanspacer{} ปญฺจปิ เจโตขิลา นตฺถิ, ตโยปิ ขิลา นตฺถิ, ราคขิโล โทสขิโล โมหขิโล นตฺถิ น สนฺติ น สํวิชฺชติ นุปลพฺภติ ปหีโน สมุจฺฉินฺโน วูปสนฺโต ปฏิปสฺสทฺโธ อภพฺพุปฺปตฺติโก ญาณคฺคินา ทฑฺโฒ.\csromanspacer{}\csromanspacer{}\textbf{กุหิญฺจี}ติ กุหิญฺจิ กิมฺหิจิ กตฺถจิ อชฺฌตฺตํ วา พหิทฺธา วา อชฺฌตฺตพหิทฺธา วาติ \textbf{ตสฺมา} มุนิ นตฺถิ ขิโล กุหิญฺจิติ.\csromanspacer{} เตนาห ภควา\csromandash{}}

\csromangathameasure{วทนฺติ เว ทุฏฺฐมนาปิ เอเก, \\ อโถปิ เว สจฺจมนา วทนฺติ. \\ วาทญฺจ ชาตํ มุนิ โน อุเปติ,}
\csromangathagroup{\csromangathastanza{วทนฺติ เว ทุฏฺฐมนาปิ เอเก, \\ อโถปิ เว สจฺจมนา วทนฺติ. \\ วาทญฺจ ชาตํ มุนิ โน อุเปติ,}}

\prose{\textbf{ตสฺมา มุนี นตฺถิ ขิโล กุหิญฺจี}ติ.}

\proseitem{16}{\csromansymbolfootnote{*}{ขุ 1. 401 ปิฏฺเฐ.}สกญฺหิ ทิฏฺฐิํ กถมจฺจเยยฺย,}

\prose{\textbf{ฉนฺทานุนีโต รุจิยา นิวิฏฺโฐ.}}

\csromangathameasure{\textbf{สยํ สมตฺตานิ ปกุพฺพมาโน,} \\ \textbf{ยถา หิ ชาเนยฺย ตถา วเทยฺย.}}
\csromangathagroup{\csromangathastanza{\textbf{สยํ สมตฺตานิ ปกุพฺพมาโน,} \\ \textbf{ยถา หิ ชาเนยฺย ตถา วเทยฺย.}}}

\prose{\textbf{สกญฺหิ ทิฏฺฐิํ กถมจฺจเยยฺยา}ติ ยํ เต ติตฺถิยา \textbf{สุนฺทริปริพฺพาชิกํ หนฺตฺวา สฺ}อมณานํ สกฺยปุตฺติยานํ อวณฺณํ ปกาสยิตฺวา เอวํ เอตํ ลาภํ ยสสกฺการํ สมฺมานํ ปจฺจาหริสฺสามาติ.\csromanspacer{} เต เอวํทิฏฺฐิกา เอวํขนฺติกา เอวํรุจิกา เอวํลทฺธิกา เอวํอชฺฌาสยา เอวํอธิปฺปายา.\csromanspacer{} เต นาสกฺขิํสุ สกํ ทิฏฺฐิํ สกํ ขนฺติํ สกํ รุจิํ สกํ ลทฺธิํ สกํ อชฺฌาสยํ สกํ อธิปฺปายํ อติกฺกมิตุํ, อถ โข เสฺวว อยโส เต ปจฺจาคโตติ.\csromanspacer{} เอวมฺปิ สกญฺหิ ทิฏฺฐิํ กถมจฺจเยยฺย.\csromanspacer{}\csromanspacer{}อถ วา “สสฺสโต โลโก อิทเมว สจฺจํ โมฆมญฺญนฺ”ติ โย โส เอวํวาโท, โส สกํ ทิฏฺฐิํ สกํ ขนฺติํ สกํ รุจิํ สกํ ลทฺธิํ สกํ อชฺฌาสยํ สกํ อธิปฺปายํ กถํ อจฺจเยยฺย อติกฺกเมยฺย สมติกฺกเมยฺย วีติวตฺเตยฺย.\csromanspacer{} ตํ กิสฺส เหตุ, ตสฺส สา ทิฏฺฐิ ตถา สมตฺตา สมาทินฺนา คหิตา ปรามฏฺฐา อภินิวิฏฺฐา อชฺโฌสิตา อธิมุตฺตาติ.\csromanspacer{} เอวมฺปิ สกญฺหิ ทิฏฺฐิํ กถมจฺจเยยฺย.\csromanspacer{}\csromanspacer{}อสสฺสโต โลโก.\csromanspacer{} อนฺตวา โลโก.\csromanspacer{} อนนฺตวา โลโก.\csromanspacer{} ตํ ชีวํ ตํ สรีรํ.}

\vspace{\tipitakaparskip}
\csromancenter{ทุฏฺฐฏฺฐกสุตฺตนิทฺเทส อญฺญํ ชีวํ อญฺญํ สรีรํ, โหติ ตถาคโต ปรํ มรณา.\csromanspacer{} น โหติ ตถาคโต ปรํ มรณา.\csromanspacer{} โหติ จ น จ โหติ ตถาคโต ปรํ มรณา.\csromanspacer{} เนว โหติ น น โหติ ตถคโต ปรํ มรณา. “อิทเมว สจฺจํ โมฆมญฺญนฺ”ติ โย โส เอวํ วาโท, โส สกํ ทิฏฺฐิํ สกํ ขนฺติํ สกํ รุจิํ สกํ ลทฺธิํ สกํ อชฺฌาสยํ สกํ อธิปฺปายํ กถํ อจฺจเยยฺย อติกฺกเมยฺย สมติกฺกเมยฺย วีติวตฺเตยฺย.\csromanspacer{} ตํ กิสฺส เหตุ, ตสฺส สา ทิฏฺฐิ ตถา สมตฺตา สมาทินฺนา คหิตา ปรามฏฺฐา อภินิวิฏฺฐา อชฺโฌสิตา อธิมุตฺตาติ เอวมฺปิ สกญฺหิ ทิฏฺฐิํ กถมจฺจเยยฺย.}
\prose{\csromanfolio{49}\textbf{ฉนฺทานุนีโต รุจิยา นิวิฏฺโฐ}ติ \textbf{ฉนฺทานุนีโต}ติ สกาย ทิฏฺฐิยา สกาย ขนฺติยา สกาย รุจิยา สกาย ลทฺธิยา ยายติ นิยฺยติ วุยฺหติ สํหรียติ, ยถา หตฺถิยาเนน วา อสฺสยาเนน วา รถยาเนน วา โคยาเนน วา อชยาเนน วา เมณฺฑยาเนน วา โอฏฺฐยาเนน วา ขรยาเนน วา ยายติ นิยฺยติ วุยฺหติ สํหรียติ.\csromanspacer{} เอวเมว สกาย ทิฏฺฐิยา สกาย ขนฺติยา สกาย รุจิยา สกาย ลทฺธิยา ยายติ นิยฺยติ วุยฺหติ สํหรียตีติ \textbf{ฉนฺทานุนีโต}.\csromanspacer{}\csromanspacer{}\textbf{รุจิยา นิวิฏฺโฐ}ติ สกาย ทิฏฺฐิยา สกาย รุจิยา สกาย ลทฺธิยา นิวิฏฺโฐ ปติฏฺฐิโต อลฺลีโน อุปคโต อชฺโฌสิโต อธิมุตฺโตติ \textbf{ฉนฺทานุนีโต} รุจิยา นิวิฏฺโฐ.}

\prose{\textbf{สยํ สมตฺตานิ ปกุพฺพมาโน}ติ สยํ สมตฺตํ กโรติ, ปริปุณฺณํ กโรติ, อโนมํ กโรติ, อคฺคํ เสฏฺฐํ วิสิฏฺฐํ ปาโมกฺขํ อุตฺตมํ ปวรํ กโรติ, “อยํ สตฺถา สพฺพญฺญู”ติ สยํ สมตฺตํ กโรติ, ปริปุณฺณํ กโรติ, อโนมํ กโรติ, อคฺคํ เสฏฺฐํ วิสิฏฺฐํ ปาโมกฺขํ อุตฺตมํ ปวรํ กโรติ, อยํ ธมฺโม สฺวากฺขาโต, อยํ คโณ สุปฺปฏิปนฺโน, อยํ ทิฏฺฐิ ภทฺทิกา, อยํ ปฏิปทา สุปญฺญตฺตา, อยํ มคฺโค นิยฺยานิโกติ สยํ สมตฺตํ กโรติ, ปริปุณฺณํ กโรติ, อโนมํ กโรติ, อคฺคํ เสฏฺฐํ วิสิฏฺฐํ ปาโมกฺขํ อุตฺตมํ ปวรํ กโรติ ชเนติ สญฺชเนติ นิพฺพตฺเตติ อภินิพฺพตฺเตตีติ สยํ สมตฺตานิ ปกุพฺพมาโน.}

\prose{\textbf{ยถา หิ ชาเนยฺย ตถา วเทยฺยา}ติ ยถา ชาเนยฺย, ตถา วเทยฺย กเถยฺย ภเณยฺย ทีปเยยฺย โวหเรยฺย “สสฺสโต โลโก อิทเมว สจฺจํ โมฆมญฺญนฺ”ติ.\csromanspacer{} ยถา ชาเนยฺย, ตถา วเทยฺย กเถยฺย ภเณยฺย ทีปเยยฺย โวหเรยฺย “อสสฺสโต โลเก ฯเปฯ เนว \csromanfolio{50}โหติ น น โหติ ตถาคโต ปรํ มรณา อิทเมว สจฺจํ โมฆมญฺญนฺ”ติ\textbf{.}\csromanspacer{} ยถา ชาเนยฺย, ตถา วเทยฺย กเถยฺย ภเณยฺย ทีปเยยฺย โวหเรยฺยาติ ยถา หิ ชาเนยฺย ตถา วเทยฺย\textbf{.}\csromanspacer{} เตนาห ภควา\csromandash{}}

\csromangathameasure{สกญฺหิ ทิฏฺฐิํ กถมจฺจเยยฺย, \\ ฉนฺทานุนีโต รุจิยา นิวิฏฺโฐ\textbf{.} \\ สยํ สมตฺตานิ ปกุพฺพมาโน, \\ ยถา หิ ชาเนยฺย ตถา วเทยฺยาติ\textbf{.}}
\csromangathagroup{\csromangathastanza{สกญฺหิ ทิฏฺฐิํ กถมจฺจเยยฺย, \\ ฉนฺทานุนีโต รุจิยา นิวิฏฺโฐ\textbf{.} \\ สยํ สมตฺตานิ ปกุพฺพมาโน, \\ ยถา หิ ชาเนยฺย ตถา วเทยฺยาติ\textbf{.}}}

\proseitem{17}{\csromansymbolfootnote{*}{ขุ 1. 401 ปิฏฺเฐปิ.}โย อตฺตโน สีลวตานิ ชนฺตุ,}

\prose{\textbf{อนานุปุฏฺโฐว}\footnote{อนานุปุฏฺโฐ จ (สฺยา)} \textbf{ปเรส ปาว}\footnote{ปาวา (สี, สฺยา)}\textbf{.}}

\csromanmatikaanchor{mtk.4}
\csromantocmarkref{subsection}{3. ทุฏฺฐฏฺฐกสุตฺตนิทฺเทส}{mtk.4}
\bookmark[dest={mtk.4},level=2]{3. ทุฏฺฐฏฺฐกสุตฺตนิทฺเทส}
\csromangathameasure{\textbf{อนริยธมฺมํ กุสลา ตมาหุ,} \\ \textbf{โย อาตุมานํ สยเมว ปาว2.}}
\csromangathagroup{\csromangathastanza{\textbf{อนริยธมฺมํ กุสลา ตมาหุ,} \\ \textbf{โย อาตุมานํ สยเมว ปาว2.}}}

\prose{\textbf{โย อตฺตโน สีลวตานิ ชนฺตู}ติ \textbf{โย}ติ \textbf{โย} ยาทิโส ยถายุตฺโต ยถาวิหิโต ยถาปกาโร ยํฐานปฺปตฺโต ยํธมฺมสมนฺนาคโต ขตฺติ\textbf{โย} วา พฺราหฺมโณ วา เวสฺโส วา สุทฺโท วา คหฏฺโฐ วา ปพฺพชิโต วา เทโว วา มนุสฺโส วา\textbf{.}\csromanspacer{}\csromanspacer{}\textbf{สีลวตานี}ติ อตฺถิ สีลญฺเจว วตญฺจ\footnote{วตฺตญฺจ (สฺยา), เอวมุปริปิ.}, อตฺถิ วตํ น สีลํ\textbf{.}\csromanspacer{} กตมํ สีลญฺเจว วตญฺจ\csromandash{}อิธ ภิกฺขุ สีลวา โหติ ปาติโมกฺขสํวรสํวุโต วิหรติ อาจารโคจรสมฺปนฺโน อณุมตฺเตสุ วชฺเชสุ ภยทสฺสาวี, สมาทาย สิกฺขติ สิกฺขาปเทสุ\textbf{.}\csromanspacer{} โย ตตฺถ สํยโม สํวโร อวีติกฺกโม, อิทํ สีลํ\textbf{.}\csromanspacer{} ยํ สมาทานํ ตํ วตํ, สํวรฏฺเฐน สีลํ, สมาทานฏฺเฐน วตํ, อิทํ วุจฺจติ สีลญฺเจว วตญฺจ\textbf{.}\csromanspacer{} กตมํ วตํ น สีลํ\csromandash{}อฏฺฐ ธุตงฺคานิ อารญฺญิกงฺคํ ปิณฺฑปาติกงฺคํ ปํสุกูลิกงฺคํ เตจีวริกงฺคํ สปทานจาริกงฺคํ ขลุปจฺฉาภตฺติกงฺคํ เนสชฺชิกงฺคํ ยถาสนฺถติกงฺคํ อิทํ วุจฺจติ วตํ, น สีลํ\textbf{.}\csromanspacer{}\csromanspacer{}วีริย สมาทานมฺปิ วุจฺจติ วตํ, น สีลํ\textbf{.} + กามํ ตโจ จ นฺหารุ\footnote{นหารุ (สี, สฺยา)} จ อฏฺฐิ จ อวสิสฺสตุ\footnote{อวสุสฺสตุ (สฺยา)} สรีเร อุปสฺสุสฺสตุ มํสโลหิตํ\textbf{.}\csromanspacer{} ยํ ตํ ปุริสถาเมน ปุริสพเลน ปุริสวีริเยน ปุริสปรกฺกเมน ปตฺตพฺพํ, น ตํ อปาปุณิตฺวา วีริยสฺส สณฺฐานํ ภวิสฺสตีติ จิตฺตํ ปคฺคณฺหาติ ปทหติ\textbf{.}\csromanspacer{} เอวรูปํ วีริยสมาทานํ อิทํ วุจฺจติ วตํ, น สีลํ\textbf{.}}

\csromanheader{2}{3. ทุฏฺฐฏฺฐกสุตฺตนิทฺเทส}
\csromangathameasure{“นาสิสฺสํ น ปิวิสฺสามิ, วิหารโต น นิกฺขเม. \\ นปิ ปสฺสํ นิปาเตสฺสํ, ตณฺหาสลฺเล อนูหเต”ติ.}
\csromangathagroup{\csromangathastanza{\csromanfolio{51}“นาสิสฺสํ น ปิวิสฺสามิ, วิหารโต น นิกฺขเม. \\ นปิ ปสฺสํ นิปาเตสฺสํ, ตณฺหาสลฺเล อนูหเต”ติ.}}

\prose{จิตฺตํ ปคฺคณฺหาติ ปทหติ.\csromanspacer{} เอวรูปมฺปิ วีริยสมาทานํ วุจฺจติ วตํ, น สีลํ.\csromanspacer{}\csromanspacer{}“น ตาวาหํ อิมํ ปลฺลงฺกํ ภินฺทิสฺสามิ, ยาว เม น อนุปาทาย อาสเวหิ จิตฺตํ วิมุจฺจิสฺสตี”ติ จิตฺตํ ปคฺคณฺหาติ ปทหติ.\csromanspacer{} เอวรูปมฺปิ วีริยสมาทานํ วุจฺจติ วตํ, น สีลํ.\csromanspacer{}\csromanspacer{}“น ตาวาหํ อิมมฺหา อาสนา วุฏฺฐหิสฺสามิ.\csromanspacer{} จงฺกมา โอโรหิสฺสามิ.\csromanspacer{} วิหารา นิกฺขมิสฺสามิ.\csromanspacer{} อฑฺฒโยคา นิกฺขมิสฺสามิ.\csromanspacer{} ปาสาทา นิกฺขมิสฺสามิ.\csromanspacer{} หมฺมิยา นิกฺขมิสฺสามิ.\csromanspacer{} คุหาย นิกฺขมิสฺสามิ.\csromanspacer{} เลณา นิกฺขมิสฺสามิ.\csromanspacer{} กุฏิยา นิกฺขมิสฺสามิ.\csromanspacer{} กูฏาคารา นิกฺขมิสฺสามิ.\csromanspacer{} อฏฺฏา นิกฺขมิสฺสามิ.\csromanspacer{} มาฬา นิกฺขมิสฺสามิ.\csromanspacer{} อุทฺทณฺฑา\footnote{อุฏฺฏณฺฑา (ก) ขุ 8. 286, 297 ปิฏฺเฐสุปิ.} นิกฺขมิสฺสามิ.\csromanspacer{} อุปฏฺฐานสาลาย นิกฺขมิสฺสามิ.\csromanspacer{} มณฺฑปา นิกฺขมิสฺสามิ.\csromanspacer{} รุกฺขมูลา นิกฺขมิสฺสามิ, ยาว เม น อนุปาทาย อาสเวหิ จิตฺตํ วิมุจฺจิสฺสตี”ติ จิตฺตํ ปคฺคณฺหาติ ปทหติ.\csromanspacer{} เอวรูปมฺปิ วีริยสมาทานํ วุจฺจติ วตํ, น สีลํ.\csromanspacer{}\csromanspacer{}“อิมสฺมิญฺเญว ปุพฺพณฺหสมยํ อริยธมฺมํ อาหริสฺสามิ สมาหริสฺสามิ อธิคจฺฉิสฺสามิ ผสฺสยิสฺสามิ สจฺฉิกริสฺสามี”ติ จิตฺตํ ปคฺคณฺหาติ ปทหติ.\csromanspacer{} เอวรูปมฺปิ วีริยสมาทานํ วุจฺจติ วตํ, น สีลํ. “อิมสฺมิญฺเญว มชฺฌนฺหิกสมยํ.\csromanspacer{} สายนฺหสมยํ.\csromanspacer{} ปุเรภตฺตํ.\csromanspacer{} ปจฺฉาภตฺตํ.\csromanspacer{} ปุริมํ ยามํ.\csromanspacer{} มชฺฌิมํ ยามํ.\csromanspacer{} ปจฺฉิมํ ยามํ.\csromanspacer{} กาเฬ.\csromanspacer{} ชุเณฺห.\csromanspacer{} วสฺเส.\csromanspacer{} เหมนฺเต.\csromanspacer{} คิเมฺห.\csromanspacer{} ปุริเม วโยขนฺเธ.\csromanspacer{} มชฺฌิเม วโยขนฺเธ.\csromanspacer{} ปจฺฉิเม วโยขนฺเธ อริยธมฺมํ อาหริสฺสามิ สมาหริสฺสามิ อธิคจฺฉิสฺสามิ ผสฺสยิสฺสามิ สจฺฉิกริสฺสามี”ติ จิตฺตํ ปคฺคณฺหาติ ปทหติ.\csromanspacer{} เอวรูปมฺปิ วีริยสมาทานํ วุจฺจติ วตํ, น สีลํ.\csromanspacer{}\csromanspacer{}\textbf{ชนฺตู}ติ สตฺโต นโร มานโว\footnote{มาณโว (ก)} โปโส ปุคฺคโล ชีโว ชาคุ ชนฺตุ อินฺทคุ มนุโชติ โย อตฺตโน สีลวตานิ ชนฺตุ.}

\prose{\textbf{อนานุปุฏฺโฐว ปเรส ปาวา}ติ \textbf{ปเรสนฺ}ติ ปเรสํ ขตฺติยานํ พฺราหฺมณานํ เวสฺสานํ สุทฺทานํ คหฏฺฐานํ ปพฺพชิตานํ เทวานํ มนุสฺสานํ.\csromanspacer{}\csromanspacer{}\textbf{อนานุปุฏฺโฐ}ติ อปุฏฺโฐ อปุจฺฉิโต อยาจิโต อนชฺเฌสิโต อปสาทิโต.\csromanspacer{}\csromanspacer{}\textbf{ปาวา}ติ อตฺตโน สีลํ วา วตํ วา สีลพฺพตํ ปาวทติ. “อหมสฺมิ สีลสมฺปนฺโน”ติ วา “วตสมฺปนฺโน”ติ วา “สีลพฺพตสมฺปนฺโน”ติ วา ชาติยา วา โคตฺเตน วา โกลปุตฺติเยน วา วณฺณโปกฺขรตาย วา ธเนน วา อชฺเฌเนน วา กมฺมายตเนน วา}

\prose{\csromanfolio{52}สิปฺปายตเนน วา วิชฺชาฏฺฐาเนน\footnote{วิชฺชฏฺฐาเนน (สฺยา)} วา สุเตน วา ปฏิภาเนน\footnote{ปฏิภาเณน (สี, สฺยา, ก)} วา อญฺญตรญฺญตเรน วา วตฺถุนา, “อุจฺจา กุลา ปพฺพชิโต”ติ วา “มหากุลา ปพฺพชิโต”ติ วา “มหาโภคกุลา ปพฺพชิโต”ติ วา “อุฬารโภคกุลา ปพฺพชิโต”ติ วา, “ญาโต ยสสฺสี สคหฏฺฐปพฺพชิตานนฺ”ติ วา “ลาภิมฺหิ จีวรปิณฺฑปาตเสนาสนคิลานปจฺจยเภสชฺชปริกฺขารานนฺ”ติ วา “สุตฺตนฺติโก”ติ วา “วินยธโร”ติ วา “ธมฺมกถิโก”ติ วา “อารญฺญิโก”ติ วา “ปิณฺฑปาติโก”ติ วา ปํสุกูลิโก”ติ วา “เตจีวริโก”ติ วา “สปทานจาริโก”ติ วา ขลุปจฺฉาภตฺติโก”ติ วา “เนสชฺชิโก”ติ วา “ยถาสนฺถติโก”ติ วา “ปฐมสฺส ฌานสฺส ลาภี”ติ วา “ทุติยสฺส ฌานสฺส ลาภี”ติ วา “ตติยสฺส ฌานสฺส ลาภี”ติ วา “จตุตฺถสฺส ฌานสฺส ลาภี”ติ วา “อากาสานญฺจายตน- สมาปตฺติยา ลาภี”ติ วา “วิญฺญาณญฺจายตนสมาปตฺติยา ลาภี”ติ วา “อากิญฺจญฺญายตนสมาปตฺติยา ลาภี”ติ วา “เนวสญฺญานาสญฺญายตน- สมาปตฺติยา ลาภี”ติ วา ปาวทติ กเถติ ภณติ ทีปยติ โวหรตีติ อนานุปุฏฺโฐว ปเรสํ ปาว.}

\prose{\textbf{อนริยธมฺมํ กุสลา ตมาหู}ติ \textbf{กุสลา}ติ เย เต ขนฺธ\textbf{กุสลา} ธาตุ\textbf{กุสลา} อายตน\textbf{กุสลา} ปฏิจฺจสมุปฺปาท\textbf{กุสลา} สติปฏฺฐาน\textbf{กุสลา} สมฺมปฺปธาน\textbf{กุสลา} อิทฺธิปาท\textbf{กุสลา} อินฺทฺริย\textbf{กุสลา} พล\textbf{กุสลา} โพชฺฌงฺค\textbf{กุสลา} มคฺค\textbf{กุสลา} ผล\textbf{กุสลา} นิพฺพาน\textbf{กุสลา}, เต \textbf{กุสลา} เอวมาหํสุ “อนริยานํ เอโส ธมฺโม, เนโส ธมฺโม อริยานํ”. “พาลานํ เอโส ธมฺโม, เนโส ธมฺโม ปณฺฑิตานํ”. “อสปฺปุริสานํ เอโส ธมฺโม, เนโส ธมฺโม สปฺปุริสานนฺ”ติ เอวมาหํสุ เอวํ กเถนฺติ เอวํ ภณนฺติ เอวํ ทีปยนฺติ เอวํ โวหรนฺตีติ อนริยธมฺมํ \textbf{กุสลา} ตมาหุ.}

\prose{\textbf{โย อาตุมานํ สยเมว ปาวา}ติ \textbf{อาตุมา} วุจฺจติ อตฺตา.\csromanspacer{}\csromanspacer{}\textbf{สยเมว ปาวา}ติ สยเมว อตฺตานํ ปาวทติ “อหมสฺมิ สีลสมฺปนฺโน”ติ วา “วตสมฺปนฺโน”ติ วา “สีลพฺพตสมฺปนฺโน”ติ วา ชาติยา วา โคตฺเตน วา โกลปุตฺติเยน วา วณฺณโปกฺขรตาย วา ธเนน วา อชฺเฌเนน วา กมฺมายตเนน วา สิปฺปายตเนน วา วิชฺชาฏฺฐาเนน วา สุเตน วา ปฏิภาเนน วา อญฺญตรญฺญตเรน วา วตฺถุนา, “อุจฺจา กุลา}

\vspace{\tipitakaparskip}
\csromancenter{ทุฏฺฐฏฺฐกสุตฺตนิทฺเทส ปพฺพชิโต”ติ วา “มหากุลา ปพฺพชิโต”ติ วา “มหาโภคกุลา ปพฺพชิโต”ติ วา “อุฬารโภคกุลา ปพฺพชิโต”ติ วา “ญาโต ยสสฺสี สคหฏฺฐปพฺพชิตานนฺ”ติ วา, “ลาภิมฺหิ จีวรปิณฺฑปาตเสนาสนคิลานปจฺจยเภสชฺชปริกฺขารานนฺ”ติ วา “สุตฺตนฺติโก”ติ วา “วินยธโร”ติ วา “ธมฺมกถิโก”ติ วา “อารญฺญิโก”ติ วา “ปิณฺฑปาติโก”ติ วา “ปํสุกูลิโก”ติ วา “เตจีวริโก”ติ วา “สปทานจาริโก”ติ วา “ขลุปจฺฉาภตฺติโก”ติ วา “เนสชฺชิโก”ติ วา “ยถาสนฺถติโก”ติ วา ปฐมสฺส ฌานสฺส ลาภี”ติ วา “ทุติยสฺส ฌานสฺส ลาภี”ติ วา “ตติยสฺส ฌานสฺส ลาภี”ติ วา “จตุตฺถสฺส ฌานสฺส ลาภี”ติ วา “อากาสานญฺจายตนสมาปตฺติยา ลาภี”ติ วา “วิญฺญาณญฺจายตนสมาปตฺติยา ลาภี”ติ วา อากิญฺจญฺญายตนสมาปตฺติยา ลาภี”ติ วา “เนวสญฺญานาสญฺญายตนสมาปตฺติยา ลาภี”ติ วา ปาวทติ กเถติ ภณติ ทีปยติ โวหรตีติ โย อาตุมานํ สยเมว ปาวาติ.\csromanspacer{} เตนาห ภควา\csromandash{}}
\vspace{0.5\baselineskip}
\csromangathameasure{โย อตฺตโน สีลวตานิ ชนฺตุ, \\ อนานุปุฏฺโฐว ปเรส ปาว. \\ อนริยธมฺมํ กุสลา ตมาหุ, \\ โย อาตุมานํ สยเมว ปาวาติ.}
\csromangathagroup{\csromangathastanza{\csromanfolio{53}โย อตฺตโน สีลวตานิ ชนฺตุ, \\ อนานุปุฏฺโฐว ปเรส ปาว. \\ อนริยธมฺมํ กุสลา ตมาหุ, \\ โย อาตุมานํ สยเมว ปาวาติ.}}

\proseitem{18}{\csromansymbolfootnote{*}{ขุ 1. 401 ปิฏฺเฐปิ.}\textbf{สนฺโต จ ภิกฺขุ อภินิพฺพุตตฺโต},}

\prose{\textbf{อิติ’หนฺติ สีเลสุ อกตฺถมาโน.}}

\csromangathameasure{\textbf{ตมริยธมฺมํ กุสลา วทนฺติ,} \\ \textbf{ยสฺสุสฺสทา นตฺถิ กุหิญฺจิ โลเก.}}
\csromangathagroup{\csromangathastanza{\textbf{ตมริยธมฺมํ กุสลา วทนฺติ,} \\ \textbf{ยสฺสุสฺสทา นตฺถิ กุหิญฺจิ โลเก.}}}

\prose{\textbf{สนฺโต จ ภิกฺขุ อภินิพฺพุตตฺโต}ติ \textbf{สนฺโต}ติ + ราคสฺส สมิตตฺตา \textbf{สนฺโต}, โทสสฺส สมิตตฺตา \textbf{สนฺโต}, โมหสฺส สมิตตฺตา \textbf{สนฺโต}, โกธสฺส.\csromanspacer{} อุปนาหสฺส.\csromanspacer{} มกฺขสฺส.\csromanspacer{} ปฬาสสฺส\footnote{ปลาสสฺส (สี, ก)}.\csromanspacer{} อิสฺสาย.\csromanspacer{} มจฺฉริยสฺส.\csromanspacer{} มายาย.\csromanspacer{} สาเฐยฺยสฺส.\csromanspacer{} ถมฺภสฺส.\csromanspacer{} สารมฺภสฺส.\csromanspacer{} มานสฺส.\csromanspacer{} อติมานสฺส.\csromanspacer{} มทสฺส.\csromanspacer{} ปมาทสฺส.\csromanspacer{} สพฺพกิเลสานํ.\csromanspacer{} สพฺพทุจฺจริตานํ.\csromanspacer{} สพฺพทรถานํ.\csromanspacer{} สพฺพปริฬาหานํ.\csromanspacer{} สพฺพสนฺตาปานํ.\csromanspacer{} สพฺพกุสลาภิสงฺขารานํ สนฺตตฺตา สมิตตฺตา วูปสมิตตฺตา วิชฺฌาตตฺตา นิพฺพุตตฺตา วิคตตฺตา ปฏิปสฺสทฺธตฺตา \textbf{สนฺโต} \csromanfolio{54}อุปสนฺโต วูปสนฺโต นิพฺพุโต ปฏิปสฺสทฺโธติ สนฺโต.\csromanspacer{}\csromanspacer{}* \textbf{ภิกฺขู}ติ สตฺตนฺนํ ธมฺมํ ภินฺนตฺตา ภิกฺขุ, สกฺกายทิฏฺฐิ ภินฺนา โหติ, วิจิกิจฺฉา ภินฺนา โหติ, สีลพฺพตปรามาโส ภินฺโน โหติ, ราโค ภินฺโน โหติ, โทโส ภินฺโน โหติ, โมโห ภินฺโน โหติ, มาโน ภินฺโน โหติ, ภินฺนาสฺส โหนฺติ ปาปกา อ\textbf{กุสลา} ธมฺมา สํกิเลสิกา โปโนภวิกา\footnote{โปโนพฺภวิกา (สฺยา, ก)} สทรา ทุกฺขวิปากา อายติํ ชาติชรามรณิยา.}

\prose{ปชฺเชน กเตน อตฺตนา, (สภิยาติ ภควา,)}

\prose{ปรินิพฺพานคโต วิติณฺณกงฺโข.}

\csromangathameasure{วิภวญฺจ ภวญฺจ วิปฺปหาย, \\ วุสิตวา ขีณปุนพฺภโว ส ภิกฺขูติ.}
\csromangathagroup{\csromangathastanza{วิภวญฺจ\footnote{วิภวํ (สี, ก) ขุ 1. 357 ปิฏฺเฐปิ.} ภวญฺจ วิปฺปหาย, \\ วุสิตวา ขีณปุนพฺภโว ส ภิกฺขูติ.}}

\prose{\textbf{สนฺโต จ ภิกฺขุ อภินิพฺพุตตฺโต}ติ ราคสฺส นิพฺพาปิตตฺตา อภินิพฺพุตตฺโต, โทสสฺส นิพฺพาปิตตฺตา อภินิพฺพุตตฺโต, โมหสฺส นิพฺพาปิตตฺตา อภินิพฺพุตตฺโต, โกธสฺส.\csromanspacer{} อุปนาหสฺส.\csromanspacer{} มกฺขสฺส.\csromanspacer{} ปฬาสสฺส.\csromanspacer{} อิสฺสาย.\csromanspacer{} มจฺฉริยสฺส.\csromanspacer{} มายาย.\csromanspacer{} สาเฐยฺยสฺส.\csromanspacer{} ถมฺภสฺส.\csromanspacer{} สารมฺภสฺส.\csromanspacer{} มานสฺส.\csromanspacer{} อติมานสฺส.\csromanspacer{} มทสฺส.\csromanspacer{} ปมาทสฺส.\csromanspacer{} สพฺพกิเลสานํ.\csromanspacer{} สพฺพทุจฺจริตานํ.\csromanspacer{} สพฺพทรถานํ.\csromanspacer{} สพฺพปริฬาหานํ.\csromanspacer{} สพฺพสนฺตาปานํ.\csromanspacer{} สพฺพากุสลาภิสงฺขารานํ นิพฺพาปิตตฺตา อภินิพฺพุตตฺโตติ สนฺโต จ ภิกฺขุ อภินิพฺพุตตฺโต.}

\prose{\textbf{อิติ’หนฺติ สีเลสุ อกตฺถมาโน}ติ \textbf{อิตี}ติ ปทสนฺธี ปทสํสคฺโค ปทปาริปูรี อกฺขรสมวาโย พฺยญฺจนสิลิฏฺฐตา ปทานุปุพฺพตา เมตํ\footnote{ปทานุปุพฺพตา เปตํ (สี, ก)} อิติ’หนฺติ.\csromanspacer{} \textbf{สีเลสุ อกตฺถมาโน}ติ อิเธกจฺโจ กตฺถี โหติ วิกตฺถี, โส กตฺถติ วิกตฺถติ. “อหมสฺมิ สีลสมฺปนฺโน”ติ วา “วตสมฺปนฺโน”ติ วา “สีลพฺพตสมฺปนฺโน”ติ วา ชาติยา วา โคตฺเตน วา โกลปุตฺติเยน วา วณฺณโปกฺขรตาย วา ฯเปฯ “เนวสญฺญานาสญฺญายตนสมาปตฺติยา ลาภี”ติ วา กตฺถติ วิกตฺถติ.\csromanspacer{} เอวํ น กตฺถติ น วิกตฺถติ, กตฺถนา อารโต วิรโต ปฏิวิรโต นิกฺขนฺโต นิสฺสโฏ วิปฺปมุตฺโต วิสญฺญุตฺโต วิมริยาทิกเตน เจตสา วิหรตีติ อิติ’หนฺติ สีเลสุ อกตฺถมาโน.}

\csromanheader{2}{3. ทุฏฺฐฏฺฐกสุตฺตนิทฺเทส}
\prose{\csromanfolio{55}\textbf{ตมริยธมฺมํ กุสลา วทนฺตี}ติ \textbf{กุสลา}ติ เย เต ขนฺธ\textbf{กุสลา} ธาตุ\textbf{กุสลา} อายตน\textbf{กุสลา} ปฏิจฺจสมุปฺปาท\textbf{กุสลา} สติปฏฺฐาน\textbf{กุสลา} สมฺมปฺปธาน\textbf{กุสลา} อิทฺธิปาท\textbf{กุสลา} อินฺทฺริย\textbf{กุสลา} พล\textbf{กุสลา} โพชฺฌงฺค\textbf{กุสลา} มคฺค\textbf{กุสลา} ผล\textbf{กุสลา} นิพฺพาน\textbf{กุสลา}, เต \textbf{กุสลา} เอวํ วทนฺติ “อริยานํ เอโส ธมฺโม, เนโส ธมฺโม อนริยานํ”. “ปณฺฑิตานํ เอโส ธมฺโม, เนโส ธมฺโม พาลานํ”. “สปฺปุริสานํ เอโส ธมฺโม, เนโส ธมฺโม อสปฺปุริสานนฺ”ติ เอวํ วทนฺติ อริยานํ เอวํ กเถนฺติ เอวํ ภณนฺติ เอวํ ทีปยนฺติ เอวํ โวหรนฺตีติ ตมริยธมฺมํ \textbf{กุสลา} วทนฺติ.}

\prose{\textbf{ยสฺสุสฺสทา นตฺถิ กุหิญฺจิ โลเก}ติ \textbf{ยสฺสา}ติ อรหโต ขีณาสวสฺส.\csromanspacer{} \textbf{อุสฺสทา}ติ สตฺตุสฺสทา ราคุสฺสโท โทสุสฺสโท โมหุสฺสโท มานุสฺสโท ทิฏฺฐุสฺสโท กิเลสุสฺสโท กมฺมุสฺสโท, ยสฺสิเม\footnote{ตสฺสิเม (สี, สฺยา)} อุสฺสทา นตฺถิ น สนฺติ น วิชฺชนฺติ นุปลพฺภนฺติ, ปหีนา สมุจฺฉินฺนา วูปสนฺตา ปฏิปสฺสทฺธา อภพฺพุปฺปตฺติกา ญาณคฺคินา ทฑฺฒา.\csromanspacer{}\csromanspacer{}\textbf{กุหิญฺจี}ติ กุหิญฺจิ กิมฺหิจิ กตฺถจิ อชฺฌตฺตํ วา พหิทฺธา วา อชฺฌตฺตพหิทฺธา วา.\csromanspacer{}\csromanspacer{}\textbf{โลเก}ติ อปายโลเก มนุสฺสโลเก เทวโลเก ขนฺธโลเก ธาตุโลเก อายตนโลเกติ ยสฺสุสฺสทา นตฺถิ กุหิญฺจิ โลเก.\csromanspacer{} เตนาห ภควา\csromandash{}}

\csromangathameasure{สนฺโต จ ภิกฺขุ อภินิพฺพุตตฺโต, \\ อิติ’หนฺติ สีเลสุ อกตฺถมาโน.}
\csromangathagroup{\csromangathastanza{สนฺโต จ ภิกฺขุ อภินิพฺพุตตฺโต, \\ อิติ’หนฺติ สีเลสุ อกตฺถมาโน.}}

\prose{ตมริยธมฺมํ \textbf{กุสลา} วทนฺติ,}

\prose{\textbf{ยสฺสุสฺสทา นตฺถิ กุหิญฺจิ โลเก}ติ.}

\proseitem{19}{\csromansymbolfootnote{*}{ขุ 1. 401 ปิฏฺเฐปิ.}\textbf{ปกปฺปิตา สงฺขตา ยสฺส ธมฺมา},}

\prose{\textbf{ปุรกฺขตา}\footnote{ปุเรกฺขตา (สี, ก)} \textbf{สนฺติ อวีวทาตา.}}

\csromangathameasure{\textbf{ยทตฺตนิ ปสฺสติ อานิสํสํ,} \\ \textbf{ตํ นิสฺสิโต กุปฺปปฏิจฺจสนฺติํ.}}
\csromangathagroup{\csromangathastanza{\textbf{ยทตฺตนิ ปสฺสติ อานิสํสํ,} \\ \textbf{ตํ นิสฺสิโต กุปฺปปฏิจฺจสนฺติํ.}}}

\prose{\textbf{ปกปฺปิตา สงฺขตา ยสฺส ธมฺมา}ติ \textbf{ปกปฺปนา}ติ เทฺว \textbf{ปกปฺปนา} ตณฺหา\textbf{ปกปฺปนา} จ ทิฏฺฐิ\textbf{ปกปฺปนา} จ ฯเปฯ อยํ ตณฺหา\textbf{ปกปฺปนา} ฯเปฯ อยํ ทิฏฺฐิ\textbf{ปกปฺปนา}.\csromanspacer{} \textbf{สงฺกตา}ติ สงฺขตา อภิสงฺขตา สณฺฐปิตาติปิ สงฺขตา.\csromanspacer{} อถ วา อนิจฺจา สงฺขตา ปฏิจฺจสมุปฺปนฺนา ขยธมฺมา วยธมฺมา วิราคธมฺมา นิโรธธมฺมาติปิ \csromanfolio{56}สงฺขตา\textbf{.}\csromanspacer{}\csromanspacer{}\textbf{ยสฺสา}ติ ทิฏฺฐิคติกสฺส\textbf{.}\csromanspacer{} \textbf{ธมฺมา} วุจฺจนฺติ ทฺวาสฏฺฐิ ทิฏฺฐิคตานีติ ปกปฺปิตา สงฺขตา ยสฺส ธมฺมา\textbf{.}}

\prose{\textbf{ปุรกฺขตา สนฺติ อวีวทาตา}ติ \textbf{ปุรกฺขตา}ติ เทฺว ปุเรกฺขารา ตณฺหาปุเรกฺขาโร จ ทิฏฺฐิปุเรกฺขาโร จ ฯเปฯ อยํ ตณฺหาปุเรกฺขาโร ฯเปฯ อยํ ทิฏฺฐิปุเรกฺขาโร\textbf{.}\csromanspacer{} ตสฺส ตณฺหาปุเรกฺขาโร อปฺปหีโน, ทิฏฺฐิปุเรกฺขาโร อปฺปฏินิสฺสฏฺโฐ, ตสฺส ตณฺหาปุเรกฺขารสฺส อปฺปหีนตฺตา, ทิฏฺฐิปุเรกฺขารสฺส อปฺปฏินิสฺสฏฺฐตฺตา โส ตณฺหํ วา ทิฏฺฐิํ วา ปุเรโต กตฺวา จรติ ตณฺหาธโช ตณฺหาเกตุ ตณฺหาธิปเตยฺโย, ทิฏฺฐิธโช ทิฏฺฐิเกตุ ทิฏฺฐาธิปเตยฺโย, ตณฺหาย วา ทิฏฺฐิยา วา ปริวาริโต จรตีติ \textbf{ปุรกฺขตา.}\csromanspacer{}\csromanspacer{}\textbf{สนฺตี}ติ สนฺติ สํวิชฺชนฺติ อตฺถิ อุปลพฺภนฺติ\textbf{.}\csromanspacer{}\csromanspacer{}\textbf{อวีวทาตา}ติ อเววทาตา อโวทาตา อปริสุทฺธา สํกิลิฏฺฐา สํกิเลสิกาติ \textbf{ปุรกฺขตา} สนฺติ อวีวทาตา\textbf{.}}

\prose{\textbf{ยทตฺตนิ ปสฺสติ อานิสํส}นฺติ \textbf{ยทตฺตนี}ติ ยํ อตฺตนิ\textbf{.}\csromanspacer{} \textbf{อตฺตา} วุจฺจติ ทิฏฺฐิคตํ\textbf{.}\csromanspacer{} อตฺตโน ทิฏฺฐิยา เทฺว อานิสํเส ปสฺสติ ทิฏฺฐธมฺมิกญฺจ อานิสํสํ สมฺปรายิกญฺจ อานิสํสํ\textbf{.}\csromanspacer{} กตโม ทิฏฺฐิยา ทิฏฺฐธมฺมิโก อานิสํโส, ยํทิฏฺฐิโก สตฺถา โหติ, ตํทิฏฺฐิกา สาวกา โหนฺติ\textbf{.}\csromanspacer{} ตํทิฏฺฐิกํ สตฺถารํ สาวกา สกฺกโรนฺติ ครุํ กโรนฺติ\footnote{ครุกโรนฺติ (สี, สฺยา)} มาเนนฺติ ปูเชนฺติ อปจิติํ กโรนฺติ, ลภติ จ ตโตนิทานํ จีวรปิณฺฑปาตเสนาสนคิลานปจฺจย- เภสชฺชปริกฺขารํ\textbf{.}\csromanspacer{} อยํ ทิฏฺฐิยา ทิฏฺฐธมฺมิโก อานิสํโส\textbf{.}\csromanspacer{}\csromanspacer{}กตโม ทิฏฺฐิยา สมฺปรายิโก อานิสํโส, อยํ ทิฏฺฐิ อลํ นาคตฺตาย วา สุปณฺณตฺตาย วา ยกฺขตฺตาย วา อสุรตฺตาย วา คนฺธพฺพตฺตาย วา มหาราชตฺตาย วา อินฺทตฺตาย วา พฺรหฺมตฺตาย วา เทวตฺตาย วา อยํ ทิฏฺฐิ สุทฺธิยา วิสุทฺธิยา ปริสุทฺธิยา, มุตฺติยา วิมุตฺติยา ปริมุตฺติยา, อิมาย ทิฏฺฐิยา สุชฺฌนฺติ วิสุชฺฌนฺติ ปริสุชฺฌนฺติ, มุจฺจนฺติ วิมุจฺจนฺติ ปริมุจฺจนฺติ, อิมาย ทิฏฺฐิยา สุชฺฌิสฺสามิ วิสุชฺฌิสฺสามิ ปริสุชฺฌิสฺสามิ, มุจฺจิสฺสามิ วิมุจฺจิสฺสามิ ปริมุจฺจิสฺสามีติ อายติํ ผลปาฏิกงฺขี โหติ\textbf{.}\csromanspacer{} อยํ ทิฏฺฐิยา สมฺปรายิโก อานิสํโส\textbf{.}\csromanspacer{}\csromanspacer{}อตฺตโน ทิฏฺฐิยา อิเม เทฺว อานิสํเส ปสฺสติ ทกฺขติ โอโลเกติ นิชฺฌายติ อุปปริกฺขตีติ ยทตฺตนิ ปสฺสติ อานิสํสํ\textbf{.}}

\csromanheader{2}{3. ทุฏฺฐฏฺฐกสุตฺตนิทฺเทส}
\prose{\csromanfolio{57}\textbf{ตํ นิสฺสิโต กุปฺปปฏิจฺจสนฺตินฺ}ติ ติสฺโส สนฺติโย อจฺจนฺตสนฺติ ตทงฺคสนฺติ สมฺมุติสนฺติ.\csromanspacer{} กตมา \textbf{อจฺจนฺตสนฺติ,} อจฺจนฺตสนฺติ วุจฺจติ อมตํ นิพฺพานํ.\csromanspacer{} โย โส สพฺพสงฺขารสมโถ สพฺพูปธิปฏินิสฺสคฺโค ตณฺหกฺขโย วิราโค นิโรโธ นิพฺพานํ, อยํ อจฺจนฺตสนฺติ.\csromanspacer{}\csromanspacer{}กตมา \textbf{ตทงฺคสนฺติ,} ปฐมํ ฌานํ สมาปนฺนสฺส นีวรณา สนฺตา โหนฺติ, ทุติยํ ฌานํ สมาปนฺนสฺส วิตกฺกวิจารา สนฺตา โหนฺติ, ตติยํ ฌานํ สมาปนฺนสฺส ปีติ สนฺตา โหติ, จตุตฺถํ ฌานํ สมาปนฺนสฺส สุขทุกฺขา สนฺตา โหนฺติ, อากาสานญฺจายตนํ สมาปนฺนสฺส รูปสญฺญา ปฏิฆสญฺญา นานตฺตสญฺญา สนฺตา โหติ, วิญฺญาณญฺจายตนํ สมาปนฺนสฺส อากาสานญฺจายตนสญฺญา สนฺตา โหติ, อากิญฺจญฺญายตนํ สมาปนฺนสฺส วิญฺญาณญฺจายตนสญฺญา สนฺตา โหติ, เนวสญฺญานาสญฺญายตนํ สมาปนฺนสฺส อากิญฺจญฺญายตนสญฺญา สนฺตา โหติ.\csromanspacer{} อยํ ตทงฺคสนฺติ.\csromanspacer{}\csromanspacer{}กตมา \textbf{สมฺมุติสนฺติ,} สมฺมุติสนฺติโย วุจฺจนฺติ ทฺวาสฏฺฐิ ทิฏฺฐิคตานิ ทิฏฺฐิสนฺติโย.\csromanspacer{} อปิ จ สมฺมุติสนฺติ อิมสฺมิํ อตฺเถ อธิปฺเปตา สนฺตีติ.\csromanspacer{}\csromanspacer{}\textbf{ตํ นิสฺสิโต กุปฺปปฏิจฺจสนฺตินฺ}ติ กุปฺปสนฺติํ ปกุปฺปสนฺติํ เอริตสนฺติํ สเมริตสนฺติํ จลิตสนฺติํ ฆฏฺฏิตสนฺติํ กปฺปิตสนฺติํ ปกปฺปิตสนฺติํ อนิจฺจํ สงฺขตํ ปฏิจฺจสมุปฺปนฺนํ ขยธมฺมํ วยธมฺมํ วิราคธมฺมํ นิโรธธมฺมํ, สนฺติํ นิสฺสิโต อสฺสิโต อลฺลีโน อุปคโต อชฺโฌสิโต อธิมุตฺโตติ ตํ นิสฺสิโต กุปฺปปฏิจฺจสนฺติ.\csromanspacer{} เตนาห ภควา\csromandash{}}

\csromangathameasure{ปกปฺปิตา สงฺขตา ยสฺส ธมฺมา, \\ ปุรกฺขตา สนฺติ อวีวทาตา. \\ ยทตฺตนิ ปสฺสติ อานิสํสํ,}
\csromangathagroup{\csromangathastanza{ปกปฺปิตา สงฺขตา ยสฺส ธมฺมา, \\ ปุรกฺขตา สนฺติ อวีวทาตา. \\ ยทตฺตนิ ปสฺสติ อานิสํสํ,}}

\prose{\textbf{ตํ นิสฺสิโต กุปฺปปฏิจฺจสนฺตินฺ}ติ.}

\proseitem{20}{\csromansymbolfootnote{*}{ขุ 1. 401 ปิฏฺเฐปิ.}\textbf{ทิฏฺฐีนิเวสา น หิ สฺวาติวตฺตา},}

\prose{\textbf{ธมฺเมสุ นิจฺเฉยฺย สมุคฺคหีตํ.}}

\csromangathameasure{\textbf{ตสฺมา นโร เตสุ นิเวสเนสุ,} \\ \textbf{นิรสฺสตี อาทิยตี จ ธมฺมํ.}}
\csromangathagroup{\csromangathastanza{\textbf{ตสฺมา นโร เตสุ นิเวสเนสุ,} \\ \textbf{นิรสฺสตี อาทิยตี จ ธมฺมํ.}}}

\prose{\textbf{ทิฏฺฐีนิเวสา น หิ สฺวาติวตฺตา}ติ \textbf{ทิฏฺฐีนิเวสา}ติ “สสฺสโต โลโก อิทเมว สจฺจํ โมฆมญฺญนฺ”ติ อภินิเวสปรามาโส ทิฏฺฐินิเวสนํ. “อสสฺสโต โลโก.\csromanspacer{} อนฺตวา โลโก.\csromanspacer{} อนนฺตวา \csromanfolio{58}โลโก.\csromanspacer{} ตํ ชีวํ ตํ สรีรํ.\csromanspacer{} อญฺญํ ชีวํ อญฺญํ สรีรํ.\csromanspacer{} โหติ ตถาคโต ปรํ มรณา.\csromanspacer{} น โหติ ตถาคโต ปรํ มรณา.\csromanspacer{} โหติ จ น จ โหติ ตถาคโต ปรํ มรณา.\csromanspacer{} เนว โหติ น น โหติ ตถาคโต ปรํ มรณา อิทเมว สจฺจํ โมฆมญฺญนฺ”ติ อภินิเวสปรามาโส ทิฏฺฐินิเวสนนฺติ.\csromanspacer{}\csromanspacer{}\textbf{ทิฏฺฐีเนเวสา น หิ สฺวาติวตฺตา}ติ ทิฏฺฐินิเวสา น หิ สฺวาติวตฺตา ทุรติวตฺตา ทุตฺตรา ทุปฺปตรา ทุสฺสมติกฺกมา ทุพฺพินิวตฺตาติ\footnote{ทุพฺพีติวตฺตาติ (สี, สฺยา, ก)} ทิฏฺฐีนิเวสา น หิ สฺวาติวตฺตา.}

\prose{\textbf{ธมฺเมสุ นิจฺเฉยฺย สมุคฺคหีตนฺ}ติ \textbf{ธมฺเมสู}ติ ทฺวาสฏฺฐิทิฏฺฐิคเตสุ.\csromanspacer{}\csromanspacer{}\textbf{นิจฺเฉยฺยา}ติ นิจฺฉินิตฺวา วินิจฺฉินิตฺวา วิจินิตฺวา ปวิจินิตฺวา ตุลยิตฺวา ตีรยิตฺวา วิภาวยิตฺวา วิภูตํ กตฺวา.\csromanspacer{}\csromanspacer{}\textbf{สมุคฺคหีตนฺ}ติ นิเวสเนสุ โอธิคฺคาโห พิลคฺคาโห วรคฺคาโห โกฏฺฐาสคฺคาโห อุจฺจยคฺคาโห สมุจฺจยคฺคาโห, อิทํ สจฺจํ ตจฺฉํ ตถํ ภูตํ ยาถาวํ อวิปรีตนฺติ คหิตํ ปรามฏฺฐํ อภินิวิฏฺฐํ อชฺโฌสิตํ อธิมุตฺตนฺติ ธมฺเมสุ นิจฺเฉยฺย สมุคฺคหีตํ.}

\prose{\textbf{ตสฺมา นโร เตสุ นิเวสเนสู}ติ \textbf{ตสฺมา}ติ \textbf{ตสฺมา} ตํการณา ตํเหตุ ตปฺปจฺจยา ตํนิทานํ.\csromanspacer{}\csromanspacer{}\textbf{นโร}ติ สตฺโต นโร มานโว โปโส ปุคฺคโล ชีโว ชาคุ ชนฺตุ อินฺทคุ มนุโช.\csromanspacer{}\csromanspacer{}\textbf{เตสุ นิเวสเนสู}ติ เตสุ ทิฏฺฐินิเวสเนสูติ \textbf{ตสฺมา} นโร เตสุ นิเวสเนสุ.}

\prose{\textbf{นิรสฺสตี อาทิยตี จ ธมฺมนฺ}ติ \textbf{นิรสฺสตี}ติ ทฺวีหิ การเณหิ นิรสฺสติ, ปรวิจฺฉินฺทนาย วา นิรสฺสติ, อนภิสมฺภุณนฺโต วา นิรสฺสติ.\csromanspacer{} กถํ ปรวิจฺฉินฺทนาย นิรสฺสติ\csromandash{}ปโร วิจฺฉินฺเทติ, โส สตฺถา น สพฺพญฺญู, ธมฺโม น สฺวากฺขาโต, คโณ น สุปฺปฏิปนฺโน, ทิฏฺฐิ น ภทฺทิกา, ปฏิปทา น สุปญฺญตฺตา, มคฺโค น นิยฺยานิโก, นตฺเถตฺถ สุทฺธิ วา วิสุทฺธิ วา ปริสุทฺธิ วา, มุตฺติ วา วิมุตฺติ วา ปริมุตฺติ วา, นตฺเถตฺถ สุชฺฌนฺติ วา วิสุชฺฌนฺติ วา ปริสุชฺฌนฺติ วา, มุจฺจนฺติ วา วิมุจฺจนฺติ วา ปริมุจฺจนฺติ วา, หีนา นิหีนา โอมกา ลามกา ฉตุกฺกา\footnote{ชตุกฺกา (สี, สฺยา)} ปริตฺตาติ เอวํ ปโร วิจฺฉินฺเทติ.\csromanspacer{} เอวํ วิจฺฉินฺทิยมาโน สตฺถารํ นิรสฺสติ, ธมฺมกฺขานํ นิรสฺสติ, คณํ นิรสฺสติ, ทิฏฺฐิํ นิรสฺสติ, ปฏิปทํ นิรสฺสติ, มคฺคํ นิรสฺสติ.\csromanspacer{} เอวํ ปรวิจฺฉินฺทนาย นิรสฺสติ.\csromanspacer{} กถํ อนภิสมฺภูณนฺโต นิรสฺสติ\csromandash{}สีลํ อนภิสมฺภุณนฺโต สีลํ นิรสฺสติ, วตํ อนภิสมฺภุณนฺโต วตํ นิรสฺสติ,}

\vspace{\tipitakaparskip}
\csromancenter{ทุฏฺฐฏฺฐกสุตฺตนิทฺเทส สีลพฺพตํ อนภิสมฺภุณนฺโต สีลพฺพตํ นิรสฺสติ.\csromanspacer{} เอวํ อนภิสมฺภูณนฺโต นิรสฺสติ.\csromanspacer{} \textbf{อาทิยตี จ ธมฺมนฺ}ติ สตฺถารํ คณฺหาติ, ธมฺมกฺขานํ คณฺหาติ, คณํ คณฺหาติ, ทิฏฺฐิํ คณฺหาติ, ปฏิปทํ คณฺหาติ, มคฺคํ คณฺหาติ ปรามสติ อภินิวิสตีติ นิรสฺสตี อาทิยตี จ ธมฺมํ.\csromanspacer{} เตนาห ภควา\csromandash{}}
\vspace{0.5\baselineskip}
\csromangathameasure{ทิฏฺฐีนิเวสา น หิ สฺวาติวตฺตา, \\ ธมฺเมสุ นิจฺเฉยฺย สมุคฺคหีตํ. \\ ตสฺมา นโร เตสุ นิเวสเนสุ, \\ นิรสฺสตี อาทิยตี จ ธมฺมนฺติ.}
\csromangathagroup{\csromangathastanza{\csromanfolio{59}ทิฏฺฐีนิเวสา น หิ สฺวาติวตฺตา, \\ ธมฺเมสุ นิจฺเฉยฺย สมุคฺคหีตํ. \\ ตสฺมา นโร เตสุ นิเวสเนสุ, \\ นิรสฺสตี อาทิยตี จ ธมฺมนฺติ.}}

\proseitem{21}{\csromansymbolfootnote{*}{ขุ 1. 402 ปิฏฺเฐปิ.}โธนสฺส หิ นตฺถิ กุหิญฺจิ โลเก,}

\prose{\textbf{ปกปฺปิตา ทิฏฺฐิ ภวาภเวสุ.}}

\prose{\textbf{มายญฺจ มานญฺจ ปหาย โธโน,}}

\csromangathameasure{ส เกน คจฺเฉยฺย, อนูปโย โส.}
\csromangathagroup{\csromangathastanza{ส เกน คจฺเฉยฺย, อนูปโย โส.}}

\prose{\textbf{โธนสฺส หิ นตฺถิ กุหิญฺจิ โลเก, ปกปฺปิตา ทิฏฺฐิ ภวาภเวสู}ติ \textbf{โธโน}ติ โธนา วุจฺจติ ปญฺญา.\csromanspacer{} ยา ปญฺญา ปชานนา วิจโย ปวิจโย ธมฺมวิจโย สลฺลกฺขณา อุปลกฺขณา ปจฺจุปลกฺขณา ปณฺฑิจฺจํ โกสลฺลํ เนปุญฺญํ เวภพฺยา จินฺตา อุปปริกฺขา ภูริ เมธา ปริณายิกา วิปสฺสนา สมฺปชญฺญํ ปโตโท ปญฺญา ปญฺญินฺทฺริยํ ปญฺญาพลํ ปญฺญาสตฺถํ ปญฺญาปาสาโท ปญฺญา-อาโลโก ปญฺญา-โอภาโส ปญฺญาปชฺโชโต ปญฺญารตนํ อโมโห ธมฺมวิจโย สมฺมาทิฏฺฐิ.\csromanspacer{} กิํการณา โธนา วุจฺจติ ปญฺญา, ตาย ปญฺญาย กายทุจฺจริตํ ธุตญฺจ โธตญฺจ สนฺโธตญฺจ นิทฺโธตญฺจ, วจีทุจฺจริตํ.\csromanspacer{} มโนทุจฺจริตํ ธุตญฺจ โธตญฺจ สนฺโธตญฺจ นิทฺโธตญฺจ.\csromanspacer{} ราโค ธุโต จ โธโต จ สนฺโธโต จ นิทฺโธโต จ, โทโส.\csromanspacer{} โมโห.\csromanspacer{} โกโธ.\csromanspacer{} อุปนาโห.\csromanspacer{} มกฺโข.\csromanspacer{} ปฬาโส ธุโต จ โธโต จ สนฺโธโต จ นิทฺโธโต จ, อิสฺสา ธุตา จ โธตา จ สนฺโธตา จ นิทฺโธตา จ, มจฺฉริยํ ธุตญฺจ โธตญฺจ สนฺโธตญฺจ นิทฺโธตญฺจ, มายา ธุตา จ โธตา จ สนฺโธตา จ นิทฺโธตา จ, สาเฐยฺยํ ธุตญฺจ โธตญฺจ สนฺโธตญฺจ นิทฺโธตญฺจ, ถมฺโภ ธุตา จ โธโต จ สนฺโธโต จ นิทฺโธโต จ, สารมฺโภ.\csromanspacer{} มาโน.\csromanspacer{} อติมาโน.\csromanspacer{} มโท.\csromanspacer{} ปมาโท ธุโต จ โธโต จ สนฺโธโต จ นิทฺโธโต จ, สพฺเพ กิเลสา.\csromanspacer{} สพฺเพ ทุจฺจริตา.\csromanspacer{} สพฺเพ ทรถา.\csromanspacer{} สพฺเพ ปริฬาหา.\csromanspacer{} สพฺเพ สนฺตาปา. \csromanfolio{60}สพฺพากุสลาภิสงฺขารา ธุตา จ โธตา จ สนฺโธตา จ นิทฺโธตา จ, ตํการณา โธนา วุจฺจติ ปญฺญา.}

\prose{อถ วา สมฺมาทิฏฺฐิยา มิจฺฉาทิฏฺฐิ ธุตา จ โธตา จ สนฺโธตา จ นิทฺโธตา จ, สมฺมาสงฺกปฺเปน มิจฺฉาสงฺกปฺโป ธุโต จ โธโต จ สนฺโธโต จ นิทฺโธโต จ, สมฺมาวาจาย มิจฺฉาวาจา ธุตา จ โธตา จ ฯเปฯ สมฺมากมฺมนฺเตน มิจฺฉากมฺมนฺโต ธุโต จ.\csromanspacer{} สมฺมา-อาชีเวน มิจฺฉา-อาชีโว ธุโต จ.\csromanspacer{} สมฺมาวายาเมน มิจฺฉาวายาโม ธุโต จ.\csromanspacer{} สมฺมาสติยา มิจฺฉาสติ ธุตา จ.\csromanspacer{} สมฺมาสมาธินา มิจฺฉาสมาธิ ธุโต จ โธโต จ สนฺโธโต จ นิทฺโธโต จ, สมฺมาญาเณน มิจฺฉาญาณํ ธุตญฺจ.\csromanspacer{} สมฺมาวิมุตฺติสา มิจฺฉาวิมุตฺติ ธุตา จ โธตา จ สนฺโธตา จ นิทฺโธตา จ.}

\prose{อถ วา อริเยน อฏฺฐงฺคิเกน มคฺเคน สพฺเพ กิเลสา.\csromanspacer{} สพฺเพ ทุจฺจริตา.\csromanspacer{} สพฺเพ ทรถา.\csromanspacer{} สพฺเพ ปริฬาหา.\csromanspacer{} สพฺเพ สนฺตาปา.\csromanspacer{} สพฺพากุสลาภิสงฺขารา ธุตา จ โธตา จ สนฺโธตา จ นิทฺโธตา จ.\csromanspacer{} อรหา อิเมหิ โธเนยฺเยหิ ธมฺเมหิ อุเปโต สมุเปโต อุปคโต สมุปคโต อุปปนฺโน สมุปปนฺโน สมนฺนาคโต, ตสฺมา อรหา โธโน.\csromanspacer{} โส ธุตราโค ธุตปาโป ธุตกิเลโส ธุตปริฬาโหติ โธโน.\csromanspacer{}\csromanspacer{}\textbf{กุหิญฺจี}ติ กุหิญฺจิ กิมฺหิจิ กตฺถจิ อชฺฌตฺตํ วา พหิทฺธา วา อชฺฌตฺตพหิทฺธา วา.\csromanspacer{} \textbf{โลเก}ติ อปายโลเก ฯเปฯ อายตนโลเก.}

\prose{\textbf{ปกปฺปิตา}ติ เทฺว ปกปฺปนา ตณฺหาปกปฺปนา จ ทิฏฺฐิปกปฺปนา จ ฯเปฯ อยํ ตณฺหาปกปฺปนา ฯเปฯ อยํ ทิฏฺฐิปกปฺปนา.\csromanspacer{}\csromanspacer{}\textbf{ภวาภเวสู}ติ ภวาภเว กมฺมภเว ปุนพฺภเว กามภเว, กมฺมภเว กามภเว ปุนพฺภเว รูปภเว, กมฺมภเว รูปภเว ปุนพฺภเว อรูปภเว, กมฺมภเว อรูปภเว ปุนพฺภเว ปุนปฺปุนภเว, ปุนปฺปุนคติยา ปุนปฺปุน-อุปปตฺติยา ปุนปฺปุนปฏิสนฺธิยา ปุนปฺปุน- อตฺตภาวาภินิพฺพตฺติยา.\csromanspacer{}\csromanspacer{}\textbf{โธนสฺส หิ นตฺถิ กุหิญฺจิ โลเก, ปกปฺปิตา ทิฏฺฐิ ภวาภเวสู}ติ โธนสฺส กุหิญฺจิ โลเก ภวาภเวสุ จ กปฺปิตา ปกปฺปิตา อภิสงฺขตา สณฺฐปิตา ทิฏฺฐิ นตฺถิ น สนฺติ น สํวิชฺชติ นุปลพฺภติ, ปหีนา สมุจฺฉินฺนา วูปสนฺตา ปฏิปสฺสทฺธา อภพฺพุปฺปตฺติกา ญาณคฺคินา ทฑฺฒาติ โธนสฺส หิ นตฺถิ กุหิญฺจิ โลเก, ปกปฺปิตา ทิฏฺฐิ ภวาภเวสุ.}

\prose{\textbf{มายญฺจ มานญฺจ ปหาย โธโน}ติ \textbf{มายา} วุจฺจติ วญฺจนิกา จริยา. * อิเธกจฺโจ กาเยน ทุจฺจริตํ จริตฺวา วาจาย ทุจฺจริตํ จริตฺวา มนสา}

\vspace{\tipitakaparskip}
\csromancenter{ทุฏฺฐฏฺฐกสุตฺตนิทฺเทส ทุจฺจริตํ จริตฺวา ตสฺส ปฏิจฺฉาทนเหตุ ปาปิกํ อิจฺฉํ ปณิทหติ\csromandash{}“มา มํ ชญฺญา”ติ อิจฺฉติ, “มา มํ ชญฺญา”ติ สงฺกปฺเปติ, “มา มํ ชญฺญา”ติ วาจํ ภาสติ, “มา มํ ชญฺญา”ติ กาเยน ปรกฺกมติ, ยา เอวรูปา มายา มายาวิตา อจฺจาสรา วญฺจนา นิกติ นิกิรณา ปริหรณา คูหนา ปริคูหนา ฉาทนา ปริจฺฉาทนา อนุตฺตานิกมฺมํ อนาวิกมฺมํ โวจฺฉาทนา ปาปกิริยา, อยํ วุจฺจติ มายา.}
\prose{\csromanfolio{61}\textbf{มาโน}ติ เอกวิเธน มาโน, ยา จิตฺตสฺส อุนฺนติ\footnote{อุณฺณติ (สฺยา, ก)}.\csromanspacer{}\csromanspacer{}ทุวิเธน มาโน, อตฺตุกฺกํสนมาโน ปรวมฺภนมาโน.\csromanspacer{}\csromanspacer{}ติวิเธน มาโน, “เสยฺโยหมสฺมี”ติ มาโน “สทิโสหมสฺมี”ติ มาโน “หีโนหมสฺมี”ติ มาโน.\csromanspacer{}\csromanspacer{}จตุพฺพิเธน มาโน, ลาเภน มานํ ชเนติ, ยเสน มานํ ชเนติ, ปสํสาย มานํ ชเนติ, สุเขน มานํ ชเนติ.\csromanspacer{}\csromanspacer{}ปญฺจวิเธน มาโน, “ลาภิมฺหิ มนาปิกานํ รูปานนฺ”ติ มานํ ชเนติ, “ลาภิมฺหิ มนาปิกานํ สทฺทานํ.\csromanspacer{} คนฺธานํ.\csromanspacer{} รสานํ.\csromanspacer{} โผฏฺฐพฺพานนฺ”ติ มานํ ชเนติ.\csromanspacer{}\csromanspacer{}ฉพฺพิเธน มาโน, จกฺขุสมฺปทาย มานํ ชเนติ, โสตสมฺปทาย.\csromanspacer{} ฆานสมฺปทาย.\csromanspacer{} ชิวฺหาสมฺปทาย.\csromanspacer{} กายสมฺปทาย.\csromanspacer{} มโนสมฺปทาย มานํ ชเนติ.\csromanspacer{}\csromanspacer{}สตฺตวิเธน มาโน, มาโน อติมาโน มานาติมาโน โอมาโน อธิมาโน อสฺมิมาโน มิจฺฉามาโน.\csromanspacer{}\csromanspacer{}อฏฺฐวิเธน มาโน, ลาเภน มานํ ชเนติ, อลาเภน โอมานํ ชเนติ, ยเสน มานํ ชเนติ, อยเสน โอมานํ ชเนติ, ปสํสาย มานํ ชเนติ, นินฺทาย โอมานํ ชเนติ, สุเขน มานํ ชเนติ, ทุกฺเขน โอมานํ ชเนติ.\csromanspacer{}\csromanspacer{}นววิเธน มาโน, เสยฺยสฺส * “เสยฺโยหมสฺมี”ติ มาโน, เสยฺยสฺส “สทิโสหมสฺมี”ติ มาโน, เสยฺยสฺส “หีโนหมสฺมี”ติ มาโน.\csromanspacer{} สทิสสฺส “เสยฺโยหมสฺมี”ติ มาโน, สทิสสฺส “สทิโสหมสฺมี”ติ มาโน, สทิสสฺส “หีโนหมสฺมี”ติ มาโน.\csromanspacer{} หีนสฺส “เสยฺโยหมสฺมี”ติ มาโน, หีนสฺส “สทิโสหมสฺมี”ติ มาโน, หีนสฺส “หีโนหมสฺมี”ติ มาโน.\csromanspacer{}\csromanspacer{}ทสวิเธน มาโน, อิเธกจฺโจ มานํ ชเนติ ชาติยา วา โคตฺเตน วา โกลปุตฺติเยน วา วณฺณโปกฺขรตาย วา ธเนน วา อชฺเฌเนน วา กมฺมายตเนน วา สิปฺปายตเนน วา วิชฺชาฏฺฐาเนน วา สุเตน วา ปฏิภาเนน วา อญฺญตรญฺญตเรน วา วตฺถุนา. + โย เอวรูโป มาโน มญฺญนา มญฺญิตตฺตํ อุนฺนติ อุนฺนโม\footnote{อุณฺณโม (สฺยา, ก)} ธโช \csromanfolio{62}สมฺปคฺคาโห เกตุกมฺยตา จิตฺตสฺส.\csromanspacer{} อยํ วุจฺจติ มาโน.\csromanspacer{}\csromanspacer{}\textbf{มายญฺจ มานญฺจ ปหาย โธโน}ติ โธโน มายญฺจ มานญฺจ ปหาย ปชหิตฺวา วิโนเทตฺวา พฺยนฺติํ กริตฺวา อนภาวํ คเมตฺวาติ มายญฺจ มานญฺจ ปหาย โธโน.}

\prose{\textbf{ส เกน คจฺเฉยฺย, อนูปโย โส}ติ \textbf{อุปยา}ติ เทฺว \textbf{อุปยา} ตณฺหูปโย จ ทิฏฺฐูปโย จ ฯเปฯ อยํ ตณฺหูปโย ฯเปฯ อยํ ทิฏฺฐูปโย.\csromanspacer{} ตสฺส ตณฺหูปโย ปหีโน, ทิฏฺฐูปโย ปฏินิสฺสฏฺโฐ.\csromanspacer{} ตณฺหูปยสฺส ปหีนตฺตา, ทิฏฺฐูปยสฺส ปฏินิสฺสฏฺฐตฺตา อนูปโย ปุคฺคโล เกน ราเคน คจฺเฉยฺย, เกน โทเสน คจฺเฉยฺย, เกน โมเหน คจฺเฉยฺย, เกน มาเนน คจฺเฉยฺย, กาย ทิฏฺฐิยา คจฺเฉยฺย, เกน อุทฺธจฺเจน คจฺเฉยฺย, กาย วิจิกิจฺฉาย คจฺเฉยฺย, เกหิ อนุสเยหิ คจฺเฉยฺย “รตฺโต”ติ วา “ทุฏฺโฐ”ติ วา “มูโฬฺห”ติ วา “วินิพทฺโธ”ติ วา “ปรามฏฺโฐ”ติ วา “วิกฺเขปคโต”ติ วา “อนิฏฺฐงฺคโต”ติ วา “ถามคโต”ติ วา.\csromanspacer{} เต อภิสงฺขารา ปหีนา, อภิสงฺขารานํ ปหีนตฺตา คติโย เกน คจฺเฉยฺย “เนรยิโก”ติ วา “ติรจฺฉานโยนิโก”ติ วา “เปตฺติวิสยิโก”ติ วา “มนุสฺโส”ติ วา “เทโว”ติ วา “รูปี”ติ วา “อรูปี”ติ วา “สญฺญี”ติ วา “อสญฺญี”ติ วา “เนวสญฺญีนาสญฺญี”ติ วา.\csromanspacer{} โส เหตุ นตฺถิ, ปจฺจโย นตฺถิ, การณํ นตฺถิ, เยน คจฺเฉยฺยาติ ส เกน คจฺเฉยฺย, อนูปโย โส.\csromanspacer{} เตนาห ภควา\csromandash{}}

\csromangathameasure{โธนสฺส หิ นตฺถิ กุหิญฺจิ โลเก, \\ ปกปฺปิตา ทิฏฺฐิ ภวาภเวสุ. \\ \textbf{มายญฺจ มานญฺจ ปหาย โธโน},}
\csromangathagroup{\csromangathastanza{โธนสฺส หิ นตฺถิ กุหิญฺจิ โลเก, \\ ปกปฺปิตา ทิฏฺฐิ ภวาภเวสุ. \\ \textbf{มายญฺจ มานญฺจ ปหาย โธโน},}}

\prose{\textbf{ส เกน คจฺเฉยฺย, อนูปโย โส}ติ.}

\proseitem{22}{\csromansymbolfootnote{*}{ขุ 1. 402 ปิฏฺเฐปิ.}อุปโย หิ ธมฺเมสุ อุเปติ วาทํ,}

\prose{\textbf{อนูปยํ เกน กถํ วเทยฺย.}}

\csromangathameasure{\textbf{อตฺตา นิรตฺตา น หิ ตสฺส อตฺถิ,} \\ \textbf{อโธสิ โส ทิฏฺฐิมิเธว สพฺพํ.}}
\csromangathagroup{\csromangathastanza{\textbf{อตฺตา นิรตฺตา น หิ ตสฺส อตฺถิ,} \\ \textbf{อโธสิ โส ทิฏฺฐิมิเธว สพฺพํ.}}}

\prose{\textbf{อุปโย หิ ธมฺเมสุ อุเปติ วาทนฺ}ติ \textbf{อุปยา}ติ เทฺว \textbf{อุปยา} ตณฺหูปโย จ ทิฏฺฐูปโย จ ฯเปฯ อยํ ตณฺหูปโย ฯเปฯ อยํ ทิฏฺฐูปโย.\csromanspacer{} ตสฺส ตณฺหูปโย อปฺปหีโน, ทิฏฺฐูปโย อปฺปฏินิสฺสฏฺโฐ.\csromanspacer{} ตณฺหูปยสฺส}

\vspace{\tipitakaparskip}
\csromancenter{ทุฏฺฐฏฺฐกสุตฺตนิทฺเทส อปฺปหีนฺ\textbf{อตฺตา}, ทิฏฺฐูปยสฺส อปฺปฏินิสฺสฏฺฐฺ\textbf{อตฺตา} ธมฺเมสุ วาทํ อุเปติ “รตฺโต”ติ วา “ทุฏฺโฐ”ติ วา “มูฏฺโฐ”ติ วา “วินิพทฺโธ”ติ วา “ปรามฏฺโฐ”ติ วา “วิกฺเขปคโต”ติ วา “อนิฏฺฐงฺคโต”ติ วา “ถามคโต”ติ วา.\csromanspacer{} เต อภิสงฺขารา อปฺปหีนา, อภิสงฺขารานํ อปฺปหีนฺ\textbf{อตฺตา} คติยา วาทํ อุเปติ “เนรยิโก”ติ วา “ติรจฺฉานโยนิโก”ติ วา “เปตฺติวิสยิโก”ติ วา “มนุสฺโส”ติ วา “เทโว”ติ วา “รุปี”ติ วา “อรูปี”ติ วา “สญฺญี”ติ วา “อสญฺญี”ติ วา “เนวสญฺญีนาสญฺญี”ติ วา วาทํ อุเปติ อุปคจฺฉติ คณฺหาติ ปรามสติ อภินิวิสตีติ อุปโย หิ ธมฺเมสุ อุเปติ วาทํ.}
\prose{\csromanfolio{63}\textbf{อนูปยํ เกน กถํ วเทยฺยา}ติ \textbf{อุปยา}ติ เทฺว \textbf{อุปยา} ตณฺหูปโย จ ทิฏฺฐูปโย จ ฯเปฯ อยํ ตณฺหูปโย ฯเปฯ อยํ ทิฏฺฐูปโย.\csromanspacer{} ตสฺส ตณฺหูปโย ปหีโน ทิฏฺฐูปโย ปฏินิสฺสฏฺโฐ.\csromanspacer{} ตณฺหูปยสฺส ปหีนฺ\textbf{อตฺตา}, ทิฏฺฐูปยสฺส ปฏินิสฺสฏฺฐฺ\textbf{อตฺตา} อนูปยํ ปุคฺคลํ เกน ราเคน วเทยฺย, เกน โทเสน วเทยฺย, เกน โมเหน วเทยฺย, เกน มาเนน วเทยฺย, กาย ทิฏฺฐิยา วเทยฺย, เกน อุทฺธจฺเจน วเทยฺย, กาย วิจิกิจฺฉาย วเทยฺย, เกหิ อนุสเยหิ วเทยฺย “รตฺโต”ติ วา “ทุฏฺโฐ”ติ วา “มูโฬฺห”ติ วา “วินิพทฺโธ”ติ วา “ปรามฏฺโฐ”ติ วา “วิกฺเขปคโต”ติ วา “อนิฏฺฐงฺคโต”ติ วา “ถามคโต”ติ วา.\csromanspacer{} เต อภิสงฺขารา ปหีนา, อภิสงฺขารานํ ปหีนฺ\textbf{อตฺตา} คติโย เกน วเทยฺย “เนรยิโก”ติ วา ฯเปฯ “เนวสญฺญีนาสญฺญี”ติ วา.\csromanspacer{} โส เหตุ นตฺถิ, ปจฺจโย นตฺถิ, การณํ นตฺถิ, เยน วเทยฺย กเถยฺย ภเณยฺย ทีปเยยฺย โวหเรยฺยาติ อนูปยํ เกน กถํ วเทยฺย.}

\prose{\textbf{อตฺตา นิรตฺตา น หิ ตสฺส อตฺถี}ติ \textbf{อตฺตา}ติ \textbf{อตฺตา}นุทิฏฺฐิ นตฺถิ.\csromanspacer{} \textbf{นิรตฺตา}ติ อุจฺเฉททิฏฺฐิ นตฺถิ.\csromanspacer{} \textbf{อตฺตา}ติ คหิตํ นตฺถิ.\csromanspacer{} \textbf{นิรตฺตา}ติ มุญฺจิตพฺพํ นตฺถิ.\csromanspacer{} ยสฺสตฺถิ คหิตํ, ตสฺสตฺถิ มุญฺจิตพฺพํ, ยสฺสตฺถิ มุญฺจิตพฺพํ, ตสฺสตฺถิ คหิตํ, คหณํ มุญฺจนา สมติกฺกนฺโต อรหา พุทฺธิปริหานิวีติวตฺโต.\csromanspacer{} โส วุฏฺฐวาโส จิณฺณจรโณ คตทฺโธ คตทิโส ชาติมรณสํสาโร นตฺถิ ตสฺส ปุนพฺภโวติ \textbf{อตฺตา} นิรฺ\textbf{อตฺตา} น หิ ตสฺส อตฺถิ.}

\prose{\textbf{อโธสิ โส ทิฏฺฐิมิเธว สพฺพนฺ}ติ ตสฺส ทฺวาสฏฺฐิ ทิฏฺฐิคตานิ ปหีนานิ สมุจฺฉินฺนานิ วูปสนฺตานิ ปฏิปสฺสทฺธานิ อภพฺพุปฺปตฺติกานิ ญาณคฺคินา ทฑฺฒานิ.\csromanspacer{} โส สพฺพทิฏฺฐิคตํ อิเธว อโธสิ ธุนิ สนฺธุนิ นิทฺธุนิ ปชหิ วิโนเทสิ \csromanfolio{64}พฺยนฺติํ อกาสิ อนภาวํ คเมสีติ อโธสิ โส ทิฏฺฐิมิเธว สพฺพํ.\csromanspacer{} เตนาห ภควา\csromandash{}}

\prose{อุปโย หิ ธมฺเมสุ อุเปติ วาทํ.}

\prose{อนูปยํ เกน กถํ วเทยฺย.}

\prose{อตฺตา นิรตฺตา น หิ ตสฺส อตฺถิ.}

\prose{อโธสิ โส ทิฏฺฐิมิเธว สพฺพนฺติ.\csromanspacer{} ทุฏฺฐฏฺฐกสุตฺตนิทฺเทโส ตติโย.}
\csromansectionrule

\csromanheader{2}{4. สุทฺธฏฺฐกสุตฺตนิทฺเทส}
\csromanheader{2}{อถ สุทฺธฏฺฐกสุตฺตนิทฺเทสํ วกฺขติ\csromandash{}}
\proseitem{23}{\csromansymbolfootnote{*}{ขุ 1. 402 ปิฏฺเฐปิ.}ปสฺสามิ สุทฺธํ ปรมํ อโรคํ,}

\prose{\textbf{ทิฏฺเฐน สํสุทฺธิ นรสฺส โหติ.}}

\csromangathameasure{\textbf{เอวาภิชานํ “ปรมนฺ”ติ ญตฺวา,} \\ \textbf{สุทฺธานุปสฺสีติ ปจฺเจติ ญาณํ.}}
\csromangathagroup{\csromangathastanza{\textbf{เอวาภิชานํ “ปรมนฺ”ติ ญตฺวา,} \\ \textbf{สุทฺธานุปสฺสีติ ปจฺเจติ ญาณํ.}}}

\prose{\textbf{ปสฺสามิ สุทฺธํ ปรมํ อโรคนฺ}ติ \textbf{ปสฺสามิ สุทฺธนฺ}ติ ปสฺสามิ สุทฺธํ, ทกฺขามิ สุทฺธํ โอโลเกมิ สุทฺธํ นิชฺฌายามิ สุทฺธํ อุปปริกฺขามิ สุทฺธํ.\csromanspacer{} \textbf{ปรมํ อโรคนฺ}ติ ปรมํ อาโรคฺยปฺปตฺตํ ตาณปฺปตฺตํ เลณปฺปตฺตํ สรณปฺปตฺตํ อภยปฺปตฺตํ อจฺจุตปฺปตฺตํ อมตปฺปตฺตํ นิพฺพานปฺปตฺตนฺติ ปสฺสามิ สุทฺธํ ปรมํ อโรคํ.}

\prose{\textbf{ทิฏฺเฐน สํสุทฺธิ นรสฺส โหตี}ติ จกฺขุวิญฺญาณํ\footnote{จกฺขุวิญฺญาเณน (สี, สฺยา)} รูปทสฺสเนน นรสฺส สุทฺธิ วิสุทฺธิ ปริสุทฺธิ, มุตฺติ วิมุตฺติ ปริมุตฺติ โหติ, นโร สุชฺฌติ วิสุชฺฌติ ปริสุชฺฌติ, มุจฺจติ วิมุจฺจติ ปริมุจฺจตีติ ทิฏฺเฐน สํสุทฺธิ นรสฺส โหติ.}

\prose{\textbf{เอวาภิชานํ “ปรมนฺ”ติ ญตฺวา}ติ เอวํ อภิชานนฺโต อาชานนฺโต วิชานนฺโต ปฏิวิชานนฺโต ปฏิวิชฺฌนฺโต “อิทํ ปรมํ อคฺคํ เสฏฺฐํ วิสิฏฺฐํ ปาโมกฺขํ อุตฺตมํ ปวรนฺ”ติ ญตฺวา ชานิตฺวา ตุลยิตฺวา ตีรยิตฺวา วิภาวยิตฺวา วิภูตํ กตฺวาติ เอวาภิชานํ “ปรมนฺ”ติ ญตฺวา.}

\csromanheader{2}{4. สุทฺธฏฺฐกสุตฺตนิทฺเทส}
\prose{\csromanfolio{65}\textbf{สุทฺธานุปสฺสีติ ปจฺเจติ ญาณนฺ}ติ โย สุทฺธํ ปสฺสติ, โส สุทฺธานุปสฺสี, \textbf{ปจฺเจติ ญาณนฺ}ติ จกฺขุวิญฺญาณํ รูปทสฺสเนน ญาณนฺติ ปจฺเจติ, มคฺโคติ ปจฺเจติ, ปโถติ ปจฺเจติ, นิยฺยานนฺติ ปจฺเจตีติ สุทฺธานุปสฺสี ปจฺเจติ ญาณํ.\csromanspacer{} เตนาห ภควา\csromandash{}}

\csromangathameasure{ปสฺสามิ สุทฺธํ ปรมํ อโรคํ, \\ ทิฏฺเฐน สํสุทฺธิ นรสฺส โหติ. \\ เอวาภิชานํ “ปรมนฺ”ติ ญตฺวา,}
\csromangathagroup{\csromangathastanza{ปสฺสามิ สุทฺธํ ปรมํ อโรคํ, \\ ทิฏฺเฐน สํสุทฺธิ นรสฺส โหติ. \\ เอวาภิชานํ “ปรมนฺ”ติ ญตฺวา,}}

\prose{\textbf{สุทฺธานุปสฺสีติ ปจฺเจติ ญาณนฺ}ติ.}

\proseitem{24}{\csromansymbolfootnote{*}{ขุ 1. 402 ปิฏฺเฐปิ.}ทิฏฺเฐน เจ สุทฺธิ นรสฺส โหติ,}

\prose{\textbf{ญาเณน วา โส ปชหาติ ทุกฺขํ.}}

\csromangathameasure{\textbf{อญฺเญน โส สุชฺฌติ โสปธีโก,} \\ \textbf{ทิฏฺฐี หิ นํ ปาว ตถา วทานํ.}}
\csromangathagroup{\csromangathastanza{\textbf{อญฺเญน โส สุชฺฌติ โสปธีโก,} \\ \textbf{ทิฏฺฐี หิ นํ ปาว ตถา วทานํ.}}}

\prose{\textbf{ทิฏฺเฐน เจ สุทฺธิ นรสฺส โหตี}ติ จกฺขุวิญฺญาณํ รูปทสฺสเนน เจ นรสฺส สุทฺธิ วิสุทฺธิ ปริสุทฺธิ, มุตฺติ วิมุตฺติ ปริมุตฺติ โหติ, นโร สุชฺฌติ วิสุชฺฌติ ปริสุชฺฌติ, มุจฺจติ วิมุจฺจติ ปริมุจฺจตีติ ทิฏฺเฐน เจ สุทฺธิ นรสฺส โหติ.}

\prose{\textbf{ญาเณน วา โส ปชหาติ ทุกฺขนฺ}ติ จกฺขุวิญฺญาณํ รูปทสฺสเนน เจ นโร ชาติทุกฺขํ ปชหติ, ชราทุกฺขํ ปชหติ, พฺยาธิทุกฺขํ ปชหติ, มรณทุกฺขํ ปชหติ, โสกปริเทวทุกฺขโทมนสฺสุปายาสทุกฺขํ ปชหตีติ ญาเณน วา โส ปชหาติ ทุกฺขํ.}

\prose{\textbf{อญฺเญน โส สุชฺฌติ โสปธีโก}ติ อญฺเญน อสุทฺธิมคฺเคน มิจฺฉาปฏิปทาย อนิยฺยานิกปเถน อญฺญตฺร สติปฏฺฐาเนหิ อญฺญตฺร สมฺมปฺปธาเนหิ อญฺญตฺร อิทฺธิปาเทหิ อญฺญตฺร อินฺทฺริเยหิ อญฺญตฺร พเลหิ อญฺญตฺร โพชฺฌงฺเคหิ อญฺญตฺร อริยา อฏฺฐงฺคิกา มคฺคานโร สุชฺฌติ วิสุชฺฌติ ปริสุชฺฌติ, มุจฺจติ วิมุจฺจติ ปริมุจฺจติ.\csromanspacer{}\csromanspacer{}\textbf{โสปธีโก}ติ สราโค สโทโส สโมโห สมาโน สตโณฺห สทิฏฺฐิ สกิเลโส ส-อุปาทาโนติ อญฺเญน โส สุชฺฌติ โสปธิโก.}

\prose{\textbf{ทีฏฺฐี หิ นํ ปาว ตถา วทานนฺ}ติ สาว ทิฏฺฐิ ตํ ปุคฺคลํ ปาวทติ, อิติ วายํ ปุคฺคโล มิจฺฉาทิฏฺฐิโก วิปรีตทสฺสโน.\csromanspacer{}\csromanspacer{}\textbf{ตถา วทานนฺ}ติ ตถา วทนฺตํ กเถนฺตํ ภณนฺตํ ทีปยนฺตํ โวหรนฺตํ, “สสฺสโต \csromanfolio{66}โลโก, อิทเมว สจฺจํ โมฆมญฺญนฺ”ติ ตถา วทนฺตํ กเถนฺตํ ภณนฺตํ ทีปยนฺตํ โวหรนฺตํ, “อสสฺสโต โลโก.\csromanspacer{} อนฺตวา โลโก.\csromanspacer{} อนนฺตวา โลโก.\csromanspacer{} ตํ ชีวํ ตํ สรีรํ.\csromanspacer{} อญฺญํ ชีวํ อญฺญํ สรีรํ.\csromanspacer{} โหติ ตถาคโต ปรํ มรณา.\csromanspacer{} น โหติ ตถาคโต ปรํ มรณา.\csromanspacer{} โหติ จ น จ โหติ ตถาคโต ปรํ มรณา.\csromanspacer{} เนว โหติ น น โหติ ตถาคโต ปรํ มรณา, อิทเมว สจฺจํ โมฆมญฺญนฺ”ติ ตถา วทนฺตํ กเถนฺตํ ภณนฺตํ ทีปยนฺตํ โวหรนฺตนฺติ ทิฏฺฐี หิ นํ ปาว ตถา วทานํ.\csromanspacer{} เต\textbf{นา}ห ภควา\csromandash{}}

\csromangathameasure{ทิฏฺเฐน เจ สุทฺธิ นรสฺส โหติ, \\ ญาเณน วา โส ปชหาติ ทุกฺขํ. \\ อญฺเญน โส สุชฺฌติ โสปธีโก, \\ ทีฏฺฐี หิ นํ ปาว ตถา วทานนฺติ.}
\csromangathagroup{\csromangathastanza{ทิฏฺเฐน เจ สุทฺธิ นรสฺส โหติ, \\ ญาเณน วา โส ปชหาติ ทุกฺขํ. \\ อญฺเญน โส สุชฺฌติ โสปธีโก, \\ ทีฏฺฐี หิ นํ ปาว ตถา วทานนฺติ.}}

\proseitem{25}{\csromansymbolfootnote{*}{ขุ 1. 402 ปิฏฺเฐปิ.}น \textbf{พฺราหฺมโณ} อญฺญโต สุทฺธิมาห,}

\prose{\textbf{ทิฏฺเฐ สุเต สีลวเต มุเต วา.}}

\csromangathameasure{\textbf{ปุญฺเญ จ ปาเป จ อนูปลิตฺโต,} \\ \textbf{อตฺตญฺชโห นยิธ ปกุพฺพมาโน.}}
\csromangathagroup{\csromangathastanza{\textbf{ปุญฺเญ จ ปาเป จ อนูปลิตฺโต,} \\ \textbf{อตฺตญฺชโห นยิธ ปกุพฺพมาโน.}}}

\prose{\textbf{น พฺราหฺมโณ อญฺญโต สุทฺธิมาห, ทิฏฺเฐ สุเต สีลวเต มุเต วา}ติ \textbf{นา}ติ ปฏิกฺเขโป, \textbf{พฺราหฺมโณ}ติ สตฺตนฺนํ ธมฺมานํ พาหิตตฺตา \textbf{พฺราหฺมโณ}, สกฺกายทิฏฺฐิ พาหิตา โหติ, วิจิกิจฺฉา พาหิตา โหติ, สีลพฺพตปรามาโส พาหิโต โหติ, ราโค พาหิโต โหติ, โทโส พาหิโต โหติ, โมโห พาหิโต โหติ, มาโน พาหิโต โหติ, พาหิตาสฺส โหนฺติ ปาปกา อกุสลา ธมฺมา สํกิเลสิกา โปโนภวิกา สทรา ทุกฺขวิปากา อายติํ ชาติชรามรณิยา\csromandash{}}

\prose{\csromansymbolfootnote{+}{ขุ 1. 358; ขุ 8. 84 ปิฏฺเฐสุปิ.}พาหิตฺวา สพฺพปาปกานิ, (สภิยาติ ภควา,)}

\prose{วิมโล สาธุสมาหิโต ฐิตตฺโต.}

\csromangathameasure{สํสารมติจฺจ เกวลี โส, \\ อสิโต ตาทิ ปวุจฺจเต ส พฺรหฺมา.}
\csromangathagroup{\csromangathastanza{สํสารมติจฺจ เกวลี โส, \\ อสิโต\footnote{อนิสฺสิโต (สฺยา)} ตาทิ ปวุจฺจเต ส พฺรหฺมา.}}

\csromanheader{2}{4. สุทฺธฏฺฐกสุตฺตนิทฺเทส}
\prose{\csromanfolio{67}\textbf{น พฺราหฺมโณ อญฺญโต สุทฺธิมาหา}ติ พฺราหฺมโณ อญฺเญน อสุทฺธิมคฺเคน มิจฺฉาปฏิปทาย อนิยฺยานิกปเถน อญฺญตฺร สติปฏฺฐาเนหิ อญฺญตฺร สมฺมปฺปธาเนหิ อญฺญตฺร อิทฺธิปาเทหิ อญฺญตฺร อินฺทฺริเยหิ อญฺญตฺร พเลหิ อญฺญตฺร โพชฺฌงฺเคหิ อญฺญตฺร อริเยน อฏฺฐงฺคิเกน มคฺเคน สุทฺธิํ วิสุทฺธิํ ปริสุทฺธิํ, มุตฺติํ วิมุตฺติํ ปริมุตฺติํ, นาห น กเถติ น ภณติ น ทีปยติ น โวหรตีติ น พฺราหฺมโณ อญฺญโต สุทฺธิมาห.}

\prose{\textbf{ทิฏฺเฐ สุเต สีลวเต มุเต วา}ติ สนฺเตเก สมณพฺราหฺมณา ทิฏฺฐิสุทฺธิกา, เต เอกจฺจานํ รูปานํ ทสฺสนํ มงฺคลํ ปจฺเจนฺติ, เอกจฺจานํ รูปานํ ทสฺสนํ อมงฺคลํ ปจฺเจนฺติ.\csromanspacer{} กตเมสํ รูปานํ ทสฺสนํ มงฺคลํ ปจฺเจนฺติ, เต กาลโต วุฏฺฐหิตฺวา อภิมงฺคลคตานิ รูปานิ ปสฺสนฺติ, จาฏกสกุณํ\footnote{วาตสกุณํ (สฺยา), จาปสกุณํ (ก)} ปสฺสนฺติ, ผุสฺสเวฬุวลฏฺฐิํ ปสฺสนฺติ, คพฺภินิตฺถิํ ปสฺสนฺติ, กุมารกํ ขนฺเธ อาโรเปตฺวา คจฺฉนฺตํ ปสฺสนฺติ, ปุณฺณฆฏํ ปสฺสนฺติ, โรหิตมจฺฉํ ปสฺสนฺติ, อาชญฺญํ ปสฺสนฺติ, อาชญฺญรถํ ปสฺสนฺติ, อุสภํ ปสฺสนฺติ, โคกปิลํ ปสฺสนฺติ, เอวรูปานํ รูปานํ ทสฺสนํ มงฺคลํ ปจฺเจนฺติ.\csromanspacer{} กตเมสํ รูปานํ ทสฺสนํ อมงฺคลํ ปจฺเจนฺติ, ปลาลปุญฺชํ ปสฺสนฺติ, ตกฺกฆฏํ ปสฺสนฺติ, ริตฺตฆฏํ ปสฺสนฺติ, นฏํ ปสฺสนฺติ, นคฺคสมณกํ ปสฺสนฺติ, ขรํ ปสฺสนฺติ, ขรยานํ ปสฺสนฺติ, เอกยุตฺตยานํ ปสฺสนฺติ, กาณํ ปสฺสนฺติ, กุณิํ ปสฺสนฺติ, ขญฺชํ ปสฺสนฺติ, ปกฺขหตํ\footnote{ปกฺขปาทํ (ก)} ปสฺสนฺติ, ชิณฺณกํ ปสฺสนฺติ, พฺยาธิกํ\footnote{พฺยาธิตํ (สี)} ปสฺสนฺติ, มตํ ปสฺสนฺติ, เอวรูปานํ รูปานํ ทสฺสนํ อมงฺคลํ ปจฺเจนฺติ.\csromanspacer{} อิเม เต สมณพฺราหฺมณา ทิฏฺฐิสุทฺธิกา.\csromanspacer{} เต ทิฏฺเฐน สุทฺธิํ วิสุทฺธิํ ปริสุทฺธิํ, มุตฺติํ วิมุตฺติํ ปริมุตฺติํ ปจฺเจนฺติ.}

\prose{สนฺเตเก สมณพฺราหฺมณา สุตสุทฺธิกา, เต เอกจฺจานํ สทฺทานํ สวนํ มงฺคลํ ปจฺเจนฺติ, เอกจฺจานํ สทฺทานํ สวนํ อมงฺคลํ ปจฺเจนฺติ.\csromanspacer{} กตเมสํ สทฺทานํ สวนํ มงฺคลํ ปจฺเจนฺติ, เต กาลโต วุฏฺฐหิตฺวา อภิมงฺคลคตานิ สทฺทานิ สุณนฺติ “วฑฺฒา”ติ วา “วฑฺฒมานา”ติ วา “ปุณฺณา”ติ วา “ผุสฺสา”ติ วา “อโสกา”ติ วา “สุมนา”ติ วา “สุนกฺขตฺตา”ติ วา “สุมงฺคลา”ติ วา “สิรี”ติ วา “สิรีวฑฺฒา”ติ วา เอวรูปานํ สทฺทานํ สวนํ มงฺคลํ ปจฺเจนฺติ.\csromanspacer{}\csromanspacer{}กตเมสํ สทฺทานํ สวนํ อมงฺคลํ ปจฺเจนฺติ, “กาโณ”ติ วา “กุณี”ติ วา “ขญฺโช”ติ วา “ปกฺขหโต”ติ วา “ชิณฺณโก”ติ วา “พฺยาธิโก”ติ วา “มโต”ติ วา “ฉินฺทนฺ”ติ วา “ภินฺทนฺ”ติ วา “ทฑฺฒนฺ”ติ วา “นฏฺฐนฺ”ติ วา “นตฺถี”ติ วา เอวรูปานํ สทฺทานํ}

\prose{\csromanfolio{68}สวนํ อมงฺคลํ ปจฺเจนฺติ.\csromanspacer{} อิเม เต สมณพฺราหฺมณา สุตสุทฺธิกา, เต สุเตน สุทฺธิํ วิสุทฺธิํ ปริสุทฺธิํ, มุตฺติํ วิมุตฺติํ ปริมุตฺติํ ปจฺเจนฺติ.}

\prose{สนฺเตเก สมณพฺราหฺมณา สีลสุทฺธิกา, เต สีลมตฺเตน สํยมมตฺเตน สํวรมตฺเตน อวีติกฺกมมตฺเตน สุทฺธิํ วิสุทฺธิํ ปริสุทฺธิํ, มุตฺติํ วิมุตฺติํ ปริมุตฺติํ ปจฺเจนฺติ.\csromanspacer{} สมณมุณฺฑิกาปุตฺโต\footnote{สมโณ มณฺฑิกาปุตฺโต (สี), ม 2. 215 ปิฏฺเฐ ปสฺสิตพฺพํ.} เอวมาห\csromandash{} จตูหิ โข อหํ คหปติ\footnote{ถปติ (สี, สฺยา) เอวมุปริปิ.} ธมฺเมหิ สมนฺนาคตํ ปุริสปุคฺคลํ ปญฺญาเปมิ สมฺปนฺนกุสลํ ปรมกุสลํ อุตฺตมปตฺติปฺปตฺตํ สมณํ อโยชฺฌํ.\csromanspacer{} กตเมหิ จตูหิ\csromandash{}อิธ คหปติ น กาเยน ปาปกํ กมฺมํ กโรติ, น ปาปิกํ วาจํ ภาสติ, น ปาปกํ สงฺกปฺปํ สงฺกปฺเปติ, น ปาปกํ อาชีวํ อาชีวติ.\csromanspacer{} อิเมหิ โข อหํ คหปติ จตูหิ ธมฺเมหิ สมนฺนาคตํ ปุริสปุคฺคลํ ปญฺญาเปมิ สมฺปนฺนกุสลํ ปรมกุสลํ อุตฺตมปตฺติปฺปตฺตํ สมณํ อโยชฺฌํ.\csromanspacer{} เอวเมวํ สนฺเตเก สมณพฺราหฺมณา สีลสุทฺธิกา, เต สีลมตฺเตน สํยมมตฺเตน สํวรมตฺเตน อวีติกฺกมมตฺเตน สุทฺธิํ วิสุทฺธิํ ปริสุทฺธิํ, มุตฺติํ วิมุตฺติํ ปริมุตฺติํ ปจฺเจนฺติ.}

\prose{สนฺเตเก สมณพฺราหฺมณา วตสุทฺธิกา, เต หตฺถิวติกา วา โหนฺติ, อสฺสวติกา วา โหนฺติ, โควติกา วา โหนฺติ, กุกฺกุรวติกา วา โหนฺติ, กากวติกา วา โหนฺติ, วาสุเทววติกา วา โหนฺติ, พลเทววติกา วา โหนฺติ, ปุณฺณภทฺทวติกา วา โหนฺติ, มณิภทฺทวติกา วา โหนฺติ, อคฺคิวติกา วา โหนฺติ, นาควติกา วา โหนฺติ, สุปณฺณวติกา วา โหนฺติ, ยกฺขวติกา วา โหนฺติ, อสุรวติกา วา โหนฺติ, คนฺธพฺพวติกา วา โหนฺติ, มหาราชวติกา วา โหนฺติ, จนฺทวติกา วา โหนฺติ, สูริยวติกา วา โหนฺติ, อินฺทวติกา วา โหนฺติ, พฺรหฺมวติกา วา โหนฺติ, เทววติกา วา โหนฺติ, ทิสาวติกา วา โหนฺติ.\csromanspacer{} อิเม เต สมณพฺราหฺมณา วตสุทฺธิกา, เต วเตน สุทฺธิํ วิสุทฺธิํ ปริสุทฺธิํ, มุตฺติํ วิมุตฺติํ ปริมุตฺติํ ปจฺเจนฺติ.}

\prose{สนฺเตเก สมณพฺราหฺมณา มุตสุทฺธิกา, เต กาลโต อุฏฺฐหิตฺวา ปถวิํ อามสนฺติ, หริตํ อามสนฺติ, โคมยํ อามสนฺติ, กจฺฉปํ อามสนฺติ, ผาลํ อกฺกมนฺติ, ติลวาหํ อามสนฺติ, ผุสฺสติลํ ขาทนฺติ}

\vspace{\tipitakaparskip}
\csromancenter{สุทฺธฏฺฐกสุตฺตนิทฺเทส ผุสฺสเตลํ มกฺเขนฺติ, ผุสฺสทนฺตกฏฺฐํ ขาทนฺติ, ผุสฺสมตฺติกาย นฺหายนฺติ, ผุสฺสสาฏกํ นิวาเสนฺติ, ผุสฺสเวฐนํ เวเฐนฺติ.\csromanspacer{} อิเม เต สมณพฺราหฺมณา มุตสุทฺธิกา, เต มุเตน สุทฺธิํ วิสุทฺธิํ ปริสุทฺธิํ, มุตฺติํ วิมุตฺติํ ปริมุตฺติํ ปจฺเจนฺติ น พฺราหฺมโณ อญฺญโต สุทฺธิมาห.}
\prose{\csromanfolio{69}\textbf{ทิฏฺเฐ สุเต สีลวเต มุเต วา}ติ พฺราหฺมโณ ทิฏฺฐสุทฺธิยาปิ สุทฺธิํ นาห, สุตสุทฺธิยาปิ สุทฺธิํ นาห, สีลสุทฺธิยาปิ สุทฺธิํ นาห, วตสุทฺธิยาปิ สุทฺธิํ นาห, มุตสุทฺธิยาปิ สุทฺธิํ นาห น กเถติ น ภณติ น ทีปยติ น โวหรตีติ น พฺราหฺมโณ อญฺญโต สุทฺธิมาห, ทิฏฺเฐ สุเต สีลวเต มุเต วา.}

\prose{\textbf{ปุญฺเญ ว ปาเป จ อนูปลิตฺโต}ติ \textbf{ปุญฺญํ} วุจฺจติ ยํ กิญฺจิ เตธาตุกํ กุสลาภิสงฺขารํ, \textbf{อปุญฺญํ} วุจฺจติ สพฺพํ อกุสลํ.\csromanspacer{} ยโต ปุญฺญาภิสงฺขาโร จ อปุญฺญาภิสงฺขาโร จ อาเนญฺชาภิสงฺขาโร จ ปหีนา โหนฺติ อุจฺฉินฺนมูลา ตาลาวตฺถุกตา อนภาวํกตา อายติํ อนุปฺปาทธมฺมา.\csromanspacer{} เอตฺตาวตา ปุญฺเญ จ ปาเป จ น ลิมฺปติ น ปลิมฺปติ น อุปลิมฺปติ, อลิตฺโต อปลิตฺโต อนูปลิตฺโต นิกฺขนฺโต นิสฺสโฏ วิปฺปมุตฺโต วิสญฺญุตฺโต วิมริยาทิกเตน เจตสา วิหรตีติ ปุญฺเญ จ ปาเป จ อนูปลิตฺโต.}

\prose{\textbf{อตฺตญฺชโห นยิธ ปกุพฺพมาโน}ติ \textbf{อตฺตญฺชโห}ติ อตฺตทิฏฺฐิชโห.\csromanspacer{} \textbf{อตฺตญฺชโห}ติ คาหํ ชโห\footnote{คาหชโห (สี, สฺยา), อตฺตคาหํ ชโห (ก)}.\csromanspacer{} \textbf{อตฺตญฺชโห}ติ ตณฺหาวเสน ทิฏฺฐิวเสน คหิตํ ปรามฏฺฐํ อภินิวิฏฺฐํ อชฺโฌสิตํ อธิมุตฺตํ.\csromanspacer{} สพฺพํ ตํ จตฺตํ โหติ วนฺตํ มุตฺตํ ปหีนํ ปฏินิสฺสฏฺฐํ.\csromanspacer{}\csromanspacer{}\textbf{นยิธ ปกุพฺพมาโน}ติ ปุญฺญาภิสงฺขารํ วา อปุญฺญาภิสงฺขารํ วา อาเนญฺชาภิสงฺขารํ วา อปกุพฺพมาโน อชนยมาโน อสญฺชนยมาโน อนิพฺพตฺตยมาโน อนภินิพฺพตฺตยมาโนติ \textbf{อตฺตญฺชโห} นยิธ ปกุพฺพมาโน.\csromanspacer{} เตนาห ภควา\csromandash{}}

\csromangathameasure{น พฺราหฺมโณ อญฺญโต สุทฺธิมาห,}
\csromangathagroup{\csromangathastanza{น พฺราหฺมโณ อญฺญโต สุทฺธิมาห,}}

\prose{\textbf{ทิฏฺเฐ สุเต สีลวเต มุเต วา}.}

\csromangathameasure{ปุญฺเญ จ ปาเป จ อนูปลิตฺโต,}
\csromangathagroup{\csromangathastanza{ปุญฺเญ จ ปาเป จ อนูปลิตฺโต,}}

\prose{\textbf{อตฺตญฺชโห นยิธ ปกุพฺพมาโน}ติ.}

\proseitem{26}{\csromanfolio{70}\csromansymbolfootnote{*}{ขุ 1. 402 ปิฏฺเฐปิ.}\textbf{ปุริมํ ปหาย อปรํ สิตาเส},}

\prose{\textbf{เอชานุคา เต น ตรนฺติ สงฺคํ.}}

\csromangathameasure{\textbf{เต อุคฺคหายนฺติ นิรสฺสชนฺติ,} \\ \textbf{กปีว สาขํ ปมุญฺจํ} \textbf{คหาย.}}
\csromangathagroup{\csromangathastanza{\textbf{เต อุคฺคหายนฺติ นิรสฺสชนฺติ,} \\ \textbf{กปีว สาขํ ปมุญฺจํ}\footnote{ปมุขํ (สี, สฺยา)} \textbf{คหาย.}}}

\csromanmatikaanchor{mtk.5}
\csromantocmarkref{subsection}{4. สุทฺธฏฺฐกสุตฺตนิทฺเทส}{mtk.5}
\bookmark[dest={mtk.5},level=2]{4. สุทฺธฏฺฐกสุตฺตนิทฺเทส}
\prose{\textbf{ปุริมํ ปหาย อปรํ สิตาเส}ติ ปุริมํ สตฺถารํ ปหาย ปรํ สตฺถารํ นิสฺสิตา, ปุริมํ ธมฺมกฺขานํ ปหาย อปรํ ธมฺมกฺขานํ นิสฺสิตา, ปุริมํ คณํ ปหาย อปรํ คณํ นิสฺสิตา, ปุริมํ ทิฏฺฐิํ ปหาย อปรํ ทิฏฺฐิํ นิสฺสิตา, ปุริมํ ปฏิปทํ ปหาย อปรํ ปฏิปทํ นิสฺสิตา, ปุริมํ มคฺคํ ปหาย อปรํ มคฺคํ นิสฺสิตา สนฺนิสฺสิตา อลฺลีนา อุปคตา อชฺโฌสิตา อธิมุตฺตาติ ปุริมํ ปหาย อปรํ สิตาเส.}

\prose{\textbf{เอชานุคา เต น ตรนฺติ สงฺคนฺ}ติ \textbf{เอชา} วุจฺจติ ตณฺหา.\csromanspacer{} โย ราโค สาราโค ฯเปฯ อภิชฺฌา โลโภ อกุสลมูลํ.\csromanspacer{} \textbf{เอชานุคา}ติ \textbf{เอชา}นุคา \textbf{เอชา}นุคตา \textbf{เอชา}นุสฏา \textbf{เอชา}ย ปนฺนา ปติตา อภิภูตา ปริยาทินฺนจิตฺตา.\csromanspacer{}\csromanspacer{}\textbf{เต น ตรนฺติ สงฺคนฺ}ติ ราคสงฺคํ โทสสงฺคํ โมหสงฺคํ มานสงฺคํ ทิฏฺฐิสงฺคํ กิเลสสงฺคํ ทุจฺจริตสงฺคํ น ตรนฺติ น อุตฺตรนฺติ น ปตรนฺติ น สมติกฺกมนฺติ น วีติวตฺตนฺตีติ \textbf{เอชา}นุคา เต น ตรนฺติ สงฺคํ.}

\prose{\textbf{เต อุคฺคหายนฺติ นิรสฺสชนฺตี}ติ สตฺถารํ คณฺหนฺติ, ตํ มุญฺจิตฺวา อญฺญํ สตฺถารํ คณฺหนฺติ.\csromanspacer{} ธมฺมกฺขานํ คณฺหนฺติ, ตํ มุญฺจิตฺวา อญฺญํ ธมฺมกฺขานํ คณฺหนฺติ.\csromanspacer{} คณํ คณฺหนฺติ, ตํ มุญฺจิตฺวา อญฺญํ คณํ คณฺหนฺติ.\csromanspacer{} ทิฏฺฐิํ คณฺหนฺติ, ตํ มุญฺจิตฺวา อญฺญํ ทิฏฺฐิํ คณฺหนฺติ.\csromanspacer{} ปฏิปทํ คณฺหนฺติ, ตํ มุญฺจิตฺวา อญฺญํ ปฏิปทํ คณฺหนฺติ.\csromanspacer{} มคฺคํ คณฺหนฺติ, ตํ มุญฺจิตฺวา อญฺญํ มคฺคํ คณฺหนฺติ, คณฺหนฺติ จ มุญฺจนฺติ จ อาทิยนฺติ จ นิรสฺสชนฺติ จาติ เต อุคฺคหายนฺติ นิรสฺสชนฺติ.}

\prose{\textbf{กปีว สาขํ ปมุญฺจํ คหายา}ติ ยถา มกฺกโฏ อรญฺเญ ปวเน จรมาโน สาขํ คณฺหาติ, ตํ มุญฺจิตฺวา อญฺญํ สาขํ คณฺหาติ.\csromanspacer{} เอวเมวํ ปุถุสมณพฺราหฺมณา ปุถุทิฏฺฐิคตานิ คณฺหนฺติ จ มุญฺจนฺติ จ อาทิยนฺติ นิรสฺสชนฺติ จาติ กปีว สาขํ ปมุญฺจํ \textbf{คหาย.}\csromanspacer{} เตนาห ภควา\csromandash{}}

\prose{\textbf{ปุริมํ ปหาย อปรํ สิตาเส},}

\prose{\textbf{เอชานุคา เต น ตรนฺติ สงฺคํ.}}

\csromangathameasure{\textbf{เต อุคฺคหายนฺติ นิรสฺสชนฺติ,} \\ \textbf{กปีว สาขํ ปมุญฺจํ คหายา}ติ.}
\csromangathagroup{\csromangathastanza{\textbf{เต อุคฺคหายนฺติ นิรสฺสชนฺติ,} \\ \textbf{กปีว สาขํ ปมุญฺจํ คหายา}ติ.}}

\csromanheader{2}{4. สุทฺธฏฺฐกสุตฺตนิทฺเทส}
\proseitem{27}{\csromanfolio{71}\csromansymbolfootnote{*}{ขุ 1. 403 ปิฏฺเฐ.}สยํ สมาทาย วตานิ ชนฺตุ,}

\prose{\textbf{อุจฺจาวจํ คจฺฉติ สญฺญสตฺโต.}}

\csromangathameasure{\textbf{วิทฺวา จ เวเทหิ สเมจฺจ ธมฺมํ,} \\ \textbf{น อุจฺจาวจํ คจฺฉติ ภูริปญฺโญ.}}
\csromangathagroup{\csromangathastanza{\textbf{วิทฺวา จ เวเทหิ สเมจฺจ ธมฺมํ,} \\ \textbf{น อุจฺจาวจํ คจฺฉติ ภูริปญฺโญ.}}}

\prose{\textbf{สยํ สมาทาย วตานิ ชนฺตู}ติ \textbf{สยํ สมาทายา}ติ สามํ สมาทาย.\csromanspacer{} \textbf{วตานี}ติ หตฺถิวตํ วา อสฺสวตํ วา โควตํ วา กุกฺกูรวตํ วา กากวตํ วา วาสุเทววตํ วา พลเทววตํ วา ปุณฺณภทฺทวตํ วา มณิภทฺทวตํ วา อคฺคิวตํ วา นาควตํ วา สุปณฺณวตํ วา ยกฺขวตํ วา อสุรวตํ วา ฯเปฯ ทิสาวตํ วา อาทาย สมาทาย อาทิยิตฺวา สมาทิยิตฺวา คณฺหิตฺวา ปรมสิตฺวา อภินิวิสิตฺวา.\csromanspacer{}\csromanspacer{}\textbf{ชนฺตู}ติ สตฺโต นโร ฯเปฯ มนุโชติ สยํ สมาทาย วตานิ ชนฺตุ.}

\prose{\textbf{อุจฺจาวจํ คจฺฉติ สญฺญสตฺโต}ติ สตฺถารโต สตฺถารํ คจฺฉติ, ธมฺมกฺขานโต ธมฺมกฺขานํ คจฺฉติ, คณโต คณํ คจฺฉติ, ทิฏฺฐิยา ทิฏฺฐิํ คจฺฉติ, ปฏิปทาย ปฏิปทํ คจฺฉติ, มคฺคโต มคฺคํ คจฺฉติ.\csromanspacer{}\csromanspacer{}\textbf{สญฺญสตฺโต}ติ กามสญฺญาย พฺยาปาทสญฺญาย วิหิํสาสญฺญาย ทิฏฺฐิสญฺญาย สตฺโต วิสตฺโต อาสตฺโต ลคฺโค ลคฺคิโต ปลิพุทฺโธ.\csromanspacer{} ยถา ภิตฺติขิเล วา นาคทนฺเต วา ภณฺฑํ สตฺตํ วิสตฺตํ อาสตฺตํ ลคฺคํ ลคฺคิตํ ปลิพุทฺธํ, เอวเมวํ กามสญฺญาย พฺยาปาทสญฺญาย วิหิํสาสญฺญาย ทิฏฺฐิสญฺญาย สตฺโต วิสตฺโต อาสตฺโต ลคฺโค ลคฺคิโต ปลิพุทฺโธติ อุจฺจาวจํ คจฺฉติ สญฺญสตฺโต.}

\prose{\textbf{วิทฺวา จ เวเทหิ สเมจฺจ ธมฺมนฺ}ติ \textbf{วิทฺวา}ติ \textbf{วิทฺวา} วิชฺชาคโต ญาณี วิภาวี เมธาวี.\csromanspacer{} \textbf{เวเทหี}ติ \textbf{เวทา} วุจฺจนฺติ จตูสุ มคฺเคสุ ญาณํ, ปญฺญา ปญฺญินฺทฺริยํ ปญฺญาพลํ ธมฺมวิจยสมฺโพชฺฌงฺโค วีมํสา วิปสฺสนา สมฺมาทิฏฺฐิ.\csromanspacer{} เตหิ เวเทหิ ชาติชรามรณสฺส อนฺตคโต อนฺตปฺปตฺโต, โกฏิคโต โกฏิปฺปตฺโต, ปริยนฺตคโต ปริยนฺตปฺปตฺโต, โวสานคโต โวสานปฺปตฺโต, ตาณคโต ตาณปฺปตฺโต, เลณคโต เลณปฺปตฺโต, สรณคโต สรณปฺปตฺโต, อภยคโต อภยปฺปตฺโต, อจฺจุตคโต อจฺจุตปฺปตฺโต, อมตคโต อมตปฺปตฺโต, นิพฺพานคโต นิพฺพานปฺปตฺโต, \textbf{เวทา}นํ วา อนฺตคโตติ \textbf{เวทคู,} เวเทหิ วา อนฺตคโตติ \textbf{เวทคู,} สตฺตนฺนํ วา ธมฺมานํ วิทิตตฺตา \textbf{เวทคู,} สกฺกายทิฏฺฐิ วิทิตา โหติ, วิจิกิจฺฉา วิทิตา โหติ, สีลพฺพตปรามาโส \csromanfolio{72}วิทิโต โหติ, ราโค วิทิโต โหติ, โทโส วิทิโต โหติ, โมโห วิทิโต โหติ, มาโน วิทิโต โหติ, วิทิตาสฺส โหนฺติ ปาปกา อกุสลา ธมฺมา สํกิเลสิกา โปโนภวิกา สทรา ทุกฺขวิปากา อายติํ ชาติชรามรณิยา.}

\vspace{\tipitakaparskip}
\csromancenter{เวทานิ วิเจยฺย เกวลานิ, (สภิยาติ ภควา,)}
\prose{สมณานํ ยานีธตฺถิ\footnote{ยานิปตฺถิ (สี, สฺยา) ขุ 1. 360; ขุ 8. 64 ปิฏฺเฐสุปิ.} พฺราหฺมณานํ.}

\csromangathameasure{สพฺพเวทนาสุ วีตราโค, \\ สพฺพํ เวทมติจฺจ เวทคู โสติ.}
\csromangathagroup{\csromangathastanza{สพฺพเวทนาสุ วีตราโค, \\ สพฺพํ เวทมติจฺจ เวทคู โสติ.}}

\prose{\textbf{วิทฺวา จ เวเทหิ สเมจฺจ ธมฺมนฺ}ติ สเมจฺจ อภิสเมจฺจ ธมฺมํ, “สพฺเพ สงฺขารา อนิจฺจา”ติ สเมจฺจ อภิสเมจฺจ ธมฺมํ, “สพฺเพ สงฺขารา ทุกฺขา”ติ สเมจฺจ อภิสเมจฺจ ธมฺมํ, “สพฺเพ ธมฺมา อนตฺตา”ติ สเมจฺจ อภิสเมจฺจ ธมฺมํ. “อวิชฺชาปจฺจยา สงฺขารา”ติ สเมจฺจ อภิสเมจฺจ ธมฺมํ, “สงฺขารปจฺจยา วิญฺญาณนฺ”ติ สเมจฺจ อภิสเมจฺจ ธมฺมํ, “วิญฺญาณปจฺจยา นามรูปนฺ”ติ. “นามรูปปจฺจยา สฬายตนนฺ”ติ. “สฬายตนปจฺจยา ผสฺโส”ติ. “ผสฺสปจฺจยา เวทนา”ติ. “เวทนาปจฺจยา ตณฺหา”ติ. “ตณฺหาปจฺจยา อุปาทานนฺ”ติ. “อุปาทานปจฺจยา ภโว”ติ. “ภวปจฺจยา ชาตี”ติ. “ชาติปจฺจยา ชรามรณนฺ”ติ สเมจฺจ อภิสเมจฺจ ธมฺมํ. “อวิชฺชานิโรธา สงฺขารนิโรโธ”ติ สเมจฺจ อภิสเมจฺจ ธมฺมํ, “สงฺขารนิโรธา วิญฺญาณนิโรโธ”ติ สเมจฺจ อภิสเมจฺจ ธมฺมํ, “วิญฺญาณนิโรธา นามรูปนิโรโธ”ติ. “นามรูปนิโรธา สฬายตนนิโรโธ”ติ. “สฬายตนนิโรธา ผสฺสนิโรโธ”ติ. “ผสฺสนิโรธา เวทนานิโรโธ”ติ. “เวทนานิโรธา ตณฺหานิโรโธ”ติ. “ตณฺหานิโรธา อุปาทานนิโรโธ”ติ.\csromanspacer{} อุปาทานนิโรธา ภวนิโรโธ”ติ. “ภวนิโรธา ชาตินิโรโธ”ติ. “ชาตินิโรธา ชรามรณนิโรโธ”ติ สเมจฺจ อภิสเมจฺจ ธมฺมํ. “อิทํ ทุกฺขนฺ”ติ สเมจฺจ อภิสเมจฺจ ธมฺมํ, “อยํ ทุกฺขสมุทโย”ติ. “อยํ ทุกฺขนิโรโธ”ติ. “อยํ ทุกฺขนิโรธคามินี ปฏิปทา”ติ สเมจฺจ อภิสเมจฺจ ธมฺมํ. “อิเม อาสวา”ติ สเมจฺจ อภิสเมจฺจ ธมฺมํ, “อยํ อาสวสมุทโย”ติ. “อยํ อาสวนิโรโธ”ติ. “อยํ อาสวนิโรธคามินี ปฏิปทา”ติ สเมจฺจ อภิสเมจฺจ ธมฺมํ. “อิเม ธมฺมา อภิญฺเญยฺยา”ติ สเมจฺจ อภิสเมจฺจ ธมฺมํ, “อิเม ธมฺมา ปริญฺเญยฺยา”ติ. “อิเม ธมฺมา ปหาตพฺพา”ติ. “อิเม ธมฺมา}

\vspace{\tipitakaparskip}
\csromancenter{สุทฺธฏฺฐกสุตฺตนิทฺเทส ภาเวตพฺพา”ติ. “อิเม ธมฺมา สจฺฉิกาตพฺพา”ติ สเมจฺจ อภิสเมจฺจ ธมฺมํ.\csromanspacer{} ฉนฺนํ ผสฺสายตนานํ สมุทยญฺจ อตฺถงฺคมญฺจ อสฺสาทญฺจ อาทีนวญฺจ นิสฺสรณญฺจ สเมจฺจ อภิสเมจฺจ ธมฺมํ, ปญฺจนฺนํ อุปาทานกฺขนฺธานํ สมุทยญฺจ อตฺถงฺคมญฺจ อสฺสาทญฺจ อาทีนวญฺจ นิสฺสรณญฺจ สเมจฺจ อภิสเมจฺจ ธมฺมํ, จตุนฺนํ มหาภูตานํ สมุทยญฺจ อตฺถงฺคมญฺจ อสฺสาทญฺจ อาทีนวญฺจ นิสฺสรณญฺจ สเมจฺจ อภิสเมจฺจ ธมฺมํ, “ยํ กิญฺจิ สมุทยธมฺมํ สพฺพํ ตํ นิโรธธมฺมนฺ”ติ สเมจฺจ อภิสเมจฺจ ธมฺมนฺติ วิทฺวา จ เวเทหิ สเมจฺจ ธมฺมํ.}
\prose{\csromanfolio{73}\textbf{น อุจฺจาวจํ คจฺฉติ ภูริปญฺโญ}ติ น สตฺถารโต สตฺถารํ คจฺฉติ, น ธมฺมกฺขานโต ธมฺมกฺขานํ คจฺฉติ, น คณโต คณํ คจฺฉติ, น ทิฏฺฐิยา ทิฏฺฐิํ คจฺฉติ, น ปฏิปทาย ปฏิปทํ คจฺฉติ, น มคฺคโต มคฺคํ คจฺฉติ.\csromanspacer{}\csromanspacer{}\textbf{ภูริปญฺโญ}ติ ภูริปญฺโญ มหาปญฺโญ ปุถุปญฺโญ หาสปญฺโญ ชวนปญฺโญ ติกฺขปญฺโญ นิพฺเพธิกปญฺโญ.\csromanspacer{} \textbf{ภูริ} วุจฺจติ ปถวี, ตาย ปถวิสมาย ปญฺญาย วิปุลาย วิตฺถตาย สมนฺนาคโตติ น อุจฺจาวจํ คจฺฉติ ภูริปญฺโญ.\csromanspacer{} เตนาห ภควา\csromandash{}}

\csromangathameasure{สยํ สมาทาย วตานิ ชนฺตุ, \\ อุจฺจาวจํ คจฺฉติ สญฺญสตฺโต. \\ วิทฺวา จ เวเทหิ สเมจฺจ ธมฺมํ,}
\csromangathagroup{\csromangathastanza{สยํ สมาทาย วตานิ ชนฺตุ, \\ อุจฺจาวจํ คจฺฉติ สญฺญสตฺโต. \\ วิทฺวา จ เวเทหิ สเมจฺจ ธมฺมํ,}}

\prose{\textbf{น อุจฺจาวจํ คจฺฉติ ภูริปญฺโญ}ติ.}

\proseitem{28}{\csromansymbolfootnote{*}{ขุ 1. 403 ปิฏฺเฐปิ.}ส สพฺพธมฺเมสุ วิเสนิภูโต,}

\prose{\textbf{ยํ กิญฺจิ ทิฏฺฐํ ว สุตํ มุตํ วา.}}

\csromangathameasure{\textbf{ตเมว ทสฺสิํ วิวฏํ จรนฺตํ,} \\ \textbf{เกนีธ โลกสฺมิ วิกปฺปเยยฺย.}}
\csromangathagroup{\csromangathastanza{\textbf{ตเมว ทสฺสิํ วิวฏํ จรนฺตํ,} \\ \textbf{เกนีธ โลกสฺมิ วิกปฺปเยยฺย.}}}

\prose{\textbf{ส สพฺพธมฺเมสุ วิเสนิภูโต, ยํ กิญฺจิ ทิฏฺฐํ ว สุตํ มุตํ วา}ติ \textbf{เสนา} วุจฺจติ มาร\textbf{เสนา}.\csromanspacer{} กายทุจฺจริตํ มาร\textbf{เสนา}, วจีทุจฺจริตํ มาร\textbf{เสนา}, มโนทุจฺจริตํ มาร\textbf{เสนา}, ราโค มาร\textbf{เสนา}, โทโส มาร\textbf{เสนา}, โมโห มาร\textbf{เสนา}, โกโธ.\csromanspacer{} อุปนาโห ฯเปฯ สพฺพากุสลาภิสงฺขารา มาร\textbf{เสนา}.\csromanspacer{} วุตฺตํ เหตํ ภควตา\csromandash{}}

\csromangathameasure{+กามา เต ปฐมา เสนา, ทุติยา อรติ วุจฺจติ. \\ ตติยา ขุปฺปิปาสา เต, จตุตฺถี ตณฺหา วุจฺจติ. \\ ปญฺจมี ถินมิทฺธํ เต, ฉฏฺฐา ภีรู ปวุจฺจติ. \\ สตฺตมี วิจิกิจฺฉา เต, มกฺโข ถมฺโภ เต อฏฺฐโม. \\ ลาโภ สิโลโก สกฺกาโร, \\ มิจฺฉาลทฺโธ จ โย ยโส. \\ โย จตฺตานํ สมุกฺกํเส, ปเร จ อวชานติ. \\ เอสา นมุจิ เต เสนา, กณฺหาสฺสาภิปฺปหารินี. \\ น นํ อสูโร ชินาติ, เชตฺวาว ลภเต สุขนฺติ.}
\csromangathagroup{\csromangathastanza{\csromansymbolfootnote{+}{ขุ 1. 342; ขุ 8. 114 ปิฏฺเฐสุปิ.}กามา เต ปฐมา เสนา, ทุติยา อรติ วุจฺจติ. \\ ตติยา ขุปฺปิปาสา เต, จตุตฺถี ตณฺหา วุจฺจติ.}\csromangathastanza{\csromanfolio{74}ปญฺจมี ถินมิทฺธํ เต, ฉฏฺฐา ภีรู ปวุจฺจติ. \\ สตฺตมี วิจิกิจฺฉา เต, มกฺโข ถมฺโภ เต อฏฺฐโม.}\csromangathastanza{ลาโภ สิโลโก สกฺกาโร, \\ มิจฺฉาลทฺโธ จ โย ยโส.}\csromangathastanza{โย จตฺตานํ สมุกฺกํเส, ปเร จ อวชานติ. \\ เอสา นมุจิ เต เสนา, กณฺหาสฺสาภิปฺปหารินี. \\ น นํ อสูโร ชินาติ, เชตฺวาว ลภเต สุขนฺติ.}}

\prose{ยโต จตูหิ อริยมคฺเคหิ สพฺพา จ มารเสนา สพฺเพ จ ปฏิเสนิกรา กิเลสา ชิตา จ ปราชิตา จ ภคฺคา วิปฺปลุคฺคา ปรมฺมุขา, โส วุจฺจติ วิเสนิภูโต.\csromanspacer{} โส ทิฏฺเฐ วิเสนิภูโต, สุเต วิเสนิภูโต, มุเต วิเสนิภูโต, วิญฺญาเต วิเสนิภูโตติ ส สพฺพธมฺเมสุ เวเสนิภูโต, ยํ กิญฺจิ ทิฏฺฐํ ว สุตํ มุตํ วา.}

\prose{\textbf{ตเมว ทสฺสิํ วิวฏํ จรนฺตนฺ}ติ ตเมว สุทฺธทสฺสิํ วิสุทฺธทสฺสิํ ปริสุทฺธทสฺสิํ โวทาตทสฺสิํ ปริโยทาตทสฺสิํ.\csromanspacer{} อถ วา สุทฺธทสฺสนํ วิสุทฺธทสฺสนํ ปริสุทฺธทสฺสนํ โวทาตทสฺสนํ ปริโยทาตทสฺสนํ.\csromanspacer{}\csromanspacer{}\textbf{วิวฏนฺ}ติ ตณฺหาฉทนํ ทิฏฺฐิฉทนํ กิเลสฉทนํ ทุจฺจริตฉทนํ อวิชฺชาฉทนํ.\csromanspacer{} ตานิ ฉทนานิ วิวฏานิ โหนฺติ วิทฺธํสิตานิ อุคฺฆาฏิตานิ สมุคฺฆาฏิตานิ ปหีนานิ สมุจฺฉินฺนานิ วูปสนฺตานิ ปฏิปสฺสทฺธานิ อภพฺพุปฺปตฺติกานิ ญาณคฺคินา ทฑฺฒานิ.\csromanspacer{}\csromanspacer{}\textbf{จรนฺตนฺ}ติ จรนฺตํ วิจรนฺตํ วิหรนฺตํ อิริยนฺตํ วตฺเตนฺตํ ปาเลนฺตํ ยเปนฺตํ ยาเปนฺตนฺติ ตเมว ทสฺสิํ วิวฏํ จรนฺตํ.}

\prose{\textbf{เกนีธ โลกสฺมิ วิกปฺปเยยฺยา}ติ \textbf{กปฺปา}ติ เทฺว \textbf{กปฺปา} ตณฺหากปฺโป จ ทิฏฺฐิกปฺโป จ ฯเปฯ อยํ ตณฺหากปฺโป ฯเปฯ อยํ ทิฏฺฐิกปฺโป.\csromanspacer{} ตสฺส ตณฺหากปฺโป ปหีโน, ทิฏฺฐิกปฺโป ปฏินิสฺสฏฺโฐ, ตณฺหากปฺปสฺส ปหีนตฺตา, ทิฏฺฐิกปฺปสฺส ปฏินิสฺสฏฺฐตฺตา เกน ราเคน กปฺเปยฺย, เกน โทเสน กปฺเปยฺย, เกน โมเหน กปฺเปยฺย, เกน มาเนน กปฺเปยฺย, กาย ทิฏฺฐิยา กปฺเปยฺย, เกน อุทฺธจฺเจน กปฺเปยฺย, กาย วิจิกิจฺฉาย กปฺเปยฺย, เกหิ อนุสเยหิ กปฺเปยฺย “รตฺโต”ติ วา “ทุฏฺโฐ”ติ วา “มูโฬฺห”ติ วา “วินิพทฺโธ”ติ วา “ปรามฏฺโฐ”ติ วา “วิกฺเขปคโต”ติ วา “อนิฏฺฐงฺคโต”ติ วา “ถามคโต”ติ วา.\csromanspacer{} เต อภิสงฺขารา ปหีนา, อภิสงฺขารานํ}

\vspace{\tipitakaparskip}
\csromancenter{สุทฺธฏฺฐกสุตฺตนิทฺเทส ปหีนตฺตา คติโย เกน กปฺเปยฺย “เนรยิโก”ติ วา “ติรจฺฉานโยนิโก”ติ วา “เปตฺติวิสยิโก”ติ วา “มนุสฺโส”ติ วา “เทโว”ติ วา “รูปี”ติ วา “อรูปี”ติ วา “สญฺญี”ติ วา “อสญฺญี”ติ วา “เนวสญฺญีนาสญฺญี”ติ วา.\csromanspacer{} โส เหตุ นตฺถิ, ปจฺจโย นตฺถิ, การณํ นตฺถิ, เยน กปฺเปยฺย วิกปฺเปยฺย วิกปฺปํ อาปชฺเชยฺย.\csromanspacer{} \textbf{โลกสฺมินฺ}ติ อปายโลเก มนุสฺสโลเก เทวโลเก ขนฺธโลเก ธาตุโลเก อายตนโลเกติ เกนีธ โลกสฺมิํ วิกปฺปเยยฺย.\csromanspacer{} เตนาห ภควา\csromandash{}}
\vspace{0.5\baselineskip}
\csromangathameasure{ส สพฺพธมฺเมสุ วิเสนิภูโต, \\ ยํ กิญฺจิ ทิฏฺฐํ ว สุตํ มุตํ วา. \\ ตเมว ทสฺสิํ วิวฏํ จรนฺตํ, \\ เกนีธ โลกสฺมิ วิกปฺปเยยฺยาติ.}
\csromangathagroup{\csromangathastanza{\csromanfolio{75}ส สพฺพธมฺเมสุ วิเสนิภูโต, \\ ยํ กิญฺจิ ทิฏฺฐํ ว สุตํ มุตํ วา. \\ ตเมว ทสฺสิํ วิวฏํ จรนฺตํ, \\ เกนีธ โลกสฺมิ วิกปฺปเยยฺยาติ.}}

\proseitem{29}{\csromansymbolfootnote{*}{ขุ 1. 403 ปิฏฺเฐปิ.}น กปฺปยนฺติ น ปุเรกฺขโรนฺติ,}

\prose{\textbf{“อจฺจนฺตสุธี”ติ น เต วทนฺติ.}}

\csromangathameasure{\textbf{อาทานคนฺถํ คถิตํ วิสชฺช,} \\ \textbf{อาสํ น กุพฺพนฺติ กุหิญฺจิ โลเก.}}
\csromangathagroup{\csromangathastanza{\textbf{อาทานคนฺถํ คถิตํ วิสชฺช,} \\ \textbf{อาสํ น กุพฺพนฺติ กุหิญฺจิ โลเก.}}}

\prose{\textbf{น กปฺปยนฺติ น ปุเรกฺขโรนฺตี}ติ \textbf{กปฺปา}ติ เทฺว \textbf{กปฺปา} ตณฺหากปฺโป จ ทิฏฺฐิกปฺโป จ ฯเปฯ อยํ ตณฺหากปฺโป ฯเปฯ อยํ ทิฏฺฐิกปฺโป.\csromanspacer{} เตสํ ตณฺหากปฺโป ปหีโน, ทิฏฺฐิกปฺโป ปฏินิสฺสฏฺโฐ, ตณฺหากปฺปสฺส ปหีนตฺตา, ทิฏฺฐิกปฺปสฺส ปฏินิสฺสฏฺฐตฺตา ตณฺหากปฺปํ วา ทิฏฺฐิกปฺปํ วา น กปฺเปนฺติ น ชเนนฺติ น สญฺชเนนฺติ น นิพฺพตฺเตนฺติ นาภินิพฺพตฺเตนฺตีติ น กปฺปยนฺติ.\csromanspacer{}\csromanspacer{}\textbf{น ปุเรกฺขโรนฺตี}ติ \textbf{ปุเรกฺขารา}ติ เทฺว \textbf{ปุเรกฺขารา} ตณฺหาปุเรกฺขาโร จ ทิฏฺฐิปุเรกฺขาโร จ ฯเปฯ อยํ ตณฺหาปุเรกฺขาโร ฯเปฯ อยํ ทิฏฺฐิปุเรกฺขาโร.\csromanspacer{} เตสํ ตณฺหาปุเรกฺขาโร ปหีโน, ทิฏฺฐิปุเรกฺขาโร ปฏินิสฺสฏฺโฐ, ตณฺหาปุเรกฺขารสฺส ปหีนตฺตา, ทิฏฺฐิปุเรกฺขารสฺส ปฏินิสฺสฏฺฐตฺตา น ตณฺหํ วา น ทิฏฺฐิํ วา ปุรโต กตฺวา จรนฺติ, น ตณฺหาธชา น ตณฺหาเกตู น ตณฺหาธิปเตยฺยา, น ทิฏฺฐิธชา น ทิฏฺฐิเกตู น ทิฏฺฐาธิปเตยฺยา, น ตณฺหาย วา น ทิฏฺฐิยา วา ปริวาริตา จรนฺตีติ น กปฺปยนฺติ น ปุเรกฺขโรนฺติ.}

\prose{\textbf{“อจฺจนฺตสุทฺธี”ติ น เต วทนฺตี}ติ อจฺจนฺตสุทฺธิํ สํสารสุทฺธิํ อกิริยทิฏฺฐิํ สสฺสตวาทํ น วทนฺติ น กเถนฺติ น ภณนฺติ น ทีปยนฺติ น โวหรนฺตีติ “อจฺจนฺตสุทฺธี”ติ น เต วทนฺติ.}

\prose{\csromanfolio{76}\textbf{อาทานคนฺถํ คถิตํ วิสชฺชา}ติ \textbf{คนฺถา}ติ * จตฺตาโร \textbf{คนฺถา} อภิชฺฌากายคนฺโถ พฺยาปาโท กายคนฺโถ สีลพฺพตปรามาโส กายคนฺโถ อิทํสจฺจาภินิเวโส กายคนฺโถ.\csromanspacer{} อตฺตโน ทิฏฺฐิยา ราโค อภิชฺฌากายคนฺโถ, ปรวาเทสุ อาฆาโต อปฺปจฺจโย พฺยาปาโท กายคนฺโถ.\csromanspacer{} อตฺตโน สีลํ วา วตํ วา สีลวตํ วา ปรามสนฺตีติ สีลพฺพตปรามาโส กายคนฺโถ.\csromanspacer{} อตฺตโน ทิฏฺฐิ อิทํสจฺจาภินิเวโส กายคนฺโถ.\csromanspacer{} กิํการณา วุจฺจติ อาทานคนฺโถ, เตหิ คนฺเถหิ รูปํ อาทิยนฺติ อุปาทิยนฺติ คณฺหนฺติ ปรามสนฺติ อภินิวิสนฺติ.\csromanspacer{} เวทนํ.\csromanspacer{} สญฺญํ.\csromanspacer{} สงฺขาเร.\csromanspacer{} วิญฺญาณํ.\csromanspacer{} คติํ.\csromanspacer{} อุปปตฺติํ.\csromanspacer{} ปฏิสนฺธิํ.\csromanspacer{} ภวํ.\csromanspacer{} สํสารวฏฺฏํ อาทิยนฺติ อุปาทิยนฺติ คณฺหนฺติ ปรามสนฺติ อภินิวิสนฺติ, ตํการณา วุจฺจติ อาทานคนฺโถ.\csromanspacer{}\csromanspacer{}\textbf{วิสชฺชา}ติ คนฺเถ โวสชฺชิตฺวา วา วิสชฺช.\csromanspacer{} อถ วา คนฺเถ คถิเต คนฺถิเต พนฺเธ วิพนฺเธ อาพนฺเธ ลคฺเค ลคฺคิเต ปลิพุทฺเธ พนฺธเน โผฏยิตฺวา\footnote{โปฏยิตฺวา (สี, ก)} วิสชฺช.\csromanspacer{} ยถา วยฺหํ วา รถํ วา สกฏํ วา สนฺทมานิกํ วา สชฺชํ กโรนฺติ วิโกเปนฺติ.\csromanspacer{} เอวเมวํ คนฺเถ โวสชฺชิตฺวา วิสชฺช.\csromanspacer{} อถ วา คนฺเถ คถิเต คนฺถิเต พนฺเธ วิพนฺเธ อาพนฺเธ ลคฺเค ลคฺคิเต ปลิพุทฺเธ พนฺธเน โปฏยิตฺวา วิสชฺชาติ อาทานคนฺถํ คถิตํ วิสชฺช.}

\prose{\textbf{อาสํ น กุพฺพนฺติ กุหิญฺจิ โลเก}ติ \textbf{อาสา} วุจฺจติ ตณฺหา.\csromanspacer{} โย ราโค สาราโค ฯเปฯ อภิชฺฌา โลโภ อกุสลมูลํ.\csromanspacer{}\csromanspacer{}\textbf{อาสํ น กุพฺพนฺตี}ติ อาสํ น กุพฺพนฺติ น ชเนนฺติ น สญฺชเนนฺติ น นิพฺพตฺเตนฺติ น อภินิพฺพตฺเตนฺติ.\csromanspacer{}\csromanspacer{}\textbf{กุหิญฺจี}ติ กุหิญฺจิ กิมฺหิจิ กตฺถจิ อชฺฌตฺตํ วา พหิทฺธา วา อชฺฌตฺตพหิทฺธา วา.\csromanspacer{}\csromanspacer{}\textbf{โลเก}ติ อปายโลเก ฯเปฯ อายตนโลเกติ อาสํ น กุพฺพนฺติ กุหิญฺจิ โลเก.\csromanspacer{} เตนาห ภควา\csromandash{}}

\csromangathameasure{น กปฺปยนฺติ น ปุเรกฺขโรนฺติ, \\ “อจฺจนฺตสุทฺธี”ติ น เต วทนฺติ.}
\csromangathagroup{\csromangathastanza{น กปฺปยนฺติ น ปุเรกฺขโรนฺติ, \\ “อจฺจนฺตสุทฺธี”ติ น เต วทนฺติ.}}

\prose{อาทานคนฺถํ คถิตํ วิสชฺช,}

\prose{\textbf{อาสํ น กุพฺพนฺติ กุหิญฺจิ โลเก}ติ.}

\proseitem{30}{\csromansymbolfootnote{+}{ขุ 1. 403 ปิฏฺเฐปิ.}สีมาติโค พฺราหฺมโณ ตสฺส นตฺถิ,}

\prose{\textbf{ญตฺวา ว ทิสฺวา ว สมุคฺคหีตํ.}}

\csromangathameasure{\textbf{น ราคราคี นวิราครตฺโต,} \\ \textbf{ตสฺสีธ นตฺถิ ปรมุคฺคหีตํ.}}
\csromangathagroup{\csromangathastanza{\textbf{น ราคราคี นวิราครตฺโต,} \\ \textbf{ตสฺสีธ นตฺถิ ปรมุคฺคหีตํ.}}}

\csromanheader{2}{4. สุทฺธฏฺฐกสุตฺตนิทฺเทส}
\prose{\csromanfolio{77}\textbf{สีมาติโค พฺราหฺมโณ ตสฺส นตฺถิ, ญตฺวา ว ทิสฺวา ว สมุคฺคหีตนฺ}ติ \textbf{สีมา}ติ จตสฺโส \textbf{สีมา}โย สกฺกายทิฏฺฐิ วิจิกิจฺฉา สีลพฺพตปรามาโส ทิฏฺฐานุสโย วิจิกิจฺฉานุสโย ตเทกฏฺฐา จ กิเลสา, อยํ ปฐมา \textbf{สีมา}.\csromanspacer{} โอฬาริกํ กามราคสญฺโญชนํ ปฏิฆสญฺโญชนํ โอฬาริโก กามราคานุสโย ปฏิฆานุสโย ตเทกฏฺฐา จ กิเลสา, อยํ ทุติยา \textbf{สีมา}.\csromanspacer{} อณุสหคตํ กามราคสญฺโญชนํ ปฏิฆสญฺโญชนํ อณุสหคโต กามราคานุสโย ปฏิฆานุสโย ตเทกฏฺฐา จ กิเลสา, อยํ ตติยา \textbf{สีมา}.\csromanspacer{} รูปราโค อรูปราโค มาโน อุทฺธจฺจํ อวิชฺชา มานานุสโย ภวราคานุสโย อวิชฺชานุสโย ตเทกฏฺฐา จ กิเลสา, อยํ จตุตฺถา \textbf{สีมา}.\csromanspacer{} ยโต จ จตูหิ อริยมคฺเคหิ อิมา จตสฺโส \textbf{สีมา}โย อติกฺกนฺโต โหติ สมติกฺกนฺโต วีติวตฺโต, โส วุจฺจติ \textbf{สีมา}ติโค.\csromanspacer{}\csromanspacer{}\textbf{พฺราหฺมโณ}ติ สตฺตนฺนํ ธมฺมานํ พาหิตตฺตา พฺราหฺมโณ, สกฺกายทิฏฺฐิ พาหิตา โหติ, วิจิกิจฺฉา พาหิตา โหติ, สีลพฺพตปรามาโส พาหิโต โหติ ฯเปฯ อสิโต ตาทิ ปวุจฺจเต สพฺรหฺมา.\csromanspacer{}\csromanspacer{}\textbf{ตสฺสา}ติ อรหโต ขีณาสวสฺส.}

\prose{\textbf{ญตฺวา}ติ ปรจิตฺตญาเณน วา ญตฺวา ปุพฺเพนิวาสานุสฺสติญาเณน วา ญตฺวา.\csromanspacer{} \textbf{ทิสฺวา}ติ มํสจกฺขุนา วา ทิสฺวา ทิพฺพจกฺขุนา วา ทิสฺวา.\csromanspacer{} \textbf{สีมาติโค พฺราหฺมโณ ตสฺส นตฺถิ, ญตฺวา ว ทิสฺวา ว สมุคฺคหีตนฺ}ติ ตสฺส อิทํ ปรมํ อคฺคํ เสฏฺฐํ วิสิฏฺฐํ\footnote{วิเสฏฺฐํ (สี, สฺยา)} ปาโมกฺขํ อุตฺตมํ ปวรนฺติ คหิตํ ปรามฏฺฐํ อภินิวิฏฺฐํ อชฺโฌสิตํ อธิมุตฺตํ นตฺถิ น สนฺติ น สํวิชฺชติ นุปลพฺภติ ปหีนํ สมุจฺฉินฺนํ วูปสนฺตํ ปฏิปสฺสทฺธํ อภพฺพุปฺปตฺติกํ ญาณคฺคินา ทฑฺฒนฺติ \textbf{สีมา}ติโค พฺราหฺมโณ ตสฺส นตฺถิ, ญตฺวา ว ทิสฺวา ว สมุคฺคหีตํ.}

\prose{\textbf{น ราคราคี น วิราครตฺโต}ติ \textbf{ราครตฺตา} วุจฺจนฺติ เย ปญฺจสุ กามคุเณสุ รตฺตา คิทฺธา คธิตา มุจฺฉิตา อชฺโฌสนฺนา ลคฺคา ลคฺคิตา ปลิพุทฺธา.\csromanspacer{} วิ\textbf{ราครตฺตา} วุจฺจนฺติ เย รูปาวจร-อรูปาวจรสมาปตฺตีสุ รตฺตา คิทฺธา คธิตา มุจฺฉิตา อชฺโฌสนฺนา ลคฺคา ลคฺคิตา ปลิพุทฺธา.\csromanspacer{} \textbf{น ราคราคี น วิราครตฺโต}ติ ยโต กามราโค จ รูปราโค จ อรูปราโค จ ปหีนา โหนฺติ อุจฺฉินฺนมูลา ตาลาวตฺถุกตา อนภาวํกตา\footnote{อนภาวกตา (สี), อนภาวํคตา (สฺยา)} อายติํ อนุปฺปาทธมฺมา.\csromanspacer{} เอตฺตาวตา น ราคราคี น วิราครตฺโต.}

\prose{\csromanfolio{78}\textbf{ตสฺสีธ นตฺถิ ปรมุคฺคหีตนฺ}ติ \textbf{ตสฺสา}ติ อรหโต ขีณาสวสฺส ตสฺส อิทํ ปรมํ อคฺคํ เสฏฺฐํ วิสิฏฺฐํ ปาโมกฺขํ อุตฺตมํ ปวรนฺติ คหิตํ ปรามฏฺฐํ อภินิวิฏฺฐํ อชฺโฌสิตํ อธิมุตฺตํ นตฺถิ น สนฺติ น สํวิชฺชติ นุปลพฺภติ ปหีนํ สมุจฺฉินฺนํ วูปสนฺตํ ปฏิปสฺสทฺธํ อภพฺพุปฺปตฺติกํ ญาณคฺคินา ทฑฺฒนฺติ ตสฺสีธ นตฺถิ ปรมุคฺคหีตํ.\csromanspacer{} เตนาห ภควา\csromandash{}}

\csromangathameasure{สีมาติโค พฺราหฺมโณ ตสฺส นตฺถิ, \\ ญตฺวา ว ทิสฺวา ว สมุคฺคหีตํ. \\ น ราคราคี น วิราครตฺโต,}
\csromangathagroup{\csromangathastanza{สีมาติโค พฺราหฺมโณ ตสฺส นตฺถิ, \\ ญตฺวา ว ทิสฺวา ว สมุคฺคหีตํ. \\ น ราคราคี น วิราครตฺโต,}}

\prose{\textbf{ตสฺสีธ นตฺถิ ปรมุคฺคหีตนฺ}ติ.\csromanspacer{} สุทฺธฏฺฐกสุตฺตนิทฺเทโส จตุตฺโถ.}
\csromansectionrule

\csromanheader{2}{5. ปรมฏฺฐกสุตฺตนิทฺเทส}
\csromanheader{2}{อถ ปรมฏฺฐกสุตฺตนิทฺเทสํ วกฺขติ\csromandash{}}
\proseitem{31}{\csromansymbolfootnote{*}{ขุ 1. 403 ปิฏฺเฐ.}\textbf{ปรมนฺติ ทิฏฺฐีสุ ปริพฺพสาโน},}

\prose{\textbf{ยทุตฺตริํ กุรุเต ชนฺตุ โลเก.}}

\csromangathameasure{\textbf{หีนาติ อญฺเญ ตโต สพฺพมาห,} \\ \textbf{ตสฺมา วิวาทานิ อวีติวตฺโต.}}
\csromangathagroup{\csromangathastanza{\textbf{หีนาติ อญฺเญ ตโต สพฺพมาห,} \\ \textbf{ตสฺมา วิวาทานิ อวีติวตฺโต.}}}

\prose{\textbf{ปรมนฺติ ทิฏฺฐีสุ ปริพฺพสาโน}ติ สนฺเตเก สมณพฺราหฺมณา ทิฏฺฐิคติกา, เต ทฺวาสฏฺฐิยา ทิฏฺฐิคตานํ อญฺญตรญฺญตรํ ทิฏฺฐิคตํ “อิทํ ปรมํ อคฺคํ เสฏฺฐํ วิสิฏฺฐํ ปาโมกฺขํ อุตฺตมํ ปวรนฺ”ติ คเหตฺวา อุคฺคเหตฺวา คณฺหิตฺวา ปรามสิตฺวา อภินิวิสิตฺวา สกาย สกาย ทิฏฺฐิยา วสนฺติ ปวสนฺติ\footnote{สํวสนฺติ (สฺยา) นตฺถิ สีหฬ-โปตฺถเก.} อาวสนฺติ ปริวสนฺติ.\csromanspacer{} ยถา อาคาริกา วา ฆเรสุ วสนฺติ, สาปตฺติกา วา อาปตฺตีสุ วสนฺติ, สกิเลสา วา กิเลเสสุ วสนฺติ.\csromanspacer{} เอวเมวํ สนฺเตเก สมณพฺราหฺมณา ทิฏฺฐิคติกา, เต ทฺวาสฏฺฐิยา ทิฏฺฐิคตานํ อญฺญตรญฺญตรํ ทิฏฺฐิคตํ “อิทํ ปรมํ อคฺคํ เสฏฺฐํ วิสิฏฺฐํ ปาโมกฺขํ อุตฺตมํ ปวรนฺ”ติ คเหตฺวา อุคฺคเหตฺวา คณฺหิตฺวา ปรามสิตฺวา อภินิวิสิตฺวา สกาย สกาย ทิฏฺฐิยา วสนฺติ ปวสนฺติ1 อาวสนฺติ ปริวสนฺตีติ ปรมนฺติ ทิฏฺฐีสุ ปริพฺพสาโน.}

\csromanheader{2}{5. ปรมฏฺฐกสุตฺตนิทฺเทส}
\prose{\csromanfolio{79}\textbf{ยทุตฺตริํ กุรุเต ชนฺตุ โลเก}ติ \textbf{ยทิ}ติ ยํ.\csromanspacer{} \textbf{อุตฺตริํ กุรุเต}ติ อุตฺตริํ กโรติ, อคฺคํ เสฏฺฐํ วิสิฏฺฐํ ปาโมกฺขํ อุตฺตมํ ปวรํ กโรติ. “อยํ สตฺถา สพฺพญฺญู”ติ อุตฺตริํ กโรติ, อคฺคํ เสฏฺฐํ วิสิฏฺฐํ ปาโมกฺขํ อุตฺตมํ ปวรํ กโรติ, “อยํ ธมฺโม สฺวากฺขาโต, อยํ คโณ สุปฺปฏิปนฺโน, อยํ ทิฏฺฐิ ภทฺทิกา, อยํ ปฏิปทา สุปญฺญตฺตา, อยํ มคฺโค นิยฺยานิโก”ติ อุตฺตริํ กโรติ, อคฺคํ เสฏฺฐํ วิสิฏฺฐํ ปาโมกฺขํ อุตฺตมํ ปวรํ กโรติ นิพฺพตฺเตติ อภินิพฺพตฺเตติ.\csromanspacer{} \textbf{ชนฺตู}ติ สตฺโต นโร ฯเปฯ มนุโช.\csromanspacer{} \textbf{โลเก}ติ อปายโลเก ฯเปฯ อายตนโลเกติ ยทุตฺตริํ กุรุเต ชนฺตุ โลเก.}

\prose{\textbf{หีนาติ อญฺเญ ตโต สพฺพมาหา}ติ อตฺตโน สตฺถารํ ธมฺมกฺขานํ คณํ ทิฏฺฐิํ ปฏิปทํ มคฺคํ ฐเปตฺวา สพฺเพ ปรปฺปวาเท ขิปติ อุกฺขิปติ ปริกฺขิปติ. “โส สตฺถา น สพฺพญฺญู, ธมฺโม น สฺวากฺขาโต, คโณ น สุปฺปฏิปนฺโน, ทิฏฺฐิ น ภทฺทิกา, ปฏิปทา น สุปญฺญตฺตา, มคฺโค น นิยฺยานิโก, นตฺเถตฺถ สุทฺธิ วา วิสุทฺธิ วา ปริสุทฺธิ วา มุตฺติ วา วิมุตฺติ วา ปริมุตฺติ วา, นตฺเถตฺถ สุชฺฌนฺติ วา วิสุชฺฌนฺติ วา ปริสุชฺฌนฺติ วา มุจฺจนฺติ วา วิมุจฺจนฺติ วา ปริมุจฺจนฺติ วา, หีนา นิหีนา โอมกา ลามกา ฉตุกฺกา ปริตฺตา”ติ เอวมาห เอวํ กเถติ เอวํ ภณติ เอวํ ทีปยติ เอวํ โวหรตีติ หีนาติ อญฺเญ ตโต สพฺพมาห.}

\prose{\textbf{ตสฺมา วิวาทานิ อวีติวตฺโต}ติ \textbf{ตสฺมา}ติ ตํการณา ตํเหตุ ตปฺปจฺจยา ตํนิทานา.\csromanspacer{} \textbf{วิวาทานี}ติ ทิฏฺฐิกลหานิ ทิฏฺฐิภณฺฑนานิ ทิฏฺฐิวิคฺคหานิ ทิฏฺฐิวิวาทานิ ทิฏฺฐิเมธคานิ จ.\csromanspacer{} \textbf{อวีติวตฺโต}ติ อนติกฺกนฺโต อสมติกฺกนฺโต อวีติวตฺโตติ \textbf{ตสฺมา} วิวาทานิ อวีติวตฺโต.\csromanspacer{} เตนาห ภควา\csromandash{}}

\prose{ปรมนฺติ ทิฏฺฐีสุ ปริพฺพสาโน,}

\prose{\textbf{ยทุตฺตริํ กุรุเต ชนฺตุ โลเก}.}

\csromangathameasure{หีนาติ อญฺเญ ตโต สพฺพมาห, \\ \textbf{ตสฺมา วิวาทานิ อวีติวตฺโต}ติ.}
\csromangathagroup{\csromangathastanza{หีนาติ อญฺเญ ตโต สพฺพมาห, \\ \textbf{ตสฺมา วิวาทานิ อวีติวตฺโต}ติ.}}

\proseitem{32}{\csromansymbolfootnote{*}{ขุ 1. 483 ปิฏฺเฐ.}ยทตฺตนี ปสฺสติ อานิสํสํ,}

\prose{\textbf{ทิฏฺเฐ สุเต สีลวเต มุเต วา.}}

\csromangathameasure{\textbf{ตเทว โส ตตฺถ สมุคฺคหาย,} \\ \textbf{นิหีนโต ปสฺสติ สพฺพมญฺญํ.}}
\csromangathagroup{\csromangathastanza{\textbf{ตเทว โส ตตฺถ สมุคฺคหาย,} \\ \textbf{นิหีนโต ปสฺสติ สพฺพมญฺญํ.}}}

\prose{\csromanfolio{80}\textbf{ยทตฺตนี ปสฺสติ อานิสํสํ, ทิฏฺเฐ สุเต สีลวเต มุเต วา}ติ \textbf{ยทตฺตนี}ติ ยํ อตฺตนิ.\csromanspacer{} \textbf{อตฺตา} วุจฺจติ ทิฏฺฐิคตํ.\csromanspacer{} อตฺตโน ทิฏฺฐิยา เทฺว อานิสํเส ปสฺสติ ทิฏฺฐธมฺมิกญฺจ อานิสํสํ สมฺปรายิกญฺจ อานิสํสํ.\csromanspacer{} กตโม ทิฏฺฐิยา ทิฏฺฐธมฺมิโก อานิสํโส, ยํทิฏฺฐิโก สตฺถา โหติ, ตํทิฏฺฐิกา สาวกา โหนฺติ.\csromanspacer{} ตํทิฏฺฐิกํ สตฺถารํ สาวกา สกฺกโรนฺติ ครุํ กโรนฺติ มาเนนฺติ ปูเชนฺติ, ลภติ จ ตโตนิทานํ จีวรปิณฺฑปาตเสนาสนคิลาน- ปจฺจยเภสชฺชปริกฺขารํ.\csromanspacer{} อยํ ทิฏฺฐิยา ทิฏฺฐธมฺมิโก อานิสํโส.\csromanspacer{} กตโม ทิฏฺฐิยา สมฺปรายิโก อานิสํโส, อยํ ทิฏฺฐิ อลํ นาคตฺตาย วา สุปณฺณตฺตาย วา ยกฺขตฺตาย วา อสุรตฺตาย วา คนฺธพฺพตฺตาย วา มหาราชตฺตาย วา อินฺทตฺตาย วา พฺรหฺมตฺตาย วา เทวตฺตาย วา, อยํ ทิฏฺฐิ อลํ สุทฺธิยา วิสุทฺธิยา ปริสุทฺธิยา, มุตฺติยา วิมุตฺติยา ปริมุตฺติยา, อิมาย ทิฏฺฐิยา สุชฺฌนฺติ วิสุชฺฌนฺติ ปริสุชฺฌนฺติ, มุจฺจนฺติ วิมุจฺจนฺติ ปริมุจฺจนฺติ, อิมาย ทิฏฺฐิยา สุชฺฌิสฺสามิ วิสุชฺฌิสฺสามิ ปริสุชฺฌิสฺสามิ, มุจฺจิสฺสามิ วิมุจฺจิสฺสามิ ปริมุจฺจิสฺสามิ, อายติํ ผลปาฏิกงฺขี โหติ.\csromanspacer{} อยํ ทิฏฺฐิยา สมฺปรายิโก อานิสํโส.\csromanspacer{} อตฺตโน ทิฏฺฐิยา อิเม เทฺว อานิสํเส ปสฺสติ, ทิฏฺฐสุทฺธิยาปิ เทฺว อานิสํเส ปสฺสติ, สุตสุทฺธิยาปิ เทฺว อานิสํเส ปสฺสติ, สีตสุทฺธิยาปิ เทฺว อานิสํเส ปสฺสติ, วตสุทฺธิยาปิ เทฺว อานิสํเส ปสฺสติ, มุตสุทฺธิยาปิ เทฺว อานิสํเส ปสฺสติ, ทิฏฺฐธมฺมิกญฺจ อานิสํสํ สมฺปรายิกญฺจ อานิสํสํ.\csromanspacer{} กตโม มุตสุทฺธิยา ทิฏฺฐธมฺมิโก อานิสํโส, ยํทิฏฺฐิโก สตฺถา โหติ, ตํทิฏฺฐิกา สาวกา โหนฺติ ฯเปฯ.\csromanspacer{} อยํ มุตสุทฺธิยา ทิฏฺฐธมฺมิโก อานิสํโส.\csromanspacer{} กตโม มุตสุทฺธิยา สมฺปรายิโก อานิสํโส, อยํ ทิฏฺฐิ อลํ นาคตฺตาย วา ฯเปฯ.\csromanspacer{} อยํ มุตสุทฺธิยา สมฺปรายิโก อานิสํโส.\csromanspacer{} มุตสุทฺธิยาปิ อิเม เทฺว อานิสํเส ปสฺสติ ทกฺขติ โอโลเกติ นิชฺฌายติ อุปปริกฺขตีติ \textbf{ยทตฺตนี} ปสฺสติ อานิสํสํ, ทิฏฺเฐ สุเต สีลวเต มุเต วา.}

\prose{\textbf{ตเทว โส ตตฺถ สมุคฺคหายา}ติ \textbf{ตเทวา}ติ ตํ ทิฏฺฐิคตํ.\csromanspacer{} \textbf{ตตฺถา}ติ สกาย ทิฏฺฐิยา สกาย ขนฺติยา สกาย รุจิยา สกาย ลทฺธิยา.\csromanspacer{} \textbf{สมุคฺคหายา}ติ อิทํ ปรมํ อคฺคํ เสฏฺฐํ วิสิฏฺฐํ ปาโมกฺขํ อุตฺตมํ ปวรนฺติ คเหตฺวา อุคฺคเหตฺวา คณฺหิตฺวา ปรามสิตฺวา อภินิวิสิตฺวาติ ตเทว โส ตตฺถ สมุคฺคหาย.}

\csromanheader{2}{5. ปรมฏฺฐกสุตฺตนิทฺเทส}
\prose{\csromanfolio{81}\textbf{นิหีนโต ปสฺสติ สพฺพมญฺญนฺ}ติ อญฺญํ สตฺถารํ ธมฺมกฺขานํ คณํ ทิฏฺฐิํ ปฏิปทํ มคฺคํ หีนโต นิหีนโต โอมกโต ลามกโต ฉตุกฺกโต ปริตฺตโต ทิสฺสติ ปสฺสติ ทกฺขติ โอโลเกติ นิชฺฌายติ อุปปริกฺขตีติ นิหีนโต ปสฺสติ สพฺพมญฺญํ.\csromanspacer{} เตนาห ภควา\csromandash{}}

\csromangathameasure{ยทตฺตนี ปสฺสติ อานิสํสํ, \\ ทิฏฺเฐ สุเต สีลวเต มุเต วา. \\ ตเทว โส ตตฺถ สมุคฺคหาย,}
\csromangathagroup{\csromangathastanza{ยทตฺตนี ปสฺสติ อานิสํสํ, \\ ทิฏฺเฐ สุเต สีลวเต มุเต วา. \\ ตเทว โส ตตฺถ สมุคฺคหาย,}}

\prose{\textbf{นิหีนโต ปสฺสติ สพฺพมญฺญนฺ}ติ.}

\proseitem{33}{\csromansymbolfootnote{*}{ขุ 1. 404 ปิฏฺเฐปิ.}ตํ วาปิ\footnote{จาปิ (สี)} \textbf{คนฺถํ กุสลา วทนฺติ,}}

\prose{\textbf{ยํ นิสฺสิโต ปสฺสติ หีนมญฺญํ.}}

\csromangathameasure{\textbf{ตสฺมา หิ ทิฏฺฐํ ว สุตํ มุตํ วา,} \\ \textbf{สีลพฺพตํ ภิกฺขุ น นิสฺสเยยฺย.}}
\csromangathagroup{\csromangathastanza{\textbf{ตสฺมา หิ ทิฏฺฐํ ว สุตํ มุตํ วา,} \\ \textbf{สีลพฺพตํ ภิกฺขุ น นิสฺสเยยฺย.}}}

\prose{\textbf{ตํ วาปิ คนฺถํ กุสลา วทนฺตี}ติ \textbf{กุสลา}ติ เย เต ขนฺธ\textbf{กุสลา} ธาตุ\textbf{กุสลา} อายตน\textbf{กุสลา} ปฏิจฺจสมุปฺปาท\textbf{กุสลา} สติปฏฺฐาน\textbf{กุสลา} สมฺมปฺปธาน\textbf{กุสลา} อิทฺธิปาท\textbf{กุสลา} อินฺทฺริย\textbf{กุสลา} พล\textbf{กุสลา} โพชฺฌงฺค\textbf{กุสลา} มคฺค\textbf{กุสลา} ผล\textbf{กุสลา} นิพฺพาน\textbf{กุสลา}, เต \textbf{กุสลา} เอวํ วทนฺติ “คนฺโถ เอโส, ลคฺคนํ เอตํ, พนฺธนํ เอตํ, ปลิโพโธ เอโส”ติ เอวํ วทนฺติ เอวํ กเถนฺติ เอวํ ภณนฺติ เอวํ ทีปยนฺติ เอวํ โวหรนฺตีติ ตํ วาปิ คนฺถํ \textbf{กุสลา} วทนฺติ.}

\prose{\textbf{ยํ นิสฺสิโต ปสฺสติ หีนมญฺญนฺ}ติ \textbf{ยํ นิสฺสิโต}ติ ยํ สตฺถารํ ธมฺมกฺขานํ คณํ ทิฏฺฐิํ ปฏิปทํ มคฺคํ นิสฺสิโต สนฺนิสฺสิโต อลฺลีโน อุปคโต อชฺโฌสิโต อธิมุตฺโต.\csromanspacer{} \textbf{ปสฺสติ หีนมญฺญนฺ}ติ อญฺญํ สตฺถารํ ธมฺมกฺขานํ คณํ ทิฏฺฐิํ ปฏิปทํ มคฺคํ หีนโต นิหีนโต โอมกโต ลามกโต ฉตุกฺกโต ปริตฺตโต ทิสฺสติ ปสฺสติ ทกฺขติ โอโลเกติ นิชฺฌายติ อุปนิชฺฌายติ อุปปริกฺขตีติ \textbf{ยํ นิสฺสิโต} ปสฺสติ หีนมญฺญํ.}

\prose{\textbf{ตสฺมา หิ ทิฏฺฐํ ว สุตํ มุตํ วา, สีลพฺพตํ ภิกฺขุ น นิสฺสเยยฺยา}ติ \textbf{ตสฺมา}ติ \textbf{ตสฺมา} ตํการณา ตํเหตุ ตปฺปจฺจยา ตํนิทานา ทิฏฺฐํ วา ทิฏฺฐสุทฺธิํ วา สุตํ วา สุตสุทฺธิํ วา มุตํ วา มุตสุทฺธิํ วา สีลํ \csromanfolio{82}วา สีลสุทฺธิํ วา วตํ วา วตสุทฺธิํ วา น นิสฺสเยยฺย น คเณยฺย น ปรามเสยฺย นาภินิเวเสยฺยาติ ตสฺมา หิ ทิฏฺฐํ ว สุตํ มุตํ วา, สีลพฺพตํ ภิกฺขุ น นิสฺสเยยฺย.\csromanspacer{} เตนาห ภควา\csromandash{}}

\csromangathameasure{ตํ วาปิ คนฺถํ กุสลา วทนฺติ, \\ ยํ นิสฺสิโต ปสฺสติ หีนมญฺญํ. \\ ตสฺมา หิ ทิฏฺฐํ ว สุตํ มุตํ วา, \\ สีลพฺพตํ ภิกฺขุ น นิสฺสเยยฺยาติ.}
\csromangathagroup{\csromangathastanza{ตํ วาปิ คนฺถํ กุสลา วทนฺติ, \\ ยํ นิสฺสิโต ปสฺสติ หีนมญฺญํ. \\ ตสฺมา หิ ทิฏฺฐํ ว สุตํ มุตํ วา, \\ สีลพฺพตํ ภิกฺขุ น นิสฺสเยยฺยาติ.}}

\proseitem{34}{\csromansymbolfootnote{*}{ขุ 1. 404 ปิฏฺเฐปิ.}ทิฏฺฐิมฺปิ โลกสฺมิํ น กปฺปเยยฺย,}

\prose{\textbf{ญาเณน วา สีลวเตน วาปิ.}}

\csromangathameasure{\textbf{สโมติ อตฺตานมนูปเนยฺย,} \\ \textbf{หีโน น มญฺเญถ วิเสสิ วาปิ.}}
\csromangathagroup{\csromangathastanza{\textbf{สโมติ อตฺตานมนูปเนยฺย,} \\ \textbf{หีโน น มญฺเญถ วิเสสิ วาปิ.}}}

\prose{\textbf{ทิฏฺฐิมฺปิ โลกสฺมิํ น กปฺปเยยฺย, ญาเณน วา สีลวเตน วาปี}ติ อฏฺฐสมาปตฺติญาเณน วา ปญฺจาภิญฺญาญาเณน วา มิจฺฉาญาเณน วา, สีเลน วา วเตน วา สีลพฺพเตน วา ทิฏฺฐิํ น กปฺปเยยฺย น ชเนยฺย น สญฺชเนยฺย น นิพฺพตฺเตยฺย น อภินิพฺพตฺเตยฺย.\csromanspacer{}\csromanspacer{}\textbf{โลกสฺมินฺ}ติ อปายโลเก ฯเปฯ อายตนโลเกติ ทิฏฺฐิมฺปิ โลกสฺมิํ น กปฺปเยยฺย, ญาเณน วา สีลวเตน วาปิ.}

\prose{\textbf{สโมติ อตฺตานมนูปเนยฺยา}ติ “สทิโสหมสฺมี”ติ อตฺตานํ น อุปเนยฺย ชาติยา วา โคตฺเตน วา โกลปุตฺติเยน วา วณฺณโปกฺขรตาย วา ธเนน วา อชฺเฌเนน วา กมฺมายตเนน วา สิปฺปายตเนน วา วิชฺชาฏฺฐาเนน วา สุเตน วา ปฏิภาเนน วา อญฺญตรญฺญตเรน วา วตฺถุนาติ สโมติ อตฺตานมนูปเนยฺย.}

\prose{\textbf{หีโน น มญฺเญถ วิเสสิ วาปี}ติ “หีโนหมสฺมี”ติ อตฺตานํ น อุปเนยฺย ชาติยา วา โคตฺเตน วา ฯเปฯ อญฺญตรญฺญตเรน วา วตฺถุนา. “เสยฺโยหมสฺมี”ติ อตฺตานํ น อุปเนยฺย ชาติยา วา โคตฺเตน วา ฯเปฯ อญฺญตรญฺญตเรน วา วตฺถุนาติ หีโน น มญฺเญถ วิเสสิ วาปิ.\csromanspacer{} เตนาห ภควา\csromandash{}}

\csromanheader{2}{5. ปรมฏฺฐกสุตฺตนิทฺเทส}
\csromangathameasure{ทิฏฺฐิมฺปิ โลกสฺมิํ น กปฺปเยยฺย, \\ ญาเณน วา สีลวเตน วาปิ. \\ สโมติ อตฺตานมนูปเนยฺย, \\ หีโน น มญฺเญถ วิเสสิ วาปีติ.}
\csromangathagroup{\csromangathastanza{\csromanfolio{83}ทิฏฺฐิมฺปิ โลกสฺมิํ น กปฺปเยยฺย, \\ ญาเณน วา สีลวเตน วาปิ. \\ สโมติ อตฺตานมนูปเนยฺย, \\ หีโน น มญฺเญถ วิเสสิ วาปีติ.}}

\proseitem{35}{\csromansymbolfootnote{*}{ขุ 1. 404 ปิฏฺเฐปิ.}\textbf{อตฺตํ ปหาย อนุปาทิยาโน},}

\prose{\textbf{ญาเณนปิ โส นิสฺสยํ โน กโรติ.}}

\csromangathameasure{\textbf{ส เว วิยตฺเตสุ น วคฺคสารี,} \\ \textbf{ทิฏฺฐิมฺปิ โส น ปจฺเจติ กิญฺจิ.}}
\csromangathagroup{\csromangathastanza{\textbf{ส เว วิยตฺเตสุ น วคฺคสารี,} \\ \textbf{ทิฏฺฐิมฺปิ โส น ปจฺเจติ กิญฺจิ.}}}

\prose{\textbf{อตฺตํ ปหาย อนุปาทิยาโน}ติ \textbf{อตฺตํ ปหายา}ติ อตฺตทิฏฺฐิํ ปหาย.\csromanspacer{} \textbf{อตฺตํ ปหายา}ติ คาหํ\footnote{อตฺตคาหํ (สี, ก)} ปหาย.\csromanspacer{} \textbf{อตฺตํ ปหายา}ติ ตณฺหาวเสน ทิฏฺฐิวเสน คหิตํ ปรามฏฺฐํ อภินิวิฏฺฐํ อชฺโฌสิตํ อธิมุตฺตํ ปหาย ปชหิตฺวา วิโนเทตฺวา พฺยนฺติํ กริตฺวา อนภาวํ คเมตฺวาติ อตฺตํ ปหาย.\csromanspacer{}\csromanspacer{}\textbf{อนุปาทิยาโน}ติ จตูหิ อุปาทาเนหิ อนุปาทิยมาโน อคณฺหมาโน อปรามาสมาโน อนภินิวิสมาโนติ อตฺตํ ปหาย อนุปาทิยาโน.}

\prose{\textbf{ญาเณนปิ โส นิสฺสยํ โน กโรตี}ติ อฏฺฐสมาปตฺติญาเณน วา ปญฺจาภิญฺญาญาเณน วา มิจฺฉาญาเณน วา ตณฺหานิสฺสยํ วา ทิฏฺฐินิสฺสยํ วา น กโรติ น ชเนติ น สญฺชเนติ น นิพฺพตฺเตติ น อภินิพฺพตฺเตตีติ ญาเณนปิ โส นิสฺสยํ โน กโรติ.}

\prose{\textbf{ส เว วิยตฺเตสุ น วคฺคสารี}ติ ส เว วิยตฺเตสุ ภินฺเนสุ เทฺวชฺฌาปนฺเนสุ เทฺวฬฺหกชาเตสุ นานาทิฏฺฐิเกสุ นานาขนฺติเกสุ นานารุจิเกสุ นานาลทฺธิเกสุ นานาทิฏฺฐินิสฺสยํ นิสฺสิเตสุ ฉนฺทาคติํ คจฺฉนฺเตสุ, โทสาคติํ คจฺฉนฺเตสุ, โมหาคติํ คจฺฉนฺเตสุ, ภยาคติํ คจฺฉนฺเตสุ น ฉนฺทาคติํ คจฺฉติ, น โทสาคติํ คจฺฉติ, น โมหาคติํ คจฺฉติ, น ภยาคติํ คจฺฉติ, น ราควเสน คจฺฉติ, น โทสวเสน คจฺฉติ, น โมหวเสน คจฺฉติ, น มานวเสน คจฺฉติ, น ทิฏฺฐิวเสน คจฺฉติ, น อุทฺธจฺจวเสน คจฺฉติ, น วิจิกิจฺฉาวเสน คจฺฉติ, น อนุสยวเสน คจฺฉติ, น วคฺเคหิ ธมฺเมหิ ยายติ นิยฺยติ วุยฺหติ สํหรียตีติ ส เว วิยตฺเตสุ น วคฺคสารี.}

\prose{\csromanfolio{84}\textbf{ทิฏฺฐิมฺปิ โส น ปจฺเจติ กิญฺจี}ติ ตสฺส ทฺวาสฏฺฐิ ทิฏฺฐิคตานิ ปหีนานิ สมุจฺฉินฺนานิ วูปสนฺตานิ ปฏิปสฺสทฺธานิ อภพฺพุปฺปตฺติกานิ ญาณคฺคินา ทฑฺฒานิ.\csromanspacer{} โส กิญฺจิ ทิฏฺฐิคตํ น ปจฺเจติ น ปจฺจาคจฺฉตีติ ทิฏฺฐิมฺปิ โส น ปจฺเจติ กิญฺจิ.\csromanspacer{} เตนาห ภควา\csromandash{}}

\csromangathameasure{อตฺตํ ปหาย อนุปาทิยาโน, \\ ญาเณนปิ โส นิสฺสยํ โน กโรติ. \\ ส เว วิยตฺเตสุ น วคฺคสารี,}
\csromangathagroup{\csromangathastanza{อตฺตํ ปหาย อนุปาทิยาโน, \\ ญาเณนปิ โส นิสฺสยํ โน กโรติ. \\ ส เว วิยตฺเตสุ น วคฺคสารี,}}

\prose{\textbf{ทิฏฺฐิมฺปิ โส น ปจฺเจติ กิญฺจี}ติ.}

\proseitem{36}{\csromansymbolfootnote{*}{ขุ 1. 404 ปิฏฺเฐปิ.}ยสฺสูภยนฺเต ปณิธีธ นตฺถิ,}

\prose{\textbf{ภวาภวาย อิธ วา หุรํ วา.}}

\csromangathameasure{\textbf{นิเวสนา ตสฺส น สนฺติ เกจิ,} \\ \textbf{ธมฺเมสุ นิจฺเฉยฺย สมุคฺคหีตํ.}}
\csromangathagroup{\csromangathastanza{\textbf{นิเวสนา ตสฺส น สนฺติ เกจิ,} \\ \textbf{ธมฺเมสุ นิจฺเฉยฺย สมุคฺคหีตํ.}}}

\prose{\textbf{ยสฺสูภยนฺเต ปณิธีธ นตฺถิ, ภวาภวาย อิธ วา หุรํ วา}ติ \textbf{ยสฺสา}ติ อรหโต ขีณาสวสฺส.\csromanspacer{} \textbf{อนฺโต}ติ\footnote{อนฺตาติ (สฺยา) ขุ 8. 44 ปิฏฺเฐปิ.} ผสฺโส เอโก \textbf{อนฺโต}, ผสฺสสมุทโย ทุติโย \textbf{อนฺโต}, อตีโต เอโก \textbf{อนฺโต}, อนาคโต ทุติโย \textbf{อนฺโต}, สุขา เวทนา เอโก \textbf{อนฺโต}, ทุกฺขา เวทนา ทุติโย \textbf{อนฺโต}, นามํ เอโก \textbf{อนฺโต}, รูปํ ทุติโย \textbf{อนฺโต}, ฉ อชฺฌตฺติกานิ อายตนานิ เอโก \textbf{อนฺโต}, ฉ พาหิรานิ อายตนานิ ทุติโย \textbf{อนฺโต}, สกฺกาโย เอโก \textbf{อนฺโต}, สกฺกายสมุทโย ทุติโย \textbf{อนฺโต}.\csromanspacer{} \textbf{ปณิธิ} วุจฺจติ ตณฺหา.\csromanspacer{} โย ราโค สาราโค ฯเปฯ อภิชฺฌา โลโภ อกุสลมูลํ.}

\prose{\textbf{ภวาภวายา}ติ ภวาภวาย กมฺมภวาย ปุนพฺภวาย กามภวาย, กมฺมภวาย กามภวาย ปุนพฺภวาย รูปภวาย, กมฺมภวาย รูปภวาย ปุนพฺภวาย อรูปภวาย, กมฺมภวาย อรูปภวาย ปุนพฺภวาย ปุนปฺปุนภวาย, ปุนปฺปุนคติยา ปุนปฺปุน-อุปปตฺติยา ปุนปฺปุนปฏิสนฺธิยา ปุนปฺปุน-อตฺตภาวาภินิพฺพตฺติยา.\csromanspacer{}\csromanspacer{}\textbf{อิธา}ติ สกตฺตภาโว.\csromanspacer{} \textbf{หุรา}ติ ปรตฺตภาโว.\csromanspacer{} \textbf{อิธา}ติ สกรูปเวทนาสญฺญาสงฺขารวิญฺญาณํ.\csromanspacer{} \textbf{หุรา}ติ ปรรูปเวทนาสญฺญาสงฺขารวิญฺญาณํ.\csromanspacer{} \textbf{อิธา}ติ ฉ อชฺฌตฺติกานิ อายตนานิ.\csromanspacer{} \textbf{หุรา}ติ ฉ พาหิรานิ อายตนานิ.\csromanspacer{} \textbf{อิธา}ติ มนุสฺสโลโก.\csromanspacer{} \textbf{หุรา}ติ}

\vspace{\tipitakaparskip}
\csromancenter{ปรมฏฺฐกสุตฺตนิทฺเทส เทวโลโก.\csromanspacer{} \textbf{อิธา}ติ กามธาตุ.\csromanspacer{} \textbf{หุรา}ติ รูปธาตุ อรูปธาตุ.\csromanspacer{} \textbf{อิธา}ติ กามธาตุ รูปธาตุ.\csromanspacer{} \textbf{หุรา}ติ อรูปธาตุ.\csromanspacer{}\csromanspacer{}\textbf{ยสฺสูภยนฺเต ปณิธีธ นตฺถิ, ภวาภวาย อิธ วา หุรํ วา}ติ ยสฺส อุโภ อนฺเต จ ภวาภวาย จ อิธ หุรํ จ ปณิธิ ตณฺหา นตฺถิ น สนฺติ น สํวิชฺชติ นุปลพฺภติ ปหีนา สมุจฺฉินฺนา วูปสนฺตา ปฏิปสฺสทฺธา อภพฺพุปฺปตฺติกา ญาณคฺคินา ทฑฺฒาติ ยสฺสูภยนฺเต ปณิธีธ นตฺถิ, ภวาภวาย อิธ วา หุรํ วา.}
\prose{\csromanfolio{85}\textbf{นิเวสนา ตสฺส น สนฺติ เกจี}ติ \textbf{นิเวสนา}ติ เทฺว \textbf{นิเวสนา} ตณฺหา\textbf{นิเวสนา} จ ทิฏฺฐิ\textbf{นิเวสนา} จ ฯเปฯ อยํ ตณฺหา\textbf{นิเวสนา} ฯเปฯ อยํ ทิฏฺฐิ\textbf{นิเวสนา}.\csromanspacer{} \textbf{ตสฺสา}ติ อรหโต ขีณาสวสฺส.\csromanspacer{}\csromanspacer{}\textbf{นิเวสนา ตสฺส น สนฺติ เกจี}ติ \textbf{นิเวสนา} ตสฺส น สนฺติ เกจิ นตฺถิ น สนฺติ น สํวิชฺชนฺติ นุปลพฺภนฺติ ปหีนา สมุจฺฉินฺนา วูปสนฺตา ปฏิปสฺสทฺธา อภพฺพุปฺปตฺติกา ญาณคฺคินา ทฑฺฒาติ \textbf{นิเวสนา} ตสฺส น สนฺติ เกจิ.}

\prose{\textbf{ธมฺเมสุ นิจฺเฉยฺย สมุคฺคหีตนฺ}ติ \textbf{ธมฺเมสู}ติ ทฺวาสฏฺฐิยา ทิฏฺฐิคเตสุ.\csromanspacer{} \textbf{นิจฺเฉยฺยา}ติ นิจฺฉินิตฺวา วินิจฺฉินิตฺวา วิจินิตฺวา ปวิจินิตฺวา ตุลยิตฺวา ตีรยิตฺวา วิภาวยิตฺวา วิภูตํ กตฺวา.\csromanspacer{} \textbf{สมุคฺคหีตนฺ}ติ โอธิคฺคาโห พิลคฺคาโห วรคฺคาโห โกฏฺฐาสคฺคาโห อุจฺจยคฺคาโห สมุจฺจยคฺคาโห, อิทํ สจฺจํ ตจฺฉํ ตถํ ภูตํ ยาถาวํ อวิปรีตนฺติ คหิตํ ปรามฏฺฐํ, อภินิวิฏฺฐํ อชฺโฌสิตํ อธิมุตฺตํ นตฺถิ น สนฺติ น สํวิชฺชติ นุปลพฺภติ, ปหีนํ สมุจฺฉินฺนํ วูปสนฺตํ ปฏิปสฺสทฺธํ อภพฺพุปฺปตฺติกํ ญาณคฺคินา ทฑฺฒนฺติ ธมฺเมสุ นิจฺเฉยฺย สมุคฺคหีตํ.\csromanspacer{} เตนาห ภควา\csromandash{}}

\csromangathameasure{ยสฺสูภยนฺเต ปณิธีธ นตฺถิ, \\ ภวาภวาย อิธ วา หุรํ วา. \\ นิเวสนา ตสฺส น สนฺติ เกจิ,}
\csromangathagroup{\csromangathastanza{ยสฺสูภยนฺเต ปณิธีธ นตฺถิ, \\ ภวาภวาย อิธ วา หุรํ วา. \\ นิเวสนา ตสฺส น สนฺติ เกจิ,}}

\prose{\textbf{ธมฺเมสุ นิจฺเฉยฺย สมุคฺคหีตนฺ}ติ.}

\proseitem{37}{\csromansymbolfootnote{*}{ขุ 1. 404 ปิฏฺเฐปิ.}ตสฺสีธ ทิฏฺเฐ ว สุเต มุเต วา,}

\prose{\textbf{ปกปฺปิตา นตฺถิ อณูปิ สญฺญา.}}

\csromangathameasure{\textbf{ตํ พฺราหฺมณํ ทิฏฺฐิมนาทิยานํ,} \\ \textbf{เกนีธ โลกสฺมิํ วิกปฺปเยยฺย.}}
\csromangathagroup{\csromangathastanza{\textbf{ตํ พฺราหฺมณํ ทิฏฺฐิมนาทิยานํ,} \\ \textbf{เกนีธ โลกสฺมิํ วิกปฺปเยยฺย.}}}

\prose{\csromanfolio{86}\textbf{ตสฺสีธ ทิฏฺเฐ ว สุเต มุเต วา, ปกปฺปิตา นตฺถิ อณูปิ สญฺญา}ติ \textbf{ตสฺสา}ติ อรหโต ขิณาสวสฺส ตสฺส ทิฏฺเฐ วา ทิฏฺฐสุทฺธิยา วา สุเต วา สุตสุทฺธิยา วา มุเต วา มุตสุทฺธิยา วา สญฺญาปุพฺพงฺคมตา สญฺญาวิกปฺปเยยฺยตา สญฺญาวิคฺคเหน สญฺญาย อุฏฺฐปิตา สมุฏฺฐปิตา กปฺปิตา ปกปฺปิตา สงฺขตา อภิสงฺขตา สณฺฐปิตา, ทิฏฺฐิ นตฺถิ น สนฺติ น สํวิชฺชติ นุปลพฺภติ, ปหีนา สมุจฺฉินฺนา วูปสนฺตา ปฏิปสฺสทฺธา อภพฺพุปฺปตฺติกา ญาณคฺคินา ทฑฺฒาติ ตสฺสีธ ทิฏฺเฐ ว สุเต มุเต วา, ปกปฺปิตา นตฺถิ อณูปิ สญฺญา.}

\prose{\textbf{ตํ พฺราหฺมณํ ทิฏฺฐิมนาทิยานนฺ}ติ พฺราหฺมโณติ สตฺตนฺนํ ธมฺมานํ พาหิตตฺตา พฺราหฺมโณ สกฺขายทิฏฺฐิ พาหิตา โหติ ฯเปฯ อสิโต ตาทิ ปวุจฺจเต ส พฺรหฺมา.\csromanspacer{} \textbf{ตํ พฺราหฺมณํ ทิฏฺฐิมนาทิยานนฺ}ติ ตํ พฺราหฺมณํ ทิฏฺฐิมนาทิยนฺตํ อคณฺหนฺตํ อปรามสนฺตํ อนภินิเวสนฺตนฺติ ตํ พฺราหฺมณํ ทิฏฺฐิมนาทิยานํ.}

\prose{\textbf{เกนีธ โลกสฺมิํ วิกปฺปเยยฺยา}ติ \textbf{กปฺปา}ติ เทฺว \textbf{กปฺปา} ตณฺหากปฺโป จ ทิฏฺฐิกปฺโป จ ฯเปฯ อยํ ตณฺหากปฺโป ฯเปฯ อยํ ทิฏฺฐิกปฺโป.\csromanspacer{} ตสฺส ตณฺหากปฺโป ปหีโน, ทิฏฺฐิกปฺโป ปฏินิสฺสฏฺโฐ, ตณฺหากปฺปสฺส ปหีนตฺตา, ทิฏฺฐิกปฺปสฺส ปฏินิสฺสฏฺฐตฺตา เกน ราเคน กปฺเปยฺย, เกน โทเสน กปฺเปยฺย, เกน โมเหน กปฺเปยฺย, เกน มาเนน กปฺเปยฺย, กาย ทิฏฺฐิยา กปฺเปยฺย เกน อุทฺธจฺเจน กปฺเปยฺย, กาย วิจิกิจฺฉาย กปฺเปยฺย, เกหิ อนุสเยหิ กปฺเปยฺย “รตฺโต”ติ วา “ทุฏฺโฐ”ติ วา “มูโฬฺห”ติ วา “วินิพทฺโธ”ติ วา “ปรามฏฺโฐ”ติ วา “วิกฺเขปคโต”ติ วา “อนิฏฺฐงฺคโต”ติ วา “ถามคโต”ติ วา.\csromanspacer{} เต อภิสงฺขารา ปหีนา, อภิสงฺขารานํ ปหีนตฺตา คติโย เกน กปฺเปยฺย “เนรยิโก”ติ วา “ติรจฺฉานโยนิโก”ติ วา “เปตฺติวิสยิโก”ติ วา “มนุสฺโส”ติ วา “เทโว”ติ วา “รูปี”ติ วา “อรูปี”ติ วา “สญฺญี”ติ วา “อสญฺญี”ติ วา “เนวสญฺญีนาสญฺญี”ติ วา.\csromanspacer{} โส เหตุ นตฺถิ, ปจฺจโย นตฺถิ, การณํ นตฺถิ, เยน กปฺเปยฺย วิกปฺเปยฺย วิกปฺปํ อาปชฺเชยฺย.\csromanspacer{}\csromanspacer{}\textbf{โลกสฺมินฺ}ติ อปายโลเก ฯเปฯ อายตนโลเกติ เกนีธ โลกสฺมิํ วิกปฺปเยยฺย.\csromanspacer{} เตนาห ภควา\csromandash{}}

\prose{ตสฺสีธ ทิฏฺเฐ ว สุเต มุเต วา,}

\prose{ปกปฺปิตา นตฺถิ อณูปิ สญฺญา.}

\prose{ตํ พฺราหฺมณํ ทิฏฺฐิมนาทิยานํ,}

\prose{\textbf{เกนีธ โลกสฺมิํ วิกปฺปเยยฺยา}ติ.}

\csromanheader{2}{5. ปรมฏฺฐกสุตฺตนิทฺเทส}
\proseitem{38}{\csromanfolio{87}\csromansymbolfootnote{*}{ขุ 1. 404 ปิฏฺเฐปิ.}น กปฺปยนฺติ น ปุเรกฺขโรนฺติ,}

\prose{\textbf{ธมฺมาปิ เตสํ น ปฏิจฺฉิตาเส.}}

\csromangathameasure{\textbf{น พฺราหฺมโณ สีลวเตน เนยฺโย,} \\ \textbf{ปารงฺคโต น ปจฺเจติ ตาที.}}
\csromangathagroup{\csromangathastanza{\textbf{น พฺราหฺมโณ สีลวเตน เนยฺโย,} \\ \textbf{ปารงฺคโต น ปจฺเจติ ตาที.}}}

\csromanmatikaanchor{mtk.6}
\csromantocmarkref{subsection}{5. ปรมฏฺฐกสุตฺตนิทฺเทส}{mtk.6}
\bookmark[dest={mtk.6},level=2]{5. ปรมฏฺฐกสุตฺตนิทฺเทส}
\prose{\textbf{น กปฺปยนฺติ น ปุรกฺขโรนฺตี}ติ \textbf{กปฺปา}ติ เทฺว \textbf{กปฺปา} ตณฺหากปฺโป จ ทิฏฺฐิกปฺโป จ.\csromanspacer{} กตโม ตณฺหากปฺโป, ยาวตา ตณฺหาสงฺขาเตน สีมกตํ มริยาทิกตํ โอธิกตํ ปริยนฺตกตํ ปริคฺคหิตํ มมายิตํ, อิทํ มมํ, เอตํ มมํ, เอตฺตกํ มมํ, เอตฺตาวตา มมํ, มม รูปา สทฺทา คนฺธา รสา โผฏฺฐพฺพา อตฺถรณา ปาวุรณา ทาสิทาสา\footnote{ทาสีทาสา (สฺยา, ก)} อเชฬกา กุกฺกุฏสูกรา หตฺถิควาสฺสวฬวา เขตฺตํ วตฺถุ หิรญฺญํ สุวณฺณํ คามนิคมราชธานิโย รฏฺฐญฺจ ชนปโท จ โกโส จ โกฏฺฐาคารญฺจ เกวลมฺปิ มหาปถวิํ ตณฺหาวเสน มมายติ, ยาวตา อฏฺฐสตตณฺหาวิจริตํ, อยํ ตณฺหากปฺโป.\csromanspacer{}\csromanspacer{}กตโม ทิฏฺฐิกปฺโป, วีสติวตฺถุกา สกฺกายทิฏฺฐิ, ทสวตฺถุกา มิจฺฉาทิฏฺฐิ, ทสวตฺถุกา อนฺตคฺคาหิกา ทิฏฺฐิ. + ยา เอวรูปา ทิฏฺฐิ ทิฏฺฐิคตํ ทิฏฺฐิคหนํ ทิฏฺฐิกนฺตารํ ทิฏฺฐิวิสูกายิกํ ทิฏฺฐิวิปฺผนฺทิตํ ทิฏฺฐิสญฺโญชนํ คาโห ปฏิคฺคาโห อภินิเวโส ปรามาโส กุมฺมคฺโค มิจฺฉาปโถ มิจฺฉตฺตํ ติตฺถายตนํ วิปริเยสคฺคาโห\footnote{วิปริยาสคฺคาโห (ก)} วิปรีตคฺคาโห วิปลฺลาสคฺคาโห มิจฺฉาคาโห, อยาถาวกสฺมิํ ยาถาวกนฺติ คาโห, ยาวตา ทฺวาสฏฺฐิ ทิฏฺฐิคตานิ อยํ ทิฏฺฐิกปฺโป.\csromanspacer{} เตสํ ตณฺหากปฺโป ปหีโน, ทิฏฺฐิกปฺโป ปฏินิสฺสฏฺโฐ, ตณฺหากปฺปสฺส ปหีนตฺตา, ทิฏฺฐิกปฺปสฺส ปฏินิสฺสฏฺฐตฺตา ตณฺหากปฺปํ วา ทิฏฺฐิกปฺปํ วา น กปฺเปนฺติ น ชเนนฺติ น สญฺชเนนฺติ น นิพฺพตฺเตนฺติ น อภินิพฺพตฺเตนฺตีติ น กปฺปยนฺติ.}

\prose{\textbf{น ปุเรกฺขโรนฺตี}ติ \textbf{ปุเรกฺขารา}ติ เทฺว \textbf{ปุเรกฺขารา} ตณฺหาปุเรกฺขาโร จ ทิฏฺฐิปุเรกฺขาโร จ ฯเปฯ อยํ ตณฺหาปุเรกฺขาโร ฯเปฯ อยํ ทิฏฺฐิปุเรกฺขาโร.\csromanspacer{} เตสํ ตณฺหาปุเรกฺขาโร ปหีโน, ทิฏฺฐิปุเรกฺขาโร ปฏินิสฺสฏฺโฐ, ตณฺหาปุเรกฺขารสฺส ปหีนตฺตา, ทิฏฺฐิปุเรกฺขารสฺส ปฏินิสฺสฏฺฐตฺตา น ตณฺหํ วา น ทิฏฺฐิํ วา ปุรโต กตฺวา จรนฺติ, น ตณฺหาธชา น ตณฺหาเกตู น ตณฺหาธิปเตยฺยา น ทิฏฺฐิธชา น ทิฏฺฐิเกตู น ทิฏฺฐาธิปเตยฺยา น ตณฺหาย วา น ทิฏฺฐิยา วา ปริวาเรตฺวา จรนฺตีติ น กปฺปยนฺติ น ปุเรกฺขโรนฺติ.}

\prose{\csromanfolio{88}\textbf{ธมฺมาปิ เตสํ น ปฏิจฺฉิตาเส}ติ \textbf{ธมฺมา} วุจฺจนฺติ ทฺวาสฏฺฐิ ทิฏฺฐิคตานิ.\csromanspacer{} \textbf{เตสนฺ}ติ เตสํ อรหนฺตานํ ขีณาสวานํ.\csromanspacer{} \textbf{น ปฏิจฺฉิตาเส}ติ “สสฺสโต โลโก, อิทเมว สจฺจํ โมฆมญฺญนฺ”ติ น ปฏิจฺฉิตาเส, “อสสฺสโต โลโก.\csromanspacer{} อนฺตวา โลโก.\csromanspacer{} อนนฺตวา โลโก.\csromanspacer{} ตํ ชีวํ ตํ สรีรํ.\csromanspacer{} อญฺญํ ชีวํ อญฺญํ สรีรํ.\csromanspacer{} โหติ ตถาคโต ปรํ มรณา.\csromanspacer{} น โหติ ตถาคโต ปรํ มรณา.\csromanspacer{} โหติ จ น จ โหติ ตถาคโต ปรํ มรณา, เนว โหติ น น โหติ ตถาคโต ปรํ มรณา, อิทเมว สจฺจํ โมฆมญฺญนฺ”ติ น ปฏิจฺฉิตาเสติ \textbf{ธมฺมา}ปิ เตสํ น ปฏิจฺฉิตาเส.}

\prose{\textbf{น พฺราหฺมโณ สีลวเตน เนยฺโย}ติ \textbf{นา}ติ ปฏิกฺเขโป.\csromanspacer{} \textbf{พฺราหฺมโณ}ติ สตฺตนฺนํ \textbf{ธมฺมา}นํ พาหิตตฺตา พฺราหฺมโณ.\csromanspacer{} สกฺกายทิฏฺฐิ พาหิตา โหติ ฯเปฯ อสิโต ตาทิ วุจฺจเต ส พฺรหฺมา.\csromanspacer{}\csromanspacer{}\textbf{น พฺราหฺมโณ สีลวเตน เนยฺโย}ติ พฺราหฺมโณ สีเลน วา วเตน วา สีลพฺพเตน วา น ยายติ น นิยฺยติ น วุยฺหติ น สํหรียตีติ น พฺราหฺมโณ สีลวเตน เนยฺโย.}

\prose{\textbf{ปารงฺคโต น ปจฺเจติ ตาที}ติ \textbf{ปารํ} วุจฺจติ อมตํ นิพฺพานํ.\csromanspacer{} โย โส สพฺพสงฺขารสมโถ สพฺพูปธิปฏินิสฺสคฺโค ตณฺหกฺขโย วิราโค นิโรโธ นิพฺพานํ, โส ปารงฺคโต ปารปฺปตฺโต อนฺตคโต อนฺตปฺปตฺโต โกฏิคโต โกฏิปฺปตฺโต (วิตฺถาโร) ชาติมรณสํสาโร, นตฺถิ ตสฺส ปุนพฺภโวติ ปารงฺคโต.\csromanspacer{} \textbf{น ปจฺเจตี}ติ * โสตาปตฺติมคฺเคน เย กิเลสา ปหี\textbf{นา}, เต กิเลเส น ปุเนติ น ปจฺเจติ น ปจฺจาคจฺฉติ.\csromanspacer{} สกทาคามิมคฺเคน เย กิเลสา ปหี\textbf{นา}, เต กิเลเส น ปุเนติ น ปจฺเจติ น ปจฺจาคจฺฉติ.\csromanspacer{} อ\textbf{นา}คามิมคฺเคน เย กิเลสา ปหี\textbf{นา}, เต กิเลเส น ปุเนติ น ปจฺเจติ น ปจฺจาคจฺฉติ.\csromanspacer{} อรหตฺตมคฺเคน เย กิเลสา ปหี\textbf{นา}, เต กิเลเส น ปุเนติ น ปจฺเจติ น ปจฺจาคจฺฉตีติ ปารงฺคโต น ปจฺเจติ.\csromanspacer{}\csromanspacer{}\textbf{ตาที}ติ อรหา ปญฺจหากาเรหิ ตาที อิฏฺฐานิฏฺเฐ ตาที, จตฺตาวีติ ตาที, ติณฺณาวีติ ตาที, มุตฺตาวีติ ตาที, ตํนิทฺเทสา ตาที.}

\prose{กถํ อรหา อิฏฺฐานิฏฺเฐ ตาที, อรหา ลาเภปิ ตาที, อลาเภปิ ตาที, ยเสปิ ตาที, อยเสปิ ตาที, ปสํสายปิ ตาที, นินฺทายปิ ตาที, สุเขปิ ตาที, ทุกฺเขปิ ตาที, เอกญฺเจ พาหํ\footnote{องฺคํ (สี)} คนฺเธน ลิมฺเปยฺยุํ,}

\vspace{\tipitakaparskip}
\csromancenter{ปรมฏฺฐกสุตฺตนิทฺเทส เอกญฺเจ พาหํ\footnote{องฺคํ (สี)} วาสิยา ตจฺเฉยฺยุํ, อมุสฺมิํ นตฺถิ ราโค, อมุสฺมิํ นตฺถิ ปฏิฆํ, อนุนยปฏิฆวิปฺปหีโน, อุคฺฆาตินิฆาติวีติวตฺโต, อนุโรธวิโรธสมติกฺกนฺโต, เอวํ อรหา อิฏฺฐานิฏฺเฐ ตาที.}
\prose{\csromanfolio{89}กถํ อรหา จตฺตาวีติ ตาที, อรหโต ราโค จตฺโต วนฺโต มุตฺโต ปหีโน ปฏินิสฺสฏฺโฐ, โทโส.\csromanspacer{} โมโห.\csromanspacer{} โกโธ.\csromanspacer{} อุปนาโห.\csromanspacer{} มกฺโข.\csromanspacer{} ปฬาโส.\csromanspacer{} อิสฺสา.\csromanspacer{} มจฺฉริยํ.\csromanspacer{} มยา.\csromanspacer{} สาเฐยฺยํ.\csromanspacer{} ถมฺโภ.\csromanspacer{} สารมฺโภ.\csromanspacer{} มาโน.\csromanspacer{} อติมาโน.\csromanspacer{} มโท.\csromanspacer{} ปมาโท.\csromanspacer{} สพฺเพ กิเลสา.\csromanspacer{} สพฺเพ ทุจฺจริตา.\csromanspacer{} สพฺเพ ทรถา.\csromanspacer{} สพฺเพ ปริฬาหา.\csromanspacer{} สพฺเพ สนฺตาปา.\csromanspacer{} สพฺพากุสลาภิสงฺขารา จตฺตา วนฺตา มุตฺตา ปหีนา ปฏินิสฺสฏฺฐา, เอวํ อรหา จตฺตาวีติ ตาที.}

\prose{กถํ อรหา ติณฺณาวีติ ตาที, อรหา กาโมฆํ ติณฺโณ ภโวฆํ ติณฺโณ ทิฏฺโฐฆํ ติณฺโณ อวิชฺโชฆํ ติณฺโณ สพฺพํ สํสารปถํ ติณฺโณ อุตฺติณฺโณ นิตฺติณฺโณ อติกฺกนฺโต สมติกฺกนฺโต วีติวตฺโต, โส วุฏฺฐวาโส จิณฺณจรโณ ชาติมรณสํสาโร, นตฺถิ ตสฺส ปุนพฺภโวติ, เอวํ อรหา ติณฺณาวีติ ตาที.}

\prose{กถํ อรหา มุตฺตาวีติ ตาที, อรหโต ราคา จิตฺตํ มุตฺตํ วิมุตฺตํ สุวิมุตฺตํ, โทสา จิตฺตํ มุตฺตํ วิมุตฺตํ สุวิมุตฺตํ, โมหา จิตฺตํ มุตฺตํ วิมุตฺตํ สุวิมุตฺตํ, โกธา.\csromanspacer{} อุปนาหา.\csromanspacer{} มกฺขา.\csromanspacer{} ปฬาสา.\csromanspacer{} อิสฺสาย.\csromanspacer{} มจฺฉริยา.\csromanspacer{} มายาย.\csromanspacer{} สาเฐยฺยา.\csromanspacer{} ถมฺภา.\csromanspacer{} สารมฺภา.\csromanspacer{} มาน.\csromanspacer{} อติมานา.\csromanspacer{} มทา.\csromanspacer{} ปมาทา.\csromanspacer{} สพฺพกิเลเสหิ.\csromanspacer{} สพฺพทุจฺจริเตหิ.\csromanspacer{} สพฺพทรเถหิ.\csromanspacer{} สพฺพปริฬาเหหิ.\csromanspacer{} สพฺพสนฺตาเปหิ.\csromanspacer{} สพฺพากุสลาภิสงฺขาเรหิ จิตฺตํ มุตฺตํ วิมุตฺตํ สุวิมุตฺตํ, เอวํ อรหา มุตฺตาวีติ ตาที.}

\prose{กถํ อรหา ตํนิทฺเทสา ตาที, อรหา สีเล สติ สีลวาติ ตํนิทฺเทสา ตาที, สทฺธาย สติ สทฺโธติ ตํนิทฺเทสา ตาที, วีริเย สติ วีริยวาติ ตํนิทฺเทสา ตาที, สติยา สติ สติมาติ ตํนิทฺเทสา ตาที, สมาธิมฺหิ สติ สมาหิโต ตํนิทฺเทสา ตาที, ปญฺญาย สติ ปญฺญวาติ ตํนิทฺเทสา ตาที, วิชฺชาย สติ เตวิชฺโชติ ตํนิทฺเทสา ตาที, อภิญฺญาย สติ ฉฬภิญฺโญติ ตํนิทฺเทสา ตาที, เอวํ อรหา ตํนิทฺเทสา ตาทีติ ปารงฺคโต น ปจฺเจติ ตาที.\csromanspacer{} เตนาห ภควา\csromandash{}}

\csromangathameasure{น กปฺปนฺติ น ปุเรกฺขโรนฺติ, \\ ธมฺมาปิ เตสํ น ปฏิจฺฉิตาเส\textbf{.} \\ น พฺราหฺมโณ สีลวเตน เนยฺโย, \\ ปารงฺคโต น ปจฺเจติ ตาทีติ\textbf{.}}
\csromangathagroup{\csromangathastanza{\csromanfolio{90}น กปฺปนฺติ น ปุเรกฺขโรนฺติ, \\ ธมฺมาปิ เตสํ น ปฏิจฺฉิตาเส\textbf{.} \\ น พฺราหฺมโณ สีลวเตน เนยฺโย, \\ ปารงฺคโต น ปจฺเจติ ตาทีติ\textbf{.}}}

\csromanheader{2}{ปรมฏฺฐกสุตฺตนิทฺเทโส ปญฺจโม\textbf{.}}
\csromansectionrule
\csromanheader{2}{6. ชราสุตฺตนิทฺเทส}
\csromanheader{2}{อถ ชราสุตฺตนิทฺเทสํ วกฺขติ\csromandash{}}
\proseitem{39}{\csromansymbolfootnote{*}{ขุ 1. 405 ปิฏฺเฐ.}อปฺปํ วต ชีวิตํ อิทํ,}

\prose{\textbf{โอรํ วสฺสสตาปิ มิยฺยติ}\footnote{มียติ (สี)}\textbf{.}}

\csromangathameasure{\textbf{โย เจปิ อติจฺจ ชีวติ,} \\ \textbf{อถ โข โส ชรสาปิ มิยฺยติ1.}}
\csromangathagroup{\csromangathastanza{\textbf{โย เจปิ อติจฺจ ชีวติ,} \\ \textbf{อถ โข โส ชรสาปิ มิยฺยติ1.}}}

\prose{\textbf{อปฺปํ วต ชีวิตํ อิทนฺ}ติ \textbf{ชีวิตนฺ}ติ + อายุฐิติ ยปนา ยาปนา อิริยนา วตฺตนา ปาลนา ชีวิตํ ชีวิตินฺทฺริยํ\textbf{.}\csromanspacer{} อปิ จ ทฺวีหิ การเณหิ อปฺปกํ ชีวิตํ โถกํ ชีวิตํ ฐิติปริตฺตตาย วา อปฺปกํ ชีวิตํ, สรสปริตฺตตาย วา อปฺปกํ ชีวิตํ\textbf{.}\csromanspacer{} กถํ ฐิติปริตฺตตาย วา อปฺปกํ ชีวิตํ, อตีเต จิตฺตกฺขเณ ชีวิตฺถ, น ชีวติ น ชีวิสฺสติ, อนาคเต จิตฺตกฺขเณ ชีวิสฺสติ, น ชีวติ น ชีวิตฺถ, ปจฺจุปฺปนฺเน จิตฺตกฺขเณ ชีวติ, น ชีวิตฺถ น ชีวิสฺสติ\textbf{.}}

\csromangathameasure{ชีวิตํ อตฺตภาโว จ, สุขทุกฺขา จ เกวลา. \\ เอกจิตฺตสมายุตฺตา, ลหุโส วตฺตเต ขโณ. \\ จุลฺลาสีติสหสฺสานิ, กปฺปา ติฏฺฐนฺติ เย มรู. \\ น เตฺวว เตปิ ชีวนฺติ, ทฺวีหิ จิตฺเตหิ สํยุตา. \\ เย นิรุทฺธา มรนฺตสฺส, ติฏฺฐมานสฺส วา อิธ. \\ สพฺเพปิ สทิสา ขนฺธา, คตา อปฺปฏิสนฺธิกา. \\ อนนฺตรา จ เย ภคฺคา, เย จ ภคฺคา อนาคตา. \\ ตทนฺตเร นิรุทฺธานํ, เวสมํ นตฺถิ ลกฺขเณ. \\ อนิพฺพตฺเตน น ชาโต, ปจฺจุปฺปนฺเนน ชีวติ. \\ จิตฺตภคฺคา มโต โลโก, ปญฺญตฺติ ปรมตฺถิยา.}
\csromangathagroup{\csromangathastanza{ชีวิตํ อตฺตภาโว จ, สุขทุกฺขา จ เกวลา. \\ เอกจิตฺตสมายุตฺตา, ลหุโส วตฺตเต ขโณ.}\csromangathastanza{จุลฺลาสีติสหสฺสานิ, กปฺปา ติฏฺฐนฺติ เย มรู. \\ น เตฺวว เตปิ ชีวนฺติ, ทฺวีหิ จิตฺเตหิ สํยุตา\footnote{สโมหิตา (สี, สฺยา, ก)}.}\csromangathastanza{เย นิรุทฺธา มรนฺตสฺส, ติฏฺฐมานสฺส วา อิธ. \\ สพฺเพปิ สทิสา ขนฺธา, คตา อปฺปฏิสนฺธิกา.}\csromangathastanza{อนนฺตรา จ เย ภคฺคา, เย จ ภคฺคา อนาคตา. \\ ตทนฺตเร นิรุทฺธานํ, เวสมํ นตฺถิ ลกฺขเณ.}\csromangathastanza{อนิพฺพตฺเตน น ชาโต, ปจฺจุปฺปนฺเนน ชีวติ. \\ จิตฺตภคฺคา มโต โลโก, ปญฺญตฺติ ปรมตฺถิยา.}}

\csromanheader{2}{6. ชราสุตฺตนิทฺเทส}
\csromangathameasure{ยถา นินฺนา ปวตฺตนฺติ, ฉนฺเทน ปริณามิตา. \\ อจฺฉินฺนธารา วตฺตนฺติ, สฬายตนปจฺจยา. \\ อนิธานคตา ภคฺคา, ปุญฺโช นตฺถิ อนาคเต. \\ นิพฺพตฺตา เย จ ติฏฺฐนฺติ, อารคฺเค สาสปูปมา. \\ นิพฺพตฺตานญฺจ ธมฺมานํ, ภงฺโค เนสํ ปุรกฺขโต. \\ ปโลกธมฺมา ติฏฺฐนฺติ, ปุราเณหิ อมิสฺสิตา. \\ อทสฺสนโต อายนฺติ, ภงฺคา คจฺฉนฺตุทสฺสนํ.}
\csromangathagroup{\csromangathastanza{\csromanfolio{91}ยถา นินฺนา ปวตฺตนฺติ, ฉนฺเทน ปริณามิตา. \\ อจฺฉินฺนธารา วตฺตนฺติ, สฬายตนปจฺจยา.}\csromangathastanza{อนิธานคตา ภคฺคา, ปุญฺโช นตฺถิ อนาคเต. \\ นิพฺพตฺตา เย จ ติฏฺฐนฺติ, อารคฺเค สาสปูปมา.}\csromangathastanza{นิพฺพตฺตานญฺจ ธมฺมานํ, ภงฺโค เนสํ ปุรกฺขโต. \\ ปโลกธมฺมา ติฏฺฐนฺติ, ปุราเณหิ อมิสฺสิตา. \\ อทสฺสนโต อายนฺติ, ภงฺคา คจฺฉนฺตุทสฺสนํ.}}

\vspace{\tipitakaparskip}
\csromancenter{วิชฺชุปฺปาโทว อากาเส, อุปฺปชฺชนฺติ วยนฺติ จาติ\csromandash{}}
\prose{เอวํ ฐิติปริตฺตตาย อปฺปกํ ชีวิตํ.}

\prose{กถํ สรสปริตฺตตาย อปฺปกํ ชีวิตํ\csromandash{}อสฺสาสูปนิพทฺธํ\footnote{อสฺสาสูปนิพนฺธํ (ก)} ชีวิตํ, ปสฺสาสูปนิพทฺธํ ชีวิตํ, อสฺสาสปสฺสาสูปนิพทฺธํ ชีวิตํ, มหาภูตูปนิพทฺธํ ชีวิตํ, กพฬีการาหารูปนิพทฺธํ ชีวิตํ, อุสฺมูปนิพทฺธํ ชีวิตํ, วิญฺญาณูปนิพทฺธํ ชีวิตํ, มูลมฺปิ อิเมสํ ทุพฺพลํ, ปุพฺพเหตูหิ อิเมสํ ทุพฺพลา, เย ปจฺจยา, เตปิ ทุพฺพลา, เยปิ ปภาวิกา, เตปิ ทุพฺพลา, สหภูปิ อิเมสํ ทุพฺพลา, สมฺปโยคาปิ อิเมสํ ทุพฺพลา, สหชาปิ อิเมสํ ทุพฺพลา, ยาปิ ปโยชิกา, สาปิ ทุพฺพลา, อญฺญมญฺญํ อิเม นิจฺจทุพฺพลา, อญฺญมญฺญํ อนวฏฺฐิตา อิเม, อญฺญมญฺญํ ปริปาตยนฺติ อิเม, อญฺญมญฺญสฺส หิ นตฺถิ ตายิตา, น จาปิ ฐเปนฺติ อญฺญมญฺญํ อิเม, โยปิ นิพฺพตฺตโก, โส น วิชฺชติ.}

\csromangathameasure{น จ เกนจิ โกจิ หายติ, \\ คนฺธพฺพา จ อิเม หิ สพฺพโส. \\ ปุริเมหิ ปภาวิกา อิเม, \\ เยปิ ปภาวิกา เต ปุเร มตา. \\ ปุริมาปิ จ ปจฺฉิมาปิ จ,}
\csromangathagroup{\csromangathastanza{น จ เกนจิ โกจิ หายติ, \\ คนฺธพฺพา จ อิเม หิ สพฺพโส. \\ ปุริเมหิ ปภาวิกา อิเม, \\ เยปิ ปภาวิกา เต ปุเร มตา.}\csromangathastanza{ปุริมาปิ จ ปจฺฉิมาปิ จ,}}

\csromanheader{2}{อญฺญมญฺญํ น กทาจิ มทฺทสํสูติ\csromandash{}}
\prose{เอวํ สรสปริตฺตตาย อปฺปกํ ชีวิตํ.}

\prose{อปิ จ จาตุมหาราชิกานํ\footnote{จาตุมฺมหาราชิกานํ (สี, สฺยา) เหฏฺฐา 33 ปิฏฺเฐปิ.} เทวานํ ชีวิตํ อุปาทาย มนุสฺสานํ อปฺปกํ ชีวิตํ ปริตฺตํ ชีวิตํ โถกํ ชีวิตํ ขณิกํ ชีวิตํ ลหุกํ ชีวิตํ อิตฺตรํ}

\prose{\csromanfolio{92}ชีวิตํ อนทฺธนิยํ ชีวิตํ นจิรฏฺฐิติกํ ชีวิตํ.\csromanspacer{} ตาวติํสานํ เทวานํ ฯเปฯ.\csromanspacer{} ยามานํ เทวานํ.\csromanspacer{} ตุสิตานํ เทวานํ.\csromanspacer{} นิมฺมานรตีนํ เทวานํ.\csromanspacer{} ปรนิมฺมิตวสวตฺตีนํ เทวานํ.\csromanspacer{} พฺรหฺมกายิกานํ เทวานํ ชีวิตํ อุปาทาย มนุสฺสานํ อปฺปกํ ชีวิตํ ปริตฺตํ ชีวิตํ โถกํ ชีวิตํ ขณิกํ ชีวิตํ ลหุกํ ชีวิตํ อิตฺตรํ ชีวิตํ อนทฺธนียํ ชีวิตํ นจิรฏฺฐิติกํ ชีวิตํ.\csromanspacer{} วุตฺตํ เหตํ ภควตา\csromandash{} * อปฺปมิทํ ภิกฺขเว มนุสฺสานํ อายุ, คมนิโย สมฺปราโย, มนฺตาย โพทฺธพฺพํ, กตฺตพฺพํ กุสลํ, จริตพฺพํ พฺรหฺมจริยํ, นตฺถิ ชาตสฺส อมรณํ.\csromanspacer{} โย ภิกฺขเว จิรํ ชีวติ, โส วสฺสสตํ, อปฺปํ วา ภิยฺโย.}

\csromangathameasure{*อปฺปมายุ มนุสฺสานํ, หีเฬยฺย นํ สุโปริโส. \\ จเรยฺยาทิตฺตสีโสว, นตฺถิ มจฺจุสฺสนาคโม. \\ อจฺจยนฺติ อโหรตฺตา, ชีวิตํ อุปรุชฺฌติ.}
\csromangathagroup{\csromangathastanza{\csromansymbolfootnote{*}{สํ 1. 109, 110 ปิฏฺเฐสุ.}อปฺปมายุ มนุสฺสานํ, หีเฬยฺย นํ สุโปริโส. \\ จเรยฺยาทิตฺตสีโสว, นตฺถิ มจฺจุสฺสนาคโม. \\ อจฺจยนฺติ อโหรตฺตา, ชีวิตํ อุปรุชฺฌติ.}}

\prose{อายุ ขิยฺยติ มจฺจานํ, กุนฺนทีนํว โอทกนฺติ\csromandash{}}

\prose{อปฺปํ วต ชีวิตํ อิทํ.}

\prose{\textbf{โอรํ วสฺสสตาปิ มิยฺยตี}ติ กลลกาเลปิ จวติ มรติ อนฺตรธายติ วิปฺปลุชฺชติ, อพฺพุทกาเลปิ จวติ มรติ อนฺตรธายติ วิปฺปลุชฺชติ, เปสิกาเลปิ จวติ มรติ อนฺตรธายติ วิปฺปลุชฺชติ, ฆนกาเลปิ จวติ มรติ อนฺตรธายติ วิปฺปลุชฺชติ, ปสาขกาเลปิ จวติ มรติ อนฺตรธายติ วิปฺปลุชฺชติ, ชาตมตฺโตปิ จวติ มรติ อนฺตรธายติ วิปฺปลุชฺชติ.\csromanspacer{} สูติฆเรปิ\footnote{ปสูติฆเร (สฺยา), สูติกฆเร (ก)} จวติ มรติ อนฺตรธายติ วิปฺปลุชฺชติ, อทฺธมาสิโกปิ จวติ มรติ อนฺตรธายติ วิปฺปลุชฺชติ, มาสิโกปิ จวติ มรติ อนฺตรธายติ วิปฺปลุชฺชติ, เทฺวมาสิโกปิ.\csromanspacer{} เตมาสิโกปิ.\csromanspacer{} จตุมาสิโกปิ.\csromanspacer{} ปญฺจมาสิโกปิ จวติ มรติ อนฺตรธายติ วิปฺปลุชฺชติ, ฉมาสิโกปิ.\csromanspacer{} สตฺตมาสิโกปิ.\csromanspacer{} อฏฺฐมาสิโกปิ.\csromanspacer{} นวมาสิโกปิ.\csromanspacer{} ทสมาสิโกปิ.\csromanspacer{} สํวจฺฉริโกปิ จวติ มรติ อนฺตรธายติ วิปฺปลุชฺชติ, เทฺววสฺสิโกปิ.\csromanspacer{} ติวสฺสิโกปิ.\csromanspacer{} จตุวสฺสิโกปิ.\csromanspacer{} ปญฺจวสฺสิโกปิ.\csromanspacer{} ฉวสฺสิโกปิ.\csromanspacer{} สตฺตวสฺสิโกปิ.\csromanspacer{} อฏฺฐวสฺสิโกปิ.\csromanspacer{} นววสฺสิโกปิ.\csromanspacer{} ทสวสฺสิโกปิ.\csromanspacer{} วีสติวสฺสิโกปิ.\csromanspacer{} ติํสวสฺสิโกปิ.\csromanspacer{} จตฺตารีสวสฺสิโกปิ.\csromanspacer{} ปญฺญาสวสฺสิโกปิ.\csromanspacer{} สฏฺฐิวสฺสิโกปิ.\csromanspacer{} สตฺตติวสฺสิโกปิ.\csromanspacer{} อสีติวสฺสิโกปิ.\csromanspacer{} น วุติวสฺสิโกปิ จวติ มรติ อนฺตรธายติ วิปฺปลุชฺชตีติ โอรํ วสฺสสตาปิ มิยฺยติ.}

\csromanheader{2}{6. ชราสุตฺตนิทฺเทส}
\prose{\csromanfolio{93}\textbf{โย เจปิ อติจฺจ ชีวตี}ติ โย วสฺสสตํ อติกฺกมิตฺวา ชีวติ, โส เอกํ วา วสฺสํ ชีวติ.\csromanspacer{} เทฺว วา วสฺสานิ ชีวติ.\csromanspacer{} ตีณิ วา วสฺสานิ ชีวติ.\csromanspacer{} จตฺตาริ วา วสฺสานิ ชีวติ.\csromanspacer{} ปญฺจ วา วสฺสานิ ชีวติ ฯเปฯ.\csromanspacer{} ทส วา วสฺสานิ ชีวติ.\csromanspacer{} วีสติ วา วสฺสานิ ชีวติ.\csromanspacer{} ติํสํ วา วสฺสานิ ชีวติ.\csromanspacer{} จตฺตารีสํ วา วสฺสานิ ชีวตีติ โส เจปิ อติจฺจ ชีวติ.\csromanspacer{} \textbf{อถ โข โส ชรสาปิ มิยฺยตี}ติ ยทา ชิณฺโณ โหติ วุทฺโธ มหลฺลโก อทฺธคโต วโย-อนุปฺปตฺโต ขณฺฑทนฺโต ปลิตเกโส วิลูนํ ขลิตสิโร\footnote{ขลิตํ สิโร (สี)} วลินํ ติลกาหตคตฺโต วงฺโก โภคฺโค ทณฺฑปรายโณ โส ชรายปิ จวติ มรติ อนฺตรธายติ วิปฺปลุชฺชติ, นตฺถิ มรณมฺหา โมกฺโข.}

\csromangathameasure{ผลานมิว ปกฺกานํ, ปาโต ปตนโต ภยํ. \\ เอวํ ชาตาน มจฺจานํ, นิจฺจํ มรณโต ภยํ. \\ ยถาปิ กุพฺภการสฺส, กตํ มตฺติกภาชนํ. \\ สพฺพํ เภทนปริยนฺตํ, เอวํ มจฺจาน ชีวิตํ. \\ ทหรา จ มหนฺตา จ, เย พาลา เย จ ปณฺฑิตา. \\ สพฺเพ มจฺจุวสํ ยนฺติ, สพฺเพ มจฺจุปรายณา. \\ เตสํ มจฺจุปเรตานํ, คจฺฉตํ ปรโลกโต. \\ น ปิตา ตายเต ปุตฺตํ, ญาตี วา ปน ญาตเก. \\ เปกฺขตญฺเญว ญาตีนํ, ปสฺส ลาลปฺปตํ ปุถุ. \\ เอกเมโกว มจฺจานํ, โควชฺโฌ วิย นิยฺยติ. \\ เอวมพฺภาหโต โลโก, มจฺจุนา จ ชราย จาติ.}
\csromangathagroup{\csromangathastanza{ผลานมิว ปกฺกานํ, ปาโต ปตนโต\footnote{ปปตโต (สี) ขุ 1. 371 ปิฏฺเฐ.} ภยํ. \\ เอวํ ชาตาน มจฺจานํ, นิจฺจํ มรณโต ภยํ.}\csromangathastanza{ยถาปิ กุพฺภการสฺส, กตํ มตฺติกภาชนํ. \\ สพฺพํ เภทนปริยนฺตํ, เอวํ มจฺจาน ชีวิตํ.}\csromangathastanza{ทหรา จ มหนฺตา จ, เย พาลา เย จ ปณฺฑิตา. \\ สพฺเพ มจฺจุวสํ ยนฺติ, สพฺเพ มจฺจุปรายณา.}\csromangathastanza{เตสํ มจฺจุปเรตานํ, คจฺฉตํ ปรโลกโต. \\ น ปิตา ตายเต ปุตฺตํ, ญาตี วา ปน ญาตเก.}\csromangathastanza{เปกฺขตญฺเญว ญาตีนํ, ปสฺส ลาลปฺปตํ ปุถุ. \\ เอกเมโกว มจฺจานํ, โควชฺโฌ วิย นิยฺยติ. \\ เอวมพฺภาหโต โลโก, มจฺจุนา จ ชราย จาติ.}}

\csromanheader{2}{อถ โข โส ชรสาปิ มิยฺยติ. เตนาห ภควา\csromandash{}}
\csromangathameasure{อปฺปํ วต ชีวิตํ อิทํ, \\ โอรํ วสฺสสตาปิ มิยฺยติ. \\ โย เจปิ อติจฺจ ชีวติ,}
\csromangathagroup{\csromangathastanza{อปฺปํ วต ชีวิตํ อิทํ, \\ โอรํ วสฺสสตาปิ มิยฺยติ. \\ โย เจปิ อติจฺจ ชีวติ,}}

\prose{\textbf{อถ โข โส ชรสาปิ มิยฺยตี}ติ.}

\proseitem{40}{\csromanfolio{94}\csromansymbolfootnote{*}{ขุ 1. 405 ปิฏฺเฐ.}\textbf{โสจนฺติ ชนา มมายิเต},}

\prose{\textbf{ญ หิ สนฺติ นิจฺจา ปริคฺคหา.}}

\csromangathameasure{\textbf{วินาภาวํ สนฺตเมวิทํ,} \\ \textbf{อิติ ทิสฺวา นาคารมาวเส.}}
\csromangathagroup{\csromangathastanza{\textbf{วินาภาวํ สนฺตเมวิทํ,} \\ \textbf{อิติ ทิสฺวา นาคารมาวเส.}}}

\prose{\textbf{โสจนฺติ ชนา มมายิเต}ติ \textbf{ชนา}ติ ขตฺติยา จ พฺราหฺมณา จ เวสฺสา จ สุทฺทา จ คหฏฺฐา จ ปพฺพชิตา จ เทวา จ มนุสฺสา จ.\csromanspacer{} \textbf{มมตฺตา}ติ เทฺว มมตฺตา ตณฺหามมตฺตญฺจ ทิฏฺฐิมมตฺตญฺจ ฯเปฯ อิทํ ตณฺหามมตฺตํ ฯเปฯ อิทํ ทิฏฺฐิมมตฺตํ.\csromanspacer{} มมายิตํ วตฺถุํ อจฺเฉทสงฺกิโนปิ โสจนฺติ, อจฺฉิชฺชนฺเตปิ โสจนฺติ, อจฺฉินฺเนปิ โสจนฺติ.\csromanspacer{} มมายิตํ วตฺถุํ วิปริณามสงฺกิโนปิ โสจนฺติ, วิปริณามนฺเตปิ โสจนฺติ, วิปริณเตปิ โสจนฺติ กิลมนฺติ ปริเทวนฺติ อุรตฺตาฬิํ กนฺทนฺติ สมฺโมหํ อาปชฺชนฺตีติ โสจนฺติ \textbf{ชนา} มมายิเต.}

\prose{\textbf{น หิ สนฺติ นิจฺจา ปริคฺคหา}ติ เทฺว ปริคฺคหา ตณฺหาปริคฺคโห จ ทิฏฺฐิปริคฺคโห จ ฯเปฯ อยํ ตณฺหาปริคฺคโห ฯเปฯ อยํ ทิฏฺฐิปริคฺคโห.\csromanspacer{} ตณฺหาปริคฺคโห อนิจฺโจ สงฺขโต ปฏิจฺจสมุปฺปนฺโน ขยธมฺโม วยธมฺโม วิราคธมฺโม นิโรธธมฺโม วิปริณามธมฺโม.\csromanspacer{} ทิฏฺฐิปริคฺคโหปิ อนิจฺโจ สงฺขโต ปฏิจฺจสมุปฺปนฺโน ขยธมฺโม วยธมฺโม วิราคธมฺโม นิโรธธมฺโม วิปริณามธมฺโม.\csromanspacer{} วุตฺตํ เหตํ ภควตา\csromandash{} + ปสฺสถ โน ตุเมฺห ภิกฺขเว ตํ ปริคฺคหํ, ยฺวายํ ปริคฺคโห นิจฺโจ ธุโว สสฺสโต อวิปริณามธมฺโม, สสฺสติสมํ ตเถว ติฏฺเฐยฺยาติ.\csromanspacer{} โน เหตํ ภนฺเต.\csromanspacer{} สาธุ ภิกฺขเว อหมฺปิ โข ตํ ภิกฺขเว ปริคฺคหํ น สมนุปสฺสามิ, ยฺวายํ ปริคฺคโห นิจฺโจ ธุโว สสฺสโต อวิปริณามธมฺโม, สสฺสติสมํ ตเถว ติฏฺเฐยฺยาติ.\csromanspacer{} ปริคฺคหา นิจฺจา ธุวา สสฺสตา อวิปริณามธมฺมา นตฺถิ น สนฺติ น สํวิชฺชนฺติ น ลพฺภนฺตีติ น หิ สนฺติ นิจฺจา ปริคฺคหา.}

\prose{\textbf{วินาภาวํ สนฺตเมวิทนฺ}ติ นานาภาเว วินาภาเว อญฺญถาภาเว สนฺเต สํวิชฺชมาเน อุปลพฺภิยมาเน.\csromanspacer{} วุตฺตํ เหตํ ภควตา\csromandash{}อลํ อานนฺท มา โสจิ มา ปริเทวิ, นนุ เอตํ อานนฺท มยา ปฏิกจฺเจว\footnote{ปฏิคจฺเจว (สี) ที 2. 119 ปิฏฺเฐปิ.} อกฺขาตํ “สพฺเพเหว ปิเยหิ มนาเปหิ นานาภาโว วินาภาโว อญฺญถาภาโว, ตํ กุเตตฺถ อานนฺท ลพฺภา.\csromanspacer{} ยํ ตํ}

\vspace{\tipitakaparskip}
\csromancenter{ชราสุตฺตนิทฺเทส ชาตํ ภูตํ สงฺขตํ ปโลกธมฺมํ\textbf{,} ตํ วต มา ปลุชฺชี”ติ\textbf{,} เนตํ ฐานํ วิชฺชติ.\csromanspacer{} ปุริมานํ ปุริมานํ ขนฺธานํ ธาตูนํ อายตนานํ วิปริณามญฺญถาภาวา ปจฺฉิมา ปจฺฉิมา ขนฺธา จ ธาตุโย จ อายตนานิ จ ปวตฺตนฺตีติ วินาภาวํ สนฺตเมวิทํ.}
\prose{\csromanfolio{95}\textbf{อิติ ทิสฺวา นาคารมาวเส}ติ \textbf{อิตี}ติ ปทสนฺทิ ปทสํสคฺโค ปทปาริปูรี อกฺขรสมวาโย พฺยญฺชนสิลิฏฺฐตา ปทานุปุพฺพตา เมตํ \textbf{อิตี}ติ.\csromanspacer{} \textbf{อิติ ทิสฺวา} ปสฺสิตฺวา ตุลยิตฺวา ตีรยิตฺวา วิภาวยิตฺวา วิภูตํ กตฺวา มมตฺเตสูติ อิติ ทิสฺวา.\csromanspacer{} \textbf{นาคารมาวเส}ติ สพฺพํ ฆราวาสปลิโพธํ ฉินฺทิตฺวา ปุตฺตทารปลิโพธํ ฉินฺทิตฺวา ญาติปลิโพธํ ฉินฺทิตฺวา มิตฺตามจฺจปลิโพธํ ฉินฺทิตฺวา สนฺนิธิปลิโพธํ ฉินฺทิตฺวา เกสมสฺสุํ โอหาเรตฺวา กาสายานิ วตฺตานิ อจฺฉาเทตฺวา อคารสฺมา อนคาริยํ ปพฺพชิตฺวา อกิญฺจนภาวํ อุปคนฺตฺวา เอโก จเรยฺย วิหเรยฺย อิริเยยฺย วตฺเตยฺย ปาเลยฺย ยเปยฺย ยาเปยฺยาติ อิติ ทิสฺวา นาคารมาวเส.\csromanspacer{} เตนาห ภควา\csromandash{}}

\csromangathameasure{โสจนฺติ ชนา มมายิเต\textbf{,} \\ น หิ สนฺติ นิจฺจา ปริคฺคหา. \\ วินาภาวํ สนฺตเมวิทํ\textbf{,} \\ \textbf{อิติ ทิสฺวา นาคารมาวเส}ติ.}
\csromangathagroup{\csromangathastanza{โสจนฺติ ชนา มมายิเต\textbf{,} \\ น หิ สนฺติ นิจฺจา ปริคฺคหา. \\ วินาภาวํ สนฺตเมวิทํ\textbf{,} \\ \textbf{อิติ ทิสฺวา นาคารมาวเส}ติ.}}

\proseitem{41}{\csromansymbolfootnote{*}{ขุ 1. 405 ปิฏฺเฐ.}มรเณนปิ ตํ ปหียติ\footnote{ปหิยฺยติ (ก)}\textbf{,}}

\prose{\textbf{ยํ ปุริโส มมิทนฺติ มญฺญติ.}}

\csromangathameasure{\textbf{เอตมฺปิ วิทิตฺวาน} \textbf{ปณฺฑิโต,} \\ \textbf{น มมตฺตาย นเมถ มามโก.}}
\csromangathagroup{\csromangathastanza{\textbf{เอตมฺปิ วิทิตฺวาน}\footnote{เอตํ ทิสฺวาน (สี, ก)} \textbf{ปณฺฑิโต,} \\ \textbf{น มมตฺตาย นเมถ มามโก.}}}

\prose{\textbf{มรเณนปิ ตํ ปหียตี}ติ \textbf{มรณนฺ}ติ + ยา เตสํ เตสํ สตฺตานํ ตมฺหา ตมฺหา สตฺตนิกายา จุติ จวนตา เภโท อนฺตรธานํ มจฺจุมรณํ กาลํกิริยา ขนฺธานํ เภโท กเฬวรสฺส นิกฺเขโป ชีวิตินฺทฺริยสฺสุปจฺเฉโท.\csromanspacer{} \textbf{ตนฺ}ติ รูปคตํ เวทนาคตํ สญฺญาคตํ สงฺขารคตํ วิญฺญาณคตํ.\csromanspacer{} \textbf{ปหียตี}ติ ปหียติ ชหียติ วิชหียติ อนฺตรธายติ วิปฺปลุชฺชติ\textbf{,} ภาสิตมฺปิ เหตํ\csromandash{}}

\csromangathameasure{*ปุพฺเพว มจฺจํ วิชหนฺติ โภคา,}
\csromangathagroup{\csromangathastanza{\csromanfolio{96}\csromansymbolfootnote{*}{ขุ 5. 120 ปิฏฺเฐ ชาตเก.}ปุพฺเพว มจฺจํ วิชหนฺติ โภคา,}}

\vspace{\tipitakaparskip}
\csromancenter{มจฺโจ ว เน ปุพฺพตรํ ชหาติ.}
\prose{อสสฺสตา โภคิโน กามกามี,}

\prose{ตฺสฺมา น โสจามหํ โสกกาเล.}

\prose{อุเทติ อาปูรติ เวติ จนฺโท,}

\prose{อตฺตํ ตเปตฺวาน ปเลติ สูริโย.}

\vspace{\tipitakaparskip}
\csromancenter{วิทิตา มยา สตฺตุก โลกธมฺมา,}
\csromancenter{ตสฺมา น โสจามหํ โสกกาเลติ.}
\prose{\textbf{มรเณนปิ ตํ ปหียติ, ยํ ปุริโส มมิทนฺติ มญฺญตี}ติ \textbf{ยนฺติ} รูปคตํ เวทนาคตํ สญฺญาคตํ สงฺขารคตํ วิญฺญาณคตํ.\csromanspacer{} \textbf{ปุริโส}ติ สงฺขา สมญฺญา ปญฺญตฺติ โวหาโร\footnote{โลกโวหาโร (สฺยา) อภิ 1. 256 ปิฏฺเฐปิ.} นามํ นามกมฺมํ นามเธยฺยํ นิรุตฺติ พฺยญฺชนํ อภิลาโป.\csromanspacer{} \textbf{มมิทนฺติ มญฺญตี}ติ ตณฺหามญฺญนาย มญฺญติ, ทิฏฺฐิมญฺญนาย มญฺญติ, มานมญฺญนาย มญฺญติ, กิเลสมญฺญนาย มญฺญติ, ทุจฺจริตมญฺญนาย มญฺญติ, ปโยคมญฺญนาย มญฺญติ, วิปากมญฺญนาย มญฺญตีติ ยํ ปุริโส มมิทนฺติ มญฺญติ.}

\prose{\textbf{เอตมฺปิ วิทิตฺวาน ปณฺฑิโต}ติ เอตํ อาทีนวํ ญตฺวา ชานิตฺวา ตุลยิตฺวา ตีรยิตฺวา วิภาวยิตฺวา วิภูตํ กตฺวา มมตฺเตสูติ เอตมฺปิ วิทิตฺวา.\csromanspacer{} \textbf{ปณฺฑิโต}ภิ ธีโร ปณฺฑิโต ปญฺญวา พุทฺธิมา ญาณี วิภาวี เมธาวีติ เอตมฺปิ วิทิตฺวาน ปณฺฑิโต.}

\prose{\textbf{น มมตฺตาย นเมถ มามโก}ติ \textbf{มมตฺตา}ติ เทฺว \textbf{มมตฺตา} ตณฺหามมตฺตญฺจ ทิฏฺฐิมมตฺตญฺจ ฯเปฯ อิทํ ตณฺหามมตฺตํ ฯเปฯ อิทํ ทิฏฺฐิมมตฺตํ.\csromanspacer{} \textbf{มามโก}ติ พุทฺธมามโก ธมฺมมามโก สํฆมามโก.\csromanspacer{} โส ภควนฺตํ มมายติ, ภควา ตํ ปุคฺคลํ ปริคฺคณฺหาติ.\csromanspacer{} วุตฺตํ เหตํ ภควตา\csromandash{}เย เต ภิกฺขเว ภิกฺขุ กุหา ถทฺธา\footnote{พทฺธา (ก) ขุ 1. 270 ปิฏฺเฐ.} ลปา สิงฺคี อุนฺนฬา อสมา หิตา, น เม เต ภิกฺขเว ภิกฺขู มามกา.\csromanspacer{} อปคตา จ เต ภิกฺขเว ภิกฺขู อิมสฺมา ธมฺมวินยา, น จ เต อิมสฺมิํ ธมฺมวินเย วุทฺธิํ วิรูฬฺหิํ เวปุลฺลํ อาปชฺชนฺติ.\csromanspacer{} เย จ โข เต ภิกฺขเว ภิกฺขู นิกฺกุหา นิลฺลปา ธีรา อตฺถทฺธา สุสมาหิตา, เต โข เม ภิกฺขเว ภิกฺขู มามกา, อนปคตา จ เต ภิกฺขเว ภิกฺขู อิมสฺมา ธมฺมวินยา, เต จ อิมสฺมิํ ธมฺมวินเย วุทฺธิํ วิรูฬฺหิํ เวปุลฺลํ อาปชฺชนฺติ.}

\csromanheader{2}{6. ชราสุตฺตนิทฺเทส}
\csromangathameasure{กุหา ถทฺธา ลปา สิงฺคี, อุนฺนฬา อสมาหิตา. \\ น เต ธมฺเม วิรูหนฺติ, สมฺมาสมฺพุทฺธเทสิเต. \\ นิกฺกุหา นิลฺลปา ธีรา, อตฺถทฺธา สุสมาหิตา. \\ เต เว ธมฺเม วิรูหนฺติ, สมฺมาสมฺพุทฺธเทสิเต.}
\csromangathagroup{\csromangathastanza{\csromanfolio{97}กุหา ถทฺธา ลปา สิงฺคี, อุนฺนฬา อสมาหิตา. \\ น เต ธมฺเม วิรูหนฺติ, สมฺมาสมฺพุทฺธเทสิเต.}\csromangathastanza{นิกฺกุหา นิลฺลปา ธีรา, อตฺถทฺธา สุสมาหิตา. \\ เต เว ธมฺเม วิรูหนฺติ, สมฺมาสมฺพุทฺธเทสิเต.}}

\prose{\textbf{น มมตฺตาย นเมถ มามโก}ติ มามโก ตณฺหามมตฺตํ ปหาย ทิฏฺฐิมมตฺตํ ปฏินิสฺสชฺชิตฺวา มมตฺตาย น นเมยฺย น โอนเมยฺย, น ตํนินฺโน อสฺส, น ตปฺโปโณ น ตปฺปพฺภาโร น ตทธิมุตฺโต น ตทธิปเตยฺโยติ น มมตฺตาย นเมถ มามโก, เตนาห ภควา\csromandash{}}

\csromangathameasure{มรเณนปิ ตํ ปหียติ, \\ ยํ ปุริโส มมิทนฺติ มญฺญนฺติ. \\ เอตมฺปิ วิทิตฺวาน ปณฺฑิโต,}
\csromangathagroup{\csromangathastanza{มรเณนปิ ตํ ปหียติ, \\ ยํ ปุริโส มมิทนฺติ มญฺญนฺติ. \\ เอตมฺปิ วิทิตฺวาน ปณฺฑิโต,}}

\prose{\textbf{น มมตฺตาย นเมถ มามโก}ติ.}

\proseitem{42}{\csromansymbolfootnote{*}{ขุ 1. 405; ขุ 10. 173 ปิฏฺเฐสุ.}สุปิเนน ยถาปิ สงฺคตํ,}

\prose{\textbf{ปฏิพุทฺโธ ปุริโส น ปสฺสติ.}}

\csromangathameasure{\textbf{เอวมฺปิ ปิยายิตํ ชนํ,} \\ \textbf{เปตํ กาลงฺกตํ} \textbf{น ปสฺสติ.}}
\csromangathagroup{\csromangathastanza{\textbf{เอวมฺปิ ปิยายิตํ ชนํ,} \\ \textbf{เปตํ กาลงฺกตํ}\footnote{กาลกตํ (สี, สฺยา)} \textbf{น ปสฺสติ.}}}

\prose{\textbf{สุปิเนน ยถาปิ สงฺคตนฺ}ติ สงฺคตํ สมาคตํ สมาหิตํ สนฺนิปติตนฺติ สุปิเนน ยถาปิ สงฺคตํ.\csromanspacer{}\csromanspacer{}\textbf{ปฏิพุทฺโธ ปุริโส น ปสฺสตี}ติ ยถา ปุริโส สุปินคโต จนฺทํ ปสฺสติ, สูริยํ ปสฺสติ, มหาสมุทฺทํ ปสฺสติ, สิเนรุํ ปพฺพตราชานํ ปสฺสติ, หตฺถิํ ปสฺสติ, อสฺสํ ปสฺสติ, รถํ ปสฺสติ, ปตฺติํ ปสฺสติ, เสนาพฺยูหํ ปสฺสติ, อารามรามเณยฺยกํ ปสฺสติ, วนรามเณยฺยกํ.\csromanspacer{} ภูมิรามเณยฺยกํ.\csromanspacer{} โปกฺขรณีรามเณยฺยกํ ปสฺสติ.\csromanspacer{} ปฏิพุทฺโธ น กิญฺจิ ปสฺสตีติ ปฏิพุทฺโธ ปุริโส \textbf{น ปสฺสติ.}}

\prose{\textbf{เอวมฺปิ ปิยายิตํ ชนนฺ}ติ \textbf{เอวนฺ}ติ โอปมฺมสมฺปฏิปาทนํ.\csromanspacer{} \textbf{ปิยายิตํ ชนนฺ}ติ มมายิตํ ชนํ มาตรํ วา ปิตรํ วา ภาตรํ วา ภคินิํ วา ปุตฺตํ วา ธีตรํ วา มิตฺตํ วา อมจฺจํ วา ญาติํ วา สาโลหิตํ วาติ เอวมฺปิ ปิยายิตํ ชนํ.}

\prose{\csromanfolio{98}\textbf{เปตํ กาลงฺกตํ น ปสฺสตี}ติ \textbf{เปโต} วุจฺจติ มโต.\csromanspacer{} กาลงฺกตํ น ปสฺสติ น ทกฺขติ นาธิคจฺฉติ น วินฺทติ น ปฏิลภตีติ เปตํ กาลงฺกตํ น ปสฺสติ.\csromanspacer{} เตนาห ภควา\csromandash{}}

\csromangathameasure{สุปิเนน ยถาปิ สงฺคตํ, \\ ปฏิพุทฺโธ ปุริโส น ปสฺสติ. \\ เอวมฺปิ ปิยายิตํ ชนํ,}
\csromangathagroup{\csromangathastanza{สุปิเนน ยถาปิ สงฺคตํ, \\ ปฏิพุทฺโธ ปุริโส น ปสฺสติ. \\ เอวมฺปิ ปิยายิตํ ชนํ,}}

\prose{\textbf{เปตํ กาลงฺกตํ น ปสฺสตี}ติ.}

\proseitem{43}{\csromansymbolfootnote{*}{ขุ 1. 405 ปิฏฺเฐ.}\textbf{ทิฏฺฐาปิ สุตาปิ เต ชนา},}

\prose{\textbf{เยสํ นามมิทํ ปวุจฺจติ.}}

\csromangathameasure{\textbf{นามํเยวา’วสิสฺสติ}, \\ \textbf{อกฺเขยฺยํ เปตสฺส ชนฺตุโน.}}
\csromangathagroup{\csromangathastanza{\textbf{นามํเยวา’วสิสฺสติ}\footnote{นามเมวา’วสิสฺสติ (สี, สฺยา)}, \\ \textbf{อกฺเขยฺยํ เปตสฺส ชนฺตุโน.}}}

\prose{\textbf{ทิฏฺฐาปิ สุตาปิ เต ชนา}ติ \textbf{ทิฏฺฐา}ติ จกฺขุวิญฺญาณาภิสมฺภูตา.\csromanspacer{} \textbf{สุตา}ติ เย โสตวิญฺญาณาภิสมฺภูตา.\csromanspacer{} \textbf{เต ชนา}ติ ขตฺติยา จ พฺราหฺมณา จ เวสฺสา จ สุทฺทา จ คหฏฺฐา จ ปพฺพชิตา จ เทวา จ มนุสฺสา จาติ \textbf{ทิฏฺฐา}ปิ สุตาปิ เต ชนา.}

\prose{\textbf{เยสํ นามมิทํ ปวุจฺจตี}ติ \textbf{เยสนฺ}ติ เยสํ ขตฺติยานํ พฺราหฺมณานํ เวสฺสานํ สุทฺทานํ คหฏฺฐานํ ปพฺพชิตานํ เทวานํ มนุสฺสานํ.\csromanspacer{} \textbf{นามนฺ}ติ สงฺขา สมญฺญา ปญฺญตฺติ โวหาโร นามํ นามกมฺมํ นามเธยฺยํ นิรุตฺติ พฺยญฺชนํ อภิลาโป.\csromanspacer{} \textbf{ปวุจฺจตี}ติ วุจฺจติ ปวุจฺจติ กถียติ ภณียติ ทีปียติ โวหรียตีติ เยสํ นามมิทํ ปวุจฺจติ.}

\prose{\textbf{นามํเยวา’วสิสฺสติ อกฺเขยฺยนฺ}ติ รูปคตํ เวทนาคตํ สญฺญาคตํ สงฺขารคตํ วิญฺญาณคตํ ปหียติ ชหียติ วิชหียติ อนฺตรธายติ วิปฺปลุชฺชติ นามํเยวา’วสิสฺสติ.\csromanspacer{} \textbf{อกฺเขยฺยนฺ}ติ อกฺขาตุํ กเถตุํ ภณิตุํ ทีปยิตุํ โวหริตุนฺติ นามํ เอวาวสิสฺสติ อกฺเขยฺยํ.\csromanspacer{} \textbf{เปตสฺส ชนฺตุโน}ติ \textbf{เปตสฺสา}ติ มตสฺส กาลงฺกตสฺส.\csromanspacer{} \textbf{ชนฺตุโน}ติ สตฺตสฺส นรสฺส มานวสฺส โปสสฺส ปุคฺคลสฺส ชีวสฺส ชาคุสฺส ชนฺตุสฺส อินฺทคุสฺส มนุชสฺสาติ เปตสฺส ชนฺตุโน.\csromanspacer{} เตนาห ภควา\csromandash{}}

\csromanheader{2}{6. ชราสุตฺตนิทฺเทส}
\csromangathameasure{ทิฏฺฐาปิ สุตาปิ เต ชนา, \\ เยสํ นามมิทํ ปวุจฺจติ. \\ นามํเยวา’วสิสฺสภิ, \\ อกฺเขยฺยํ เปตสฺส ชนฺตุโนติ.}
\csromangathagroup{\csromangathastanza{\csromanfolio{99}ทิฏฺฐาปิ สุตาปิ เต ชนา, \\ เยสํ นามมิทํ ปวุจฺจติ. \\ นามํเยวา’วสิสฺสภิ, \\ อกฺเขยฺยํ เปตสฺส ชนฺตุโนติ.}}

\proseitem{44}{\csromansymbolfootnote{*}{ขุ 1. 405 ปิฏฺเฐ.}โสกปฺปริเทวมจฺฉรํ,}

\prose{\textbf{น ปชหนฺติ คิธา มมายิเต.}}

\csromangathameasure{\textbf{ตสฺมา มุนโย ปริคฺคหํ,} \\ \textbf{หิตฺวา อจริํสุ เขมทสฺสิโน.}}
\csromangathagroup{\csromangathastanza{\textbf{ตสฺมา มุนโย ปริคฺคหํ,} \\ \textbf{หิตฺวา อจริํสุ เขมทสฺสิโน.}}}

\prose{\textbf{โสกปฺปริเทวมจฺฉรํ, น ปชหนฺติ คิทฺธา มมายิเต}ติ \textbf{โสโก}ติ ญาติพฺยสเนน วา ผุฏฺฐสฺส โภคพฺยสเนน วา ผุฏฺฐสฺส โรคพฺยสเนน วา ผุฏฺฐสฺส สีลพฺยสเนน วา ผุฏฺฐสฺส ทิฏฺฐิพฺยสเนน วา ผุฏฺฐสฺส อญฺญตรญฺญตเรน พฺยสเนน สมนฺนาคตสฺส อญฺญตรญฺญตเรน ทุกฺขธมฺเมน ผุฏฺฐสฺส \textbf{โสโก} โสจนา โสจิตตฺตํ อนฺโต\textbf{โสโก} อนฺโตปริ\textbf{โสโก} อนฺโตทาโห อนฺโตปริทาโห\footnote{อนฺโตฑาโห อนฺโตปริฑาโห (สฺยา)} เจตโส ปริชฺฌายนา โทมนสฺสํ โสกสลฺลํ.\csromanspacer{} \textbf{ปริเทโว}ติ ญาติพฺยสเนน วา ผุฏฺฐสฺส ฯเปฯ ทิฏฺฐิพฺยสเนน วา ผุฏฺฐสฺส อญฺญตรญฺญตเรน พฺยสเนน สมนฺนาคตสฺส อญฺญตรญฺญตเรน ทุกฺขธมฺเมน ผุฏฺฐสฺส อาเทโว ปริเทโว อาเทวนา ปริเทวนา อาเทวิตตฺตํ ปริเทวิตตฺตํ วาจา ปลาโป วิปฺปลาโป ลาลปฺโป ลาลปฺปายนา ลาลปฺปายิตตฺตํ.\csromanspacer{}\csromanspacer{}\textbf{มจฺฉริยนฺ}ติ ปญฺจ มจฺฉริยานิ อาวาสมจฺฉริยํ กุลมจฺฉริยํ ลาภมจฺฉริยํ วณฺณมจฺฉริยํ ธมฺมมจฺฉริยํ.\csromanspacer{} ยํ เอวรูปํ มจฺฉริยํ มจฺฉรายนา มจฺฉรายิตตฺตํ เววิจฺฉํ กทริยํ กฏุกญฺจุกตา อคฺคหิตตฺตํ จิตฺตสฺส.\csromanspacer{} อิทํ วุจฺจติ มจฺฉริยํ.\csromanspacer{}\csromanspacer{}อปิ จ ขนฺธมจฺฉริยมฺปิ มจฺฉริยํ, ธาตุมจฺฉริยมฺปิ มจฺฉริยํ, อายตนมจฺฉริยมฺปิ มจฺฉริยํ คาโห, อิทํ วุจฺจติ มจฺฉริยํ.\csromanspacer{} เคโธ วุจฺจติ ตณฺหา.\csromanspacer{} โย ราโค สาราโค ฯเปฯ อภิชฺฌา โลโภ อกุสลมูลํ.\csromanspacer{} \textbf{มมตฺตา}ติ เทฺว มมตฺตา ตณฺหามมตฺตญฺจ ทิฏฺฐิมมตฺตญฺจ ฯเปฯ อิทํ ตณฺหามมตฺตํ ฯเปฯ อิทํ ทิฏฺฐิมมตฺตํ.\csromanspacer{} มมายิตํ วตฺถุํ อจฺเฉทสงฺกิโนปิ โสจนฺติ, อจฺฉิชฺชนฺเตปิ โสจนฺติ, อจฺฉินฺเนปิ โสจนฺติ, มมายิตํ วตฺถุํ วิปริณามสงฺกิโนปิ โสจนฺติ, วิปริณามนฺเตปิ โสจนฺติ, วิปริณเตปิ โสจนฺติ, มมายิตํ วตฺถุํ อจฺเฉทสงฺกิโนปิ ปริเทวนฺติ, \csromanfolio{100}อจฺฉิชฺชนฺเตปิ ปริเทวนฺติ, อจฺฉินฺเนปิ ปริเทวนฺติ.\csromanspacer{} มมายิตํ วตฺถุํ วิปริณามสงฺกิโนปิ ปริเทวนฺติ, วิปริณามนฺเตปิ ปริเทวนฺติ, วิปริณเตปิ ปริเทวนฺติ, มมายิตํ วตฺถุํ รกฺขนฺติ โคเปนฺติ ปริคฺคณฺหนฺติ มมายนฺติ มจฺฉรายนฺติ, มมายิตสฺมิํ วตฺถุสฺมิํ โสกํ น ชหนฺติ, ปริเทวํ น ชหนฺติ, มจฺฉริยํ น ชหนฺติ, เคธํ น ชหนฺติ นปฺปชหนฺติ น วิโนเทนฺติ น พฺยนฺติํ กโรนฺติ น อนภาวํ คเมนฺตีติ โสกปฺปริเทวมจฺฉรํ, นปฺปชหนฺติ คิทฺธา มมายิเต.}

\prose{\textbf{ตสฺมา มุนโย ปริคฺคหํ, หิตฺวา อจริํสุ เขมทสฺสิโน}ติ \textbf{ตสฺมา}ติ \textbf{ตสฺมา} ตํการณา ตํเหตุ ตปฺปจฺจยา ตํนิทานา เอตํ อาทีนวํ สมฺปสฺสมานา มมตฺเตสูติ \textbf{ตสฺมา}.\csromanspacer{} \textbf{มุนโย}ติ \textbf{โมนํ} วุจฺจติ ญาณํ.\csromanspacer{} ยา ปญฺญา ปชานนา ฯเปฯ อโมโห ธมฺมวิจโย สมฺมาทิฏฺฐิ.\csromanspacer{} เตน ญาเณน สมนฺนาคตา มุนโย โมนปฺปตฺตา.\csromanspacer{} ตีณิ โมเนยฺยานิ กายโมเนยฺยํ วจีโมเนยฺยํ มโนโมเนยฺยํ ฯเปฯ สงฺคชาลมติจฺจ โส มุนิ.\csromanspacer{} \textbf{ปริคฺคโห}ติ เทฺว ปริคฺคหา ตณฺหาปริคฺคโห จ ทิฏฺฐิปริคฺคโห จ ฯเปฯ อยํ ตณฺหาปริคฺคโห ฯเปฯ อยํ ทิฏฺฐิปริคฺคโห.\csromanspacer{} \textbf{มุนโย} ตณฺหาปริคฺคหํ ปริจฺจชิตฺวา ทิฏฺฐิปริคฺคหํ ปฏินิสฺสชฺชิตฺวา จชิตฺวา ปชหิตฺวา วิโนเทตฺวา พฺยนฺติํ กริตฺวา อนภาวํ คเมตฺวา อจริํสุ วิหริํสุ อิริยิํสุ วตฺติํสุ ปาลิํสุ ยปิํสุ ยาปิํสุ.\csromanspacer{}\csromanspacer{}\textbf{เขมทสฺสิโน}ติ \textbf{เขมํ} วุจฺจติ อมตํ นิพฺพานํ.\csromanspacer{} โย โส สพฺพสงฺขารสมโถ สพฺพูปธิปฏินิสฺสคฺโค ตณฺหกฺขโย วิราโค นิโรโธ นิพฺพานํ.\csromanspacer{} \textbf{เขมทสฺสิโน}ติ เขมทสฺสิโน ตาณทสฺสิโน เลณทสฺสิโน สรณทสฺสิโน อภยทสฺสิโน อจฺจุตทสฺสิโน อมตทสฺสิโน นิพฺพานทสฺสิโนติ \textbf{ตสฺมา} มุนโย ปริคฺคหํ, ยิตฺวา อจริํสุ เขมทสฺสิโน.\csromanspacer{} เตนาห ภควา\csromandash{}}

\csromangathameasure{โสกปฺปริเทวมจฺฉรํ, \\ น ชหนฺติ คิทฺธา มมายิเต.}
\csromangathagroup{\csromangathastanza{โสกปฺปริเทวมจฺฉรํ, \\ น ชหนฺติ คิทฺธา มมายิเต.}}

\prose{ตสฺมา มุนโย ปริคฺคหํ,}

\prose{หิตฺวา อจริํสุ เขมทสฺสิโนติ.}

\proseitem{45}{\csromansymbolfootnote{*}{ขุ 1. 406 ปิฏฺเฐ.}ปติลีนจรสฺส ภิกฺขุโน,}

\prose{\textbf{ภชมานสฺส วิวิตฺตมาสนํ.}}

\csromanmatikaanchor{mtk.7}
\csromantocmarkref{subsection}{6. ชราสุตฺตนิทฺเทส}{mtk.7}
\bookmark[dest={mtk.7},level=2]{6. ชราสุตฺตนิทฺเทส}
\csromangathameasure{\textbf{สามคฺคิยมาหุ ตสฺส ตํ,} \\ \textbf{โย อตฺตานํ ภวเน น ทสฺสเย.}}
\csromangathagroup{\csromangathastanza{\textbf{สามคฺคิยมาหุ ตสฺส ตํ,} \\ \textbf{โย อตฺตานํ ภวเน น ทสฺสเย.}}}

\csromanheader{2}{6. ชราสุตฺตนิทฺเทส}
\prose{\csromanfolio{101}\textbf{ปติลีนจรสฺส ภิกฺขุโน}ติ ปติลีนจรา วุจฺจนฺติ สตฺต เสกฺขา\footnote{เสขา (สี, สฺยา, ก)}, อรหา ปติลีโน.\csromanspacer{} กิํการณา ปติลีนจรา วุจฺจนฺติ สตฺต เสกฺขา, เต ตโต ตโต จิตฺตํ ปติลีเนนฺตา ปติกุเฏนฺตา ปติวฏฺเฏนฺตา สนฺนิรุทฺธนฺตา\footnote{สนฺนิรุมฺเภนฺตา (สี)} สนฺนิคฺคณฺหนฺตา สนฺนิวาเรนฺตา รกฺขนฺตา โคเปนฺตา จรนฺติ วิจรนฺติ วิหรนฺติ อิริยนฺติ วตฺเตนฺติ ปาเลนฺติ ยเปนฺติ ยาเปนฺติ, จกฺขุทฺวาเร จิตฺตํ ปติลีเนนฺตา ปติกุเฏนฺตา ปติวฏฺเฏนฺตา สนฺนิรุธนฺตา สนฺนิคฺคณฺหนฺตา สนฺนิวาเรนฺตา รกฺขนฺตา โคเปนฺตา จรนฺติ วิจรนฺติ วิหรนฺติ อิริยนฺติ วตฺเตนฺติ ปาเลนฺติ ยเปนฺติ ยาเปนฺติ.\csromanspacer{} โสตทฺวาเร จิตฺตํ ฯเปฯ.\csromanspacer{} ฆานทฺวาเร จิตฺตํ.\csromanspacer{} ชิวฺหาทฺวาเร จิตฺตํ.\csromanspacer{} กายทฺวาเร จิตฺตํ.\csromanspacer{} มโนทฺวาเร จิตฺตํ ปติลีเนนฺตา ปติกุเฏนฺตา ปติวฏฺเฏนฺตา สนฺนิรุทฺธนฺตา สนฺนิคฺคณฺหนฺตา สนฺนิวาเรนฺตา รกฺขนฺตา โคเปนฺตา จรนฺติ วิจรนฺติ วิหรนฺติ อิริยนฺติ วตฺเตนฺติ ปาเลนฺติ ยเปนฺติ ยาเปนฺติ.\csromanspacer{} ยถา กุกฺกุฏปตฺตํ วา นฺหารุททฺทุลํ วา อคฺคิมฺหิ ปกฺขิตฺตํ ปติลียติ ปติกุฏติ ปติวฏฺฏติ น สมฺปสาริยติ.\csromanspacer{} เอวเมวํ ตโต ตโต จิตฺตํ ปติลีเนนฺตา ปติกุเฏนฺตา ปติวฏฺเฏนฺตา สนฺนิรุทฺธนฺตา สนฺนิคฺคณฺหนฺตา สนฺนิวาเรนฺตา รกฺขนฺตา โคเปนฺตา จรนฺติ วิจรนฺติ วิหรนฺติ อิริยนฺติ วตฺเตนฺติ ปาเลนฺติ ยเปนฺติ ยาเปนฺติ.\csromanspacer{}\csromanspacer{}จกฺขุทฺวาเร จิตฺตํ.\csromanspacer{} โสตทฺวาเร จิตฺตํ.\csromanspacer{} ฆานทฺวาเร จิตฺตํ.\csromanspacer{} ชิวฺหาทฺวาเร จิตฺตํ.\csromanspacer{} กายทฺวาเร จิตฺตํ.\csromanspacer{} มโนทฺวาเร จิตฺตํ ปติลีเนนฺตา ปติกุเฏนฺตา ปติวฏฺเฏนฺตา สนฺนิรุทฺธนฺตา สนฺนิคฺคณฺหนฺตา สนฺนิวาเรนฺตา รกฺขนฺตา โคเปนฺตา จรนฺติ วิจรนฺติ วิหรนฺติ อิริยนฺติ วตฺเตนฺติ ปาเลนฺติ ยเปนฺติ ยาเปนฺติ, ตํการณา ปติลีนจรา วุจฺจนฺติ สตฺต เสกฺขา.\csromanspacer{} \textbf{ภิกฺขุโน}ติ ปุถุชฺชนกลฺยาณกสฺส วา ภิกฺขุโน, เสกฺขสฺส วา ภิกฺขุโนติ ปติลีนจรสฺส ภิกฺขุโน.}

\prose{\textbf{ภชมานสฺส วิวิตฺตมาสนนฺ}ติ \textbf{อาสนํ} วุจฺจติ ยตฺถ นิสีทนฺติ.\csromanspacer{} มญฺโจ ปีฐํ ภิสิ ตฏฺฏิกา จมฺมขณฺโฑ ติณสนฺถาโร ปณฺณสนฺถาโร ปลาลสนฺถาโร.\csromanspacer{} ตํ \textbf{อาสนํ} อสปฺปายรูปทสฺสเนน ริตฺตํ วิวิตฺตํ ปวิวิตฺตํ, อสปฺปายสทฺทสฺสวเนน ริตฺตํ วิวิตฺตํ ปวิวิตฺตํ, อสปฺปายคนฺธฆายเนน.\csromanspacer{} อสปฺปายรสสายเนน.\csromanspacer{} อสปฺปายโผฏฺฐพฺพผุสเนน.\csromanspacer{} อสปฺปาเยหิ ปญฺจหิ กามคุเณหิ ริตฺตํ วิวิตฺตํ ปวิวิตฺตํ, ตํ วิวิตฺตํ \textbf{อาสนํ} ภชโต สมฺภชโต เสวโต นิเสวโต สํเสวโต ปฏิเสวโตติ ภชมานสฺส วิวิตฺตมฺ\textbf{อาสนํ}.}

\prose{\csromanfolio{102}\textbf{สามคฺคิยมาหุ ตสฺส ตํ, โย อตฺตานํ ภวเน น ทสฺสเย}ติ \textbf{สามคฺคิโย}ติ ติสฺโส \textbf{สามคฺคิโย} คณสามคฺคี ธมฺมสามคฺคี อนภินิพฺพตฺติสามคฺคี.\csromanspacer{} กตมา คณสามคฺคี, พหู เจปิ ภิกฺขู สมคฺคา สมฺโมทมานา อวิวทมานา ขีโรทกีภูตา อญฺญมญฺญํ ปิยจกฺขูหิ สมฺปสฺสนฺตา วิหรนฺติ, อยํ คณสามคฺคี.\csromanspacer{}\csromanspacer{}กตมา ธมฺมสามคฺคี, จตฺตาโร สติปฏฺฐานา จตฺตาโร สมฺมปฺปธานา จตฺตาโร อิทฺธิปาทา ปญฺจินฺทฺริยานิ ปญฺจ พลานิ สตฺต โพชฺฌงฺคา อริโย อฏฺฐิงฺคิโก มคฺโค.\csromanspacer{} เต เอกโต ปกฺขนฺทนฺติ ปสีทนฺติ สมฺปติฏฺฐนฺติ วิมุจฺจนฺติ.\csromanspacer{} น เตสํ ธมฺมานํ วิวาโท ปวิวาโท อตฺถิ, อยํ ธมฺมสามคฺคี.\csromanspacer{}\csromanspacer{}กตมา อนภินิพฺพตฺติสามคฺคี, พหู เจปิ ภิกฺขู อนุปาทิเสสาย นิพฺพานธาตุยา ปรินิพฺพายนฺติ, น เตสํ\footnote{เตน (สี)} นิพฺพานธาตุยา อูนตฺตํ วา ปุณฺณตฺตํ วา ปญฺญายติ, อยํ อนภินิพฺพตฺติสามคฺคี.\csromanspacer{}\csromanspacer{}\textbf{ภวเน}ติ เนรยิกานํ นิรโย ภวนํ, ติรจฺฉานโยนิกานํ ติรจฺฉานโยนิ ภวนํ, เปตฺติวิสยิกานํ เปตฺติวิสโย ภวนํ, มนุสฺสานํ มนุสฺสโลโก ภวนํ, เทวานํ เทวโลโก ภวนฺติ.\csromanspacer{} \textbf{สามคฺคิยมาหุ ตสฺส ตํ, โย อตฺตานํ ภวเน น ทสฺสเย}ติ ตสฺเสสา สามคฺคี เอตํ ฉนฺนํ เอตํ ปติรูปํ เอตํ อนุจฺฉวิกํ เอตํ อนุโลมํ.\csromanspacer{} โย เอวํ ปฏิจฺฉนฺเน นิรเย อตฺตานํ น ทสฺเสยฺย, ติรจฺฉานโยนิยํ อตฺตานํ น ทสฺเสยฺย, เปตฺติวิสเย อตฺตานํ น ทสฺเสยฺย, มนุสฺสโลเก อตฺตานํ น ทสฺเสยฺย, เทวโลเก อตฺตานํ น ทสฺเสยฺยาติ เอวมาหํสุ เอวํ กเถนฺติ เอวํ ภณนฺติ เอวํ ทีปยนฺติ เอวํ โวหรนฺตีติ สามคฺคิยมาหุ ตสฺส ตํ, โย อตฺตานํ ภวเน น ทสฺสเย.\csromanspacer{} เตนาห ภควา\csromandash{}}

\csromangathameasure{ปติลีนจรสฺส ภิกฺขุโน, \\ ภชมานสฺส วิวิตฺตมาสนํ.}
\csromangathagroup{\csromangathastanza{ปติลีนจรสฺส ภิกฺขุโน, \\ ภชมานสฺส วิวิตฺตมาสนํ.}}

\prose{สามคฺคิยมาหุ ตสฺส ตํ,}

\prose{โย อตฺตานํ ภวเน น ทสฺสเยติ.}

\proseitem{46}{\csromansymbolfootnote{*}{ขุ 1. 406 ปิฏฺเฐ.}สพฺพตฺถ มุนี อนิสฺสิโต,}

\prose{\textbf{น ปิยํ กุพฺพติ โนปิ อปฺปิยํ.}}

\csromangathameasure{\textbf{ตสฺมิํ ปริเทวมจฺฉรํ,} \\ \textbf{ปณฺเณ วาริ ยถา น ลิมฺปติ.}}
\csromangathagroup{\csromangathastanza{\textbf{ตสฺมิํ ปริเทวมจฺฉรํ,} \\ \textbf{ปณฺเณ วาริ ยถา น ลิมฺปติ.}}}

\csromanheader{2}{6. ชราสุตฺตนิทฺเทส}
\prose{\csromanfolio{103}\textbf{สพฺพตฺถ มุนี อนิสฺสิโต}ติ \textbf{สพฺพํ} วุจฺจติ ทฺวาทสายตนานิ จกฺขุญฺเจว รูปา จ, โสตญฺจ สทฺทา จ, ฆานญฺจ คนฺธา จ, ชิวฺหา จ รสา จ, กาโย จ โผฏฺฐพฺพา จ, มโน จ ธมฺมา จ.\csromanspacer{} \textbf{มุนี}ติ \textbf{โมนํ} วุจฺจติ ญาณํ ฯเปฯ สงฺคชาลมติจฺจ โส มุนิ.\csromanspacer{} \textbf{อนิสฺสิโต}ติ เทฺว นิสฺสยา ตณฺหานิสฺสโย จ ทิฏฺฐินิสฺสโย จ ฯเปฯ อยํ ตณฺหานิสฺสโย ฯเปฯ อยํ ทิฏฺฐินิสฺสโย.\csromanspacer{} มุนิ ตณฺหานิสฺสยํ ปหาย ทิฏฺฐินิสฺสยํ ปฏินิสฺสชฺชิตฺวา จกฺขุํ อนิสฺสิโต, โสตํ อนิสฺสิโต, ฆานํ อนิสฺสิโต, ชิวฺหํ อนิสฺสิโต, กายํ อนิสฺสิโต, มนํ อนิสฺสิโต, รูเป.\csromanspacer{} สทฺเท.\csromanspacer{} คนฺเธ.\csromanspacer{} รเส.\csromanspacer{} โผฏฺฐพฺเพ.\csromanspacer{} ธมฺเม.\csromanspacer{} กุลํ.\csromanspacer{} คณํ.\csromanspacer{} อาวาสํ.\csromanspacer{} ลาภํ.\csromanspacer{} ยสํ.\csromanspacer{} ปสํสํ.\csromanspacer{} สุขํ.\csromanspacer{} จีวรํ.\csromanspacer{} ปิณฺฑปาตํ.\csromanspacer{} เสนาสนํ.\csromanspacer{} คิลานปจฺจยเภสชฺชปริกฺขารํ.\csromanspacer{} กามธาตุํ.\csromanspacer{} รูปธาตุํ.\csromanspacer{} อรูปธาตุํ.\csromanspacer{} กามภวํ.\csromanspacer{} รูปภวํ.\csromanspacer{} อรูปภวํ.\csromanspacer{} สญฺญาภวํ.\csromanspacer{} อสญฺญาภวํ.\csromanspacer{} เนวสญฺญานาสญฺญาภวํ.\csromanspacer{} เอกโวการภวํ.\csromanspacer{} จตุโวการภวํ.\csromanspacer{} ปญฺจโวการภวํ.\csromanspacer{} อตีตํ.\csromanspacer{} อนาคตํ.\csromanspacer{} ปจฺจุปฺปนฺนํ.\csromanspacer{} ทิฏฺฐํ.\csromanspacer{} สุตํ.\csromanspacer{} มุตํ.\csromanspacer{} วิญฺญาตํ.\csromanspacer{} สพฺเพ ธมฺเม อนิสฺสิโต อนลฺลีโน อนุปคโต อนชฺโฌสิโต อนธิมุตฺโต นิกฺขนฺโต นิสฺสโฏ วิปฺปมุตฺโต วิสญฺญุตฺโต วิมริยาทิกเตน เจตสา วิหรตีติ สพฺพตฺถ มุนิ อนิสฺสิโต.}

\prose{\textbf{น ปิยํ กุพฺพติ โนปิ อปฺปิยนฺติ ปิยา}ติ เทฺว ปิยา สตฺตา วา สงฺขารา วา.\csromanspacer{} กตเม สตฺตา ปิยา, อิธ ย’สฺส เต โหนฺติ อตฺถกามา หิตกามา ผาสุกามา โยคกฺเขมกามา มาตา วา ปิตา วา ภาตา วา ภคินี วา ปุตฺตา วา ธีตรา วา มิตฺตา วา อมจฺจา วา ญาตี วา สาโลหิตา วา, อิเม สตฺตา ปิยา.\csromanspacer{} กตเม สงฺขารา ปิยา, มนาปิกา รูปา มนาปิกา สทฺทา มนาปิกา คนฺธา มนาปิกา รสา มนาปิกา โผฏฺฐพฺพา, อิเม สงฺขารา ปิยา.\csromanspacer{} \textbf{อปฺปิยา}ติ เทฺว อปฺปิยา สตฺตา วา สงฺขารา วา.\csromanspacer{} กตเม สตฺตา อปฺปิยา, อิธ ย’สฺส เต โหนฺติ อนตฺถกามา อหิตกามา อผาสุกามา อโยคกฺเขมกามา ชีวิตา โวโรเปตุกามา, อิเม สตฺตา อปฺปิยา.\csromanspacer{}\csromanspacer{}กตเม สงฺขารา อปฺปิยา, อมนาปิกา รูปา อมนาปิกา สทฺทา อมนาปิกา คนฺธา อมนาปิกา รสา อมนาปิกา โผฏฺฐพฺพา, อิเม สงฺขารา อปฺปิยา.\csromanspacer{}\csromanspacer{}\textbf{น ปิยํ กุพฺพติ โนปิ อปฺปิยนฺ}ติ “อยํ เม สตฺโต ปิโย, อิเม จ สงฺขารา มนาปา”ติ ราควเสน ปิยํ น กโรติ, “อยํ เม สตฺโต อปฺปิโย, อิเม จ สงฺขารา อมนาปา”ติ}

\prose{\csromanfolio{104}ปฏิฆวเสน อปฺปิยํ น กโรติ น ชเนติ น สญฺชเนติ น นิพฺพตฺเตติ นาภินิพฺพตฺเตตีติ น ปิยํ กุพฺพติ โนปิ อปฺปิยํ.}

\prose{\textbf{ตสฺมิํ ปริเทวมจฺฉรํ, ปณฺเณ วาริ ยถา น ลิมฺปตี}ติ \textbf{ตสฺมินฺ}ติ ตสฺมิํ ปุคฺคเล อรหนฺเต ขีณาสเว.\csromanspacer{} \textbf{ปริเทโว}ติ ญาติพฺยสเนน วา ผุฏฺฐสฺส โภคพฺยสเนน วา ผุฏฺฐสฺส โรคพฺยสเนน วา ผุฏฺฐสฺส สีลพฺยสเนน วา ผุฏฺฐสฺส ทิฏฺฐิพฺยสเนน วา ผุฏฺฐสฺส อญฺญตรญฺญตเรน พฺยสเนน สมนฺนาคตสฺส อญฺญตรญฺญตเรน ทุกฺขธมฺเมน ผุฏฺฐสฺส อาเทโว ปริเทโว อาเทวนา ปริเทวนา อาเทวิตตฺตํ ปริเทวิตตฺตํ วาจา ปลาโป วิปฺปลาโป ลาลปฺโป ลาลปฺปายนา ลาลปฺปายิตตฺตํ.\csromanspacer{}\csromanspacer{}\textbf{มจฺฉริยนฺ}ติ ปญฺจ มจฺฉริยานิ อาวาสมจฺฉริยํ กุลมจฺฉริยํ ลาภมจฺฉริยํ วณฺณมจฺฉริยํ ธมฺมมจฺฉริยํ.\csromanspacer{} ยํ เอวรูปํ มจฺฉริยํ มจฺฉรายนา มจฺฉรายิตตฺตํ เววิจฺฉํ กทริยํ กตุกญฺจุกตา อคฺคหิตตฺตํ จิตฺตสฺส, อิทํ วุจฺจติ มจฺฉริยํ.\csromanspacer{} อปิ จ ขนฺธมจฺฉริยมฺปิ มจฺฉริยํ, ธาตุมจฺฉริยมฺปิ มจฺฉริยํ, อายตนมจฺฉริยมฺปิ มจฺฉริยํ คาโห, อิทํ วุจฺจติ มจฺฉริยํ.}

\prose{\textbf{ปณฺเณ วาริ ยถา น ลิมฺปตี}ติ \textbf{ปณฺณํ} วุจฺจติ ปทุมปตฺตํ.\csromanspacer{} \textbf{วาริ} วุจฺจติ อุทกํ.\csromanspacer{} ยถา วาริ ปทุมปตฺตํ น ลิมฺปติ น ปลิมฺปติ, น อุปลิมฺปติ อลิตฺตํ อปลิตฺตํ อนุปลิตฺตํ.\csromanspacer{} เอวเมวํ ตสฺมิํ ปุคฺคเล อรหนฺเต ขีณาสเว ปริเทโว มจฺฉริยญฺจ น ลิมฺปติ น ปลิมฺปติ น อุปลิมฺปติ, อลิตฺตา อปลิตฺตา อนุปลิตฺตา.\csromanspacer{} โส จ ปุคฺคโล อรหนฺโต เตหิ กิเลเสหิ น ลิปฺปติ น ปลิปฺปติ น อุปลิปฺปติ, อลิตฺโต อปลิตฺโต อนุปลิตฺโต นิกฺขนฺโต นิสฺสโฏ วิปฺปมุตฺโต วิสญฺญุตฺโต วิมริยาทิกเตน เจตสา วิหรตีติ ตสฺมิํ ปริเทวมจฺฉรํ, ปณฺเณ วาริ ยถา น ลิมฺปติ.\csromanspacer{} เตนาห ภควา\csromandash{}}

\csromangathameasure{สพฺพตฺถ มุนี อนิสฺสิโต, น ปิยํ กุพฺพติ โนปิ อปฺปิยํ.}
\csromangathagroup{\csromangathastanza{สพฺพตฺถ มุนี อนิสฺสิโต, น ปิยํ กุพฺพติ โนปิ อปฺปิยํ.}}

\prose{\textbf{ตสฺมิํ ปริเทวมจฺฉรํ, ปณฺเณ วาริ ยถา น ลิมฺปตี}ติ.}

\proseitem{47}{\csromansymbolfootnote{*}{ขุ 1. 406 ปิฏฺเฐ.}\textbf{อุทพินฺทุ ยถาปิ โปกฺขเร}, ปทุเม วาริ ยถา น ลิมฺปติ.}

\csromangathameasure{เอวํ มุนิ โนปลิมฺปติ, ยทิทํ ทิฏฺฐสุตํ มุเตสุ วา.}
\csromangathagroup{\csromangathastanza{เอวํ มุนิ โนปลิมฺปติ, ยทิทํ ทิฏฺฐสุตํ มุเตสุ\footnote{ทิฏฺฐสุเต มุเตสุ (สี), ทิฏฺฐสุตมุเตสุ (ก)} วา.}}

\prose{\textbf{อุทพินฺทุ ยถาปิ โปกฺขเร}ติ อุทพินฺทุ วุจฺจติ อุทกเถโว.\csromanspacer{} โปกฺขรํ วุจฺจติ ปทุมปตฺตํ.\csromanspacer{} ยถา อุทพินฺทุ ปทุมปตฺเต น ลิมฺปติ น ปลิมฺปติ น}

\vspace{\tipitakaparskip}
\csromancenter{ชราสุตฺตนิทฺเทส อุปลิมฺปติ, อลิตฺตํ อปลิตฺตํ อนุปลิตฺตนฺติ อุทพินฺทุ ยถาปิ โปกฺขเร.\csromanspacer{}\csromanspacer{}\textbf{ปทุเม วาริ ยถา น ลิมฺปตี}ติ ปทุมํ วุจฺจติ ปทุมปุปฺผํ.\csromanspacer{} วาริ วุจฺจติ อุทกํ.\csromanspacer{} ยถา วาริ ปทุมปุปฺผํ น ปลิมฺปติ น ปลิมฺปติ น อุปลิมฺปติ อลิตฺตํ อปลิตฺตํ อนุปลิตฺตนฺติ ปทุเม วาริ ยถา น ลิมฺปติ.}
\prose{\csromanfolio{105}\textbf{เอวํ มุนิ โนปลิมฺปติ, ยทิทํ ทิฏฺฐสุตํ มุเตสุ วา}ติ \textbf{เอว}นฺติ โอปมฺมสมฺปฏิปาทนํ.\csromanspacer{} \textbf{มุนี}ติ \textbf{โมนํ} วุจฺจติ ญาณํ ฯเปฯ สงฺคชาลมติจฺจ โส มุนิ.\csromanspacer{} \textbf{เลปา}ติ เทฺว เลปา ตณฺหาเลโป จ ทิฏฺฐิเลโป จ ฯเปฯ อยํ ตณฺหาเลโป ฯเปฯ อยํ ทิฏฺฐิเลโป.\csromanspacer{} มุนิ ตณฺหาเลปํ ปหาย ทิฏฺฐิเลปํ ปฏินิสฺสชฺชิตฺวา ทิฏฺเฐ น ลิมฺปติ, สุเต น ลิมฺปติ, มุเต น ลิมฺปติ วิญฺญาเต น ลิมฺปติ น ปลิมฺปติ, น อุปลิมฺปติ, อลิตฺโต อปลิตฺโต อนุปลิตฺโต นิกฺขนฺโต นิสฺสโฏ วิปฺปมุตฺโต วิสญฺญุตฺโต วิมริยาทิกเตน เจตสา วิหรตีติ \textbf{เอว}ํ มุนิ โนปลิมฺปติ, ยทิทํ ทิฏฺฐสุตํ มุเตสุ วา.\csromanspacer{} เตนาห ภควา\csromandash{}}

\csromangathameasure{อุทพินฺทุ ยถาปิ โปกฺขเร, \\ ปฺทุเม วาริ ยถา น ลิมฺปติ.}
\csromangathagroup{\csromangathastanza{อุทพินฺทุ ยถาปิ โปกฺขเร, \\ ปฺทุเม วาริ ยถา น ลิมฺปติ.}}

\prose{เอวํ มุนิ โนปลิมฺปติ,}

\prose{ยฺทิทํ ทิฏฺฐสุตํ มุเตสุ วาติ.}

\proseitem{48}{\csromansymbolfootnote{*}{ขุ 1. 406 ปิฏฺเฐ.}โธโน น หิ เตน มญฺญติ,}

\prose{\textbf{ยทิทํ ทิฏฺฐสุตํ มุเตสุ วา.}}

\csromangathameasure{\textbf{นาญฺเญน วิสุทฺธิมิจฺฉติ,} \\ \textbf{น หิ โส รชฺชติ โน วิรชฺชติ.}}
\csromangathagroup{\csromangathastanza{\textbf{นาญฺเญน วิสุทฺธิมิจฺฉติ,} \\ \textbf{น หิ โส รชฺชติ โน วิรชฺชติ.}}}

\prose{\textbf{โธโน น หิ เตน มญฺญติ, ยทิทํ ทิฏฺฐสุตํ มุเตสุ วา}ติ \textbf{โธโน}ติ \textbf{โธนา} วุจฺจติ ปญฺญา.\csromanspacer{} ยา ปญฺญา ปชานนา ฯเปฯ อโมโห ธมฺมวิจโย สมฺมาทิฏฺฐิ.\csromanspacer{} กิํ การณา \textbf{โธนา} วุจฺจติ ปญฺญา.\csromanspacer{} ตาย ปญฺญาย กายทุจฺจริตํ ธุตญฺจ โธตญฺจ สนฺโธตญฺจ นิทฺโธตญฺจ, วจีทุจฺจริตํ มโนทุจฺจริตํ ธุตญฺจ โธตญฺจ สนฺโธตญฺจ นิทฺโธตญฺจ, ราโค ธุโต จ โธโต จ สนฺโธโต จ นิทฺโธโต จ, โทโส.\csromanspacer{} โมโห.\csromanspacer{} โกโธ.\csromanspacer{} อุปนาโห.\csromanspacer{} มกฺโข.\csromanspacer{} ปฬาโส.\csromanspacer{} อิสฺสา.\csromanspacer{} มจฺฉริยํ.\csromanspacer{} มายา.\csromanspacer{} สาเฐยฺยํ.\csromanspacer{} ถมฺโภ.\csromanspacer{} สารมฺโภ.\csromanspacer{} มาโน.\csromanspacer{} อติมาโน.\csromanspacer{} มโท.\csromanspacer{} ปมาโท.\csromanspacer{} สพฺเพ กิเลสา.\csromanspacer{} สพฺเพ ทุจฺจริตา.\csromanspacer{} สพฺเพ ทรถา.\csromanspacer{} สพฺเพ ปริฬาหา.\csromanspacer{} สพฺเพ สนฺตาปา.\csromanspacer{} สพฺพากุสลาภิสงฺขารา ธุตา จ โธตา จ สนฺโธตา จ นิทฺโธตา จ.\csromanspacer{} ตํการณา \textbf{โธนา} วุจฺจติ ปญฺญา.}

\prose{\csromanfolio{106}อถ วา สมฺมาทิฏฺฐิยา มิจฺฉาทิฏฺฐิ ธุตา จ โธตา จ สนฺโธตา จ นิทฺโธตา จ, สมฺมาสงฺกปฺเปน มิจฺฉาสงฺกปฺโป ธุโต จ โธโต จ สนฺโธโต จ นิทฺโธโต จ, สมฺมาวาจาย มิจฺฉาวาจา ธุตา จ.\csromanspacer{} สมฺมากมฺมนฺเตน มิจฺฉากมฺมนฺโต ธุโต จ.\csromanspacer{} สมฺมา-อาชีเวน มิจฺฉา-อาชีโว ธุโต จ.\csromanspacer{} สมฺมาวายาเมน มิจฺฉาวายาโม ธุโต จ.\csromanspacer{} สมฺมาสติยา มิจฺฉาสติ ธุตา จ.\csromanspacer{} สมฺมาสมาธินา มิจฺฉาสมาธิ ธุโต จ.\csromanspacer{} สมฺมาญาเณน มิจฺฉาญาณํ ธุตญฺจ.\csromanspacer{} สมฺมาวิมุตฺติยา มิจฺฉาวิมุตฺติ ธุตา จ โธตา จ สนฺโธตา จ นิทฺโธตา จ.}

\prose{อถ วา อริเยน อฏฺฐงฺคิเกน มคฺเคน สพฺเพ กิเลสา.\csromanspacer{} สพฺเพ ทุจฺจริตา.\csromanspacer{} สพฺเพ ทรถา.\csromanspacer{} สพฺเพ ปริฬาหา.\csromanspacer{} สพฺเพ สนฺตาปา.\csromanspacer{} สพฺพากุสลาภิสงฺขารา ธุตา จ โธตา จ สนฺโธตา จ นิทฺโธตา จ.\csromanspacer{} อรหา อิเมหิ โธเนหิ ธมฺเมหิ อุเปโต สมุเปโต อุปคโต สมุปคโต อุปปนฺโน สมุปปนฺโน สมนฺนาคโต.\csromanspacer{} ตสฺมา อรหา โธโน.\csromanspacer{} โส ธุตราโค ธุตปาโป ธุตกิเลโส ธุตปริฬาโหติ โธโน.}

\prose{\textbf{โธโน น หิ เตน มญฺญติ, ยทิทํ ทิฏฺฐสุตํ มุเตสุ วา}ติ โธโน ทิฏฺฐํ น มญฺญติ, ทิฏฺฐสมิํ น มญฺญติ, ทิฏฺฐโต น มญฺญติ, “ทิฏฺฐา เม”ติ น มญฺญติ.\csromanspacer{} สุตํ น มญฺญติ, สุตสฺมิํ น มญฺญติ, สุตโต น มญฺญติ, “สุตํ เม”ติ น มญฺญติ.\csromanspacer{} มุตํ น มญฺญติ, มุตสฺมิํ น มญฺญติ, มุตโต น มญฺญติ, “มุตํ เม”ติ น มญฺญติ.\csromanspacer{} วิญฺญาตํ น มญฺญติ, วิญฺญาตสฺมิํ น มญฺญติ, วิญฺญาตโต น มญฺญติ, “วิญฺญาตํ เม”ติ น มญฺญติ.\csromanspacer{} วุตฺตมฺปิ เหตํ ภควตา\csromandash{} * “อสฺมี”ติ ภิกฺขเว มญฺญิตเมตํ, “อยมหมสฺมี”ติ มญฺญิตเมตํ. “ภวิสฺสนฺ”ติ มญฺญิตเมตํ, “น ภวิสฺสนฺ”ติ มญฺญิตเมตํ. “รูปี ภวิสฺสนฺ”ติ มญฺญิตเมตํ, “อรูปี ภวิสฺสนฺ”ติ มญฺญิตเมตํ. “สญฺญี ภวิสฺสนฺ”ติ มญฺญิตเมตํ, “อสญฺญี ภวิสฺสนฺ”ติ มญฺญิตเมตํ, “เนวสญฺญีนาสญฺญี ภวิสฺสนฺ”ติ มญฺญิตเมตํ.\csromanspacer{} มญฺญิตํ\footnote{มญฺญิตํ หิ (สี)} ภิกฺขเว โรโค, มญฺญิตํ คณฺโฑ, มญฺญิตํ สลฺลํ, มญฺญิตํ อุปทฺทโว.\csromanspacer{} ตสฺมาติห ภิกฺขเว อมญฺญมาเนน เจตสา วิหริสฺสามาติ เอวญฺหิ โว ภิกฺขเว สิกฺขิตพฺพนฺติ โธโน น หิ เตน มญฺญติ, ยทิทํ ทิฏฺฐสุตํ มุเตสุ วา.}

\csromanheader{2}{7. ติสฺสเมตฺเตยฺยสุตฺตนิทฺเทส}
\prose{\csromanfolio{107}\textbf{นาญฺเญน วิสุทฺธิมิจฺฉตี}ติ โธโน อญฺเญน อสุทฺธิมคฺเคน มิจฺฉาปฏิปทาย อนิยฺยานิกปเถน อญฺญตฺร สติปฏฺฐาเนหิ อญฺญตฺร สมฺมปฺปธาเนหิ อญฺญตฺร อิทฺธิปาเทหิ อญฺญตฺร อินฺทฺริเยหิ อญฺญตฺร พเลหิ อญฺญตฺร โพชฺฌงฺเคหิ อญฺญตฺร อริยา อฏฺฐงฺคิกา มคฺคา สุทฺธิํ วิสุทฺธิํ ปริสุทฺธิํ มุตฺติํ วิมุตฺติํ ปริมุตฺติํ น อิจฺฉติ น สาทิยติ น ปตฺเถติ น ปิเหติ นาภิชปฺปตีติ นาญฺเญน วิสุทฺธิมิจฺฉติ.}

\prose{\textbf{น หิ โส รชฺชติ โน วิรชฺชตี}ติ สพฺเพ พาลปุถุชฺชนา รชฺชนฺติ, ปุถุชฺชนกลฺยณกํ อุปาทาย สตฺต เสกฺขา วิรชฺชนฺติ, อรหา เนว รชฺชติ โน วิรชฺชติ.\csromanspacer{} วิรตฺโต โส ขยา ราคสฺส วีตราคตฺตา, ขยา โทสสฺส วีตโทสตฺตา, ขยา โมหสฺส วีตโมหตฺตา, โส วุฏฺฐวาโส จิณฺณจรโณ ฯเปฯ.\csromanspacer{} ชาติชรามรณสํสาโร, นตฺถิ ตสฺส ปุนพฺภโวติ น หิ โส รชฺชติ โน วิรชฺชติ.\csromanspacer{} เตนาห ภควา\csromandash{}}

\csromangathameasure{โธโน น หิ เตน มญฺญติ, \\ ยทิทํ ทิฏฺฐสุตํ มุเตสุ วา. \\ นาญฺเญน วิสุทฺธิมิจฺฉติ,}
\csromangathagroup{\csromangathastanza{โธโน น หิ เตน มญฺญติ, \\ ยทิทํ ทิฏฺฐสุตํ มุเตสุ วา. \\ นาญฺเญน วิสุทฺธิมิจฺฉติ,}}

\prose{\textbf{น หิ โส รชฺชติ โน วิรชฺชตี}ติ.\csromanspacer{} ชราสุตฺตนิทฺเทโส ฉฏฺโฐ.}
\csromansectionrule

\csromanheader{2}{7. ติสฺสเมตฺเตยฺยสุตฺตนิทฺเทส}
\csromanheader{2}{อถ ติสฺสเมตฺเตยฺยสุตฺตนิทฺเทสํ วกฺขติ\csromandash{}}
\proseitem{49}{\csromansymbolfootnote{*}{ขุ 1. 406 ปิฏฺเฐ.}เมถุนมนุยุตฺตสฺส, (อิจฺจายสฺมา ติสฺโส เมตฺเตยฺโย,)}

\prose{\textbf{วิฆาตํ พฺรูหิ มาริส.}}

\csromangathameasure{\textbf{สุตฺวาน ตว สาสนํ,} \\ \textbf{วิเวเก สิกฺขิสฺสามเส.}}
\csromangathagroup{\csromangathastanza{\textbf{สุตฺวาน ตว สาสนํ,} \\ \textbf{วิเวเก สิกฺขิสฺสามเส.}}}

\prose{\textbf{เมถุนมนุยุตฺตสฺสา}ติ + \textbf{เมถุนธมฺโม} นาม โย โส อสทฺธมฺโม คามธมฺโม วสลธมฺโม ทุฏฺฏุลฺโล โอทกนฺติโก รหสฺโส ทฺวยํทฺวยสมาปตฺติ.\csromanspacer{} กิํ การณา วุจฺจติ \textbf{เมถุนธมฺโม}, อุภินฺนํ รตฺตานํ สารตฺตานํ อวสฺสุตานํ ปริยุฏฺฐิตานํ ปริยาทินฺนจิตฺตานํ อุภินฺนํ สทิสานํ ธมฺโมติ, ตํการณา วุจฺจติ \textbf{เมถุนธมฺโม}.\csromanspacer{} ยถา อุโภ กลหการกา \csromanfolio{108}เมถุนกาติ วุจฺจติ.\csromanspacer{} อุโภ ภณฺฑนการกา เมถุนกาติ วุจฺจนฺติ.\csromanspacer{} อุโภ ภสฺสการกา เมถุนกาติ วุจฺจนฺติ.\csromanspacer{} อุโภ วิวาทการกา เมถุนกาติ วุจฺจนฺติ.\csromanspacer{} อุโภ อธิกรณการกา เมถุนกาติ วุจฺจนฺติ.\csromanspacer{} อุโภ วาทิโน เมถุนกาติ วุจฺจนฺติ.\csromanspacer{} อุโภ สลฺลาปกา เมถุนกาติ วุจฺจนฺติ, เอวเมวํ อุภินฺนํ รตฺตานํ สารตฺตานํ อวสฺสุตานํ ปริยุฏฺฐิตานํ ปริยาทินฺนจิตฺตานํ อุภินฺนํ สทิสานํ ธมฺโมติ ตํการณา วุจฺจติ เมถุนธมฺโม.}

\prose{\textbf{เมถุนมนุยุตฺตสฺสา}ติ เมถุนธมฺเม ยุตฺตสฺส ปยุตฺตสฺส อายุตฺตสฺส สมายุตฺตสฺส ตจฺจริตสฺส ตพฺพหุลสฺส ตคฺครุกสฺส ตนฺนินฺนสฺส ตปฺโปณสฺส ตปฺปพฺภารสฺส ตทธิมุตฺตสฺส ตทธิปเตยฺยสฺสาติ เมถุนมนุยุตฺตสฺส.}

\prose{\textbf{อิจฺจายสฺมา ติสฺโส เมตฺเตยฺโย}ติ \textbf{อิจฺจา}ติ ปทส\textbf{นฺธี} ปทสํสคฺโค ปทปาริปูรี อกฺขรสมวาโย พฺยญฺชนสิลิฏฺฐตา ปทานุปุพฺพตาเปตํ \textbf{อิจฺจา}ติ.\csromanspacer{}\csromanspacer{}\textbf{อายสฺมา}ติ ปิยวจนํ ครุวจนํ สคารววจนํ สปฺปติสฺสวจนเมตํ อายสฺมาติ.\csromanspacer{}\csromanspacer{}\textbf{ติสฺโส}ติ ตสฺส เถรสฺส นามํ สงฺขา สมญฺญา ปญฺญตฺติ โวหาโร นามํ นามกมฺมํ นามเธยฺยํ นิรุตฺติ พฺยญฺชนํ อภิลาโป.\csromanspacer{}\csromanspacer{}\textbf{เมตฺเตยฺโย}ติ ตสฺส เถรสฺส โคตฺตํ สงฺขา สมญฺญา ปญฺญตฺติ โวหาโรติ \textbf{อิจฺจา}ยสฺมา \textbf{ติสฺโส} \textbf{เมตฺเตยฺโย}.}

\prose{\textbf{วิฆาตํ พฺรูหิ มาริสา}ติ \textbf{วิฆาตนฺ}ติ วิฆาตํ อุปฆาตํ ปีฬนํ ฆฏฺฏนํ อุปทฺทวํ อุปสคฺคํ พฺรูหิ อาจิกฺข เทเสหิ ปญฺญเปหิ ปฏฺฐเปหิ วิวร วิภช อุตฺตานีกโรหิ\footnote{อุตฺตานิํ กโรหิ (ก)} ปกาเสหิ.\csromanspacer{}\csromanspacer{}\textbf{มาริสา}ติ ปิยวจนํ ครุวจนํ สคารววจนํ สปฺปติสฺสวจนเมตํ มาริสาติ วิฆาตํ พฺรูหิ มาริส.}

\prose{\textbf{สุตฺวาน ตว สาสนนฺ}ติ ตุยฺหํ วจนํ พฺยปฺปถํ เทสนํ อนุสาสนํ อนุสิฏฺฐิํ สุตฺวา สุณิตฺวา อุคฺคเหตฺวา อุปธารยิตฺวา อุปลกฺขยิตฺวาติ สุตฺวาน ตว สาสนํ.}

\prose{\textbf{วิเวเก สิกฺขิสฺสามเส}ติ \textbf{วิเวโก}ติ * ตโย วิเวกา กาย\textbf{วิเวโก} จิตฺต\textbf{วิเวโก} อุปธิ\textbf{วิเวโก}.\csromanspacer{} กตโม กาย\textbf{วิเวโก}\csromandash{}อิธ ภิกฺขุ วิวิตฺตํ เสนาสนํ ภชติ อรญฺญํ รุกฺขมูลํ ปพฺพตํ กนฺทรํ คิริคุหํ สุสานํ วนปตฺถํ อพฺโภกาสํ ปลาสปุญฺชํ, กาเยน วิวิตฺโต วิหรติ.\csromanspacer{} โส เอโก คจฺฉติ, เอโก ติฏฺฐติ, เอโก นิสีทติ, เอโก เสยฺยํ กปฺเปติ, เอโก คามํ ปิณฺฑาย ปวิสติ, เอโก ปฏิกฺกมติ, เอโก รโห}

\vspace{\tipitakaparskip}
\csromancenter{ติสฺสเมตฺเตยฺยสุตฺตนิทฺเทส นิสีทติ, เอโก จงฺกมํ อธิฏฺฐาติ, เอโก จรติ, เอโก วิหรติ อิริยติ วตฺตติ ปาเลติ ยเปติ ยาเปติ, อยํ กายวิเวโก.}
\prose{\csromanfolio{109}กตโม จิตฺตวิเวโก\csromandash{}ปฐมํ ฌานํ สมาปนฺนสฺส นีวรเณหิ จิตฺตํ วิวิตฺตํ โหติ, ทุติยํ ฌานํ สมาปนฺนสฺส วิตกฺกวิจาเรหิ จิตฺตํ วิวิตฺตํ โหติ, ตติยํ ฌานํ สมาปนฺนสฺส ปีติยา จิตฺตํ วิวิตฺตํ โหติ, จตุตฺถํ ฌานํ สมาปนฺนสฺส สุขทุกฺเขหิ จิตฺตํ วิวิตฺตํ โหติ, อากาสานญฺจายตนํ สมาปนฺนสฺส รูปสญฺญาย ปฏิฆสญฺญาย นานตฺตสญฺญาย จิตฺตํ วิวิตฺตํ โหติ, วิญฺญาณญฺจายตนํ สมาปนฺนสฺส อากาสานญฺจายตนสญฺญาย จิตฺตํ วิวิตฺตํ โหติ, อากิญฺจญฺญายตนํ สมาปนฺนสฺส วิญฺญาณญฺจายตนสญฺญาย จิตฺตํ วิวิตฺตํ โหติ, เนวสญฺญานาสญฺญายตนํ สมาปนฺนสฺส อากิญฺจญฺญายตนสญฺญาย จิตฺตํ วิวิตฺตํ โหติ.\csromanspacer{} โสตาปนฺนสฺส สกฺกายทิฏฺฐิยา วิจิกิจฺฉาย สีลพฺพตปรามาสา ทิฏฺฐานุสยา วิจิกิจฺฉานุสยา ตเทกฏฺเฐหิ จ กิเลเสหิ จิตฺตํ วิวิตฺตํ โหติ.\csromanspacer{} สกทาคามิสฺส โอฬาริกา กามราคสญฺโญชนา ปฏิฆสญฺโญชนา โอฬาริกา กามราคานุสยา ปฏิฆานุสยา ตเทกฏฺเฐหิ จ กิเลเสหิ จิตฺตํ วิวิตฺตํ โหติ, อนาคามิสฺส อณุสหคตา กามราคสญฺโญชนา ปฏิฆสญฺโญชนา อณุสหคตา กามราคานุสยา ปฏิฆานุสยา ตเทกฏฺเฐหิ จ กิเลเสหิ จิตฺตํ วิวิตฺตํ โหติ, อรหโต รูปราคา อรูปราคา มานา อุทฺธจฺจา อวิชฺชาย มานานุสยา ภวราคานุสยา อวิชฺชานุสยา ตเทกฏฺเฐหิ จ กิเลเสหิ พหิทฺธา จ สพฺพนิมิตฺเตหิ จิตฺตํ วิวิตฺตํ โหติ.\csromanspacer{} อยํ จิตฺตวิเวโก.}

\prose{กตโม \textbf{อุปธิ}วิเวโก\csromandash{}\textbf{อุปธิ} วุจฺจนฺติ กิเลสา จ ขนฺธา จ อภิสงฺขารา จ.\csromanspacer{} อุปธิวิเวโก วุจฺจติ อมตํ นิพฺพานํ.\csromanspacer{} โย โส สพฺพสงฺขารสมโถ สพฺพูปธิปฏินิสฺสคฺโค ตณฺหกฺขโย วิราโค นิโรโธ นิพฺพานํ, อยํ \textbf{อุปธิ}วิเวโก.\csromanspacer{} กายวิเวโก จ วิเวกฏฺฐกายานํ เนกฺขมฺมาภิรตานํ, จิตฺตวิเวโก จ ปริสุทฺธจิตฺตานํ ปรมโวทานปตฺตานํ, \textbf{อุปธิ}วิเวโก จ นิรูปธีนํ ปุคฺคลานํ วิสงฺขารคตานํ.\csromanspacer{}\csromanspacer{}\textbf{วิเวเก สิกฺขิสฺสามเส}ติ โส เถโร ปกติยา สิกฺขิตสิกฺโข.\csromanspacer{} อปิ จ ธมฺมเทสนํ อุปาทาย ธมฺมเทสนํ สาเวนฺโต\footnote{ยาจนฺโต (สี, สฺยา)} เอวมาห วิเวเก สิกฺขิสฺสามเสติ.\csromanspacer{} เตนาห เถโร ติสฺสเมตฺเตยฺโย\csromandash{}}

\prose{\csromanfolio{110}เมถุนมนุยุตฺตสฺส, (อิจฺจายสฺมา ติสฺโส เมตฺเตยฺโย,)}

\prose{วิฆาตํ พฺรูหิ มาริส.}

\csromangathameasure{สุตฺวาน ตว สาสนํ, \\ วิเวเก สิกฺขิสฺสามเสติ.}
\csromangathagroup{\csromangathastanza{สุตฺวาน ตว สาสนํ, \\ วิเวเก สิกฺขิสฺสามเสติ.}}

\proseitem{50}{\csromansymbolfootnote{*}{ขุ 1. 406 ปิฏฺเฐปิ.}เมถุนมนุยุตฺตสฺส, (\textbf{เมตฺเตยฺยา}ติ \textbf{ภควา},)}

\prose{\textbf{มุสฺสเต วาปิ สาสนํ.}}

\csromangathameasure{\textbf{มิจฺฉา จ ปฏิปชฺชติ,} \\ \textbf{เอตํ ตสฺมิํ อนาริยํ.}}
\csromangathagroup{\csromangathastanza{\textbf{มิจฺฉา จ ปฏิปชฺชติ,} \\ \textbf{เอตํ ตสฺมิํ อนาริยํ.}}}

\prose{\textbf{เมถุนมนุยุตฺตสฺสา}ติ + \textbf{เมถุนธมฺโม} นาม โย โส อสทฺธมฺโม คามธมฺโม วสลธมฺโม ทุฏฺฐุลฺโล โอทกนฺติโก รหสฺโส ทฺวยํทฺวยสมาปตฺติ.\csromanspacer{} กิํการณา วุจฺจติ \textbf{เมถุนธมฺโม}, อุภินฺนํ รตฺตานํ สารตฺตานํ อวสฺสุตานํ ปริยุฏฺฐิตานํ ปริยาทินฺนจิตฺตานํ อุภินฺนํ สทิสานํ ธมฺโมติ ตํการณา วุจฺจติ \textbf{เมถุนธมฺโม}.\csromanspacer{} ยถา อุโภ กลหการกา เมถุนกาติ วุจฺจนฺติ.\csromanspacer{} อุโภ ภณฺฑนการกา เมถุนกาติ วุจฺจนฺติ.\csromanspacer{} อุโภ ภสฺสการกา เมถุนกาติ วุจฺจนฺติ.\csromanspacer{} อุโภ วิวาทการกา เมถุนกาติ วุจฺจนฺติ.\csromanspacer{} อุโภ อธิกรณการกา เมถุนกาติ วุจฺจนฺติ.\csromanspacer{} อุโภ วาทิโน เมถุนกาติ วุจฺจนฺติ.\csromanspacer{} อุโภ สลฺลาปกา เมถุนกาติ วุจฺจนฺติ, เอวเมวํ อุภินฺนํ รตฺตานํ สารตฺตานํ อวสฺสุตานํ ปริยุฏฺฐิตานํ ปริยาทินฺนจิตฺตานํ อุภินฺนํ สทิสานํ ธมฺโมติ ตํการณา วุจฺจติ \textbf{เมถุนธมฺโม}.}

\prose{\textbf{เมถุนมนุยุตฺตสฺสา}ติ เมถุนธมฺเม ยุตฺตสฺส ปยุตฺตสฺส อายุตฺตสฺส สมายุตฺตสฺส ตจฺจริตสฺส ตพฺพหุลสฺส ตคฺครุกสฺส ตนฺนินฺนสฺส ตปฺโปณสฺส ตปฺปพฺภารสฺส ตทธิมุตฺตสฺส ตทธิปเตยฺยสฺสาติ เมถุนมนุยุตฺตสฺส.}

\prose{\textbf{เมตฺเตยฺยา}ติ \textbf{ภควา} ตํ เถรํ โคตฺเตน อาลปติ.\csromanspacer{}\csromanspacer{}\textbf{ภควา}ติ คารวาธิวจนํ.\csromanspacer{} อปิ จ ภคฺคราโคติ \textbf{ภควา}, ภคฺคโทโสติ \textbf{ภควา}, ภคฺคโมโหติ \textbf{ภควา}, ภคฺคมาโนติ \textbf{ภควา}, ภคฺคิทิฏฺฐีติ \textbf{ภควา}, ภคฺคกณฺฑโกติ\footnote{ภคฺคกณโกติ (สี, สฺยา)} \textbf{ภควา}, ภคฺคกิเลโสติ \textbf{ภควา}, ภชิ วิภชิ ปวิภชิ ธมฺมรตนนฺติ \textbf{ภควา}, ภวานํ อนฺตกโรติ \textbf{ภควา}, ภาวิตกาโย, ภาวิตสีโล, ภาวิตจิตฺโต, ภาวิตปญฺโญติ \textbf{ภควา}, ภชิ วา \textbf{ภควา} อรญฺญวนปตฺถานิ\footnote{อรญฺเญ วนปตฺถานิ (สี)} ปนฺตานิ เสนาสนานิ อปฺปสทฺทานิ อปฺปนิคฺโฆสานิ}

\vspace{\tipitakaparskip}
\csromancenter{ติสฺสเมตฺเตยฺยสุตฺตนิทฺเทส วิชนวาตานิ มนุสฺสาราหสฺเสยฺยกานิ\footnote{มนุสฺสราหเสยฺยกานิ (สี, สฺยา)} ปฏิสลฺลานสารุปฺปานีติ ภควา, ภาคี วา ภควา จีวรปิณฺฑปาตเสนาสนคิลานปจฺจยเภสชฺชปริกฺขารานนฺติ ภควา, ภาคี วา ภควา อตฺถรสสฺส ธมฺมรสสฺส วิมุตฺติรสสฺส, อธีสีลสฺส อธิจิตฺตสฺส อธิปญฺญายาติ ภควา, ภาคี วา ภควา จตุนฺนํ ฌานานํ จตุนฺนํ อปฺปมญฺญานํ จตุนฺนํ อรูปสมาปตฺตีนนฺติ ภควา, ภาคี วา ภควา อฏฺฐนฺนํ วิโมกฺขานํ อฏฺฐนฺนํ อภิภายตนานํ นวนฺนํ อนุปุพฺพวิหารสมาปตฺตีนนฺติ ภควา, ภาคี วา ภควา ทสนฺนํ สญฺญาภาวนานํ ทสนฺนํ กสิณสมาปตฺตีนํ อานาปานสฺสติสมาธิสฺส อสุภสมาปตฺติยาติ ภควา, ภาคี วา ภควา จตุนฺนํ สติปฏฺฐานานํ จตุนฺนํ สมฺมปฺปธานานํ จตุนฺนํ อิทฺธิปาทานํ ปญฺจนฺนํ อินฺทฺริยานํ ปญฺจนฺนํ พลานํ สตฺตนฺนํ โพชฺฌงฺคานํ อริยสฺส อฏฺฐงฺคิกสฺส มคฺคสฺสาติ ภควา, ภาคี วา ภควา ทสนฺนํ ตถาคตพลานํ จตุนฺนํ เวสารชฺชานํ จตุนฺนํ ปฏิสมฺภิทานํ ฉนฺนํ อภิญฺญานํ\footnote{อภิญฺญาณานํ (สี)} ฉนฺนํ พุทฺธธมฺมานนฺติ ภควา.\csromanspacer{} ภควาติ เนตํ นามํ มาตรา กตํ, น ปิตรา กตํ, น ภาตรา กตํ, น ภคินิยา กตํ, น มิตฺตามจฺเจหิ กตํ, น ญาติสาโลหิเตหิ กตํ, น สมณพฺราหฺมเณปิ กตํ, น เทวตาหิ กตํ, วิโมกฺขนฺติกเมตํ พุทฺธานํ ภควนฺตานํ โพธิยา มูเล สห สพฺพญฺญุตญฺญาณสฺส ปฏิลาภา สจฺฉิกา ปญฺญตฺติ ยทิทํ ภควาติ เมตฺเตยฺยาติ ภควา.}
\prose{\csromanfolio{111}\textbf{มุสฺสเต วาปิ สาสนนฺ}ติ ทฺวีหิ การเณหิ สาสนํ มุสฺสติ ปริยตฺติสาสนมฺปิ มุสฺสติ, ปฏิปตฺติสาสนมฺปิ มุสฺสติ.\csromanspacer{} กตมํ ปริยตฺติสาสนํ, ยํ ตสฺส ปริยาปุฏํ สุตฺตํ เคยฺยํ เวยฺยากรณํ คาถา อุทานํ อิติวุตฺตกํ ชาตกํ อพฺภุตธมฺมํ เวทลฺลํ, อิทํ ปริยตฺติสาสนํ.\csromanspacer{} ตมฺปิ มุสฺสติ สมฺมุสฺสติ สมฺปมุสฺสติ ปริพาหิโร โหตีติ เอวมฺปิ มุสฺสเต วาปิ สาสนํ.}

\prose{กตมํ ปฏิปตฺติสาสนํ.\csromanspacer{} สมฺมาปฏิปทา อนุโลมปฏิปทา อปจฺจนีกปฏิปทา อนฺวตฺถปฏิปทา ธมฺมานุธมฺมปฏิปทา สีเลสุ ปริปูรการิตา อินฺทฺริเยสุ คุตฺตทฺวารตา โภชเน มตฺตญฺญุตา ชาคริยานุโยโค สติสมฺปชญฺญํ จตฺตาโร สติปฏฺฐานา จตฺตาโร สมฺมปฺปธานา จตฺตาโร อิทฺธิปาทา ปญฺจินฺทฺริยานิ ปญฺจ พลานิ สตฺต โพชฺฌงฺคา อริโย อฏฺฐงฺคิโก มคฺโค, อิทํ ปฏิปตฺติสาสนํ.\csromanspacer{} ตมฺปิ มุสฺสติ สมฺมุสฺสติ ปมุสฺสติ สมฺปมุสฺสติ ปริพาหิโร โหตีติ เอวมฺปิ มุสฺสเต วาปิ สาสนํ.}

\prose{\csromanfolio{112}\textbf{มิจฺฉา จ ปฏิปชฺชตี}ติ ปาณมฺปิ หนติ, อทินฺนมฺปิ อาทิยติ, สนฺธิมฺปิ ฉินฺทติ, นิลฺโลปมฺปิ หรติ, เอกาคาริกมฺปิ กโรติ, ปริปนฺเถปิ ติฏฺฐติ, ปรทารมฺปิ คจฺฉติ, มุสาปิ ภณตีติ มิจฺฉา จ ปฏิปชฺชติ.}

\prose{\textbf{เอตํ ตสฺมิํ อนาริยนฺ}ติ เอตํ ตสฺมิํ ปุคฺคเล อนริยธมฺโม พาลธมฺโม มูฬฺหธมฺโม อญฺญาณธมฺโม อมราวิกฺเขปธมฺโม ยทิทํ มิจฺฉาปฏิปทาติ เอตํ ตสฺมิํ อนาริยํ.\csromanspacer{} เตนาห ภควา\csromandash{}}

\prose{เมถุนมนุยุตฺตสฺส, (เมตฺเตยฺยาติ ภควา,)}

\prose{มุสฺสเต วาปิ สาสนํ.}

\csromangathameasure{มิจฺฉา จ ปฏิปชฺชติ, \\ \textbf{เอตํ ตสฺมิํ อนาริยนฺ}ติ.}
\csromangathagroup{\csromangathastanza{มิจฺฉา จ ปฏิปชฺชติ, \\ \textbf{เอตํ ตสฺมิํ อนาริยนฺ}ติ.}}

\proseitem{51}{\csromansymbolfootnote{*}{ขุ 1. 407 ปิฏฺเฐปิ.}เอโก ปุพฺเพ จริตฺวาน, เมถุนํ โย นิเสวติ.}

\csromangathameasure{ยานํ ภนฺตํว ตํ โลเก, หีนมาหุ ปุถุชฺชนํ.}
\csromangathagroup{\csromangathastanza{ยานํ ภนฺตํว ตํ โลเก, หีนมาหุ ปุถุชฺชนํ.}}

\prose{\textbf{เอโก ปุพฺเพ จริตฺวานา}ติ ทฺวีหิ การเณหิ เอโก ปุพฺเพ จริตฺวาน ปพฺพชฺชาสงฺขาเตน วา คณาววสฺสคฺคฏฺเฐน วา.\csromanspacer{} กถํ ปพฺพชฺชาสงฺขาเตน เอโก ปุพฺเพ จริตฺวาน.\csromanspacer{} สพฺพํ ฆราวาสปลิโพธํ ฉินฺทิตฺวา ปุตฺตทารปลิโพธํ ฉินฺทิตฺวา ญาติปลิโพธํ ฉินฺทิตฺวา มิตฺตามจฺจปลิโพธํ ฉินฺทิตฺวา สนฺนิธิปลิโพธํ ฉินฺทิตฺวา เกสมสฺสุํ โอหาเรตฺวา กาสายานิ วตฺถานิ อจฺฉาเทตฺวา อคารสฺมา อนคาริยํ ปพฺพชิตฺวา อกิญฺจนภาวํ อุปคนฺตฺวา เอโก จรติ วิหรติ อิริยติ วตฺตติ ปาเลติ ยเปติ ยาเปติ เอวํ ปพฺพชฺชาสงฺขาเตน เอโก ปุพฺเพ จริตฺวาน.}

\prose{กถํ คณาววสฺสคฺคฏฺเฐน เอโก ปุพฺเพ จริตฺวาน.\csromanspacer{} โส เอวํ ปพฺพชิโต สมาโน เอโก อรญฺญวนปตฺถานิ ปนฺตานิ เสนาสนานิ ปฏิเสวติ อปฺปสทฺทานิ อปฺปนิคฺโฆสานิ วิชนวาตานิ มนุสฺสราหสฺเสยฺยกานิ ปฏิสลฺลานสารุปฺปานิ.\csromanspacer{} โส เอโก คจฺฉติ, เอโก ติฏฺฐติ, เอโก นิสีทติ, เอโก เสยฺยํ กปฺเปติ, เอโก คามํ ปิณฺฑาย ปวิสติ, เอโก ปฏิกฺกมติ, เอโก รโห นิสีทติ, เอโก จงฺกมํ อธิฏฺฐาติ, เอโก จรติ วิหรติ อิริยติ วตฺตติ ปาเลติ ยเปติ ยาเปติ, เอวํ คณาววสฺสคฺคฏฺเฐน เอโก ปุพฺเพ จริตฺวาน.}

\prose{\textbf{เมถุนํ โย นิเสวตี}ติ \textbf{เมถุนธมฺโม} นาม โย โส อสทฺธมฺโม ฯเปฯ ตํการณา วุจฺจติ \textbf{เมถุนธมฺโม}.\csromanspacer{}\csromanspacer{}\textbf{เมถุนํ โย}}

\vspace{\tipitakaparskip}
\csromancenter{ติสฺสเมตฺเตยฺยสุตฺตนิทฺเทส \textbf{นีเสวตี}ติ โย อปเรน สมเยน พุทฺธํ ธมฺมํ สํฆํ สิกฺขํ ปจฺจกฺขาย หีนายาวตฺติตฺวา เมถุนํ ธมฺมํ เสวติ นิเสวติ สํเสวติ ปฏิเสวตีติ เมถุนํ โย เนเสวติ.}
\prose{\csromanfolio{113}\textbf{ยานํ ภนฺตํว ตํ โลเก}ติ \textbf{ยานนฺ}ติ หตฺถิยานํ อสฺสยานํ โคยานํ อชยานํ เมณฺฑยานํ โอฏฺฐยานํ ขรยานํ ภนฺตํ อทนฺตํ อการิตํ อวินีตํ อุปฺปถํ คณฺหาติ, วิสมํ ขาณุมฺปิ ปาสาณมฺปิ อภิรุหติ, ยานมฺปิ อาโรหนกมฺปิ ภญฺชติ, ปปาเตปิ ปปตติ.\csromanspacer{}\csromanspacer{}ยถา ตํ ภนฺตํ ยานํ อทนฺตํ อการิตํ อวินีตํ อุปฺปถํ คณฺหาติ, เอวเมวํ โส วิพฺภนฺตโก ภนฺตยานปฏิภาโค อุปฺปถํ คณฺหาติ, มิจฺฉาทิฏฺฐิํ คณฺหาติ ฯเปฯ มิจฺฉาสมาธิํ คณฺหาติ.\csromanspacer{}\csromanspacer{}ยถา ตํ ภนฺตํ ยานํ อทนฺตํ อการิตํ อวินีตํ วิสมํ ขาณุมฺปิ ปาสาณมฺปิ อภิรุหติ, เอวเมตํ โส วิพฺภนฺตโก ภนฺตยานปฏิภาโค วิสมํ กายกมฺมํ อภิรุหติ, วิสมํ วจีกมฺมํ อภิรุหติ, วิสมํ มโนกมฺมํ อภิรุหติ, วิสมํ ปาณาติปาตํ อภิรุหติ, วิสมํ อทินฺนาทานํ อภิรุหติ, วิสมํ กาเมสุมิจฺฉาจารํ อภิรุหติ, วิสมํ มุสาวาทํ อภิรุหติ, วิสมํ ปิสุณวาจํ อภิรุหติ, วิสมํ ผรุสวาจํ อภิรุหติ, วิสมํ สมฺผปฺปลาปํ อภิรุหติ, วิสมํ อภิชฺฌํ อภิรุหติ, วิสมํ พฺยาปาทํ อภิรุหติ, วิสมํ มิจฺฉาทิฏฺฐิํ อภิรุหติ, วิสเม สงฺขาเร อภิรุหติ, วิสเม ปญฺจ กามคุเณ อภิรุหติ, วิสเม นีวรเณ อภิรุหติ.\csromanspacer{}\csromanspacer{}ยถา ตํ ภนฺตํ ยานํ อทนฺตํ อการิตํ อวินีตํ ยานมฺปิ อาโรหนกมฺปิ ภญฺชติ, เอวเมวํ โส วิพฺภนฺตโก ภนฺตยานปฏิภาโค นิรเย อตฺตานํ ภญฺชติ, ติรจฺฉานโยนิยํ อตฺตานํ ภญฺชติ, เปตฺติวิสเย อตฺตานํ ภญฺชติ, มนุสฺสโลเก อตฺตานํ ภญฺชติ, เทวโลเก อตฺตานํ ภญฺชติ.\csromanspacer{}\csromanspacer{}ยถา ตํ ภนฺตํ ยานํ อทนฺตํ อการิตํ อวินีตํ ปปาเต ปปตติ, เอวเมวํ โส วิพฺภนฺตโก ภนฺตยานปฏิภาโค ชาติปปาตมฺปิ ปปตติ, ชราปปาตมฺปิ ปปตติ, พฺยาธิปปาตมฺปิ ปปตติ, มรณปปาตมฺปิ ปปตติ, โสกปริเทวทุกฺขโทมนสฺสุปายาสปปาตมฺปิ ปปตติ.\csromanspacer{} \textbf{โลเก}ติ อปายโลเก มนุสฺสโลเกติ ยานํ ภนฺตํว ตํ โลเก.}

\prose{\textbf{หีนมาหุ ปุถุชฺชน}นฺติ \textbf{ปุถุชฺชนา}ติ เกนฏฺเฐน \textbf{ปุถุชฺชนา}.\csromanspacer{} ปุถุกิเลเส ชเนนฺตีติ \textbf{ปุถุชฺชนา}, ปุถุ อวิหตสกฺกายทิฏฺฐิกาติ \textbf{ปุถุชฺชนา}, ปุถุ สตฺถารานํ มุขุลฺโลกิกาติ \textbf{ปุถุชฺชนา}, ปุถุ สพฺพคตีหิ อวุฏฺฐิตาติ \csromanfolio{114}ปุถุชฺชนา, ปุถุ นานาภิสงฺขาเร\footnote{นานาภิสงฺขาเรหิ (สฺยา) อุปริ 192 ปิฏฺเฐปิ.} อภิสงฺขโรนฺตีติ ปุถุชฺชนา, ปุถุ นานา-โอเฆหิ วุยฺหนฺตีติ ปุถุชฺชนา, ปุถุ นานาสนฺตาเปหิ สนฺตปนฺตีติ ปุถุชฺชนา, ปุถุ นานาปริฬาเหหิ ปริทยฺหนฺตีติ ปุถุชฺชนา, ปุถุ ปญฺจสุ กามคุเณสุ รตฺตา คิทฺธา คธิตา มุจฺฉิตา อชฺโฌสนฺนา\footnote{อชฺโฌปนฺนา (สี, สฺยา)} ลคฺคา ลคฺคิตา ปลิพุทฺธาติ ปุถุชฺชนา, ปุถุ ปญฺจหิ นีวรเณหิ อาวุตา นิวุตา โอวุตา ปิหิตา ปฏิจฺฉนฺนา ปฏิกุชฺชิตาติ ปุถุชฺชนา.\csromanspacer{}\csromanspacer{}\textbf{หีนมาหุ ปุถุชฺชนนฺ}ติ ปุถุชฺชนํ หีนํ นิหีนํ โอมกํ ลามกํ ฉตุกฺกํ ปริตฺตนฺติ เอวมาหํสุ เอวํ กเถนฺติ เอวํ ภณนฺติ เอวํ ทีปยนฺติ เอวํ โวหรนฺตีติ หีนมาหุ ปุถุชฺชนํ.\csromanspacer{} เตนาห ภควา\csromandash{}}

\csromangathameasure{เอโก ปุพฺเพ จริตฺวาน, เมถุนํ โย นิเสวติ. \\ ยานํ ภนฺตํว ตํ โลเก, หีนมาหุ ปุถุชฺชนนฺติ.}
\csromangathagroup{\csromangathastanza{เอโก ปุพฺเพ จริตฺวาน, เมถุนํ โย นิเสวติ. \\ ยานํ ภนฺตํว ตํ โลเก, หีนมาหุ ปุถุชฺชนนฺติ.}}

\proseitem{52}{\csromansymbolfootnote{*}{ขุ 1. 407 ปิฏฺเฐ.}\textbf{ยโส กิตฺติ จ ยา ปุพฺเพ, หายเต วาปิ ตสฺส สา}.}

\csromangathameasure{เอตมฺปิ ทิสฺวา สิกฺเขถ, เมถุนํ วิปฺปหาตเว.}
\csromangathagroup{\csromangathastanza{เอตมฺปิ ทิสฺวา สิกฺเขถ, เมถุนํ วิปฺปหาตเว.}}

\prose{\textbf{ยโส กิตฺติ จ ยา ปุพฺเพ, หายเต วาปิ ตสฺส สา}ติ กตโม \textbf{ยโส.}\csromanspacer{} อิเธกจฺโจ ปุพฺเพ สมณภาเว สกฺกโต โหติ ครุกโต มานิโต ปูชิโต อปจิโต, ลาภี จีวรปิณฺฑปาตเสนาสนคิลานปจฺจยเภสชฺชปริกฺขารานํ, อยํ \textbf{ยโส.}\csromanspacer{}\csromanspacer{}กตมา \textbf{กิตฺติ} อิเธกจฺโจ ปุพฺเพ สมณภาเว \textbf{กิตฺติ}วณฺณคโต โหติ, ปณฺฑิโต วิยตฺโต เมธาวี พหุสฺสุโต จิตฺตกถี กลฺยาณปฏิภาโน “สุตฺตนฺติโก”ติ วา “วินยธโร”ติ วา “ธมฺมกถิโก”ติ วา “อารญฺญิโก”ติ วา “ปิณฺฑปาติโก”ติ วา “ปํสุกูลิโก”ติ วา “เตจีวริโก”ติ วา “สปทานจาริโก”ติ วา”ขลุปจฺฉาภตฺติโก”ติ วา “เนสชฺชิโก”ติ วา “ยถาสนฺถติโก”ติ วา “ปฐมสฺส ฌานสฺส ลาภี”ติ วา “ทุติยสฺส ฌานสฺส ลาภี”ติ วา “ตติยสฺส ฌานสฺส ลาภี”ติ วา “จตุตฺถสฺส ฌานสฺส ลาภี”ติ วา “อากาสานญฺจายตนสมาปตฺติยา ลาภี”ติ วา “วิญฺญาณญฺจายตนสมาปตฺติยา ลาภี”ติ วา “อากิญฺจญฺญายตนสมาปตฺติยา ลาภี”ติ วา “เนวสญฺญานาสญฺญายตนสมาปตฺติยา ลาภี”ติ วา, อยํ กิตฺตีติ ยโส \textbf{กิตฺติ} จ ยา ปุพฺเพ.}

\csromanheader{2}{7. ติสฺสเมตฺเตยฺยสุตฺตนิทฺเทส}
\prose{\csromanfolio{115}\textbf{หายเต วาปิ ตสฺส สา}ติ ตสฺส อปเรน สมเยน พุทฺธํ ธมฺมํ สํฆํ สิกฺขํ ปจฺจกฺขาย หีนายาวตฺตสฺส โส จ ยโส สา จ กิตฺติ หายติ ปริหายติ ปริธํสติ ปริปตติ อนฺตรธายติ วิปฺปลุชฺชตีติ ยโส กิตฺติ จ ยา ปุพฺเพ, หายเต วาปิ ตสฺส สา.}

\prose{\textbf{เอตมฺปิ ทิสฺวา สิกฺเขถ, เมถุนํ วิปฺปหาตเว}ติ \textbf{เอตนฺ}ติ ปุพฺเพ สมณภาเว ยโส กิตฺติ จ, อปรภาเค พุทฺธํ ธมฺมํ สํฆํ สิกฺขํ ปจฺจกฺขาย หีนายาวตฺตสฺส อยโส จ อกิตฺติ จ, เอตํ สมฺปตฺติํ วิปตฺติํ.\csromanspacer{}\csromanspacer{}\textbf{ทิสฺวา}ติ ปสฺสิตฺวา ตุลยิตฺวา ตีรยิตฺวา วิภาวยิตฺวา วิภูตํ กตฺวาติ เอตมฺปิ ทิสฺวา.\csromanspacer{}\csromanspacer{}\textbf{สิกฺเขถา}ติ ติสฺโส สิกฺขา อธิสีลสิกฺขา อธิจิตฺตสิกฺขา อธิปญฺญาสิกฺขา.\csromanspacer{} กตมา อธิสีลสิกฺขา.\csromanspacer{}\csromanspacer{}อิธ ภิกฺขุ สีลวา โหติ ปาติโมกฺขสํวรสํวุโต วิหรติ อาจารโคจรสมฺปนฺโน อณุมตฺเตสุ วชฺเชสุ ภยทสฺสาวี สมาทาย สิกฺขติ สิกฺขาปเทสุ.\csromanspacer{} ขุทฺทโก สีลกฺขนฺโธ, มหนฺโต สีลกฺขนฺโธ, สีลํ ปติฏฺฐา อาทิ จรณํ สํยโม สํวโร มุขํ ปมุขํ กุสลานํ ธมฺมานํ สมาปตฺติยา.\csromanspacer{} อยํ อธิสีลสิกฺขา.}

\prose{กตมา อธิจิตฺตสิกฺขา.\csromanspacer{} อิธ ภิกฺขุ วิวิจฺเจว กาเมหิ วิวิจฺจ อกุสเลหิ ธมฺเมหิ สวิตกฺกํ สวิจารํ วิเวกชํ ปีติสุขํ ปฐมํ ฌานํ อุปสมฺปชฺช วิหรติ ฯเปฯ ทุติยํ ฌานํ.\csromanspacer{} ตติยํ ฌานํ.\csromanspacer{} จตุตฺถํ ฌานํ อุปสมฺปชฺช วิหรติ.\csromanspacer{} อยํ อธิจิตฺตสิกฺขา.}

\prose{กตมา อธิปญฺญาสิกฺขา.\csromanspacer{} อิธ ภิกฺขุ ปญฺญวา โหติ อุทยตฺถคามินิยา ปญฺญาย สมนฺนาคโต อริยาย นิพฺเพธิกาย สมฺมา ทุกฺขกฺขยคามินิยา.\csromanspacer{} โส “อิทํ ทุกฺขนฺ”ติ ยถาภูตํ ปชานาติ, “อยํ ทุกฺข สมุทโย”ติ ยถาภูตํ ปชานาติ, “อยํ ทุกฺขนิโรโธ”ติ ยถาภูตํ ปชานาติ,”อยํ ทุกฺขนิโรธคามินี ปฏิปทา”ติ ยถาภูตํ ปชานาติ, “อิเม อาสวา”ติ ยถาภูตํ ปชานาติ, “อยํ อาสวสมุทโย”ติ ยถาภูตํ ปชานาติ, “อยํ อาสวนิโรโธ”ติ ยถาภูตํ ปชานาติ, “อยํ อาสวนิโรธคามินี ปฏิปทา”ติ ยถาภูตํ ปชานาติ.\csromanspacer{} อยํ อธิปญฺญาสิกฺขา.\csromanspacer{}\csromanspacer{}\textbf{เมถุนธมฺโม} นาม โย โส อสทฺธมฺโม ฯเปฯ ตํการณา วุจฺจติ เมถุนธมฺโม.}

\prose{\csromanfolio{116}\textbf{เอตมฺปิ ทิสฺวา สิกฺเขถ, เมถุนํ วิปฺปหาตเว}ติ เมถุนธมฺมสฺส ปหานาย วูปสมาย ปฏินิสฺสคฺคาย ปฏิปสฺสทฺธิยา อธิสีลมฺปิ สิกฺเขยฺย, อธิจิตฺตมฺปิ สิกฺเขยฺย, อธิปญฺญมฺปิ สิกฺเขยฺย.\csromanspacer{} อิมา ติสฺโส สิกฺขาโย อาวชฺชนฺโต สิกฺเขยฺย, ชานนฺโต สิกฺเขยฺย, ปสฺสนฺโต สิกฺเขยฺย, ปจฺจเวกฺขนฺโต สิกฺเขยฺย, จิตฺตํ อธิฏฺฐหนฺโต สิกฺเขยฺย, สทฺธาย อธิมุจฺจนฺโต สิกฺเขยฺย, วีริยํ ปคฺคณฺหนฺโต สิกฺเขยฺย, สติํ อุปฏฺฐเปนฺโต สิกฺเขยฺย, จิตฺตํ สมาทหนฺโต สิกฺเขยฺย, ปญฺญาย ปชานนฺโต สิกฺเขยฺย, อภิญฺเญยฺยํ อภิชานนฺโต สิกฺเขยฺย, ปริญฺเญยฺยํ ปริชานนฺโต สิกฺเขยฺย, ปหาตพฺพํ ปชหนฺโต สิกฺเขยฺย, ภาเวตพฺพํ ภาเวนฺโต สิกฺเขยฺย, สจฺฉิกาตพฺพํ สจฺฉิกโรนฺโต สิกฺเขยฺย อาจเรยฺย สมาจเรยฺย สมาทาย วตฺเตยฺยาติ เอตมฺปิ ทิสฺวา สิกฺเขถ, เมถุนํ วิปฺปหาตเว.\csromanspacer{} เตนาห ภควา\csromandash{}}

\csromangathameasure{ยโส กิตฺติ จ ยา ปุพฺเพ, หายเต วาปิ ตสฺส สา.}
\csromangathagroup{\csromangathastanza{ยโส กิตฺติ จ ยา ปุพฺเพ, หายเต วาปิ ตสฺส สา.}}

\prose{\textbf{เอตมฺปิ ทิสฺวา สิกฺเขถ, เมถุนํ วิปฺปหาตเว}ติ.}

\proseitem{53}{\csromansymbolfootnote{*}{ขุ 1. 407; ขุ 10. 173 ปิฏฺเฐสุ.}สงฺกปฺเปหิ ปเรโต โส, กปโณ วิย ฌายติ.}

\csromangathameasure{สุตฺวา ปเรสํ นิคฺโฆสํ, มงฺกุ โหติ ตถาวิโธ.}
\csromangathagroup{\csromangathastanza{สุตฺวา ปเรสํ นิคฺโฆสํ, มงฺกุ โหติ ตถาวิโธ.}}

\prose{\textbf{สงฺกปฺเปหิ ปเรโต โส, กปโณ วิย ฌายตี}ติ กามสงฺกปฺเปน พฺยาปาทสงฺกปฺเปน วิหิํสาสงฺกปฺเปน ทิฏฺฐิสงฺกปฺเปน ผุฏฺโฐ ปเรโต สโมหิโต สมนฺนาคโต ปิหิโต กปโณ วิย มนฺโท วิย โมมูโห วิย ฌายติ ปชฺฌายติ นิชฺฌายติ อปชฺฌายติ\footnote{อวชฺฌายติ (สฺยา)}.\csromanspacer{} ยถา อุลูโก รุกฺขสาขายํ มูสิกํ มคยมาโน ฌายติ ปชฺฌายติ นิชฺฌายติ อปชฺฌายติ.\csromanspacer{} ยถา โกตฺถุ นทีตีเร มจฺเฉ มคยมาโน ฌายติ ปชฺฌายติ นิชฺฌายติ อปชฺฌายติ.\csromanspacer{} ยถา พิฬาโร สนฺธิสมลสงฺกฏิเร มูสิกํ มคยมาโน ฌายติ ปชฺฌายติ นิชฺฌายติ อปชฺฌายติ.\csromanspacer{} ยถา คทฺรโภ วหจฺฉินฺโน สนฺธิสมลสงฺกฏิเร ฌายติ ปชฺฌายติ นิชฺฌายติ อปชฺฌายติ, เอวเมวํ โส วิพฺภนฺตโก กามสงฺกปฺเปน พฺยาปาทสงฺกปฺเปน วิหิํสาสงฺกปฺเปน ทิฏฺฐิสงฺกปฺเปน ผุฏฺโฐ ปเรโต สโมหิโต สมนฺนาคโต ปิหิโต กปโณ วิย มนฺโท วิย โมมูโห วิย ฌายติ ปชฺฌายติ นิชฺฌายติ อปชฺฌายตีติ สงฺกปฺเปติ ปเรโต โส, กปโณ วิย ฌายติ.}

\csromanheader{2}{7. ติสฺสเมตฺเตยฺยสุตฺตนิทฺเทส}
\prose{\csromanfolio{117}\textbf{สุตฺวา ปเรสํ นิคฺโฆสํ, มงฺกุ โหติ ตถาวิโธ}ติ \textbf{ปเรสนฺ}ติ อุปชฺฌายา วา อาจริยา วา สมานุปชฺฌายกา วา สมานาจริยกา วา มิตฺตา วา สนฺทิฏฺฐา วา สมฺภตฺตา วา สหายา วา โจเทนฺติ “ตสฺส เต อาวุโส อลาภา, ตสฺส เต ทุลฺลทฺธํ, ยํ ตฺวํ เอวรูปํ อุฬารํ สตฺถารํ ลภิตฺวา เอวํ สฺวากฺขาเต ธมฺมวินเย ปพฺพชิตฺวา เอวรูปํ อริยคณํ ลภิตฺวา หีนสฺส เมถุนธมฺมสฺส การณา พุทฺธํ ธมฺมํ สํฆํ สิกฺขํ ปจฺจกฺขาย หีนายาวตฺโตสิ, สทฺธาปิ นาม เต นาโหสิ กุสเลสุ ธมฺเมสุ, หิรีปิ นาม เต นาโหสิ กุสเลสุ ธมฺเมสุ, โอตฺตปฺปมฺปิ นาม เต นาโหสิ กุสเลสุ ธมฺเมสุ, วีริยมฺปิ นาม เต นาโหสิ กุสเลสุ ธมฺเมสุ, สติปิ นาม เต นาโหสิ กุสเลสุ ธมฺเมสุ, ปญฺญาปิ นาม เต นาโหสิ กุสเลสุ ธมฺเมสู”ติ.\csromanspacer{} เตสํ วจนํ พฺยปฺปถํ เทสนํ อนุสาสนํ อนุสิฏฺฐิํ สุตฺวา สุณิตฺวา อุคฺคเหตฺวา อุปธารยิตฺวา อุปลกฺขยิตฺวา มงฺกุ โหติ, ปีฬิโต ฆฏฺฏิโต พฺยาธิโต โทมนสฺสิโต โหติ.\csromanspacer{}\csromanspacer{}\textbf{ตถาวิโธ}ติ ตฺ\textbf{อถา}วิโธ ตาทิโส ตสฺสณฺฐิโต ตปฺปกาโร ตปฺปฏิภาโค โย โส วิพฺภนฺตโกติ สุตฺวา ปเรสํ นิคฺโฆสํ มงฺกุ โหติ ตฺ\textbf{อถา}วิโธ.\csromanspacer{} เตนาห ภควา\csromandash{}}

\csromangathameasure{สงฺกปฺเปหิ ปเรโต โส, กปโณ วิย ฌายติ.}
\csromangathagroup{\csromangathastanza{สงฺกปฺเปหิ ปเรโต โส, กปโณ วิย ฌายติ.}}

\prose{\textbf{สุตฺวา ปเรสํ นิคฺโฆสํ, มงฺกุ โหติ ตถาวิโธ}ติ.}

\proseitem{54}{\csromansymbolfootnote{*}{ขุ 1. 407 ปิฏฺเฐ.}\textbf{อถ สตฺถานิ กุรุเต, ปรวาเทหิ โจทิโต}.}

\csromangathameasure{เอส ขฺวสฺส มหาเคโธ, โมสวชฺชํ ปคาหติ.}
\csromangathagroup{\csromangathastanza{เอส ขฺวสฺส มหาเคโธ, โมสวชฺชํ ปคาหติ\footnote{สํคาหติ (ก)}.}}

\prose{\textbf{อถ สตฺถานิ กุรุเต, ปรวาเทหิ โจทิโต}ติ \textbf{อถา}ติ ปทสนฺธี ปทสํสคฺโค ปทปาริปูรี อกฺขรสมวาโย พฺยญฺชนสิลิฏฺฐตา ปทานุปุพฺพตาเปตํ \textbf{อถา}ติ.\csromanspacer{}\csromanspacer{}\textbf{สตฺถานี}ติ ตีณิ สตฺถานิ กายสตฺถํ วจีสตฺถํ มโนสตฺถํ.\csromanspacer{} ติวิธํ กายทุจฺจริตํ กายสตฺถํ, จตุพฺพิธํ วจีทุจฺจริตํ วจีสตฺถํ, ติวิธํ มโนทุจฺจริตํ มโนสตฺถํ.\csromanspacer{}\csromanspacer{}\textbf{ปรวาเทหิ โจทิโต}ติ อุปชฺฌาเยหิ วา อาจริเยหิ วา สมานุปชฺฌายเกหิ วา สมานาจริยเกหิ วา มิตฺเตหิ วา สนฺทิฏฺเฐหิ วา สมฺภตฺเตหิ วา สหาเยหิ วา โจทิโต สมฺปชานมุสา ภาสติ “อภิรโต อหํ ภนฺเต อโหสิํ ปพฺพชฺชาย, มาตา เม โปเสตพฺพา, เตนมฺหิ วิพฺภนฺโต”ติ \csromanfolio{118}ภณติ. “ปิตา เม โปเสตพฺโพ, เตนมฺหิ วิพฺภนฺโต”ติ ภณติ. “ภาตา เม โปเสตพฺโพ.\csromanspacer{} ภคินี เม โปเสตพฺพา.\csromanspacer{} ปุตฺโต เม โปเสตพฺโพ.\csromanspacer{} ธีตา เม โปเสตพฺพา.\csromanspacer{} มิตฺตา เม โปเสตพฺพา.\csromanspacer{} อมจฺจา เม โปเสตพฺพา.\csromanspacer{} ญาตกา เม โปเสตพฺพา.\csromanspacer{} สาโลหิตา เม โปเสตพฺพา, เตนมฺหิ วิพฺภนฺโต”ติ ภณติ วจีสตฺถํ กโรติ สงฺกโรติ ชเนติ สญฺชเนติ นิพฺพตฺเตติ อภินิพฺพตฺเตตีติ อถ สตฺถานิ กุรุเต, ปรวาเทหิ โจทิโต.}

\prose{\textbf{เอส ขฺวสฺส มหาเคโธ}ติ เอโส ตสฺส มหาเคโธ มหาวนํ มหาคหนํ มหากนฺตาโร มหาวิสโม มหากุฏิโล มหาปงฺโก มหาปลิโป มหาปลิโพโธ มหาพนฺธนํ ยทิทํ สมฺปชานมุสาวาโทติ เอส ขฺวสฺส มหาเคโธ.}

\prose{\textbf{โมสวชฺชํ ปคาหตี}ติ โมสวชฺชํ วุจฺจติ มุสาวาโท.\csromanspacer{} อิเธกจฺโจ สภคฺคโต วา ปริสคฺคโต วา ญาติมชฺฌคโต วา ปูคมชฺฌคโต วา ราชกุลมชฺฌคโต วา อภินีโต สกฺขิปุฏฺโฐ “เอหมฺโภ ปุริส ยํ ชานาสิ, ตํ วเทหี”ติ โส อชานํ วา อาห “ชานามี”ติ, ชานํ วา อาห “น ชานามี”ติ, อปสฺสํ วา อาห “ปสฺสามี”ติ, ปสฺสํ วา อาห “น ปสฺสามี”ติ.\csromanspacer{} อิติ อตฺตเหตุ วา ปรเหตุ วา อามิสกิญฺจิกฺขเหตุ วา สมฺปชานมุสา ภาสติ, อิทํ วุจฺจติ โมสวชฺชํ.}

\prose{อปิ จ ตีหากาเรหิ มุสาวาโท โหติ, * ปุพฺเพวสฺส โหติ “มุสา ภณิสฺสนฺ”ติ, ภณนฺตสฺส โหติ “มุสา ภณามี”ติ, ภณิตสฺส โหติ “มุสา มยา ภณิตนฺ”ติ อิเมหิ ตีหากาเรหิ มุสาวาโท โหติ.\csromanspacer{}\csromanspacer{}อปิ จ จตูหากาเรหิ มุสาวาโท โหติ, ปุพฺเพวสฺส โหติ “มุสา ภณิสฺสนฺ”ติ, ภณนฺตสฺส โหติ “มุสา ภณามี”ติ, ภณิตสฺส โหติ “มุสา มยา ภณิตนฺ”ติ วินิธาย ทิฏฺฐิํ.\csromanspacer{} อิเมหิ จตูหากาเรหิ มุสาวาโท โหติ.\csromanspacer{}\csromanspacer{}อปิ จ ปญฺจหากาเรหิ.\csromanspacer{} ฉหากาเรหิ.\csromanspacer{} สตฺตหากาเรหิ.\csromanspacer{} อฏฺฐหากาเรหิ มุสาวาโท โหติ, ปุพฺเพวสฺส โหติ “มุสา ภณิสฺสนฺ”ติ, ภณนฺตสฺส โหติ “มุสา ภณามี”ติ, ภณิตสฺส โหติ “มุสา มยา ภณิตนฺ”ติ วินิธาย ทิฏฺฐิํ, วินิธาย ขนฺติํ, วินิธาย รุจิํ, วินิธาย สญฺญํ, วินิธาย ภาวํ.\csromanspacer{} อิเมหิ อฏฺฐหากาเรหิ มุสาวาโท โหติ.\csromanspacer{}\csromanspacer{}\textbf{โมสวชฺชํ ปคาหตี}ติ}

\vspace{\tipitakaparskip}
\csromancenter{ติสฺสเมตฺเตยฺยสุตฺตนิทฺเทส โมสวชฺชํ ปคาหติ โอคาหติ อชฺโฌคาหติ ปวิสตีติ โมสวชฺชํ ปคาหติ.\csromanspacer{} เตนาห ภควา\csromandash{}}
\vspace{0.5\baselineskip}
\csromangathameasure{อถ สตฺถานิ กุรุเต, ปรวาเทหิ โจทิโต. \\ เอส ขฺวสฺส มหาเคโธ, โมสวชฺชํ ปคาหตีติ.}
\csromangathagroup{\csromangathastanza{\csromanfolio{119}อถ สตฺถานิ กุรุเต, ปรวาเทหิ โจทิโต. \\ เอส ขฺวสฺส มหาเคโธ, โมสวชฺชํ ปคาหตีติ.}}

\proseitem{55}{\csromansymbolfootnote{*}{ขุ 1. 407 ปิฏฺเฐ.}\textbf{ปณฺฑิโตติ สมญฺญาโต}, เอกจริยํ\footnote{เอกจฺจริยํ (ก)} \textbf{อธิฏฺฐิโต.}}

\csromangathameasure{ส จาปิ เมถุเน ยุตฺโต, มนฺโทว ปริกิสฺสติ.}
\csromangathagroup{\csromangathastanza{ส จาปิ เมถุเน ยุตฺโต, มนฺโทว ปริกิสฺสติ.}}

\prose{\textbf{ปณฺฑิโตติ สมญฺญาโต}ติ อิเธกจฺโจ ปุพฺเพ สมณภาเว กิตฺติ วณฺณคโต โหติ ปณฺฑิโต วิยตฺโต เมธาวี พหุสฺสุโต จิตฺตกถี กลฺยาณปฏิภาโน “สุตฺตนฺติโก”ติ วา “วินยธโร”ติ วา “ธมฺมกถิโก”ติ วา ฯเปฯ “เนวสญฺญานาสญฺญายตนสมาปตฺติยา ลาภี”ติ วา เอวํ ญาโต โหติ ปญฺญาโต สมญฺญาโต โหตีติ ปณฺฑิโตติ สมญฺญาโต.}

\prose{\textbf{เอกจริยํ อธิฏฺฐิโต}ติ ทฺวีหิ การเณหิ เอกจริยํ อธิฏฺฐิโต ปพฺพชฺชาสงฺขาเตน คณาววสฺสคฺคฏฺเฐน วา.\csromanspacer{} กถํ ปพฺพชฺชาสงฺขาเตน เอกจริยํ \textbf{อธิฏฺฐิโต.}\csromanspacer{} สพฺพํ ฆราวาสปลิโพธํ ฉินฺทิตฺวา ฯเปฯ เอวํ ปพฺพชฺชาสงฺขาเตน เอกจริยํ \textbf{อธิฏฺฐิโต.}\csromanspacer{} กถํ คณาววสฺสคฺคฏฺเฐน เอกจริยํ \textbf{อธิฏฺฐิโต.}\csromanspacer{} โส เอวํ ปพฺพชิโต สมาโน เอโก อรญฺญวนปตฺถานิ ปนฺตานิ ฯเปฯ เอวํ คณาววสฺสคฺคฏฺเฐน เอกจริยํ อธิฏฺฐิโตติ เอกจริยํ \textbf{อธิฏฺฐิโต.}}

\prose{\textbf{ส จาปิ เมถุเน ยุตฺโต}ติ \textbf{เมถุนธมฺโม} นาม โย โส อสทฺธมฺโม คามธมฺโม ฯเปฯ ตํการณา วุจฺจติ \textbf{เมถุนธมฺโม}.\csromanspacer{}\csromanspacer{}\textbf{ส จาปิ เมถุเน ยุตฺโต}ติ โส อปเรน สมเยน พุทฺธํ ธมฺมํ สํฆํ สิกฺขํ ปจฺจกฺขาย หีนายาวตฺติตฺวา เมถุนธมฺเม ยุตฺโต\footnote{ยุตฺโต สํยุตฺโต (สี)} ปยุตฺโต อายุตฺโต สมายุตฺโตติ ส จาปิ เมถุเน ยุตฺโต.}

\prose{\textbf{มนฺโทว ปริกิสฺสตี}ติ กปโณ วิย มนฺโท วิย โมมูโห วิย กิสฺสติ ปริกิสฺสติ ปริกิลิสฺสติ ปาณมฺปิ หนติ, อทินฺนมฺปิ อาทิยติ, สนฺธิมฺปิ ฉินฺทติ, นิลฺโลปมฺปิ หรติ, เอกาคาริกมฺปิ กโรติ, ปริปนฺเถปิ ติฏฺฐติ, ปรทารมฺปิ คจฺฉติ, มุสาปิ ภณติ.\csromanspacer{} เอวมฺปิ กิสฺสติ ปริกิสฺสติ ปริกิลิสฺสติ.\csromanspacer{} ตเมนํ ราชาโน คเหตฺวา วิวิธา กมฺมการณา \csromanfolio{120}กาเรนฺติ, * กสาหิปิ ตาเฬนฺติ, เวตฺเตหิปิ ตาเฬนฺติ, อทฺธทณฺฑเกหิปิ ตาเฬนฺติ, หตฺถมฺปิ ฉินฺทินฺติ, ปาทมฺปิ ฉินฺทนฺติ, หตฺถปาทมฺปิ ฉินฺทนฺติ, กณฺณมฺปิ ฉินฺทนฺติ, นาสมฺปิ ฉินฺทนฺติ, กณฺณนาสมฺปิ ฉินฺทนฺติ, พิลงฺคถาลิกมฺปิ กโรนฺติ, สงฺขมุณฺฑิกมฺปิ กโรนฺติ, ราหุมุขมฺปิ กโรนฺติ, โชติมาลิกมฺปิ กโรนฺติ, หตฺถปชฺโชติกมฺปิ กโรนฺติ, เอรกวตฺติกมฺปิ กโรนฺติ, จีรกวาสิกมฺปิ กโรนฺติ, เอเณยฺยกมฺปิ กโรนฺติ, พฬิสมํสิกมฺปิ กโรนฺติ, กหาปณิกมฺปิ กโรนฺติ, ขาราปตจฺฉิกมฺปิ1 กรนฺติ, ปลิฆปริวตฺติกมฺปิ กโรนฺติ, ปลาลปีฐกมฺปิ กโรนฺติ, ตตฺเตนปิ เตเลน โอสิญฺจนฺติ, สุนเขหิปิ ขาทาเปนฺติ, ชีวนฺตมฺปิ สูเล อุตฺตาเสนฺติ, อสินาปิ สีสํ ฉินฺทนฺติ.\csromanspacer{} เอวมฺปิ กิสฺสติ ปริกิสฺสติ ปริกิลิสฺสติ.}

\prose{อถ วา กามตณฺหาย อภิภูโต ปริยาทินฺนจิตฺโต โภเค ปริเยสนฺโต นาวาย มหาสมุทฺทํ ปกฺขนฺทติ, สีตสฺส ปุรกฺขโต อุณฺหสฺส ปุรกฺขโต ฑํส มกส วาตาตป สรีสป สมฺผสฺเสหิ ปีฬิยมาโน ขุปฺปิปาสาย มิยฺยมาโน ติคุมฺภํ คจฺฉติ, ตกฺโกลํ คจฺฉติ, ตกฺกสีลํ คจฺฉติ, กาลมุขํ คจฺฉติ, ปุรปูรํ คจฺฉติ, เวสุงฺคํ คจฺฉติ, เวราปถํ คจฺฉติ, ชวํ คจฺฉติ, ตามลิํ\footnote{กมลิํ (สฺยา), ตํมลิํ (ก) 3. ปีนํ (สฺยา)} คจฺฉติ, วงฺคํ คจฺฉติ, เอฬพนฺธนํ คจฺฉติ, สุวณฺณกูฏํ คจฺฉติ, สุวณฺณภูมิํ คจฺฉติ, ตมฺพปาณิํ คจฺฉติ, สุปฺปารกํ คจฺฉติ, ภารุกจฺฉํ คจฺฉติ, สุรฏฺฐํ คจฺฉติ, ภงฺคโลกํ คจฺฉติ, ภงฺคณํ คจฺฉติ, ปรมภงฺคณํ คจฺฉติ, โยนํ คจฺฉติ, ปรมโยนํ3 คจฺฉนฺติ, วินกํ\footnote{นวกํ (สี)} คจฺฉติ, มูลปทํ คจฺฉติ, มรุกนฺตารํ คจฺฉติ, ชณฺณุปถํ คจฺฉติ, อชปถํ คจฺฉติ, เมณฺฑปถํ คจฺฉติ, สงฺกุปถํ คจฺฉติ, ฉตฺตปถํ คจฺฉติ, วํสปถํ คจฺฉติ, สกุณปถํ คจฺฉติ, มูสิกปถํ คจฺฉติ, ทริปถํ คจฺฉติ, เวตฺตาจารํ คจฺฉติ.\csromanspacer{} เอวมฺปิ กิสฺสติ ปริกิสฺสติ ปริกิลิสฺสติ.}

\prose{คเวสนฺโต น วินฺทติ, อลาภมูลกมฺปิ ทุกฺขํ โทมนสฺสํ ปฏิสํเวเทติ.\csromanspacer{} เอวมฺปิ กิสฺสติ ปริกิสฺสติ ปริกิลิสฺสติ.}

\prose{คเวสนฺโต วินฺทติ, ลทฺธาปิ อารกฺขมูลกมฺปิ ทุกฺขํ โทมนสฺสํ ปฏิสํเวเทติ, “กินฺติ เม โภเค เนว ราชาโน หเรยฺยุํ, น โจรา หเรยฺยุํ, น อคฺคิ ทเหยฺย, น อุทกํ วเหยฺย, น อปิยา ทายาทา หเรยฺยุนฺ”ติ ตสฺส เอวํ อารกฺขโต โคปยโต เต โภคา วิปฺปลุชฺชนฺติ, โส}

\vspace{\tipitakaparskip}
\csromancenter{ติสฺสเมตฺเตยฺยสุตฺตนิทฺเทส วิปฺปโยคมูลกมฺปิ ทุกฺขํ โทมนสฺสํ ปฏิสํเวเทติ.\csromanspacer{} เอวมฺปิ กิสฺสติ ปริกิสฺสติ ปริกิลิสฺสตีติ ส จาปิ เมถุเน ยุตฺโต มนฺโทว ปริกิสฺสติ.\csromanspacer{} เตนาห ภควา\csromandash{}}
\vspace{0.5\baselineskip}
\csromangathameasure{ปณฺฑิโตติ สมญฺญาโต, เอกจริยํ อธิฏฺฐิโต. \\ ส จาปิ เมถุเน ยุตฺโต, มนฺโทว ปริกิสฺสตีติ.}
\csromangathagroup{\csromangathastanza{\csromanfolio{121}ปณฺฑิโตติ สมญฺญาโต, เอกจริยํ อธิฏฺฐิโต. \\ ส จาปิ เมถุเน ยุตฺโต, มนฺโทว ปริกิสฺสตีติ.}}

\proseitem{56}{\csromansymbolfootnote{*}{ขุ 1. 407 ปิฏฺเฐ.}เอตมาทีนวํ ญตฺวา, มุนิ ปุพฺพาปเร อิธ.}

\csromangathameasure{เอกจริยํ ทฬฺหํ กยิรา, น นิเสเวถ เมถุนํ.}
\csromangathagroup{\csromangathastanza{เอกจริยํ ทฬฺหํ กยิรา, น นิเสเวถ เมถุนํ.}}

\prose{\textbf{เอตมาทีนวํ ญตฺวา, มุนิ ปุพฺพาปเร อิธา}ติ \textbf{เอตนฺ}ติ ปุพฺเพ สมณภาเว ยโส จ กิตฺติ จ.\csromanspacer{} อปรภาเค พุทฺธํ ธมฺมํ สํฆํ สิกฺขํ ปจฺจกฺขาย หีนายาวตฺตสฺส อยโส จ อกิตฺติ จ, เอตํ สมฺปตฺติํ วิปตฺติญฺจ.\csromanspacer{} \textbf{ญตฺวา}ติ ชานิตฺวา ตุลยิตฺวา ตีรยิตฺวา วิภาวยิตฺวา วิภูตํ กตฺวา.\csromanspacer{} \textbf{มุนี}ติ \textbf{โมนํ} วุจฺจติ ญาณํ.\csromanspacer{} ยา ปญฺญา ปชานนา ฯเปฯ สงฺคชาลมติจฺจ โส มุนิ.\csromanspacer{} \textbf{อิธา}ติ อิมิสฺสา ทิฏฺฐิยา อิมิสฺสา ขนฺติยา อิมิสฺสา รุจิยา อิมสฺมิํ อาทาเย อิมสฺมิํ ธมฺเม อิมสฺมิํ วินเย อิมสฺมิํ ธมฺมวินเย อิมสฺมิํ ปาวจเน อิมสฺมิํ พฺรหฺมจริเย อิมสฺมิํ สตฺถุสาสเน อิมสฺมิํ อตฺตภาเว อิมสฺมิํ มนุสฺสโลเกติ เอตมาทีนวํ ญตฺวา, มุนิ ปุพฺพาปเร อิธ.}

\prose{\textbf{เอกจริยํ ทฬฺหํ กยิรา}ติ ทฺวีหิ การเณหิ เอกจริยํ ทฬฺหํ กเรยฺย ปพฺพชฺชาสงฺขาเตน วา คณาววสฺสคฺคฏฺเฐน วา.\csromanspacer{} กถํ ปพฺพชฺชาสงฺขาเตน เอกจริยํ ทฬฺหํ กเรยฺย.\csromanspacer{} สพฺพํ ฆราวาสปลิโพธํ ฉินฺทิตฺวา ปุตฺตทารปลิโพธํ ฉินฺทิตฺวา ญาติปลิโพธํ ฉินฺทิตฺวา มิตฺตามจฺจปลิโพธํ ฉินฺทิตฺวา สนฺนิธิปลิโพธํ ฉินฺทิตฺวา เกสมสฺสุํ โอหาเรตฺวา กาสายานิ วตฺถานิ อจฺฉาเทตฺวา อคารสฺมา อนคาริยํ ปพฺพชิตฺวา อกิญฺจนภาวํ อุปคนฺตฺวา เอโก จเรยฺย วิหเรยฺย อิริเยยฺย วตฺเตยฺย ปาเลยฺย ยเปยฺย ยาเปยฺย.\csromanspacer{} เอวํ ปพฺพชฺชาสงฺขาเตน เอกจริยํ ทฬฺหํ กเรยฺย.}

\prose{กถํ คณาววสฺสคฺคฏฺเฐน เอกจริยํ ทฬฺหํ กเรยฺย.\csromanspacer{} โส เอวํ ปพฺพชิโต สมาโน เอโก อรญฺญวนปตฺถานิ ปนฺตานิ เสนาสนานิ ปฏิเสเวยฺย อปฺปสทฺทานิ อปฺปนิคฺโฆสานิ วิชนวาตานิ มนุสฺสราหสฺเสยฺยกานิ ปฏิสลฺลานสารุปฺปานิ.\csromanspacer{} โส เอโก คจฺเฉยฺย, เอโก ติฏฺเฐยฺย, เอโก นิสีเทยฺย, เอโก เสยฺยํ กปฺเปยฺย, เอโก คามํ ปิณฺฑาย ปวิเสยฺย, \csromanfolio{122}เอโก ปฏิกฺกเมยฺย, เอโก รโห นิสีเทยฺย, เอโก จงฺกมํ อธิฏฺเฐยฺย, เอโก จเรยฺย, วิหเรยฺย อิริเยยฺย วตฺเตยฺย ปาเลยฺย, ยเปยฺย ยาเปยฺย.\csromanspacer{} เอวํ คณาววสฺสคฺคฏฺเฐน.\csromanspacer{}\csromanspacer{}\textbf{เอกจริยํ ทฬฺหํ กเรยฺยา}ติ เอกจริยํ ทฬฺหํ กเรยฺย, ถิรํ กเรยฺย, ทฬฺหํ สมาทาโน อสฺส, อวฏฺฐิตสมาทาโน อสฺส กุสเลสุ ธมฺเมสูติ เอกจริยํ ทฬฺหํ กยิรา.}

\prose{\textbf{น นิเสเวถ เมถุนนฺ}ติ \textbf{เมถุนธมฺโม} นาม โย โส อสทฺธมฺโม คามธมฺโม.\csromanspacer{} ตํการณา วุจฺจติ \textbf{เมถุนธมฺโม}.\csromanspacer{} เมถุนธมฺมํ น เสเวยฺย น นิเสเวยฺย น สํเสเวยฺย น ปฏิเสเวยฺย น จเรยฺย น สมาจเรยฺย น สมาทาย วตฺเตยฺยาติ น นิเสเวถ เมถุนํ.\csromanspacer{} เตนาห ภควา\csromandash{}}

\csromangathameasure{เอตมาทีนวํ ญตฺวา, มุนิ ปุพฺพาปเร อิธ. \\ เอกจริยํ ทฬฺหํ กยิรา, น นิเสเวถ เมถุนนฺติ.}
\csromangathagroup{\csromangathastanza{เอตมาทีนวํ ญตฺวา, มุนิ ปุพฺพาปเร อิธ. \\ เอกจริยํ ทฬฺหํ กยิรา, น นิเสเวถ เมถุนนฺติ.}}

\proseitem{57}{\csromansymbolfootnote{*}{ขุ 1. 407 ปิฏฺเฐ.}วิเวกญฺเญว สิกฺเขถ, เอตํ \textbf{อริยา}นมุตฺตมํ.}

\csromangathameasure{น เตน เสฏฺโฐ มญฺเญถ, ส เว นิพฺพานสนฺติเก.}
\csromangathagroup{\csromangathastanza{น เตน เสฏฺโฐ มญฺเญถ, ส เว นิพฺพานสนฺติเก.}}

\prose{\textbf{วิเวกญฺเญว สิกฺเขถา}ติ \textbf{วิเวโก}ติ ตโย วิเวกา กาย\textbf{วิเวโก} จิตฺต\textbf{วิเวโก} อุปธิ\textbf{วิเวโก}.\csromanspacer{} กตโม กาย\textbf{วิเวโก} ฯเปฯ อยํ อุปธิ\textbf{วิเวโก}.\csromanspacer{} กาย\textbf{วิเวโก} จ วิเวกฏฺฐกายานํ เนกฺขมฺมาภิรตานํ.\csromanspacer{} จิตฺต\textbf{วิเวโก} จ ปริสุทฺธจิตฺตานํ ปรมโวทานปฺปตฺตานํ.\csromanspacer{} อุปธิ\textbf{วิเวโก} จ นิรูปธีนํ ปุคฺคลานํ วิสงฺขารคตานํ.\csromanspacer{}\csromanspacer{}\textbf{สิกฺขา}ติ ติสฺโส สิกฺขา อธิสีลสิกฺขา อธิจิตฺตสิกฺขา อธิปญฺญาสิกฺขา ฯเปฯ อยํ อธิปญฺญาสิกฺขา.\csromanspacer{}\csromanspacer{}\textbf{วิเวกญฺเญว สิกฺเขถา}ติ วิเวกญฺเญว สิกฺเขยฺย อาจเรยฺย สมาจเรยฺย สมาทาย วตฺเตยฺยาติ วิเวกญฺเญว สิกฺเขถ.}

\prose{\textbf{เอตํ อริยานมุตฺตมนฺ}ติ \textbf{อริยา} วุจฺจนฺติ พุทฺธา จ พุทฺธสาวกา จ ปจฺเจกพุทฺธา จ \textbf{อริยา}นํ เอตํ อคฺคํ เสฏฺฐํ วิสิฏฺฐํ ปาโมกฺขํ อุตฺตมํ ปวรํ ยทิทํ วิเวกจฺ\textbf{อริยา}ติ เอตํ \textbf{อริยา}นมุตฺตมํ.}

\prose{\textbf{น เตน เสฏฺโฐ มญฺเญถา}ติ กายวิเวกจฺ\textbf{อริยา}ย อุนฺนติํ น กเรยฺย, อุนฺนมํ น กเรยฺย, มานํ น กเรยฺย, ถามํ น กเรยฺย, ถมฺภํ น กเรยฺย, น เตน มานํ ชเนยฺย, น เตน ถทฺโธ อสฺส ปตฺถทฺโธ ปคฺคหิตสิโรติ เตน เสฏฺโฐ น มญฺเญถ.}

\csromanheader{2}{7. ติสฺสเมตฺเตยฺยสุตฺตนิทฺเทส}
\prose{\csromanfolio{123}\textbf{ส เว นิพฺพานสนฺติเก}ติ โส นิพฺพานสฺส สนฺติเก สามนฺตา อาสนฺเน อวิทูเร อุปกฏฺเฐติ ส เว นิพฺพานสนฺติเก.\csromanspacer{} เตนาห ภควา\csromandash{}}

\csromangathameasure{วิเวกญฺเญว สิกฺเขถ, เอตํ อริยานมุตฺตมํ. \\ น เตน เสฏฺโฐ มญฺเญถ, ส เว นิพฺพานสนฺติเกติ.}
\csromangathagroup{\csromangathastanza{วิเวกญฺเญว สิกฺเขถ, เอตํ อริยานมุตฺตมํ. \\ น เตน เสฏฺโฐ มญฺเญถ, ส เว นิพฺพานสนฺติเกติ.}}

\proseitem{58}{\csromansymbolfootnote{*}{ขุ 1. 407 ปิฏฺเฐ.}\textbf{ริตฺตสฺส มุนิโน จรโต}, กาเมสุ อนเปกฺขิโน.}

\csromangathameasure{โอฆติณฺณสฺส ปิหยนฺติ, กาเมสุ คธิตา ปชา.}
\csromangathagroup{\csromangathastanza{โอฆติณฺณสฺส ปิหยนฺติ, กาเมสุ คธิตา ปชา.}}

\prose{\textbf{ริตฺตสฺส มุนิโน จรโต}ติ ริตฺตสฺส วิวิตฺตสฺส ปวิวิตฺตสฺส, กายทุจฺจริเตน ริตฺตสฺส วิวิตฺตสฺส ปวิวิตฺตสฺส.\csromanspacer{} วจีทุจฺจริเตน.\csromanspacer{} มโนทุจฺจริเตน.\csromanspacer{} ราเคน.\csromanspacer{} โทเสน.\csromanspacer{} โมเหน.\csromanspacer{} โกเธน.\csromanspacer{} อุปนาเหน.\csromanspacer{} มกฺเขน.\csromanspacer{} ปฬาเสน.\csromanspacer{} อิสฺสาย.\csromanspacer{} มจฺฉริเยน.\csromanspacer{} มายาย.\csromanspacer{} สาเฐยฺเยน.\csromanspacer{} ถมฺเภน.\csromanspacer{} สารมฺเภน.\csromanspacer{} มาเนน.\csromanspacer{} อติมาเนน.\csromanspacer{} มเทน.\csromanspacer{} ปมาเทน.\csromanspacer{} สพฺพกิเลเสหิ สพฺพทุจฺจริเตหิ.\csromanspacer{} สพฺพทรเถหิ.\csromanspacer{} สพฺพปริฬาเหหิ.\csromanspacer{} สพฺพสนฺตาเปหิ.\csromanspacer{} สพฺพากุสลาภิสงฺขาเรหิ ริตฺตสฺส วิวิตฺตสฺส ปวิวิตฺตสฺส.\csromanspacer{}\csromanspacer{}\textbf{มุนิโน}ติ \textbf{โมนํ} วุจฺจติ ญาณํ ฯเปฯ สงฺคชาลมติจฺจ โส มุนิ.\csromanspacer{}\csromanspacer{}\textbf{จรโต}ติ จรโต วิหรโต อิริยโต วตฺตโต ปาลยโต ยปยโต ยาปยโตติ ริตฺตสฺส มุนิโน จรโต.}

\prose{\textbf{กาเมสุ อนเปกฺขิโน}ติ \textbf{กามา}ติ อุทฺทานโต เทฺว \textbf{กามา} วตฺถุ \textbf{กามา} จ กิเลส\textbf{กามา} จ ฯเปฯ อิเม วุจฺจนฺติ วตฺถุ\textbf{กามา} ฯเปฯ อิเม วุจฺจนฺติ กิเลส\textbf{กามา}.\csromanspacer{} วตฺถุกาเม ปริชานิตฺวา กิเลสกาเม ปหาย ปชหิตฺวา วิโนเทตฺวา พฺยนฺติํ กริตฺวา อนภาวํ คเมตฺวา กาเมสุ อนเปกฺขมาโน จตฺตกาโม วนฺตกาโม มุตฺตกาโม ปหีนกาโม ปฏินิสฺสฏฺฐกาโม, วีตราโค จตฺตราโค วนฺตราโค มุตฺตราโค ปหีนราโค ปฏินิสฺสฏฺฐราโค นิจฺฉาโต นิพฺพุโต สีติภูโต\footnote{สีติภูโต (สี)} สุขปฺปฏิสํเวที พฺรหฺมภูเตน อตฺตนา วิหรตีติ กาเมสุ อนเปกฺขิโน.}

\prose{\textbf{โอฆติณฺณสฺส ปิหยนฺติ, กาเมสุ คธิตา ปชา}ติ \textbf{ปชา}ติ สตฺตาธิวจนํ, \textbf{ปชา} กาเมสุ รตฺตา คิทฺธา คธิตา มุจฺฉิตา อชฺโฌสนฺนา ลคฺคา ลคฺคิตา ปลิพุทฺธา.\csromanspacer{} เต กาโมฆํ ติณฺณสฺส ภโวฆํ ติณฺณสฺส ทิฏฺโฐฆํ ติณฺณสฺส อวิชฺโชฆํ ติณฺณสฺส สพฺพสงฺขารปถํ ติณฺณสฺส อุตฺติณฺณสฺส นิตฺติณฺณสฺส อติกฺกนฺตสฺส สมติกฺกนฺตสฺส วีติวตฺตสฺส ปารํ คตสฺส ปารํ \csromanfolio{124}ปตฺตสฺส อนฺตํ คตสฺส อนฺตํ ปตฺตสฺส โกฏิํ คตสฺส โกฏิํ ปตฺตสฺส ปริยนฺตํ คตสฺส ปริยนฺตํ ปตฺตสฺส โวสานํ คตสฺส โวสานํ ปตฺตสฺส ตาณํ คตสฺส ตาณํ ปตฺตสฺส เลณํ คตสฺส เลณํ ปตฺตสฺส สรณํ คตสฺส สรณํ ปตฺตสฺส อภยํ คตสฺส อภยํ ปตฺตสฺส อจฺจุตํ คตสฺส อจฺจุตํ ปตฺตสฺส อมตํ คตสฺส อมตํ ปตฺตสฺส นิพฺพานํ คตสฺส นิพฺพานํ ปตฺตสฺส อิจฺฉนฺติ สาทิยนฺติ ปตฺถยนฺติ ปิหยนฺติ อภิชาปฺปนฺติ.\csromanspacer{} ยถา อิณายิกา อานณฺยํ\footnote{อาณณฺยํ (ฏฺฐ)} ปตฺเถนฺติ ปิหยนฺติ.\csromanspacer{} ยถา อาพาธิกา อาโรคฺยํ ปตฺเถนฺติ ปิหยนฺติ.\csromanspacer{} ยถา พนฺธนพทฺธา พนฺธนโมกฺขํ ปตฺเถนฺติ ปิหยนฺติ.\csromanspacer{} ยถา ทาสา ภุชิสฺสํ ปตฺเถนฺติ ปิหยนฺติ.\csromanspacer{} ยถา กนฺตารทฺธานปกฺขนฺทา\footnote{กนฺตารทฺธานปกฺขนฺตา (สี), กนฺตารทฺธานปกฺขนฺนา (สฺยา)} เขมนฺตภูมิํ ปตฺเถนฺติ ปิหยนฺติ.\csromanspacer{} เอวเมวํ ปชา กาเมสุ รตฺตา คิทฺธา คธิตา มุจฺฉิตา อชฺโฌสนฺนา ลคฺคา ลคฺคิตา ปลิพุทฺธา.\csromanspacer{} เต กาโมฆํ ติณฺณสฺส ภโวฆํ ติณฺณสฺส ฯเปฯ นิพฺพานํ คตสฺส นิพฺพานํ ปตฺตสฺส อิจฺฉนฺติ สาทิยนฺติ ปตฺถยนฺติ ปิหยนฺติ อภิชปฺปนฺตีติ โอฆติณฺณสฺส ปิหยนฺติ กาเมสุ คธิตา ปชา.\csromanspacer{} เตนาห ภควา\csromandash{}}

\csromangathameasure{ริตฺตสฺส มุนิโน จรโต, กาเมสุ อนเปกฺขิโน. \\ โอฆติณฺณสฺส ปิหยนฺติ, กาเมสุ คธิตา ปชาติ. ตสฺสเมตฺเตยฺยสุตฺตนิทฺเทโส สตฺตโม.}
\csromangathagroup{\csromangathastanza{ริตฺตสฺส มุนิโน จรโต, กาเมสุ อนเปกฺขิโน.}\csromangathastanza{โอฆติณฺณสฺส ปิหยนฺติ, กาเมสุ คธิตา ปชาติ.\csromanspacer{} ตสฺสเมตฺเตยฺยสุตฺตนิทฺเทโส สตฺตโม.}}

\csromanheader{2}{8. ปสูรสุตฺตนิทฺเทส}
\csromanheader{2}{อถ ปสูรสุตฺตนิทฺเทสํ วกฺขติ\csromandash{}}
\proseitem{59}{\csromansymbolfootnote{*}{ขุ 1. 407 ปิฏฺเฐ.}อิเธว สุทฺธิํ อิติ วาทยนฺติ,}

\prose{\textbf{นาญฺเญสุ ธมฺเมสุ วิสุทฺธิมาหุ.}}

\csromangathameasure{\textbf{ยํ นิสฺสิตา ตตฺถ สุตํ วทานา,} \\ \textbf{ปจฺเจกสจฺเจสุ ปุถู นิวิฏฺฐา.}}
\csromangathagroup{\csromangathastanza{\textbf{ยํ นิสฺสิตา ตตฺถ สุตํ วทานา,} \\ \textbf{ปจฺเจกสจฺเจสุ ปุถู นิวิฏฺฐา.}}}

\prose{\textbf{อิเธว สุทฺธิํ อิติ วาทยนฺตี}ติ อิเธว สุทฺธิํ วิสุทฺธิํ ปริสุทฺธิํ, มุตฺติํ วิมุตฺติํ ปริมุตฺติํ วทนฺติ กเถนฺติ ภณนฺติ ทีปยนฺติ โวหรนฺติ, “สสฺสโต โลโก, อิทเมว สจฺจํ โมฆมญฺญนฺ”ติ สุทฺธิํ วิสุทฺธิํ ปริสุทฺธิํ, มุตฺติํ วิมุตฺติํ ปริมุตฺติํ วทนฺติ กเถนฺติ ภณนฺติ ทีปยนฺติ โวหรนฺติ, “อสสฺสโต}

\vspace{\tipitakaparskip}
\csromancenter{ปสูรสุตฺตนิทฺเทส โลโก.\csromanspacer{} อนฺตวา โลโก.\csromanspacer{} อนนฺตวา โลโก.\csromanspacer{} ตํ ชีวํ ตํ สรีรํ.\csromanspacer{} อญฺญํ ชีวํ อญฺญํ สรีรํ.\csromanspacer{} โหติ ตถาคโต ปรํ มรณา.\csromanspacer{} น โหติ ตถาคโต ปรํ มรณา.\csromanspacer{} โหติ จ น จ โหติ ตถาคโต ปรํ มรณา.\csromanspacer{} เนว โหติ น น โหติ ตถาคโต ปรํ มรณา, อิทเมว สจฺจํ โมฆมญฺญนฺ”ติ สุทฺธิํ วิสุทฺธิํ ปริสุทฺธิํ, มุตฺติํ วิมุตฺติํ ปริมุตฺติํ วทนฺติ กเถนฺติ ภณนฺติ ทีปยนฺติ โวหรนฺตีติ อิเธว สุทฺธิํ อิติ วาทยนฺติ.}
\csromanmatikaanchor{mtk.8}
\csromantocmarkref{subsection}{7. ติสฺสเมตฺเตยฺยสุตฺตนิทฺเทส}{mtk.8}
\bookmark[dest={mtk.8},level=2]{7. ติสฺสเมตฺเตยฺยสุตฺตนิทฺเทส}
\prose{\csromanfolio{125}\textbf{นาญฺเญสุ ธมฺเมสุ วิสุทฺธิมาหู}ติ อตฺตโน สตฺถารํ ธมฺมกฺขานํ คณํ ทิฏฺฐิํ ปฏิปทํ มคฺคํ ฐเปตฺวา สพฺเพ ปรวาเท ขิปนฺติ อุกฺขิปนฺติ ปริกฺขิปนฺติ, “โส สตฺถา น สพฺพญฺญู, ธมฺโม น สฺวากฺขาโต, คโณ น สุปฺปฏิปนฺโน, ทิฏฺฐิ น ภทฺทิกา, ปฏิปทา น สุปญฺญตฺตา, มคฺโค น นิยฺยานิโก, นตฺเถตฺถ สุทฺธิ วา วิสุทฺธิ วา ปริสุทฺธิ วา, มุตฺติ วา วิมุตฺติ วา ปริมุตฺติ วา, น ตตฺถ สุชฺฌนฺติ วา วิสุชฺฌนฺติ วา ปริสุชฺฌนฺติ วา, มุจฺจนฺติ วา วิมุจฺจนฺติ วา ปริมุจฺจนฺติ วา, หีนา นิหีนา โอมกา ลามกา ฉตุกฺกา ปริตฺตา”ติ เอวมาหํสุ เอวํ วทนฺติ เอวํ กเถนฺติ เอวํ ภณนฺติ เอวํ ทีปยนฺติ เอวํ โวหรนฺตีติ นาญฺเญสุ ธมฺเมสุ วิสุทฺธิมาหุ.}

\prose{\textbf{ยํ นิสฺสิตา ตตฺถ สุภํ วทานา}ติ ยํ \textbf{นิสฺสิตา}ติ สตฺถารํ ธมฺมกฺขานํ คณํ ทิฏฺฐิํ ปฏิปทํ มคฺคํ \textbf{นิสฺสิตา} สนฺ\textbf{นิสฺสิตา}\footnote{ปติฏฺฐิตา (สี), อานิสฺสิตา (ก) เหฏฺฐา 70 ปิฏฺเฐปิ.} อลฺลีนา, อุปคตา อชฺโฌสิตา อธิมุตฺตา.\csromanspacer{}\csromanspacer{}\textbf{ตตฺถา}ติ สกาย ทิฏฺฐิยา สกาย ขนฺติยา สกาย รุจิยา สกาย ลทฺธิยา.\csromanspacer{}\csromanspacer{}\textbf{สุตํ วทานา}ติ สุภวาทา โสภนวาทา ปณฺฑิตวาทา ถิรวาทา\footnote{ธีรวาทา (สฺยา)} ญายวาทา เหตุวาทา ลกฺขณาวาทา การณวาทา ฐานวาทา สกาย ลทฺธิยาติ ยํ \textbf{นิสฺสิตา} ตตฺถ สุภํ วทานา.}

\prose{\textbf{ปจฺเจกสจฺเจสุ ปุถู นิวิฏฺฐา}ติ ปุถู สมณพฺราหฺมณา ปุถู ปจฺเจกสจฺเจสุ นิวิฏฺฐา ปติฏฺฐิตา อลฺลีนา อุปคตา อชฺโฌสิตา อธิมุตฺตา “สสฺสโต โลโก, อิทเมว สจฺจํ โมฆมญฺญนฺ”ติ นิวิฏฺฐา ปติฏฺฐิตา อลฺลีนา อุปคตา อชฺโฌสิตา อธิมุตฺตา. “อสสฺสโต โลโก ฯเปฯ เนว โหติ น น โหติ ตถาคโต ปรํ มรณา, อิทเมว สจฺจํ โมฆมญฺญนฺ”ติ นิวิฏฺฐา ปติฏฺฐิตา อลฺลีนา อุปคตา อชฺโฌสิตา อธิมุตฺตาติ ปจฺเจกสจฺเจสุ ปุถู นิวิฏฺฐา.\csromanspacer{} เตนาห ภควา\csromandash{}}

\csromangathameasure{อิเธว สุทฺธิํ อิติ วาทยนฺติ, \\ นาญฺเญสุ ธมฺเมสุ วิสุทฺธิมาหุ. \\ ยํ นิสฺสิตา ตตฺถ สุภํ วทานา, \\ ปจฺเจกสจฺเจสุ ปุถู นิวิฏฺฐาติ.}
\csromangathagroup{\csromangathastanza{\csromanfolio{126}อิเธว สุทฺธิํ อิติ วาทยนฺติ, \\ นาญฺเญสุ ธมฺเมสุ วิสุทฺธิมาหุ. \\ ยํ นิสฺสิตา ตตฺถ สุภํ วทานา, \\ ปจฺเจกสจฺเจสุ ปุถู นิวิฏฺฐาติ.}}

\proseitem{60}{\csromansymbolfootnote{*}{ขุ 1. 408 ปิฏฺเฐ.}เต วาทกามา ปริสํ วิคยฺห,}

\prose{\textbf{พาลํ ทหนฺตี มิถุ อญฺญมญฺญํ.}}

\csromangathameasure{\textbf{วทนฺติ เต อญฺญสิตา กโถชฺชํ,} \\ \textbf{ปสํสกามา กุสลาวทานา.}}
\csromangathagroup{\csromangathastanza{\textbf{วทนฺติ เต อญฺญสิตา กโถชฺชํ,} \\ \textbf{ปสํสกามา กุสลาวทานา.}}}

\prose{\textbf{เต วาทกามา ปริสํ วิคยฺหา}ติ \textbf{เต วาทกามา}ติ \textbf{เต วาทกามา} วาทตฺถิกา วาทาธิปฺปายา วาทปุเรกฺขารา วาทปริเยสนํ จรนฺตา.\csromanspacer{}\csromanspacer{}\textbf{ปริสํ วิคยฺหา}ติ ขตฺติยปริสํ พฺราหฺมณปริสํ คหปติปริสํ สมณปริสํ วิคยฺห โอคยฺห อชฺโฌคาเหตฺวา ปวิสิตฺวาติ \textbf{เต วาทกามา} ปริสํ วิคยฺห.}

\prose{\textbf{พาลํ ทหนฺตี มิถุ อญฺญมญฺญ}นฺติ \textbf{มิถู}ติ เทฺว ชนา เทฺว กลหการกา เทฺว ภณฺฑนการกา เทฺว ภสฺสการกา เทฺว วิวาทการกา เทฺว อธิกรณการกา เทฺว วาทิโน เทฺว สลฺลาปกา เต อญฺญมญฺญํ พาลโต หีนโต นิหีนโต โอมกโต ลามกโต ฉตุกฺกโต ปริตฺตโต ทหนฺติ ปสฺสนฺติ ทกฺขนฺติ โอโลเกนฺติ นิชฺฌายนฺติ อุปปริกฺขนฺตีติ พาลํ ทหนฺตี มิถุ อญฺญมญฺญํ.}

\prose{\textbf{วทนฺติ เต อญฺญสิตา กโถชฺชนฺ}ติ อญฺญํ สตฺถารํ ธมฺมกฺขานํ คณํ ทิฏฺฐิํ ปฏิปทํ มคฺคํ นิสฺสิตา สนฺนิสฺสิตา อลฺลีนา อุปคตา อชฺโฌสิตา อธิมุตฺตา.\csromanspacer{}\csromanspacer{}\textbf{กโถชฺชํ} วุจฺจติ กลโห ภณฺฑนํ วิคฺคโห วิวาโท เมธคํ.\csromanspacer{} อถ วา \textbf{กโถชฺช}นฺติ อโนชวนฺตี สา กถา.\csromanspacer{} \textbf{กโถชฺชํ} วทนฺติ, กลหํ วทนฺติ, ภณฺฑนํ วทนฺติ, วิคฺคหํ วทนฺติ, วิวาทํ วทนฺติ, เมธคํ วทนฺติ กเถนฺติ ภณนฺติ ทีปยนฺติ โวหรนฺตีติ วทนฺติ เต อญฺญสิตา \textbf{กโถชฺช}ํ.}

\prose{\textbf{ปสํสกามา กุสลาวทานา}ติ \textbf{ปสํสกามา}ติ \textbf{ปสํสกามา} ปสํสตฺถิกา ปสํสาธิปฺปายา ปสํสปุเรกฺขารา ปสํสปริเยสนํ จรนฺตา.\csromanspacer{} \textbf{กุสลาวทานา}ติ กุสลวาทา ปณฺฑิตวาทา ถิรวาทา ญายวาทา เหตุวาทา ลกฺขณวาทา การณวาทา ฐานวาทา สกาย ลทฺธิยาติ \textbf{ปสํสกามา} กุสลาวทานา.\csromanspacer{} เตนาห ภควา\csromandash{}}

\csromanheader{2}{8. ปสูรสุตฺตนิทฺเทส}
\csromangathameasure{เต วาทกามา ปริสํ วิคยฺห, \\ ปาลํ ทหนฺตี มิถุ อญฺญมญฺญํ. \\ วทนฺติ เต อญฺญสิตา กโถชฺชํ, \\ ปสํสกามา กุสลาวทานาติ.}
\csromangathagroup{\csromangathastanza{\csromanfolio{127}เต วาทกามา ปริสํ วิคยฺห, \\ ปาลํ ทหนฺตี มิถุ อญฺญมญฺญํ. \\ วทนฺติ เต อญฺญสิตา กโถชฺชํ, \\ ปสํสกามา กุสลาวทานาติ.}}

\proseitem{61}{\csromansymbolfootnote{*}{ขุ 1. 408 ปิฏฺเฐ.}\textbf{ยุตฺโต กถายํ ปริสาย มชฺเฌ},}

\prose{\textbf{ปสํสมิจฺฉํ วินิฆาติ โหติ.}}

\csromangathameasure{\textbf{อปาหตสฺมิํ ปน มงฺกุ โหติ,} \\ \textbf{นินฺทาย โส กุปฺปติ รนฺธเมสี.}}
\csromangathagroup{\csromangathastanza{\textbf{อปาหตสฺมิํ ปน มงฺกุ โหติ,} \\ \textbf{นินฺทาย โส กุปฺปติ รนฺธเมสี.}}}

\prose{\textbf{ยุตฺโต กถายํ ปริสาย มชฺเฌ}ติ ขตฺติยปริสาย วา พฺราหฺมณปริสาย วา คหปติปริสาย วา สมณปริสาย วา มชฺเฌ อตฺตโน กถายํ ยุตฺโต ปยุตฺโต อายุตฺโต สมายุตฺโต สมฺปยุตฺโต กเถตุนฺติ ยุตฺโต กถายํ ปริสาย มชฺเฌ.}

\prose{\textbf{ปสํสมิจฺฉํ วินิฆาติ โหตี}ติ \textbf{ปสํสมิจฺฉนฺ}ติ ปสํสํ โถมนํ กิตฺติํ วณฺณหาริยํ อิจฺฉนฺโต สาทิยนฺโต ปตฺถยนฺโต ปิหยนฺโต อภิชปฺปนฺโต.\csromanspacer{} \textbf{วินิฆาติ โหตี}ติ ปุพฺเพว สลฺลาปา กถํกถี วินิฆาตี โหติ “ชโย นุ โข เม ภวิสฺสติ, ปราชโย นุ โข เม ภวิสฺสติ, กถํ นิคฺคหํ กริสฺสามิ, กถํ ปฏิกมฺมํ กริสฺสามิ, กถํ วิเสสํ กริสฺสามิ, กถํ ปฏิวิเสสํ กริสฺสามิ, กถํ อาเวฐิยํ\footnote{อาเวธิยํ (สฺยา)} กริสฺสามิ, กตํ นิพฺเพฐิยํ\footnote{นิพฺเพธิยํ (สฺยา, ก)} กริสฺสามิ, กถํ เฉทํ กริสฺสามิ, กถํ มณฺฑลํ กริสฺสามี”ติ เอวํ ปุพฺเพว สลฺลาปา กถํกถี วินิฆาตี โหตีติ ปสํสมิจฺฉํ วินิฆาติ โหติ.}

\prose{\textbf{อปาหตสฺมิํ ปน มงฺกุ โหตี}ติ เย เต ปญฺหวีมํสกา ปริสา ปาริสชฺชา ปาสาริกา\footnote{ปาสนิกา (สฺยา)}, เต อปหรนฺติ, “อตฺถาปคตํ ภณิตนฺ”ติ อตฺถโต อปหรนฺติ, “พฺยญฺชนาปคตํ ภณิตนฺ”ติ พฺยญฺชนโต อปหรนฺติ, “อตฺถพฺยญฺชนาปคตํ ภณิตนฺ”ติ อตฺถพฺยญฺชนโต อปหรนฺติ, อตฺโถ เต ทุนฺนีโต, พฺยญฺชนํ เต ทุโรปิตํ, อตฺถพฺยญฺชนํ เต ทุนฺนีตํ ทุโรปิตํ, นิคฺคโห เต อกโต, ปฏิกมฺมํ เต ทุกฺกฏํ, วิเสโส เต อกโต, ปฏิวิเสโส เต ทุกฺกโฏ, อาเวฐิยา เต อกตา, นิพฺเพฐิยา เต ทุกฺกฏา, เฉโท เต อกโต, มณฺฑลํ เต ทุกฺกฏํ, วิสมกถํ ทุกฺกถิตํ ทุพฺภณิตํ ทุลฺลปิตํ ทุรุตฺตํ ทุพฺภาสิตนฺติ อปหรนฺติ.\csromanspacer{}\csromanspacer{}\textbf{อปาหตสฺมิํ ปน มงฺกุ} \csromanfolio{128}\textbf{โหตี}ติ อปาหตสฺมิํ มงฺกุ โหติ ปีฬิโต ฆฏฺฏิโต พฺยาธิโต โทมนสฺสิโต \textbf{โหตี}ติ อปาหตสฺมิํ ปน มงฺกุ โหติ\textbf{.}}

\prose{\textbf{นินฺทาย โส กุปฺปติ รนฺธเมสี}ติ นินฺทาย ครหาย อกิตฺติยา อวณฺณหาริกาย กุปฺปติ พฺยาปชฺชติ ปติตฺถียติ, โกปญฺจ โทสญฺจ อปฺปจฺจยญฺจ ปาตุกโรตีติ นินฺทาย โส กุปฺปติ\textbf{.}\csromanspacer{}\csromanspacer{}\textbf{รนฺธเมสี}ติ วิรนฺธเมสี อปรทฺธเมสี ขลิตเมสี คฬิตเมสี วิวรเมสีติ นินฺทาย โส กุปฺปติ รนฺธเมสี\textbf{.}\csromanspacer{} เตนาห ภควา\csromandash{}}

\csromangathameasure{ยุตฺโต กถายํ ปริสาย มชฺเฌ, \\ ปสํสมิจฺฉํ วินิฆาติ โหติ\textbf{.} \\ อปาหตสฺมิํ ปน มงฺกุ โหติ,}
\csromangathagroup{\csromangathastanza{ยุตฺโต กถายํ ปริสาย มชฺเฌ, \\ ปสํสมิจฺฉํ วินิฆาติ โหติ\textbf{.} \\ อปาหตสฺมิํ ปน มงฺกุ โหติ,}}

\prose{\textbf{นินฺทาย โส กุปฺปติ รนฺธเมสี}ติ\textbf{.}}

\proseitem{62}{\csromansymbolfootnote{*}{ขุ 1. 408 ปิฏฺเฐ.}ยมสฺส วาทํ ปริหีนมาหุ,}

\prose{\textbf{อปาหตํ ปญฺหวิมํสกาเส}\footnote{ปญฺหวิมํสกา เย (สฺยา)}\textbf{.}}

\prose{\textbf{ปริเทวติ โสจติ หีนวาโท.}}

\prose{\textbf{“อุปจฺจคา มนฺ}”\textbf{ติ อนุตฺถุนาติ.}}

\prose{\textbf{ยมสฺส วาทํ ปริหีนมาหู}ติ ยํ ตสฺส วาทํ หีนํ นิหีนํ ปริหีนํ ปริหาปิตํ น ปริปูริตํ เอวมาหํสุ เอวํ กเถนฺติ เอวํ ภณนฺติ เอวํ ทีปยนฺติ เอวํ โวหรนฺตีติ ยมสฺส วาทํ ปริหีนมาหุ\textbf{.}}

\prose{\textbf{อปาหตํ ปญฺหวิมํสกาเส}ติ เย เต ปญฺหวีมํสกา ปริสา ปาริสชฺชา ปาสาริกา, เต อปหรนฺติ, “อตฺถาปคตํ ภณิตนฺ”ติ อตฺถโต อปหรนฺติ, “พฺยญฺชนาปคตํ ภณิตนฺ”ติ พฺยญฺชนโต อปหรนฺติ, “อตฺถพฺยญฺชนาปคตํ ภณิตนฺ”ติ อตฺถพฺยญฺชนโต อปหรนฺติ, อตฺโถ เต ทุนฺนีโต, พฺยญฺชนํ เต ทุโรปิตํ, อตฺถพฺยญฺชนํ เต ทุนฺนีตํ ทุโรปิตํ, นิคฺคโห เต อกโต, ปฏิกมฺมํ เต ทุกฺกฏํ, วิเสโส เต อกโต, ปฏิวิเสโส เต ทุกฺกโฏ, อาเวฐิยา เต อกตา, นิพฺเพฐิยา เต ทุกฺกฏา, เฉโท เต อกโต, มณฺฑลํ เต ทุกฺกฏํ, วิสมกถํ ทุกฺกถิตํ ทุพฺภณิตํ ทุลฺลปิตํ ทุรุตฺตํ ทุพฺภาสิตนฺติ อปหรนฺตีติ อปาหตํ ปญฺหวิมํสกาเส\textbf{.}}

\csromanheader{2}{8. ปสูรสุตฺตนิทฺเทส}
\prose{\csromanfolio{129}\textbf{ปริเทวติ โสจติ หีนวาโท}ติ \textbf{ปริเทวตี}ติ อญฺญํ มยา อาวชฺชิตํ อญฺญํ จินฺติตํ, อญฺญํ อุปธาริตํ, อญฺญํ อุปลกฺขิตํ, โส มหาปกฺโข มหาปริโส มหาปริวาโร, ปริสา จายํ วคฺคา น สมคฺคา, สมคฺคาย ปริสาย เหตุ กถาสลฺลาโป ปุน ภญฺชิสฺสามีติ ยา เอวรูปา\footnote{โย เอวรูโป (สฺยา)} วาจา ปลาโป วิปฺปลาโป ลาลปฺโป ลาลปฺปายนา ลาลปฺปายิตตฺตนฺติ ปริเทวติ.\csromanspacer{}\csromanspacer{}\textbf{โสจตี}ติ “ตสฺส ชโย”ติ โสจติ, “มยฺหํ ปราชโย”ติ โสจติ, “ตสฺส ลาโภ”ติ โสจติ, “มยฺหํ อลาโภ”ติ โสจติ, “ตสฺส ยโส”ติ โสจติ. “มยฺหํ อยโส”ติ โสจติ, “ตสฺส ปสํสา”ติ โสจติ, “มยฺหํ นินฺทา”ติ โสจติ, “ตสฺส สุขนฺ”ติ โสจติ, “มยฺหํ ทุกฺขนฺ”ติ โสจติ, โสสกฺกโต ครุกโต มานิโต ปูชิโต อปจิโต ลาภี จีวรปิณฺฑปาตเสนาสนคิลานปจฺจยเภสชฺชปริกฺขารานํ, “อหมสฺมิ อสกฺกโต อครุกโต อมานิโต อปูชิโต อนปจิโต น ลาภี จีวรปิณฺฑปาตเสนาสนคิลานปจฺจยเภสชฺชปริกฺขารานนฺ”ติ โสจติ กิลมติ ปริเทวติ อุรตฺตาฬิํ กนฺทติ สมฺโมหํ อาปชฺชตีติ ปริเทวติ โสจติ.\csromanspacer{}\csromanspacer{}\textbf{หีนวาโท}ติ หีนวาโท นิหีนวาโท ปริหีนวาโท ปริหาปิตวาโท น ปริปูรวาโทติ ปริเทวติ โสจติ หีนวาโท.}

\prose{\textbf{“อุปจฺจคา มนฺ”ติ อนุตฺถุนาตี}ติ โส มํ วาเทน วาทํ อจฺจคา อุปจฺจคา อติกฺกนฺโต สมติกฺกนฺโต วีติวตฺโตติ เอวมฺปิ อุปจฺจคา มนฺติ.\csromanspacer{}\csromanspacer{}อถ วา มํ วาเทน วาทํ อภิภวิตฺวา อชฺโฌตฺถริตฺวา ปริยาทิยิตฺวา มทฺทยิตฺวา จรติ วิหรติ อิริยติ วตฺตติ ปาเลติ ยเปติ ยาเปตีติ เอวมฺปิ อุปจฺจคา มนฺติ.\csromanspacer{}\csromanspacer{}\textbf{อนุตฺถุนา} วุจฺจติ วาจาปลาโป วิปฺปลาโป ลาลปฺโป ลาลปฺปายนา ลาลปฺปายิตตฺตนฺติ อุปจฺจคา มนฺติ อนุตฺถุนาติ.\csromanspacer{} เตนาห ภควา\csromandash{}}

\csromangathameasure{ยมสฺส วาทํ ปริหีนมาหุ, \\ อปาหตํ ปญฺหวิมํสกาเส. \\ \textbf{ปริเทวติ โสจติ หีนวาโท},}
\csromangathagroup{\csromangathastanza{ยมสฺส วาทํ ปริหีนมาหุ, \\ อปาหตํ ปญฺหวิมํสกาเส. \\ \textbf{ปริเทวติ โสจติ หีนวาโท},}}

\prose{\textbf{“อุปจฺจคา มนฺ”ติ อนุตฺถุนาตี}ติ.}

\proseitem{63}{\csromanfolio{130}\csromansymbolfootnote{*}{ขุ 1. 408 ปิฏฺเฐ.}\textbf{เอเต วิวาทา สมเณสุ ชาตา},}

\prose{\textbf{เอเตสุ อุคฺฆาติ นิฆาติ โหติ.}}

\csromangathameasure{\textbf{เอตมฺปิ ทิสฺวา วิรเม กโถชฺชํ,} \\ \textbf{น ห’ญฺญทตฺถ’ตฺถิ ปสํสลาภา.}}
\csromangathagroup{\csromangathastanza{\textbf{เอตมฺปิ ทิสฺวา วิรเม กโถชฺชํ,} \\ \textbf{น ห’ญฺญทตฺถ’ตฺถิ ปสํสลาภา.}}}

\prose{\textbf{เอเต วิวาทา สมเณสุ ชาตา}ติ \textbf{สมณา}ติ เย เกจิ อิโต พหิทฺธา ปริพฺพชูปคตา ปริพฺพชสมาปนฺนา, เอเต ทิฏฺฐิกลหา ทิฏฺฐิภณฺฑนา ทิฏฺฐิวิคฺคหา ทิฏฺฐิวิวาทา ทิฏฺฐิเมธคา สมเณสุ ชาตา สญฺชาตา นิพฺพตฺตา อภินิพฺพตฺตา ปาตุภูตาติ เอเต วิวาทา สมเณสุ ชาตา.}

\prose{\textbf{เอเตสุ อุคฺฆาตินิฆาติ} \textbf{โหตี}ติ ชยปราชโย โหติ ลาภาลาโภ โหติ, ยสายโส โหติ, นินฺทาปสํสา โหติ, สุขทุกฺขํ โหติ, โสมนสฺสโทมนสฺสํ โหติ, อิฏฺฐานิฏฺฐํ โหติ, อนุนยปฏิฆํ โหติ, อุคฺฆาติตนิคฺฆาติตํ โหติ, อนุโรธวิโรโธ โหติ, ชเยน จิตฺตํ อุคฺฆาติตํ โหติ, ปราชเยน จิตฺตํ นิคฺฆาติตํ โหติ, ลาเภน จิตฺตํ อุคฺฆาติตํ โหติ, อลาเภน จิตฺตํ นิคฺฆาติตํ โหติ, ยเสน จิตฺตํ อุคฺฆาติตํ โหติ, อยเสน จิตฺตํ นิคฺฆาติตํ โหติ, ปสํสาย จิตฺตํ อุคฺฆาติตํ โหติ, นินฺทาย จิตฺตํ นิคฺฆาติตํ โหติ, สุเขน จิตฺตํ อุคฺฆาติตํ โหติ, ทุกฺเขน จิตฺตํ นิคฺฆาติตํ โหติ, โสมนสฺเสน จิตฺตํ อุคฺฆาติตํ โหติ, โทมนสฺเสน จิตฺตํ นิคฺฆาติตํ โหติ, อุนฺนติยา\footnote{อุณฺณติยา (สฺยา, ก)} จิตฺตํ อุคฺฆาติตํ โหติ, โอนติยา\footnote{โอณติยา (สฺยา, ก)} จิตฺตํ นิคฺฆาติตํ \textbf{โหตี}ติ เอเตสุ อุคฺฆาติ นิฆาติ โหติ.}

\prose{\textbf{เอตมฺปิ ทิสฺวา วิรเม กโถชฺชนฺ}ติ \textbf{เอตมฺปิ ทิสฺวา}ติ เอตํ อาทีนวํ ทิสฺวา ปสฺสิตฺวา ตุลยิตฺวา ตีรยิตฺวา วิภาวยิตฺวา วิภูตํ กตฺวา ทิฏฺฐิกลเหสุ ทิฏฺฐิภณฺฑเนสุ ทิฏฺฐิวิคฺคเหสุ ทิฏฺฐิวิวาเทสุ ทิฏฺฐิเมธเคสูติ.\csromanspacer{}\csromanspacer{}\textbf{เอตมฺปิ ทิสฺวา วิรเม กโถชฺชนฺ}ติ กโถชฺชํ วุจฺจติ กลโห ภณฺฑนํ วิคฺคโห วิวาโท เมธคํ.\csromanspacer{} อถ วา \textbf{กโถชฺชนฺ}ติ อโนชวนฺตี สา กถา กโถชฺชํ น กเรยฺย, กลหํ น กเรยฺย, ภณฺฑนํ น กเรยฺย, วิคฺคหํ น กเรยฺย, วิวาทํ น กเรยฺย, เมธคํ น กเรยฺย, กลหภณฺฑนวิคฺคหวิวาทเมธคํ ปชเหยฺย วิโนเทยฺย พฺยนฺติํ กเรยฺย อนภาวํ คเมยฺย, กลหภณฺฑนวิคฺคหวิวาทเมธคา อารโต อสฺส วิรโต ปฏิวิรโต}

\vspace{\tipitakaparskip}
\csromancenter{ปสูรสุตฺตนิทฺเทส นิกฺขนฺโต นิสฺสโฏ วิปฺปยุตฺโต วิสญฺญุตฺโต วิมริยาทิกเตน เจตสา วิหเรยฺยาติ เอตมฺปิ ทิสฺวา วิรเม กโถชฺชํ.}
\prose{\csromanfolio{131}\textbf{น ห’ญฺญทตฺถ’ตฺถิ ปสํสลาภา}ติ ปสํสลาภา อญฺโญ อตฺโถ นตฺถิ อตฺตตฺโถ วา ปรตฺโถ วา อุภยตฺโถ วา ทิฏฺฐธมฺมิโก วา อตฺโถ, สมฺปรายิโก วา อตฺโถ, อุตฺตาโน วา อตฺโถ, คมฺภีโร วา อตฺโถ, คูโฬฺห วา อตฺโถ, ปฏิจฺฉนฺโน วา อตฺโถ, เนยฺโย วา อตฺโถ, นีโต วา อตฺโถ, อนวชฺโช วา อตฺโถ, นิกฺกิเลโส วา อตฺโถ, โวทาโน วา อตฺโถ, ปรมตฺโถ วา อตฺโถ, นตฺถิ น สนฺติ น สํวิชฺชนฺติ นุปลพฺภนฺตีติ น ห’ญฺญทตฺถ’ตฺถิ ปสํสลาภา.\csromanspacer{} เตนาห ภควา\csromandash{}}

\csromangathameasure{เอเต วิวาทา สมเณสุ ชาตา, \\ เอเตสุ อุคฺฆาติ นิฆาติ โหติ. \\ เอตมฺปิ ทิสฺวา วิรเม กโถชฺชํ,}
\csromangathagroup{\csromangathastanza{เอเต วิวาทา สมเณสุ ชาตา, \\ เอเตสุ อุคฺฆาติ นิฆาติ โหติ. \\ เอตมฺปิ ทิสฺวา วิรเม กโถชฺชํ,}}

\prose{\textbf{น ห’ญฺญทตฺถ’ตฺถิ ปสํสลาภา}ติ.}

\proseitem{64}{\csromansymbolfootnote{*}{ขุ 1. 408 ปิฏฺเฐ.}ปสํสิโต วา ปน ตตฺถ โหติ,}

\prose{\textbf{อกฺขาย วาทํ ปริสาย มชฺเฌ.}}

\csromangathameasure{\textbf{โส} \textbf{หสฺสตี อุนฺนมตี} \textbf{จ เตน,} \\ \textbf{ปปฺปุยฺย ตมตฺถํ ยถา มโน อหุ.}}
\csromangathagroup{\csromangathastanza{\textbf{โส}\footnote{โส ตํ (สี)} \textbf{หสฺสตี อุนฺนมตี}\footnote{อุณฺณมตี (สฺยา, ก)} \textbf{จ เตน,} \\ \textbf{ปปฺปุยฺย ตมตฺถํ ยถา มโน อหุ.}}}

\prose{\textbf{ปสํสิโต วา ปน ตตฺถ โหตี}ติ \textbf{ตตฺถา}ติ สกาย ทิฏฺฐิยา สกาย ขนฺติยา สกาย รุจิยา สกาย ลทฺธิยา ปสํสิโต โถมิโต กิตฺติโต วณฺณิโต โหตีติ ปสํสิโต วา ปน ตตฺถ โหติ.}

\prose{\textbf{อกฺขาย วาทํ ปริสาย มชฺเฌ}ติ ขตฺติยปริสาย วา พฺราหฺมณปริสาย วา คหปติปริสาย วา สมณปริสาย วา มชฺเฌ อตฺตโน วาทํ อกฺขาย อาจิกฺขิตฺวา อนุวาทํ อกฺขาย อาจิกฺขิตฺวา ถมฺภยิตฺวา พฺรูหยิตฺวา ทีปยิตฺวา โชตยิตฺวา โวหริตฺวา ปริคฺคณฺหิตฺวาติ อกฺขาย วาทํ ปริสาย มชฺเฌ.}

\prose{\textbf{โส หสฺสตี อุนฺนมตี จ เตนา}ติ โส เตน ชยตฺเถน ตุฏฺโฐ โหติ หฏฺโฐ ปหฏฺโฐ อตฺตมโน ปริปุณฺณสงฺกปฺโป.\csromanspacer{} อถ วา ทนฺตวิทํสกํ หสมาโน.\csromanspacer{} \textbf{โส หสฺสตี อุนฺนมตี จ เตนา}ติ โส เตน \csromanfolio{132}ชยตฺเถน อุนฺนโต โหติ อุนฺนโม ธโช สมฺปคฺคาโห เกตุกมฺยตา จิตฺตสฺสาติ โส หสฺสตี อุนฺนมตี จ เตน.}

\prose{\textbf{ปปฺปุยฺย ตมตฺถํ ยถา มโน อหู}ติ ตํ ชยตฺถํ ปปฺปุยฺย ปาปุณิตฺวา อธิคนฺตฺวา วินฺทิตฺวา ปฏิลภิตฺวา.\csromanspacer{}\csromanspacer{}\textbf{ยถา มโน อหู}ติ ยถา มโน อหุ, ยถา จิตฺโต อหุ, ยถา สงฺกปฺโป อหุ, ยถา วิญฺญาโณ อหูติ ปปฺปุยฺย ตมตฺถํ มโน อหุ.\csromanspacer{} เตนาห ภควา\csromandash{}}

\csromangathameasure{ปสํสิโต วา ปน ตตฺถ โหติ, \\ อกฺขาย วาทํ ปริสาย มชฺเฌ. \\ โส หสฺสตี อุนฺนมตี จ เตน,}
\csromangathagroup{\csromangathastanza{ปสํสิโต วา ปน ตตฺถ โหติ, \\ อกฺขาย วาทํ ปริสาย มชฺเฌ. \\ โส หสฺสตี อุนฺนมตี จ เตน,}}

\prose{\textbf{ปปฺปุยฺย ตมตฺถํ ยถา มโน อหู}ติ.}

\proseitem{65}{\csromansymbolfootnote{*}{ขุ 1. 408 ปิฏฺเฐ.}ยา อุนฺนตี สา’สฺส วิฆาตภูมิ,}

\prose{\textbf{มานาติมานํ วทเต ปเนโส.}}

\csromangathameasure{\textbf{เอตมฺปิ ทิสฺวา น วิวาทเยถ,} \\ \textbf{น หิ เตน สุทฺธิํ กุสลา วทนฺติ.}}
\csromangathagroup{\csromangathastanza{\textbf{เอตมฺปิ ทิสฺวา น วิวาทเยถ,} \\ \textbf{น หิ เตน สุทฺธิํ กุสลา วทนฺติ.}}}

\prose{\textbf{ยา อุนฺนตี สา’สฺส วิฆาตภูมี}ติ ยา อุนฺนติ อุนฺนโม ธโช สมฺปคฺคาโห เกตุกมฺยตา จิตฺตสฺสาติ ยา อุนฺนติ.\csromanspacer{} \textbf{สา’สฺส วิฆาตภูมี}ติ สา ตสฺส วิฆาตภูมิ อุปฆาตภูมิ ปีฬนภูมิ ฆฏฺฏนภูมิ อุปทฺทวภูมิ อุปสคฺคภูมีติ ยา อุนฺนตี สา’สฺส วิฆาตภูมิ.}

\prose{\textbf{มานาติมานํ วทเต ปเนโส}ติ โส ปุคฺคโล มานญฺจ วทติ อติมานญฺจ วทตีติ มานาติมานํ วทเต ปเนโส.}

\prose{\textbf{เอตมฺปิ ทิสฺวา น วิวาทเยถา}ติ เอตํ อาทีนวํ ทิสฺวา ปสฺสิตฺวา ตุลยิตฺวา ตีรยิตฺวา วิภาวยิตฺวา วิภูตํ กตฺวา ทิฏฺฐิกลเหสุ ทิฏฺฐิภณฺฑเนสุ ทิฏฺฐิวิคฺคเหสุ ทิฏฺฐิวิวาเทสุ ทิฏฺฐิเมธเคสูติ เอตมฺปิ ทิสฺวา.\csromanspacer{} \textbf{น วิวาทเยถา}ติ น กลหํ กเรยฺย, น ภณฺฑนํ กเรยฺย, น วิคฺคหํ กเรยฺย, น วิวาทํ กเรยฺย, น เมธคํ กเรยฺย, กลหภณฺฑนวิคฺคหวิวาทเมธคํ ปชเหยฺย วิโนเทยฺย พฺยนฺติํ กเรยฺย อนภาวํ คเมยฺย, กลหภณฺฑนวิคฺคหวิวาทเมธคา อารโต อสฺส วิรโต ปฏิวิรโต นิกฺขนฺโต นิสฺสโฏ วิปฺปยุตฺโต วิสญฺญุตฺโต วิมริยาทิกเตน เจตสา วิหเรยฺยาติ เอตมฺปิ ทิสฺวา น วิวาทเยถ.}

\csromanheader{2}{8. ปสูรสุตฺตนิทฺเทส}
\prose{\csromanfolio{133}\textbf{น หิ เตน สุทฺธิํ กุสลา วทนฺตี}ติ \textbf{กุสลา}ติ เย เต ขนฺธ\textbf{กุสลา} ธาตุ\textbf{กุสลา} อายตน\textbf{กุสลา} ปฏิจฺจสมุปฺปาท\textbf{กุสลา} สติปฏฺฐาน\textbf{กุสลา} สมฺมปฺปธาน\textbf{กุสลา} อิทฺธิปาท\textbf{กุสลา} อินฺทฺริย\textbf{กุสลา} พล\textbf{กุสลา} โพชฺฌงฺค\textbf{กุสลา} มคฺค\textbf{กุสลา} ผล\textbf{กุสลา} นิพฺพาน\textbf{กุสลา}, เต \textbf{กุสลา} ทิฏฺฐิกลเหน ทิฏฺฐิภณฺฑเนน ทิฏฺฐิวิคฺคเหน ทิฏฺฐิวิวาเทน ทิฏฺฐิเมธเคน สุทฺธิํ วิสุทฺธิํ ปริสุทฺธิํ, มุตฺติํ วิมุตฺติํ ปริมุตฺติํ น วทนฺติ น กเถนฺติ น ภณนฺติ น ทีปยนฺติ น โวหรนฺตีติ น หิ เตน สุทฺธิํ \textbf{กุสลา} วทนฺติ.\csromanspacer{} เตนาห ภควา\csromandash{}}

\csromangathameasure{ยา อุนฺนตี สา’สฺส วิฆาตภูมิ, \\ มานาติมานํ วทเต ปเนโส. \\ เอตมฺปิ ทิสฺวา น วิวาทเยถ,}
\csromangathagroup{\csromangathastanza{ยา อุนฺนตี สา’สฺส วิฆาตภูมิ, \\ มานาติมานํ วทเต ปเนโส. \\ เอตมฺปิ ทิสฺวา น วิวาทเยถ,}}

\prose{\textbf{น หิ เตน สุทฺธิํ กุสลา วทนฺตี}ติ.}

\proseitem{66}{\csromansymbolfootnote{*}{ขุ 1. 408 ปิฏฺเฐ.}\textbf{สูโร ยถา ราชขาทาย ปุฏฺโฐ},}

\prose{\textbf{อภิคชฺช’เมติ ปฏิสูร’มิจฺฉํ.}}

\csromangathameasure{\textbf{เยเนว โส, เตน ปเลหิ สูร,} \\ \textbf{ปุพฺเพว นตฺถิ ยทิทํ ยุธาย.}}
\csromangathagroup{\csromangathastanza{\textbf{เยเนว โส, เตน ปเลหิ สูร,} \\ \textbf{ปุพฺเพว นตฺถิ ยทิทํ ยุธาย.}}}

\prose{\textbf{สูโร ยถา ราชขาทาย ปุฏฺโฐ}ติ \textbf{สูโร}ติ \textbf{สูโร} วีโร วิกฺกนฺโต อภีรู อฉมฺภี อนุตฺราสี อปลายี.\csromanspacer{}\csromanspacer{}\textbf{ราชขาทาย ปุฏฺโฐ}ติ ราชขาทนีเยน ราชโภชนีเยน ปุฏฺโฐ โปสิโต อาปาทิโต วฑฺฒิโตติ \textbf{สูโร} ยถา ราชขาทาย ปุฏฺโฐ.}

\prose{\textbf{อภิคชฺช’เมติ ปฏิสูร’มิจฺฉนฺ}ติ โส คชฺชนฺโต อุคฺคชฺชนฺโต อภิคชฺชนฺโต เอติ อุเปติ อุปคจฺฉติ ปฏิสูรํ ปฏิปุริสํ ปฏิสตฺตุํ ปฏิมลฺลํ อิจฺฉนฺโต สาทิยนฺโต ปตฺถยนฺโต ปิหยนฺโต อภิชปฺปนฺโตติ อภิคชฺช’เมติ ปฏิสูร’มิจฺฉํ.}

\prose{\textbf{เยเนว โส, เตน ปเลหิ สูรา}ติ เยเนว โส ทิฏฺฐิคติโก, เตน ปเลหิ, เตน วช, เตน คจฺฉ, เตน อติกฺกม, โส ตุยฺหํ ปฏิ\textbf{สูโร} ปฏิปุริโส ปฏิสตฺตุ ปฏิมลฺโลติ เยเนว โส, เตน ปเลหิ สูร.}

\prose{\textbf{ปุพฺเพว นตฺถิ ยทิทํ ยุธายา}ติ ปุพฺเพว โพธิยา มูเล เย ปฏิเสนิกรา กิเลสา ปฏิโลมกรา ปฏิกณฺฑกกรา ปฏิปกฺขกรา, เต นตฺถิ \csromanfolio{134}น สนฺติ น สํวิชฺชนฺติ นุปลพฺภนฺติ, ปหีนา สมุจฺฉินฺนา วูปสนฺตา ปฏิปสฺสทฺธา อภพฺพุปฺปตฺติกา ญาณคฺคินา ทฑฺฒา.\csromanspacer{}\csromanspacer{}\textbf{ยทิทํ ยุธายา}ติ ยทิทํ ยุทฺธตฺถาย กลหตฺถาย ภณฺฑนตฺถาย วิคฺคหตฺถาย วิวาทตฺถาย เมธคตฺถายาติ ปุพฺเพว นตฺถิ ยทิทํ ยุธาย.\csromanspacer{} เตนาห ภควา\csromandash{}}

\csromangathameasure{สูโร ยถา ราชขาทาย ปุฏฺโฐ, \\ อภิคชฺช’เมติ ปฏิสูร’มิจฺฉํ. \\ เยเนว โส, เตน ปเลหิ สูร, \\ ปุพฺเพว นตฺถิ ยทิทํ ยุธายาติ.}
\csromangathagroup{\csromangathastanza{สูโร ยถา ราชขาทาย ปุฏฺโฐ, \\ อภิคชฺช’เมติ ปฏิสูร’มิจฺฉํ. \\ เยเนว โส, เตน ปเลหิ สูร, \\ ปุพฺเพว นตฺถิ ยทิทํ ยุธายาติ.}}

\proseitem{67}{\csromansymbolfootnote{*}{ขุ 1. 409 ปิฏฺเฐ.}เย ทิฏฺฐิมุคฺคยฺห วิวาทยนฺติ,}

\prose{\textbf{“อิทเมว สจฺจนฺ”ติ จ วาทยนฺติ.}}

\csromangathameasure{\textbf{เต ตฺวํ วทสฺสู น หิ เตธ อตฺถิ,} \\ \textbf{วาทมฺหิ ชาเต ปฏิเสนิกตฺตา.}}
\csromangathagroup{\csromangathastanza{\textbf{เต ตฺวํ วทสฺสู น หิ เตธ อตฺถิ,} \\ \textbf{วาทมฺหิ ชาเต ปฏิเสนิกตฺตา.}}}

\prose{\textbf{เย ทิฏฺฐิมุคฺคยฺห วิวาทยนฺตี}ติ เย ทฺวาสฏฺฐิทิฏฺฐิคตานํ อญฺญตรญฺญตรํ ทิฏฺฐิคตํ คเหตฺวา คณฺหิตฺวา อุคฺคณฺหิตฺวา ปรามสิตฺวา อภินิวิสิตฺวา วิวาทยนฺติ, กลหํ กโรนฺติ, ภณฺฑนํ กโรนฺติ, วิคฺคหํ กโรนฺติ, วิวาทํ กโรนฺติ, เมธคํ กโรนฺติ “น ตฺวํ อิมํ ธมฺมวินยํ อาชานาสิ, อหํ อิมํ ธมฺมวินยํ อาชานามิ, กิํ ตฺวํ อิมํ ธมฺมวินยํ อาชานิสฺสสิ, มิจฺฉาปฏิปนฺโน ตฺวมสิ, อหมสฺมิ สมฺมาปฏิปนฺโน, สหิตํ เม, อสหิตํ เต, ปุเร วจนียํ ปจฺฉา อวจ, ปจฺฉา วจนียํ ปุเร อวจ, อธิจิณฺณํ เต วิปราวตฺตํ, อาโรปิโต เต วาโท, นิคฺคหิโต ตฺวมสิ, จร วาทปฺปโมกฺขาย, นิพฺเพเฐหิ วา สเจ ปโหสี”ติ เย ทิฏฺฐิมุคฺคยฺห วิวาทยนฺติ.}

\prose{\textbf{“อิทเมว สจฺจนฺ”ติ จ วาทยนฺตี}ติ “สสฺสโต โลโก, อิทเมว สจฺจํ โมฆมญฺญนฺ”ติ วาทยนฺติ กเถนฺติ ภณนฺติ ทีปยนฺติ โวหรนฺติ.\csromanspacer{} อสสฺสโต โลโก ฯเปฯ เนว โหติ น น โหติ ตถาคโต ปรํ มรณา, อิทเมว สจฺจํ โมฆมญฺญนฺ”ติ วาทยนฺติ กเถนฺติ ภณนฺติ ทีปยนฺติ โวหรนฺตีติ “อิทเมว สจฺจนฺ”ติ จ วาทยนฺติ.}

\prose{\textbf{เต ตฺวํ วทสฺสู น หิ เตธ อตฺถิ, วาทมฺหิ ชาเต ปฏิเสนิกตฺตา}ติ เต ตฺวํ ทิฏฺฐิคติเก วทสฺสุ วาเทน วาทํ นิคฺคเหน นิคฺคหํ ปฏิกมฺเมน ปฏิกมฺมํ วิเสเสน วิเสสํ ปฏิวิเสเสน ปฏิวิเสสํ อาเวฐิยาย}

\vspace{\tipitakaparskip}
\csromancenter{ปสูรสุตฺตนิทฺเทส อาเวฐิยํ นิพฺเพฐิยาย นิพฺเพฐิยํ เฉเทน เฉทํ มณฺฑเลน มณฺฑลํ, เต ตุยฺหํ ปฏิสูรา ปฏิปุริสา ปฏิสตฺตู ปฏิมลฺลาติ เต ตฺวํ วทสฺสู น หิ เตธ อตฺถิ.\csromanspacer{} \textbf{วาทมฺหิ ชาเต ปฏิเสนิกตฺตา}ติ วาเท ชาเต สญฺชาเต นิพฺพตฺเต อภินิพฺพตฺเต ปาตุภูเตเยว ปฏิเสนิกตฺตา\footnote{ปฏิเสนิกตา (ก), เอวํ เสเสสุ ตีสุ ปเทสุปิ.} ปฏิโลมกตฺตา ปฏิกณฺฑกกตฺตา ปฏิปกฺขกตฺตา กลหํ กเรยฺยุํ, ภณฺฑนํ กเรยฺยุํ, วิคฺคหํ กเรยฺยุํ, เมธคํ กเรยฺยุํ, เต นตฺถิ น สนฺติ น สํวิชฺชนฺติ นุปลพฺภนฺติ ปหีนา ฯเปฯ ญาณคฺคินา ทฑฺฒาติ เต ตฺวํ วทสฺสู น หิ เตธ อตฺถิ, วาทมฺหิ ชาเต ปฏิเสนิกตฺตา.\csromanspacer{} เตนาห ภควา\csromandash{}}
\vspace{0.5\baselineskip}
\csromangathameasure{\textbf{เย} ทิฏฺฐิมุคฺคยฺห วิวาทยนฺติ, \\ “อิทเมว สจฺจนฺ”ติ จ วาทยนฺติ. \\ เต ตฺวํ วทสฺสู น หิ เตธ อตฺถิ, \\ \textbf{วาทมฺหิ ชาเต ปฏิเสนิกตฺตา}ติ.}
\csromangathagroup{\csromangathastanza{\csromanfolio{135}\textbf{เย} ทิฏฺฐิมุคฺคยฺห วิวาทยนฺติ, \\ “อิทเมว สจฺจนฺ”ติ จ วาทยนฺติ. \\ เต ตฺวํ วทสฺสู น หิ เตธ อตฺถิ, \\ \textbf{วาทมฺหิ ชาเต ปฏิเสนิกตฺตา}ติ.}}

\proseitem{68}{\csromansymbolfootnote{*}{ขุ 1. 409 ปิฏฺเฐ.}วิเสนิกตฺวา ปน เย จรนฺติ,}

\prose{\textbf{ทิฏฺฐีหิ ทิฏฺฐิํ อวิรุชฺฌามานา.}}

\csromangathameasure{\textbf{เตสุ ตฺวํ กิํ ลเภโถ ปสูร,} \\ \textbf{เยสีธ นตฺถิ ปรมุคฺคหีตํ.}}
\csromangathagroup{\csromangathastanza{\textbf{เตสุ ตฺวํ กิํ ลเภโถ ปสูร,} \\ \textbf{เยสีธ นตฺถิ ปรมุคฺคหีตํ.}}}

\prose{\textbf{วิเสนิกตฺวา ปน เย จรนฺตี}ติ เสนา วุจฺจติ มารเสนา.\csromanspacer{} กายทุจฺจริตํ มารเสนา, วจีทุจฺจริตํ มารเสนา, มโนทุจฺจริตํ มารเสนา, โลโภ มารเสนา, โทโส มารเสนา, โมโห มารเสนา, โกโธ.\csromanspacer{} อุปนาโห.\csromanspacer{} มกฺโข.\csromanspacer{} ปฬาโส.\csromanspacer{} อิสฺสา.\csromanspacer{} มจฺฉริยํ.\csromanspacer{} มายา.\csromanspacer{} สาเฐยฺยํ.\csromanspacer{} ถมฺโภ.\csromanspacer{} สารมฺโภ.\csromanspacer{} มาโน.\csromanspacer{} อติมาโน.\csromanspacer{} มโท.\csromanspacer{} ปมาโท.\csromanspacer{} สพฺเพ กิเลสา.\csromanspacer{} สพฺเพ ทุจฺจริตา.\csromanspacer{} สพฺเพ ทรถา.\csromanspacer{} สพฺเพ ปริฬาหา.\csromanspacer{} สพฺเพ สนฺตาปา.\csromanspacer{} สพฺพากุสลาภิสงฺขารา มารเสนา.\csromanspacer{} วุตฺตํ เหตํ ภควตา\csromandash{}}

\csromangathameasure{+กามา เต ปฐมา เสนา, ทุติยา อรติ วุจฺจติ ฯเปฯ. \\ น นํ อสูโร ชินาติ, เชตฺวาว ลภเต สุขนฺติ.}
\csromangathagroup{\csromangathastanza{\csromansymbolfootnote{+}{ขุ 1. 342; ขุ 8. 114 ปิฏฺเฐสุปิ.}กามา เต ปฐมา เสนา, ทุติยา อรติ วุจฺจติ ฯเปฯ. \\ น นํ อสูโร ชินาติ, เชตฺวาว ลภเต สุขนฺติ.}}

\prose{ยโต จตูหิ อริยมคฺเคหิ สพฺพา จ มารเสนา สพฺเพ จ ปฏิเสนิกรา กิเลสา ชิตา จ ปราชิตา จ ภคฺคา วิปฺปลุคฺคา ปรมฺมุขา, เตน วุจฺจติ วิเสนิกตฺวาติ.\csromanspacer{}\csromanspacer{}\textbf{เย}ติ อรหนฺโต ขีณาสวา.\csromanspacer{} \textbf{จรนฺตี}ติ จรนฺติ \csromanfolio{136}วิหรนฺติ อิริยนฺติ วตฺเตนฺติ ปาเลนฺติ ยเปนฺติ ยาเปนฺตีติ วิเสนิกตฺวา ปน เย จรนฺติ.}

\prose{\textbf{ทิฏฺฐีหิ ทิฏฺฐิํ อวิรุชฺฌมานา}ติ เยสํ ทฺวาสฏฺฐิ ทิฏฺฐิคตานิ ปหีนานิ สมุจฺฉินฺนานิ วูปสนฺตานิ ปฏิปสฺสทฺธานิ อภพฺพุปฺปตฺติกานิ ญาณคฺคินา ทฑฺฒานิ\textbf{,} เต ทิฏฺฐีหิ ทิฏฺฐิํ อวิรุชฺฌมานา อปฺปฏิวิรุชฺฌมานา อปฺปหียมานา อปฺปฏิหญฺญมานา อปฺปฏิหตมานาติ ทิฏฺฐีหิ ทิฏฺฐิํ อวิรุชฺฌมานา.}

\prose{\textbf{เตสุ ตฺวํ กิํ ลเภโถ ปสูรา}ติ เตสุ อรหนฺเตสุ ขีณาสเวสุ กิํ ลเภโถ ปฏิสูรํ ปฏิปุริสํ ปฏิสตฺตุํ ปฏิมลฺลนฺติ เตสุ ตฺวํ กิํ ลเภโถ ปสูร.}

\prose{\textbf{เยสีธ นตฺถิ ปรมุคฺคหีตนฺ}ติ เยสํ อรหนฺตานํ ขีณาสวานํ “อิทํ ปรมํ อคฺคํ เสฏฺฐํ วิสิฏฺฐํ ปาโมกฺขํ อุตฺตมํ ปวรนฺ”ติ คหิตํ ปรามฏฺฐํ อภินิวิฏฺฐํ อชฺโฌสิตํ อธิมุตฺตํ นตฺถิ น สนฺติ น สํวิชฺชติ นุปลพฺภติ\textbf{,} ปหีนํ สมุจฺฉินฺนํ วูปสนฺตํ ปฏิปสฺสทฺธํ อภพฺพุปฺปตฺติกํ ญาณคฺคินา ทฑฺฒนฺติ เยสีธ นตฺถิ ปรมุคฺคหีตํ.\csromanspacer{} เตนาห ภควา\csromandash{}}

\prose{วิเสนิกตฺวา ปน เย จรนฺติ\textbf{,}}

\prose{\textbf{ทิฏฺฐีหิ ทิฏฺฐิํ อวิรุชฺฌมานา}.}

\prose{เตสุ ตฺวํ กิํ ลเภโถ ปสูร\textbf{,}}

\prose{\textbf{เยสีธ นตฺถิ ปรมุคฺคหีตนฺ}ติ.}

\proseitem{69}{\csromansymbolfootnote{*}{ขุ 1. 409 ปิฏฺเฐ.}\textbf{อถ ตฺวํ ปวิตกฺกมาคมา}\footnote{ปวิตกฺกมาคม (สี), สวิตกฺกมาคมา (ก)}\textbf{,}}

\prose{\textbf{มนสา ทิฏฺฐิคตานิ จินฺตยนฺโต.}}

\csromangathameasure{\textbf{โธเนน ยุคํ สมาคมา,} \\ \textbf{น หิ ตฺวํ สกฺขสิ สมฺปยาตเว.}}
\csromangathagroup{\csromangathastanza{\textbf{โธเนน ยุคํ สมาคมา,} \\ \textbf{น หิ ตฺวํ สกฺขสิ สมฺปยาตเว.}}}

\prose{\textbf{อถ ตฺวํ ปวิตกฺกมาคมา}ติ \textbf{อถา}ติ ปทสนฺธิ ปทสํสคฺโค ปทปาริปูรี อกฺขรสมวาโย พฺยญฺชนสิลิฏฺฐตา ปทานุปุพฺพตาเปตํ \textbf{อถา}ติ.\csromanspacer{}\csromanspacer{}\textbf{ปวิตกฺกมาคมา}ติ ตกฺเกนฺโต วิตกฺเกนฺโต สงฺกปฺเปนฺโต ชโย นุ โข เม ภวิสฺสติ\textbf{,} ปราชโย นุ โข เม ภวิสฺสติ\textbf{,} กถํ นิคฺคหํ กริสฺสามิ\textbf{,} กถํ ปฏิกมฺมํ กริสฺสามิ\textbf{,} กถํ วิเสสํ กริสฺสามิ\textbf{,} กถํ ปฏิวิเสสํ กริสฺสามิ\textbf{,} กถํ อาเวฐิยํ กริสฺสามิ\textbf{,} กถํ นิพฺเพฐิยํ}

\vspace{\tipitakaparskip}
\csromancenter{ปสูรสุตฺตนิทฺเทส กริสฺสามิ, กถํ เฉทํ กริสฺสามิ, กถํ มณฺฑลํ กริสฺสามิ, เอวํ ตกฺเกนฺโต วิตกฺเกนฺโต สงฺกปฺเปนฺโต อาคโตสิ อุปคโตสิ สมฺปตฺโตสิ มยา สทฺธิํ สมาคโตสีติ อถ ตฺวํ ปวิตกฺกมาคมา.}
\prose{\csromanfolio{137}\textbf{มนสา ทิฏฺฐิคตานิ จินฺตยนฺโต}ติ \textbf{มโน}ติ * ยํ จิตฺตํ \textbf{มโน} มานสํ หทยํ ปณฺฑรํ \textbf{มโน} มนายตนํ มนินฺทฺริยํ วิญฺญาณํ วิญฺญาณกฺขนฺโธ ตชฺชา \textbf{มโน}วิญฺญาณธาตุ.\csromanspacer{} จิตฺเตน ทิฏฺฐิํ จินฺเตนฺโต วิจินฺเตนฺโต “สสฺสโต โลโก”ติ วา, “อสสฺสโต โลโก”ติ วา ฯเปฯ “เนว โหติ น น โหติ ตถาคโต ปรํ มรณา”ติ วาติ มนสา ทิฏฺฐิคตานิ จินฺตยนฺโต.}

\prose{\textbf{โธเนน ยุคํ สมาคมา, น หิ ตฺวํ สกฺขสิ สมฺปยาตเว}ติ \textbf{โธนา} วุจฺจติ ปญฺญา.\csromanspacer{} ยา ปญฺญา ปชานนา ฯเปฯ อโมโห ธมฺมวิจโย สมฺมาทิฏฺฐิ กิํการณา \textbf{โธนา} วุจฺจติ ปญฺญา, ตาย ปญฺญาย กายทุจฺจริตํ ธุตญฺจ โธตญฺจ สนฺโธตญฺจ นิทฺโธตญฺจ, วจีทุจฺจริตํ ฯเปฯ สพฺพากุสลาภิสงฺขารา ธุตา จ โธตา จ สนฺโธตา จ นิทฺโธตา จ.\csromanspacer{} อถ วา สมฺมาทิฏฺฐิยา มิจฺฉาทิฏฺฐิ.\csromanspacer{} สมฺมาสงฺกปฺเปน มิจฺฉาสงฺกปฺโป ฯเปฯ สมฺมาวิมุตฺติยา มิจฺฉาวิมุตฺติ ธุตา จ โธตา จ สนฺโธตา จ นิทฺโธตา จ.\csromanspacer{} อถ วา อริเยน อฏฺฐงฺคิเกน มคฺเคน สพฺเพ กิเลสา.\csromanspacer{} สพฺเพ ทุจฺจริตา.\csromanspacer{} สพฺเพ ทรถา.\csromanspacer{} สพฺเพ ปริฬาหา.\csromanspacer{} สพฺเพ สนฺตาปา.\csromanspacer{} สพฺพากุสลาภิสงฺขารา ธุตา จ โธตา จ สนฺโธตา จ นิทฺโธตา จ.\csromanspacer{} ภควา อิเมหิ โธเนยฺเยหิ ธมฺเมหิ อุเปโต สมุเปโต อุปคโต สมุปคโต อุปปนฺโน สมุปปนฺโน สมนฺนาคโต, ตสฺมา ภควา โธโน.\csromanspacer{} โส ธุตราโค ธุตปาโป ธุตกิเลโส ธุตปริฬาโหติ โธโนติ.}

\prose{\textbf{โธเนน ยุคํ สมาคมา, น หิ ตฺวํ สกฺขสิ สมฺปยาตเว}ติ ปสูโร ปริพฺพาชโก น ปฏิพโล โธเนน พุทฺเธน ภควตา สทฺธิํ ยุคํ สมาคมํ สมาคนฺตฺวา ยุคคฺคาหํ คณฺหิตฺวา สากจฺเฉตุํ สลฺลปิตุํ สากจฺฉํ สมาปชฺชิตุํ.\csromanspacer{} ตํ กิสฺส เหตุ, ปสูโร ปริพฺพาชโก หีโน นิหีโน โอมโก ลามโก ฉตุกฺโก ปริตฺโต.\csromanspacer{} โส หิ ภควา อคฺโค จ เสฏฺโฐ จ วิสิฏฺโฐ จ ปาโมกฺโข จ อุตฺตโม จ ปวโร จ.\csromanspacer{} ยถา สโส น ปฏิพโล มตฺเตน มาตงฺเคน สทฺธิํ ยุคํ สมาคมํ สมาคนฺตฺวา ยุคคฺคาหํ คณฺหิตุํ.\csromanspacer{}\csromanspacer{}ยถา โกตฺถุโก น ปฏิพโล สีเหน มิครญฺญา สทฺธิํ ยุคํ สมาคมํ สมาคนฺตฺวา ยุคคฺคาหํ คณฺหิตุํ.\csromanspacer{}\csromanspacer{}ยถา วจฺฉโก ตรุณโก \csromanfolio{138}เธนุปโก น ปฏิพโล อุสเภน จลกกุนา\footnote{พลกฺกกุนา (สฺยา)} สทฺธิํ ยุคํ สมาคมํ สมาคนฺตฺวา ยุคคฺคาหํ คณฺหิตุํ.\csromanspacer{}\csromanspacer{}ยถา ธงฺโก น ปฏิพโล ครุเฬน เวนเตยฺเยน สทฺธิํ ยุคํ สมาคมํ สมาคนฺตฺวา ยุคคฺคาหํ คณฺหิตุํ.\csromanspacer{}\csromanspacer{}ยถา จณฺฑาโล น ปฏิพโล รญฺญา จกฺกวตฺตินา สทฺธิํ ยุคํ สมาคมํ สมาคนฺตฺวา ยุคคฺคาหํ คณฺหิตุํ.\csromanspacer{}\csromanspacer{}ยถา ปํสุปิสาจโก น ปฏิพโล อินฺเทน เทวรญฺญา สทฺธิํ ยุคํ สมาคมํ สมาคนฺตฺวา ยุคคฺคาหํ คณฺหิตุํ, เอวเมวํ ปสูโร ปริพฺพาชโก น ปฏิพโล โธเนน พุทฺเธน ภควตา สทฺธิํ ยุคํ สมาคมํ สมาคนฺตฺวา ยุคคฺคาหํ คณฺหิตฺวา สากจฺเฉตุํ สลฺลปิตุํ สากจฺฉํ สมาปชฺชิตุํ.\csromanspacer{} ตํ กิสฺส เหตุ, ปสูโร ปริพฺพาชโก หีนปญฺโญ นิหีนปญฺโญ โอมกปญฺโญ ลามกปญฺโญ ฉตุกฺกปญฺโญ ปริตฺตปญฺโญ.\csromanspacer{} โส หิ ภควา มหาปญฺโญ ปุถุปญฺโญ หาสปญฺโญ ชวนปญฺโญ ติกฺขปญฺโญ นิพฺเพธิกปญฺโญ ปญฺญาปเภทกุสโล ปภินฺนญาโณ อธิคตปฏิสมฺภิโท จตุเวสารชฺชปฺปตฺโต ทสพลธารี ปุริสาสโภ ปุริสสีโห ปุริสนาโค ปุริสาชญฺโญ ปุริสโธรโยฺห อนนฺตญาโณ อนนฺตเตโช อนนฺตยโส อฑฺโฒ มหทฺธโน ธนวา เนตา วิเนตา อนุเนตา ปญฺญาเปตา นิชฺฌาเปตา เปกฺเขตา ปสาเทตา.\csromanspacer{} โส หิ ภควา อนุปฺปนฺนสฺส มคฺคสฺส อุปฺปาเทตา อสญฺชาตสฺส มคฺคสฺส สญฺชเนตา อนกฺขาตสฺส มคฺคสฺส อกฺขาตา มคฺคญฺญู มคฺควิทู มคฺคโกวิโท มคฺคานุคา จ ปนสฺส เอตรหิ สาวกา วิหรนฺติ ปจฺฉา สมนฺนาคตา.}

\prose{โส หิ ภควา ชานํ ชานาติ ปสฺสํ ปสฺสติ จกฺขุภูโต ญาณภูโต ธมฺมภูโต พฺรหฺมภูโต วตฺตา ปวตฺตา อตฺถสฺส นินฺเนตา อมตสฺส ทาตา ธมฺมสฺสามี ตถาคโต, นตฺถิ ตสฺส ภควโต อญฺญาตํ อทิฏฺฐํ อวิทิตํ อสจฺฉิกตํ อผสฺสิตํ\footnote{อผุสิตํ (สฺยา) ขุ 8. 175; ขุ 9. 375 ปิฏฺเฐสุปิ.} ปญฺญาย.\csromanspacer{} อตีตํ อนาคตํ ปจฺจุปฺปนฺนํ อุปาทาย สพฺเพ ธมฺมา สพฺพากาเรน พุทฺธสฺส ภควโต ญาณมุเข อาปาถํ อาคจฺฉนฺติ, ยํ กิญฺจิ เนยฺยํ นาม อตฺถิ ชานิตพฺพํ อตฺตตฺโถ วา ปรตฺโถ วา อุภยตฺโถ วา ทิฏฺฐธมฺมิโก วา อตฺโถ สมฺปรายิโก วา อตฺโถ อุตฺตาโน วา อตฺโถ คมฺภีโร วา อตฺโถ คูโฬฺห วา อตฺโถ ปฏิจฺฉนฺโน วา อตฺโถ เนยฺโย วา อตฺโถ นีโต วา อตฺโถ อนวชฺโช วา อตฺโถ นิกฺกิเลโส วา อตฺโถ โวทาโน วา อตฺโถ ปรมตฺโถ วา อตฺโถ, สพฺพํ ตํ อนฺโตพุทฺธญาเณ ปริวตฺตติ.}

\csromanheader{2}{8. ปสูรสุตฺตนิทฺเทส}
\prose{\csromanfolio{139}สพฺพํ กายกมฺมํ พุทฺธสฺส ภควโต ญาณานุปริวตฺติ, สพฺพํ วจีกมฺมํ ญาณานุปริวตฺติ, สพฺพํ มโนกมฺมํ ญาณานุปริวตฺติ.\csromanspacer{} อตีเต พุทฺธสฺส ภควโต อปฺปฏิหตํ ญาณํ, อนาคเต อปฺปฏิหตํ ญาณํ, ปจฺจุปฺปนฺเน อปฺปฏิหตํ ญาณํ.\csromanspacer{} ยาวตกํ เนยฺยํ, ตาวตกํ ญาณํ, ยาวตกํ ญาณํ, ตาวตกํ เนยฺยํ, เนยฺยปริยนฺติกํ ญาณํ, ญาณปริยนฺติกํ เนยฺยํ, เนยฺยํ อติกฺกมิตฺวา ญาณํ นปฺปวตฺตติ, ญาณํ อติกฺกมิตฺวา เนยฺยปโถ นตฺถิ.\csromanspacer{} อญฺญมญฺญปริยนฺตฏฺฐายิโน เต ธมฺมา.\csromanspacer{} ยถา ทฺวินฺนํ สมุคฺคปฏลานํ สมฺมา ผุสิตานํ เหฏฺฐิมํ สมุคฺคปฏลํ อุปริมํ นาติวตฺตติ, อุปริมํ สมุคฺคปฏลํ เหฏฺฐิมํ นาติวตฺตติ, อญฺญมญฺญปริยนฺตฏฺฐายิโน.\csromanspacer{} เอวเมวํ พุทฺธสฺส ภควโต เนยฺยญฺจ ญาณญฺจ อญฺญมญฺญปริยนฺตฏฺฐายิโน.\csromanspacer{} ยาวตกํ เนยฺยํ ตาวตกํ ญาณํ, ยาวตกํ ญาณํ, ตาวตกํ เนยฺยํ, เนยฺยปริยนฺติกํ ญาณํ, ญาณปริยนฺติกํ เนยฺยํ, เนยฺยํ อติกฺกมิตฺวา ญาณํ นปฺปวตฺตติ, ญาณํ อติกฺกมิตฺวา เนยฺยปโถ นตฺถิ.\csromanspacer{} อญฺญมญฺญปริยนฺตฏฺฐายิโน เต ธมฺมา.\csromanspacer{} สพฺพธมฺเมสุ พุทฺธสฺส ภควโต ญาณํ ปวตฺตติ.}

\prose{สพฺเพ ธมฺมา พุทฺธสฺส ภควโต อาวชฺชนปฏิพทฺธา อากงฺขปฏิพทฺธา มนสิการปฏิพทฺธา จิตฺตุปฺปาทปฏิพทฺธา.\csromanspacer{} สพฺพสตฺเตสุ พุทฺธสฺส ภควโต ญาณํ ปวตฺตติ, สพฺเพสญฺจ สตฺตานํ ภควา อาสยํ ชานาติ อนุสยํ ชานาติ, จริตํ ชานาติ, อธิมุตฺติํ ชานาติ.\csromanspacer{} อปฺปรชกฺเข มหารชกฺเข ติกฺขินฺทฺริเย มุทินฺทฺริเย สฺวากาเร ทฺวากาเร สุวิญฺญาปเย ทุวิญฺญาปเย ภพฺพาภพฺเพ สตฺเต ปชานาติ.\csromanspacer{} สเทวโก โลโก สมารโก สพฺรหฺมโก สสฺสมณพฺราหฺมณี ปชา สเทวมนุสฺสา อนฺโตพุทฺธญาเณ ปริวตฺตติ.}

\prose{ยถา เย เกจิ มจฺฉกจฺฉปา อนฺตมโส ติมิติมิงฺคลํ อุปาทาย อนฺโตมหาสมุทฺเท ปริวตฺตนฺติ.\csromanspacer{} เอวเมวํ สเทวโก โลโก สมารโก สพฺรหฺมโก สสฺสมพฺราหฺมณี ปชา สเทวมนุสฺสา อนฺโตพุทฺธญาเณ ปริวตฺตติ.\csromanspacer{}\csromanspacer{}ยถา เย เกจิ ปกฺขี อนฺตมโส ครุฬํ เวนเตยฺยํ อุปาทาย อากาสสฺส ปเทเส ปริวตฺตนฺติ.\csromanspacer{} เอวเมวํ เยปิ เต สาริปุตฺตสมา ปญฺญาย, เตปิ พุทฺธญาณสฺส ปเทเส ปริวตฺตนฺติ.\csromanspacer{} พุทฺธญาณํ เทวมนุสฺสานํ ปญฺญํ ผริตฺวา อภิภวิตฺวา ติฏฺฐติ.}

\prose{เยปิ เต ขตฺติยปณฺฑิตา พฺราหฺมณปณฺฑิตา คหปติปณฺฑิตา สมณปณฺฑิตา นิปุณา กตปรปฺปวาทา วาลเวธิรูปา โวภินฺทนฺตา\footnote{เต ภินฺทนฺตา (ก)} มญฺเญ จรนฺติ}

\prose{\csromanfolio{140}ปญฺญาคเตน ทิฏฺฐิคตานิ.\csromanspacer{} เต ปเญฺห อภิสงฺขริตฺวา อภิสงฺขริตฺวา ตถาคเต อุปสงฺกมิตฺวา ปุจฺฉนฺติ คูฬฺหานิ จ ปฏิจฺฉนฺนานิ จ, กถิตา วิสชฺชิตาวเต ปญฺหา ภควตา โหนฺติ นิทฺทิฏฺฐการณา, อุปกฺขิตฺตกา จ เต ภควโต สมฺปชฺชนฺติ.\csromanspacer{} อถ โข ภควาว ตตฺถ อติโรจติ ยทิทํ ปญฺญายาติ โธเนน ยุคํ สมาคมา น หิ ตฺวํ สกฺขติ สมฺปยาตเว.\csromanspacer{} เตนาห ภควา\csromandash{}}

\csromangathameasure{อถ ตฺวํ ปวิตกฺกมาคมา, มนสา ทิฏฺฐิคตานิ จินฺตยนฺโต. \\ โธเนน ยุคํ สมาคมา, น หิ ตฺวํ สกฺขสิ สมฺปยาตเวติ.}
\csromangathagroup{\csromangathastanza{อถ ตฺวํ ปวิตกฺกมาคมา, มนสา ทิฏฺฐิคตานิ จินฺตยนฺโต. \\ โธเนน ยุคํ สมาคมา, น หิ ตฺวํ สกฺขสิ สมฺปยาตเวติ.}}

\csromanheader{2}{ปสูรสุตฺตนิทฺเทโส อฏฺฐโม.}
\csromansectionrule
\csromanheader{2}{9. มาคณฺฑิยสุตฺตนิทฺเทส}
\csromanheader{2}{อถ มาคณฺฑิยสุตฺตนิทฺเทสํ วกฺขติ\csromandash{}}
\proseitem{70}{\csromansymbolfootnote{*}{ขุ 1. 409 ปิฏฺเฐ.}ทิสฺวาน ตณฺหํ อรติํ รคญฺจ,}

\csromanmatikaanchor{mtk.9}
\csromantocmarkref{subsection}{8. ปสูรสุตฺตนิทฺเทส}{mtk.9}
\bookmark[dest={mtk.9},level=2]{8. ปสูรสุตฺตนิทฺเทส}
\prose{\textbf{นาโหสิ ฉนฺโท อปิ เมถุนสฺมิํ.}}

\csromangathameasure{\textbf{กิเมวิทํ มุตฺตกรีสปุณฺณํ,} \\ \textbf{ปาทาปิ นํ สมฺผุสิตุํ น อิจฺเฉ.}}
\csromangathagroup{\csromangathastanza{\textbf{กิเมวิทํ มุตฺตกรีสปุณฺณํ,} \\ \textbf{ปาทาปิ นํ สมฺผุสิตุํ น อิจฺเฉ.}}}

\prose{\textbf{ทิสฺวาน ตณฺหํ อรติํ รคญฺจ, นาโหสิ ฉนฺโท อปิ เมถุนสฺมินฺ}ติ ตณฺหญฺจ อรติญฺจ รคญฺจ มารธีตโร ทิสฺวา ปสฺสิตฺวา เมถุนธมฺเม ฉนฺโท วา ราโค วา เปมํ วา นาโหสีติ ทิสฺวาน ตณฺหํ อรติํ รคญฺจ, นาโหสิ ฉนฺโท อปิ เมถุนสฺมิํ.}

\prose{\textbf{กิเมวิทํ มุตฺตกรีสปุณฺณํ, ปาทาปิ นํ สมฺผุสิตุํ น อิจฺเฉ}ติ กิเมวิทํ สรีรํ มุตฺตปุณฺณํ กรีสปุณฺณํ เสมฺหปุณฺณํ รุหิรปุณฺณํ อฏฺฐิสงฺฆาตนฺหารุสมฺพนฺธํ รุธิรมํสาวเลปนํ จมฺมวินทฺธํ ฉวิยา ปฏิจฺฉนฺนํ ฉิทฺทาวฉิทฺทํ อุคฺฆรนฺตํ ปคฺฆรนฺตํ กิมิสํฆนิเสวิตํ นานากลิมลปริปูรํ ปาเทน อกฺกมิตุํ น อิจฺเฉยฺย, กุโต ปน สํวาโส วา สมาคโม วาติ กิเมวิทํ มุตฺตกรีสปุณฺณํ, ปาทาปิ นํ สมฺผุสิตุํ น อิจฺเฉ.\csromanspacer{} อนจฺฉริยญฺเจตํ มนุสฺโส ทิพฺเพ กาเม ปตฺถยนฺโต มานุสเก กาเม น อิจฺเฉยฺย, มานุสเก วา กาเม}

\vspace{\tipitakaparskip}
\csromancenter{มาคณฺฑิยสุตฺตนิทฺเทส ปตฺถยนฺโต ทิพฺเพ กาเม อิจฺเฉยฺย, ยํ ตฺวํ อุโภปิ น อิจฺฉสิ น สาทิยสิ น ปตฺเถสิ น ปิเหสิ นาภิชปฺปสิ, กิํ เต ทสฺสนํ กตมาย ตฺวํ ทิฏฺฐิยา สมนฺนาคโตติ ปุจฺฉตีติ.\csromanspacer{} เตนาห \textbf{ภควา}\csromandash{}}
\vspace{0.5\baselineskip}
\csromangathameasure{ทิสฺวาน ตณฺหํ อรติํ รคญฺจ, \\ นาโหสิ ฉนฺโท อปิ เมถุนสฺมิํ. \\ กิเมวิทํ มุตฺตกรีสปุณฺณํ, \\ ปาทาปิ นํ สมฺผุสิตุํ น อิจฺเฉติ.}
\csromangathagroup{\csromangathastanza{\csromanfolio{141}ทิสฺวาน ตณฺหํ อรติํ รคญฺจ, \\ นาโหสิ ฉนฺโท อปิ เมถุนสฺมิํ. \\ กิเมวิทํ มุตฺตกรีสปุณฺณํ, \\ ปาทาปิ นํ สมฺผุสิตุํ น อิจฺเฉติ.}}

\proseitem{71}{\csromansymbolfootnote{*}{ขุ 1. 409 ปิฏฺเฐ.}เอตาทิสํ เจ รตนํ น อิจฺฉสิ,}

\prose{\textbf{นาริํ นรินฺเทหิ พหูหิ ปตฺถิตํ.}}

\csromangathameasure{\textbf{ทิฏฺฐิคตํ สีลวตํ นุ ชีวิตํ,} \\ \textbf{ภวูปปตฺติญฺจ วเทสิ กีทิสํ.}}
\csromangathagroup{\csromangathastanza{\textbf{ทิฏฺฐิคตํ สีลวตํ นุ ชีวิตํ,} \\ \textbf{ภวูปปตฺติญฺจ วเทสิ กีทิสํ.}}}

\proseitem{72}{\csromansymbolmark{*}อิทํ วทามีติ น ตสฺส โหติ, (\textbf{มาคณฺฑิยา}ติ\footnote{มาคนฺทิยาติ (สี, สฺยา)} \textbf{ภควา,)}}

\prose{\textbf{ธมฺเมสุ นิจฺเฉยฺย สมุคฺคหีตํ.}}

\csromangathameasure{\textbf{ปสฺสญฺจ ทิฏฺฐีสุ อนุคฺคหาย,} \\ \textbf{อชฺฌตฺตสนฺติํ ปจินํ อทสฺสํ.}}
\csromangathagroup{\csromangathastanza{\textbf{ปสฺสญฺจ ทิฏฺฐีสุ อนุคฺคหาย,} \\ \textbf{อชฺฌตฺตสนฺติํ ปจินํ อทสฺสํ.}}}

\prose{\textbf{อิทํ วทามีติ น ตสฺส โหตี}ติ \textbf{อิทํ วทามี}ติ อิทํ วทามิ, เอตํ วทามิ, เอตฺตกํ วทามิ, เอตฺตาวตา วทามิ, อิทํ ทิฏฺฐิคตํ วทามิ “สสฺสโต โลโก”ติ วา ฯเปฯ “เนว โหติ น น โหติ ตถาคโต ปรํ มรณา”ติ วา.\csromanspacer{}\csromanspacer{}\textbf{น ตสฺส โหตี}ติ น มยฺหํ โหติ, “เอตฺตาวตา วทามี”ติ น ตสฺส โหตีติ \textbf{อิทํ วทามี}ติ น ตสฺส โหติ.}

\prose{\textbf{มาคณฺฑิยา}ติ \textbf{ภควา} ตํ พฺราหฺมณํ นาเมน อาลปติ.\csromanspacer{} \textbf{ภควา}ติ คารวาธิวจนํ ฯเปฯ สจฺฉิกา ปญฺญตฺติ ยทิทํ \textbf{ภควา}ติ \textbf{มาคณฺฑิยา}ติ \textbf{ภควา}.}

\prose{\textbf{ธมฺเมสุ นิจฺเฉยฺย สมุคฺคหีตนฺ}ติ \textbf{ธมฺเมสู}ติ ทฺวาสฏฺฐิยา ทิฏฺฐิคเตสุ.\csromanspacer{} \textbf{นิจฺเฉยฺยา}ติ นิจฺฉินิตฺวา วินิจฺฉินิตฺวา วิจินิตฺวา ปวิจินิตฺวา ตุลยิตฺวา ตีรยิตฺวา วิภาวยิตฺวา วิภูตํ กตฺวา.\csromanspacer{} \textbf{สมุคฺคหีตนฺ}ติ โอธิคฺคาโห พิลคฺคาโห วรคฺคาโห โกฏฺฐาสคฺคาโห อุจฺจยคฺคาโห สมุจฺจยคฺคาโห, “อิทํ สจฺจํ ตจฺฉํ ตถํ ภูตํ ยาถาวํ อวิปรีตนฺ”ติ คหิตํ ปรามฏฺฐํ อภินิวิฏฺฐํ \csromanfolio{142}อชฺโฌสิตํ อธิมุตฺตํ นตฺถิ น สนฺติ น สํวิชฺชติ นุปลพฺภติ, ปหีนํ สมุจฺฉินฺนํ วูปสนฺตํ ปฏิปสฺสทฺธํ อภพฺพุปฺปตฺติกํ ญาณคฺคินา ทฑฺฒนฺติ ธมฺเมสุ นิจฺเฉยฺย สมุคฺคหีตํ.}

\prose{\textbf{ปสฺสญฺจ ทิฏฺฐีสุ อนุคฺคหายา}ติ ทิฏฺฐีสุ อาทีนวํ ปสฺสนฺโต ทิฏฺฐิโย น คณฺหามิ น ปรามสามิ นาภินิวิสามิ.\csromanspacer{} อถ วา น คณฺหิตพฺพา น ปรามสิตพฺพา นาภินิวิสิตพฺพาติ เอวมฺปิ ปสฺสญฺจ ทิฏฺฐีสุ อนุคฺคหาย.}

\prose{อถ วา “สสฺสโต โลโก, อิทเมว สจฺจํ โมฆมญฺญนฺ”ติ ทิฏฺฐิคตเมตํ ทิฏฺฐิคหนํ ทิฏฺฐิกนฺตาโร ทิฏฺฐิวิสูกายิกํ ทิฏฺฐิวิปฺผนฺทิตํ ทิฏฺฐิสญฺโญชนํ สทุกฺขํ สวิฆาตํ ส-อุปายาสํ สปริฬาหํ น นิพฺพิทาย น วิราคาย น นิโรธาย น อุปสมาย น อภิญฺญาย น สมฺโพธาย น นิพฺพานาย สํวตฺตตีติ ทิฏฺฐีสุ อาทีนวํ ปสฺสนฺโต ทิฏฺฐิโย น คณฺหามิ น ปรามสามิ นาภินิวิสามิ.\csromanspacer{} อถ วา น คณฺหิตพฺพา น ปรามสิตพฺพา นาภินิวิตพฺพาติ เอวมฺปิ ปสฺสญฺจ ทิฏฺฐีสุ อนุคฺคหาย.}

\prose{อถ วา “อสสฺสโต โลโก, อนฺตวา โลโก, อนนฺตวา โลโก, ตํ ชีวํ ตํ สรีรํ, อญฺญํ ชีวํ อญฺญํ สรีรํ, โหติ ตถาคโต ปรํ มรณา, น โหติ ตถาคโต ปรํ มรณา, โหติ จ น จ โหติ ตถาคโต ปรํ มรณา, เนว โหติ น น โหติ ตถาคโต ปรํ มรณา, อิทเมว สจฺจํ โมฆมญฺญนฺ”ติ ทิฏฺฐิคตเมตํ ทิฏฺฐิคหนํ ทิฏฺฐิกนฺตาโร ทิฏฺฐิวิสูกายิกํ ทิฏฺฐิวิปฺผนฺทิตํ ทิฏฺฐิสญฺโญชนํ สทุกฺขํ สวิฆาตํ ส-อุปายาสํ สปริฬาหํ น นิพฺพิทาย น วิราคาย น นิโรธาย น อุปสมาย น อภิญฺญาย น สมฺโพธาย น นิพฺพานาย สํวตฺตตีติ ทิฏฺฐีสุ อาทีนวํ ปสฺสนฺโต ทิฏฺฐิโย น คณฺหามิ น ปรามสามิ นาภินิวิสามิ.\csromanspacer{} อถ วา น คณฺหิตพฺพา น ปรามสิตพฺพา นาภินิวิสิตพฺพาติ เอวมฺปิ ปสฺสญฺจ ทิฏฺฐีสุ อนุคฺคหาย.}

\prose{อถ วา อิมา ทิฏฺฐิโย เอวํคหิตา เอวํปรามฏฺฐา เอวํคติกา ภวิสฺสนฺติ เอวํอภิสมฺปรายาติ ทิฏฺฐีสุ อาทีนวํ ปสฺสนฺโต ทิฏฺฐิโย น คณฺหามิ น ปรามสามิ นาภินิวิสามิ.\csromanspacer{} อถ วา น คณฺหิตพฺพา น ปรามสิตพฺพา นาภินิวิสิตพฺพาติ เอวมฺปิ ปสฺสญฺจ ทิฏฺฐีสุ อนุคฺคหาย.}

\prose{อถ วา อิมา ทิฏฺฐิโย นิรยสํวตฺตนิกา ติรจฺฉานโยนิสํวตฺตนิกา เปตฺติวิสยสํวตฺตนิกาติ ทิฏฺฐีสุ อาทีนวํ ปสฺสนฺโต ทิฏฺฐิโย น}

\vspace{\tipitakaparskip}
\csromancenter{มาคณฺฑิยสุตฺตนิทฺเทส คณฺหามิ น ปรามสามิ นาภินิวิสามิ.\csromanspacer{} อถ วา น คณฺหิตพฺพา น ปรามสิตพฺพา นาภินิวิสิตพฺพาติ เอวมฺปิ ปสฺสญฺจ ทิฏฺฐีสุ อนุคฺคหาย.}
\prose{\csromanfolio{143}อถ วา อิมา ทิฏฺฐิโย อนิจฺจา สงฺขตา ปฏิจฺจสมุปฺปนฺนา ขยธมฺมา วยธมฺมา วิราคธมฺมา นิโรธธมฺมาติ ทิฏฺฐีสุ อาทีนวํ ปสฺสนฺโต ทิฏฺฐิโย น คณฺหามิ น ปรามสามิ นาภินิวิสามิ.\csromanspacer{} อถ วา น คณฺหิตพฺพา น ปรามสิตพฺพา นาภินิวิสิตพฺพาติ เอวมฺปิ ปสฺสญฺจ ทิฏฺฐีสุ อนุคฺคหาย.}

\prose{\textbf{อชฺฌตฺตสนฺติํ ปจินํ อทสฺสนฺ}ติ อชฺฌตฺตสนฺติํ อชฺฌตฺตํ ราคสฺส สนฺติํ, โทสสฺส สนฺติํ, โมหสฺส สนฺติํ.\csromanspacer{} โกธสฺส.\csromanspacer{} อุปนาหสฺส.\csromanspacer{} มกฺขสฺส.\csromanspacer{} ปฬาสสฺส.\csromanspacer{} อิสฺสาย.\csromanspacer{} มจฺฉริยสฺส.\csromanspacer{} มายาย.\csromanspacer{} สาเฐยฺยสฺส.\csromanspacer{} ถมฺภสฺส.\csromanspacer{} สารมฺภสฺส.\csromanspacer{} มานสฺส.\csromanspacer{} อติมานสฺส.\csromanspacer{} มทสฺส.\csromanspacer{} ปมาทสฺส.\csromanspacer{} สพฺพกิเลสานํ.\csromanspacer{} สพฺพทุจฺจริตานํ.\csromanspacer{} สพฺพทรถานํ.\csromanspacer{} สพฺพปริฬาหานํ.\csromanspacer{} สพฺพสนฺตาปานํ.\csromanspacer{} สพฺพากุสลาภิสงฺขารานํ สนฺติํ อุปสนฺติํ วูปสนฺติํ นิพฺพุติํ ปฏิปสฺสทฺธิํ สนฺติํ.\csromanspacer{} \textbf{ปจินนฺ}ติ ปจินนฺโต วิจินนฺโต ปวิจินนฺโต ตุลยนฺโต ตีรยนฺโต วิภาวยนฺโต วิภูตํ กโรนฺโต”สพฺเพ สงฺขารา อนิจฺจา”ติ ปจินนฺโต วิจินนฺโต ปวิจินนฺโต ตุลยนฺโต ตีรยนฺโต วิภาวยนฺโต วิภูตํ กโรนฺโต “สพฺเพ สงฺขารา ทุกฺขา”ติ. “สพฺเพ ธมฺมา อนตฺตา”ติ ปจินนฺโต วิจินนฺโต ปวิจินนฺโต, “ยํ กิญฺจิ สมุทยธมฺมํ, สพฺพํ ตํ นิโรธธมฺมนฺ”ติ ปจินนฺโต วิจินนฺโต ปวิจินนฺโต ตุลยนฺโต ตีรยนฺโต วิภาวยนฺโต วิภูตํ กโรนฺโต.\csromanspacer{} \textbf{อทสฺสนฺ}ติ อทสฺสํ อทกฺขิํ อปสฺสิํ ปฏิวิชฺฌินฺติ อชฺฌตฺตสนฺติํ ปจินํ อทสฺสํ.\csromanspacer{} เตนาห ภควา\csromandash{}}

\prose{อิทํ วทามีติ น ตสฺส โหติ, (มาคณฺฑิยาติ ภควา,)}

\prose{ธมฺเมสุ นิจฺเฉยฺย สมุคฺคหีตํ.}

\prose{ปสฺสญฺจ ทิฏฺฐีสุ อนุคฺคหาย,}

\prose{\textbf{อชฺฌตฺตสนฺติํ ปจินํ อทสฺสนฺ}ติ.}

\proseitem{73}{\csromansymbolfootnote{*}{ขุ 1. 410 ปิฏฺเฐ.}วินิจฺฉยา ยานิ ปกปฺปิตานิ, (อิติ มาคณฺฑิโย,)}

\prose{\textbf{เต เว มุนี พฺรูสิ อนุคฺคหาย.}}

\csromangathameasure{\textbf{อชฺฌตฺตสนฺตีติ ยเมตมตฺถํ,} \\ \textbf{กถํ นุ ธีเรหิ ปเวทิตํ ตํ.}}
\csromangathagroup{\csromangathastanza{\textbf{อชฺฌตฺตสนฺตีติ ยเมตมตฺถํ,} \\ \textbf{กถํ นุ ธีเรหิ ปเวทิตํ ตํ.}}}

\prose{\textbf{วินิจฺฉยา ยานิ ปกปฺปิตานี}ติ วินิจฺฉยา วุจฺจนฺติ ทฺวาสฏฺฐิ ทิฏฺฐิคตานิ ทิฏฺฐิวินิจฺฉยา.\csromanspacer{}\csromanspacer{}\textbf{ปกปฺปิตานี}ติ กปฺปิตา ปกปฺปิตา อภิสงฺขตา สณฺฐปิตาติปิ \csromanfolio{144}ปกปฺปิตา.\csromanspacer{} อถ วา อนิจฺจา สงฺขตา ปฏิจฺจสมุปฺปนฺนา ขยธมฺมา วยธมฺมา วิราคธมฺมา นิโรธธมฺมา วิปริณามธมฺมาติปิ ปกปฺปิตาติ วินิจฺฉยา ยานิ ปกปฺปิตานิ.}

\prose{\textbf{อิติ มาคณฺฑิโย}ติ \textbf{อิตี}ติ ปทสนฺธิ ปทสํสคฺโค ปทปาริปูรี อกฺขรสมวาโย พฺยญฺชนสิลิฏฺฐตา ปทานุปุพฺพตาเปตํ \textbf{อิตี}ติ.\csromanspacer{}\csromanspacer{}\textbf{มาคณฺฑิโย}ติ ตสฺส พฺราหฺมณสฺส นามํ สงฺขา สมญฺญา ปญฺญตฺติ โวหาโร อิติ \textbf{มาคณฺฑิโย}ติ.}

\prose{\textbf{เต เว มุนี พฺรูสิ อนุคฺคหาย, อชฺฌตฺตสนฺตีติ ยเมตมตฺถนฺติ เตเว}ติ ทฺวาสฏฺฐิ ทิฏฺฐิคตานิ.\csromanspacer{} \textbf{มุนี}ติ \textbf{โมนํ} วุจฺจติ ญาณํ ฯเปฯ สงฺคชาลมติจฺจ โส มุนีติ.\csromanspacer{} \textbf{อนุคฺคหายา}ติ ทิฏฺฐีสุ อาทีนวํ ปสฺสนฺโต ทิฏฺฐิโย น คณฺหามิ น ปรามสามิ นาภินิวิสามีติ จ ภณสิ, อชฺฌตฺตสนฺตีติ จ ภณสิ.\csromanspacer{}\csromanspacer{}\textbf{ยเมตมตฺถนฺ}ติ ยํ ปรมตฺถนฺติ เต เว มุนี พฺรูสิ อนุคฺคหาย อชฺฌตฺตสนฺตีติ ยเมตมตฺถํ.}

\prose{\textbf{กถํ นุ ธีเรหิ ปเวทิตํ ตนฺติกถํ นู}ติ ปทํ สํสยปุจฺฉา วิมติปุจฺฉา เทฺวฬฺหกปุจฺฉา อเนกํสปุจฺฉา, เอวํ นุ โข นนุ โข กิํ นุ โข กถํ นุ โขติ กถํ นุ.\csromanspacer{}\csromanspacer{}\textbf{ธีเรหี}ติ ธีเรหิ ปณฺฑิเตหิ ปญฺญวนฺเตหิ\footnote{ปญฺญาวนฺเตหิ (สี, สฺยา)} พุทฺธิมนฺเตหิ ญาณีหิ วิภาวีหิ เมธาวีหิ.\csromanspacer{} \textbf{ปเวทิตนฺ}ติ เวทิตํ ปเวทิตํ อาจิกฺขิตํ เทสิตํ ปญฺญาปิตํ ปฏฺฐปิตํ วิวฏํ วิภตฺตํ อุตฺตานีกตํ ปกาสิตนฺติ กถํ นุ ธีเรหิ ปเวทิตํ ตํ.\csromanspacer{} เตนาห โส พฺราหฺมโณ\csromandash{}}

\prose{วินิจฺฉยา ยานิ ปกปฺปิตานิ, (อิติ \textbf{มาคณฺฑิโย},)}

\prose{เต เว มุนี พฺรูสิ อนุคฺคหาย.}

\prose{อชฺฌตฺตสนฺตีติ ยเมตมตฺถํ,}

\prose{กถํ นุ ธีเรหิ ปเวทิตํ ตนฺติ.}

\proseitem{74}{\csromansymbolfootnote{*}{ขุ 1. 410 ปิฏฺเฐ.}น ทิฏฺฐิยา น สุติยา น ญาเณน, (มาคณฺฑิยาติ ภควา,)}

\prose{\textbf{สีลพฺพเตนาปิ น สุทฺธิมาห.}}

\csromangathameasure{\textbf{อทิฏฺฐิยา อสฺสุติยา อญาณา,} \\ \textbf{อสีลตา อพฺพตา โนปิ เตน.} \\ \textbf{เอเต จ นิสฺสชฺช อนุคฺคหาย,} \\ \textbf{สนฺโต อนิสฺสย ภวํ น ชปฺเป.}}
\csromangathagroup{\csromangathastanza{\textbf{อทิฏฺฐิยา อสฺสุติยา อญาณา,} \\ \textbf{อสีลตา อพฺพตา โนปิ เตน.} \\ \textbf{เอเต จ นิสฺสชฺช อนุคฺคหาย,} \\ \textbf{สนฺโต อนิสฺสย ภวํ น ชปฺเป.}}}

\csromanheader{2}{9. มาคณฺฑิยสุตฺตนิทฺเทส}
\prose{\csromanfolio{145}\textbf{น ทิฏฺฐิยา น สุติยา น ญาเณนา}ติ ทิฏฺเฐนาปิ สุทฺธิํ วิสุทฺธิํ ปริสุทฺธิํ, มุตฺติํ วิมุตฺติํ ปริมุตฺติํ นาห น กเถสิ น ภณสิ น ทีปยสิ น โวหรสิ, สุเตนาปิ สุทฺธิํ วิสุทฺธิํ ปริสุทฺธิํ, มุตฺติํ วิมุตฺติํ ปริมุตฺติํ นาห น กเถสิ น ภณสิ น ทีปยสิ น โวหรสิ, ทิฏฺฐสุเตนาปิ สุทฺธิํ วิสุทฺธิํ ปริสุทฺธิํ, มุตฺติํ วิมุตฺติํ ปริมุตฺติํ นาห น กเถสิ น ภณสิ น ทีปยสิ น โวหรสิ, ญาเณนาปิ สุทฺธิํ วิสุทฺธิํ ปริสุทฺธิํ, มุตฺติํ วิมุตฺติํ ปริมุตฺติํ นาห นกเถสิ น ภณสิ น ทีปยสิ น โวหรสีติ น ทิฏฺฐิยา น สุติยา น ญาเณน.}

\prose{\textbf{มาคณฺฑิยา}ติ \textbf{ภควา} ตํ พฺราหฺมณํ นาเมน อาลปติ.\csromanspacer{} \textbf{ภควา}ติ คารวาธิวจนํ ฯเปฯ สจฺฉิกา ปญฺญตฺติ ยทิทํ \textbf{ภควา}ติ \textbf{มาคณฺฑิยา}ติ \textbf{ภควา}.}

\prose{\textbf{สีลพฺพเตนาปิ น สุทฺธิมาหา}ติ สีเลนาปิ สุทฺธิํ วิสุทฺธิํ ปริสุทฺธิํ, มุตฺติํ วิมุตฺติํ ปริมุตฺติํ นาห น กเถสิ น ภณสิ น ทีปยสิ น โวหรสิ, วเตนาปิ สุทฺธิํ วิสุทฺธิํ ปริสุทฺธิํ, มุตฺติํ วิมุตฺติํ ปริมุตฺติํ นาห น กเถสิ น ภณสิ น ทีปยสิ น โวหรสิ, สีลพฺพเตนาปิ สุทฺธิํ วิสุทฺธิํ ปริสุทฺธิํ, มุตฺติํ วิมุตฺติํ ปริมุตฺติํ นาห น กเถสิ น ภณสิ น ทีปยสิ น โวหรสีติ สีลพฺพเตนาปิ น สุทฺธิมาห.}

\prose{\textbf{อทิฏฺฐิยา อสฺสุติยา อญาณา, อสีลตา อพฺพตา โนปิ เตนา}ติ ทิฏฺฐิปิ อิจฺฉิตพฺพา ทสวตฺถุกา สมฺมาทิฏฺฐิ “อตฺถิ ทินฺนํ, อตฺถิ ยิฏฺฐํ, อตฺถิ หุตํ, อตฺถิ สุกตทุกฺกฏานํ\footnote{สุกฏทุกฺกฏานํ (สี) อภิ 1. 241 ปิฏฺเฐปิ.} กมฺมานํ ผลํ วิปาโก, อตฺถิ อยํ โลโก, อตฺถิ ปโร โลโก, อตฺถิ มาตา, อตฺถิ ปิตา, อตฺถิ สตฺตา โอปปาติกา, อตฺถิ โลเก สมณพฺราหฺมณา สมฺมคฺคตา\footnote{สมคฺคตา (ก)} สมฺมาปฏิปนฺนา, เย อิมญฺจ โลกํ ปรญฺจ โลกํ สยํ อภิญฺญา สจฺฉิกตฺวา ปเวเทนฺตี”ติ, สวนมฺปิ อิจฺฉิตพฺพํ ปรโต โฆโส, สุตฺตํ เคยฺยํ เวยฺยากรณํ คาถา อุทานํ อิติวุตฺตกํ ชาตกํ อพฺภุตธมฺมํ เวทลฺลํ, ญาณมฺปิ อิจฺฉิตพฺพํ กมฺมสฺสกตญาณํ สจฺจานุโลมิกญาณํ\footnote{กมฺมสฺสกตํ ญาณํ สจฺจานุโลมิกํ ญาณํ (สี, ก) ญาณวิภงฺเคปิ.} อภิญฺญาญาณํ สมาปตฺติญาณํ, สีลมฺปิ อิจฺฉิตพฺพํ ปาติโมกฺขสํวโร, วตมฺปิ อิจฺฉิตพฺพํ อฏฺฐ ธุตงฺคานิ อารญฺญิกงฺคํ ปิณฺฑปาติกงฺคํ ปํสุกูลิกงฺคํ เตจีวริกงฺคํ สปทานจาริกงฺคํ ขลุปจฺฉาภตฺติกงฺคํ เนสชฺชิกงฺคํ ยถาสนฺถติกงฺคนฺติ.}

\prose{\csromanfolio{146}\textbf{อทิฏฺฐิยา อสฺสุติยา อญาณา, อสีลตา อพฺพตา โนปิ เตนา}ติ นาปิ สมฺมาทิฏฺฐิมตฺเตน, นาปิ สวนมตฺเตน, นาปิ ญาณมตฺเตน, นาปิ สีลมตฺเตน, นาปิ วตมตฺเตน อชฺฌตฺตสนฺติํ ปตฺโต โหติ, นาปิ วินา \textbf{เอเต}หิ ธมฺเมหิ อชฺฌตฺตสนฺติํ ปาปุณาติ.\csromanspacer{} อปิ จ สมฺภารา อิเม ธมฺมา โหนฺติ อชฺฌตฺตสนฺติํ ปาปุณิตุํ อธิคนฺตุํ ผสฺสิตุํ สจฺฉิกาตุนฺติ อทิฏฺฐิยา อสฺสุติยา อญาณา, อสีลตา อพฺพตา โนปิ เตน.}

\prose{\textbf{เอเต จ นิสฺสชฺช อนุคฺคหายา}ติ \textbf{เอเต}ติ กณฺหปกฺขิกานํ ธมฺมานํ สมุคฺฆาตโต ปหานํ อิจฺฉิตพฺพํ, เตธาตุเกสุ กุสเลสุ ธมฺเมสุ อตมฺมยตา\footnote{อกมฺมยถา (สี, ก)} อิจฺฉิตพฺพา, ยโต กณฺหปกฺขิยา ธมฺมา สมุคฺฆาตปหาเนน ปหีนา โหนฺติ อุจฺฉินฺนมูลา ตาลาวตฺถุกตา อนภาวํกตา อายติํ อนุปฺปาทธมฺมา, เตธาตุเกสุ จ กุสเลสุ ธมฺเมสุ อตมฺมยตา โหติ, เอตฺตาวตาปิ น คณฺหาติ น ปรามสติ นาภินิวิสติ.\csromanspacer{} อถ วา น คณฺหิตพฺพา น ปรามสิตพฺพา นาภินิวิสิตพฺพาติ เอวมฺปิ \textbf{เอเต} จ นิสฺสชฺช อนุคฺคหาย.\csromanspacer{} ยโต ตณฺหา จ ทิฏฺฐิ จ มาโน จ ปหีนา โหนฺติ อุจฺฉินฺนมูลา ตาลาวตฺถุกตา อนภาวํกตา อายติํ อนุปฺปาทธมฺมาติ, เอตฺตาวตาปิ น คณฺหาติ น ปรามสติ นาภินิวิสตีติ เอวมฺปิ \textbf{เอเต} จ นิสฺสชฺช อนุคฺคหาย.}

\prose{ยโต ปุญฺญาภิสงฺขาโรจ อปุญฺญาภิสงฺขาโร จ อาเนญฺชาภิสงฺขาโร จ ปหีนา โหนฺติ อุจฺฉินฺนมูลา ตาลาวตฺถุกตา อนภาวํกตา อายติํ อนุปฺปาทธมฺมาติ, เอตฺตาวตาปิ น คณฺหาติ น ปรามสติ นาภินิวิสตีติ เอวมฺปิ \textbf{เอเต} จ นิสฺสชฺช อนุคฺคหาย.}

\prose{\textbf{สนฺโต อนิสฺสาย ภวํ น ชปฺเป}ติ \textbf{สนฺโต}ติ ราคสฺส สมิตตฺตา \textbf{สนฺโต}, โทสสฺส สมิตตฺตา \textbf{สนฺโต}, โมหสฺส สมิตตฺตา \textbf{สนฺโต}, โกธสฺส.\csromanspacer{} อุปนาหสฺส.\csromanspacer{} มกฺขสฺส.\csromanspacer{} ปฬาสสฺส.\csromanspacer{} อิสฺสาย.\csromanspacer{} มจฺฉริยสฺส.\csromanspacer{} มายาย.\csromanspacer{} สาเฐยฺยสฺส.\csromanspacer{} ถมฺภสฺส.\csromanspacer{} สารมฺภสฺส.\csromanspacer{} มานสฺส.\csromanspacer{} อติมานสฺส.\csromanspacer{} มทสฺส.\csromanspacer{} ปมาทสฺส.\csromanspacer{} สพฺพกิเลสานํ.\csromanspacer{} สพฺพทุจฺจริตานํ.\csromanspacer{} สพฺพทรถานํ.\csromanspacer{} สพฺพปริฬาหานํ.\csromanspacer{} สพฺพสนฺตาปานํ.\csromanspacer{} สพฺพากุสลาภิสงฺขารานํ สนฺตตฺตา สมิตตฺตา วูปสมิตตฺตา วิชฺฌาตตฺตา นิพฺพุตตฺตา วิคตตฺตา ปฏิปสฺสทฺธตฺตา \textbf{สนฺโต} อุป\textbf{สนฺโต} วูป\textbf{สนฺโต} นิพฺพุโต ปฏิปสฺสทฺโธติ \textbf{สนฺโต}.}

\csromanheader{2}{9. มาคณฺฑิยสุตฺตนิทฺเทส}
\prose{\csromanfolio{147}\textbf{อนิสฺสายา}ติ เทฺว นิสฺสยา ตณฺหานิสฺสโย จ ทิฏฺฐินิสฺสโย จ ฯเปฯ อยํ ตณฺหานิสฺสโย ฯเปฯ อยํ ทิฏฺฐินิสฺสโย.\csromanspacer{} ตณฺหานิสฺสยํ ปหาย ทิฏฺฐินิสฺสยํ ปฏินิสฺสชฺชิตฺวา จกฺขุํ อนิสฺสาย, โสตํ อนิสฺสาย, ฆานํ อนิสฺสาย, ชิวฺหํ อนิสฺสาย, กายํ อนิสฺสาย, มนํ อนิสฺสาย.\csromanspacer{} รูเป.\csromanspacer{} สทฺเท.\csromanspacer{} คนฺเธ.\csromanspacer{} รเส.\csromanspacer{} โผฏฺฐพฺเพ.\csromanspacer{} กุลํ.\csromanspacer{} คณํ.\csromanspacer{} อาวาสํ.\csromanspacer{} ลาภํ.\csromanspacer{} ยสํ.\csromanspacer{} ปสํสํ.\csromanspacer{} สุขํ.\csromanspacer{} จีวรํ.\csromanspacer{} ปิณฺฑปาตํ.\csromanspacer{} เสนาสนํ.\csromanspacer{} คิลานปจฺจยเภสชฺชปริกฺขารํ.\csromanspacer{} กามธาตุํ.\csromanspacer{} รูปธาตุํ.\csromanspacer{} อรูปธาตุํ.\csromanspacer{} กามภวํ.\csromanspacer{} รูปภวํ.\csromanspacer{} อรูปภวํ.\csromanspacer{} สญฺญาภวํ.\csromanspacer{} อสญฺญาภวํ.\csromanspacer{} เนวสญฺญานาสญฺญาภวํ.\csromanspacer{} เอกโวการภวํ.\csromanspacer{} จตุโวการภวํ.\csromanspacer{} ปญฺจโวการภวํ.\csromanspacer{} อตีตํ.\csromanspacer{} อนาคตํ.\csromanspacer{} ปจฺจุปฺปนฺนํ.\csromanspacer{} ทิฏฺฐสุตมุตวิญฺญาตพฺเพ ธมฺเม อนิสฺสาย อคฺคณฺหิตฺวา อปรามสิตฺวา อนภินิวิสิตฺวาติ สนฺโต อนิสฺสาย.\csromanspacer{}\csromanspacer{}\textbf{ภวํ น ชปฺเป}ติ กามภวํ น ชปฺเปยฺย, รูปภวํ น ชปฺเปยฺย, อรูปภวํ น ชปฺเปยฺย นปฺปชปฺเปยฺย น อภิชปฺเปยฺยาติ สนฺโต อนิสฺสาย ภวํ น ชปฺเป.\csromanspacer{} เตนาห ภควา\csromandash{}}

\prose{น ทิฏฺฐิยา น สุติยา น ญาเณน, (มาคณฺฑิยาติ ภควา,)}

\prose{\textbf{สีลพฺพเตนาปิ น สุทฺธิมาห.}}

\csromangathameasure{\textbf{อทิฏฺฐิยา อสฺสุติยา อญาณา,} \\ \textbf{อสีลตา อพฺพตา โนปิ เตน.} \\ เอเต จ นิสฺสชฺช อนุคฺคหาย, \\ สนฺโต อนิสฺสาย ภวํ น ชปฺเปติ.}
\csromangathagroup{\csromangathastanza{\textbf{อทิฏฺฐิยา อสฺสุติยา อญาณา,} \\ \textbf{อสีลตา อพฺพตา โนปิ เตน.} \\ เอเต จ นิสฺสชฺช อนุคฺคหาย, \\ สนฺโต อนิสฺสาย ภวํ น ชปฺเปติ.}}

\proseitem{75}{\csromansymbolfootnote{*}{ขุ 1. 410 ปิฏฺเฐ.}โน เจ กิร ทิฏฺฐิยา น สุติยา น ญาเณน, (อิติ มาคณฺฑิโย,)}

\prose{\textbf{สีลพฺพเตนาปิ น สุทฺธิมาห.}}

\csromangathameasure{\textbf{อทิฏฺฐิยา อสฺสุติยา อญาณา,} \\ \textbf{อสีลตา อพฺพตา โนปิ เตน.} \\ \textbf{มญฺญามหํ โมมุหเมว ธมฺมํ,} \\ \textbf{ทิฏฺฐิยา เอเก ปจฺเจนฺติ สุทฺธิํ.}}
\csromangathagroup{\csromangathastanza{\textbf{อทิฏฺฐิยา อสฺสุติยา อญาณา,} \\ \textbf{อสีลตา อพฺพตา โนปิ เตน.} \\ \textbf{มญฺญามหํ โมมุหเมว ธมฺมํ,} \\ \textbf{ทิฏฺฐิยา เอเก ปจฺเจนฺติ สุทฺธิํ.}}}

\prose{\textbf{โน เจ กิร ทิฏฺฐิยา น สุติยา น ญาเณนา}ติ ทิฏฺฐิยาปิ สุทฺธิํ วิสุทฺธิํ ปริสุทฺธิํ, มุตฺติํ วิมุตฺติํ ปริมุตฺติํ นาห น กเถสิ น ภณสิ น ทีปยสิ น โวหรสิ, สุเตนาปิ สุทฺธิํ วิสุทฺธิํ ปริสุทฺธิํ, มุตฺติํ วิมุตฺติํ ปริมุตฺติํ.\csromanspacer{} ทิฏฺฐสุเตนาปิ สุทฺธิํ วิสุทฺธิํ ปริสุทฺธิํ, มุตฺติํ วิมุตฺติํ ปริมุตฺติํ.\csromanspacer{} ญาเณนาปิ สุทฺธิํ วิสุทฺธิํ ปริสุทฺธิํ, มุตฺติํ วิมุตฺติํ ปริมุตฺติํ นาห น กเถสิ \csromanfolio{148}น ภณสิ น ทีปยสิ น โวหรสีติ โน เจ กิร ทิฏฺฐิยา น สุติยา น ญาเณน\textbf{.}}

\prose{\textbf{อิติ มาคณฺฑิโย}ติ \textbf{อิตี}ติ ปทสนฺธิ ฯเปฯ.\csromanspacer{} \textbf{มาคณฺฑิโย}ติ ตสฺส พฺราหฺมณสฺส นามํ ฯเปฯ อิติ \textbf{มาคณฺฑิโย}.}

\prose{\textbf{สีลพฺพเตนาปิ น สุทฺธิมาหา}ติ สีเลนาปิ สุทฺธิํ วิสุทฺธิํ ปริสุทฺธิํ ฯเปฯ วเตนาปิ สุทฺธิํ วิสุทฺธิํ ปริสุทฺธิํ ฯเปฯ สีลพฺพเตนาปิ สุทฺธิํ วิสุทฺธิํ ปริสุทฺธิํ, มุตฺติํ วิมุตฺติํ ปริมุตฺติํ นาห น กเถสิ น ภณสิ น ทีปยสิ นโวหรสีติ สีลพฺพเตนาปิ น สุทฺธิมาห.}

\prose{\textbf{อทิฏฺฐิยา อสฺสุติยา อญาณา, อสีลตา อพฺพตา โนปิ เตนา}ติ ทิฏฺฐิปิ อิจฺฉิตพฺพาติ เอวํ ภณสิ, สวนมฺปิ อิจฺฉิตพฺพนฺติ เอวํ ภณสิ, ญาณมฺปิ อิจฺฉิตพฺพนฺติ เอวํ ภณสิ, น สกฺโกสิ เอกํเสน อนุชานิตุํ, นปิ สกฺโกสิ เอกํเสน ปฏิกฺขิปิตุนฺติ อทิฏฺฐิยา อสฺสุติยา อญาณา, อสีลตา อพฺพตา โนปิ เตน.}

\prose{\textbf{มญฺญามหํ โมมุหเมว ธมฺมนฺ}ติ โมมูหธมฺโม อยํ ตุยฺหํ พาลธมฺโม มูฬฺหธมฺโม อญฺญาณธมฺโม อมราวิกฺเขปธมฺโมติ เอวํ มญฺญามิ เอวํ ชานามิ เอวํ อาชานามิ เอวํ วิชานามิ เอวํ ปฏิวิชานามิ เอวํ ปฏิวิชฺฌามีติ มญฺญามหํ โมมุหเมว ธมฺมํ.}

\prose{\textbf{ทิฏฺฐิยา เอเก ปจฺเจนฺติ สุทฺธินฺ}ติ สุทฺธิทิฏฺฐิยา เอเก สมณพฺราหฺมณา สุทฺธิํ วิสุทฺธิํ ปริสุทฺธิํ, มุตฺติํ วิมุตฺติํ ปริมุตฺติํ ปจฺเจนฺติ, “สสฺสโต โลโก, อิทเมว สจฺจํ โมฆมญฺญนฺ”ติ ทิฏฺฐิยา เอเก สมณพฺราหฺมณา สุทฺธิํ วิสุทฺธิํ ปริสุทฺธิํ, มุตฺติํ วิมุตฺติํ ปริมุตฺติํ ปจฺเจนฺติ, “อสสฺสโต โลโก ฯเปฯ เนว โหติ น น โหติ ตถาคโต ปรํ มรณา, อิทเมว สจฺจํ โมฆมญฺญนฺ”ติ ทิฏฺฐิยา เอเก สมณพฺราหฺมณา สุทฺธิํ วิสุทฺธิํ ปริสุทฺธิํ, มุตฺติํ วิมุตฺติํ ปริมุตฺติํ ปจฺเจนฺตีติ ทิฏฺฐิยา เอเก ปจฺเจนฺติ สุทฺธิํ.\csromanspacer{} เตนาห โส พฺราหฺมโณ\csromandash{}}

\prose{โน เจ กิร ทิฏฺฐิยา น สุติยา น ญาเณน, (อิติ \textbf{มาคณฺฑิโย},)}

\prose{สีลพฺพเตนาปิ น สุทฺธิมาห.}

\prose{อทิฏฺฐิยา อสฺสุติยา อญาณา,}

\prose{อสีลตา อพฺพตา โนปิ เตน.}

\prose{มญฺญามหํ โมมุหเมว ธมฺมํ,}

\prose{\textbf{ทิฏฺฐิยา เอเก ปจฺเจนฺติ สุทฺธินฺ}ติ.}

\csromanheader{2}{9. มาคณฺฑิยสุตฺตนิทฺเทส}
\proseitem{76}{\csromanfolio{149}\csromansymbolmark{*}\textbf{ทิฏฺฐิญฺจ}\footnote{ทิฏฺฐีสุ (สี, สฺยา, ก) ขุ 1. 410 ปิฏฺเฐปิ ปสฺสิตพฺพํ.} \textbf{นิสฺสายนุปุจฺฉมาโน, (มาคณฺฑิยาติ ภควา,)}}

\prose{\textbf{สมุคฺคหีเตสุ ปโมหมาคา}\footnote{สโมหมาคา (ก)}\textbf{.}}

\csromangathameasure{\textbf{อิโต จ นาทฺทกฺขิ อณุมฺปิ สญฺญํ,} \\ \textbf{ตสฺมา ตุวํ โมมุหโต ทหาสิ.}}
\csromangathagroup{\csromangathastanza{\textbf{อิโต จ นาทฺทกฺขิ อณุมฺปิ สญฺญํ,} \\ \textbf{ตสฺมา ตุวํ โมมุหโต ทหาสิ.}}}

\prose{\textbf{ทิฏฺฐิญฺจ นิสฺสายนุปุจฺฉมาโน}ติ มาคณฺฑิโย พฺราหฺมโณ ทิฏฺฐิํ นิสฺสาย ทิฏฺฐิํ ปุจฺฉติ, ลคฺคนํ นิสฺสาย ลคฺคนํ ปุจฺฉติ, พนฺธนํ นิสฺสาย พนฺธนํ ปุจฺฉติ, ปลิโพธํ นิสฺสาย ปลิโพธํ ปุจฺฉติ\textbf{.}\csromanspacer{}\csromanspacer{}\textbf{อนุปุจฺฉมาโน}ติ ปุนปฺปุนํ ปุจฺฉตีติ ทิฏฺฐิญฺจ นิสฺสายนุปุจฺฉมาโน\textbf{.}}

\prose{\textbf{มาคณฺฑิยา}ติ \textbf{ภควา} ตํ พฺราหฺมณํ นาเมน อาลปติ\textbf{.}\csromanspacer{} \textbf{ภควา}ติ คารวาธิวจนํ ฯเปฯ สจฺฉิกา ปญฺญตฺติ ยทิทํ \textbf{ภควา}ติ \textbf{มาคณฺฑิยา}ติ \textbf{ภควา.}}

\prose{\textbf{สมุคฺคหีเตสุ ปโมหมาคา}ติ ยา สา ทิฏฺฐิ ตยา คหิตา ปรามฏฺฐา อภินิวิฏฺฐา อชฺโฌสิตา อธิมุตฺตา, ตาเยว ตฺวํ ทิฏฺฐิยา มูโฬฺหสิ ปมูโฬฺหสิ สมฺมูโฬฺหสิ โมหํ อาคโตสิ ปโมหํ อาคโตสิ สมฺโมหํ อาคโตสิ อนฺธการํ ปกฺขนฺโทสีติ สมุคฺคหีเตสุ ปโมหมาคา\textbf{.}}

\prose{\textbf{อิโต จ นาทฺทกฺขิ อณุมฺปิ สญฺญนฺ}ติ อิโต อชฺฌตฺตสนฺติโต วา ปฏิปทาโต วา ธมฺมเทสนาโต วา ยุตฺตสญฺญํ ปตฺตสญฺญํ ลกฺขณสญฺญํ การณสญฺญํ ฐานสญฺญํ น ปฏิลภติ, กุโต ญาณนฺติ, เอวมฺปิ อิโต จ นาทฺทกฺขิ อณุมฺปิ สญฺญํ\textbf{.}\csromanspacer{}\csromanspacer{}อถ วา อนิจฺจํ วา อนิจฺจสญฺญานุโลมํ วา ทุกฺขํ วา ทุกฺขสญฺญานุโลมํ วา อนตฺตํ วา อนตฺตสญฺญานุโลมํ วา สญฺญุปฺปาทมตฺตํ วา สญฺชานิตมตฺตํ วา น ปฏิลภติ กุโต ญาณนฺติ, เอวมฺปิ อิโต จ นาทฺทกฺขิ อณุมฺปิ สญฺญํ\textbf{.}}

\prose{\textbf{ตสฺมา ตุวํ โมมุหโต ทหาสี}ติ \textbf{ตสฺมา}ติ \textbf{ตสฺมา} ตํการณา ตํเหตุ ตปฺปจฺจยา ตํนิทานา โมมูหธมฺมโต พาลธมฺมโต มูฬฺหธมฺมโต อญฺญาณธมฺมโต อมราวิกฺเขปธมฺมโต ทหาสิ ปสฺสสิ ทกฺขสิ โอโลเกสิ นิชฺฌายสิ อุปปริกฺขสีติ \textbf{ตสฺมา} ตุวํ โมมุหโต ทหาสิ\textbf{.}\csromanspacer{} เตนาห \textbf{ภควา}\csromandash{}}

\prose{\csromanfolio{150}ทิฏฺฐิญฺจ นิสฺสายนุปุจฺฉมาโน, (มาคณฺฑิยาติ ภควา,)}

\prose{สมุคฺคหีเตสุ ปโมหมาคา.}

\csromangathameasure{อิโต จ นาทฺทกฺขิ อณุมฺปิ สญฺญํ, \\ ตสฺมา ตุวํ โมมุหโต ทหาสีติ.}
\csromangathagroup{\csromangathastanza{อิโต จ นาทฺทกฺขิ อณุมฺปิ สญฺญํ, \\ ตสฺมา ตุวํ โมมุหโต ทหาสีติ.}}

\proseitem{77}{\csromansymbolfootnote{*}{ขุ 1. 410 ปิฏฺเฐ.}สโม วิเสสี อุท วา นิหีโน,}

\prose{\textbf{โย มญฺญติ โส วิวเทถ เตน.}}

\csromangathameasure{\textbf{ตีสุ วิธาสุ อวิกมฺปมาโน,} \\ \textbf{สโม วิเสสีติ น ตสฺส โหติ.}}
\csromangathagroup{\csromangathastanza{\textbf{ตีสุ วิธาสุ อวิกมฺปมาโน,} \\ \textbf{สโม วิเสสีติ น ตสฺส โหติ.}}}

\prose{\textbf{สโม วิเสสี อุท วา นิหีโน, โย มญฺญติ โส วิวเทถ เตนา}ติ “สทิโสหมสฺมี”ติ วา “เสยฺโยหมสฺมี”ติ วา “หีโนหมสฺมี”ติ วา โย มญฺญติ, โส เตน มาเนน ตาย ทิฏฺฐิยา เตน วา ปุคฺคเลน กลหํ กเรยฺย ภณฺฑนํ กเรยฺย วิคฺคหํ กเรยฺย วิวาทํ กเรยฺย เมธคํ กเรยฺย “น ตฺวํ อิมํ ธมฺมวินยํ อาชานาสิ, อหํ อิมํ ธมฺมวินยํ อาชานามิ, กิํ ตฺวํ อิมํ ธมฺมวินยํ อาชานิสฺสสิ, มิจฺฉาปฏิปนฺโน ตฺวมสิ, อหมสฺมิ สมฺมาปฏิปนฺโน, สหิตํ เม, อสหิตํ เต.\csromanspacer{} ปุเร วจนียํ ปจฺฉา อวจ, ปจฺฉา วจนียํ ปุเร อวจ, อธิจิณฺณํ เต วิปราวตฺตํ, อาโรปิโต เต วาโท, นิคฺคหิโต ตฺวมสิ, จร วาทปฺปโมกฺขาย, นิพฺเพเฐหิ วา สเจ ปโหสี”ติ สโม วิเสสี อุท วา นิหีโน, โย มญฺญติ โส วิวเทถ เตน.}

\prose{\textbf{ตีสุ วิธาสุ อวิกมฺปมาโน, สโม วิเสสีติ น ตสฺส โหตี}ติ ยสฺเสตา ติสฺโส วิธา ปหีนา สมุจฺฉินฺนา วูปสนฺตา ปฏิปสฺสทฺธา อภพฺพุปฺปตฺติกา ญาณคฺคินา ทฑฺฒา.\csromanspacer{} โส ตีสุ วิธาสุ น กมฺปติ น วิกมฺปติ, อวิกมฺปมานสฺส ปุคฺคลสฺส “สทิโสหมสฺมี”ติ วา “เสยฺโยหมสฺมี”ติ วา “หีโนหมสฺมี”ติ วา \textbf{น โหตี}ติ ตีสุ วิธาสุ อวิกมฺปมาโน สโม วิเสสีติ น ตสฺส โหติ.\csromanspacer{} เตนาห ภควา\csromandash{}}

\prose{สโม วิเสสี อุท วา นิหีโน,}

\prose{\textbf{โย มญฺญติ โส วิวเทถ เตน.}}

\csromangathameasure{\textbf{ตีสุ วิธาสุ อวิกมฺปมาโน,} \\ สโม วิเสสีติ น ตสฺส โหตีติ.}
\csromangathagroup{\csromangathastanza{\textbf{ตีสุ วิธาสุ อวิกมฺปมาโน,} \\ สโม วิเสสีติ น ตสฺส โหตีติ.}}

\csromanheader{2}{9. มาคณฺฑิยสุตฺตนิทฺเทส}
\proseitem{78}{\csromanfolio{151}\csromansymbolfootnote{*}{ขุ 1. 410 ปิฏฺเฐ.}สจฺจนฺติ โส \textbf{พฺราหฺมโณ} กิํ วเทยฺย,}

\prose{\textbf{มุสาติ วา โส วิวเทถ เกน.}}

\csromangathameasure{\textbf{ยสฺมิํ สมํ วิสมํ วาปิ นตฺถิ,} \\ \textbf{ส เกน วาทํ ปฏิสํยุเชยฺย.}}
\csromangathagroup{\csromangathastanza{\textbf{ยสฺมิํ สมํ วิสมํ วาปิ นตฺถิ,} \\ \textbf{ส เกน วาทํ ปฏิสํยุเชยฺย.}}}

\prose{\textbf{สจฺจนฺติ โส พฺราหฺมโณ กิํ วเทยฺยา}ติ \textbf{พฺราหฺมโณ}ติ สตฺตนฺนํ ธมฺมานํ พาหิตตฺตา \textbf{พฺราหฺมโณ} ฯเปฯ อสิโต ตาทิ ปวุจฺจเต ส พฺรหฺมา.\csromanspacer{} \textbf{สจฺจนฺติ โส พฺราหฺมโณ กิํ วเทยฺยา}ติ “สสฺสโต โลโก, อิทเมว สจฺจํ โมฆมญฺญนฺ”ติ \textbf{พฺราหฺมโณ} กิํ วเทยฺย กิํ กเถยฺย กิํ ภเณยฺย กิํ ทีปเยยฺย กิํ โวหเรยฺย. “อสสฺสโต โลโก ฯเปฯ เนว โหติ น น โหติ ตถาคโต ปรํ มรณา, อิทเมว สจฺจํ โมฆมญฺญนฺ”ติ \textbf{พฺราหฺมโณ} กิํ วเทยฺย กิํ กเถยฺย กิํ ภเณยฺย กิํ ทีปเยยฺย กิํ โวหเรยฺยาติ สจฺจนฺติ โส \textbf{พฺราหฺมโณ} กิํ วเทยฺย.}

\prose{\textbf{มุสาติ วา โส วิวเทถ เกนา}ติ \textbf{พฺราหฺมโณ} มยฺหํว สจฺจํ, ตุยฺหํ มุสาติ เกน มาเนน, กาย ทิฏฺฐิยา, เกน วา ปุคฺคเลน กลหํ กเรยฺย ภณฺฑนํ กเรยฺย วิคฺคหํ กเรยฺย วิวาทํ กเรยฺย เมธคํ กเรยฺย “น ตฺวํ อิมํ ธมฺมวินยํ อาชานาสิ ฯเปฯ นิพฺเพเฐหิ วา สเจ ปโหสี”ติ มุสาติ วา โส วิวเทถ เกน.}

\prose{\textbf{ยสฺมิํ สมํ วิสมํ วาปิ นตฺถี}ติ \textbf{ยสฺมินฺ}ติ ยสฺมิํ ปุคฺคเล อรหนฺเต ขีณาสเว “สทิโสหมสฺมี”ติ มาโน นตฺถิ. “เสยฺโยหมสฺมี”ติ มาโน นตฺถิ, “หีโนหมสฺมี”ติ โอมาโน นตฺถิ น สนฺติ น สํวิชฺชติ นุปลพฺภติ, ปหีนํ สมุจฺฉินฺนํ วูปสนฺตํ ปฏิปสฺสทฺธํ อภพฺพุปฺปตฺติกํ ญาณคฺคินา ทฑฺฒนฺติ ยสฺมิํ สมํ วิสมํ วาปิ นตฺถิ.}

\prose{\textbf{ส เกน วาทํ ปฏิสํยุเชยฺยา}ติ โส เกน มาเนน กาย ทิฏฺฐิยา เกน วา ปุคฺคเลน วาทํ ปฏิสญฺโญเชยฺย ปฏิพเลยฺย กลหํ กเรยฺย ภณฺฑนํ กเรยฺย วิคฺคหํ กเรยฺย วิวาทํ กเรยฺย เมธคํ กเรยฺย “น ตฺวํ อิมํ ธมฺมวินยํ อาชานาสิ ฯเปฯ นิพฺเพเฐหิ วา สเจ ปโหสี”ติ ส เกน วาทํ ปฏิสํยุเชยฺย.\csromanspacer{} เตนาห ภควา\csromandash{}}

\prose{สจฺจนฺติ โส \textbf{พฺราหฺมโณ} กิํ วเทยฺย,}

\prose{\textbf{มุสาติ วา โส วิวเทถ เกน.}}

\prose{\textbf{ยสฺมิํ สมํ วิสมํ วาปิ นตฺถิ,}}

\prose{\textbf{ส เกน วาทํ ปฏิสํยุเชยฺยา}ติ.}

\proseitem{79}{\csromanfolio{152}\csromansymbolfootnote{*}{ขุ 1. 411; สํ 2. 8 ปิฏฺเฐสุ.}โอกํ ปหาย อนิเกตสารี,}

\prose{\textbf{คาเม อกุพฺพํ มุนิ สนฺถวานิ}\footnote{สนฺธวานิ (ก)}\textbf{.}}

\csromangathameasure{\textbf{กาเมหิ ริตฺโต อปุรกฺขราโน,} \\ \textbf{กถํ น วิคฺคยฺห ชเนน กยิรา.}}
\csromangathagroup{\csromangathastanza{\textbf{กาเมหิ ริตฺโต อปุรกฺขราโน,} \\ \textbf{กถํ น วิคฺคยฺห ชเนน กยิรา.}}}

\prose{\csromansymbolmark{*}อถ โข หาลิทฺทกานิ\footnote{หลิทฺทกานี (สี)} คหปติ เยนายสฺมา มหากจฺจาโน เตนุปสงฺกมิ, อุปสงฺกมิตฺวา อายสฺมนฺตํ มหากจฺจานํ อภิวาเทตฺวา เอกมนฺตํ นิสีทิ, เอกมนฺตํ นิสินฺโน โข หาลิทฺทกานิ คหปติ อายสฺมนฺตํ มหากจฺจานํ เอตทโวจ\csromandash{}วุตฺตมิทํ ภนฺเต กจฺจาน ภควตา อฏฺฐกวคฺคิเย มาคณฺฑิยปเญฺห\csromandash{}}

\prose{“โอกํ ปหาย อนิเกตสารี,}

\prose{\textbf{คาเม อกุพฺพํ มุนิ สนฺถวานิ.}}

\csromangathameasure{กาเมหิ ริตฺโต อปุเรกฺขราโน, \\ กถํ น วิคฺคยฺห ชเนน กยิรา”ติ\textbf{.}}
\csromangathagroup{\csromangathastanza{กาเมหิ ริตฺโต อปุเรกฺขราโน, \\ กถํ น วิคฺคยฺห ชเนน กยิรา”ติ\textbf{.}}}

\prose{อิมสฺส นุ โข ภนฺเต กจฺจาน ภควตา สํขิตฺเตน ภาสิตสฺส กถํ วิตฺถาเรน อตฺโถ ทฏฺฐพฺโพติ\textbf{.}}

\prose{รูปธาตุ โข คหปติ วิญฺญาณสฺส โอโก รูปธาตุราควินิพทฺธญฺจ\footnote{ราควินิพนฺธญฺจ (สฺยา, ก)} ปน วิญฺญาณํ “โอกสารี”ติ วุจฺจติ\textbf{.}\csromanspacer{} เวทนาธาตุ โข คหปติ\textbf{.}\csromanspacer{} สญฺญาธาตุ โข คหปติ\textbf{.}\csromanspacer{} สงฺขารธาตุ โข คหปติ วิญฺญาณสฺส โอโก สงฺขารธาตุราควินิพทฺธญฺจ ปน วิญฺญาณํ “โอกสารี”ติ วุจฺจติ\textbf{.}\csromanspacer{} เอวํ โข คหปติ โอกสารี โหติ\textbf{.}}

\prose{กถญฺจ คหปติ อโนกสารี โหติ, รูปธาตุยา โข คหปติ โย ฉนฺโท โย ราโค ยา นนฺที ยา ตณฺหา เย อุปายุปาทานา\footnote{อุปยุปาทานา (ก)} เจตโส อธิฏฺฐานาภินิเวสานุสยา, เต ตถาคตสฺส ปหีนา อุจฺฉินฺนมูลา ตาลาวตฺถุกตา อนภาวํกตา อายติํ อนุปฺปาทธมฺมา, ตสฺมา ตถาคโต “อโนกสารี”ติ วุจฺจติ\textbf{.}\csromanspacer{} เวทนาธาตุยา โข คหปติ\textbf{.}\csromanspacer{} สญฺญาธาตุยา โข คหปติ\textbf{.}\csromanspacer{} สงฺขารธาตุยา โข คหปติ\textbf{.}\csromanspacer{} วิญฺญาณธาตุยา โข คหปติ โย ฉนฺโท โย ราโค ยา นนฺที ยา ตณฺหา เย อุปายุปาทาน เจตโส อธิฏฺฐานาภินิเวสานุสยา,}

\vspace{\tipitakaparskip}
\csromancenter{มาคณฺฑิยสุตฺตนิทฺเทส เต ตถาคตสฺส ปหีนา อุจฺฉินฺนมูลา ตาลาวตฺถุกตา อนภาวํกตา อายติํ อนุปฺปาทธมฺมา, ตสฺมา ตถาคโต “อโนกสารี”ติ วุจฺจติ.\csromanspacer{} เอวํ โข คหปติ อโนกสารี โหติ.}
\prose{\csromanfolio{153}กถญฺจ คหปติ นิเกตสารี โหติ, รูปนิมิตฺตนิเกตวิสารวินิพนฺธา โข คหปติ “นิเกตสารี”ติ วุจฺจติ.\csromanspacer{} สทฺทนิมิตฺต.\csromanspacer{} คนฺธนิมิตฺต.\csromanspacer{} รสนิมิตฺต.\csromanspacer{} โผฏฺฐพฺพนิมิตฺต.\csromanspacer{} ธมฺมนิมิตฺตนิเกตวิสารวินิพนฺธา โข คหปติ “นิเกตสารี”ติ วุจฺจติ.\csromanspacer{} เอวํ โข คหปติ นิเกตสารี โหติ.}

\prose{กถญฺจ คหปติ อนิเกตสารี โหติ, รูปนิมิตฺตนิเกตวิสารวินิพนฺธา โข คหปติ ตถาคตสฺส ปหีนา อุจฺฉินฺนมูลา ตาลาวตฺถุกตา อนภาวํกตา อายติํ อนุปฺปาทธมฺมา, ตสฺมา ตถาคโต “อนิเกตสารี”ติ วุจฺจติ.\csromanspacer{} สทฺทนิมิตฺต.\csromanspacer{} คนฺธนิมิตฺต.\csromanspacer{} รสนิมิตฺต.\csromanspacer{} โผฏฺฐพฺพนิมิตฺต.\csromanspacer{} ธมฺมนิมิตฺตนิเกตวิสารวินิพนฺธา โข คหปติ ตถาคตสฺส ปหีนา อุจฺฉินฺนมูลา ตาลาวตฺถุกตา อนภาวํกตา อายติํ อนุปฺปาทธมฺมา, ตสฺมา ตถาคโต “อนิเกตสารี”ติ วุจฺจติ.\csromanspacer{} เอวํ โข คหปติ อนิเกตสารี โหติ.}

\prose{กถญฺจ คหปติ คาเม สนฺถวชาโต โหติ, อิธ คหปติ เอกจฺโจ ภิกฺขุ คิหีหิ สํสฏฺโฐ วิหรติ สหนนฺที สหโสกี สุขิเตสุ สุขิโต ทุกฺขิเตสุ ทุกฺขิโต อุปฺปนฺเนสุ กิจฺจกรณีเยสุ อตฺตนา โวโยคํ อาปชฺชติ.\csromanspacer{} เอวํ โข คหปติ คาเม สนฺถวชาโต โหติ.}

\prose{กถญฺจ คหปติ คาเม น สนฺถวชาโต โหติ, อิธ คหปติ เอกจฺโจ ภิกฺขุ คิหีหิ อสํสฏฺโฐ วิหรติ น สหนนฺที น สหโสกี น สุขิเตสุ สุขิโต น ทุกฺขิเตสุ ทุกฺขิโต อุปฺปนฺเนสุ กิจฺจกรณีเยสุ น อตฺตนา โวโยคํ อาปชฺชติ.\csromanspacer{} เอวํ โข คหปติ คาเม น สนฺถวชาโต โหติ.}

\prose{กถญฺจ คหปติ กาเมหิ อริตฺโต โหติ, อิธ คหปติ เอกจฺโจ ภิกฺขุ กาเมสุ อวีตราโค โหติ อวิคตจฺฉนฺโท อวิคตเปโม อวิคตปิปาโส อวิคตปริฬาโห อวิคตตโณฺห.\csromanspacer{} เอวํ โข คหปติ กาเมหิ อริตฺโต โหติ.}

\prose{\csromanfolio{154}กถญฺจ คหปติ กาเมหิ ริตฺโต โหติ, อิธ คหปติ เอกจฺโจ ภิกฺขุ กาเมสุ วีตราโค โหติ วิคตจฺฉนฺโท วิคตเปโม วิคตปิปาโส วิคตปริฬาโห วิคตตโณฺห.\csromanspacer{} เอวํ โข คหปติ กาเมหิ ริตฺโต โหติ.}

\prose{กถญฺจ คหปติ ปุรกฺขราโน โหติ, อิธ คหปติ เอกจฺจสฺส ภิกฺขุโน เอวํ โหติ * “เอวํรูโป สิยํ อนาคตมทฺธานนฺ”ติ “เอวํเวทโน สิยํ.\csromanspacer{} เอวํสญฺโญ สิยํ.\csromanspacer{} เอวํสงฺขาโร สิยํ.\csromanspacer{} เอวํวิญฺญาโณ สิยํ อนาคตมทฺธานนฺ”ติ เอวํ โข คหปติ ปุรกฺขราโน โหติ.}

\prose{กถญฺจ คหปติ อปุรกฺขราโน โหติ, อิธ คหปติ เอกจฺจสฺส ภิกฺขุโน น เอวํ โหติ “เอวํรูโป สิยํ อนาคตมทฺธานนฺ”ติ “เอวํเวทโน สิยํ.\csromanspacer{} เอวํสญฺโญ สิยํ.\csromanspacer{} เอวํสงฺขาโร สิยํ.\csromanspacer{} เอวํวิญฺญาโณ สิยํ อนาคตมทฺธานนฺ”ติ เอวํ โข คหปติ อปุรกฺขราโน โหติ.}

\prose{กถญฺจ คหปติ กถํ วิคฺคยฺห ชเนน กตฺตา โหติ, อิธ คหปติ เอกจฺโจ เอวรูปิํ กถํ กตฺตา โหติ “น ตฺวํ อิมํ ธมฺมวินยํ อาชานาสิ, อหํ อิมํ ธมฺมวินยํ อาชานามิ, กิํ ตฺวํ อิมํ ธมฺมวินยํ อาชานิสฺสสิ, มิจฺฉาปฏิปนฺโน ตฺวมสิ, อหมสฺมิ สมฺมาปฏิปนฺโน, สหิตํ เม, อสหิตํ เต, ปุเร วจนียํ ปจฺฉา อวจ, ปจฺฉา วจนียํ ปุเร อวจ, อธิจิณฺณํ เต วิปราวตฺตํ, อาโรปิโต เต วาโท, นิคฺคหิโต ตฺวมสิ, จรวาทปฺปโมกฺขาย, นิพฺเพเฐหิ วา สเจ ปโหสี”ติ.\csromanspacer{} เอวํ โข คหปติ กถํ วิคฺคยฺห ชเนน กตฺตา โหติ.}

\prose{กถญฺจ คหปติ กถํ น วิคฺคยฺห ชเนน กตฺตา โหติ, อิธ คหปติ เอกจฺโจ น เอวรูปิํ กถํ กตฺตา โหติ “น ตฺวํ อิมํ ธมฺมวินยํ อาชานาสิ ฯเปฯ นิพฺเพเฐหิ วา สเจ ปโหสี”ติ.\csromanspacer{} เอวํ โข คหปติ น วิคฺคยฺห ชเนน กตฺตา โหติ.\csromanspacer{} อิติ โข คหปติ ยํ ตํ วุตฺตํ ภควตา อฏฺฐกวคฺคิเย มาคณฺฑิยปเญฺห\csromandash{}}

\csromangathameasure{“โอกํ ปหาย อนิเกตสารี, \\ คาเม อกุพฺพํ มุนิ สนฺถวานิ. \\ กาเมหิ ริตฺโต อปุรกฺขราโน, \\ กถํ น วิคฺคยฺห ชเนน กริยา”ติ.}
\csromangathagroup{\csromangathastanza{“โอกํ ปหาย อนิเกตสารี, \\ คาเม อกุพฺพํ มุนิ สนฺถวานิ. \\ กาเมหิ ริตฺโต อปุรกฺขราโน, \\ กถํ น วิคฺคยฺห ชเนน กริยา”ติ.}}

\csromanheader{2}{9. มาคณฺฑิยสุตฺตนิทฺเทส}
\prose{\csromanfolio{155}อิมสฺส โข คหปติ ภควตา สํขิตฺเตน ภาสิตสฺส เอวํ วิตฺถาเรน อตฺโถ ทฏฺฐพฺโพติ.\csromanspacer{} เตนาห ภควา\csromandash{}}

\csromangathameasure{โอกํ ปหาย อนิเกตสารี, \\ คาเม อกุพฺพํ มุนิ สนฺถวานิ. \\ กาเมหิ ริตฺโต อปุรกฺขราโน, \\ กถํ น วิคฺคยฺห ชเนน กริยาติ.}
\csromangathagroup{\csromangathastanza{โอกํ ปหาย อนิเกตสารี, \\ คาเม อกุพฺพํ มุนิ สนฺถวานิ. \\ กาเมหิ ริตฺโต อปุรกฺขราโน, \\ กถํ น วิคฺคยฺห ชเนน กริยาติ.}}

\proseitem{80}{\csromansymbolfootnote{*}{ขุ 1. 411 ปิฏฺเฐ.}\textbf{เยหิ วิวิตฺโต วิจเรยฺย โลเก},}

\prose{\textbf{น ตานิ อุคฺคยฺห วเทยฺย นาโค.}}

\csromangathameasure{\textbf{เอลมฺพุชํ กณฺฑกวาริชํ} \textbf{ยถา,} \\ \textbf{ชเลน ปงฺเกน จนูปลิตฺตํ.} \\ \textbf{เอวํ มุนี สนฺติวาโท อคิทฺโธ,} \\ \textbf{กาเม จ โลเก จ อนูปลิตฺโต.}}
\csromangathagroup{\csromangathastanza{\textbf{เอลมฺพุชํ กณฺฑกวาริชํ}\footnote{กณฺฑกํ วาริชํ (สี)} \textbf{ยถา,} \\ \textbf{ชเลน ปงฺเกน จนูปลิตฺตํ.} \\ \textbf{เอวํ มุนี สนฺติวาโท อคิทฺโธ,} \\ \textbf{กาเม จ โลเก จ อนูปลิตฺโต.}}}

\prose{\textbf{เยหิ วิวิตฺโต วิจเรยฺย โลเก}ติ \textbf{เยหี}ติ เยหิ ทิฏฺฐิคเตหิ.\csromanspacer{} \textbf{วิวิตฺโต}ติ กายทุจฺจริเตน ริตฺโต วิวิตฺโต ปวิวิตฺโต, วจีทุจฺจริเตน.\csromanspacer{} มโนทุจฺจริเตน.\csromanspacer{} ราเคน ฯเปฯ สพฺพากุสลาภิสงฺขาเรหิ ริตฺโต วิวิตฺโต ปวิวิตฺโต.\csromanspacer{}\csromanspacer{}\textbf{วิจเรยฺยา}ติ วิจเรยฺย วิหเรยฺย อิริเยยฺย วตฺเตยฺย ปาเลยฺย ยเปยฺย ยาเปยฺย.\csromanspacer{}\csromanspacer{}\textbf{โลเก}ติ มนุสฺสโลเกติ เยหิ วิวิตฺโต วิจเรยฺย โลเก.}

\prose{\textbf{น ตานิ อุคฺคยฺห วเทยฺย นาโค}ติ \textbf{นาโค}ติ อาคุํ น กโรตีติ \textbf{นาโค}, น คจฺฉตีติ \textbf{นาโค}, นาคจฺฉตีติ \textbf{นาโค}.\csromanspacer{} กถํ อาคุํ น กโรตีติ \textbf{นาโค}, \textbf{อาคู} วุจฺจนฺติ ปาปกา อกุสลา ธมฺมา สํกิเลสิกา โปโนภวิกา สทรา ทุกฺขวิปากา อายติํ ชาติชรามรณิยา.}

\prose{\csromansymbolfootnote{+}{ขุ 1. 359; ขุ 8. 83 ปิฏฺเฐสุปิ.}อาคุํ น กโรติ กิญฺจิ โลเก, (สภิยาติ ภควา,)}

\prose{สพฺพสญฺโญเค วิสชฺช พนฺธนานิ.}

\csromangathameasure{สพฺพตฺถ น สชฺชตี วิมุตฺโต, \\ นาโค ตาทิ ปวุจฺจเต ตถตฺตา.}
\csromangathagroup{\csromangathastanza{สพฺพตฺถ น สชฺชตี วิมุตฺโต, \\ นาโค ตาทิ ปวุจฺจเต ตถตฺตา.}}

\prose{เอวํ อาคุํ น กโรตีติ \textbf{นาโค}.}

\prose{\csromanfolio{156}กถํ น คจฺฉตีติ นาโค, น ฉนฺทาคติํ คจฺฉติ, น โทโสคติํ คจฺฉติ, น โมหาคติํ คจฺฉติ, น ภยาคติํ คจฺฉติ, น ราควเสน คจฺฉติ, น โทสวเสน คจฺฉติ, น โมหวเสน คจฺฉติ, น มานวเสน คจฺฉติ, น ทิฏฺฐิวเสน คจฺฉติ, น อุทฺธจฺจวเสน คจฺฉติ, น วิจิกิจฺฉาวเสน คจฺฉติ, นานุสยวเสน คจฺฉติ, น วคฺเคหิ ธมฺเมหิ ยายติ นียติ วุยฺหติ สํหรียติ.\csromanspacer{} เอวํ น คจฺฉตีติ นาโค.}

\prose{กถํ นาคจฺฉตีติ นาโค, โสตาปตฺติมคฺเคน เย กิเลสา ปหีนา, เต กิเลเส น ปุเนติ น ปจฺเจติ น ปจฺฉาคจฺฉติ.\csromanspacer{} สกทาคามิมคฺเคน.\csromanspacer{} อนาคามิมคฺเคน.\csromanspacer{} อรหตฺตมคฺเคน เย กิเลสา ปหีนา.\csromanspacer{} เต กิเลเส น ปุเนติ น ปจฺเจติ น ปจฺจาคจฺฉติ.\csromanspacer{} เอวํ นาคจฺฉตีติ นาโค.}

\prose{\textbf{น ตานิ อุคฺคยฺห วเทยฺย นาโค}ติ นาโค น ตานิ ทิฏฺฐิคตานิ คเหตฺวา อุคฺคเหตฺวา คณิตฺวา ปรามสิตฺวา อภินิวิสิตฺวา วเทยฺย กเถยฺย ภเณยฺย ทีปเยยฺย โวหเรยฺย, “สสฺสโต โลโก ฯเปฯ เนว โหติ น น โหติ ตถาคโต ปรํ มรณา, อิทเมว สจฺจํ โมฆมญฺญนฺ”ติ วเทยฺย กเถยฺย ภเณยฺย ทีปเยยฺย โวหเรยฺยาติ น ตานิ อุคฺคยฺห วเทยฺย นาโค.}

\prose{\textbf{เอลมฺพุชํ กณฺฑกวาริชํ ยถา, ชเลน ปงฺเกน จนูปลิตฺตนฺ}ติ เอลํ วุจฺจติ อุทกํ, อมฺพุชํ วุจฺจติ ปทุมํ, กณฺฑโก วุจฺจติ ขรทณฺโฑ, วาริ วุจฺจติ อุทกํ, วาริชํ วุจฺจติ ปทุมํ วาริสมฺภวํ, ชลํ วุจฺจติ อุทกํ, ปงฺโก วุจฺจติ กทฺทโม, ยถา ปทุมํ วาริชํ วาริสมฺภวํ ชเลน จ ปงฺเกน จ น ลิมฺปติ น ปลิมฺปติ น อุปลิมฺปติ, อลิตฺตํ อสํลิตฺตํ อนุปลิตฺตนฺติ เอลมฺพุชํ กณฺฑกวาริชํ ยถา, ชเลน ปงฺเกน จนูปลิตฺตํ.}

\prose{\textbf{เอวํ มุนี สนฺติวาโท อคิทฺโธ, กาเม จ โลเก จ อนุปลิตฺโต}ติ \textbf{เอวนฺ}ติ โอปมฺมสํปฏิปาทนํ.\csromanspacer{} \textbf{มุนี}ติ \textbf{โมนํ} วุจฺจติ ญาณํ ฯเปฯ สงฺคชาลมติจฺจ โส มุนิ.\csromanspacer{} \textbf{สนฺติวาโท}ติ สนฺติวาโท มุนี ตาณวาโท เลณวาโท สรณวาโท อภยวาโท อจฺจุตวาโท อมตวาโท นิพฺพานวาโทติ เอวํ มุนิ สนฺติวาโท.\csromanspacer{}\csromanspacer{}\textbf{อคิทฺโธ}ติ \textbf{เคโธ} วุจฺจติ ตณฺหา.\csromanspacer{} โย ราโค สาราโค ฯเปฯ อภิชฺฌา โลโภ อกุสลมูลํ.\csromanspacer{} ยสฺเสโส \textbf{เคโธ} ปหีโน สมุจฺฉินฺโน วูปสนฺโต ปฏิปสฺสทฺโธ อภพฺพุปฺปตฺติโก ญาณคฺคินา ทฑฺโฒ, โส วุจฺจติ อคิทฺโธ.\csromanspacer{} โส รูเป อคิทฺโธ, สทฺเท.\csromanspacer{} คนฺเธ.\csromanspacer{} รเส.\csromanspacer{} โผฏฺฐพฺเพ.\csromanspacer{} กุเล.\csromanspacer{} คเณ.\csromanspacer{} อาวาเส.\csromanspacer{} ลาเภ.\csromanspacer{} ยเส.\csromanspacer{} ปสํสาย.\csromanspacer{} สุเข.\csromanspacer{} จีวเร.}

\vspace{\tipitakaparskip}
\csromancenter{มาคณฺฑิยสุตฺตนิทฺเทส ปิณฺฑปาเต.\csromanspacer{} เสนาสเน.\csromanspacer{} คิลานปจฺจยเภสชฺชปริกฺขาเร.\csromanspacer{} กามธาตุยา.\csromanspacer{} รูปธาตุยา.\csromanspacer{} อรูปธาตุยา.\csromanspacer{} กามภเว.\csromanspacer{} รูปภเว.\csromanspacer{} อรูปภเว.\csromanspacer{} สญฺญาภเว.\csromanspacer{} อสญฺญาภเว.\csromanspacer{} เนวสญฺญานาสญฺญาภเว.\csromanspacer{} เอกโวการภเว.\csromanspacer{} จตุโวการภเว.\csromanspacer{} ปญฺจโวการภเว.\csromanspacer{} อตีเต.\csromanspacer{} อนาคเต.\csromanspacer{} ปจฺจุปฺปนฺเน.\csromanspacer{} ทิฏฺฐ สุต มุต วิญฺญาตพฺเพสุ ธมฺเมสุ อคิทฺโธ อคธิโต อมุจฺฉิโต อนชฺโฌสนฺโน\footnote{อนชฺฌาปนฺโน (สี), อนชฺโฌปนฺโน (สฺยา)} วีตเคโธ วิคตเคโธ จตฺตเคโธ วนฺตเคโธ มุตฺตเคโธ ปหีนเคโธ ปฏินิสฺสฏฺฐเคโธ, วิตราโค วิคตราโค จตฺตราโค วนฺตราโค มุตฺตราโค ปหีนราโค ปฏินิสฺสฏฺฐราโค นิจฺฉาโต นิพฺพุโต สีติภูโต สุขปฺปฏิสํเวที พฺรหฺมภูเตน อตฺตนา วิหรตีติ เอวํ มุนิ สนฺติวาโท อคิทฺโธ.}
\prose{\csromanfolio{157}\textbf{กาเม จ โลเก จ อนูปลิตฺโต}ติ \textbf{กามา}ติ อุทานโต เทฺว \textbf{กามา} วตฺถุ\textbf{กามา} จ กิเลส\textbf{กามา} จ ฯเปฯ อิเม วุจฺจนฺติ วตฺถุ\textbf{กามา} ฯเปฯ อิเม วุจฺจนฺติ กิเลส\textbf{กามา}.\csromanspacer{} \textbf{โลเก}ติ อปายโลเก มนุสฺสโลเก เทวโลเก, ขนฺธโลเก ธาตุโลเก อายตนโลเก.\csromanspacer{} \textbf{เลปา}ติ เทฺว เลปา ตณฺหาเลโป จ ทิฏฺฐิเลโป จ ฯเปฯ อยํ ตณฺหาเลโป ฯเปฯ อยํ ทิฏฺฐิเลโป.\csromanspacer{} มุนิ ตณฺหาเลปํ ปหาย ทิฏฺฐิเลปํ ปฏินิสฺสชฺชิตฺวา กาเม จ โลเก จ น ลิมฺปติ น ปลิมฺปติ น อุปลิมฺปติ, อลิตฺโต อปลิตฺโต อนุปลิตฺโต นิกฺขนฺโต นิสฺสโฏ วิปฺปมุตฺโต วิสญฺญุตฺโต วิมริยาทิกเตน เจตสา วิหรตีติ เอวํ มุนิ สนฺติวาโท อคิทฺโธ, กาเม จ โลเก จ อนูปลิตฺโต, เตนาห ภควา\csromandash{}}

\csromangathameasure{เยหิ วิวิตฺโต วิจเรยฺย โลเก, \\ น ตานิ อุคฺคยฺห วเทยฺย นาโค. \\ เอลมฺพุชํ กณฺฑกวาริชํ ยถา, \\ ชเลน ปงฺเกน จนูปลิตฺตํ. \\ เอวํ มุนิ สนฺติวาโท อคิทฺโธ,}
\csromangathagroup{\csromangathastanza{เยหิ วิวิตฺโต วิจเรยฺย โลเก, \\ น ตานิ อุคฺคยฺห วเทยฺย นาโค. \\ เอลมฺพุชํ กณฺฑกวาริชํ ยถา, \\ ชเลน ปงฺเกน จนูปลิตฺตํ.}\csromangathastanza{เอวํ มุนิ สนฺติวาโท อคิทฺโธ,}}

\prose{\textbf{กาเม จ โลเก จ อนูปลิตฺโต}ติ.}

\csromanhangingbegin
\hangingheaditem{81}{น เวทคู ทิฏฺฐิยายโก\footnote{ทิฏฺฐิยา (สี, สฺยา) ขุ 1. 411 ปิฏฺเฐปิ.} \textbf{น มุติยา,}}
\hangingbody{\textbf{ส มานเมติ น หิ ตมฺมโย โส.}}
\hangingbody{\textbf{น กมฺมุนา โนปิ สุเตน เนยฺโย,}}
\hangingbody{\textbf{อนูปนีโต ส นิเวสเนสุ.}}
\csromanhangingend

\csromanmatikaanchor{mtk.10}
\csromantocmarkref{subsection}{9. มาคณฺฑิยสุตฺตนิทฺเทส}{mtk.10}
\bookmark[dest={mtk.10},level=2]{9. มาคณฺฑิยสุตฺตนิทฺเทส}
\prose{\csromanfolio{158}\textbf{น เวทคู ทิฏฺฐิยายโก น มุติยา, สมานเมตี}ติ \textbf{นา}ติ ปฏิกฺเขโป. * \textbf{เวทคู}ติ \textbf{เวโท} วุจฺจติ จตูสุ มคฺเคสุ ญาณํ, ปญฺญา ปญฺญินฺทฺริยํ ปญฺญาพลํ ธมฺมวิจยสมฺโพชฺฌงฺโค วีมํสา วิปสฺส\textbf{นา} สมฺมาทิฏฺฐิ.\csromanspacer{} เตหิ เวเทหิ ชาติชรามรณสฺส อนฺตคโต อนฺตปฺปตฺโต, โกฏิคโต โกฏิปฺปตฺโต ปริยนฺตคโต ปริยนฺตปฺปตฺโต, โวสานคโต โวสานปฺปตฺโต, ตาณคโต ตาณปฺปตฺโต, เลณคโต เลณปฺปตฺโต, สรณคโต สรณปฺปตฺโต, อภยคโต อภยปฺปตฺโต, อจฺจุตคโต อจฺจุตปฺปตฺโต, อมตคโต อมตปฺปตฺโต, นิพฺพานคโต นิพฺพานปฺปตฺโต, เวทานํ วา อนฺตํ คโตติ เวทคู, เวทหิ วา อนฺตํ คโตติ เวทคู, สตฺตนฺนํ วา ธมฺมานํ วิทิตตฺตา เวทคู, สกฺกายทิฏฺฐิ วิทิตา โหติ, วิจิกิจฺฉา วิทิตา โหติ, สีลพฺพตปรามาโส วิทิโต โหติ, ราโค วิทิโต โหติ, โทโส วิทิโต โหติ, โมโห วิทิโต โหติ, มาโน วิทิโต โหติ, วิทิตาสฺส โหนฺติ ปาปกา อกุสลา ธมฺมา สํกิเลสิกา โปโนภวิกา สทรา ทุกฺขวิปากา อายติํ ชาติชรามรณิยา.}

\vspace{\tipitakaparskip}
\csromancenter{\csromansymbolfootnote{+}{ขุ 1. 360 ปิฏฺเฐปิ.}เวทานิ วิเจยฺย เกวลานิ, (สภิยาติ ภควา,)}
\prose{สมณานํ ยานีธตฺถิ พฺราหฺมณานํ.}

\csromangathameasure{สพฺพเวท\textbf{นา}สุ วีตราโค, \\ สพฺพํ เวทมติจฺจ เวทคู โสติ.}
\csromangathagroup{\csromangathastanza{สพฺพเวท\textbf{นา}สุ วีตราโค, \\ สพฺพํ เวทมติจฺจ เวทคู โสติ.}}

\prose{\textbf{น ทิฏฺฐิยา}ติ ตสฺส ทฺวาสฏฺฐิ ทิฏฺฐิคตานิ ปหี\textbf{นา}นิ สมุจฺฉินฺ\textbf{นา}นิ วูปสนฺตานิ ปฏิปสฺสทฺธานิ อภพฺพุปฺปตฺติกานิ ญาณคฺคิ\textbf{นา} ทฑฺฒานิ.\csromanspacer{} โส ทิฏฺฐิยา น ยายติ น นียติ น วุยฺหติ น สํหรียติ.\csromanspacer{} นปิ ตํ ทิฏฺฐิคตํ สารโต ปจฺเจติ น ปจฺจาคจฺฉตีติ น เวทคู ทิฏฺฐิยา.\csromanspacer{}\csromanspacer{}\textbf{น มุติยา}ติ มุตรูเปน วา ปรโต โฆเสน วา มหาชนสมฺมุติยา วา มานํ เนติ น อุเปติ น อุปคจฺฉติ น คณฺหาติ น ปรามสติ \textbf{นา}ภินิวิสตีติ น เวทคู ทิฏฺฐิยายโก น มุติยา ส มานเมติ.}

\prose{\textbf{น หิ ตมฺมโย โส}ติ น ตณฺหาวเสน ทิฏฺฐิวเสน ตมฺมโย โหติ ตปฺปรโม ตปฺปรายโณ.\csromanspacer{} ยโต ตณฺหา จ ทิฏฺฐิ จ มาโน จสฺส ปหี\textbf{นา} โหนฺติ อุจฺฉินฺนมูลา ตาลาวตฺถุกตา อนภาวํกตา อายติํ อนุปฺปาทธมฺมา.\csromanspacer{} เอตฺตาวตา น ตมฺมโย โหติ น ตปฺปรโม น ตปฺปรายโณติ ส มานเมติ น หิ ตมฺมโย โส.}

\csromanheader{2}{9. มาคณฺฑิยสุตฺตนิทฺเทส}
\prose{\csromanfolio{159}\textbf{น กมฺมุนา โนปิ สุเตน เนยฺโย}ติ \textbf{น} \textbf{กมฺมุนา}ติ ปุญฺญาภิสงฺขาเร\textbf{น} วา อปุญฺญาภิสงฺขาเร\textbf{น} วา อาเนญฺชาภิสงฺขาเร\textbf{น} วา \textbf{น} ยายติ \textbf{น} นียติ \textbf{น} วุยฺหติ \textbf{น} สํหรียตีติ \textbf{น} \textbf{กมฺมุนา}.\csromanspacer{}\csromanspacer{}\textbf{โนปิ สุเตน เนยฺโย}ติ สุตสุทฺธิยา วา ปรโต โฆเส\textbf{น} วา มหาช\textbf{น}สมฺมุติยา วา \textbf{น} ยายติ \textbf{น} นียติ \textbf{น} วุยฺหติ \textbf{น} สํหรียตีติ \textbf{น} \textbf{กมฺมุนา} โนปิ สุเต\textbf{น} เนยฺโย.}

\prose{\textbf{อนูปนีโต ส นิเวสเนสู}ติ \textbf{อุปยา}ติ เทฺว \textbf{อุปยา} ตณฺหูปโย จ ทิฏฺฐูปโย จ ฯเปฯ อยํ ตณฺหูปโย ฯเปฯ อยํ ทิฏฺฐูปโย.\csromanspacer{} ตสฺส ตณฺหูปโย ปหีโน, ทิฏฺฐูปโย ปฏินิสฺสฏฺโฐ.\csromanspacer{} ตณฺหูปยสฺส ปหี\textbf{น}ตฺตา, ทิฏฺฐูปยสฺส ปฏินิสฺสฏฺฐตฺตา โส นิเวสเนสุ อนุปนีโต อนุปลิตฺโต อนุปคโต อ\textbf{น}ชฺโฌสิโต อ\textbf{น}ธิมุตฺโต นิกฺขนฺโต นิสฺสโฏ วิปฺปมุตฺโต วิสญฺญุตฺโต วิมริยาทิกเต\textbf{น} เจตสา วิหรตีติ อนุปนีโต ส นิเวสเนสุ.\csromanspacer{} เตนาห ภควา\csromandash{}}

\csromangathameasure{น เวทคู ทิฏฺฐิยายโก \textbf{น} มุติยา, \\ ส มา\textbf{น}เมติ \textbf{น} หิ ตมฺมโย โส.}
\csromangathagroup{\csromangathastanza{น เวทคู ทิฏฺฐิยายโก \textbf{น} มุติยา, \\ ส มา\textbf{น}เมติ \textbf{น} หิ ตมฺมโย โส.}}

\prose{\textbf{น กมฺมุนา โนปิ สุเตน เนยฺโย},}

\prose{อนุปนีโต ส นิเวสเนสูติ.}

\proseitem{82}{\csromansymbolmark{*}\textbf{สญฺญาวิรตฺตสฺส น สนฺติ คนฺถา},}

\prose{\textbf{ปญฺญาวิมุตฺตสฺส น สนฺติ โมหา.}}

\csromangathameasure{\textbf{สญฺญญฺจ ทิฏฺฐิญฺจ เย อคฺคเหสุํ,} \\ \textbf{เต ฆฏฺฏยนฺตา} \textbf{วิจรนฺติ โลเก.}}
\csromangathagroup{\csromangathastanza{\textbf{สญฺญญฺจ ทิฏฺฐิญฺจ เย อคฺคเหสุํ,} \\ \textbf{เต ฆฏฺฏยนฺตา}\footnote{ฆฏฺฏมานา (สี, ก) ขุ 1. 411 ปิฏฺเฐปิ.} \textbf{วิจรนฺติ โลเก.}}}

\prose{\textbf{สญฺญาวิรตฺตสฺส น สนฺติ คนฺถา}ติ โย สมถปุพฺพงฺคมํ อริยมคฺคํ ภาเวติ, ตสฺส อาทิโต อุปาทาย คนฺถา วิกฺขมฺภิตา โหนฺติ, อรหตฺเต ปตฺเต อรหโต คนฺถา จ โมหา จ นีวรณา จ กามสญฺญา พฺยาปาทสญฺญา วิหิํสาสญฺญา ทิฏฺฐิสญฺญา จ ปหีนา โหนฺติ อุจฺฉินฺ\textbf{น}มูลา ตาลาวตฺถุกตา อ\textbf{น}ภาวํกตา อายติํ อนุปฺปาทธมฺมาติ สญฺญาวิรตฺตสฺส \textbf{น} สนฺติ คนฺถา.}

\prose{\textbf{ปญฺญาวิมุตฺตสฺส น สนฺติ โมหา}ติ โย วิปสฺสนาปุพฺพงฺคมํ อริยมคฺคํ ภาเวติ, ตสฺส อาทิโต อุปาทาย โมหา วิกฺขมฺภิตา โหนฺติ, \csromanfolio{160}อรหตฺเต ปตฺเต อรหโต โมหา จ คนฺถา จ นีวรณา จ กามสญฺญา พฺยาปาทสญฺญา วิหิํสาสญฺญา ทิฏฺฐิสญฺญา จ ปหีนา โหนฺติ อุจฺฉินฺนมูลา ตาลาวตฺถุกตา อนภาวํกตา อายติํ อนุปฺปาทธมฺมาติ ปญฺญาวิมุตฺตสฺส น สนฺติ โมหา.}

\prose{\textbf{สญฺญญฺจ ทิฏฺฐิญฺจ เย อคฺคเหสุํ, เต ฆฏฺฏยนฺตา วิจรนฺติ โลเก}ติ เย สญฺญํ คณฺหนฺติ กามสญฺญํ พฺยาปาทสญฺญํ วิหิํสาสญฺญํ, เต สญฺญาวเสน ฆฏฺเฏนฺติ สํฆฏฺเฏนฺติ.\csromanspacer{} ราชาโนปิ ราชูหิ วิวทนฺติ, ขตฺติยาปิ ขตฺติเยหิ วิวทนฺติ, พฺราหฺมณาปิ พฺราหฺมเณหิ วิวทนฺติ, คหปตีปิ คหปตีหิ วิวทนฺติ, มาตาปิ ปุตฺเตน วิวทติ, ปุตฺโตปิ มาตรา วิวทติ, ปิตาปิ ปุตฺเตน วิวทติ, ปุตฺโตปิ ปิตรา วิวทติ, ภาตาปิ ภาตรา วิวทติ, ภคินีปิ ภคินิยา วิวทติ, ภาตาปี ภคินิยา วิวทติ, ภคินีปิ ภาตรา วิวทติ, สหาโยปิ สหาเยน วิวทติ.\csromanspacer{} เต ตตฺถ กลหวิคฺคหวิวาทมาปนฺนา อญฺญมญฺญํ ปาณีหิปิ อุปกฺกมนฺติ, เลฑฺฑูหิปิ\footnote{เลฏฺฏูหิปิ (ก)} อุปกฺกมนฺติ, ทณฺเฑหิปิ อุปกมนฺติ, สตฺเถหิปิ อุปกฺกมนฺติ.\csromanspacer{} เต ตตฺถ มรณมฺปิ นิคจฺฉนฺติ มรณมตฺตมฺปิ ทุกฺขํ.\csromanspacer{} เย ทิฏฺฐิํ คณฺหนฺติ “สสฺสโต โลโก”ติ วา ฯเปฯ “เนว โหติ น น โหติ ตถาคโต ปรํ มรณา”ติ วา, เต ทิฏฺฐิวเสน ฆฏฺเฏนฺติ สํฆฏฺเฏนฺติ, สตฺถารโต สตฺถารํ ฆฏฺเฏนฺติ, ธมฺมกฺขานโต ธมฺมกฺขานํ ฆฏฺเฏนฺติ, คณโต คณํ ฆฏฺเฏนฺติ, ทิฏฺฐิยา ทิฏฺฐิ ฆฏฺเฏนฺติ, ปฏิปทาย ปฏิปทํ ฆฏฺเฏนฺติ, มคฺคโต มคฺคํ ฆฏฺเฏนฺติ.}

\prose{อถ วา เต วิวทนฺติ, กลหํ กโรนฺติ, ภณฺฑนํ กโรนฺติ, วิคฺคหํ กโรนฺติ, วิวาทํ กโรนฺติ, เมธคํ กโรนฺติ “น ตฺวํ อิมํ ธมฺมวินยํ อาชานาสิ ฯเปฯ นิพฺเพเฐหิ วา สเจ ปโหสี”ติ.\csromanspacer{} เตสํ อภิสงฺขารา อปฺปหีนา, อภิสงฺขารานํ อปฺปหีนตฺตา คติยา ฆฏฺเฏนฺติ, นิรเย ฆฏฺเฏนฺติ, ติรจฺฉานโยนิยา ฆฏฺเฏนฺติ, เปตฺติวิสเย ฆฏฺเฏนฺติ, มนุสฺสโลเก ฆฏฺเฏนฺติ, เทวโลเก ฆฏฺเฏนฺติ, คติยา คติํ.\csromanspacer{} อุปปตฺติยา อุปปตฺติํ.\csromanspacer{} ปฏิสนฺธิยา ปฏิสนฺธิํ.\csromanspacer{} ภเวน ภวํ.\csromanspacer{} สํสาเรน สํสารํ.\csromanspacer{} วฏฺเฏน วฏฺฏํ ฆฏฺเฏนฺติ สํฆฏฺเฏนฺติ วทนฺติ วิจรนฺติ วิหรนฺติ อิริยนฺติ วตฺเตนฺติ ปาเลนฺติ ยเปนฺติ ยาเปนฺติ.\csromanspacer{} \textbf{โลเก}ติ อปายโลเก ฯเปฯ อายตนโลเกติ สญฺญญฺจ ทิฏฺฐิญฺจ เย อคฺคเหสุํ, เต ฆฏฺฏยนฺตา วิจรนฺติ โลเก.\csromanspacer{} เตนาห ภควา\csromandash{}}

\csromanheader{2}{9. มาคณฺฑิยสุตฺ\textbf{ตนฺ}อิทฺเทส}
\csromangathameasure{สญฺญาวิรตฺตสฺส น สนฺติ คนฺถา, \\ ปญฺญาวิมุตฺตสฺส น สนฺติ โมหา. \\ สญฺญญฺจ ทิฏฺฐิญฺจ เย อคฺคเหสุํ, \\ เก ฆฏฺฏยนฺตา วิจรนฺติ โลเกติ. มาคณฺฑิยสุตฺ\textbf{ตนฺ}อิทฺเทโส นวโม.}
\csromangathagroup{\csromangathastanza{\csromanfolio{161}สญฺญาวิรตฺตสฺส น สนฺติ คนฺถา, \\ ปญฺญาวิมุตฺตสฺส น สนฺติ โมหา. \\ สญฺญญฺจ ทิฏฺฐิญฺจ เย อคฺคเหสุํ, \\ เก ฆฏฺฏยนฺตา วิจรนฺติ โลเกติ.\csromanspacer{} มาคณฺฑิยสุตฺ\textbf{ตนฺ}อิทฺเทโส นวโม.}}

\csromanheader{2}{10. ปุราเภทสุตฺ\textbf{ตนฺ}อิทฺเทส}
\csromanheader{2}{อถ ปุราเภทสุตฺ\textbf{ตนฺ}อิทฺเทสํ วกฺขติ\csromandash{}}
\proseitem{83}{\csromansymbolfootnote{*}{ขุ 1. 411 ปิฏฺเฐ.}\textbf{กถํทสฺสี กถํสีโล, อุปสนฺโตติ วุจฺจติ}.}

\csromangathameasure{ตํ เม โคตม ปพฺรูหิ, ปุจฺฉิโต อุตฺตมํ นรํ.}
\csromangathagroup{\csromangathastanza{ตํ เม โคตม ปพฺรูหิ, ปุจฺฉิโต อุตฺตมํ นรํ.}}

\prose{\textbf{กถํทสฺสี กถํสีโล, อุปสนฺโตติ วุจฺจติ}ติ \textbf{กถํทสฺสี}ติ กีทิเสน ทสฺสเนน สมนฺนาคโต กิํสณฺฐิเตน กิํปกาเรน กิํปฏิภาเคนาติ \textbf{กถํทสฺสี}.\csromanspacer{}\csromanspacer{}\textbf{กถํสีโล}ติ กีทิเสน สีเลน สมนฺนาคโต กิํสณฺฐิเตน กิํปกาเรน กิํปฏิภาเคนาติ \textbf{กถํทสฺสี} กถํสีโล.\csromanspacer{}\csromanspacer{}\textbf{อุปสนฺโตติ วุจฺจตี}ติ สนฺโต อุปสนฺโต วูปสนฺโต นิพฺพุโต ปฏิปสฺสทฺโธติ วุจฺจติ ปวุจฺจติ กถียติ ภณียติ ทีปียติ โวหรียติ.\csromanspacer{}\csromanspacer{}กถํทสฺสีติ อธิปญฺญํ ปุจฺฉติ, กถํสีโลติ อธิสีลํ ปุจฺฉติ.\csromanspacer{} อุปสนฺโตติ อธิจิตฺตํ ปุจฺฉตีติ \textbf{กถํทสฺสี} กถํสีโล อุปสนฺโตติ วุจฺจติ.}

\prose{\textbf{ตํ เม โคตม ปพฺรูหี}ติ \textbf{ตนฺ}ติ ยํ ปุจฺฉามิ, ยํ ยาจามิ, ยํ อชฺเฌสามิ, ยํ ปสาเทมิ.\csromanspacer{} \textbf{โคตมา}ติ โส นิมฺมิโต พุทฺธํ ภควนฺตํ โคตฺเตน อาลปติ.\csromanspacer{} \textbf{ปพฺรูหี}ติ พฺรูหิ อาจิกฺข เทเสหิ ปญฺญเปหิ ปฏฺฐเปหิ วิวร วิภช อุตฺตานีกโรหิ ปกาเสหีติ ตํ เม โคตม ปพฺรูหิ.}

\prose{\textbf{ปุจฺฉิโต อุตฺตมํ นรนฺ}ติ \textbf{ปุจฺฉิโต}ติ ปุฏฺโฐ \textbf{ปุจฺฉิโต} ยาจิโต อชฺเฌสิโต ปสาทิโต.\csromanspacer{} \textbf{อุตฺตมํ นรนฺ}ติ อคฺคํ เสฏฺฐํ วิสิฏฺฐํ ปาโมกฺขํ อุตฺตมํ ปวรํ นรนฺติ \textbf{ปุจฺฉิโต} อุตฺตมํ นรํ.\csromanspacer{} เตนาห โส นิมฺมิโต\csromandash{}}

\prose{\textbf{กถํทสฺสี กถํสีโล, อุปสนฺโตติ วุจฺจติ}.}

\csromangathameasure{ตํ เม โคตม ปพฺรูหิ, ปุจฺฉิโต อุตฺตมํ นรนฺติ.}
\csromangathagroup{\csromangathastanza{ตํ เม โคตม ปพฺรูหิ, ปุจฺฉิโต อุตฺตมํ นรนฺติ.}}

\proseitem{84}{\csromanfolio{162}\csromansymbolfootnote{*}{ขุ 1. 411 ปิฏฺเฐ.}\textbf{วีตตโณฺห ปุราเภทา}, (อิติ \textbf{ภควา},)}

\prose{\textbf{ปุพฺพมนฺต’มนิสฺสิโต.}}

\csromangathameasure{\textbf{เวมชฺเฌ นุปสงฺเขยฺโย,} \\ \textbf{ตสฺส นตฺถิ ปุรกฺขตํ.}}
\csromangathagroup{\csromangathastanza{\textbf{เวมชฺเฌ นุปสงฺเขยฺโย,} \\ \textbf{ตสฺส นตฺถิ ปุรกฺขตํ.}}}

\prose{\textbf{วีตตโณฺห ปุราเภทา}ติ ปุรา กายสฺส เภทา, ปุรา อตฺตภาวสฺส เภทา ปุรา กเฬวรสฺส นิกฺเขปา ปุรา ชีวิตินฺทฺริยสฺส อุปจฺเฉทา วีตตโณฺห วิคตตโณฺห จตฺตตโณฺห วนฺตตโณฺห มุตฺตตโณฺห ปหีนตโณฺห ปฏินิสฺสฏฺฐตโณฺห, วีตราโค วิคตราโค จตฺตราโค วนฺตราโค มุตฺตราโค ปหีนราโค ปฏินิสฺสฏฺฐราโค, นิจฺฉาโต นิพฺพุโต สีติภูโต สุขปฺปฏิสํเวที พฺรหฺมภูเตน อตฺตนา วิหรติ.}

\prose{\textbf{ภควา}ติ คารวาธิวจนํ.\csromanspacer{} อปิ จ ภคฺคราโคติ \textbf{ภควา}, ภคฺคโทโสติ \textbf{ภควา}, ภคฺคโมโหติ \textbf{ภควา}, ภคฺคมาโนติ \textbf{ภควา}, ภคฺคทิฏฺฐีติ \textbf{ภควา}, ภคฺคตโณฺหติ \textbf{ภควา}, ภคฺคกิเลโสติ \textbf{ภควา}, ภชิ วิภชิ ปวิภชิ ธมฺมรตนนฺติ \textbf{ภควา}, ภวานํ อนฺตกโรติ \textbf{ภควา}, ภาวิตกาโย ภาวิตสีโล ภาวิตจิตฺโต ภาวิตปญฺโญติ \textbf{ภควา}, ภชิ วา \textbf{ภควา} อรญฺญวนปตฺถานิ ปนฺตานิ เสนาสนานิ อปฺปสทฺทานิ อปฺปนิคฺโฆสานิ วิชนวาตานิ มนุสฺสราหสฺเสยฺยกานิ ปฏิสลฺลานสารุปฺปานีติ \textbf{ภควา}.\csromanspacer{} ภาคี วา \textbf{ภควา} จีวรปิณฺฑปาตเสนาสนคิลานปจฺจย- เภสชฺชปริกฺขรานนฺติ \textbf{ภควา}.\csromanspacer{} ภาคี วา \textbf{ภควา} อตฺถรสสฺส ธมฺมรสสฺส วิมุตฺติรสสฺส อธิสีลสฺส อธิจิตฺตสฺส อธิปญฺญายาติ \textbf{ภควา}.\csromanspacer{} ภาคี วา \textbf{ภควา} จตุนฺนํ ฌานานํ จตุนฺนํ อปฺปมญฺญานํ จตุนฺนํ อรูปสมาปตฺตีนนฺติ \textbf{ภควา}.\csromanspacer{} ภาคี วา \textbf{ภควา} อฏฺฐนฺนํ วิโมกฺขานํ อฏฺฐนฺนํ อภิภายตนานํ นวนฺนํ อนุปุพฺพวิหารสมาปตฺตีนนฺติ \textbf{ภควา}.\csromanspacer{} ภาคี วา \textbf{ภควา} ทสนฺนํ ปญฺญาภาวนานํ ทสนฺนํ กสิณสมาปตฺตีนํ อานาปานสฺสติสมาธิสฺส อสุภสมาปตฺติยาติ \textbf{ภควา}.\csromanspacer{} ภาคี วา \textbf{ภควา} จตุนฺนํ สติปฏฺฐานานํ จตุนฺนํ สมฺมปฺปธานานํ จตุนฺนํ อิทฺธิปาทานํ ปญฺจนฺนํ อินฺทฺริยานํ ปญฺจนฺนํ พลานํ สตฺตนฺนํ โพชฺฌงฺคานํ อริยสฺส อฏฺฐงฺคิกสฺส มคฺคสฺสาติ \textbf{ภควา}.\csromanspacer{} ภาคี วา \textbf{ภควา} ทสนฺนํ ตถาคตพลานํ จตุนฺนํ เวสารชฺชานํ จตุนฺนํ ปฏิสมฺภิทานํ ฉนฺนํ อภิญฺญานํ ฉนฺนํ พุทฺธธมฺมานนฺติ \textbf{ภควา}.\csromanspacer{} \textbf{ภควา}ติ เนตํ นามํ มาตรา กตํ, น ปิตรา กตํ, น ภาตรา กตํ, น ภคินิยา กตํ, น มิตฺตามจฺเจหิ กตํ, น ญาติสาโลหิเตหิ}

\vspace{\tipitakaparskip}
\csromancenter{ปุราเภทสุตฺตนิทฺเทส กตํ, น สมณพฺราหฺมเณหิ กตํ, น เทวตาหิ กตํ, วิโมกฺขนฺติกเมตํ พุทฺธานํ ภควนฺตานํ โพธิยา มูเล สห สพฺพญฺญุตญาณสฺส ปฏิลาภา สจฺฉิกา ปญฺญาตฺติ ยทิทํ ภควาติ วีตตโณฺห ปุราเภทาติ ภควา.}
\prose{\csromanfolio{163}\textbf{ปุพฺพมนฺตมนิสฺสิโต}ติ ปุพฺพนฺโต วุจฺจติ อตีโต อทฺธา.\csromanspacer{} อตีตํ อทฺธานํ อารพฺภ ตณฺหา ปหีนา, ทิฏฺฐิ ปฏินิสฺสฏฺฐา, ตณฺหาย ปหีนตฺตา, ทิฏฺฐิยา ปฏินิสฺสฏฺฐตฺตา เอวมฺปิ ปุพฺพมนฺตมนิสฺสิโต.\csromanspacer{}\csromanspacer{}อถ วา * “เอวํรูโป อโหสิํ อตีตมทฺธานนฺ”ติ ตตฺถ นนฺทิํ น สมนฺนาเนติ, “เอวํ เวทโน อโหสิํ.\csromanspacer{} เอวํสญฺโญ อโหสิํ.\csromanspacer{} เอวํสงฺขาโร อโหสิํ.\csromanspacer{} เอวํวิญฺญาโณ อโหสิํ อตีตมทฺธานนฺ”ติ ตตฺถ นนฺทิํ น สมนฺนาเนติ, เอวมฺปิ ปุพฺพมนฺตมนิสฺสิโต.\csromanspacer{}\csromanspacer{}อถ วา อิติ เม จกฺขุ\footnote{จกฺขุํ (สี, ก) ม 3. 236 ปิฏฺเฐปิ.} อโหสิ อตีตมทฺธานํ, อิติ รูปาติ, ตตฺถ น ฉนฺทราคปฏิพทฺธํ โหติ วิญฺญาณํ, น ฉนฺทราคปฏิพทฺธตฺตา วิญฺญาณสฺส น ตทภินนฺทติ, น ตทภินนฺทนฺโต เอวมฺปิ ปุพฺพมนฺตมนิสฺสิโต.\csromanspacer{} อิติ เม โสตํ อโหสิ อตีตมทฺธานํ, อิติ สทฺทาติ.\csromanspacer{} อิติ เม ฆานํ อโหสิ อตีตมทฺธานํ, อิติ คนฺธาติ.\csromanspacer{} อิติ เม ชิวฺหา อโหสิ อตีตมทฺธานํ, อิติ รสาติ.\csromanspacer{} อิติ เม กาโย อโหสิ อตีตมทฺธานํ, อิติ โผฏฺฐพฺพาติ.\csromanspacer{} อิติ เม มโน อโหสิ อตีตมทฺธานํ, อิติ ธมฺมาติ ตตฺถ น ฉนฺทราคปฏิพทฺธํ โหติ วิญฺญาณํ.\csromanspacer{} น ฉนฺทราคปฏิพทฺธตฺตา วิญฺญาณสฺส น ตทภินนฺทติ, น ตทภินนฺทนฺโต เอวมฺปิ ปุพฺพมนฺตมนิสฺสิโต.\csromanspacer{}\csromanspacer{}อถ วา ยานิ ตานิ ปุพฺเพ มาตุคาเมน สทฺธิํ หสิตลปิตกีฬิตานิ น ตทสฺสาเทติ, น ตํ นิกาเมติ, น จ เตน วิตฺติํ อาปชฺชติ เอวมฺปิ ปุพฺพมนฺตมนิสฺสิโต.}

\prose{\textbf{เวมชฺเฌ นุปสงฺเขยฺโย}ติ \textbf{เวมชฺฌํ} วุจฺจติ ปจฺจุปฺปนฺโน อทฺธา.\csromanspacer{} ปจฺจุปฺปนฺนํ อทฺธานํ อารพฺภ ตณฺหา ปหีนา, ทิฏฺฐิ ปฏินิสฺสฏฺฐา, ตณฺหาย ปหีนตฺตา, ทิฏฺฐิยา ปฏินิสฺสฏฺฐตฺตา รตฺโตติ นุปสงฺเขยฺโย, ทุฏฺโฐติ นุปสงฺเขยฺโย, มูโฬฺหติ นุปสงฺเขยฺโย, วินิพทฺโธติ นุปสงฺเขยฺโย, ปรามฏฺโฐติ นุปสงฺเขยฺโย, วิกฺเขปคโตติ นุปสงฺเขยฺโย, อนิฏฺฐงฺคโตติ นุปสงฺเขยฺโย, ถามคโตติ นุปสงฺเขยฺโย.\csromanspacer{} เต อภิสงฺขารา ปหีนา, อภิสงฺขารานํ ปหีนตฺตา คติยา นุปสงฺเขยฺโย “เนรยิโก”ติ วา “ติรจฺฉานโยนิโก”ติ วา “เปตฺติวิสยิโก”ติ วา “มนุสฺโส”ติ \csromanfolio{164}วา “เทโว”ติ วา “รูปี”ติ วา “อรูปี”ติ วา “สญฺญี”ติ วา “อสญฺญี”ติ วา “เนวสญฺญีนาสญฺญี”ติ วา.\csromanspacer{} โส เหตุ นตฺถิ, ปจฺจโย นตฺถิ, การณํ นตฺถิ, เยน สงฺขํ คจฺเฉยฺยาติ เวมชฺเฌ นุปสงฺเขยฺโย.}

\prose{\textbf{ตสฺส นตฺถิ ปุรกฺขตนฺ}ติ \textbf{ตสฺสา}ติ อรหโต ขีณาสวสฺส.\csromanspacer{} \textbf{ปุเรกฺขารา}ติ เทฺว ปุเรกฺขารา ตณฺหาปุเรกฺขาโร จ ทิฏฺฐิปุเรกฺขาโร จ ฯเปฯ อยํ ตณฺหาปุเรกฺขาโร ฯเปฯ อยํ ทิฏฺฐิปุเรกฺขาโร.\csromanspacer{} ตสฺส ตณฺหาปุเรกฺขาโร ปหีโน, ทิฏฺฐิปุเรกฺขาโร ปฏินิสฺสฏฺโฐ, ตณฺหาปุเรกฺขารสฺส ปหีนตฺตา, ทิฏฺฐิปุเรกฺขารสฺส ปฏินิสฺสฏฺฐตฺตา น ตณฺหํ วา ทิฏฺฐิํ วา ปุรโต กตฺวา จรติ, น ตณฺหาธโช, น ตณฺหาเกตุ, น ตณฺหาธิปเตยฺโย, น ทิฏฺฐิธโช, น ทิฏฺฐิเกตุ, น ทิฏฺฐาธิปเตยฺโย, น ตณฺหาย วา ทิฏฺฐิยา วา ปริวาริโต จรติ, เอวมฺปิ ตสฺส นตฺถิ ปุรกฺขตํ.\csromanspacer{}\csromanspacer{}อถ วา * “เอวํรูโป สิยํ อนาคตมทฺธานนฺ”ติ ตตฺถ นนฺทิํ น สมนฺนาเนติ, “เอวํเวทโน สิยํ.\csromanspacer{} เอวํสญฺโญ สิยํ.\csromanspacer{} เอวํสงฺขาโร สิยํ.\csromanspacer{} เอวํวิญฺญาโณ สิยํ อนาคตมทฺธานนฺ”ติ ตตฺถ นนฺทิํ น สมนฺนาเนติ.\csromanspacer{} เอวมฺปิ ตสฺส นตฺถิ ปุรกฺขตํ.\csromanspacer{}\csromanspacer{}อถ วา + “อิติ เม จกฺขุ สิยา อนาคตมทฺธานํ, อิติ รูปา”ติ อปฺปฏิลทฺธสฺส ปฏิลาภาย จิตฺตํ น ปณิทหติ, เจตโส อปฺปณิธานปฺปจฺจยา น ตทภินนฺทติ, น ตทภินนฺทนฺโต เอวมฺปิ ตสฺส นตฺถิ ปุรกฺขตํ. “อิติ เม โสตํ สิยา อนาคตมทฺธานํ, อิติ สทฺทา”ติ. “อิติ เม ฆานํ สิยา อนาคตมทฺธานํ, อิติ คนฺธา”ติ. “อิติ เม ชิวฺหา สิยา อนาคตมทฺธานํ, อิติ รสา”ติ. “อิติ เม กาโย สิยา อนาคตมทฺธานํ, อิติ โผฏฺฐพฺพา”ติ. “อิติ เม มโน สิยา อนาคตมทฺธานํ, อิติ ธมฺมา”ติ อปฺปฏิลทฺธสฺส ปฏิลาภาย จิตฺตํ น ปณิทหติ, เจตโส อปฺปณิธานปฺปจฺจยา น ตทภินนฺทติ, น ตทภินนฺทนฺโต เอวมฺปิ ตสฺส นตฺถิ ปุรกฺขตํ.\csromanspacer{}\csromanspacer{}อถ วา “อิมินาหํ สีเลน วา วเตน วา ตเปน วา พฺรหฺมจริเยน วา เทโว วา ภวิสฺสามิ เทวญฺญตโร”ติ วา อปฺปฏิลทฺธสฺส ปฏิลาภาย จิตฺตํ น ปณิทหติ, เจตโส อปฺปณิธานปฺปจฺจยา น ตทภินนฺทติ, น ตทภินนฺทนฺโต เอวมฺปิ ตสฺส นตฺถิ ปุรกฺขตํ.\csromanspacer{} เตนาห ภควา\csromandash{}}

\vspace{\tipitakaparskip}
\csromancenter{วีตตโณฺห ปุราเภทา, (อิติ ภควา,)}
\prose{ปุพฺพมนฺตมนิสฺสิโต.}

\csromangathameasure{เว มชฺเฌ นุปสงฺเขยฺโย, \\ \textbf{ตสฺส นตฺถิ ปุรกฺขตนฺ}ติ.}
\csromangathagroup{\csromangathastanza{เว มชฺเฌ นุปสงฺเขยฺโย, \\ \textbf{ตสฺส นตฺถิ ปุรกฺขตนฺ}ติ.}}

\csromanheader{2}{10. ปุราเภทสุตฺตนิทฺเทส}
\proseitem{85}{\csromanfolio{165}\csromansymbolfootnote{*}{ขุ 1. 412 ปิฏฺเฐ.}อกฺโกธโน อสนฺตาสี, อวิกตฺถี อกุกฺกุโจ.}

\csromangathameasure{มนฺตภาณี อนุทฺธโต, ส เว วาจายโต มุนิ.}
\csromangathagroup{\csromangathastanza{มนฺตภาณี อนุทฺธโต, ส เว วาจายโต มุนิ.}}

\prose{\textbf{อโกธโน อสนฺตาสี}ติ \textbf{อกฺโกธโน}ติ ยํ หิ โข วุตฺตํ.\csromanspacer{} อปิ จ โกโธ ตาว วตฺตพฺโพ.\csromanspacer{} ทสหากาเรหิ โกโธ ชายติ + “อนตฺถํ เม อจรี”ติ โกโธ ชายติ, “อนตฺถํ เม จรตี”ติ “โกโธ ชายติ, “อนตฺถํ เม จริสฺสตี”ติ โกโธ ชายติ. “ปิยสฺส เม มนาปสฺส อนตฺถํ อจริ.\csromanspacer{} อนตฺถํ จรติ.\csromanspacer{} อนตฺถํ จริสฺสตี”ติ โกโธ ชายติ, “อปฺปิยสฺส เม อมนาปสฺส อตฺถํ อจริ.\csromanspacer{} อตฺถํ จรติ.\csromanspacer{} อตฺถํ จริสฺสตี”ติ โกโธ ชายติ, อฏฺฐาเน วา ปน โกโธ ชายติ.\csromanspacer{} โย เอวรูโป จิตฺตสฺส อาฆาโต ปฏิฆาโต ปฏิฆํ ปฏิวิโรโธ โกโป ปโกโป สมฺปโกโป โทโส ปโทโส สมฺปโทโส จิตฺตสฺส พฺยาปตฺติ มโนปโทโส โกโธ กุชฺฌนา กุชฺฌิตตฺตํ โทโส ทุสฺสนา ทุสฺสิตตฺตํ พฺยาปตฺติ พฺยาปชฺชนา พฺยาปชฺชิตตฺตํ วิโรโธ ปฏิวิโรโธ จณฺฑิกฺกํ อสุโรโป\footnote{อสฺสุโรโป (สี, ก)} อนตฺตมนตา จิตฺตสฺส.\csromanspacer{} อยํ วุจฺจติ โกโธ.}

\prose{อปิ จ โกธสฺส อธิมตฺตปริตฺตตา เวทิตพฺพา, อตฺถิ กญฺจิ\footnote{กิญฺจิ (ก)} กาลํ โกโธ จิตฺตาวิลกรณมตฺโต โหติ, น จ ตาว มุขกุลานวิกุลาโน โหติ.\csromanspacer{} อตฺถิ กญฺจิ กาลํ โกโธ มุขกุลานวิกุลานมตฺโต โหติ, น จ ตาว หนุสญฺโจปโน โหติ.\csromanspacer{} อตฺถิ กญฺจิ กาลํ โกโธ หนุสญฺโจปนมตฺโต โหติ, น จ ตาว ผรุสวาจํ นิจฺฉารโณ\footnote{ผรุสวาจนิจฺฉารโณ (สฺยา) ขุ 8. 68 ปิฏฺเฐปิ.} โหติ.\csromanspacer{} อตฺถิ กญฺจิ กาลํ โกโธ ผรุสวาจํ นิจฺฉารณมตฺโต โหติ, น จ ตาว ทิสาวิทิสานุวิโลกโน โหติ.\csromanspacer{} อตฺถิ กญฺจิ กาลํ โกโธ ทิสาวิทิสานุวิโลกนมตฺโต โหติ, น จ ตาว ทณฺฑสตฺถปรามสโน โหติ.\csromanspacer{} อตฺถิ กญฺจิ กาลํ โกโธ ทณฺฑสตฺถปรามสนมตฺโต โหติ, น จ ตาว ทณฺฑสตฺถ-อพฺภุกฺกิรโณ โหติ.\csromanspacer{} อตฺถิ กญฺจิ กาลํ โกโธ ทณฺฑสตฺถพฺภุกฺกิรมตฺโต โหติ, น จ ตาว ทณฺฑสตฺถ-อภินิปาตโน โหติ.\csromanspacer{} อตฺถิ กญฺจิ กาลํ โกโธ ทณฺฑสตฺถ-อภินิปาตมตฺโต โหติ, น จ ตาว ฉินฺนวิจฺฉินฺนกรโณ โหติ.\csromanspacer{} อตฺถิ กญฺจิ กาลํ โกโธ ฉินฺนวิจฺฉิกรณมตฺโต โหติ, น จ ตาว สมฺภญฺชนปลิภญฺชโน โหติ.\csromanspacer{} อตฺถิ กญฺจิ กาลํ \csromanfolio{166}โกโธ สมฺภญฺชนปลิภญฺชนมตฺโต โหติ, น จ ตาว องฺคมงฺค- อปกฑฺฒโน โหติ, อตฺถิ กญฺจิ กาลํ โกโธ องฺคมงฺค-อปกฑฺฒนมตฺโต โหติ, น จ ตาว ชีวิตาโวโรปโน\footnote{ชีวิตปนาสโน (สฺยา) ขุ 8. 61 ปิฏฺเฐปิ.} โหติ.\csromanspacer{} อตฺถิ กญฺจิ กาลํ โกโธ ชีวิตาโวโรปนมตฺโต โหติ, น จ ตาว สพฺพจาคปริจฺจาคาย สณฺฐิโต โหติ.\csromanspacer{} ยโต โกโธ ปรปุคฺคลํ ฆาเตตฺวา อตฺตานํ ฆาเตติ, เอตฺตาวตา โกโธ ปรมุสฺสทคโต ปรมเวปุลฺลปฺปตฺโต โหติ.\csromanspacer{} ยสฺส โส โกโธ ปหีโน สมุจฺฉินฺโน วูปสนฺโต ปฏิปสฺสทฺโธ อภพฺพุปฺปตฺติโก ญาณคฺคินา ทฑฺโฒ, โส วุจฺจติ อกฺโกธโน โกธสฺส ปหีนตฺตา อกฺโกธโน, โกธวตฺถุสฺส ปริญฺญาตตฺตา อกฺโกธโน, โกธเหตุสฺส อุปจฺฉินฺนตฺตา อกฺโกธโนติ อกฺโกธโน.}

\prose{\textbf{อสนฺตาสี}ติ อิเธกจฺโจ ตาสี โหติ อุตฺตาสี ปริตฺตาสี, โส ตสติ อุตฺตสติ ปริตฺตสติ ภายติ สนฺตาสํ อาปชฺชติ.\csromanspacer{} กุลํ วา น ลภามิ, คณํ วา น ลภามิ, อาวาสํ วา น ลภามิ, ลาภํ วา น ลภามิ, ยสํ วา น ลภามิ, ปสํสํ วา น ลภามิ, สุขํ วา น ลภามิ, จีวรํ วา น ลภามิ, ปิณฺฑปาตํ วา น ลภามิ, เสนาสนํ วา น ลภามิ, คิลานปจฺจยเภสชฺชปริกฺขารํ วา น ลภามิ, คิลานุปฏฺฐากํ วา น ลภามิ, “อปฺปญฺญาโตมฺหี”ติ ตสติ อุตฺตสติ ปริตฺตสติ ภายติ สนฺตาสํ อาปชฺชติ.}

\prose{อิธ ภิกฺขุ อสนฺตาสี โหติ อนุตฺตาสี อปริตฺตาสี, โส น ตสติ น อุตฺตสติ น ปริตฺตสติ น ภายติ น สนฺตาสํ อาปชฺชติ.\csromanspacer{} กุลํ วา น ลภามิ, คณํ วา น ลภามิ, อาวาสํ วา น ลภามิ, ลาภํ วา น ลภามิ, ยสํ วา น ลภามิ, ปสํสํ วา น ลภามิ, สุขํ วา น ลภามิ, จีวรํ วา น ลภามิ, ปิณฺฑปาตํ วา น ลภามิ, เสนาสนํ วา น ลภามิ, คิลานปจฺจยเภสชฺชปริกฺขารํ วา น ลภามิ, คิลานุปฏฺฐากํ วา น ลภามิ, “อปฺปญฺญาโตมฺหี”ติ น ตสติ น อุตฺตสติ น ปริตฺตสติ น ภายติ น สนฺตาสํ อาปชฺชตีติ อกฺโกธโน อสนฺตาสี.}

\prose{\textbf{อวิกตฺถี อกุกฺกุโจ}ติ อิเธกจฺโจ กตฺถี โหติ วิกตฺถี, โส กตฺถติ วิกตฺถติ “อหมสฺมิ สีลสมฺปนฺโน”ติ วา “วตสมฺปนฺโน”ติ วา “สีลพฺพตสมฺปนฺโน”ติ วา ชาติยา วา โคตฺเตน วา โกลปุตฺติเยน วา}

\vspace{\tipitakaparskip}
\csromancenter{ปุราเภทสุตฺตนิทฺเทส วณฺณโปกฺขรตาย วา ธเนน วา อชฺเฌเนน วา กมฺมายตเนน วา สิปฺปายตเนน วา วิชฺชาฏฺฐาเนน วา สุเตน วา ปฏิภาเนน วา อญฺญตรญฺญตเรน วา วตฺถุนา “อุจฺจา กุลา ปพฺพชิโต”ติ วา “มหากุลา ปพฺพชิโต”ติ วา “มหาโภคกุลา ปพฺพชิโต”ติ วา “อุฬารโภคกุลา ปพฺพชิโต”ติ วา “ญาโต ยสสฺสี คหฏฺฐปพฺพชิตานนฺ”ติ วา “ลาภิมฺหิ จีวรปิณฺฑปาตเสนาสนคิลานปจฺจยเภสชฺชปริกฺขารานนฺ”ติ วา, “สุตฺตนฺติโก”ติ วา “วินยธโร”ติ วา “ธมฺมกถิโก”ติ วา, “อารญฺญิโก”ติ วา “ปิณฺฑปาติโก”ติ วา “ปํสุกูลิโก”ติ วา “เตจีวริโก”ติ วา “สปทานจาริโก”ติ วา “ขลุปจฺฉาภตฺติโก”ติ วา “เนสชฺชิโก”ติ วา “ยถาสนฺถติโก”ติ วา, “ปฐมสฺส ฌานสฺส ลาภี”ติ วา “ทุติยสฺส ฌานสฺส ลาภี”ติ วา “ตติยสฺส ฌานสฺส ลาภี”ติ วา “จตุตฺถสฺส ฌานสฺส ลาภี”ติ วา “อากาสานญฺจายตนสมาปตฺติยา.\csromanspacer{} วิญฺญาณญฺจายตนสมาปตฺติยา.\csromanspacer{} อากิญฺจญฺญายตนสมาปตฺติยา.\csromanspacer{} เนวสญฺญานาสญฺญายตนสมาปตฺติยา ลาภี”ติ วา กตฺถติ วิกตฺถติ.\csromanspacer{} เอวํ น กตฺถติ น วิกตฺถติ, กตฺถนา วิกตฺถนา อารโต วิรโต ปฏิวิรโต, นิกฺขนฺโต นิสฺสโฏ วิปฺปมุตฺโต วิสญฺญุตฺโต วิมริยาทิกเตน เจตสา วิหรตีติ อวิกตฺถี.}
\prose{\csromanfolio{167}\textbf{อกุกฺกุจฺโจ}ติ \textbf{กุกฺกุจฺจนฺ}ติ หตฺถกุกฺกุจฺจมฺปิ กุกฺกุจฺจํ ปาทกุกฺกุจฺจมฺปิ กุกฺกุจฺจํ หตฺถปาทกุกฺกุจฺจมฺปิ กุกฺกุจฺจํ, อกปฺปิเย กปฺปิยสญฺญิตา, กปฺปิเย อกปฺปิยสญฺญิตา, วิกาเล กาลสญฺญิตา, กาเล วิกาลสญฺญิตา, อวชฺเช วชฺชสญฺญิตา, วชฺเช อวชฺชสญฺญิตา, ยํ เอวรูปํ กุกฺกุจฺจํ กุกฺกุจฺจายนา กุกฺกุจฺจายิตตฺตํ เจตโส วิปฺปฏิสาโร มโนวิเลโข.\csromanspacer{} อิทํ วุจฺจติ กุกฺกุจฺจํ.}

\prose{อปิ จ ทฺวีหิ การเณหิ อุปฺปชฺชติ กุกฺกุจฺจํ เจตโส วิปฺปฏิสาโร มโนวิเลโข กตตฺตา จ อกตตฺตา จ.\csromanspacer{} กถํ กตตฺตา จ อกตตฺตา จ อุปฺปชฺชติ กุกฺกุจฺจํ เจตโส วิปฺปฏิสาโร มโนวิเลโข\csromandash{}“กตํ เม กายทุจฺจริตํ, อกตํ เม กายสุจริตนฺ”ติ อุปฺปชฺชติ กุกฺกุจฺจํ เจตโส วิปฺปฏิสาโร มโนวิเลโข, “กตํ เม วจีทุจฺจริตํ, อกตํ เม วจีสุจริตํ.\csromanspacer{} กตํ เม มโนทุจฺจริตํ, อกตํ เม มโนสุจริตนฺ”ติ อุปฺปชฺชติ กุกฺกุจฺจํ เจตโส วิปฺปฏิสาโร มโนวิเลโข, “กโต เม ปาณาติปาโต, อกตา เม ปาณาติปาตา เวรมณี”ติ อุปฺปชฺชติ \csromanfolio{168}กุกฺกุจฺจํ เจตโส วิปฺปฏิสาโร มโนวิเลโข, “กตํ เม อทินฺนาทานํ, อกตา เม อทินฺนาทานา เวรมณี”ติ อุปฺปชฺชติ กุกฺกุจฺจํ เจตโส วิปฺปฏิสาโร มโนวิเลโข, “กโต เม กาเมสุมิจฺฉาจาโร, อกตา เม กาเมสุมิจฺฉาจารา เวรมณี”ติ อุปฺปชฺชติ กุกฺกุจฺจํ เจตโส วิปฺปฏิสาโร มโนวิเลโข. “กโต เม มุสาวาโท, อกตา เม มุสาวาทา เวรมณี”ติ. “กตา เม ปิสุณา วาจา, อกตา เม ปิสุณาย วาจาย เวรมณี”ติ. “กตา เม ผรุสา วาจา, อกตา เม ผรุสาย วาจาย เวรมณี”ติ. “กโต เม สมฺผปฺปลาโป, อกตา เม สมฺผปฺปลาปา เวรมณี”ติ. “กตา เม อภิชฺฌา, อกตา เม อนภิชฺฌา”ติ. “กโต เม พฺยาปาโท, อกโต เม อพฺยาปาโท”ติ. “กตา เม มิจฺฉาทิฏฺฐิ, อกตา เม สมฺมาทิฏฺฐี”ติ อุปฺปชฺชติ กุกฺกุจฺจํ เจตโส วิปฺปฏิสาโร มโนวิเลโข.\csromanspacer{} เอวํ กตตฺตา จ อกตตฺตา จ อุปฺปชฺชติ กุกฺกุจฺจํ เจตโส วิปฺปฏิสาโร มโนวิเลโข.}

\prose{อถ วา “สีเลสุ’มฺหิ น ปริปูรการี”ติ อุปฺปชฺชติ กุกฺกุจฺจํ เจตโส วิปฺปฏิสาโร มโนวิเลโข, “อินฺทฺริเยสุ’มฺหิ อคุตฺตทฺวาโร”ติ. “โภชเน อมตฺตญฺญุ’มฺหี”ติ. “ชาคริยํ อนนุยุตฺโต’มฺหี”ติ. “น สติสมฺปชญฺเญน สมนฺนาคโต’มฺหี”ติ. “อภาวิตา เม จตฺตาโร สติปฏฺฐานา”ติ. “อภาวิตา เม จตฺตาโร สมฺมปฺปธานา”ติ. “อภาวิตา เม จตฺตาโร อิทฺธิปาทา”ติ. “อภาวิตานิ เม ปญฺจินฺทฺริยานี”ติ. “อภาวิตานิ เม ปญฺจ พลานี”ติ. “อภาวิตา เม สตฺต โพชฺฌงฺคา”ติ. “อภาวิโต เม อริโย อฏฺฐงฺคิโก มคฺโค”ติ. “ทุกฺขํ เม อปริญฺญาตนฺ”ติ. “สมุทโย เม อปฺปหีโน”ติ. “มคฺโค เม อภาวิโต”ติ. “นิโรโธ เม อสจฺฉิกโต”ติ อุปฺปชฺชติ กุกฺกุจฺจํ เจตโส วิปฺปฏิสาโร มโนวิเลโข.\csromanspacer{} ยสฺเสตํ กุกฺกุจฺจํ ปหีนํ สมุจฺฉินฺนํ วูปสนฺตํ ปฏิปสฺสทฺธํ อภพฺพุปฺปตฺติกํ ญาณคฺคินา ทฑฺฒํ.\csromanspacer{} โส วุจฺจติ อกุกฺกุจฺโจติ อวิกตฺถี อกุกฺกุโจ.}

\prose{\textbf{มนฺตภาณี อนุทฺธโต}ติ \textbf{มนฺตา} วุจฺจติ ปญฺญา.\csromanspacer{} ยา ปญฺญา ปชานนา ฯเปฯ อโมโห ธมฺมวิจโย สมฺมาทิฏฺฐิ.\csromanspacer{} มนฺตาย ปริคฺคเหตฺวา ปริคฺคเหตฺวา วาจํ ภาสติ พหุมฺปิ กเถนฺโต พหุมฺปิ ภณนฺโต พหุมฺปิ ทีปยนฺโต พหุมฺปิ โวหรนฺโต ทุกฺกถิตํ ทุพฺภณิตํ ทุลฺลปิตํ ทุรุตฺตํ ทุพฺภาสิตํ วาจํ น ภาสตีติ มนฺตภาณี.\csromanspacer{}\csromanspacer{}\textbf{อนุทฺธโต}ติ ตตฺถ กตมํ อุทฺธจฺจํ, ยํ จิตฺตสฺส}

\vspace{\tipitakaparskip}
\csromancenter{ปุราเภทสุตฺตนิทฺเทส อุทฺธจฺจํ อวูปสโม เจตโส วิกฺเขโป ภนฺตตฺตํ จิตฺตสฺส, อิทํ วุจฺจติ อุทฺธจฺจํ.\csromanspacer{} ยสฺเสตํ อุทฺธจฺจํ ปหีนํ สมุจฺฉินฺนํ วูปสนฺตํ ปฏิปสฺสทฺธํ อภพฺพุปฺปตฺติกํ ญาณคฺคินา ทฑฺฒํ, โส วุจฺจติ อนุทฺธโตติ มนฺตภาณี อนุทฺธโต.}
\prose{\csromanfolio{169}\textbf{ส เว วาจายโต มุนี}ติ \textbf{*} อิธ ภิกฺขุ มุสาวาทํ ปหาย มุสาวาทา ปฏิวิรโต โหติ สจฺจวาที สจฺจสนฺโธ เถโต ปจฺจยิโก อวิสํวาทโก โลกสฺส, ปิสุณํ วาจํ ปหาย ปิสุณาย วาจาย ปฏิวิรโต โหติ อิโต สุตฺวา น อมุตฺร อกฺขาตา อิเมสํ เภทาย, อมุตฺร วา สุตฺวา น อิเมสํ อกฺขาตา อมูสํ เภทาย, อิติ ภินฺนานํ วา สนฺธาตา, สหิตานํ วา อนุปฺปทาตา, สมคฺคาราโม สมคฺครโต สมคฺคนนฺที สมคฺคกรณิํ วาจํ ภาสิตา โหติ.\csromanspacer{} ผรุสํ วาจํ ปหาย ผรุสาย วาจาย ปฏิวิรโต โหติ, ยา สา วาจา เนลา กณฺณสุขา เปมนียา หทยงฺคมา โปรี พหุชนกนฺตา พหุชนมนาปา, ตถารูปิํ วาจํ ภาสิตา โหติ.\csromanspacer{} สมฺผปฺปลาปํ ปหาย สมฺผปฺปลาปา ปฏิวิรโต โหติ กาลวาที ภูตวาที อตฺถวาที ธมฺมวาที วินยวาที นิธานวติํ วาจํ ภาสิตา โหติ กาเลน สาปเทสํ ปริยนฺตวติํ อตฺถสํหิตํ.\csromanspacer{} จตูหิ วจีสุจริเตหิ สมนฺนาคโต จตุทฺโทสาปคตํ วาจํ ภาสติ, พาตฺติํสาย ติรจฺฉานกถาย อารโต อสฺส วิรโต ปฏิวิรโต นิกฺขนฺโต นิสฺสโฏ วิปฺปมุตฺโต วิสญฺญุตฺโต วิมริยาทิกเตน เจตสา วิหรติ.}

\prose{\csromansymbolfootnote{+}{อํ 3. 359 ปิฏฺเฐ.}ทส กถาวตฺถูนิ กเถสิ.\csromanspacer{} เสยฺยถิทํ, อปฺปิจฺฉกถํ กเถติ, สนฺตุฏฺฐีกถํ กเถติ, ปวิเวกกถํ.\csromanspacer{} อสํสคฺคกถํ.\csromanspacer{} วีริยารมฺภกถํ.\csromanspacer{} สีลกถํ.\csromanspacer{} สมาธิกถํ.\csromanspacer{} ปญฺญากถํ.\csromanspacer{} วิมุตฺติกถํ.\csromanspacer{} วิมุตฺติญาณทสฺสนกถํ.\csromanspacer{} สติปฏฺฐานกถํ.\csromanspacer{} สมฺมปฺปธานกถํ.\csromanspacer{} อิทฺธิปาทกถํ.\csromanspacer{} อินฺทฺริยกถํ.\csromanspacer{} พลกถํ.\csromanspacer{} โพชฺฌงฺคกถํ.\csromanspacer{} มคฺคกถํ.\csromanspacer{} ผลกถํ.\csromanspacer{} นิพฺพานกถํ กเถติ.\csromanspacer{}\csromanspacer{}\textbf{วาจายโต}ติ ยตฺโต ปริยตฺโต คุตฺโต โคปิโต รกฺขิโต วูปสนฺโต.\csromanspacer{} \textbf{มุนี}ติ \textbf{โมนํ} วุจฺจติ ญาณํ.\csromanspacer{} ยา ปญฺญา ปชานนา ฯเปฯ อโมโห ธมฺมวิจโย สมฺมาทิฏฺฐิ ฯเปฯ สงฺคชาลมติจฺจ โส มุนีติ ส เว วาจายโต มุนิ.\csromanspacer{} เตนาห ภควา\csromandash{}}

\csromangathameasure{อกฺโกธโน อสนฺตาสี, อวิกตฺถี อกุกฺกุโจ. \\ มนฺตภาณี อนุทฺธโต, ส เว วาจายโต มุนีติ.}
\csromangathagroup{\csromangathastanza{อกฺโกธโน อสนฺตาสี, อวิกตฺถี อกุกฺกุโจ. \\ มนฺตภาณี อนุทฺธโต, ส เว วาจายโต มุนีติ.}}

\proseitem{86}{\csromanfolio{170}\csromansymbolfootnote{*}{ขุ 1. 412 ปิฏฺเฐ.}\textbf{นิราสตฺติ อนาคเต}, อตีตํ นานุโสจติ.}

\csromangathameasure{วิเวกทสฺสี ผสฺเสสุ, ทิฏฺฐีสุ จ น นียติ.}
\csromangathagroup{\csromangathastanza{วิเวกทสฺสี ผสฺเสสุ, ทิฏฺฐีสุ จ น นียติ.}}

\prose{\textbf{นิราสตฺติ อนาคเต}ติ \textbf{อาสตฺติ} วุจฺจติ ตณฺหา.\csromanspacer{} โย ราโค สาราโค ฯเปฯ อภิชฺฌา โลโภ อกุสลมูลํ.\csromanspacer{} ยสฺเสสา \textbf{อาสตฺติ} ตณฺหา ปหีนา สมุจฺฉินฺนา วูปสนฺตา ปฏิปสฺสทฺธา อภพฺพุปฺปตฺติกา ญาณคฺคินา ทฑฺฒาติ เอวมฺปิ นิรฺ\textbf{อาสตฺติ} อนาคเต.\csromanspacer{}\csromanspacer{}อถ วา + “เอวํรูโป สิยํ อนาคตมทฺธานนฺ”ติ ตตฺถ นนฺทิํ น สมนฺนาเนติ, “เอวํเวทโน สิยํ.\csromanspacer{} เอวํสญฺโญ สิยํ.\csromanspacer{} เอวํสงฺขาโร สิยํ.\csromanspacer{} เอวํวิญฺญาโณ สิยํ อนาคตมทฺธานนฺ”ติ ตตฺถ นนฺทิํ น สมนฺนาเนติ.\csromanspacer{} เอวมฺปิ นิรฺ\textbf{อาสตฺติ} อนาคเต.\csromanspacer{}\csromanspacer{}อถ วา อิติ เม จกฺขุ สิยา อนาคตมทฺธานํ, อิติ รูปา”ติ อปฺปฏิลทฺธสฺส ปฏิลาภาย จิตฺตํ น ปณิทหติ, เจตโส อปฺปณิธานปฺปจฺจยา นตทภินนฺทติ, น ตทภินนฺทนฺโต เอวมฺปิ นิรฺ\textbf{อาสตฺติ} อนาคเต.\csromanspacer{} อิติ เม โสตํ สิยา อนาคตมทฺธานํ, อิติ สทฺทา”ติ ฯเปฯ “อิติ เม มโน สิยา อนาคตมทฺธานํ, อิติ ธมฺมา”ติ อปฺปฏิลทฺธสฺส ปฏิลาภาย จิตฺตํ น ปณิทหติ, เจตโส อปฺปณิธานปฺปจฺจยา น ตทภินนฺทติ, น ตทภินนฺทนฺโต เอวมฺปิ นิรฺ\textbf{อาสตฺติ} อนาคเต.\csromanspacer{}\csromanspacer{}อถ วา “อิมินาหํ สีเลน วา วเตน วา ตเปน วา พฺราหฺมจริเยน วา เทโว วา ภวิสฺสามิ เทวญฺญตโร วา”ติ อปฺปฏิลทฺธสฺส ปฏิลาภาย จิตฺตํ น ปณิทหติ, เจตโส อปฺปณิธานปฺปจฺจยา น ตทภินนฺทติ, น ตทภินนฺทนฺโต เอวมฺปิ นิรฺ\textbf{อาสตฺติ} อนาคเต.}

\prose{\textbf{อตีตํนานุโสจตี}ติ วิปริณตํ วา วตฺถุํ น โสจติ, วิปริณตสฺมิํ วา วตฺถุสฺมิํ น โสจติ, “จกฺขุ เม วิปริณตนฺ”ติ น โสจติ, “โสตํ เม.\csromanspacer{} ฆานํ เม.\csromanspacer{} ชิวฺหา เม.\csromanspacer{} กาโย เม.\csromanspacer{} รูปา เม.\csromanspacer{} สทฺทา เม.\csromanspacer{} คนฺธา เม.\csromanspacer{} รสา เม.\csromanspacer{} โผฏฺฐพฺพา เม.\csromanspacer{} กุลํ เม.\csromanspacer{} คโณ เม.\csromanspacer{} อาวาโส เม.\csromanspacer{} ลาโภ เม.\csromanspacer{} ยโส เม.\csromanspacer{} ปสํสา เม.\csromanspacer{} สุขํ เม.\csromanspacer{} จีวรํ เม.\csromanspacer{} ปิณฺฑปาโต เม.\csromanspacer{} เสนาสนํ เม.\csromanspacer{} คิลานปจฺจยเภสชฺชปริกฺขาโร เม.\csromanspacer{} มาตา เม.\csromanspacer{} ปิตา เม.\csromanspacer{} ภาตา เม.\csromanspacer{} ภคินี เม.\csromanspacer{} ปุตฺโต เม.\csromanspacer{} ธีตา เม.\csromanspacer{} มิตฺตา เม.\csromanspacer{} อมจฺจา เม.\csromanspacer{} ญาตกา เม.\csromanspacer{} สาโลหิตา เม วิปริณตา”ติ น โสจติ น กิลมติ น ปริเทวติ น อุรตฺตาฬิํ กนฺทติ น สมฺโมหํ อาปชฺชตีติ อตีตํ นานุโสจติ.}

\csromanheader{2}{10. ปุราเภทสุตฺตนิทฺเทส}
\prose{\csromanfolio{171}\textbf{วิเวกทสฺสี ผสฺเสสู}ติ จกฺขุสมฺผสฺโส โสตสมฺผสฺโส ฆานสมฺผสฺโส ชิวฺหาสมฺผสฺโส กายสมฺผสฺโส มโนสมฺผสฺโส อธิวจนสมฺผสฺโส ปฏิฆสมฺผสฺโส, สุขเวทนีโย ผสฺโส ทุกฺขเวทนีโย ผสฺโส อทุกฺขมสุขเวทนีโย ผสฺโส, กุสโล ผสฺโส อกุสโล ผสฺโส อพฺยากโต ผสฺโส, กามาวจโร ผสฺโส รูปาวจโร ผสฺโส อรูปาวจโร ผสฺโส, สุญฺญโต ผสฺโส อนิมิตฺโต ผสฺโส อปฺปณิหิโต ผสฺโส, โลกิโย ผสฺโส โลกุตฺตโร ผสฺโส, อตีโต ผสฺโส อนาคโต ผสฺโส ปจฺจุปฺปนฺโน ผสฺโส, โย เอวรูโป ผสฺโส ผุสนา สมฺผุสนา สมฺผุสิตตฺตํ, อยํ วุจฺจติ ผสฺโส.}

\prose{\textbf{วิเวกทสฺสี ผสฺเสสู}ติ จกฺขุสมฺผสํ วิวิตฺตํ ปสฺสติ อตฺเตน วา อตฺตนิเยน วา นิจฺเจน วา ธุเวน วา สสฺสเตน วา อวิปริณามธมฺเมน วา, โสตสมฺผสฺสํ วิวิตฺตํ ปสฺสติ.\csromanspacer{} ฆานสมฺผสฺสํ วิวิตฺตํ ปสฺสติ.\csromanspacer{} ชิวฺหาสมฺผสฺสํ วิวิตฺตํ ปสฺสติ.\csromanspacer{} กายสมฺผสฺสํ วิวิตฺตํ ปสฺสติ.\csromanspacer{} มโนสมฺผสฺสํ วิวิตฺตํ ปสฺสติ.\csromanspacer{} อธิวจนสมฺผสฺสํ วิวิตฺตํ ปสฺสติ.\csromanspacer{} ปฏิฆสมฺผสฺสํ วิวิตฺตํ ปสฺสติ.\csromanspacer{} สุขเวทนียํ ผสฺสํ.\csromanspacer{} ทุกฺขเวทนียํ ผสฺสํ.\csromanspacer{} อทุกฺขมสุขเวทนียํ ผสฺสํ.\csromanspacer{} กุสลํ ผสฺสํ.\csromanspacer{} อกุสลํ ผสฺสํ.\csromanspacer{} อพฺยากตํ ผสฺสํ.\csromanspacer{} กามาวจรํ ผสฺสํ.\csromanspacer{} รูปาวจรํ ผสฺสํ.\csromanspacer{} อรูปาวจรํ ผสฺสํ.\csromanspacer{} โลกิยํ ผสฺสํ วิวิตฺตํ ปสฺสติ อตฺเตน วา อตฺตนิเยน วา นิจฺเจน วา ธุเวน วา สสฺสเตน วา อวิปริณามธมฺเมน วา.}

\prose{อถ วา อตีตํ ผสฺสํ อนาคเตหิ จ ปจฺจุปฺปนฺเนหิ จ ผสฺเสหิ วิวิตฺตํ ปสฺสติ, อนาคตํ ผสฺสํ อตีเตหิ จ ปจฺจุปฺปนฺเนหิ จ ผสฺเสหิ วิวิตฺตํ ปสฺสติ, ปจฺจุปฺปนฺนํ ผสฺสํ อตีเตหิ จ อนาคเตหิ จ ผสฺเสหิ วิวิตฺตํ ปสฺสติ.\csromanspacer{}\csromanspacer{}อถ วา เย เต ผสฺสา อริยา อนาสวา โลกุตฺตรา สุญฺญตปฏิสญฺญุตฺตา, เต ผสฺเส วิวิตฺเต ปสฺสติ ราเคน โทเสน โมเหน, โกเธน อุปนาเหน มกฺเขน ปฬาเสน อิสฺสาย มจฺฉริเยน มายาย สาเฐยฺเยน ถมฺเภน สารมฺเภน มาเนน อติมาเนน มเทน ปมาเทน สพฺพกิเลเสหิ สพฺพทุจฺจริเตหิ สพฺพทรเถหิ สพฺพปริฬาเหหิ สพฺพสนฺตาเปหิ สพฺพากุสลาภิสงฺขาเรหิ วิวิตฺเต ปสฺสตีติ วิเวกทสฺสี ผสฺเสสุ.}

\prose{\textbf{ทิฏฺฐีสุ จ น นียตี}ติ ตสฺส ทฺวาสฏฺฐิ ทิฏฺฐิคตานิ ปหีนานิ สมุจฺฉินฺนานิ วูปสนฺตานิ ปฏิปสฺสทฺธานิ อภพฺพุปฺปตฺติกานิ ญาณคฺคินา ทฑฺฒานิ.\csromanspacer{} โส ทิฏฺฐิยา}

\prose{\csromanfolio{172}น ยายติ น นียติ น วุยฺหติ น สํหรียติ.\csromanspacer{} นปิ ตํ ทิฏฺฐิคตํ สารโต ปจฺเจติ น ปจฺจาคจฺฉตีติ ทิฏฺฐีสุ จ น นียติ.\csromanspacer{} เตนาห ภควา\csromandash{}}

\csromangathameasure{นิราสตฺติ อนาคเต, อตีตํ นานุโสจติ. \\ วิเวกทสฺสี ผสฺเสสุ, ทิฏฺฐีสุ จ น นียตีติ.}
\csromangathagroup{\csromangathastanza{นิราสตฺติ อนาคเต, อตีตํ นานุโสจติ. \\ วิเวกทสฺสี ผสฺเสสุ, ทิฏฺฐีสุ จ น นียตีติ.}}

\proseitem{87}{\csromansymbolfootnote{*}{ขุ 1. 412 ปิฏฺเฐ.}\textbf{ปติลีโน อกุหโก}, อปิหาลุ อมจฺฉรี.}

\csromangathameasure{อปฺปคมฺโภ อเชคุจฺโฉ, เปสุเณยฺย จ โน ยุโต.}
\csromangathagroup{\csromangathastanza{อปฺปคมฺโภ อเชคุจฺโฉ, เปสุเณยฺย จ โน ยุโต.}}

\prose{\textbf{ปติลีโน อกุหโก}ติ \textbf{ปติลีโน}ติ ราคสฺส ปหีนตฺตา \textbf{ปติลีโน}, โทสสฺส ปหีนตฺตา \textbf{ปติลีโน}, โมหสฺส ปหีนตฺตา ปติลิโน, โกธสฺส.\csromanspacer{} อุปนาหสฺส.\csromanspacer{} มกฺขสฺส.\csromanspacer{} ปฬาสสฺส.\csromanspacer{} อิสฺสาย.\csromanspacer{} มจฺฉริยสฺส ฯเปฯ สพฺพากุสลา- ภิสงฺขารานํ ปหีนตฺตา \textbf{ปติลีโน}.\csromanspacer{} วุตฺตํ เหตํ ภควตา\csromandash{} + กถญฺจ ภิกฺขเว ภิกฺขุ \textbf{ปติลีโน} โหติ, อิธ ภิกฺขเว ภิกฺขุโน อสฺมิมาโน ปหีโน โหติ อุจฺฉินฺนมูโล ตาลาวตฺถุกโต อนภาวํกโต อายติํ อนุปฺปาทธมฺโม.\csromanspacer{} เอวํ โข ภิกฺขเว ภิกฺขุ \textbf{ปติลีโน} โหตีติ \textbf{ปติลีโน}.}

\prose{\textbf{อกุหโก}ติ ++ ตีณิ กุหนวตฺถูนิ ปจฺจยปฏิเสวนสงฺขาตํ กุหนวตฺถุ อิริยาปถสงฺขาตํ กุหนวตฺถุ สามนฺตชปฺปนสงฺขาตํ กุหนวตฺถุ.}

\prose{กตมํ ปจฺจยปฏิเสวนสงฺขาตํ กุหนวตฺถุ, อิธ คหปติกา ภิกฺขุํ นิมนฺเตนฺติ จีวรปิณฺฑปาตเสนาสนคิลานปจฺจยเภสชฺช- ปริกฺขาเรหิ, โส ปาปิจฺโฉ อิจฺฉาปกโต อตฺถิโก จีวรปิณฺฑปาตเสนาสน- คิลานปจฺจยเภสชฺชปริกฺขารานํ ภิยฺโยกมฺยตํ อุปาทาย จีวรํ ปจฺจกฺขาติ, ปิณฺฑปาตํ ปจฺจกฺขาติ, เสนาสนํ ปจฺจกฺขาติ, คิลานปจฺจยเภสชฺชปริกฺขารํ ปจฺจกฺขาติ.\csromanspacer{} โส เอวมาห “กิํ สมณสฺส มหคฺเฆน จีวเรน, เอตํ สารุปฺปํ, ยํ สมโณ สุสานา วา สงฺการกูฏา วา ปาปณิกา วา นนฺตกานิ อุจฺจินิตฺวา สงฺฆาฏิํ กตฺวา ธาเรยฺย.\csromanspacer{} กิํ สมณสฺส มหคฺเฆน ปิณฺฑปาเตน, เอตํ สารุปฺปํ, ยํ สมโณ อุญฺฉาจริยาย ปิณฺฑิยาโลเปน ชีวิกํ กปฺเปยฺย.\csromanspacer{} กิํ สมณสฺส มหคฺเฆน เสนาสเนน, เอตํ สารุปฺปํ, ยํ สมโณ รุกฺขมูลิโก วา อสฺส โสสานิโก วา อพฺโภกาสิโก วา.\csromanspacer{} กิํ สมณสฺส มหคฺเฆน คิลานปจฺจยเภสฺสชฺชปริกฺขาเรน, เอตํ สารุปฺปํ, ยํ สมโณ ปูติมุตฺเตน วา หริตกีขณฺเฑน วา โอสธํ กเรยฺยา”ติ.\csromanspacer{} ตทุปาทาย ลูขํ จีวรํ ธาเรติ, ลูขํ ปิณฺฑปาตํ ปริภุญฺชติ,}

\vspace{\tipitakaparskip}
\csromancenter{ปุราเภทสุตฺตนิทฺเทส ลูขํ เสนาสนํ ปฏิเสวติ, ลูขํ คิลานปจฺจยเภสชฺชปริกฺขารํ ปฏิเสวติ, ตเมนํ คหปติกา เอวํ ชานนฺติ “อยํ สมโณ อปฺปิจฺโฉ สนฺตุฏฺโฐ ปวิวิตฺโต อสํสฏฺโฐ อารทฺธวีริโย ธุตวาโท”ติ ภิยฺโย ภิยฺโย นิมนฺเตนฺติ จีวรปิณฺฑปาตเสนาสนคิลานคิลานปจฺจยเภสชฺช- ปริกฺขาเรหิ.\csromanspacer{} โส เอวมาห “ติณฺณํ สมฺมุขีภาวา สทฺโธ กุลปุตฺโต พหุํ ปุญฺญํ ปสวติ.\csromanspacer{} สทฺธาย สมฺมุขีภาวา สทฺโธ กุลปุตฺโต พหุํ ปุญฺญํ ปสวติ, เทยฺยธมฺมสฺส สมฺมุขีภาวา สทฺโธ กุลปุตฺโต พหุํ ปุญฺญํ ปสวติ, ทกฺขิเณยฺยานํ สมฺมุขีภาวา สทฺโธ กุลปุตฺโต พหุํ ปสวติ.\csromanspacer{} ตุมฺหากญฺเจวายํ สทฺธา อตฺถิ, เทยฺยธมฺโม จ สํวิชฺชติ, อหญฺจ ปฏิคฺคาหโก, สเจหํ น ปฏิคฺคเหสฺสามิ, เอวํ ตุเมฺห ปุญฺเญน ปริพาหิรา ภวิสฺสถ, น มยฺหํ อิมินา อตฺโถ, อปิ จ ตุมฺหากํเยว อนุกมฺปาย ปฏิคฺคณฺหามี”ติ ตทุปาทาย พหุมฺปิ จีวรํ ปฏิคฺคณฺหาติ, พหุมฺปิ ปิณฺฑปาตํ ปฏิคฺคณฺหาติ, พหุมฺปิ เสนาสนํ ปฏิคฺคณฺหาติ, พหุมฺปิ คิลานปจฺจยเภสชฺชปริกฺขารํ ปฏิคฺคณฺหาติ.\csromanspacer{} ยา เอวรูปา ภากุฏิกา ภากุฏิยํ กุหนา กุหายนา กุหิตตฺตํ, อิทํ ปจฺจยปฏิเสวนสงฺขาตํ กุหนวตฺถุ.}
\prose{\csromanfolio{173}กตมํ อิริยาปถสงฺขาตํ กุหนวตฺถุ, อิเธกจฺโจ ปาปิจฺโฉ อิจฺฉาปกโต สมฺภาวนาธิปฺปาโย “เอวํ มํ ชโน สมฺภาเวสฺสตี”ติ คมนํ สณฺฐเปติ, ฐานํ สณฺฐเปติ, นิสชฺชํ สณฺฐเปติ, สยนํ สณฺฐเปติ, ปณิธาย คจฺฉติ, ปณิธาย ติฏฺฐติ, ปณิธาย นิสีทติ, ปณิธาย เสยฺยํ กปฺเปติ, สมาหิโต วิย คจฺฉติ, สมาหิโต วิย ติฏฺฐติ, สมาหิโต วิย นิสีทติ, สมาหิโต วิย เสยฺยํ กปฺเปติ, อาปาถกชฺฌายีว โหติ.\csromanspacer{} ยา เอวรูปา อิริยาปถสฺส อาฐปนา\footnote{อฏฺฐปนา (สี)} ฐปนา สณฺฐปนา ภากุฏิกา ภากุฏิยํ กุหนา กุหายนา กุหิตตฺตํ, อิทํ อิริยาปถสงฺขาตํ กุหนวตฺถุ.}

\prose{กตมํ สามนฺตชปฺปนสงฺขาตํ กุหนวตฺถุ, อิเธกจฺโจ ปาปิจฺโฉ อิจฺฉาปกโต สมฺภาวนาธิปฺปาโย “เอวํ มํ ชาโน สมฺโภเวสฺสตี”ติ อริยธมฺมสนฺนิสฺสิตํ วาจํ ภาสติ, “โย เอวรูปํ จีวรํ ธาเรติ, โส สมโณ มเหสกฺโข”ติ ภณติ. “โย เอวรูปํ ปตฺตํ ธาเรติ.\csromanspacer{} โลหถาลกํ ธาเรติ.\csromanspacer{} ธมฺมกรณํ ธาเรติ.\csromanspacer{} ปริสาวนํ ธาเรติ.\csromanspacer{} กุญฺจิกํ ธาเรติ.\csromanspacer{} อุปาหนํ ธาเรติ.\csromanspacer{} กายพนฺธนํ ธาเรติ.\csromanspacer{} อาโยคํ ธาเรติ, โส \csromanfolio{174}สมโณ มเหสกฺโข”ติ ภณติ. “ยสฺส เอวรูโป อุปชฺโฌโย, โส สมโณ มเหสกฺโข”ติ ภณติ. “ยสฺส เอวรูโป อาจริโย.\csromanspacer{} เอวรูปา สมานุปชฺฌายกา.\csromanspacer{} สมานาจริยกา.\csromanspacer{} มิตฺตา.\csromanspacer{} สนฺทิฏฺฐา.\csromanspacer{} สมฺภตฺตา.\csromanspacer{} สหายา, โส สมโณ มเหสกฺโข”ติ ภณติ. “โย เอวรูเป วิหาเร วสติ, โส สมโณ มเหสกฺโข”ติ ภณติ. “โย เอวรูเป อฑฺฒโยเค วสติ.\csromanspacer{} ปาสาเท วสติ.\csromanspacer{} หมฺมิเย วสติ.\csromanspacer{} คุหายํ วสติ.\csromanspacer{} เลเณ วสติ.\csromanspacer{} กุฏิยา วสติ.\csromanspacer{} กูฏาคาเร วสติ.\csromanspacer{} อฏฺเฏ วสติ.\csromanspacer{} มาเฬ วสติ.\csromanspacer{} อุทฺทณฺเฑ วสติ.\csromanspacer{} อุปฏฺฐานสาลายํ วสติ.\csromanspacer{} มณฺฑเป วสติ.\csromanspacer{} รูกฺขมูเล วสติ, โส สมโณ มเหสกฺโข”ติ ภณติ.}

\prose{อถ วา โกรชิกโกรชิโก\footnote{โกรญฺชิกโกรญฺชิโก (สี) ขุ 8. 274 ปิฏฺเฐ.} ภากุฏิกภากุฏิโก กุหกกุหโก ลปกลปโก มุขสมฺภาวิโก “อยํ สมโณ อิมาสํ เอวรูปานํ สนฺตานํ วิหารสมาปตฺตีนํ ลาภี”ติ ตาทิสํ คมฺภีรํ คูฬฺหํ นิปุณํ ปฏิจฺฉนฺนํ โลกุตฺตรํ สุญฺญตาปฏิสํยุตฺตํ กถํ กเถสิ.\csromanspacer{} ยา เอวรูปา ภากุฏิกา ภากุฏิยํ กุหนา กุหายนา กุหิตตฺตํ, อิทํ สามนฺตชปฺปนสงฺขาตํ กุหนวตฺถุ.\csromanspacer{} ยสฺสิมานิ ตีณิ กุหนวตฺถูนิ ปหีนานิ สมุจฺฉินฺนานิ วูปสนฺตานิ ปฏิปสฺสทฺธานิ อภพฺพุปฺปตฺติกานิ ญาณคฺคินา ทฑฺฒานิ, โส วุจฺจติ อกุหโกติ ปติลีโน อกุหโก.}

\prose{\textbf{อปิหาลุ อมจฺฉรี}ติ \textbf{ปิหา} วุจฺจติ ตณฺหา.\csromanspacer{} โย ราโค สาราโค ฯเปฯ อภิชฺฌา โลโภ อกุสลมูลํ.\csromanspacer{} ยสฺเสสา \textbf{ปิหา} ตณฺหา ปหีนา สมุจฺฉินฺนา วูปสนฺตา ปฏิปสฺสทฺธา อภพฺพุปฺปตฺติกา ญาณคฺคินา ทฑฺฒา, โส วุจฺจติ อ\textbf{ปิหา}ลุ.\csromanspacer{} โส รูเป น ปิเหติ.\csromanspacer{} สทฺเท.\csromanspacer{} คนฺเธ.\csromanspacer{} รเส.\csromanspacer{} โผฏฺฐพฺเพ.\csromanspacer{} กุลํ.\csromanspacer{} คณํ.\csromanspacer{} อาวาสํ.\csromanspacer{} ลาภํ.\csromanspacer{} ยสํ.\csromanspacer{} ปสํสํ.\csromanspacer{} สุขํ.\csromanspacer{} จีวรํ.\csromanspacer{} ปิณฺฑปาตํ.\csromanspacer{} เสนาสนํ.\csromanspacer{} คิลานปจฺจยเภสชฺชปริกฺขารํ.\csromanspacer{} กามธาตุํ.\csromanspacer{} รูปธาตุํ.\csromanspacer{} อรูปธาตุํ.\csromanspacer{} กามภวํ.\csromanspacer{} รูปภวํ.\csromanspacer{} อรูปภวํ.\csromanspacer{} สญฺญาภวํ.\csromanspacer{} อสญฺญาภวํ.\csromanspacer{} เนวสญฺญานาสญฺญาภวํ.\csromanspacer{} เอกโวการภวํ.\csromanspacer{} จตุโวการภวํ.\csromanspacer{} ปญฺจโวการภวํ.\csromanspacer{} อตีตํ.\csromanspacer{} อนาคตํ.\csromanspacer{} ปจฺจุปฺปนฺนํ.\csromanspacer{} ทิฏฺฐสุตมุตวิญฺญาตพฺเพ ธมฺเม น ปิเหติ น อิจฺฉติ น สาทิยติ น ปตฺเถติ นาภิชปฺปตีติ อ\textbf{ปิหา}ลุ.\csromanspacer{}\csromanspacer{}\textbf{อมจฺฉรี}ติ ปญฺจ มจฺฉริยานิ อาวาสมจฺฉริยํ กุลมจฺฉริยํ ลาภมจฺฉริยํ วณฺณมจฺฉริยํ ธมฺมมจฺฉริยํ.\csromanspacer{} ยํ เอวรูปํ มจฺเฉรํ มจฺฉรายนา มจฺฉรายิตตฺตํ เววิจฺฉํ กทริยํ}

\vspace{\tipitakaparskip}
\csromancenter{ปุราเภทสุตฺตนิทฺเทส กฏุกญฺจุกตา อคฺคหิตตฺตํ จิตฺตสฺส, อิทํ วุจฺจติ มจฺฉริยํ.\csromanspacer{} อปิ จ ขนฺธมจฺฉริยมฺปิ มจฺฉริยํ, ธาตุมจฺฉริยมฺปิ มจฺฉริยํ, อายตนมจฺฉริยมฺปิ มจฺฉริยํ คาโห, อิทํ วุจฺจติ มจฺฉริยํ.\csromanspacer{} ยสฺเสตํ มจฺฉริยํ ปหีนํ สมุจฺฉินฺนํ วูปสนฺตํ ปฏิปสฺสทฺธํ อภพฺพุปฺปตฺติกํ ญาณคฺคินา ทฑฺฒํ, โส วุจฺจติ อมจฺฉรีติ อปิหาลุ อมจฺฉรี.}
\prose{\csromanfolio{175}\textbf{อปฺปคพฺโภ อเชคุจฺโฉ}ติ \textbf{ปาคพฺภิยนฺ}ติ ตีณิ ปาคพฺภิยานิ กายิกํ ปาคพฺภิยํ วาจสิกํ ปาคพฺภิยํ เจตสิกํ ปาคพฺภิยํ.\csromanspacer{} กตมํ กายิกํ ปาคพฺภิยํ.\csromanspacer{} อิเธกจฺโจ สํฆคโตปิ กายิกํ ปาคพฺภิยํ ทสฺเสติ, คณคโตปิ กายิกํ ปาคพฺภิยํ ทสฺเสติ, โภชนสาลายมฺปิ กายิกํ ปาคพฺภิยํ ทสฺเสติ, ชนฺตาฆเรปิ กายิกํ ปาคพฺภิยํ ทสฺเสติ, อุทกติตฺเตปิ กายิกํ ปาคพฺภิยํ ทสฺเสติ, อนฺตรฆรํ ปวิสนฺโตปิ กายิกํ ปาคพฺภิยํ ทสฺเสติ, อนฺตรฆรํ ปวิฏฺโฐปิ กายิกํ ปาคพฺภิยํ ทสฺเสติ.}

\prose{กถํ สํฆคโต กายิกํ ปาคพฺภิยํ ทสฺเสติ, อิเธกจฺโจ สํฆคโต อจิตฺตีการกโต\footnote{อจิตฺติการกโต (สฺยา, ก)} เถเร ภิกฺขู ฆฏฺฏยนฺโตปิ ติฏฺฐติ, ฆฏฺฏยนฺโตปิ นิสีทติ, ปุรโตปิ ติฏฺฐติ, ปุรโตปิ นิสีทติ, อุจฺเจปิ อาสเน นิสีทติ, สสีสํ ปารุปิตฺวาปิ นิสีทติ, ฐิตโกปิ ภณติ, พาหาวิกฺเขปโกปิ ภณติ.\csromanspacer{} เอวํ สํฆคโต กายิกํ ปาคพฺภิยํ ทสฺเสติ.}

\prose{กถํ คณคโต กายิกํ ปาคพฺภิยํ ทสฺเสติ, อิเธกจฺโจ คณคโต อจิตฺตีการกโต เถรานํ ภิกฺขูนํ อนุปาหนานํ จงฺกมนฺตานํ ส- อุปาหโน จงฺกมติ, นีเจ จงฺกเม จงฺกมนฺตานํ อุจฺเจ จงฺกเม จงฺกมติ, ฉมาย จงฺกมนฺตานํ จงฺกเม จงฺกมติ, ฆฏฺฏยนฺโตปิ ติฏฺฐติ, ฆฏฺฏยนฺโตปิ นิสีทติ, ปุรโตปิ ติฏฺฐติ, ปุรโตปิ นิสีทติ, อุจฺเจปิ อาสเน นิสีทติ, สสีสํ ปารุปิตฺวา นิสีทติ, ฐิตโกปิ ภณติ, พาหาวิกฺเขปโกปิ ภณติ, เอวํ คณคโต กายิกํ ปาคพฺภิยํ ทสฺเสติ.}

\prose{กถํ โภชนสาลายํ กายิกํ ปาคพฺภิยํ ทสฺเสติ, อิเธกจฺโจ โภชนสาลายํ อจิตฺตีการกโต เถเร ภิกฺขู อนุปขชฺช นิสีทติ, นเวปิ ภิกฺขู อาสเนน ปฏิพาหติ, ฆฏฺฏยนฺโตปิ ติฏฺฐติ, ฆฏฺฏยนฺโตปิ นิสีทติ, ปุรโตปิ ติฏฺฐติ, ปุรโตปิ นิสีทติ, อุจฺเจปิ อาสเน นิสีทติ, สสีสํ ปารุปิตฺวาปิ นิสีทติ, ฐิตโกปิ ภณติ, \csromanfolio{176}พาหาวิกฺเขปโกปิ ภณติ.\csromanspacer{} เอวํ โภชนสาลายํ กายิตํ ปาคพฺภิยํ ทสฺเสติ.}

\prose{กถํ ชนฺตาฆเร กายิกํ ปาคพฺภิยํ ทสฺเสติ, อิเธกจฺโจ ชนฺตาฆเร อจิตฺตีการกโต เถเร ภิกฺขู ฆฏฺฏยนฺโตปิ ติฏฺฐติ, ฆฏฺฏยนฺโตปิ นิสีทติ, ปุรโตปิ ติฏฺฐติ, ปุรโตปิ นิสีทติ, อุจฺเจปิ อาสเน นิสีทติ, อนาปุจฺฉมฺปิ อนชฺฌิฏฺโฐปิ กฏฺฐํ ปกฺขิปติ, ทฺวารมฺปิ ปิทหติ, พาหาวิกฺเขปโกปิ ภณติ.\csromanspacer{} เอวํ ชนฺตาฆเร กายิกํ ปาคพฺภิยํ ทสฺเสติ.}

\prose{กถํ อุทกติตฺเถ กายิกํ ปาคพฺภิยํ ทสฺเสติ, อิเธกจฺโจ อุทกติตฺเถ อจิตฺตีการกโต เถเร ภิกฺขู ฆฏฺฏยนฺโตปิ โอตรติ, ปุรโตปิ โอตรติ, ฆฏฺฏยนฺโตปิ นฺหายติ\footnote{นหายติ (สี)}, ปุรโตปิ นฺหายติ, อุปริโตปิ นฺหายติ, ฆฏฺฏยนฺโตปิ อุตฺตรติ, ปุรโตปิ อุตฺตรติ, อุปริโตปิ อุตฺตรติ.\csromanspacer{} เอวํ อุทกติตฺเถ กายิกํ ปาคพฺภิยํ ทสฺเสติ.}

\prose{กถํ อนฺตรฆรํ ปวิสนฺโต กายิกํ ปาคพฺภิยํ ทสฺเสติ, อิเธกจฺโจ อนฺตรฆรํ ปวิสนฺโต อจิตฺตีการกโต เถเร ภิกฺขู ฆฏฺฏยนฺโตปิ คจฺฉติ, ปุรโตปิ คจฺฉติ, โวกฺกมฺมาปิ เถรานํ ภิกฺขูนํ ปุรโต ปุรโต คจฺฉติ.\csromanspacer{} เอวํ อนฺตรฆรํ ปวิสนฺโต กายิกํ ปาคพฺภิยํ ทสฺเสติ.}

\prose{กถํ อนฺตรฆรํ ปวิฏฺโฐ กายิกํ ปาคพฺภิยํ ทสฺเสติ, อิเธกจฺโจ อนฺตรฆรํ ปวิฏฺโฐ “น ปวิส\footnote{ปวิสถ (สี) เอวมญฺเญสุ ปททฺวเยสุปิ.} ภนฺเต”ติ วุจฺจมาโน ปวิสติ, “น ติฏฺฐ ภนฺเต”ติ วุจฺจมาโน ติฏฺฐติ, “น นิสีท ภนฺเต”ติ วุจฺจมาโน นิสีทติ, อโนกาสมฺปิ ปวิสติ, อโนกาเสปิ ติฏฺฐติ, อโนกาเสปิ นิสีทติ, ยานิปิ ตานิ โหนฺติ กุลานํ โอวรกานิ คูฬฺหานิ จ ปฏิจฺฉนฺนานิ จ, ยตฺถ กุลิตฺถิโย กุลธีตโร กุลสุณฺหโย กุลกุมาริโย นิสีทนฺติ, ตตฺถปิ สหสา ปวิสติ, กุมารกสฺสปิ สิรํ ปรามสติ.\csromanspacer{} เอวํ อนฺตรฆรํ ปวิฏฺโฐ กายิกํ ปาคพฺภิยํ ทสฺเสติ.\csromanspacer{} อิทํ กายิกํ ปาคพฺภิยํ ทสฺเสติ.}

\prose{กตมํ วาจสิกํ ปาคพฺภิยํ ทสฺเสติ, อิเธกจฺโจ สํฆคโตปิ วาจสิกํ ปาคพฺภิยํ ทสฺเสติ, คณคโถปิ วาจสิกํ ปาคพฺภิยํ ทสฺเสติ, อนฺตรฆรํ ปวิฏฺโฐปิ วาจสิกํ ปาคพฺภิยํ ทสฺเสติ.}

\csromanheader{2}{10. ปุราเภทสุตฺตนิทฺเทส}
\prose{\csromanfolio{177}กถํ สํฆคโต วาจสิกํ ปาคพฺภิยํ ทสฺเสติ, อิเธกจฺโจ สํฆคโต อจิตฺตีการกโต เถเร ภิกฺขู อนาปุจฺฉํ วา อนชฺฌิฏฺโฐ วา อารามคตานํ ภิกฺขูนํ ธมฺมํ ภณติ, ปญฺหํ วิสชฺเชติ, ปาติโมกฺขํ อุทฺทิสติ, ฐิตโกปิ ภณติ, พาหาวิกฺเขปโกปิ ภณติ.\csromanspacer{} เอวํ สํฆคโต วาจสิกํ ปาคพฺภิยํ ทสฺเสติ.}

\prose{กถํ คณคโต วาจสิกํ ปาคพฺภิยํ ทสฺเสติ, อิเธกจฺโจ คณคโต อจิตฺตีการกโต เถเร ภิกฺขู อนาปุจฺฉํ วา อนชฺฌิฏฺโฐ วา อารามคตานํ ภิกฺขูนํ ธมฺมํ ภณติ, ปญฺหํ วิสชฺเชติ, ฐิตโกปิ ภณติ, พาหาวิกฺเขปโกปิ.\csromanspacer{} อารามคตานํ ภิกฺขุนีนํ อุปาสกานํ อุปาสิกานํ ธมฺมํ ภณติ, ปญฺหํ วิสชฺเชติ, ฐิตโกปิ ภณติ, พาหาวิกฺเขปโกปิ ภณติ.\csromanspacer{} เอวํ คณคโต วาจสิกํ ปาคพฺภิยํ ทสฺเสติ.}

\prose{กถํ อนฺตรฆรํ ปวิฏฺโฐ วาจสิกํ ปาคพฺภิยํ ทสฺเสติ, อิเธกจฺโจ อนฺตรฆรํ ปวิฏฺโฐ อิตฺถิํ วา กุมาริํ วา เอวมาห “อิตฺถํนาเม อิตฺถํโคตฺเต กิํ อตฺถิ, ยาคุ อตฺถิ, ภตฺตํ อตฺถิ, ขาทนียํ อตฺถิ.\csromanspacer{} กิํ ปิวิสฺสาม, กิํ ภุญฺชิสฺสาม, กิํ ขาทิสฺสาม.\csromanspacer{} กิํ วา อตฺถิ, กิํ วา เม ทสฺสาถา”ติ วิปฺปลปติ.\csromanspacer{} ยา เอวรูปา วาจา ปลาโป วิปฺปลาโป ลาลปฺโป ลาลปฺปนา ลาลปฺปิตตฺตํ.\csromanspacer{} เอวํ อนฺตรฆรํ ปวิฏฺโฐ วาจสิกํ ปาคพฺภิยํ ทสฺเสติ.\csromanspacer{} อิทํ วาจสิกํ ปาคพฺภิยํ.}

\prose{กตมํ เจตสิกํ ปาคพฺภิยํ, อิเธกจฺโจ น อุจฺจา กุลา ปพฺพชิโต สมาโน อุจฺจา กุลา ปพฺพชิเตน สทฺธิํ สทิสํ อตฺตานํ ทหติ จิตฺเตน, น มหากุลา ปพฺพชิโต สมาโน มหากุลา ปพฺพชิเตน สทฺธิํ สทิสํ อตฺตานํ ทหติ จิตฺเตน, น มหาโภคกุลา ปพฺพชิโต สมาโน มหาโภคกุลา ปพฺพชิเตน สทฺธิํ สทิสํ อตฺตานํ ทหติ จิตฺเตน, น อุฬารโภคกุลา ปพฺพชิโต สมาโน.\csromanspacer{} น สุตฺตนฺติโก สมาโน สุตฺตนฺติเกน สทฺธิํ สทิสํ อตฺตานํ ทหติ จิตฺเตน.\csromanspacer{} น วินยธโร สมาโน.\csromanspacer{} น ธมฺมกถิโก สมาโน.\csromanspacer{} น อารญฺญิโก สมาโน.\csromanspacer{} น ปิณฺฑปาติโก สมาโน.\csromanspacer{} น ปํสุกูลิโก สมาโน.\csromanspacer{} น เตจีวริโก สมาโน.\csromanspacer{} น สปทานจาริโก สมาโน.\csromanspacer{} น ขลุปจฺฉาภตฺติโก สมาโน.\csromanspacer{} น เนสชฺชิโก สมาโน.\csromanspacer{} น ยถาสนฺถติโก สมาโน.\csromanspacer{} น ปฐมสฺส ฌานสฺส ลาภี สมาโน ปฐมสฺส ฌานสฺส ลาภินา สทฺธิํ สทิสํ อตฺตานํ ทหติ จิตฺเตน ฯเปฯ.\csromanspacer{} น เนวสญฺญานาสญฺญายตนสมาปตฺติยา}

\prose{\csromanfolio{178}ลาภี สมาโน เนวสญฺญานาสญฺญายตนสมาปตฺติยา ลาภินา สทฺธิํ สทิสํ อตฺตานํ ทหติ จิตฺเตน.\csromanspacer{} อิทํ เจตสิกํ ปาคพฺภิยํ.\csromanspacer{} ยสฺสิมานิ ตีณิ ปาคพฺภิยานิ ปหีนานิ สมุจฺฉินฺนานิ วูปสนฺตานิ ปฏิปสฺสทฺธานิ อภพฺพุปฺปตฺติกานิ ญาณคฺคินา ทฑฺฒานิ.\csromanspacer{} โส วุจฺจติ อปฺปคพฺโภติ อปฺปคพฺโภ.}

\prose{\textbf{อเชคุจฺโฉ}ติ อตฺถิ ปุคฺคโล เชคุจฺโฉ, อตฺถิ อเชคุจฺโฉ.\csromanspacer{} กตโม จ ปุคฺคโล เชคุจฺโฉ, อิเธกจฺโจ ปุคฺคโล ทุสฺสีโล โหติ ปาปธมฺโม อสุจิสงฺกสฺสรสมาจาโร ปฏิจฺฉนฺนกมฺมนฺโต อสฺสมโณ สมณปฏิญฺโญ อพฺรหฺมจารี พฺรหฺมจาริปฏิญฺโญ อนฺโตปูติ อวสฺสุโต กสมฺพุชาโต, อยํ วุจฺจติ ปุคฺคโล เชคุจฺโฉ.\csromanspacer{}\csromanspacer{}อถ วา โกธโน โหติ อุปายาสพหุโล, อปฺปมฺปิ วุตฺโต สมาโน อภิสชฺชติ กุปฺปติ พฺยาปชฺชติ ปติตฺถียติ, โกปญฺจ โทสญฺจ อปฺปจฺจยญฺจ ปาตุกโรติ, อยํ วุจฺจติ ปุคฺคโล เชคุจฺโฉ.\csromanspacer{}\csromanspacer{}อถ วา โกธโน โหติ อุปนาหี, มกฺขี โหติ ปฬาสี, อิสฺสุกี โหติ มจฺฉรี, สโฐ โหติ มายาวี, ถทฺโธ โหติ อติมานี, ปาปิจฺโฉ โหติ มิจฺฉาทิฏฺฐิ\footnote{มิจฺฉาทิฏฺฐี (สี)}, สนฺทิฏฺฐิปรามาสี โหติ อาธานคฺคาหี ทุปฺปฏินิสฺสคฺคี, อยํ วุจฺจติ ปุคฺคโล เชคุจฺโฉ.}

\prose{กตโม จ ปุคฺคโล อเชคุจฺโฉ, อิธ ภิกฺขุ สีลวา โหติ, ปาติโมกฺขสํวรสํวุโต วิหรติ อาจารโคจรสมฺปนฺโน อณุมตฺเตสุ วชฺเชสุ ภยทสฺสาวี, สมาทาย สิกฺขติ สิกฺขาปเทสุ, อยํ วุจฺจติ ปุคฺคโล อเชคุจฺโฉ.\csromanspacer{}\csromanspacer{}อถ วา อกฺโกธโน โหติ อนุปายาสพหุโล, พหุมฺปิ วุตฺโต สมาโน น อภิสชฺชติ น กุปฺปติ น พฺยาปชฺชติ น ปติตฺถียติ, น โกปญฺจ โทสญฺจ อปฺปจฺจยญฺจ ปาตุกโรติ, อยํ วุจฺจติ ปุคฺคโล อเชคุจฺโฉ.\csromanspacer{}\csromanspacer{}อถ วา อกฺโกธโน โหติ อนุปนาหี, อมกฺขี โหติ อปฬาสี, อนิสฺสุกี โหติ อมจฺฉรี, อสโฐ โหติ อมายาวี, อถทฺโธ โหติ อนติมานี, น ปาปิจฺโฉ โหติ น มิจฺฉาทิฏฺฐิ, อสนฺทิฏฺฐิปรามาสี โหติ อนาธานคฺคาหี สุปฺปฏินิสฺสคฺคี, อยํ วุจฺจติ ปุคฺคโล อเชคุจฺโฉ.\csromanspacer{} สพฺเพ พาลปุถุชฺชนา เชคุจฺฉา, ปุถุชฺชนกลฺยณกํ อุปาทาย อฏฺฐ อริยปุคฺคลา อเชคุจฺฉาติ อปฺปคพฺโภ อเชคุจฺโฉ.}

\prose{\textbf{เปสุเณยฺเย จ โน ยุโต}ติ \textbf{เปสุญฺญนฺ}ติ อิเธกจฺโจ ปิสุณวาโจ โหติ อิโต สุตฺวา อมุตฺร อกฺขาตา อิเมสํ เภทาย, อมุตฺร วา}

\vspace{\tipitakaparskip}
\csromancenter{ปุราเภทสุตฺตนิทฺเทส สุตฺวา อิเมสํ อกฺขาตา อมูสํ เภทาย, อิติ สมคฺคานํ วา เภตฺตา\footnote{เภโท (ก)} ภินฺนานํ วา อนุปฺปทาตา, วคฺคาราโม วคฺครโต วคฺคนนฺที วคฺคกรณิํ วาจํ ภาสิตา โหติ.\csromanspacer{} อิทํ วุจฺจติ เปสุญฺญํ.}
\prose{\csromanfolio{179}อปิ จ ทฺวีหิ การเณหิ เปสุญฺญํ อุปสํหรติ ปิยกมฺยตาย วา เภทาธิปฺปาเยน\footnote{เภทาธิปฺปาโย (พหูสุ)} วา.\csromanspacer{} กถํ ปิยกมฺยตาย เปสุญฺญํ อุปสํหรติ, “อิมสฺส ปิโย ภวิสฺสามิ, มนาโป ภวิสฺสามิ, วิสฺสาสิโก ภวิสฺสามิ, อพฺภนฺตริโก ภวิสฺสามิ, สุหทโย ภวิสฺสามี”ติ เอวํ ปิยกมฺยตาย เปสุญฺญํ อุปสํหรติ.\csromanspacer{}\csromanspacer{}กถํ เภทาธิปฺปาเยน เปสุญฺญํ อุปสํหรติ, ‘กถํ อิเม นานา อสฺสุ, วินา อสฺสุ, วคฺคา อสฺสุ, เทฺวธา อสฺสุ, เทฺวชฺฌา อสฺสุ, เทฺว ปกฺขา อสฺสุ, ภิชฺเชยฺยุํ น สมาคจฺเฉยฺยุํ, ทุกฺขํ น ผาสุ\footnote{อผาสุํ (สี)} วิหเรยฺยุนฺ”ติ เอวํ เภทาธิปฺปาเยน เปสุญฺญํ อุปสํหรติ.\csromanspacer{} ยสฺเสตํ เปสุญฺญํ ปหีนํ สมุจฺฉินฺนํ วูปสนฺตํ ปฏิปสฺสทฺธํ อภพฺพุปฺปตฺติกํ ญาณคฺคินา ทฑฺฒํ.\csromanspacer{} โส เปสุญฺเญ โน ยุโต น ยุตฺโต น ปยุตฺโต น สมฺมายุตฺโตติ เปสุเณยฺเย จ โน ยุโต.\csromanspacer{} เตนาห ภควา\csromandash{}}

\csromangathameasure{ปติลีโน อกุหโก, อปิหาลุ อมจฺฉรี. \\ อปฺปคพฺโภ อเชคุจฺโฉ, เปสุเณยฺเย จ โน ยุโตติ.}
\csromangathagroup{\csromangathastanza{ปติลีโน อกุหโก, อปิหาลุ อมจฺฉรี. \\ อปฺปคพฺโภ อเชคุจฺโฉ, เปสุเณยฺเย จ โน ยุโตติ.}}

\proseitem{88}{\csromansymbolfootnote{*}{ขุ 1. 412 ปิฏฺเฐ.}\textbf{สาติเยสุ อนสฺสาวี}, อติมาเน จ โน ยุโต.}

\csromangathameasure{สโณฺห จ ปฏิภานวา, น สทฺโธ น วิรชฺชติ.}
\csromangathagroup{\csromangathastanza{สโณฺห จ ปฏิภานวา, น สทฺโธ น วิรชฺชติ.}}

\prose{\textbf{สาติเยสุ อนสฺสาวี}ติ \textbf{สาติยา} วุจฺจนฺติ ปญฺจ กามคุณา.\csromanspacer{} กิํการณา \textbf{สาติยา} วุจฺจนฺติ ปญฺจ กามคุณา.\csromanspacer{} เยภุยฺเยน เทวมนุสฺสา ปญฺจ กามคุเณ อิจฺฉนฺติ สาติยนฺติ ปตฺถยนฺติ ปิหยนฺติ อภิชปฺปนฺติ, ตํการณา \textbf{สาติยา} วุจฺจนฺติ ปญฺจ กามคุณา.\csromanspacer{} เยสํ เอสา \textbf{สาติยา} ตณฺหา อปฺปหีนา, เตสํ จกฺขุโต รูปตณฺหา สวติ อาสวติ\footnote{ปสวภิ (สฺยา)} สนฺทติ ปวตฺตติ.\csromanspacer{} โสตโต สทฺทตณฺหา.\csromanspacer{} ฆานโต คนฺธตณฺหา.\csromanspacer{} ชิวฺหาโต รสตณฺหา.\csromanspacer{} กายโต โผฏฺฐพฺพตณฺหา.\csromanspacer{} มนโต ธมฺมตณฺหา สวติ อาสวติ สนฺทติ ปวตฺตติ.\csromanspacer{} เยสํ เอสา \textbf{สาติยา} ตณฺหา ปหีนา สมุจฺฉินฺนา วูปสนฺตา ปฏิปสฺสทฺธา อภพฺพุปฺปตฺติกา ญาณคฺคินา ทฑฺฒา, เตสํ จกฺขุโต รูปตณฺหา น สวติ นาสวติ\footnote{น ปสวติ (สฺยา)} น สนฺทติ น \csromanfolio{180}ปวตฺตติ.\csromanspacer{} โสตโต สทฺทตณฺหา ฯเปฯ.\csromanspacer{} มนโต ธมฺมตณฺหา น สวติ นาสวติ น สนฺทติ น ปวตฺตตีติ สาติเยสุ อนสฺสาวี.}

\prose{\textbf{อติมาเน จ โน ยุโต}ติ กตโม อติมาโน, อิเธกจฺโจ ปรํ อติมญฺญติ ชาติยา วา โคตฺเตน วา ฯเปฯ อญฺญตรญฺญตเรน วา วตฺถุนา, โย เอวรูโป มาโน มญฺญนา มญฺญิตตฺตํ อุนฺนติ อุนฺนโม ธโช สมฺปคฺคาโห เกตุกมฺยตา จิตฺตสฺส.\csromanspacer{} อยํ วุจฺจติ อติมาโน.\csromanspacer{} ยสฺเสโส อติมาโน ปหีโน สมุจฺฉินฺโน วูปสนฺโต ปฏิปสฺสทฺโธ อภพฺพุปฺปตฺติโก ญาณคฺคินา ทฑฺโฒ.\csromanspacer{} โส อติมาเน จ โน ยุโต น ยุตฺโต นปฺปยุตฺโต น สมฺมายุตฺโตติ อติมาเน จ โน ยุโต.}

\prose{\textbf{สโณฺห จ ปฏิภานวา}ติ \textbf{สโณฺห}ติ สเณฺหน กายกมฺเมน สมนฺนาคโตติ \textbf{สโณฺห}.\csromanspacer{} สเณฺหน วจีกมฺเมน.\csromanspacer{} สเณฺหน มโนกมฺเมน สมนฺนาคโตติ \textbf{สโณฺห}, สเณฺหหิ สติปฏฺฐาเนหิ สมนฺนาคโตติ \textbf{สโณฺห}, สเณฺหหิ สมฺมปฺปธาเนหิ.\csromanspacer{} สเณฺหหิ อิทฺธิปาเทหิ.\csromanspacer{} สเณฺหหิ อินฺทฺริเยหิ.\csromanspacer{} สเณฺหหิ พเลหิ.\csromanspacer{} สเณฺหหิ โพชฺฌงฺเคหิ สมนฺนาคโตติ \textbf{สโณฺห}.\csromanspacer{} สเณฺหน อริเยน อฏฺฐงฺคิเกน มคฺเคน สมนฺนาคโตติ \textbf{สโณฺห}.}

\prose{\textbf{ปฏิภานวา}ติ ตโย ปฏิภานวนฺโต ปริยตฺติปฏิภานวา ปริปุจฺฉาปฏิภานวา อธิคมปฏิภานวา.\csromanspacer{} กตโม ปริยตฺติปฏิภานวา, อิเธกจฺจสฺสา ปกติยา ปริยาปุฏํ โหติ สุตฺตํ เคยฺยํ เวยฺยากรณํ คาถา อุทานํ อิติวุตฺตกํ ชาตกํ อพฺภุตธมฺมํ เวทลฺลํ, ตสฺส ปริยตฺติํ นิสฺสาย ปฏิภายติ, อยํ ปริยตฺติปฏิภานวา.\csromanspacer{}\csromanspacer{}กตโม ปริปุจฺฉาปฏิภานวา, อิเธกจฺโจ ปริปุจฺฉิตา\footnote{ปริปุจฺฉิตํ (สี), ปริปุจฺฉโก (สฺยา)} โหติ อตฺตตฺเถ จ ญายตฺเถ จ ลกฺขเณ จ การเณ จ ฐานาฐาเน จ.\csromanspacer{} ตสฺส ตํ ปริปุจฺฉํ นิสฺสาย ปฏิภายติ, อยํ ปริปุจฺฉาปฏิภานวา.\csromanspacer{}\csromanspacer{}กตโม อธิคมปฏิภานวา, อิเธกจฺจสฺส อธิคตา โหนฺติ จตฺตาโร สติปฏฺฐานา จตฺตาโร สมฺมปฺปธานา จตฺตาโร อิทฺธิปาทา ปญฺจินฺทฺริยานิ ปญฺจ พาลานิ สตฺต โพชฺฌงฺคา อริโย อฏฺฐงฺคิโก มคฺโค จตฺตาโร อริยมคฺคา จตฺตาริ สามญฺญผลานิ จตสฺโส ปฏิสมฺภิทาโย ฉ อภิญฺญาโย, ตสฺส อตฺโถ ญาโต, ธมฺโม ญาโต, นิรุตฺติ ญาตา, อตฺเถ ญาเต อตฺโถ ปฏิภายติ, ธมฺเม ญาเต ธมฺโม ปฏิภายติ, นิรุตฺติยา ญาตาย นิรุตฺติ ปฏิภายติ, อิเมสุ}

\vspace{\tipitakaparskip}
\csromancenter{ปุราเภทสุตฺตนิทฺเทส ตีสุ ญาเณสุ ญาณํ ปฏิภานปฏิสมฺภิทา.\csromanspacer{} อิมาย ปฏิภานปฏิสมฺภิทาย อุเปโต สมุเปโต อุปคโต สมุปคโต อุปปนฺโน สมุปปนฺโน สมนฺนาคโต, โส วุจฺจติ ปฏิภานวา.\csromanspacer{} ยสฺส ปริยตฺติ นตฺถิ, ปริปุจฺฉา นตฺถิ, อธิคโม นตฺถิ, กิํ ตสฺส ปฏิภายิสฺสตีติ สโณฺห จ ปฏิภานวา.}
\prose{\csromanfolio{181}\textbf{น สทฺโธ น วิรชฺชตี}ติ \textbf{น สทฺโธ}ติ สามํ สยํ อภิญฺญาตํ อตฺตปจฺจกฺขํ ธมฺมํ, น กสฺสจิ สทฺธหติ อญฺญสฺส สมณสฺส วา พฺราหฺมณสฺส วา เทวสฺส วา มารสฺส วา พฺรหฺมุโน วา. “สพฺเพ สงฺขารา อนิจฺจา”ติ สามํ สยํ อภิญฺญาตํ ฯเปฯ “สพฺเพ สงฺขารา ทุกฺขา”ติ. “สพฺเพ ธมฺมา อนตฺตา”ติ. “อวิชฺชาปจฺจยา สงฺขารา”ติ ฯเปฯ “ชาติปจฺจยา ชรามรณนฺ”ติ. “อวิชฺชานิโรธา สงฺขารนิโรโธ”ติ ฯเปฯ “ชาตินิโรธา ชรามรณนิโรโธ”ติ. “อิทํ ทุกฺขนฺ”ติ ฯเปฯ “อยํ ทุกฺขนิโรธคามินี ปฏิปทา”ติ. “อิเม อาสวา”ติ ฯเปฯ “อยํ อาสวนิโรธคามินี ปฏิปทา”ติ. “อิเม ธมฺมา อภิญฺเญยฺยา”ติ ฯเปฯ “อิเม ธมฺมา สจฺฉิกาตพฺพา”ติ สามํ สยํ อภิญฺญาตํ ฯเปฯ ฉนฺนํ ผสฺสายตนานํ สมุทยญฺจ อตฺถงฺคมญฺจ อสฺสาทญฺจ อาทีนวญฺจ นิสฺสรณญฺจ.\csromanspacer{} ปญฺจนฺนํ อุปาทานกฺขนฺธานํ สมุทยญฺจ ฯเปฯ จตุนฺนํ มหาภูตานํ สมุทยญฺจ อตฺถงฺคมญฺจ อสฺสาทญฺจ อาทีนวญฺจ นิสฺสรณญฺจ สามํ สยํ อภิญฺญาตํ ฯเปฯ “ยํ กิญฺจิ สมุทยธมฺมํ, สพฺพํ ตํ นิโรธธมฺมนฺ”ติ สามํ สยํ อภิญฺญาตํ อตฺตปจฺจกฺขํ ธมฺมํ, น กสฺสจิ สทฺทหติ อญฺญสฺส สมณสฺส วา พฺราหฺมณสฺส วา เทวสฺส วา มารสฺส วา พฺรหฺมุโน วา\footnote{พฺรหฺมุโน วา ฯเปฯ (สี, ก)}.}

\prose{วุตฺตํ เหตํ ภควตา\csromandash{} * สทฺทหสิ ตฺวํ สาริปุตฺต สทฺธินฺทฺริยํ ภาวิตํ พหุลีกตํ อมโตคธํ โหติ อมตปรายณํ อมตปริโยสานํ.\csromanspacer{} วีริยินฺทฺริยํ.\csromanspacer{} สตินฺทฺริยํ.\csromanspacer{} สมาธินฺทฺริยํ.\csromanspacer{} ปญฺญินฺทฺริยํ ภาวิตํ พหุลีกตํ อมโต คธํ โหติ อมตปรายณํ อมตปริโยสานนฺติ.}

\prose{น ขฺวาหํ เอตฺถ ภนฺเต ภควโต สทฺธาย คจฺฉามิ สทฺธินฺทฺริยํ.\csromanspacer{} วีริยินฺทฺริยํ.\csromanspacer{} สตินฺทฺริยํ.\csromanspacer{} สมาธินฺทฺริยํ.\csromanspacer{} ปญฺญินฺทฺริยํ ภาวิตํ พหุลีกตํ อมโตคธํ โหติ อมตรายณํ อมตปริโยสานํ.\csromanspacer{} เยสํ เหตํ ภนฺเต อญฺญาตํ อสฺส อทิฏฺฐํ อวิทิตํ อสจฺฉิกตํ อผสฺสิตํ ปญฺญาย, เต ตตฺถ ปเรสํ สทฺธาย คจฺเฉยฺยุํ สทฺธินฺทฺริยํ ภาวิตํ พหุลีกตํ อมโตคธํ โหติ อมตปรายณํ อมตปริโยสานํ.\csromanspacer{} วีริยินฺทฺริยํ.\csromanspacer{} สตินฺทฺริยํ. \csromanfolio{182}สมาธินฺทฺริยํ.\csromanspacer{} ปญฺญินฺทฺริยํ ภาวิตํ พหุลีกตํ อมโตคธํ โหติ อมตปรายณํ อมตปริโยสานํ.\csromanspacer{} เยสญฺจ โข เอตํ ภนฺเต ญาตํ ทิฏฺฐํ วิทิตํ สจฺฉิกตํ ผสฺสิตํ ปญฺญาย, นิกฺกงฺขา เต ตตฺถ นิพฺพิจิกิจฺฉา, สทฺธินฺทฺริยํ.\csromanspacer{} วีริยินฺทฺริยํ.\csromanspacer{} สตินฺทฺริยํ.\csromanspacer{} สมาธินฺทฺริยํ.\csromanspacer{} ปญฺญินฺทฺริยํ ภาวิตํ พหุลีกตํ อมโตคธํ โหติ อมตปรายณํ อมตปริโยสานํ.\csromanspacer{} มยฺหญฺจ โข เอตํ ภนฺเต ญาตํ ทิฏฺฐํ วิทิตํ สจฺฉิกตํ ผสฺสิตํ ปญฺญาย, นิกฺกงฺโขหํ ตตฺถ นิพฺพิจิกิจฺโฉ.\csromanspacer{} สทฺธินฺทฺริยํ.\csromanspacer{} วีริยินฺทฺริยํ.\csromanspacer{} สตินฺทฺริยํ.\csromanspacer{} สมาธินฺทฺริยํ.\csromanspacer{} ปญฺญินฺทฺริยํ ภาวิตํ พหุลีกตํ อมโตคธํ โหติ อมตปรายณํ อมตปริโยสานนฺติ.}

\prose{สาธุ สาธุ สาริปุตฺต เยสํ เหตํ สาริปุตฺต อญฺญาตํ อสฺส อทิฏฺฐํ อวิทิตํ อสจฺฉิกตํ อผสฺสิตํ ปญฺญาย, เต ตตฺถ ปเรสํ สทฺธาย คจฺเฉยฺยุํ สทฺธินฺทฺริยํ ฯเปฯ ปญฺญินฺทฺริยํ ภาวิตํ พหุลิกตํ อมโตคธํ โหติ อมตปรายณํ อมตปริโยสานนฺติ.}

\csromangathameasure{*อสฺสทฺโธ อกตญฺญู จ, สนฺธิจฺเฉโท จ โย นโร. \\ หตาวกาโส วนฺตาโส, ส เว อุตฺตมโปริโสติ.}
\csromangathagroup{\csromangathastanza{\csromansymbolfootnote{*}{ขุ 1. 27 ธมฺมปเท.}อสฺสทฺโธ อกตญฺญู จ, สนฺธิจฺเฉโท จ โย นโร. \\ หตาวกาโส วนฺตาโส, ส เว อุตฺตมโปริโสติ.}}

\prose{\textbf{น สทฺโธ น วิรชฺชตี}ติ สพฺเพ พาลปุถุชฺชนา รชฺชนฺติ ปุถุชฺชนกลฺยาณกํ อุปาทาย สตฺต เสกฺขา วิรชฺชนฺติ.\csromanspacer{} อรหา เนว รชฺชติ โน วิรชฺชติ, วิรตฺโต โส ขยา ราคสฺส วีตราคตฺตา, ขยา โทสสฺส วีตโทสตฺตา, ขยา โมหสฺส วีตโมหตฺตา, โส วุฏฺฐวาโส จิณฺณจรโณ ฯเปฯ.\csromanspacer{} ชาติมรณสํสาโร นตฺถิ ตสฺส ปุนพฺภโวติ น สทฺโธ น วิรชฺชติ.\csromanspacer{} เตนาห ภควา\csromandash{}}

\csromangathameasure{สาติเยสุ อนสฺสาวี, อติมาเน จ โน ยุโต. \\ สโณฺห จ ปฏิภานวา, น สทฺโธ น วิรชฺชตีติ.}
\csromangathagroup{\csromangathastanza{สาติเยสุ อนสฺสาวี, อติมาเน จ โน ยุโต. \\ สโณฺห จ ปฏิภานวา, น สทฺโธ น วิรชฺชตีติ.}}

\proseitem{89}{\csromansymbolfootnote{+}{ขุ 1. 412 ปิฏฺเฐ.}ลาภกมฺยา น สิกฺขติ, อลาเภ จ น กุปฺปติ.}

\csromangathameasure{อวิรุทฺโธ จ ตณฺหาย, รเสสุ นานุคิชฺฌติ.}
\csromangathagroup{\csromangathastanza{อวิรุทฺโธ จ ตณฺหาย, รเสสุ\footnote{รเส จ (สี, สฺยา)} นานุคิชฺฌติ.}}

\prose{\textbf{ลาภกมฺยา น สิกฺขติ, อลาเภ จ น กุปฺปตี}ติ กถํ ลาภกมฺยา สิกฺขติ, อิธ ภิกฺขเว ภิกฺขุ ภิกฺขุํ ปสฺสติ ลาภิํ จีวรปิณฺฑปาตเสนาสนคิลานปจฺจยเภสชฺชปริกฺขารานํ, ตสฺส เอวํ โหติ “เกน}

\vspace{\tipitakaparskip}
\csromancenter{ปุราเภทสุตฺตนิทฺเทส นุ โข อยมายสฺมา ลาภี จีวรปิณฺฑปาตเสนาสนคิลานปจฺจยเภสชฺช- ปริกฺขารานนฺ”ติ.\csromanspacer{} ตสฺส เอวํ โหติ “อยํ โข อายสฺมา สุตฺตนฺติโก, เตนายมายสฺมา ลาภี จีวรปิณฺฑปาตเสนาสนคิลานปจฺจยเภสชฺช- ปริกฺขารานนฺ”ติ.\csromanspacer{} โส ลาภเหตุ ลาภปจฺจยา ลาภการณา ลาภาภินิพฺพตฺติยา ลาภํ ปริปาเจนฺโต สุตฺตนฺตํ ปริยาปุณาติ.\csromanspacer{} เอวมฺปิ ลาภกมฺยา สิกฺขติ.}
\prose{\csromanfolio{183}อถ วา ภิกฺขุ ภิกฺขุํ ปสฺสติ ลาภิํ จีวรปิณฺฑปาตเสนาสน- คิลานปจฺจยเภสชฺชปริกฺขารานํ, ตสฺส เอวํ โหติ “เกน นุ โข อยมายสฺมา ลาภี จีวรปิณฺฑปาตเสนาสนคิลานปจฺจยเภสชฺช- ปริกฺขารานนฺ”ติ.\csromanspacer{} ตสฺส เอวํ โหติ “อยํ โข อายสฺมา วินยธโร ฯเปฯ ธมฺมกถิโก.\csromanspacer{} อาภิธมฺมิโก.\csromanspacer{} เตนายมายสฺมา ลาภี จีวรปิณฺฑปาตเสนาสนคิลานปจฺจยเภสชฺชปริกฺขารานนฺ”ติ.\csromanspacer{} โส ลาภเหตุ ลาภปจฺจยา ลาภการณา ลาภาภินิพฺพตฺติยา ลาภํ ปริปาเจนฺโต อภิธมฺมํ ปริยาปุณาติ.\csromanspacer{} เอวมฺปิ ลาภกมฺยา สิกฺขติ.}

\prose{อถ วา ภิกฺขุ ภิกฺขุํ ปสฺสติ ลาภิํ จีวรปิณฺฑปาตเสนาสน- คิลานปจฺจยเภสชฺชปริกฺขารานํ, ตสฺส เอวํ โหติ “เกน นุ โข อยมายสฺมา ลาภี จีวรปิณฺฑปาตเสนาสนคิลานปจฺจยเภสชฺช- ปริกฺขารานนฺ”ติ.\csromanspacer{} ตสฺส เอวํ โหติ “อยํ โข อายสฺมา อารญฺญิโก.\csromanspacer{} ปิณฺฑปาติโก.\csromanspacer{} ปํสุกูลิโก.\csromanspacer{} เตจีวริโก.\csromanspacer{} สปทานจาริโก.\csromanspacer{} ขลุปจฺฉาภตฺติโก.\csromanspacer{} เนสชฺชิโก.\csromanspacer{} ยถาสนฺถติโก, เตนายมายสฺมา ลาภี จีวรปิณฺฑปาตเสนาสนคิลานปจฺจยเภสชฺชปริกฺขารานนฺ”ติ.\csromanspacer{} โส ลาภเหตุ ลาภปจฺจยา ลาภการณา ลาภาภินิพฺพตฺติยา ลาภํ ปริปาเจนฺโต อารญฺญิโก โหติ ฯเปฯ ยถาสนฺถติโก โหติ.\csromanspacer{} เอวมฺปิ ลาภกมฺยา สิกฺขติ.}

\prose{กถํ น ลาภกมฺยา สิกฺขติ, อิธ ภิกฺขุ น ลาภเหตุ น ลาภปจฺจยา น ลาภการณา น ลาภาภินิพฺพตฺติยา น ลาภํ ปริปาเจนฺโต ยาวเทว อตฺตทมตฺถาย อตฺตสมตฺถาย อตฺตปรินิพฺพาปนตฺถาย สุตฺตนฺตํ ปริยาปุณาติ, วินยํ ปริยาปุณาติ, อภิธมฺมํ ปริยาปุณาติ.\csromanspacer{} เอวมฺปิ น ลาภกมฺยา สิกฺขติ.}

\prose{\csromanfolio{184}อถ วา ภิกฺขุ น ลาภเหตุ น ลาภปจฺจยา น ลาภการณา น ลาภาภินิพฺพตฺติยา น ลาภํ ปริปาเจนฺโต ยาวเทว อปฺปิจฺฉญฺเญว\footnote{อปฺปิจฺฉํเยว (สี)} นิสฺสาย สนฺตุฏฺฐิญฺเญว นิสฺสาย สลฺเลขญฺเญว นิสฺสาย ปวิเวกญฺเญว นิสฺสาย อิทมตฺถิตญฺเญว\footnote{อิทมตฺถิกตญฺเญว (สี)} นิสฺสาย อารญฺญิโก โหติ.\csromanspacer{} ปิณฺฑปาติโก โหติ.\csromanspacer{} ปํสุกูลิโก โหติ.\csromanspacer{} เตจีวริโก โหติ.\csromanspacer{} สปทานจาริโก โหติ.\csromanspacer{} ขลุปจฺฉาภตฺติโก โหติ.\csromanspacer{} เนสชฺชิโก โหติ.\csromanspacer{} ยถาสนฺถติโก โหติ.\csromanspacer{} เอวมฺปิ น ลาภกมฺยา สิกฺขตีติ ลาภกมฺยา น สิกฺขติ.}

\prose{\textbf{อลาเภ จ น กุปฺปตี}ติ กถํ อลาเภ กุปฺปติ, อิเธกจฺโจ “กุลํ วา น ลภามิ, คณํ วา น ลภามิ, อาวาสํ วา น ลภามิ, ลาภํ วา น ลภามิ, ยสํ วา น ลภามิ, ปสํสํ วา น ลภามิ, สุขํ วา น ลภามิ, จีวรํ วา น ลภามิ, ปิณฺฑปาตํ วา น ลภามิ, เสนาสนํ วา น ลภามิ, คิลานปจฺจยเภสชฺชปริกฺขารํ วา น ลภามิ, คิลานุปฏฺฐากํ วา น ลาภามิ อปฺปญฺญโตมฺหี”ติ กุปฺปติ พฺยาปชฺชติ ปติตฺถียติ, โกปญฺจ โทสญฺจ อปฺปจฺจยญฺจ ปาตุกโรติ.\csromanspacer{} เอวํ อลาเภ กุปฺปติ.}

\prose{กถํ อลาเภ น กุปฺปติ, อิธ ภิกฺขุ “กุลํ วา น ลภามิ, คณํ วา น ลภามิ ฯเปฯ อปฺปญฺญาโตมฺหี”ติ น กุปฺปติ น พฺยาปชฺชติ น ปติตฺถียติ, น โกปญฺจ โทสญฺจ อปฺปจฺจยญฺจ ปาตุกโรติ.\csromanspacer{} เอวํ อลาเภ น กุปฺปตีติ ลาภกมฺยา น สิกฺขติ, อลาเภ จ น กุปฺปติ.}

\prose{\textbf{อวิรุทฺโธ จ ตณฺหาย, รเสสุ นานุคิชฺฌตี}ติ \textbf{วิรุทฺโธ}ติ * โย จิตฺตสฺส อาฆาโต ปฏิฆาโต ปฏิฆํ ปฏิวิโรโธ โกโป ปโกโป สมฺปโกโป โทโส ปโทโส สมฺปโทโส จิตฺตสฺส พฺยาปตฺติ มโนปโทโส โกโธ กุชฺฌนา กุชฺฌิตตฺตํ โทโส ทุสฺสนา ทุสฺสิตตฺตํ พฺยาปตฺติ พฺยาปชฺชนา พฺยาปชฺชิตตฺตํ วิโรโธ ปฏิวิโรโธ จณฺฑิกฺกํ อสุโรโป อนตฺตมนตา จิตฺตสฺส, อยํ วุจฺจติ วิโรโธ.\csromanspacer{} ยสฺเสโส วิโรโธ ปหีโน สมุจฺฉินฺโน วูปสนฺโต ปฏิปสฺสทฺโธ อภพฺพุปฺปตฺติโก ญาณคฺคินา ทฑฺโฒ, โส วุจฺจติ อ\textbf{วิรุทฺโธ}.\csromanspacer{}\csromanspacer{}\textbf{ตณฺหา}ติ รูปตณฺหา สทฺทตณฺหา คนฺธตณฺหา รสตณฺหา โผฏฺฐพฺพตณฺหา ธมฺมตณฺหา.\csromanspacer{}\csromanspacer{}\textbf{รโส}ติ มูลรโส ขนฺธรโส ตจรโส ปตฺตรโส ปุปฺผรโส ผลรโส}

\vspace{\tipitakaparskip}
\csromancenter{ปุราเภทสุตฺตนิทฺเทส อมฺพิลํ มธุรํ ติตฺตกํ กฏุกํ โลณิกํ ขาริกํ ลมฺพิกํ\footnote{ลปิลํ (สี) ลมฺพิลํ (สฺยา), ลพิลํ (ก); อภิ 1. 167; อายตนวิภงฺเค 73 ปิฏฺเฐ.} กสาโว สาทุ อสาทุ สีตํ อุณฺหํ.\csromanspacer{} สนฺเตเก สมณพฺราหฺมณา รสคิทฺธา, เต ชิวฺหคฺเคน รสคฺคานิ ปริเยสนฺตา อาหิณฺฑนฺติ.\csromanspacer{} เต อมฺพิลํ ลภิตฺวา อนมฺพิลํ ปริเยสนฺติ, อนมฺพิลํ ลภิตฺวา อมฺพิลํ ปริเยสนฺติ, มธุรํ ลภิตฺวา อมธุรํ ปริเยสนฺติ, อมธุรํ ลภิตฺวา มธุรํ ปริเยสนฺติ, ติตฺตกํ ลภิตฺวา อติตฺตกํ ปริเยสนฺติ, อติตฺตกํ ลภิตฺวา ติตฺตกํ ปริเยสนฺติ, กฏุกํ ลภิตฺวา อกฏุกํ ปริเยสนฺติ, อกฏุกํ ลภิตฺวา กฏุกํ ปริเยสนฺติ, โลณิกํ ลภิตฺวา อโลณิกํ ปริเยสนฺติ, อโลณิกํ ลภิตฺวา โลณิกํ ปริเยสนฺติ.\csromanspacer{} ขาริกํ ลภิตฺวา อขาริกํ ปริเยสนฺติ, อขาริกํ ลภิตฺวา ขาริกํ ปริเยสนฺติ, ลมฺพิกํ ลภิตฺวา กสาวํ ปริเยสนฺติ, กสาวํ ลภิตฺวา ลมฺพิกํ ปริเยสนฺติ.\csromanspacer{} สาทุํ ลภิตฺวา อสาทุํ ปริเยสนฺติ, อสาทุํ ลภิตฺวา สาทุํ ปริเยสนฺติ, สีตํ ลภิตฺวา อุณฺหํ ปริเยสนฺติ, อุณฺหํ ลภิตฺวา สีตํ ปริเยสนฺติ.\csromanspacer{} เต ยํ ยํ ลภิตฺวา เตน เตน น สนฺตุสฺสนฺติ, อปราปรํ ปริเยสนฺติ.\csromanspacer{} มนาปิเกสุ รเสสุ รตฺตา คิทฺธา คธิตา มุจฺฉิตา อชฺโฌสนฺนา ลคฺคา ลคฺคิตา ปลิพุทฺธา.\csromanspacer{} ยสฺเสสา รสตณฺหา ปหีนา สมุจฺฉินฺนา วูปสนฺตา ปฏิปสฺสทฺธา อภพฺพุปฺปตฺติกา ญาณคฺคินา ทฑฺฒา.\csromanspacer{} โส ปฏิสงฺขา โยนิโส อาหารํ อาหาเรติ เนว ทวาย น มทาย น มณฺฑนาย น วิภูสนาย ยาวเทว อิมสฺส กายสฺส ฐิติยา ยาปนาย วิหิํสูปรติยา พฺรหฺมจริยานุคฺคหาย อิติ ปุราณญฺจ เวทนํ ปฏิหงฺขามิ, นวญฺจ เวทนํ น อุปฺปาเทสฺสามิ, ยาตฺรา จ เม ภวิสฺสติ อนวชฺชตา จ ผาสุวิหาโร จ.}
\prose{\csromanfolio{185}ยถา วนํ อาลิมฺเปยฺย ยาวเทว โรปนตฺถาย.\csromanspacer{} ยถา วา ปน อกฺขํ อพฺภญฺเชยฺย ยาวเทว ภารสฺส นิตฺถรณตฺถาย.\csromanspacer{} ยถา วา ปน ปุตฺตมํสํ อาหารํ อาหเรยฺย ยาวเทว กนฺตารสฺส นิตฺถรณตฺถาย.\csromanspacer{} เอวเมวํ ภิกฺขุ ปฏิสงฺขา โยนิโส อาหารํ อาหาเรติ เนว ทวาย ฯเปฯ ผาสุวิหาโร จาติ รสตณฺหํ ปชหติ วิโนเทติ พฺยนฺติํ กโรติ อนภาวํ คเมติ, รสตณฺหาย อารโต อสฺส วิรโต ปฏิวิรโต นิกฺขนฺโต นิสฺสโฏ วิปฺปมุตฺโต วิสญฺญุตฺโต วิมริยาทิกเตน เจตสา วิหรตีติ อวิรุทฺโธ จ ตณฺหาย, รเสสุ นานุคิชฺฌติ.\csromanspacer{} เตนาห ภควา\csromandash{}}

\csromangathameasure{ลาภกมฺยา น สิกฺขติ, อลาเภ จ น กุปฺปติ. \\ อวิรุทฺโธ จ ตณฺหาย, รเสสุ นานุคิชฺฌตีติ.}
\csromangathagroup{\csromangathastanza{\csromanfolio{186}ลาภกมฺยา น สิกฺขติ, อลาเภ จ น กุปฺปติ. \\ อวิรุทฺโธ จ ตณฺหาย, รเสสุ นานุคิชฺฌตีติ.}}

\proseitem{90}{\csromansymbolfootnote{*}{ขุ 1. 412 ปิฏฺเฐ.}\textbf{อุเปกฺขโก สทา สโต}, น โลเก มญฺญเต สมํ.}

\csromangathameasure{น วิเสสี น นีเจยฺโย, ตสฺส โน สนฺติ อุสฺสทา.}
\csromangathagroup{\csromangathastanza{น วิเสสี น นีเจยฺโย, ตสฺส โน สนฺติ อุสฺสทา.}}

\prose{\textbf{อุเปกฺขโก สทา สโต}ติ \textbf{อุเปกฺขโก}ติ ฉฬงฺคุเปกฺขาย สมนฺนาคโต. + จกฺขุนา รูปํ ทิสฺวา เนว สุมโน โหติ น ทุมฺมโน, \textbf{อุเปกฺขโก} วิหรติ สโต สมฺปชาโน.\csromanspacer{} โสเตน สทฺทํ สุตฺวา ฯเปฯ.\csromanspacer{} มนสา ธมฺมํ วิญฺญาย เนว สุมโน โหติ น ทุมฺมโน, \textbf{อุเปกฺขโก} วิหรติ สโต สมฺปชาโน.\csromanspacer{} จกฺขุนา รูปํ ทิสฺวา มนาปํ นาภิคิชฺฌติ นาภิหํสติ1, น ราคํ ชเนติ, ตสฺส ฐิโตว กาโย โหติ, ฐิตํ จิตฺตํ อชฺฌตฺตํ สุสณฺฐิตํ สุวิมุตฺตํ, จกฺขุนา โข ปเนว รูปํ ทิสฺวา อมนาปํ น มงฺกุ โหติ อปฺปติฏฺฐิตจิตฺโต\footnote{อปฺปติฏฺฐีน จิตฺโต (สฺยา), อปฺปติฏฺฐนจิตฺต (ก) 3. อาทินมนโส (สฺยา)} อลีนมนโส3 อพฺยาปนฺนเจตโส, ตสฺส ฐิโตว กาโย โหติ, ฐิตํ จิตฺตํ อชฺฌตฺตํ สุสณฺฐิตํ สุวิมุตฺตํ.\csromanspacer{} โสเตน สทฺทํ สุตฺวา.\csromanspacer{} ฆาเนน คนฺธํ ฆายิตฺวา.\csromanspacer{} ชิวฺหาย รสํ สายิตฺวา.\csromanspacer{} กาเยน โผฏฺฐพฺพํ ผุสิตฺวา.\csromanspacer{} มนสา ธมฺมํ วิญฺญาย มนาปํ นาภิคิชฺฌติ นาภิหํสติ น ราคํ ชเนติ, ตสฺส ฐิโตว กาโย โหติ, ฐิตํ จิตฺตํ อชฺฌตฺตํ สุสณฺฐิตํ สุวิมุตฺตํ.\csromanspacer{} มนสา โข ปเนว ธมฺมํ วิญฺญาย อมนาปํ น มงฺกุ โหติ อปฺปติฏฺฐิตจิตฺโต อลีนมนโส อพฺยาปนฺนเจตโส, ตสฺส ฐิโตว กาโย โหติ, ฐิตํ จิตฺตํ อชฺฌตฺตํ สุสณฺฐิตํ สุวิมุตฺตํ.}

\prose{จกฺขุนา รูปํ ทิสฺวา มนาปามนาเปสุ รูเปสุ ตสฺส ฐิโตว กาโย โหติ, ฐิตํ จิตฺตํ อชฺฌตฺตํ สุสณฺฐิตํ สุวิมุตฺตํ.\csromanspacer{} โสเตน สทฺทํ สุตฺวา ฯเปฯ.\csromanspacer{} มนสา ธมฺมํ วิญฺญาย มนาปามนาเปสุ ธมฺเมสุ ตสฺส ฐิโตว กาโย โหติ, ฐิตํ จิตฺตํ อชฺฌตฺตํ สุสณฺฐิตํ สุวิมุตฺตํ.}

\prose{จกฺขุนา รูปํ ทิสฺวา รชนีเย น รชฺชติ, ทุสฺสนีเย\footnote{โทสนีเย (พหูสุ)} น ทุสฺสติ, โมหนีเย น มุยฺหติ, โกปนีเย น กุปฺปติ, มทนีเย น มชฺชติ, กิเลสนีเย น กิลิสฺสติ.\csromanspacer{} โสเตน สทฺทํ สุตฺวา ฯเปฯ.\csromanspacer{} มนสา ธมฺมํ วิญฺญาย รชนีเย น รชฺชติ, ทุสฺสนีเย น ทุสฺสติ, โมหนีเย น มุยฺหติ, โกปนีเยน น กุปฺปติ, มทนีเย น มชฺชติ, กิเลสนีเย น กิลิสฺสติ.}

\vspace{\tipitakaparskip}
\csromancenter{ปุราเภทสุตฺตนิทฺเทส ทิฏฺเฐ ทิฏฺฐมตฺโต, สุเต สุตมตฺโต, มุเต มุตมตฺโต, วิญฺญาเต วิญฺญาตมตฺโต.\csromanspacer{} ทิฏฺเฐ น ลิมฺปติ, สุเต น ลิมฺปติ, มุเต น ลิมฺปติ, วิญฺญาเต น ลิมฺปติ.\csromanspacer{} ทิฏฺเฐ อนูปโย อนปาโย อนิสฺสิโต อปฺปฏิพทฺโธ วิปฺปมุตฺโต วิสญฺญุตฺโต วิมริยาทิกเตน เจตสา วิหรติ.\csromanspacer{} สุเต.\csromanspacer{} มุเต.\csromanspacer{} วิญฺญาเต อนูปโย อนปาโย อนิสฺสิโต อปฺปฏิพทฺโธ วิปฺปมุตฺโต วิสญฺญุตฺโต วิมริยาทิกเตน เจตสา วิหรติ.}
\prose{\csromanfolio{187}สํวิชฺชติ อรหโต จกฺขุ, ปสฺสติ อรหา จกฺขุนา รูปํ, ฉนฺทราโค อรหโต นตฺถิ, สุวิมุตฺตจิตฺโต อรหา.\csromanspacer{} สํวิชฺชติ อรหโต โสตํ, สุณาติ อรหา โสเตน สทฺทํ, ฉนฺทราโค อรหโต นตฺถิ, สุวิมุตฺตจิตฺโต อรหา.\csromanspacer{} สํวิชฺชติ อรหโต ฆานํ, ฆายติ อรหา ฆาเนน คนฺธํ, ฉนฺทราโค อรหโต นตฺถิ, สุวิมุตฺตจิตฺโต อรหา.\csromanspacer{} สํวิชฺชติ อรหโต ชิวฺหา, สายติ อรหา ชิวฺหาย รสํ ฯเปฯ.\csromanspacer{} สํวิชฺชติ อรหโต กาโย, ผุสติ อรหา กาเยน โผฏฺฐพฺพํ ฯเปฯ.\csromanspacer{} สํวิชฺชติ อรหโต มโน, วิชานาติ อรหา มนสา ธมฺมํ, ฉนฺทราโค อรหโต นตฺถิ, สุวิมุตฺตจิตฺโต อรหา.}

\prose{จกฺขุ รูปารามํ รูปรตํ รูปสมฺมุทิตํ, ตํ อรหโต ทนฺตํ คุตฺตํ รกฺขิตํ สํวุตํ, ตสฺส จ สํวราย ธมฺมํ เทเสติ.\csromanspacer{} โสตํ สทฺทารามํ ฯเปฯ.\csromanspacer{} ฆานํ คนฺธารามํ.\csromanspacer{} ชิวฺหา รสารามา รสรตา รสสมฺมุทิตา, สา อรหโต ทนฺตา คุตฺตา รกฺขิตา สํวุตา, ตสฺสา จ สํวราย ธมฺมํ เทเสติ.\csromanspacer{} กาโย โผฏฺฐพฺพาราโม ฯเปฯ.\csromanspacer{} มโน ธมฺมาราโม ธมฺมรโต ธมฺมสมฺมุทิโต, โส อรหโต ทนฺโต คุตฺโต รกฺขิโต สํวุโต, ตสฺส จ สํวราย ธมฺมํ เทเสติ.}

\csromangathameasure{ทนฺตํ นยนฺติ สมิติํ, ทนฺตํ ราชาภิรูหติ. \\ ทนฺโต เสฏฺโฐ มนุสฺเสสุ, โย’ติวากฺยํ ติติกฺขติ. \\ วรมสฺสตรา ทนฺตา, อาชานียา จ สินฺธวา. \\ กุญฺชรา จ มหานาคา, อตฺตทนฺโต ตโต วรํ. \\ น หิ เอเตหิ ยาเนหิ, คจฺเฉยฺย อคตํ ทิสํ. \\ ยถา’ตฺตนา สุทนฺเตน, ทนฺโต ทนฺเตน คจฺฉติ. \\ วิธาสุ น วิกมฺปนฺติ, วิปฺปมุตฺตา ปุนพฺภวา. \\ ทนฺตภูมิมนุปฺปตฺตา, เต โลเก วิชิตาวิโน. \\ ยสฺสินฺทฺริยานิ ภาวิตานิ, \\ อชฺฌตฺตํ พหิทฺธา จ สพฺพโลเก. \\ นิพฺพิชฺฌ อิมํ ปรญฺจ โลกํ, \\ กาลํ กงฺขติ ภาวิโต ส ทนฺโตติ.}
\csromangathagroup{\csromangathastanza{ทนฺตํ นยนฺติ สมิติํ, ทนฺตํ ราชาภิรูหติ. \\ ทนฺโต เสฏฺโฐ มนุสฺเสสุ, โย’ติวากฺยํ ติติกฺขติ.}\csromangathastanza{วรมสฺสตรา ทนฺตา, อาชานียา จ\footnote{อาชานิยาว (สฺยา) ขุ 1. 59; ขุ 8. 66 ปิฏฺเฐสุปิ.} สินฺธวา. \\ กุญฺชรา จ มหานาคา, อตฺตทนฺโต ตโต วรํ.}\csromangathastanza{น หิ เอเตหิ ยาเนหิ, คจฺเฉยฺย อคตํ ทิสํ. \\ ยถา’ตฺตนา สุทนฺเตน, ทนฺโต ทนฺเตน คจฺฉติ.}\csromangathastanza{\csromanfolio{188}วิธาสุ น วิกมฺปนฺติ, วิปฺปมุตฺตา ปุนพฺภวา. \\ ทนฺตภูมิมนุปฺปตฺตา, เต โลเก วิชิตาวิโน.}\csromangathastanza{ยสฺสินฺทฺริยานิ ภาวิตานิ\footnote{วิภาวิตานิ (สี)}, \\ อชฺฌตฺตํ พหิทฺธา จ\footnote{วิภาวิตานิ (สี)} สพฺพโลเก. \\ นิพฺพิชฺฌ อิมํ\footnote{วิภาวิตานิ (สี)} ปรญฺจ โลกํ, \\ กาลํ กงฺขติ ภาวิโต ส ทนฺโตติ.}}

\prose{\textbf{อุเปกฺขโก สทา}ติ สทา สพฺพทา สพฺพกาลํ นิจฺจกาลํ ธุวกาลํ ฯเปฯ ปจฺฉิเม วโยขนฺเธ.\csromanspacer{} \textbf{สโต}ติ จตูหิ การเณหิ สโต, กาเย กายานุปสฺสนาสติปฏฺฐานํ ภาเวนฺโต สโต.\csromanspacer{} เวทนาสุ.\csromanspacer{} จิตฺเต.\csromanspacer{} ธมฺเมสุ ธมฺมานุปสฺสนาสติปฏฺฐานํ ภาเวนฺโต สโต ฯเปฯ โส วุจฺจติ สโตติ อุเปกฺขโก สทา สโต.}

\prose{\textbf{น โลเก มญฺญเต สมนฺ}ติ “สทิโสหมสฺมี”ติ มานํ น ชเนติ ชาติยา วา โคตฺเตน วา ฯเปฯ อญฺญตรญฺญตเรน วา วตฺถุนาติ น โลเก มญฺญเต สมํ.}

\prose{\textbf{น วิเสสี น นีเจยฺโย}ติ “เสยฺโยหมสฺมี”ติ อติมานํ น ชเนติ ชาติยา วา โคตฺเตน วา ฯเปฯ อญฺญตรญฺญตเรน วา วตฺถุนา. “หีโนหมสฺมี”ติ โอมานํ น ชเนติ ชาติยา วา โคตฺเตน วา ฯเปฯ อญฺญตรญฺญตเรน วา วตฺถุนาติ น วิเสสี น นีเจยฺโย.}

\prose{\textbf{ตสฺส โน สนฺติ อุสฺสทา}ติ \textbf{ตสฺสา}ติ อรหโต ขีณาสวสฺส.\csromanspacer{} \textbf{อุสฺสทา}ติ สตฺตุสฺสทา ราคุสฺสโท โทสุสฺสโท โมหุสฺสโท มานุสฺสโท ทิฏฺฐุสฺสโท กิเลสุสฺสโท กมฺมุสฺสโท.\csromanspacer{} ตสฺสิเม อุสฺสทา นตฺถิ น สนฺติ น สํวิชฺชนฺติ นุปลพฺภนฺติ, ปหีนา สมุจฺฉินฺนา วูปสนฺตา ปฏิปสฺสทฺธา อภพฺพุปฺปตฺติกา ญาณคฺคินา ทฑฺฒาติ ตสฺส โน สนฺติ อุสฺสทา.\csromanspacer{} เตนาห ภควา\csromandash{}}

\csromangathameasure{อุเปกฺขโก สทา สโต, น โลเก มญฺญเต สมํ. \\ น วิเสสี น นีเจยฺโย, ตสฺส โน สนฺติ อุสฺสทาติ.}
\csromangathagroup{\csromangathastanza{อุเปกฺขโก สทา สโต, น โลเก มญฺญเต สมํ. \\ น วิเสสี น นีเจยฺโย, ตสฺส โน สนฺติ อุสฺสทาติ.}}

\csromanheader{2}{10. ปุราเภทสุตฺตนิทฺเทส}
\proseitem{91}{\csromanfolio{189}\csromansymbolfootnote{*}{ขุ 1. 412 ปิฏฺเฐ.}ยสฺส นิสฺสยตา\footnote{นิสฺสยนา (ก)} \textbf{นตฺถิ, ญตฺวา ธมฺมํ อนิสฺสิโต.}}

\csromangathameasure{ภวาย วิภวาย วา, ตณฺหา ยสฺส น วิชฺชติ.}
\csromangathagroup{\csromangathastanza{ภวาย วิภวาย วา, ตณฺหา ยสฺส น วิชฺชติ.}}

\prose{\textbf{ยสฺส นิสฺสยตา นตฺถี}ติ \textbf{ยสฺสา}ติ อรหโต ขีณาสวสฺส.\csromanspacer{} \textbf{นิสฺสยา}ติ เทฺว นิสฺสยา \textbf{ตณฺหา}นิสฺสโย จ ทิฏฺฐินิสฺสโย จ ฯเปฯ อยํ \textbf{ตณฺหา}นิสฺสโย ฯเปฯ อยํ ทิฏฺฐินิสฺสโย.\csromanspacer{} ตสฺส \textbf{ตณฺหา}นิสฺสโย ปหีโน, ทิฏฺฐินิสฺสโย ปฏินิสฺสฏฺโฐ, \textbf{ตณฺหา}นิสฺสยสฺส ปหีนตฺตา ทิฏฺฐินิสฺสยสฺส ปฏินิสฺสฏฺฐตฺตา นิสฺสยตา ยสฺส นตฺถิ น สนฺติ น สํวิชฺชติ นุปลพฺภติ, ปหีนา สมุจฺฉินฺนา วูปสนฺตา ปฏิปสฺสทฺธา อภพฺพุปฺปตฺติกา ญาณคฺคินา ทฑฺฒาติ ยสฺส นิสฺสยตา นตฺถิ.}

\prose{\textbf{ญตฺวา ธมฺมํ อนิสฺสิโต}ติ \textbf{ญตฺวา}ติ \textbf{ญตฺวา} ชานิตฺวา ตุลยิตฺวา ตีรยิตฺวา วิภาวยิตฺวา วิภูตํ กตฺวา, “สพฺเพ สงฺขารา อนิจฺจา”ติ \textbf{ญตฺวา} ชานิตฺวา ตุลยิตฺวา ตีรยิตฺวา วิภาวยิตฺวา วิภูตํ กตฺวา, “สพฺเพ สงฺขารา ทุกฺขา”ติ. “สพฺเพ ธมฺมา อนตฺตา”ติ ฯเปฯ “ยํ กิญฺจิ สมุทยธมฺมํ, สพฺพํ ตํ นิโรธธมฺมนฺ”ติ \textbf{ญตฺวา} ชานิตฺวา ตุลยิตฺวา ตีรยิตฺวา วิภาวยิตฺวา วิภูตํ กตฺวา.\csromanspacer{}\csromanspacer{}\textbf{อนิสฺสิโต}ติ เทฺว นิสฺสยา \textbf{ตณฺหา}นิสฺสโย จ ทิฏฺฐินิสฺสโย จ ฯเปฯ อยํ \textbf{ตณฺหา}นิสฺสโย ฯเปฯ อยํ ทิฏฺฐินิสฺสโย.\csromanspacer{} ตณฺหานิสฺสยํ ปหาย ทิฏฺฐินิสฺสยํ ปฏินิสฺสชฺชิตฺวา จกฺขุํ อนิสฺสิโต.\csromanspacer{} โสตํ อนิสฺสิโต.\csromanspacer{} ฆานํ อนิสฺสิโต.\csromanspacer{} ชิวฺหํ อนิสฺสิโต.\csromanspacer{} กายํ อนิสฺสิโต.\csromanspacer{} มนํ อนิสฺสิโต.\csromanspacer{} รูเป.\csromanspacer{} สทฺเท.\csromanspacer{} คนฺเธ.\csromanspacer{} รเส.\csromanspacer{} โผฏฺฐพฺเพ.\csromanspacer{} กุลํ.\csromanspacer{} คณํ.\csromanspacer{} อาวาสํ ฯเปฯ.\csromanspacer{} ทิฏฺฐสุตมุตวิญฺญาตพฺเพ ธมฺเม อนิสฺสิโต อนลฺลีโน อนุปคโต อนชฺโฌสิโต อนธิมุตฺโต นิกฺขนฺโต นิสฺสโฏ วิปฺปมุตฺโต วิสญฺญุตฺโต วิมริยาทิกเตน เจตสา วิหรตีติ \textbf{ญตฺวา} ธมฺมํ อนิสฺสิโต.}

\prose{\textbf{ภวาย วิภวาย วา, ตณฺหา ยสฺส น วิชฺชตี}ติ \textbf{ตณฺหา}ติ รูป\textbf{ตณฺหา} สทฺท\textbf{ตณฺหา} คนฺธ\textbf{ตณฺหา} รส\textbf{ตณฺหา} โผฏฺฐพฺพ\textbf{ตณฺหา} ธมฺม\textbf{ตณฺหา}.\csromanspacer{} \textbf{ยสฺสา}ติ อรหโต ขีณาสวสฺส.\csromanspacer{} \textbf{ภวายา}ติ ภวทิฏฺฐิยา.\csromanspacer{} \textbf{วิภวายา}ติ วิภวทิฏฺฐิยา.\csromanspacer{} \textbf{ภวายา}ติ สสฺสตทิฏฺฐิยา.\csromanspacer{} \textbf{วิภวายา}ติ อุจฺเฉททิฏฺฐิยา.\csromanspacer{} \textbf{ภวายา}ติ ปุนปฺปุนภวาย ปุนปฺปุนคติยา ปุนปฺปุน-อุปปตฺติยา ปุนปฺปุนปฏิสนฺธิยา ปุนปฺปุน-อตฺตภาวาภินิพฺพตฺติยา.\csromanspacer{} ตณฺหา ยสฺส นตฺถิ น \csromanfolio{190}สนฺติ น สํวิชฺชติ นุปลพฺภติ, ปหี\textbf{นา} สมุจฺฉินฺ\textbf{นา} วูปสนฺตา ปฏิปสฺสทฺธา อภพฺพุปฺปตฺติกา ญาณคฺคิ\textbf{นา} ทฑฺฒาติ ภวาย วิภวาย วา, ตณฺหา ยสฺส น วิชฺชติ.\csromanspacer{} เต\textbf{นา}ห ภควา\csromandash{}}

\csromangathameasure{ยสฺส นิสฺสยตา นตฺถิ, ญตฺวา ธมฺมํ อนิสฺสิโต. \\ ภวาย วิภวาย วา, ตณฺหา ยสฺส น วิชฺชตีติ.}
\csromangathagroup{\csromangathastanza{ยสฺส นิสฺสยตา นตฺถิ, ญตฺวา ธมฺมํ อนิสฺสิโต. \\ ภวาย วิภวาย วา, ตณฺหา ยสฺส น วิชฺชตีติ.}}

\proseitem{92}{\csromansymbolfootnote{*}{ขุ 1. 412 ปิฏฺเฐ.}\textbf{ตํ พฺรูมิ อุปสนฺโต}ติ, กาเมสุ อนเปกฺขินํ.}

\csromangathameasure{คนฺถา ตสฺส น วิชฺชนฺติ, อตรี โส วิสตฺติกํ.}
\csromangathagroup{\csromangathastanza{คนฺถา ตสฺส น วิชฺชนฺติ, อตรี โส วิสตฺติกํ.}}

\prose{\textbf{ตํ พฺรูมิ อุปสนฺโต}ติ อุปสนฺโต วูปสนฺโต นิพฺพุโต ปฏิปสฺสทฺโธติ ตํ พฺรูมิ ตํ กเถมิ ตํ ภณามิ ตํ ทีปยามิ ตํ โวหรามีติ ตํ พฺรูมิ อุปสนฺโตติ.}

\prose{\textbf{กาเมสุ อนเปกฺขินนฺ}ติ \textbf{กามา}ติ อุทฺทานโต เทฺว \textbf{กามา} วตฺถุ\textbf{กามา} จ กิเลส\textbf{กามา} จ ฯเปฯ อิเม วุจฺจนฺติ วตฺถุ\textbf{กามา} ฯเปฯ อิเม วุจฺจนฺติ กิเลส\textbf{กามา}.\csromanspacer{} วตฺถุกาเม ปริชานิตฺวา กิเลสกาเม ปหาย ปชหิตฺวา วิโนเทตฺวา พฺยนฺติํ กริตฺวา อนภาวํ คเมตฺวา กาเมสุ อนเปกฺขิโน วีตกาโม จตฺตกาโม วนฺตกาโม มุตฺตกาโม ปหีนกาโม ปฏินิสฺสฏฺฐกาโม, กาเมสุ วีตราโค วิคตราโค จตฺตราโค วนฺตราโค มุตฺตราโค ปหีนราโค ปฏินิสฺสฏฺฐราโค นิจฺฉาโต นิพฺพุโต สีติภูโต สุขปฺปฏิสํเวที พฺรหฺมภูเตน อตฺตนา วิหรตีติ กาเมสุ อนเปกฺขินํ.}

\prose{\textbf{คนฺถา ตสฺส น วิชฺชนฺตี}ติ \textbf{คนฺถา}ติ จตฺตาโร \textbf{คนฺถา} อภิชฺฌากายคนฺโถ พฺยาปาโท กายคนฺโถ สีลพฺพตปรามาโส กายคนฺโถ อิทํสจฺจาภินิเวโส กายคนฺโถ.\csromanspacer{} อตฺตโน ทิฏฺฐิยา ราโค อภิชฺฌากายคนฺโถ, ปรวาเทสุ อาฆาโต อปฺปจฺจโย พฺยาปาโท กายคนฺโถ, อตฺตโน สีลํ วา วตํ วา สีลพฺพตํ วา ปรามาโส สีลพฺพตปรามาโส กายคนฺโถ, อตฺตโน ทิฏฺฐิ อิทํสจฺจาภินิเวโส กายคนฺโถ.\csromanspacer{}\csromanspacer{}\textbf{ตสฺสา}ติ อรหโต ขีณาสวสฺส.\csromanspacer{} \textbf{คนฺถา ตสฺส น วิชฺชนฺตี}ติ \textbf{คนฺถา} ตสฺส นตฺถิ น สนฺติ น สํวิชฺชนฺติ นุปลพฺภนฺติ, ปหีนา สมุจฺฉินฺนา วูปสนฺตา ปฏิปสฺสทฺธา อภพฺพุปฺปตฺติกา ญาณคฺคินา ทฑฺฒาติ \textbf{คนฺถา} ตสฺส น วิชฺชนฺติ.}

\prose{\textbf{อตรี โส วิสตฺติกนฺ}ติ \textbf{วิสตฺติกา} วุจฺจติ ตณฺหา.\csromanspacer{} โย ราโค สาราโค ฯเปฯ อภิชฺฌา โลโภ อกุสลมูลํ.\csromanspacer{} \textbf{วิสตฺติกา}ติ เกนฏฺเฐน}

\vspace{\tipitakaparskip}
\csromancenter{ปุราเภทสุตฺตนิทฺเทส วิสตฺติกา.\csromanspacer{} วิสตาติ วิสตฺติกา, วิสาลาติ วิสตฺติกา, วิสฏาติ วิสตฺติกา, วิสมาติ วิสตฺติกา, วิสกฺกตีติ วิสตฺติกา, วิสํหรตีติ วิสตฺติกา, วิสํวาทิกาติ วิสตฺติกา, วิสมูลาติ วิสตฺติกา, วิสผลาติ วิสตฺติกา, วิสปริโภคาติ วิสตฺติกา, วิสาลา วา ปน สา ตณฺหา รูเป.\csromanspacer{} สทฺเท.\csromanspacer{} คนฺเธ.\csromanspacer{} รเส.\csromanspacer{} โผฏฺฐพฺเพ.\csromanspacer{} กุเล.\csromanspacer{} คเณ.\csromanspacer{} อาวาเส ฯเปฯ.\csromanspacer{} ทิฏฺฐสุตมุตวิญฺญาตพฺเพสุ ธมฺเมสุ วิสฏา วิตฺถตาติ วิสตฺติกา.\csromanspacer{}\csromanspacer{}\textbf{อตรี โส วิสตฺติกนฺ}ติ โส อิมํ วิสตฺติกํ ตณฺหํ อตริ อุตฺตริ ปตริ สมติกฺกมิ วีติวตฺตีติ อตรี โส วิสตฺติกํ.\csromanspacer{} เต\textbf{นา}ห ภควา\csromandash{}}
\vspace{0.5\baselineskip}
\csromangathameasure{ตํ พฺรูมิ อุปสนฺโตติ, กาเมสุ อนเปกฺขินํ. \\ คนฺถา ตสฺส น วิชฺชนฺติ, อตรี โส วิสตฺติกนฺติ.}
\csromangathagroup{\csromangathastanza{\csromanfolio{191}ตํ พฺรูมิ อุปสนฺโตติ, กาเมสุ อนเปกฺขินํ. \\ คนฺถา ตสฺส น วิชฺชนฺติ, อตรี โส วิสตฺติกนฺติ.}}

\proseitem{93}{\csromansymbolfootnote{*}{ขุ 1. 412 ปิฏฺเฐ.}น ตสฺส ปุตฺตา ปสโว, เขตฺตํ วตฺถุญฺจ วิชฺชติ.}

\csromangathameasure{อตฺตา วาปิ นิรตฺตา วา, น ตสฺมิํ อุปลพฺภติ.}
\csromangathagroup{\csromangathastanza{อตฺตา วาปิ นิรตฺตา วา, น ตสฺมิํ อุปลพฺภติ.}}

\prose{\textbf{น ตสฺส ปุตฺตา ปสโว, เขตฺตํ วตฺถุญฺจ วิชฺชตี}ติ \textbf{นา}ติ ปฏิกฺเขโป.\csromanspacer{} \textbf{ตสฺสา}ติ อรหโต ขีณาสวสฺส.\csromanspacer{} \textbf{ปุตฺตา}ติ จฺ\textbf{อตฺตา}โร ปุตฺตา อตฺตโช ปุตฺโต, เขตฺตโช ปุตฺโต, ทินฺนโก ปุตฺโต, อนฺเตวาสิโก ปุตฺโต.\csromanspacer{} \textbf{ปสโว}ติ อเชฬกา กุกฺกุฏสูกรา หตฺถิคาวาสฺสวฬวา.\csromanspacer{} \textbf{เขตฺตนฺ}ติ สาลิเขตฺตํ วีหิเขตฺตํ มุคฺคเขตฺตํ มาสเขตฺตํ ยวเขตฺตํ โคธุมเขตฺตํ ติลเขตฺตํ.\csromanspacer{} \textbf{วตฺถุ}นฺติ ฆรวตฺถุํ โกฏฺฐวตฺถุํ ปุเรวตฺถุํ ปจฺฉาวตฺถุํ อารามวตฺถุํ วิหารวตฺถุํ.\csromanspacer{}\csromanspacer{}\textbf{น ตสฺส ปุตฺตา ปสโว, เขตฺตํ วตฺถุญฺจ วิชฺชตี}ติ ตสฺส ปุตฺตปริคฺคโห วา ปสุปริคฺคโห วา เขตฺตปริคฺคโห วา วตฺถุปริคฺคโห วา นตฺถิ น สนฺติ น สํวิชฺชนฺติ นุปลพฺภนฺติ, ปหี\textbf{นา} สมุจฺฉินฺ\textbf{นา} วูปสนฺตา ปฏิปสฺสทฺธา อภพฺพุปฺปตฺติกา ญาณคฺคิ\textbf{นา} ทฑฺฒาติ น ตสฺส ปุตฺตา ปสโว, เขตฺตํ วตฺถุญฺจ วิชฺชติ.}

\prose{\textbf{อตฺตา วาปิ นิรตฺตา วา, น ตสฺมิํ อุปลพฺภตี}ติ \textbf{อตฺตา}ติ อตฺตทิฏฺฐิ.\csromanspacer{} \textbf{นิรตฺตา}ติ อุจฺเฉททิฏฺฐิ.\csromanspacer{} \textbf{อตฺตา}ติ คหิตํ นตฺถิ.\csromanspacer{} \textbf{นิรตฺตา}ติ มุญฺจิตพฺพํ นตฺถิ.\csromanspacer{} ยสฺส นตฺถิ คหิตํ, ตสฺส นตฺถิ มุญฺจิตพฺพํ.\csromanspacer{} ยสฺส นตฺถิ มุญฺจิตพฺพํ, ตสฺส นตฺถิ คหิตํ.\csromanspacer{} คาหมุญฺจนสมติกฺกนฺโต อรหา วุทฺธิปริหานิวีติวตฺโต, โส วุฏฺฐวาโส จิณฺณจรโณ ฯเปฯ.\csromanspacer{} ชาติมรณสํสาโร นตฺถิ ตสฺส ปุนพฺภโวติ \textbf{อตฺตา} วาปิ นิรฺ\textbf{อตฺตา} วา, น ตสฺมิํ อุปลพฺภติ.\csromanspacer{} เต\textbf{นา}ห ภควา\csromandash{}}

\csromangathameasure{น ตสฺส ปุตฺตา ปสโว, เขตฺตํ วตฺถุญฺจ วิชฺชติ. \\ อตฺตา วาปิ นิรตฺตา วา, น ตสฺมิํ อุปลพฺภตีติ.}
\csromangathagroup{\csromangathastanza{\csromanfolio{192}น ตสฺส ปุตฺตา ปสโว, เขตฺตํ วตฺถุญฺจ วิชฺชติ. \\ อตฺตา วาปิ นิรตฺตา วา, น ตสฺมิํ อุปลพฺภตีติ.}}

\proseitem{94}{\csromansymbolfootnote{*}{ขุ 1. 412 ปิฏฺเฐ.}\textbf{เยน นํ วชฺชุํ ปุถุชฺชนา, อโถ สมณพฺราหฺมณา}.}

\csromangathameasure{ตํ ตสฺส อปุรกฺขตํ, ตสฺมา วาเทสุ เน’ชติ.}
\csromangathagroup{\csromangathastanza{ตํ ตสฺส อปุรกฺขตํ, ตสฺมา วาเทสุ เน’ชติ.}}

\prose{\textbf{เยน นํ วชฺชุํ ปุถุชฺชนา, อโถ สมณพฺราหฺมณา}ติ \textbf{ปุถุชฺชนา}ติ ปุถุ กิเลเส ชเนนฺตีติ \textbf{ปุถุชฺชนา}, ปุถุ อวิหตสกฺกายทิฏฺฐิกาติ \textbf{ปุถุชฺชนา}, ปุถุ สตฺถารานํ มุขุลฺโลกิกาติ\footnote{มุขุลฺโลกกาติ (สี) เหฏฺฐา 113 ปิฏฺเฐปิ.} \textbf{ปุถุชฺชนา}, ปุถุ สพฺพคตีหิ อวุฏฺฐิตาติ \textbf{ปุถุชฺชนา}, ปุถุ นานาภิสงฺขาเร อภิสงฺขโรนฺตีติ \textbf{ปุถุชฺชนา}, ปุถุ นานา-โอเฆหิ วุยฺหนฺตีติ \textbf{ปุถุชฺชนา}, ปุถุ นานาสนฺตาเปหิ สนฺตปฺเปนฺตีติ \textbf{ปุถุชฺชนา}, ปุถุ นานาปริฬาเหหิ ปริฑยฺหนฺตีติ \textbf{ปุถุชฺชนา}, ปุถุ ปญฺจสุ กามคุเณสุ รตฺตา คิทฺธา คธิตา มุจฺฉิตา อชฺโฌสนฺนา ลคฺคา ลคฺคิตา ปลิพุทฺธาติ \textbf{ปุถุชฺชนา}, ปุถุ ปญฺจหิ นีวรเณหิ อาวุตา นิวุตา โอวุตา ปิหิตา ปฏิจฺฉนฺนา ปฏิกุชฺชิตาติ \textbf{ปุถุชฺชนา}.\csromanspacer{} \textbf{สมณา}ติ เย เกจิ อิโต พหิทฺธา ปริพฺพชูปคตา ปริพฺพชสมาปนฺนา.\csromanspacer{} \textbf{พฺราหฺมณา}ติ เย เกจิ โภวาทิกา.\csromanspacer{}\csromanspacer{}\textbf{เยน นํ วชฺชุํ ปุถุชฺชนา, อโถ สมณพฺราหฺมณา}ติ \textbf{ปุถุชฺชนา} เยน ตํ ราเคน วเทยฺยุํ, เยน โทเสน วเทยฺยุํ, เยน โมเหน วเทยฺยุํ, เยน มาเนน วเทยฺยุํ, ยาย ทิฏฺฐิยา วเทยฺยุํ, เยน อุทฺธจฺเจน วเทยฺยุํ, ยาย วิจิกิจฺฉาย วเทยฺยุํ, เยหิ อนุสเยหิ วเทยฺยุํ “รตฺโต”ติ วา “ทุฏฺโฐ”ติ วา “มูโฬฺห”ติ วา “วินิพทฺโธ”ติ วา “ปรามฏฺโฐ”ติ วา “วิกฺเขปคโต”ติ วา “อนิฏฺฐงฺคโต”ติ วา “ถามคโต”ติ วา.\csromanspacer{} เต อภิสงฺขารา ปหีนา, อภิสงฺขารานํ ปหีนตฺตา คติยา\footnote{คติโย (สฺยา)} เยน ตํ วเทยฺยุํ “เนรยิโก”ติ วา “ติรจฺฉานโยนิโก”ติ วา “เปตฺติวิสยิโก”ติ วา “มนุสฺโส”ติ วา “เทโว”ติ วา “รูปี”ติ วา “อรูปี”ติ วา “สญฺญี”ติ วา “อสญฺญี”ติ วา “เนวสญฺญีนาสญฺญี”ติ วา.\csromanspacer{} โส เหตุ นตฺถิ, ปจฺจโย นตฺถิ, การณํ นตฺถิ.\csromanspacer{} เยน นํ วเทยฺยุํ กเถยฺยุํ ภเณยฺยุํ ทีปเยยฺยุํ โวหเรยฺยุนฺติ เยน นํ วชฺชุํ \textbf{ปุถุชฺชนา}, อโถ สมณพฺราหฺมณา.}

\prose{\textbf{ตํ ตสฺส อปุรกฺขตนฺ}ติ \textbf{ตสฺสา}ติ อรหโต ขีณาสวสฺส \textbf{ปุเรกฺขารา}ติ เทฺว \textbf{ปุเรกฺขารา} ตณฺหาปุเรกฺขาโร จ ทิฏฺฐิปุเรกฺขาโร จ ฯเปฯ}

\vspace{\tipitakaparskip}
\csromancenter{ปุราเภทสุตฺตนิทฺเทส อยํ ตณฺหาปุเรกฺขาโร ฯเปฯ อยํ ทิฏฺฐิปุเรกฺขาโร.\csromanspacer{} ตสฺส ตณฺหาปุเรกฺขาโร ปหีโน, ทิฏฺฐิปุเรกฺขาโร ปฏินิสฺสฏฺโฐ, ตณฺหาปุเรกฺขารสฺส ปหีนตฺตา, ทิฏฺฐิปุเรกฺขารสฺส ปฏินิสฺสฏฺฐตฺตา น ตณฺหํ วา ทิฏฺฐิํ วา ปุรโต กตฺวา จรติ, น ตณฺหาธโช น ตณฺหาเกตุ น ตณฺหาธิปเตยฺโย น ทิฏฺฐิธโช น ทิฏฺฐิเกตุ น ทิฏฺฐาธิปเตยฺโย น ตณฺหาย วา น ทิฏฺฐิยา วา ปริวาริโต จรตีติ ตํ ตสฺส อปุรกฺขตํ.}
\prose{\csromanfolio{193}\textbf{ตสฺมา วาเทสุ เน’ชตี}ติ \textbf{ตสฺมา}ติ \textbf{ตสฺมา} ตํการณา ตํเหตุ ตปฺปจฺจยา ตนฺนิทานา วาเทสุ อุปวาเทสุ นินฺทาย ครหาย อกิตฺติยา อวณฺณหาริกาย เน’ชติ น อิญฺชติ น จลติ น เวธติ นปฺปเวธติ น สมฺปเวธตีติ \textbf{ตสฺมา} วาเทสุ เน’ชติ.\csromanspacer{} เตนาห ภควา\csromandash{}}

\csromangathameasure{เยน นํ วชฺชุํ ปุถุชฺชนา, อโถ สมณพฺราหฺมณา. \\ ตํ ตสฺส อปุรกฺขตํ, ตสฺมา วาเทสุ เน’ชตีติ.}
\csromangathagroup{\csromangathastanza{เยน นํ วชฺชุํ ปุถุชฺชนา, อโถ สมณพฺราหฺมณา. \\ ตํ ตสฺส อปุรกฺขตํ, ตสฺมา วาเทสุ เน’ชตีติ.}}

\proseitem{95}{\csromansymbolfootnote{*}{ขุ 1. 412 ปิฏฺเฐ.}\textbf{วีตเคโธ อมจฺฉรี}, น อุสฺเสสุ วทเต มุนิ.}

\csromangathameasure{น สเมสุ น โอเมสุ, กปฺปํ เน’ติ อกปฺปิโย.}
\csromangathagroup{\csromangathastanza{น สเมสุ น โอเมสุ, กปฺปํ เน’ติ อกปฺปิโย.}}

\prose{\textbf{วีตเคโธ อมจฺฉรี}ติ \textbf{เคโธ} วุจฺจติ ตณฺหา.\csromanspacer{} โย ราโค สาราโค ฯเปฯ อภิชฺฌา โลโภ อกุสลมูลํ.\csromanspacer{} ยสฺเสโส \textbf{เคโธ} ปหีโน สมุจฺฉินฺโน วูปสนฺโต ปฏิปสฺสทฺโธ อพพฺพุปฺปตฺติโก ญาณคฺคินา ทฑฺโฒ.\csromanspacer{} โส วุจฺจติ วีต\textbf{เคโธ}.\csromanspacer{} โส รูเป อคิทฺโธ ฯเปฯ ทิฏฺฐสุตมุตวิญฺญาตพฺเพสุ ธมฺเมสุ อคิทฺโธ อคธิโต อมุจฺฉิโต อนชฺโฌสิโต วีต\textbf{เคโธ} วิคต\textbf{เคโธ} จตฺต\textbf{เคโธ} วนฺต\textbf{เคโธ} มุตฺต\textbf{เคโธ} ปหีน\textbf{เคโธ} ปฏินิสฺสฏฺฐ\textbf{เคโธ} นิจฺฉาโต ฯเปฯ พฺรหฺมภูเตน อตฺตนา วิหรตีติ วีต\textbf{เคโธ}.\csromanspacer{}\csromanspacer{}\textbf{อมจฺฉรี}ติ \textbf{มจฺฉริยนฺ}ติ ปญฺจ มจฺฉริยานิ อาวาสมจฺฉริยํ กุลมจฺฉริยํ ลาภมจฺฉริยํ วณฺณมจฺฉริยํ ธมฺมมจฺฉริยํ.\csromanspacer{} ยํ เอวรูปํ ฯเปฯ คาโห, อิทํ วุจฺจติ มจฺฉริยํ.\csromanspacer{} ยสฺเสตํ มจฺฉริยํ ปหีนํ สมุจฺฉินฺนํ วูปสนฺตํ ปฏิปสฺสทฺธํ อภพฺพุปฺปตฺติกํ ญาณคฺคินา ทฑฺฒํ, โส วุจฺจติ อมจฺฉรีติ วีต\textbf{เคโธ} อมจฺฉรี.}

\prose{\textbf{น อุสฺเสสุ วทเต มุนิ.}\csromanspacer{}\textbf{ น สเมสุ น โอเมสู}ติ \textbf{มุนี}ติ \textbf{โมนํ} วุจฺจติ ญาณํ ฯเปฯ สงฺคชาลมติจฺจ โส มุนิ. “เสยฺโยหมสฺมี”ติ วา “สทิโสหมสฺมี”ติ วา “หีโนหมสฺมี”ติ วา น วทติ น กเถติ น ภณติ น ทีปยติ น โวหรตีติ น อุสฺเสสุ วทเต มุนิ.\csromanspacer{} น สเมสุ น โอเมสุ.}

\prose{\csromanfolio{194}\textbf{กปฺปํ เน’ติ อกปฺปิโย}ติ \textbf{กปฺปา}ติ เทฺว \textbf{กปฺปา} ตณฺหากปฺโป จ ทิฏฺฐิกปฺโป จ ฯเปฯ อยํ ตณฺหากปฺโป ฯเปฯ อยํ ทิฏฺฐิกปฺโป.\csromanspacer{} ตสฺส ตณฺหากปฺโป ปหีโน, ทิฏฺฐิกปฺโป ปฏินิสฺสฏฺโฐ, ตณฺหากปฺปสฺส ปหีนตฺตา, ทิฏฺฐิกปฺปสฺส ปฏินิสฺสฏฺฐตฺตา ตณฺหากปฺปํ วา ทิฏฺฐิกปฺปํ วา เน’ติ น อุเปติ น อุปคจฺฉติ น คณฺหาติ น ปรามสติ นาภินิวิสตีติ กปฺปํ เน’ติ.\csromanspacer{} \textbf{อกปฺปิโย}ติ \textbf{กปฺปา}ติ เทฺว\textbf{กปฺปา} ตณฺหากปฺโป จ ทิฏฺฐิกปฺโป จ ฯเปฯ อยํ ตณฺหากปฺโป ฯเปฯ อยํ ทิฏฺฐิกปฺโป.\csromanspacer{} ตสฺส ตณฺหากปฺโป ปหีโน, ทิฏฺฐิกปฺโป ปฏินิสฺสฏฺโฐ, ตสฺส ตณฺหากปฺปสฺส ปหีนตฺตา ทิฏฺฐิกปฺปสฺส ปฏินิสฺสฏฺฐตฺตา ตณฺหากปฺปํ วา ทิฏฺฐิกปฺปํ วา น กปฺเปติ น ชเนติ น สญฺชเนติ น นิพฺพตฺเตติ นาภินิพฺพตฺเตตีติ กปฺปํ เน’ติ อกปฺปิโย.\csromanspacer{} เตนาห ภควา\csromandash{}}

\csromangathameasure{วีตเคโธ อมจฺฉรี, น อุสฺเสสุ วทเต มุนิ. \\ น สเมสุ น โอเมสุ, กปฺปํ เน’ติ อกปฺปิโยติ.}
\csromangathagroup{\csromangathastanza{วีตเคโธ อมจฺฉรี, น อุสฺเสสุ วทเต มุนิ. \\ น สเมสุ น โอเมสุ, กปฺปํ เน’ติ อกปฺปิโยติ.}}

\proseitem{96}{\csromansymbolfootnote{*}{ขุ 1. 412 ปิฏฺเฐ.}ยสฺส โลเก สกํ นตฺถิ, อสตา จ น โสจติ.}

\csromangathameasure{ธมฺเมสุ จ น คจฺฉติ, ส เว สนฺโตติ วุจฺจติ.}
\csromangathagroup{\csromangathastanza{ธมฺเมสุ จ น คจฺฉติ, ส เว สนฺโตติ วุจฺจติ.}}

\prose{\textbf{ยสฺส โลเก สกํ นตฺถี}ติ \textbf{ยสฺสา}ติ อรหโต ขีณาสวสฺส.\csromanspacer{} \textbf{โลเก สกํ นตฺถี}ติ ตสฺส มยฺหํ วา อิทํ, ปเรสํ วา อิทนฺติ กิญฺจิ รูปคตํ เวทนาคตํ สญฺญาคตํ สงฺขารคตํ วิญฺญาณคตํ คหิตํ ปรามฏฺฐํ อภินิวิฏฺฐํ อชฺโฌสิตํ อธิมุตฺตํ นตฺถิ น สนฺติ ฯเปฯ ญาณคฺคินา ทฑฺฒนฺติ ยสฺส โลเก สกํ นตฺถิ.\csromanspacer{}\csromanspacer{}\textbf{อสตา จ น โสจตี}ติ วิปริณตํ วา วตฺถุํ น โสจติ, วิปริณตสฺมิํ วา วตฺถุสฺมิํ น โสจติ. “จกฺขุ เม วิปริณตนฺ”ติ น โสจติ, “โสตํ เม.\csromanspacer{} ฆานํ เม.\csromanspacer{} ชิวฺหา เม.\csromanspacer{} กาโย เม.\csromanspacer{} มโน เม.\csromanspacer{} รูปา เม.\csromanspacer{} สทฺทา เม.\csromanspacer{} คนฺธา เม.\csromanspacer{} รสา เม.\csromanspacer{} โผฏฺฐพฺพา เม.\csromanspacer{} กุลํ เม.\csromanspacer{} คโณ เม.\csromanspacer{} อาวาโส เม.\csromanspacer{} ลโภ เม.\csromanspacer{} ยโส เม.\csromanspacer{} ปสํสา เม.\csromanspacer{} สุขํ เม.\csromanspacer{} จีวรํ เม.\csromanspacer{} ปิณฺฑปาโต เม.\csromanspacer{} เสนาสนํ เม.\csromanspacer{} คิลานปจฺจยเภสชฺชปริกฺขาโร เม.\csromanspacer{} มาตา เม.\csromanspacer{} ปิตา เม.\csromanspacer{} ภาตา เม.\csromanspacer{} ภคินี เม.\csromanspacer{} ปุตฺโต เม.\csromanspacer{} ธีตา เม.\csromanspacer{} มิตฺตาเม.\csromanspacer{} อมจฺจา เม.\csromanspacer{} ญาตกา เม.\csromanspacer{} สาโลหิตา เม วิปริณตา”ติ น โสจติ น กิลมติ น ปริเทวติ น อุรตฺตาฬิํ กนฺทติ น สมฺโมหํ อาปชฺชตีติ เอวมฺปิ อสตา จ น โสจติ.}

\prose{อถ วา อสนฺตาย\footnote{อสตาย (สี), อสาตาย (สฺยา)} ทุกฺขาย เวทนาย ผุฏฺโฐ ปเรโต สโมหิโต สมนฺนาคโต น โสจติ น กิลมติ น ปริเทวติ น}

\vspace{\tipitakaparskip}
\csromancenter{ปุราเภทสุตฺตนิทฺเทส อุรตฺตาฬิํ กนฺทติ น สมฺโมหํ อาปชฺชติ.\csromanspacer{} จกฺขุโรเคน ผุฏฺโฐ ปเรโต สโมหิโต สมนฺนาคโต น โสจติ น กิลมติ น ปริเทวติ น อุรตฺตาฬิํ กนฺทติ น สมฺโมหํ อาปชฺชติ.\csromanspacer{} โสตโรเคน.\csromanspacer{} ฆานโรเคน.\csromanspacer{} ชิวฺหาโรเคน.\csromanspacer{} กายโรเคน.\csromanspacer{} สีสโรเคน.\csromanspacer{} กณฺณโรเคน.\csromanspacer{} มุขโรเคน.\csromanspacer{} ทนฺตโรเคน.\csromanspacer{} กาเสน.\csromanspacer{} สาเสน.\csromanspacer{} ปินาเสน.\csromanspacer{} ทาเหน.\csromanspacer{} ชเรน.\csromanspacer{} กุจฺฉิโรเคน.\csromanspacer{} มุจฺฉาย.\csromanspacer{} ปกฺขนฺทิกาย.\csromanspacer{} สูเลน.\csromanspacer{} วิสูจิกาย.\csromanspacer{} กุฏฺเฐน.\csromanspacer{} คณฺเฑน.\csromanspacer{} กิลาเสน.\csromanspacer{} โสเสน.\csromanspacer{} อปมาเรน.\csromanspacer{} ททฺทุยา.\csromanspacer{} กณฺฑุยา.\csromanspacer{} กจฺฉุยา.\csromanspacer{} รขสาย.\csromanspacer{} วิตจฺฉิกาย.\csromanspacer{} โลหิเตน.\csromanspacer{} ปิตฺเตน.\csromanspacer{} มธุเมเหน.\csromanspacer{} อํสาย.\csromanspacer{} ปิฬกาย.\csromanspacer{} ภคนฺทลาย\footnote{ภคนฺทเลน (ก)}.\csromanspacer{} ปิตฺตสมุฏฺฐาเนน อาพาเธน.\csromanspacer{} เสมฺหสมุฏฺฐาเนน อาพาเธน.\csromanspacer{} วาตสมุฏฺฐาเนน อาพาเธน.\csromanspacer{} สนฺนิปาติเกน อาพาเธน.\csromanspacer{} อุตุปริณามเชน อาพาเธน.\csromanspacer{} วิสมปริหารเชน อาพาเธน.\csromanspacer{} โอปกฺกมิเกน อาพาเธน.\csromanspacer{} กมฺมวิปากเชน อาพาเธน.\csromanspacer{} สีเตน.\csromanspacer{} อุเณฺหน.\csromanspacer{} ชิฆจฺฉาย.\csromanspacer{} ปิปาสาย.\csromanspacer{} qอํสมกสวาตาตปสรีสปสมฺผสฺเสหิ ผุฏฺโฐ ปเรโต สโมหิโต สมนฺนาคโต น โสจติ น กิลมติ น ปริเทวติ น อุรตฺตาฬิํ กนฺทติ น สมฺโมหํ อาปชฺชตีติ เอวมฺปิ อสตา จ น โสจติ.}
\prose{\csromanfolio{195}อถ วา อสนฺเต อสํวิชฺชมาเน อนุปลพฺภมาเน\footnote{อนุปลพฺภิยมาเน (สฺยา, ก)} “อโห วต เม ตํ นตฺถิ, สิยา วต เม ตํ, ตํ วตาหํ น จ ลภามี”ติ น โสจติ น กิลมติ น ปริเทวติ น อุรตฺตาฬิํ กนฺทติ น สมฺโมหํ อาปชฺชตีติ เอวมฺปิ อสตา จ น โสจติ.\csromanspacer{} \textbf{ธมฺเมสุ จ น คจฺฉตี}ติ น ฉนฺทาคติํ คจฺฉติ, น โทสาคติํ คจฺฉติ, น โมหาคติํ คจฺฉติ, น ภยาคติํ คจฺฉติ, น ราควเสน คจฺฉติ, น โทสวเสน คจฺฉติ, น โมหวเสน คจฺฉติ, น มานวเสน คจฺฉติ, น ทิฏฺฐิวเสน คจฺฉติ, น อุทฺธจฺจวเสน คจฺฉติ, น วิจิกิจฺฉาวเสน คจฺฉติ, น อนุสยวเสน คจฺฉติ, น จ วคฺเคหิ ธมฺเมหิ ยายติ นียติ วุยฺหติ สํหรียตีติ ธมฺเมสุ จ น คจฺฉติ.}

\prose{\textbf{ส เว สนฺโตติ วุจฺจติ}ติ โส สนฺโต อุปสนฺโต วูปสนฺโต นิพฺพุโต ปฏิปสฺสทฺโธติ วุจฺจติ ปวุจฺจติ กถียติ ภณียติ ทีปียติ โวหรียตีติ ส เว สนฺโตติ วุจฺจติ.\csromanspacer{} เตนาห ภควา\csromandash{}}

\csromangathameasure{ยสฺส โลเก สกํ นตฺถิ, อสตา จ น โสจติ. \\ ธมฺเมสุ จ น คจฺฉติ, ส เว สนฺโตติ วุจฺจตีติ. ปุราเภทสุตฺตนิทฺเทโส ทสโม.}
\csromangathagroup{\csromangathastanza{\csromanfolio{196}ยสฺส โลเก สกํ นตฺถิ, อสตา จ น โสจติ.}\csromangathastanza{ธมฺเมสุ จ น คจฺฉติ, ส เว สนฺโตติ วุจฺจตีติ.\csromanspacer{} ปุราเภทสุตฺตนิทฺเทโส ทสโม.}}

\csromanheader{2}{11. กลหวิวาทสุตฺตนิทฺเทส}
\csromanheader{2}{อถ กลหวิวาทสุตฺตนิทฺเทสํ วกฺขติ\csromandash{}}
\proseitem{97}{\csromansymbolfootnote{*}{ขุ 1. 413 ปิฏฺเฐ.}\textbf{กุโตปหูตา กลหา วิวาทา},}

\prose{\textbf{ปริเทวโสกา สหมจฺฉรา จ.}}

\csromangathameasure{\textbf{มานาติมานา สหเปสุณา จ,} \\ \textbf{กุโตปหูตา เต ตทิงฺฆ พฺรูหิ.}}
\csromangathagroup{\csromangathastanza{\textbf{มานาติมานา สหเปสุณา จ,} \\ \textbf{กุโตปหูตา เต ตทิงฺฆ พฺรูหิ.}}}

\prose{\textbf{กุโตปหูตา กลหา วิวาทา}ติ \textbf{กลโห}ติ เอเกน อากาเรน \textbf{กลโห}.\csromanspacer{} \textbf{วิวาโท}ติปิ ตญฺเญว.\csromanspacer{} โย \textbf{กลโห}, โส วิวาโท, โย วิวาโท, โส \textbf{กลโห}.\csromanspacer{} อถ วา อปเรน อากาเรน วิวาโท วุจฺจติ กลหสฺส ปุพฺพภาโค วิวาโท, ราชาโนปิ ราชูหิ วิวทนฺติ, ขตฺติยาปิ ขตฺติเยหิ วิวทนฺติ, พฺราหฺมณาปิ พฺราหฺมเณหิ วิวทนฺติ, คหปตีปิ คหปตีหิ วิวทนฺติ, มาตาปิ ปุตฺเตน วิวทติ, ปุตฺโตปิ มาตรา วิวทติ, ปิตาปิ ปุตฺเตน วิวทติ, ปุตฺโตปิ ปิตรา วิวทติ, ภาตาปิ ภาตรา วิวทติ, ภาตาปิ ภคินิยา วิวทติ, ภคินีปิ ภาตรา วิวทติ, สหาโยปิ สหาเยน วิวทติ, อยํ วิวาโท.\csromanspacer{}\csromanspacer{}กตโม \textbf{กลโห}, อาคาริกา ทณฺฑปสุตา กาเยน วาจาย กลหํ กโรนฺติ, ปพฺพชิตา อาปตฺติํ อาปชฺชนฺตา กาเยน วาจาย กลหํ กโรนฺติ.\csromanspacer{} อยํ \textbf{กลโห}.}

\prose{\textbf{กุโตปหูตา กลหา วิวาทา}ติ กลหา จ วิวาทา จ กุโตปหูตา กุโตชาตา กุโตสญฺชาตา กุโตนิพฺพตฺตา กุโต-อภินิพฺพตฺตา กุโตปาตุภูตา กิํนิทานา กิํสมุทยา กิํชาติกา กิํปภวาติ กลหสฺส จ วิวาทสฺส จ มูลํ ปุจฺฉติ, เหตุํ ปุจฺฉติ, นิทานํ ปุจฺฉติ, สมฺภวํ ปุจฺฉติ, ปภวํ ปุจฺฉติ, สมุฏฺฐานํ ปุจฺฉติ, อาหารํ ปุจฺฉติ, อารมฺมณํ ปุจฺฉติ, ปจฺจยํ ปุจฺฉติ, สมุทยํ ปุจฺฉติ,}

\vspace{\tipitakaparskip}
\csromancenter{กลหวิวาทสุตฺตนิทฺเทส ปปุจฺฉติ ยาจติ อชฺเฌสติ\footnote{อชฺโฌสติ (สี)} ปสาเทตีติ กุโตปหูตา กลหา วิวาทา.}
\prose{\csromanfolio{197}\textbf{ปริเทวโสกา สหมจฺฉรา จา}ติ \textbf{ปริเทโว}ติ ญาติพฺยสเนน วา ผุฏฺฐสฺส โภคพฺยสเนน วา ผุฏฺฐสฺส โรคพฺยสเนน วา ผุฏฺฐสฺส สีลพฺยสเนน วา ผุฏฺฐสฺส ทิฏฺฐิพฺยสเนน วา ผุฏฺฐสฺส อญฺญตรญฺญตเรน วา พฺยสเนน สมนฺนาคตสฺส อญฺญตรญฺญตเรน วา ทุกฺขธมฺเมน ผุฏฺฐสฺส อาเทโว \textbf{ปริเทโว} อาเทวนา ปริเทวนา อาเทวิตตฺตํ ปริเทวิตตฺตํ วาจา ปลาโป วิปฺปลาโป ลาลปฺโป ลาลปฺปายนา ลาลปฺปายิตตฺตํ.\csromanspacer{}\csromanspacer{}\textbf{โสโก}ติ ญาติพฺยสเนน วา ผุฏฺฐสฺส โภค โรค สีล ทิฏฺฐิพฺยสเนน วา ผุฏฺฐสฺส อญฺญตรญฺญตเรน วา พฺยสเนน สมนฺนาคตสฺส อญฺญตรญฺญตเรน วา ทุกฺขธมฺเมน ผุฏฺฐสฺส โสโก โสจนา โสจิตตฺตํ อนฺโตโสโก อนฺโตปริโสโก อนฺโตทาโห อนฺโตปริทาโห เจตโส ปริชฺฌายนา โทมนสฺสํ โสกสลฺลํ.\csromanspacer{}\csromanspacer{}\textbf{มจฺฉรนฺ}ติ ปญฺจ มจฺฉริยานิ อาวาสมจฺฉริยํ กุลมจฺฉริยํ ลาภมจฺฉริยํ วณฺณมจฺฉริยํ ธมฺมมจฺฉริยํ.\csromanspacer{} ยํ เอวรูปํ มจฺฉริยํ มจฺฉรายนํ มจฺฉรายิตตฺตํ เววิจฺฉํ กทริยํ กฏุกญฺจุกตา อคฺคหิตตฺตํ จิตฺตสฺส, อิทํ วุจฺจติ มจฺฉริยํ.\csromanspacer{} อปิ จ ขนฺธมจฺฉริยมฺปิ มจฺฉริยํ, ธาตุมจฺฉริยมฺปิ มจฺฉริยํ, อายตนมจฺฉริยมฺปิ มจฺฉริยํ, คาโห อิทํ วุจฺจติ มจฺฉริยนฺติ ปริเทวโสกา สหมจฺฉรา จ.}

\prose{\textbf{มานาติมานา สหเปสุณา จา}ติ \textbf{มาโน}ติ อิเธกจฺโจ มานํ ชเนติ ชาติยา วา โคตฺเตน วา โกลปุตฺติเยน วา วณฺณโปกฺขรตาย วา ธเนน วา อชฺเฌเนน วา กมฺมายตเนน วา สิปฺปายตเนน วา วิชฺชาฏฺฐาเนน วา สุเตน วา ปฏิภาเนน วา อญฺญตรญฺญตเรน วา วตฺถุนา.\csromanspacer{} \textbf{อติมาโน}ติ อิเธกจฺโจ ปรํ อติมญฺญติ ชาติยา วา โคตฺเตน วา ฯเปฯ อญฺญตรญฺญตเรน วา วตฺถุนา.\csromanspacer{} \textbf{เปสุญฺญนฺ}ติ อิเธกจฺโจ ปิสุณาวาโจ โหติ อิโต สุตฺวา อมุตฺร อกฺขาตา อิเมสํ เภทาย, อมุตฺร วา สุตฺวา อิเมสํ อกฺขาตา อมูสํ เภทาย.\csromanspacer{} อิติ สมคฺคานํ วา เภตฺตา, ภินฺนานํ วา อนุปฺปทาตา, วคฺคาราโม วคฺครโต วคฺคนนฺที วคฺคกรณิํ วาจํ ภาสิตา โหติ, อิทํ วุจฺจติ เปสุญฺญํ.\csromanspacer{} อปิ จ ทฺวีหิ การเณหิ เปสุญฺญํ อุปสํหรติ ปิยกมฺยตาย วา เภทาธิปฺปาเยน วา.\csromanspacer{} กถํ ปิยกมฺยตาย เปสุญฺญํ อุปสํหรติ, “อิมสฺส \csromanfolio{198}ปิโย ภวิสฺสามิ, มนาโป ภวิสฺสามิ, วิสฺสาสิโก ภวิสฺสามิ, อพฺภนฺตริโก ภวิสฺสามิ, สุหทโย ภวิสฺสามี”ติ เอวํ ปิยกมฺยตาย เปสุญฺญํ อุปสํหรหิ.\csromanspacer{}\csromanspacer{}กถํ เภทาธิปฺปาเยน เปสุญฺญํ อุปสํหรติ, “กถํ อิเม นานา อสฺสุ, วินา อสฺสุ, วคฺคา อสฺสุ, ทฺวิธา อสฺสุ, เทฺวชฺฌา อสฺสุ, เทฺว ปกฺขา อสฺสุ ภิชฺเชยฺยุํ น สมาคจฺเฉยฺยุํ, ทุกฺขํ น ผาสุ วิหเรยฺยุนฺ”ติ เอวํ เภทาธิปฺปาเยน เปสุญฺญํ อุปสํหรตีติ มานาติมานา สหเปสุณา จ.}

\prose{\textbf{กุโตปหูตา เต ตทิงฺฆ พฺรูหี}ติ กลโห จ วิวาโท จ ปริเทโว จ โสโก จ มจฺฉริยญฺจ มาโน จ อติมาโน จ เปสุญฺญญฺจาติ อิเม อฏฺฐ กิเลสา กุโตปหูตา กุโตชาตา กุโตสญฺชาตา กุโตนิพฺพตฺตา กุโต-อภินิพฺพตฺตา กุโตปาตุภูตา กิํนิทานา กิํสมุทยา กิํชาติกา กิํปภวาติ อิเมสํ อฏฺฐนฺนํ กิเลสานํ มูลํ ปุจฺฉติ, เหตุํ ปุจฺฉติ, นิทานํ ปุจฺฉติ, สมฺภวํ ปุจฺฉติ, ปภวํ ปุจฺฉติ, สมุฏฺฐานํ ปุจฺฉติ, อาหารํ ปุจฺฉติ, อารมฺมณํ ปุจฺฉติ, ปจฺจยํ ปุจฺฉติ, สมุทยํ ปุจฺฉติ ปปุจฺฉติ ยาจติ อชฺเฌสติ ปสาเทตีติ กุโตปหูตา เต ตทิงฺฆ พฺรูหีติ อิงฺฆ พฺรูหิ อาจิกฺข เทเสหิ ปญฺญเปหิ ปฏฺฐเปหิ วิวร วิภช อุตฺตานีกโรหิ ปกาเสหีติ กุโตปหูตา เต ตทิงฺฆ พฺรูหิ.\csromanspacer{} เตนาห โส นิมฺมิโต\csromandash{}}

\csromangathameasure{กุโตปหูตา กลหา วิวาทา, \\ \textbf{ปริเทวโสกา สหมจฺฉรา จ.} \\ \textbf{มานาติมานา สหเปสุณา จ,}}
\csromangathagroup{\csromangathastanza{กุโตปหูตา กลหา วิวาทา, \\ \textbf{ปริเทวโสกา สหมจฺฉรา จ.} \\ \textbf{มานาติมานา สหเปสุณา จ,}}}

\prose{\textbf{กุโตปหูตา เต ตทิงฺฆ พฺรูหี}ติ.}

\csromanmatikaanchor{mtk.11}
\csromantocmarkref{subsection}{10. ปุราเภทสุตฺตนิทฺเทส}{mtk.11}
\bookmark[dest={mtk.11},level=2]{10. ปุราเภทสุตฺตนิทฺเทส}
\proseitem{98}{\csromansymbolfootnote{*}{ขุ 1. 413 ปิฏฺเฐ.}ปิยปฺปหูตา กลหา วิวาทา,}

\prose{\textbf{ปริเทวโสกา สหมจฺฉรา จ.}}

\csromangathameasure{\textbf{มานาติมานา สหเปสุณา จ,} \\ \textbf{มจฺเฉรยุตฺตา กลหา วิวาทา.} \\ \textbf{วิวาทชาเตสุ จ เปสุณานิ.}}
\csromangathagroup{\csromangathastanza{\textbf{มานาติมานา สหเปสุณา จ,} \\ \textbf{มจฺเฉรยุตฺตา กลหา วิวาทา.} \\ \textbf{วิวาทชาเตสุ จ เปสุณานิ.}}}

\prose{\textbf{ปิยปฺปหูตา กลหา วิวาทา, ปริเทวโสกา สหมจฺฉรา จา}ติ \textbf{ปิยา}ติ เทฺว \textbf{ปิยา} สตฺตา วา สงฺขารา วา.\csromanspacer{} กตเม สตฺตา \textbf{ปิยา}, อิธ}

\vspace{\tipitakaparskip}
\csromancenter{กลหวิวาทสุตฺตนิทฺเทส ยสฺส เต โหนฺติ อตฺถกามา หิตกามา ผาสุกามา โยคกฺเขมกามา มาตา วา ปิตา วา ภาตา วา ภคินี วา ปุตฺโต วา ธีตา วา มิตฺตา วา อมจฺจา วา ญาตี วา สาโลหิตา วา, อิเม สตฺตา ปิยา.\csromanspacer{} กตเม สงฺขารา ปิยา, มนาปิกา รูปา มนาปิกา สทฺทา มนาปิกา คนฺธา มนาปิกา รสา มนาปิกา โผฏฺฐพฺพา, อิเม สงฺขารา ปิยา.}
\prose{\csromanfolio{199}ปิยํ วตฺถุํ อจฺเฉทสงฺกิโนปิ กลหํ กโรนฺติ, อจฺฉิชฺชนฺเตปิ กลหํ กโรนฺติ, อจฺฉินฺเนปิ กลหํ กโรนฺติ.\csromanspacer{} ปิยํ วตฺถุํ วิปริณามสงฺกิโนปิ กลหํ กโรนฺติ, วิปริณามนฺเตปิ กลหํ กโรนฺติ, วิปริณเตปิ กลหํ กโรนฺติ.\csromanspacer{}\csromanspacer{}ปิยํ วตฺถุํ อจฺเฉทสงฺกิโนปิ วิวทนฺติ, อจฺฉิชฺชนฺเตปิ วิวทนฺติ, อจฺฉินฺเนปิ วิวทนฺติ.\csromanspacer{} ปิยํ วตฺถุํ วิปริณามสงฺกิโนปิ วิวทนฺติ, วิปริณามนฺเตปิ วิวทนฺติ, วิปริณเตปิ วิวทนฺติ.\csromanspacer{}\csromanspacer{}ปิยํ วตฺถุํ อจฺเฉทสงฺกิโนปิ ปริเทวนฺติ, อจฺฉิชฺชนฺเตปิ ปริเทวนฺติ, อจฺฉินฺเนปิ ปริเทวนฺติ.\csromanspacer{} ปิยํ วตฺถุํ วิปริณามสงฺกิโนปิ ปริเทวนฺติ, วิปริณามนฺเตปิ ปริเทวนฺติ, วิปริณเตปิ ปริเทวนฺติ.\csromanspacer{}\csromanspacer{}ปิยํ วตฺถุํ อจฺเฉทสงฺกิโนปิ โสจนฺติ, อจฺฉิชฺชนฺเตปิ โสจนฺติ, อจฺฉินฺเนปิ โสจนฺติ.\csromanspacer{} ปิยํ วตฺถุํ วิปริณามสงฺกิโนปิ โสจนฺติ, วิปริณามนฺเตปิ โสจนฺติ, วิปริณเตปิ โสจนฺติ.\csromanspacer{} ปิยํ วตฺถุํ รกฺขนฺติ โคเปนฺติ ปริคฺคณฺหนฺติ มมายนฺติ มจฺฉรายนฺติ.}

\prose{\textbf{มานาติมานา สหเปสุณา จา}ติ ปิยํ วตฺถุํ นิสฺสาย มานํ ชเนนฺติ, ปิยํ วตฺถุํ นิสฺสาย อติมานํ ชเนนฺติ.\csromanspacer{} กถํ ปิยํ วตฺถุํ นิสฺสาย มานํ ชเนนฺติ, “มยํ ลาภิโน มนาปิกานํ รูปานํ สทฺทานํ คนฺธานํ รสานํ โผฏฺฐพฺพานนฺ”ติ เอวํ ปิยํ วตฺถุํ นิสฺสาย มานํ ชเนนฺติ.\csromanspacer{} กถํ ปิยํ วตฺถุํ นิสฺสาย อติมานํ ชเนนฺติ, “มยํ ลาภิโน มนาปิกานํ รูปานํ สทฺทานํ คนฺธานํ รสานํ โผฏฺฐพฺพานํ, อิเม ปนญฺเญ น ลาภิโน มนาปิกานํ รูปานํ สทฺทานํ คนฺธานํ รสานํ โผฏฺฐพฺพานนฺ”ติ เอวํ ปิยํ วตฺถุํ นิสฺสาย อติมานํ ชเนนฺติ.\csromanspacer{} \textbf{เปสุญฺญนฺ}ติ อิเธกจฺโจ ปิสุณวาโจ โหติ อิโต สุตฺวา อมุตฺร อกฺขาตา อิเมสํ เภทาย ฯเปฯ เอวํ เภทาธิปฺปาเยน เปสุญฺญํ อุปสํหรตีติ มานาติมานา สหเปสุณา จ.}

\prose{\textbf{มจฺเฉรยุตฺตา กลหา วิวาทา}ติ กลโห จ วิวาโท จ ปริเทโว จ โสโก จ มาโน จ อติมาโน จ เปสุญฺญญฺจาติ อิเม สตฺต กิเลสา มจฺฉริเย ยุตฺตา ปยุตฺตา อายุตฺตา สมายุตฺตาติ มจฺเฉรยุตฺตา กลหา วิวาทา.}

\prose{\csromanfolio{200}\textbf{วิวาทชาเตสุ จ เปสุณานี}ติ วิวาเท ชาเต สญฺชาเต นิพฺพตฺเต อภินิพฺพตฺเต ปาตุภูเต เปสุญฺญํ อุปสํหรนฺติ อิโต สุตฺวา อมุตฺร อกฺขายนฺติ อิเมสํ เภทาย, อมุตฺร วา สุตฺวา อิเมสํ อกฺขายนฺติ อมูสํ เภทาย.\csromanspacer{} อิติ สมคฺคานํ วา เภตฺตาโร, ภินฺนานํ วา อนุปฺปทาตาโร, วคฺคารามา วคฺครตา วคฺคนนฺที วคฺคกรณิํ วาจํ ภาสิตาโร โหนฺติ, อิทํ วุจฺจติ เปสุญฺญํ.\csromanspacer{}\csromanspacer{}อปิ จ ทฺวีหิ การเณหิ เปสุญฺญํ อุปสํหรนฺติ ปิยกมฺยตาย วา เภทาธิปฺปาเยน วา.\csromanspacer{} กถํ ปิยกมฺยตาย เปสุญฺญํ อุปสํหรนฺติ, “อิมสฺส ปิยา ภวิสฺสาม, มนาปา ภวิสฺสาม, วิสฺสาสิกา ภวิสฺสาม, อพฺภนฺตริกา ภวิสฺสาม, สุหทยา ภวิสฺสามา”ติ เอวํ ปิยกมฺยตาย เปสุญฺญํ อุปสํหรนฺติ.\csromanspacer{} กถํ เภทาธิปฺปาเยน เปสุญฺญํ อุปสํหรนฺติ. “กถํ อิเม นานา อสฺสุ, วินา อสฺสุ, วคฺคา อสฺสุ, เทฺวธา อสฺสุ, เทฺวชฺฌา อสฺสุ, เทฺว ปกฺขา อสฺสุ ภิชฺเชยฺยุํ น สมาคจฺเฉยฺยุํ, ทุกฺขํ น ผาสุ วิหเรยฺยุนฺ”ติ เอวํ เภทาธิปฺปาเยน เปสุญฺญํ อุปสํหรนฺตีติ วิวาทชาเตสุ จ เปสุณานิ.\csromanspacer{} เตนาห ภควา\csromandash{}}

\csromangathameasure{ปิยปฺปหูตา กลหา วิวาทา, \\ ปริเทวโสกา สหมจฺฉรา จ. \\ มนาติมานา สหเปสุณา จ, \\ มจฺเฉรยุตฺตา กลหา วิวาทา. \\ วิวาทชาเตสุ จ เปสุณานีติ.}
\csromangathagroup{\csromangathastanza{ปิยปฺปหูตา กลหา วิวาทา, \\ ปริเทวโสกา สหมจฺฉรา จ. \\ มนาติมานา สหเปสุณา จ, \\ มจฺเฉรยุตฺตา กลหา วิวาทา.}\csromangathastanza{วิวาทชาเตสุ จ เปสุณานีติ.}}

\proseitem{99}{\csromansymbolfootnote{*}{ขุ 1. 413 ปิฏฺเฐ.}\textbf{ปิยา สุ โลกสฺมิํ กุโตนิทานา},}

\prose{\textbf{เย จาปิ โลภา วิจรนฺติ โลเก.}}

\csromangathameasure{\textbf{อาสา จ นิฏฺฐา จ กุโตนิทานา,} \\ \textbf{เย สมฺปรายาย นรสฺส โหนฺติ.}}
\csromangathagroup{\csromangathastanza{\textbf{อาสา จ นิฏฺฐา จ กุโตนิทานา,} \\ \textbf{เย สมฺปรายาย นรสฺส โหนฺติ.}}}

\prose{\textbf{ปิยา สุ โลกสฺมิํ กุโตนิทานา}ติ ปิยา กุโตนิทานา กุโตชาตา กุโตสญฺชาตา กุโตนิพฺพตฺตา กุโต-อภินิพฺพตฺตา กุโตปาตุภูตา, กิํนิทานา กิํสมุทยา กิํชาติกา กิํปภวาติ ปิยานํ มูลํ ปุจฺฉติ ฯเปฯ สมุทยํ ปุจฺฉติ ปปุจฺฉติ ยาจติ อชฺเฌสติ ปสาเทตีติ ปิยา สุ โลกสฺมิํ กุโตนิทานา.}

\prose{\textbf{เย จาปิ โลภา วิจรนฺติ โลเก}ติ \textbf{เย จาปี}ติ ขตฺติยา จ พฺราหฺมณา จ เวสฺสา จ สุทฺทา จ คหฏฺฐา จ ปพฺพชิตา จ เทวา จ}

\vspace{\tipitakaparskip}
\csromancenter{กลหวิวาทสุตฺตนิทฺเทส มนุสฺสา จ.\csromanspacer{} \textbf{โลภา}ติ โย โลโภ ลุพฺภนา ลุพฺภิตตฺตํ สาราโค สารชฺชนา สารชฺชิตตฺตํ อภิชฺฌา โลโภ อกุสลมูลํ.\csromanspacer{} \textbf{วิจรนฺตี}ติ วิจรนฺติ วิหรนฺติ อิริยนฺติ วตฺตนิ ปาเลนฺติ ยเปนฺติ ยาเปนฺติ.\csromanspacer{} \textbf{โลเก}ติ อปายโลเก มนุสฺสโลเก เทวโลเก ขนฺธโลเก ธาตุโลเก อายตนโลเกติ เย จาปิ \textbf{โลภา วิจรนฺติ โลเก.}}
\prose{\csromanfolio{201}\textbf{อาสา จ นิฏฺฐา จ กุโตนิทานา}ติ อาสา จ นิฏฺฐา จ กุโต นิทานา กุโตชาตา กุโตสญฺชาตา กุโตนิพฺพตฺตา กุโต-อภินิพฺพตฺตา กุโตปาตุภูตา, กิํนิทานา กิํสมุทยา กิํชาติกา กิํปภวาติ อาสาย จ นิฏฺฐาย จ มูลํ ปุจฺฉติ ฯเปฯ สมุทยํ ปุจฺฉติ ปปุจฺฉติ ยาจติ อชฺเฌสติ ปสาเทตีติ อาสา จ นิฏฺฐา จ กุโตนิทานา.\csromanspacer{} \textbf{เย สมฺปรายาย นรสฺส โหนฺตี}ติ เย นรสฺส ปรายณา โหนฺติ ทีปา โหนฺติ ตาณา โหนฺติ เลณา โหนฺติ สรณา โหนฺติ นิฏฺฐา ปรายณา โหนฺตีติ เย สมฺปรายาย นรสฺส โหนฺติ.\csromanspacer{} เตนาห โส นิมฺมิโต\csromandash{}}

\csromangathameasure{ปิยา สุ โลกสฺมิํ กุโตนิทานา, \\ \textbf{เย จาปิ} \textbf{โลภา วิจรนฺติ โลเก.} \\ \textbf{อาสา จ นิฏฺฐา จ กุโตนิทานา},}
\csromangathagroup{\csromangathastanza{ปิยา สุ โลกสฺมิํ กุโตนิทานา, \\ \textbf{เย จาปิ} \textbf{โลภา วิจรนฺติ โลเก.} \\ \textbf{อาสา จ นิฏฺฐา จ กุโตนิทานา},}}

\prosecont{\textbf{เย สมฺปรายาย นรสฺส โหนฺตี}ติ.}

\proseitem{100}{\csromansymbolfootnote{*}{ขุ 1. 413 ปิฏฺเฐ.}ฉนฺทานิทานานิ ปิยานิ โลเก,}

\prose{\textbf{เย จาปิ}\footnote{เย วาปิ (สฺยา)} \textbf{โลภา วิจรนฺติ โลเก.}}

\csromangathameasure{\textbf{อาสา จ นิฏฺฐา จ อิโตนิทานา,} \\ \textbf{เย สมฺปรายาย นรสฺส โหนฺติ.}}
\csromangathagroup{\csromangathastanza{\textbf{อาสา จ นิฏฺฐา จ อิโตนิทานา,} \\ \textbf{เย สมฺปรายาย นรสฺส โหนฺติ.}}}

\prose{\textbf{ฉนฺทนิทานิ ปิยานิ โลเก}ติ \textbf{ฉนฺโท}ติ โย กาเมสุ กามจฺ\textbf{ฉนฺโท} กามราโค กามนนฺที กามตณฺหา กามสฺเนโห กามปริฬาโห กามมุจฺฉา กามชฺโฌสานํ กาโมโฆ กามโยโค กามุปาทานํ กามจฺฉนฺทนีวรณํ.\csromanspacer{}\csromanspacer{}อปิ จ ปญฺจ ฉนฺทา ปริเยสนจฺ\textbf{ฉนฺโท} ปฏิลาภจฺ\textbf{ฉนฺโท} ปริโภคจฺ\textbf{ฉนฺโท} สนฺนิธิจฺ\textbf{ฉนฺโท} วิสชฺชนจฺ\textbf{ฉนฺโท}.\csromanspacer{} กตโม ปริเยสนจฺ\textbf{ฉนฺโท}, อิเธกจฺโจ อชฺโฌสิโตเยว อตฺถิโก ฉนฺทชาโต รูเป ปริเยสติ.\csromanspacer{} สทฺเท.\csromanspacer{} คนฺธ.\csromanspacer{} รเส.\csromanspacer{} โผฏฺฐพฺเพ ปริเยสติ, \csromanfolio{202}อยํ ปริเยสนจฺฉนฺโท.\csromanspacer{} กตโม ปฏิลาภจฺฉนฺโท, อิเธกจฺโจ อชฺโฌสิโตเยว อตฺถิโก ฉนฺทชาโต รูเป ปฏิลภติ.\csromanspacer{} สทฺเท.\csromanspacer{} คนฺเธ.\csromanspacer{} รเส.\csromanspacer{} โผฏฺฐพฺเพ ปฏิลภติ, อยํ ปฏิลาภจฺฉนฺโท.\csromanspacer{} กตโม ปริโภคจฺฉนฺโท, อิเธกจฺโจ อชฺโฌสิโตเยว อตฺถิโก ฉนฺทชาโต รูเป ปริภุญฺชติ.\csromanspacer{} สทฺเท.\csromanspacer{} คนฺเธ.\csromanspacer{} รเส.\csromanspacer{} โผฏฺฐพฺเพ ปริภุญฺชติ, อยํ ปริโภคจฺฉนฺโท.\csromanspacer{} กตโม สนฺนิธิจฺฉนฺโท, อิเธกจฺโจ อชฺโฌสิโตเยว อตฺถิโก ฉนฺทชาโต ธนสนฺนิจยํ กโรติ “อาปทาสุ ภวิสฺสตี”ติ, อยํ สนฺนิธิจฺฉนฺโท.\csromanspacer{} กตโม วิสชฺชนจฺฉนฺโท, อิเธกจฺโจ อชฺโฌสิโตเยว อตฺถิโก ฉนฺทชาโต ธนํ วิสชฺเชติ หตฺถาโรหานํ อสฺสาโรหานํ รถิกานํ ธนุคฺคหานํ ปตฺติกานํ “อิเม มํ รกฺขิสฺสนฺติ โคปิสฺสนฺติ สมฺปริวาริสฺสนฺตี”ติ, อยํ วิสชฺชนจฺฉนฺโท.\csromanspacer{} \textbf{ปิยานี}ติ เทฺว ปิยา สตฺตา วา สงฺขารา วา ฯเปฯ อิเม สตฺตา ปิยา ฯเปฯ อิเม สงฺขารา ปิยา.\csromanspacer{} \textbf{ฉนฺทานิทานานิ ปิยานิ โลเก}ติ ปิยา ฉนฺทนิทานา ฉนฺทสมุทยา ฉนฺทชาติกา ฉนฺทปภวาติ ฉนฺทานิทานานิ ปิยานิ โลเก.}

\prose{\textbf{เย จาปิ โลภา วิจรนฺติ โลเก}ติ \textbf{เย จาปี}ติ ขตฺติยา จ พฺราหฺมณา จ เวสฺสา จ สุทฺทา จ คหฏฺฐา จ ปพฺพชิตา จ เทวา จ มนุสฺสา จ.\csromanspacer{} \textbf{โลภา}ติ โย โลโภ ลุพฺภนา ลุพฺภิตตฺตํ สาราโค สารชฺชนฺ\textbf{อา สา}รชฺชิตตฺตํ อภิชฺฌา โลโภ อกุสลมูลํ.\csromanspacer{} \textbf{วิจรนฺตี}ติ วิจรนฺติ วิหรนฺติ อิริยนฺติ วตฺตนฺติ ปาเลนฺติ ยเปนฺติ ยาเปนฺติ.\csromanspacer{} \textbf{โลเก}ติ อปายโลเก ฯเปฯ อายกนโลเกติ เย จาปิ โลภา วิจรนฺติ โลเก.}

\prose{\textbf{อาสา จ นิฏฺฐา จ อิโตนิทานา}ติ \textbf{อาสา} วุจฺจติ ตณฺหา.\csromanspacer{} โย ราโค สาราโค ฯเปฯ อภิชฺฌา โลโภ อกุสลมูลํ.\csromanspacer{} \textbf{นิฏฺฐา}ติ อิเธกจฺโจ รูเป ปริเยสนฺโต รูปํ ปฏิลภติ, รูปนิฏฺโฐ โหติ, สทฺเท.\csromanspacer{} คนฺเธ.\csromanspacer{} รเส.\csromanspacer{} โผฏฺฐพฺเพ.\csromanspacer{} กุลํ.\csromanspacer{} คณํ.\csromanspacer{} อาวาสํ.\csromanspacer{} ลาภํ.\csromanspacer{} ยสํ.\csromanspacer{} ปสํสํ.\csromanspacer{} สุขํ.\csromanspacer{} จีวรํ.\csromanspacer{} ปิณฺฑปาตํ.\csromanspacer{} เสนาสนํ.\csromanspacer{} คิลานปจฺจยเภสชฺชปริกฺขารํ.\csromanspacer{} สุตฺตนฺตํ.\csromanspacer{} วินยํ.\csromanspacer{} อภิธมฺมํ.\csromanspacer{} อารญฺญิกงฺคํ.\csromanspacer{} ปิณฺฑปาติกงฺคํ.\csromanspacer{} ปํสุกูลิกงฺคํ.\csromanspacer{} เตจีวริกงฺคํ.\csromanspacer{} สปทานจาริกงฺคํ.\csromanspacer{} ขลุปจฺฉาภตฺติกงฺคํ.\csromanspacer{} เนสชฺชิกงฺคํ.\csromanspacer{} ยถาสนฺถติกงฺคํ.\csromanspacer{} ปฐมํ ฌานํ.\csromanspacer{} ทุติยํ ฌานํ.\csromanspacer{} ตติยํ ฌานํ.\csromanspacer{} จตุตฺถํ ฌานํ.\csromanspacer{} อากฺ\textbf{อาสา}นญฺจายตนสมาปตฺติํ.\csromanspacer{} วิญฺญาณญฺจายตนสมาปตฺติํ.}

\vspace{\tipitakaparskip}
\csromancenter{กลหวิวาทสุตฺตนิทฺเทส อากิญฺจญฺญายตนสมาปตฺติํ.\csromanspacer{} เนวสญฺญานาสญฺญายตนสมาปตฺติํ ปริเยสนฺโต เนวสญฺญานาสญฺญายตนสมาปตฺติํ ปฏิลภติ, เนวสญฺญานาสญฺญายตนสมาปตฺตินิฏฺโฐ โหติ.}
\vspace{0.5\baselineskip}
\csromangathameasure{อาสาย กสเต เขตฺตํ, พีชํ อาสาย วปฺปติ. \\ อาสาย วาณิชา ยนฺติ, สมุทฺทํ ธนหารกา.}
\csromangathagroup{\csromangathastanza{\csromanfolio{203}อาสาย กสเต เขตฺตํ, พีชํ อาสาย วปฺปติ. \\ อาสาย วาณิชา ยนฺติ, สมุทฺทํ ธนหารกา.}}

\prose{ยาย อาสาย ติฏฺฐามิ, สา เม อาสา สมิชฺฌตีติ\csromandash{}}

\prose{อาสาย สมิทฺธิ วุจฺจเต นิฏฺฐา.\csromanspacer{}\csromanspacer{}\textbf{อาสา จ นิฏฺฐา จ อิโตนิทานา}ติ อาสา จ นิฏฺฐา จ อิโต ฉนฺทนิทานา ฉนฺทสมุทยา ฉนฺทชาติกา ฉนฺทปภวาติ อาสา จ นิฏฺฐา จ อิโตนิทานา.}

\prose{\textbf{เย สมฺปรายาย นรสฺส โหนฺตี}ติ เย นรสฺส ปรายณา โหนฺติ ทีปา โหนฺติ ตาณา โหนฺติ เลณา โหนฺติ สรณา โหนฺติ นิฏฺฐา ปรายณา โหนฺตีติ เย สมฺปรายาย นรสฺส โหนฺติ.\csromanspacer{} เตนาห ภควา\csromandash{}}

\csromangathameasure{ฉนฺทานิทานานิ ปิยานิ โลเก, \\ เย จาปิ โลภา วิจรนฺติ โลเก. \\ \textbf{อาสา จ นิฏฺฐา จ อิโตนิทานา},}
\csromangathagroup{\csromangathastanza{ฉนฺทานิทานานิ ปิยานิ โลเก, \\ เย จาปิ โลภา วิจรนฺติ โลเก. \\ \textbf{อาสา จ นิฏฺฐา จ อิโตนิทานา},}}

\prose{\textbf{เย สมฺปรายาย นรสฺส โหนฺตี}ติ.}

\proseitem{101}{\csromansymbolfootnote{*}{ขุ 1. 411 ปิฏฺเฐ.}\textbf{ฉนฺโท นุ โลกสฺมิํ กุโตนิทาโน},}

\prose{\textbf{วินิจฺฉยา จาปิ}\footnote{วาปิ (สี, สฺยา)}\textbf{ กุโตปหูตา}.}

\csromangathameasure{\textbf{โกโธ โมสวชฺชญฺจ กถํกถา จ,} \\ \textbf{เย จาปิ ธมฺมา สมเณน วุตฺตา.}}
\csromangathagroup{\csromangathastanza{\textbf{โกโธ โมสวชฺชญฺจ กถํกถา จ,} \\ \textbf{เย จาปิ ธมฺมา สมเณน วุตฺตา.}}}

\prose{\textbf{ฉนฺโท นุ โลกสฺมิํ กุโตนิทาโน}ติ ฉนฺโท กุโตนิทาโน กุโตชาโต กุโตสญฺชาโต กุโตนิพฺพตฺโต กุโต-อภินิพฺพตฺโต กุโตปาตุภูโต, กิํนิทาโน กิํสมุทโย กิํชาติโก กิํปภโวติ ฉนฺทสฺส มูลํ ปุจฺฉติ ฯเปฯ สมุทยํ ปุจฺฉติ ปปุจฺฉติ ยาจติ อชฺเฌสติ ปสาเทตีติ ฉนฺโท นุ โลกสฺมิํ กุโตนิทาโน.}

\prose{\textbf{วินิจฺฉยา จาปิ กุโตปหูตา}ติ วินิจฺฉยา กุโตปหูตา กุโตชาตา กุโตสญฺชาตา กุโตนิพฺพตฺตา กุโต-อภินิพฺพตฺตา \csromanfolio{204}กุโตปาตุภูตา, กิํนิทานา กิํสมุทยา กิํชาติกา กิํปภวาติ \textbf{วินิจฺฉยา}นํ มูลํ ปุจฺฉติ ฯเปฯ สมุทยํ ปุจฺฉติ ปปุจฺฉติ ยาจติ อชฺเฌสติ ปสาเทตีติ \textbf{วินิจฺฉยา} จาปิ กุโตปหูตา.}

\prose{\textbf{โกโธ โมสวชฺชญฺจ กถํกถา จา}ติ \textbf{โกโธ}ติ โย เอวรูโป จิตฺตสฺส อาฆาโต ปฏิฆาโต ปฏิฆํ ปฏิวิโรโธ โกโป ปโกโป สมฺปโกโป โทโส ปโทโส สมฺปโทโส จิตฺตสฺส พฺยาปตฺติ มโนปโทโส \textbf{โกโธ} กุชฺฌนา กุชฺฌิตตฺตํ โทโส ทุสฺสนา ทุสฺสิตตฺตํ พฺยาปตฺติ พฺยาปชฺชนา พฺยาปชฺชิตตฺตํ วิโรโธ ปฏิวิโรโธ จณฺฑิกฺกํ อสุโรโป อนตฺตมนตา จิตฺตสฺส, \textbf{โมสวชฺชํ} วุจฺจติ มุสาวาโท.\csromanspacer{} \textbf{กถํกถา} วุจฺจติ วิจิกิจฺฉาติ \textbf{โกโธ} โมสวชฺชญฺจ กถํกถา จ.}

\prose{\textbf{เย จาปิ ธมฺมา สมเณน วุตฺตา}ติ \textbf{เย จาปี}ติ เย โกเธน จ โมสวชฺเชน จ กถํกถาย จ สหคตา สหชาตา สํสฏฺฐา สมฺปยุตฺตา เอกุปฺปาทา เอกนิโรธา เอกวตฺถุกา เอการมฺมณา, อิเม วุจฺจนฺติ เย จาปิ ธมฺมา.\csromanspacer{}\csromanspacer{}อถ วา เย เต กิเลสา อญฺญชาติกา อญฺญวิหิตกา, อิเม วุจฺจนฺติ เย จาปิ ธมฺมา.\csromanspacer{} \textbf{สมเณน วุตฺตา}ติ สมเณน สมิตปาเปน พฺราหฺมเณน พาหิตปาปธมฺเมน ภิกฺขุนา ภินฺนกิเลสมูเลน สพฺพากุสลมูลพนฺธนา ปมุตฺเตน วุตฺตา ปวุตฺตา อาจิกฺขิตา เทสิตา ปญฺญปิตา ปฏฺฐปิตา วิวฏา วิภตฺตา อุตฺตานีกตา ปกาสิตาติ เย จาปิ ธมฺมา สมเณน วุตฺตา.\csromanspacer{} เตนาห โส นิมฺมิโต\csromandash{}}

\csromangathameasure{ฉนฺโท นุ โลกสฺมิํ กุโตนิทาโน, \\ วินิจฺฉยา จาปิ กุโตปหูตา.}
\csromangathagroup{\csromangathastanza{ฉนฺโท นุ โลกสฺมิํ กุโตนิทาโน, \\ วินิจฺฉยา จาปิ กุโตปหูตา.}}

\prose{โกโธ โมสวชฺชญฺจ กถํกถา จ,}

\prose{\textbf{เย จาปิ ธมฺมา สมเณน วุตฺตา}ติ.}

\proseitem{102}{\csromansymbolfootnote{*}{ขุ 1. 413 ปฺปิฏฺเฐ.}\textbf{สาตํ อสาตนฺติ ยมาหุ โลเก},}

\prose{\textbf{ตมูปนิสฺสาย ปโหติ ฉนฺโท.}}

\csromangathameasure{\textbf{รูเปสุ ทิสฺวา วิภวํ ภวญฺจ,} \\ \textbf{วินิจฺฉยํ กุพฺพติ} \textbf{ชนฺตุ โลเก.}}
\csromangathagroup{\csromangathastanza{\textbf{รูเปสุ ทิสฺวา วิภวํ ภวญฺจ,} \\ \textbf{วินิจฺฉยํ กุพฺพติ}\footnote{กูรุเต (สฺยา)} \textbf{ชนฺตุ โลเก.}}}

\prose{\textbf{สาตํ อสาตนฺติ ยมาหุ โลเก}ติ \textbf{สาตนฺ}ติ สุขา จ เวทนา อิฏฺฐญฺจ วตฺถุ\footnote{วตฺถุํ (สี, ก)}.\csromanspacer{} \textbf{อสาตนฺ}ติ ทุกฺขา จ เวทนา อนิฏฺฐญฺจ วตฺถุ.\csromanspacer{} \textbf{ยมาหุ}}

\vspace{\tipitakaparskip}
\csromancenter{กลหวิวาทสุตฺตนิทฺเทส \textbf{โลเก}ติ ยํ อาหํสุ ยํ กเถนฺติ ยํ ภณนฺติ ยํ ทีเปนฺติ ยํ โวหรนฺตีติ สาตํ อสาตนฺติ ยมาหุ \textbf{โลเก}.}
\prose{\csromanfolio{205}\textbf{ตมูปนิสฺสาย ปโหติ ฉนฺโท}ติ สาตาสาตํ นิสฺสาย สุขทุกฺขํ นิสฺสาย โสมนสฺสโทมนสฺสํ นิสฺสาย อิฏฺฐานิฏฺฐํ นิสฺสาย อนุนยปฏิฆํ นิสฺสาย ฉนฺโท ปโหติ ปภวติ ชายติ สญฺชายติ นิพฺพตฺตติ อภินิพฺพตฺตตีติ ตมูปนิสฺสาย ปโหติ ฉนฺโท.}

\prose{\textbf{รูเปสุ ทิสฺวา วิภวํ ภวญฺจา}ติ \textbf{รูเปสู}ติ จตฺตาโร จ มหาภูตา จตุนฺนญฺจ มหาภูตานํ อุปาทาย รูปํ.\csromanspacer{} กตโม รูปานํ ภโว, โย รูปานํ ภโว ชาติ สญฺชาติ นิพฺพตฺติ อภินิพฺพตฺติ ปาตุภาโว, อยํ รูปานํ ภโว.\csromanspacer{} กตโม รูปานํ วิภโว, โย รูปานํ ขโย วโย เภโท ปริเภโท อนิจฺจตา อนฺตรธานํ, อยํ รูปานํ วิภโว.\csromanspacer{}\csromanspacer{}\textbf{รูเปสุ ทิสฺวา วิภวํ ภวญฺจา}ติ รูเปสุ ภวญฺจ วิภวญฺจ ทิสฺวา ปสฺสิตฺวา ตุลยิตฺวา ตีรยิตฺวา วิภาวยิตฺวา วิภูตํ กตฺวาติ รูเปสุ ทิสฺวา วิภวํ ภวญฺจ.}

\prose{\textbf{วินิจฺฉยํ กุพฺพติ ชนฺตุ โลเก}ติ \textbf{วินิจฺฉยา}ติ เทฺว \textbf{วินิจฺฉยา} ตณฺหาวินิจฺฉโย จ ทิฏฺฐิวินิจฺฉโย จ.\csromanspacer{} กถํ ตณฺหาวินิจฺฉยํ กโรติ, อิเธกจฺจสฺส อนุปฺปนฺนา เจว โภคา น อุปฺปชฺชนฺติ, อุปฺปนฺนา จ โภคา ปริกฺขยํ คจฺฉนฺติ, ตสฺส เอวํ โหติ “เกน นุ โข เม อุปาเยน อนุปฺปนฺนา เจว โภคา น อุปฺปชฺชนฺติ, อุปฺปนฺนา จ โภคา ปริกฺขยํ คจฺฉนฺตี”ติ, ตสฺส ปน เอวํ โหติ “สุราเมรยมชฺชปฺปมาทฏฺฐานานุโยคํ อนุยุตฺตสฺส เม อนุปฺปนฺนา เจว โภคา น อุปฺปชฺชนฺติ, อุปฺปนฺนา จ โภคา ปริกฺขยํ คจฺฉนฺติ.\csromanspacer{} วิกาลวิสิขาจริยานุโยคํ อนุยุตฺตสฺส เม อนุปฺปนฺนา เจว โภคา น อุปฺปชฺชนฺติ, อุปฺปนฺนา จ โภคา ปริกฺขยํ คจฺฉนฺติ.\csromanspacer{} สมชฺชาภิจรณํ อนุยุตฺตสฺส เม.\csromanspacer{} ชุตปฺปมาทฏฺฐานานุโยคํ อนุยุตฺตสฺส เม.\csromanspacer{} ปาปมิตฺตานุโยคํ อนุยุตฺตสฺส เม อนุปฺปนฺนา เจว โภคา น อุปฺปชฺชนฺติ, อุปฺปนฺนา จ โภคา ปริกฺขยํ คจฺฉนฺติ.\csromanspacer{} อาลสฺยานุโยคํ อนุยุตฺตสฺส เม อนุปฺปนฺนา เจว โภคา น อุปฺปชฺชนฺติ, อุปฺปนฺนา จ โภคา ปริกฺขยํ คจฺฉนฺตี”ติ เอวํ ญาณํ กตฺวา ฉ โภคานํ อปายมุขานิ น เสวติ, ฉ โภคานํ อายมุขานิ เสวติ.\csromanspacer{} เอวมฺปิ ตณฺหาวินิจฺฉยํ กโรติ.}

\prose{\csromanfolio{206}อถ วา กสิยา วา วณิชฺชาย วา โครกฺเขน วา อิสฺสตฺเถน\footnote{อิสฺสตฺเตน (ก-สี, ก) อิสุ + สตฺถ.} วา ราชโปริเสน วา สิปฺปญฺญตเรน วา ปฏิปชฺชติ.\csromanspacer{} เอวมฺปิ ตณฺหาวินิจฺฉยํ กโรติ.\csromanspacer{} กถํ ทิฏฺฐิวินิจฺฉยํ กโรติ, จกฺขุสฺมิํ อุปฺปนฺเน ชานาติ “อตฺตา เม อุปฺปนฺโน”ติ, จกฺขุสฺมิํ อนฺตรหิเต ชานาติ “อตฺตา เม อนฺตรหิโต, วิคโต เม อตฺตา”ติ.\csromanspacer{} เอวมฺปิ ทิฏฺฐิวินิจฺฉยํ กโรติ.\csromanspacer{} โสตสฺมิํ.\csromanspacer{} ฆานสฺมิํ.\csromanspacer{} ชิวฺหาย.\csromanspacer{} กายสฺมิํ.\csromanspacer{} รูปสฺมิํ.\csromanspacer{} สทฺทสฺมิํ.\csromanspacer{} คนฺธสฺมิํ.\csromanspacer{} รสสฺมิํ.\csromanspacer{} โผฏฺฐพฺพสฺมิํ อุปฺปนฺเน ชานาติ “อตฺตา เม อุปฺปนฺโน”ติ, โผฏฺฐพฺพสฺมิํ อนฺตรหิเต ชานาติ “อตฺตา เม อนฺตรหิโต, วิคโต เม อตฺตา”ติ.\csromanspacer{} เอวมฺปิ ทิฏฺฐิวินิจฺฉยํ กโรติ ชเนติ สญฺชเนติ นิพฺพตฺเตติ อภินิพฺพตฺเตติ.\csromanspacer{} \textbf{ชนฺตู}ติ สตฺโต นโร มานโว ฯเปฯ มนุโช.\csromanspacer{} \textbf{โลเก}ติ อปายโลเก ฯเปฯ อายตนโลเกติ วินิจฺฉยํ กุพฺพติ ชนฺตุ โลเก.\csromanspacer{} เตนาห ภควา\csromandash{}}

\csromangathameasure{สาตํ อสาตนฺติ ยมาหุ โลเก, \\ ตมูปนิสฺสาย ปโหติ ฉนฺโท. \\ รูเปสุ ทิสฺวา วิภวํ ภวญฺจ, \\ วินิจฺฉยํ กุพฺพติ ชนฺตุ โลเกติ.}
\csromangathagroup{\csromangathastanza{สาตํ อสาตนฺติ ยมาหุ โลเก, \\ ตมูปนิสฺสาย ปโหติ ฉนฺโท. \\ รูเปสุ ทิสฺวา วิภวํ ภวญฺจ, \\ วินิจฺฉยํ กุพฺพติ ชนฺตุ โลเกติ.}}

\proseitem{103}{\csromansymbolfootnote{*}{ขุ 1. 414 ปิฏฺเฐ.}โกโธ โมสวชฺชญฺจ \textbf{กถํกถา} จ,}

\prose{\textbf{เอเตปิ ธมฺมา ทฺวยเมว สนฺเต.}}

\csromangathameasure{\textbf{กถํกถี ญาณปถาย สิกฺเข,} \\ \textbf{ญตฺวา ปวุตฺตา สมเณน ธมฺมา.}}
\csromangathagroup{\csromangathastanza{\textbf{กถํกถี ญาณปถาย สิกฺเข,} \\ \textbf{ญตฺวา ปวุตฺตา สมเณน ธมฺมา.}}}

\prose{\textbf{โกโธ โมสวชฺชญฺจ กถํกถา จา}ติ \textbf{โกโธ}ติ โย เอวรูโป จิตฺตสฺส อาฆาโต ปฏิฆาโต ฯเปฯ.\csromanspacer{} \textbf{โมสวชฺชํ} วุจฺจติ มุสาวาโท, \textbf{กถํกถา} วุจฺจติ วิจิกิจฺฉา.\csromanspacer{} อิฏฺฐํ วตฺถุํ นิสฺสายปิ \textbf{โกโธ} ชายติ, อนิฏฺฐํ วตฺถุํ นิสฺสายปิ \textbf{โกโธ} ชายติ.\csromanspacer{} อิฏฺฐํ วตฺถุํ นิสฺสายปิ มุสาวาโท อุปฺปชฺชติ, อนิฏฺฐํ วตฺถุํ นิสฺสายปิ มุสาวาโท อุปฺปชฺชติ.\csromanspacer{} อิฏฺฐํ วตฺถุํ นิสฺสายปิ \textbf{กถํกถา} อุปฺปชฺชติ, อนิฏฺฐํ วตฺถุํ นิสฺสายปิ \textbf{กถํกถา} อุปฺปชฺชติ.}

\prose{กถํ อนิฏฺฐํ วตฺถุํ นิสฺสาย \textbf{โกโธ} ชายติ, ปกติยา อนิฏฺฐํ วตฺถุํ นิสฺสาย \textbf{โกโธ} ชายติ “อนตฺถํ เม อจรี”ติ \textbf{โกโธ} ชายติ, “อนตฺถํ เม จรตี”ติ \textbf{โกโธ} ชายติ, “อนตฺถํ เม}

\vspace{\tipitakaparskip}
\csromancenter{กลหวิวาทสุตฺตนิทฺเทส จริสฺสตี”ติ โกโธ ชายติ, “ปิยสฺส เม มนาปสฺส อนตฺถํ อจริ.\csromanspacer{} อนตฺถํ จรติ.\csromanspacer{} อนตฺถํ จริสฺสตี”ติ โกโธ ชายติ. “อปฺปิยสฺส เม อมนาปสฺส อตฺถํ อจริ.\csromanspacer{} อตฺถํ จรติ.\csromanspacer{} อตฺถํ จริสฺสตี”ติ โกโธ ชายติ.\csromanspacer{} เอวํ อนิฏฺฐํ วตฺถุํ นิสฺสาย โกโธ ชายติ.}
\prose{\csromanfolio{207}กถํ อิฏฺฐํ วตฺถุํ นิสฺสาย โกโธ ชายติ, อิฏฺฐํ วตฺถุํ อจฺเฉทสงฺกิโนปิ โกโธ ชายติ, อจฺฉิชฺชนฺเตปิ โกโธ ชายติ, อจฺฉินฺเนปิ โกโธ ชายติ.\csromanspacer{} อิฏฺฐํ วตฺถุํ วิปริณามสงฺกิโนปิ โกโธ ชายติ, วิปริณามนฺเตปิ โกโธ ชายติ, วิปริณเตปิ โกโธ ชายติ.\csromanspacer{} เอวํ อิฏฺฐํ วตฺถุํ นิสฺสาย โกโธ ชายติ.}

\prose{กถํ อนิฏฺฐํ วตฺถุํ นิสฺสาย มุสาวาโท อุปฺปชฺชติ, อิเธกจฺโจ อนฺทุพนฺธเนน\footnote{อทฺทุพนฺธเนน (สฺยา, ก)} วา พทฺโธ\footnote{พนฺโธ (สฺยา, ก)}, ตสฺส พนฺธนสฺส โมกฺขตฺถาย สมฺปชานมุสา ภาสติ.\csromanspacer{} รชฺชุพนฺธเนน วา พทฺโธ.\csromanspacer{} สงฺขลิกพนฺธเนน วา พทฺโธ.\csromanspacer{} เวตฺตพนฺธเนน วา พทฺโธ.\csromanspacer{} ลตาพนฺธเนน วา พทฺโธ.\csromanspacer{} ปกฺเขปพนฺธเนน วา พทฺโธ.\csromanspacer{} ปริกฺเขปพนฺธเนน วา พทฺโธ.\csromanspacer{} คาม นิคม นคร รฏฺฐพนฺธเนน วา พทฺโธ.\csromanspacer{} ชนปทพนฺธเนน วา พทฺโธ, ตสฺส พนฺธนสฺส โมกฺขตฺถาย สมฺปชานมุสา ภาสติ.\csromanspacer{} เอวํ อนิฏฺฐํ วตฺถุํ นิสฺสาย มุสาวาโท อุปฺปชฺชตีติ.}

\prose{กถํ อิฏฺฐํ วตฺถุํ นิสฺสาย มุสาวาโท อุปฺปชฺชติ, อิเธกจฺโจ มนาปิกานํ\footnote{มนาปานํ (สี)} รูปานํ เหตุ สมฺปชานมุสา ภาสติ.\csromanspacer{} มนาปิกานํ สทฺทานํ.\csromanspacer{} คนฺธานํ.\csromanspacer{} รสานํ.\csromanspacer{} โผฏฺฐพฺพานํ เหตุ.\csromanspacer{} จีวรเหตุ.\csromanspacer{} ปิณฺฑปาตเหตุ.\csromanspacer{} เสนาสนเหตุ.\csromanspacer{} คิลานปจฺจยเภสชฺชปริกฺขารเหตุ สมฺปชานมุสา ภาสติ.\csromanspacer{} เอวํ อิฏฺฐํ วตฺถุํ นิสฺสาย มุสาวาโท อุปฺปชฺชติ.}

\prose{กถํ อนิฏฺฐํ วตฺถุํ นิสฺสาย กถํกถา อุปฺปชฺชติ, “มุจฺจิสฺสามิ\footnote{มุญฺจิสฺสามิ (สี)} นุ โข จกฺขุโรคโต, น นุ โข มุจฺจิสฺสามิ จกฺขุโรคโต.\csromanspacer{} มุจฺจิสฺสามิ นุ โข โสตโรคโต.\csromanspacer{} ฆานโรคโต.\csromanspacer{} ชิวฺหาโรคโต.\csromanspacer{} กายโรคโต.\csromanspacer{} สีสโรคโต.\csromanspacer{} กณฺณโรคโต.\csromanspacer{} มุขโรคโต.\csromanspacer{} มุจฺจิสฺสามิ นุ โข ทนฺตโรคโต, น นุ โข มุจฺจิสฺสามิ ทนฺตโรคโต”ติ เอวํ อนิฏฺฐํ วตฺถุํ นิสฺสาย กถํกถา อุปฺปชฺชติ.}

\prose{\csromanfolio{208}กถํ อิฏฺฐํ วตฺถุํ นิสฺสาย กถํกถา อุปฺปชฺชติ, “ลภิสฺสามิ นุ โข มนาปิเก\footnote{มนาปิเย (สี, ก)} รูเป, น นุ โข ลภิสฺสามิ มนาปิเก รูเป.\csromanspacer{} ลภิสฺสามิ นุ โข มนาปิเก สทฺเท.\csromanspacer{} คนฺเธ.\csromanspacer{} รเส.\csromanspacer{} โผฏฺฐพฺเพ.\csromanspacer{} กุลํ.\csromanspacer{} คณํ.\csromanspacer{} อาวาสํ.\csromanspacer{} ลาภํ.\csromanspacer{} ยสํ.\csromanspacer{} ปสํสํ.\csromanspacer{} สุขํ.\csromanspacer{} จีวรํ.\csromanspacer{} ปิณฺฑปาตํ.\csromanspacer{} เสนาสนํ.\csromanspacer{} คิลานปจฺจยเภสชฺชปริกฺขารนฺ”ติ.\csromanspacer{} เอวํ อิฏฺฐํ วตฺถุํ นิสฺสาย กถํกถา อุปฺปชฺชตีติ โกโธ โมสวชฺชญฺจ กถํกถา จ.}

\prose{\textbf{เอเตปิ ธมฺมา ทฺวยเมว สนฺเต}ติ สาตาสาเต สนฺเต, สุขทุกฺเข สนฺเต, โสมนสฺสโทมนสฺเส สนฺเต, อิฏฺฐานิฏฺเฐ สนฺเต, อนุนยปฏิเฆ สนฺเต สํวิชฺชมาเน อตฺถิ อุปลพฺภมาเนติ เอเตปิ ธมฺมา ทฺวยเมว สนฺเต.}

\prose{\textbf{กถํกถี ญาณปถาย สิกฺเข}ติ ญาณมฺปิ ญาณปโถ, ญาณสฺส อารมฺมณมฺปิ ญาณปโถ, ญาณสหภุโนปิ ธมฺมา ญาณปโถ.\csromanspacer{} ยถา อริยมคฺโค อริยปโถ, เทวมคฺโค เทวปโถ, พฺรหฺมมคฺโค พฺรหฺมปโถ, เอวเมวํ ญาณมฺปิ ญาณปโถ, ญาณสฺส อารมฺมณมฺปิ ญาณปโถ, ญาณสหภุโนปิ ธมฺมา ญาณปโถ.}

\prose{\textbf{สิกฺเข}ติ ติสฺโส สิกฺขา อธิสีลสิกฺขา อธิจิตฺตสิกฺขา อธิปญฺญาสิกฺขา.\csromanspacer{} กตมา อธิสีลสิกฺขา, อิธ ภิกฺขุ สีลวา โหติ, ปาติโมกฺขสํวรสํวุโต วิหรติ อาจารโคจรสมนฺโน อณุมตฺเตสุ วชฺเชสุ ภยทสฺสาวี, สมาทาย สิกฺขติ สิกฺขาปเทสุ, ขุทฺทโก สีลกฺขนฺโธ มหนฺโต สีลกฺขนฺโธ สีลํ ปติฏฺฐา อาทิ จรณํ สํยโม สํวโร มุขํ ปมุขํ กุสลานํ ธมฺมานํ สมาปตฺติยา.\csromanspacer{} อยํ อธิสีลสิกฺขา.\csromanspacer{} กตมา อธิจิตฺตสิกฺขา, อิธ ภิกฺขุ วิวิจฺเจว กาเมหิ ฯเปฯ จตุตฺถํ ฌานํ อุปสมฺปชฺช วิหรติ, อยํ อธิจิตฺตสิกฺขา.\csromanspacer{}\csromanspacer{}กตมา อธิปญฺญาสิกฺขา, อิธ ภิกฺขุ ปญฺญวา โหติ อุทยตฺถคามินิยา ปญฺญาย สมนฺนาคโต อริยาย นิพฺเพธิกาย สมฺมาทุกฺขกฺขยคามินิยา. “โส อิทํ ทุกฺขนฺ”ติ ยถาภูตํ ปชานาติ ฯเปฯ “อยํ ทุกฺขนิโรธคามินี ปฏิปทา”ติ ยถาภูตํ ปชานาติ, “อิเม อาสวา”ติ ยถาภูตํ ปชานาติ ฯเปฯ “อยํ อาสวนิโรธคามินี ปฏิปทา”ติ ยถาภูตํ ปชานาติ, อยํ อธิปญฺญาสิกฺขา.}

\csromanheader{2}{11. กลหวิวาทสุตฺตนิทฺเทส}
\prose{\csromanfolio{209}\textbf{กถํกถี ญาณปถาย สิกฺเข}ติ กถํกถี ปุคฺคโล สกงฺโข สวิเลโข สเทฺวฬฺหโก สวิจิกิจฺโฉ ญาณาธิคมาย ญาณผุสนาย ญาณสจฺฉิกิริยาย, อธิสีลมฺปิ สิกฺเขยฺย, อธิจิตฺตมฺปิ สิกฺเขยฺย, อธิปญฺญมฺปิ สิกฺเขยฺย, อิมา ติสฺโส สิกฺขาโย อาวชฺชนฺโต สิกฺเขยฺย, ชานนฺโต สิกฺเขยฺย, ปสฺสนฺโต สิกฺเขยฺย, ปจฺจเวกฺขนฺโต สิกฺเขยฺย, จิตฺตํ อธิฏฺฐหนฺโต สิกฺเขยฺย, สทฺธาย อธิมุจฺจนฺโต สิกฺเขยฺย, วีริยํ ปคฺคณฺหนฺโต สิกฺเขยฺย, สติํ อุปฏฺฐหนฺโต สิกฺเขยฺย, จิตฺตํ สมาทหนฺโต สิกฺเขยฺย, ปญฺญาย ปชานนฺโต สิกฺเขยฺย, อภิญฺเญยฺยํ อภิชานนฺโต สิกฺเขยฺย, ปริญฺเญยฺยํ ปริชานนฺโต สิกฺเขยฺย, ปหาตพฺพํ ปชหนฺโต สิกฺเขยฺย, ภาเวตพฺพํ ภาเวนฺโต สิกฺเขยฺย, สจฺฉิกาตพฺพํ สจฺฉิกโรนฺโต สิกฺเขยฺย อาจเรยฺย สมาจเรยฺย สมาทาย วตฺเตยฺยาติ กถํกถี ญาณปถาย สิกฺเข.}

\prose{\textbf{ญตฺวา ปวุตฺตา สมเณน ธมฺมา}ติ \textbf{ญตฺวา}ติ \textbf{ญตฺวา} ชานิตฺวา ตุลยิตฺวา ตีรยิตฺวา วิภาวยิตฺวา วิภูตํ กตฺวา วุตฺตา ปวุตฺตา อาจิกฺขิตา เทสิตา ปญฺญปิตา ปฏฺฐปิตา วิวฏา วิภตฺตา อุตฺตานีกตา\footnote{อุตฺตานิํ กตา (ก)} ปกาสิตา, “สพฺเพ สงฺขารา อนิจฺจา”ติ \textbf{ญตฺวา} ชานิตฺวา ตุลยิตฺวา ตีรยิตฺวา วิภาวยิตฺวา วิภูตํ กตฺวา วุตฺตา ปวุตฺตา อาจิกฺขิตา เทสิตา ปญฺญปิตา ปฏฺฐปิตา วิวฏา วิภตฺตา อุตฺตานีกตา ปกาสิตา, “สพฺเพสงฺขารา ทุกฺขา’ติ. “สพฺเพ ธมฺมา อนตฺตา”ติ. “อวิชฺชาปจฺจยา สงฺขารา”ติ ฯเปฯ “ชาติปจฺจยา ชรามรณนฺ”ติ. “อวิชฺชานิโรธา สงฺขารนิโรโธ”ติ ฯเปฯ “ชาตินิโรธา ชรามรณนิโรโธ”ติ. “อิทํ ทุกฺขนฺ”ติ ฯเปฯ “อยํ ทุกฺขนิโรธคามินี ปฏิปทา”ติ. “อิเม อาสวา”ติ ฯเปฯ “อยํ อาสวนิโรธคามินี ปฏิปทา”ติ. “อิเม ธมฺมา อภิญฺเญยฺยา”ติ. “อิเม ธมฺมา ปริญฺเญยฺยา”ติ. “อิเม ธมฺมา ปหาตพฺพา”ติ. “อิเม ธมฺมา ภาเวตพฺพา”ติ.\csromanspacer{} อิเม ธมฺมา สจฺฉิกาตพฺพา”ติ.\csromanspacer{} ฉนฺนํ ผสฺสายตนานํ สมุทยญฺจ อตฺถงฺคมญฺจ อสฺสาทญฺจ อาทีนวญฺจ นิสฺสรณญฺจ.\csromanspacer{} ปญฺจนฺนํ อุปาทานกฺขนฺธานํ.\csromanspacer{} จตุนฺนํ มหาภูตานํ. “ยํ กิญฺจิ สมุทยธมฺมํ, สพฺพํ ตํ นิโรธธมฺมนฺ”ติ \textbf{ญตฺวา} ชานิตฺวา ตุลยิตฺวา ตีรยิตฺวา วิภาวยิตฺวา วิภูตํ กตฺวา วุตฺตา ปวุตฺตา อาจิกฺขิตา เทสิตา ปญฺญปิตา ปฏฺฐปิตา วิวฏา วิภตฺตา อุตฺตานีกตา ปกาสิตา.}

\prose{\csromanfolio{210}วุตฺตํ เหตํ ภควตา\csromandash{} * อภิญฺญายาหํ ภิกฺขเว ธมฺมํ เทเสมิ, โน อนภิญฺญาย, สนิทานาหํ ภิกฺขเว ธมฺมํ เทเสมิ, โน อนิทานํ, สปฺปาฏิหาริยาหํ ภิกฺขเว ธมฺมํ เทเสมิ, โน อปฺปาฏิหาริยํ.\csromanspacer{} ตสฺส มยฺหํ ภิกฺขเว อภิญฺญาย ธมฺมํ เทสยโต, โน อนภิญฺญาย, สนิทานํ ธมฺมํ เทสยโต, โน อนิทานํ, สปฺปาฏิหาริยํ ธมฺมํ เทสยโต, โน อปฺปาฏิหาริยํ กรณีโย โอวาโท, กรณียา อนุสาสนี.\csromanspacer{} อลญฺจ ปน ภิกฺขเว โว ตุฏฺฐิยา อลํ ปาโมชฺชาย อลํ โสมนสฺสาย “สมฺมาสมฺพุทฺโธ ภควา, สฺวากฺขาโต ธมฺโม, สุปฺปฏิปนฺโน สํโฆ”ติ ฯเปฯ.\csromanspacer{} อิมสฺมิํ จ ปน เวยฺยากรณสฺมิํ ภญฺญมาเน ทสสหสฺสี โลกธาตุ อกมฺปิตฺถาติ ญตฺวา ปวุตฺตา สมเณน ธมฺมา.\csromanspacer{} เตนาห ภควา\csromandash{}}

\csromangathameasure{โกโธ โมสวชฺชญฺจ กถํกถา จ, \\ เอเตปิ ธมฺมา ทฺวยเมว สนฺเต. \\ กถํกถี ญาณปถาย สิกฺเข, \\ ญตฺวา ปวุตฺตา สมเณน ธมฺมาติ.}
\csromangathagroup{\csromangathastanza{โกโธ โมสวชฺชญฺจ กถํกถา จ, \\ เอเตปิ ธมฺมา ทฺวยเมว สนฺเต. \\ กถํกถี ญาณปถาย สิกฺเข, \\ ญตฺวา ปวุตฺตา สมเณน ธมฺมาติ.}}

\proseitem{104}{\csromansymbolfootnote{+}{ขุ 1. 414 ปิฏฺเฐ.}\textbf{สาตํ อสาตญฺจ กุโตนิทานา},}

\prose{\textbf{กิสฺมิํ อสนฺเต น ภวนฺติ เหเต.}}

\csromangathameasure{\textbf{วิภวํ ภวญฺจาปิ ยเมตมตฺถํ,} \\ \textbf{เอตํ เม ปพฺรูหิ ยโตนิทานํ.}}
\csromangathagroup{\csromangathastanza{\textbf{วิภวํ ภวญฺจาปิ ยเมตมตฺถํ,} \\ \textbf{เอตํ เม ปพฺรูหิ ยโตนิทานํ.}}}

\prose{\textbf{สาตํ อสาตญฺจ กุโตนิทานา}ติ สาตา อสาตา กุโตนิทานา กุโตชาตา กุโตสญฺชาตา กุโตนิพฺพตฺตา กุโต-อภินิพฺพตฺตา กุโตปาตุภูตา, กิํนิทานา กิํสมุทยา กิํชาติกา กิํปภวาติ สาตาสาตานํ มูลํ ปุจฺฉติ ฯเปฯ สมุทยํ ปุจฺฉติ ปปุจฺฉติ ยาจติ อชฺเฌสติ ปสาเทตีติ สาตํ อสาตญฺจ กุโตนิทานา.}

\prose{\textbf{กิสฺมิํ อสนฺเต น ภวนฺติ เหเต}ติ กิสฺมิํ อสนฺเต อสํวิชฺชมาเน นตฺถิ อนุปลพฺภมาเน สาตาสาตา น ภวนฺติ นปฺปภวนฺติ น ชายนฺติ น สญฺชายนฺติ น นิพฺพตฺตนฺติ น อภินิพฺพตฺตนฺตีติ กิสฺมิํ อสนฺเต น ภวนฺติ เหเต.}

\prose{\textbf{วิภวํ ภวญฺจาปิ ยเมตมตฺถนฺ}ติ กตโม สาตาสาตานํ ภโว, โย สาตาสาตานํ ภโว ปภโว ชาติ สญฺชาติ นิพฺพตฺติ อภินิพฺพตฺติ ปาตุภาโว, อยํ สาตาสาตานํ ภโว.\csromanspacer{} กตโม สาตาสาตานํ วิภโว, โย สาตาสาตานํ ขโย วโย เภโท ปริเภโท}

\vspace{\tipitakaparskip}
\csromancenter{กลหวิวาทสุตฺตนิทฺเทส อนิจฺจตา อนฺตรธานํ, อยํ สาตาสาตานํ วิภโว.\csromanspacer{}\csromanspacer{}\textbf{ยเมตมตฺถนฺ}ติ ยํ ปรมตฺถนฺติ วิภวํ ภวญฺจาปิ ยเมตมตฺถํ.}
\prose{\csromanfolio{211}\textbf{เอตํ เม ปพฺรูหิ ยโตนิทานนฺ}ติ \textbf{เอตนฺ}ติ ยํ ปุจฺฉามิ, ยํ ยาจามิ, ยํ อชฺเฌสามิ, ยํ ปสาเทมิ.\csromanspacer{} \textbf{ปพฺรูหี}ติ พฺรูหิ วเทหิ อาจิกฺข เทเสหิ ปญฺญเปหิ ปฏฺฐเปหิ วิวร วิภช อุตฺตานีกโรหิ ปกาเสหีติ เอตํ เม ปพฺรูหิ.\csromanspacer{} \textbf{ยโตนิทานนฺ}ติ ยํนิทานํ ยํสมุทยํ ยํชาติกํ ยํปภวนฺติ เอตํ เม ปพฺรูหิ ยโตนิทานํ.\csromanspacer{} เตนาห โส นิมฺมิโต\csromandash{}}

\csromangathameasure{สาตํ อสาตญฺจ กุโตนิทานา, \\ กิสฺมิํ อสนฺเต น ภวนฺติ เหเต. \\ \textbf{วิภวํ ภวญฺจาปิ ยเมตมตฺถํ,}}
\csromangathagroup{\csromangathastanza{สาตํ อสาตญฺจ กุโตนิทานา, \\ กิสฺมิํ อสนฺเต น ภวนฺติ เหเต. \\ \textbf{วิภวํ ภวญฺจาปิ ยเมตมตฺถํ,}}}

\prose{\textbf{เอตํ เม ปพฺรูหิ ยโตนิทานนฺ}ติ.}

\proseitem{105}{\csromansymbolfootnote{*}{ขุ 1. 414 ปิฏฺเฐ.}ผสฺสนิทานํ สาตํ อสาตํ,}

\prose{\textbf{ผสฺเส อสนฺเต น ภวนฺติ เหเต.}}

\csromangathameasure{\textbf{วิภวํ ภวญฺจาปิ ยเมตมตฺถํ,} \\ \textbf{เอตํ เต ปพฺรูมิ อิโตนิทานํ.}}
\csromangathagroup{\csromangathastanza{\textbf{วิภวํ ภวญฺจาปิ ยเมตมตฺถํ,} \\ \textbf{เอตํ เต ปพฺรูมิ อิโตนิทานํ.}}}

\prose{\textbf{ผสฺสนิทานํ สาตํ อสาตนฺ}ติ สุขเวทนียํ ผสฺสํ ปฏิจฺจ อุปฺปชฺชติ สุขา เวทนา.\csromanspacer{} ยา ตสฺเสว สุขเวทนียสฺส ผสฺสสฺส นิโรธา ยํ ตชฺชํ เวทยิตํ สุขเวทนียํ ผสฺสํ ปฏิจฺจ ปฏิจฺจ อุปฺปนฺนา สุขา เวทนา, สา นิรุชฺฌติ, สา วูปสมฺมติ.\csromanspacer{} ทุกฺขเวทนียํ ผสฺสํ ปฏิจฺจ อุปฺปชฺชติ ทุกฺขา เวทนา.\csromanspacer{} ยา ตสฺเสว ทุกฺขเวทนียสฺส ผสฺสสฺส นิโรธา ยํ ตชฺชํ เวทยิตํ ทุกฺขเวทนียํ ผสฺสํ ปฏิจฺจ อุปฺปนฺนา ทุกฺขา เวทนา, สา นิรุชฺฌติ, สา วูปสมฺมติ.\csromanspacer{} อทุกฺขมสุขเวทนียํ ผสฺสํ ปฏิจฺจ อุปฺปชฺชติ อทุกฺขมสุขา เวทนา.\csromanspacer{} ยา ตสฺเสว อทุกฺขมสุขเวทนียสฺส ผสฺสสฺส นิโรธา ยํ ตชฺชํ เวทยิตํ อทุกฺขมสุขเวทนียํ ผสฺสํ ปฏิจฺจ อุปฺปนฺนา อทุกฺขมสุขา เวทนา, สา นิรุชฺฌติ วูปสมฺมติ.\csromanspacer{}\csromanspacer{}\textbf{ผสฺสนิทานํ สาตํ อสาตนฺ}ติ สาตาสาตา ผสฺสนิทานา ผสฺสสมุทยา ผสฺสชาติกา ผสฺสปฺปภวาติ ผสฺสนิทานํ สาตํ อสาตํ.}

\prose{\textbf{ผสฺเส อสนฺเต น ภวนฺติ เหเต}ติ ผสฺเส อสนฺเต อสํวิชฺชมาเน นตฺถิ อนุปลพฺภมาเน สาตาสาตา น ภวนฺติ นปฺปภวนฺติ น ชายนฺติ น สญฺชายนฺติ น นิพฺพตฺตนฺติ นาภินิพฺพตฺตนฺติ น ปาตุภวนฺตีติ ผสฺเส อสนฺเต น ภวนฺติ เหเต.}

\prose{\csromanfolio{212}\textbf{วิภวํ ภวญฺจาปิ ยเมตมตฺถนฺ}ติ ภวทิฏฺฐิปิ ผสฺสนิทานา, วิภวทิฏฺฐิปิ ผสฺสนิทานา.\csromanspacer{} \textbf{ยเมตมตฺถนฺ}ติ ยํ ปรมตฺถนฺติ วิภวํ ภวญฺจาปิ ยมฺ\textbf{เอต}มตฺถํ.}

\prose{\textbf{เอตํ เต ปพฺรูมิ อิโตนิทานนฺ}ติ \textbf{เอต}นฺติ ยํ ปุจฺฉสิ, ยํ ยาจสิ, ยํ อชฺเฌสสิ, ยํ ปสาเทสิ.\csromanspacer{}\csromanspacer{}\textbf{ปพฺรูมี}ติ พฺรูมิ อาจิกฺขามิ เทเสมิ ปญฺญเปมิ ปฏฺฐเปมิ วิวรามิ วิภชามิ อุตฺตานีกโรมิ ปกาเสมีติ \textbf{เอต}ํ เต ปพฺรูมิ.\csromanspacer{} \textbf{อิโตนิทานนฺ}ติ อิโต ผสฺสนิทานํ ผสฺสสมุทยํ ผสฺสชาติกํ ผสฺสปฺปภวนฺติ \textbf{เอต}ํ เต ปพฺรูมิ อิโตนิทานํ.\csromanspacer{} เตนาห ภควา\csromandash{}}

\csromangathameasure{ผสฺสนิทานํ สาตํ อสาตํ, \\ ผสฺเส อสนฺเต น ภวนฺติ เหเต.}
\csromangathagroup{\csromangathastanza{ผสฺสนิทานํ สาตํ อสาตํ, \\ ผสฺเส อสนฺเต น ภวนฺติ เหเต.}}

\prose{วิภวํ ภวญฺจาปิ ยมฺ\textbf{เอต}มตฺถํ,}

\prose{\textbf{เอตํ เต ปพฺรูมิ อิโตนิทานนฺ}ติ.}

\proseitem{106}{\csromansymbolfootnote{*}{ขุ 1. 414 ปิฏฺเฐ.}\textbf{ผสฺโส นุ โลกสฺมิํ กุโตนิทาโน},}

\prose{\textbf{ผริคฺคหา จาปิ}\footnote{วาปิ (สี, สฺยา)} \textbf{กุโตปหูตา.}}

\csromangathameasure{\textbf{กิสฺมิํ อสนฺเต น มมตฺตมตฺถิ,} \\ \textbf{กิสฺมิํ วิภูเต น ผุสนฺติ ผสฺสา.}}
\csromangathagroup{\csromangathastanza{\textbf{กิสฺมิํ อสนฺเต น มมตฺตมตฺถิ,} \\ \textbf{กิสฺมิํ วิภูเต น ผุสนฺติ ผสฺสา.}}}

\prose{\textbf{ผสฺโส นุ โลกสฺมิํ กุโตนิทาโน}ติ ผสฺโส กุโตนิทาโน กุโตชาโต กุโตสญฺชาโต กุโตนิพฺพตฺโต กุโต-อภินิพฺพตฺโต กุโตปาตุภูโต, กิํนิทาโน กิํสมุทโย กิํชาติโก กิํปภโวติ ผสฺสสฺส มูลํ ปุจฺฉติ, เหตุํ ปุจฺฉติ ฯเปฯ สมุทยํ ปุจฺฉติ ปปุจฺฉติ ยาจติ อชฺเฌสติ ปสาเทตีติ ผสฺโส นุ โลกสฺมิํ กุโตนิทาโน.}

\prose{\textbf{ปริคฺคหา จาปิ กุโตปหูตา}ติ ปริคฺคหา กุโตปหูตา กุโตชาตา กุโตสญฺชาตา กุโตนิพฺพตฺตา กุโต-อภินิพฺพตฺตา กุโตปาตุภูตา, กิํนิทานา กิํสมุทยา กิํชาติกา กิํปภวาติ ปริคฺคหานํ มูลํ ปุจฺฉติ, เหตุํ ปุจฺฉติ ฯเปฯ สมุทยํ ปุจฺฉติ ปปุจฺฉติ ยาจติ อชฺเฌสติ ปสาเทตีติ ปริคฺคหา จาปิ \textbf{กุโตปหูตา.}}

\csromanheader{2}{11. กลหวิวาทสุตฺตนิทฺเทส}
\prose{\csromanfolio{213}\textbf{กิสฺมิํ อสนฺเต น มมตฺตมตฺถี}ติ กิสฺมิํ อสนฺเต อสํวิชฺชมาเน นตฺถิ อนุปลพฺภมาเน มมตฺตา นตฺถิ น สนฺติ น สํวิชฺชนฺติ นุปลพฺภนฺติ, ปหีนา สมุจฺฉินฺนา วูปสนฺตา ปฏิปสฺสทฺโธ อภพฺพุปฺปตฺติกา ญาณคฺคินา ทฑฺฒาติ กิสฺมิํ อสนฺเต น มมตฺตมตฺถิ.}

\prose{\textbf{กิสฺมิํ วิภูเต น ผุสนฺติ ผสฺสา}ติ กิสฺมิํ วิภูเต วิภวิเต อติกฺกนฺเต สมติกฺกนฺเต วีติวตฺเต ผสฺสา น ผุสนฺตีติ กิสฺมิํ วิภูเต น ผุสนฺติ ผสฺสา.\csromanspacer{} เตนาห โส นิมฺมิโต\csromandash{}}

\csromangathameasure{ผสฺโส นุ โลกสฺมิํ กุโตนิทาโน, \\ ปริคฺคหา จาปิ กุโตปหูตา.}
\csromangathagroup{\csromangathastanza{ผสฺโส นุ โลกสฺมิํ กุโตนิทาโน, \\ ปริคฺคหา จาปิ กุโตปหูตา.}}

\prose{กิสฺมิํ อสนฺเต น มมตฺตมตฺถิ,}

\prose{\textbf{กิสฺมิํ วิภูเต น ผุสนฺติ ผสฺสา}ติ.}

\proseitem{107}{\csromansymbolfootnote{*}{ขุ 1. 414 ปิฏฺเฐ.}\textbf{นามญฺจ รูปญฺจ ปฏิจฺจ ผสฺโส},}

\prose{\textbf{อิจฺฉานิทานานิ ปริคฺคหานิ.}}

\csromangathameasure{\textbf{อิจฺฉาย’สนฺตฺยา น มมตฺตมตฺถิ,} \\ \textbf{รูเป วิภูเต น ผุสนฺติ ผสฺสา.}}
\csromangathagroup{\csromangathastanza{\textbf{อิจฺฉาย’สนฺตฺยา น มมตฺตมตฺถิ,} \\ \textbf{รูเป วิภูเต น ผุสนฺติ ผสฺสา.}}}

\prose{\textbf{นามญฺจ รูปญฺจ ปฏิจฺจ ผสฺโส}ติ + จกฺขุญฺจ ปฏิจฺจ รูเป จ อุปฺปชฺชติ จกฺขุวิญฺญาณํ, ติณฺณํ สงฺคติ ผสฺโส.\csromanspacer{} จกฺขุ จ\footnote{จกฺขุญฺจ (พหูสุ)} รูปา จ รูปสฺมิํ จกฺขุสมฺผสฺสํ ฐเปตฺวา สมฺปยุตฺตกา ธมฺมา นามสฺมิํ, เอวมฺปิ นามญฺจ รูปญฺจ ปฏิจฺจ ผสฺโส.\csromanspacer{} โสตญฺจ ปฏิจฺจ สทฺเท จ อุปฺปชฺชติ โสตวิญฺญาณํ, ติณฺณํ สงฺคติ ผสฺโส.\csromanspacer{} โสตญฺจ สทฺทา จ รูปสฺมิํ โสตสมฺผสฺสํ ฐเปตฺวา สมฺปยุตฺตกา ธมฺมา นามสฺมิํ, เอวมฺปิ นามญฺจ รูปญฺจ ปฏิจฺจ ผสฺโส.\csromanspacer{} ฆานญฺจ ปฏิจฺจ คนฺเธ จ อุปฺปชฺชติ ฆานวิญฺญาณํ, ติณฺณํ สงฺคติ ผสฺโส.\csromanspacer{} ฆานญฺจ คนฺธา จ รูปสฺมิํ ฆานสมฺผสฺสํ ฐเปตฺวา สมฺปยุตฺตกา ธมฺมา นามสฺมิํ, เอวมฺปิ นามญฺจ รูปญฺจ ปฏิจฺจ ผสฺโส.\csromanspacer{} ชิวฺหญฺจ ปฏิจฺจ รเส จ อุปฺปชฺชติ ชิวฺหาวิญฺญาณํ, ติณฺณํ สงฺคติ ผสฺโส.\csromanspacer{} ชิวฺหา จ รสา จ รูปสฺมิํ ชิวฺหาสมฺผสฺสํ ฐเปตฺวา สมฺปยุตฺตกา ธมฺมา นามสฺมิํ, เอวมฺปิ นามญฺจ รูปญฺจ ปฏิจฺจ ผสฺโส.\csromanspacer{} กายญฺจ ปฏิจฺจ โผฏฺฐพฺเพ จ อุปฺปชฺชติ กายวิญฺญาณํ, ติณฺณํ สงฺคติ ผสฺโส.\csromanspacer{} กาโย จ โผฏฺฐพฺพา จ รูปสฺมิํ กายสมฺผสฺสํ ฐเปตฺวา สมฺปยุตฺตกา ธมฺมา นามสฺมิํ.\csromanspacer{} เอวมฺปิ นามญฺจ รูปญฺจ ปฏิจฺจ ผสฺโส.\csromanspacer{} มนญฺจ ปฏิจฺจ ธมฺเม จ อุปฺปชฺชติ มโนวิญฺญาณํ, ติณฺณํ สงฺคติ ผสฺโส. \csromanfolio{214}วตฺถุ รูปํ รูปสฺมิํ, ธมฺมา รูปิโน รูปสฺมิํ มโนสมฺผสฺสํ ฐเปตฺวา สมฺปยุตฺตกา ธมฺมา นามสฺมิํ, เอวมฺปิ นามญฺจ รูปญฺจ ปฏิจฺจ ผสฺโส.}

\prose{\textbf{อิจฺฉานิทานานิ ปริคฺคหานี}ติ \textbf{อิจฺฉา} วุจฺจติ ตณฺหา.\csromanspacer{} โย ราโค สาราโค ฯเปฯ อภิชฺฌา โลโภ อกุสลมูลํ.\csromanspacer{} \textbf{ปริคฺคหา}ติ \textbf{เทฺว} ปริคฺคหา ตณฺหาปริคฺคโห จ ทิฏฺฐิปริคฺคโห จ ฯเปฯ อยํ ตณฺหาปริคฺคโห ฯเปฯ อยํ ทิฏฺฐิปริคฺคโห.\csromanspacer{} \textbf{อิจฺฉานิทานานิ ปริคฺคหานี}ติ ปริคฺคหา \textbf{อิจฺฉา}นิทานา อิจฺฉเหตุกา \textbf{อิจฺฉา}ปจฺจยา \textbf{อิจฺฉา}การณา \textbf{อิจฺฉา}ปภวาติ \textbf{อิจฺฉา}นิทานานิ ปริคฺคหานิ.}

\prose{\textbf{อิจฺฉาย’สนฺตฺยา น มมตฺตมตฺถี}ติ \textbf{อิจฺฉา} วุจฺจติ ตณฺหา.\csromanspacer{} โย ราโค สาราโค ฯเปฯ อภิชฺฌา โลโภ อกุสลมูลํ.\csromanspacer{} \textbf{มมตฺตา}ติ \textbf{เทฺว} มมตฺตา ตณฺหามมตฺตญฺจ ทิฏฺฐิมมตฺตญฺจ ฯเปฯ อิทํ ตณฺหามมตฺตํ ฯเปฯ อิทํ ทิฏฺฐิมมตฺตํ.\csromanspacer{} \textbf{อิจฺฉาย’สนฺตฺยา น มมตฺตมตฺถี}ติ \textbf{อิจฺฉา}ย อสนฺตฺยา อสํวิชฺชมานาย นตฺถิ อนุปลพฺภมานาย มมตฺตา นตฺถิ น สนฺติ น สํวิชฺชนฺติ นุปลพฺภนฺติ, ปหีนา สมุจฺฉินฺนา วูปสนฺตา ปฏิปสฺสทฺธา อภพฺพุปฺปตฺติกา ญาณคฺคินา ทฑฺฒาติ \textbf{อิจฺฉา}ย’สนฺตฺยา น มมตฺตมตฺถิ.}

\prose{\textbf{รูเป วิภูเต น ผุสนฺติ ผสฺสา}ติ \textbf{รูเป}ติ จตฺตาโร จ มหาภูตา จตุนฺนญฺจ มหาภูตานํ อุปาทาย รูปํ.\csromanspacer{} \textbf{รูเป วิภูเต}ติ จตูหากาเรหิ\footnote{จตูหิ การเณหิ (สฺยา)} รูปํ วิภูตํ โหติ ญาตวิภูเตน\footnote{ญาณวิภูเตน (สี), ตทฏฺฐกถายํ ปน ญาตวีติวตฺเตนาติ ทิสฺสติ.} ตีรณวิภูเตน ปหานวิภูเตน สมติกฺกมวิภูเตน.\csromanspacer{} กถํ ญาตวิภูเตน รูปํ วิภูตํ โหติ, รูปํ ชานาติ “ยํ กิญฺจิ รูปํ, สพฺพํ รูปํ จตฺตาริ จ มหาภูตานิ จตุนฺนญฺจ มหาภูตานํ อุปาทาย รูปนฺ”ติ ชานาติ ปสฺสติ.\csromanspacer{} เอวํ ญาตวิภูเตน รูปํ วิภูตํ โหติ.}

\prose{กถํ ตีรณวิภูเตน รูปํ วิภูตํ โหติ, เอวํ ญาตํ กตฺวา รูปํ ตีเรติ, อนิจฺจโต ทุกฺขโต โรคโต คณฺฑโต สลฺลโต อฆโต อาพาธโต ปรโต ปโลกโต อีติโต อุปทฺทวโต ภยโต อุปสคฺคโต จลโต ปภงฺคุโต อธุวโต อตาณโต อเลณโต อสรณโต ริตฺตโต ตุจฺฉโต สุญฺญโต อนตฺตโต อาทีนวโต วิปริณามธมฺมโต อสารกโต อฆมูลโต}

\vspace{\tipitakaparskip}
\csromancenter{กลหวิวาทสุตฺตนิทฺเทส วธกโต วิภวโต สาสวโต สงฺขตโต มารามิสโต ชาติธมฺมโต ชราธมฺมโต พฺยาธิธมฺมโต มรณธมฺมโต โสกปริเทวทุกฺขโทมนสฺสุปายาสธมฺมโต สํกิเลสิกธมฺมโต สมุทยโต อตฺถงฺคมโต อสฺสาทโต อาทีนวโต นิสฺสรณโต ตีเรติ.\csromanspacer{} เอวํ ตีรณวิภูเตน รูปํ วิภูตํ โหติ.}
\prose{\csromanfolio{215}กถํ ปหานวิภูเตน รูปํ วิภูตํ โหติ, เอวํ ตีรยิตฺวา รูเป ฉนฺทราคํ ปชหติ วิโนเทติ พฺยนฺติํ กโรติ อนภาวํ คเมติ, วุตฺตํเหตํ ภควตา\csromandash{} * โย ภิกฺขเว รูปสฺมิํ ฉนฺทราโค, ตํ ปชหถ.\csromanspacer{} เอวํ ตํ รูปํ ปหีนํ ภวิสฺสติ อุจฺฉินฺนมูลํ ตาลาวตฺถุกตํ อนภาวํกตํ อายติํ อนุปฺปาทธมฺมนฺติ.\csromanspacer{} เอวํ ปหานวิภูเตน รูปํ วิภูตํ โหติ.}

\prose{กถํ สมติกฺกมวิภูเตน รูปํ วิภูตํ โหติ, จตสฺโส อรูปสมาปตฺติโย ปฏิลทฺธสฺส รูปา วิภูตา โหนฺติ วิภาวิตา อติกฺกนฺตา สมติกฺกนฺตา วีติวตฺตา.\csromanspacer{} เอวํ สมติกฺกมวิภูเตน รูปํ วิภูตํ โหติ.\csromanspacer{} อิเมหิ จตูหิ การเณหิ รูปํ วิภูตํ โหติ.}

\prose{\textbf{รูเป วิภูเต น ผุสนฺติ ผสฺสา}ติ รูเป วิภูเต วิภาวิเต อติกฺกนฺเต สมติกฺกนฺเต วีติวตฺเต ปญฺจ ผสฺสา น ผุสนฺติ จกฺขุสมฺผสฺโส โสตสมฺผสฺโส ฆานสมฺผสฺโส ชิวฺหาสมฺผสฺโส กายสมฺผสฺโสติ รูเป วิภูเต น ผุสนฺติ ผสฺสา.\csromanspacer{} เตนาห ภควา\csromandash{}}

\csromangathameasure{นามญฺจ รูปญฺจ ปฏิจฺจ ผสฺโส, \\ อิจฺฉานิทานานิ ปริคฺคหานิ. \\ อิจฺฉาย’สนฺตฺยา น มมตฺตมตฺถิ,}
\csromangathagroup{\csromangathastanza{นามญฺจ รูปญฺจ ปฏิจฺจ ผสฺโส, \\ อิจฺฉานิทานานิ ปริคฺคหานิ. \\ อิจฺฉาย’สนฺตฺยา น มมตฺตมตฺถิ,}}

\prose{\textbf{รูเป วิภูเต น ผุสนฺติ ผสฺสา}ติ.}

\proseitem{108}{\csromansymbolfootnote{+}{ขุ 1. 414 ปิฏฺเฐ.}\textbf{กถํ สเมตสฺส วิโภติ รูป}ํ,}

\prose{\textbf{สุขํ ทุขญฺจาปิ}\footnote{ทุกฺขํ วาปิ (สฺยา)} \textbf{กถํ วิโภติ.}}

\csromangathameasure{\textbf{เอตํ เม ปพฺรูหิ ยถา วิโภติ,} \\ \textbf{ตํ ชานิยามาติ} \textbf{เม มโน อหุ.}}
\csromangathagroup{\csromangathastanza{\textbf{เอตํ เม ปพฺรูหิ ยถา วิโภติ,} \\ \textbf{ตํ ชานิยามาติ}\footnote{ชานิสฺสามาติ (สี, ก)} \textbf{เม มโน อหุ.}}}

\prose{\textbf{กถํ สเมตสฺส วิโภติ รูป}นฺติ \textbf{กถํ สเมตสฺสา}ติ กถํ สเมตสฺส กถํ ปฏิปนฺนสฺส กถํ อิริยนฺตสฺส กถํ วตฺเตนฺตสฺส กถํ ปาเลนฺตสฺส กถํ ยเปนฺตสฺส กถํ ยาเปนฺตสฺส รูปํ วิโภติ วิภาวียติ \csromanfolio{216}อติกฺกมียติ สมติกฺกมียติ\footnote{วิภาวิยฺยติ อติกฺกมิยฺยติ สมติกฺกมิยฺยติ (พหูสุ)} วีติวตฺตียตีติ กถํ สเมตสฺส วิโภติ รูปํ.}

\prose{\textbf{สุขํ ทุขญฺจาปิ กถํ วิโภตี}ติ สุขญฺจ ทุกฺขญฺจ กถํ วิโภติ วิภาวียติ อติกฺกมียติ สมติกฺกมียติ วีติวตฺตียตีติ สุขํ ทุขญฺจาปิ กถํ วิโภติ.}

\prose{\textbf{เอตํ เม ปพฺรูหิ ยถา วิโภตี}ติ \textbf{เอตนฺ}ติ ยํ ปุจฺฉามิ, ยํ ยาจามิ, ยํ อชฺเฌสามิ, ยํ ปสาเทมีติ เอตํ.\csromanspacer{} \textbf{เม ปพฺรูหี}ติ เม ปพฺรูหิ อาจิกฺข เทเสหิ ปญฺญเปหิ ปฏฺฐเปหิ วิวร วิภช อุตฺตานีกโรหิ ปกาเสหีติ เอตํ เม ปพฺรูหิ.\csromanspacer{} \textbf{ยถา วิโภตี}ติ ยถา วิโภติ วิภาวียติ อติกฺกมียติ สมติกฺกมียติ วีติวตฺตียตีติ เอตํ เม ปพฺรูหิ ยถา วิโภติ.}

\prose{\textbf{ตํ ชานิยามาติ เม มโน อหู}ติ \textbf{ตํ ชานิยามา}ติ ตํ ชาเนยฺยาม อาชาเนยฺยาม วิชาเนยฺยาม ปฏิวิชาเนยฺยาม ปฏิวิชฺเฌยฺยามาติ ตํ ชานิยาม.\csromanspacer{} \textbf{อิติ เม มโน อหู}ติ อิติ เม มโน อหุ, อิติ เม จิตฺตํ อหุ, อิติ เม สงฺกปฺโป อหุ, อิติ เม วิญฺญาณํ อหูติ ตํ ชานิยาม อิติ เม มโน อหุ.\csromanspacer{} เตนาห โส นิมฺมิโต\csromandash{}}

\csromangathameasure{กถํ สเมตสฺส วิโภติ รูปํ,}
\csromangathagroup{\csromangathastanza{กถํ สเมตสฺส วิโภติ รูปํ,}}

\prose{สุขํ ทุขญฺจาปิ กถํ วิโภติ.}

\prose{เอตํ เม ปพฺรูหิ ยถา วิโภติ,}

\prose{\textbf{ตํ ชานิยามาติ เม มโน อหู}ติ.}

\proseitem{109}{\csromansymbolfootnote{*}{ขุ 1. 414 ปิฏฺเฐ.}\textbf{น สญฺญสญฺญี น วิสญฺญสญฺญี},}

\prose{\textbf{โนปิ อสญฺญี น วิภูตสญฺญี.}}

\csromangathameasure{\textbf{เอวํ สเมตสฺส วิโภติ รูปํ,} \\ \textbf{สญฺญานิทานา หิ ปปญฺจสงฺขา.}}
\csromangathagroup{\csromangathastanza{\textbf{เอวํ สเมตสฺส วิโภติ รูปํ,} \\ \textbf{สญฺญานิทานา หิ ปปญฺจสงฺขา.}}}

\prose{\textbf{น สญฺญสญฺญี น วิสญฺญสญฺญี}ติ สญฺญสญฺญิโน วุจฺจนฺติ เย ปกติสญฺญาย ฐิตา, นปิ โส ปกติสญฺญาย ฐิโต.\csromanspacer{} วิสญฺญสญฺญิโน วุจฺจนฺติ อุมฺมตฺตกา เย จ ขิตฺตจิตฺตา\footnote{อุกฺขิตฺตจิตฺตา (สฺยา)}, นปิ โส อุมฺมตฺตโก, โนปิ ขิตฺตจิตฺโตติ น สญฺญสญฺญี น วิสญฺญสญฺญี.}

\csromanheader{2}{11. กลหวิวาทสุตฺตนิทฺเทส}
\prose{\csromanfolio{217}\textbf{โนปิ อสญฺญี น วิภูตสญฺญี}ติ อสญฺญิโน วุจฺจนฺติ นิโรธสมาปนฺนา เย จ อสญฺญสตฺตา.\csromanspacer{} นปิ โส นิโรธสมาปนฺโน, นปิ อสญฺญสตฺโต.\csromanspacer{} วิภูตสญฺญิโน วุจฺจนฺติ เย จตุนฺนํ อรูปสมาปตฺตีนํ ลาภิโน.\csromanspacer{} นปิ โส จตุนฺนํ อรูปสมาปตฺตีนํ ลาภีติ โนปิ อสญฺญี น วิภูตสญฺญี.}

\prose{\textbf{เอวํ สเมตสฺส วิโภติ รูปนฺ}ติ อิธ ภิกฺขุ สุขสฺส จ ปหานา ฯเปฯ จตุตฺถํ ฌานํ อุปสมฺปชฺช วิหรติ, โส เอวํ สมาหิเต จิตฺเต ปริสุทฺเธ ปริโยทาเต อนงฺคเณ วิคตูปกฺกิเลเส มุทุภูเต กมฺมนิเย ฐิเต อาเนญฺชปฺปตฺเต อากาสานญฺจายตนสมาปตฺติปฏิลาภตฺถาย จิตฺตํ อภินีหรติ อภินินฺนาเมติ อารุปฺปมคฺคสมงฺคีติ เอวํ สเมตสฺส เอวํ ปฏิปนฺนสฺส เอวํ อิริยนฺตสฺส เอวํ วตฺเตนฺตสฺส เอวํ ปาเลนฺตสฺส เอวํ ยเปนฺตสฺส เอวํ ยาเปนฺตสฺส รูปํ วิโภติ วิภาวียติ อติกฺกมียติ สมติกฺกมียติ วีติวตฺตียตีติ เอวํ สเมตสฺส วิโภติ รูปํ.}

\prose{\textbf{สญฺญานิทานา หิ ปปญฺจสงฺขา}ติ ปปญฺจาเยว ปปญฺจสงฺขา.\csromanspacer{} ตณฺหาปปญฺจสงฺขา ทิฏฺฐิปปญฺจสงฺขา มานปปญฺจสงฺขา สญฺญานิทานา สญฺญาสมุทยา สญฺญาชาติกา สญฺญาปภวาติ สญฺญานิทานา หิ ปปญฺจสงฺขา.\csromanspacer{} เตนาห ภควา\csromandash{}}

\csromangathameasure{น สญฺญสญฺญี น วิสญฺญสญฺญี, \\ \textbf{โนปิ อสญฺญี น วิภูตสญฺญี}. \\ เอวํ สเมตสฺส วิโภติ รูปํ,}
\csromangathagroup{\csromangathastanza{น สญฺญสญฺญี น วิสญฺญสญฺญี, \\ \textbf{โนปิ อสญฺญี น วิภูตสญฺญี}. \\ เอวํ สเมตสฺส วิโภติ รูปํ,}}

\prose{\textbf{สญฺญานิทานา หิ ปปญฺจสงฺขา}ติ.}

\proseitem{110}{\csromansymbolfootnote{*}{ขุ 1. 415; ขุ 8. 71 ปิฏฺเฐสุปิ.}\textbf{ยํ ตํ อปุจฺฉิมฺห อกิตฺตยี โน},}

\prose{\textbf{อญฺญํ ตํ ปุจฺฉาม ตทิงฺฆ พฺรูหิ.}}

\csromangathameasure{\textbf{เอตฺตาวต’คฺคํ นุ วทนฺติ เหเก,} \\ \textbf{ยกฺขสฺส สุทฺธิํ อิธ ปณฺฑิตาเส.} \\ \textbf{อุทาหุ อญฺญมฺปิ วทนฺติ เอตฺโต.}}
\csromangathagroup{\csromangathastanza{\textbf{เอตฺตาวต’คฺคํ นุ วทนฺติ เหเก,} \\ \textbf{ยกฺขสฺส สุทฺธิํ อิธ ปณฺฑิตาเส.} \\ \textbf{อุทาหุ อญฺญมฺปิ วทนฺติ เอตฺโต.}}}

\prose{\textbf{ยํ ตํ อปุจฺฉิมฺห อกิตฺตยี โน}ติ ยํ ตํ อปุจฺฉิมฺห อยาจิมฺห อชฺเฌสิมฺห ปสาทยิมฺห.\csromanspacer{} \textbf{อกิตฺตยี โน}ติ กิตฺติตํ ปกิตฺติตํ อาจิกฺขิตํ เทสิตํ ปญฺญปิตํ ปฏฺฐปิตํ วิวฏํ วิภตฺตํ อุตฺตานีกตํ ปกาสิตนฺติ ยํ ตํ อปุจฺฉิมฺห อกิตฺตยี โน.}

\prose{\textbf{อญฺญํ ตํ ปุจฺฉาม ตทิงฺฆ พฺรูหี}ติ อญฺญํ ตํ ปุจฺฉาม, อญฺญํ ตํ ยาจาม, อญฺญํ ตํ อชฺเฌสาม, อญฺญํ ตํ ปสาเทม, อุตฺตริํ ตํ ปุจฺฉาม.\csromanspacer{}\csromanspacer{}\textbf{ตทิงฺฆ พฺรูหี}ติ \csromanfolio{218}อิงฺฆ พฺรูหิ อาจิกฺข เทเสหิ ปญฺญเปหิ ปฏฺฐเปหิ วิวร วิภช อุตฺตานีกโรหิ ปกาเสหีติ อญฺญํ ตํ ปุจฺฉาม ตทิงฺฆ พฺรูหิ.}

\prose{\textbf{เอตฺตาวต’คฺคํ นุ วทนฺติ เหเก, ยกฺขสฺส สุทฺธิํ อิธ ปณฺฑิตาเส}ติ เอเก สมณพฺราหฺมณา เอตา อรูปสมาปตฺติโย อคฺคํ เสฏฺฐํ วิสิฏฺฐํ ปาโมกฺขํ อุตฺตมํ ปวรํ วทนฺติ กเถนฺติ ภณนฺติ ทีปยนฺติ โวหรนฺติ.\csromanspacer{} \textbf{ยกฺขสฺสา}ติ สตฺตสฺส นรสฺส มานวสฺส โปสสฺส ปุคฺคลสฺส ชีวสฺส ชาคุสฺส ชนฺตุสฺส อินฺทคุสฺส มนุชสฺส.\csromanspacer{}\csromanspacer{}\textbf{สุทฺธินฺ}ติ สุทฺธิํ วิสุทฺธิํ ปริสุทฺธิํ, มุตฺติํ วิมุตฺติํ ปริมุตฺติํ.\csromanspacer{}\csromanspacer{}\textbf{อิธ ปณฺฑิตาเส}ติ อิธ ปณฺฑิตวาทา ถิรวาทา ญายวาทา เหตุวาทา ลกฺขณวาทา การณวาทา ฐานวาทา สกาย ลทฺธิยาติ เอตฺตาวต’คฺคํ นุ วทนฺติ เหเก, ยกฺขสฺส สุทฺธิํ อิธ ปณฺฑิตาเส.}

\prose{\textbf{อุทาหุ อญฺญมฺปิ วทนฺติ เอตฺโต}ติ อุทาหุ เอเก สมณพฺราหฺมณา เอตา อรูปสมาปตฺติโย อติกฺกมิตฺวา สมติกฺกมิตฺวา วีติวตฺเตตฺวา เอตฺโต อรูปสมาปตฺติโต อญฺญํ อุตฺตริํ ยกฺขสฺส สุทฺธิํ วิสุทฺธิํ ปริสุทฺธิํ มุตฺติํ วิมุตฺติํ ปริมุตฺติํ วทนฺติ กเถนฺติ ภณนฺติ ทีปยนฺติ โวหรนฺตีติ อุทาหุ อญฺญมฺปิ วทนฺติ เอตฺโต.\csromanspacer{} เตนาห โส นิมฺมิโต\csromandash{}}

\csromangathameasure{ยํ ตํ อปุจฺฉิมฺห อกิตฺตยี โน, \\ อญฺญํ ตํ ปุจฺฉาม ตทิงฺฆ พฺรูหิ.}
\csromangathagroup{\csromangathastanza{ยํ ตํ อปุจฺฉิมฺห อกิตฺตยี โน, \\ อญฺญํ ตํ ปุจฺฉาม ตทิงฺฆ พฺรูหิ.}}

\prose{เอตฺตาวต’คฺคํ นุ วทนฺติ เหเก,}

\prose{\textbf{ยกฺขสฺส สุทฺธิํ อิธ ปณฺฑิตาเส.}}

\prose{\textbf{อุทาหุ อญฺญมฺปิ วทนฺติ เอตฺโต}ติ.}

\proseitem{111}{\csromansymbolfootnote{*}{ขุ 1. 415 ปิฏฺเฐ.}เอตฺตาวต’คฺคมฺปิ วทนฺติ เหเก,}

\prose{\textbf{ยกฺขสฺส สุทฺธิํ อิธ ปณฺฑิตาเส.}}

\csromangathameasure{\textbf{เตสํ ปเนเก สมยํ วทนฺติ,} \\ \textbf{อนุปาทิเสเส กุสลาวทานา.}}
\csromangathagroup{\csromangathastanza{\textbf{เตสํ ปเนเก สมยํ วทนฺติ,} \\ \textbf{อนุปาทิเสเส กุสลาวทานา.}}}

\prose{\textbf{เอตฺตาวตคฺคมฺปิ วทนฺติ เหเก, ยกฺขสฺส สุทฺธิํ อิธ ปณฺฑิตาเส}ติ \textbf{สนฺเตเก} สมณพฺราหฺมณา สสฺสตวาทา เอตา อรูปสมาปตฺติโย อคฺคํ เสฏฺฐํ วิสิฏฺฐํ ปาโมกฺขํ อุตฺตมํ ปวรํ วทนฺติ กเถนฺติ ภณนฺติ ทีปยนฺติ โวหรนฺติ.\csromanspacer{} \textbf{ยกฺขสฺสา}ติ สตฺตสฺส นรสฺส มานวสฺส โปสสฺส ปุคฺคลสฺส ชีวสฺส ชาคุสฺส ชนฺตุสฺส อินฺทคุสฺส มนุชสฺส.\csromanspacer{}\csromanspacer{}\textbf{สุทฺธินฺ}ติ สุทฺธิํ วิสุทฺธิํ ปริสุทฺธิํ, มุตฺติํ วิมุตฺติํ ปริมุตฺติํ.\csromanspacer{}\csromanspacer{}\textbf{อิธ ปณฺฑิตาเส}ติ อิธ ปณฺฑิตวาทา ถิรวาทา ญายวาทา}

\vspace{\tipitakaparskip}
\csromancenter{กลหวิวาทสุตฺตนิทฺเทส เหตุวาทา ลกฺขณวาทา การณวาท, ฐานวาทา สกาย ลทฺธิยาติ เอตฺตาวต’คฺคมฺปิ วทนฺติ เหเก, ยกฺขสฺส สุทฺธิํ อิธ ปณฺฑิตาเส.}
\prose{\csromanfolio{219}\textbf{เตสํ ปเนเก สมยํ วทนฺติ, อนุปาทิเสเส กุสลาวทานา}ติ เตสํเยว สมณพฺราหฺมณานํ เอเก สมณพฺราหฺมณา อุจฺเฉทวาทา ภวตชฺชิตา วิภวํ อภินนฺทนฺติ, เต สตฺตสฺส สมํ อุปสมํ วูปสมํ นิโรธํ ปฏิปสฺสทฺธินฺติ วทนฺติ, ยโต กิํ โภ อยํ อตฺตา กายสฺส เภทา อุจฺฉิชฺชติ วินสฺสติ น โหติ ปรํ มรณา, เอตฺตาวตา อนุปาทิเสโสติ.\csromanspacer{} \textbf{กุสลาวทานา}ติ กุสลวาทา ปณฺฑิตวาทา ถิรวาทา ญายวาทา เหตุวาทา ลกฺขณวาทา การณวาทา ฐานวาทา สกาย ลทฺธิยาติ เตสํ ปเนเก สมยํ วทนฺติ, อนุปาทิเสเส กุสลาวทานา.\csromanspacer{} เตนาห ภควา\csromandash{}}

\csromangathameasure{เอตฺตาวต’คฺคมฺปิ วทนฺติ เหเก, \\ ยกฺขสฺส สุทฺธิํ อิธ ปณฺฑิตาเส.}
\csromangathagroup{\csromangathastanza{เอตฺตาวต’คฺคมฺปิ วทนฺติ เหเก, \\ ยกฺขสฺส สุทฺธิํ อิธ ปณฺฑิตาเส.}}

\prose{เตสํ ปเนเก สมยํ วทนฺติ,}

\prose{อนุปาทิเสเส กุสลาวทานาติ.}

\proseitem{112}{\csromansymbolfootnote{*}{ขุ 1. 415 ปิฏฺเฐ.}\textbf{เอเต จ ญ ตฺวา อุปนิสฺสิตา}ติ,}

\prose{\textbf{ญตฺวา มุนี นิสฺสเย โส วิมํสี.}}

\csromangathameasure{\textbf{ญตฺวา วิมุตฺโต น วิวาท’เมติ,} \\ \textbf{ภวาภวาย น สเมติ ธีโร.}}
\csromangathagroup{\csromangathastanza{\textbf{ญตฺวา วิมุตฺโต น วิวาท’เมติ,} \\ \textbf{ภวาภวาย น สเมติ ธีโร.}}}

\prose{\textbf{เอเต จ ญตฺวา อุปนิสฺสิตา}ติ \textbf{เอเต}ติ ทิฏฺฐิคติเก.\csromanspacer{} \textbf{อุปนิสฺสิตา}ติ “สสฺสตทิฏฺฐินิสฺสิตา”ติ ญตฺวา, “อุจฺเฉททิฏฺฐินิสฺสิตา”ติ ญตฺวา, “สสฺสตุจฺเฉททิฏฺฐินิสฺสิตา”ติ ญตฺวา ชานิตฺวา ตุลยิตฺวา ตีรยิตฺวา วิภาวยิตฺวา วิภูตํ กตฺวาติ \textbf{เอเต} จ ญตฺวา อุปนิสฺสิตาติ.}

\csromanmatikaanchor{mtk.12}
\csromanmatikaanchor{mtk.13}
\csromantocmarkref{subsection}{11. กลหวิวาทสุตฺตนิทฺเทส}{mtk.12}
\bookmark[dest={mtk.12},level=2]{11. กลหวิวาทสุตฺตนิทฺเทส}
\csromantocmarkref{subsection}{12. จูฬวิยูหสุตฺตนิทฺเทส}{mtk.13}
\bookmark[dest={mtk.13},level=2]{12. จูฬวิยูหสุตฺตนิทฺเทส}
\prose{\textbf{ญตฺวา มุนี นิสฺสเย โส วิมํสี}ติ \textbf{มุนี}ติ \textbf{โมนํ} วุจฺจติ ญาณํ ฯเปฯ สงฺคชาลมติจฺจ โส มุนิ.\csromanspacer{} มุนิ “สสฺสตทิฏฺฐินิสฺสิตา”ติ ญตฺวา, “อุจฺเฉททิฏฺฐินิสฺสิตา”ติ ญตฺวา, “สสฺสตุจฺเฉททิฏฺฐินิสฺสิตา”ติ ญตฺวา ชานิตฺวา ตุลยิตฺวา ตีรยิตฺวา วิภาวยิตฺวา วิภูตํ กตฺวา.\csromanspacer{}\csromanspacer{}\textbf{โส วิมํสี}ติ ปณฺฑิโต ปญฺญวา พุทฺธิมา ญาณี วิภาวี เมธาวีติ ญตฺวา มุนิ นิสฺสเย โส วิมํสี.\csromanspacer{}\csromanspacer{}\textbf{ญตฺวา วิมุตฺโต น วิวาท เมตี}ติ ญตฺวา ชานิตฺวา ตุลยิตฺวา ตีรยิตฺวา วิภาวยิตฺวา วิภูตํ กตฺวา.\csromanspacer{} \textbf{วิมุตฺโต}ติ มุตฺโต วิมุตฺโต ปริมุตฺโต สุวิมุตฺโต อจฺจนฺต-อนุปาทาวิโมกฺเขน, “สพฺเพ สงฺขารา อนิจฺจา”ติ ญตฺวา ชานิตฺวา ตุลยิตฺวา ตีรยิตฺวา วิภาวยิตฺวา วิภูตํ \csromanfolio{220}กตฺวา มุตฺโต วิมุตฺโต ปริมุตฺโต สุวิมุตฺโต อจฺจนฺต-อนุปาทาวิโมกฺเขน, “สพฺเพ สงฺขารา ทุกฺขา”ติ. “สพฺเพ ธมฺมา อนตฺตา”ติ ฯเปฯ “ยํ กิญฺจิ สมุทยธมฺมํ, สพฺพํ ตํ นิโรธธมฺมนฺ”ติ ญตฺวา ชานิตฺวา ตุลยิตฺวา ตีรยิตฺวา วิภาวยิตฺวา วิภูตํ กตฺวา มุตฺโต วิมุตฺโต ปริมุตฺโต สุวิมุตฺโต อจฺจนฺต-อนุปาทาวิโมกฺเขนาติ ญตฺวา วิมุตฺโต.\csromanspacer{}\csromanspacer{}\textbf{น วิวาทเมตี}ติ น กลหํ กโรติ.\csromanspacer{} น ภณฺฑนํ กโรติ.\csromanspacer{} น วิคฺคหํ กโรติ, น วิวาทํ กโรติ, น เมธคํ กโรติ, วุตฺตํ เหตํ ภควตา\csromandash{} * เอวํ วิมุตฺตจิตฺโต โข อคฺคิเวสฺสน ภิกฺขุ น เกนจิ สํวทติ, น เกนจิ วิวทติ.\csromanspacer{} ยญฺจ โลเก วุตฺตํ, เตน จ โวหรติ อปรามสนฺติ ญตฺวา วิมุตฺโต น วิวาทเมติ.}

\prose{\textbf{ภวาภวาย น สเมติ ธีโร}ติ \textbf{ภวาภวายา}ติ ภวาย กมฺมภวาย ปุนพฺภวาย กามภวาย, กมฺมภวาย กามภวาย ปุนพฺภวาย รูปภวาย, กมฺมภวาย รูปภวาย ปุนพฺภวาย อรูปภวาย, กมฺมภวาย อรูปภวาย ปุนพฺภวาย ปุนปฺปุนพฺภวาย ปุนปฺปุนคติยา ปุนปฺปุน-อุปปตฺติยา ปุนปฺปุนปฏิสนฺธิยา ปุนปฺปุน-อตฺตภาวาย ปุนปฺปุนาภินิพฺพตฺติยา น สเมติ น สมาคจฺฉติ น คณฺหาติ น ปรามสติ นาภินิวิสติ.\csromanspacer{} \textbf{ธีโร}ติ ธีโร ปณฺฑิโต ปญฺญวา พุทฺธิมา ญาณี วิภาวี เมธาวีติ ภวาภวาย น สเมติ ธีโร.\csromanspacer{} เตนาห ภควา\csromandash{}}

\csromangathameasure{เอเต จ ญตฺวา อุปนิสฺสิตาติ, \\ ญตฺวา มุนี นิสฺสเย โส วิมํสี. \\ ญตฺวา วิมุตฺโต น วิวาทเมติ,}
\csromangathagroup{\csromangathastanza{เอเต จ ญตฺวา อุปนิสฺสิตาติ, \\ ญตฺวา มุนี นิสฺสเย โส วิมํสี. \\ ญตฺวา วิมุตฺโต น วิวาทเมติ,}}

\prose{\textbf{ภวาภวาย น สเมติ ธีโร}ติ.}

\prose{กลหวิวาทสุตฺตนิทฺเทโส เอกาทสโม.}
\csromansectionrule

\csromanheader{2}{12. จูฬพฺยูหสุตฺตนิทฺเทส}
\csromanheader{2}{อถ จูฬพฺยูหสุตฺตนิทฺเทสํ\footnote{จุฬวิยูหสุตฺตนิทฺเทสํ (สี, สฺยา)} วกฺขติ\csromandash{}}
\proseitem{113}{\csromansymbolfootnote{+}{ขุ 1. 415 ปิฏฺเฐ.}สกํ สกํ ทิฏฺฐิปริพฺพสานา,}

\prose{\textbf{วิคฺคยฺห นานา กุสลา วทนฺติ.}}

\csromangathameasure{\textbf{โย เอวํ ชานาติ} \textbf{ส เวทิ ธมฺมํ,} \\ \textbf{อิทํ ปฏิกฺโกส’มเกวลี โส.}}
\csromangathagroup{\csromangathastanza{\textbf{โย เอวํ ชานาติ}\footnote{เอวํ ปชานาติ (สี)} \textbf{ส เวทิ ธมฺมํ,} \\ \textbf{อิทํ ปฏิกฺโกส’มเกวลี โส.}}}

\csromanheader{2}{12. จูฬพฺยูหสุตฺตนิทฺเทส}
\prose{\csromanfolio{221}\textbf{สกํ สกํ ทิฏฺฐิปริพฺพสานา}ติ สนฺเตเก สมณพฺราหฺมณา ทิฏฺฐิคติกา, เต ทฺวาสฏฺฐิยา ทิฏฺฐิคตานํ อญฺญตรญฺญตรํ ทิฏฺฐิคตํ คเหตฺวา อุคฺคเหตฺวา คณฺหิตฺวา ปรามสิตฺวา อภินิวิสิตฺวา สกาย สกาย ทิฏฺฐิยา วสนฺติ สํวสนฺติ อาวสนฺติ ปริวสนฺติ, ยถา อคาริกา ฆเรสุ วสนฺติ สาปตฺติกา วา อาปตฺตีสุ วสนฺติ, สกิเลสา วา กิเลเสสุ วสนฺติ.\csromanspacer{} เอวเมวํ สนฺเตเก สมณพฺราหฺมณา ทิฏฺฐิคติกา, เต ทฺวาสฏฺฐิยา ทิฏฺฐิคตานํ อญฺญตรญฺญตรํ ทิฏฺฐิคตํ คเหตฺวา อุคฺคเหตฺวา คณฺหิตฺวา ปรามสิตฺวา อภินิวิสิตฺวา สกาย สกาย ทิฏฺฐิยา วสนฺติ สํวสนฺติ อาวสนฺติ ปริวสนฺตีติ สกํ สกํ ทิฏฺฐิปริพฺพสานา.}

\prose{\textbf{วิคฺคยฺห นานา กุสลา วทนฺตี}ติ \textbf{วิคฺคยฺห}ติ คเหตฺวา อุคฺคเหตฺวา คณฺหิตฺวา ปรามสิตฺวา อภินิวิสิตฺวา นานา วทนฺติ วิวิธํ วทนฺติ อญฺโญญฺญํ วทนฺติ ปุถุ\footnote{ปุถุํ (สี)} วทนฺติ, น เอกํ วทนฺติ กเถนฺติ ภณนฺติ ทีปยนฺติ โวหรนฺติ.\csromanspacer{} \textbf{กุสลา}ติ กุสลวาทา ปณฺฑิตวาทา ถิรวาทา ญายวาทา เหตุวาทา ลกฺขณวาทา การณวาทา ฐานวาทา สกาย ลทฺธิยาติ \textbf{วิคฺคยฺห} นานา กุสลา วทนฺติ.}

\prose{\textbf{โย เอวํ ชานาติ สเวทิ ธมฺมนฺ}ติ โย อิมํ\footnote{อิทํ (สี, ก)} ธมฺมํ ทิฏฺฐิํ ปฏิปทํ มคฺคํ ชานาติ.\csromanspacer{} โส ธมฺมํ เวทิ อญฺญาสิ อปสฺสิ ปฏิวิชฺฌีติ โย เอวํ ชานาติ ส เวทิ ธมฺมํ.}

\prose{\textbf{อิทํ ปฏิกฺโกส’มเกวลี โส}ติ โย อิมํ ธมฺมํ ทิฏฺฐิํ ปฏิปทํ มคฺคํ ปฏิกฺโกสติ.\csromanspacer{} อเกวลี โส อสมตฺโต โส อปริปุณฺโณ โส หีโน นิหีโน โอมโก ลามโก ฉตุกฺโก ปริตฺโตติ อิทํ ปฏิกฺโกส’มเกวลี โส.\csromanspacer{} เตนาห โส นิมฺมิโต\csromandash{}}

\prose{\textbf{สกํ สกํ ทิฏฺฐิปริพฺพสานา},}

\prose{วิคฺคยฺห นานา กุสลา วทนฺติ.}

\csromangathameasure{โย เอวํ ชานาติ ส เวทิ ธมฺมํ,}
\csromangathagroup{\csromangathastanza{โย เอวํ ชานาติ ส เวทิ ธมฺมํ,}}

\prose{\textbf{อิทํ ปฏิกฺโกส’มเกวลี โส}ติ.}

\proseitem{114}{\csromansymbolfootnote{*}{ขุ 1. 415 ปิฏฺเฐ.}เอวมฺปิ \textbf{วิคฺคยฺห} วิวาทยนฺติ,}

\prose{\textbf{พาโล ปโร อกฺกุสโลติ จาหุ.}}

\csromangathameasure{\textbf{สจฺโจ นุ วาโท กตโม อิเมสํ,} \\ \textbf{สพฺเพว หีเม กุสลาวทานา.}}
\csromangathagroup{\csromangathastanza{\textbf{สจฺโจ นุ วาโท กตโม อิเมสํ,} \\ \textbf{สพฺเพว หีเม กุสลาวทานา.}}}

\prose{\csromanfolio{222}\textbf{เอวมฺปิ วิคฺคยฺห วิวาทยนฺตี}ติ เอวํ คเหตฺวา อุคฺคเหตฺวา คณฺหิตฺวา ปรามสิตฺวา อภินิวิสิตฺวา วิวาทยนฺติ กลหํ กโรนฺติ, ภณฺฑนํ กโรนฺติ, วิคฺคหํ กโรนฺติ, วิวาทํ กโรนฺติ, เมธคํ กโรนฺติ. “น ตฺวํ อิมํ ธมฺมวินยํ อาชานาสิ ฯเปฯ นิพฺเพเฐหิ วา สเจ ปโหสี”ติ วิคฺคยฺห วิวาทยนฺติ.}

\prose{\textbf{พาโล ปโร อกฺกุสโลติ จาหู}ติ ปโร พาโล หีโน นิหีโน โอมโก ลามโก ฉตุกฺโก ปริตฺโต อกุสโล อวิทฺวา อวิชฺชาคโต อญฺญานี อวิภาวี ทุปฺปญฺโญติ เอวมาหํสุ เอวํ กเถนฺติ เอวํ ภณนฺติ เอวํ ทีปยนฺติ เอวํ โวหรนฺตีติ พาโล ปโร อกฺกุสโลติ จาหุ.}

\prose{\textbf{สจฺโจ นุ วาโท กตโม อิเมสนฺ}ติ อิเมสํ สมณพฺราหฺมณานํ วาโท กตโม สจฺโจ ตจฺโฉ ตโถ ภูโต ยาถาโว อวิปรีโตติ สจฺโจ นุ วาโท กตโม อิเมสํ.}

\prose{\textbf{สพฺเพว หีเม กุสลาวทานา}ติ สพฺเพวิเม สมณพฺราหฺมณา กุสลวาทา ปณฺฑิตวาทา ถิรวาทา ญายวาทา เหตุวาทา ลกฺขณวาทา การณวาทา ฐานวาทา สกาย ลทฺธิยาติ สพฺเพว หีเม กุสลาวทานา.\csromanspacer{} เตนาห โส นิมฺมิโต\csromandash{}}

\csromangathameasure{เอวมฺปิ วิคฺคยฺห วิวาทยนฺติ, \\ พาโล ปโร อกฺกุสโลติ จาหุ. \\ สจฺโจ นุ วาโท กตโม อิเมสํ, \\ \textbf{สพฺเพว หีเม กุสลาวทานา}ติ.}
\csromangathagroup{\csromangathastanza{เอวมฺปิ วิคฺคยฺห วิวาทยนฺติ, \\ พาโล ปโร อกฺกุสโลติ จาหุ. \\ สจฺโจ นุ วาโท กตโม อิเมสํ, \\ \textbf{สพฺเพว หีเม กุสลาวทานา}ติ.}}

\proseitem{115}{\csromansymbolfootnote{*}{ขุ 1. 416 ปิฏฺเฐ.}ปรสฺส เจ ธมฺมมนานุชานํ,}

\prose{\textbf{พาโล’มโก โหติ นิหีนปญฺโญ.}}

\csromangathameasure{\textbf{สพฺเพว พาลาสุ นิหีนปญฺญา,} \\ \textbf{สพฺเพวิเม ทิฏฺฐิปริพฺพสนา.}}
\csromangathagroup{\csromangathastanza{\textbf{สพฺเพว พาลาสุ นิหีนปญฺญา,} \\ \textbf{สพฺเพวิเม ทิฏฺฐิปริพฺพสนา.}}}

\prose{\textbf{ปรสฺส เจ ธมฺมมนานุชานนฺ}ติ ปรสฺส ธมฺมํ ทิฏฺฐิํ ปฏิปทํ มคฺคํ อนานุชานนฺโต อนานุปสฺสนฺโต อนานุมนนฺโต อนานุมญฺญนฺโต อนานุโมทนฺโตติ ปรสฺส เจ ธมฺมมนานุชานํ.}

\csromanheader{2}{12. จูฬพฺยูหสุตฺตนิทฺเทส}
\prose{\csromanfolio{223}\textbf{พาโล’มโก โหติ นิหีนปญฺโญ}ติ ปโร พาโล โหติ หีโน นิหีโน โอมโก ลามโก ฉตุกฺโก ปริตฺโต หีนปญฺโญ นิหีนปญฺโญ โอมกปญฺโญ ลามกปญฺโญ ฉตุกฺกปญฺโญ ปริตฺตปญฺโญติ พาโล’มโก โหติ นิหีนปญฺโญ.}

\prose{\textbf{สพฺเพว พาลา สุนิหีนปญฺญา}ติ สพฺเพวิเม สมณพฺราหฺมณา พาลา หีนา นิหีนา โอมกา ลามกา ฉตุกฺกา ปริตฺตา\textbf{,} สพฺเพว หีนปญฺญา นิหีนปญฺญา โอมกปญฺญา ลามกปญฺญา ฉตุกฺกปญฺญา ปริตฺตปญฺญาติ สพฺเพว พาลา สุนิหีนปญฺญา.}

\prose{\textbf{สพฺเพวิเม ทิฏฺฐิปริพฺพสานา}ติ สพฺเพวิเม สมณพฺราหฺมณา ทิฏฺฐิคติกา\textbf{,} เต ทฺวาสฏฺฐิยา ทิฏฺฐิคตานํ อญฺญตรญฺญตรํ ทิฏฺฐิคตํ คเหตฺวา อุคฺคเหตฺวา คณฺหิตฺวา ปรามสิตฺวา อภินิวิสิตฺวา สกาย สกาย ทิฏฺฐิยา วสนฺติ สํวสนฺติ อาวสนฺติ ปริวสนฺติ.\csromanspacer{} ยถา อคาริกา วา ฆเรสุ วสนฺติ\textbf{,} สาปตฺติกา วา อาปตฺตีสุ วสนฺติ\textbf{,} สกิเลสา วา กิเลเสสุ วสนฺติ\textbf{,} เอวเมวํ สพฺเพวิเม สมณพฺราหฺมณา ทิฏฺฐิคติกา ฯเปฯ ปริวสนฺตีติ สพฺเพวิเม ทิฏฺฐิปริพฺพสานา.\csromanspacer{} เตนาห ภควา\csromandash{}}

\csromangathameasure{ปรสฺส เจ ธมฺมมนานุชานํ\textbf{,} \\ \textbf{พาโล’มโก โหติ นิหีนปญฺโญ}. \\ \textbf{สพฺเพว พาลา สุนิหีนปญฺญา,}}
\csromangathagroup{\csromangathastanza{ปรสฺส เจ ธมฺมมนานุชานํ\textbf{,} \\ \textbf{พาโล’มโก โหติ นิหีนปญฺโญ}. \\ \textbf{สพฺเพว พาลา สุนิหีนปญฺญา,}}}

\prose{\textbf{สพฺเพวิเม ทิฏฺฐิปริพฺพสานา}ติ.}

\csromanhangingbegin
\hangingheaditem{116}{\textbf{สนฺทิฏฺฐิยา เจว นวีวทาตา}\footnote{เจวนเววทาตา (สี), เจ ปน วิวทาตา (สฺยา) ขุ 1. 416 ปิฏฺเฐปิ.}\textbf{,}}
\hangingbody{\textbf{สํสุทฺธปญฺญา กุสลา มุตีมา.}}
\hangingbody{\textbf{น เตสํ โกจิ ปริหีนปญฺโญ,}}
\hangingbody{\textbf{ทิฏฺฐี หิ เตสมฺปิ ตถา สมตฺตา.}}
\csromanhangingend

\prose{\textbf{สนฺทิฏฺฐิยา เจว นวีวทาตา}ติ สกาย ทิฏฺฐิยา สกาย ขนฺติยา สกาย รุจิยา สกาย ลทฺธิยา อนวีวทาตา อโวทาตา อปริโยทาตา สํกิลิฏฺฐา สํกิเลสิกาติ สนฺทิฏฺฐิยา เจว นวีวทาตา.}

\prose{\textbf{สํสุทฺธปญฺญา กุสลา มุตีมา}ติ สุทฺธปญฺญา วิสุทฺธปญฺญา ปริสุทฺธปญฺญา โวทาตปญฺญา ปริโยทาตปญฺญา.\csromanspacer{} อถ วา สุทฺธทสฺสนา วิสุทฺธทสฺสนา \csromanfolio{224}ปริสุทฺธทสฺส\textbf{นา} โวทาตทสฺส\textbf{นา} ปริโยทาตทสฺส\textbf{นา}ติ สํสุทฺธปญฺญา.\csromanspacer{} \textbf{กุสลา}ติ กุสลา ปณฺฑิตา ปญฺญวนฺโต อิทฺธิมนฺโต ญาณิโน วิภาวิโน เมธาวิโนติ สํสุทฺธปญฺญา กุสลา.\csromanspacer{} \textbf{มุตีมา}ติ มุติมา ปณฺฑิตา ปญฺญวนฺโต อิทฺธิมนฺโต ญาณิโน วิภาวิโน เมธาวิโนติ สํสุทฺธปญฺญา กุสลา มุตีมา.}

\prose{\textbf{เตสํ น โกจิ ปริหีนปญฺโญ}ติ เตสํ สมณพฺราหฺมณานํ น โกจิ หีนปญฺโญ นิหีนปญฺโญ โอมกปญฺโญ ลามกปญฺโญ ฉตุกฺกปญฺโญ ปริตฺตปญฺโญ อตฺถิ.\csromanspacer{} สพฺเพว เสฏฺฐปญฺญา วิสิฏฺฐปญฺญา ปาโมกฺขปญฺญา อุตฺตมปญฺญา ปวรปญฺญาติ เตสํ น โกจิ ปริหีนปญฺโญ.}

\prose{\textbf{ทิฏฺฐี หิ เตสมฺปิตถา สมตฺตา}ติ เตสํ สมณพฺราหฺมณานํ ทิฏฺฐิ ตถา สมตฺตา สมาทินฺ\textbf{นา} คหิตา ปรามฏฺฐา อภินิวิฏฺฐา อชฺโฌสิตา อธิมุตฺตาติ ทิฏฺฐี หิ เตสมฺปิ ตถา สมตฺตา.\csromanspacer{} เต\textbf{นา}ห ภควา\csromandash{}}

\csromangathameasure{สนฺทิฏฺฐิยา เจว นวีวทาตา, \\ สํสุทฺธปญฺญา กุสลา มุตีมา. \\ \textbf{เตสํ น โกจิ ปริหีนปญฺโญ}, \\ \textbf{ทิฏฺฐี หิ เตสมฺปิ ตถา สมตฺตา}ติ.}
\csromangathagroup{\csromangathastanza{สนฺทิฏฺฐิยา เจว นวีวทาตา, \\ สํสุทฺธปญฺญา กุสลา มุตีมา. \\ \textbf{เตสํ น โกจิ ปริหีนปญฺโญ}, \\ \textbf{ทิฏฺฐี หิ เตสมฺปิ ตถา สมตฺตา}ติ.}}

\csromanhangingbegin
\hangingheaditem{117}{น วาหเมตํ ตถิยนฺติ\footnote{ตถิวนฺติ (สฺยา) ขุ 1. 416 ปิฏฺเฐปิ.} \textbf{พฺรูมิ,}}
\hangingbody{\textbf{ยมาหุ พาลา มิถุ อญฺญมญฺญํ.}}
\hangingbody{\textbf{สกํ สกํ ทิฏฺฐิมกํสุ สจฺจํ,}}
\hangingbody{\textbf{ตสฺมา หิ พาโลติ ปรํ ทหนฺติ.}}
\csromanhangingend

\prose{\textbf{น วาหเมตํ ตถิยนฺติ พฺรูมี}ติ \textbf{นา}ติ ปฏิกฺเขโป.\csromanspacer{} \textbf{เอตนฺ}ติ ทฺวาสฏฺฐิทิฏฺฐิคตานิ \textbf{นา}หํ เอตํ ตจฺฉํ ตถํ ภูตํ ยาถาวํ อวิปรีตนฺติ พฺรูมิ อาจิกฺขามิ เทเสมิ ปญฺญเปมิ ปฏฺฐเปมิ วิวรามิ วิภชามิ อุตฺตานีกโรมิ ปกาเสมีติ น วาหเมตํ ตถิยนฺติ พฺรูมิ.}

\prose{\textbf{ยมาหุ พาลา มิถุ อญฺญมญฺญนฺ}ติ \textbf{มิถู}ติ เทฺว ช\textbf{นา} เทฺว กลหการกา เทฺว ภณฺฑนการกา เทฺว ภสฺสการกา เทฺว วิวาทการกา เทฺว อธิกรณการกา เทฺว วาทิโน เทฺว สลฺลปกา, เต อญฺญมญฺญํ “พาโล\footnote{พาลโต (สี, ก) เอวมญฺเญสุ ฉปฺปเทสุปิ โตปจฺจยนฺตวเสน.} หีโน นิหีโน โอมโก ลามโก ฉตุกฺโก ปริตฺโต”ติ}

\vspace{\tipitakaparskip}
\csromancenter{จูฬพฺยูหสุตฺตนิทฺเทส เอวมาหํสุ เอวํ กเถนฺติ เอวํ ภณนฺติ เอวํ ทีปยนฺติ เอวํ โวหรนฺตีติ ยมาตุ พาลา มิถุ อญฺญมญฺญํ.}
\prose{\csromanfolio{225}\textbf{สกํ สกํ ทิฏฺฐิมกํสุ สจฺจนฺ}ติ “สสฺสโต โลโก, อิทเมว สจฺจํ โมฆมญฺญนฺ”ติ, สกํ สกํ ทิฏฺฐิมกํสุ สจฺจํ, “อสสฺสโต โลโก, อิทเมวสจฺจํ โมฆมญฺญนฺ”ติ ฯเปฯ “เนว โหติ น น โหติ ตถาคโต ปรํ มรณา, อิทเมว สจฺจํ โมฆมญฺญนฺ”ติ สกํ สกํ ทิฏฺฐิมกํสุ สจฺจํ.}

\prose{\textbf{ตสฺมา หิ พาโลติ ปรํ ทหนฺตี}ติ \textbf{ตสฺมา}ติ \textbf{ตสฺมา} ตํการณา ตํเหตุ ตปฺปจฺจยา ตํนิทานา ปรํ “พาโล หีโน นิหีโน โอมโก ลามโก ฉตุกฺโก ปริตฺโต”ติ ทหนฺติ ปสฺสนฺติ ทกฺขนฺติ โอโลเกนฺติ นิชฺฌายนฺติ อุปปริกฺขนฺตีติ \textbf{ตสฺมา} หิ พาโลติ ปรํ ทหนฺติ.\csromanspacer{} เตนาห ภควา\csromandash{}}

\csromangathameasure{น วาหเมตํ ตถิยนฺติ พฺรูมิ, \\ ยมาหุ พาลา มิถุ อญฺญมญฺญํ.}
\csromangathagroup{\csromangathastanza{น วาหเมตํ ตถิยนฺติ พฺรูมิ, \\ ยมาหุ พาลา มิถุ อญฺญมญฺญํ.}}

\prose{สกํ สกํ ทิฏฺฐิมกํสุ สจฺจํ,}

\prose{\textbf{ตสฺมา หิ พาโลติ ปรํ ทหนฺตี}ติ.}

\csromanhangingbegin
\hangingheaditem{118}{\textbf{ยมาหุ สจฺจํ ตถิยนฺติ เอเก},}
\hangingbody{\textbf{ตมาหุ อญฺเญปิ}1 \textbf{ตุจฺฉํ มุสาติ.}}
\hangingbody{\textbf{เอวมฺปิ วิคฺคยฺห วิวาทยนฺติ,}}
\hangingbody{\textbf{กสฺมา น เอกํ สมณา วทนฺติ.}}
\csromanhangingend

\prose{\textbf{ยมาหุ สจฺจํ ตถิยนฺติ เอเก}ติ ยํ ธมฺมํ ทิฏฺฐิํ ปฏิปทํ มคฺคํ เอเก สมณพฺราหฺมณา “อิทํ สจฺจํ ตจฺฉํ ตถํ ภูตํ ยาถาวํ อวิปรีตนฺ”ติ เอวมาหํสุ เอวํ กเถนฺติ เอวํ ภณนฺติ เอวํ ทีปยนฺติ เอวํ โวหรนฺตีติ ยมาหุ สจฺจํ ตถิยนฺติ เอเก.}

\prose{\textbf{ตมาหุ อญฺเญปิ ตุจฺฉํ มุสาตี}ติ ตเมว ธมฺมํ ทิฏฺฐิํ ปฏิปทํ มคฺคํ เอเก สมณพฺราหฺมณา “ตุจฺฉํ เอตํ, มุสา เอตํ, อภูตํ เอตํ, อลิกํ เอตํ, อยาถาวํ เอตนฺ”ติ เอวมาหํสุ เอวํ กเถนฺติ เอวํ ภณนฺติ เอวํ ทีปยนฺติ เอวํ โวหรนฺตีติ ตมาหุ อญฺเญปิ \textbf{ตุจฺฉํ มุสาติ.}}

\prose{\csromanfolio{226}\textbf{เอวมฺปิ วิคฺคยฺห วิวาทยนฺตี}ติ เอวํ คเหตฺวา อุคฺคเหตฺวา คณฺหิตฺวา ปรามสิตฺวา อภินิวิสิตฺวา วิวาทยนฺติ, กลหํ กโรนฺติ ภณฺฑนํ กโรนฺติ วิคฺคหํ กโรนฺติ, วิวาทํ กโรนฺติ, เมธคํ กโรนฺติ “น ตฺวํ อิมํ ธมฺมวินยํ อาชานาสิ ฯเปฯ นิพฺเพเฐหิ วา สเจ ปโหสี”ติ เอวมฺปิ วิคฺคยฺห วิวาทยนฺติ.}

\prose{\textbf{กสฺมา น เอกํ สมณา วทนฺตี}ติ \textbf{กสฺมา}ติ \textbf{กสฺมา} กิํการณา กิํเหตุ กิํปจฺจยา กิํนิทานา กิํสมุทยา กิํชาติกา กิํปภวา น เอกํ วทนฺติ นานา วทนฺติ วิวิธํ วทนฺติ อญฺโญญฺญํ\footnote{อญฺโญญฺเญ (ก)} วทนฺติ ปุถุ วทนฺติ กเถนฺติ ภณนฺติ ทีปยนฺติ โวหรนฺตีติ \textbf{กสฺมา} น เอกํ สมณา วทนฺติ.\csromanspacer{} เตนาห โส นิมฺมิโต\csromandash{}}

\csromangathameasure{ยมาหุ สจฺจํ ตถิยนฺติ เอเก, \\ ตมาหุ อญฺเญปิ ตุจฺฉํ มุสาติ. \\ เอวมฺปิ วิคฺคยฺห วิวาทยนฺติ,}
\csromangathagroup{\csromangathastanza{ยมาหุ สจฺจํ ตถิยนฺติ เอเก, \\ ตมาหุ อญฺเญปิ ตุจฺฉํ มุสาติ. \\ เอวมฺปิ วิคฺคยฺห วิวาทยนฺติ,}}

\prose{\textbf{กสฺมา น เอกํ สมณา วทนฺตี}ติ.}

\proseitem{119}{\csromansymbolfootnote{*}{ขุ 1. 416 ปิฏฺเฐ.}เอกํ หิ สจฺจํ น ทุตียมตฺถิ,}

\prose{\textbf{ยสฺมิํ ปชา โน วิวเท ปชานํ.}}

\csromangathameasure{\textbf{นานา เต สจฺจานิ สยํ ถุนนฺติ,} \\ \textbf{ตสฺมา น เอกํ สมณา วทนฺติ.}}
\csromangathagroup{\csromangathastanza{\textbf{นานา เต สจฺจานิ สยํ ถุนนฺติ,} \\ \textbf{ตสฺมา น เอกํ สมณา วทนฺติ.}}}

\prose{\textbf{เอกํ หิ สจฺจํ น ทุตียมตฺถี}ติ เอกํ สจฺจํ วุจฺจติ ทุกฺขนิโรโธ นิพฺพานํ.\csromanspacer{} โย โส สพฺพสงฺขารสมโถ สพฺพูปธิปฏินิสฺสคฺโค ตณฺหกฺขโย วิราโค นิโรโธ นิพฺพานํ.\csromanspacer{} อถ วา เอกํ สจฺจํ วุจฺจติ มคฺคสจฺจํ นิยฺยานสจฺจํ ทุกฺขนิโรธคามินี ปฏิปทา อริโย อฏฺฐงฺคิโก มคฺโค.\csromanspacer{} เสยฺยถิทํ.\csromanspacer{} สมฺมาทิฏฺฐิ สมฺมาสงฺกปฺโป สมฺมาวาจา สมฺมากมฺมนฺโต สมฺมา-อาชีโว สมฺมาวายาโม สมฺมาสติ สมฺมาสมาธีติ เอกํ หิ สจฺจํ น ทุตียมตฺถิ.}

\prose{\textbf{ยสฺมิํ ปชา โน วิวเท ปชานนฺ}ติ \textbf{ยสฺมินฺ}ติ ยสฺมิํ สจฺเจ.\csromanspacer{} \textbf{ปชา}ติ สตฺตาธิวจนํ.\csromanspacer{} \textbf{ปชานนฺ}ติ\footnote{ปชานํ (สี, ก), ปชา (สฺยา)} ยํ สจฺจํ ปชานนฺตา อาชานนฺตา วิชานนฺตา ปฏิวิชานนฺตา ปฏิวิชฺฌนฺตา น กลหํ กเรยฺยุํ, น ภณฺฑนํ กเรยฺยุํ, น วิคฺคหํ กเรยฺยุํ, น วิวาทํ กเรยฺยุํ, น เมธคํ กเรยฺยุํ, กลหํ ภณฺฑนํ วิคฺคหํ วิวาทํ เมธคํ ปชเหยฺยุํ วิโนเทยฺยุํ พฺยนฺติํ กเรยฺยุํ\footnote{พฺยนฺตีกเรยฺยุํ (สี, สฺยา)} อนภาวํ คเมยฺยุนฺติ ยสฺมิํ ปชา โน วิวเท ปชานํ.}

\csromanheader{2}{12. จูฬพฺยูหสุตฺตนิทฺเทส}
\prose{\csromanfolio{227}\textbf{นานา เต สจฺจานิ สยํ ถุนนฺตี}ติ นานา เต สจฺจานิ สยํ ถุนนฺติ วทนฺติ กเถนฺติ ภณนฺติ ทีปยนฺติ โวหรนฺติ “สสฺสโต โลโก, อิทเมว สจฺจํ โมฆมญฺญนฺ”ติ สยํ ถุนนฺติ วทนฺติ กเถนฺติ ภณนฺติ ทีปยนฺติ โวหรนฺติ, “อสสฺสโต โลโก ฯเปฯ เนว โหติ น น โหติ ตถาคโต ปรํ มรณา, อิทเมว สจฺจํ โมฆมญฺญนฺ”ติ สยํ ถุนนฺติ วทนฺติ กเถนฺติ ภณนฺติ ทีปยนฺติ โวหรนฺตีติ นานา เต สจฺจานิ สยํ ถุนนฺติ.}

\prose{\textbf{ตสฺมา น เอกํ สมณา วทนฺตี}ติ \textbf{ตสฺมา}ติ \textbf{ตสฺมา} ตํการณา ตํเหภุ ตปฺปจฺจยา ตํนิทานา น เอกํ วทนฺติ นานา วทนฺติ วิวิธํ วทนฺติ อญฺโญญฺญํ วทนฺติ ปุถุ วทนฺติ กเถนฺติ ภณนฺติ ทีปยนฺติ โวหรนฺตีติ \textbf{ตสฺมา} น เอกํ สมณา วทนฺติ.\csromanspacer{} เตนาห ภควา\csromandash{}}

\csromangathameasure{เอกํ หิ สจฺจํ น ทุตียมตฺถิ, \\ ยสฺมิํ ปชา โน วิวเท ปชานํ.}
\csromangathagroup{\csromangathastanza{เอกํ หิ สจฺจํ น ทุตียมตฺถิ, \\ ยสฺมิํ ปชา โน วิวเท ปชานํ.}}

\prose{นานา เต สจฺจานิ สยํ ถุนนฺติ,}

\prose{\textbf{ตสฺมา น เอกํ สมณา วทนฺตี}ติ.}

\proseitem{120}{\csromansymbolfootnote{*}{ขุ 1. 416 ปิฏฺเฐ.}\textbf{กสฺมา นุ สจฺจานิ วทนฺติ นานา},}

\prose{\textbf{ปวาทิยาเส กุสลาวทานา.}}

\csromangathameasure{\textbf{สจฺจานิ สุตานิ พหูนิ นานา,} \\ \textbf{อุทาหุ เต ตกฺกมนุสฺสรนฺติ.}}
\csromangathagroup{\csromangathastanza{\textbf{สจฺจานิ สุตานิ พหูนิ นานา,} \\ \textbf{อุทาหุ เต ตกฺกมนุสฺสรนฺติ.}}}

\prose{\textbf{กสฺมา นุ สจฺจานิ วทนฺติ นานา}ติ \textbf{กสฺมา}ติ \textbf{กสฺมา} กิํการณา กิํเหตุ กิํปจฺจยา กิํนิทานา สจฺจานิ นานา\footnote{นานานิ (ก)} วทนฺติ, วิวิธานิ วทนฺติ, อญฺโญญฺญานิ วทนฺติ, ปุถูนิ วทนฺติ กเถนฺติ ภณนฺติ ทีปยนฺติ โวหรนฺตีติ \textbf{กสฺมา} นุ สจฺจานิ วทนฺติ นานา.}

\prose{\textbf{ปวาทิยาเส กุสลาวทานา}ติ \textbf{ปวาทิยาเส}ติ วิปฺปวทนฺตีติปิ \textbf{ปวาทิยาเส}.\csromanspacer{} อถ วา สกํ สกํ ทิฏฺฐิคตํ ปวทนฺติ กเถนฺติ ภณนฺติ ทีปยนฺติ โวหรนฺติ “สสฺสโต โลโก, อิทเมว สจฺจํ โมฆมญฺญนฺ”ติ ปวทนฺติ กเถนฺติ ภณนฺติ ทีปยนฺติ โวหรนฺติ, “อสสฺสโต โลโก ฯเปฯ เนว โหติ น น โหติ ตถาคโต ปรํ มรณา, อิทเมว สจฺจํ โมฆมญฺญนฺ”ติ ปวทนฺติ กเถนฺติ ภณนฺติ ทีปยนฺติ โวหรนฺติ. \csromanfolio{228}\textbf{กุสลาวทานา}ติ กุสลวาทา ปณฺฑิตวาทา ถิรวาทา ญายวาทา เหตุวาทา ลกฺขณวาทา การณวาทา ฐานวาทา สกาย ลทฺธิยาติ ปวาทิยาเส กุสลาวทานา.}

\prose{\textbf{สจฺจานิ สุตานิ พหูนิ นานา}ติ สจฺจานิ สุตานิ พหุกานิ นานานิ วิวิธานิ อญฺโญญฺญานิ ปุถูนีติ สจฺจานิ สุตานิ พหูนิ นานา.}

\prose{\textbf{อุทาหุ เต ตกฺกมนุสฺสรนฺตี}ติ อุทาหุ ตกฺเกน สงฺกปฺเปน ยายนฺติ นียนฺติ วุยฺหนฺติ สํหรียนฺตีติ เอวมฺปิ อุทาหุ เต ตกฺกมนุสฺสรนฺติ.\csromanspacer{} อถ วา ตกฺกปริยาหตํ วีมํสานุจริตํ สยํ ปฏิภานํ วทนฺติ กเถนฺติ ภณนฺติ ทีปยนฺติ โวหรนฺตีติ เอวมฺปิ อุทาหุ เต ตกฺกมนุสฺสรนฺติ.\csromanspacer{} เตนาห โส นิมฺมิโต\csromandash{}}

\csromangathameasure{กสฺมา นุ สจฺจานิ วทนฺติ นานา, \\ ปวาทิยาเส กุสลาวทานา. \\ \textbf{สจฺจานิ สุตานิ พหูนิ นานา},}
\csromangathagroup{\csromangathastanza{กสฺมา นุ สจฺจานิ วทนฺติ นานา, \\ ปวาทิยาเส กุสลาวทานา. \\ \textbf{สจฺจานิ สุตานิ พหูนิ นานา},}}

\prose{\textbf{อุทาหุ เต ตกฺกมนุสฺสรนฺตี}ติ.}

\proseitem{121}{\csromansymbolfootnote{*}{ขุ 1. 416 ปิฏฺเฐ.}\textbf{น เหว สจฺจานิ พหูนิ นานา},}

\prose{\textbf{อญฺญตฺร สญฺญาย นิจฺจานิ โลเก.}}

\csromangathameasure{\textbf{ตกฺกญฺจ ทิฏฺฐีสุ ปกปฺปยิตฺวา,} \\ \textbf{สจฺจํ มุสาติ ทฺวยธมฺมมาหุ.}}
\csromangathagroup{\csromangathastanza{\textbf{ตกฺกญฺจ ทิฏฺฐีสุ ปกปฺปยิตฺวา,} \\ \textbf{สจฺจํ มุสาติ ทฺวยธมฺมมาหุ.}}}

\prose{\textbf{น เหว สจฺจานิ พหูนิ นานา}ติ น เหว สจฺจานิ พหุกานิ นานานิ วิวิธานิ อญฺโญญฺญานิ ปุถูนีติ น เหว สจฺจานิ พหูนิ นานา.}

\prose{\textbf{อญฺญตฺร สญฺญาย นิจฺจานิ โลเก}ติ อญฺญตฺร สญฺญาย นิจฺจคฺคาหา เอกญฺเญว สจฺจํ โลเก กถียติ ภณียติ ทีปียติ โวหรียติ.\csromanspacer{} ทุกฺขนิโรโธ นิพฺพานํ, โย โส สพฺพสงฺขารสมโถ สพฺพูปธิปฏินิสฺสคฺโค ตณฺหกฺขโย วิราโค นิโรโธ นิพฺพานํ.\csromanspacer{} อถ วา เอกํ สจฺจํ วุจฺจติ มคฺคสจฺจํ นิยฺยานสจฺจํ ทุกฺขนิโรธคามินี ปฏิปทา อริโย อฏฺฐงฺคิโก มคฺโค.\csromanspacer{} เสยฺยถิทํ, สมฺมาทิฏฺฐิ ฯเปฯ สมฺมาสมาธีติ อญฺญตฺร สญฺญาย นิจฺจานิ โลเก.}

\prose{\textbf{ตกฺกญฺจ ทิฏฺฐีสุ ปกปฺปยิตฺวา, สจฺจํ มุสาติ ทฺวยธมฺมมาหู}ติ ตกฺกํ วิตกฺกํ สงฺกปฺปํ ตกฺกยิตฺวา วิตกฺกยิตฺวา สงฺกปฺปยิตฺวา ทิฏฺฐิคตานิ ชเนนฺติ}

\vspace{\tipitakaparskip}
\csromancenter{จูฬพฺยูหสุตฺตนิทฺเทส สญฺชเนนฺติ นิพฺพตฺเตนฺติ อภินิพฺพตฺเตนฺติ ทิฏฺฐิคตานิ ชเนตฺวา สญฺชเนตฺวา นิพฺพตฺเตตฺวา อภินิพฺพตฺเตตฺวา “มยฺหํ สจฺจํ, ตุยฺหํ มุสา”ติ เอวมาหํสุ เอวํ กเถนฺติ เอวํ ภณนฺติ เอวํ ทีปยนฺติ เอวํ โวหรนฺตีติ ตกฺกญฺจ ทิฏฺฐีสุ ปกปฺปยิตฺวา, สจฺจํ มุสาติ ทฺวยธมฺมมาหุ.\csromanspacer{} เตนาห ภควา\csromandash{}}
\vspace{0.5\baselineskip}
\csromangathameasure{น เหว สจฺจานิ พหูนิ นานา, \\ อญฺญตฺร สญฺญาย นิจฺจานิ โลเก. \\ ตกฺกญฺจ ทิฏฺฐีสุ ปกปฺปยิตฺวา, \\ สจฺจํ มุสาติ ทฺวยธมฺมมาหูติ.}
\csromangathagroup{\csromangathastanza{\csromanfolio{229}น เหว สจฺจานิ พหูนิ นานา, \\ อญฺญตฺร สญฺญาย นิจฺจานิ โลเก. \\ ตกฺกญฺจ ทิฏฺฐีสุ ปกปฺปยิตฺวา, \\ สจฺจํ มุสาติ ทฺวยธมฺมมาหูติ.}}

\proseitem{122}{\csromansymbolfootnote{*}{ขุ 1. 417 ปิฏฺเฐ.}ทิฏฺเฐ สุเต สีลวเต มุเต วา,}

\prose{\textbf{เอเต จ}\footnote{เอเตสุ (สี)} \textbf{นิสฺสาย วิมานทสฺสี.}}

\csromangathameasure{\textbf{วินิจฺฉเย ฐตฺวา ปหสฺสมาโน,} \\ \textbf{พาโล ปโร อกฺกุสโลติ จาห.}}
\csromangathagroup{\csromangathastanza{\textbf{วินิจฺฉเย ฐตฺวา ปหสฺสมาโน,} \\ \textbf{พาโล ปโร อกฺกุสโลติ จาห.}}}

\prose{\textbf{ทิฏฺเฐ สุเต สีลวเต มุเต วา, เอเต จ นิสฺสาย วิมานทสฺสี}ติ ทิฏฺฐํ วา ทิฏฺฐสุทฺธิํ วา สุตํ วา สุตสุทฺธิํ วา สีลํ วา สีลสุทฺธิํ วา วตํ วา วตสุทฺธิํ วา มุตํ วา มุตสุทฺธิํ วา นิสฺสาย อุปนิสฺสาย คณฺหิตฺวา ปรามสิตฺวา อภินิวิสิตฺวาติ ทิฏฺเฐ สุเต สีลวเต มุเต วา, เอเต จ นิสฺสาย.\csromanspacer{} \textbf{วิมานทสฺสี}ติ น สมฺมาเนตีติปิ วิมานทสฺสี.\csromanspacer{} อถ วา โทมนสฺสํ ชเนตีติปิ วิมานทสฺสีติ ทิฏฺเฐ สุเต สีลวเต มุเต วา, เอเต จ \textbf{นิสฺสาย วิมานทสฺสี.}}

\prose{\textbf{วินิจฺฉเย ฐตฺวา ปหสฺสมาโน}ติ วินิจฺฉยา วุจฺจนฺติ ทฺวาสฏฺฐิ ทิฏฺฐิคตานิ.\csromanspacer{} ทิฏฺฐิวินิจฺฉเย วินิจฺฉยทิฏฺฐิยา ฐตฺวา ปติฏฺฐหิตฺวา คณฺหิตฺวา ปรามสิตฺวา อภินิวิสิตฺวาติ วินิจฺฉเย ฐตฺวา.\csromanspacer{} \textbf{ปหสฺสมาโน}ติ ตุฏฺโฐ โหติ หฏฺโฐ ปหฏฺโฐ อตฺตมโน ปริปุณฺณสงฺกปฺโป.\csromanspacer{} อถ วา ทนฺตวิทํสกํ ปหสฺสมาโนติ วินิจฺฉเย ฐตฺวา ปหสฺสมาโน.}

\prose{\textbf{พาโล ปโร อกฺกุสโลติ จาหา}ติ ปโร “พาโล หีโน นิหีโน โอมโก ลามโก ฉตุกฺโก ปริตฺโต อกุสโล อวิทฺวา อวิชฺชาคโต อญฺญาณี อวิภาวี อเมธาวี ทุปฺปญฺโญ”ติ เอวมาห เอวํ กเถติ เอวํ ภณติ เอวํ ทีปยติ เอวํ โวหรตีติ พาโล ปโร อกฺกุสโลติ จาห.\csromanspacer{} เตนาห ภควา\csromandash{}}

\csromangathameasure{ทิฏฺเฐ สุเต สีลวเต มุเต วา, \\ เอเต จ นิสฺสาย วิมานทสฺสี. \\ วินิจฺฉเย ฐตฺวา ปหสฺสมาโน, \\ พาโล ปโร อกฺกุสโลติ จาหาติ.}
\csromangathagroup{\csromangathastanza{\csromanfolio{230}ทิฏฺเฐ สุเต สีลวเต มุเต วา, \\ เอเต จ นิสฺสาย วิมานทสฺสี. \\ วินิจฺฉเย ฐตฺวา ปหสฺสมาโน, \\ พาโล ปโร อกฺกุสโลติ จาหาติ.}}

\proseitem{123}{\csromansymbolfootnote{*}{ขุ 1. 417 ปิฏฺเฐ.}เยเนว พาโลติ ปรํ ทหาติ,}

\prose{\textbf{เตนา’ตุมานํ กุสโลติ จาห.}}

\csromangathameasure{\textbf{สย’มตฺตนา โส กุสลาวทาโน,} \\ \textbf{อญฺญํ วิมาเนติ ตเทว ปาว.}}
\csromangathagroup{\csromangathastanza{\textbf{สย’มตฺตนา โส กุสลาวทาโน,} \\ \textbf{อญฺญํ วิมาเนติ ตเทว ปาว.}}}

\prose{\textbf{เยเนว พาโลติ ปรํ ทหาตี}ติ เยเนว เหตุนา เยน ปจฺจเยน เยน การเณน เยน ปภเวน ปรํ พาลโต หีนโต นิหีนโต โอมกโต ลามกโต ฉตุกฺกโต ปริตฺตโต ทหติ ปสฺสติ ทกฺขติ โอโลเกติ นิชฺฌายติ อุปปริกฺขตีติ เยเนว พาโลติ ปรํ ทหาติ.}

\prose{\textbf{เตนา’ตุมานํ กุสโลติ จาหา}ติ \textbf{อาตุมาโน} วุจฺจติ อตฺตา.\csromanspacer{} โสปิ เตเนว เหตุนา เตน ปจฺจเยน เตน การเณน เตน ปภเวน อตฺตานํ อหมสฺมิ กุสโล ปณฺฑิโต ปญฺญวา พุทฺธิมา ญาณี วิภาวี เมธาวีติ เตนา’ตุมานํ กุสโลติ จาห.}

\prose{\textbf{สย’มตฺตนา โส กุสลาวทาโน}ติ สยเมว อตฺตานํ กุสลวาโท ปณฺฑิตวาโท ถิรวาโท ญายวาโท เหตุวาโท ลกฺขณวาโท การณวาโท ฐานวาโท สกาย ลทฺธิยาติ สย’มตฺตนา โส กุสลาวทาโน.}

\prose{\textbf{อญฺญํ วิมาเนติ ตเทว ปาวา}ติ น สมฺมาเนตีติปิ อญฺญํ วิมาเนติ.\csromanspacer{} อถ วา โทมนสฺสํ ชเนตีติปิ อญฺญํ วิมาเนติ.\csromanspacer{} \textbf{ตเทว ปาวา}ติ ตเทว ตํ ทิฏฺฐิคตํ ปาวทติ, “อิติปายํ ปุคฺคโล มิจฺฉาทิฏฺฐิโก วิปรีตทสฺสโน”ติ อญฺญํ วิมาเนติ ตเทว ปาวท.\csromanspacer{} เตนาห ภควา\csromandash{}}

\prose{เยเนว พาโลติ ปรํ ทหาติ,}

\prose{\textbf{เตนา’ตุมานํ กุสโลติ จาห.}}

\csromangathameasure{\textbf{สย’มตฺตนา โส กุสลาวทาโน,}}
\csromangathagroup{\csromangathastanza{\textbf{สย’มตฺตนา โส กุสลาวทาโน,}}}

\prose{\textbf{อญฺญํ วิมาเนติ ตเทว ปาวา}ติ.}

\csromanheader{2}{12. จูฬพฺยูหสุตฺตนิทฺเทส}
\proseitem{124}{\csromanfolio{231}\csromansymbolfootnote{*}{ขุ 1. 417 ปิฏฺเฐ.}\textbf{อติสารทิฏฺฐิยา โส สมตฺโต},}

\prose{\textbf{มาเนน มตฺโต ปริปุณฺณมานี.}}

\csromangathameasure{\textbf{สยเมว สามํ มนสาภิสิตฺโต,} \\ \textbf{ทิฏฺฐี หิ สา ตสฺส ตถา สมตฺตา.}}
\csromangathagroup{\csromangathastanza{\textbf{สยเมว สามํ มนสาภิสิตฺโต,} \\ \textbf{ทิฏฺฐี หิ สา ตสฺส ตถา สมตฺตา.}}}

\prose{\textbf{อติสารทิฏฺฐิยา โส สมตฺโต}ติ อติสารทิฏฺฐิโย วุจฺจนฺติ ทฺวาสฏฺฐิ ทิฏฺฐิคตานิ.\csromanspacer{} กิํการณา อติสารทิฏฺฐิโย วุจฺจนฺติ ทฺวาสฏฺฐิ ทิฏฺฐิคตานิ, สพฺพา ตา ทิฏฺฐิโย การณาติกฺกนฺตา ลกฺขณาติกฺกนฺตา ฐานาติกฺกนฺตา, ตํการณา อติสารทิฏฺฐิโย วุจฺจนฺติ ทฺวาสฏฺฐิ ทิฏฺฐิคตานิ.\csromanspacer{} สพฺเพปิ ติตฺถิยา อติสารทิฏฺฐิยา\footnote{สพฺพาปิ ทิฏฺฐิโย อติสารทิฏฺฐิโย (สี, ก) 2. สพฺพาปิ ทิฏฺฐิโย วุจฺจนฺติ อติสารทิฏฺฐิโย (สี, ก)}.\csromanspacer{} กิํการณา สพฺเพปิ ติตฺถิยา อติสารทิฏฺฐิยา2, เต\footnote{ตา (ก)} อญฺญมญฺญํ อติกฺกมิตฺวา สมติกฺกมิตฺวา วีติวตฺติตฺวา ทิฏฺฐิคตานิ ชเนนฺติ สญฺชเนนฺติ นิพฺพตฺเตนฺติ อภินิพฺพตฺเตนฺติ, ตํการณา สพฺเพปิ ติตฺถิยา อติสารทิฏฺฐิยา2.\csromanspacer{} \textbf{อติสารทิฏฺฐิยา โส สมตฺโต}ติ อติสารทิฏฺฐิยา สมตฺโต ปริปุณฺโณ อโนโมติ อติสารทิฏฺฐิยา โส สมตฺโต.}

\prose{\textbf{มาเนน มตฺโต ปริปุณฺณมานี}ติ สกาย ทิฏฺฐิยา ทิฏฺฐิมาเนน มตฺโต ปมตฺโต อุมฺมตฺโต อติมตฺโตติ มาเนน มตฺโต.\csromanspacer{} \textbf{ปริปุณฺณมานี}ติ ปริปุณฺณมานี สมตฺตมานี อโนมมานีติ มาเนน มตฺโต ปริปุณฺณมานี.}

\prose{\textbf{สยเมว สามํ มนสาภิสิตฺโต}ติ สยเมว อตฺตานํ จิตฺเตน อภิสิญฺจติ “อหมสฺมิ กุสโล ปณฺฑิโต ปญฺญวา พุทฺธิมา ญาณี วิภาวี เมธาวี”ติ สยเมว สามํ มนสาภิสิตฺโต.}

\prose{\textbf{ทิฏฺฐี หิ สา ตสฺส ตถา สมตฺตา}ติ ตสฺส สา ทิฏฺฐิ ตถา สมตฺตา สมาทินฺนา คหิตา ปรามฏฺฐา อภินิวิฏฺฐา อชฺโฌสิตา อธิมุตฺตาติ ทิฏฺฐี หิ สา ตสฺส ตถา สมตฺตา.\csromanspacer{} เตนาห ภควา\csromandash{}}

\prose{\textbf{อติสารทิฏฺฐิยา โส สมตฺโต},}

\prose{มเนน มตฺโต ปริปุณฺณมานี.}

\csromangathameasure{\textbf{สยเมว สามํ มนสาภิสิตฺโต,} \\ \textbf{ทิฏฺฐี หิ สา ตสฺส ตถา สมตฺตา}ติ.}
\csromangathagroup{\csromangathastanza{\textbf{สยเมว สามํ มนสาภิสิตฺโต,} \\ \textbf{ทิฏฺฐี หิ สา ตสฺส ตถา สมตฺตา}ติ.}}

\proseitem{125}{\csromanfolio{232}\csromansymbolfootnote{*}{ขุ 1. 417 ปิฏฺเฐ.}\textbf{ปรสฺส เจ หิ วจสา นิหีโน},}

\prose{\textbf{ตุโม สหา โหติ นิหีนปญฺโญ.}}

\csromangathameasure{\textbf{อถ เจ สยํ เวทคู โหติ ธีโร,} \\ \textbf{น โกจิ พาโล สมเณสุ อตฺถิ.}}
\csromangathagroup{\csromangathastanza{\textbf{อถ เจ สยํ เวทคู โหติ ธีโร,} \\ \textbf{น โกจิ พาโล สมเณสุ อตฺถิ.}}}

\prose{\textbf{ปรสฺส เจ หิ วจสา นิหีโน}ติ ปรสฺส เจ วาจาย วจเนน นินฺทิตการณา ครหิตการณา อุปวทิตการณา ปโร พาโล โหติ หีโน นิหีโน โอมโก ลามโก ฉตุกฺโก ปริตฺโตติ ปรสฺส เจ หิ วจสา นิหีโน.\csromanspacer{} \textbf{ตุโม สหา โหติ นิหีนปญฺโญ}ติ โสปิ เตเนว สห โหติ หีนปญฺโญ นิหีนปญฺโญ โอมกปญฺโญ ลามกปญฺโญ ฉตุกฺกปญฺโญ ปริตฺตปญฺโญติ ตุโม สหา โหติ นิหีนปญฺโญ.}

\prose{\textbf{อถ เจ สยํ เวทคู โหติ ธีโร}ติ อถ เจ สยํ เวทคู โหติ ธีโร ปณฺฑิโต ปญฺญวา พุทฺธิมา ญาณี วิภาวี เมธาวีติ อถ เจ สยํ เวทคู โหติ ธีโร.}

\prose{\textbf{น โกจิ พาโล สมเณสุ อตฺถี}ติ สมเณสุ น โกจิ พาโล หีโน นิหีโน โอมโก ลามโก ฉตุกฺโก ปริตฺโต อตฺถิ, สพฺเพว เสฏฺฐปญฺญา\footnote{อคฺคปญฺญา เสฏฺฐปญฺญา (สฺยา)} วิสิฏฺฐปญฺญา ปาโมกฺขปญฺญา อุตฺตมปญฺญา ปวรปญฺญาติ น โกจิ พาโล สมเณสุ อตฺถิ.\csromanspacer{} เตนาห ภควา\csromandash{}}

\prose{\textbf{ปรสฺส เจ หิ วจสา นิหีโน},}

\prose{\textbf{ตุโม สหา โหติ นิหีนปญฺโญ.}}

\csromangathameasure{\textbf{อถ เจ สยํ เวทคู โหติ ธีโร,}}
\csromangathagroup{\csromangathastanza{\textbf{อถ เจ สยํ เวทคู โหติ ธีโร,}}}

\prose{\textbf{น โกจิ พาโล สมเณสุ อตฺถี}ติ.}

\proseitem{126}{\csromansymbolmark{*}อญฺญํ อิโต ยา’ภิวทนฺติ ธมฺมํ,}

\prose{\textbf{อปรทฺธา สุทฺธิ’มเกวลี เต.}}

\csromangathameasure{\textbf{เอวมฺปิ ติตฺถฺยา ปุถุโส วทนฺติ,} \\ \textbf{สนฺทิฏฺฐิราเคน หิ เต’ภิรตฺตา.}}
\csromangathagroup{\csromangathastanza{\textbf{เอวมฺปิ ติตฺถฺยา ปุถุโส วทนฺติ,} \\ \textbf{สนฺทิฏฺฐิราเคน หิ เต’ภิรตฺตา.}}}

\prose{\textbf{อญฺญํ อิโต ยา’ภิวทนฺติ ธมฺมํ, อปรทฺธา สุทฺธิ’มเกวลี เต}ติ อิโต อญฺญํ ธมฺมํ ทิฏฺฐิํ ปฏิปทํ มคฺคํ เย อภิวทนฺติ, เต สุทฺธิมคฺคํ วิสุทฺธิมคฺคํ ปริสุทฺธิมคฺคํ โวทาตมคฺคํ ปริโยทาตมคฺคํ วิรทฺธา อปรทฺธา ขลิตา คลิตา อญฺญาย}

\vspace{\tipitakaparskip}
\csromancenter{จูฬพฺยูหสุตฺตนิทฺเทส อปรทฺธา อเกวลี เต, อสมตฺตา เต, อปริปุณฺณา เต, หีนา นิหีนา โอมกา ลามกา ฉตุกฺกา ปริตฺตาติ อญฺญํ อิโต ยา’ภิวทนฺติ ธมฺมํ, อปรทฺทา สุทฺธิ’มเกวลี เต.}
\prose{\csromanfolio{233}\textbf{เอวมฺปิ ติตฺถฺยา ปุถุโส วทนฺติ}ติ \textbf{ติตฺถํ} วุจฺจติ ทิฏฺฐิคตํ.\csromanspacer{} ติตฺถิยา วุจฺจนฺติ ทิฏฺฐิคติกา.\csromanspacer{} ปุถุทิฏฺฐิยา ปุถุทิฏฺฐิคตานิ วทนฺติ กเถนฺติ ภณนฺติ ทีปยนฺติ โวหรนฺตีติ เอวมฺปิ ติตฺถฺยา ปุถุโส วทนฺติ.}

\prose{\textbf{สนฺทิฏฺฐิราเคน หิ เต’ภิรตฺตา}ติ สกาย ทิฏฺฐิยา ทิฏฺฐิราเคน รตฺตา อภิรตฺตาติ สนฺทิฏฺฐิราเคน หิ เต’ภิรตฺตา.\csromanspacer{} เตนาห ภควา\csromandash{}}

\csromangathameasure{อญฺญํ อิโต ยา’ภิวทนฺติ ธมฺมํ, \\ อปรทฺธา สุทฺธิ’มเกวลี เต.}
\csromangathagroup{\csromangathastanza{อญฺญํ อิโต ยา’ภิวทนฺติ ธมฺมํ, \\ อปรทฺธา สุทฺธิ’มเกวลี เต.}}

\prose{\textbf{เอวมฺปิ ติตฺถฺยา ปุถุโส วทนฺติ},}

\prose{\textbf{สนฺทิฏฺฐิราเคน หิ เต’ภิรตฺตา}ติ.}

\proseitem{127}{\csromansymbolfootnote{*}{ขุ 1. 417 ปิฏฺเฐ.}อิเธว สุทฺธิํ อิติ วาทยนฺติ,}

\prose{\textbf{นาญฺเญสุ ธมฺเมสุ วิสุทฺธิมาหุ.}}

\csromangathameasure{\textbf{เอวมฺปิ ติตฺถฺยา ปุถุโส นิวิฏฺฐา,} \\ \textbf{สกายเน ตตฺถ ทฬฺหํ วทานา.}}
\csromangathagroup{\csromangathastanza{\textbf{เอวมฺปิ ติตฺถฺยา ปุถุโส นิวิฏฺฐา,} \\ \textbf{สกายเน ตตฺถ ทฬฺหํ วทานา.}}}

\prose{\textbf{อิเธว สุทฺธิํ อิติ วาทยนฺตี}ติ อิธ สุทฺธิํ วิสุทฺธิํ ปริสุทฺธิํ, มุตฺติํ วิมุตฺติํ ปริมุตฺติํ วทนฺติ กเถนฺติ ภณนฺติ ทีปยนฺติ โวหรนฺติ, “สสฺสโต โลโก, อิทเมว สจฺจํ โมฆมญฺญนฺ”ติ อิธ สุทฺธิํ วิสุทฺธิํ ปริสุทฺธิํ มุตฺติํ วิมุตฺติํ ปริมุตฺติํ วทนฺติ กเถนฺติ ภณนฺติ ทีปยนฺติ โวหรนฺติ, “อสสฺสโต โลโก ฯเปฯ เนว โหติ น น โหติ ตถาคโต ปรํ มรณา, อิทเมว สจฺจํ โมฆมญฺญนฺ”ติ อิธ สุทฺธิํ วิสุทฺธิํ ปริสุทฺธิํ, มุตฺติํ วิมุตฺติํ ปริมุตฺติํ วทนฺติ กเถนฺติ ภณนฺติ ทีปยนฺติ โวหรนฺตีติ อิเธว สุทฺธิํ อิติ วาทยนฺติ.}

\prose{\textbf{นาญฺเญสุ ธมฺเมสุ วิสุทฺธิมาหู}ติ อตฺตโน สตฺถารํ ธมฺมกฺขานํ คณํ ทิฏฺฐิํ ปฏิปทํ มคฺคํ ฐเปตฺวา สพฺเพ ปรวาเท ขิปนฺติ อุกฺขิปนฺติ ปริกฺขิปนฺติ, โส สตฺถา น สพฺพญฺญู, ธมฺโม น สฺวากฺขาโต, คโณ น สุปฺปฏิปนฺโน, ทิฏฺฐิ น ภทฺทิกา, ปฏิปทา น สุปญฺญตฺตา, มคฺโค น นิยฺยานิโก, นตฺเถตฺถ สุทฺธิ วา วิสุทฺธิ วา ปริสุทฺธิ วา, มุตฺติ วา วิมุตฺติ วา ปริมุตฺติ วา, นตฺเถตฺถ สุชฺฌนฺติ วา วิสุชฺฌนฺติ วา ปริสุชฺฌนฺติ วา, มุจฺจนฺติ วา วิมุจฺจนฺติ วา ปริมุจฺจนฺติ วา, \csromanfolio{234}หีนา นิหีนา โอมกา ลามกา ฉตุกฺกา ปริตฺตาติ เอวมาหํสุ เอวํ กเถนฺติ เอวํ ภณนฺติ เอวํ ทีปยนฺติ เอวํ โวหรนฺตีติ นาญฺเญสุ ธมฺเมสุ วิสุทฺธิมาหุ.}

\prose{\textbf{เอวมฺปิ ติตฺถฺยา ปุถุโส นิวิฏฺฐา}ติ \textbf{ติตฺถํ} วุจฺจติ ทิฏฺฐิคตํ, ติตฺถิยา วุจฺจนฺติ ทิฏฺฐิคติกา, ปุถุทิฏฺฐิยา\footnote{ปุถุติตฺถิยา (สี, ก) ปุริมคาถาย ปาฐเภโท นตฺถิ.} ปุถุทิฏฺฐิคเตสุ นิวิฏฺฐา ปติฏฺฐิตา อลฺลีนา อุปคตา อชฺโฌสิตา อธิมุตฺตาติ เอวมฺปิ ติตฺถฺยา ปุถุโส นิวิฏฺฐา.}

\prose{\textbf{สกายเน ตตฺถ ทฬฺหํ วทานา}ติ ธมฺโม สกายนํ ทิฏฺฐิ สกายนํ ปฏิปทา สกายนํ มคฺโค สกายนํ, สกายเน ทฬฺหวาทา ถิรวาทา พลิกวาทา อวฏฺฐิตวาทาติ สกายเน ตตฺถ ทฬฺหํ วทานา.\csromanspacer{} เตนาห ภควา\csromandash{}}

\csromangathameasure{อิเธว สุทฺธิํ อิติ วาทยนฺติ, \\ นาญฺเญสุ ธมฺเมสุ วิสุทฺธิมาหุ.}
\csromangathagroup{\csromangathastanza{อิเธว สุทฺธิํ อิติ วาทยนฺติ, \\ นาญฺเญสุ ธมฺเมสุ วิสุทฺธิมาหุ.}}

\prose{\textbf{เอวมฺปิ ติตฺถฺยา ปุถุโส นิวิฏฺฐา},}

\prose{\textbf{สกายเน ตตฺถ ทฬฺหํ วทานา}ติ.}

\csromanhangingbegin
\hangingheaditem{128}{\textbf{สกายเน วาปิ ทฬฺหํ วทาโน},}
\hangingbody{\textbf{กเมตฺถ}2 \textbf{พาโลติ ปรํ ทเหยฺย.}}
\hangingbody{\textbf{สยํว}3 \textbf{โส เมธคมาวเหยฺย,}}
\hangingbody{\textbf{ปรํ วทํ พาลมสุทฺธิธมฺมํ.}}
\csromanhangingend

\prose{\textbf{สกายเน วาปิ ทฬฺหํ วทาโน}ติ ธมฺโม สกายนํ ทิฏฺฐิ สกายนํ ปฏิปทา สกายนํ มคฺโค สกายนํ, สกายเน ทฬฺหวาโท ถิรวาโท พลิกวาโท อวฏฺฐิตวาโทติ สกายเน วาปิ ทฬฺหํ วทาโน.}

\prose{\textbf{กเมตฺถ พาโลติ ปรํ ทเหยฺยา}ติ \textbf{เอตฺถา}ติ สกาย ทิฏฺฐิยา สกาย ขนฺติยา สกาย รุจิยา สกาย ลทฺธิยา ปรํ พาลโต หีนโต นิหีนโต โอมกโต ลามกโต ฉตุกฺกโต ปริตฺตโต กํ ทเหยฺย, กํ ปสฺเสยฺย, กํ ทกฺเขยฺย, กํ โอโลเกยฺย, กํ นิชฺฌาเยยฺย, กํ อุปปริกฺเขยฺยาติ กเมตฺถ \textbf{พาโลติ ปรํ ทเหยฺย.}}

\csromanheader{2}{12. จูฬพฺยูหสุตฺตนิทฺเทส}
\prose{\csromanfolio{235}\textbf{สยํว โส เมธคมาวเหยฺย, ปรํ วทํ พาลมสุทฺธิธมฺมนฺ}ติ ปโรพาโล หีโน นิหีโน โอมโก ลามโก ฉตุกฺโก ปริตฺโต อสุทฺธิธมฺโม อวิสุทฺธิธมฺโม อปริสุทฺธิธมฺโม อโวทาตธมฺโมติ เอวํ วทนฺโต เอวํ กเถนฺโต เอวํ ภณนฺโต เอวํ ทีปยนฺโต เอวํ โวหรนฺโต สยเมว กลหํ ภณฺฑนํ วิคฺคหํ วิวาทํ เมธคํ อาวเหยฺย สมาวเหยฺย อาหเรยฺย สมาหเรยฺย อากฑฺเฒยฺย สมากฑฺเฒยฺย คเณฺหยฺย ปรามเสยฺย อภินิวิเสยฺยาติ สยํว โส เมธคมาวเหยฺย, ปรํ วทํ พาลมสุทฺธิธมฺมํ.\csromanspacer{} เตนาห ภควา\csromandash{}}

\csromangathameasure{สกายเน วาปิ ทฬฺหํ วทาโน, \\ กเมตฺถ พาโลติ ปรํ ทเหยฺย. \\ สยํว โส เมธคมาวเหยฺย. \\ ปรํ วทํ พาลมสุทฺธิธมฺมนฺติ.}
\csromangathagroup{\csromangathastanza{สกายเน วาปิ ทฬฺหํ วทาโน, \\ กเมตฺถ พาโลติ ปรํ ทเหยฺย. \\ สยํว โส เมธคมาวเหยฺย. \\ ปรํ วทํ พาลมสุทฺธิธมฺมนฺติ.}}

\proseitem{129}{\csromansymbolfootnote{*}{ขุ 1. 418 ปิฏฺเฐ.}วินิจฺฉเย ฐตฺวา สยํ ปมาย,}

\prose{\textbf{อุทฺธํ ส}\footnote{อุทฺธํ โส (สฺยา)} \textbf{โลกสฺมิํ วิวาทเมติ.}}

\csromangathameasure{\textbf{หิตฺวาน สพฺพานิ วินิจฺฉยานิ,} \\ \textbf{น เมธคํ กุพฺพติ ชนฺตุ โลเก.}}
\csromangathagroup{\csromangathastanza{\textbf{หิตฺวาน สพฺพานิ วินิจฺฉยานิ,} \\ \textbf{น เมธคํ กุพฺพติ ชนฺตุ โลเก.}}}

\prose{\textbf{วินิจฺฉเย ฐตฺวา สยํ ปมายา}ติ วินิจฺฉยา วุจฺจนฺติ ทฺวาสฏฺฐิ ทิฏฺฐิคตานิ.\csromanspacer{} วินิจฺฉยทิฏฺฐิยา ฐตฺวา ปติฏฺฐหิตฺวา คณฺหิตฺวา ปรามสิตฺวา อภินิวิสิตฺวาติ วินิจฺฉเย ฐตฺวา.\csromanspacer{} \textbf{สยํ ปมายา}ติ สยํ ปมาย ปมินิตฺวา “อยํ สตฺถา สพฺพญฺญู”ติ สยํ ปมาย ปมินิตฺวา “อยํ ธมฺโม สฺวากฺขาโต, อยํ คโณ สุปฺปฏิปนฺโน, อยํ ทิฏฺฐิ ภทฺทิกา, อยํ ปฏิปทา สุปญฺญตฺตา, อยํ มคฺโค นิยฺยานิโก”ติ สยํ ปมาย ปมินิตฺวาติ วินิจฺฉเย ฐตฺวา สยํ ปมาย.}

\prose{\textbf{อุทฺธํ ส โลกสฺมิํ วิวาทเมตี}ติ อุทฺธํ โส วุจฺจติ อนาคตํ.\csromanspacer{} อตฺตโน วาทํ อุทฺธํ ฐเปตฺวา สยเมว กลหํ ภณฺฑนํ วิคฺคหํ วิวาทํ เมธคํ เอติ อุเปติ อุปคจฺฉติ คณฺหาติ ปรามสติ อภินิวิสตีติ เอวมฺปิ อุทฺธํ ส \textbf{โลกสฺมิํ วิวาทเมติ.}\csromanspacer{} อถ วา อญฺเญน อุทฺธํ วาเทน สทฺธิํ กลหํ กโรติ, ภณฺฑนํ กโรติ, วิคฺคหํ กโรติ, วิวาทํ กโรติ, เมธคํ กโรติ. “น ตฺวํ อิมํ ธมฺมวินยํ อาชานาสิ ฯเปฯ นิพฺเพเฐหิ วา สเจ ปโหสี”ติ เอวมฺปิ อุทฺธํ ส \textbf{โลกสฺมิํ วิวาทเมติ.}}

\prose{\csromanfolio{236}\textbf{หิตฺวาน สพฺพานิ วินิจฺฉยานี}ติ วินิจฺฉยา วุจฺจนฺติ ทฺวาสฏฺฐิ ทิฏฺฐิคตานิ\textbf{.}\csromanspacer{} ทิฏฺฐิวินิจฺฉยา สพฺเพ วินิจฺฉเย\footnote{สพฺพา วินิจฺฉิตทิฏฺฐิโย (สฺยา), สพฺพา วินิจฺฉยทิฏฺฐิโย (ก)} หิตฺวา จชิตฺวา ปริจฺจชิตฺวา ชหิตฺวา ปชหิตฺวา วิโนเทตฺวา พฺยนฺติํ กริตฺวา อนภาวํ คเมตฺวาติ หิตฺวาน สพฺพานิ วินิจฺฉยานิ\textbf{.}}

\prose{\textbf{น เมธคํ กุพฺพติ ชนฺตุ โลเก}ติ น กลหํ กโรติ, น ภณฺฑนํ กโรติ, น วิคฺคหํ กโรติ, น วิวาทํ กโรติ, น เมธคํ กโรติ\textbf{.}\csromanspacer{} วุตฺตํ เหตํ ภควตา\csromandash{} * เอวํ วิมุตฺตจิตฺโต โข อคฺคิเวสฺสน ภิกฺขุ น เกนจิ สํวทติ, น เกนจิ วิวทติ\textbf{.}\csromanspacer{} ยญฺจ โลเก วุตฺตํ, เตน จ โวหรติ อปรามสนฺติ\textbf{.}\csromanspacer{}\csromanspacer{}\textbf{ชนฺตู}ติ สตฺโต นโร มานโว โปโส ปุคฺคโล ชีโว ชาคุ ชนฺตุ อินฺทคุ มนุโช\textbf{.}\csromanspacer{} \textbf{โลเก}ติ อปายโลเก ฯเปฯ อายตนโลเกติ น เมธคํ กุพฺพติ ชนฺตุ โลเกติ\textbf{.}\csromanspacer{} เตนาห ภควา\csromandash{}}

\csromangathameasure{วินิจฺฉเย ฐตฺวา สยํ ปมาย, \\ อุทฺธํส โลกสฺมิํ วิวาทเมติ\textbf{.} \\ หิตฺวาน สพฺพานิ วินิจฺฉยานิ,}
\csromangathagroup{\csromangathastanza{วินิจฺฉเย ฐตฺวา สยํ ปมาย, \\ อุทฺธํส โลกสฺมิํ วิวาทเมติ\textbf{.} \\ หิตฺวาน สพฺพานิ วินิจฺฉยานิ,}}

\prose{\textbf{น เมธคํ กุพฺพติ ชนฺตุ โลเก}ติ\textbf{.}\csromanspacer{} จูฬวิยูหสุตฺตนิทฺเทโส ทฺวาทสโม\textbf{.}}
\csromansectionrule

\csromanheader{2}{13. มหาวิยูหสุตฺตนิทฺเทส}
\csromanheader{2}{อถ มหาวิยูหสุตฺตนิทฺเทสํ วกฺขติ\csromandash{}}
\csromanhangingbegin
\hangingheaditem{130}{\textbf{เย เกจิเม ทิฏฺฐิปริพฺพสานา},}
\hangingbody{\textbf{อิทเมว สจฺจนฺติ จ วาทยนฺติ}2\textbf{.}}
\hangingbody{\textbf{สพฺเพว เต นินฺทมนฺวานยนฺติ,}}
\hangingbody{\textbf{อโถ ปสํสมฺปิ ลภนฺติ ตตฺถ.}}
\csromanhangingend

\prose{\textbf{เย เกจิเม ทิฏฺฐิปริพฺพสานา}ติ \textbf{เย เกจี}ติ สพฺเพน สพฺพํ สพฺพถา สพฺพํ อเสสํ นิสฺเสสํ ปริยาทิยนวจนเมตํ \textbf{เย เกจี}ติ\textbf{.}\csromanspacer{} \textbf{ทิฏฺฐิปริพฺพสานา}ติ สนฺเตเก สมณพฺราหฺมณา ทิฏฺฐิคติกา, เต ทฺวาสฏฺฐิยา ทิฏฺฐิคตานํ อญฺญตรญฺญตรํ ทิฏฺฐิคตํ คเหตฺวา อุคฺคเหตฺวา คณฺหิตฺวา}

\vspace{\tipitakaparskip}
\csromancenter{มหาวิยูหสุตฺตนิทฺเทส ปรามสิตฺวา อภินิวิสิตฺวา สกาย สกาย ทิฏฺฐิยา วสนฺติ สํวสนฺติ อาวสนฺติ ปริวสนฺติ.\csromanspacer{} ยถา อคาริกา วา ฆเรสุ วสนฺติ, สาปตฺติกา วา อาปตฺตีสุ วสนฺติ, สกิเลสา วา กิเลเสสุ วสนฺติ, เอวเมวํ สนฺเตเก ฯเปฯ ปริวสนฺตีติ เย เกจิเม ทิฏฺฐิปริพฺพสานา.}
\prose{\csromanfolio{237}\textbf{อิทเมว สจฺจนฺติ จ วาทยนฺตี}ติ “สสฺสโต โลโก, อิทเมว สจฺจํ โมฆมญฺญนฺ”ติ วทนฺติ กเถนฺติ ภณนฺติ ทีปยนฺติ โวหรนฺติ, “อสสฺสโต โลโก ฯเปฯ เนว โหติ น น โหติ ตถาคโต ปรํ มรณา, อิทเมว สจฺจํ โมฆมญฺญนฺ”ติ วทนฺติ กเถนฺติ ภณนฺติ ทีปยนฺติ โวหรนฺตีติ อิทเมว สจฺจนฺติ จ วาทยนฺติ.}

\prose{\textbf{สพฺเพว เต นินฺทมนฺวานยนฺตี}ติ สพฺเพว เต สมณพฺราหฺมณา นินฺทเมว อเนฺวนฺติ, ครหเมว อเนฺวนฺติ, อกิตฺติเมว อเนฺวนฺติ, สพฺเพ นินฺทิตาเยว โหนฺติ, ครหิตาเยว โหนฺติ, อกิตฺติตาเยว โหนฺตีติ สพฺเพว เต นินฺทมนฺวานยนฺติ.}

\prose{\textbf{อโถ ปสํสมฺปิ ลภนฺติ ตตฺถา}ติ ตตฺถ สกาย ทิฏฺฐิยา สกาย ขนฺติยา สกาย รุจิยา สกาย ลทฺธิยา ปสํสํ โถมนํ กิตฺติํ วณฺณหาริกํ ลภนฺติ ปฏิลภนฺติ อุปคจฺฉนฺติ วินฺทนฺตีติ อโถ ปสํสมฺปิ ลภนฺติ ตตฺถ.\csromanspacer{} เตนาห โส นิมฺมิโต\csromandash{}}

\prose{เย เกจิเม ทิฏฺฐิปริพฺพสานา.}

\prose{อิทเมว สจฺจนฺติ จ วาทยนฺติ.}

\csromangathameasure{สพฺเพว เต นินฺทมนฺวานยนฺติ,}
\csromangathagroup{\csromangathastanza{สพฺเพว เต นินฺทมนฺวานยนฺติ,}}

\prose{\textbf{อโถ ปสํสมฺปิ ลภนฺติ ตตฺถา}ติ.}

\proseitem{131}{\csromansymbolfootnote{*}{ขุ 1. 418 ปิฏฺเฐ.}อปฺปํ หิ เอตํ น อลํ สมาย,}

\prose{\textbf{ทุเว วิวาทสฺส ผลานิ พฺรูมิ.}}

\csromangathameasure{\textbf{เอตมฺปิ ทิสฺวา น วิวาทเยถ,} \\ \textbf{เขมา’ภิปสฺสํ อวิวาทภูมิํ.}}
\csromangathagroup{\csromangathastanza{\textbf{เอตมฺปิ ทิสฺวา น วิวาทเยถ,} \\ \textbf{เขมา’ภิปสฺสํ อวิวาทภูมิํ.}}}

\prose{\textbf{อปฺปํ หิ เอตํ น อลํ สมายา}ติ \textbf{อปฺปํ หิ เอตนฺ}ติ อปฺปกํ เอตํ, โอมกํ เอตํ, โถกกํ เอตํ, ลามกํ เอตํ, ฉตุกฺกํ เอตํ, ปริตฺตกํ เอตนฺติ อปฺปํ หิ เอตํ.\csromanspacer{} \textbf{น อลํ สมายา}ติ นาลํ ราคสฺส สมาย, โทสสฺส สมาย, โมหสฺส \csromanfolio{238}สมาย, โกธสฺส.\csromanspacer{} อุปนาหสฺส.\csromanspacer{} มกฺขสฺส.\csromanspacer{} ปฬาสสฺส.\csromanspacer{} อิสฺสาย.\csromanspacer{} มจฺฉริยสฺส.\csromanspacer{} มายาย.\csromanspacer{} สาเฐยฺยสฺส.\csromanspacer{} ถมฺภสฺส.\csromanspacer{} สารมฺภสฺส.\csromanspacer{} มานสฺส.\csromanspacer{} อติมานสฺส.\csromanspacer{} มทสฺส.\csromanspacer{} ปมาทสฺส.\csromanspacer{} สพฺพกิเลสานํ.\csromanspacer{} สพฺพทุจฺจริตานํ.\csromanspacer{} สพฺพทรถานํ.\csromanspacer{} สพฺพปริฬาหานํ.\csromanspacer{} สพฺพสนฺตาปานํ.\csromanspacer{} สพฺพากุสลาภิสงฺขารานํ สมาย อุปสมาย วูปสมาย นิพฺพานาย ปฏินิสฺสคฺคาย ปฏิปสฺสทฺธิยาติ อปฺปํ หิ เอตํ น อลํ สมาย.}

\prose{\textbf{ทุเว วิวาทสฺส ผลานิ พฺรูมี}ติ ทิฏฺฐิกลหสฺส ทิฏฺฐิภณฺฑนสฺส ทิฏฺฐิวิคฺคหสฺส ทิฏฺฐิวิวาทสฺส ทิฏฺฐิเมธคสฺส เทฺว ผลานิ โหนฺติ.\csromanspacer{} ชยปราชโย โหติ, ลาภาลาโภ โหติ, ยสายโส โหติ, นินฺทาปสํโส โหติ, สุขทุกฺขํ โหติ, โสมนสฺสโทมนสฺสํ โหติ, อิฏฺฐานิฏฺฐํ โหติ, อนุนยปฏิฆํ โหติ, อุคฺฆาตินิคฺฆาติ โหติ, อนุโรธวิโรโธ โหติ.\csromanspacer{} อถ วา ตํ กมฺมํ นิรยสํวตฺตนิกํ ติรจฺฉานโยนิสํวตฺตนิกํ เปตฺติวิสยสํวตฺตนิกนฺติ พฺรูมิ อาจิกฺขามิ เทเสมิ ปญฺญเปมิ ปฏฺฐเปมิ วิวรามิ วิภชามิ อุตฺตานีกโรมิ ปกาเสมีติ ทุเว วิวาทสฺส ผลานิ พฺรูมิ.}

\prose{\textbf{เอตมฺปิ ทิสฺวา น วิวาทเยถา}ติ \textbf{เอตมฺปิ ทิสฺวา}ติ เอตํ อาทีนวํ ทิสฺวา ปสฺสิตฺวา ตุลยิตฺวา ตีรยิตฺวา วิภาวยิตฺวา วิภูตํ กตฺวา ทิฏฺฐิกลเหสุ ทิฏฺฐิภณฺฑเนสุ ทิฏฺฐิวิคฺคเหสุ ทิฏฺฐิวิวาเทสุ ทิฏฺฐิเมธเคสูติ \textbf{เอตมฺปิ ทิสฺวา}.\csromanspacer{} \textbf{น วิวาทเยถา}ติ น กลหํ กเรยฺย, น ภณฺฑนํ กเรยฺย, น วิคฺคหํ กเรยฺย, น วิวาทํ กเรยฺย, น เมธคํ กเรยฺย, กลหํ ภณฺฑนํ วิคฺคหํ วิวาทํ เมธคํ ปชเหยฺยวิโนเทยฺย พฺยนฺติํ กเรยฺย อนภาวํ คเมยฺย, กลหา ภณฺฑนา วิคฺคหา วิวาทา เมธคา อารโต อสฺส วิรโต ปฏิวิรโต นิกฺขนฺโต นิสฺสโฏ วิปฺปมุตฺโต วิสญฺญุตฺโต วิมริยาทิกเตน เจตสา วิหเรยฺยาติ \textbf{เอตมฺปิ ทิสฺวา} น วิวาทเยถ.}

\prose{\textbf{เขมา’ภิปสฺสํ อวิวาทภูมิ}นฺติ \textbf{อวิวาทภูมิ} วุจฺจติ อมตํ นิพฺพานํ.\csromanspacer{} โย โส สพฺพสงฺขารสมโถ สพฺพูปธิปฏินิสฺสคฺโค ตณฺหกฺขโย วิราโค นิโรโธ นิพฺพานํ, เอตํ \textbf{อวิวาทภูมิ}ํ เขมโต ตาณโต เลณโต สรณโต อภยโต อจฺจุตโต อมตโต นิพฺพานโต ปสฺสนฺโต ทกฺขนฺโต โอโลเกนฺโต นิชฺฌายนฺโต อุปปริกฺขนฺโตติ เขมา’ภิปสฺสํ \textbf{อวิวาทภูมิ}ํ.\csromanspacer{} เตนาห ภควา\csromandash{}}

\csromanheader{2}{13. มหาวิยูหสุตฺตนิทฺเทส}
\csromangathameasure{อปฺปํ หิ เอตํ น อลํ สมาย, \\ ทุเว วิวาทสฺส ผลานิ พฺรูมิ, \\ เอตมฺปิ ทิสฺวา น วิวาทเยถ, \\ เขมา’ภิปสฺสํ อวิวาทภูมินฺติ.}
\csromangathagroup{\csromangathastanza{\csromanfolio{239}อปฺปํ หิ เอตํ น อลํ สมาย, \\ ทุเว วิวาทสฺส ผลานิ พฺรูมิ, \\ เอตมฺปิ ทิสฺวา น วิวาทเยถ, \\ เขมา’ภิปสฺสํ อวิวาทภูมินฺติ.}}

\proseitem{132}{\csromansymbolfootnote{*}{ขุ 1. 418 ปิฏฺเฐ.}\textbf{ยา กาจิมา สมฺมุติโย ปุถุชฺชา},}

\prose{\textbf{สพฺพาว เอตา น อุเปติ วิทฺวา.}}

\csromangathameasure{\textbf{อนูปโย โส อุปยํ กิเมยฺย,} \\ \textbf{ทิฏฺเฐ สุเต ขนฺติมกุพฺพมาโน.}}
\csromangathagroup{\csromangathastanza{\textbf{อนูปโย โส อุปยํ กิเมยฺย,} \\ \textbf{ทิฏฺเฐ สุเต ขนฺติมกุพฺพมาโน.}}}

\prose{\textbf{ยา กาจิมา สมฺมุติโย ปุถุชฺชา}ติ \textbf{ยา กาจี}ติ สพฺเพน สพฺพํ สพฺพถา สพฺพํ อเสสํ นิสฺเสสํ ปริยาทิยนวจนเมตํ \textbf{ยา กาจี}ติ.\csromanspacer{} \textbf{สมฺมุติโย}ติ สมฺมุติโย วุจฺจนฺติ ทฺวาสฏฺฐิ ทิฏฺฐิคตานิ ทิฏฺฐิสมฺมุติโย.\csromanspacer{} \textbf{ปุถุชฺชา}ติ ปุถุชฺชเนหิ ชนิตา สมฺมุติโยติ ปุถุชฺชา, ปุถุ นานาชเนหิ ชนิตา วา สมฺมุติโยติ ปุถุชฺชาติ ยา กาจิมา สมฺมุติโย ปุถุชฺชา.}

\prose{\textbf{สพฺพาว เอตา น อุเปติ วิทฺวา}ติ วิทฺวา วิชฺชาคโต ญาณี วิภาวี เมธาวี สพฺพาว เอตา ทิฏฺฐิสมฺมุติโย เนติ น อุเปติ น อุปคจฺฉติ น คณฺหาติ น ปรามสติ นาภินิวิสตีติ สพฺพาว เอตา น อุเปติ วิทฺวา.}

\prose{\textbf{อนุปโย โส อุปยํ กิเมยฺยา}ติ \textbf{อุปโย}ติ เทฺว อุปยา ตณฺหูปโย จ ทิฏฺฐูปโย จ ฯเปฯ อยํ ตณฺหูปโย ฯเปฯ อยํ ทิฏฺฐูปโย.\csromanspacer{} ตสฺส ตณฺหูปโย ปหีโน, ทิฏฺฐูปโย ปฏินิสฺสฏฺโฐ, ตณฺหูปยสฺส ปหีนตฺตา, ทิฏฺฐูปยสฺส ปฏินิสฺสฏฺฐตฺตา อนูปโย ปุคฺคโล กิํ รูปํ อุเปยฺย อุปคจฺเฉยฺย คเณฺหยฺย ปรามเสยฺย อภินิเวเสยฺย “อตฺตา เม”ติ.\csromanspacer{} กิํ เวทนํ.\csromanspacer{} กิํ สญฺญํ.\csromanspacer{} กิํ สงฺขาเร.\csromanspacer{} กิํ วิญฺญาณํ.\csromanspacer{} กิํ คติํ.\csromanspacer{} กิํ อุปปตฺติํ.\csromanspacer{} กิํ ปฏิสนฺธิํ.\csromanspacer{} กิํ ภวํ.\csromanspacer{} กิํ สํสารํ.\csromanspacer{} กิํ วฏฺฏํ อุเปยฺย อุปคจฺเฉยฺย คเณฺหยฺย ปรามเสยฺย อภินิเวเสยฺยาติ อนฺ\textbf{อุปโย} โส อุปยํ กิเมยฺย.}

\prose{\textbf{ทิฏฺเฐ สุเต ขนฺติมกุพฺพมาโน}ติ ทิฏฺเฐ วา ทิฏฺฐสุทฺธิยา วา สุเต วา สุตสุทฺธิยา วา มุเต วา มุตสุทฺธิยา วา ขนฺติํ อกุพฺพมาโน ฉนฺทํ อกุพฺพมาโน เปมํ อกุพฺพมาโน ราคํ อกุพฺพมาโน อชนยมาโน อสญฺชนยมาโน อนิพฺพตฺตยมาโน อนภินิพฺพตฺตยมาโนติ ทิฏฺเฐ สุเต ขนฺติมกุพฺพมาโน.\csromanspacer{} เตนาห ภควา\csromandash{}}

\csromangathameasure{ยา กาจิมา สมฺมุติโย ปุถุชฺชา, \\ สพฺพาว เอตา น อุเปติ วิทฺวา. \\ อนูปโย โส อุปยํ กิเมยฺย, \\ ทิฏฺเฐ สุเต ขนฺติมกุพฺพมาโนติ.}
\csromangathagroup{\csromangathastanza{\csromanfolio{240}ยา กาจิมา สมฺมุติโย ปุถุชฺชา, \\ สพฺพาว เอตา น อุเปติ วิทฺวา. \\ อนูปโย โส อุปยํ กิเมยฺย, \\ ทิฏฺเฐ สุเต ขนฺติมกุพฺพมาโนติ.}}

\proseitem{133}{\csromansymbolfootnote{*}{ขุ 1. 418 ปิฏฺเฐ.}สีลุตฺตมา สญฺญเมนาหุ สุทฺธิํ,}

\prose{\textbf{วตํ สมาทาย อุปฏฺฐิตาเส.}}

\csromangathameasure{\textbf{อิเธว สิกฺเขม อถสฺส สุทฺธิํ,} \\ \textbf{ภวูปนีตา กุสลาวทานา.}}
\csromangathagroup{\csromangathastanza{\textbf{อิเธว สิกฺเขม อถสฺส สุทฺธิํ,} \\ \textbf{ภวูปนีตา กุสลาวทานา.}}}

\prose{\textbf{สีลุตฺตมา สญฺญเมนาหุ สุทฺธินฺ}ติ สนฺเตเก สมณพฺราหฺมณา สีลุตฺตมวาทา, เต สีลมตฺเตน สญฺญมมตฺเตน สํวรมตฺเตน อวีติกฺกมมตฺเตน สุทฺธิํ วิสุทฺธิํ ปริวิสุทฺธิํ, มุตฺติํ วิมุตฺติํ ปริวิมุตฺติํ อาหํสุ\footnote{อาหุ (สฺยา)} วทนฺติ กเถนฺติ ภณนฺติ ทีปยนฺติ โวหรนฺติ.}

\prose{\csromansymbolfootnote{+}{ม 2. 215 ปิฏฺเฐ.}สมณมุณฺฑิกาปุตฺโต เอวมาห\csromandash{}จตูหิ โข อหํ คหปติ ธมฺเมหิ สมนฺนาคตํ ปุริสปุคฺคลํ ปญฺญเปมิ สมฺปนฺนกุสลํ ปรมกุสลํ อุตฺตมปตฺติปตฺตํ สมณํ อโยชฺฌํ.\csromanspacer{} กตเมหิ จตูหิ, อิธ คหปติ น กาเยน ปาปกมฺมํ กโรติ, น ปาปิกํ\footnote{ปาปกํ (สี)} วาจํ ภาสติ, น ปาปกํ สงฺกปฺปํ สงฺกปฺเปติ, น ปาปกํ อาชีวํ อาชีวติ, อิเมหิ โข อหํ คหปติ จตูหิ ธมฺเมหิ สมนฺนาคตํ ปุริสปุคฺคลํ ปญฺญเปมิ สมฺปนฺนกุสลํ ปรมกุสลํ อุตฺตมปตฺติปตฺตํ สมณํ อโยชฺฌํ.\csromanspacer{} เอวเมวํ สนฺเตเก สมณพฺราหฺมณา สีลุตฺตมวาทา, เต สีลมตฺเตน สญฺญมมตฺเตน สํวรมตฺเตน อวีติกฺกมมตฺเตน สุทฺธิํ วิสุทฺธิํ ปริสุทฺธิํ, มุตฺติํ วิมุตฺติํ ปริมุตฺติํ อาหํสุ วทนฺติ กเถนฺติ ภณนฺติ ทีปยนฺติ โวหรนฺตีติ สีลุตฺตมา สญฺญเมนาหุ สุทฺธิํ.}

\prose{\textbf{วตํ สมาทาย อุปฏฺฐิตาเส}ติ \textbf{วต}นฺติ หตฺถิ\textbf{วต}ํ วา อสฺส\textbf{วต}ํ วา โค\textbf{วต}ํ วา กุกฺกุร\textbf{วต}ํ วา กาก\textbf{วต}ํ วา วาสุเทว\textbf{วต}ํ วา พลเทว\textbf{วต}ํ วา ปุณฺณภทฺท\textbf{วต}ํ วา มณิภทฺท\textbf{วต}ํ วา อคฺคิ\textbf{วต}ํ วา นาค\textbf{วต}ํ วา สุปณฺณ\textbf{วต}ํ วา ยกฺข\textbf{วต}ํ วา อสุร\textbf{วต}ํ วา คนฺธพฺพ\textbf{วต}ํ วา มหาราช\textbf{วต}ํ วา จนฺท\textbf{วต}ํ วา สูริย\textbf{วต}ํ วา อินฺท\textbf{วต}ํ วา พฺรหฺม\textbf{วต}ํ วา เทว\textbf{วต}ํ วา ทิสา\textbf{วต}ํ วา อาทาย สมาทาย อาทิยิตฺวา สมาทิยิตฺวา คณฺหิตฺวา ปรามสิตฺวา}

\vspace{\tipitakaparskip}
\csromancenter{มหาวิยูหสุตฺตนิทฺเทส อภินิวิสิตฺวา อุปฏฺฐิตา ปจฺจุปฏฺฐิตา อลฺลีนา อุปคตา อชฺโฌสิตา อธิมุตฺตาติ วตํ สมาทาย อุปฏฺฐิตาเส.}
\prose{\csromanfolio{241}\textbf{อิเธว สิกฺเขม อถสฺส สุทฺธินฺ}ติ \textbf{อิธา}ติ สกาย ทิฏฺฐิยา สกาย ขนฺติยา สกาย รุจิยา สกาย ลทฺธิยา.\csromanspacer{} \textbf{สิกฺเขมา}ติ สิกฺเขม อาจเรม สมาจเรม สมาทาย วตฺเตมาติ อิเธว สิกฺเขม.\csromanspacer{} \textbf{อถสฺส สุทฺธินฺ}ติ อถสฺส สุทฺธิํ วิสุทฺธิํ ปริสุทฺธิํ, มุตฺติํ วิมุตฺติํ ปริมุตฺตินฺติ อิเธว สิกฺเขม อถสฺส สุทฺธิํ.}

\prose{\textbf{ภวูปนีตา กุสลาวทานา}ติ \textbf{ภวูปนีตา}ติ \textbf{ภวูปนีตา} ภวูปคตา ภวชฺโฌสิตา ภวาธิมุตฺตาติ \textbf{ภวูปนีตา}.\csromanspacer{} \textbf{กุสลาวทานา}ติ กุสลวาทา ปณฺฑิตวาทา ถิรวาทา ญายวาทา เหตุวาทา ลกฺขณวาทา การณวาทา ฐานวาทา สกาย ลทฺธิยาติ \textbf{ภวูปนีตา} กุสลาวทานา.\csromanspacer{} เตนาห ภควา\csromandash{}}

\csromangathameasure{สีลุตฺตมา สญฺญเมนาหุ สุทฺธิํ, \\ วตํ สมาทาย อุปฏฺฐิตาเส.}
\csromangathagroup{\csromangathastanza{สีลุตฺตมา สญฺญเมนาหุ สุทฺธิํ, \\ วตํ สมาทาย อุปฏฺฐิตาเส.}}

\prose{อิเธว สิกฺเขม อถสฺส สุทฺธิํ,}

\prose{\textbf{ภวูปนีตา กุสลาวทานา}ติ.}

\csromanhangingbegin
\hangingheaditem{134}{สเจ จุโต สีลวตโต โหติ,}
\hangingbody{\textbf{ปเวธตี}1 \textbf{กมฺม วิราธยิตฺวา.}}
\hangingbody{\textbf{ปชปฺปตี}2 \textbf{ปตฺถยตี จ สุทฺธิํ,}}
\hangingbody{\textbf{สตฺถาว หีโน ปวสํ ฆรมฺหา.}}
\csromanhangingend

\prose{\textbf{สเจ จุโต สีลวตโต โหตี}ติ ทฺวีหิ การเณหิ สีลวตโต จวติ, ปรวิจฺฉินฺทนาย วา จวติ, อนภิสมฺภุณนฺโต วา จวติ.\csromanspacer{} กถํ ปรวิจฺฉินฺทนาย จวติ, ปโร วิจฺฉินฺทติ โส สตฺถา น สพฺพญฺญู, ธมฺโม น สฺวากฺขาโต, คโณ น สุปฺปฏิปนฺโน, ทิฏฺฐิ น ภทฺทิกา, ปฏิปทา น สุปญฺญตฺตา, มคฺโค น นิยฺยานิโก, นตฺเถตฺถ สุทฺธิ วา วิสุทฺธิ วา ปริสุทฺธิ วา, มุตฺติ วา วิมุตฺติ วา ปริมุตฺติ วา, นตฺเถตฺถ สุชฺฌนฺติ วา วิสุชฺฌนฺติ วา ปริสุชฺฌนฺติ วา, มุจฺจนฺติ วา วิมุจฺจนฺติ วา ปริมุจฺจนฺติ วา, หีนา นิหีนา โอมกา ลามกา ฉตุกฺกา ปริตฺตาติ เอวํ ปโร วิจฺฉินฺทติ.\csromanspacer{} เอวํ วิจฺฉินฺทิยมาโน สตฺถารา จวติ, ธมฺมกฺขานา จวติ, คณา จวติ, ทิฏฺฐิยา จวติ, ปฏิปทาย จวติ, มคฺคโต จวติ.\csromanspacer{} เอวํ \csromanfolio{242}ปรวิจฺฉินฺทนาย จวติ.\csromanspacer{} กถํ อนภิสมฺภุณนฺโต จวติ, สีลํ อนภิสมฺภุณนฺโต สีลโต จวติ, วตํ อนภิสมฺภุณนฺโต วตโต จวติ, สีลพฺพตํ อนภิสมฺภุณนฺโต สีลพฺพตโต จวติ.\csromanspacer{} เอวํ อนภิสมฺภุณนฺโต จวตีติ สเจ จุโต สีลวตโต โหติ.}

\prose{\textbf{ปเวธติ กมฺม วิราธยิตฺวา}ติ \textbf{ปเวธตี}ติ สีลํ วา วตํ วา สีลพฺพตํ วา “วิรทฺธํ มยา, อปรทฺธํ มยา, ขลิตํ มยา, คลิตํ มยา, อญฺญาย อปรทฺโธ อหนฺ”ติ เวธติ ปเวธติ สมฺ\textbf{ปเวธตี}ติ ปเวธติ.\csromanspacer{}\csromanspacer{}\textbf{กมฺม วิราธยิตฺวา}ติ ปุญฺญาภิสงฺขารํ วา อปุญฺญาภิสงฺขารํ วา อาเนญฺชาภิสงฺขารํ วา “วิรทฺธํ มยา, อปรทฺธํ มยา, ขลิตํ มยา, คลิตํ มยา, อญฺญาย อปรทฺโธ อหนฺ”ติ เวธติ ปเวธติ สมฺ\textbf{ปเวธตี}ติ ปเวธติ กมฺม วิราธยิตฺวา.}

\prose{\textbf{ปชปฺปตี ปตฺถยตี จ สุทฺธิ}นฺติ \textbf{ปชปฺปตี}ติ สีลํ วา ชปฺปติ, วตํ วา ชปฺปติ, สีลพฺพตํ วา ชปฺปติ ปชปฺปติ อภิชปฺปตีติ ปชปฺปติ.\csromanspacer{} \textbf{ปตฺถยตี จ สุทฺธินฺ}ติ สีลสุทฺธิํ วา ปตฺเถติ, วตสุทฺธิํ วา ปตฺเถติ, สีลพฺพตสุทฺธิํ วา ปตฺเถติ ปิเหติ อภิชปฺปตีติ \textbf{ปชปฺปตี} ปตฺถยตี จ สุทฺธิํ.}

\prose{\textbf{สตฺถาว หีโน ปวสํ ฆรมฺหา}ติ ยถา ปุริโส ฆรโต นิกฺขนฺโต สตฺเถน ปวสํ วสนฺโต สตฺถา โอหีโน.\csromanspacer{} ตํ วา สตฺถํ อนุพนฺธติ, สกํ วา ฆรํ ปจฺจาคจฺฉติ.\csromanspacer{} เอวเมวํ โส ทิฏฺฐิคติโก ตํ วา สตฺถารํ คณฺหาติ, อญฺญํ วา สตฺถารํ คณฺหาติ.\csromanspacer{} ตํ วา ธมฺมกฺขานํ คณฺหาติ, อญฺญํ วา ธมฺมกฺขานํ คณฺหาติ.\csromanspacer{} ตํ วา คณํ คณฺหาติ, อญฺญํ วา คณํ คณฺหาติ.\csromanspacer{} ตํ วา ทิฏฺฐิํ คณฺหาติ, อญฺญํ วา ทิฏฺฐิํ คณฺหาติ.\csromanspacer{} ตํ วา ปฏิปทํ คณฺหาติ, อญฺญํ วา ปฏิปทํ คณฺหาติ.\csromanspacer{} ตํ วา มคฺคํ คณฺหาติ, อญฺญํ วา มคฺคํ คณฺหาติ ปรามสติ อภินิวิสตีติ สตฺถาว หีโน ปวสํ ฆรมฺหา.\csromanspacer{} เตนาห ภควา\csromandash{}}

\csromangathameasure{สเจ จุโต สีลวตโต โหติ, \\ ปเวธตี กมฺม วิราธยิตฺวา.}
\csromangathagroup{\csromangathastanza{สเจ จุโต สีลวตโต โหติ, \\ ปเวธตี กมฺม วิราธยิตฺวา.}}

\prose{\textbf{ปชปฺปตี ปตฺถยตี จ สุทฺธิ}ํ,}

\prose{\textbf{สตฺถาว หีโน ปวสํ ฆรมฺหา}ติ.}

\csromanheader{2}{13. มหาวิยูหสุตฺตนิทฺเทส}
\proseitem{135}{\csromanfolio{243}\csromansymbolfootnote{*}{ขุ 1. 419 ปิฏฺเฐ.}สีลพฺพตํ วาปิ ปหาย สพฺพํ,}

\prose{\textbf{กมฺมญฺจ สาวชฺชนวชฺชเมตํ.}}

\csromangathameasure{\textbf{สุทฺธิํ อสุทฺธินฺติ อปตฺถยาโน,} \\ \textbf{วิรโต จเร สนฺติมนุคฺคหาย.}}
\csromangathagroup{\csromangathastanza{\textbf{สุทฺธิํ อสุทฺธินฺติ อปตฺถยาโน,} \\ \textbf{วิรโต จเร สนฺติมนุคฺคหาย.}}}

\prose{\textbf{สีลพฺพตํ วาปิ ปหาย สพฺพนฺ}ติ สพฺพา สีลสุทฺธิโย ปหาย ปชหิตฺวา วิโนเทตฺวา พฺยนฺติํ กริตฺวา อนภาวํ คเมตฺวา, สพฺพา วตสุทฺธิโย ปหาย ปชหิตฺวา วิโนเทตฺวา พฺยนฺติํ กริตฺวา อนภาวํ คเมตฺวา, สพฺพา สีลพฺพตสุทฺธิโย ปหาย ปชหิตฺวา วิโนเทตฺวา พฺยนฺติํ กริตฺวา อนภาวํ คเมตฺวาติ สีลพฺพตํ วาปิ ปหาย สพฺพํ.}

\prose{\textbf{กมฺมญฺจ สาวชฺชนวชฺชเมตนฺ}ติ \textbf{สาวชฺชกมฺมํ} วุจฺจติ กณฺหํ กณฺหวิปากํ.\csromanspacer{} \textbf{อนวชฺชกมฺมํ} วุจฺจติ สุกฺกํ สุกฺกวิปากํ.\csromanspacer{} สาวชฺชญฺจ กมฺมํ, อนวชฺชญฺจ กมฺมํ ปหาย ปชหิตฺวา วิโนเทตฺวา พฺยนฺติํ กริตฺวา อนภาวํ คเมตฺวาติ กมฺมญฺจ สาวชฺชนวชฺชเมตํ.}

\prose{\textbf{สุทฺธิํ อสุทฺธินฺติ อปตฺถยาโน}ติ \textbf{อสุทฺธินฺ}ติ อสุทฺธิํ ปตฺเถนฺติ.\csromanspacer{} อกุสเล ธมฺเม ปตฺเถนฺติ, \textbf{สุทฺธินฺ}ติ สุทฺธิํ ปตฺเถนฺติ.\csromanspacer{} ปญฺจ กามคุเณ ปตฺเถนฺติ, อสุทฺธิํ ปตฺเถนฺติ, อกุสเล ธมฺเม ปตฺเถนฺติ, ปญฺจ กามคุเณ ปตฺเถนฺติ, สุทฺธิํ ปตฺเถนฺติ, ทฺวาสฏฺฐิ ทิฏฺฐิคตานิ ปตฺเถนฺติ, อสุทฺธิํ ปตฺเถนฺติ, อกุสเล ธมฺเม ปตฺเถนฺติ, ปญฺจ กามคุเณ ปตฺเถนฺติ, ทฺวาสฏฺฐิ ทิฏฺฐิคตานิ ปตฺเถนฺติ, สุทฺธิํ ปตฺเถนฺติ.\csromanspacer{} เตธาตุเก กุสเล ธมฺเม ปตฺเถนฺติ, อสุทฺธิํ ปตฺเถนฺติ, อกุสเล ธมฺเม ปตฺเถนฺติ.\csromanspacer{} ปญฺจ กามคุเณ ปตฺเถนฺติ, ทฺวาสฏฺฐิ ทิฏฺฐิคตานิ ปตฺเถนฺติ.\csromanspacer{} เตธาตุเก กุสเล ธมฺเม ปตฺเถนฺติ, สุทฺธิํ ปตฺเถนฺติ, ปุถุชฺชนกลฺยาณกา\footnote{กลฺยาณปุถุชฺชนา (สฺยา) เอวมีทิเสสุ ฐาเนสุ.} นิยามาวกฺกนฺติํ\footnote{นิยามาวตฺตนฺติํ (ก)} ปตฺเถนฺติ, เสกฺขา อคฺคธมฺมํ อรหตฺตํ ปตฺเถนฺติ, อรหตฺเต ปตฺเต อรหา เนว อกุสเล ธมฺเม ปตฺเถติ, นปิ ปญฺจ กามคุเณ ปตฺเถติ, นปิ ทฺวาสฏฺฐิ ทิฏฺฐิคตานิ ปตฺเถติ, นปิ เตธาตุเก กุสเล ธมฺเม ปตฺเถติ, นปิ นิยามาวกฺกนฺติํ ปตฺเถติ, นปิ อคฺคธมฺมํ อรหตฺตํ ปตฺเถติ.\csromanspacer{} ปตฺถนา สมติกฺกนฺโต อรหา วุทฺธิปาริหานิวีติวตฺโต\footnote{วุทฺธิปาริหานิํ วีติวตฺโต (สี)}, โส วุฏฺฐวาโส จิณฺณจรโณ ฯเปฯ.\csromanspacer{} ชาติชรามรณสํสาโร, นตฺถิ ตสฺส ปุนพฺภโวติ สุทฺธิํ \textbf{อสุทฺธินฺ}ติ อปตฺถยาโน.}

\prose{\csromanfolio{244}\textbf{วิรโต จเร สนฺติมนุคฺคหายา}ติ \textbf{วิรโต}ติ สุทฺธิ-อสุทฺธิยา อารโต อสฺส \textbf{วิรโต} ปฏิ\textbf{วิรโต} นิกฺขนฺโต นิสฺสโฏ วิปฺปมุตฺโต วิสญฺญุตฺโต, วิมริยาทิกเตน เจตสา วิหรตีติ \textbf{วิรโต}.\csromanspacer{} \textbf{จเร}ติ จเรยฺย วิจเรยฺย วิหเรยฺย อิริเยยฺย วตฺเตยฺย ปาเลยฺย ยเปยฺย ยาเปยฺยาติ \textbf{วิรโต} จเร.\csromanspacer{} \textbf{สนฺติมนุคฺคหายา}ติ \textbf{สนฺติโย} วุจฺจนฺติ ทฺวาสฏฺฐิ ทิฏฺฐิคตานิ.\csromanspacer{} ทิฏฺฐิ\textbf{สนฺติโย} อคณฺหนฺโต อปรามสนฺโต อนภินิวิสนฺโตติ \textbf{วิรโต} จเร สนฺติมนุคฺคหาย.\csromanspacer{} เตนาห ภควา\csromandash{}}

\csromangathameasure{สีลพฺพตํ วาปิ ปหาย สพฺพํ, \\ กมฺมญฺจ สาวชฺชนวชฺชเมตํ. \\ สุทฺธิํ อสุทฺธินฺติ อปตฺถยาโน,}
\csromangathagroup{\csromangathastanza{สีลพฺพตํ วาปิ ปหาย สพฺพํ, \\ กมฺมญฺจ สาวชฺชนวชฺชเมตํ. \\ สุทฺธิํ อสุทฺธินฺติ อปตฺถยาโน,}}

\prose{\textbf{วิรโต จเร สนฺติมนุคฺคหายา}ติ.}

\proseitem{136}{\csromansymbolfootnote{*}{ขุ 1. 419 ปิฏฺเฐ.}\textbf{ตมูปนิสฺสาย ชิคุจฺฉิตํ วา},}

\prose{\textbf{อถ วาปิ ทิฏฺฐิํ ว สุตํ มุตํ วา.}}

\csromangathameasure{\textbf{อุทฺธํสรา สุทฺธิ’มนุตฺถุนนฺติ,} \\ \textbf{อวีตตณฺหาเส ภวาภเวสุ.}}
\csromangathagroup{\csromangathastanza{\textbf{อุทฺธํสรา สุทฺธิ’มนุตฺถุนนฺติ,} \\ \textbf{อวีตตณฺหาเส ภวาภเวสุ.}}}

\prose{\textbf{ตมูปนิสฺสาย ชิคุจฺฉิตํ วา}ติ สนฺเตเก สมณพฺราหฺมณา ตโปชิคุจฺฉวาทา ตโปชิคุจฺฉสารา ตโปชิคุจฺฉนิสฺสิตา อานิสฺสิตา อลฺลีนา อุปคตา อชฺโฌสิตา อธิมุตฺตาติ ตมูปนิสฺสาย ชิคุจฺฉิตํ วา.}

\prose{\textbf{อถ วาปิ ทิฏฺฐํ ว สุตํ มุตํ วา}ติ ทิฏฺฐํ วา ทิฏฺฐสุทฺธิํ วา สุตํ วา สุตสุทฺธิํ วา มุตํ วา มุตสุทฺธิํ วา นิสฺสาย อุปนิสฺสาย คณฺหิตฺวา ปรามสิตฺวา อภินิวิสิตฺวาติ อถ วาปิ ทิฏฺฐํ ว สุตํ มุตํ วา.}

\prose{\textbf{อุทฺธํสรา สุทฺธิมนุตฺถุนนฺตี}ติ สนฺเตเก สมณพฺราหฺมณา อุทฺธํสราวาทา.\csromanspacer{} กตเม เต สมณพฺราหฺมณา อุทฺธํสราวาทา, เย เต สมณพฺราหฺมณา อจฺจนฺตสุทฺธิกา สํสารสุทฺธิกา อกิริยทิฏฺฐิกา สสฺสตวาทา, อิเม เต สมณพฺราหฺมณา อุทฺธํสราวาทา, เต สํสาเร สุทฺธิํ วิสุทฺธิํ ปริสุทฺธิํ, มุตฺติํ วิมุตฺติํ ปริมุตฺติํ ถุนนฺติ วทนฺติ กเถนฺติ ภณนฺติ ทีปยนฺติ โวหรนฺตีติ อุทฺธํสรา สุทฺธิมนุตฺถุนนฺติ.}

\prose{\textbf{อวีตตณฺหาเส ภวาภเวสู}ติ \textbf{ตณฺหา}ติ รูป\textbf{ตณฺหา} สทฺท\textbf{ตณฺหา} คนฺธ\textbf{ตณฺหา} รส\textbf{ตณฺหา} โผฏฺฐพฺพ\textbf{ตณฺหา} ธมฺม\textbf{ตณฺหา}.\csromanspacer{} \textbf{ภวาภเวสูติ} ภวาภเว กมฺมภเว ปุนพฺภเว กามภเว, กมฺมภเว กามภเว ปุนพฺภเว}

\vspace{\tipitakaparskip}
\csromancenter{มหาวิยูหสุตฺตนิทฺเทส รูปภเว, กมฺมภเว รูปภเว ปุนพฺภเว อรูปภเว, กมฺมภเว อรูปภเว ปุนพฺภเว ปุนปฺปุนพฺภเว, ปุนปฺปุนคติยา ปุนปฺปุน- อุปปตฺติยา ปุนปฺปุนปฏิสนฺธิยา ปุนปฺปุน-อตฺตภาวาภินิพฺพตฺติยา อวีตตณฺหา อวิคตตณฺหา อจตฺตตณฺหา อวนฺตตณฺหา อมุตฺตตณฺหา อปฺปหีนตณฺหา อปฺปฏินิสฺสฏฺฐตณฺหาติ อวีตตณฺหาเส ภวาภเวสุ.\csromanspacer{} เตนาห ภควา\csromandash{}}
\vspace{0.5\baselineskip}
\csromangathameasure{ตมูปนิสฺสาย ชิคุจฺฉิตํ วา, \\ อถ วาปิ ทิฏฺฐิํ ว สุตํ มุตํ วา. \\ อุทฺธํสรา สุทฺธิมนุตฺถุนนฺติ, \\ อวีตตณฺหาเส ภวาภเวสูติ.}
\csromangathagroup{\csromangathastanza{\csromanfolio{245}ตมูปนิสฺสาย ชิคุจฺฉิตํ วา, \\ อถ วาปิ ทิฏฺฐิํ ว สุตํ มุตํ วา. \\ อุทฺธํสรา สุทฺธิมนุตฺถุนนฺติ, \\ อวีตตณฺหาเส ภวาภเวสูติ.}}

\proseitem{137}{\csromansymbolfootnote{*}{ขุ 1. 419 ปิฏฺเฐ.}ปตฺถยมานสฺส หิ ชปฺปิตานิ,}

\prose{\textbf{ปเวธิตํ วาปิ ปกปฺปิเตสุ.}}

\csromangathameasure{\textbf{จุตูปปาโต อิธ ยสฺส นตฺถิ,} \\ \textbf{ส เกน เวเธยฺย กุหิํ ว ชปฺเป.}}
\csromangathagroup{\csromangathastanza{\textbf{จุตูปปาโต อิธ ยสฺส นตฺถิ,} \\ \textbf{ส เกน เวเธยฺย กุหิํ ว ชปฺเป.}}}

\prose{\textbf{ปตฺถยมานสฺส หิ ชปฺปิตานี}ติ \textbf{ปตฺถนา} วุจฺจติ ตณฺหา, โย ราโค สาราโค ฯเปฯ อภิชฺฌา โลโภ อกุสลมูลํ.\csromanspacer{} \textbf{ปตฺถยมานสฺสา}ติ ปตฺถยมานสฺส อิจฺฉมานสฺส สาทิยมานสฺส ปิหยมานสฺส อภิชปฺปยมานสฺสาติ ปตฺถยมานสฺส หิ.\csromanspacer{} \textbf{ชปฺปิตานี}ติ \textbf{ชปฺปนา} วุจฺจติ ตณฺหา.\csromanspacer{} โย ราโค สาราโค ฯเปฯ อภิชฺฌา โลโภ อกุสลมูลนฺติ ปตฺถยมานสฺส หิ ชปฺปิตานิ.}

\prose{\textbf{ปเวธิตํ วาปิ ปกปฺปิเตสู}ติ \textbf{ปกปฺปนา}ติเทฺว \textbf{ปกปฺปนา} ตณฺหา\textbf{ปกปฺปนา} จ ทิฏฺฐิ\textbf{ปกปฺปนา} จ ฯเปฯ อยํ ตณฺหา\textbf{ปกปฺปนา} ฯเปฯ อยํ ทิฏฺฐิ\textbf{ปกปฺปนา}.\csromanspacer{} \textbf{ปเวธิตํ วาปิ ปกปฺปิเตสู}ติ ปกปฺปิตํ วตฺถุํ อจฺเฉทสงฺกิโนปิ เวเธนฺติ, อจฺฉิชฺชนฺเตปิ เวเธนฺติ, อจฺฉินฺเนปิ เวเธนฺติ.\csromanspacer{} ปกปฺปิตํ วตฺถุํ วิปริณามสงฺกิโนปิ เวเธนฺติ, วิปริณมนฺเตปิ เวเธนฺติ, วิปริณเตปิ เวเธนฺติ ปเวเธนฺติ สมฺปเวเธนฺตีติ ปเวธิตํ วาปิ ปกปฺปิเตสุ.}

\prose{\textbf{จุตูปปาโต อิธ ยสฺส นตฺถี}ติ \textbf{ยสฺสา}ติ อรหโต ขีณาสวสฺส.\csromanspacer{} ยสฺส คมนํ อาคมนํ คมนาคมนํ กาลํ คติ ภวาภโว จุติ จ อุปปตฺติ จ นิพฺพตฺติ จ เภโท จ ชาติ จ ชรามรณญฺจ นตฺถิ น สนฺติ น สํวิชฺชนฺติ นุปลพฺภนฺติ, ปหีนา สมุจฺฉินฺนา วูปสนฺตา ปฏิปสฺสทฺธา อภพฺพุปฺปตฺติกา ญาณคฺคินา ทฑฺฒาติ จุตูปปาโต อิธ ยสฺส นตฺถิ.}

\prose{\csromanfolio{246}\textbf{ส เกน เวเธยฺย กุหิํ ว ชปฺเป}ติ โส เกน ราเคน เวเธยฺย, เกน โทเสน เวเธยฺย, เกน โมเหน เวเธยฺย, เกน มาเนน เวเธยฺย, กาย ทิฏฺฐิยา เวเธยฺย, เกน อุทฺธจฺเจน เวเธยฺย, กาย วิจิกิจฺฉาย เวเธยฺย, เกหิ อนุสเยหิ เวเธยฺย “รตฺโต”ติ วา “ทุฏฺโฐ”ติ วา “มูโฬฺห”ติ วา “วินิพทฺโธ”ติ วา “ปรามฏฺโฐ”ติ วา “วิกฺเขปคโต”ติ วา “อนิฏฺฐงฺคโต”ติ วา “ถามคโต”ติ วา.\csromanspacer{} เต อภิสงฺขารา ปหีนา, อภิสงฺขารานํ ปหีนตฺตา คติยา เกน เวเธยฺย “เนรยิโก”ติ วา “ติรจฺฉนโยนิโก”ติ วา “เปตฺติวิสยิโก”ติ วา “มนุสฺโส”ติ วา “เทโว”ติ วา “รูปี”ติ วา “อรูปี”ติ วา “สญฺญี”ติ วา “อสญฺญี”ติ วา “เนวสญฺญีนาสญฺญี”ติ วา.\csromanspacer{} โส เหตุ นตฺถิ, ปจฺจโย นตฺถิ, การณํ นตฺถิ, เยน เวเธยฺย ปเวเธยฺย สมฺปเวเธยฺยาติ ส เกน เวเธยฺย.\csromanspacer{} \textbf{กุหิํ ว ชปฺเป}ติ กุหิํ วา ชปฺเปยฺย, กิมฺหิ ชปฺเปยฺย, กตฺถ ชปฺเปยฺย ปชปฺเปยฺย อภิชปฺเปยฺยาติ ส เกน เวเธยฺย กุหิํ ว ชปฺเป.\csromanspacer{} เตนาห ภควา\csromandash{}}

\csromangathameasure{ปตฺถยมานสฺส หิ ชปฺปิตานิ, \\ ปเวธิตํ วาปิ ปกปฺปิเตสุ. \\ จุตูปปาโต อิธ ยสฺส นตฺถิ,}
\csromangathagroup{\csromangathastanza{ปตฺถยมานสฺส หิ ชปฺปิตานิ, \\ ปเวธิตํ วาปิ ปกปฺปิเตสุ. \\ จุตูปปาโต อิธ ยสฺส นตฺถิ,}}

\prose{\textbf{ส เกน เวเธยฺย กุหิํ ว ชปฺเป}ติ.}

\proseitem{138}{\csromansymbolfootnote{*}{ขุ 1. 419 ปิฏฺเฐ.}\textbf{ยมาหุ ธมฺมํ ปรมนฺติ เอเก},}

\prose{\textbf{ตเมว หีนนฺติ ปนาหุ อญฺเญ.}}

\csromangathameasure{\textbf{สจฺโจ นุ วาโท กตโม อิเมสํ,} \\ \textbf{สพฺเพว หีเม กุสลาวทานา.}}
\csromangathagroup{\csromangathastanza{\textbf{สจฺโจ นุ วาโท กตโม อิเมสํ,} \\ \textbf{สพฺเพว หีเม กุสลาวทานา.}}}

\prose{\textbf{ยมาหุ ธมฺมํ ปรมนฺติ เอเก}ติ ยํ ธมฺมํ ทิฏฺฐิํ ปฏิปทํ มคฺคํ เอเก สมณพฺราหฺมณา “อิทํ ปรมํ อคฺคํ เสฏฺฐํ วิสิฏฺฐํ ปาโมกฺขํ อุตฺตมํ ปวรนฺ”ติ เอวมาหํสุ เอวํ กเถนฺติ เอวํ ภณนฺติ เอวํ ทีปยนฺติ เอวํ โวหรนฺตีติ ยมาหุ ธมฺมํ ปรมนฺติ เอเก.}

\prose{\textbf{ตเมว หีนนฺติ ปนาหุ อญฺเญ}ติ ตเมว ธมฺมํ ทิฏฺฐิํ ปฏิปทํ มคฺคํ เอเก สมณพฺราหฺมณา “หีนํ เอตํ, นิหีนํ เอตํ, โอมกํ เอตํ, ลามกํ เอตํ, ฉตุกฺกํ เอตํ, ปริตฺตกํ เอตนฺ”ติ เอวมาหํสุ เอวํ กเถนฺติ เอวํ ภณนฺติ เอวํ ทีปยนฺติ เอวํ โวหรนฺตีติ ตเมว หีนนฺติ ปนาหุ อญฺเญ.}

\csromanheader{2}{13. มหาวิยูหสุตฺตนิทฺเทส}
\prose{\csromanfolio{247}\textbf{สจฺโจ นุ วาโท กตโม อิเมสนฺ}ติ อิเมสํ สมณพฺราหฺมณานํ วาโท กตโม สจฺโจ ตจฺโฉ ตโถ ภูโต ยาถาโว อวิปรีโตติ สจฺโจ นุ วาโท กตโม อิเมสํ.}

\prose{\textbf{สพฺเพว หีเม กุสลาวทานา}ติ สพฺเพวิเม สมณพฺราหฺมณา กุสลวาทา ปณฺฑิตวาทา ถิรวาทา ญายวาทา เหตุวาทา ลกฺขณวาทา การณวาทา ฐานวาทา สกาย ลทฺธิยาติ สพฺเพว หีเม กุสลาวทานา.\csromanspacer{} เตนาห โส นิมฺมิโต\csromandash{}}

\csromangathameasure{ยมาหุ ธมฺมํ ปรมนฺติ เอเก, \\ ตเมว หีนนฺติ ปนาหุ อญฺเญ. \\ สจฺโจ นุ วาโท กตโม อิเมสํ, \\ \textbf{สพฺเพว หีเม กุสลาวทานา}ติ.}
\csromangathagroup{\csromangathastanza{ยมาหุ ธมฺมํ ปรมนฺติ เอเก, \\ ตเมว หีนนฺติ ปนาหุ อญฺเญ. \\ สจฺโจ นุ วาโท กตโม อิเมสํ, \\ \textbf{สพฺเพว หีเม กุสลาวทานา}ติ.}}

\proseitem{139}{\csromansymbolfootnote{*}{ขุ 1. 419 ปิฏฺเฐ.}สกํ หิ ธมฺมํ ปริปุณฺณมาหุ,}

\prose{\textbf{อญฺญสฺส ธมฺมํ ปน หีนมาหุ.}}

\csromangathameasure{\textbf{เอวมฺปิ วิคฺคยฺห วิวาทยนฺติ,} \\ \textbf{สกํ สกํ สมฺมุติมาหุ สจฺจํ.}}
\csromangathagroup{\csromangathastanza{\textbf{เอวมฺปิ วิคฺคยฺห วิวาทยนฺติ,} \\ \textbf{สกํ สกํ สมฺมุติมาหุ สจฺจํ.}}}

\prose{\textbf{สกํ หิ ธมฺมํ ปริปุณฺณมาหู}ติ สกํ ธมฺมํ ทิฏฺฐิํ ปฏิปทํ มคฺคํ เอเก สมณพฺราหฺมณา “อิทํ สมตฺตํ ปริปุณฺณํ อโนมนฺ”ติ เอวมาหํสุ ฯเปฯ เอวํ โวหรนฺตีติ สกํ หิ ธมฺมํ ปริปุณฺณมาหุ.}

\prose{\textbf{อญฺญสฺส ธมฺมํ ปน หีนมาหู}ติ อญฺญสฺส ธมฺมํ ทิฏฺฐิํ ปฏิปทํ มคฺคํ เอเก สมณพฺราหฺมณา “หีนํ เอตํ, นิหีนํ เอตํ, โอมกํ เอตํ, ลามกํ เอตํ, ฉตุกฺกํ เอตํ, ปริตฺตกํ เอตนฺ”ติ เอวมาหํสุ เอวํ กเถนฺติ เอวํ ภณนฺติ เอวํ ทีปยนฺติ เอวํ โวหรนฺตีติ อญฺญสฺส ธมฺมํ ปน หีนมาหุ.}

\prose{\textbf{เอวมฺปิ วิคฺคยฺห วิวาทยนฺตี}ติ เอวํ คเหตฺวา อุคฺคเหตฺวา คณฺหิตฺวา ปรามสิตฺวา อภินิวิสิตฺวา วิวาทยนฺติ, กลหํ กโรนฺติ, ภณฺฑนํ กโรนฺติ, วิคฺคหํ กโรนฺติ, วิวาทํ กโรนฺติ, เมธคํ กโรนฺติ “น ตฺวํ อิมํ ธมฺมวินยํ อาชานาสิ ฯเปฯ นิพฺเพเฐหิ วา สเจ ปโหสี”ติ เอวมฺปิ วิคฺคยฺห วิวาทยนฺติ.}

\prose{\textbf{สกํ สกํ สมฺมุติมาหุ สจฺจนฺ}ติ “สสฺสโต โลโก, อิทเมว สจฺจํ โมฆมญฺญนฺ”ติ สกํ สกํ สมฺมุติมาหุ สจฺจํ. “อสสฺสโต \csromanfolio{248}โลโก ฯเปฯ เนว โหติ น น โหติ ตถาคโต ปรํ มรณา, อิทเมว สจฺจํ โมฆมญฺญนฺ”ติ สกํ สกํ สมฺมุติมาหุ สจฺจํ.\csromanspacer{} เตนาห ภควา\csromandash{}}

\csromangathameasure{สกํ หิ ธมฺมํ ปริปุณฺณมาหุ, \\ อญฺญสฺส ธมฺมํ ปน หีนมาหุ. \\ เอวมฺปิ วิคฺคยฺห วิวาทยนฺติ, \\ สกํ สกํ สมฺมุติมาหุ สจฺจนฺติ.}
\csromangathagroup{\csromangathastanza{สกํ หิ ธมฺมํ ปริปุณฺณมาหุ, \\ อญฺญสฺส ธมฺมํ ปน หีนมาหุ. \\ เอวมฺปิ วิคฺคยฺห วิวาทยนฺติ, \\ สกํ สกํ สมฺมุติมาหุ สจฺจนฺติ.}}

\proseitem{140}{\csromansymbolfootnote{*}{ขุ 1. 419 ปิฏฺเฐ.}\textbf{ปรสฺส เจ วมฺภยิเตน หีโน},}

\prose{\textbf{น โกจิ ธมฺเมสุ วิเสสิ อสฺส.}}

\csromangathameasure{\textbf{ปุถู หิ อญฺญสฺส วทนฺติ ธมฺมํ,} \\ \textbf{นิหีนโต สมฺหิ ทฬฺหํ วทานา.}}
\csromangathagroup{\csromangathastanza{\textbf{ปุถู หิ อญฺญสฺส วทนฺติ ธมฺมํ,} \\ \textbf{นิหีนโต สมฺหิ ทฬฺหํ วทานา.}}}

\prose{\textbf{ปรสฺส เจ วมฺภยิเตน หีโน}ติ ปรสฺส เจ วมฺภยิตการณา นินฺทิตการณา ครหิตการณา อุปวทิตการณา ปโร พาโล โหติ หีโน นิหีโน โอมโก ลามโก ฉตุกฺโก ปริตฺโตติ ปรสฺส เจ วมฺภยิเตน หีโน.}

\prose{\textbf{น โกจิ ธมฺเมสุ วิเสสิ อสฺสา}ติ ธมฺเมสุ น โกจิ อคฺโค เสฏฺโฐ วิสิฏฺโฐ ปาโมกฺโข อุตฺตโม ปวโร อสฺสาติ น โกจิ ธมฺเมสุ วิเสสิ อสฺส.}

\prose{\textbf{ปุถู หิ อญฺญสฺส วทนฺติ ธมฺมํ นิหีนโต}ติ พหุกาปิ พหูนํ ธมฺมํ วทนฺติ อุปวทนฺติ นินฺทนฺติ ครหนฺติ หีนโต นิหีนโต โอมกโต ลามกโต ฉตุกฺกโต ปริตฺตโต, พหุกาปิ เอกสฺส ธมฺมํ วทนฺติ อุปวทนฺติ นินฺทนฺติ ครหนฺติ หีนโต นิหีนโต โอมกโต ลามกโต ฉตุกฺกโต ปริตฺตโต, เอโกปิ พหูนํ ธมฺมํ วทติ อุปวทติ นินฺทติ ครหติ หีนโต นิหีนโต โอมกโต ลามกโต ฉตุกฺกโต ปริตฺตโต, เอโกปิ เอกสฺส ธมฺมํ วทติ อุปวทติ นินฺทติ ครหติ หีนโต นิหีนโต โอมกโต ลามกโต ฉตุกฺกโต ปริตฺตโตติ ปุถู หิ อญฺญสฺส วทนฺติ ธมฺมํ.}

\prose{\textbf{นิหีนโต สมฺหิ ทฬฺหํ วทานา}ติ ธมฺโม สกายนํ ทิฏฺฐิ สกายนํ ปฏิปทา สกายนํ มคฺโค สกายนํ, สกายเน ทฬฺหวาทา ถิรวาทา}

\vspace{\tipitakaparskip}
\csromancenter{มหาวิยูหสุตฺตนิทฺเทส พลิกวาทา อฏฺฐิตวาทาติ นิหีนโต สมฺหิ ทฬฺหํ วทานา.\csromanspacer{} เตนาห ภควา\csromandash{}}
\vspace{0.5\baselineskip}
\csromangathameasure{ปรสฺส เจ วมฺภยิเตน หีโน, \\ น โกจิ ธมฺเมสุ วิเสสิ อสฺส. \\ ปุถู หิ อญฺญสฺส วทนฺติ ธมฺมํ, \\ นิหีนโต สมฺหิ ทฬฺหํ วทานาติ.}
\csromangathagroup{\csromangathastanza{\csromanfolio{249}ปรสฺส เจ วมฺภยิเตน หีโน, \\ น โกจิ ธมฺเมสุ วิเสสิ อสฺส. \\ ปุถู หิ อญฺญสฺส วทนฺติ ธมฺมํ, \\ นิหีนโต สมฺหิ ทฬฺหํ วทานาติ.}}

\csromanhangingbegin
\hangingheaditem{141}{สทฺธมฺมปูชาปิ\footnote{สทฺธมฺมปูชา จ (สี, สฺยา)} \textbf{เนสํ ตเถว,}}
\hangingbody{\textbf{ยถา ปสํสนฺติ สกายนานิ.}}
\hangingbody{\textbf{สพฺเพว วาทา}2 \textbf{ตถิยา3} \textbf{ภเวยฺยุํ,}}
\hangingbody{\textbf{สุทฺธี หิ เนสํ ปจฺจตฺตเมว.}}
\csromanhangingend

\prose{\textbf{สทฺธมฺมปูชาปิ เนสํ ตเถวา}ติ กตมา สทฺธมฺมปูชา, สกํ สตฺถารํ สกฺกโรติ ครุํ กโรติ มาเนติ ปูเชติ “อยํ สตฺถา สพฺพญฺญู”ติ, อยํ สทฺธมฺมปูชา.\csromanspacer{} สกํ ธมฺมกฺขานํ สกํ คณํ สกํ ทิฏฺฐิํ สกํ ปฏิปทํ สกํ มคฺคํ สกฺกโรติ ครุํ กโรติ มาเนติ ปูเชติ “อยํ มคฺโค นิยฺยานิโก”ติ, อยํ สทฺธมฺมปูชา.\csromanspacer{} \textbf{สทฺธมฺมปูชาปิ เนสํ ตเถวา}ติ สทฺธมฺมปูชา ตถา ตจฺฉา ภูตา ยถาวา อวิปรีตาติ สทฺธมฺมปูชาปิ เนสํ ตเถว.}

\prose{\textbf{ยถา ปสํสนฺติ สกายนานี}ติ ธมฺโม สกายนํ ทิฏฺฐิ สกายนํ ปฏิปทา สกายนํ มคฺโค สกายนํ, สกายนานิ ปสํสนฺติ โถเมนฺติ กิตฺเตนฺติ วณฺเณนฺตีติ ยถา ปสํสนฺติ สกายนานิ.}

\prose{\textbf{สพฺเพว วาทา ตถิยา ภเวยฺยุนฺ}ติ สพฺเพว วาทา ตถา ตจฺฉา ภูตา ยาถาวา อวิปรีตา ภเวยฺยุนฺติ สพฺเพว วาทา \textbf{ตถิยา} ภเวยฺยุํ.}

\prose{\textbf{สุทฺธี หิ เนสํ ปจฺจตฺตเมวา}ติ ปจฺจตฺตเมว เตสํ สมณพฺราหฺมณานํ สุทฺธิ วิสุทฺธิ ปริสุทฺธิ มุตฺติ วิมุตฺติ ปริมุตฺตีติ สุทฺธี หิ เนสํ ปจฺจตฺตเมว.\csromanspacer{} เตนาห ภควา\csromandash{}}

\prose{สทฺธมฺมปูชาปิ \textbf{เนสํ ตเถว,}}

\prose{\textbf{ยถา ปสํสนฺติ สกายนานิ.}}

\prose{\textbf{สพฺเพว วาทา} \textbf{ตถิยา} \textbf{ภเวยฺยุํ,}}

\prose{\textbf{สุทฺธี หิ เนสํ ปจฺจตฺตเมวา}ติ.}

\proseitem{142}{\csromanfolio{250}\csromansymbolfootnote{*}{ขุ 1. 420 ปิฏฺเฐ.}น พฺราหฺมณสฺส ปรเนยฺยมตฺถิ,}

\prose{\textbf{ธมฺเมสุ นิจฺเฉยฺย สมุคฺคหีตํ.}}

\csromangathameasure{\textbf{ตสฺมา วิวาทานิ อุปาติวตฺโต,} \\ \textbf{น หิ เสฏฺฐโต ปสฺสติ ธมฺมมญฺญํ.}}
\csromangathagroup{\csromangathastanza{\textbf{ตสฺมา วิวาทานิ อุปาติวตฺโต,} \\ \textbf{น หิ เสฏฺฐโต ปสฺสติ ธมฺมมญฺญํ.}}}

\prose{\textbf{น พฺราหฺมณสฺส ปรเนยฺยมตฺถี}ติ \textbf{นา}ติ ปฏิกฺเขโป.\csromanspacer{} \textbf{พฺราหฺมโณ}ติ สตฺตนฺนํ ธมฺมานํ พาหิตตฺตา พฺราหฺมโณ ฯเปฯ อสิโต ตาทิ ปวุจฺจเต ส พฺรหฺมา.\csromanspacer{}\csromanspacer{}\textbf{น พฺราหฺมณสฺส ปรเนยฺยมตฺถี}ติ พฺราหฺมณสฺส ปรเนยฺยตา นตฺถิ, พฺราหฺมโณ น ปรเนยฺโย, น ปรปตฺติโย, น ปรปจฺจโย, น ปรปฏิพทฺธคู\footnote{ปรปฏิพนฺธคู (ก)} ชา\textbf{นา}ติ ปสฺสติ อสมฺมูโฬฺห สมฺปชาโน ปฏิสฺสโต. “สพฺเพ สงฺขารา อนิจฺจา”ติ พฺราหฺมณสฺส ปรเนยฺยตา นตฺถิ, พฺราหฺมโณ น ปรเนยฺโย, น ปรปตฺติโย, น ปรปจฺจโย, น ปรปฏิพทฺธคู ชา\textbf{นา}ติ ปสฺสติ อสมฺมูโฬฺห สมฺปชาโน ปฏิสฺสโต. “สพฺเพ สงฺขารา ทุกฺขา”ติ ฯเปฯ “ยํ กิญฺจิ สมุทยธมฺมํ, สพฺพํ ตํ นิโรธธมฺมนฺ”ติ พฺราหฺมณสฺส ปรเนยฺยตา นตฺถิ, พฺราหฺมโณ น ปรเนยฺโย, น ปรปตฺติโย, น ปรปจฺจโย, น ปรปฏิพทฺธคู ชา\textbf{นา}ติ ปสฺสติ อสมฺมูโฬฺห สมฺปชาโน ปฏิสฺสโตติ น พฺราหฺมณสฺส ปรเนยฺยมตฺถิ.}

\prose{\textbf{ธมฺเมสุ นิจฺเฉยฺย สมุคฺคหีต}นฺติ \textbf{ธมฺเมสู}ติ ทฺวาสฏฺฐิยา ทิฏฺฐิคเตสุ.\csromanspacer{} \textbf{นิจฺเฉยฺยา}ติ นิจฺฉินิตฺวา วินิจฺฉินิตฺวา วิจินิตฺวา ปวิจินิตฺวา ตุลยิตฺวา ตีรยิตฺวา วิภาวยิตฺวา วิภูตํ กตฺวา.\csromanspacer{} \textbf{สมุคฺคหีต}นฺติ โอธิคฺคาโห พิลคฺคาโห วรคฺคาโห โกฏฺฐาสคฺคาโห อุจฺจยคฺคาโห สมุจฺจยคฺคาโห อิทํ สจฺจํ ตถํ ภูตํ ยาถาวํ อวิปรีตนฺติ คหิตํ ปรามฏฺฐํ อภินิวิฏฺฐํ อชฺโฌสิตํ อธิมุตฺตํ นตฺถิ น สนฺติ น สํวิชฺชติ นุปลพฺภติ, ปหีนํ สมุจฺฉินฺนํ วูปสนฺตํ ปฏิปสฺสทฺธํ อภพฺพุปฺปตฺติกํ ญาณคฺคิ\textbf{นา} ทฑฺฒนฺติ ธมฺเมสุ นิจฺเฉยฺย สมุคฺคหีตํ.}

\prose{\textbf{ตสฺมา วิวาทานิ อุปาติวตฺโต}ติ \textbf{ตสฺมา}ติ \textbf{ตสฺมา} ตํการณา ตํเหตุ ตปฺปจฺจยา ตํนิทา\textbf{นา} ทิฏฺฐิกลหานิ ทิฏฺฐิภณฺฑ\textbf{นา}นิ ทิฏฺฐิวิคฺคหานิ ทิฏฺฐิวิวาทานิ ทิฏฺฐิเมธคานิ อุปาติวตฺโต อติกฺกนฺโต สมติกฺกนฺโต วีติวตฺโตติ \textbf{ตสฺมา} วิวาทานิ อุปาติวตฺโต.}

\prose{\textbf{น หิ เสฏฺฐโต ปสฺสติ ธมฺมมญฺญนฺ}ติ อญฺญํ สตฺถารํ ธมฺมกฺขานํ คณํ ทิฏฺฐิํ ปฏิปทํ มคฺคํ อญฺญตฺร สติปฏฺฐาเนหิ อญฺญตฺร สมฺมปฺปธาเนหิ อญฺญตฺร}

\vspace{\tipitakaparskip}
\csromancenter{มหาวิยูหสุตฺตนิทฺเทส อิทฺธิปาเทหิ อญฺญตฺร อินฺทฺริเยหิ อญฺญตฺร พเลหิ อญฺญตฺร โพชฺฌงฺเคหิ อญฺญตฺร อริยา อฏฺฐงฺคิกา มคฺคา อคฺคํ เสฏฺฐํ วิสิฏฺฐํ ปาโมกฺขํ อุตฺตมํ ปวรํ ธมฺมํ น ปสฺสติ น ทกฺขติ น โอโลเกติ น นิชฺฌายติ น อุปปริกฺขตีติ น หิ เสฏฺฐโต ปสฺสติ ธมฺมมญฺญํ.\csromanspacer{} เตนาห ภควา\csromandash{}}
\vspace{0.5\baselineskip}
\csromangathameasure{น พฺราหฺมณสฺส ปรเนยฺยมตฺถิ, \\ ธมฺเมสุ นิจฺเฉยฺย สมุคฺคหีตํ. \\ ตสฺมา วิวาทานิ อุปาติวตฺโต, \\ น หิ เสฏฺฐโต ปสฺสติ ธมฺมมญฺญนฺติ.}
\csromangathagroup{\csromangathastanza{\csromanfolio{251}น พฺราหฺมณสฺส ปรเนยฺยมตฺถิ, \\ ธมฺเมสุ นิจฺเฉยฺย สมุคฺคหีตํ. \\ ตสฺมา วิวาทานิ อุปาติวตฺโต, \\ น หิ เสฏฺฐโต ปสฺสติ ธมฺมมญฺญนฺติ.}}

\proseitem{143}{\csromansymbolfootnote{*}{ขุ 1. 420 ปิฏฺเฐ.}\textbf{ชานามิ ปสฺสามิ ตเถว เอต}ํ,}

\prose{\textbf{ทิฏฺฐิยา เอเก ปจฺเจนฺติ สุทฺธิํ.}}

\csromangathameasure{\textbf{อทกฺขิ เจ กิญฺหิ ตุมสฺส เตน,} \\ \textbf{อติสิตฺวา อญฺเญน วทนฺติ สุทฺธิํ.}}
\csromangathagroup{\csromangathastanza{\textbf{อทกฺขิ เจ กิญฺหิ ตุมสฺส เตน,} \\ \textbf{อติสิตฺวา อญฺเญน วทนฺติ สุทฺธิํ.}}}

\prose{\textbf{ชานามิ ปสฺสามิ ตเถว เอต}นฺติ \textbf{ชานามี}ติ ปรจิตฺตญาเณน\footnote{ปรจิตฺตวิชานนญาเณน (สี) อฏฺฐกถา โอโลเกตพฺพา.} วา ชานามิ, ปุพฺเพนิวาสานุสฺสติญาเณน วา ชานามิ.\csromanspacer{} \textbf{ปสฺสามี}ติ มํสจกฺขุนา วา ปสฺสามิ, ทิพฺเพน จกฺขุนา วา ปสฺสามิ.\csromanspacer{} \textbf{ตเถว เอตนฺ}ติ เอตํ ตถํ ตจฺฉํ ภูตํ ยาถาวํ อวิปรีตนฺติ ชานามิ ปสฺสามิ ตเถว เอตํ.}

\prose{\textbf{ทิฏฺฐิยา เอเก ปจฺเจนฺติ สุทฺธินฺ}ติ ทิฏฺฐิยา เอเก สมณพฺราหฺมณา สุทฺธิํ วิสุทฺธิํ ปริสุทฺธิํ, มุตฺติํ วิมุตฺติํ ปริมุตฺติํ ปจฺเจนฺติ, “สสฺสโต โลโก อิทเมว สจฺจํ โมฆมญฺญนฺ”ติ ทิฏฺฐิยา เอเก สมณพฺราหฺมณา สุทฺธิํ วิสุทฺธิํ ปริสุทฺธิํ, มุตฺติํ วิมุตฺติํ ปริมุตฺติํ ปจฺเจนฺติ. “อสสฺสโต โลโก ฯเปฯ เนว โหติ น น โหติ ตถาคโต ปรํ มรณา, อิทเมว สจฺจํ โมฆมญฺญนฺ”ติ ทิฏฺฐิยา เอเก สมณพฺราหฺมณา สุทฺธิํ วิสุทฺธิํ ปริสุทฺธิํ, มุตฺติํ วิมุตฺติํ ปริมุตฺติํ ปจฺเจนฺตีติ ทิฏฺฐิยา เอเก ปจฺเจนฺติ สุทฺธิํ.}

\prose{\textbf{อทกฺขิ เจ กิญฺหิ ตุมสฺส เตนา}ติ \textbf{อทกฺขี}ติ ปรจิตฺตญาเณน วา อทกฺขิ, ปุพฺเพนิวาสานุสฺสติญาเณน วา อทกฺขิ, มํสจกฺขุนา วา อทกฺขิ, ทิพฺเพน จกฺขุนา วา \textbf{อทกฺขี}ติ อทกฺขิ เจ.\csromanspacer{} \textbf{กิญฺหิ ตุมสฺส เตนา}ติ ตสฺส เตน ทสฺสเนน กิํ กตํ, น ทุกฺขปริญฺญา อตฺถิ, น สมุทยสฺส ปหานํ อตฺถิ, น มคฺคภาวนา อตฺถิ, น ผลสจฺฉิกิริยา อตฺถิ, น ราคสฺส สมุจฺเฉทปหานํ อตฺถิ, น โทสสฺส สมุจฺเฉทปหานํ อตฺถิ, น โมหสฺส สมุจฺเฉทปหานํ \csromanfolio{252}อตฺถิ, น กิเลสานํ สมุจฺเฉทปหานํ อตฺถิ, น สํสารวฏฺฏสฺส อุปจฺเฉโท อตฺถีติ อทกฺขิ เจ กิญฺหิ ตุมสฺส เตน.}

\prose{\textbf{อติสิตฺวา อญฺเญน วทนฺติ สุทฺธินฺ}ติ เต ติตฺถิยา สุทฺธิมคฺคํ วิสุทฺธิมคฺคํ ปริสุทฺธิมคฺคํ โวทาตมคฺคํ ปริโวทาตมคฺคํ อติกฺกมิตฺวา สมติกฺกมิตฺวา วีติวตฺติตฺวา อญฺญตฺร สติปฏฺฐาเนหิ อญฺญตฺร สมฺมปฺปธาเนหิ อญฺญตฺร อิทฺธิปาเทหิ อญฺญตฺร อินฺทฺริเยหิ อญฺญตฺร พเลหิ อญฺญตฺร โพชฺฌงฺเคหิ อญฺญตฺร อริยา อฏฺฐงฺคิกา มคฺคา สุทฺธิํ วิสุทฺธิํ ปริสุทฺธิํ มุตฺติํ วิมุตฺติํ ปริมุตฺติํ วทนฺติ กเถนฺติ ภณนฺติ ทีปยนฺติ โวหรนฺตีติ เอวมฺปิ อติสิตฺวา อญฺเญน วทนฺติ สุทฺธิํ.}

\prose{อถ วา พุทฺธา จ พุทฺธสาวกา จ ปจฺเจกพุทฺธา จ เตสํ ติตฺถิยานํ อสุทฺธิมคฺคํ อวิสุทฺธิมคฺคํ อปริสุทฺธิมคฺคํ อโวทาตมคฺคํ อปริโวทาตมคฺคํ อติกฺกมิตฺวา สมติกฺกมิตฺวา วีติวตฺติตฺวา จตูหิ สติปฏฺฐาเนหิ จตูหิ สมฺมปฺปธาเนหิ จตูหิ อิทฺธิปาเทหิ ปญฺจหิ อินฺทฺริเยหิ ปญฺจหิ พเลหิ สตฺตหิ โพชฺฌงฺเคหิ อริเยน อฏฺฐงฺคิเกน มคฺเคน สุทฺธิํ วิสุทฺธิํ ปริสุทฺธิํ, มุตฺติํ วิมุตฺติํ ปริมุตฺติํ วทนฺติ กเถนฺติ ภณนฺติ ทีปยนฺติ โวหรนฺตีติ เอวมฺปิ อติสิตฺวา อญฺเญน วทนฺติ สุทฺธิํ.\csromanspacer{} เตนาห ภควา\csromandash{}}

\csromangathameasure{ชานามิ ปสฺสามิ ตเถว เอตํ, \\ ทิฏฺฐิยา เอเก ปจฺเจนฺติ สุทฺธิํ. \\ อทกฺขิ เจ กิญฺหิ ตุมสฺส เตน,}
\csromangathagroup{\csromangathastanza{ชานามิ ปสฺสามิ ตเถว เอตํ, \\ ทิฏฺฐิยา เอเก ปจฺเจนฺติ สุทฺธิํ. \\ อทกฺขิ เจ กิญฺหิ ตุมสฺส เตน,}}

\prose{\textbf{อติสิตฺวา อญฺเญน วทนฺติ สุทฺธินฺ}ติ.}

\proseitem{144}{\csromansymbolfootnote{*}{ขุ 1. 420 ปิฏฺเฐ.}ปสฺสํ นโร ทกฺขติ นามรูปํ,}

\prose{\textbf{ทิสฺวาน วา ญสฺสติ}\footnote{ญายติ (ก)} \textbf{ตานิเมว.}}

\csromangathameasure{\textbf{กามํ พหุํ ปสฺสตุ อปฺปกํ วา,} \\ \textbf{น หิ เตน สุทฺธิํ กุสลา วทนฺติ.}}
\csromangathagroup{\csromangathastanza{\textbf{กามํ พหุํ ปสฺสตุ อปฺปกํ วา,} \\ \textbf{น หิ เตน สุทฺธิํ กุสลา วทนฺติ.}}}

\prose{\textbf{ปสฺสํ นโร ทกฺขติ นามรูปนฺ}ติ ปสฺสํ นโร ทกฺขติ ปรจิตฺตญาเณน วา ปสฺสนฺโต ปุพฺเพนิวาสานุสฺสติญาเณน วา ปสฺสนฺโต มํสจกฺขุนา วา ปสฺสนฺโต ทิพฺเพน จกฺขุนา วา ปสฺสนฺโต นามรูปํเยว ทกฺขติ นิจฺจโต สุขโต อตฺตโต, น เตสํ ธมฺมานํ สมุทยํ วา อตฺถงฺคมํ วา อสฺสาทํ วา อาทีนวํ วา นิสฺสรณํ วา ทกฺขตีติ ปสฺสํ นโร ทกฺขติ นามรูปํ.}

\csromanheader{2}{13. มหาวิยูหสุตฺตนิทฺเทส}
\prose{\csromanfolio{253}\textbf{ทิสฺวาน วา ญสฺสติ ตานิเมวา}ติ \textbf{ทิสฺวา}ติ ปรจิตฺตญาเณน วา \textbf{ทิสฺวา} ปุพฺเพนิวาสานุสฺสติญาเณน วา \textbf{ทิสฺวา} มํสจกฺขุนา วา \textbf{ทิสฺวา} ทิพฺเพน จกฺขุนา วา \textbf{ทิสฺวา} นามรูปํเยว \textbf{ทิสฺวา} ญสฺสติ นิจฺจโต สุขโต อตฺตโต, น เตสํ ธมฺมานํ สมุทยํ วา อตฺถงฺคมํ วา อสฺสาทํ วา อาทีนวํ วา นิสฺสรณํ วา ญสฺสตีติ \textbf{ทิสฺวา}น วา ญสฺสติ ตานิเมว.}

\prose{\textbf{กามํ พหุํ ปสฺสตุ อปฺปกํ วา}ติ กามํ พหุกํ วา ปสฺสนฺโต นามรูปํ อปฺปกํ วา นิจฺจโต สุขโต อตฺตโตติ กามํ พหุํ ปสฺสตุ อปฺปกํ วา.}

\prose{\textbf{น หิ เตน สุทฺธิํ กุสลา วทนฺตี}ติ \textbf{กุสลา}ติ เย เต ขนฺธ\textbf{กุสลา} ธาตุ\textbf{กุสลา} อายตน\textbf{กุสลา} ปฏิจฺจสมุปฺปาท\textbf{กุสลา} สติปฏฺฐาน\textbf{กุสลา} สมฺมปฺปธาน\textbf{กุสลา} อิทฺธิปาท\textbf{กุสลา} อินฺทฺริย\textbf{กุสลา} พล\textbf{กุสลา} โพชฺฌงฺค\textbf{กุสลา} มคฺค\textbf{กุสลา} ผล\textbf{กุสลา} นิพฺพาน\textbf{กุสลา}, เต \textbf{กุสลา} ปรจิตฺตญาเณน วา ปุพฺเพนิวาสานุสฺสติญาเณน วา มํสจกฺขุนา วา ทิพฺเพน จกฺขุนา วา นามรูปทสฺสเนน สุทฺธิํ วิสุทฺธิํ ปริสุทฺธิํ มุตฺติํ วิมุตฺติํ ปริมุตฺติํ น วทนฺติ น กเถนฺติ น ภณนฺติ น ทีปยนฺติ น โวหรนฺตีติ น หิ เตน สุทฺธิํ \textbf{กุสลา} วทนฺติ.\csromanspacer{} เตนาห ภควา\csromandash{}}

\csromangathameasure{ปสฺสํ นโร ทกฺขติ นามรูปํ, \\ ทิสฺวาน วา ญสฺสติ ตานิเมว.}
\csromangathagroup{\csromangathastanza{ปสฺสํ นโร ทกฺขติ นามรูปํ, \\ ทิสฺวาน วา ญสฺสติ ตานิเมว.}}

\prose{\textbf{กามํ พหุํ ปสฺสตุ อปฺปกํ วา},}

\prose{\textbf{น หิ เตน สุทฺธิํ กุสลา วทนฺตี}ติ.}

\proseitem{145}{\csromansymbolfootnote{*}{ขุ 1. 420 ปิฏฺเฐ.}นิวิสฺสวาที น หิ สุพฺพินาโล,}

\prose{\textbf{ปกปฺปิตา ทิฏฺฐิปุรกฺขราโน.}}

\csromangathameasure{\textbf{ยํ นิสฺสิโต ตตฺถ สุภํ วทาโน,} \\ \textbf{สุทฺธิํ วโท ตตฺถ ตถทฺทสา โส.}}
\csromangathagroup{\csromangathastanza{\textbf{ยํ นิสฺสิโต ตตฺถ สุภํ วทาโน,} \\ \textbf{สุทฺธิํ วโท ตตฺถ ตถทฺทสา โส.}}}

\prose{\textbf{นิวิสฺสวาที น หิ สุพฺพินาโย}ติ “สสฺสโต โลโก, อิทเมว สจฺจํ โมฆมญฺญนฺ”ติ นิวิสฺสวาที, “อสสฺสโต โลโก ฯเปฯ เนว โหติ น น โหติ ตถาคโต ปรํ มรณา, อิทเมว สจฺจํ โมฆมญฺญนฺ”ติ นิวิสฺสวาที.\csromanspacer{} \textbf{น หิ สุพฺพินาโย}ติ นิวิสฺสวาที ทุพฺพินโย ทุปฺปญฺญาปโย\footnote{ทุปฺปญฺญาปิโย (สี) เอวมีทิเสสุ ฐาเนสุ.} \csromanfolio{254}ทุนฺนิชฺฌาปโย ทุปฺเปกฺขาปโย ทุปฺปสาทโยติ นิวิสฺสวาที น หิ สุพฺพิ\textbf{นา}โย.}

\prose{\textbf{ปกปฺปิตา ทิฏฺฐิปุรกฺขราโน}ติ กปฺปิตา ปกปฺปิตา อภิสงฺขตา สณฺฐปิตา ทิฏฺฐิํ ปุเรกฺขโต กตฺวา จรติ.\csromanspacer{} ทิฏฺฐิธโช ทิฏฺฐิเกตุ ทิฏฺฐาธิปเตยฺโย ทิฏฺฐิยา ปริวาริโต จรตีติ ปกปฺปิตา ทิฏฺฐิปุรกฺขราโน.}

\prose{\textbf{ยํ นิสฺสิโต ตตฺถ สุภํ วทาโน}ติ \textbf{ยํ นิสฺสิโต}ติ ยํ สตฺถารํ ธมฺมกฺขานํ คณํ ทิฏฺฐิํ ปฏิปทํ มคฺคํ นิสฺสิโต อานิสฺสิโต อลฺลีโน อุปคโต อชฺโฌสิโต อธิมุตฺโตติ \textbf{ยํ นิสฺสิโต}.\csromanspacer{} \textbf{ตตฺถา}ติ สกาย ทิฏฺฐิยา สกาย ขนฺติยา สกาย รุจิยา สกาย ลทฺธิยา.\csromanspacer{} \textbf{สุภํ วทาโน}ติ สุภวาโท โสภนวาโท ปณฺฑิตวาโท ถิรวาโท ญายวาโท เหตุวาโท ลกฺขณวาโท การณวาโท ฐานวาโท สกาย ลทฺธิยาติ \textbf{ยํ นิสฺสิโต} ตตฺถ สุภํ วทาโน.}

\prose{\textbf{สุทฺธิํวโท ตตฺถ ตถทฺทสา โส}ติ สุทฺธิวาโท วิสุทฺธิวาโท ปริสุทฺธิวาโท โวทตวาโท ปริโวทาตวาโท.\csromanspacer{} อถ วา สุทฺธิทสฺสโน วิสุทฺธิทสฺสโน ปริสุทฺธิทสฺสโน โวทาตทสฺสโน ปริโวทาตทสฺสโนติ สุทฺธิํ วาโท.\csromanspacer{} \textbf{ตตฺถา}ติ สกาย ทิฏฺฐิยา สกาย ขนฺติยา สกาย รุจิยา สกาย ลทฺธิยา ตถํ ตจฺฉํ ภูตํ ยาถาวํ อวิปรีตนฺติ อทฺทสฺส อทกฺขิ อปสฺสิ ปฏิวิชฺฌีติ สุทฺธิํ วาโท ตตฺถ ตถทฺทสา โส.\csromanspacer{} เต\textbf{นา}ห ภควา\csromandash{}}

\csromangathameasure{นิวิสฺสวาที น หิ สุพฺพิ\textbf{นา}โย,}
\csromangathagroup{\csromangathastanza{นิวิสฺสวาที น หิ สุพฺพิ\textbf{นา}โย,}}

\prose{\textbf{ปกปฺปิตา ทิฏฺฐิปุรกฺขราโน}.}

\prose{\textbf{ยํ นิสฺสิโต ตตฺถ สุภํ วทาโน},}

\prose{\textbf{สุทฺธิํ วโท ตตฺถ ตถทฺทสา โส}ติ.}

\proseitem{146}{\csromansymbolfootnote{*}{ขุ 1. 420 ปิฏฺเฐ.}\textbf{น พฺราหฺมโณ กปฺปมุเปติ สงฺขา},}

\prose{\textbf{น ทิฏฺฐิสารี นปิ ญาณพนฺธุ.}}

\csromangathameasure{\textbf{ญตฺวา จ โส สมฺมุติโย ปุถุชฺชา,} \\ \textbf{อุเปกฺขตี อุคฺคหณนฺติ มญฺเญ.}}
\csromangathagroup{\csromangathastanza{\textbf{ญตฺวา จ โส สมฺมุติโย ปุถุชฺชา,} \\ \textbf{อุเปกฺขตี อุคฺคหณนฺติ มญฺเญ.}}}

\prose{\textbf{น พฺราหฺมโณ กปฺปมุเปติ สงฺขา}ติ \textbf{นา}ติ ปฏิกฺเขโป.\csromanspacer{} \textbf{พฺราหฺมโณ}ติ สตฺตนฺนํ ธมฺมานํ พาหิตตฺตา พฺราหฺมโณ ฯเปฯ อสิโต ตาทิ ปวุจฺจเต ส พฺรหฺมา.\csromanspacer{} \textbf{กปฺปา}ติ เทฺว กปฺปา ตณฺหากปฺโป จ ทิฏฺฐิกปฺโป จ ฯเปฯ อยํ ตณฺหากปฺโป ฯเปฯ อยํ ทิฏฺฐิกปฺโป.\csromanspacer{} \textbf{สงฺขา} วุจฺจติ ญาณํ.\csromanspacer{} ยา ปญฺญา ปชาน\textbf{นา} ฯเปฯ}

\vspace{\tipitakaparskip}
\csromancenter{มหาวิยูหสุตฺตนิทฺเทส อโมโห ธมฺมวิจโย สมฺมาทิฏฺฐิ.\csromanspacer{} \textbf{น พฺราหฺมโณ กปฺปมุเปติ สงฺขา}ติ พฺราหฺมโณ สงฺขาย ชานิตฺวา ตุลยิตฺวา ตีรยิตฺวา วิภาวยิตฺวา วิภูตํ กตฺวา, “สพฺเพ สงฺขารา อนิจฺจา, สพฺเพ สงฺขารา ทุกฺขา ฯเปฯ ยํ กิญฺจิ สมุทยธมฺมํ, สพฺพํ ตํ นิโรธธมฺมนฺ”ติ สงฺขาย ชานิตฺวา ตุลยิตฺวา ตีรยิตฺวา วิภาวยิตฺวา วิภูตํ กตฺวา ตณฺหากปฺปํ วา ทิฏฺฐิกปฺปํ วา เนติ น อุเปติ น อุปคจฺฉติ น คณฺหาติ น ปรามสติ นาภินิวิสตีติ น พฺราหฺมโณ กปฺปมุเปติ สงฺขา.}
\prose{\csromanfolio{255}\textbf{น ทิฏฺฐิสารี นปิ ญาณพนฺธู}ติ ตสฺส ทฺวาสฏฺฐิ ทิฏฺฐิคตานิ ปหีนานิ สมุจฺฉินฺนานิ วูปสนฺตานิ ปฏิปสฺสทฺธานิ อภพฺพุปฺปตฺติกานิ ญาณคฺคินา ทฑฺฒานิ.\csromanspacer{} โส ทิฏฺฐิยา น ยายติ น นียติ น วุยฺหติ น สํหรียติ นปิ ตํ ทิฏฺฐิคตํ สารโต ปจฺเจติ น ปจฺจาคจฺฉตีติ น ทิฏฺฐิสารี.\csromanspacer{} \textbf{นปิ ญาณพนฺธู}ติ อฏฺฐสมาปตฺติญาเณน วา ปญฺจาภิญฺญาญาเณน วา ตณฺหาพนฺธุํ วา ทิฏฺฐิพนฺธุํ วา น กโรติ น ชเนติ น สญฺชเนติ น นิพฺพตฺเตติ นาภินิพฺพตฺเตตีติ น ทิฏฺฐิสารี นปิ ญาณพนฺธุ.}

\prose{\textbf{ญตฺวา จ โส สมฺมุติโย ปุถุชฺชา}ติ \textbf{ญตฺวา}ติ \textbf{ญตฺวา} ชานิตฺวา ตุลยิตฺวา ตีรยิตฺวา วิภาวยิตฺวา วิภูตํ กตฺวา. “สพฺเพ สงฺขารา อนิจฺจา”ติ \textbf{ญตฺวา} ชานิตฺวา ตุลยิตฺวา ตีรยิตฺวา วิภาวยิตฺวา วิภูตํ กตฺวา, “สพฺเพ สงฺขารา ทุกฺขา”ติ ฯเปฯ “ยํ กิญฺจิ สมุทยธมฺมํ, สพฺพํ ตํ นิโรธธมฺมนฺ”ติ \textbf{ญตฺวา} ชานิตฺวา ตุลยิตฺวา ตีรยิตฺวา วิภาวยิตฺวา วิภูตํ กตฺวาติ \textbf{ญตฺวา} จ โส.\csromanspacer{} \textbf{สมฺมุติโย} วุจฺจนฺติ ทฺวาสฏฺฐิ ทิฏฺฐิคตานิ ทิฏฺฐิสมฺมุติโย.\csromanspacer{} \textbf{ปุถุชฺชา}ติ ปุถุชฺชเนหิ ชนิตา วา ตา สมฺมุติโยติ ปุถุชฺชา.\csromanspacer{} ปุถุ นานาชเนหิ ชนิตา วา สมฺมุติโยติ ปุถุชฺชาติ \textbf{ญตฺวา} จ โส สมฺมุติโย ปุถุชฺชา.}

\prose{\textbf{อุเปกฺขตี อุคฺคหณนฺติ มญฺเญ}ติ อญฺเญ ตณฺหาวเสน ทิฏฺฐิวเสน คณฺหนฺติ ปรามสนฺติ อภินิวิสนฺติ.\csromanspacer{} อรหา อุเปกฺขติ น คณฺหาติ น ปรามสติ นาภินิวิสตีติ อุเปกฺขตี อุคฺคหณนฺติ มญฺเญ.\csromanspacer{} เตนาห ภควา\csromandash{}}

\csromangathameasure{\textbf{น พฺราหฺมโณ กปฺปมุเปติ สงฺขา},}
\csromangathagroup{\csromangathastanza{\textbf{น พฺราหฺมโณ กปฺปมุเปติ สงฺขา},}}

\prose{น ทิฏฺฐิสารี นปิ ญาณพนฺธุ.}

\prose{\textbf{ญตฺวา จ โส สมฺมุติโย ปุถุชฺชา},}

\prose{\textbf{อุเปกฺขตี อุคฺคหณนฺติ มญฺเญ}ติ.}

\proseitem{147}{\csromanfolio{256}\csromansymbolfootnote{*}{ขุ 1. 420 ปิฏฺเฐ.}\textbf{วิสฺสชฺช คนฺถานิ มุนีธ โลเก},}

\prose{\textbf{วิวาทชาเตสุ น วคฺคสารี.}}

\csromangathameasure{\textbf{สนฺโต อสนฺเตสุ อุเปกฺขโก โส,} \\ \textbf{อนุคฺคโห อุคฺคหณนฺติ มญฺเญ.}}
\csromangathagroup{\csromangathastanza{\textbf{สนฺโต อสนฺเตสุ อุเปกฺขโก โส,} \\ \textbf{อนุคฺคโห อุคฺคหณนฺติ มญฺเญ.}}}

\prose{\textbf{วิสฺสชฺช คนฺถานิ มุนีธ โลเก}ติ \textbf{คนฺถา}ติ จตฺตาโร \textbf{คนฺถา} อภิชฺฌากายคนฺโถ พฺยาปาโท กายคนฺโถ สีลพฺพตปรามาโส กายคนฺโถ อิทํสจฺจาภินิเวโส กายคนฺโถ.\csromanspacer{} อตฺตโน ทิฏฺฐิยา ราโค อภิชฺฌากายคนฺโถ, ปรวาเทสุ อาฆาโต อปฺปจฺจโย พฺยาปาโท กายคนฺโถ, อตฺตโน สีลํ วา วตํ วา สีลวตํ วา ปรามสติ สีลพฺพตปรามาโส กายคนฺโถ, อตฺตโน ทิฏฺฐิ อิทํสจฺจาภินิเวโส กายคนฺโถ.\csromanspacer{} \textbf{วิสฺสชฺชา}ติ คนฺเถ โวสฺสชฺชิตฺวา วิสฺสชฺช.\csromanspacer{} อถ วา คนฺเถ คถิเต คนฺถิเต\footnote{คถิเต คณฺฐิเต (พหูสุ) 204 ปิฏฺเฐ อฏฺฐกถา โอโลเกตพฺพา.} พนฺเธ วิพนฺเธ อาพนฺเธ\footnote{พทฺเธ วิพทฺเธ อาพทฺเธ (สี)} ลคฺเค ลคฺคิเต ปลิพุทฺเธ พนฺธเน โผฏยิตฺวา วิสฺสชฺช.\csromanspacer{} ยถา วยฺหํ วา รถํ วา สกฏํ วา สนฺทมานิกํ วา สชฺชํ วิสฺสชฺชํ กโรนฺติ วิโกเปนฺติ, เอวเมวํ คนฺเถ โวสฺสชฺชิตฺวา วิสฺสชฺช.\csromanspacer{} อถ วา คนฺเถ คถิเต คณฺฐิเต พนฺเธ วิพนฺเธ อาพนฺเธ ลคฺเค ลคฺคิเต ปลิพุทฺเธ พนฺธเน โผฏยิตฺวา วิสฺสชฺช.\csromanspacer{} \textbf{มุนี}ติ \textbf{โมนํ}วุจฺจติ ญาณํ ฯเปฯ สงฺคชาลมติจฺจ โส มุนิ.\csromanspacer{} \textbf{อิธา}ติ อิมิสฺสา ทิฏฺฐิยา ฯเปฯ อิมสฺมิํ มนุสฺสโลเกติ วิสฺสชฺช \textbf{คนฺถา}นิ มุนีธ โลเก.}

\prose{\textbf{วิวาทชาเตสุ น วคฺคสารี}ติ วิวาทชาเตสุ สญฺชาเตสุ นิพฺพตฺเตสุ อภินิพฺพตฺเตสุ ปาตุภูเตสุ\footnote{วิวาเท ชาเต สญฺชาเต นิพฺพตฺเต อภินิพฺพตฺเต ปาตุภูเต (สี, ก)} ฉนฺทาคติํ คจฺฉนฺเตสุ โทสาคติํ คจฺฉนฺเตสุ ภยาคติํ คจฺฉนฺเตสุ โมหาคติํ คจฺฉนฺเตสุ น ฉนฺทาคติํ คจฺฉติ, น โทสาคติํ คจฺฉติ, น ภยาคติํ คจฺฉติ, น โมหาคติํ คจฺฉติ, น ราควเสน คจฺฉติ, น โทสวเสน คจฺฉติ, น โมหวเสน คจฺฉติ, น มานวเสน คจฺฉติ, น ทิฏฺฐิวเสน คจฺฉติ, น อุทฺธจฺจวเสน คจฺฉติ, น วิจิกิจฺฉาวเสน คจฺฉติ, น อนุสยวเสน คจฺฉติ, น วคฺเคหิ ธมฺเมหิ ยายติ นียติ วุยฺหติ สํหรียตีติ วิวาทชาเตสุ น วคฺคสารี.}

\prose{\textbf{สนฺโต อสนฺเตสุ อุเปกฺขโก โส}ติ \textbf{สนฺโต}ติ ราคสฺส สนฺตตฺตา \textbf{สนฺโต}, โทสสฺส สนฺตตฺตา \textbf{สนฺโต}, โมหสฺส สนฺตตฺตา \textbf{สนฺโต} ฯเปฯ สพฺพากุสลาภิสงฺขารานํ สนฺตตฺตา สมิตตฺตา วูปสมิตตฺตา วิชฺฌาตตฺตา}

\vspace{\tipitakaparskip}
\csromancenter{มหาวิยูหสุตฺตนิทฺเทส นิพฺพุตตฺตา วิคตตฺตา ปฏิปสฺสทฺธตฺตา สนฺโต อุปสนฺโต วูปสนฺโต นิพฺพุโต ปฏิปสฺสทฺโธติ สนฺโต.\csromanspacer{} \textbf{อสนฺเตสู}ติ อสนฺเตสุ อนุปสนฺเตสุ อวูปสนฺเตสุ อนิพฺพุเตสุ อปฺปฏิปสฺสทฺเธสูติ สนฺโต อสนฺเตสุ.\csromanspacer{} \textbf{อุเปกฺขโก โส}ติ * อรหา ฉฬงฺคุเปกฺขาย สมนฺนาคโต จกฺขุนา รูปํ ทิสฺวา เนว สุมโน โหติ น ทุมฺมโน อุเปกฺขโก วิหรติ สโต สมฺปชาโน.\csromanspacer{} โสเตน สทฺทํ สุตฺวา ฯเปฯ กาลํ กงฺขติ ภาวิโต ส ทนฺโตติ สนฺโต อสนฺเตสุ อุเปกฺขโก โส.}
\prose{\csromanfolio{257}\textbf{อนุคฺคโห อุคฺคหณนฺติ มญฺเญ}ติ อญฺเญ ตณฺหาวเสน ทิฏฺฐิวเสน คณฺหนฺเต ปรามสนฺเต อภินิวิสนฺเต อรหา อุเปกฺขติ น คณฺหาติ น ปรามสติ นาภินิวิสตีติ อนุคฺคโห อุคฺคหณนฺติ มญฺเญ.\csromanspacer{} เตนาห ภควา\csromandash{}}

\csromangathameasure{วิสฺสชฺช คนฺถานิ มุนีธ โลเก, \\ วิวาทชาเตสุ น วคฺคสารี. \\ สนฺโต อสนฺเตสุ อุเปกฺขโก โส,}
\csromangathagroup{\csromangathastanza{วิสฺสชฺช คนฺถานิ มุนีธ โลเก, \\ วิวาทชาเตสุ น วคฺคสารี. \\ สนฺโต อสนฺเตสุ อุเปกฺขโก โส,}}

\prose{\textbf{อนุคฺคโห อุคฺคหณนฺติ มญฺเญ}ติ.}

\proseitem{148}{\csromansymbolfootnote{+}{ขุ 1. 420 ปิฏฺเฐ.}ปุพฺพาสเว หิตฺวา นเว อกุพฺพํ,}

\prose{\textbf{น ฉนฺทคู โนปิ นิวิสฺสวาที.}}

\csromangathameasure{\textbf{ส วิปฺปมุตฺโต ทิฏฺฐิคเตหิ ธีโร,} \\ \textbf{น ลิมฺปติ โลเก อนตฺตครหี.}}
\csromangathagroup{\csromangathastanza{\textbf{ส วิปฺปมุตฺโต ทิฏฺฐิคเตหิ ธีโร,} \\ \textbf{น ลิมฺปติ โลเก อนตฺตครหี.}}}

\prose{\textbf{ปุพฺพาสเว หิตฺวา นเว อกุพฺพนฺ}ติ \textbf{ปุพฺพาสวา} วุจฺจนฺติ อตีตา รูปเวทนาสญฺญาสงฺขารวิญฺญาณา, อตีเต สงฺขาเร อารพฺภ เย กิเลสา อุปฺปชฺเชยฺยุํ, เต กิเลเส หิตฺวา จชิตฺวา ปริจฺจชิตฺวา ปชหิตฺวา วิโนเทตฺวา พฺยนฺติํ กริตฺวา อนภาวํ คเมตฺวาติ ปุพฺพาสเว หิตฺวา.\csromanspacer{} \textbf{นเว อกุพฺพนฺ}ติ \textbf{นวา} วุจฺจนฺติ ปจฺจุปฺปนฺนา รูปเวทนาสญฺญาสงฺขารวิญฺญาณา, ปจฺจุปฺปนฺเน สงฺขาเร อารพฺภ ฉนฺทํ\footnote{ขนฺติํ (ก)} อกุพฺพมาโน เปมํ อกุพฺพมาโน ราคํ อกุพฺพมาโน อชนยมาโน อสญฺชนยมาโน อนิพฺพตฺตยมาโน อนภินิพฺพตฺตยมาโนติ ปุพฺพาสเว หิตฺวา นเว อกุพฺพํ.}

\prose{\textbf{น ฉนฺทคู โนปิ นิวิสฺสวาที}ติ ฉนฺทาคติํ คจฺฉติ, น โทสาคติํ คจฺฉติ, น โมหาคติํ คจฺฉติ, น ภยาคติํ คจฺฉติ, น ราควเสน \csromanfolio{258}คจฺฉติ, น โทสวเสน คจฺฉติ, น โมหวเสน คจฺฉติ, น มานวเสน คจฺฉติ, น ทิฏฺฐิวเสน คจฺฉติ, น อุทฺธจฺจวเสน คจฺฉติ, น วิจิกิจฺฉาวเสน คจฺฉติ, น อนุสยวเสน คจฺฉติ, น วคฺเคหิ ธมฺเมหิ ยายติ นียติ วุยฺหติ น สํหรียตีติ น ฉนฺทคู.\csromanspacer{} \textbf{โนปิ นิวิสฺสวาที}ติ “สสฺสโต โลโก, อิทเมว สจฺจํ โมฆมญฺญนฺ”ติ น นิวิสฺสวาที ฯเปฯ “เนว โหติ น น โหติ ตถาคโต ปรํ มรณา, อิทเมว สจฺจํ โมฆมญฺญนฺ”ติ น นิวิสฺสวาทีติ น ฉนฺทคู โนปิ นิวิสฺสวาที.}

\prose{\textbf{ส วิปฺปมุตฺโต ทิฏฺฐิคเตหิ ธีโร}ติ ตสฺส ทฺวาสฏฺฐิ ทิฏฺฐิคตานิ ปหีนานิ สมุจฺฉินฺนานิ วูปสนฺตานิ ปฏิปสฺสทฺธานิ อภพฺพุปฺปตฺติกานิ ญาณคฺคินา ทฑฺฒานิ.\csromanspacer{} โส ทิฏฺฐิคเตหิ วิปฺปมุตฺโต วิสญฺญุตฺโต, วิมริยาทิกเตน เจตสา วิหรติ.\csromanspacer{} \textbf{ธีโร}ติ ธีโร ปณฺฑิโต ปญฺญวา พุทฺธิมา ญาณี วิภาวี เมธาวีติ ส วิปฺปมุตฺโต ทิฏฺฐิคเตหิ ธีโร.}

\prose{\textbf{น ลิมฺปติ โลเก อนตฺตครหี}ติ \textbf{เลปา}ติ เทฺว \textbf{เลปา} ตณฺหาเลโป จ ทิฏฺฐิเลโป จ ฯเปฯ อยํ ตณฺหาเลโป ฯเปฯ อยํ ทิฏฺฐิเลโป.\csromanspacer{} ตสฺส ตณฺหาเลโป ปหีโน, ทิฏฺฐิเลโป ปฏินิสฺสฏฺโฐ, ตสฺส ตณฺหาเลปสฺส ปหีนตฺตา, ทิฏฺฐิเลปสฺส ปฏินิสฺสฏฺฐตฺตา อปายโลเก น ลิมฺปติ, มนุสฺสโลเก น ลิมฺปติ, เทวโลเก น ลิมฺปติ, ขนฺธโลเก น ลิมฺปติ.\csromanspacer{} ธาตุโลเก น ลิมฺปติ, อายตนโลเก น ลิมฺปติ น ปลิมฺปติ น อุปลิมฺปติ, อลิตฺโต อปลิตฺโต อนุปลิตฺโต นิกฺขนฺโต นิสฺสโฏ วิปฺปมุตฺโต วิสญฺญุตฺโต, วิมริยาทิกเตน เจตสา วิหรตีติ น ลิมฺปติ โลเก.}

\prose{\textbf{อนตฺตครหี}ติ ทฺวีหิ การเณหิ อตฺตานํ ครหติ กตตฺตา จ อกตตฺตา จ.\csromanspacer{} กถํ กตตฺตา จ อกตตฺตา จ อตฺตานํ ครหติ, “กตํ เม กายทุจฺจริตํ, อกตํ เม กายสุจริตนฺ”ติ อตฺตานํ ครหติ. “กตํ เม วจีทุจฺจริตํ ฯเปฯ กตํ เม มโนทุจฺจริตํ.\csromanspacer{} กโต เม ปาณาติปาโต ฯเปฯ กตา เม มิจฺฉาทิฏฺฐิ, อกตา เม สมฺมาทิฏฺฐี”ติ อตฺตานํ ครหติ.\csromanspacer{} เอวํ กตตฺตา จ อกตตฺตา จ อตฺตานํ ครหติ.}

\prose{อถ วา “สีเลสุ’มฺหิ น ปริปูรการี”ติ อตฺตานํ ครหติ. “อินฺทฺริเยสุ’มฺหิ อคุตฺตทฺวาโร”ติ. “โภชเน’มฺหิ อมตฺตญฺญู”ติ. “ชาคริย’มฺหิ อนนุยุตฺโต”ติ. “น สติสมฺปชญฺเญนา’มฺหิ สมนฺนาคโต”ติ. “อภาวิตา}

\vspace{\tipitakaparskip}
\csromancenter{มหาวิยูหสุตฺตนิทฺเทส เม จตฺตาโร สติปฏฺฐานา”ติ. “อภาวิตา เม จตฺตาโร สมฺมปฺปธานา”ติ. “อภาวิตา เม จตฺตาโร อิทฺธิปาทา”ติ. “อภาวิตานิ เม ปญฺจินฺทฺริยานี”ติ. “อภาวิตานิ เม ปญฺจ พลานี”ติ. “อภาวิตา เม สตฺต โพชฺฌงฺคา”ติ. “อภาวิโต เม อริโย อฏฺฐงฺคิโก มคฺโค”ติ. “ทุกฺขํ เม อปริญฺญาตนฺ”ติ. “ทุกฺขสมุทโย เม อปฺปหีโน”ติ. “มคฺโค เม อภาวิโต”ติ. “นิโรโธ เม อสจฺฉิกโต”ติ อตฺตานํ ครหติ.\csromanspacer{} เอวํ กตตฺตา จ อกตตฺตา จ อตฺตานํ ครหติ.\csromanspacer{} เอวํ อตฺตครหี.\csromanspacer{} ตยิทํ กมฺมํ อกุพฺพมาโน อชนยมาโน อสญฺชนยมาโน อนิพฺพตฺตยมาโน อนภินิพฺพตฺตยมาโน อนตฺตครหีติ น ลิมฺปติ โลเก อนตฺตครหี.\csromanspacer{} เตนาห ภควา\csromandash{}}
\vspace{0.5\baselineskip}
\csromangathameasure{ปุพฺพาสเว หิตฺวา นเว อกุพฺพํ, \\ น ฉนฺทคู โนปิ นิวิสฺสวาที. \\ ส วิปฺปมุตฺโต ทิฏฺฐิคเตหิ ธีโร, \\ น ลิมฺปติ โลเก อนตฺตครหีติ.}
\csromangathagroup{\csromangathastanza{\csromanfolio{259}ปุพฺพาสเว หิตฺวา นเว อกุพฺพํ, \\ น ฉนฺทคู โนปิ นิวิสฺสวาที. \\ ส วิปฺปมุตฺโต ทิฏฺฐิคเตหิ ธีโร, \\ น ลิมฺปติ โลเก อนตฺตครหีติ.}}

\proseitem{149}{\csromansymbolfootnote{*}{ขุ 1. 421 ปิฏฺเฐ.}ส สพฺพธมฺเมสุ วิเสนิภูโต,}

\prose{\textbf{ยํ กิญฺจิ ทิฏฺฐํ ว สุตํ มุตํ วา.}}

\csromangathameasure{\textbf{ส ปนฺนภาโร มุนิ วิปฺปมุตฺโต,}}
\csromangathagroup{\csromangathastanza{\textbf{ส ปนฺนภาโร มุนิ วิปฺปมุตฺโต,}}}

\vspace{\tipitakaparskip}
\csromancenter{\textbf{น กปฺปิโย นุปรโต น ปตฺถิโย. (อิติ ภควา,)}}
\prose{\textbf{ส สพฺพธมฺเมสุ วิเสนิภูโต, ยํ กิญฺจิ ทิฏฺฐํ ว สุตํ มุตํ วา}ติ \textbf{เสนา} วุจฺจติ มาร\textbf{เสนา}, กายทุจฺจริตํ มาร\textbf{เสนา}, วจีทุจฺจริตํ มาร\textbf{เสนา}, มโนทุจฺจริตํ มาร\textbf{เสนา}, ราโค.\csromanspacer{} โทโส.\csromanspacer{} โมโห.\csromanspacer{} โกโธ.\csromanspacer{} อุปนาโห.\csromanspacer{} มกฺโข.\csromanspacer{} ปฬาโส.\csromanspacer{} อิสฺสา.\csromanspacer{} มจฺฉริยํ.\csromanspacer{} มายา.\csromanspacer{} สาเฐยฺยํ.\csromanspacer{} ถมฺโภ.\csromanspacer{} สารมฺโภ.\csromanspacer{} มาโน.\csromanspacer{} อติมาโน.\csromanspacer{} มโท.\csromanspacer{} ปมาโท.\csromanspacer{} สพฺเพ กิเลสา.\csromanspacer{} สพฺเพ ทุจฺจริตา.\csromanspacer{} สพฺเพ ทรถา.\csromanspacer{} สพฺเพ ปริฬาหา.\csromanspacer{} สพฺเพ สนฺตาปา.\csromanspacer{} สพฺพากุสลาภิสงฺขารา มาร\textbf{เสนา}.\csromanspacer{} วุตฺตํ เหตํ ภควตา\csromandash{}}

\csromangathameasure{+กามา เต ปฐมา เสนา, ทุติยา อรติ วุจฺจติ ฯเปฯ. \\ น นํ อสูโร ชินาติ, เชตฺวาว ลภเต สุขนฺติ.}
\csromangathagroup{\csromangathastanza{\csromansymbolfootnote{+}{ขุ 1. 342; ขุ 8. 114 ปิฏฺเฐสุ.}กามา เต ปฐมา เสนา, ทุติยา อรติ วุจฺจติ ฯเปฯ. \\ น นํ อสูโร ชินาติ, เชตฺวาว ลภเต สุขนฺติ.}}

\prose{ยโต จตูหิ อริยมคฺเคหิ สพฺพา จ มาร\textbf{เสนา} สพฺเพ จ ปฏิเสนิกรา กิเลสา ชิตา จ ปราชิตา จ ภคฺคา วิปฺปลุคฺคา ปรมฺมุขา, โส วุจฺจติ วิเสนิภูโต.\csromanspacer{} โส ทิฏฺเฐ วิเสนิภูโต, สุเต.\csromanspacer{} มุเต.\csromanspacer{} วิญฺญาเต วิเสนิภูโตติ ส สพฺพธมฺเมสุ วิเสนิภูโต, ยํ กิญฺจิ ทิฏฺฐํ ว สุตํ มุตํ วา.}

\prose{\csromanfolio{260}\textbf{ส ปนฺนภาโร มุนิ วิปฺปมุตฺโต}ติ \textbf{ภารา}ติ ตโย \textbf{ภารา} ขนฺธภาโร กิเลสภาโร อภิสงฺขารภาโร.\csromanspacer{} กตโม ขนฺธภาโร, ปฏิสนฺธิยา รูปํ เวทนา สญฺญา สงฺขารา วิญฺญาณํ, อยํ ขนฺธภาโร.\csromanspacer{} กตโม กิเลสภาโร, ราโค โทโส โมโห ฯเปฯ สพฺพากุสลาภิสงฺขารา, อยํ กิเลสภาโร.\csromanspacer{} กตโม อภิสงฺขารภาโร, ปุญฺญาภิสงฺขาโร อปุญฺญาภิสงฺขาโร อาเนญฺชาภิสงฺขาโร, อยํ อภิสงฺขารภาโร.\csromanspacer{} ยโต ขนฺธภาโร จ กิเลสภาโร จ อภิสงฺขารภาโร จ ปหีนา โหนฺติ อุจฺฉินฺนมูลา ตาลาวตฺถุกตา อนภาวํกตา อายติํ อนุปฺปาทธมฺมา, โส วุจฺจติ ปนฺนภาโร ปติตภาโร โอโรปิตภาโร สโมโรปิตภาโร นิกฺขิตฺตภาโร ปฏิปสฺสทฺธภาโร.}

\prose{\csromansymbolfootnote{*}{ขุ 8. 73 ปิฏฺเฐ.}มุนีติ \textbf{โมนํ} วุจฺจติ ญาณํ.\csromanspacer{} ยา ปญฺญา ปชานนา วิจโย ปวิจโย ธมฺมวิจโย สลฺลกฺขณา อุปลกฺขณา ปจฺจุปลกฺขณา ปณฺฑิจฺจํ โกสลฺลํ เนปุญฺญํ เวภพฺยา จินฺตา อุปปริกฺขา ภูริ เมธา ปรินายิกา วิปสฺสนา สมฺปชญฺญํ ปโตโท ปญฺญา ปญฺญินฺทฺริยํ ปญฺญาพลํ ปญฺญาสตฺถํ ปญฺญาปาสาโท ปญฺญา อาโลโก ปญฺญา-โอภาโส ปญฺญาปชฺโชโต ปญฺญารตนํ อโมโห ธมฺมวิจโย สมฺมาทิฏฺฐิ.\csromanspacer{} เตน ญาเณน สมนฺนาคโต มุนิ โมนปฺปตฺโต.}

\prose{ตีณิ โมเนยฺยานิ กายโมเนยฺยํ วจีโมเนยฺยํ มโนโมเนยฺยํ, กตมํ กายโมเนยฺยํ, ติวิธานํ กายทุจฺจริตานํ ปหานํ กายโมเนยฺยํ, ติวิธํ กายสุจริตํ กายโมเนยฺยํ, กายารมฺมเณ ญาณํ กายโมเนยฺยํ, กายปริญฺญา กายโมเนยฺยํ, ปริญฺญาสหคโต มคฺโค กายโมเนยฺยํ, กาเย ฉนฺทราคสฺส ปหานํ กายโมเนยฺยํ, กายสงฺขารนิโรโธ จตุตฺถชฺฌานสมาปตฺติ กายโมเนยฺยํ, อิทํ กายโมเนยฺยํ.}

\prose{กตมํ วจีโมเนยฺยํ จตุพฺพิธานํ วจีทุจฺจริตานํ ปหานํ วจีโมเนยฺยํ, จตุพฺพิธํ วจีสุจริตํ วจีโมเนยฺยํ, วาจารมฺมเณ ญาณํ วจีโมเนยฺยํ, วาจาปริญฺญา วจีโมเนยฺยํ, ปริญฺญาสหคโต มคฺโค วจีโมเนยฺยํ, วาจาย ฉนฺทราคสฺส ปหานํ วจีโมเนยฺยํ, วจีสงฺขารนิโรโธ ทุติยชฺฌานสมาปตฺติ วจีโมเนยฺยํ, อิทํ วจีโมเนยฺยํ.}

\prose{กตมํ มโนโมเนยฺยํ, ติวิธานํ มโนทุจฺจริตานํ ปหานํ มโนโมเนยฺยํ, ติวิธํ มโนสุจริตํ มโนโมเนยฺยํ, จิตฺตารมฺมเณ ญาณํ}

\vspace{\tipitakaparskip}
\csromancenter{มหาวิยูหสุตฺตนิทฺเทส มโนโมเนยฺยํ, จิตฺตปริญฺญา มโนโมเนยฺยํ, ปริญฺญาสหคโต มคฺโค มโนโมเนยฺยํ, จิตฺเต ฉนฺทราคสฺส ปหานํ มโนโมเนยฺยํ, จิตฺตสงฺขารนิโรโธ สญฺญาเวทยิตนิโรธสมาปตฺติ มโนโมเนยฺยํ, อิทํ มโนโมเนยฺยํ.}
\prose{\csromanfolio{261}กายมุนิํ วาจามุนิํ, มโนมุนิมนาสวํ.}

\prose{มุนิํ โมเนยฺยสมฺปนฺนํ, อาหุ สพฺพปฺปหายินํ.}

\prose{กายมุนิํ วาจามุนิํ, มโนมุนิมนาสวํ.}

\prose{มุนิํ โมเนยฺยสมฺปนฺนํ, อาหุ นินฺหาตปาปกนฺติ\footnote{นินหาตปาปกนฺติ (สี) อํ 1. 277; ขุ 1. 234; ขุ 8. 73 ปิฏฺเฐสุ.}.}

\prose{อิเมหิ ตีหิ โมเนยฺเยหิ ธมฺเมหิ สมนฺนาคตา ฉ มุนิโน\footnote{มุนโย (สี, สฺยา, ก)} อคารมุนิโน อนคารมุนิโน เสขมุนิโน อเสขมุนิโน ปจฺเจกมุนิโน มุนิมุนิโนติ\footnote{มุนิมุนิโน (สี, สฺยา, ก) 44 ปิฏฺเฐปิ.}.\csromanspacer{} กตเม อคารมุนิโน, เย เต อคาริกา ทิฏฺฐปทา วิญฺญาตสาสนา, อิเม อคารมุนิโน.\csromanspacer{} กตเม อนคารมุนิโน, เย เต ปพฺพชิตา ทิฏฺฐปทา วิญฺญาตสาสนา, อิเม อนคารมุนิโน.\csromanspacer{} สตฺต เสขา เสขมุนิโน.\csromanspacer{} อรหนฺโต อเสขมุนิโน.\csromanspacer{} ปจฺเจกพุทฺธา ปจฺเจกมุนิโน.\csromanspacer{} มุนิมุนิโน วุจฺจนฺติ ตถคตา อรหนฺโต สมฺมาสมฺพุทฺธา.}

\csromanmatikaanchor{mtk.14}
\csromantocmarkref{subsection}{13. มหาวิยูหสุตฺตนิทฺเทส}{mtk.14}
\bookmark[dest={mtk.14},level=2]{13. มหาวิยูหสุตฺตนิทฺเทส}
\prose{น โมเนน มุนี โหติ, มูฬฺหรูโป อวิทฺทสุ.}

\prose{โย จ ตุลํว ปคฺคยฺห, วรมาทาย ปณฺฑิโต.}

\prose{ปาปานิ ปริวชฺเชติ, ส มุนี เตน โส มุนิ.}

\prose{โย มุนาติ อุโภ โลเก, “มุนิ” เตน ปวุจฺจติ.}

\prose{อสตญฺจ สตญฺจ ญตฺวา ธมฺมํ, อชฺฌตฺตํ พหิทฺธา จ สพฺพโลเก.}

\prose{เทวมนุสฺเสหิ ปูชิโต โย, สงฺคชาลมติจฺจ โส มุนิ.}

\prose{\textbf{วิปฺปมุตฺโต}ติ มุนิโน ราคา จิตฺตํ มุตฺตํ วิมุตฺตํ สุวิมุตฺตํ.\csromanspacer{} โทสา จิตฺตํ.\csromanspacer{} โมหา จิตฺตํ มุตฺตํ วิมุตฺตํ สุวิมุตฺตํ ฯเปฯ สพฺพากุสลาภิสงฺขาเรหิ จิตฺตํ มุตฺตํ วิมุตฺตํ สุวิมุตฺตนฺติ ส ปนฺนภาโร มุนิ วิปฺปมุตฺโต.}

\prose{\textbf{น กปฺปิโย นุปรโต น ปตฺถิโยติ ภควา}ติ \textbf{กปฺปา}ติ เทฺว \textbf{กปฺปา} ตณฺหากปฺโป จ ทิฏฺฐิกปฺโป จ ฯเปฯ อยํ ตณฺหากปฺโป ฯเปฯ อยํ ทิฏฺฐิกปฺโป.\csromanspacer{} ตสฺส ตณฺหากปฺโป ปหีโน, ทิฏฺฐิกปฺโป ปฏินิสฺสฏฺโฐ, ตณฺหากปฺปสฺส \csromanfolio{262}ปหีนตฺตา\textbf{,} ทิฏฺฐิกปฺปสฺส ปฏินิสฺสฏฺฐตฺตา ตณฺหากปฺปํ วา ทิฏฺฐิกปฺปํวา น กปฺเปติ น ชเนติ น สญฺชเนติ น นิพฺพตฺเตติ นาภินิพฺพตฺเตตีติ น กปฺปิโย\textbf{.}\csromanspacer{} \textbf{นูปรโต}ติ สพฺเพ พาลปุถุชฺชนา รชฺชนฺติ\textbf{,} ปุถุชฺชนกลฺยาณกํ อุปาทาย สตฺต เสขา อปฺปตฺตสฺส ปตฺติยา อนธิคตสฺส อธิคมาย อสจฺฉิกตสฺส สจฺฉิกิริยาย อารมนฺติ วิรมนฺติ ปฏิวิรมนฺติ\textbf{,} อรหา อารโต วิรโต ปฏิวิรโต นิกฺขนฺโต นิสฺสโฏ วิปฺปมุตฺโต วิสญฺญุตฺโถ\textbf{,} วิมริยาทิกเตน เจตสา วิหรตีติ น กปฺปิโย นูปรโต\textbf{.}\csromanspacer{} \textbf{น ปตฺถิโย}ติ \textbf{ปตฺถนา} วุจฺจติ ตณฺหา\textbf{.}\csromanspacer{} โย ราโค สาราโค ฯเปฯ อภิชฺฌา โลโภ อกุสลมูลํ\textbf{.}\csromanspacer{} ยสฺเสสา \textbf{ปตฺถนา} ตณฺหา ปหีนา สมุจฺฉินฺนา วูปสนฺตา ปฏิปสฺสทฺธา อภพฺพุปฺปตฺติกา ญาณคฺคินา ทฑฺฒา\textbf{,} โส วุจฺจติ น ปตฺถิโย\textbf{.}}

\prose{\textbf{ภควา}ติ คารวาธิวจนํ.\csromanspacer{} อปิ จ ภคฺคราโคติ \textbf{ภควา}, ภคฺคโทโสติ \textbf{ภควา}, ภคฺคโมโหติ \textbf{ภควา}, ภคฺคมาโนติ \textbf{ภควา}, ภคฺคทิฏฺฐีติ \textbf{ภควา}, ภคฺคกณฺฑโกติ \textbf{ภควา}, ภคฺคกิเลโสติ \textbf{ภควา}, ภชิ วิภชิ ปวิภชิ ธมฺมรตนนฺติ \textbf{ภควา}, ภวานํ อนฺตกโรติ \textbf{ภควา}, ภาวิตกาโย, ภาวิตสีโล, ภาวิตจิตฺโต, ภาวิตปญฺโญติ \textbf{ภควา}, ภชิ วา \textbf{ภควา} อรญฺญวนปตฺถานิ ปนฺตานิ เสนาสนานิ อปฺปสทฺทานิ อปฺปนิคฺโฆสานิ วิชนวาตานิ มนุสฺสราหสฺเสยฺยกานิ ปฏิสลฺลานสารุปฺปานีติ \textbf{ภควา}, ภาคี วา \textbf{ภควา} จีวรปิณฺฑปาตเสนาสนคิลานปจฺจยเภสชฺช- ปริกฺขารานนฺติ \textbf{ภควา}.\csromanspacer{} ภาคี วา \textbf{ภควา} อตฺถรสสฺส ธมฺมรสสฺส วิมุตฺติรสสฺส อธิสีลสฺส อธิจิตฺตสฺส อธิปญฺญายาติ \textbf{ภควา}, ภาคี วา \textbf{ภควา} จตุนฺนํ ฌานานํ จตุนฺนํ อปฺปมญฺญานํ จตุนฺนํ อรูปสมาปตฺตีนนฺติ \textbf{ภควา}, ภาคี วา \textbf{ภควา} อฏฺฐนฺนํ วิโมกฺขานํ อฏฺฐนฺนํ อภิภายตนานํ นวนฺนํ อนุปุพฺพวิหารสมาปตฺตีนนฺติ \textbf{ภควา}, ภาคี วา \textbf{ภควา} ทสนฺนํ สญฺญาภาวนานํ ทสนฺนํ กสิณสมาปตฺตีนํ อานาปานสฺสติสมาธิสฺส อสุภสมาปตฺติยาติ \textbf{ภควา}, ภาคี วา \textbf{ภควา} จตุนฺนํ สติปฏฺฐานานํ จตุนฺนํ สมฺมปฺปธานานํ จตุนฺนํ อิทฺธิปาทานํ ปญฺจนฺนํ อินฺทฺริยานํ ปญฺจนฺนํ พลานํ สตฺตนฺนํ โพชฺฌงฺคานํ อริยสฺส อฏฺฐงฺคิกสฺส มคฺคสฺสาติ \textbf{ภควา}, ภาคี วา \textbf{ภควา} ทสนฺนํ ตถาคตพลานํ จตุนฺนํ เวสารชฺชานํ จตุนฺนํ ปฏิสมฺภิทานํ ฉนฺนํ อภิญฺญานํ ฉนฺนํ พุทฺธธมฺมานนฺติ \textbf{ภควา}.\csromanspacer{} \textbf{ภควา}ติ เนตํ นามํ มาตรา กตํ, น ปิตรา กตํ, น ภาตรา กตํ, น ภคินิยา กตํ, น มิตฺตามจฺเจหิ กตํ, น}

\vspace{\tipitakaparskip}
\csromancenter{ตุวฏกสุตฺตนิทฺเทส ญาติสาโลหิเตหิ กตํ\textbf{,} น สมณพฺราหฺมเณหิ กตํ\textbf{,} น เทวตาหิ กตํ\textbf{,} วิโมกฺขนฺติกเมตํ พุทฺธานํ ภควนฺตานํ โพธิยา มูเล สห สพฺพญฺญุตญฺญาณสฺส ปฏิลาภา สจฺฉิกา ปญฺญตฺติ ยทิทํ ภควาติ น กปฺปิโย นูปรโต น ปตฺถิโย อิติ ภควา\textbf{.}\csromanspacer{} เตนาห ภควา\csromandash{}}
\vspace{0.5\baselineskip}
\csromanmatikaanchor{mtk.15}
\csromantocmarkref{subsection}{14. ตุวฏฺฏกสุตฺตนิทฺเทส}{mtk.15}
\bookmark[dest={mtk.15},level=2]{14. ตุวฏฺฏกสุตฺตนิทฺเทส}
\csromangathameasure{ส สพฺพธมฺเมสุ วิเสนิภูโต\textbf{,} \\ ยํ กิญฺจิ ทิฏฺฐํ ว สุตํ มุตํ วา\textbf{.} \\ ส ปนฺนภาโร มุนิ วิปฺปมุตฺโต\textbf{,}}
\csromangathagroup{\csromangathastanza{\csromanfolio{263}ส สพฺพธมฺเมสุ วิเสนิภูโต\textbf{,} \\ ยํ กิญฺจิ ทิฏฺฐํ ว สุตํ มุตํ วา\textbf{.} \\ ส ปนฺนภาโร มุนิ วิปฺปมุตฺโต\textbf{,}}}

\vspace{\tipitakaparskip}
\csromancenter{น กปฺปิโย นุปรโต น ปตฺถิโย\textbf{.} (อิติ ภควาติ\textbf{,}) มหาพฺยูหสุตฺตนิทฺเทโส เตรสโม\textbf{.}}
\csromansectionrule
\csromanheader{2}{14. ตุวฏกสุตฺตนิทฺเทส}
\csromanheader{2}{อถ ตุวฏกสุตฺตนิทฺเทสํ วกฺขติ\csromandash{}}
\proseitem{150}{\csromansymbolmark{*}ปุจฺฉามิ ตํ อาทิจฺจพนฺธุ\footnote{อาทิจฺจพนฺธู (สี, สฺยา)}\textbf{,}}

\prose{\textbf{วิเวกํ สนฺติปทญฺจ มเหสิ}\footnote{มเหสิํ (สี, สฺยา)}\textbf{.}}

\csromangathameasure{\textbf{กถํ ทิสฺวา นิพฺพาติ ภิกฺขุ,} \\ \textbf{อนุปาทิยาโน โลกสฺมิํ กิญฺจิ.}}
\csromangathagroup{\csromangathastanza{\textbf{กถํ ทิสฺวา นิพฺพาติ ภิกฺขุ,} \\ \textbf{อนุปาทิยาโน โลกสฺมิํ กิญฺจิ.}}}

\prose{\textbf{ปุจฺฉามิ ตํ อาทิจฺจพนฺธู}ติ \textbf{ปุจฺฉา}ติ + ติสฺโส \textbf{ปุจฺฉา} อทิฏฺฐโชตนา \textbf{ปุจฺฉา} ทิฏฺฐสํสนฺทนา \textbf{ปุจฺฉา} วิมติจฺเฉทนา \textbf{ปุจฺฉา.}\csromanspacer{} กตมา อทิฏฺฐโชตนา \textbf{ปุจฺฉา,} ปกติยา ลกฺขณํ อญฺญาตํ โหติ อทิฏฺฐํ อตุลิตํ อตีริตํ อวิภูตํ อวิภาวิตํ\textbf{,} ตสฺส ญาณาย ทสฺสนาย ตุลนาย ตีรณาย วิภาวนาย ปญฺหํ ปุจฺฉติ\textbf{,} อยํ อทิฏฺฐโชตนา \textbf{ปุจฺฉา.}\csromanspacer{}\csromanspacer{}กตมา ทิฏฺฐสํสนฺทนา \textbf{ปุจฺฉา,} ปกติยา ลกฺขณํ ญาตํ โหติ ทิฏฺฐํ ตุลิตํ ตีริตํ วิภูตํ วิภาวิตํ\textbf{,} อญฺเญหิ ปณฺฑิเตหิ สทฺธิํ สํสนฺทนตฺถาย ปญฺหํ ปุจฺฉติ\textbf{,} อยํ ทิฏฺฐสํสนฺทนา \textbf{ปุจฺฉา.}\csromanspacer{}\csromanspacer{}กตมา วิมติจฺเฉทนา \textbf{ปุจฺฉา,} ปกติยา สํสยปกฺขนฺโท\footnote{สํสยปกฺขนฺโน (สี, สฺยา) * ขุ 1. 42 ปิฏฺเฐ.} โหติ\textbf{,} วิมติปกฺขนฺโท เทฺวฬฺหกชาโต “เอวํ นุ โข\textbf{,} น นุ โข\textbf{,} กิํ นุ โข\textbf{,} กถํ นุ โข”ติ\textbf{,} โส วิมติจฺเฉทนตฺถาย ปญฺหํ ปุจฺฉติ\textbf{,} อยํ วิมติจฺเฉทนา \textbf{ปุจฺฉา.}\csromanspacer{} อิมา ติสฺโส \textbf{ปุจฺฉา.}}

\prose{\csromanfolio{264}อปราปิ ติสฺโส ปุจฺฉา มนุสฺสปุจฺฉา อมนุสฺสปุจฺฉา นิมฺมิตปุจฺฉา.\csromanspacer{} กตมา มนุสฺสปุจฺฉา, มนุสฺสา พุทฺธํ ภควนฺตํ อุปสงฺกมิตฺวา ปญฺหํ ปุจฺฉนฺติ, ภิกฺขู ปุจฺฉนฺติ, ภิกฺขุนิโย ปุจฺฉนฺติ, อุปาสกา ปุจฺฉนฺติ, อุปาสิกาโย ปุจฺฉนฺติ, ราชาโน ปุจฺฉนฺติ, ขตฺติยา ปุจฺฉนฺติ, พฺราหฺมณา ปุจฺฉนฺติ, เวสฺสา ปุจฺฉนฺติ, สุทฺทา ปุจฺฉนฺติ, คหฏฺฐา ปุจฺฉนฺติ, ปพฺพชิตา ปุจฺฉนฺติ, อยํ มนุสฺสปุจฺฉา.\csromanspacer{}\csromanspacer{}กตมา อมนุสฺสปุจฺฉา, อมนุสฺสา พุทฺธํ ภควนฺตํ อุปสงฺกมิตฺวา ปญฺหํ ปุจฺฉนฺติ, นาคา ปุจฺฉนฺติ, สุปณฺณา ปุจฺฉนฺติ, ยกฺขา ปุจฺฉนิตฺ, อสุรา ปุจฺฉนฺติ, คนฺธพฺพา ปุจฺฉนฺติ, มหาราชาโน ปุจฺฉนฺติ, อินฺทา ปุจฺฉนฺติ, พฺรหฺมาโน ปุจฺฉนฺติ, เทวตาโย ปุจฺฉนฺติ, อยํ อมนุสฺสปุจฺฉา.\csromanspacer{}\csromanspacer{}กตมา นิมฺมิตปุจฺฉา, ยํ ภควา รูปํ อภินิมฺมินาติ มโนมยํ สพฺพงฺคปจฺจงฺคํ อหีนินฺทฺริยํ, ตํ โส นิมฺมิโต พุทฺธํ ภควนฺตํ อุปสงฺกมิตฺวา ปญฺหํ ปุจฺฉติ, ภควา ตสฺส\footnote{นตฺถิ สี-สฺยา-โปตฺถเกสุ.} วิสชฺเชติ, อยํ นิมฺมิตปุจฺฉา.\csromanspacer{} อิมา ติสฺโส ปุจฺฉา.}

\prose{อปราปิ ติสฺโส ปุจฺฉา อตฺตตฺถปุจฺฉา ปรตฺถปุจฺฉา อุภยตฺถปุจฺฉา.\csromanspacer{}\csromanspacer{}อปราปิ ติสฺโส ปุจฺฉา ทิฏฺฐธมฺมิกตฺถปุจฺฉา สมฺปรายิกตฺถปุจฺฉา ปรมตฺถปุจฺฉา.\csromanspacer{} อปราปิ ติสฺโส ปุจฺฉา อนวชฺชตฺถปุจฺฉา นิกฺกิเลสตฺถปุจฺฉา\footnote{นิกฺเขปตฺถปุจฺฉา (สี, ก)} โวทานตฺถปุจฺฉา.\csromanspacer{}\csromanspacer{}อปราปิ ติสฺโส ปุจฺฉา อตีตปุจฺฉา อนาคตปุจฺฉา ปจฺจุปฺปนฺนปุจฺฉา.\csromanspacer{}\csromanspacer{}อปราปิ ติสฺโส ปุจฺฉา อชฺฌตฺตปุจฺฉา พหิทฺธาปุจฺฉา อชฺฌตฺตพหิทฺธาปุจฺฉา.\csromanspacer{}\csromanspacer{}อปราปิ ติสฺโส ปุจฺฉา กุสลปุจฺฉา อกุสลปุจฺฉา อพฺยากตปุจฺฉา.\csromanspacer{}\csromanspacer{}อปราปิ ติสฺโส ปุจฺฉา ขนฺธปุจฺฉา ธาตุปุจฺฉา อายตนปุจฺฉา.\csromanspacer{}\csromanspacer{}อปราปิ ติสฺโส ปุจฺฉา สติปฏฺฐานปุจฺฉา สมฺมปฺปธานปุจฺฉา อิทฺธิปาทปุจฺฉา.\csromanspacer{}\csromanspacer{}อปราปิ ติสฺโส ปุจฺฉา อินฺทฺริยปุจฺฉา พลปุจฺฉา โพชฺฌงฺคปุจฺฉา.\csromanspacer{}\csromanspacer{}อปราปิ ติสฺโส ปุจฺฉา มคฺคปุจฺฉา ผลปุจฺฉา นิพฺพานปุจฺฉา.}

\prose{\textbf{ปุจฺฉามิ ตนฺ}ติ ปุจฺฉามิ ตํ, ยาจามิ ตํ, อชฺเฌสามิ ตํ, ปสาเทมิ ตํ, “กถยสฺสุ เม”ติ ปุจฺฉามิ ตํ, \textbf{อาทิจฺจพนฺธู}ติ \textbf{อาทิจฺโจ} วุจฺจติ สูริโย\footnote{สุริโย (สี, สฺยา)}.\csromanspacer{} สูริโย โคตโม โคตฺเตน, ภควาปิ โคตโม โคตฺเตน, ภควา สูริยสฺส โคตฺตญาตโก โคตฺตพนฺธุ, ตสฺมา พุทฺโธ อาทิจฺจพนฺธูติ ปุจฺฉามิ ตํ อาทิจฺจพนฺธุ.}

\prose{\textbf{วิเวกํ สนฺติปทญฺจ มเหสี}ติ \textbf{วิเวกา}ติ ตโย \textbf{วิเวกา} กายวิเวโก จิตฺตวิเวโก อุปธิวิเวโก.\csromanspacer{} กตโม กายวิเวโก,}

\vspace{\tipitakaparskip}
\csromancenter{ตุวฏกสุตฺตนิทฺเทส อิธ ภิกฺขุ วิวิตฺตํ เสนาสนํ ภชติ อรญฺญํ รุกฺขมูลํ ปพฺพตํ กนฺทรํ คิริคุหํ สุสานํ วนปตฺถํ อพฺโภกาสํ ปลาลปุญฺชํ.\csromanspacer{} กาเยน วิวิตฺเตน วิหรติ โส เอโก คจฺฉติ, เอโก ติฏฺฐติ, เอโก นิสีทติ, เอโก เสยฺยํ กปฺเปติ, เอโก คามํ ปิณฺฑาย ปวิสติ, เอโก ปฏิกฺกมติ, เอโก รโห นิสีทติ, เอโก จงฺกมํ อธิฏฺฐาติ, เอโก จรติ วิหรติ อิริยติ วตฺตติ ปาเลติ ยเปติ ยาเปติ.\csromanspacer{} อยํ กายวิเวโก.}
\prose{\csromanfolio{265}กตโม จิตฺตวิเวโก, ปฐมํ ฌานํ สมาปนฺนสฺส นีวรเณหิ จิตฺตํ วิวิตฺตํ โหติ, ทุติยํ ฌานํ สมาปนฺนสฺส วิตกฺกวิจาเรหิ จิตฺตํ วิวิตฺตํ โหติ, ตติยํ ฌานํ สมาปนฺนสฺส ปีติยา จิตฺตํ วิวิตฺตํ โหติ, จตุตฺถํ ฌานํ สมาปนฺนสฺส สุขทุกฺเขหิ จิตฺตํ วิวิตฺตํ โหติ.\csromanspacer{} อากาสานญฺจายตนํ สมาปนฺนสฺส รูปสญฺญาย ปฏิฆสญฺญาย นานตฺตสญฺญาย จิตฺตํ วิวิตฺตํ โหติ, วิญฺญาณญฺจายตนํ สมาปนฺนสฺส อากาสานญฺจยตนสญฺญาย จิตฺตํ วิวิตฺตํ โหติ, อากิญฺจญฺญายตนํ สมาปนฺนสฺส วิญฺญาณญฺจายตนสญฺญาย จิตฺตํ วิวิตฺตํ โหติ, เนวสญฺญานาสญฺญายตนํ สมาปนฺนสฺส อากิญฺจญฺญายตนสญฺญาย จิตฺตํ วิวิตฺตํ โหติ.\csromanspacer{} โสตาปนฺนสฺส สกฺกายทิฏฺฐิยา วิจิกิจฺฉาย สีลพฺพตปรามาสา ทิฏฺฐานุสยา วิจิกิจฺฉานุสยา, ตเทกฏฺเฐหิ จ กิเลเสหิ จิตฺตํ วิวิตฺตํ โหติ, สกทาคามิสฺส โอฬาริกา กามราคสญฺโญชนา ปฏิฆสญฺโญชนา โอฬาริกา กามราคานุสยา ปฏิฆานุสยา, ตเทกฏฺเฐหิ จ กิเลเสหิ จิตฺตํ วิวิตฺตํ โหติ, อนาคามิสฺส อณุสหคตา กามราคสญฺโญชนา ปฏิฆสญฺโญชนา อณุสหคตา กามราคานุสยา ปฏิฆานุสยา, ตเทกฏฺเฐหิ จ กิเลเสหิ จิตฺตํ วิวิตฺตํ โหติ, อรหโต รูปราคา อรูปราคา มานา อุทฺธจฺจา อวิชฺชาย มานานุสยา ภวราคานุสยา อวิชฺชานุสยา, ตเทกฏฺเฐหิ จ กิเลเสหิ พหิทฺธา จ สพฺพนิมิตฺเตหิ จิตฺตํ วิวิตฺตํ โหติ.\csromanspacer{} อยํ จิตฺตวิเวโก.}

\prose{กตโม \textbf{อุปธิ}วิเวโก, \textbf{อุปธิ} วุจฺจนฺติ กิเลสา จ ขนฺธา จ อภิสงฺขารา จ.\csromanspacer{} อุปธิวิเวโก วุจฺจติ อมตํ นิพฺพานํ.\csromanspacer{} โย โส สพฺพสงฺขารสมโถ สพฺพูปธิปฏินิสฺสคฺโค ตณฺหกฺขโย วิราโค นิโรโธ นิพฺพานํ, อยํ \textbf{อุปธิ}วิเวโก.\csromanspacer{} กายวิเวโก จ วิเวกฏฺฐกายานํ\footnote{วูปกฏฺฐกายานํ (สี)} เนกฺขมฺมาภิรตานํ, จิตฺตวิเวโก จ ปริสุทฺธจิตฺตานํ ปรมโวทานปตฺตานํ, \textbf{อุปธิ}วิเวโก \csromanfolio{266}จ นิรุปธีนํ ปุคฺคลานํ วิสงฺขารคตานํ\textbf{.}\csromanspacer{} \textbf{สนฺตี}ติ เอเกน อากาเรน สนฺติปิ สนฺติปทมฺปิ ตํเยว อมตํ นิพฺพานํ\textbf{.}\csromanspacer{} โย โส สพฺพสงฺขารสมโถ สพฺพูปธิปฏินิสฺสคฺโค ตณฺหกฺขโย คิราโค นิโรโธ นิพฺพานํ\textbf{.}\csromanspacer{} วุตฺตํ เหตํ ภควตา\csromandash{} * สนฺตเมตํ ปทํ, ปณีตเมตํ ปทํ, ยทิทํ สพฺพสงฺขารสมโถ สพฺพูปธิปฏินิสฺสคฺโค ตณฺหกฺขโย วิราโค นิโรโธ นิพฺพานนฺติ\textbf{.}\csromanspacer{} อถ อปเรน อากาเรน เย ธมฺมา สนฺตาธิคมาย สนฺติผุสนาย สนฺติสจฺฉิกิริยาย สํวตฺตนฺติ\textbf{.}\csromanspacer{} เสยฺยถิทํ, จตฺตาโร สติปฏฺฐานา จตฺตาโร สมฺมปฺปธานา จตฺตาโร อิทฺธิปาทา ปญฺจินฺทฺริยานิ ปญฺจ พลานิ สตฺต โพชฺฌงฺคา อริโย อฏฺฐงฺคิโก มคฺโค, อิเม วุจฺจนฺติ สนฺติปทํ ตาณปทํ เลณปทํ สรณปทํ อภยปทํ อจฺจุตปทํ อมตปทํ นิพฺพานปทํ\textbf{.}}

\prose{\textbf{มเหสี}ติ กิํ มเหสิ ภควา.\csromanspacer{} มหนฺตํ สีลกฺขนฺธํ เอสี คเวสี ปริเยสีติ มเหสิ, มหนฺตํ สมาธิกฺขนฺธํ.\csromanspacer{} มหนฺตํ ปญฺญากฺขนฺธํ.\csromanspacer{} มหนฺตํ วิมุตฺติกฺขนฺธํ.\csromanspacer{} มหนฺตํ วิมุตฺติญาณทสฺสนกฺขนฺธํ เอสี คเวสี ปริเยสีติ มเหสิ, มหโต ตโมกายสฺส ปทาลนํ, มหโต วิปลฺลาสสฺส เภทนํ, มหโต ตณฺหาสลฺลสฺส อพฺพหนํ,1 มหโต ทิฏฺฐิสํฆาตสฺส วินิเวฐนํ, มหโต มานธชสฺส ปปาตนํ\footnote{ปวาหนํ (สฺยา)} มหโต อภิสงฺขารสฺส วูปสมํ, มหโต โอฆสฺส นิตฺถรณํ, มหโต ภารสฺส นิกฺเขปนํ, มหโต สํสารวฏฺฏสฺส อุปจฺเฉทํ, มหโต สนฺตาปสฺส นิพฺพาปนํ, มหโต ปริฬาหสฺส ปฏิปสฺสทฺธิํ, มหโต ธมฺมธชสฺส อุสฺสาปนํ เอสี คเวสี ปริเยสีติ มเหสิ, มหนฺเต สติปฏฺฐาเน มหนฺเต สมฺมปฺปธาเน มหนฺเต อิทฺธิปาเท มหนฺตานิ อินฺทฺริยานิ มหนฺตานิ พลานิ มหนฺเต โพชฺฌงฺเค มหนฺตํ อริยํ อฏฺฐงฺคิกํ มคฺคํ มหนฺตํ ปรมตฺถํ อมตํ นิพฺพานํ เอสี คเวสี ปริเยสีติ มเหสิ, มเหสกฺเขหิ วา สตฺเตหิ เอสิโต คเวสิโต ปริเยสิโต “กหํ พุทฺโธ กหํ ภควา กหํ เทวเทโว กหํ นราสโภ”ติ \textbf{มเหสี}ติ วิเวกํ สนฺติปทญฺจ มเหสิ.}

\prose{\textbf{กถํ ทิสฺวา นิพฺพาติ ภิกฺขู}ติ กถํ ทิสฺวา ปสฺสิตฺวา ตุลยิตฺวา ตีรยิตฺวา วิภาวยิตฺวา วิภูตํ กตฺวา อตฺตโน ราคํ นิพฺพาเปติ, โทสํ นิพฺพาเปติ, โมหํ นิพฺพาเปติ, โกธํ.\csromanspacer{} อุปนาหํ.\csromanspacer{} มกฺขํ.\csromanspacer{} ปฬาสํ.\csromanspacer{} อิสฺสํ.\csromanspacer{} มจฺฉริยํ.\csromanspacer{} มายํ.\csromanspacer{} สาเฐยฺยํ.\csromanspacer{} ถมฺภํ.\csromanspacer{} สารมฺภํ.\csromanspacer{} มานํ.\csromanspacer{} อติมานํ.\csromanspacer{} มทํ.\csromanspacer{} ปมาทํ.\csromanspacer{} สพฺเพ กิเลเส.\csromanspacer{} สพฺเพ ทุจฺจริเต.\csromanspacer{} สพฺเพ ทรเถ.\csromanspacer{} สพฺเพ ปริฬาเห.\csromanspacer{} สพฺเพ}

\vspace{\tipitakaparskip}
\csromancenter{ตุวฏกสุตฺตนิทฺเทส สนฺตาเป\textbf{.}\csromanspacer{} สพฺพากุสลาภิสงฺขาเร นิพฺพาเปติ สเมติ อุปสเมติ วูปสเมติ ปฏิปสฺสมฺเภติ\textbf{.}\csromanspacer{}\csromanspacer{}\textbf{ภิกฺขู}ติ ปุถุชฺชนกลฺยาณโก วา ภิกฺขุ เสโข วา ภิกฺขูติ กถํ ทิสฺวา นิพฺพาติ ภิกฺขุ\textbf{.}}
\prose{\csromanfolio{267}\textbf{อนุปาทิยาโน โลกสฺมิํ กิญฺจี}ติ จตูหิ อุปาทาเนหิ อนุปาทิยมาโน อคณฺหยมาโน อปรามสมาโน อนภินิวิสมาโน\textbf{.}\csromanspacer{} \textbf{โลกสฺมินฺ}ติ อปายโลเก มนุสฺสโลเก เทวโลเก, ขนฺธโลเก ธาตุโลเก อายตนโลเก\textbf{.}\csromanspacer{} \textbf{กิญฺจี}ติ กิญฺจิ รูปคตํ เวทนาคตํ สญฺญาคตํ สงฺขารคตํ วิญฺญาณคตนฺติ อนุปาทิยาโน โลกสฺมิํ กิญฺจิ\textbf{.}\csromanspacer{} เตนาห โส นิมฺมิโต\csromandash{}}

\csromangathameasure{ปุจฺฉามิ ตํ อาทิจฺจพนฺธุ, \\ วิเวกํ สนฺติปทญฺจ มเหสิ\textbf{.} \\ กถํ ทิสฺวา นิพฺพาติ ภิกฺขุ,}
\csromangathagroup{\csromangathastanza{ปุจฺฉามิ ตํ อาทิจฺจพนฺธุ, \\ วิเวกํ สนฺติปทญฺจ มเหสิ\textbf{.} \\ กถํ ทิสฺวา นิพฺพาติ ภิกฺขุ,}}

\prose{\textbf{อนุปาทิยาโน โลกสฺมิํ กิญฺจี}ติ\textbf{.}}

\csromanhangingbegin
\hangingheaditem{151}{มูลํ ปปญฺจสงฺขาย, (อิติ \textbf{ภควา},)}
\hangingbody{\textbf{มนฺตา อสฺมีติ สพฺพมุปรุนฺเธ}1\textbf{.}}
\hangingbody{\textbf{ยา กาจิ ตณฺหา อชฺฌตฺตํ,}}
\hangingbody{\textbf{ตาสํ วินยา สทา สโต สิกฺเข.}}
\csromanhangingend

\prose{\textbf{มูลํ ปปญฺจสงฺขาย, (อิติ ภควา,) มนฺตา อสฺมีติ สพฺพมุปรุนฺเธ}ติ ปปญฺจาเยว ปปญฺจสงฺขา, ตณฺหาปปญฺจสงฺขา ทิฏฺฐิปปญฺจสงฺขา\textbf{.}\csromanspacer{} กตมํ ตณฺหาปปญฺจสฺส มูลํ, อวิชฺชามูลํ อโยนิโสมนสิกาโร มูลํ, อสฺมิมาโน มูลํ, อหิริกํ มูลํ, อโนตฺตปฺปํ มูลํ, อุทฺธจฺจํ มูลํ, อิทํ ตณฺหาปปญฺจสฺส มูลํ\textbf{.}\csromanspacer{} กตมํ ทิฏฺฐิปปญฺจสฺส มูลํ, อวิชฺชามูลํ, อโยนิโสมนสิกาโร มูลํ, อสฺมิมาโน มูลํ, อหิริกํ มูลํ, อโนตฺตปฺปํ มูลํ, อุทฺธจฺจํ มูลํ, อิทํ ทิฏฺฐิปปญฺจสฺส มูลํ\textbf{.}}

\prose{\textbf{ภควา}ติ คารวาธิวจนํ\textbf{.}\csromanspacer{} อปิ จ ภคฺคราโคติ \textbf{ภควา}, ภคฺคโทโสติ \textbf{ภควา}, ภคฺคโมโหติ \textbf{ภควา}, ภคฺคมาโนติ \textbf{ภควา}, ภคฺคทิฏฺฐีติ \textbf{ภควา}, ภคฺคกณฺฑโกติ \textbf{ภควา}, ภคฺคกิเลโสติ \textbf{ภควา}, ภชิ วิภชิ ปวิภชิ ธมฺมรตนนฺติ \textbf{ภควา}, ภวานํ อนฺตกโรติ \textbf{ภควา}, ภาวิตกาโย ภาวิตสีโล ภาวิตจิตฺโต ภาวิตปญฺโญติ \textbf{ภควา}, ภชิ วา \csromanfolio{268}ภควา อรญฺญวนปตฺถานิ ปนฺตานิ เสนาสนานิ อปฺปสทฺทานิ อปฺปนิคฺโฆสานิ วิชนวาตานิ มนุสฺสราหสฺเสยฺยกานิ ปฏิสลฺลานสารุปฺปานีติ ภควา, ภาคี วา ภควา จีวรปิณฺฑปาตเสนาสนคิลานปจฺจยเภสชฺช- ปริกฺขารานนฺติ ภควา, ภาคี วา ภควา อตฺถรสสฺส ธมฺมรสสฺส วิมุตฺติรสสฺส อธิสีลสฺส อธิจิตฺตสฺส อธิปญฺญายาติ ภควา, ภาคี วา ภควา จตุนฺนํ ฌานานํ จตุนฺนํ อปฺปมญฺญานํ จตุนฺนํ อรูปสมาปตฺตีนนฺติ ภควา, ภาคี วา ภควา อฏฺฐนฺนํ วิโมกฺขานํ อฏฺฐนฺนํ อภิภายตนานํ นวนฺนํ อนุปุพฺพวิหารสมาปตฺตีนนฺติ ภควา, ภาคี วา ภควา ทสนฺนํ สญฺญาภาวนานํ ทสนฺนํ กสิณสมาปตฺตีนํ อานาปานสฺสติสมาธิสฺส อสุภสมาปตฺติยาติ ภควา, ภาคี วา ภควา จตุนฺนํ สติปฏฺฐานานํ จตุนฺนํ สมฺมปฺปธานานํ จตุนฺนํ อิทฺธิปาทานํ ปญฺจนฺนํ อินฺทฺริยานํ ปญฺจนฺนํ พลานํ สตฺตนฺนํ โพชฺฌงฺคานํ อริยสฺส อฏฺฐงฺคิสฺส มคฺคสฺสาติ ภควา, ภาคี วา ภควา ทสนฺนํ ตถาคตพลานํ จตุนฺนํ เวสารชฺชานํ จตุนฺนํ ปฏิสมฺภิทานํ ฉนฺนํ อภิญฺญานํ ฉนฺนํ พุทฺธธมฺมานนฺติ ภควา, ภควาติ เนตํ นามํ มาตรา กตํ, น ปิตรา กตํ, น ภาตรา กตํ, น ภคินิยา กตํ, น มิตฺตามจฺเจหิ กตํ, น ญาติสาโลหิเตหิ กตํ, น สมณพฺราหฺมเณหิ กตํ, น เทวตาหิ กตํ, วิโมกฺขนฺติกเมตํ พุทฺธานํ ภควนฺตานํ โพธิยา มูเล สห สพฺพญฺญุตญาณสฺส ปฏิลาภา สจฺฉิกา ปญฺญตฺติ ยทิทํ ภควาติ มูลํ ปปญฺจสงฺขาย อิติ ภควา.}

\prose{\textbf{มนฺตา อสฺมีติ สพฺพมุปรุนฺเธ}ติ \textbf{มนฺตา} วุจฺจติ ปญฺญา.\csromanspacer{} ยา ปญฺญา ปชานนา ฯเปฯ อโมโห ธมฺมวิจโย สมฺมาทิฏฺฐิ.\csromanspacer{} \textbf{อสฺมี}ติ รูเป “อสฺมี”ติ มาโน “อสฺมี”ติ ฉนฺโท “อสฺมี”ติ อนุสโย.\csromanspacer{} เวทนาย.\csromanspacer{} สญฺญาย.\csromanspacer{} สงฺขาเรสุ.\csromanspacer{} วิญฺญาเณ “อสฺมี”ติ มาโน “อสฺมี”ติ ฉนฺโท “อสฺมี”ติ อนุสโยติ มูลํ ปปญฺจสงฺขาย อิติ ภควา.\csromanspacer{} \textbf{มนฺตา อสฺมีติ สพฺพมุปรุนฺเธ}ติ ปปญฺจสงฺขาย มูลญฺจ อสฺมิมานญฺจ \textbf{มนฺตา}ย สพฺพํ รุนฺเธยฺย อุปรุนฺเธยฺย นิโรเธยฺย วูปสเมยฺย อตฺถงฺคเมยฺย ปฏิปสฺสมฺเภยฺยาติ มูลํ ปปญฺจสงฺขาย อิติ ภควา, \textbf{มนฺตา} อสฺมีติ สพฺพมุปรุนฺเธ.}

\prose{\textbf{ยา กาจิ ตณฺหา อชฺฌตฺตนฺ}ติ \textbf{ยา กาจี}ติ สพฺเพน สพฺพํ สพฺพถา สพฺพํ อเสสํ นิสฺเสสํ ปริยาทิยนวจนเมตํ \textbf{ยา กาจี}ติ.\csromanspacer{} \textbf{ตณฺหา}ติ รูปตณฺหา ฯเปฯ ธมฺมตณฺหา.\csromanspacer{} \textbf{อชฺฌตฺตนฺ}ติ อชฺฌตฺตสมุฏฺฐานา วา\footnote{อชฺฌตฺตํ สมุฏฺฐาติ (สฺยา)} สา ตณฺหาติ}

\vspace{\tipitakaparskip}
\csromancenter{ตุวฏกสุตฺตนิทฺเทส อชฺฌตฺตํ.\csromanspacer{} อถ วา \textbf{อชฺฌตฺติกํ} วุจฺจติ จิตฺตํ.\csromanspacer{} ยํ จิตฺตํ มโน มานสํ หทยํ ปณฺฑรํ มโน มนายตนํ มนินฺทฺริยํ วิญฺญาณํ วิญฺญาณกฺขนฺโธ ตชฺชา มโนวิญฺญาณธาตุ.\csromanspacer{} จิตฺเตน สา ตณฺหา สหคตา สหชาตา สํสฏฺฐา สมฺปยุตฺตา เอกุปฺปาทา เอกนิโรธา เอกวตฺถุกา เอการมฺมณาติปิ อชฺฌตฺตนฺติ ยา กาจิ ตณฺหา อชฺฌตฺตํ.}
\prose{\csromanfolio{269}\textbf{ตาสํ วินยา สทา สโต สิกฺเข}ติ \textbf{สทา}ติ \textbf{สทา} สพฺพทา สพฺพกาลํ นิจฺจกาลํ ธุวกาลํ สตตํ สมิตํ อพฺโพกิณฺณํ โปงฺขานุโปงฺขํ อุทกูมิคชาตํ อวีจิ สนฺตติ สหิตํ ผุสิตํ ปุเรภตฺตํ ปจฺฉาภตฺตํ ปุริมํ ยามํ มชฺฌิมํ ยามํ ปจฺฉิมํ ยามํ กาเฬ ชุเณฺห วสฺเส เหมนฺเต คิเมฺห ปุริเม วโยขนฺเธ มชฺฌิเม วโยขนฺเธ ปจฺฉิเม วโยขนฺเธ.\csromanspacer{}\csromanspacer{}\textbf{สโต}ติ จตูหิ การเณหิ สโต.\csromanspacer{} กาเย กายานุปสฺสนาสติปฏฺฐานํ ภาเวนฺโต สโต, เวทนาสุ ฯเปฯ.\csromanspacer{} จิตฺเต ฯเปฯ.\csromanspacer{} ธมฺเมสุ ธมฺมานุปสฺสนาสติปฏฺฐานํ ภาเวนฺโต สโต, อปเรหิปิ จตูหิ การเณหิ สโต, อสติ ปริวชฺชนาย สโต, สติกรณียานํ ธมฺมานํ กตตฺตา สโต, สติปฏิปกฺขานํ ธมฺมานํ หตตฺตา สโต, สตินิมิตฺตานํ อสมฺมุฏฺฐตฺตา สโต.\csromanspacer{} อปเรหิปิ จตูหิ การเณหิ สโต, สติยา สมนฺนาคตตฺตา สโต, สติยา วสิตตฺตา สโต, สติยา ปาคุญฺญตาย สโต, สติยา อปจฺโจโรหนตาย สโต.\csromanspacer{} อปเรหิปิ จตูหิ การเณหิ สโต, สตฺตตฺตา สโต, สนฺตตฺตา สโต, สมิตตฺตา สโต, สนฺตธมฺมสมนฺนาคตตฺตา สโต, พุทฺธานุสฺสติยา สโต, ธมฺมานุสฺสติยา สโต, สํฆานุสฺสติยา สโต, สีลานุสฺสติยา สโต, จาคานุสฺสติยา สโต, เทวตานุสฺสติยา สโต, อานาปานสฺสติยา สโต, มรณสฺสติยา สโต, กายคตาสติยา สโต, อุปสมานุสฺสติยา สโต, ยา สติ อนุสฺสติ ปฏิสฺสติ สติ สรณตา ธารณตา อปิลาปนตา อสมฺมุสฺสนตา สติ สตินฺทฺริยํ สติพลํ สมฺมาสติ สติสมฺโพชฺฌงฺโค เอกายนมคฺโค, อยํ วุจฺจติ สติ.\csromanspacer{} อิมาย สติยา อุเปโต สมุเปโต อุปคโต สมุปคโต อุปปนฺโน สมุปปนฺโน สมนฺนาคโต โส วุจฺจติ สโต.}

\prose{\textbf{สิกฺเข}ติ ติสฺโส สิกฺขา อธิสีลสิกฺขา อธิจิตฺตสิกฺขา อธิปญฺญาสิกฺขา.\csromanspacer{} กตมา อธิสีลสิกฺขา, อิธ ภิกฺขุ สีลวา โหติ, ปาติโมกฺขสํวรสํวุโต วิหรติ อาจารโคจรสมฺปนฺโน อณุมตฺเตสุ \csromanfolio{270}วชฺเชสุ ภยทสฺสาวี, สมาทาย สิกฺขติ สิกฺขาปเทสุ, ขุทฺทโก สีลกฺขนฺโธ.\csromanspacer{} มหนฺโต สีลกฺขนฺโธ.\csromanspacer{} สีลํ ปติฏฺฐา อาทิ จรณํ สํยโม สํวโร มุขํ ปมุขํ กุสลานํ ธมฺมานํ สมาปตฺติยา.\csromanspacer{} อยํ อธิสีลสิกฺขา.}

\prose{กตมา อธิจิตฺตสิกฺขา, อิธ ภิกฺขุ วิวิจฺเจว กาเมหิ วิวิจฺจ อกุสเลหิ ธมฺเมหิ สวิตกฺกํ สวิจารํ วิเวกชํ ปีติสุขํ ปฐมํ ฌานํ อุปสมฺปชฺช วิหรติ.\csromanspacer{} วิตกฺกวิจารานํ วูปสมา อชฺฌตฺตํ สมฺปสาทนํ เจตโส เอโกทิภาวํ อวิตกฺกํ อวิจารํ สมาธิชํ ปีติสุขํ ทุติยํ ฌานํ อุปสมฺปชฺช วิหรติ.\csromanspacer{} ปีติยา จ วิราคา อุเปกฺขโก จ วิหรติ สโต จ สมฺปชาโน สุขญฺจ กาเยน ปฏิสํเวเทติ, ยํ ตํ อริยา อาจิกฺขนฺติ “อุเปกฺขโก สติมา สุขวิหารี”ติ ตติยํ ฌานํ อุปสมฺปชฺช วิหรติ.\csromanspacer{} สุขสฺส จ ปหานา ทุกฺขสฺส จ ปหานา ปุพฺเพว โสมนสฺสโทมนสฺสานํ อตฺถงฺคมา อทุกฺขมสุขํ อุเปกฺขาสติปาริสุทฺธิํ จตุตฺถํ ฌานํ อุปสมฺปชฺช วิหรติ.\csromanspacer{} อยํ อธิจิตฺตสิกฺขา.}

\prose{กตมา อธิปญฺญาสิกฺขา, อิธ ภิกฺขุ ปญฺญวา โหติ อุทยตฺถคามินิยา ปญฺญาย สมนฺนาคโต อริยาย นิพฺเพธิกาย สมฺมาทุกฺขกฺขยคามินิยา.\csromanspacer{} โส “อิทํ ทุกฺขนฺ”ติ ยถาภูตํ ปชานาติ, “อยํ ทุกฺขสมุทโย”ติ ยถาภูตํ ปชานาติ, “อยํ ทุกฺขนิโรโธ”ติ ยถาภูตํ ปชานาติ, “อยํ ทุกฺขนิโรธคามินี ปฏิปทา”ติ ยถาภูตํ ปชานาติ, “อิเม อาสวา”ติ ยถาภูตํ ปชานาติ, “อยํ อาสวสมุทโย”ติ ยถาภูตํ ปชานาติ, “อยํ อาสวนิโรโธ”ติ ยถาภูตํ ปชานาติ, “อยํ อาสวนิโรธคามินี ปฏิปทา”ติ ยถาภูตํ ปชานาติ.\csromanspacer{} อยํ อธิปญฺญาสิกฺขา.}

\prose{\textbf{ตาสํ วินยา สทา สโต สิกฺเข}ติ ตาสํ ตณฺหานํ วินยาย ปฏิวินยาย ปหานาย วูปสมาย ปฏินิสฺสคฺคาย ปฏิปสฺสทฺธิยา อธิสีลมฺปิ สิกฺเขยฺย, อธิจิตฺตมฺปิ สิกฺเขยฺย, อธิปญฺญมฺปิ สิกฺเขยฺย.\csromanspacer{} อิมา ติสฺโส สิกฺขาโย อาวชฺชนฺโต สิกฺเขยฺย, ปชานนฺโต สิกฺเขยฺย, ปสฺสนฺโต สิกฺเขยฺย, ปจฺจเวกฺขนฺโต สิกฺเขยฺย, จิตฺตํ อธิฏฺฐหนฺโต สิกฺเขยฺย, สทฺธาย อธิมุจฺจนฺโต สิกฺเขยฺย, วีริยํ ปคฺคณฺหนฺโต สิกฺเขยฺย, สติํ อุปฏฺฐเปนฺโต สิกฺเขยฺย, จิตฺตํ สมาทหนฺโต สิกฺเขยฺย, ปญฺญาย ปชานนฺโต สิกฺเขยฺย, อภิญฺเญยฺยํ อภิชานนฺโต สิกฺเขยฺย, ปริญฺเญยฺยํ ปริชานนฺโต สิกฺเขยฺย, ปหาตพฺพํ ปชหนฺโต สิกฺเขยฺย, ภาเวตพฺพํ ภาเวนฺโต สิกฺเขยฺย,}

\vspace{\tipitakaparskip}
\csromancenter{ตุวฏกสุตฺตนิทฺเทส สจฺฉิกาตพฺพํ สจฺฉิกโรนฺโต สิกฺเขยฺย อาจเรยฺย สมาจเรยฺย สมาทาย วตฺเตยฺยาติ ตาสํ วินยา สทา สโต สิกฺเข.\csromanspacer{} เตนาห ภควา\csromandash{}}
\prose{\csromanfolio{271}มูลํ ปปญฺจสงฺขาย, (อิติ ภควา,)}

\prose{มนฺตา อสฺมีติ สพฺพมุปรุนฺเธ.}

\csromangathameasure{ยา กาจิ ตณฺหา อชฺฌตฺตํ, \\ ตาสํ วินยา สทา สโต สิกฺเขติ.}
\csromangathagroup{\csromangathastanza{ยา กาจิ ตณฺหา อชฺฌตฺตํ, \\ ตาสํ วินยา สทา สโต สิกฺเขติ.}}

\proseitem{152}{\csromansymbolfootnote{*}{ขุ 1. 421 ปิฏฺเฐ.}ยํ กิญฺจิ ธมฺมมภิชญฺญา, อชฺฌตฺตํ อถ วาปิ พหิทฺธา.}

\csromangathameasure{น เตน ถามํ กุพฺเพถ, น หิ สา นิพฺพุติ สตํ วุตฺตา.}
\csromangathagroup{\csromangathastanza{น เตน ถามํ กุพฺเพถ, น หิ สา นิพฺพุติ สตํ วุตฺตา.}}

\prose{\textbf{ยํ กิญฺจิ ธมฺมมภิชญฺญา อชฺฌตฺตนฺ}ติ ยํ กิญฺจิ อตฺตโน คุณํ ชาเนยฺย กุสเล วา ธมฺเม อพฺยากเต วา ธมฺเม.\csromanspacer{} กตเม อตฺตโน คุณา, อุจฺจา กุลา ปพฺพชิโต วา อสฺสํ\footnote{อสฺส (สฺยา)}, มหาโภคกุลา ปพฺพชิโต วา อสฺสํ, อุฬารโภคกุลา ปพฺพชิโต วา อสฺสํ, ญาโต ยสสฺสี สคหฏฺฐปพฺพชิตานนฺติ วา อสฺสํ, ลาภิมฺหิ จีวร ปิณฺฑปาต เสนาสน คิลานปจฺจย เภสชฺช ปริกฺขารานนฺติ วา อสฺสํ, สุตฺตนฺติโก วา อสฺสํ, วินยธโร วา อสฺสํ, ธมฺมกถิโก วา อสฺสํ, อารญฺญิโก วา อสฺสํ, ปิณฺฑปาติโก วา อสฺสํ, ปํสุกูลิโก วา อสฺสํ, เตจีวริโก วา อสฺสํ, สปทานจาริโก วา อสฺสํ, ขลุปจฺฉาภตฺติโก วา อสฺสํ, เนสชฺชิโก วา อสฺสํ, ยถาสนฺถติโก วา อสฺสํ, ปฐมสฺส ฌานสฺส ลาภีติ วา อสฺสํ, ทุติยสฺส ฌานสฺส ลาภีติ วา อสฺสํ, ตติยสฺส ฌานสฺส ลาภีติ วา อสฺสํ, จตุตฺถสฺส ฌานสฺส ลาภีติ วา อสฺสํ, อากาสานญฺจายตนสมาปตฺติยา ลาภีติ วา อสฺสํ, วิญฺญาณญฺจายตนสมาปตฺติยา.\csromanspacer{} อากิญฺจญฺญายตนสมาปตฺติยา.\csromanspacer{} เนวสญฺญานาสญฺญายตนสมาปตฺติยา ลาภีติ วา อสฺสํ.\csromanspacer{} อิเม วุจฺจนฺติ อตฺตโน คุณา.\csromanspacer{} ยํ กิญฺจิ อตฺตโน คุณํ ชาเนยฺย อาชาเนยฺย วิชาเนยฺย ปฏิวิชาเนยฺย ปฏิวิชฺเฌยฺยาติ ยํ กิญฺจิ ธมฺมมภิชญฺญา อชฺฌตฺตํ.\csromanspacer{} \textbf{อถ วาปิ พหิทฺธา}ติ อุปชฺฌายสฺส วา อาจริยสฺส วา เต คุณา อสฺสูติ\footnote{อสฺสูติ อชฺฌตฺตํ (พหูสุ)} อถ วาปิ พหิทฺธา.}

\prose{\csromanfolio{272}\textbf{น เตน ถามํ กุพฺเพถา}ติ อตฺตโน วา คุเณน ปเรสํ วา คุเณน ถามํ น กเรยฺย, ถมฺภํ น กเรยฺย, มานํ น กเรยฺย, อุนฺนติํ น กเรยฺย, อุนฺนมํ น กเรยฺย, น เตน มานํ ชเนยฺย, น เตน ถทฺโธ อสฺส ปตฺถทฺโธ ปคฺคหิตสิโรติ น เตน ถามํ กุพฺเพถ.}

\prose{\textbf{น หิ สา นิพฺพุติ สตํ วุตฺตา}ติ สตานํ สนฺตานํ สปฺปุริสานํ พุทฺธานํ พุทฺธสาวกานํ ปจฺเจกพุทฺธานํ สา นิพฺพุตีติ น วุตฺตา น ปวุตฺตา น อาจิกฺขิตา น เทสิตา น ปญฺญปิตา น ปฏฺฐปิตา น วิวฏา น วิภตฺตา น อุตฺตานีกตา นปฺปกาสิตาติ น หิ สา นิพฺพุติ สตํ วุตฺตา.\csromanspacer{} เตนาห ภควา\csromandash{}}

\csromangathameasure{ยํ กิญฺจิ ธมฺมมภิชญฺญา, อชฺฌตฺตํ อถ วาปิ พหิทฺธา. \\ น เตน ถามํ กุพฺเพถ, น หิ สา นิพฺพุติ สตํ วุตฺตาติ.}
\csromangathagroup{\csromangathastanza{ยํ กิญฺจิ ธมฺมมภิชญฺญา, อชฺฌตฺตํ อถ วาปิ พหิทฺธา. \\ น เตน ถามํ กุพฺเพถ, น หิ สา นิพฺพุติ สตํ วุตฺตาติ.}}

\proseitem{153}{\csromansymbolfootnote{*}{ขุ 1. 421 ปิฏฺเฐ.}เสยฺโย น เตน มญฺเญยฺย,}

\prose{\textbf{นีเจยฺโย อถ วาปิ สริกฺโข.}}

\csromangathameasure{\textbf{ผุฏฺโฐ อเนกรูเปหิ,} \\ \textbf{นาตุมานํ วิกปฺปยํ ติฏฺเฐ.}}
\csromangathagroup{\csromangathastanza{\textbf{ผุฏฺโฐ อเนกรูเปหิ,} \\ \textbf{นาตุมานํ วิกปฺปยํ ติฏฺเฐ.}}}

\prose{\textbf{เสยฺโย น เตน มญฺเญยฺยา}ติ “เสยฺโยหมสฺมี”ติ อติมานํ น ชเนยฺย ชาติยา วา โคตฺเตน วา โกลปุตฺติเยน วา วณฺณโปกฺขรตาย วา ธเนน วา อชฺเฌเนน วา กมฺมายตเนน วา สิปฺปายตเนน วา วิชฺชาฏฺฐาเนน วา สุเตน วา ปฏิภาเนน วา อญฺญตรญฺญตเรน วา วตฺถุนาติ เสยฺโย น เตน มญฺเญยฺย.}

\prose{\textbf{นีเจยฺโย อถ วาปิ สริกฺโข}ติ “หีโนหมสฺมี”ติ โอมานํ น ชเนยฺย ชาติยา วา โคตฺเตน วา ฯเปฯ อญฺญตรญฺญตเรน วา วตฺถุนา “สทิโสหมสฺมี”ติ มานํ น ชเนยฺย ชาติยา วา โคตฺเตน วา โกลปุตฺติเยน วา วณฺณโปกฺขรตาย วา ธเนน วา อชฺเฌเนน วา กมฺมายตเนน วา สิปฺปายตเนน วา วิชฺชาฏฺฐาเนน วา สุเตน วา ปฏิภาเนน วา อญฺญตรญฺญตเรน วา วตฺถุนาติ นีเจยฺโย อถ วาปิ สริกฺโข.}

\prose{\textbf{ผุฏฺโฐ อเนกรูเปหี}ติ อเนกวิเธหิ อากาเรหิ ผุฏฺโฐ ปเรโต สโมหิโต สมนฺนาคโตติ ผุฏฺโฐ อเนกรูเปติ.}

\csromanheader{2}{14. ตุวฏกสุตฺตนิทฺเทส}
\prose{\csromanfolio{273}\textbf{นาตุมานํ วิกปฺปยํ ติฏฺเฐ}ติ \textbf{อาตุมา} วุจฺจติ อตฺตา.\csromanspacer{} \textbf{อตฺตา}นํ กปฺเปนฺโต วิกปฺเปนฺโต วิกปฺปํ อาปชฺชนฺโต น ติฏฺเฐยฺยาติ นฺ\textbf{อาตุมา}นํ วิกปฺปยํ ติฏฺเฐ.\csromanspacer{} เตนาห ภควา\csromandash{}}

\csromangathameasure{เสยฺโย น เตน มญฺเญยฺย, \\ นีเจยฺโย อถ วาปิ สริกฺโข. \\ ผุฏฺโฐ อเนกรูเปหิ, \\ มฺ\textbf{อาตุมา}นํ วิกปฺปยํ ติฏฺเฐติ.}
\csromangathagroup{\csromangathastanza{เสยฺโย น เตน มญฺเญยฺย, \\ นีเจยฺโย อถ วาปิ สริกฺโข. \\ ผุฏฺโฐ อเนกรูเปหิ, \\ มฺ\textbf{อาตุมา}นํ วิกปฺปยํ ติฏฺเฐติ.}}

\csromanhangingbegin
\hangingheaditem{154}{\textbf{อชฺฌตฺตเมวุ’ปสเม},}
\hangingbody{\textbf{น อญฺญโต ภิกฺขุ สนฺติเมเสยฺย.}}
\hangingbody{\textbf{อชฺฌตฺตํ อุปสนฺตสฺส,}}
\hangingbody{\textbf{นตฺถิ อตฺตา}1 \textbf{กุโต นิรตฺตา}2 \textbf{วา.}}
\csromanhangingend

\prose{\textbf{อชฺฌตฺตเมวุ’ปสเม}ติ อชฺฌตฺตํ ราคํ สเมยฺย, โทสํ สเมยฺย, โมหํ สเมยฺย, โกธํ.\csromanspacer{} อุปนาหํ.\csromanspacer{} มกฺขํ.\csromanspacer{} ปฬาสํ.\csromanspacer{} อิสฺสํ.\csromanspacer{} มจฺฉริยํ.\csromanspacer{} มายํ.\csromanspacer{} สาเฐยฺยํ.\csromanspacer{} ถมฺภํ.\csromanspacer{} สารมฺภํ.\csromanspacer{} มานํ.\csromanspacer{} อติมานํ.\csromanspacer{} มทํ.\csromanspacer{} ปมาทํ.\csromanspacer{} สพฺเพ กิเลเส.\csromanspacer{} สพฺเพ ทุจฺจริเต.\csromanspacer{} สพฺเพ ทรเถ.\csromanspacer{} สพฺเพ ปริฬาเห.\csromanspacer{} สพฺเพ สนฺตาเป.\csromanspacer{} สพฺพากุสลาภิสงฺขาเร สเมยฺย อุปสเมยฺย วูปสเมยฺย นิพฺพาเปยฺย ปฏิปสฺสมฺเภยฺยาติ อชฺฌตฺตเมวุ’ปสเม.}

\prose{\textbf{น อญฺญโต ภิกฺขุ สนฺติเมเสยฺยา}ติ อญฺญโต อสุทฺธิมคฺเคน มิจฺฉาปฏิปทาย อนิยฺยานปเตน อญฺญตฺร สติปฏฺฐาเนหิ อญฺญตฺร สมฺมปฺปธาเนหิ อญฺญตฺร อิทฺธิปาเทหิ อญฺญตฺร อินฺทฺริเยหิ อญฺญตฺร พเลหิ อญฺญตฺร โพชฺฌงฺเคหิ อญฺญตฺร อริยา อฏฺฐงฺคิกา มคฺคา สนฺติํ อุปสนฺติํ วูปสนฺติํ นิพฺพุติํ ปฏิปสฺสทฺธิํ น เอเสยฺย น คเวเสยฺย น ปริเยเสยฺยาติ น อญฺญโต ภิกฺขุ สนฺติเมเสยฺย.}

\prose{\textbf{อชฺฌตฺตํ อุปสนฺตสฺสา}ติ อชฺฌตฺตํ ราคํ สนฺตสฺส โทสํ สนฺตสฺส โมหํ สนฺตสฺส ฯเปฯ สพฺพากุสลาภิสงฺขาเร สนฺตสฺส อุปสนฺตสฺส วูปสนฺตสฺส นิพฺพุตสฺส ปฏิปสฺสทฺธิยาติ อชฺฌตฺตํ อุปสนฺตสฺส.}

\prose{\textbf{นตฺถิ อตฺตา กุโต นิรตฺตา วา}ติ \textbf{นตฺถี}ติ ปฏิกฺเขโป.\csromanspacer{} \textbf{อตฺตา}ติ อตฺตทิฏฺฐิ นตฺถิ.\csromanspacer{} \textbf{นิรตฺตา}ติ อุจฺเฉททิฏฺฐิ นตฺถิ.\csromanspacer{} \textbf{อตฺตา}ติ คหิตํ นตฺถิ.\csromanspacer{} \textbf{นิรตฺตา}ติ มุญฺจิตพฺพํ นตฺถิ.\csromanspacer{} ยสฺสตฺถิ คหิตํ, ตสฺสตฺถิ มุญฺจิตพฺพํ.\csromanspacer{} ยสฺสตฺถิ มุญฺจิตพฺพํ, \csromanfolio{274}ตสฺส คหิตํ, คาหํ มุญฺจนํ สมติกฺกนฺโต อรหา วุทฺธิปาริหานิวีติวตฺโต, โส วุฏฺฐวาโส จิณฺณจรโณ ฯเปฯ.\csromanspacer{} ชาติมรณสํสาโร, นตฺถิ ตสฺส ปุนพฺภโวติ นตฺถิ อตฺตา กุโต นิรตฺตา วา.\csromanspacer{} เตนาห ภควา\csromandash{}}

\csromangathameasure{อชฺฌตฺตเมวุ’ปสเม, \\ น อญฺญโต ภิกฺขุ สนฺติเมเสยฺย. \\ อชฺฌตฺตํ อุปสนฺตสฺส, \\ นตฺถิ อตฺตา กุโต นิรตฺตา วาติ.}
\csromangathagroup{\csromangathastanza{อชฺฌตฺตเมวุ’ปสเม, \\ น อญฺญโต ภิกฺขุ สนฺติเมเสยฺย. \\ อชฺฌตฺตํ อุปสนฺตสฺส, \\ นตฺถิ อตฺตา กุโต นิรตฺตา วาติ.}}

\csromanhangingbegin
\hangingheaditem{155}{มชฺเฌ ยถา สมุทฺทสฺส,}
\hangingbody{อู\textbf{มิ โน ชายตี ฐิโต โหติ.}}
\hangingbody{\textbf{เอวํ ฐิโต อเนชสฺส,}}
\hangingbody{\textbf{อุสฺสทํ ภิกฺขุ น กเรยฺย}1 \textbf{กุหิญฺจิ.}}
\csromanhangingend

\prose{\textbf{มชฺเฌ ยถา สมุทฺทสฺส, อูมิ โน ชายตี ฐิโต โหตี}ติ สมุทฺโท จตุราสีติโยชนสหสฺสานิ อุพฺเพเธน คมฺภีโร.\csromanspacer{} เหฏฺฐา จตฺตารีสโยชนสหสฺสานิ อุทกํ มจฺฉกจฺฉเปหิ กมฺปติ.\csromanspacer{} อุปริ จตฺตารีสโยชนสหสฺสานิ อุทกํ วาเตหิ กมฺปติ.\csromanspacer{} มชฺเฌ จตฺตารีสโยชนสหสฺสานิ อุทกํ น กมฺปติ น วิกมฺปติ น จลติ น เวธติ นปฺปเวธติ น สมฺปเวธติ.\csromanspacer{} อเนริโต อฆฏฺฏิโต อจลิโต อลุฬิโต อภนฺโต วูปสนฺโต, ตตฺร อูมิ โน ชายติ, ฐิโต โหติ สมุทฺโทติ, เอวมฺปิ มชฺเฌ ยถา สมุทฺทสฺส, อู\textbf{มิ โน ชายตี ฐิโต โหติ.}}

\prose{อถ วา สตฺตนฺนํ ปพฺพตานํ อนฺตริกาสุ สตฺต สีทนฺตรา มหาสมุทฺทา\footnote{สีทนฺตรสมุทฺโท (สฺยา)}, ตตฺร อุทกํ น กมฺปติ น วิกมฺปติ น จลติ น เวธติ นปฺปเวธติ น สมฺปเวธติ.\csromanspacer{} อเนริโต อฆฏฺฏิโต อจลิโต อลุฬิโต อภนฺโต วูปสนฺโต, ตตฺร อูมิ โน ชายติ ฐิโต โหติ สมุทฺโทติ, เอวมฺปิ มชฺเฌ ยถา สมุทฺทสฺส, อู\textbf{มิ โน ชายตี ฐิโต โหติ.}}

\prose{\textbf{เอวํ ฐิโต อเนชสฺสา}ติ \textbf{เอวนฺ}ติ โอปมฺมสมฺปฏิปาทนํ.\csromanspacer{} \textbf{ฐิโต}ติ ลาเภปิ น กมฺปติ, อลาเภปิ น กมฺปติ, ยเสปิ น กมฺปติ, อยเสปิ น กมฺปติ, ปสํสายปิ น กมฺปติ, นินฺทายปิ น กมฺปติ, สุเขปิ น กมฺปติ, ทุกฺเขปิ น กมฺปติ น วิกมฺปติ น จลติ น เวธติ นปฺปเวธติ น สมฺปเวธตีติ เอวํ ฐิโต.\csromanspacer{} \textbf{อเนชสฺสา}ติ \textbf{เอชา} วุจฺจติ ตณฺหา.\csromanspacer{} โย}

\vspace{\tipitakaparskip}
\csromancenter{ตุวฏกสุตฺตนิทฺเทส ราโค สาราโค ฯเปฯ อภิชฺฌา โลโภ อกุสลมูลํ.\csromanspacer{} ยสฺเสสา เอชา ตณฺหา ปหีนา อุจฺฉินฺนา สมุจฺฉินฺนา วูปสนฺตา ปฏิปสฺสทฺธา อภพฺพุปฺปตฺติกา ญาณคฺคินา ทฑฺฒา, โส วุจฺจติ อเนโช.\csromanspacer{} เอชาย ปหีนตฺตา อเนโช, โส ลาเภปิ น อิญฺชติ, อลาเภปิ น อิญฺชติ, ยเสปิ น อิญฺชติ, อยเสปิ น อิญฺชติ, ปสํสายปิ น อิญฺชติ, นินฺทายปิ น อิญฺชติ, สุเขปิ น อิญฺชติ,ทุกฺเขปิ น อิญฺชติ, น จลติ น เวธติ นปฺปเวธติ น สมฺปเวธตีติ เอวํ ฐิโต อเนชสฺส.}
\prose{\csromanfolio{275}\textbf{อุสฺสทํ ภิกฺขุ น กเรยฺย กุหิญฺจี}ติ \textbf{อุสฺสทา}ติ สตฺตฺ\textbf{อุสฺสทา}.\csromanspacer{} ราคุสฺสทํ\footnote{ราคุสฺสโท (สฺยา)} โทสุสฺสทํ โมหุสฺสทํ มานุสฺสทํ ทิฏฺฐุสฺสทํ กิเลสุสฺสทํ กมฺมุสฺสทํ น กเรยฺย น ชเนยฺย น สญฺชเนยฺย น นิพฺพตฺเตยฺย นาภินิพฺพตฺเตยฺย.\csromanspacer{} \textbf{กุหิญฺจี}ติ กุหิญฺจิ กิสฺมิญฺจิ กตฺถจิ อชฺฌตฺตํ วา พหิทฺธา วา อชฺฌตฺตพหิทฺธา วาติ อุสฺสทํ ภิกฺขุ น กเรยฺย กุหิญฺจิ.\csromanspacer{} เตนาห ภควา\csromandash{}}

\csromangathameasure{มชฺเฌ ยถา สมุทฺทสฺส, \\ อูมิ โน ชายตี ฐิโต โหติ. \\ เอวํ ฐิโต อเนชสฺส,}
\csromangathagroup{\csromangathastanza{มชฺเฌ ยถา สมุทฺทสฺส, \\ อูมิ โน ชายตี ฐิโต โหติ. \\ เอวํ ฐิโต อเนชสฺส,}}

\prose{\textbf{อุสฺสทํ ภิกฺขุ น กเรยฺย กุหิญฺจี}ติ.}

\proseitem{156}{\csromansymbolfootnote{*}{ขุ 1. 422 ปิฏฺเฐ.}อกิตฺตยี วิวฏจกฺขุ,}

\prose{\textbf{สกฺขิธมฺมํ ปริสฺสยวินยํ.}}

\csromangathameasure{\textbf{ปฏิปทํ วเทหิ ภทฺทนฺเต,} \\ \textbf{ปาติโมกฺขํ อถ วาปิ สมาธิํ.}}
\csromangathagroup{\csromangathastanza{\textbf{ปฏิปทํ วเทหิ ภทฺทนฺเต,} \\ \textbf{ปาติโมกฺขํ อถ วาปิ สมาธิํ.}}}

\prose{\textbf{อกิตฺตยี วิวฏจกฺขู}ติ \textbf{อกิตฺตยี}ติ อกิตฺตยิ ปริกิตฺตยิ อาจิกฺขิ เทเสสิ ปญฺญเปสิ ปฏฺฐเปสิ วิวริ วิภชิ อุตฺตานิมกาสิ ปกาเสสีติ อกิตฺตยิ\footnote{อกิตฺตยีติ กิตฺติตํ อาจิกฺขิตํ เทสิตํ ปญฺญปิตํ ปฏฺฐปิตํ วิวฏํ วิภตฺตํ อุตฺตานีกตํ ปกาสิตนฺติ อกิตฺตยิ (สี, ก)}.\csromanspacer{} \textbf{วิวฏจกฺขู}ติ + ภควา ปญฺจหิ จกฺขูหิ วิวฏจกฺขุ, มํสจกฺขุนาปิ วิวฏจกฺขุ, ทิพฺเพน จกฺขุนาปิ วิวฏจกฺขุ, ปญฺญาจกฺขุนาปิ วิวฏจกฺขุ, พุทฺธจกฺขุนาปิ วิวฏจกฺขุ, สมนฺตจกฺขุนาปิ วิวฏจกฺขุ.}

\prose{กถํ ภควา มํสจกฺขุนาปิ วิวฏจกฺขุ, มํสจกฺขุมฺหิปิ ภควโต ปญฺจ วณฺณา สํวิชฺชนฺติ, นีโล จ วณฺโณ, ปีตโก จ วณฺโณ, โลหิตโก \csromanfolio{276}จ วณฺโณ, กโณฺห จ วณฺโณ, โอทาโต จ วณฺโณ.\csromanspacer{} อกฺขิโลมานิ จ ภควโต, ยตฺถ จ อกฺขิโลมานิ ปติฏฺฐิตานิ, ตํ นีลํ โหติ สุนีลํ ปาสาทิกํทสฺสเนยฺยํ อุมาปุปฺผสมานํ\footnote{อุมฺมาปุปฺผสมานํ (สี, ก), อุมฺมารปุปฺผสมานํ (สฺยา)}.\csromanspacer{} ตสฺส ปรโต ปีตกํ โหติ สุปีตกํ สุวณฺณวณฺณํ ปาสาทิกํ ทสฺสเนยฺยํ กณิการปุปฺผ- สมานํ.\csromanspacer{} อุภโต จ อกฺขิกูฏานิ ภควโต โลหิตกานิ โหนฺติ สุโลหิตกานิ ปาสาทิกานิ ทสฺสเนยฺยานิ อินฺทโคปกสมานานิ.\csromanspacer{} มชฺเฌ กณฺหํ โหติ สุกณฺหํ อลูขํ สินิทฺธํ ปาสาทิกํ ทสฺสเนยฺยํ อทฺทาริฏฺฐกสมานํ\footnote{อฬาริฏฺฐกสมานํ (สฺยา), ภทฺทาริฏฺฐกสมานํ (ฏฺฐ 331)}.\csromanspacer{} ตสฺส ปรโต โอทาตํ โหติ สุ-โอทาตํ เสตํ ปณฺฑรํ ปาสาทิกํ ทสฺสเนยฺยํ โอสธิตารกสมานํ.\csromanspacer{} เตน ภควา ปากติเกน มํสจกฺขุนา อตฺตภาวปริยาปนฺเนน ปุริมสุจริตกมฺมา- ภินิพฺพตฺเตน สมนฺตา โยชนํ ปสฺสติ ทิวา เจว รตฺติญฺจ.\csromanspacer{} ยทา หิปิ จตุรงฺคสมนฺนาคโต อนฺธกาโร โหติ.\csromanspacer{} สูริโย วา อตฺถงฺคโต โหติ, กาฬปกฺโข จ อุโปสโถ โหติ, ติพฺโพ จ วนสณฺโฑ โหติ, มหา จ กาฬเมโฆ\footnote{อกาลเมโฆ (สฺยา)} อพฺภุฏฺฐิโต โหติ.\csromanspacer{} เอวรูเปปิ จตุรงฺคสมนฺนาคเต อนฺธกาเร สมนฺตา โยชนํ ปสฺสติ.\csromanspacer{} นตฺถิ โส กุฏฺโฏ\footnote{กุฑฺโฑ (สี)} วา กวาฏํ วา ปากาโร วา ปพฺพโต วา คจฺโฉ วา ลตา วา อาวรณํ รูปานํ ทสฺสนาย.\csromanspacer{} เอกญฺเจ ติลผลํ นิมิตฺตํ กตฺวา ติลวาเห ปกฺขิเปยฺย.\csromanspacer{} ตญฺเญว ติลผลํ อุทฺธเรยฺย.\csromanspacer{} เอวํ ปริสุทฺธํ ภควโต ปากติกํ มํสจกฺขุ, เอวํ ภควา มํสจกฺขุนาปิ วิวฏจกฺขุ.}

\prose{กถํ ภควา ทิพฺเพน จกฺขุนาปิ วิวฏจกฺขุ, ภควา ทิพฺเพน จกฺขุนา วิสุทฺเธน อติกฺกนฺตมานุสเกน สตฺเต ปสฺสติ จวมาเน อุปปชฺชมาเน หีเน ปณีเต สุวณฺเณ ทุพฺพณฺเณ สุคเต ทุคฺคเต, ยถากมฺมูปเค สตฺเต ปชานาติ “อิเม วต โภนฺโต สตฺตา กายทุจฺจริเตน สมนฺนาคตา วจีทุจฺจริเตน สมนฺนาคตา มโนทุจฺจริเตน สมนฺนาคตา อริยานํ อุปวาทกา มิจฺฉทิฏฺฐิกา มิจฺฉาทิฏฺฐิกมฺมสมาทานา, เต กายสฺส เภทา ปรํ มรณา อปายํ ทุคฺคติํ วินิปาตํ นิรยํ อุปปนฺนา, อิเม วา ปน โภนฺโต สตฺตา กายสุจริเตน สมนฺนาคตา วจีสุจริเตน สมนฺนาคตา มโนสุจริเตน สมนฺนาคตา อริยานํ อนุปวาทกา สมฺมาทิฏฺฐิกา สมฺมาทิฏฺฐิกมฺมสมาทานา, เต กายสฺส เภทา ปรํ มรณา สุคติํ สคฺคํ โลกํ}

\vspace{\tipitakaparskip}
\csromancenter{ตุวฏกสุตฺตนิทฺเทส อุปปนฺนา”ติ อิติ ทิพฺเพน จกฺขุนา วิสุทฺเธน อติกฺกนฺตมานุสเกน สตฺเต ปสฺสติ จวมาเน อุปปชฺชมาเน หีเน ปณีเต สุวณฺเณ ทุพฺพณฺเณ สุคเต ทุคฺคเต, ยถากมฺมูปเค สตฺเต ปชานาติ.\csromanspacer{} อากงฺขมาโน จ ภควา เอกมฺปิ โลกธาตุํ ปสฺเสยฺย, เทฺวปิ โลกธาตุโย ปสฺเสยฺย, ติสฺโสปิ โลกธาตุโย ปสฺเสยฺย, จตสฺโสปิ โลกธาตุโย ปสฺเสยฺย, ปญฺจปิ โลกธาตุโย ปสฺเสยฺย, ทสปิ โลกธาตุโย ปสฺเสยฺย, วีสมฺปิ โลกธาตุโย ปสฺเสยฺย, ติํสมฺปิ โลกธาตุโย ปสฺเสยฺย, จตฺตาลีสมฺปิ โลกธาตุโย ปสฺเสยฺย, ปญฺญาสมฺปิ โลกธาตุโย ปสฺเสยฺย, สตมฺปิ โลกธาตุโย ปสฺเสยฺย, สหสฺสิมฺปิ จูฬนิกํ โลกธาตุํ ปสฺเสยฺย, ทฺวิสหสฺสิมฺปิ มชฺฌิมิกํ โลกธาตุํ ปสฺเสยฺย, ติสหสฺสิมฺปิ โลกธาตุํ ปสฺเสยฺย, มหาสหสฺสิมฺปิ\footnote{ติสหสฺสิํ มหาสหสฺสิมฺปิ (สี, ก)} โลกธาตุํ ปสฺเสยฺย.\csromanspacer{} ยาวตกํ วา ปน อากงฺเขยฺย, ตาวตกํ ปสฺเสยฺย.\csromanspacer{} เอวํ ปริสุทฺธํ ภควโต ทิพฺพจกฺขุ.\csromanspacer{} เอวํ ภควา ทิพฺเพน จกฺขุนาปิ วิวฏจกฺขุ.}
\prose{\csromanfolio{277}กถํ ภควา ปญฺญาจกฺขุนาปิ วิวฏจกฺขุ, ภควา มหาปญฺโญ ปุถุปญฺโญ หาสปญฺโญ ชวนปญฺโญ ติกฺขปญฺโญ นิพฺเพธิกปญฺโญ ปญฺญาปเภทกุสโล ปภินฺนญาโณ อธิคตปฏิสมฺภิโท จตุเวสารชฺชปฺปตฺโต ทสพลธารี ปุริสาสโภ ปุริสสีโห ปุริสนาโค ปุริสาชญฺโญ ปุริสโธรโยฺห อนนฺตญาโณ อนนฺตเตโช อนนฺตยโส อฑฺโฒ มหทฺธโน ธนวา, เนตา วิเนตา อนุเนตา, ปญฺญาเปตา อนิชฺฌาเปตา, เปกฺเขตา ปสาเทตา, โส หิ ภควา อนุปฺปนฺนสฺส มคฺคสฺส อุปฺปาเทตา, อสญฺชาตสฺส มคฺคสฺส สญฺชเนตา, อนกฺขาตสฺส มคฺคสฺส อกฺขาตา, มคฺคญฺญู มคฺควิทู มคฺคโกวิโท, มคฺคานุคา จ ปน เอตรหิ สาวกา วิหรนฺติ ปจฺฉา สมนฺนาคตา.}

\prose{โส หิ ภควา ชานํ ชานาติ, ปสฺสํ ปสฺสติ, จกฺขุภูโต ญาณภูโต ธมฺมภูโต พฺรหฺมภูโต วตฺตา ปวตฺตา อตฺถสฺส นินฺเนตา อมตสฺส ทาตา ธมฺมสฺสามี ตถาคโต, นตฺถิ ตสฺส ภควโต อญฺญาตํ อทิฏฺฐํ อวิทิตํ อสจฺฉิกตํ อผสฺสิตํ ปญฺญาย.\csromanspacer{} อตีตํ อนาคตํ ปจฺจุปฺปนฺนํ\footnote{อตีตานาคตปจฺจุปฺปนฺนํ (สฺยา) 138; ขุ 8. 175; ขุ 9. 375 ปิฏฺเฐสุปิ.} อุปาทาย สพฺเพ ธมฺมา สพฺพากาเรน พุทฺธสฺส ภควโต ญาณมุเข อาปาถํ อาคจฺฉนฺติ.\csromanspacer{} ยํ กิญฺจิ เนยฺยํ นาม อตฺถิ ชานิตพฺพํ\footnote{อตฺถิ ธมฺมํ ชานิตพฺพํ (สี, ก)} อตฺตตฺโถ วา ปรตฺโถ วา อุภยตฺโถ \csromanfolio{278}วา ทิฏฺฐธมฺมิโก วา อตฺโถ สมฺปรายิโก วา อตฺโถ อุตฺตาโน วา อตฺโถ คมฺภีโร วา อตฺโถ คูโฬฺห วา อตฺโถ ปฏิจฺฉนฺโน วา อตฺโถ เนยฺโย วา อตฺโถ นีโต วา อตฺโถ อนวชฺโช วา อตฺโถ นิกฺกิเลโส วา อตฺโถ โวทาโน วา อตฺโถ ปรมตฺโถ วา อตฺโถ, สพฺพํ ตํ\footnote{ปรมตฺโถ วา, สพฺพนฺตํ (สฺยา)} อนฺโตพุทฺธญาเณ ปริวตฺตติ.}

\prose{สพฺพํ กายกมฺมํ พุทฺธสฺส ภควโต ญาณานุปริวตฺติ, สพฺพํ วจีกมฺมํ.\csromanspacer{} สพฺพํ มโนกมฺมํ.\csromanspacer{} อตีเต พุทฺธสฺส ภควโต อปฺปฏิหตํ ญาณํ, อนาคเต.\csromanspacer{} ปจฺจุปฺปนฺเน อปฺปฏิหตํ ญาณํ.\csromanspacer{} ยาวตกํ เนยฺยํ, ตาวตกํ ญาณํ, ยาวตกํ ญาณํ, ตาวตกํ เนยฺยํ, เนยฺยปริยนฺติกํ ญาณํ, ญาณปริยนฺติกํ เนยฺยํ, เนยฺยํ อติกฺกมิตฺวา ญาณํ นปฺปวตฺตติ, ญาณํ อติกฺกมิตฺวา เนยฺยปโถ นตฺถิ, อญฺญมญฺญปริยนฺตฏฺฐายิโน เต ธมฺมา, ยถา ทฺวินฺนํ สมุคฺคปฏลานํ สมฺมาผุสิตานํ เหฏฺฐิมํ สมุคฺคปฏลํ อุปริมํ นาติวตฺตติ, อุปริมํ สมุคฺคปฏลํ เหฏฺฐิมํ นาติวตฺตติ, อญฺญมญฺญปริยนฺตฏฺฐายิโน.\csromanspacer{} เอวเมวํ พุทฺธสฺส ภควโต เนยฺยญฺจ ญาณญฺจ อญฺญมญฺญปริยนฺตฏฺฐยิโน.\csromanspacer{} ยาวตกํ เนยฺยํ, ตาวตกํ ญาณํ, ยาวตกํ ญาณํ, ตาวตกํ เนยฺยํ, เนยฺยปริยนฺติกํ ญาณํ, ญาณปริยนฺติกํ เนยฺยํ, เนยฺยํ อติกฺกมิตฺวา ญาณํ นปฺปวตฺตติ, ญาณํ อติกฺกมิตฺวา เนยฺยปโถ นตฺถิ, อญฺญมญฺญปริยนฺตฏฺฐายิโน เต ธมฺมา.}

\prose{สพฺพธมฺเมสุ พุทฺธสฺส ภควโต ญาณํ ปวตฺตติ, สพฺเพ ธมฺมา พุทฺธสฺส ภควโต อาวชฺชนปฏิพทฺธา อากงฺขปฏิพทฺธา มนสิการปฏิพทฺธา จิตฺตุปฺปาทปฏิพทฺธา.\csromanspacer{} สพฺพสตฺเตสุ พุทฺธสฺส ภควโต ญาณํ ปวตฺตติ, สพฺเพสํ สตฺตานํ ภควา อาสยํ ชานาติ, อนุสยํ ชานาติ, จริตํ ชานาติ, อธิมุตฺติํ ชานาติ, อปฺปรชกฺเข มหารชกฺเข ติกฺขินฺทฺริเย มุทินฺทฺริเย สฺวากาเร ทฺวากาเร สุวิญฺญาปเย ทุวิญฺญาปเย ภพฺพาภพฺเพ สตฺเต ปชานาติ, สเทวโก โลโก สมารโก สพฺรหฺมโก สสฺสมณพฺราหฺมณี ปชา สเทวมนุสฺสา อนฺโตพุทฺธญาเณ ปริวตฺตติ\footnote{ปริวตฺตนฺติ (ก) ปสูรสุตฺตนิทฺเทเส อวสาเนปิ.}.}

\prose{ยถา เย เกจิ มจฺฉกจฺฉปา อนฺตมโส ติมิ ติมิงฺคลํ อุปาทาย อนฺโตมหาสมุทฺเท ปริวตฺตนฺติ, เอวเมวํ สเทวโก โลโก สมารโก สพฺรหฺมโก สสฺสมณพฺราหฺมณี ปชา สเทวมนุสฺสา อนฺโตพุทฺธญาเณ ปริวตฺตติ2.\csromanspacer{}\csromanspacer{}ยถา เย เกจิ ปกฺขี อนฺตมโส ครุฬํ เวนเตยฺยํ อุปาทาย อากาสสฺส ปเทเส ปริวตฺตนฺติ, เอวเมวํ เยปิ}

\vspace{\tipitakaparskip}
\csromancenter{ตุวฏกสุตฺตนิทฺเทส เต สาริปุตฺตสมา ปญฺญาย, เตปิ พุทฺธญาณสฺส ปเทเส ปริวตฺตนฺติ.\csromanspacer{} พุทฺธญาณํ เทวมนุสฺสานํ ปญฺญํ ผริตฺวา อภิภวิตฺวา ติฏฺฐติ.\csromanspacer{} เยปิ เต ขตฺติยปณฺฑิตา พฺราหฺมณปณฺฑิตา คหปติปณฺฑิตา สมณปณฺฑิตา นิปุณา กตปรปฺปวาทา วาลเวธิรูปา โวภินฺทนฺตา มญฺเญ จรนฺติ ปญฺญาคเตน ทิฏฺฐิคตานิ.\csromanspacer{} เต ปญฺหํ อภิสงฺขริตฺวา อภิสงฺขริตฺวา ตถาคตํ อุปสงฺกมิตฺวา ปุจฺฉนฺติ คูฬฺหานิ จ ปฏิจฺฉนฺนานิ จ.\csromanspacer{} กถิตา วิสชฺชิตา จ เต ปญฺหา ภควตา โหนฺติ นิทฺทิฏฺฐการณา.\csromanspacer{} อุปกฺขิตฺตกา จ เต ภควโต สมฺปชฺชนฺติ.\csromanspacer{} อถ โข ภควา ตตฺถ อติโรจติ ยทิทํ ปญฺญายาติ, เอวํ ภควา ปญฺญาจกฺขุนาปิ วิวฏจกฺขุ.}
\prose{\csromanfolio{279}\csromansymbolfootnote{*}{ขุ 8.177; ขุ 9. 116 ปิฏฺเฐสุปิ.}กถํ ภควา พุทฺธจกฺขุนาปิ วิวฏจกฺขุ, ภควา พุทฺธจกฺขุนา โลกํ โวโลเกนฺโต อทฺทส สตฺเต อปฺปรชกฺเข มหารชกฺเข ติกฺขินฺทฺริเย มุทินฺทฺริเย สฺวากาเร ทฺวากาเร สุวิญฺญาปเย ทุวิญฺญาปเย, อปฺเปกจฺเจ ปรโลกวชฺชภยทสฺสาวิโน วิหรนฺเต, อปฺเปกจฺเจ นปรโลกวชฺชภยทสฺสาวิโน วิหรนฺเต, เสยฺยถาปิ นาม อุปฺปลินิยํ วา ปทุมินิยํ วา ปุณฺฑรีกินิยํ วา อปฺเปกจฺจานิ อุปฺปลานิ วา ปทุมานิ วา ปุณฺฑรีกานิ วา อุทเก ชาตานิ, อุทเก สํวฑฺฒานิ, อุทกา อนุคฺคตานิ, อนฺโตนิมุคฺคโปสีนิ, อปฺเปกจฺจานิ อุปฺปลานิ วา ปทุมานิ วา ปุณฺฑรีกานิ วา อุทเก ชาตานิ, อุทเก สํวฑฺฒานิ, สโมทกํ ฐิตานิ, อเปกจฺจานิ อุปฺปลานิ วา ปทุมานิ วา ปุณฺฑรีกานิ วา อุทเก ชาตานิ, อุทเก สํวฑฺฒานิ, อุทกา อจฺจุคฺคมฺม ติฏฺฐนฺติ อนุปลิตฺตานิ อุทเกน.\csromanspacer{} เอวเมวํ ภควา พุทฺธจกฺขุนา โลกํ โวโลเกนฺโต อทฺทส สตฺเต อปฺปรชกฺเข มหารชกฺเข ติกฺขินฺทฺริเย มุทินฺทฺริเย สฺวากาเร ทฺวากาเร สุวิญฺญาปเย ทุวิญฺญาปเย, อปฺเปกจฺเจ ปรโลกวชฺชภยทสฺสาวิโน วิหรนฺเต, อปฺเปกจฺเจ นปรโลกวชฺชภยทสฺสาวิโน วิหรนฺเต.\csromanspacer{} ชานาติ ภควา “อยํ ปุคฺคโล ราคจริโต, อยํ โทสจริโต, อยํ โมหจริโต, อยํ วิตกฺกจริโต, อยํ สทฺธาจริโต, อยํ ญาณจริโต”ติ.\csromanspacer{} ราคจริตสฺส ภควา ปุคฺคลสฺส อสุภกถํ กเถติ.\csromanspacer{} โทสจริตสฺส ภควา ปุคฺคลสฺส เมตฺตาภาวนํ อาจิกฺขติ.\csromanspacer{} โมหจริตสฺส ภควา ปุคฺคลสฺส อุทฺเทเส ปริปุจฺฉาย กาเลน ธมฺมสฺสวเน กาเลน ธมฺมสากจฺฉาย ครุสํวาเส นิเวเสติ.\csromanspacer{} วิตกฺกจริตสฺส ภควา ปุคฺคลสฺส อานาปานสฺสติํ อาจิกฺขติ.\csromanspacer{} สทฺธาจริตสฺส ภควา ปุคฺคลสฺส ปสาทนียํ นิมิตฺตํ อาจิกฺขติ \csromanfolio{280}พุทฺธสุโพธิํ\footnote{พุทฺธสุพุทฺธตํ (ก)} ธมฺมสุธมฺมตํ สํฆสุปฺปฏิปตฺติํ สีลานิ จ อตฺตโน.\csromanspacer{} ญาณจริตสฺส ภควา ปุคฺคลสฺส อาจิกฺขติ วิปสฺสนานิมิตฺตํ อนิจฺจาการํ ทุกฺขาการํ อนตฺตาการํ.}

\csromangathameasure{*“เสเล ยถา ปพฺพตมุทฺธนิฏฺฐิโต, \\ ยถาปิ ปสฺเส ชนตํ สมนฺตโต. \\ ตถูปมํ ธมฺมมยํ สุเมธ, \\ ปาสาทมารุยฺห \textbf{สมนฺตจกฺขุ}. \\ โสกาวติณฺณํ ชนตมเปตโสโก, \\ อเวกฺขสฺสุ ชาติชราภิภูตนฺ”ติ.}
\csromangathagroup{\csromangathastanza{\csromansymbolfootnote{*}{ที 2. 33; ม 1. 225; ม 2. 292; สํ 1. 139 ปิฏฺเฐสุปิ.}“เสเล ยถา ปพฺพตมุทฺธนิฏฺฐิโต, \\ ยถาปิ ปสฺเส ชนตํ สมนฺตโต. \\ ตถูปมํ ธมฺมมยํ สุเมธ, \\ ปาสาทมารุยฺห \textbf{สมนฺตจกฺขุ}. \\ โสกาวติณฺณํ ชนตมเปตโสโก, \\ อเวกฺขสฺสุ ชาติชราภิภูตนฺ”ติ.}}

\prose{เอวํ ภควา พุทฺธจกฺขุนาปิ วิวฏจกฺขุ.}

\prose{กถํ ภควา \textbf{สมนฺตจกฺขุ}นาปิ วิวฏจกฺขุ, \textbf{สมนฺตจกฺขุ} วุจฺจติ สพฺพญฺญุตญาณํ.\csromanspacer{} ภควา สพฺพญฺญุตญาเณน อุเปโต สมุเปโต อุปคโต สมุปคโต อุปปนฺโน สมุปปนฺโน สมนฺนาคโต\csromandash{}}

\csromangathameasure{+“น ตสฺส อทิฏฺฐมิธตฺถิ กิญฺจิ, \\ อโถ อวิญฺญาตมชานิตพฺพํ. \\ สพฺพํ อภิญฺญาสิ ยทตฺถิ เนยฺยํ, \\ ตถาคโต เตน สมนฺตจกฺขู”ติ.}
\csromangathagroup{\csromangathastanza{\csromansymbolfootnote{+}{ขุ 8. 93, 178; ขุ 9. 127, 227 ปิฏฺฐาทีสุ.}“น ตสฺส อทิฏฺฐมิธตฺถิ กิญฺจิ, \\ อโถ อวิญฺญาตมชานิตพฺพํ. \\ สพฺพํ อภิญฺญาสิ ยทตฺถิ เนยฺยํ, \\ ตถาคโต เตน สมนฺตจกฺขู”ติ.}}

\prose{เอวํ ภควา \textbf{สมนฺตจกฺขุ}นาปิ วิวฏจกฺขูติ อกิตฺตยิ วิวฏจกฺขุ.}

\prose{\textbf{สกฺขิธมฺมํ ปริสฺสยวินยนฺ}ติ \textbf{สกฺขิธมฺมนฺ}ติ น อิติหิติหํ, น อิติกิราย, น ปรมฺปราย, น ปิฏกสมฺปทาย, น ตกฺกเหตุ, น นยเหตุ, น อาการปริวิตกฺเกน, น ทิฏฺฐินิชฺฌานกฺขนฺติยา สามํ สยมภิญฺญาตํ อตฺตปจฺจกฺขํ ธมฺมนฺติ สกฺขิธมฺมํ.\csromanspacer{}\csromanspacer{}\textbf{ปริสฺสยวินยนฺ}ติ \textbf{ปริสฺสยา}ติ เทฺว \textbf{ปริสฺสยา} ปากฏ\textbf{ปริสฺสยา} จ ปฏิจฺฉนฺน\textbf{ปริสฺสยา} จ.\csromanspacer{} กตเม ปากฏ\textbf{ปริสฺสยา}, สีหา พฺยคฺฆา ทีปี อจฺฉา ตรจฺฉา โกกา มหิํสา\footnote{โคมหิสา (สฺยา) 10 ปิฏฺเฐปิ.} หตฺถี อหี วิจฺฉิกา สตปที โจรา วา อสฺสุ มานวา วา กตกมฺมา วา อกตกมฺมา วา จกฺขุโรโค โสตโรโค ฆานโรโค ชิวฺหาโรโค กายโรโค สีสโรโค กณฺณโรโค มุขโรโค ทนฺตโรโค กาโส สาโส ปินาโส ฑาโห\footnote{qอโห (สี, สฺยา) 10 ปิฏฺเฐปิ.} ชโร กุจฺฉิโรโค มุจฺฉา ปกฺขนฺทิกา สูลา วิสูจิกา กุฏฺฐํ คณฺโฑ กิลาโส โสโส}

\vspace{\tipitakaparskip}
\csromancenter{ตุวฏกสุตฺตนิทฺเทส อปมาโร ททฺทุ กณฺฑุ กจฺฉุ รขสา วิตจฺฉิกา โลหิตปิตฺตํ\footnote{โลหิตํ ปิตฺตํ (สฺยา, ก)} มธุเมโห อํสา ปิฬกา ภคนฺทลา ปิตฺตสมุฏฺฐานา อาพาธา เสมฺหสมุฏฺฐานา อาพาธา วาตสมุฏฺฐานา อาพาธา สนฺนิปาติกา อาพาธา อุตุปริณามชา อาพาธา วิสมปริหารชา อาพาธา โอปกฺกมิกา อาพาธา กมฺมวิปากชา อาพาธา สีตํ อุณฺหํ ชิฆจฺฉา ปิปาสา อุจฺจาโร ปสฺสาโว ฑํสมกสวาตาตปสรีสปสมฺผสฺสา อิติ วา, อิเม วุจฺจนฺติ ปากฏปริสฺสยา.}
\prose{\csromanfolio{281}กตเม ปฏิจฺฉนฺนปริสฺสยา, กายทุจฺจริตํ วจีทุจฺจริตํ มโนทิจฺจริตํ กามจฺฉนฺทนีวรณํ พฺยาปาทนีวรณํ ถินมิทฺธนีวรณํ อุทฺธจฺจกุกฺกุจฺจนีวรณํ วิจิกิจฺฉานีวรณํ, ราโค โทโส โมโห โกโธ อุปนาโห มกฺโข ปฬาโส อิสฺสา มจฺฉริยํ มายา สาเฐยฺยํ ถมฺโภ สารมฺโภ มาโน อติมาโน มโท ปมาโท สพฺเพ กิเลสา สพฺเพ ทุจฺจริตา สพฺเพ ทรถา สพฺเพ ปริฬาหา สพฺเพ สนฺตาปา สพฺพากุสลาภิสงฺขารา, อิเม วุจฺจนฺติ ปฏิจฺฉนฺนปริสฺสยา.}

\prose{\textbf{ปริสฺสยา}ติ เกนฏฺเฐน ปริสฺสยา, ปริสหนฺตีติ ปริสฺสยา.\csromanspacer{} ปริหานาย สํวตฺตนฺตีติ ปริสฺสยา.\csromanspacer{} ตตฺราสยาติ ปริสฺสยา.\csromanspacer{} กถํ ปริสหนฺตีติ ปริสฺสยา, เต ปริสฺสยา ตํ ปุคฺคลํ สหนฺติ ปริสหนฺติ อภิภวนฺติ อชฺโฌตฺถรนฺติ ปริยาทิยนฺติ มทฺทนฺติ, เอวํ ปริสหนฺตีติ ปริสฺสยา.}

\prose{กถํ ปริหานาย สํวตฺตนฺตีติ ปริสฺสยา, เต ปริสฺสยา กุสลานํ ธมฺมานํ อนฺตรายาย ปริหานาย สํวตฺตนฺติ.\csromanspacer{} กตเมสํ กุสลานํ ธมฺมานํ สมฺมาปฏิปทาย อนุโลมปฏิปทาย อปจฺจนีกปฏิปทาย อวิรุทฺธปฏิปทาย อนฺวตฺถปฏิปทาย ธมฺมานุธมฺมปฏิปทาย สีเลสุ ปริปูรการิตาย อินฺทฺริเยสุ คุตฺตทฺวารตาย โภชเน มตฺตญฺญุตาย ชาคริยานุโยคสฺส สติสมฺปชญฺญสฺส จตุนฺนํ สติปฏฺฐานานํ ภาวนานุโยคสฺส จตุนฺนํ สมฺมปฺปธานานํ จตุนฺนํ อิทฺธิปาทานํ ปญฺจนฺนํ อินฺทฺริยานํ ปญฺจนฺนํ พลานํ สตฺตนฺนํ โพชฺฌงฺคานํ อริยสฺส อฏฺฐงฺคิกสฺส มคฺคสฺส ภาวนานุโยคสฺส อิเมสํ กุสลานํ ธมฺมานํ อนฺตรายาย ปริหานาย สํวตฺตนฺติ, เอวมฺปิ ปริหานาย สํวตฺตนฺตีติ ปริสฺสยา.}

\prose{\csromanfolio{282}กถํ ตตฺราสยาติ ปริสฺสยา, ตตฺเถเต ปาปกา อกุสลา ธมฺมา อุปฺปชฺชนฺติ อตฺตภาวสนฺนิสฺสยา, ยถา พิเล พิลาสยา ปาณา สยนฺติ, ทเก ทกาสยา ปาณา สยนฺติ, วเน วนาสยา ปาณา สยนฺติ, รุกฺเข รุกฺขาสยา ปาณา สยนฺติ, เอวเมวํ ตตฺเถเต ปาปกา อกุสลา ธมฺมา อุปฺปชฺชนฺติ อตฺตภาวสนฺนิสฺสยาติ, เอวมฺปิ ตตฺราสยาติ ปริสฺสยา.}

\prose{วุตฺตํ เหตํ ภควตา\csromandash{} * สานฺเตวาสิโก ภิกฺขเว ภิกฺขุ สาจริยโก ทุกฺขํ น ผาสุ วิหรติ.\csromanspacer{} กถญฺจ ภิกฺขเว ภิกฺขุ สานฺเตวาสิโก สาจริยโก ทุกฺขํ น ผาสุ วิหรติ, อิธ ภิกฺขเว ภิกฺขุโน จกฺขุนา รูปํ ทิสฺวา อุปฺปชฺชนฺติ ปาปกา อกุสลา ธมฺมา สรสงฺกปฺปา สญฺโญชนิยา, ตฺยสฺส อนฺโต วสนฺติ อนฺวาสฺสวนฺติ ปาปกา อกุสลา ธมฺมาติ ตสฺมา “สานฺเตวาสิโก”ติ วุจฺจติ, เต นํ สมุทาจรนฺติ.\csromanspacer{} สมุทาจรนฺติ\footnote{สมุทาจาเรนฺติ (สี)} นํ ปาปกา อกุสลา ธมฺมาติ ตสฺมา “สาจริยโก”ติ วุจฺจติ.}

\prose{ปุน จปรํ ภิกฺขเว ภิกฺขุโน โสเตน สทฺทํ สุตฺวา ฯเปฯ ฆาเนน คนฺธํ ฆายิตฺวา.\csromanspacer{} ชิวฺหาย รสํ สายิตฺวา.\csromanspacer{} กาเยน โผฏฺฐพฺพํ ผุสิตฺวา.\csromanspacer{} มนสา ธมฺมํ วิญฺญาย อุปฺปชฺชนฺติ ปาปกา อกุสลา ธมฺมา สรสงฺกปฺปา สญฺโญชนิยา, ตฺยสฺส อนฺโต วสนฺติ อนฺวาสฺสวนฺติ ปาปกา อกุสลา ธมฺมาติ ตสฺมา “สานฺเตวาสิโก”ติ วุจฺจติ, เต นํ สมุทาจรนฺติ.\csromanspacer{} สมุทาจรนฺติ นํ ปาปกา อกุสลา ธมฺมาติ ตสฺมา “สาจริยโก”ติ วุจฺจติ.\csromanspacer{} เอวํ โข ภิกฺขเว ภิกฺขุ สานฺเตวาสิโก สาจริยโก ทุกฺขํ น ผาสุ วิหรตีติ.\csromanspacer{} เอวมฺปิ ตตฺราสยาติ ปริสฺสยา.}

\prose{วุตฺตํ เหตํ ภควตา\csromandash{}ตโยเม ภิกฺขเว อนฺตรามลา อนฺตรา-อมิตฺตา อนฺตราสปตฺตา อนฺตราวธกา อนฺตราปจฺจตฺถิกา.\csromanspacer{} กตเม ตโย, โลโภ ภิกฺขเว อนฺตรามโล\footnote{อนฺตรามลํ (สี, ก); อํ 2. 471; ขุ 1. 252; ขุ 8. 251 ปิฏฺเฐสุปิ.} อนฺตรา-อมิตฺโต อนฺตราสปตฺโต อนฺตราวธโก อนฺตราปจฺจตฺถิโก, โทโส ภิกฺขเว ฯเปฯ โมโห ภิกฺขเว อนฺตรามโล อนฺตรา-อมิตฺโต อนฺตราสปตฺโต อนฺตราวธโก อนฺตราปจฺจตฺถิโก.\csromanspacer{} อิเม โข ภิกฺขเว ตโย อนฺตรามลา อนฺตรา-อมิตฺตา อนฺตราสปตฺตา อนฺตราวธกา อนฺตราปจฺจตฺถิกาติ.}

\csromanheader{2}{14. ตุวฏกสุตฺตนิทฺเทส}
\csromangathameasure{*“อนตฺถชนโน โลโภ, โลโภ จิตฺตปฺปโกปโน. \\ ภยมนฺตรโต ชาตํ, ตํ ชโน นาวพุชฺฌติ. \\ ลุทฺโธ อตฺถํ น ชานาติ, ลุทฺโธ ธมฺมํ น ปสฺสติ. \\ อนฺธนฺตมํ ตทา โหติ, ยํ โลโภ สหเต นรํ. \\ อนตฺถชนโน โทโส, โทโส จิตฺตปฺปโกปโน. \\ ภยมนฺตรโต ชาตํ, ตํ ชโน นาวพุชฺฌติ. \\ กุทฺโธ อตฺถํ น ชานาติ, กุทฺโธ ธมฺมํ น ปสฺสติ. \\ อนฺธนฺตมํ ตทา โหติ, ยํ โทโส สหเต นรํ. \\ อนตฺถชนโน โมโห, โมโห จิตฺตปฺปโกปโน. \\ ภยมนฺตรโต ชาตํ, ตํ ชโน นาวพุชฺฌติ. \\ มูโฬฺห อตฺถํ น ชานาติ, มูโฬฺห ธมฺมํ น ปสฺสติ. \\ อนฺธนฺตมํ ตทา โหติ, ยํ โมโห สหเต นรนฺ”ติ.}
\csromangathagroup{\csromangathastanza{\csromanfolio{283}\csromansymbolfootnote{*}{สํ 1. 70 ปิฏฺเฐปิ.}“อนตฺถชนโน โลโภ, โลโภ จิตฺตปฺปโกปโน. \\ ภยมนฺตรโต ชาตํ, ตํ ชโน นาวพุชฺฌติ.}\csromangathastanza{ลุทฺโธ อตฺถํ น ชานาติ, ลุทฺโธ ธมฺมํ น ปสฺสติ. \\ อนฺธนฺตมํ\footnote{อนฺธตมํ (สฺยา); อํ 2. 471; ขุ 1. 252; ขุ 8. 251 ปิฏฺเฐสุปิ.} ตทา โหติ, ยํ โลโภ สหเต นรํ.}\csromangathastanza{อนตฺถชนโน โทโส, โทโส จิตฺตปฺปโกปโน. \\ ภยมนฺตรโต ชาตํ, ตํ ชโน นาวพุชฺฌติ.}\csromangathastanza{กุทฺโธ อตฺถํ น ชานาติ, กุทฺโธ ธมฺมํ น ปสฺสติ. \\ อนฺธนฺตมํ ตทา โหติ, ยํ โทโส สหเต นรํ.}\csromangathastanza{อนตฺถชนโน โมโห, โมโห จิตฺตปฺปโกปโน. \\ ภยมนฺตรโต ชาตํ, ตํ ชโน นาวพุชฺฌติ.}\csromangathastanza{มูโฬฺห อตฺถํ น ชานาติ, มูโฬฺห ธมฺมํ น ปสฺสติ. \\ อนฺธนฺตมํ ตทา โหติ, ยํ โมโห สหเต นรนฺ”ติ.}}

\prose{เอวมฺปิ ตตฺราสยาติ ปริสฺสยา.}

\prose{วุตฺตํ เหตํ ภควตา\csromandash{} * ตโย โข มหาราช ปุริสสฺส ธมฺมา อชฺฌตฺตํ อุปฺปชฺชมานา อุปฺปชฺชนฺติ อหิตาย ทุกฺขาย อผาสุวิหาราย.\csromanspacer{} กตเม ตโย, โลโภ โข มหาราช ปุริสสฺส ธมฺโม อชฺฌตฺตํ อุปฺปชฺชมาโน อุปฺปชฺชติ อหิตาย ทุกฺขาย อผาสุวิหาราย, โทโส โข มหาราช ฯเปฯ โมโห โข มหาราช ปุริสสฺส ธมฺโม อชฺฌตฺตํ อุปฺปชฺชมาโน อุปฺปชฺชติ อหิตาย ทุกฺขาย อผาสุวิหาราย.\csromanspacer{} อิเม โข มหาราช ตโย ปุริสสฺส ธมฺมา อชฺฌตฺตํ อุปฺปชฺชมานา อุปฺปชฺชนฺติ อหิตาย ทุกฺขาย อผาสุวิหาราย.}

\csromangathameasure{+“โลโภ โทโส จ โมโห จ, ปุริสํ ปาปเจตสํ. \\ หิํสนฺติ อตฺตสมฺภูตา, ตจสารํว สมฺผลนฺ”ติ.}
\csromangathagroup{\csromangathastanza{\csromansymbolfootnote{+}{ขุ 1. 226; ขุ 8. 252 ปิฏฺเฐสุปิ.}“โลโภ โทโส จ โมโห จ, ปุริสํ ปาปเจตสํ. \\ หิํสนฺติ อตฺตสมฺภูตา, ตจสารํว สมฺผลนฺ”ติ.}}

\prose{เอวมฺปิ ตตฺราสยาติ ปริสฺสยา.}

\csromanheader{2}{วุตฺตํ เหตํ ภควตา\csromandash{}}
\csromangathameasure{*“ราโค จ โทโส จ อิโต นิทานา,}
\csromangathagroup{\csromangathastanza{\csromanfolio{284}\csromansymbolfootnote{*}{สํ 1. 209; ขุ 1. 320; ขุ 8. 252; ขุ 10. 126 ปิฏฺเฐสุปิ.}“ราโค จ โทโส จ อิโต นิทานา,}}

\prose{อรตี รตี โลมหํโส อิโตชา.}

\prose{อิโต สมุฏฺฐาย มโนวิตกฺกา,}

\vspace{\tipitakaparskip}
\csromancenter{กุมารกา ธงฺกมิโวสฺสชนฺตี”ติ.}
\prose{เอวมฺปิ ตตฺราสยาติ ปริสฺสยา.}

\prose{\textbf{ปริสฺสยวินยนฺ}ติ ปริสฺสยวินยํ ปริสฺสยปฺปหานํ ปริสฺสยวูปสมํ ปริสฺสยปฏินิสฺสคฺคํ ปริสฺสยปฏิปสฺสทฺธิํ อมตํ นิพฺพานนฺติ สกฺขิธมฺมํ ปริสฺสยวินยํ.}

\prose{\textbf{ปฏิปทํ วเทหิ ภทฺทนฺเต}ติ ปฏิปทํ วเทหิ สมฺมาปฏิปทํ อนุโลมปฏิปทํ อปจฺจนีกปฏิปทํ อวิรุทฺธปฏิปทํ อนฺวตฺถปฏิปทํ ธมฺมานุธมฺมปฏิปทํ สีเลสุ ปริปูรการิตํ อินฺทฺริเยสุ คุตฺตทฺวารตํ โภชเน มตฺตญฺญุตํ ชาคริยานุโยคํ สติสมฺปชญฺญํ จตฺตาโร สติปฏฺฐาเน จตฺตาโร สมฺมปฺปธาเน จตฺตาโร อิทฺธิปาเท ปญฺจินฺทฺริยานิ ปญฺจ พลานิ สตฺต โพชฺฌงฺเค อริยํ อฏฺฐงฺคิกํ มคฺคํ นิพฺพานญฺจ นิพฺพานคามินิญฺจ ปฏิปทํ วเทหิ อาจิกฺข เทเสหิ ปญฺญเปหิ ปฏฺฐเปหิ วิวร วิภช อุตฺตานีกโรหิ ปกาเสหีติ ปฏิปทํ วเทหิ.\csromanspacer{}\csromanspacer{}\textbf{ภทฺทนฺเต}ติ โส นิมฺมิโต พุทฺธํ ภควนฺตํ อาลปติ.\csromanspacer{} อถ วา ยํ ตฺวํ ธมฺมํ อาจิกฺขสิ เทเสสิ ปญฺญเปสิ ปฏฺฐเปสิ วิวริ วิภชิ อุตฺตานี-อกาสิํ ปกาเสสิ, สพฺพํ ตํ สุนฺทรํ ภทฺทกํ กลฺยาณํ อนวชฺชํ เสวิตพฺพนฺติ ปฏิปทํ วเทหิ ภทฺทนฺเต.}

\prose{\textbf{ปาติโมกฺขํ อถ ปาปิ สมาธิ}นฺติ \textbf{ปาติโมกฺขนฺ}ติ สีลํ ปติฏฺฐา อาทิ จรณํ สํยโม สํวโร มุขํ ปมุขํ กุสลานํ ธมฺมานํ สมาปตฺติยา.\csromanspacer{} \textbf{อถ วาปิ สมาธินฺ}ติ ยา จิตฺตสฺส ฐิติ สณฺฐิติ อวฏฺฐิติ อวิสาหาโร อวิกฺเขโป อวิสาหตมานสตา สมโถ สมาธินฺทฺริยํ สมาธิพลํ สมฺมาสมาธีติ ปาติโมกฺขํ อถ วาปิ สมาธิํ.\csromanspacer{} เตนาห โส นิมฺมิโต\csromandash{}}

\prose{อกิตฺตยี วิวฏจกฺขุ,}

\vspace{\tipitakaparskip}
\csromancenter{สกฺขิธมฺมํ ปริสฺสยวินยํ.}
\csromancenter{\textbf{ปฏิปทํ วเทหิ ภทฺทนฺเต},}
\csromancenter{ปาติโมกฺขํ \textbf{อถ วาปิ สมาธินฺ}ติ.}
\csromanheader{2}{14. ตุวฏกสุตฺตนิทฺเทส}
\proseitem{157}{\csromanfolio{285}\csromansymbolfootnote{*}{ขุ 1. 422 ปิฏฺเฐปิ.}จกฺขูหิ เนว โลล’สฺส,}

\prose{\textbf{คามกถาย อาวรเย โสตํ.}}

\csromangathameasure{\textbf{รเส จ นานุคิชฺเฌยฺย,} \\ \textbf{น จ มมาเยถ กิญฺจิ โลกสฺมิํ.}}
\csromangathagroup{\csromangathastanza{\textbf{รเส จ นานุคิชฺเฌยฺย,} \\ \textbf{น จ มมาเยถ กิญฺจิ โลกสฺมิํ.}}}

\prose{\textbf{จกฺขูหิ เนว โลล’สฺสา}ติ กถํ จกฺขุโลโลติ, อิเธกจฺโจ จกฺขุโลลิเยน สมนฺนาคโต โหติ “อทิฏฺฐํ ทกฺขิตพฺพํ, ทิฏฺฐํ สมติกฺกมิตพฺพนฺ”ติ อาราเมน อารามํ, อุยฺยาเนน อุยฺยานํ, คาเมน คามํ, นิคเมน นิคมํ, นคเรน นครํ, รฏฺเฐน รฏฺฐํ, ชนปเทน ชนปทํ ทีฆจาริกํ อนวฏฺฐิตจาริกํ\footnote{อนวตฺถิตจาริกํ (สี, สฺยา) อุปริ 400 ปิฏฺเฐปิ.} อนุยุตฺโต จ โหติ รูปสฺส ทสฺสนาย, เอวมฺปิ จกฺขุโลโล โหติ.}

\prose{อถ วา ภิกฺขุ อนฺตรฆรํ ปวิฏฺโฐ วีถิํ ปฏิปนฺโน อสํวุโต คจฺฉติ หตฺถิํ โอโลเกนฺโต, อสฺสํ โอโลเกนฺโต, รถํ โอโลเกนฺโต, ปตฺติํ โอโลเกนฺโต, อิตฺถิโย โอโลเกนฺโต, ปุริเส โอโลเกนฺโต, กุมารเก โอโลเกนฺโต, กุมาริกาโย โอโลเกนฺโต, อนฺตราปณํ โอโลเกนฺโต, ฆรมุขานิ โอโลเกนฺโต, อุทฺธํ โอโลเกนฺโต, อโธ โอโลเกนฺโต, ทิสาวิทิสํ วิเปกฺขมาโน คจฺฉติ, เอวมฺปิ จกฺขุโลโล โหติ.}

\prose{อถ วา ภิกฺขุ จกฺขุนา รูปํ ทิสฺวา นิมิตฺตคฺคาหี โหติ อนุพฺยญฺชนคฺคาหี, ยตฺวาธิกรณเมนํ จกฺขุนฺทฺริยํ อสํวุตํ วิหรนฺตํ อภิชฺฌาโทมนสฺสา ปาปกา อกุสลา ธมฺมา อนฺวาสฺสเวยฺยุํ, ตสฺส สํวราย น ปฏิปชฺชติ, น รกฺขติ จกฺขุนฺทฺริยํ, จกฺขุนฺทฺริเย น สํวรํ อาปชฺชติ.\csromanspacer{} เอวมฺปิ จกฺขุโลโล โหติ.}

\prose{ยถา วา ปเนเก โภนฺโต สมณพฺราหฺมณา สทฺธาเทยฺยานิ โภชนานิ ภุญฺชิตฺวา เต เอวรูปํ วิสูกทสฺสนํ อนุยุตฺตา วิหรนฺติ.\csromanspacer{} เสยฺยถิทํ, นจฺจํ คีตํ วาทิตํ เปกฺขํ อกฺขานํ ปาณิสฺสรํ เวตาฬํ กุมฺภถูณํ\footnote{กุมฺภถูนํ (สี, สฺยา, ก); ที 1. 6; ขุ 8. 285 ปิฏฺเฐสุปิ.} โสภนกํ\footnote{โสภน ครกํ (สี, สฺยา)} จณฺฑาลํ วํสํ โธวนํ หตฺถิยุทฺธํ อสฺสยุทฺธํ มหิํสยุทฺธํ\footnote{มหิสยุทฺธํ (สี, สฺยา)} อุสภยุทฺธํ อชยุทฺธํ เมณฺฑยุทฺธํ กุกฺกุฏยุทฺธํ วฏฺฏกยุทฺธํ ทณฺฑยุทฺธํ มุฏฺฐิยุทฺธํ นิพฺพุทฺธํ อุยฺโยธิกํ พลคฺคํ เสนาพฺยูหํ อนีกทสฺสนํ อิติ วา, เอวมฺปิ จกฺขุโลโล โหติ.}

\prose{\csromanfolio{286}กถํ น จกฺขุโลโล โหติ, อิธ ภิกฺขุ อนฺตรฆรํ ปวิฏฺโฐ วีถิํ ปฏิปนฺโน สํวุโต คจฺฉติ น หตฺถิํ โอโลเกนฺโต, น อสฺสํ โอโลเกนฺโต, น รถํ โอโลเกนฺโต, น ปตฺติํ โอโลเกนฺโต, น อิตฺถิโย โอโลเกนฺโต, น ปุริเส โอโลเกนฺโต, น กุมารเก โอโลเกนฺโต, น กุมาริกาโย โอโลเกนฺโต, น อนฺตราปณํ โอโลเกนฺโต, น ฆรมุขานิ โอโลเกนฺโต, น อุทฺธํ โอโลเกนฺโต, น อโธ โอโลเกนฺโต, น ทิสาวิทิสาวิเปกฺขมาโน คจฺฉติ, เอวมฺปิ น จกฺขุโลโล โหติ.}

\prose{อถ วา ภิกฺขุ จกฺขุนา รูปํ ทิสฺวาน นิมิตฺตคฺคาหี โหติ นานุพฺยญฺชนคฺคาหี.\csromanspacer{} ยตฺวาธิกรณเมนํ จกฺขุนฺทฺริยํ อสํวุตํ วิหรนฺตํ อภิชฺฌาโทมนสฺสา ปาปกา อกุสลา ธมฺมา อนฺวาสฺสเวยฺยุํ, ตสฺส สํวราย ปฏิปชฺชติ, รกฺขติ จกฺขุนฺทฺริยํ, จกฺขุนฺทฺริเย สํวรํ อาปชฺชติ, เอวมฺปิ น จกฺขุโลโล โหติ.}

\prose{ยถา วา ปเนเก โภนฺโต สมณพฺราหฺมณา สทฺธาเทยฺยานิ โภชนานิ ภุญฺชิตฺวา เต เอวรูปํ วิสูกทสฺสนํ อนนุยุตฺตา วิหรนฺติ.\csromanspacer{} เสยฺยถิทํ, นจฺจํ คีตํ วาทิตํ เปกฺขํ อกฺขานํ ฯเปฯ อนีกทสฺสนํ อิติ วา, เอวรูปา วิสูกทสฺสนา ปฏิวิรโต โหติ, เอวมฺปิ น จกฺขุโลโล โหติ.}

\prose{\textbf{จกฺขูหิ เนว โลล’สฺสา}ติ จกฺขุโลลิยํ ปชเหยฺย วิโนเทยฺยํ พฺยนฺติํ กเรยฺย อนภาวํ คเมยฺย, จกฺขุโลลิยา อารโต อสฺส วิรโต ปฏิวิรโต นิกฺขนฺโต นิสฺสโฏ วิปฺปมุตฺโต วิสญฺญุตฺโต วิมริยาทิกเตน เจตสา วิหเรยฺยาติ จกฺขูหิ เนว โลลสฺส.}

\prose{\textbf{คามกถาย อาวรเย โสตนฺ}ติ คามกถา วุจฺจติ พาตฺติํส ติรจฺฉานกถา.\csromanspacer{} เสยฺยถิทํ, ราชกถา โจรกถา มหามตฺตกถา เสนากถา ภยกถา ยุทฺธกถา อนฺนกถา ปานกถา วตฺถกถา ยานกถา สยนกถา มาลากถา คนฺธกถา ญาติกถา คามกถา นิคมกถา นครกถา ชนปทกถา อิตฺถิกถา\footnote{อิตฺถิกถา ปุริสกถา (พหูสุ) ที 1. 7; ม 2. 181; สํ 3. 368; อํ 3. 358; วิ 2. 213; ม-ฏฺฐ 3. 156 ปิฏฺเฐสุ ปสฺสิตพฺพํ.} สูรกถา วิสิขากถา กุมฺภฏฺฐานกถา ปุพฺพเปตกถา นานตฺตกถา โลกกฺขายิกา สมุทฺทกฺขายิกา อิติภวาภวกถา อิติ วา.\csromanspacer{} \textbf{คามกถาย อาวรเย}}

\vspace{\tipitakaparskip}
\csromancenter{ตุวฏกสุตฺตนิทฺเทส \textbf{โสตนฺ}ติ คามกถาย โสตํ อาวเรยฺย นิวาเรยฺย สํวเรยฺย รกฺเขยฺย โคเปยฺย ปิทเหยฺย ปจฺฉินฺเทยฺยาติ คามกถาย อาวรเย โสตํ.}
\prose{\csromanfolio{287}\textbf{รเส จ นานุคิชฺเฌยฺยา}ติ \textbf{รเสจา}ติ มูลรโส ขนฺธรโส ตจรโส ปตฺตรโส ปุปฺผรโส ผลรโส อมฺพิลํ มธุรํ ติตฺตกํ กฏุกํ โลณิกํ ขาริกํ ลมฺพิกํ กสาโว สาทุ อสาทุ สีตํ อุณฺหํ.\csromanspacer{} สนฺเตเก สมณพฺราหฺมณา รสคิทฺธา, เต ชิวฺหคฺเคน รสคฺคานิ ปริเยสนฺตา อาหิณฺฑนฺติ, เต อมฺพิลํ ลภิตฺวา อนมฺพิลํ ปริเยสนฺติ, อนมฺพิลํ ลภิตฺวา อมฺพิลํ ปริเยสนฺติ ฯเปฯ สีตํ ลภิตฺวา อุณฺหํ ปริเยสนฺติ.\csromanspacer{} อุณฺหํ ลภิตฺวา สีตํ ปริเยสนฺติ.\csromanspacer{} เต ยํ ยํ ลภิตฺวา เตน เตน น ตุสฺสนฺติ, อปราปรํ ปริเยสนฺติ, มนาปิเกสุ รเสสุ รตฺตา คิทฺธา คธิตา มุจฺฉิตา อชฺโฌสนฺนา ลคฺคา ลคฺคิตา ปลิพุทฺธา.\csromanspacer{} ยสฺเสสา รสตณฺหา ปหีนา สมุจฺฉินฺนา ฯเปฯ ญาณคฺคินา ทฑฺฒา, โส ปฏิสงฺขา โยนิโส อาหารํ อาหาเรติ เนว ทวาย ฯเปฯ อนวชฺชตา จ ผาสุวิหาโร จาติ.}

\prose{ยถา วณํ อาลิมฺเปยฺย ยาวเทว อารุหณตฺถาย\footnote{อารุหนตฺถาย (ก), โรปนตฺถาย เหฏฺฐา 185 ปิฏฺเฐ.}.\csromanspacer{} ยถา วา ปน อกฺขํ อมฺภญฺเชยฺย ยาวเทว ภรสฺส นิตฺถรณตฺถาย.\csromanspacer{} ยถา วา ปุตฺตมํสํ อาหารํ อาหาเรยฺย ยาวเทว กนฺตารสฺส นิตฺถรณตฺถาย, เอวเมวํ ภิกฺขุ ปฏิสงฺขา โยนิโส อาหารํ อาหาเรติ เนว ทวาย ฯเปฯ อนวชฺชตา จ ผาสุวิหาโร จาติ รสตณฺหํ ปชเหยฺย วิโนเทยฺย พฺยนฺติํ กเรยฺย อนภาวํ คเมยฺย, รสตณฺหาย อารโต อสฺส วิรโต ปฏิวิรโต นิกฺขนฺโต นิสฺสโฏ วิปฺปมุตฺโต วิสญฺญุตฺโต วิมริยาทิกเตน เจตสา วิหเรยฺยาติ รเส จ นานุคิชฺเฌยฺย.}

\prose{\textbf{น จ มมาเยถ กิญฺจิ โลกสฺมิ}นฺติ \textbf{มมตฺตา}ติ เทฺว \textbf{มมตฺตา} ตณฺหามมตฺตญฺจ ทิฏฺฐิมมตฺตญฺจ ฯเปฯ อิทํ ตณฺหามมตฺตํ ฯเปฯ อิทํ ทิฏฺฐิมมตฺตํ.\csromanspacer{} ตณฺหามมตฺตํ ปหาย ทิฏฺฐิมมตฺตํ ปฏินิสฺสชฺชิตฺวา จกฺขุํ น มมาเยยฺย น คเณฺหยฺย น ปรามเสยฺย นาภินิวิเสยฺย.\csromanspacer{} โสตํ.\csromanspacer{} ฆานํ.\csromanspacer{} ชิวฺหํ.\csromanspacer{} กายํ.\csromanspacer{} รูเป.\csromanspacer{} สทฺเท.\csromanspacer{} คนฺเธ.\csromanspacer{} รเส.\csromanspacer{} โผฏฺฐพฺเพ.\csromanspacer{} กุลํ.\csromanspacer{} คณํ.\csromanspacer{} อาวาสํ.\csromanspacer{} ลาภํ.\csromanspacer{} ยสํ.\csromanspacer{} ปสํสํ.\csromanspacer{} สุขํ.\csromanspacer{} จีวรํ.\csromanspacer{} ปิณฺฑปาตํ.\csromanspacer{} เสนาสนํ.\csromanspacer{} คิลานปจฺจยเภสชฺชปริกฺขารํ.\csromanspacer{} กามธาตุํ.\csromanspacer{} รูปธาตุํ.\csromanspacer{} อรูปธาตุํ.\csromanspacer{} กามภวํ.\csromanspacer{} รูปภวํ.\csromanspacer{} อรูปภวํ.\csromanspacer{} สญฺญาภวํ.\csromanspacer{} อสญฺญาภวํ.\csromanspacer{} เนวสญฺญานาสญฺญาภวํ.\csromanspacer{} เอกโวการภวํ.\csromanspacer{} จตุโวการภวํ. \csromanfolio{288}ปญฺจโวการภวํ.\csromanspacer{} อตีตํ.\csromanspacer{} อนาคตํ.\csromanspacer{} ปจฺจุปฺปนฺนํ.\csromanspacer{} ทิฏฺฐสุตมุตวิญฺญาตพฺเพ ธมฺเม น มมาเยยฺย น คเณฺหยฺย น ปรามเสยฺย นาภินิวิเสยฺย.\csromanspacer{} \textbf{กิญฺจี}ติ กิญฺจิ รูปคตํ เวทนาคตํ สญฺญาคตํ สงฺขารคตํ วิญฺญาณคตํ.\csromanspacer{} \textbf{โลกสฺมินฺ}ติ อปายโลเก ฯเปฯ อายตนโลเกติ น จ มมาเยถ กิญฺจิ โลกสฺมิํ.\csromanspacer{} เตนาห ภควา\csromandash{}}

\csromangathameasure{จกฺขู หิ เนว โลล’สฺส, \\ คามกถาย อาวรเย โสตํ. \\ รเส จ นานุคิชฺเฌยฺย, \\ น จ มมาเยถ กิญฺจิ โลกสฺมินฺติ.}
\csromangathagroup{\csromangathastanza{จกฺขู หิ เนว โลล’สฺส, \\ คามกถาย อาวรเย โสตํ. \\ รเส จ นานุคิชฺเฌยฺย, \\ น จ มมาเยถ กิญฺจิ โลกสฺมินฺติ.}}

\proseitem{158}{\csromansymbolfootnote{*}{ขุ 1. 422 ปิฏฺเฐปิ.}ผสฺเสน ยทา ผุฏฺฐสฺส,}

\prose{\textbf{ปริเทวํ ภิกฺขุ น กเรยฺย กุหิญฺจิ.}}

\csromangathameasure{\textbf{ภวญฺจ นาภิชปฺเปยฺย,} \\ \textbf{เภรเวสุ จ น สมฺปเวเธยฺย.}}
\csromangathagroup{\csromangathastanza{\textbf{ภวญฺจ นาภิชปฺเปยฺย,} \\ \textbf{เภรเวสุ จ น สมฺปเวเธยฺย.}}}

\prose{\textbf{ผสฺเสน ยทา ผุฏฺฐสฺสา}ติ \textbf{ผสฺโส}ติ โรค\textbf{ผสฺโส} โรคผสฺเสน ปุฏฺโฐ ปเรโต สโมหิโต สมนฺนาคโต อสฺส, จกฺขุโรเคน ผุฏฺโฐ ปเรโต สโมหิโต สมนฺนาคโต อสฺส, โสตโรเคน.\csromanspacer{} ฆานโรเคน.\csromanspacer{} ชิวฺหาโรเคน.\csromanspacer{} กายโรเคน.\csromanspacer{} สีสโรเคน.\csromanspacer{} กณฺณโรเคน.\csromanspacer{} มุขโรเคน.\csromanspacer{} ทนฺตโรเคน.\csromanspacer{} กาเสน.\csromanspacer{} สาเสน.\csromanspacer{} ปินาเสน.\csromanspacer{} ทาเหน.\csromanspacer{} ชเรน.\csromanspacer{} กุจฺฉิโรเคน.\csromanspacer{} มุจฺฉาย.\csromanspacer{} ปกฺขนฺทิกาย.\csromanspacer{} สูลาย.\csromanspacer{} วิสูจิกาย.\csromanspacer{} กุฏฺเฐน.\csromanspacer{} คณฺเฑน.\csromanspacer{} กิลาเสน.\csromanspacer{} โสเสน.\csromanspacer{} อปมาเรน.\csromanspacer{} ททฺทุยา.\csromanspacer{} กณฺฑุยา.\csromanspacer{} กจฺฉุยา.\csromanspacer{} รขสาย.\csromanspacer{} วิตจฺฉิกาย.\csromanspacer{} โลหิเตน.\csromanspacer{} ปิตฺเตน.\csromanspacer{} มธุเมเหน.\csromanspacer{} อํสาย.\csromanspacer{} ปิฬกาย.\csromanspacer{} ภคนฺทลาย.\csromanspacer{} ปิตฺตสมุฏฺฐาเนน อาพาเธน.\csromanspacer{} เสมฺหสมุฏฺฐาเนน อาพาเธน.\csromanspacer{} วาตสมุฏฺฐาเนน อาพาเธน.\csromanspacer{} สนฺนิปาติเกน อาพาเธน.\csromanspacer{} อุตุปริณามเชน อาพาเธน.\csromanspacer{} วิสมปริหารเชน อาพาเธน.\csromanspacer{} โอปกฺกมิเกน อาพาเธน.\csromanspacer{} กมฺมวิปากเชน อาพาเธน.\csromanspacer{} สีเตน.\csromanspacer{} อุเณฺหน.\csromanspacer{} ชิฆจฺฉาย.\csromanspacer{} ปิปาสาย.\csromanspacer{} อุจฺจาเรน.\csromanspacer{} ปสฺสาเวน.\csromanspacer{} qอํสมกสวาตาตปสรีสปสมฺผสฺเสหิ ผุฏฺโฐ ปเรโต สโมหิโต สมนฺนาคโต อสฺสาติ ผสฺเสน ยทา ผุฏฺฐสฺส.}

\prose{\textbf{ปริเทวํ ภิกฺขุ น กเรยฺย กุหิญฺจี}ติ อาเทวํ ปริเทวํ อาเทวนํ ปริเทวนํ อาเทวิตตฺตํ ปริเทวิตตฺตํ วาจา ปลาปํ วิปฺปลาปํ ลาลปฺปํ ลาลปฺปายนํ ลาลปฺปายิตตฺตํ น กเรยฺย น ชเนยฺย น สญฺชเนยฺย}

\vspace{\tipitakaparskip}
\csromancenter{ตุวฏกสุตฺตนิทฺเทส น นิพฺพตฺเตยฺย นาภินิพฺพตฺเตยฺย.\csromanspacer{}\csromanspacer{}\textbf{กุหิญฺจี}ติ กุหิญฺจิ กิมฺหิจิ กตฺถจิ อชฺฌตฺตํ วา พหิทฺธา วา อชฺฌตฺตพหิทฺธา วาติ ปริเทวํ ภิกฺขุ น กเรยฺย กุหิญฺจิ.}
\prose{\csromanfolio{289}\textbf{ภวญฺจ นาภิชปฺเปยฺยา}ติ กามภวํ น ชปฺเปยฺย, รูปภวํ น ชปฺเปยฺย, อรูปภวํ น ชปฺเปยฺย น ปชปฺเปยฺย นาภิชปฺเปยฺยาติ ภวญฺจ นาภิชปฺเปยฺย.}

\prose{\textbf{เภรเวสุ จ น สมฺปเวเธยฺยา}ติ \textbf{เภรวา}ติ เอเกนากาเรน ภยมฺปิ เภรวมฺปิ ตญฺเญว.\csromanspacer{} วุตฺตํ เหตํ ภควตา\csromandash{} * เอตํ นูน ตํ ภยเภรวํ อาคจฺฉตีติ.\csromanspacer{} พหิทฺธารมฺมณํ วุตฺตํ, สีหา พฺยคฺฆา ทีปี อจฺฉา ตรจฺฉา โกกา มหิํสา อสฺสา หตฺถี อหิวิจฺฉิกา สตปที, โจรา วา อสฺสุ มานวา วา กตกมฺมา วา อกตกมฺมา วา.\csromanspacer{} อถาปเรน อากาเรน ภยํ วุจฺจติ อชฺฌตฺติกํ จิตฺตสมุฏฺฐานํ ภยํ ภยานกํ ฉมฺภิตตฺตํ โลมหํโส เจตโส อุพฺเพโค อุตฺราโส ชาติภยํ ชราภยํ พฺยาธิภยํ มรณภยํ ราชภยํ โจรภยํ อคฺคิภยํ อุทกภยํ อตฺตานุวาทภยํ ปรานุวาทภยํ ทณฺฑภยํ ทุคฺคติภยํ อูมิภยํ กุมฺภีลภยํ อาวฏฺฏภยํ สุสุกาภยํ\footnote{สุํสุกาภยํ (สฺยา)} อาชีวิกภยํ อสิโลกภยํ ปริสาย สารชฺชภยํ มทนภยํ ทุคฺคติภยํ ภยํ ภยานกํ ฉมฺภิตตฺตํ โลมหํโส เจตโส อุพฺเพโค อุตฺราโส.\csromanspacer{} \textbf{เภรเวสุ จ น สมฺปเวเธยฺยา}ติ เภรเว ปสฺสิตฺวา วา สุณิตฺวา วา น เวเธยฺย น ปเวเธยฺย น สมฺปเวเธยฺย น ตเสยฺย น อุตฺตเสยฺย น ปริตเสยฺย น ภาเยยฺย น สนฺตาสํ อาปชฺเชยฺย, อภีรู อสฺส อฉมฺภี อนุตฺราสี อปลายี, ปหีนภยเภรโว วิคตโลมหํโส วิหเรยฺยาติ เภรเวสุ จ น สมฺปเวเธยฺย.\csromanspacer{} เตนาห ภควา\csromandash{}}

\csromangathameasure{ผสฺเสน ยทา ผุฏฺฐสฺส, \\ ปริเทวํ ภิกฺขุ น กเรยฺย กุหิญฺจิ. \\ ภวญฺจ นาภิชปฺเปยฺย,}
\csromangathagroup{\csromangathastanza{ผสฺเสน ยทา ผุฏฺฐสฺส, \\ ปริเทวํ ภิกฺขุ น กเรยฺย กุหิญฺจิ. \\ ภวญฺจ นาภิชปฺเปยฺย,}}

\prose{\textbf{เภรเวสุ จ น สมฺปเวเธยฺยา}ติ.}

\proseitem{159}{\csromansymbolfootnote{+}{ขุ 1. 422 ปิฏฺเฐปิ.}อนฺนานมโถ ปานานํ,}

\prose{\textbf{ขาทนียานมโถปิ วตฺถานํ.}}

\csromangathameasure{\textbf{ลทฺธา น สนฺนิธิํ กยิรา,} \\ \textbf{น จ ปริตฺตเส ตานิ อลภมาโน.}}
\csromangathagroup{\csromangathastanza{\textbf{ลทฺธา น สนฺนิธิํ กยิรา,} \\ \textbf{น จ ปริตฺตเส ตานิ อลภมาโน.}}}

\prose{\csromanfolio{290}\textbf{อนฺนานมโถ ปานานํ, ขาทนียานมโถปิ วตฺถาน}นฺติ \textbf{อนฺนานนฺ}ติ โอทโน กุมฺมาโส สตฺตุ มจฺโฉ มํสํ.\csromanspacer{} \textbf{ปานานนฺ}ติ * อฏฺฐ ปานานิ อมฺพปานํ ชมฺพุปานํ โจจปานํ โมจปานํ มธุปานํ มุทฺทิกปานํ สาลูกปานํ ผารุสกปานํ.\csromanspacer{} อปรานิปิ อฏฺฐ ปานานิ โกสมฺพปานํ โกลปานํ พทรปานํ ฆตปานํ เตลปานํ ปโยปานํ ยาคุปานํ รสปานํ.\csromanspacer{}\csromanspacer{}\textbf{ขาทนียานนฺ}ติ ปิฏฺฐขชฺชกํ ปูวขชฺชกํ มูลขชฺชกํ ตจขชฺชกํ ปตฺตขชฺชกํ ปุปฺผขชฺชกํ ผลขชฺชกํ.\csromanspacer{} \textbf{วตฺถานนฺ}ติ ฉ จีวรานิ โขมํ กปฺปาสิกํ โกเสยฺยํ กมฺพลํ สาณํ ภงฺคนฺติ อนฺนนมโถ ปานานํ ขาทนียานมโถปิ วตฺถานํ.}

\prose{\textbf{ลทฺธา น สนฺนิธิํ กยิรา}ติ \textbf{ลทฺธา}ติ \textbf{ลทฺธา} ลภิตฺวา อธิคนฺตฺวา วินฺทิตฺวา ปฏิลภิตฺวา น กุหนาย น ลปนาย น เนมิตฺติกตาย น นิปฺเปสิกตาย น ลาเภน ลาภํ นิชิคีสนตาย\footnote{นิชิคิํสนตาย (สี, สฺยา) ขุ 8. 258 ปิฏฺเฐปิ.} น ทารุทาเนน น เวฬุทาเนน น ปตฺตทาเนน น ปุปฺผทาเนน น ผลทาเนน น สินานทาเนน น จุณฺณทาเนน น มตฺติกาทาเนน น ทนฺตกฏฺฐทาเนน น มุโขทกทาเนน น จาฏุกมฺยตาย น มุคฺคสูปฺยตาย\footnote{น มุคฺคสุปฺปตาย (สี), น มุคฺคสูปตาย (สฺยา)} น ปาริภฏฺยตาย น ปีฐมทฺทิกตาย\footnote{ปิฏฺฐิมํสิกตาย (สี), ปิฏฺฐิมทฺทิกตาย (ก)} น วตฺถุวิชฺชาย น ติรจฺฉานวิชฺชาย น องฺควิชฺชาย น นกฺขตฺตวิชฺชาย น ทูตคมเนน น ปหิณคมเนน น ชงฺฆเปสนิเยน น เวชฺชกมฺเมน น นวกมฺเมน น ปิณฺฑปฏิปิณฺฑเกน น ทานานุปฺปทาเนน ธมฺเมน สเมน \textbf{ลทฺธา} ลภิตฺวา อธิคนฺตฺวา วินฺทิตฺวา ปฏิลภิตฺวาติ \textbf{ลทฺธา}.\csromanspacer{} \textbf{น สนฺนิธิํ กยิรา}ติ อนฺนสนฺนิธิํ ปานสนฺนิธิํ วตฺถสนฺนิธิํ ยานสนฺนิธิํ สยนสนฺนิธิํ คนฺธสนฺนิธิํ อามิสสนฺนิธิํ น กเรยฺย น ชเนยฺย น สญฺชเนยฺย น นิพฺพตฺเตยฺย นาภินิพฺพตฺเตยฺยาติ \textbf{ลทฺธา} น สนฺนิธิํ กยิรา.}

\prose{\textbf{น จ ปริตฺตเส ตานิ อลภมาโน}ติ อนฺนํ วา น ลภามิ, ปานํ วา น ลภามิ, วตฺถํ วา น ลภามิ, กุลํ วา น ลภามิ, คณํ วา น ลภามิ, อาวาสํ วา น ลภามิ, ลาภํ วา น ลภามิ, ยสํ วา น ลภามิ, ปสํสํ วา น ลภามิ, สุขํ วา น ลภามิ, จีวรํ วา น ลภามิ, ปิณฺฑปาตํ วา น ลภามิ, เสนาสนํ วา น ลภามิ, คิลานปจฺจย เภสชฺชปริกฺขารํ วา น ลภามิ, คิลานุปฏฺฐากํ วา น ลภามิ “อปฺปญฺญาโตมฺหี”ติ น ตเสยฺย น อุตฺตเสยฺย น ปริตฺตเสยฺย น ภาเยยฺย น สนฺตาสํ}

\vspace{\tipitakaparskip}
\csromancenter{ตุวฏกสุตฺตนิทฺเทส อาปชฺเชยฺย, อภีรู อสฺส อฉมฺภี อนุตฺราสี อปลายี, ปหีนภยเภรโว วิคตโลมหํโส วิหเรยฺยาติ น จ ปริตฺตเส ตานิ อลภมาโน.\csromanspacer{} เตนาห ภควา\csromandash{}}
\vspace{0.5\baselineskip}
\csromangathameasure{อนฺนานมโถ ปานานํ, \\ ขาทนียานมโถปิ วตฺถานํ. \\ ลทฺธา น สนฺนิธิํ กยิรา, \\ น จ ปริตฺตเส ตานิ อลภมาโนติ.}
\csromangathagroup{\csromangathastanza{\csromanfolio{291}อนฺนานมโถ ปานานํ, \\ ขาทนียานมโถปิ วตฺถานํ. \\ ลทฺธา น สนฺนิธิํ กยิรา, \\ น จ ปริตฺตเส ตานิ อลภมาโนติ.}}

\proseitem{160}{\csromansymbolfootnote{*}{ขุ 1. 422 ปิฏฺเฐ.}ฌายี น ปาทโลล’สฺส,}

\prose{\textbf{วิรเม กุกฺกุจฺจา นปฺปมชฺเชยฺย.}}

\csromangathameasure{\textbf{อถาสเนสุ สยเนสุ,} \\ \textbf{อปฺปสทฺเทสุ ภิกฺขุ วิหเรยฺย.}}
\csromangathagroup{\csromangathastanza{\textbf{อถาสเนสุ สยเนสุ,} \\ \textbf{อปฺปสทฺเทสุ ภิกฺขุ วิหเรยฺย.}}}

\prose{\textbf{ฌายี น ปาทโลล’สฺสา}ติ \textbf{ฌายี}ติ ปฐเมนปิ ฌาเนน \textbf{ฌายี}, ทุติเยนปิ ฌาเนน \textbf{ฌายี}, ตติเยนปิ ฌาเนน \textbf{ฌายี}, จตุตฺเถนปิ ฌาเนน \textbf{ฌายี}, สวิตกฺกสวิจาเรนปิ ฌาเนน \textbf{ฌายี}, อวิตกฺกวิจารมตฺเตนปิ ฌาเนน \textbf{ฌายี}, อวิตกฺก-อวิจาเรนปิ ฌาเนน \textbf{ฌายี}, สปฺปีติเกนปิ ฌาเนน \textbf{ฌายี}, นิปฺปีติเกนปิ ฌาเนน \textbf{ฌายี}, ปีติสหคเตนปิ ฌาเนน \textbf{ฌายี}, สาตสหคเตนปิ ฌาเนน \textbf{ฌายี}, สุขสหคเตนปิ ฌาเนน \textbf{ฌายี}, อุเปกฺขาสหคเตนปิ ฌาเนน \textbf{ฌายี}, สุญฺญเตนปิ ฌาเนน \textbf{ฌายี}, อนิมิตฺเตนปิ ฌาเนน \textbf{ฌายี}, อปฺปณิหิเตนปิ ฌาเนน \textbf{ฌายี}, โลกิเยนปิ ฌาเนน \textbf{ฌายี}, โลกุตฺตเรนปิ ฌาเนน \textbf{ฌายี} ฌานรโต เอกตฺตมนุยุตฺโต ปรมตฺถครุโกติ\footnote{สทตฺถครุโกติ (สี, สฺยา)} \textbf{ฌายี}.}

\prose{\textbf{น ปาทโลล’สฺสา}ติ กถํ ปาทโลโล โหติ, อิเธกจฺโจ ปาทโลลิเยน สมนฺนาคโต โหติ, อาราเมน อารามํ, อุยฺยาเนน อุยฺยานํ, คาเมน คามํ, นิคเมน นิคมํ, นคเรน นครํ, รฏฺเฐน รฏฺฐํ, ชนปเทน ชนปทํ ทีฆจาริกํ อนวฏฺฐิตจาริกํ อนุยุตฺโต วิหรติ.\csromanspacer{} เอวมฺปิ ปาทโลโล โหติ.}

\prose{อถ วา ภิกฺขุ อนฺโตปิสํฆาราเม ปาทโลลิเยน สมนฺนาคโต โหติ, น อตฺถเหตุ น การณเหตุ อุทฺธโต อวูปสนฺตจิตฺโต ปริเวณโต ปริเวณํ คจฺฉติ, วิหารโต วิหารํ คจฺฉติ, \csromanfolio{292}อฑฺฒโยคโต อฑฺฒโยคํ คจฺฉติ, ปาสาทโต ปาสาทํ คจฺฉติ, หมฺมิยโต หมฺมิยํ คจฺฉติ, คุหาย คุหํ คจฺฉติ, เลณโต เลณํ คจฺฉติ, กุฏิโต กุฏิํ คจฺฉติ, กูฏาคารโต กูฏาคารํ คจฺฉติ, อฏฺฏโต อฏฺฏํ คจฺฉติ, มาฬโต มาฬํ คจฺฉติ, อุทฺทณฺฑโต อุทฺทณฺฑํ คจฺฉติ, อุปฏฺฐานสาลโต อุปฏฺฐานสาลํ คจฺฉติ, มณฺฑลมาฬโต มณฺฑลมาฬํ คจฺฉติ, รุกฺขมูลโต รุกฺขมูลํ คจฺฉติ, ยตฺถ วา ปน ภิกฺขู นิสีทนฺติ, ตหิํ คจฺฉติ, ตตฺถ เอกสฺส วา ทุติโย โหติ.\csromanspacer{} ทฺวินฺนํ วา ตติโย โหติ, ติณฺณํ วา จตุตฺโถ โหติ, ตตฺถ พหุํ สมฺผปฺปลาปํ ปลปติ.\csromanspacer{} เสยฺยถิทํ.\csromanspacer{} ราชกถํ โจรกถํ มหามตฺตกถํ เสนากถํ ภยกถํ ยุทฺธกถํ อนฺนกถํ ปานกถํ วตฺถกถํ ยานกถํ สยนกถํ มาลากถํ คนฺธกถํ ญาติกถํ คามกถํ นิคมกถํ นครกถํ ชนปทกถํ อิตฺถิกถํ สูรกถํ วิสิขากถํ กุมฺภฏฺฐานกถํ ปุพฺพเปตกถํ นานตฺตกถํ โลกกฺขายิกํ สมุทฺทกฺขายิกํ อิติ ภวาภวกถํ อิติ วา, เอวมฺปิ ปาทโลโล โหติ.}

\prose{\textbf{น ปาทโลล’สฺสา}ติ ปาทโลลิยํ ปชเหยฺย วิโนเทยฺย พฺยนฺติํ กเรยฺย อนภาวํ คเมยฺย, ปาทโลลิยา อารโต อสฺส วิรโต ปฏิวิรโต นิกฺขนฺโต นิสฺสโฏ วิปฺปมุตฺโต วิสญฺญุตฺโต วิมริยาทิกเตน เจตสา วิหเรยฺย จเรยฺย วิจเรยฺย อิริเยยฺย วตฺเตยฺย ปาเลยฺย ยเปยฺย ยาเปยฺย ปฏิสลฺลานาราโม อสฺส ปฏิสลฺลานรโต อชฺฌตฺตํ เจโตสมถมนุยุตฺโต อนิรากตชฺฌาโน วิปสฺสนาย สมนฺนาคโต พฺรูเหตา สุญฺญาคารานํ ฌายี ฌานรโต เอกตฺตมนุยุตฺโต ปรมตฺถครุโกติ ฌายี น ปาทโลลสฺส.}

\prose{\textbf{วิรเม กุกฺกุจฺจา นปฺปมชฺเชยฺยา}ติ \textbf{กุกฺกุจฺจนฺ}ติ หตฺถกุกฺกุจฺจมฺปิ กุกฺกุจฺจํ, ปาทกุกฺกุจฺจมฺปิ กุกฺกุจฺจํ, หตฺถปาทกุกฺกุจฺจมฺปิ กุกฺกุจฺจํ, อกปฺปิเย กปฺปิยสญฺญิตา, กปฺปิเย อกปฺปิยสญฺญิตา,วิกาเล กาลสญฺญิตา, กาเล วิกาลสญฺญิตา, อวชฺเช วชฺชสญฺญิตา, วชฺเช อวชฺชสญฺญิตา, ยํ เอวรูปํ กุกฺกุจฺจํ กุกฺกุจฺจายนา กุกฺกุจฺจายิตตฺตํ เจตโส วิปฺปฏิสาโร มโนวิเลโข, อิทํ วุจฺจติ กุกฺกุจฺจํ.}

\prose{อปิ จ ทฺวีหิ กาเรเณหิ อุปฺปชฺชติ กุกฺกุจฺจํ เจตโส วิปฺปฏิสาโร มโนวิเลโข กตตฺตา จ อกตตฺตา จ.\csromanspacer{} กถํ กตตฺตา จ อกตตฺตา จ อุปฺปชฺชติ กุกฺกุจฺจํ เจตโส วิปฺปฏิสาโร มโนวิเลโข\csromandash{}“กตํ เม กายทุจฺจริตํ, อกตํ เม กายสุจริตนฺ”ติ อุปฺปชฺชติ กุกฺกุจฺจํ}

\vspace{\tipitakaparskip}
\csromancenter{ตุวฏกสุตฺตนิทฺเทส เจตโส วิปฺปฏิสาโร มโนวิเลโข. “กตํ เม วจีทุจฺจริตํ, อกตํ เม วจีสุจริตนฺ”ติ อุปฺปชฺชติ กุกฺกุจฺจํ เจตโส วิปฺปฏิสาโร มโนวิเลโข. “กตํ เม มโนทุจฺจริตํ, อกตํ เม มโนสุจริตํ.\csromanspacer{} กโต เม ปาณาติปาโต, อกตา เม ปาณาติปาตา เวรมณี “ติ อุปฺปชฺชติ กุกฺกุจฺจํ ฯเปฯ มโนวิเลโข, “กตํ เม อทินฺนาทานํ.\csromanspacer{} กโต เม กาเมสุมิจฺฉาจาโร.\csromanspacer{} กโต เม มุสาวาโท.\csromanspacer{} กตา เม ปิสุณา วาจา.\csromanspacer{} กตา เม ผรุสา วาจา.\csromanspacer{} กโต เม สมฺผปฺปลาโป.\csromanspacer{} กตา เม อภิชฺฌา.\csromanspacer{} กโต เม พฺยาปาโท.\csromanspacer{} กตา เม มิจฺฉาทิฏฺฐิ, อกตา เม สมฺมาทิฏฺฐี”ติ อุปฺปชฺชติ กุกฺกุจฺจํ เจตโส วิปฺปฏิสาโร มโนวิเลโข.\csromanspacer{} เอวํ กตตฺตา จ อกตตฺตา จ อุปฺปชฺชติ กุกฺกุจฺจํ เจตโส วิปฺปฏิสาโร มโนวิเลโข.}
\prose{\csromanfolio{293}อถ วา “สีเลสุ’มฺหิ น ปริปูรการี”ติ อุปฺปชฺชติ กุกฺกุจฺจํ เจตโส วิปฺปฏิสาโร มโนวิเลโข, “อินฺทฺริเยสุ’มฺหิ อคุตฺตทฺวาโร”ติ. “โภชเน อมตฺตญฺญุ’มฺหี”ติ. “ชาคริยํ อนนุยุตฺโต’มฺหี”ติ. “น สติสมฺปชญฺเญน สมนฺนาคโต’มฺหี”ติ. “อภาวิตา เม จตฺตาโร สติปฏฺฐานา”ติ. “อภาวิตา เม จตฺตาโร สมฺมปฺปธานา”ติ. “อภาวิตา เม จตฺตาโร อิทฺธิปาทา”ติ. “อภาวิตานิ เม ปญฺจินฺทฺริยานี”ติ. “อภาวิตานิ เม ปญฺจ พลานี”ติ. “อภาวิตา เม สตฺต โพชฺฌงฺคา”ติ. “อภาวิโต เม อริโย อฏฺฐงฺคิโก มคฺโค”ติ. “ทุกฺขํ เม อปริญฺญาตนฺ”ติ. “สมุทโย เม อปฺปหีโน”ติ. “มคฺโค เม อภาวิโต”ติ. “นิโรโธ เม อสจฺฉิกโต”ติ อุปฺปชฺชติ กุกฺกุจฺจํ เจตโส วิปฺปฏิสาโร มโนวิเลโข.\csromanspacer{} \textbf{วิรเม กุกฺกุจฺจา}ติ กุกฺกุจฺจา อารเมยฺย วิรเมยฺย ปฏิวิรเมยฺย กุกฺกุจฺจํ ปชเหยฺย วิโนเทยฺย พฺยนฺติํ กเรยฺย อนภาวํ คเมยฺย.\csromanspacer{} กุกฺกุจฺจา อารโต อสฺส วิรโต ปฏิวิรโต นิกฺขนฺโต นิสฺสโฏ วิปฺปมุตฺโต วิสญฺญุตฺโต วิมริยาทิกเตน เจตสา วิหเรยฺยาติ วิรเม กุกฺกุจฺจา.}

\prose{\textbf{นปฺปมชฺเชยฺยา}ติ สกฺกจฺจการี อสฺส สาตจฺจการี อฏฺฐิตการี อโนลีนวุตฺติโก อนิกฺขิตฺตจฺฉนฺโท อนิกฺขิตฺตธุโร อปฺปมาโท กุสเลสุ ธมฺเมสุ. “กทาหํ อปริปูรํ วา สีลกฺขนฺธํ ปริปูเรยฺยํ, ปริปูรํ วา สีลกฺขนฺธํ ตตฺถ ตตฺถ ปญฺญาย อนุคฺคเณฺหยฺยนฺ”ติ โย ตตฺถ ฉนฺโท จ วายาโม จ อุสฺสาโห จ อุสฺโสฬฺหี จ ถาโม จ อปฺปฏิวานี จ สติ จ สมฺปชญฺญญฺจ อาตปฺปํ ปธานํ อธิฏฺฐานํ อนุโยโค อปฺปมาโท กุสเลสุ \csromanfolio{294}ธมฺเมสุ. “กทาหํ อปริปูรํ วา สมาธิกฺขนฺธํ.\csromanspacer{} ปญฺญากฺขนฺธํ.\csromanspacer{} วิมุตฺติกฺขนฺธํ.\csromanspacer{} วิมุตฺติญาณทสฺสนกฺขนฺธํ.\csromanspacer{} กทาหํ อปริญฺญาตํ วา ทุกฺขํ ปริชาเนยฺยํ, อปฺปหีเน วา กิเลเส ปชเหยฺยํ, อภาวิตํ วา มคฺคํ ภาเวยฺยํ, อสจฺฉิกตํ วา นิโรธํ สจฺฉิกเรยฺยนฺ”ติ โย ตตฺถ ฉนฺโท จ วายาโม จ อุสฺสาโห จ อุสฺโสฬฺหี จฺ\textbf{อ ถา}โม จ อปฺปฏิวานี จ สติ จ สมฺปชญฺญญฺจ \textbf{อาตปฺปํ} ปธานํ อธิฏฺฐานํ อนุโยโค อปฺปมาโท กุสเลสุ ธมฺเมสูติ วิรเม กุกฺกุจฺจา นปฺปมชฺเชยฺย.}

\prose{\textbf{อถาสเนสุ สยเนสุ, อปฺปสทฺเทสุ ภิกฺขุ วิหเรยฺยา}ติ \textbf{อถา}ติ ปทสนฺธิ ฯเปฯ.\csromanspacer{} \textbf{อาสนํ} วุจฺจติ ยตฺถ นิสีทติ มญฺโจ ปีฐํ ภิสิ ตฏฺฏิกา จมฺมขณฺโฑ ติณสนฺถาโร ปณฺณสนฺถาโร ปลาลสนฺถาโร, \textbf{สยนํ} วุจฺจติ เสนาสนํ วิหาโร อฑฺฒโยโค ปาสาโท หมฺมิยํ คุหาติ \textbf{อถา}สเนสุ สยเนสุ.}

\prose{\textbf{อปฺปสทฺเทสุ ภิกฺขุ วิหเรยฺยา}ติ อปฺปสทฺเทสุ อปฺปนิคฺโฆเสสุ วิชนวาเตสุ มนุสฺสาราหสฺเสยฺยเกสุ ปฏิสลฺลานสารุปฺเปสุ เสนาสเนสุ จเรยฺย วิจเรยฺย วิหเรยฺย อิริเยยฺย วตฺเตยฺย ปาเลยฺย ยเปยฺย ยาเปยฺยาติ \textbf{อถา}สเนสุ สยเนสุ, อปฺปสทฺเทสุ ภิกฺขุ วิหเรยฺย.\csromanspacer{} เตนาห ภควา\csromandash{}}

\csromangathameasure{ฌายี น ปาทโลล’สฺส, \\ วิรเม กุกฺกุจฺจา นปฺปมชฺเชยฺย.}
\csromangathagroup{\csromangathastanza{ฌายี น ปาทโลล’สฺส, \\ วิรเม กุกฺกุจฺจา นปฺปมชฺเชยฺย.}}

\prose{อถาสเนสุ สยเนสุ,}

\prose{\textbf{อปฺปสทฺเทสุ ภิกฺขุ วิหเรยฺยา}ติ.}

\proseitem{161}{\csromansymbolfootnote{*}{ขุ 1. 422 ปิฏฺเฐปิ.}นิทฺทํ น พหุลีกเรยฺย,}

\prose{\textbf{ชาคริยํ ภเชยฺย อาตาปี.}}

\csromangathameasure{\textbf{ตนฺทิํ มายํ หสฺสํ ขิฑฺฑํ,} \\ \textbf{เมถุนํ วิปฺปชเห สวิภูสํ.}}
\csromangathagroup{\csromangathastanza{\textbf{ตนฺทิํ มายํ หสฺสํ ขิฑฺฑํ,} \\ \textbf{เมถุนํ วิปฺปชเห สวิภูสํ.}}}

\prose{\textbf{นิทฺทํ น พหุลีกเรยฺยา}ติ รตฺตินฺทิวํ ฉโกฏฺฐาสํ ถาเรตฺวา ปญฺจโกฏฺฐาสํ ปฏิชคฺเคยฺย, เอกโกฏฺฐาสํ นิปฺปชฺเชยฺยาติ นิทฺทํ น พหุลีกเรยฺย.}

\prose{\textbf{ชาคริยํ ภเชยฺย อาตาปี}ติ อิธ ภิกฺขุ ทิวสํ จงฺกเมน นิสชฺชาย อาวรณีเยหิ ธมฺเมหิ จิตฺตํ ปริโสเธยฺย.\csromanspacer{} รตฺติยา ปฐมํ ยามํ จงฺกเมน นิสชฺชาย อาวรณีเยหิ ธมฺเมหิ จิตฺตํ ปริโสเธยฺย, รตฺติยา มชฺฌิมํ}

\vspace{\tipitakaparskip}
\csromancenter{ตุวฏกสุตฺตนิทฺเทส ยามํ ทกฺขิเณน ปสฺเสน สีหเสยฺยํ กปฺเปยฺย ปาเทปาทํ อจฺจาธาย สโต สมฺปชาโน อุฏฺฐานสญฺญํ มนสิกริตฺวา, รตฺติยา ปจฺฉิมํ ยามํ ปจฺจุฏฺฐาย จงฺกเมน นิสชฺชาย อาวรณีเยหิ ธมฺเมหิ จิตฺตํ ปริโสเธยฺย.}
\prose{\csromanfolio{295}\textbf{ชาคริยํ ภเชยฺยา}ติ ชาคริยํ ภเชยฺย สมฺภเชยฺย เสเวยฺย นิเสเวยฺย สํเสเวยฺย ปฏิเสเวยฺยาติ ชาคริยํ ภเชยฺย.}

\prose{\textbf{อาตาปี}ติ \textbf{อาตปฺปํ} วุจฺจติ วีริยํ\footnote{วิริยํ (สี, สฺยา)}.\csromanspacer{} โย เจตสิโก วีริยารมฺโภ นิกฺกโม ปรกฺกโม อุยฺยาโม วายาโม อุสฺสาโห อุสฺโสฬฺหี ถาโม ฐิติ อสิถิลปรกฺกมตา อนิกฺขิตฺตจฺฉนฺทตา อนิกฺขิตฺตธุรตา ธุรสมฺปคฺคโห วีริยํ วีริยินฺทฺริยํ วีริยพลํ สมฺมาวายาโม.\csromanspacer{} อิมินา อาตาเปน อุเปโต สมุเปโต อุปคโต สมุปคโต อุปปนฺโน สมุปปนฺโน สมนฺนาคโต, โส วุจฺจติ อาตาปีติ ชาคริยํ ภเชยฺย อาตาปี.}

\prose{\textbf{ตนฺทิํ มายํ หสฺสํ ขิฑฺฑํ, เมถุนํ วิปฺปชเห สวิภูสนฺ}ติ \textbf{ตนฺที}ติ\footnote{ตนฺทินฺติ (สี)} \textbf{ตนฺที} ตนฺทิยนา ตนฺทิยิตตฺตํ\footnote{นตฺถิ อตฺตทณฺฑสุตฺตนิทฺเทเส ขุทฺทกวตฺถุวิภงฺเคปิ.} ตนฺทิมนกตา อาลสฺยํ อาลสฺยายนา อาลสฺยายิตตฺตํ.\csromanspacer{} \textbf{มายา} วุจฺจติ วญฺจนิกา จริยา.\csromanspacer{} อิเธกจฺโจ กาเยน ทุจฺจริตํ จริตฺวา.\csromanspacer{} วาจาย ฯเปฯ.\csromanspacer{} มนสา ทุจฺจริตํ จริตฺวา ตสฺส ปฏิจฺฉาทนเหตุ ปาปิกํ อิจฺฉํ ปณิทหติ. “มา มํ ชญฺญา”ติ\footnote{ชญฺญุภิ (ก) อภิ 2. 372 ปิฏฺเฐ.} อิจฺฉติ, “มา มํ ชญฺญา”ติ สงฺกปฺเปติ, “มา มํ ชญฺญา”ติ วาจํ ภาสติ, “มา มํ ชญฺญา”ติ กาเยน ปรกฺกมติ.\csromanspacer{} ยา เอวรูปา มายาวิตา อจฺจสรา วญฺจนา นิกติ นิกิรณา ปริหรณา คูหนา ปริคูหนา ฉาทนา ปริจฺฉาทนา\footnote{ปฏิจฺฉาทนา (สี, ก) 59 ปิฏฺเฐปิ.} อนุตฺตานีกมฺมํ อนาวิกมฺมํ โวจฺฉาทนา ปาปกิริยา.\csromanspacer{} อยํ วุจฺจติ มายา.\csromanspacer{}\csromanspacer{}\textbf{หสฺสนฺ}ติ อิเธกจฺโจ อติเวลํ ทนฺตวิทํสกํ หสติ.\csromanspacer{} วุตฺตํ เหตํ ภควตา\csromandash{} * โกมารกมิทํ ภิกฺขเว อริยสฺส วินเย ยทิทํ อติเวลํ ทนฺตวิทํสกํ หสิตนฺติ.}

\prose{\textbf{ขิฑฺฑา}ติ เทฺว ขิฑฺฑา กายิกา จ ขิฑฺฑา วาจสิกา จ ขิฑฺฑา, กตมา กายิกา ขิฑฺฑา, หตฺตีหิปิ กีฬนฺติ, อสฺเสหิปิ กีฬนฺติ, รเถหิปิ กีฬนฺติ, ธนูหิปิ กีฬนฺติ, อฏฺฐปเทหิปิ กีฬนฺติ, ทสปเทหิปิ กีฬนฺติ, อากาเสปิ กีฬนฺติ, ปริหารปเถปิ กีฬนฺติ, สนฺติกายปิ กีฬนฺติ, ขลิกายปิ กีฬนฺติ, ฆฏิกายปิ กีฬนฺติ, สลากหตฺเถนปิ กีฬนฺติ, อกฺเขนปิ กีฬนฺติ, ปงฺกจีเรนปิ \csromanfolio{296}กีฬนฺติ, วงฺกเกนปิ กีฬนฺติ, โมกฺขจิกายปิ กีฬนฺติ, จิงฺคุลเกนปิ กีฬนฺติ, ปตฺตาฬฺหเกนปิ กีฬนฺติ, รถเกนปิ กีฬนฺติ, ธนุเกนปิ กีฬนฺติ, อกฺขริกายปิ กีฬนฺติ, มเนสิกายปิ กีฬนฺติ, ยถาวชฺเชนปิ กีฬนฺติ, อยํ กายิกา ขิฑฺฑา.\csromanspacer{}\csromanspacer{}กตมา วาจสิกา ขิฑฺฑา, มุขเภริกํ\footnote{มุขเภริยํ (สฺยา)} มุขาลมฺพรํ มุขฑิณฺฑิมกํ\footnote{มุขเทณฺฑิมกํ (สี, สฺยา)} มุขวลิมกํ มุขเภรุฬกํ\footnote{มุขเภรุลกํ (สี)} มุขททฺทริกํ นาฏกํ ลาปํ คีตํ ทวกมฺมํ, อยํ วาจสิกา ขิฑฺฑา.}

\prose{เมถุนธมฺโม นาม โย โส อสทฺธมฺโม คามธมฺโม วสลธมฺโม ทุฏฺฐุลฺโล โอทกนฺติโก รหสฺโส ทฺวยํทฺวยสมาปตฺติ.\csromanspacer{} กิํการณา วุจฺจติ เมถุนธมฺโม.\csromanspacer{} อุภินฺนํ รตฺตานํ สารตฺตานํ อวสฺสุตานํ ปริยุฏฺฐิตานํ ปริยาทินฺนจิตฺตานํ อุภินฺนํ สทิสานํ ธมฺโมติ ตํการณา วุจฺจติ เมถุนธมฺโม.\csromanspacer{} ยถา อุโภ กลหการกา เมถุนกาติ วุจฺจนฺติ, อุโภ ภณฺฑนการกา.\csromanspacer{} อุโภ ภสฺสการกา.\csromanspacer{} อุโภ อธิกรณการกา.\csromanspacer{} อุโภ วิวาทการกา.\csromanspacer{} อุโภ วาทิโน.\csromanspacer{} อุโภ สลฺลาปกา เมถุนกาติ วุจฺจนฺติ.\csromanspacer{} เอวเมวํ อุภินฺนํ รตฺตานํ สารตฺตานํ อวสฺสุตานํ ปริยุฏฺฐิตานํ ปริยาทินฺนจิตฺตานํ อุภินฺนํ สทิสานํ ธมฺโมติ ตํการณา วุจฺจติ เมถุนธมฺโม.}

\prose{\textbf{วิภูสา}ติ เทฺว วิภูสา.\csromanspacer{} อตฺถิ อคาริยสฺส วิภูสา, อตฺถิ ปพฺพชิตสฺส วิภูสา.\csromanspacer{} กตมา อคาริยสฺส วิภูสา, เกสา จ มสฺสุ จ มาลา จ คนฺธา จ วิเลปนา จ อาภรณา จ ปิลนฺธนา จ วตฺถญฺจ สยนาสนญฺจ เวฐนญฺจ อุจฺฉาทนํ ปริมทฺทนํ นฺหาปนํ สมฺพาหนํ อาทาสํ อญฺชนํ มาลาวิเลปนํ มุขจุณฺณกํ มุขเลปํ หตฺถิพนฺธํ สิขาพนฺธํ\footnote{วสิกฺขาพนฺธํ (สฺยา)} ทณฺฑนาฬิยํ ขคฺคํ ฉตฺตํ จิตฺรา อุปาหนา อุณฺหีสํ มณิํ วาฬพีชนิํ โอทาตานิ วตฺถานิ ทีฆทสานิ อิติ วา อยํ อคาริยสฺส วิภูสา.\csromanspacer{}\csromanspacer{}กตมา ปพฺพชิตสฺส วิภูสา, จีวรมณฺฑนา ปตฺตมณฺฑนา เสนาสนมณฺฑนา อิมสฺส วา ปูติกายสฺส พาหิรานํ วา ปริกฺขารานํ มณฺฑนา วิภูสนา เกฬนา ปริเกฬนา เคธิตตา เคธิตตฺตํ จปลตา จาปลฺยํ, อยํ ปพฺพชิตสฺส วิภูสา.}

\prose{\textbf{ตนฺทิํ มายํ หสฺสํ ขิฑฺฑํ, เมถุนํ วิปฺปชเห สวิภูสนฺ}ติ ตนฺทิญฺจ มายญฺจ หสฺสญฺจ ขิฑฺฑญฺจ เมถุนธมฺมญฺจ สวิภูสํ สปริวารํ สปริภณฺฑํ สปริกฺขารํ ปชเหยฺย วิโนเทยฺย พฺยนฺติํ กเรยฺย อนภาวํ คเมยฺยาติ ตนฺทิํ มายํ หสฺสํ ขิฑฺฑํ, เมถุนํ วิปฺปชเห สวิภูสํ.\csromanspacer{} เตนาห ภควา\csromandash{}}

\csromanheader{2}{14. ตุวฏกสุตฺตนิทฺเทส}
\csromangathameasure{นิทฺทํ น พหุลีกเรยฺย, \\ ชาคริยํ ภเชยฺย อาตาปี. \\ ตนฺทิํ มายํ หสฺสํ ขิฑฺฑํ, \\ เมถนํ วิปฺปชเห สวิภูสนฺติ.}
\csromangathagroup{\csromangathastanza{\csromanfolio{297}นิทฺทํ น พหุลีกเรยฺย, \\ ชาคริยํ ภเชยฺย อาตาปี. \\ ตนฺทิํ มายํ หสฺสํ ขิฑฺฑํ, \\ เมถนํ วิปฺปชเห สวิภูสนฺติ.}}

\proseitem{162}{\csromansymbolfootnote{*}{ขุ 1. 423 ปิฏฺเฐ.}อาถพฺพณํ สุปินํ ลกฺขณํ,}

\prose{\textbf{โน วิทเห อโถปิ นกฺขตฺตํ.}}

\csromangathameasure{\textbf{วิรุตญฺจ คพฺภกรณํ,} \\ \textbf{ติกิจฺฉํ มามโก น เสเวยฺย.}}
\csromangathagroup{\csromangathastanza{\textbf{วิรุตญฺจ คพฺภกรณํ,} \\ \textbf{ติกิจฺฉํ มามโก น เสเวยฺย.}}}

\prose{\textbf{อาถพฺพณํ สุปินํ ลกฺขณํ, โน วิทเห อโถปิ นกฺขตฺตนฺ}ติ อาถพฺพณิกา อาถพฺพณํ ปโยเชนฺติ, นคเร วา รุทฺเธ\footnote{รุนฺเธ (ก)} สงฺคาเม วา ปจฺจุปฏฺฐิเต ปรเสนปจฺจตฺถิเกสุ ปจฺจามิตฺเตสุ อีติํ อุปฺปาเทนฺติ, อุปทฺทวํ อุปฺปาเทนฺติ, โรคํ อุปฺปาเทนฺติ, ปชฺชรกํ กโรนฺติ, สูลํ กโรนฺติ, วิสูจิกํ กโรนฺติ, ปกฺขนฺทิกํ กโรนฺติ.\csromanspacer{} เอวํ อาถพฺพณิกา อาถพฺพณํ ปโยเชนฺติ.}

\prose{สุปินปาฐกา สุปินํ อาทิสนฺติ “โย ปุพฺพณฺหสมยํ สุปินํ ปสฺสติ, เอวํ วิปาโก โหติ.\csromanspacer{} โย มชฺฌนฺหิกสมยํ\footnote{มชฺฌนฺติกสมยํ (พหูสุ)} สุปินํ ปสฺสติ, เอวํ วิปาโก โหติ.\csromanspacer{} โย สายนฺหสมยํ สุปินํ ปสฺสติ, เอวํ วิปาโก โหติ.\csromanspacer{} โย ปุริเม ยาเม.\csromanspacer{} โย มชฺฌิเม ยาเม.\csromanspacer{} โย ปจฺฉิเม ยาเม.\csromanspacer{} โย ทกฺขิเณน ปสฺเสน นิปนฺโน.\csromanspacer{} โย วาเมน ปสฺเสน นิปนฺโน.\csromanspacer{} โย อุตฺตานํ นิปนฺโน.\csromanspacer{} โย อวกุชฺช นิปนฺโน.\csromanspacer{} โย จนฺทํ ปสฺสติ.\csromanspacer{} โย สูริยํ ปสฺสติ.\csromanspacer{} โย มหาสมุทฺทํ ปสฺสติ.\csromanspacer{} โย สิเนรุํ ปพฺพตราชานํ ปสฺสติ.\csromanspacer{} โย หตฺถิํ ปสฺสติ.\csromanspacer{} โย อสฺสํ ปสฺสติ.\csromanspacer{} โย รถํ ปสฺสติ.\csromanspacer{} โย ปตฺติํ ปสฺสติ.\csromanspacer{} โย เสนาพฺยูหํ ปสฺสติ.\csromanspacer{} โย อารามรามเณยฺยกํ ปสฺสติ.\csromanspacer{} โย วนรามเณยฺยกํ ปสฺสติ.\csromanspacer{} โย ภูมิรามเณยฺยกํ ปสฺสติ.\csromanspacer{} โย โปกฺขรณีรามเณยฺยกํ\footnote{โปกฺขรณิรามเณยฺยกํ (สฺยา)} ปสฺสติ, เอวํ วิปาโก โหตี”ติ.\csromanspacer{} เอวํ สุปินปาฐกา สุปินํ อาทิสนฺติ.}

\prose{ลกฺขณปาฐกา ลกฺขณํ อาทิสนฺติ, มณิลกฺขณํ ทณฺฑลกฺขณํ วตฺถลกฺขณํ อสิลกฺขณํ อุสุลกฺขณํ ธนุลกฺขณํ อาวุธลกฺขณํ อิตฺถิลกฺขณํ ปุริสลกฺขณํ กุมาริกาลกฺขณํ กุมารลกฺขณํ ทาสิลกฺขณํ ทาสลกฺขณํ หตฺถิลกฺขณํ อสฺสลกฺขณํ มหิํสลกฺขณํ อุสภลกฺขณํ \csromanfolio{298}โคณลกฺขณํ อชลกฺขณํ เมณฺฑลกฺขณํ กุกฺกุฏลกฺขณํ วฏฺฏลกฺขณํ โคธาลกฺขณํ กณฺณิกาลกฺขณํ กจฺฉปลกฺขณํ มิคลกฺขณํ อิติ วาติ, เอวํ ลกฺขณปาฐกา ลกฺขณํ อาทิสนฺติ.}

\prose{นกฺขตฺตปาฐกา นกฺขตฺตํ อาทิสนฺติ, อฏฺฐวีสติ นกฺขตฺตานิ อิมินา นกฺขตฺเตน ฆรปฺปเวโส กตฺตพฺโพ, อิมินา นกฺขตฺเตน มกุฏํ พนฺธิตพฺพํ, อิมินา นกฺขตฺเตน วาเรยฺยํ กาเรตพฺพํ, อิมินา นกฺขตฺเตน พีชนีหาโร\footnote{พีชนิหาโร (สฺยา, ก)} กตฺตพฺโพ, อิมินา นกฺขตฺเตน สํวาโส\footnote{ฆรวาโส (สฺยา) 3. มิคจกฺกํ. มิคจกฺกปาฐกา (สฺยา)} คนฺตพฺโพติ.\csromanspacer{} เอวํ นกฺขตฺตปาฐกา นกฺขตฺตํ อาทิสนฺติ.}

\prose{\textbf{อาถพฺพณํ สุปินํ ลกฺขณํ, โน วิทเห อโถปิ นกฺขตฺตนฺ}ติ อาถพฺพณญฺจ สุปินญฺจ ลกฺขณญฺจ นกฺขตฺตญฺจ โน วิทเหยฺย น จเรยฺย น สมาจเรยฺย น สมาทาย วตฺเตยฺย.\csromanspacer{} อถ วา น คเณฺหยฺย น อุคฺคเณฺหยฺย น ธาเรยฺย น อุปธาเรยฺย น อุปลกฺเขยฺย นปฺปโยเชยฺยาติ อาถพฺพณํ สุปินํ ลกฺขณํ, โน วิทเห อโถปิ นกฺขตฺตํ.}

\prose{\textbf{วิรุตญฺจ คพฺภกรณํ, ติกิจฺฉํ มามโก น เสเวยฺยา}ติ \textbf{วิรุตํ} วุจฺจติ มิควากฺกํ.\csromanspacer{} มิควากฺกปาฐกา3 มิควากฺกํ อาทิสนฺติ, สกุนฺตานํ วา จตุปฺปทานํ วา รุตํ\footnote{รุทํ (สฺยา)} วสฺสิตํ ชานนฺตีติ.\csromanspacer{} เอวํ มิควากฺกปาฐกา มิควากฺกํ อาทิสนฺติ.\csromanspacer{}\csromanspacer{}คพฺภกรณียา คพฺภํ สณฺฐาเปนฺติ, ทฺวีหิ การเณหิ คพฺโภ น สณฺฐาติ ปาณเกหิ วา วาตกุปฺเปหิ วา, ปาณกานํ วา วาตกุปฺปานํ วา ปฏิฆาตาย โอสธํ เทนฺตีติ.\csromanspacer{} เอวํ คพฺภกรณียา คพฺภํ สณฺฐาเปนฺติ.\csromanspacer{}\csromanspacer{}\textbf{ติกิจฺฉา}ติ ปญฺจ ติกิจฺฉา สาลากิยํ สลฺลกตฺติยํ กายติกิจฺฉํ ภูติยํ โกมารภจฺจํ.\csromanspacer{} \textbf{มามโก}ติ พุทฺธมามโก ธมฺมมามโก สํฆมามโก, โส วา ภควนฺตํ มมายติ, ภควา วา ตํ ปุคฺคลํ ปริคฺคณฺหาติ.\csromanspacer{} วุตฺตํ เหตํ ภควตา\csromandash{} * เย เต ภิกฺขเว ภิกฺขู กุหา ถทฺธา ลปา สิงฺคี อุนฺนฬา อสมาหิตา, น เม เต ภิกฺขเว ภิกฺขู มามกา, อปคตา จ เต ภิกฺขเว ภิกฺขู อิมสฺมา ธมฺมวินยา, น จ เต ภิกฺขู อิมสฺมิํ ธมฺมวินเย วุทฺธิํ วิรูฬฺหิํ เวปุลฺลํ อาปชฺชนฺติ.\csromanspacer{} เย จ โข ภิกฺขเว ภิกฺขู นิกฺกุหา นิลฺลปา ธีรา อตฺถทฺธา สุสมาหิตา, เต โข เม ภิกฺขเว ภิกฺขู มามกา, น จ อปคตา เต ภิกฺขเว ภิกฺขู อิมสฺมา ธมฺมวินยา, เต จ ภิกฺขู อิมสฺมิํ ธมฺมวินเย วุทฺธิํ วิรูฬฺหิํ เวปุลฺลํ อาปชฺชนฺติ.}

\csromanheader{2}{14. ตุวฏกสุตฺตนิทฺเทส}
\csromangathameasure{กุหา ถทฺธา ลปา สิงฺคี, อุนฺนฬา อสมาหิตา. \\ น เต ธมฺเม วิรูหนฺติ, สมฺมาสมฺพุทฺธเทสิเต. \\ นิกฺกุหา นิลฺลปา ธีรา, อตฺถทฺธา สุสมาหิตา. \\ เต เว ธมฺเม วิรูหนฺติ, สมฺมาสมฺพุทฺธเทสิเต.}
\csromangathagroup{\csromangathastanza{\csromanfolio{299}กุหา ถทฺธา ลปา สิงฺคี, อุนฺนฬา อสมาหิตา. \\ น เต ธมฺเม วิรูหนฺติ, สมฺมาสมฺพุทฺธเทสิเต.}\csromangathastanza{นิกฺกุหา นิลฺลปา ธีรา, อตฺถทฺธา สุสมาหิตา. \\ เต เว ธมฺเม วิรูหนฺติ, สมฺมาสมฺพุทฺธเทสิเต.}}

\prose{\textbf{วิรุตญฺจ คพฺภกรณํ, ติกิจฺฉํ มามโก น เสเวยฺยา}ติ วิรุตญฺจ คพฺภกรณญฺจ ติกิจฺฉญฺจ มามโก น เสเวยฺย น นิเสเวยฺย น สํเสเวยฺย นปฺปฏิเสเวยฺย น จเรยฺย น สมาจเรยฺย น สมาทาย วตฺเตยฺย.\csromanspacer{} อถ วา น คเณฺหยฺย น อุคฺคเณฺหยฺย น ธาเรยฺย น อุปธาเรยฺย น อุปลกฺเขยฺย นปฺปโยเชยฺยาติ วิรุตญฺจ คพฺภกรณํ, ติกิจฺฉํ มามโก น เสเวยฺย.\csromanspacer{} เตนาห ภควา\csromandash{}}

\csromangathameasure{อาถพฺพณํ สุปินํ ลกฺขณํ, \\ โน วิทเห อโถปิ นกฺขตฺตํ.}
\csromangathagroup{\csromangathastanza{อาถพฺพณํ สุปินํ ลกฺขณํ, \\ โน วิทเห อโถปิ นกฺขตฺตํ.}}

\prose{วิรุตญฺจ คพฺภกรณํ,}

\prose{ติกิจฺฉํ มามโก น เสเวยฺยาติ.}

\proseitem{163}{\csromansymbolfootnote{*}{ขุ 1. 423 ปิฏฺเฐปิ.}นินฺทาย นปฺปเวเธยฺย,}

\prose{\textbf{น อุนฺนเมยฺย ปสํสิโต ภิกฺขุ.}}

\csromangathameasure{\textbf{โลภํ สห มจฺฉริเยน,} \\ \textbf{โกธํ เปสุณิยญฺจ ปนุเทยฺย.}}
\csromangathagroup{\csromangathastanza{\textbf{โลภํ สห มจฺฉริเยน,} \\ \textbf{โกธํ เปสุณิยญฺจ ปนุเทยฺย.}}}

\prose{\textbf{นินฺทาย นปฺปเวเธยฺยา}ติ อิเธกจฺเจ ภิกฺขู นินฺทนฺติ ชาติยา วา โคตฺเตน วา โกลปุตฺติเยน วา วณฺณโปกฺขรตาย วา ธเนน วา อชฺเฌเนน วา กมฺมายตเนน วา สิปฺปายตเนน วา วิชฺชาฏฺฐาเนน วา สุเตน วา ปฏิภาเนน วา อญฺญตรญฺญตเรน วา วตฺถุนา นินฺทนฺติ ครหนฺติ อุปวทนฺติ, นินฺทิโต ครหิโต อุปวทิโต, นินฺทาย ครหาย อุปวาเทน อกิตฺติยา อวณฺณหาริกาย น เวเธยฺย นปฺปเวเธยฺย น สมฺปเวเธยฺย น ตเสยฺย น อุตฺตเสยฺย น ปริตเสยฺย น ภาเยยฺย น สนฺตาสํ อาปชฺเชยฺย, อภีรู อสฺส อฉมฺภี อนุตฺราสี อปลายี, ปหีนภยเภรโว วิคตโลมหํโส วิหเรยฺยาติ นินฺทาย นปฺปเวเธยฺย.}

\prose{\textbf{น อุนฺนเมยฺย ปสํสิโต ภิกฺขู}ติ อิเธกจฺเจ ภิกฺขู ปสํสนฺติ ชาติยา วา โคตฺเตน วา โกลปุตฺติเยน วา วณฺณโปกฺขรตาย วา ธเนน วา อชฺเฌเนน วา กมฺมายตเนน วา สิปฺปายตเนน วา วิชฺชาฏฺฐาเนน \csromanfolio{300}วา สุเตน วา ปฏิภาเนน วา อญฺญตรญฺญตเรน วา วตฺถุนา ปสํสนฺติ โถเมนฺติ กิตฺเตนฺติ วณฺเณนฺติ, ปสํสิโต โถมิโต กิตฺติโต วณฺณิโต, ปสํสาย โถมเนน กิตฺติยา วณฺณหาริกาย อุนฺนติํ น กเรยฺย อุนฺนมํ น กเรยฺย มานํ น กเรยฺย ถมฺภํ น กเรยฺย น เตน มานํ ชเนยฺย น เตน ถทฺโธ อสฺส ปตฺถทฺโธ ปคฺคหิตสิโรติ น อุนฺนเมยฺย ปสํสิโต ภิกฺขุ.}

\prose{\textbf{โลภํ สห มจฺฉริเยน, โกธํ เปสุณิยญฺจ ปนุเทยฺยา}ติ \textbf{โลโภ}ติ โย \textbf{โลโภ} ลุพฺภนา ลุพฺภิตตฺตํ สาราโค สารชฺชนา สารชฺชิตตฺตํ อภิชฺฌา \textbf{โลโภ} อกุสลมูลํ.\csromanspacer{} \textbf{มจฺฉริยนฺ}ติ ปญฺจ มจฺฉริยานิ, อาวาสมจฺฉริยํ ฯเปฯ คาโห วุจฺจติ มจฺฉริยํ.\csromanspacer{} \textbf{โกโธ}ติ โย จิตฺตสฺส อาฆาโต ปฏิฆาโต ปฏิฆํ ปฏิวิโรโธ โกโป ปโกโป สมฺปโกโป โทโส ปโทโส สมฺปโทโส จิตฺตสฺส พฺยาปตฺติ มโนปโทโส โกโธ กุชฺฌนา กุชฺฌิตตฺตํ โทโส ทุสฺสนา ทุสฺสิตตฺตํ พฺยาปตฺติ พฺยาปชฺชนา พฺยาปชฺชิตตฺตํ, วิโรโธ ปฏิวิโรโธ จณฺฑิกฺกํ อสุโรโป อนตฺตมนตา จิตฺตสฺส.\csromanspacer{} \textbf{เปสุญฺญนฺ}ติ อิเธกจฺโจ ปิสุณวาโจ โหติ, อิโต สุตฺวา อมุตฺร อกฺขาตา อิเมสํ เภทาย, อมุตฺร วา สุตฺวา อิเมสํ อกฺขาตา อมูสํ เภทาย, อิติ สมคฺคานํ วา เภตฺตา ภินฺนานํ วา อนุปฺปทาตา, วคฺคาราโม วคฺครโต วคฺคนนฺที วคฺคกรณิํ วาจํ ภาสิตา โหติ.\csromanspacer{} อิทํ วุจฺจติ เปสุญฺญํ.}

\prose{อปิ จ ทฺวีหิ การเณหิ เปสุญฺญํ อุปสํหรติ ปิยกมฺยตาย วา เภทาธิปฺปาเยน วา.\csromanspacer{} กถํ ปิยกมฺยตาย เปสุญฺญํ อุปสํหรติ, อิมสฺส ปิโย ภวิสฺสามิ, มนาโป ภวิสฺสามิ, วิสฺสาสิโก ภวิสฺสามิ, อพฺภนฺตริโก ภวิสฺสามิ, สุหทโย ภวิสฺสามีติ.\csromanspacer{} เอวํ ปิยกมฺยตาย เปสุญฺญํ อุปสํหรติ.\csromanspacer{}\csromanspacer{}กถํ เภทาธิปฺปาเยน เปสุญฺญํ อุปสํหรติ, “กถํ อิเม นานา อสฺสุ, วินา อสฺสุ, วคฺคา อสฺสุ, ทฺวิธา อสฺสุ, เทฺวชฺฌา อสฺสุ, เทฺว ปกฺขา อสฺสุ, ภิชฺเฌยฺยุํ น สมาคจฺเฉยฺยุํ, ทุกฺขํ น ผาสุ วิหเรยฺยุนฺ”ติ.\csromanspacer{} เอวํ เภทาธิปฺปาเยน เปสุญฺญํ อุปสํหรติ.\csromanspacer{}\csromanspacer{}\textbf{โลภํ สห มจฺฉริเยน, โกธํ เปสุณิยญฺจ ปนุเทยฺยา}ติ โลภญฺจ มจฺฉริยญฺจ โกธญฺจ เปสุญฺญญฺจ นุเทยฺย ปนุเทยฺย ปชเหยฺย วิโนเทยฺย พฺยนฺติํ กเรยฺย อนภาวํ คเมยฺยาติ โลภํ สห มจฺฉริเยน, โกธํ เปสุณิยญฺจ ปนุเทยฺย.\csromanspacer{} เตนาห ภควา\csromandash{}}

\csromanheader{2}{14. ตุวฏกสุตฺตนิทฺเทส}
\csromangathameasure{นินฺทาย นปฺปเวเธยฺย, \\ น อุนฺนเมยฺย ปสํสิโต ภิกฺขุ. \\ โลภํ สห มจฺฉริเยน, \\ โกธํ เปสุณิยญฺจ ปนุเทยฺยาติ.}
\csromangathagroup{\csromangathastanza{\csromanfolio{301}นินฺทาย นปฺปเวเธยฺย, \\ น อุนฺนเมยฺย ปสํสิโต ภิกฺขุ. \\ โลภํ สห มจฺฉริเยน, \\ โกธํ เปสุณิยญฺจ ปนุเทยฺยาติ.}}

\proseitem{164}{\csromansymbolfootnote{*}{ขุ 1. 423 ปิฏฺเฐปิ.}กยวิกฺกเย น ติฏฺเฐยฺย,}

\prose{\textbf{อุปวาทํ ภิกฺขุ น กเรยฺย กุหิญฺจิ.}}

\csromangathameasure{\textbf{คาเม จ นาภิสชฺเชยฺย,} \\ \textbf{ลาภกมฺยา ชนํ น ลปเยยฺย.}}
\csromangathagroup{\csromangathastanza{\textbf{คาเม จ นาภิสชฺเชยฺย,} \\ \textbf{ลาภกมฺยา ชนํ น ลปเยยฺย.}}}

\prose{\textbf{กยวิกฺกเย น ติฏฺเฐยฺยา}ติ เย กยวิกฺกยา วินเย ปฏิกฺขิตฺตา, น เต อิมสฺมิํ อตฺเถ อธิปฺเปตา.\csromanspacer{} กถํ กยวิกฺกเย ติฏฺฐติ, ปญฺจนฺนํ สทฺธิํ ปตฺตํ วา จีวรํ วา อญฺญํ วา กิญฺจิ ปริกฺขารํ วญฺจนิยํ วา กโรนฺโต อุทยํ วา ปตฺถยนฺโต ปริวตฺเตติ.\csromanspacer{} เอวํ กยวิกฺกเย ติฏฺฐติ.\csromanspacer{}\csromanspacer{}กถํ กยวิกฺกเย น ติฏฺฐติ, ปญฺจนฺนํ สทฺธิํ ปตฺตํ วา จีวรํ วา อญฺญํ วา กิญฺจิ ปริกฺขารํ น วญฺจนิยํ วา กโรนฺโต น อุทยํ วา ปตฺถยนฺโต ปริวตฺเตติ.\csromanspacer{} เอวํ กยวิกฺกเยน ติฏฺฐติ.\csromanspacer{}\csromanspacer{}\textbf{กยวิกฺกเย น ติฏฺเฐยฺยา}ติ กยวิกฺกเย น ติฏฺเฐยฺย น สนฺติฏฺเฐยฺย, กยวิกฺกยํ ปชเหยฺย วิโนเทยฺย พฺยนฺติํ กเรยฺย อนภาวํ คเมยฺย, กยวิกฺกยา อารโต อสฺส วิรโต ปฏิวิรโต นิกฺขนฺโต นิสฺสโฏ วิปฺปมุตฺโต วิสญฺญุตฺโต วิมริยาทิกเตน เจตสา วิหเรยฺยาติ กยวิกฺกเย น ติฏฺเฐยฺย.}

\prose{\textbf{อุปวาทํ ภิกฺขุ น กเรยฺย กุหิญฺจี}ติ กตเม อุปวาทกราติ กิเลสา, สนฺเตเก สมณพฺราหฺมณา อิทฺธิมนฺโต ทิพฺพจกฺขุกา ปรจิตฺตวิทุโน, เต ทูรโตปิ ปสฺสนฺติ, อาสนฺนาปิ น ทิสฺสนฺติ, เจตสาปิ จิตฺตํ ปชานนฺติ, เทวตาปิ โข สนฺติ อิทฺธิมนฺตินิโย ทิพฺพจกฺขุกา ปรจิตฺตวิทุนิโย, ตา ทูรโตปิ ปสฺสนฺติ, อาสนฺนาปิ น ทิสฺสนฺติ, เจตสาปิ จิตฺตํ ปชานนฺติ.\csromanspacer{} เต โอฬาริเกหิ วา กิเลเสหิ มชฺฌิเมหิ วา กิเลเสหิ สุขุเมหิ วา กิเลเสหิ อุปวเทยฺยุํ.\csromanspacer{} กตเม โอฬาริกา กิเลสา กายทุจฺจริตํ วจีทุจฺจริตํ มโนทุจฺจริตํ, อิเม วุจฺจนฺติ โอฬาริกา กิเลสา.\csromanspacer{} กตเม มชฺฌิมา กิเลสา, กามวิตกฺโก พฺยาปาทวิตกฺโก วิหิํสาวิตกฺโก, อิเม วุจฺจนฺติ มชฺฌิมา กิเลสา.\csromanspacer{} กตเม สุขุมา กิเลสา, ญาติวิตกฺโก ชนปทวิตกฺโก อมรวิตกฺโก ปรานุทยตา ปฏิสญฺญุตฺโต \csromanfolio{302}วิตกฺโก ลาภสกฺการสิโลกปฏิสญฺญุตฺโต วิตกฺโก อนวญฺญตฺติปฏิสญฺญุตฺโต วิตกฺโก, อิเม วุจฺจนฺติ สุขุมา กิเลสา.\csromanspacer{} เต โอฬาริเกหิ วา กิเลเสหิ มชฺฌิเมหิ วา กิเลเสหิ สุขุเมหิ วา กิเลเสหิ น อุปวเทยฺย อุปวาทํ น กเรยฺย, อุปวาทกเร กิเลเส น กเรยฺย น ชเนยฺย น สญฺชเนยฺย น นิพฺพตฺเตยฺย นาภินิพฺพตฺเตยฺย, อุปวาทกเร กิเลเส ปชเหยฺย วิโนเทยฺย พฺยนฺติํ กเรยฺย อนภาวํ คเมยฺย, อุปวาทกเรหิ กิเลเสหิ อารโต อสฺส วิรโต ปฏิวิรโต นิกฺขนฺโต นิสฺสโฏ วิปฺปมุตฺโต วิสญฺญุตฺโต วิมริยาทิกเตน เจตสา วิหเรยฺย.\csromanspacer{} \textbf{กุหิญฺจี}ติ กุหิญฺจิ กิมฺหิจิ กตฺถจิ อชฺฌตฺตํ วา พหิทฺธา วา อชฺฌตฺตพหิทฺธา วาติ อุปวาทํ ภิกฺขุ น กเรยฺย กุหิญฺจิ.}

\prose{\textbf{คาเม จ นาภิสชฺเชยฺยา}ติ กถํ คาเม สชฺชติ, อิธ ภิกฺขุ คาเม คิหีหิ สํสฏฺโฐ วิหรติ สหนนฺที สหโสกี สุขิเตสุ สุขิโต ทุกฺขิเตสุ ทุกฺขิโต อุปฺปนฺเนสุ กิจฺจกรณีเยสุ อตฺตนา โวโยคํ อาปชฺชติ.\csromanspacer{} เอวมฺปิ คาเม สชฺชติ.}

\prose{อถ วา ภิกฺขุ ปุพฺพณฺหสมยํ นิวาเสตฺวา ปตฺตจีวรมาทาย คามํ วา นิคมํ วา ปิณฺฑาย ปวิสติ อรกฺขิเตเนว กาเยน อรกฺขิตาย วาจาย อรกฺขิเตน จิตฺเตน อนุปฏฺฐิตาย สติยา อสํวุเตหิ อินฺทฺริเยหิ.\csromanspacer{} โส ตตฺร ตตฺร สชฺชติ, ตตฺร ตตฺร คณฺหาติ, ตตฺร ตตฺร พชฺฌติ, ตตฺร ตตฺร อนยพฺยสนํ อาปชฺชติ.\csromanspacer{} เอวมฺปิ คาเม สชฺชติ.}

\prose{กถํ คาเม น สชฺชติ, อิธ ภิกฺขุ คาเม คิหีหิ อสํสฏฺโฐ วิหรติ น สหนนฺที น สหโสกี น สุขิเตสุ สุขิโต น ทุกฺขิเตสุ ทุกฺขิโต อุปฺปนฺเนสุ กิจฺจกรณีเยสุ น อตฺตนา โวโยคํ อาปชฺชติ.\csromanspacer{} เอวมฺปิ คาเม น สชฺชติ.}

\prose{อถ วา ภิกฺขุ ปุพฺพณฺหสมยํ นิวาเสตฺวา ปตฺตจีวรมาทาย คามํ วา นิคมํ วา ปิณฺฑาย ปวิสติ รกฺขิเตเนว กาเยน รกฺขิตาย วาจาย รกฺขิเตน จิตฺเตน อุปฏฺฐิตาย สติยา สํวุเตหิ อินฺทฺริเยหิ.\csromanspacer{} โส ตตฺร ตตฺร น สชฺชติ, ตตฺร ตตฺร น คณฺหาติ, ตตฺร ตตฺร น พชฺฌติ, ตตฺร ตตฺร น อนยพฺยสนํ อาปชฺชติ.\csromanspacer{} เอวมฺปิ คาเม น สชฺชติ.\csromanspacer{}\csromanspacer{}\textbf{คาเม จ นาภิสชฺเชยฺยา}ติ คาเม น สชฺเชยฺย น คเณฺหยฺย น พชฺเฌยฺย น ปลิพชฺเฌยฺย,}

\vspace{\tipitakaparskip}
\csromancenter{ตุวฏกสุตฺตนิทฺเทส อคิทฺโธ อสฺส อคธิโต อมุจฺฉิโต อนชฺโฌสนฺโน วีตเคโธ วิคตเคโธ จตฺตเคโธ ฯเปฯ พฺรหฺมภูเตน อตฺตนา วิหเรยฺยาติ คาเม จ นาภิสชฺเชยฺย.}
\prose{\csromanfolio{303}\textbf{ลาภกมฺยา ชนํ น ลปเยยฺยา}ติ กตมา ลปนา, ลาภสกฺการสิโลกสนฺนิสฺสิตสฺส ปาปิจฺฉสฺส อิจฺฉาปกตสฺส อามิสจกฺขุกสฺส โลกธมฺมครุกสฺส ยา ปเรสํ อาลปนา ลปนา สลฺลปนา อุลฺลปนา สมุลฺลปนา อุนฺนหนา สมุนฺนหนา อุกฺกาจนา สมุกฺกาจนา\footnote{อุกฺกาปนา สมุกฺกาปนา (สฺยา)} อนุปิยภาณิตา จาตุกมฺยตา มุคฺคสูปฺยตา ปาริภฏฺยตา ปรปิฏฺฐิมํสิกตา\footnote{ปิฏฺฐิมํสิกตา (ก) วิสุทฺธิมคฺเค สีลนิทฺเทสวณฺณนา โอโลเกตพฺพา.} ยา ตตฺถ สณฺหวาจตา สขิลวาจตา สิถิลวาจตา\footnote{เมตฺตวาจกตา (สฺยา)} อผรุสวาจตา.\csromanspacer{} อยํ วุจฺจติ ลปนา.}

\prose{อปิ จ ทฺวีหิ การเณหิ ชนํ ลปติ, อตฺตานํ วา นีจํ ฐเปนฺโต ปรํ อุจฺจํ ฐเปนฺโต ชนํ ลปติ, อตฺตานํ วา อุจฺจํ ฐเปนฺโต ปรํ นีจํ ฐเปนฺโต ชนํ ลปติ.\csromanspacer{} กถํ อตฺตานํ นีจํ ฐเปนฺโต ปรํ อุจฺจํ ฐเปนฺโต ชนํ ลปติ, ตุเมฺห เม พหูปการา, อหํ ตุเมฺห นิสฺสาย ลภามิ จีวรปิณฺฑปาตเสนาสนคิลานปจฺจยเภสชฺชปริกฺขารํ, เยปิ เม อญฺเญ ทาตุํ วา กาตุํ วา มญฺญนฺติ ตุเมฺห นิสฺสาย ตุเมฺห สมฺปสฺสนฺตา, ยมฺปิ เม ปุราณํ มาตาเปตฺติกํ นามเธยฺยํ, ตมฺปิ เม อนฺตรหิตํ.\csromanspacer{} ตุเมฺหหิ อหํ ญายามิ “อสุกสฺส กุลูปโก อสุกาย กุลูปโก”ติ.\csromanspacer{} เอวํ อตฺตานํ นีจํ ฐเปนฺโต ปรํ อุจฺจํ ฐเปนฺโต ชนํ ลปติ.}

\prose{กถํ อตฺตานํ อุจฺจํ ฐเปนฺโต ปรํ นีจํ ฐเปนฺโต ชนํ ลปติ, อหํ ตุมฺหากํ พหูปกาโร, ตุเมฺห มํ อาคมฺม พุทฺธํ สรณํ คเต, ธมฺมํ สรณํ คเต, สํฆํ สรณํ คตา, ปาณาติปาตา ปฏิวิรตา, อทินฺนาทานา ปฏิวิรตา, กาเมสุมิจฺฉาจารา ปฏิวิรตา, มุสาวาทา ปฏิวิรตา, สุราเมรยมชฺชปฺปมาทฏฺฐานา ปฏิวิรตา, อหํ ตุมฺหากํ อุทฺเทสํ เทมิ, ปริปุจฺฉํ เทมิ, อุโปสถํ อาจิกฺขามิ, นวกมฺมํ อธิฏฺฐามิ, อถ ปน ตุเมฺห มํ อุชฺฌิตฺวา อญฺเญ สกฺกโรถ ครุํ กโรถ มาเนถ ปูเชถาติ.\csromanspacer{} เอวํ อตฺตานํ อุจฺจํ ฐเปนฺโต ปรํ นีจํ ฐเปนฺโต ชนํ ลปติ.}

\prose{\csromanfolio{304}\textbf{ลาภกมฺยา ชนํ น ลปเยยฺยา}ติ ลาภเหตุ ลาภปจฺจยา ลาภการณา ลาภาภินิพฺพตฺติยา ลาภํ ปริปาเจนฺโต ชนํ น ลปเยยฺย ลปนํ ปชเหยฺย วิโนเทยฺย พฺยนฺติํ กเรยฺย อนภาวํ คเมยฺย, ลปนา อารโต อสฺส วิรโต ปฏิวิรโต นิกฺขนฺโต นิสฺสโฏ วิปฺปมุตฺโต วิสญฺญุตฺโต วิมริยาทิกเตน เจตสา วิหเรยฺยาติ ลาภกมฺยา ชนํ น ปลปเยยฺย.\csromanspacer{} เตนาห ภควา\csromandash{}}

\csromangathameasure{กยวิกฺกเย น ติฏฺเฐยฺย, \\ อุปวาทํ ภิกฺขุ น กเรยฺย กุหิญฺจิ. \\ คาเม จ นาภิสชฺเชยฺย,}
\csromangathagroup{\csromangathastanza{กยวิกฺกเย น ติฏฺเฐยฺย, \\ อุปวาทํ ภิกฺขุ น กเรยฺย กุหิญฺจิ. \\ คาเม จ นาภิสชฺเชยฺย,}}

\prose{\textbf{ลาภกมฺยา ชนํ น ลปเยยฺยา}ติ.}

\csromanhangingbegin
\hangingheaditem{165}{น จ กตฺถิโก\footnote{กตฺถิตา (สฺยา, ก) ขุ 1. 423 ปิฏฺเฐปิ.} \textbf{สิยา ภิกฺขุ,}}
\hangingbody{\textbf{น จ วาจํ ปยุตฺตํ ภาเสยฺย.}}
\hangingbody{\textbf{ปาคพฺภิยํ น สิกฺเขยฺย,}}
\hangingbody{\textbf{กถํ วิคฺคาหิกํ น กถเยยฺย.}}
\csromanhangingend

\prose{\textbf{น จ กตฺถิโก สิยา ภิกฺขู}ติ * อิเธกจฺโจ กตฺถี โหติ วิกตฺถี.\csromanspacer{} โส กตฺถติ วิกตฺถติ “อหมสฺมิ สีลสมฺปนฺโน”ติ วา “วตสมฺปนฺโน”ติ วา “สีลพฺพตสมฺปนฺโน”ติ วา ชาติยา วา โคตฺเตน วา โกลปุตฺติเยน วา วณฺณโปกฺขรตาย วา ธเนน วา อชฺเฌเนน วา กมฺมายตเนน วา สิปฺปายตเนน วา วิชฺชาฏฺฐาเนน วา สุเตน วา ปฏิภาเนน วา อญฺญตรญฺญตเรน วา วตฺถุนา “อุจฺจา กุลา ปพฺพชิโต”ติ วา “มหาโภคกุลา ปพฺพชิโต”ติ วา “อุฬารโภคกุลา ปพฺพชิโต”ติ วา ฯเปฯ “สุตฺตนฺติโก”ติ วา “วินยธโร”ติ วา “ธมฺมกถิโก”ติ วา “อารญฺญิโก”ติ วา ฯเปฯ “เนวสญฺญานาสญฺญายตนสมาปตฺติยา ลาภี”ติ วา กตฺถติ วิกตฺถติ.\csromanspacer{} เอวํ น กตฺเถยฺย น วิกตฺเถยฺย, กตฺถํ ปชเหยฺย วิโนเทยฺย พฺยนฺติํ กเรยฺย อนภาวํ คเมยฺย, กตฺถนา อารโต อสฺส วิรโต ปฏิวิรโต นิกฺขนฺโต นิสฺสโฏ วิปฺปมุตฺโต วิสญฺญุตฺโต วิมริยาทิกเตน เจตสา วิหเรยฺยาติ น จ กตฺถิโก สิยา ภิกฺขุ.}

\csromanheader{2}{14. ตุวฏกสุตฺตนิทฺเทส}
\prose{\csromanfolio{305}\textbf{น จ วาจํ ปยุตฺตํ ภาเสยฺยา}ติ กตมา ปยุตฺตวาจา, อิเธกจฺโจ จีวรปยุตฺตํ วาจํ ภาสติ, ปิณฺฑปาตปยุตฺตํ วาจํ ภาสติ, เสนาสนปยุตฺตํ วาจํ ภาสติ, คิลานปจฺจยเภสชฺชปริกฺขารปยุตฺตํ วาจํ ภาสติ, อยมฺปิ วุจฺจติ ปยุตฺตวาจา.}

\prose{อถ วา จีวรเหตุ ปิณฺฑปาตเหตุ เสนาสนเหตุ คิลานปจฺจยเภสชฺชปริกฺขารเหตุ สจฺจมฺปิ ภณติ, มุสาปิ ภณติ.\csromanspacer{} ปิสุณมฺปิ ภณติ, อปิสุณมฺปิ ภณติ.\csromanspacer{} ผรุสมฺปิ ภณติ, อผรุสมฺปิ ภณติ.\csromanspacer{} สมฺผปฺปลาปมฺปิ ภณติ, อสมฺผปฺปลาปมฺปิ ภณติ.\csromanspacer{} มนฺตาปิ วาจํ ภาสติ, อยมฺปิ วุจฺจติ ปยุตฺตวาจา.\csromanspacer{}\csromanspacer{}อถ วา ปสนฺนจิตฺโต ปเรสํ ธมฺมํ เทเสติ “อโห วต เม ธมฺมํ สุเณยฺยุํ, สุตฺวาว ธมฺมํ ปสีเทยฺยุํ, ปสนฺนา จ เม ปสนฺนาการํ กเรยฺยุนฺ”ติ.\csromanspacer{} อยํ วุจฺจติ ปยุตฺตวาจา.\csromanspacer{} \textbf{น จ วาจํ ปยุตฺตํ ภาเสยฺยา}ติ อนฺตมโส ธมฺมเทสนํ วาจํ อุปาทาย ปยุตฺตวาจํ น ภาเสยฺย น กเถยฺย น ภเณยฺย น ทีเปยฺย น โวหเรยฺย, ปยุตฺตํ วาจํ ปชเหยฺย วิโนเทยฺย พฺยนฺติํ กเรยฺย อนภาวํ คเมยฺย, ปยุตฺตวาจาย อารโต อสฺส วิรโต ปฏิวิรโต นิกฺขนฺโต นิสฺสโฏ วิปฺปมุตฺโต วิสญฺญุตฺโต วิมริยาทิกเตน เจตสา วิหเรยฺยาติ น จ วาจํ ปยุตฺตํ ภาเสยฺย.}

\prose{\textbf{ปาคพฺภิยํ น สิกฺเขยฺยา}ติ \textbf{ปาคพฺภิยนฺ}ติ ตีณิ ปาคพฺภิยานิ กายิกํ ปาคพฺภิยํ วาจสิกํ ปาคพฺภิยํ เจตสิกํ ปาคพฺภิยํ.\csromanspacer{} กตมํ กายิกํ ปาคพฺภิยํ, อิเธกจฺโจ สํฆคโตปิ กายิกํ ปาคพฺภิยํ ทสฺเสติ, คณคโตปิ กายิกํ ปาคพฺภิยํ ทสฺเสติ, โภชนสาลายมฺปิ กายิกํ ปาคพฺภิยํ ทสฺเสติ, ชนฺตาฆเรปิ กายิกํ ปาคพฺภิยํ ทสฺเสติ, อุทกติตฺเถปิ กายิกํ ปาคพฺภิยํ ทสฺเสติ, อนฺตรฆรํ ปวิสนฺโตปิ กายิกํ ปาคพฺภิยํ ทสฺเสติ, อนฺตรฆรํ ปวิฏฺโฐปิ กายิกํ ปาคพฺภิยํ ทสฺเสติ.}

\prose{กถํ สํฆคโต กายิกํ ปาคพฺภิยํ ทสฺเสติ, อิเธกจฺโจ สํฆคโต อจิตฺตีการกโต เถเร ภิกฺขู ฆฏฺฏยนฺโตปิ ติฏฺฐติ, ฆฏฺฏยนฺโตปิ นิสีทติ, ปุรโตปิ ติฏฺฐติ, ปุรโตปิ นิสีทติ, อุจฺเจปิ อาสเน นิสีทติ, สสีสํ ปารุปิตฺวา นิสีทติ, ฐิตโกปิ ภณติ, พาหาวิกฺเขปโกปิ ภณติ.\csromanspacer{} เอวํ สํฆคโต กายิกํ ปาคพฺภิยํ ทสฺเสติ.}

\prose{\csromanfolio{306}กถํ คณคโต กายิกํ ปาคพฺภิยํ ทสฺเสติ, อิเธกจฺโจ คณคโต อจิตฺตีการกโต เถรานํ ภิกฺขูนํ อนุปาหนานํ จงฺกมนฺตานํ ส- อุปาหโน จงฺกมติ, นีเจ จงฺกเม จงฺกมนฺตานํ อุจฺเจ จงฺกเม จงฺกมติ, ฉมาย จงฺกมนฺตานํ จงฺกเม จงฺกมติ, ฆฏฺฏยนฺโตปิ ติฏฺฐติ, ฆฏฺฏยนฺโตปิ นิสีทติ, ปุรโตปิ ติฏฺฐติ, ปุรโตปิ นิสีทติ, อุจฺเจปิ อาสเน นิสีทติ, สสีสํ ปารุปิตฺวาปิ นิสีทติ, ฐิโตปิ ภณติ, พาหาวิกฺเขปโกปิ ภณติ.\csromanspacer{} เอวํ คณคโต กายิกํ ปาคพฺภิยํ ทสฺเสติ.}

\prose{กถํ โภชนสาลายํ กายิกํ ปาคพฺภิยํ ทสฺเสติ, อิเธกจฺโจ โภชนสาลายํ อจิตฺตีการกโต เถเร ภิกฺขู อนุปขชฺช นิสีทติ, นเวปิ ภิกฺขู อาสเนน ปฏิพาหติ, ฆฏฺฏยนฺโตปิ ติฏฺฐติ, ฆฏฺฏยนฺโตปิ นิสีทติ, ปุรโตปิ ติฏฺฐติ, ปุรโตปิ นิสีทติ, อุจฺเจปิ อาสเน นิสีทติ, สสีสํ ปารุปิตฺวา นิสีทติ, ฐิตโกปิ ภณติ, พาหาวิกฺเขปโกปิ ภณติ.\csromanspacer{} เอวํ โภชนสาลายํ กายิกํ ปาคพฺภิยํ ทสฺเสติ.}

\prose{กถํ ชนฺตาฆเร กายิกํ ปาคพฺภิยํ ทสฺเสติ, อิเธกจฺโจ ชนฺตาฆเร อจิตฺตีการกโต เถเร ภิกฺขู ฆฏฺฏยนฺโตปิ ติฏฺฐติ, ฆฏฺฏยนฺโตปิ นิสีทติ, ปุรโตปิ ติฏฺฐติ, ปุรโตปิ นิสีทติ, อุจฺเจปิ อาสเน นิสีทติ, อนาปุจฺฉาปิ กฏฺฐํ ปกฺขิปติ, อนาปุจฺฉาปิ ทฺวารํ ปิทหติ.\csromanspacer{} เอวํ ชนฺตาฆเร กายิกํ ปาคพฺภิยํ ทสฺเสติ.}

\prose{กถํ อุทกติตฺเถ กายิกํ ปาคพฺภิยํ ทสฺเสติ, อิเธกจฺโจ อุทกติตฺเถ อจิตฺตีการกโต เถเร ภิกฺขู ฆฏฺฏยนฺโตปิ โอตรติ, ปุรโตปิ โอตรติ, ฆฏฺฏยนฺโตปิ นฺหายติ, ปุรโตปิ นฺหายติ, อุปริโตปิ นฺหายติ, ฆฏฺฏยนฺโตปิ อุตฺตรติ, ปุรโตปิ อุตฺตรติ, อุปริโตปิ อุตฺตรติ.\csromanspacer{} เอวํ อุทกติตฺเถ กายิกํ ปาคพฺภิยํ ทสฺเสติ.}

\prose{กถํ อนฺตรฆรํ ปวิสนฺโต กายิกํ ปาคพฺภิยํ ทสฺเสติ, อิเธกจฺโจ อนฺตรฆรํ ปวิสนฺโต อจิตฺตีการกโต เถเร ภิกฺขู ฆฏฺฏยนฺโตปิ คจฺฉติ, ปุรโตปิ คจฺฉติ, โวกฺกมฺมปิ เถรานํ ภิกฺขูนํ ปุรโต คจฺฉติ.\csromanspacer{} เอวํ อนฺตรฆรํ ปวิสนฺโต กายิกํ ปาคพฺภิยํ ทสฺเสติ.}

\prose{กถํ อนฺตรฆรํ ปวิฏฺโฐ กายิกํ ปาคพฺภิยํ ทสฺเสติ, อิเธกจฺโจ อนฺตรฆรํ ปวิฏฺโฐ “น ปวิส ภนฺเต”ติ วุจฺจมาโน ปวิสติ, “น ติฏฺฐ ภนฺเต”ติ}

\vspace{\tipitakaparskip}
\csromancenter{ตุวฏกสุตฺตนิทฺเทส วุจฺจมาโน ติฏฺฐติ, “น นิสีท ภนฺเต”ติ วุจฺจมาโน นิสีทติ, อโนกาสมฺปิ ปวิสติ, อโนกาเสปิ ติฏฺฐติ, อโนกาเสปิ นิสีทติ, ยานิปิ ตานิ โหนฺติ กุลานํ โอวรกานิ คูฬฺหานิ จ ปฏิจฺฉนฺนานิ จ, ยตฺถ กุลิตฺถิโย กุลธีตโร กุลสุณฺหาโย กุลกุมาริโย นิสีทนฺติ, ตตฺถปิ สหสา ปวิสติ, กุมารกสฺสปิ สีสํ ปรามสติ.\csromanspacer{} เอวํ อนฺตรฆรํ ปวิฏฺโฐ กายิกํ ปาคพฺภิยํ ทสฺเสติ.\csromanspacer{} อิทํ กายิกํ ปาคพฺภิยํ.}
\prose{\csromanfolio{307}กตมํ วาจสิกํ ปาคพฺภิยํ, อิเธกจฺโจ สํฆคโตปิ วาจสิกํ ปาคพฺภิยํ ทสฺเสติ, คณคโตปิ วาจสิกํ ปาคพฺภิยํ ทสฺเสติ, อนฺตรฆรํ ปวิฏฺโฐปิ วาจสิกํ ปาคพฺภิยํ ทสฺเสติ.\csromanspacer{} กถํ สํฆคโต วาจสิกํ ปาคพฺภิยํ ทสฺเสติ, อิเธกจฺโจ สํฆคโต อจิตฺตีการกโต เถเร ภิกฺขู อนาปุจฺฉํ วา อนชฺฌิฏฺโฐ วา อารามคตานํ ภิกฺขูนํ ธมฺมํ ภณติ, ปญฺหํ วิสชฺเชติ, ปาติโมกฺขํ อุทฺทสติ, ฐิตโกปิ ภณติ, พาหาวิกฺเขปโกปิ ภณติ.\csromanspacer{} เอวํ สํฆคโต วาจสิกํ ปาคพฺภิยํ ทสฺเสติ.}

\prose{กถํ คณคโต วาจสิกํ ปาคพฺภิยํ ทสฺเสติ, อิเธกจฺโจ คณคโต อจิตฺตีการกโต เถเร ภิกฺขู อนาปุจฺฉํ วา อนชฺฌิฏฺโฐ วา อารามคตานํ ภิกฺขูนํ ธมฺมํ ภณติ, ปญฺหํ วิสชฺเชติ, ฐิตโกปิ ภณติ, พาหาวิกฺเขปโกปิ ภณติ, อารามคตานํ ภิกฺขุนีนํ อุปาสกานํ อุปาสิกานํ ธมฺมํ ภณติ, ปญฺหํ วิสชฺเชติ, ฐิตโกปิ ภณติ, พาหาวิกฺเขปโกปิ ภณติ.\csromanspacer{} เอวํ คณคโต วาจสิกํ ปาคพฺภิยํ ทสฺเสติ.}

\prose{กถํ อนฺตรฆรํ ปวิฏฺโฐ วาจสิกํ ปาคพฺภิยํ ทสฺเสติ, อิเธกจฺโจ อนฺตรฆรํ ปวิฏฺโฐ อิตฺถิํ วา กุมาริํ วา เอวมาห “อิตฺถนฺนาเม อิตฺถํโคตฺเต กิํ อตฺถิ, ยาคุ อตฺถิ, ภตฺตํ อตฺถิ, ขาทนียํ อตฺถิ, กิํ ปิวิสฺสาม, กิํ ภุญฺชิสฺสาม, กิํ ขาทิสฺสาม, กิํ วา อตฺถิ, กิํ วา เม ทสฺสถา”ติ วิปฺปลปติ.\csromanspacer{} เอวํ อนฺตรฆรํ ปวิฏฺโฐ วาจสิกํ ปาคพฺภิยํ ทสฺเสติ.\csromanspacer{} อิทํ วาจสิกํ ปาคพฺภิยํ.}

\prose{กตมํ เจตสิกํ ปาคพฺภิยํ, อิเธกจฺโจ น อุจฺจา กุลา ปพฺพชิโต สมาโน อุจฺจา กุลา ปพฺพชิเตน สทฺธิํ สทิสํ อตฺตานํ กโรติ จิตฺเตน, น มหาโภคกุลา ปพฺพชิโต สมาโน มหาโภคกุลา ปพฺพชิเตน สทฺธิํ สทิสํ อตฺตานํ กโรติ จิตฺเตน, น อุฬารโภคกุลา \csromanfolio{308}ปพฺพชิโต สมาโน อุฬารโภคกุลา ปพฺพชิเตน สทฺธิํ สทิสํ อตฺตานํ กโรติ จิตฺเตน, น สุตฺตนฺติโก สมาโน สุตฺตนฺติเกน สทฺธิํ สทิสํ อตฺตานํ กโรติ จิตฺเตน, น วินยธโร สมาโน.\csromanspacer{} น ธมฺมกถิโก สมาโน.\csromanspacer{} น อารญฺญิโก สมาโน.\csromanspacer{} น ปิณฺฑปาติโก สมาโน.\csromanspacer{} น ปํสุกูลิโก สมาโน.\csromanspacer{} น เตจีวริโก สมาโน.\csromanspacer{} น สปทานจาริโก สมาโน.\csromanspacer{} น ขลุปจฺฉาภตฺติโก สมาโน.\csromanspacer{} น เนสชฺชิโก สมาโน.\csromanspacer{} น ยถาสนฺถติโก สมาโน ยถาสนฺถติเกน สทฺธิํ สทิสํ อตฺตานํ กโรติ จิตฺเตน, น ปฐมสฺส ฌานสฺส ลาภี สมาโน ปฐมสฺส ฌานสฺส ลาภินา สทฺธิํ สทิสํ อตฺตานํ กโรติ จิตฺเตน, น ทุติยสฺส ฌานสฺส.\csromanspacer{} น ตติยสฺส ฌานสฺส.\csromanspacer{} น จตุตฺถสฺส ฌานสฺส ลาภี สมาโน.\csromanspacer{} น อากาสานญฺจายตนสมาปตฺติยา ลาภี สมาโน.\csromanspacer{} น วิญฺญาณญฺจายตนสมาปตฺติยา.\csromanspacer{} น อากิญฺจญฺญายตนสมาปตฺติยา.\csromanspacer{} น เนวสญฺญานาสญฺญายตนสมาปตฺติยา ลาภี สมาโน เนวสญฺญานาสญฺญายตนสมาปตฺติยา ลาภินา สทฺธิํ สทิสํ อตฺตานํ กโรติ จิตฺเตน.\csromanspacer{} อิทํ เจตสิกํ ปาคพฺภิยํ.\csromanspacer{}\csromanspacer{}\textbf{น สิกฺเขยฺยา}ติ ปาคพฺภิยํ น สิกฺเขยฺย น จเรยฺย น อาจเรยฺย น สมาทาย วตฺเตยฺย, ปาคพฺภิยํ ปชเหยฺย วิโนเทยฺย พฺยนฺติํ กเรยฺย อนภาวํ คเมยฺย, ปาคพฺภิยา อารโต อสฺส วิรโต ปฏิวิรโต นิกฺขนฺโต นิสฺสโฏ วิปฺปมุตฺโต วิสญฺญุตฺโต วิมริยาทิกเตน เจตสา วิหเรยฺยาติ ปาคพฺภิยํ น สิกฺเขยฺย.}

\prose{\textbf{กถํ วิคฺคาหิกํ น กถเยยฺยา}ติ กตมา วิคฺคาหิกา กถา, อิเธกจฺโจ เอวรูปิํ กถํ กตฺตา โหติ “น ตฺวํ อิมํ ธมฺมวินยํ อาชานาสิ ฯเปฯ นิพฺเพเฐหิ วา สเจ ปโหสี”ติ.\csromanspacer{}\csromanspacer{}วุตฺตํ เหตํ ภควตา * \csromandash{} วิคฺคาหิกาย โข โมคฺคลฺลาน กถาย สติ กถาพาหุลฺลํ ปาฏิกงฺขํ, กถาพาหุลฺเล สติ อุทฺธจฺจํ, อุทฺธตสฺส อสํวโร, อสํวุตสฺส อารา จิตฺตํ สมาธิมฺหาติ.\csromanspacer{} \textbf{กถํ วิคฺคาหิกํ น กถเยยฺยา}ติ วิคฺคาหิกํ กถํ น กเถยฺย น ภเณยฺย น ทีเปยฺย น โวหเรยฺย, วิคฺคาหิกํ กถํ ปชเหยฺย วิโนเทยฺย พฺยนฺติํ กเรยฺย อนภาวํ คเมยฺย, วิคฺคาหิกกถาย อารโต อสฺส วิรโต ปฏิวิรโต นิกฺขนฺโต นิสฺสโฏ วิปฺปมุตฺโต วิสญฺญุตฺโต วิมริยาทิกเตน เจตสา วิหเรยฺยาติ กถํ วิคฺคาหิกํ น กถเยยฺย.\csromanspacer{} เตนาห ภควา\csromandash{}}

\csromanheader{2}{14. ตุวฏกสุตฺตนิทฺเทส}
\csromangathameasure{น จ กตฺถิโก สิยา ภิกฺขุ, \\ น จ วาจํ ปยุตฺตํ ภาเสยฺย. \\ ปาคพฺภิยํ น สิกฺเขยฺย, \\ กถํ วิคฺคาหิกํ น กถเยยฺยาติ.}
\csromangathagroup{\csromangathastanza{\csromanfolio{309}น จ กตฺถิโก สิยา ภิกฺขุ, \\ น จ วาจํ ปยุตฺตํ ภาเสยฺย. \\ ปาคพฺภิยํ น สิกฺเขยฺย, \\ กถํ วิคฺคาหิกํ น กถเยยฺยาติ.}}

\proseitem{166}{\csromansymbolfootnote{*}{ขุ 1. 423 ปิฏฺเฐปิ.}โมสวชฺเช น นิยฺเยถ, สมฺปชาโน สฐานิ น กยิรา.}

\csromangathameasure{อถ ชีวิเตน ปญฺญาย, สีลพฺพเตน นาญฺญมติมญฺเญ.}
\csromangathagroup{\csromangathastanza{อถ ชีวิเตน ปญฺญาย, สีลพฺพเตน นาญฺญมติมญฺเญ.}}

\prose{\textbf{โมสวชฺเช น นิยฺเยถา}ติ โมสวชฺชํ วุจฺจติ มุสาวาโท.\csromanspacer{} อิเธกจฺโจ สภคฺคโต\footnote{สภาคคฺคโต (สี, ก)} วา ปริสคฺคโต\footnote{ปริสคโต (สี, ก) 118 ปิฏฺเฐ นตฺถิ ปาฐนานาตฺตํ.} วา ฯเปฯ อามิสกิญฺจิกฺขเหตุ วา สมฺปชานมุสา ภาสิตา โหติ, อิทํ วุจฺจติ โมสวชฺชํ.\csromanspacer{}\csromanspacer{}อปิ จ ตีหากาเรหิ มุสาวาโท โหติ, ปุพฺเพวสฺส โหติ “มุสา ภณิสฺสนฺ”ติ, ภณนฺตสฺส โหติ “มุสา ภณามี”ติ, ภณิตสฺส โหติ “มุสา มยา ภณิตนฺ”ติ, อิเมหิ ตีหากาเรหิ มุสาวาโท โหติ.\csromanspacer{}\csromanspacer{}อปิ จ จตูหากาเรหิ ปญฺจหากาเรหิ.\csromanspacer{} ฉหากาเรหิ.\csromanspacer{} สตฺตหากาเรหิ.\csromanspacer{} อฏฺฐหากาเรหิ มุสาวาโท โหติ, ปุพฺเพวสฺส โหติ “มุสา ภณิสฺสนฺ”ติ, ภณนฺตสฺส โหติ “มุสา ภณามี”ติ, ภณิตสฺส โหติ “มุสา มยา ภณิตนฺ”ติ, วินิธาย ทิฏฺฐิํ วินิธาย ขนฺติํ วินิธาย รุจิํ วินิธาย สญฺญํ, วินิธาย ภาวํ, อิเมหิ อฏฺฐหากาเรหิ มุสาวาโท โหติ.\csromanspacer{}\csromanspacer{}\textbf{โมสวชฺเช น นิยฺเยถา}ติ โมสวชฺเช น ยาเยยฺย น นิยฺยาเยยฺย น วเหยฺย\footnote{น วุเยฺหยฺย (สี, ก), นตฺถิ สฺยา-โปตฺถเก.} น สํหเรยฺย, โมสวชฺชํ ปชเหยฺย วิโนเทยฺย พฺยนฺติํ กเรยฺย อนภาวํ คเมยฺย, โมสวชฺชา อารโต อสฺส วิรโต ปฏิวิรโต นิกฺขนฺโต นิสฺสโฏ วิปฺปมุตฺโต วิสญฺญุตฺโต วิมริยาทิกเตน เจตสา วิหเรยฺยาติ โมสวชฺเช น นิยฺเยถ.}

\prose{\textbf{สมฺปชาโน สฐานิ น กยิรา}ติ กตมํ สาเฐยฺยํ, อิเธกจฺโจ สโฐ โหติ ปริสโฐ, ยํ ตตฺถ สฐํ สฐตา สาเฐยฺยํ, กกฺกรตา กกฺกริยํ, ปริกฺขตฺตตา ปาริกฺขตฺติยํ, อิทํ วุจฺจติ สาเฐยฺยํ.\csromanspacer{}\csromanspacer{}\textbf{สมฺปชาโน สฐานิ น กยิรา}ติ สมฺปชาโน หุตฺวา สาเฐยฺยํ น กเรยฺย น ชเนยฺย น สญฺชเนยฺย น นิพฺพตฺเตยฺย นาภินิพฺพตฺเตยฺย, สาเฐยฺยํ ปชเหยฺย วิโนเทยฺย พฺยนฺติํ กเรยฺย อนภาวํ คเมยฺย, สาเฐยฺยา อารโต \csromanfolio{310}อสฺส วิรโต ปฏิวิรโต นิกฺขนฺโต นิสฺสโฏ วิปฺปมุตฺโต วิสญฺญุตฺโต วิมริยาทิกเตน เจตสา วิหเรยฺยาติ สมฺปชาโน สฐานิ น กยิรา\textbf{.}}

\prose{\textbf{อถ ชีวิเตน ปญฺญาย, สีลพฺพเตน นาญฺญมติมญฺเญ}ติ \textbf{อถา}ติ ปทสนฺธิ ฯเปฯ ปทานุปุพฺพตาเปตํ \textbf{อถา}ติ.\csromanspacer{} อิเธกจฺโจ ลูขชีวิตํ ชีวนฺโต ปรํ ปณีตชีวิตํ ชีวนฺตํ อติมญฺญติ “กิํ ปนายํ พหุลาชีโว สพฺพํ สํภกฺเขติ.\csromanspacer{} เสยฺยถิทํ, มูลพีชํ ขนฺธพีชํ ผฬุพีชํ อคฺคพีชํ พีชพีชเมว ปญฺจมํ, อสนิวิจกฺกทนฺตกูฏสมณปฺปธาเนนา”ติ\footnote{อสนีว จกฺกํ ทนฺตกูฏํ สมณปฺปธาเนนาติ (สี)}.\csromanspacer{} โส ตาย ลูขชีวิตาย ปรํ ปณีตชีวิตํ ชีวนฺตํ อติมญฺญติ.}

\prose{อิเธกจฺโจ ปณีตชีวิตํ ชีวนฺโต ปรํ ลูขชีวิตํ ชีวนฺตํ อติมญฺญติ “กิํ ปนายํ อปฺปปุญฺโญ อปฺเปสกฺโข, น ลาภี จีวรปิณฺฑปาตเสนาสนคิลานปจฺจยเภสชฺชปริกฺขารานนฺ”ติ, โส ตาย ปณีตชีวิตาย ปรํ ลูขชีวิตํ ชีวนฺตํ อติมญฺญติ.\csromanspacer{}\csromanspacer{}อิเธกจฺโจ ปญฺญาสมฺปนฺโน โหติ, โส ปุฏฺโฐ ปญฺหํ วิสชฺเชติ.\csromanspacer{} ตสฺส เอวํ โหติ “อหมสฺมิ ปญฺญาสมฺปนฺโน, อิเม ปนญฺเญ น ปญฺญาสมฺปนฺนา”ติ, โส ตาย ปญฺญาสมฺปทาย ปรํ อติมญฺญติ.\csromanspacer{} อิเธกจฺโจ สีลสมฺปนฺโน โหติ ปาติโมกฺขสํวรสํวุโต วิหรติ อาจารโคจรสมฺปนฺโน อณุมตฺเตสุ วชฺเชสุ ภยทสฺสาวี, สมาทาย สิกฺขติ สิกฺขาปเทสุ, ตสฺส เอวํ โหติ “อหมสฺมิ สีลสมฺปนฺโน, อิเม ปนญฺเญ ภิกฺขู ทุสฺสีลา ปาปธมฺมา”ติ.\csromanspacer{} โส ตาย สีลสมฺปทาย ปรํ อติมญฺญติ.\csromanspacer{}\csromanspacer{}อิเธกจฺโจ วตสมฺปนฺโน โหติ อารญฺญิโก วา ปิณฺฑปาติโก วา ปํสุกูลิโก วา เตจีวริโก วา สปทานจาริโก วา ขลุปจฺฉาภตฺติโก วา เนสชฺชิโก วา ยฺ\textbf{อถา}สนฺถติโก วา, ตสฺส เอวํ โหติ “อหมสฺมิ วต สมฺปนฺโน, อิเม ปนญฺเญ น วตสมฺปนฺนา”ติ, โส ตาย วตสมฺปทาย ปรํ อติมญฺญติ.\csromanspacer{} \textbf{อถ ชีวิเตน ปญฺญาย, สีลพฺพเตน นาญฺญมติมญฺเญ}ติ ลูขชีวิตาย วา ปณีตชีวิตาย วา ปญฺญาสมฺปทาย วา สีลสมฺปทาย วา วตสมฺปทาย วา ปรํ นาติมญฺเญยฺย นาวชาเนยฺย, น เตน มานํ ชเนยฺย, น เตน ถทฺโธ อสฺส ปตฺถทฺโธ ปคฺคหิตสิโรติ อถ ชีวิเตน ปญฺญาย, สีลพฺพเตน นาญฺญมติมญฺเญ.\csromanspacer{} เตนาห ภควา\csromandash{}}

\csromangathameasure{โมสวชฺเช น นิยฺเยถ, สมฺปชาโน สฐานิ น กยิรา.}
\csromangathagroup{\csromangathastanza{โมสวชฺเช น นิยฺเยถ, สมฺปชาโน สฐานิ น กยิรา.}}

\prose{\textbf{อถ ชีวิเตน ปญฺญาย, สีลพฺพเตน นาญฺญมติมญฺเญ}ติ.}

\csromanheader{2}{14. ตุวฏกสุตฺตนิทฺเทส}
\csromanhangingbegin
\hangingheaditem{167}{\csromanfolio{311}สุตฺวา \textbf{รุสิโต}\footnote{ทูสิโต (สี, สฺยา) ขุ 1. 423 ปิฏฺเฐปิ.} \textbf{พหุํ วาจํ,}}
\hangingbody{\textbf{สมณานํ วา ปุถุชนานํ}2\textbf{.}}
\hangingbody{\textbf{ผรุเสน เน น ปฏิวชฺชา,}}
\hangingbody{\textbf{น หิ สนฺโต ปฏิเสนิํ กโรนฺติ}3\textbf{.}}
\csromanhangingend

\prose{\textbf{สุตฺวา รุสิโต พหุํ วาจํ, สมณานํ ปุถุชนานนฺ}ติ \textbf{รุสิโต}ติ ทูสิโต ขุํสิโต ฆฏฺฏิโต วมฺภิโต ครหิโต อุปวทิโต\textbf{.}\csromanspacer{} \textbf{สมณานนฺ}ติ เย เกจิ อิโต พหิทฺธา ปริพฺพชูปคตา ปริพฺพชสมาปนฺนา\textbf{.}\csromanspacer{} \textbf{ปุถุชนานนฺ}ติ ขตฺติยา จ พฺราหฺมณา จ เวสฺสา จ สุทฺทา จ คหฏฺฐา จ ปพฺพชิตา จ เทวา จ มนุสฺสา จ, เต พหุกาหิ วาจาหิ อนิฏฺฐาหิ อกนฺตาหิ อมนาปาหิ อกฺโกเสยฺยุํ ปริภาเสยฺยุํ โรเสยฺยุํ วิโรเสยฺยุํ หิํเสยฺยุํ วิหิํเสยฺยุํ เหเฐยฺยุํ วิเหเฐยฺยุํ ฆาเตยฺยุํ อุปฆาเตยฺยุํ อุปฆาตํ กเรยฺยุํ, เตสํ พหุํ วาจํ อนิฏฺฐํ อกนฺตํ อมนาปํ สุตฺวา สุณิตฺวา อุคฺคหิตฺวา อุปธารยิตฺวา อุปลกฺขยิตฺวาติ สุตฺวา \textbf{รุสิโต} \textbf{พหุํ วาจํ,} สมณานํ วา ปุถุชนานํ\textbf{.}}

\prose{\textbf{ผรุเสน เน น ปฏิวชฺชา}ติ \textbf{ผรุเสนา}ติ ผรุเสน กกฺขเฬน น ปฏิวชฺชา นปฺปฏิภเณยฺย, อกฺโกสนฺตํ น ปจฺจกฺโกเสยฺย, โรสนฺตํ นปฺปฏิโรเสยฺย, ภณฺฑนํ นปฺปฏิภณฺเฑยฺย, น กลหํ กเรย, น ภณฺฑนํ กเรยฺย, น วิคฺคหํ กเรยฺย, น วิวาทํ กเรยฺย, น เมธคํ กเรยฺย, กลหภณฺฑนวิคฺคหวิวาทเมธคํ ปชเหยฺย วิโนเทยฺย พฺยนฺติํ กเรยฺย อนภาวํ คเมยฺย, กลหภณฺฑนวิคฺคหวิวาทเมธคา อารโต อสฺส วิรโต ปฏิวิรโต นิกฺขนฺโต นิสฺสโฏ วิปฺปมุตฺโต วิสญฺญุตฺโต วิมริยาทิกเตน เจตสา วิหเรยฺยาติ ผรุเสน เน น ปฏิวชฺชา\textbf{.}}

\prose{\textbf{น หิ สนฺโต ปฏิเสนิํ กโรนฺตี}ติ \textbf{สนฺโต}ติ ราคสฺส สนฺตตฺตา \textbf{สนฺโต}, โทสสฺส\textbf{.}\csromanspacer{} โมหสฺส\textbf{.}\csromanspacer{} โกธสฺส\textbf{.}\csromanspacer{} อุปนาหสฺส ฯเปฯ สพฺพากุสลาภิสงฺขารานํ สนฺตตฺตา สมิตตฺตา วูปสมิตตฺตา วิชฺฌาตตฺตา นิพฺพุตตฺตา วิคตตฺตา ปฏิปสฺสทฺธตฺตา \textbf{สนฺโต} อุป\textbf{สนฺโต} วูป\textbf{สนฺโต} นิพฺพุโต ปฏิปสฺสทฺโธติ \textbf{สนฺโต.}\csromanspacer{} \textbf{น หิ สนฺโต ปฏิเสนิํ กโรนฺตี}ติ \textbf{สนฺโต} ปฏิเสนิํ ปฏิมลฺลํ ปฏิกณฺฏกํ\footnote{ปฏิพนฺธนํ (สี)} ปฏิปกฺขํ น กโรนฺติ น ชเนนฺติ น \csromanfolio{312}สญฺชเนนฺติ นิพฺพตฺเตนฺติ นาภินิพฺพตฺเตนฺตีติ น หิ สนฺโต ปฏิเสนิํ กโรนฺติ.\csromanspacer{} เตนาห ภควา\csromandash{}}

\csromangathameasure{สุตฺวา รุสิโต พหุํ วาจํ, \\ สมณานํ วา ปุถุชนานํ. \\ ผรุเสน เน น ปฏิวชฺชา, \\ น หิ สนฺโต ปฏิเสนิํ กโรนฺตีติ.}
\csromangathagroup{\csromangathastanza{สุตฺวา รุสิโต พหุํ วาจํ, \\ สมณานํ วา ปุถุชนานํ. \\ ผรุเสน เน น ปฏิวชฺชา, \\ น หิ สนฺโต ปฏิเสนิํ กโรนฺตีติ.}}

\proseitem{168}{\csromansymbolfootnote{*}{ขุ 1. 423 ปิฏฺเฐปิ.}เอตญฺจ ธมฺมมญฺญาย,}

\prose{\textbf{วิจินํ ภิกฺขุ สทา สโต สิกฺเข.}}

\csromangathameasure{\textbf{สนฺตีติ นิพฺพุติํ ญ ตฺวา,} \\ \textbf{สาสเน โคตมสฺส นปฺปมชฺเชยฺย.}}
\csromangathagroup{\csromangathastanza{\textbf{สนฺตีติ นิพฺพุติํ ญ ตฺวา,} \\ \textbf{สาสเน โคตมสฺส นปฺปมชฺเชยฺย.}}}

\prose{\textbf{เอตญฺจ ธมฺมมญฺญายา}ติ \textbf{เอตนฺ}ติ อาจิกฺขิตํ เทสิตํ ปญฺญปิตํ ปฏฺฐปิตํ วิวฏํ วิภตฺตํ อุตฺตานีกตํ ปกาสิตํ ธมฺมํ อญฺญาย ชานิตฺวา ตุลยิตฺวา ตีรยิตฺวา วิภาวยิตฺวา วิภูตํ กตฺวาติ เอวมฺปิ เอตญฺจ ธมฺมมญฺญาย.\csromanspacer{}\csromanspacer{}อถ วา สมญฺจ วิสมญฺจ, ปถญฺจ วิปถญฺจ, สาวชฺชญฺจ อนวชฺชญฺจ, หีนญฺจ ปณีตญฺจ, กณฺหญฺจ สุกฺกญฺจ, วิญฺญูครหิตญฺจ วิญฺญูปสตฺถญฺจ ธมฺมํ อญฺญาย ชานิตฺวา ตุลยิตฺวา ตีรยิตฺวา วิภาวยิตฺวา วิภูตํ กตฺวาติ เอวมฺปิ เอตญฺจ ธมฺมมญฺญาย.\csromanspacer{}\csromanspacer{}อถ วา สมฺมาปฏิปทํ อนุโลมปฏิปทํ อปจฺจนีกปฏิปทํ อวิรุทฺธปฏิปทํ อนฺวตฺถปฏิปทํ ธมฺมานุธมฺมปฏิปทํ สีเลสุ ปริปูรการิตํ อินฺทฺริเยสุ คุตฺตทฺวารตํ โภชเน มตฺตญฺญุตํ ชาคริยานุโยคํ สติสมฺปชญฺญํ, จตฺตาโร สติปฏฺฐาเน จตฺตาโร สมฺมปฺปธาเน จตฺตาโร อิทฺธิปาเท ปญฺจินฺทฺริยานิ ปญฺจ พลานิ สตฺต โพชฺฌงฺเค อริยํ อฏฺฐงฺคิกํ มคฺคํ, นิพฺพานญฺจ นิพฺพานคามินิญฺจ ปฏิปทํ ธมฺมํ อญฺญาย ชานิตฺวา ตุลยิตฺวา ตีรยิตฺวา วิภาวยิตฺวา วิภูตํ กตฺวาติ เอวมฺปิ ตญฺจ ธมฺมมญฺญาย.}

\prose{\textbf{วิจินํ ภิกฺขุ สทา สโต สิกฺเข}ติ \textbf{วิจินนฺ}ติ \textbf{วิจินนฺ}โต ป\textbf{วิจินนฺ}โต ตุลยนฺโต ตีรยนฺโต วิภาวยนฺโต วิภูตํ กโรนฺโต “สพฺเพ สงฺขารา อนิจฺจา”ติ ฯเปฯ “ยํ กิญฺจิ สมุทยธมฺมํ, สพฺพํ ตํ นิโรธธมฺมนฺ”ติ \textbf{วิจินนฺ}โต ป\textbf{วิจินนฺ}โต ตุลยนฺโต ตีรยนฺโต วิภาวยนฺโต วิภูตํ กโรนฺโตติ วิจินํ ภิกฺขุ.\csromanspacer{} \textbf{สทา}ติ สทา สพฺพทา สพฺพกาลํ ฯเปฯ ปจฺฉิเม วโยขนฺเธ.\csromanspacer{} \textbf{สโต}ติ จตูหิ การเณหิ สโต, กาเย กายานุปสฺสนาสติปฏฺฐานํ ภาเวนฺโต สโต ฯเปฯ โส วุจฺจติ สโต.}

\vspace{\tipitakaparskip}
\csromancenter{ตุวฏกสุตฺตนิทฺเทส \textbf{สิกฺเข}ติ ติสฺโส สิกฺขาโย อธิสีลสิกฺขา อธิจิตฺตสิกฺขา อธิปญฺญาสิกฺขา ฯเปฯ อยํ อธิปญฺญาสิกฺขา.\csromanspacer{} อิมา ติสฺโส สิกฺขาโย อาวชฺชนฺโต สิกฺเขยฺย ฯเปฯ สิกฺเขยฺย อาจเรยฺย สมาจเรยฺย สมาทาย วตฺเตยฺยาติ วิจินํ ภิกฺขุ สทา สโต สิกฺเข.}
\prose{\csromanfolio{313}\textbf{สนฺตีติ นิพฺพุติํ ญตฺวา}ติ ราคสฺส นิพฺพุติํ สนฺตีติ ญตฺวา, โทสสฺส.\csromanspacer{} โมหสฺส ฯเปฯ สพฺพากุสลาภิสงฺขารานํ นิพฺพุติํ สนฺตีติ ญตฺวา ชานิตฺวา ตุลยิตฺวา ตีรยิตฺวา วิภาวยิตฺวา วิภูตํ กตฺวาติ สนฺตีติ นิพฺพุติํ ญตฺวา.}

\prose{\textbf{สาสเน โคตมสฺส นปฺปมชฺเชยฺยา}ติ โคตมสฺส สาสเน พุทฺธสาสเน ชินสาสเน ตถาคตสาสเน เทวสาสเน อรหนฺตสาสเน.\csromanspacer{} \textbf{นปฺปมชฺเชยฺยา}ติ สกฺกจฺจการี อสฺส สาตจฺจการี อฏฺฐิตการี อโนลีนวุตฺติโก อนิกฺขิตฺตจฺฉนฺโท อนิกฺขิตฺตธุโร กุสเลสุ ธมฺเมสุ, กทาหํ อปริปูรํ วา สีลกฺขนฺธํ ปริปูเรยฺยํ ฯเปฯ อปริปูรํ วา สมาธิกฺขนฺธํ.\csromanspacer{} ปญฺญากฺขนฺธํ.\csromanspacer{} วิมุตฺติกฺขนฺธํ.\csromanspacer{} วิมุตฺติญาณทสฺสนกฺขนฺธํ.\csromanspacer{} กทาหํ อปริญฺญาตํ วา ทุกฺขํ ปริชาเนยฺยํ.\csromanspacer{} อปฺปหีเน วา กิเลเส ปชเหยฺยํ, อภาวิตํ วา มคฺคํ ภาเวยฺยํ, อสจฺฉิกตํ วา นิโรธํ สจฺฉิกเรยฺยนฺติ โย ตตฺถ ฉนฺโท จ วายาโม จ อุสฺสาโห จ อุสฺโสฬฺหี จ ถาโม จ อปฺปฏิวานี จ สติ จ สมฺปชญฺญญฺจ อาตปฺปํ ปธานํ อธิฏฺฐานํ อนุโยโค อปฺปมาโท กุสเลสุ ธมฺเมสูติ สาสเน โคตมสฺส นปฺปมชฺเชยฺย.\csromanspacer{} เตนาห ภควา\csromandash{}}

\csromangathameasure{เอตญฺจ ธมฺมมญฺญาย, \\ วิจินํ ภิกฺขุ สทา สโต สิกฺเข. \\ \textbf{สนฺตีติ นิพฺพุติํ ญตฺวา},}
\csromangathagroup{\csromangathastanza{เอตญฺจ ธมฺมมญฺญาย, \\ วิจินํ ภิกฺขุ สทา สโต สิกฺเข. \\ \textbf{สนฺตีติ นิพฺพุติํ ญตฺวา},}}

\prose{\textbf{สาสเน โคตมสฺส นปฺปมชฺเชยฺยา}ติ.}

\proseitem{169}{\csromansymbolfootnote{*}{ขุ 1. 424 ปิฏฺเฐปิ.}อภิภู หิ โส อนภิภูโต,}

\prose{\textbf{สกฺขิธมฺม’มนีติห’มทฺทสิ.}}

\csromangathameasure{\textbf{ตสฺมา หิ ตสฺส ภควโต สาสเน,}}
\csromangathagroup{\csromangathastanza{\textbf{ตสฺมา หิ ตสฺส ภควโต สาสเน,}}}

\vspace{\tipitakaparskip}
\csromancenter{\textbf{อปฺปมตฺโต สทา นมสฺส’มนุสิกฺเข. (อิติ ภควา,)}}
\prose{\csromanfolio{314}\textbf{อภิภู หิ โส อนภิภูโต}ติ \textbf{อภิภู}ติ รูปาภิภู\footnote{อภิภูตรูปา (สฺยา) เอวมญฺเญสุ ปญฺจปเทสุปิ.} สทฺทาภิภู คนฺธาภิภู รสาภิภู โผฏฺฐพฺพาภิภู ธมฺมาภิภู อนฺ\textbf{อภิภู}โต เกหิจิ กิเลเสหิ, อภิโภสิ เน ปาปเก\footnote{อภิภู หิ เน หีเน ปาปเก (สี, ก), อภิภู หิ ปาปเก (สฺยา)} อกุสเล ธมฺเม สํกิเลสิเก โปโนภวิเก สทเร ทุกฺขวิปาเก อายติํ ชาติชรามรณิเยติ \textbf{อภิภู} หิ โส อนฺ\textbf{อภิภู}โต.}

\prose{\textbf{สกฺขิธมฺม’มนีติห’มทฺทสี}ติ \textbf{สกฺขิธมฺมนฺ}ติ น อิติหิติหํ น อิติกิริยาย น ปรมฺปราย น ปิฏกสมฺปทาย น ตกฺกเหตุ น นยเหตุ น อาการปริวิตกฺเกน น ทิฏฺฐินิชฺฌานกฺขนฺติยา สามํ สยมภิญฺญาตํ อตฺตปจฺจกฺขธมฺมํ อทฺทสิ อทฺทกฺขิ อปสฺสิ ปฏิวิชฺฌีติ สกฺขิธมฺมมนีติหมทฺทสิ.}

\prose{\textbf{ตสฺมา หิ ตสฺส ภควโต สาสเน}ติ \textbf{ตสฺมา}ติ \textbf{ตสฺมา} ตํการณา ตํเหตุ ตปฺปจฺจยา ตํนิทานา.\csromanspacer{} \textbf{ตสฺส ภควโต สาสเน}ติ ตสฺส ภควโต สาสเน โคตมสาสเน พุทฺธสาสเน ชินสาสเน ตถาคตสาสเน เทวสาสเน อรหนฺตสาสเนติ \textbf{ตสฺมา} ตสฺส ภควโต สาสเน.}

\prose{\textbf{อปฺปมตฺโต สทา นมสฺส’มนุสิกฺเข. (อิติ ภควา,)}ติ \textbf{อปฺปมตฺโต}ติ สกฺกจฺจการี ฯเปฯ อปฺปมาโท กุสเลสุ ธมฺเมสุ.\csromanspacer{} \textbf{สทา}ติ สทา สพฺพกาลํ ฯเปฯ ปจฺฉิเม วโยขนฺเธ.\csromanspacer{} \textbf{นมสฺสนฺ}ติ กาเยน วา นมสฺสมาโน วาจาย วา นมสฺสมาโน จิตฺเตน วา นมสฺสมาโน อนฺวตฺถปฏิปตฺติยา วา นมสฺสมาโน ธมฺมานุธมฺมปฏิปตฺติยา วา นมสฺสมาโน สกฺกุรุมาโน ครุกุรุมาโน\footnote{สกฺการมาโน ครุการมาโน (สฺยา)} มานยมาโน ปูชยมาโน อปจยมาโน.\csromanspacer{} \textbf{อนุสิกฺเข}ติ ติสฺโส สิกฺขาโย อธิสีลสิกฺขา อธิจิตฺตสิกฺขา อธิปญฺญาสิกฺขา ฯเปฯ อยํ อธิปญฺญาสิกฺขา.\csromanspacer{} อิมา ติสฺโส สิกฺขาโย อาวชฺชนฺโต สิกฺเขยฺย ฯเปฯ สจฺฉิกาตพฺพํ สจฺฉิกโรนฺโต สิกฺเขยฺย จเรยฺย อาจเรยฺย สมาจเรยฺย สมาทาย วตฺเตยฺย.\csromanspacer{} \textbf{ภควา}ติ คารวาธิวจนํ ฯเปฯ สจฺฉิกา ปญฺญตฺติ, ยทิทํ \textbf{ภควา}ติ \textbf{อปฺปมตฺโต} สทา นมสฺส’มนุสิกฺเข. (อิติ \textbf{ภควา},) เตนาห \textbf{ภควา}\csromandash{}}

\csromanheader{2}{15. \textbf{อตฺตทณฺฑสุตฺตนิทฺเทส}}
\csromangathameasure{อภิภู หิ โส อนภิภูโต, \\ สกฺขิธมฺม’มนีติห’มทฺทสิ. \\ ตสฺมา หิ ตสฺส ภควโต สาสเน,}
\csromangathagroup{\csromangathastanza{\csromanfolio{315}อภิภู หิ โส อนภิภูโต, \\ สกฺขิธมฺม’มนีติห’มทฺทสิ. \\ ตสฺมา หิ ตสฺส ภควโต สาสเน,}}

\vspace{\tipitakaparskip}
\csromancenter{อปฺปมตฺโต สทา นมสฺส’มนุสิกฺเข. (อิติ ภควาติ,) ตุวฏกสุตฺตนิทฺเทโส\footnote{ตุวฏฺฏสุตฺตนิทฺเทโส (ก) สุตฺตนิปาเตปิ.} จุทฺทสโม.}
\csromansectionrule
\csromanheader{2}{15. \textbf{อตฺตทณฺฑสุตฺตนิทฺเทส}}
\csromanheader{2}{อถ \textbf{อตฺตทณฺฑสุตฺตนิทฺเทส}ํ วกฺขติ\csromandash{}}
\proseitem{170}{\csromansymbolfootnote{*}{ขุ 1. 424 ปิฏฺเฐปิ.}อตฺต\textbf{ทณฺฑา} ภยํ ชาตํ, ชนํ ปสฺสถ เมธคํ.}

\csromangathameasure{สํเวคํ กิตฺตยิสฺสามิ, ยถา สํวิชิตํ มยา.}
\csromangathagroup{\csromangathastanza{สํเวคํ กิตฺตยิสฺสามิ, ยถา สํวิชิตํ มยา.}}

\prose{\textbf{อตฺตทณฺฑา ภยํ ชาตนฺ}ติ \textbf{ทณฺฑา}ติ ตโย \textbf{ทณฺฑา} กายทณฺโฑ วจีทณฺโฑ มโนทณฺโฑ.\csromanspacer{} ติวิธํ กายทุจฺจริตํ กายทณฺโฑ, จตุพฺพิธํ วจีทุจฺจริตํ วจีทณฺโฑ, ติวิธํ มโนทุจฺจริตํ มโนทณฺโฑ.\csromanspacer{} \textbf{ภยนฺ}ติ เทฺว ภยานิ ทิฏฺฐธมฺมิกญฺจ ภยํ สมฺปรายิกญฺจ ภยํ.\csromanspacer{} กตมํ ทิฏฺฐธมฺมิกํ ภยํ, อิเธกจฺโจ กาเยน ทุจฺจริตํ จรติ, วาจาย ทุจฺจริตํ จรติ, มนสา ทุจฺจริตํ จรติ, ปาณมฺปิ หนติ, อทินฺนมฺปิ อาทิยติ, สนฺธิมฺปิ ฉินฺทติ, นิลฺโลปมฺปิ หรติ, เอกาคาริกมฺปิ กโรติ, ปริปนฺเถปิ ติฏฺฐติ, ปรทารมฺปิ คจฺฉติ, มุสาปิ ภณติ, ตเมนํ คเหตฺวา รญฺโญ ทสฺเสนฺติ “อยํ เทว โจโร อาคุจารี, อิมสฺส ยํ อิจฺฉสิ ตํ ทณฺฑํ ปเณหี”ติ.\csromanspacer{} ตเมนํ ราชา ปริภาสติ, โส ปริภาสปจฺจยา ภยมฺปิ อุปฺปาเทติ, ทุกฺขํ โทมนสฺสํ\footnote{ทุกฺขโทมนสฺสํ (สฺยา)} ปฏิสํเวเทติ.\csromanspacer{} เอตํ ภยํ ทุกฺขํ โทมนสฺสํ กุโต ตสฺส, อตฺตทณฺฑโต ชาตํ สญฺชาตํ นิพฺพตฺตํ อภินิพฺพตฺตํ ปาตุภูตํ.}

\prose{เอตฺตเกนปิ ราชา น ตุสฺสติ, ตเมนํ ราชา พนฺธาเปติ อนฺทุพนฺธเนน วา รชฺชุพนฺธเนน วา สงฺขลิกพนฺธเนน วา เวตฺตพนฺธเนน วา ลตาพนฺธเนน วา ปกฺเขปพนฺธเนน วา\footnote{เปกฺขพนฺธเนน วา (สฺยา) ขุ 8. 235 ปิฏฺเฐปิ.} ปริกฺเขปพนฺธเนน วา คามพนฺธเนน วา นิคมพนฺธเนน วา นครพนฺธเนน วา รฏฺฐพนฺธเนน วา ชนปทพนฺธเนน วา อนฺตมโส สวจนียมฺปิ กโรติ “น เต ลพฺภา อิโต ปกฺกมิตุนฺ”ติ.\csromanspacer{} โส พนฺธนปจฺจยาปิ ทุกฺขํ \csromanfolio{316}โทมนสฺสํ ปฏิสํเวเทติ.\csromanspacer{} เอตํ ภยํ ทุกฺขํ โทมนสฺสํ กุโต ตสฺส, อตฺตทณฺฑโต ชาตํ สญฺชาตํ นิพฺพตฺตํ อภินิพฺพตฺตํ ปาตุภูตํ.}

\prose{เอตฺตเกนปิ ราชา น ตุสฺสติ, ราชา ตสฺส ธนํ อาหราเปติ สตํ วา สหสฺสํ วา สตสหสฺสํ วา.\csromanspacer{} โส ธนชานิปจฺจยาปิ ทุกฺขํ โทมนสฺสํ ปฏิสํเวเทติ.\csromanspacer{} เอตํ ภยํ ทุกฺขํ โทมนสฺสํ กุโต ตสฺส, อตฺตทณฺฑโต ชาตํ สญฺชาตํ นิพฺพตฺตํ อภินิพฺพตฺตํ ปาตุภูตํ.}

\prose{เอตฺตเกนปิ ราชา น ตุสฺสติ.\csromanspacer{} ตเมนํ ราชา วิวิธา กมฺมการณา การาเปติ, * กสาหิปิ ตาเฬติ, เวตฺเตหิปิ ตาเฬติ, อฑฺฒทณฺฑเกหิปิ ตาเฬติ, หตฺถมฺปิ ฉินฺทติ, ปาทมฺปิ ฉินฺทติ, หตฺถปาทมฺปิ ฉินฺทติ, กณฺณมฺปิ ฉินฺทติ, นาสมฺปิ ฉินฺทติ, กณฺณนาสมฺปิ ฉินฺทติ, พิลงฺคถาลิกมฺปิ กโรติ, สงฺขมุณฺฑิกมฺปิ กโรติ, ราหุมุขมฺปิ กโรติ, โชติมาลิกมฺปิ กโรติ, หตฺถปชฺโชติกมฺปิ กโรติ, เอรกวตฺติกมฺปิ กโรติ, จีรกวาสิกมฺปิ กโรติ, เอเณยฺยกมฺปิ กโรติ, พฬิสมํสิกมฺปิ กโรติ, กหาปณิกมฺปิ กโรติ, ขาราปตจฺฉิกมฺปิ กโรติ, ปลิฆปริวตฺตกมฺปิ กโรติ, ปลาลปีฐกมฺปิ กโรติ, ตตฺเตนปิ เตเลน โอสิญฺจติ, สุนเขหิปิ ขาทาเปหิ, ชีวนฺตมฺปิ สูเล อุตฺตาเสติ, อสินาปิ สีสํ ฉินฺทติ, โส กมฺมการณปจฺจยาปิ ทุกฺขํ โทมนสฺสํ ปฏิสํเวเทติ.\csromanspacer{} เอตํ ภยํ ทุกฺขํ โทมนสฺสํ กุโต ตสฺส, อตฺตทณฺฑโต ชาตํ สญฺชาตํ นิพฺพตฺตํ อภินิพฺพตฺตํ ปาตุภูตํ.\csromanspacer{} ราชา อิเมสํ จตุนฺนํ ทณฺฑานํ อิสฺสโร.}

\prose{โส สเกน กมฺเมน กายสฺส เภทา ปรํ มรณา อปายํ ทุคฺคติํ วินิปาตํ นิรยํ อุปปชฺชติ, ตเมนํ นิรยปาลา ปญฺจวิธพนฺธนํ นาม กมฺมการณํ กาเรนฺติ, ตตฺตํ อโยขิลํ หตฺเถ คเมนฺติ, ตตฺตํ อโยขิลํ ทุติเย หตฺเถ คเมนฺติ, ตตฺตํ อโยขิลํ ปาเท คเมนฺติ, ตตฺตํ อโยขิลํ ทุติเย ปาเท คเมนฺติ, ตตฺตํ อโยขิลํ มชฺเฌ อุรสฺมิํ คเมนฺติ.\csromanspacer{} โส ตตฺถ ทุกฺขา ติพฺพา\footnote{ติปฺปา (สฺยา)} กฏุกา เวทนา เวเทติ, น จ ตาว กาลํกโรติ, ยาว น ตํ ปาปกมฺมํ พฺยนฺตีโหติ.\csromanspacer{} เอตํ ภยํ ทุกฺขํ โทมนสฺสํ กุโต ตสฺส, อตฺตทณฺฑโต ชาตํ สญฺชาตํ นิพฺพตฺตํ อภินิพฺพตฺตํ ปาตุภูตํ.}

\csromanheader{2}{15. อตฺตทณฺฑสุตฺตนิทฺเทส}
\prose{\csromanfolio{317}ตเมนํ นิรยปาลา สํเวเสตฺวา\footnote{สํเวสิตฺวา (สฺยา) ม 3. 221; ขุ 8. 237 ปิฏฺเฐสุปิ.} กุฐารีหิ\footnote{กุธารีหิ (สฺยา, ก) 3. อุทฺธปาทํ (สี)} ตจฺเฉนฺติ.\csromanspacer{} โส ตตฺถ ทุกฺขา ติพฺพา กฏุกา เวทนา เวเทติ, น จ ตาว กาลํ กโรติ ยาว น ตํ ปาปกมฺมํ พฺยนฺตีโหติ.\csromanspacer{} ตเมนํ นิรยปาลา อุทฺธํปาทํ3 อโธสิรํ คเหตฺวา วาสีหิ ตจฺเฉนฺติ, ตเมนํ นิรยปาลา รเถ โยเชตฺวา อาทิตฺตาย ปถวิยา สมฺปชฺชลิตาย สโชติภูตาย สาเรนฺติปิ ปจฺจาสาเรนฺติปิ\footnote{หาเรนฺติปิ ปจฺจาหาเรนฺติปิ (สี, ก) 5. ...ปริยตฺโต (สฺยา, ก)} ฯเปฯ.\csromanspacer{} ตเมนํ นิรยปาลา มหนฺตํ องฺคารปพฺพตํ อาทิตฺตํ สมฺปชฺชลิตํ สโชติภูตํ อาโรเปนฺติปิ โอโรเปนฺติปิ ฯเปฯ.\csromanspacer{} ตเมนํ นิรยปาลา อุทฺธํปาทํ อโธสิรํ คเหตฺวา ตตฺตาย โลหกุมฺภิยา ปกฺขิปนฺติ อาทิตฺตาย สมฺปชฺชลิตาย สโชติภูตาย, โส ตตฺถ เผณุทฺเทหกํ ปจฺจติ, โส ตตฺถ เผณุทฺเทหกํ ปจฺจมาโน สกิมฺปิ อุทฺธํ คจฺฉติ, สกิมฺปิ อโธ คจฺฉติ, สกิมฺปิ ติริยํ คจฺฉติ.\csromanspacer{} โส ตตฺถ ทุกฺขา ติพฺพา กฏุกา เวทนา เวเทติ, น จ ตาว กาลํ กโรติ, ยาว น ตํ ปาปกมฺมํ พฺยนฺตีโหติ.\csromanspacer{} เอตํ ภยํ ทุกฺขํ โทมนสฺสํ กุโต ตสฺส, อตฺตทณฺฑโต ชาตํ สญฺชาตํ นิพฺพตฺตํ อภินิพฺพตฺตํ ปาตุภูตํ.\csromanspacer{} ตเมนํ นิรยปาลา มหานิรเย ปกฺขิปนฺติ.\csromanspacer{} โส โข ปน มหานิรโย\csromandash{}}

\csromangathameasure{*จตุกฺกณฺโณ จตุทฺวาโร, วิภตฺโต ภาคโส มิโต. \\ อโยปาการปริยนฺโต5, อยสา ปฏิกุชฺชิโต. \\ ตสฺส อโยมยา ภูมิ, ชลิตา เตชสา ยุตา. \\ สมนฺตา โยชนสตํ, ผริตฺวา ติฏฺฐติ สพฺพทา. \\ กทริยาตปนา โฆรา, อจฺจิมนฺโต ทุราสทา. \\ โลมหํสนรูปา จ, เภสฺมา ปฏิภยา ทุขา. \\ ปุรตฺถิมาย ภิตฺติยา, อจฺจิกฺขนฺโธ สมุฏฺฐิโต. \\ ทหนฺโต ปาปกมฺมนฺเต, ปจฺฉิมาย ปฏิหญฺญติ. \\ ปจฺฉิมาย จ ภิตฺติยา, อจฺจิกฺขนฺโธ สมุฏฺฐิโต. \\ ทหนฺโต ปาปกมฺมนฺเต, ปุรตฺถิมาย ปฏิหญฺญติ. \\ อุตฺตราย จ ภิตฺติยา, อจฺจิกฺขนฺโธ สมุฏฺฐิโต. \\ ทหนฺโต ปาปกมฺมนฺเต, ทกฺขิณาย ปฏิหญฺญติ. \\ ทกฺขิณาย จ ภิตฺติยา, อจฺจิกฺขนฺโธ สมุฏฺฐิโต. \\ ทหนฺโต ปาปกมฺมนฺเต, อุตฺตราย ปฏิหญฺญติ. \\ เหฏฺฐโต จ สมุฏฺฐาย, อจฺจิกฺขนฺโธ ภยานโก. \\ ทหนฺโต ปาปกมฺมนฺเต, ฉทนสฺมิํ ปฏิหญฺญติ. \\ ฉทนมฺหา สมุฏฺฐาย, อจฺจิกฺขนฺโธ ภยานโก. \\ ทหนฺโต ปาปกมฺมนฺเต, ภูมิยํ ปฏิหญฺญติ. \\ อโยกปาลมาทิตฺตํ, สนฺตตฺตํ ชลิตํ ยถา. \\ เอวํ อวีจินิรโย, เหฏฺฐา อุปริ ปสฺสโต. \\ ตตฺถ สตฺตา มหาลุทฺทา, มหากิพฺพิสการิโน. \\ อจฺจนฺตปาปกมฺมนฺตา, ปจฺจนฺติ น จ มิยฺยเร. \\ ชาตเวทสโม กาโย, เตสํ นิรยวาสินํ. \\ ปสฺส กมฺมานํ ทฬฺหตฺตํ, น ภสฺมา โหติ นปี มสิ. \\ ปุรตฺถิเมนปิ ธาวนฺติ, ตโต ธาวนฺติ ปจฺฉิมํ. \\ อุตฺตเรนปิ ธาวนฺติ, ตโต ธาวนฺติ ทกฺขิณํ. \\ ยํ ยํ ทิสํ ปธาวนฺติ, ตํ ตํ ทฺวารํ ปิธียติ. \\ อภินิกฺขมิตาสา เต, สตฺตา โมกฺขคเวสิโน. \\ น เต ตโต นิกฺขมิตุํ, ลภนฺติ กมฺมปจฺจยา. \\ เตสญฺจ ปาปกมฺมนฺตํ, อวิปกฺกํ กตํ พหุนฺติ.}
\csromangathagroup{\csromangathastanza{\csromansymbolfootnote{*}{อํ 1. 140; ขุ 2. 134, 152 ปิฏฺฐาทีสุปิ.}จตุกฺกณฺโณ จตุทฺวาโร, วิภตฺโต ภาคโส มิโต. \\ อโยปาการปริยนฺโต5, อยสา ปฏิกุชฺชิโต.}\csromangathastanza{ตสฺส อโยมยา ภูมิ, ชลิตา เตชสา ยุตา. \\ สมนฺตา โยชนสตํ, ผริตฺวา ติฏฺฐติ สพฺพทา.}\csromangathastanza{กทริยาตปนา โฆรา, อจฺจิมนฺโต ทุราสทา. \\ โลมหํสนรูปา จ, เภสฺมา ปฏิภยา ทุขา.}\csromangathastanza{ปุรตฺถิมาย ภิตฺติยา, อจฺจิกฺขนฺโธ สมุฏฺฐิโต. \\ ทหนฺโต ปาปกมฺมนฺเต, ปจฺฉิมาย ปฏิหญฺญติ.}\csromangathastanza{ปจฺฉิมาย จ ภิตฺติยา, อจฺจิกฺขนฺโธ สมุฏฺฐิโต. \\ ทหนฺโต ปาปกมฺมนฺเต, ปุรตฺถิมาย ปฏิหญฺญติ.}\csromangathastanza{อุตฺตราย จ ภิตฺติยา, อจฺจิกฺขนฺโธ สมุฏฺฐิโต. \\ ทหนฺโต ปาปกมฺมนฺเต, ทกฺขิณาย ปฏิหญฺญติ.}\csromangathastanza{\csromanfolio{318}ทกฺขิณาย จ ภิตฺติยา, อจฺจิกฺขนฺโธ สมุฏฺฐิโต. \\ ทหนฺโต ปาปกมฺมนฺเต, อุตฺตราย ปฏิหญฺญติ.}\csromangathastanza{เหฏฺฐโต จ สมุฏฺฐาย, อจฺจิกฺขนฺโธ ภยานโก. \\ ทหนฺโต ปาปกมฺมนฺเต, ฉทนสฺมิํ ปฏิหญฺญติ.}\csromangathastanza{ฉทนมฺหา สมุฏฺฐาย, อจฺจิกฺขนฺโธ ภยานโก. \\ ทหนฺโต ปาปกมฺมนฺเต, ภูมิยํ ปฏิหญฺญติ.}\csromangathastanza{อโยกปาลมาทิตฺตํ, สนฺตตฺตํ ชลิตํ ยถา. \\ เอวํ อวีจินิรโย, เหฏฺฐา อุปริ ปสฺสโต.}\csromangathastanza{ตตฺถ สตฺตา มหาลุทฺทา, มหากิพฺพิสการิโน. \\ อจฺจนฺตปาปกมฺมนฺตา, ปจฺจนฺติ น จ มิยฺยเร\footnote{มียเร (สี)}.}\csromangathastanza{ชาตเวทสโม กาโย, เตสํ นิรยวาสินํ. \\ ปสฺส กมฺมานํ ทฬฺหตฺตํ, น ภสฺมา โหติ นปี มสิ.}\csromangathastanza{ปุรตฺถิเมนปิ ธาวนฺติ, ตโต ธาวนฺติ ปจฺฉิมํ. \\ อุตฺตเรนปิ ธาวนฺติ, ตโต ธาวนฺติ ทกฺขิณํ.}\csromangathastanza{ยํ ยํ ทิสํ ปธาวนฺติ\footnote{ทิสมฺปิ ธาวนฺติ (สฺยา)}, ตํ ตํ ทฺวารํ ปิธียติ\footnote{ปิถียติ (สี, สฺยา)}. \\ อภินิกฺขมิตาสา เต, สตฺตา โมกฺขคเวสิโน.}\csromangathastanza{น เต ตโต นิกฺขมิตุํ, ลภนฺติ กมฺมปจฺจยา. \\ เตสญฺจ ปาปกมฺมนฺตํ, อวิปกฺกํ กตํ พหุนฺติ.}}

\prose{เอตํ ภยํ ทุกฺขํ โทมนสฺสํ กุโต ตสฺส, อตฺตทณฺฑโต ชาตํ สญฺชาตํ นิพฺพตฺตํ อภินิพฺพตฺตํ ปาตุภูตํ.\csromanspacer{} ยานิ จ เนรยิกานิ ทุกฺขานิ, ยานิ จ ติรจฺฉานโยนิกานิ ทุกฺขานิ, ยานิ จ เปตฺติวิสยิกานิ ทุกฺขานิ, ยานิ จ มานุสิกานิ ทุกฺขานิ, ตานิ กุโต ชาตานิ กุโต สญฺชาตานิ กุโต นิพฺพตฺตานิ กุโต อภินิพฺพตฺตานิ กุโต ปาตุภูตานิ, อตฺตทณฺฑโต ชาตานิ สญฺชาตานิ นิพฺพตฺตานิ อภินิพฺพตฺตานิ ปาตุภูตานีติ อตฺตทณฺฑา ภยํ ชาตํ.}

\prose{\textbf{ชนํ ปสฺสถ เมธค}นฺติ \textbf{ชนนฺ}ติ ขตฺติยา จ พฺราหฺมณา จ เวสฺสา จ สุทฺทา จ คหฏฺฐา จ ปพฺพชิตา จ เทวา จ มนุสฺสา จ เมธคํ ชนํ กลหํ ชนํ วิรุทฺธํ ชนํ ปฏิวิรุทฺธํ ชนํ อาหตํ ชนํ ปจฺจาหตํ ชนํ อาฆาติตํ ชนํ}

\vspace{\tipitakaparskip}
\csromancenter{อตฺตทณฺฑสุตฺตนิทฺเทส ปจฺจาฆาติตํ ชนํ ปสฺสถ ทกฺขถ โอโลเกถ นิชฺฌาเยถ อุปปริกฺขถาติ ชนํ ปสฺสถ เมธคํ.}
\prose{\csromanfolio{319}\textbf{สํเวคํ กิตฺตยิสฺสามี}ติ สํเวคํ อุพฺเพคํ อุตฺราสํ ภยํ ปีฬนํ ฆฏฺฏนํ อุปทฺทวํ อุปสคฺคํ.\csromanspacer{} \textbf{กิตฺตยิสฺสามี}ติ ปกิตฺตยิสฺสามิ อาจิกฺขิสฺสามิ เทเสสฺสามิ ปญฺญเปสฺสามิ ปฐเปสฺสามิ วิวริสฺสามิ วิภชิสฺสามิ อุตฺตานีกริสฺสามิ ปกาสิสฺสามีติ สํเวคํ กิตฺตยิสฺสามิ.}

\prose{\textbf{ยถา สํวิชิตํ มยา}ติ ยถา มยา อตฺตนาเยว อตฺตานํ สํเวชิโต อุพฺเพชิโต สํเวคมาปาทิโตติ ยถา สํวิชิตํ มยา.\csromanspacer{} เตนาห ภควา\csromandash{}}

\csromangathameasure{อตฺตทณฺฑา ภยํ ชาตํ, ชนํ ปสฺสถ เมธคํ. \\ สํเวคํ กิตฺตยิสฺสามิ, ยถา สํวิชิตํ มยาติ.}
\csromangathagroup{\csromangathastanza{อตฺตทณฺฑา ภยํ ชาตํ, ชนํ ปสฺสถ เมธคํ. \\ สํเวคํ กิตฺตยิสฺสามิ, ยถา สํวิชิตํ มยาติ.}}

\proseitem{171}{\csromansymbolfootnote{*}{ขุ 1. 424 ปิฏฺเฐปิ.}\textbf{ผนฺทมานํ ปชํ ทิสฺวา}, มจฺเฉ อปฺโปทเก ยถา.}

\csromangathameasure{อญฺญมญฺเญหิ พฺยารุทฺเธ, ทิสฺวา มํ ภยมาวิสิ.}
\csromangathagroup{\csromangathastanza{อญฺญมญฺเญหิ พฺยารุทฺเธ, ทิสฺวา มํ ภยมาวิสิ.}}

\prose{\textbf{ผนฺทมานํ ปชํ ทิสฺวา}ติ \textbf{ปชา}ติ สตฺตาธิวจนํ.\csromanspacer{} ปชํ ตณฺหาผนฺทนาย ผนฺทมานํ ทิฏฺฐิผนฺทนาย ผนฺทมานํ กิเลสผนฺทนาย ผนฺทมานํ ทุจฺจริตผนฺทาย ผนฺทมานํ ปโยคผนฺทนาย ผนฺทมานํ วิปากผนฺทนาย ผนฺทมานํ รตฺตํ ราเคน ผนฺทมานํ ทุฏฺฐํ โทเสน ผนฺทมานํ มูฬฺหํ โมเหน ผนฺทมานํ วินิพทฺธํ มาเนน ผนฺทมานํ ปรามฏฺฐํ ทิฏฺฐิยา ผนฺทมานํ วิกฺเขปคตํ อุทฺธจฺเจน ผนฺทมานํ อนิฏฺฐงฺคตํ วิจิกิจฺฉาย ผนฺทมานํ ถามคตํ อนุสเยหิ ผนฺทมานํ ลาเภน ผนฺทมานํ อลาเภน ผนฺทมานํ ยเสน ผนฺทมานํ อยเสน ผนฺทมานํ ปสํสาย ผนฺทมานํ นินฺทาย ผนฺทมานํ สุเขน ผนฺทมานํ ทุกฺเขน ผนฺทมานํ ชาติยา ผนฺทมานํ ชราย ผนฺทมานํ พฺยาธินา ผนฺทมานํ มรเณน ผนฺทมานํ โสกปริเทวทุกฺขโทมนสฺสุปายาเสหิ ผนฺทมานํ เนรยิเกน ทุกฺเขน ผนฺทมานํ ติรจฺฉานโยนิเกน ทุกฺเขน ผนฺทมานํ เปตฺติวิสยิเกน ทุกฺเขน ผนฺทมานํ มานุสิเกน ทุกฺเขน ผนฺทมานํ คพฺโภกฺกนฺติมูลเกน ทุกฺเขน.\csromanspacer{} คพฺภฏฺฐิติมูลเกน\footnote{คพฺเภฐิติมูลเกน (ก) 35 ปิฏฺเฐปิ.} ทุกฺเขน.\csromanspacer{} คพฺภวุฏฺฐานมูลเกน ทุกฺเขน.\csromanspacer{} ชาตสฺสูปนิพนฺธเกน ทุกฺเขน.\csromanspacer{} ชาตสฺส ปราเธยฺยเกน ทุกฺเขน.\csromanspacer{} อตฺตูปกฺกเมน ทุกฺเขน.\csromanspacer{} ปรูปกฺกเมน ทุกฺเขน.\csromanspacer{} ทุกฺขทุกฺเขน.\csromanspacer{} สงฺขารทุกฺเขน.\csromanspacer{} วิปริณามทุกฺขเน.\csromanspacer{} จกฺขุโรเคน ทุกฺเขน.\csromanspacer{} โสตโรเคน.\csromanspacer{} ฆานโรเคน.\csromanspacer{} ชิวฺหาโรเคน.\csromanspacer{} กายโรเคน. \csromanfolio{320}สีสโรเคน.\csromanspacer{} กณฺณโรเคน.\csromanspacer{} มุขโรเคน.\csromanspacer{} ทนฺตโรเคน.\csromanspacer{} กาเสน.\csromanspacer{} สาเสน.\csromanspacer{} ปินาเสน.\csromanspacer{} qอาเหน.\csromanspacer{} ชเรน.\csromanspacer{} กุจฺฉิโรเคน.\csromanspacer{} มุจฺฉาย.\csromanspacer{} ปกฺขนฺทิกาย.\csromanspacer{} สูลาย.\csromanspacer{} วิสูจิกาย.\csromanspacer{} กุฏฺเฐน.\csromanspacer{} คณฺเฑน.\csromanspacer{} กิลาเสน.\csromanspacer{} โสเสน.\csromanspacer{} อปมาเรน.\csromanspacer{} ททฺทุยา.\csromanspacer{} กณฺฑุยา.\csromanspacer{} กจฺฉุยา.\csromanspacer{} รขสาย.\csromanspacer{} วิตจฺฉิกาย.\csromanspacer{} โลหิเตน.\csromanspacer{} ปิตฺเตน.\csromanspacer{} มธุเมเหน.\csromanspacer{} อํสาย.\csromanspacer{} ปีฬกาย.\csromanspacer{} ภคนฺทลาย.\csromanspacer{} ปิตฺตสมุฏฺฐาเนน อาพาเธน.\csromanspacer{} เสมฺหสมุฏฺฐาเนน อาพาเธน.\csromanspacer{} วาตสมุฏฺฐาเนน อาพาเธน.\csromanspacer{} สนฺนิปาติเกน อาพาเธน.\csromanspacer{} อุตุปริณามเชน อาพาเธน.\csromanspacer{} วิสมปริหารเชน อาพาเธน.\csromanspacer{} โอปกฺกมิเกน อาพาเธน.\csromanspacer{} กมฺมวิปากเชน อาพาเธน.\csromanspacer{} สีเตน.\csromanspacer{} อุเณฺหน.\csromanspacer{} ชิฆจฺฉาย.\csromanspacer{} ปิปาสาย.\csromanspacer{} อุจฺจาเรน.\csromanspacer{} ปสฺสาเวน.\csromanspacer{} qอํสมกสวาตาตปสรีสปสมฺผสฺเสน ทุกฺขเน.\csromanspacer{} มาตุมรเณน ทุกฺเขน.\csromanspacer{} ปิตุมรเณน ทุกฺเขน.\csromanspacer{} ภาตุมรเณน ทุกฺเขน.\csromanspacer{} ภคินิมรเณน ทุกฺเขน.\csromanspacer{} ปุตฺตมรเณน ทุกฺเขน.\csromanspacer{} ธีตุมรเณน ทุกฺเขน.\csromanspacer{} ญาภิมรเณน ทุกฺเขน.\csromanspacer{} โภคพฺยสเนน ทุกฺเขน.\csromanspacer{} โรคพฺยสเนน ทุกฺขเน.\csromanspacer{} สีลพฺยสเนน ทุกฺเขน.\csromanspacer{} ทิฏฺฐิพฺยสเนน ทุกฺเขน ผนฺทมานํ สมฺผนฺทมานํ วิปฺผนฺทมานํ เวธมานํ ปเวธมานํ สมฺปเวธมานํ.\csromanspacer{} \textbf{ทิสฺวา}ติ \textbf{ทิสฺวา} ปสฺสิตฺวา ตุลยิตฺวา ตีรยิตฺวา วิภาวยิตฺวา วิภูตํ กตฺวาติ ผนฺทมานํ ปชํ \textbf{ทิสฺวา}.}

\prose{\textbf{มจฺเฉ อปฺโปทเก ยถา}ติ ยถา มจฺฉา อปฺโปทเก อุทกปริยาทาเน กาเกหิ วา กุลเลหิ วา พลากาหิ วา ปริปาติยมานา อุกฺขิปิยมานา ขชฺชมานา ผนฺทนฺติ สมฺผนฺทนฺติ วิปฺผนฺทนฺติ เวธนฺติ ปเวธนฺติ สมฺปเวธนฺติ, เอวเมวํ ปชา ตณฺหาผนฺทนาย ผนฺทนฺติ ฯเปฯ ทิฏฺฐิพฺยสเนน ทุกฺเขน ผนฺทนฺติ สมฺผนฺทนฺติ วิปฺผนฺทนฺติ เวธนฺติ ปเวธนฺติ สมฺปเวธนฺตีติ มจฺเฉ อปฺโปทเก ยถา.}

\prose{\textbf{อญฺญมญฺเญหิ พฺยารุทฺเธ}ติ อญฺญมญฺญํ สตฺตา วิรุทฺธา ปฏิวิรุทฺธา อาหตา ปจฺจาหตา อาฆาติตา ปจฺจาฆาติตา.\csromanspacer{} ราชาโนปิ ราชูหิ วิวทนฺติ, ขตฺติยาปิ ขตฺติเยหิ วิวทนฺติ, พฺราหฺมณาปิ พฺราหฺมเณหิ วิวทนฺติ, คหปตีปิ คหปตีหิ วิวทนฺติ, มาตาปิ ปุตฺเตน วิวทติ, ปุตฺโตปิ มาตรา วิวทติ, ปิตาปิ ปุตฺเตน วิวทติ, ปุตฺโตปิ ปิตรา วิวทติ, ภาตาปิ ภาตรา วิวทติ, ภคินีปิ ภคินิยา วิวทติ, ภาตาปิ ภคินิยา วิวทติ, ภคินิปิ ภาตรา วิวทติ, สหาโยปิ สหาเยน วิวทติ.\csromanspacer{} เต ตตฺถ กลหวิคฺคหวิวาทาปนฺนา อญฺญมญฺญํ ปาณีหิปิ อุปกฺกมนฺติ เลฑฺฑูหิปิ อุปกฺกมนฺติ, ทณฺเฑหิปิ อุปกฺกมนฺติ, สตฺเถหิปิ อุปกฺกมนฺติ, เต ตตฺถ มรณมฺปิ นิคจฺฉนฺติ มรณมตฺตมฺปิ ทุกฺขนฺติ อญฺญมญฺเญหิ พฺยารุทฺเธ.}

\csromanheader{2}{15. อตฺตทณฺฑสุตฺตนิทฺเทส}
\prose{\csromanfolio{321}\textbf{ทิสฺวา มํ ภยมาวิสี}ติ \textbf{ทิสฺวา}ติ \textbf{ทิสฺวา} ปสฺสิตฺวา ตุลยิตฺวา ตีรยิตฺวา วิภาวยิตฺวา วิภูตํ กตฺวา ภยํ ปีฬนํ ฆฏฺฏนํ อุปทฺทโว อุปสคฺโค อาวิสีติ\footnote{อาวีสีติ (สี), อาวิสตีติ (สฺยา)} \textbf{ทิสฺวา} มํ ภยมาวิสิ.\csromanspacer{} เตนาห ภควา\csromandash{}}

\csromangathameasure{ผนฺทมานํ ปชํ ทิสฺวา, มจฺเฉ อปฺโปทเก ยถา. \\ อญฺญมญฺเญหิ พฺยารุทฺเธ, ทิสฺวา มํ ภยมาวิสีติ.}
\csromangathagroup{\csromangathastanza{ผนฺทมานํ ปชํ ทิสฺวา, มจฺเฉ อปฺโปทเก ยถา. \\ อญฺญมญฺเญหิ พฺยารุทฺเธ, ทิสฺวา มํ ภยมาวิสีติ.}}

\proseitem{172}{\csromansymbolfootnote{*}{ขุ 1. 424 ปิฏฺเฐปิ.}\textbf{สมนฺตมสาโร โลโก}, \textbf{ทิสฺวา} สพฺพา สเมริตา.}

\csromangathameasure{อิจฺฉํ ภวนมตฺตโน, นาทฺทสาสิํ อโนสิตํ.}
\csromangathagroup{\csromangathastanza{อิจฺฉํ ภวนมตฺตโน, นาทฺทสาสิํ อโนสิตํ.}}

\prose{\textbf{สมนฺตมสาโร โลโก}ติ \textbf{โลโก}ติ นิรย\textbf{โลโก} ติรจฺฉานโยนิ\textbf{โลโก} เปตฺติวิสย\textbf{โลโก} มนุสฺส\textbf{โลโก} เทว\textbf{โลโก}, ขนฺธ\textbf{โลโก} ธาตุ\textbf{โลโก} อายตน\textbf{โลโก}, อยํ \textbf{โลโก} ปโร \textbf{โลโก} พฺรหฺม\textbf{โลโก} เทว\textbf{โลโก} อยํ วุจฺจติ \textbf{โลโก}.\csromanspacer{} นิรย\textbf{โลโก} อสาโร นิสฺสาโร สาราปคโต นิจฺจสารสาเรน วา สุขสารสาเรน วา อตฺตสารสาเรน วา นิจฺเจน วา ธุเวน วา สสฺสเตน วา อวิปริณามธมฺเมน วา.\csromanspacer{} ติรจฺฉานโยนิ\textbf{โลโก}.\csromanspacer{} เปตฺติวิสย\textbf{โลโก}.\csromanspacer{} มนุสฺส\textbf{โลโก}.\csromanspacer{} เทว\textbf{โลโก}.\csromanspacer{} ขนฺธ\textbf{โลโก}.\csromanspacer{} ธาตุ\textbf{โลโก}.\csromanspacer{} อายตน\textbf{โลโก}.\csromanspacer{} อยํ \textbf{โลโก}.\csromanspacer{} ปโร \textbf{โลโก}.\csromanspacer{} พฺรหฺม\textbf{โลโก}.\csromanspacer{} เทว\textbf{โลโก} อสาโร นิสฺสาโร สาราปคโต นิจฺจสารสาเรน วา สุขสารสาเรน วา อตฺตสารสาเรน วา นิจฺเจน วา ธุเวน วา สสฺสเตน วา อวิปริณามธมฺเมน วา.}

\prose{ยถา ปน นโฬ อสาโร นิสฺสาโร สาราปคโต.\csromanspacer{} ยถา เอรณฺโฑ อสาโร นิสฺสาโร สาราปคโต.\csromanspacer{} ยถา อุทุมฺพโร อสาโร นิสฺสาโร สาราปคโต.\csromanspacer{} ยถา เสตกจฺโฉ อสาโร นิสฺสาโร สาราปคโต.\csromanspacer{} ยถา ปาริภทฺทโก อสาโร นิสฺสาโร สาราปคโต.\csromanspacer{} ยถา เผณปิณฺโฑ อสาโร นิสฺสาโร สาราปคโต.\csromanspacer{} ยถา อุทกปุพฺพุฬํ\footnote{พุพฺพุลกํ (สี), ปุพฺพุฬกํ (สฺยา)} อสารํ นิสฺสารํ สาราปคตํ.\csromanspacer{} ยถา มรีจิ อสารา นิสฺสารา สาราปคตา.\csromanspacer{} ยถา กทลิกฺขนฺโธ อสาโร นิสฺสาโร สาราปคโต.\csromanspacer{} ยถา มายา อสารา นิสฺสารา สาราปคตา, เอวเมวํ นิรย\textbf{โลโก} อสาโร นิสฺสาโร สาราปคโต \csromanfolio{322}นิจฺจสารสาเรน วา สุขสารสาเรน วา อตฺตสารสาเรน วา นิจฺเจน วา ธุเวน วา สสฺสเตน วา อวิปริณามธมฺเมน วา.}

\prose{ติรจฺฉานโยนิโลโก.\csromanspacer{} เปตฺติวิสยโลโก.\csromanspacer{} มนุสฺสโลโก.\csromanspacer{} เทวโลโก อสาโร นิสฺสาโร สาราปคโต นิจฺจสารสาเรน วา สุขสารสาเรน วา อตฺตสารสาเรน วา นิจฺเจน วา ธุเวน วา สสฺสเตน วา อวิปริณามธมฺเมน วา.\csromanspacer{} ขนฺธโลโก.\csromanspacer{} ธาตุโลโก.\csromanspacer{} อายตนโลโก.\csromanspacer{} อยํ โลโก.\csromanspacer{} ปโรโลโก.\csromanspacer{} พฺรหฺมโลโก.\csromanspacer{} เทวโลโก อสาโร นิสฺสาโร สาราปคโต นิจฺจสารสาเรน วา สุขสารสาเรน วา อตฺตสารสาเรน วา นิจฺเจน วา ธุเวน วา สสฺสเตน วา อวิปริณามธมฺเมน วาติ สมนฺตมสาโร โลโก.}

\prose{\textbf{ทิสา สพฺพา สเมริตา}ติ เย ปุรตฺถิมาย ทิสาย สงฺขารา, เตปิ เอริตา สเมริตา จลิตา ฆฏฺฏิตา อนิจฺจตาย ชาติยา อนุคตา ชราย อนุสฏา พฺยาธินา อภิภูตา มรเณน อพฺภาหตา ทุกฺเข ปติฏฺฐิตา อตาณา อเลณา อสรณา อสรณีภูตา.\csromanspacer{} เย ปจฺฉิมาย ทิสาย สงฺขารา.\csromanspacer{} เย อุตฺตราย ทิสาย สงฺขารา.\csromanspacer{} เย ทกฺขิณาย ทิสาย สงฺขารา.\csromanspacer{} เย ปุรตฺถิมาย อนุทิสาย สงฺขารา.\csromanspacer{} เย ปจฺฉิมาย อนุทิสาย สงฺขารา.\csromanspacer{} เย อุตฺตราย อนุทิสาย สงฺขารา.\csromanspacer{} เย ทกฺขิณาย อนุทิสาย สงฺขารา.\csromanspacer{} เย เหฏฺฐิมาย ทิสาย สงฺขารา.\csromanspacer{} เย อุปริมาย ทิสาย สงฺขารา.\csromanspacer{} เย ทสสุ ทิสาสุ สงฺขารา, เตปิ เอริตา สเมริตา จลิตา ฆฏฺฏิตา อนิจฺจตาย ชาติยา อนุคตา ชราย อนุสฏา พฺยาธินา อภิภูตา มรเณน อพฺภาหตา ทุกฺเข ปติฏฺฐิตา อตาณา อเลณา อสรณ อสรณีภูตา.\csromanspacer{} ภาสิตมฺปิ เจตํ\csromandash{}}

\csromangathameasure{*กิญฺจาปิ เต ตํ ชลตี วิมานํ, \\ โอภาสยํ อุตฺตริยํ ทิสาย. \\ รูเป รณํ ทิสฺวา สทา ปเวธิตํ, \\ ตสฺมา น รูเป รมตี สุเมโธ. \\ +มจฺจุนาพฺภาหโต โลโก, ชราย ปริวาริโต. \\ ตณฺหาสลฺเลน โอติณฺโณ, อิจฺฉาธูมายิโต สทา.}
\csromangathagroup{\csromangathastanza{\csromansymbolfootnote{*}{สํ 1. 150 ปิฏฺเฐ.}กิญฺจาปิ เต ตํ ชลตี วิมานํ, \\ โอภาสยํ อุตฺตริยํ ทิสาย. \\ รูเป รณํ ทิสฺวา สทา ปเวธิตํ, \\ ตสฺมา น รูเป รมตี สุเมโธ. \\ \csromansymbolfootnote{+}{สํ 1. 37 ปิฏฺเฐ.}มจฺจุนาพฺภาหโต โลโก, ชราย ปริวาริโต. \\ ตณฺหาสลฺเลน โอติณฺโณ, อิจฺฉาธูมายิโต\footnote{อิจฺฉาธุมายิโก (สฺยา)} สทา.}}

\csromanheader{2}{15. อตฺตทณฺฑสุตฺตนิทฺเทส}
\csromangathameasure{*สพฺโพ อาทีปิโต โลโก, สพฺโพ โลโก ปธูปิโต. \\ สพฺโพ ปชฺชลิโต โลโก, สพฺโพ โลโก ปกมฺปิโตติ ทิสา สพฺพา สเมริตา.}
\csromangathagroup{\csromangathastanza{\csromanfolio{323}\csromansymbolfootnote{*}{สํ 1. 135 ปิฏฺเฐ.}สพฺโพ อาทีปิโต โลโก, สพฺโพ โลโก ปธูปิโต. \\ สพฺโพ ปชฺชลิโต โลโก, สพฺโพ โลโก ปกมฺปิโตติ ทิสา สพฺพา สเมริตา.}}

\prose{\textbf{อิจฺฉํ ภวนมตฺตโน}ติ อตฺตโน ภวนํ ตาณํ เลณํ สรณํ คติํ ปรายณํ อิจฺฉนฺโต สาทิยนฺโต ปตฺถยนฺโต ปิหยนฺโต อภิชปฺปนฺโตติ อิจฺฉํ ภวนมตฺตโน.\csromanspacer{} \textbf{นาทฺทสาสิํ อโนสิตนฺ}ติ อชฺโฌสิตํเยว อทฺทสํ, อนชฺโฌสิตํ นาทฺทสํ.\csromanspacer{} สพฺพํ โยพฺพญฺญํ ชราย โอสิตํ, สพฺพํ อาโรคฺยํ พฺยาธินา โอสิตํ, สพฺพํ ชีวิตํ มรเณน โอสิตํ, สพฺพํ ลาภํ อลาเภน โอสิตํ, สพฺพํ ยสํ อยเสน โอสิตํ, สพฺพํ ปสํสํ นินฺทาย โอสิตํ, สพฺพํ สุขํ ทุกฺเขน โอสิตํ.}

\csromangathameasure{+ลาโภ อลาโภ ยโส อยโส จ, \\ นินฺทา ปสํสา จ สุขํ ทุขญฺจ. \\ เอเต อนิจฺจา มนุเชสุ ธมฺมา, \\ อสสฺสตา วิปริณามธมฺมาติ.}
\csromangathagroup{\csromangathastanza{\csromansymbolfootnote{+}{ขุ 5. 107 ปิฏฺเฐ ชาตเก.}ลาโภ อลาโภ ยโส อยโส จ, \\ นินฺทา ปสํสา จ สุขํ ทุขญฺจ. \\ เอเต อนิจฺจา มนุเชสุ ธมฺมา, \\ อสสฺสตา วิปริณามธมฺมาติ.}}

\csromanheader{2}{นาทฺทสาสิํ อโนสิตํ. เตนาห ภควา\csromandash{}}
\csromangathameasure{สมนฺตมสาโร โลโก, ทิสา สพฺพา สเมริตา. \\ อิจฺฉํ ภวนมตฺตโน, นาทฺทสาสิํ อโนสิตนฺติ.}
\csromangathagroup{\csromangathastanza{สมนฺตมสาโร โลโก, ทิสา สพฺพา สเมริตา. \\ อิจฺฉํ ภวนมตฺตโน, นาทฺทสาสิํ อโนสิตนฺติ.}}

\csromanhangingbegin
\hangingheaditem{173}{\textbf{โอสาเน เตฺวว พฺยารุทฺเธ}, ทิสฺวา เม อรตี อหุ.}
\hangingbody{\textbf{อเถตฺถ สลฺลมทฺทกฺขิํ, ทุทฺทสํ หทยสฺสิตํ.}}
\csromanhangingend

\prose{\textbf{โอสาเน เตฺวว พฺยารุทฺเธ}ติ \textbf{โอสาเน เตฺววา}ติ สพฺพํ โยพฺพญฺญํ ชรา โอสาเปติ, สพฺพํ อาโรคฺยํ พฺยาธิ โอสาเปติ, สพฺพํ ชีวิตํ มรณํ โอสาเปติ, สพฺพํ ลาภํ อลาโภ โอสาเปติ, สพฺพํ ยสํ อยโส โอสาเปติ, สพฺพํ ปสํสํ นินฺทา โอสาเปติ, สพฺพํ สุขํ ทุกฺขํ โอสาเปตีติ โอสาเน เตฺวว.\csromanspacer{}\csromanspacer{}\textbf{พฺยารุทฺเธ}ติ โยพฺพญฺญกามา สตฺตา ชราย ปฏิวิรุทฺธา, อาโรคฺยกามา สตฺตา พฺยาธินา ปฏิวิรุทฺธา, ชีวิตุกามา สตฺตา มรเณน ปฏิวิรุทฺธา, ลาภกามา สตฺตา อลาเภน ปฏิวิรุทฺธา, ยสกามา สตฺตา อยเสน ปฏิวิรุทฺธา, ปสํสกามา สตฺตา นินฺทาย ปฏิวิรุทฺธา, สุขกามา สตฺตา ทุกฺเขน ปฏิวิรุทฺธา อาหตา ปจฺจาหตา อาฆาติตา ปจฺจาฆาติตาติ โอสาเน เตฺวว พฺยารุทฺเธ.}

\prose{\csromanfolio{324}\textbf{ทิสฺวา เม อรตี อหู}ติ \textbf{ทิสฺวา}ติ \textbf{ทิสฺวา} ปสฺสิตฺวา ตุลยิตฺวา ตีรยิตฺวา วิภาวยิตฺวา วิภูตํ กตฺวาติ \textbf{ทิสฺวา}.\csromanspacer{} \textbf{เม อรตี}ติ ยา อรติ ยา อนภิรติ ยา อนภิรมนา ยา อุกฺกณฺฐิตา ยา ปริตสิตา.\csromanspacer{} \textbf{อหู}ติ \textbf{ทิสฺวา} เม อรตี อหุ.}

\prose{\textbf{อเถตฺถ สลฺลมทฺทกฺขินฺ}ติ \textbf{อถา}ติ ปทสนฺธิ ฯเปฯ ปทานุปุพฺพตาเปตํ \textbf{อถา}ติ.\csromanspacer{} \textbf{เอตฺถา}ติ สตฺเตสุ.\csromanspacer{} \textbf{สลฺล}นฺติ สตฺตสลฺลานิ ราคสลฺลํ โทสสลฺลํ โมหสลฺลํ มานสลฺลํ ทิฏฺฐิสลฺลํ โสกสลฺลํ กถํกฺ\textbf{อถา}สลฺลํ.\csromanspacer{} \textbf{อทฺทกฺขินฺ}ติ อทฺทสํ อทกฺขิํ อปสฺสิํ ปฏิวิชฺฌินฺติ อเถตฺถ สลฺลมทฺทกฺขิํ.}

\prose{\textbf{ทุทฺทสํ หทยสฺสิตนฺ}ติ \textbf{ทุทฺทสนฺ}ติ ทุทฺทสํ ทุทฺทกฺขํ ทุปฺปสฺสํ ทุพฺพุชฺฌํ ทุรนุพุชฺฌํ ทุปฺปฏิวิชฺฌนฺติ ทุทฺทสํ.\csromanspacer{} \textbf{หทยสฺสิตนฺ}ติ หทยํ วุจฺจติ จิตฺตํ.\csromanspacer{} ยํ จิตฺตํ มโน มานสํ หทยํ ปณฺฑรํ มโน มนายตนํ มนินฺทฺริยํ วิญฺญาณํ วิญฺญาณกฺขนฺโธ ตชฺชา มโนวิญฺญาณธาตุ.\csromanspacer{} \textbf{หทยสฺสิตนฺ}ติ หทยนิสฺสิตํ จิตฺตสิตํ จิตฺตนิสฺสิตํ จิตฺเตน สหคตํ สหชาตํ สํสฏฺฐํ สมฺปยุตฺตํ เอกุปฺปาทํ เอกนิโรธํ เอกวตฺถุกํ เอการมฺมณนฺติ ทุทฺทสํ หทยสฺสิตํ.\csromanspacer{} เตนาห ภควา\csromandash{}}

\csromangathameasure{โอสาเน เตฺวว พฺยารุทฺเธ, ทิสฺวา เม อรตี อหุ. \\ อเถตฺถ สลฺลมทฺทกฺขิํ, ทุทฺทสํ หทยสฺสิตนฺติ.}
\csromangathagroup{\csromangathastanza{โอสาเน เตฺวว พฺยารุทฺเธ, ทิสฺวา เม อรตี อหุ. \\ อเถตฺถ สลฺลมทฺทกฺขิํ, ทุทฺทสํ หทยสฺสิตนฺติ.}}

\csromanhangingbegin
\hangingheaditem{174}{เยนสลฺเลน โอติณฺโณ, ทิสา สพฺพา วิธาวติ.}
\hangingbody{\textbf{ตเมว สลฺลมพฺพุยฺห, น ธาวติ น สีทติ.}}
\csromanhangingend

\prose{\textbf{เยน สลฺเลน โอติณฺโณ, ทิสา สพฺพา วิธาวตี}ติ \textbf{สลฺลนฺ}ติ สตฺต สลฺลานิ ราคสลฺลํ โทสสลฺลํ โมหสลฺลํ มานสลฺลํ ทิฏฺฐิสลฺลํ โสกสลฺลํ กถํกฺ\textbf{อถา}สลฺลํ.\csromanspacer{}\csromanspacer{}กตมํ ราคสลฺลํ, โย ราโค สาราโค อนุนโย อนุโรโธ นนฺทิราโค จิตฺตสฺส สาราโค ฯเปฯ อภิชฺฌา โลโภ อกุสลมูลํ, อิทํ ราคสลฺลํ.}

\prose{กตมํ โทสสลฺลํ. “อนตฺถํ เม อจรี”ติ อาฆาโต ชายติ, “อนตฺถํ เม จรตี”ติ อาฆาโต ชายติ, “อนตฺถํ เม จริสฺสตี”ติ อาฆาโต ชายติ ฯเปฯ จณฺฑิกฺกํ อสุโรโป อนตฺตมนตา จิตฺตสฺส, อิทํ โทสสลฺลํ.}

\prose{กตมํ โมหสลฺลํ, ทุกฺเข อญฺญาณํ ฯเปฯ ทุกฺขนิโรธคามินิยา ปฏิปทาย อญฺญาณํ, ปุพฺพนฺเต อญฺญาณํ, อปรนฺเต อญฺญาณํ, ปุพฺพนฺตาปรนฺเต}

\vspace{\tipitakaparskip}
\csromancenter{อตฺตทณฺฑสุตฺตนิทฺเทส อญฺญาณํ, อิทปฺปจฺจยตาปฏิจฺจสมุปฺปนฺเนสุ ธมฺเมสุ อญฺญาณํ.\csromanspacer{} ยํ เอวรูปํ อทสฺสนํ อนภิสมโย อนนุโพโธ อสมฺโพโธ อปฺปฏิเวโธ อสงฺคาหณา อปริโยคาหณา\footnote{อสงฺคาหตา อปริโยคาหตา (สฺยา), อสงฺคาหนา อปริโยคาหนา (ก)} อสมเปกฺขนา อปจฺจเวกฺขนา อปจฺจกฺขกมฺมํ ทุมฺเมชฺฌํ พาลฺยํ โมโห ปโมโห สมฺโมโห อวิชฺชา อวิชฺโชโฆ อวิชฺชาโยโค อวิชฺชานุสโย อวิชฺชาปริยุฏฺฐานํ อวิชฺชาลงฺคี โมโห อกุสลมูลํ, อิทํ โมหสลฺลํ.}
\prose{\csromanfolio{325}กตมํ มานสลฺลํ, “เสยฺโยหมสฺมี”ติ มาโน, “สทิโสหมสฺมี”ติ มาโน, “หีโนหมสฺมี”ติ มาโน, โย เอวรูโป มาโน มญฺญนา มญฺญิตตฺตํ อุนฺนติ อุนฺนโม ธโช สมฺปคฺคาโห เกตุกมฺยตา จิตฺตสฺส, อิทํ มานสลฺลํ.}

\prose{กตมํ ทิฏฺฐิสลฺลํ, วีสติวตฺถุกา สกฺกายทิฏฺฐิ, ทสวตฺถุกา มิจฺฉาทิฏฺฐิ, ทสวตฺถุกา อนฺตคฺคาหิกา ทิฏฺฐิ.\csromanspacer{} ยา เอวรูปา ทิฏฺฐิ ทิฏฺฐิคตํ ทิฏฺฐิคหนํ ทิฏฺฐิกนฺตาโร ทิฏฺฐิวิสูกายิกํ ทิฏฺฐิวิปฺผนฺทิตํ ทิฏฺฐิสญฺโญชนํ คาโห ปฏิคฺคาโห อภินิเวโส ปรามาโส กุมฺมคฺโค\footnote{กุมคฺโค (ก) 87 ปิฏฺเฐปิ.} มิจฺฉาปโถ มิจฺฉตฺตํ ติตฺถายตนํ วิปริเยสคฺคาโห วิปรีตคฺคาโห วิปลฺลาสคฺคาโห มิจฺฉาคาโห อยาถาวกสฺมิํ “ยาถาวกนฺ”ติ คาโห ยาวตา ทฺวาสฏฺฐิ ทิฏฺฐิคตานิ, อิทํ ทิฏฺฐิสลฺลํ.}

\prose{กตมํ โสกสลฺลํ, ญาติพฺยสเนน วา ผุฏฺฐสฺส โรคพฺยสเนน วา ผุฏฺฐสฺส โภคพฺยสเนน วา ผุฏฺฐสฺส สีลพฺยสเนน วา ผุฏฺฐสฺส ทิฏฺฐิพฺยสเนน วา ผุฏฺฐสฺส อญฺญตรญฺญตเรน พฺยสเนน สมนฺนาคตสฺส อญฺญตรญฺญตเรน ทุกฺขธมฺเมน ผุฏฺฐสฺส โสโก โสจนา โสจิตตฺตํ อนฺโตโสโก อนฺโตปริโสโก อนฺโตฑาโห อนฺโตปริฑาโห เจตโส ปริชฺฌายนา โทมนสฺสํ, อิทํ โสกสลฺลํ.}

\prose{กตมํ กถํกถาสลฺลํ, ทุกฺเข กงฺขา, ทุกฺขสมุทเย กงฺขา, ทุกฺขนิโรเธ กงฺขา, ทุกฺขนิโรธคามินิยา ปฏิปทาย กงฺขา, ปุพฺพนฺเต กงฺขา, อปรนฺเต กงฺขา, ปุพฺพนฺตาปรนฺเต กงฺขา, อิทปฺปจฺจยตาปฏิจฺจสมุปฺปนฺเนสุ ธมฺเมสุ กงฺขา.\csromanspacer{} ยา เอวรูปา กงฺขา กงฺขายนา กงฺขายิตตฺตํ วิมติ วิจิกิจฺฉา เทฺวฬฺหกํ เทฺวธาปโถ สํสโย อเนกํสคฺคาโห อาสปฺปนา ปริสปฺปนา อปริโยคาหณา ฉมฺภิตตฺตํ จิตฺตสฺส มโนวิเลโข, อิทํ กถํกถาสลฺลํ.}

\prose{\csromanfolio{326}\textbf{เยน สลฺเลน โอติณฺโณ, ทิสา สพฺพา วิธาวตี}ติ ราคสลฺเลน โอติณฺโณ วิทฺโธ ผุฏฺโฐ ปเรโต สโมหิโต สมนฺนาคโต กาเยน ทุจฺจริตํ จรติ, วาจาย ทุจฺจริตํ จรติ, มนสา ทุจฺจริตํ จรติ, ปาณมฺปิ หนติ, อทินฺนมฺปิ อาทิยติ, สนฺธิมฺปิ ฉินฺทติ, นิลฺโลปมฺปิ หรติ, เอกาคาริกมฺปิ กโรติ, ปริปนฺเถปิ ติฏฺฐติ, ปรทารมฺปิ คจฺฉติ, มุสาปิ ภณติ, เอวมฺปิ ราคสลฺเลน โอติณฺโณ วิทฺโธ ผุฏฺโฐ ปเรโต สโมหิโต สมนฺนาคโต ธาวติ วิธาวติ สนฺธาวติ สํสรติ.\csromanspacer{}\csromanspacer{}อถ วา ราคสลฺเลน โอติณฺโณ วิทฺโธ ผุฏฺโฐ ปเรโต สโมหิโต สมนฺนาคโต โภเค ปริเยสนฺโต นาวาย มหาสมุทฺทํ ปกฺขนฺทติ.\csromanspacer{} สีตสฺส ปุรกฺขโต อุณฺหสฺส ปุรกฺขโต ฑํสมกส วาตาตป สรีสป สมฺผสฺเสหิ ปีฬิยมาโน\footnote{ริสฺสมาโน (สี, สฺยา) 116 ปิฏฺเฐ-ปาฐเภโท นตฺถิ.} ขุปฺปิปาสาย มิยฺยมาโน ติคุมฺพํ คจฺฉติ, ตกฺโกลํ คจฺฉติ, ตกฺกสีลํ คจฺฉติ, กาฬมุขํ คจฺฉติ, ปุรปูรํ คจฺฉติ, เวสุงฺคํ คจฺฉติ, เวราปถํ คจฺฉติ, ชวํ คจฺฉติ, ตามลิํ คจฺฉติ, วงฺกํ คจฺฉติ, เอฬพนฺธนํ คจฺฉติ, สุวณฺณกูฏํ คจฺฉติ, สุวณฺณภูมิํ คจฺฉติ, ตมฺพปณฺณิํ คจฺฉติ, สุปฺปารกํ\footnote{สุปฺปาทกํ (ก), สุปฺปารํ (สฺยา)} คจฺฉติ, ภารุกจฺฉํ คจฺฉติ, สุรฏฺฐํ คจฺฉติ, ภงฺคโลกํ\footnote{องฺคโลกํ (สี), องฺคเณกํ (สฺยา)} คจฺฉติ, ภงฺคณํ\footnote{คงฺคณํ (สี, สฺยา) เหฏฺฐา 120 ปิฏฺเฐปิ.} คจฺฉติ, ปรมภงฺคณํ คจฺฉติ, โยนํ คจฺฉติ, ปรมโยนํ คจฺฉติ, วินกํ คจฺฉติ, มูลปทํ คจฺฉติ, มรุกนฺตารํ คจฺฉติ, ชณฺณุปถํ คจฺฉติ, อชปถํ คจฺฉติ, เมณฺฑปถํ คจฺฉติ, สงฺกุปถํ คจฺฉติ, ฉตฺตปถํ คจฺฉติ, วํสปถํ คจฺฉติ, สกุณปถํ คจฺฉติ, มูสิกปถํ คจฺฉติ, ทริปถํ คจฺฉติ, เวตฺตาจารํ คจฺฉติ, ปริเยสนฺโต น ลภติ, อลาภมูลกมฺปิ ทุกฺขํ โทมนสฺสํ ปฏิสํเวเทติ, ปริเยสนฺโต ลภติ, ลทฺธา อารกฺขมูลกมฺปิ ทุกฺขํ โทมนสฺสํ ปฏิสํเวเทติ “กินฺติ เม โภเค เนว ราชาโน หเรยฺยุํ, น โจรา หเรยฺยุํ, น อคฺคิ ทเหยฺย, น อุทกํ วเหยฺย, น อปฺปิยา ทายาทา หเรยฺยุนฺ”ติ, ตสฺส เอวํ อารกฺขโต โคปยโต เต โภคา วิปฺปลุชฺชนฺติ, โส วิปฺปโยคมูลกมฺปิ ทุกฺขํ โทมนสฺสํ ปฏิสํเวเทติ.\csromanspacer{} เอวมฺปิ ราคสลฺเลน โอติณฺโณ วิทฺโธ ผุฏฺโฐ ปเรโต สโมหิโต สมนฺนาคโต ธาวติ วิธาวติ สนฺธาวติ สํสรติ.}

\prose{โทสสลฺเลน.\csromanspacer{} โมหสลฺเลน.\csromanspacer{} มานสลฺเลน โอติณฺโณ วิทฺโธ ผุฏฺโฐ ปเรโต สโมหิโต สมนฺนาคโต กาเยน ทุจฺจริตํ จรติ,}

\vspace{\tipitakaparskip}
\csromancenter{อตฺตทณฺฑสุตฺตนิทฺเทส วาจาย ทุจฺจริตํ จรติ, มนสา ทุจฺจริตํ จรติ, ปาณมฺปิ หนติ, อทินฺนมฺปิ อาทิยติ, สนฺธิมฺปิ ฉินฺทติ, นิลฺโลปมฺปิ หรติ, เอกาคาริกมฺปิ กโรติ, ปริปนฺเถปิ ติฏฺฐติ, ปรทารมฺปิ คจฺฉติ, มุสาปิ ภณติ.\csromanspacer{} เอวํ มานสลฺเลน โอติณฺโณ วิทฺโธ ผุฏฺโฐ ปเรโต สโมหิโต สมนฺนาคโต ธาวติ วิธาวติ สนฺธาวติ สํสรติ.}
\prose{\csromanfolio{327}ทิฏฺฐิสลฺเลน โอติณฺโณ วิทฺโธ ผุฏฺโฐ ปเรโต สโมหิโต สมนฺนาคโต อเจลโก โหติ มุตฺตาจาโร หตฺถาปเลขโน\footnote{หตฺถาวเลขโน (สฺยา) ที 1. 157; ม 1. 111; ม 2. 5; อํ 1. 300 ปิฏฺเฐสุปิ.}, น เอหิภทนฺติโก, น ติฏฺฐภทนฺติโก, นาภิหฏํ, น อุทฺทิสฺสกตํ, น นิมนฺตนํ สาทิยติ, โส น กุมฺภิมุขา ปฏิคฺคณฺหาติ, น กโฬปิมุขา\footnote{ขโฬปิมุขา (สี)} ปฏิคฺคณฺหาติ, น เอฬกมนฺตรํ, น ทณฺฑมนฺตรํ, น มุสลมนฺตรํ, น ทฺวินฺนํ ภุญฺชมานานํ, น คพฺภินิยา, น ปายมานาย, น ปุริสนฺตรคตาย, น สํกิตฺตีสุ, น ยตฺถ สา อุปฏฺฐิโต โหติ, น ยตฺถ มกฺขิกา สณฺฑสณฺฑจารินี, น มจฺฉํ, น มํสํ, น สุรํ, น เมรยํ, น ถุโสทกํ ปิวติ, โส เอกาคาริโก วา โหติ เอกาโลปิโก, ทฺวาคาริโก วา โหติ ทฺวาโลปิโก ฯเปฯ สตฺตาคาริโก วา โหติ สตฺตาโลปิโก, เอกิสฺสาปิ ทตฺติยา ยาเปติ, ทฺวีหิปิ ทตฺตีหิ ยาเปติ ฯเปฯ สตฺตหิปิ ทตฺตีหิ ยาเปติ, เอกาหิกมฺปิ อาหารํ อาหาเรติ, ทฺวีหิกมฺปิ อาหารํ อาหาเรติ ฯเปฯ สตฺตาหิกมฺปิ อาหารํ อาหเรติ.\csromanspacer{} อิติ เอวรูปํ อฑฺฒมาสิกมฺปิ ปริยายภตฺตโภชนานุโยคมนุยุตฺโต วิหรติ.\csromanspacer{} เอวมฺปิ ทิฏฺฐิสลฺเลน โอติณฺโณ วิทฺโธ ผุฏฺโฐ ปเรโต สโมหิโต สมนฺนาคโต ธาวติ วิธาวติ สนฺธาวติ สํสรติ.}

\prose{อถ วา ทิฏฺฐิสลฺเลน โอติณฺโณ วิทฺโธ ผุฏฺโฐ ปเรโต สโมหิโต สมนฺนาคโต, โส สากภกฺโข วา โหติ, สามากภกฺโข วา โหติ, นีวารภกฺโข วา โหติ, ททฺทุลภกฺโข วา โหติ, หฏภกฺโข วา โหติ, กณภกฺโข วา โหติ, อาจามภกฺโข วา โหติ, ปิญฺญากภกฺโข วา โหติ, ติลภกฺโข วา โหติ, ติณภกฺโข วา โหติ, โคมยภกฺโข วา โหติ, วนมูลผลาหาโร ยาเปติ ปวตฺตผลโภชโน.\csromanspacer{} โส สาณานิปิ ธาเรติ, มสาณานิปิ ธาเรติ, ฉวทุสฺสานิปิ ธาเรติ, ปํสุกูลานิปิ ธาเรติ, ติรีฏานิปิ ธาเรติ, อชินมฺปิ ธาเรติ, อชินกฺขิปมฺปิ ธาเรติ, กุสจีรมฺปิ ธาเรติ, \csromanfolio{328}วากจีรมฺปิ ธาเรติ, ผลกจีรมฺปิ ธาเรติ, เกสกมฺพลมฺปิ ธาเรติ, อุลูกปกฺขมฺปิ ธาเรติ, เกสมสฺสุโลจโกปิ โหติ, เกสมสฺสุโลจนานุโยคมนุยุตฺโต วิหรติ, อุพฺภฏฺฐโกปิ โหติ อาสนปฏิกฺขิตฺโต, อุกฺกุฏิโกปิ โหติ อุกฺกุฏิกปฺปธานมนุยุตฺโต, กณฺฏกาปสฺสยิโก โหติ, กณฺฏกาปสฺสเย เสยฺยํ กปฺเปติ, ผลกเสยฺยมฺปิ กปฺเปติ, ถณฺฑิลเสยฺยมฺปิ กปฺเปติ, เอกาปสฺสยิโก โหติ รโชชลฺลธโร, อพฺโภกาสิโกปิ โหติ ยถาสนฺถติโก, เวกฏิโกปิ โหติ วิกฏโภชนานุโยคมนุยุตฺโต, อปานโกปิ โหติ อปานกตฺตมนุยุตฺโต, สายตติยกมฺปิ อุทโกโรหนานุโยคมนุยุตฺโต วิหรติ.\csromanspacer{} อิติ เอวรูปํ อเนกวิหิตํ กายสฺส อาตาปนปริตาปนานุโยคมนุยุตฺโต วิหรติ.\csromanspacer{} เอวมฺปิ ทิฏฺฐิสลฺเลน โอติณฺโณ วิทฺโธ ผุฏฺโฐ ปเรโต สโมหิโต สมนฺนาคโต ธาวติ วิธาวติ สนฺธาวติ สํสรติ.}

\prose{โสกสลฺเลน โอติณฺโณ วิทฺโธ ผุฏฺโฐ ปเรโต สโมหิโต สมนฺนาคโต โสจติ กิลมติ ปริเทวติ อุรตฺตาฬิํ กนฺทติ สมฺโมหํ อาปชฺชติ.\csromanspacer{} วุตฺตํ เหตํ ภควตา\csromandash{} * ภูตปุพฺพํ พฺราหฺมณ อิมิสฺสาเยว สาวตฺถิยา อญฺญตริสฺสา อิตฺถิยา มาตา กาลมกาสิ, สา ตสฺสา กาลกิริยาย อุมฺมตฺติกา ขิตฺตจิตฺตา รถิยาย รถิยํ, สิงฺฆาฏเกน สิงฺฆาฏกํ อุปสงฺกมิตฺวา เอวมาห “อปิ เม มาตรํ อทฺทสฺสถ, อปิ เม มาตรํ อทฺทสฺสถา”ติ.}

\prose{ภูตปุพฺพํ พฺราหฺมณ อิมิสฺสาเยว สาวตฺถิยา อญฺญตริสฺสา อิตฺถิยา ปิตา กาลมกาสิ.\csromanspacer{} ภาตา กาลมกาสิ.\csromanspacer{} ภคินี กาลมกาสิ.\csromanspacer{} ปุตฺโต กาลมกาสิ.\csromanspacer{} ธีตา กาลมกาสิ.\csromanspacer{} สามิโก กาลมกาสิ, สา ตสฺส กาลกิริยาย อุมฺมตฺติกา ขิตฺตจิตฺตา รถิยาย รถิยํ, สิงฺฆาฏเกน สิงฺฆาฏกํ อุปสงฺกมิตฺวา เอวมาห “อปิ เม สามิกํ อทฺทสฺสถ, อปิ เม สามิกํ อทฺทสฺสถา”ติ.}

\prose{ภูตปุพฺพํ พฺราหฺมณ อิมิสฺสาเยว สาวตฺถิยา อญฺญตรสฺส ปุริสสฺส มาตา กาลมกาสิ, โส ตสฺสา กาลกิริยาย อุมฺมตฺตโก ขิตฺตจิตฺโต รถิยาย รถิยํ, สิงฺฆาฏเกน สิงฺฆาฏกํ อุปสงฺกมิตฺวา เอวมาห “อปิ เม มาตรํ อทฺทสฺสถ, อปิ เม มาตรํ อทฺทสฺสถา”ติ.}

\csromanheader{2}{15. อตฺตทณฺฑสุตฺตนิทฺเทส}
\prose{\csromanfolio{329}ภูตปุพฺพํ พฺราหฺมณ อิมิสฺสาเยว สาวตฺถิยา อญฺญตรสฺส ปุริสสฺส ปิตา กาลมกาสิ.\csromanspacer{} ภาตา กาลมกาสิ.\csromanspacer{} ภคินี กาลมกาสิ.\csromanspacer{} ปุตฺโต กาลมกาสิ.\csromanspacer{} ธีตา กาลมกาสิ.\csromanspacer{} ปชาปติ กาลมกาสิ, โส ตสฺสา กาลกิริยาย อุมฺมตฺตโก ขิตฺตจิตฺโต รถิยาย รถิยํ สิงฺฆาถาฏเกน สิงฺฆาฏกํ อุปสงฺกมิตฺวา เอวมาห “อปิ เม ปชาปติํ อทฺทสฺสถ, อปิ เม ปชาปติํ อทฺทสฺสถา”ติ.}

\prose{ภูตปุพฺพํ พฺราหฺมณ อิมิสฺสาเยว สาวตฺถิยา อญฺญตรา อิตฺถี ญาติกุลํ อคมาสิ “ตสฺสา เต ญาตกา สามิกํ อจฺฉินฺทิตฺวา อญฺญสฺส ทาตุกามา, สา จ ตํ น อิจฺฉติ.\csromanspacer{} อถ โข สา อิตฺถี สามิกํ เอตทโวจ “อิเม มํ อยฺยปุตฺต ญาตกา ตว อจฺฉินฺทิตฺวา อญฺญสฺส ทาตุกามา, อุโภ มยํ มริสฺสามา”ติ.\csromanspacer{} อถ โข โส ปุริโส ตํ อิตฺถิํ ทฺวิธา เฉตฺวา อตฺตานํ โอปาเตสิ “อุโภ เปจฺจ ภวิสฺสามา”ติ.\csromanspacer{} เอวํ โสกสลฺเลน โอติณฺโณ วิทฺโธ ผุฏฺโฐ ปเรโต สโมหิโต สมนฺนาคโต ธาวติ วิธาวติ สนฺธาวติ สํสรติ.}

\prose{กถํกถาสลฺเลน โอติณฺโณ วิทฺโธ ผุฏฺโฐ ปเรโต สโมหิโต สมนฺนาคโต สํสยปกฺขนฺโท โหติ วิมติปกฺขนฺโท เทฺวฬฺหกชาโต “อโหสิํ นุ โข อหํ อตีตมทฺธานํ, นนุ โข อโหสิํ อตีตมทฺธานํ, กิํ นุ โข อโหสิํ อตีตมทฺธานํ, กถํ นุ โข อโหสิํ อตีตมทฺธานํ, กิํ หุตฺวา กิํ อโหสิํ นุ โข อตีตมทฺธานํ.\csromanspacer{} ภวิสฺสามิ นุ โข อหํ อนาคตมทฺธานํ, นนุ โข ภวิสฺสามิ อนาคตมทฺธานํ, กิํ นุ โข ภวิสฺสามิ อนาคตมทฺธานํ, กถํ นุ โข ภวิสฺสามิ อนาคตมทฺธานํ, กิํ หุตฺวา กิํ ภวิสฺสามิ นุ โข อนาคตมทฺธานํ.\csromanspacer{} เอตรหิ วา ปจฺจุปฺปนฺนํ อทฺธานํ อชฺฌตฺตํ กถํกถี โหติ, อหํ นุ โขสฺมิ, โน นุ โขสฺมิ, กิํ นุ โขสฺมิ, กถํ นุ โขสฺมิ, อยํ นุ โข สตฺโต กุโต อาคโต, โส กุหิํ คามี ภวิสฺสตี”ติ.\csromanspacer{} เอวํ กถํกถาสลฺเลน โอติณฺโณ วิทฺโธ ผุฏฺโฐ ปเรโต สโมหิโต สมนฺนาคโต ธาวติ วิธาวติ สนฺธาวติ สํสรติ.}

\prose{เต สลฺเล อภิสงฺขโรติ, เต สลฺเล อภิสงฺขโรนฺโต สลฺลาภิสงฺขารวเสน ปุรตฺถิมํ ทิสํ ธาวติ, ปจฺฉิมํ ทิสํ ธาวติ, อุตฺตรํ ทิสํ ธาวติ, ทกฺขิณํ ทิสํ ธาวติ.\csromanspacer{} เต สลฺลาภิสงฺขารา อปฺปหีนา, สลฺลาภิสงฺขารานํ อปฺปหีนตฺตา คติยา ธาวติ, นิรเย ธาวติ, ติรจฺฉานโยนิยา ธาวติ, เปตฺติวิสเย ธาวติ, มนุสฺสโลเก ธาวติ, เทวโลเก ธาวติ,}

\prose{\csromanfolio{330}คติยา คติํ, อุปปติตฺยา อุปปตฺติํ, ปฏิสนฺธิยา ปฏิสนฺธิํ, ภเวน ภวํ, สํสาเรน สํสารํ, วฏฺเฏน วฏฺฏํ ธาวติ วิธาวติ สนฺธาวติ สํสรตีติ เยน สลฺเลน โอติณฺโณ, ทิสา สพฺพา วิธาวติ.}

\prose{\textbf{ตเมว สลฺลมพฺพุยฺห, น ธาวติ น สีทตี}ติ ตเมว ราคสลฺลํ โทสสลฺลํ โมหสลฺลํ มานสลฺลํ ทิฏฺฐิสลฺลํ โสกสลฺลํ กถํกถาสลฺลํ อพฺพุยฺห อพฺพุหิตฺวา อุทฺธริตฺวา สมุทฺธริตฺวา อุปฺปาฏยิตฺวา สมุปฺปาฏยิตฺวา\footnote{อุปฺปาทยิตฺวา สมุปฺปาทยิตฺวา (สฺยา, ก)} ปชหิตฺวา วิโนเทตฺวา พฺยนฺติํ กริตฺวา อนภาวํ คเมตฺวา เนว ปุรตฺถิมํ ทิสํ ธาวติ, น ปจฺฉิมํ ทิสํ ธาวติ, น อุตฺตรํ ทิสํ ธาวติ, น ทกฺขิณํ ทิสํ ธาวติ, เต สลฺลาภิสงฺขารา ปหีนา, สลฺลาภิสงฺขารานํ ปหีนตฺตา คติยา น ธาวติ, นิรเย น ธาวติ, ติรจฺฉานโยนิยา น ธาวติ, เปตฺติวิสเย น ธาวติ, มนุสฺสโลเก น ธาวติ, เทวโลเก น ธาวติ, น คติยา คติํ, น อุปปตฺติยา อุปปตฺติํ, น ปฏิสนฺธิยา ปฏิสนฺธิํ, น ภเวน ภวํ, น สํสาเรน สํสารํ, น วฏฺเฏน วฏฺฏํ ธาวติ วิธาวติ สนฺธาวติ สํสรตีติ ตเมว สลฺลมพฺพุยฺห น ธาวติ.\csromanspacer{}\csromanspacer{}\textbf{น สีทตี}ติ กาโมเฆ น สีทติ, ภโวเฆ น สีทติ, ทิฏฺโฐเฆ น สีทติ, อวิชฺโชเฆ น สีทติ น สํสีทติ น โอสีทติ น อวสีทติ น คจฺฉติ น อวคจฺฉตีติ ตเมว สลฺลมพฺพุยฺห, น ธาวติ น สีทติ.\csromanspacer{} เตนาห ภควา\csromandash{}}

\csromangathameasure{เยน สลฺเลน โอติณฺโณ, ทิสา สพฺพา วิธาวติ.}
\csromangathagroup{\csromangathastanza{เยน สลฺเลน โอติณฺโณ, ทิสา สพฺพา วิธาวติ.}}

\prose{\textbf{ตเมว สลฺลมพฺพุยฺห, น ธาวติ น สีทตี}ติ.}

\proseitem{175}{\csromansymbolfootnote{*}{ขุ 1. 424 ปิฏฺเฐปิ.}ตตฺถ \textbf{สิกฺขา}นุคียนฺติ,}

\prose{\textbf{ยานิ โลเก คธิตานิ.}}

\csromangathameasure{\textbf{น เตสุ ปสุโต สิยา,} \\ \textbf{นิพฺพิชฺฌ สพฺพโส กาเม.} \\ \textbf{สิกฺเข นิพฺพานมตฺตโน.}}
\csromangathagroup{\csromangathastanza{\textbf{น เตสุ ปสุโต สิยา,} \\ \textbf{นิพฺพิชฺฌ สพฺพโส กาเม.} \\ \textbf{สิกฺเข นิพฺพานมตฺตโน.}}}

\prose{\textbf{ตตฺถ สิกฺขานุคียนฺติ, ยานิ โลเก คธิตานี}ติ \textbf{สิกฺขา}ติ หตฺถิ\textbf{สิกฺขา} อสฺส\textbf{สิกฺขา} รถ\textbf{สิกฺขา} ธนุ\textbf{สิกฺขา} สาลากิยํ สลฺลกตฺติยํ กายติกิจฺฉํ ภูติยํ โกมารภจฺจํ\footnote{โกมารติกิจฺฉํ (สฺยา), โกมารสจฺจํ (ก)}.\csromanspacer{} อนุคียนฺตีติ คียนฺติ นิคฺคียนฺติ กถียนฺติ ภณียนฺติ ทีปียนฺติ โวหรียนฺติ.\csromanspacer{} อถ วา คียนฺติ คณฺหียนฺติ}

\vspace{\tipitakaparskip}
\csromancenter{อตฺตทณฺฑสุตฺตนิทฺเทส อุคฺคณฺหียนฺติ ธารียนฺติ อุปธารียนฺติ อุปลกฺขียนฺติ คธิตปฏิลาภาย.\csromanspacer{} \textbf{คธิตา} วุจฺจนฺติ ปญฺจ กามคุณา จกฺขุวิญฺเญยฺยา รูปา อิฏฺฐา กนฺตา มนาปา ปิยรูปา กามูปสญฺหิตา รชนียา.\csromanspacer{} โสตวิญฺเญยฺยา สทฺทา ฯเปฯ.\csromanspacer{} ฆานวิญฺเญยฺยา คนฺธา.\csromanspacer{} ชิวฺหาวิญฺเญยฺยา รสา.\csromanspacer{} กายวิญฺเญยฺยา โผฏฺฐพฺพา อิฏฺฐา กนฺตา มนาปา ปิยรูปา กามูปสญฺหิตา รชนียา.\csromanspacer{} กิํการณา คธิตา วุจฺจนฺติ ปญฺจ กามคุณา, เยภุยฺเยน เทวมนุสฺสา ปญฺจ กามคุเณ อิจฺฉนฺติ สาทิยนฺติ ปตฺถยนฺติ ปิหยนฺติ อภิชปฺปนฺติ, ตํการณา คธิตา วุจฺจนฺติ ปญฺจ กามคุณา.\csromanspacer{} \textbf{โลโก}ติ มนุสฺสโลเกติ ตตฺถ \textbf{สิกฺขา}นุคียนฺติ, ยานิ โลเก คธิตานิ.}
\prose{\csromanfolio{331}\textbf{น เตสุ ปสุโต สิยา}ติ ตาสุ วา \textbf{สิกฺขา}สุ, เตสุ วา ปญฺจสุ กามคุเณสุ น ปสุโต สิยา, น ตนฺนินฺโน อสฺส, น ตปฺโปโณ, น ตปฺปพฺภาโร, น ตทธิมุตฺโต, น ตทธิปเตยฺโยติ น เตสุ ปสุโต สิยา.}

\prose{\textbf{นิพฺพิชฺฌ สพฺพโส กาเม}ติ \textbf{นิพฺพิชฺฌา}ติ ปฏิวิชฺฌิตฺวา. “สพฺเพ สงฺขารา อนิจฺจา”ติ ปฏิวิชฺฌิตฺวา, “สพฺเพ สงฺขารา ทุกฺขา”ติ ปฏิวิชฺฌิตฺวา ฯเปฯ “ยํ กิญฺจิ สมุทยธมฺมํ, สพฺพํ ตํ นิโรธธมฺมนฺ”ติ ปฏิวิชฺฌิตฺวา.\csromanspacer{}\csromanspacer{}\textbf{สพฺพโส}ติ สพฺเพน สพฺพํ สพฺพถา สพฺพํ อเสสํ นิสฺเสสํ ปริยาทิยนวจนเมตํ สพฺพโสติ.\csromanspacer{} \textbf{กามา}ติ อุทฺทานโต เทฺว กามา วตฺถุกามา จ กิเลสกามา จ ฯเปฯ อิเม วุจฺจนฺติ วตฺถุกามา ฯเปฯ อิเม วุจฺจนฺติ กิเลสกามาติ นิพฺพิชฺฌ สพฺพโส กาเม.}

\prose{\textbf{สิกฺเข นิพฺพานมตฺตโน}ติ \textbf{สิกฺขา}ติ ติสฺโส \textbf{สิกฺขา} อธิสีล\textbf{สิกฺขา} อธิจิตฺต\textbf{สิกฺขา} อธิปญฺญา\textbf{สิกฺขา} ฯเปฯ อยํ อธิปญฺญา\textbf{สิกฺขา}.\csromanspacer{}\csromanspacer{}\textbf{นิพฺพานมตฺตโน}ติ อตฺตาโน ราคสฺส นิพฺพาปนาย โทสสฺส นิพฺพาปนาย โมหสฺส นิพฺพาปนาย ฯเปฯ สพฺพากุสลาภิสงฺขารานํ สมาย อุปสมาย วูปสมาย นิพฺพาปนาย ปฏินิสฺสคฺคาย ปฏิปสฺสทฺธิยา อธิสีลมฺปิ สิกฺเขยฺย อธิจิตฺตมฺปิ สิกฺเขยฺย อธิปญฺญมฺปิ สิกฺเขยฺย อิมา ติสฺโส \textbf{สิกฺขา}โย อาวชฺชนฺโต สิกฺเขยฺย ชานนฺโต สิกฺเขยฺย ฯเปฯ สจฺฉิกาตพฺพํ สจฺฉิกโรนฺโต สิกฺเขยฺย อาจเรยฺย สมาจเรยฺย สมาทาย วตฺเตยฺยาติ สิกฺเข นิพฺพานมตฺตโน.\csromanspacer{} เตนาห ภควา\csromandash{}}

\csromangathameasure{ตตฺถ สิกฺขานุคียนฺติ, \\ ยานิ โลเก คธิตานิ. \\ น เตสุ ปสุโต \textbf{สิยา}, \\ นิพฺพิชฺฌ สพฺพโส กาเม. \\ สิกฺเข นิพฺพานมตฺตโนติ.}
\csromangathagroup{\csromangathastanza{\csromanfolio{332}ตตฺถ สิกฺขานุคียนฺติ, \\ ยานิ โลเก คธิตานิ. \\ น เตสุ ปสุโต \textbf{สิยา}, \\ นิพฺพิชฺฌ สพฺพโส กาเม.}\csromangathastanza{สิกฺเข นิพฺพานมตฺตโนติ.}}

\proseitem{176}{\csromansymbolfootnote{*}{ขุ 1. 424 ปิฏฺเฐปิ.}\textbf{สจฺโจ สิยา อปฺปคพฺโภ}, อมาโย ริตฺตเปสุโณ.}

\csromangathameasure{อกฺโกธโน โลภปาปํ, เววิจฺฉํ วิตเร มุนิ.}
\csromangathagroup{\csromangathastanza{อกฺโกธโน โลภปาปํ, เววิจฺฉํ วิตเร มุนิ.}}

\prose{\textbf{สจฺโจ สิยา อปฺปคพฺโภ}ติ \textbf{สจฺโจ} \textbf{สิยา}ติ สจฺจวาจาย สมนฺนาคโต \textbf{สิยา}, สมฺมาทิฏฺฐิยา สมนฺนาคโต \textbf{สิยา}, อริเยน อฏฺฐงฺคิเกน มคฺเคน สมนฺนาคโต \textbf{สิยา}ติ \textbf{สจฺโจ} \textbf{สิยา}.\csromanspacer{}\csromanspacer{}\textbf{อปฺปคพฺโภ}ติ ตีณิ ปาคพฺภิยานิ กายิกํ ปาคพฺภิยํ วาจสิกํ ปาคพฺภิยํ เจตสิกํ ปาคพฺภิยํ ฯเปฯ.\csromanspacer{} อิทํ เจตสิกํ ปาคพฺภิยํ.\csromanspacer{} ยสฺสิมานิ ตีณิ ปาคพฺภิยานิ ปหีนานิ สมุจฺฉินฺนานิ วูปสนฺตานิ ปฏิปสฺสทฺธานิ อภพฺพุปฺปตฺติกานิ ญาณคฺคินา ทฑฺฒานิ, โส วุจฺจติ อปฺปคพฺโภติ \textbf{สจฺโจ} \textbf{สิยา} อปฺปคพฺโภ.}

\prose{\textbf{อมาโย ริตฺตเปสุโณ}ติ มายา วุจฺจติ วญฺจนิกา จริยา.\csromanspacer{} อิเธกจฺโจ กาเยน ทุจฺจริตํ จริตฺวา วาจาย ทุจฺจริตํ จริตฺวา มนสา ทุจฺจริตํ จริตฺวา ตสฺส ปฏิจฺฉาทนเหตุ ปาปิกํ อิจฺฉํ ปณิทหติ “มา มํ ชญฺญา”ติ อิจฺฉติ, “มา มํ ชญฺญา”ติ สงฺกปฺเปติ, “มา มํ ชญฺญา”ติ วาจํ ภาสติ, “มา มํ ชญฺญา”ติ กาเยน ปรกฺกมติ.\csromanspacer{} ยา เอวรูปา มายา มายาวิตา อจฺจสรา วญฺจนา นิกติ นิกิรณา ปริหรณา คูหนา ปริคูหนา ฉาทนา ปริจฺฉาทนา อนุตฺตานีกมฺมํ อนาวิกมฺมํ โวจฺฉาทนา ปาปกิริยา, อยํ วุจฺจติ มายา.\csromanspacer{} ยสฺเสสา มายา ปหีนา สมุจฺฉินฺนา วูปสนฺตา ปฏิปสฺสทฺธา อภพฺพุปฺปตฺติกา ญาณคฺคินา ทฑฺฒา, โส วุจฺจติ อมาโย.\csromanspacer{}\csromanspacer{}\textbf{ริตฺตเปสุโณ}ติ เปสุญฺญนฺติ อิเธกจฺโจ ปิสุณวาโจ โหติ ฯเปฯ เอวํ เภทาธิปฺปาเยน เปสุญฺญํ อุปสํหรติ, ยสฺเสตํ เปสุญฺญํ ปหีนํ สมุจฺฉินฺนํ วูปสนฺตํ ปฏิปสฺสทฺธํ อภพฺพุปฺปตฺติกํ ญาณคฺคินา ทฑฺฒํ, โส วุจฺจติ ริตฺตเปสุโณ วิวิตฺตเปสุโณ ปวิวิตฺตเปสุโณติ อมาโย ริตฺตเปสุโณ.}

\vspace{\tipitakaparskip}
\csromancenter{\textbf{อกฺโกธโน โลภปาปํ, เววิจฺฉํ วิตเร มุนี}ติ \textbf{อกฺโกธโน}ติ หิ โข วุตฺตํ, อปิ จ โกโธ ตาว วตฺตพฺโพ.\csromanspacer{} ทสหากาเรหิ}
\csromancenter{อตฺตทณฺฑสุตฺตนิทฺเทส โกโธ ชายติ “อนตฺถํ เม อจรี”ติ โกโธ ชายติ ฯเปฯ ยสฺเสโส โกโธ ปหีโน สมุจฺฉินฺโน วูปสนฺโต ปฏิปสฺสทฺโธ อภพฺพุปฺปตฺติโก ญาณคฺคินา ทฑฺโฒ, โส วุจฺจติ อกฺโกธโน.\csromanspacer{} โกธสฺส ปหีนตฺตา อกฺโกธโน โกธวตฺถุสฺส ปริญฺญาตตฺตา อกฺโกธโน โกธเหตุสฺส อุปจฺฉินฺนตฺตา อกฺโกธโน.\csromanspacer{} \textbf{โลโภ}ติ โย โลโภ ลุพฺภนา ลุพฺภิตตฺตํ ฯเปฯ อภิชฺฌา โลโภ อกุสลมูลํ.\csromanspacer{} \textbf{เววิจฺฉํ} วุจฺจติ ปญฺจ มจฺฉริยานิ อาวาสมจฺฉริยํ ฯเปฯ คาโห วุจฺจติ มจฺฉริยํ.\csromanspacer{} \textbf{มุนี}ติ โมนํ วุจฺจติ ญาณํ ฯเปฯ สงฺคชาลมติจฺจ โส มุนิ.\csromanspacer{}\csromanspacer{}\textbf{อกฺโกธโน โลภปาปํ, เววิจฺฉํ วิตเร มุนี}ติ มุนิ โลภปาปญฺจ เววิจฺฉญฺจ อตริ อุตฺตริ ปตริ สมติกฺกมิ วีติวตฺตยีติ\footnote{วีติว ตฺตีติ (ก)} อกฺโกธโน โลภปาปํ, เววิจฺฉํ วิตเร มุนิ.\csromanspacer{} เตนาห ภควา\csromandash{}}
\vspace{0.5\baselineskip}
\csromangathameasure{สจฺโจ สิยา อปฺปมคพฺโภ, อมาโย ริตฺตเปสุโณ.}
\csromangathagroup{\csromangathastanza{\csromanfolio{333}สจฺโจ สิยา อปฺปมคพฺโภ, อมาโย ริตฺตเปสุโณ.}}

\prosecont{\textbf{อกฺโกธโน โลภปาปํ, เววิจฺฉํ วิตเร มุนี}ติ.}

\proseitem{177}{\csromansymbolfootnote{*}{ขุ 1. 424 ปิฏฺเฐปิ.}นิทฺทํ ตนฺทิํ สเห ถีนํ, ปมาเทน น สํวเส.}

\csromangathameasure{อติมาเน น ติฏฺเฐยฺย, นิพฺพานมนโส นโร.}
\csromangathagroup{\csromangathastanza{อติมาเน น ติฏฺเฐยฺย, นิพฺพานมนโส นโร.}}

\prose{\textbf{นิทฺทํ ตนฺทิํ สเห ถีนนฺ}ติ \textbf{นิทฺทา}ติ ยา กายสฺส อกลฺยตา อกมฺมญฺญตา โอนาโห ปริโยนาโห อนฺโตสโมโรโธ มิทฺธํ สุปฺปํ ปจลายิกา สุปฺปํ สุปฺปนา สุปฺปิตตฺตํ.\csromanspacer{} \textbf{ตนฺทินฺ}ติ ยา ตนฺที\footnote{ตนฺทิ (สี, สฺยา, ก)} ตนฺทิยนา ตนฺทิยิตตฺตํ ตนฺทิมนกตา อาลสฺยํ อาลสฺยายนา อาลสฺยายิตตฺตํ.\csromanspacer{} \textbf{ถีนนฺ}ติ ยา จิตฺตสฺส อกลฺยาตา อกมฺมญฺญตา โอลียนา สลฺลียนา ลีนํ ลียนา ลียิตตฺตํ, ถินํ ถิยนา ถิยิตตฺตํ จิตฺตสฺส.\csromanspacer{}\csromanspacer{}\textbf{นิทฺทํ ตนฺทิํ สเห ถีนนฺ}ติ นิทฺทญฺจ ตนฺทิญฺจ ถินญฺจ สเห สเหยฺย ปริสเหยฺย อภิภเวยฺย อชฺโฌตฺตเรยฺย ปริยาทิเยยฺย มทฺเทยฺยาติ นิทฺทํ ตนฺทิํ สเห ถีนํ.}

\prose{\textbf{ปมาเทน น สํวเส}ติ ปมาโท วตฺตพฺโพ.\csromanspacer{} กายทุจฺจริเต วา วจีทุจฺจริเต วา มโนทุจฺจริเต วา ปญฺจสุ วา กามคุเณสุ จิตฺตสฺส โวสคฺโค โวสคฺคานุปฺปาทนํ\footnote{โวสฺสคฺโค โวสฺสคฺคานุปฺปาทนํ (พหูสุ)} วา กุสลานํ วา ธมฺมานํ ภาวนาย อสกฺกจฺจกิริยตา อสาตจฺจกิริยตา อนฏฺฐิตกิริยตา โอลีนวุตฺติตา นิกฺขิตฺตจฺฉนฺทตา นิกฺขิตฺตธุรตา อนาเสวนา อภาวนา อพหุลีกมฺมํ อนธิฏฺฐานํ อนนุโยโค ปมาโท, โย เอวรูโป ปมาโท ปมชฺชนา \csromanfolio{334}ปมชฺชิตตฺตํ, อยํ วุจฺจติ ปมาโท.\csromanspacer{} \textbf{ปมาเทน น สํวเส}ติ ปมาเทน น วเสยฺย น สํวเสยฺย น อาวเสยฺย น ปริวเสยฺย ปมาทํ ปชเหยฺย วิโนเทยฺย พฺยนฺติํ กเรยฺย อนภาวํ คเมยฺย, ปมาทา อารโต อสฺส วิรโต ปฏิวิรโต นิกฺขนฺโต นิสฺสโฏ วิปฺปมุตฺโต วิสญฺญุตฺโต วิมริยาทิกเตน เจตสา วิหเรยฺยาติ ปมาเทน น สํวเส.}

\prose{\textbf{อติมาเน น ติฏฺเฐยฺยา}ติ \textbf{อติมาโน}ติ อิเธกจฺโจ ปรํ อติมญฺญติ ชาติยา วา โคตฺเตน วา ฯเปฯ อญฺญตรญฺญตเรน วา วตฺถุนา.\csromanspacer{} โย เอวรูโป มาโน มญฺญนา มญฺญิตตฺตํ อุนฺนติ อุนฺนโม ธโช สมฺปคฺคาโห เกตุกมฺยตา จิตฺตสฺส, อยํ วุจฺจติ \textbf{อติมาโน}.\csromanspacer{}\csromanspacer{}\textbf{อติมาเน น ติฏฺเฐยฺยา}ติ อติมาเน น ติฏฺเฐยฺย น สํติฏฺเฐยฺย อติมานํ ปชเหยฺย วิโนเทยฺย พฺยนฺติํ กเรยฺย อนภาวํ คเมยฺย, อติมานา อารโต อสฺส วิรโต ปฏิวิรโต นิกฺขนฺโต นิสฺสโฏ วิปฺปมุตฺโต วิสญฺญุตฺโต วิมริยาทิกเตน เจตสา วิหเรยฺยาติ อติมาเน น ติฏฺเฐยฺย.}

\prose{\textbf{นิพฺพานมนโส นโร}ติ อิเธกจฺโจ ทานํ เทนฺโต สีลํ สมาทิยนฺโต อุโปสถกมฺมํ กโรนฺโต ปานียํ ปริโภชนียํ อุปฏฺฐเปนฺโต ปริเวณํ สมฺมชฺชนฺโต เจติยํ วนฺทนฺโต เจติเย คนฺธมาลํ อาโรเปนฺโต เจติยํ ปทกฺขิณํ กโรนฺโต ยํ กิญฺจิ เตธาตุกํ กุสลาภิสงฺขารํ อภิสงฺขโรนฺโต น คติเหตุ น อุปปตฺติเหตุ น ปฏิสนฺธิเหตุ น ภวเหตุ น สํสารเหตุ น วฏฺฏเหตุ สพฺพํ ตํ วิสํโยคาธิปฺปาโย นิพฺพานนินฺโน นิพฺพานโปณา นิพฺพานปพฺภาโร อภิสงฺขโรตีติ เอวมฺปิ นิพฺพานมนโส นโร.\csromanspacer{}\csromanspacer{}อถ วา สพฺพสงฺขารธาตุยา จิตฺตํ ปฏิวาเปตฺวา\footnote{ปฏิวาเสตฺวา (ก)} อมตาย ธาตุยา จิตฺตํ อุปสํหรติ “เอตํ สนฺตํ เอตํ ปณีตํ ยทิทํ สพฺพสงฺขารสมโถ สพฺพูปธิปฏินิสฺสคฺโค ตณฺหกฺขโย วิราโค นิโรโธ นิพฺพานนฺ”ติ เอวมฺปิ นิพฺพานมนโส นโร.}

\csromangathameasure{น ปณฺฑิตา อุปธิสุขสฺส เหตุ, \\ ททนฺติ ทานานิ ปุนพฺภวาย. \\ กามญฺจ เต อุปธิปริกฺขยาย, \\ ททนฺติ ทานํ อปุนพฺภวาย.}
\csromangathagroup{\csromangathastanza{น ปณฺฑิตา อุปธิสุขสฺส เหตุ, \\ ททนฺติ ทานานิ ปุนพฺภวาย. \\ กามญฺจ เต อุปธิปริกฺขยาย, \\ ททนฺติ ทานํ อปุนพฺภวาย.}}

\csromanheader{2}{15. อตฺตทณฺฑสุตฺตนิทฺเทส}
\csromangathameasure{น ปณฺฑิตา อุปธิสุขสฺส เหตุ, \\ ภาเวนฺติ ฌานานิ ปุนพฺภวาย. \\ กามญฺจ เต อุปธิปริกฺขยาย, \\ ภาเวนฺติ ฌานํ อปุนพฺภวาย. \\ เต นิพฺพุติํ อาสิสมานสา ททนฺติ, \\ ตนฺนินฺนจิตฺตา ตทธิมุตฺตา. \\ นชฺโช ยถา สาครมชฺฌุเปตา, \\ ภวนฺติ นิพฺพานปรายณา เตติ.}
\csromangathagroup{\csromangathastanza{\csromanfolio{335}น ปณฺฑิตา อุปธิสุขสฺส เหตุ, \\ ภาเวนฺติ ฌานานิ ปุนพฺภวาย. \\ กามญฺจ เต อุปธิปริกฺขยาย, \\ ภาเวนฺติ ฌานํ อปุนพฺภวาย.}\csromangathastanza{เต นิพฺพุติํ อาสิสมานสา\footnote{อาสิํสมานา (สี), อภิมานา (สฺยา)} ททนฺติ, \\ ตนฺนินฺนจิตฺตา\footnote{อาสิํสมานา (สี), อภิมานา (สฺยา)} ตทธิมุตฺตา. \\ นชฺโช ยถา สาครมชฺฌุเปตา\footnote{อาสิํสมานา (สี), อภิมานา (สฺยา)}, \\ ภวนฺติ นิพฺพานปรายณา เตติ.}}

\csromanheader{2}{นิพฺพานมนโส นโร. เตนาห ภควา\csromandash{}}
\csromangathameasure{นิทฺทํ ตนฺทิํ สเห ถีนํ, ปมาเทน น สํวเส. \\ อติมาเน น ติฏฺเฐยฺย, นิพฺพานมนโส นโรติ.}
\csromangathagroup{\csromangathastanza{นิทฺทํ ตนฺทิํ สเห ถีนํ, ปมาเทน น สํวเส. \\ อติมาเน น ติฏฺเฐยฺย, นิพฺพานมนโส นโรติ.}}

\proseitem{178}{\csromansymbolfootnote{*}{ขุ 1. 425 ปิฏฺเฐปิ.}โมสวชฺเช น นิยฺเยถ\footnote{นีเยถ (ก)}\textbf{, รูเป สฺเนหํ น กุพฺพเย.}}

\csromangathameasure{มานญฺจ ปริชาเนยฺย, สาหสา วิรโต จเร.}
\csromangathagroup{\csromangathastanza{มานญฺจ ปริชาเนยฺย, สาหสา วิรโต จเร.}}

\prose{\textbf{โมสวชฺเช น นิยฺเยถา}ติ \textbf{โมสวชฺชํ} วุจฺจติ มุสาวาโท.\csromanspacer{} อิเธกจฺโจ สภคฺคโต วา ปริสคฺคโต วา ญาติมชฺฌคโต วา ปูคมชฺฌคโต วา ราชกุลมชฺฌคโต วา อภินีโต สกฺขิปุฏฺโฐ เอหมฺโภ ปุริส ยํ ชานาสิ, ตํ วเทหีติ.\csromanspacer{} โส อชานํ วา อาห “ชานามี”ติ, ชานํ วา อาห “น ชานามี”ติ, อปสฺสํ วา อาห “ปสฺสามี”ติ, ปสฺสํ วา อาห “น ปสฺสามี”ติ, อิติ อตฺตเหตุ วา ปรเหตุ วา อามิสกิญฺจิกฺขเหตุ วา สมฺปชานมุสา ภาสิตา โหติ, อิทํ วุจฺจติ \textbf{โมสวชฺชํ}.\csromanspacer{}\csromanspacer{}อปิ จ ตีหากาเรหิ.\csromanspacer{} จตูหากาเรหิ.\csromanspacer{} ปญฺจหากาเรหิ.\csromanspacer{} ฉหากาเรหิ.\csromanspacer{} สตฺตหากาเรหิ.\csromanspacer{} อฏฺฐหากาเรหิ ฯเปฯ.\csromanspacer{} อิเมหิ อฏฺฐหากาเรหิ มุสาวาโท โหติ.\csromanspacer{}\csromanspacer{}\textbf{โมสวชฺเช น นิยฺเยถา}ติ โมสวชฺเช น ยาเยยฺย น นิยฺยาเยยฺย น วเหยฺย น สํหเรยฺย, \textbf{โมสวชฺชํ} ปชเหยฺย วิโนเทยฺย พฺยนฺติํ กเรยฺย อนภาวํ คเมยฺย, โมสวชฺชา อารโต อสฺส วิรโต ปฏิวิรโต นิกฺขนฺโต นิสฺสโฏ วิปฺปมุตฺโต วิสญฺญุตฺโต วิมริยาทิกเตน เจตสา วิหเรยฺยาติ โมสวชฺเช น นิยฺเยถ.}

\prose{\csromanfolio{336}\textbf{รูเป สฺเนหํ น กุพฺพเย}ติ \textbf{รูปนฺ}ติ จตฺตาโร จ มหาภูตา จตุนฺนญฺจ มหาภูตานํ อุปาทาย รูปํ.\csromanspacer{}\csromanspacer{}\textbf{รูเป สฺเนหํ น กุพฺพเย}ติ รูเป สฺเนหํ น กเรยฺย, ฉนฺทํ น กเรยฺย, เปมํ น กเรยฺย, ราคํ น กเรยฺย น ชเนยฺย น สญฺชเนยฺย น นิพฺพตฺเตยฺย นาภินิพฺพตฺเตยฺยาติ รูเป สฺเนหํ น กุพฺพเย.}

\prose{\textbf{มานญฺจ ปริชาเนยฺยา}ติ \textbf{มาโน}ติ เอกวิเธน \textbf{มาโน}, ยา จิตฺตสฺส อุนฺนติ.\csromanspacer{}\csromanspacer{}ทุวิเธน \textbf{มาโน}, อตฺตุกฺกํสน\textbf{มาโน}, ปรวมฺภน\textbf{มาโน}.\csromanspacer{}\csromanspacer{}ติวิเธน \textbf{มาโน}, “เสยฺโยหมสฺมี”ติ \textbf{มาโน}, “สทิโสหมสฺมี”ติ \textbf{มาโน}, “หีโนหมสฺมี”ติ \textbf{มาโน}.\csromanspacer{}\csromanspacer{}จตุวิเธน \textbf{มาโน}, ลาเภน มานํ ชเนติ, ยเสน มานํ ชเนติ, ปสํสาย มานํ ชเนติ, สุเขน มานํ ชเนติ.\csromanspacer{}\csromanspacer{}ปญฺจวิเธน \textbf{มาโน}, “ลาภิมฺหิ มนาปิกานํ รูปานนฺ”ติ มานํ ชเนติ, “ลาภิมฺหิ มนาปิกานํ สทฺทานํ.\csromanspacer{} คนฺธานํ.\csromanspacer{} รสานํ.\csromanspacer{} โผฏฺฐพฺพานนฺ”ติ มานํ ชเนติ.\csromanspacer{}\csromanspacer{}ฉพฺพิเธน \textbf{มาโน}, จกฺขุสมฺปทาย มานํ ชเนติ, โสตสมฺปทาย.\csromanspacer{} ฆานสมฺปทาย.\csromanspacer{} ชิวฺหาสมฺปทาย.\csromanspacer{} กายสมฺปทาย.\csromanspacer{} มโนสมฺปทาย มานํ ชเนติ.\csromanspacer{}\csromanspacer{}สตฺตวิเธน \textbf{มาโน}, อติ\textbf{มาโน} มานาติ\textbf{มาโน} โอ\textbf{มาโน} สทิส\textbf{มาโน} อธิ\textbf{มาโน} อสฺมิ\textbf{มาโน} มิจฺฉา\textbf{มาโน}.\csromanspacer{}\csromanspacer{}อฏฺฐวิเธน \textbf{มาโน}, ลาเภน มานํ ชเนติ, อลาเภน โอมานํ ชเนติ, ยเสน มานํ ชเนติ, อยเสน โอมานํ ชเนติ, ปสํสาย มานํ ชเนติ, นินฺทาย โอมานํ ชเนติ, สุเขน มานํ ชเนติ, ทุกฺเขน โอมานํ ชเนติ.\csromanspacer{}\csromanspacer{}นววิเธน \textbf{มาโน}, เสยฺยสฺส “เสยฺโยหมสฺมี”ติ \textbf{มาโน}, เสยฺยสฺส “สทิโสหมสฺมี”ติ \textbf{มาโน}, เสยฺยสฺส “หีโนหมสฺมี”ติ \textbf{มาโน}, สทิสสฺส “เสยฺโยหมสฺมี”ติ \textbf{มาโน}, สทิสสฺส สทิโสหมสฺมี”ติ \textbf{มาโน}, สทิสสฺส “หีโนหมสฺมี”ติ \textbf{มาโน}, หีนสฺส “เสยฺโยหมสฺมี”ติ \textbf{มาโน}, หีนสฺส ‘สทิโสหมสฺมี”ติ \textbf{มาโน}, หีนสฺส “หีโนหมสฺมี”ติ \textbf{มาโน}.\csromanspacer{}\csromanspacer{}ทสวิเธน \textbf{มาโน}, อิเธกจฺโจ มานํ ชเนติ ชาติยา วา โคตฺเตน วา ฯเปฯ อญฺญตรญฺญตเรน วา วตฺถุนา, โย เอวรูโป \textbf{มาโน} มญฺญนา มญฺญิตตฺตํ อุนฺนติ อุนฺนโม ธโชสมฺปคฺคาโห เกตุกมฺยตา จิตฺตสฺส, อยํ วุจฺจติ \textbf{มาโน}.}

\prose{\textbf{มานญฺจ ปริชาเนยฺยา}ติ มานํ ตีหิ ปริญฺญาหิ ปริชาเนยฺย ญาตปริญฺญาย ตีรณปริญฺญาย ปหานปริญฺญาย.\csromanspacer{} กตมา ญาตปริญฺญา, มานํ ชานาติ อยํ เอกวิเธน \textbf{มาโน}, ยา จิตฺตสฺส อุนฺนติ.\csromanspacer{} อยํ ทุวิเธน \textbf{มาโน} อตฺตุกฺกํสน\textbf{มาโน} ปรวมฺภน\textbf{มาโน} ฯเปฯ อยํ ทสวิเธน}

\vspace{\tipitakaparskip}
\csromancenter{อตฺตทณฺฑสุตฺตนิทฺเทส มาโน, อิเธกจฺโจ มานํ ชเนติ ชาติยา วา โคตฺเตน วา ฯเปฯ อญฺญตรญฺญตเรน วา วตฺถุนาติ ชานาติ ปสฺสติ, อยํ ญาตปริญฺญา\textbf{.}}
\prose{\csromanfolio{337}กตมา ตีรณปริญฺญา, เอตํ ญาตํ กตฺวา มานํ ตีเรติ อนิจฺจโต ทุกฺขโต ฯเปฯ นิสฺสรณโต ตีเรติ, อยํ ตีรณปริญฺญา\textbf{.}}

\prose{กตมา ปหานปริญฺญา, เอวํ ตีรยิตฺวา มานํ ปชหติ วิโนเทติ พฺยนฺติํ กโรติ อนภาวํ คเมติ\textbf{.}\csromanspacer{} อยํ ปหานปริญฺญา\textbf{.}\csromanspacer{}\csromanspacer{}\textbf{มานญฺจ ปริชาเนยฺยา}ติ มานํ อิมาหิ ตีหิ ปริญฺญาหิ ปริชาเนยฺยาติ มานญฺจ ปริชาเนยฺย\textbf{.}}

\prose{\textbf{สาหสา วิรโต จเร}ติ กตมา สาหสา จริยา, รตฺตสฺส ราคจริยา สาหสา จริยา\textbf{.}\csromanspacer{} ทุฏฺฐสฺส โทสจริยา สาหสา จริยา, มูฬฺหสฺส โมหจริยา สาหสา จริยา, วินิพทฺธสฺส มานจริยา สาหสา จริยา, ปรามฏฺฐสฺส ทิฏฺฐิจริยา สาหสา จริยา, วิกฺเขปคตสฺส อุทฺธจฺจจริยา สาหสา จริยา, อนิฏฺฐงฺคตสฺส วิจิกิจฺฉาจริยา สาหสา จริยา, ถามคตสฺส อนุสยจริยา สาหสา จริยา, อยํ สาหสา จริยา\textbf{.}\csromanspacer{}\csromanspacer{}\textbf{สาหสา วิรโต จเร}ติ สาหสา จริยาย อารโต อสฺส วิรโต ปฏิวิรโต นิกฺขนฺโต นิสฺสโฏ วิปฺปมุตฺโต วิสญฺญุตฺโต วิมริยาทิกเตน เจตสา วิหเรยฺย จเรยฺย วิจเรยฺย อิริเยยฺย วตฺเตยฺย ปาเลยฺย ยเปยฺย ยาเปยฺยาติ สาหสา วิรโต จเร\textbf{.}\csromanspacer{} เตนาห ภควา\csromandash{}}

\csromangathameasure{โมสวชฺเช น นิยฺเยถ, รูป สฺเนหํ น กุพฺพเย. \\ มานญฺจ ปริชาเนยฺย, สาหสา วิรโต จเรติ.}
\csromangathagroup{\csromangathastanza{โมสวชฺเช น นิยฺเยถ, รูป สฺเนหํ น กุพฺพเย. \\ มานญฺจ ปริชาเนยฺย, สาหสา วิรโต จเรติ.}}

\csromanhangingbegin
\hangingheaditem{179}{ปุราณํ นาภินนฺเทยฺย, นเว ขนฺติํ น กุพฺพเย\footnote{ขนฺติมกุพฺพเย (พหูสุ)}\textbf{.}}
\hangingbody{\textbf{หียมาเน น โสเจยฺย, อากาสํ2} \textbf{น สิโต สิยา.}}
\csromanhangingend

\prose{\textbf{ปุราณํ นาภินนฺเทยฺยา}ติ ปุราณํ วุจฺจติ อตีตา รูปเวทนาสญฺญาสงฺขาร วิญฺญาณา\textbf{.}\csromanspacer{} อตีเต สงฺขาเร ตณฺหาวเสน ทิฏฺฐิวเสน นาภินนฺเทยฺย นาภิวเทยฺย น อชฺโฌเสยฺย, อภินนฺทานํ อภิวทนํ อชฺโฌสานํ คาหํ \csromanfolio{338}ปรามาสํ อภินิเวสํ ปชเหยฺย วิโนเทยฺย พฺยนฺติํ กเรยฺย อนภาวํ คเมยฺยาติ ปุราณํ นาภินนฺเทยฺย.}

\prose{\textbf{นเว ขนฺติํ น กุพฺพเย}ติ นวา วุจฺจติ ปจฺจุปฺปนฺนา รูปเวทนาสญฺญาสงฺขารวิญฺญาณา.\csromanspacer{} ปจฺจุปฺปนฺเน สงฺขาเร ตณฺหาวเสน ทิฏฺฐิวเสน ขนฺติํ น กเรยฺย ฉนฺทํ น กเรยฺย เปมํ น กเรยฺย ราคํ น กเรยฺย น ชเนยฺย น สญฺชเนยฺย น นิพฺพตฺเตยฺย นาภินิพฺพตฺเตยฺยาติ นเว ขนฺติํ น กุพฺพเย.}

\prose{\textbf{หียมาเน น โสเจยฺยา}ติ หียมาเน หายมาเน ปริหายมาเน เวมาเน วิคจฺฉมาเน อนฺตรธายมาเน น โสเจยฺย น กิลเมยฺย น ปรามเสยฺย น ปริเทเวยฺย น อุรตฺตาฬิํ กนฺเทยฺย น สมฺโมหํ อาปชฺเชยฺย, จกฺขุสฺมิํ หียมาเน หายมาเน ปริหายมาเน เวมาเน วิคจฺฉมาเน อนฺตรธายมาเน.\csromanspacer{} โสตสฺมิํ ฯเปฯ ฆานสฺมิํ.\csromanspacer{} ชิวฺหาย.\csromanspacer{} กายสฺมิํ.\csromanspacer{} รูปสฺมิํ.\csromanspacer{} สทฺทสฺมิํ.\csromanspacer{} คนฺธสฺมิํ.\csromanspacer{} รสสฺมิํ.\csromanspacer{} โผฏฺฐพฺพสฺมิํ.\csromanspacer{} กุลสฺมิํ.\csromanspacer{} คณสฺมิํ.\csromanspacer{} อาวาสสฺมิํ.\csromanspacer{} ลาภสฺมิํ.\csromanspacer{} ยสสฺมิํ.\csromanspacer{} ปสํสาย.\csromanspacer{} สุขสฺมิํ.\csromanspacer{} จีวรสฺมิํ.\csromanspacer{} ปิณฺฑปาตสฺมิํ.\csromanspacer{} เสนาสนสฺมิํ.\csromanspacer{} คิลานปจฺจยเภสชฺชปริกฺขารสฺมิํ หียมาเน หายมาเน ปริหายมาเน เวมาเน วิคจฺฉมาเน อนฺตรธายมาเน น โสเจยฺย น กิลเมยฺย น ปรามเสยฺย น ปริเทเวยฺย น อุรตฺตาฬิํ กนฺเทยฺย น สมฺโมหํ อาปชฺเชยฺยาติ หียมาเน น โสเจยฺย.}

\prose{\textbf{อากาสํ น สิโต สิยา}ติ อากาสํ วุจฺจติ ตณฺหา.\csromanspacer{} โย ราโค สาราโค ฯเปฯ อภิชฺฌา โลโภ อกุสลมูลํ.\csromanspacer{} กิํการณา อากาสํ วุจฺจติ ตณฺหา.\csromanspacer{} ยาย ตณฺหาย รูปํ อากสฺสติ สมากสฺสติ คณฺหาติ ปรามสติ อภินิวิสติ.\csromanspacer{} เวทนํ.\csromanspacer{} สญฺญํ.\csromanspacer{} สงฺขาเร.\csromanspacer{} วิญฺญาณํ.\csromanspacer{} คติํ.\csromanspacer{} อุปปตฺติํ.\csromanspacer{} ปฏิสนฺธิํ.\csromanspacer{} ภวํ.\csromanspacer{} สํสารํ.\csromanspacer{} วฏฺฏํ อากสฺสติ สมากสฺสติ คณฺหาติ ปรามสติ อภินิวิสติ, ตํการณา อากาสํ วุจฺจติ ตณฺหา.\csromanspacer{}\csromanspacer{}\textbf{อากาสํ น สิโต สิยา}ติ ตณฺหานิสฺสิโต น สิยา, ตณฺหํ ปชเหยฺย วิโนเทยฺย พฺยนฺติํ กเรยฺย อนภาวํ คเมยฺย, ตณฺหาย อารโต อสฺส วิรโต ปฏิวิรโต นิกฺขนฺโต นิสฺสโฏ วิปฺปมุตฺโต วิสญฺญุตฺโต วิมริยาทิกเตน เจตสา วิหเรยฺยาติ อากาสํ น สิโต สิยา.\csromanspacer{} เตนาห ภควา\csromandash{}}

\csromangathameasure{ปุราณํ นาภินนฺเทยฺย, นเว ขนฺติํ น กุพฺพเย. \\ หียมาเน น โสเจยฺย, อากาสํ น สิโต สิยาติ.}
\csromangathagroup{\csromangathastanza{ปุราณํ นาภินนฺเทยฺย, นเว ขนฺติํ น กุพฺพเย. \\ หียมาเน น โสเจยฺย, อากาสํ น สิโต สิยาติ.}}

\csromanheader{2}{15. อตฺตทณฺฑสุตฺตนิทฺเทส}
\proseitem{180}{\csromanfolio{339}\csromansymbolfootnote{*}{ขุ 1. 425 ปิฏฺเฐปิ.}เคธํ พฺรูมิ มโหโฆติ, อาชวํ\footnote{อาจมํ (สฺยา, ก)} \textbf{พฺรูมิ ชปฺปนํ.}}

\csromangathameasure{อารมฺมณํ ปกมฺปนํ, กามปงฺโก ทุรจฺจโย.}
\csromangathagroup{\csromangathastanza{อารมฺมณํ ปกมฺปนํ\footnote{ปกปฺปนํ (สฺยา, ก)}, กามปงฺโก ทุรจฺจโย.}}

\prose{\textbf{เคธํ พฺรูมิ มโหโฆตี}ติ เคโธ วุจฺจติ ตณฺหา.\csromanspacer{} โย ราโค สาราโค ฯเปฯ อภิชฺฌา โลโภ อกุสลมูลํ.\csromanspacer{} มโหโฆ วุจฺจติ ตณฺหา.\csromanspacer{} โย ราโค สาราโค ฯเปฯ อภิชฺฌา โลโภ อกุสลมูลํ.\csromanspacer{} \textbf{เคธํ พฺรูมิ มโหโฆตี}ติ เคธํ “มโหโฆ”ติ พฺรูมิ อาจิกฺขามิ เทเสมิ ปญฺญาเปมิ ปฏฺฐเปมิ วิวรามิ วิภชามิ อุตฺตานีกโรมิ ปกาเสมีติ เคธํ พฺรุมิ มโหโฆติ.}

\prose{\textbf{อาชวํ พฺรูมิ ชปฺปนนฺ}ติ อาชวา วุจฺจติ ตณฺหา.\csromanspacer{} โย ราโค สาราโค ฯเปฯ อภิชฺฌา โลโภ อกุสลมุลํ.\csromanspacer{} \textbf{ชปฺปนา}ปิ วุจฺจติ ตณฺหา.\csromanspacer{} โย ราโค สาราโค ฯเปฯ อภิชฺฌา โลโภ อกุสลมูลํ.\csromanspacer{} \textbf{อาชวํ พฺรูมิ ชปฺปนนฺ}ติ อาชวํ “ชปฺปนา”ติ พฺรูมิ อาจิกฺขามิ ฯเปฯ อุตฺตานีกโรมิ ปกาเสมีติ อาชวํ \textbf{พฺรูมิ ชปฺปนํ.}}

\prose{\textbf{อารมฺมณํ ปกมฺปนนฺ}ติ อารมฺมณมฺปิ วุจฺจติ ตณฺหา.\csromanspacer{} โย ราโค สาราโค ฯเปฯ อภิชฺฌา โลโภ อกุสลมูลํ.\csromanspacer{} ปกมฺปนาปิ วุจฺจติ ตณฺหา.\csromanspacer{} โย ราโค สาราโค ฯเปฯ อภิชฺฌา โลโภ อกุสลมูลนฺติ อารมฺมณํ ปกมฺปนํ.}

\prose{\textbf{กามปงฺโก ทุรจฺจโย}ติ กามปงฺโก กามกทฺทโม กามกิเลโส กามปลิโป กามปลิโรโธ\footnote{กามปลิโพโธ (สฺยา, ก)} ทุรจฺจโย ทุรติวตฺโต ทุตฺตโร ทุปฺปตโร ทุสฺสมติกฺกโม ทุพฺพีติวตฺโตติ กามปงฺโก ทุรจฺจโย.\csromanspacer{} เตนาห ภควา\csromandash{}}

\prose{เคธํ พฺรูมิ มโหโฆติ, อาชวํ \textbf{พฺรูมิ ชปฺปนํ.}}

\csromangathameasure{อารมฺมณํ ปกมฺปนํ, กามปงฺโก ทุรจฺจโยติ.}
\csromangathagroup{\csromangathastanza{อารมฺมณํ ปกมฺปนํ, กามปงฺโก ทุรจฺจโยติ.}}

\proseitem{181}{\csromansymbolmark{*}สจฺจา อโวกฺกมํ มุนิ, ถเล ติฏฺฐติ พฺราหฺมโณ.}

\csromangathameasure{สพฺพํ โส ปฏินิสฺสชฺช, ส เว สนฺโตติ วุจฺจติ.}
\csromangathagroup{\csromangathastanza{สพฺพํ โส ปฏินิสฺสชฺช, ส เว สนฺโตติ วุจฺจติ.}}

\prose{\textbf{สจฺจา อโวกฺกมํ มุนี}ติ สจฺจวาจาย อโวกฺกมนฺโต สมฺมาทิฏฺฐิยา อโวกฺกมนฺโต อริยา อฏฺฐงฺคิกา มคฺคา อโวกฺกมนฺโต.\csromanspacer{} \textbf{มุนี}ติ โมนํ วุจฺจติ ญาณํ ฯเปฯ สงฺคชาลมติจฺจ โส มุนีติ สจฺจา อโวกฺกมํ มุนิ.}

\prose{\csromanfolio{340}\textbf{ถเล ติฏฺฐติ พฺราหฺมโณ}ติ ถลํ วุจฺจติ อมตํ นิพฺพานํ.\csromanspacer{} โย โส สพฺพสงฺขารสมโถ สพฺพูปธิปฏินิสฺสคฺโค ตณฺหกฺขโย วิราโค นิโรโธ นิพฺพานํ.\csromanspacer{}\csromanspacer{}\textbf{พฺราหฺมโณ}ติ สตฺตนฺนํ ธมฺมานํ พาหิตตฺตา พฺราหฺมโณ ฯเปฯ อสิโต ตาทิ ปวุจฺจเต ส พฺรหฺมา.\csromanspacer{}\csromanspacer{}\textbf{ถเล ติฏฺฐติ พฺราหฺมโณ}ติ ถเล ติฏฺฐติ, ทีเป ติฏฺฐติ, ตาเณ ติฏฺฐติ, เลเณ ติฏฺฐติ, สรเณ ติฏฺฐติ, อภเย ติฏฺฐติ, อจฺจุเต ติฏฺฐติ, อมเต ติฏฺฐติ, นิพฺพาเน ติฏฺฐตีติ ถเล ติฏฺฐติ พฺราหฺมโณ.}

\prose{\textbf{สพฺพํ โส ปฏินิสฺสชฺชา}ติ สพฺพํ วุจฺจติ ทฺวาทสายตนานิ จกฺขุ เจว รูปา จ ฯเปฯ มโน เจว ธมฺมา จ, ยโต อชฺฌตฺติกพาหิเรสุ อายตเนสุ ฉนฺทราโค ปหีโน โหติ อุจฺฉินฺนมูโล ตาลาวตฺถุกโต อนภาวํกโต อายติํ อนุปฺปาทธมฺโม.\csromanspacer{} เอตฺตาวตาปิ สพฺพํ จตฺตํ โหติ วนฺตํ มุตฺตํ ปหีนํ ปฏินิสฺสฏฺฐํ, ยโต ตณฺหา จ ทิฏฺฐิ จ มาโน จ ปหีนา โหนฺติ อุจฺฉินฺนมูลา ตาลาวตฺถุกตา อนภาวํกตา อายติํ อนุปฺปาทธมฺมา.\csromanspacer{} เอตฺตาวตาปิ สพฺพํ จตฺตํ โหติ วนฺตํ มุตฺตํ ปหีนํ ปฏินิสฺสฏฺฐํ, ยโต ปุญฺญาภิสงฺขาโร จ อปุญฺญาภิสงฺขาโร จ อาเนญฺชาภิสงฺขาโร จ ปหีนา โหนฺติ อุจฺฉินฺนมูลา ตาลาวตฺถุกตา อนภาวํกตา อายติํ อนุปฺปาทธมฺมา.\csromanspacer{} เอตฺตาวตาปิ สพฺพํ จตฺตํ โหติ วนฺตํ มุตฺตํ ปหีนํ ปฏินิสฺสฏฺฐนฺติ สพฺพํ โส ปฏินิสฺสชฺช.}

\prose{\textbf{ส เว สนฺโตติ วุจฺจตี}ติ โส สนฺโต อุปสนฺโต วูปสนฺโต นิพฺพุโต ปฏิปสฺสทฺโธติ วุจฺจติ กถียติ ภณียติ ทีปียติ โวหรียตีติ ส เว สนฺโตติ วุจฺจติ.\csromanspacer{} เตนาห ภควา\csromandash{}}

\csromangathameasure{สจฺจํ อโวกฺกมํ มุนิ, ถเล ติฏฺฐติ พฺราหฺมโณ. \\ สพฺพํ โส ปฏินิสฺสชฺช, ส เว สนฺโตติ วุจฺจตีติ.}
\csromangathagroup{\csromangathastanza{สจฺจํ อโวกฺกมํ มุนิ, ถเล ติฏฺฐติ พฺราหฺมโณ. \\ สพฺพํ โส ปฏินิสฺสชฺช, ส เว สนฺโตติ วุจฺจตีติ.}}

\csromanhangingbegin
\hangingheaditem{182}{\textbf{ส เว วิทฺวา}\footnote{วิทฺธา (สฺยา) ขุ 1. 425 ปิฏฺเฐปิ.} \textbf{ส เวทคู, ญตฺวา ธมฺมํ อนิสฺสิโต.}}
\hangingbody{\textbf{สมฺมา โส โลเก อิริยาโน, น ปิเหตี ธ กสฺสจิ.}}
\csromanhangingend

\prose{\textbf{ส เว วิทฺวา ส เวทคู}ติ \textbf{วิทฺวา}ติ \textbf{วิทฺวา} วิชฺชาคโต ญาณี วิภาวี เมธาวี.\csromanspacer{} \textbf{เวทคู}ติ \textbf{เวทา} วุจฺจนฺติ จตูสุ มคฺเคสุ ญาณํ ฯเปฯ สพฺพเวทนาสุ วีตราโค สพฺพเวทมติจฺจ เวทคู โสติ ส เว \textbf{วิทฺวา} ส เวทคู.}

\csromanheader{2}{15. อตฺตทณฺฑสุตฺตนิทฺเทส}
\prose{\csromanfolio{341}\textbf{ญตฺวา ธมฺมํ อนิสฺสิโต}ติ \textbf{ญตฺวา}ติ \textbf{ญตฺวา} ชานิตฺวา ตุลยิตฺวา ตีรยิตฺวา วิภาวยิตฺวา วิภูตํ กตฺวา. “สพฺเพ สงฺขารา อนิจฺจา”ติ \textbf{ญตฺวา} ชานิตฺวา ตุลยิตฺวา ตีรยิตฺวา วิภาวยิตฺวา วิภูตํ กตฺวา. “สพฺเพ สงฺขารา ทุกฺขา”ติ ฯเปฯ “ยํ กิญฺจิ สมุทยธมฺมํ, สพฺพํ ตํ นิโรธธมฺมนฺ”ติ \textbf{ญตฺวา} ชานิตฺวา ตุลยิตฺวา ตีรยิตฺวา วิภาวยิตฺวา วิภูตํ กตฺวา.\csromanspacer{} \textbf{อนิสฺสิโต}ติ เทฺว นิสฺสยา ตณฺหานิสฺสโย จ ทิฏฺฐินิสฺสโย จ ฯเปฯ อยํ ตณฺหานิสฺสโย ฯเปฯ อยํ ทิฏฺฐินิสฺสโย.\csromanspacer{} ตณฺหานิสฺสยํ ปหาย ทิฏฺฐินิสฺสยํ ปฏินิสฺสชฺชิตฺวา จกฺขุํ อนิสฺสิโต.\csromanspacer{} โสตํ อนิสฺสิโต.\csromanspacer{} ฆานํ อนิสฺสิโต ฯเปฯ ทิฏฺฐสุตมุตวิญฺญาตพฺเพ ธมฺเม อนิสฺสิโต อนลฺลีโน อนุปคโต อนชฺโฌสิโต อนธิมุตฺโต นิกฺขนฺโต นิสฺสโฏ วิปฺปมุตฺโต วิสญฺญุตฺโต วิมริยาทิกเตน เจตสา วิหรตีติ \textbf{ญตฺวา} ธมฺมํ อนิสฺสิโต.}

\prose{\textbf{สมฺมา โส โลเก อิริยาโน}ติ ยโต อชฺฌตฺติกพาหิเรสุ อายตเนสุ ฉนฺทราโค ปหีโน โหติ อุจฺฉินฺนมูโล ตาลาวตฺถุกโต อนภาวํกโต อายติํ อนุปฺปาทธมฺโม.\csromanspacer{} เอตฺตาวตาปิ สมฺมา โส โลเก จรติ วิหรติ อิริยติ วตฺตติ ปาเลติ ยเปติ ยาเปติ ฯเปฯ ยโต ปุญฺญาภิสงฺขาโร จ อปุญฺญาภิสงฺขาโร จ อาเนญฺชาภิสงฺขาโร จ ปหีนา โหนฺติ อุจฺฉินฺนมูลา ตาลาวตฺถุกตา อนภาวํกตา อายติํ อนุปฺปาทธมฺมา.\csromanspacer{} เอตฺตาวตาปิ สมฺมา โส โลเก จรติ วิหรติ อิริยติ วตฺตติ ปาเลติ ยเปติ ยาเปตีติ สมฺมา โส โลเก อิริยาโน.}

\prose{\textbf{น ปิเหตีธ กสฺสจี}ติ \textbf{วิหา} วุจฺจติ ตณฺหา.\csromanspacer{} โย ราโค สาคาโร ฯเปฯ อภิชฺฌา โลโภ อกุสลมูลํ.\csromanspacer{} ยสฺเสสา ปิหา ตณฺหา ปหีนา สมุจฺฉินฺนา วูปสนฺตา ปฏิปสฺสทฺธา อภพฺพุปฺปตฺติกา ญาณคฺคินา ทฑฺฒา.\csromanspacer{} โส กสฺสจิ น ปิเหติ ขตฺติยสฺส วา พฺราหฺมณสฺส วา เวสฺสสฺส วา สุทฺทสฺส วา คหฏฺฐสฺส วา ปพฺพชีตสฺส วา เทวสฺส วา มนุสฺสสฺส วาติ น ปิเหตีธ กสฺสจิ.\csromanspacer{} เตนาห ภควา\csromandash{}}

\csromangathameasure{ส เว วิทฺวา ส เวทคู, ญตฺวา ธมฺมํ อนิสฺสิโต. \\ สมฺมา โส โลเก อิริยาโน, น ปิเหตีธ กสฺสจีติ.}
\csromangathagroup{\csromangathastanza{ส เว วิทฺวา ส เวทคู, ญตฺวา ธมฺมํ อนิสฺสิโต. \\ สมฺมา โส โลเก อิริยาโน, น ปิเหตีธ กสฺสจีติ.}}

\proseitem{183}{\csromansymbolfootnote{*}{ขุ 1. 425 ปิฏฺเฐปิ.}โยธ กาเม อจฺจตริ, สงฺคํ โลเก ทุรจฺจยํ.}

\csromangathameasure{น โส โสจติ นาชฺเฌติ, ฉินฺนโสโต อพนฺธโน.}
\csromangathagroup{\csromangathastanza{น โส โสจติ นาชฺเฌติ, ฉินฺนโสโต อพนฺธโน.}}

\prose{\csromanfolio{342}\textbf{โยธ กาเม อจฺจตริ, สงฺคํ โลเก ทุรจฺจยนฺ}ติ \textbf{โย}ติ \textbf{โย} ยาทิโส ยถายุตฺโต ยถาวิหิโต ยถาปกาโร ยํฐานปฺปตฺโต ยํธมฺมสมนฺนาคโต ขตฺติ\textbf{โย} วา พฺราหฺมโณ วา เวสฺโส วา สุทฺโท วา คหฏฺโฐ วา ปพฺพชิโต วา เทโว วา มนุสฺโส วา.\csromanspacer{} \textbf{กามา}ติ อุทฺทานโต เทฺว กามา วตฺถุกามา จ กิเลสกามา จ ฯเปฯ อิเม วุจฺจนฺติ วตฺถุกามา ฯเปฯ อิเม วุจฺจนฺติ กิเลสกามา.\csromanspacer{} \textbf{สงฺคา}ติ สตฺต สงฺคา ราคสงฺโค โทสสงฺโค โมหสงฺโค มานสงฺโค ทิฏฺฐิสงฺโค กิเลสสงฺโค ทุจฺจริตสงฺโค.\csromanspacer{} \textbf{โลโก}ติ อปายโลเก มนุสฺสโลเก เทวโลเก ขนฺธโลเก ธาตุโลเก อายตนโลเก.\csromanspacer{} \textbf{สงฺคํ โลเก ทุรจฺจยนฺ}ติ \textbf{โย} กาเม จ สงฺเค จ โลเก ทุรจฺจเย ทุรติวตฺเต ทุตฺตเร ทุปฺปตเร ทุสฺสมติกฺกเม ทุพฺพีติวตฺเต อตริ อุตฺตริ ปตริ สมติกฺกมิ วีติวตฺตยีติ \textbf{โย}ธ กาเม อจฺจตริ, สงฺคํ โลเก ทุรจฺจยํ.}

\prose{\textbf{น โส โสจติ นาชฺเฌตี}ติ วิปริณตํ วา วตฺถุํ น โสจติ.\csromanspacer{} วิปริณตสฺมิํ วา วตฺถุสฺมิํ น โสจติ. “จกฺขุ เม วิปริณตนฺ”ติ น โสจติ.\csromanspacer{} โสตํ เม.\csromanspacer{} ฆานํ เม.\csromanspacer{} ชิวฺหา เม.\csromanspacer{} กา\textbf{โย} เม.\csromanspacer{} รูปา เม.\csromanspacer{} สทฺทา เม.\csromanspacer{} คนฺธา เม.\csromanspacer{} รสา เม.\csromanspacer{} โผฏฺฐพฺพา เม.\csromanspacer{} กุลํ เม.\csromanspacer{} คโณ เม.\csromanspacer{} อาวาโส เม.\csromanspacer{} ลาโภ เม.\csromanspacer{} ยโส เม.\csromanspacer{} ปสํสา เม.\csromanspacer{} สุขํ เม.\csromanspacer{} จีวรํ เม.\csromanspacer{} ปิณฺฑปาโต เม.\csromanspacer{} เสนาสนํ เม.\csromanspacer{} คิลานปจฺจยเภสชฺชปริกฺขารา เม.\csromanspacer{} มาตา เม.\csromanspacer{} ปิตา เม.\csromanspacer{} ภาตา เม.\csromanspacer{} ภคินี เม.\csromanspacer{} ปุตฺโต เม.\csromanspacer{} ธีตา เม.\csromanspacer{} มิตฺตา เม.\csromanspacer{} อมจฺจา เม.\csromanspacer{} ญาตี เม. “สาโลหิตา เม วิปริณตา”ติ น โสจติ น กิลมติ น ปริเทวติ น อุรตฺตาฬิํ กนฺทติ น สมฺโมหํ อาปชฺชตีติ น โสจติ.\csromanspacer{} \textbf{นาชฺเฌตี}ติ นชฺเฌติ น อชฺเฌติ น อุปนิชฺฌายติ น นิชฺฌายติ น ปชฺฌายติ.\csromanspacer{} อถ วา น ชายติ น ชิยฺยติ น มิยฺยติ น จวติ น อุปปชฺชตีติ นาชฺเฌตีติ น โส โสจติ นาชฺเฌติ.}

\prose{\textbf{ฉินฺนโสโต อพนฺธโน}ติ \textbf{โสตํ} วุจฺจติ ตณฺหา.\csromanspacer{} โย ราโค สาราโค ฯเปฯ อภิชฺฌา โลโภ อกุสลมูลํ.\csromanspacer{} ยสฺเสสา โสตา ตณฺหา ปหีนา สมุจฺฉินฺนา ฯเปฯ ญาณคฺคินา ทฑฺฒา, โส วุจฺจติ ฉินฺนโสโต.\csromanspacer{} \textbf{อพนฺธโน}ติ ราคพนฺธนํ โทสพนฺธนํ โมหพนฺธนํ มานพนฺธนํ ทิฏฺฐิพนฺธนํ กิเลสพนฺธนํ ทุจฺจริตพนฺธนํ, ยสฺเสเต พนฺธนา ปหีนา สมุจฺฉินฺนา ฯเปฯ}

\vspace{\tipitakaparskip}
\csromancenter{อตฺตทณฺฑสุตฺตนิทฺเทส ญาณคฺคินา ทฑฺฒา, โส วุจฺจติ อพนฺธโนติ ฉินฺนโสโต อพนฺธโน.\csromanspacer{} เตนาห ภควา\csromandash{}}
\vspace{0.5\baselineskip}
\csromangathameasure{โยธ กาเม อจฺจตริ, สงฺคํ โลเก ทุรจฺจยํ. \\ น โส โสจติ นาชฺเฌติ, ฉินฺนโสโต อพนฺธโนติ.}
\csromangathagroup{\csromangathastanza{\csromanfolio{343}โยธ กาเม อจฺจตริ, สงฺคํ โลเก ทุรจฺจยํ. \\ น โส โสจติ นาชฺเฌติ, ฉินฺนโสโต อพนฺธโนติ.}}

\proseitem{184}{\csromansymbolfootnote{*}{ขุ 1. 425 ปิฏฺเฐปิ.}ยํ ปุพฺเพ ตํ วิโสเสหิ, ปจฺฉา เต มาหุ กิญฺจนํ.}

\csromangathameasure{มชฺเฌ เจ โน คเหสฺสสิ, อุปสนฺโต จริสฺสสิ.}
\csromangathagroup{\csromangathastanza{มชฺเฌ เจ โน คเหสฺสสิ, อุปสนฺโต จริสฺสสิ.}}

\prose{\textbf{ยํ ปุพฺเพ ตํ วิโสเสหี}ติ อตีเต สงฺขาเร อารพฺภ เย กิเลสา อุปฺปชฺเชยฺยุํ, เต กิเลเส โสเสหิ วิโสเสหิ สุกฺขาเปหิ วิสุกฺขาเปหิ อพีชํ กโรหิ ปชห วิโนเทหิ พฺยนฺติํ กโรหิ อนภาวํ คเมหีติ เอวมฺปิ ยํ ปุพฺเพ ตํ วิโสเสหิ.\csromanspacer{} อถ วา เย อตีตา กมฺมาภิสงฺขารา อวิปกฺกวิปากา, เต กมฺมาภิสงฺขาเร โสเสหิ วิโสเสหิ สุกฺขาเปหิ วิสุกฺขาเปหิ อพีชํ กโรหิ ปชห วิโนเทหิ พฺยนฺติํ กโรหิ อนภาวํ คเมหีติ เอวมฺปิ ยํ ปุพฺเพ ตํ วิโสเสหิ.}

\prose{\textbf{ปจฺฉา เต มาหุ กิญฺจนนฺ}ติ ปจฺฉา วุจฺจติ อนาคตํ.\csromanspacer{} อนาคเต สงฺขาเร อารพฺภ ยานิ อุปฺปชฺเชยฺยุํ ราคกิญฺจนํ โทสกิญฺจนํ โมหกิญฺจนํ มานกิญฺจนํ ทิฏฺฐิกิญฺจนํ กิเลสกิญฺจนํ ทุจฺจริตกิญฺจนํ, อิมานิ กิญฺจนานิ\footnote{อิเม กิญฺจนา (ก) ปสฺส ที 3. 182 ปิฏฺเฐ.} ตุยฺหํ มา อหุ มา อกาสิ มา ชเนสิ มา สญฺชเนสิ มา นิพฺพตฺเตสิ มา อภินิพฺพตฺเตสิ ปชห วิโนเทหิ พฺยนฺติํ กโรหิ อนภาวํ คเมหีติ ปจฺฉา เต มาหุ กิญฺจนํ.}

\prose{\textbf{มชฺเฌ เจ โน คเหสฺสสี}ติ มชฺฌํ วุจฺจติ ปจฺจุปฺปนฺนา รูปเวทนาสญฺญาสงฺขารวิญฺญาณา, ปจฺจุปฺปเน สงฺขาเร ตณฺหาวเสน ทิฏฺฐิวเสน น คเหสฺสสิ น อุคฺคเหสฺสสิ น คณฺหิสฺสสิ น ปรามสิสฺสสิ นาภินนฺทิสฺสสิ นาภิจริสฺสสิ น อชฺโฌสิสฺสสิ อภินนฺทนํ อภิวทนํ อชฺโฌสานํ คาหํ ปรามาสํ อภินิเวสํ ปชหิสฺสสิ วิโนเทสฺสสิ พฺยนฺติํ กริสฺสสิ อนภาวํ คเมสฺสสีติ มชฺเฌ เจ โน คเหสฺสสิ.}

\prose{\textbf{อุปสนฺโต จริสฺสสี}ติ ราคสฺส สนฺตตฺตา สมิตตฺตา อุปสมิตตฺตา, โทสสฺส สนฺตตฺตา สมิตตฺตา อุปสมิตตฺตา ฯเปฯ สพฺพากุสลาภิสงฺขารานํ สนฺตตฺตา สมิตตฺตา อุปสมิตตฺตา วูปสมิตตฺตา วิชฺฌาตตฺตา นิพฺพุตตฺตา วิคตตฺตา ปฏิปสฺสทฺธตฺตา สนฺโต อุปสนฺโต วูปสนฺโต นิพฺพุโต \csromanfolio{344}ปฏิปสฺสทฺโธ จริสฺสสิ วิหริสฺสสิ อิริยิสฺสสิ วตฺติสฺสสิ ปาลิสฺสสิ ยปิสฺสสิ ยาปิสฺสสีติ อุปสนฺโต จริสฺสสิ.\csromanspacer{} เตนาห ภควา\csromandash{}}

\csromangathameasure{ยํ ปุพฺเพ ตํ วิโสเสหิ, ปจฺฉา เต มาหุ กิญฺจนํ. \\ มชฺเฌ เจ โน คเหสฺสสิ, อุปสนฺโต จริสฺสสีติ.}
\csromangathagroup{\csromangathastanza{ยํ ปุพฺเพ ตํ วิโสเสหิ, ปจฺฉา เต มาหุ กิญฺจนํ. \\ มชฺเฌ เจ โน คเหสฺสสิ, อุปสนฺโต จริสฺสสีติ.}}

\proseitem{185}{\csromansymbolfootnote{*}{ขุ 1. 425 ปิฏฺเฐปิ.}สพฺพโส นามรูปสฺมิํ, ยสฺส นตฺถิ มมายิตํ.}

\csromangathameasure{อสตา จ น โสจติ, ส เว โลเก น ชียติ.}
\csromangathagroup{\csromangathastanza{อสตา จ น โสจติ, ส เว โลเก น ชียติ.}}

\prose{\textbf{สพฺพโส นามรูปสฺมิํ, ยสฺส นตฺถิ มมายิตนฺ}ติ \textbf{สพฺพโส}ติ สพฺเพน สพฺพํ สพฺพถา สพฺพํ อเสสํ นิสฺเสสํ ปริยาทิยนวจนเมตํ \textbf{สพฺพโส}ติ.\csromanspacer{} \textbf{นามนฺ}ติ จตฺตาโร อรูปิโน ขนฺธา.\csromanspacer{} \textbf{รูปนฺ}ติ จตฺตาโร จ มหาภูตา จตุนฺนญฺจ มหาภูตานํ อุปาทาย รูปํ.\csromanspacer{} \textbf{ยสฺสา}ติ อรหโต ขีณาสวสฺส.\csromanspacer{} \textbf{มมายิตนฺ}ติ เทฺว มมตฺตา ตณฺหามมตฺตญฺจ ทิฏฺฐิมมตฺตญฺจ ฯเปฯ อิทํ ตณฺหามมตฺตํ ฯเปฯ อิทํ ทิฏฺฐิมมตฺตํ.\csromanspacer{}\csromanspacer{}\textbf{สพฺพโส นามรูปสฺมิํ, ยสฺส นตฺถิ มมายิตนฺ}ติ \textbf{สพฺพโส} นามรูปสฺมิํ มมตฺตา ยสฺส นตฺถิ น สนฺติ น สํวิชฺชนฺติ นุปลพฺภนฺติ, ปหีนา สมุจฺฉินฺนา วูปสนฺตา ปฏิปสฺสทฺธา อภพฺพุปฺปตฺติกา ญาณคฺคินา ทฑฺฒาติ \textbf{สพฺพโส} นามรูปสฺมิํ, ยสฺส นตฺถิ มมายิตํ.}

\prose{\textbf{อสตา จ น โสจตี}ติ วิปริณตํ วา วตฺถุํ น โสจติ, วิปริณตสฺมิํ วา วตฺถุสฺมิํ น โสจติ “จกฺขุ เม วิปริณตนฺ”ติ น โสจติ.\csromanspacer{} โสตํ เม.\csromanspacer{} ฆานํ เม.\csromanspacer{} ชิวฺหา เม.\csromanspacer{} กาโย เม.\csromanspacer{} รูปา เม.\csromanspacer{} สทฺทา เม.\csromanspacer{} คนฺธา เม.\csromanspacer{} รสา เม.\csromanspacer{} โผฏฺฐพฺพา เม.\csromanspacer{} กุลํ เม.\csromanspacer{} คโณ เม.\csromanspacer{} อาวาโส เม.\csromanspacer{} ลาโภ เม ฯเปฯ “สาโลหิตา เม วิปริณตา”ติ น โสจติ น กิลมติ น ปริเทวติ น อุรตฺตาฬิํ กนฺทติ น สมฺโมหํ อาปชฺชตีติ เอวมฺปิ อสตา จ น โสจติ.}

\prose{อถ วา อสตาย ทุกฺขาย เวทนาย ผุฏฺโฐ ปเรโต สโมหิโต สมนฺนาคโต น โสจติ น กิลมติ น ปริเทวติ น อุรตฺตาฬิํ กนฺทติ น สมฺโมหํ อาปชฺชตีติ เอวมฺปิ อสตา จ น โสจติ.\csromanspacer{}\csromanspacer{}อถ วา จกฺขุโรเคน ผุฏฺโฐ ปเรโต ฯเปฯ.\csromanspacer{} qอํสมกสวาตาตปสรีสปสมฺผสฺเสน ผุฏฺโฐ ปเรโต สโมหิโต สมนฺนาคโต น โสจติ น กิลมติ น ปริเทวติ น อุรตฺตาฬิํ กนฺทติ น สมฺโมหํ อาปชฺชตีติ เอวมฺปิ อสตา จ น โสจติ.\csromanspacer{} อถ วา อสนฺเต อสํวิชฺชมาเน อนุปลพฺภมาเน “อหุ วต เม, ตํ วต เม}

\vspace{\tipitakaparskip}
\csromancenter{อตฺตทณฺฑสุตฺตนิทฺเทส นตฺถิ, สิยา วต เม, ตํ วตาหํ น ลภามี”ติ น โสจติ น กิลมติ น ปริเทวติ น อุรตฺตาฬิํ กนฺทติ น สมฺโมหํ อาปชฺชตีติ เอวมฺปิ อสตา จ น โสจติ.}
\prose{\csromanfolio{345}\textbf{ส เว โลเก น ชียตี}ติ ยสฺส “มยฺหํ วา อิทํ, ปเรสํ วา อิทนฺ”ติ กิญฺจิ รูปคตํ เวทนาคตํ สญฺญาคตํ สงฺขารคตํ วิญฺญาณคตํ คหิตํ ปรามฏฺฐํ อภินิวิฏฺฐํ อชฺโฌสิตํ อธิมุตฺตํ อตฺถิ, ตสฺส ชานิ อตฺถิ, ภาสิตมฺปิ เหตํ\csromandash{}}

\csromangathameasure{ชีโน รถ’สฺสํ2 มณิกุณฺฑเล จ, \\ ปุตฺเต จ ทาเร จ ตเถว ชีโน. \\ สพฺเพสุ โภเคสุ อเสสเกสุ, \\ กสฺมา น สนฺตปฺปสิ โสกกาเล. \\ ปุพฺเพว มจฺจํ วิชหนฺติ โภคา, \\ มจฺโจ ธเน ปุพฺพตรํ ชหาสิ. \\ อสสฺสตา3 โภคิโน กามกามี, \\ ตสฺมา น โสจามหํ โสกกาเล. \\ อุเทติ อาปูรติ เวติ จนฺโท, \\ อตฺถํ ตเปตฺวาน ปเลติ สูริโย. \\ วิทิตา มยา สตฺตุก โลกธมฺมา, \\ ตสฺมา น โสจามหํ โสกกาเลติ.}
\csromangathagroup{\csromangathastanza{ชีโน\footnote{ชิณฺเณ (สี), ชินฺโน (สฺยา) 2. รถสฺเส (สี) ปญฺจกนิปาเต (อาทิมฺหิ) มณิกุณฺฑลชาตกฏฺฐกถา โอโลเกตพฺพา. 3. อสสฺสกา (สี), อสฺสกา (สฺยา)} รถ’สฺสํ2 มณิกุณฺฑเล จ, \\ ปุตฺเต จ ทาเร จ ตเถว ชีโน. \\ สพฺเพสุ โภเคสุ อเสสเกสุ, \\ กสฺมา น สนฺตปฺปสิ โสกกาเล.}\csromangathastanza{ปุพฺเพว มจฺจํ วิชหนฺติ โภคา, \\ มจฺโจ ธเน ปุพฺพตรํ ชหาสิ. \\ อสสฺสตา3 โภคิโน กามกามี, \\ ตสฺมา น โสจามหํ โสกกาเล.}\csromangathastanza{อุเทติ อาปูรติ เวติ จนฺโท, \\ อตฺถํ ตเปตฺวาน\footnote{อตฺถํ คมิตฺวาน (พหูสุ) ขุ 5. 120 ปิฏฺเฐ.} ปเลติ สูริโย. \\ วิทิตา มยา สตฺตุก โลกธมฺมา, \\ ตสฺมา น โสจามหํ โสกกาเลติ.}}

\prose{ยสฺส “มยฺหํ วา อิทํ, ปเรสํ วา อิทนฺ”ติ กิญฺจิ รูปคตํ เวทนาคตํ สญฺญาคตํ สงฺขารคตํ วิญฺญาณคตํ คหิตํ ปรามฏฺฐํ อภินิวิฏฺฐํ อชฺโฌสิตํ อธิมุตฺตํ นตฺถิ, ตสฺส ชานิ นตฺถิ.\csromanspacer{} ภาสิตมฺปิ เหตํ “นนฺทสิ สมณา”ติ.\csromanspacer{} กิํ ลทฺธา อาวุโสติ.\csromanspacer{} เตน หิ สมณ โสจสีติ.\csromanspacer{} กิํ ชียิตฺถ อาวุโสติ.\csromanspacer{} เตน หิ สมณ เนว นนฺทสิ น จ โสจสีติ.\csromanspacer{} เอวมาวุโสติ\csromandash{}}

\csromangathameasure{*จิรสฺสํ วต ปสฺสาม, พฺราหฺมณํ ปรินิพฺพุตํ. \\ อนนฺทิํ อนฆํ ภิกฺขุํ, ติณฺณํ โลเก วิสตฺติกนฺติ.}
\csromangathagroup{\csromangathastanza{\csromanfolio{346}\csromansymbolfootnote{*}{สํ 1. 1, 53 ปิฏฺเฐสุ.}จิรสฺสํ วต ปสฺสาม, พฺราหฺมณํ ปรินิพฺพุตํ. \\ อนนฺทิํ อนฆํ ภิกฺขุํ, ติณฺณํ โลเก วิสตฺติกนฺติ.}}

\csromanheader{2}{ส เว โลเก น ชียติ. เตนาห ภควา\csromandash{}}
\csromangathameasure{สพฺพโส นามรูปสฺมิํ, ยสฺส นตฺถิ มมายิตํ. \\ อสตา จ น โสจติ, ส เว โลเก น ชียตีติ.}
\csromangathagroup{\csromangathastanza{สพฺพโส นามรูปสฺมิํ, ยสฺส นตฺถิ มมายิตํ. \\ อสตา จ น โสจติ, ส เว โลเก น ชียตีติ.}}

\proseitem{186}{\csromansymbolfootnote{+}{ขุ 1. 425 ปิฏฺเฐปิ. ++ สํ 1. 294 ปิฏฺเฐ.}ยสฺส นตฺถิ “อิทํ เม”ติ, ปเรสํ วาปิ กิญฺจนํ.}

\csromangathameasure{มมตฺตํ โส อสํวินฺทํ, “นตฺถิ เม”ติ น โสจติ.}
\csromangathagroup{\csromangathastanza{มมตฺตํ โส อสํวินฺทํ, “นตฺถิ เม”ติ น โสจติ.}}

\prose{\textbf{ยสฺส นตฺถิ “อิทํ เม”ติ, ปเรสํ วาปิ กิญฺจนนฺ}ติ \textbf{ยสฺสา}ติ อรหโต ขีณาสวสฺส.\csromanspacer{} ยสฺส “มยฺหํ วา อิทํ, ปเรสํ วา อิทนฺ”ติ กิญฺจิ รุปคตํ เวทนาคตํ สญฺญาคตํ สงฺขารคตํ วิญฺญาณคตํ คหิตํ ปรามฏฺฐํ อภินิวิฏฺฐํ อชฺโฌสิตํ อธิมุตฺตํ นตฺถิ น สนฺติ น สํวิชฺชติ นุปลพฺภติ, ปหีนํ สมุจฺฉินฺนํ วูปสนฺตํ ปฏิปสฺสทฺธํ อภพฺพุปฺปตฺติกํ ญาณคฺคินา ทฑฺฒนฺติ เอวมฺปิ ยสฺส นตฺถิ “อิทํ เม”ติ, ปเรสํ วาปิ กิญฺจนํ.}

\prose{วุตฺตํ เหตํ ภควตา\csromandash{} ++ นายํ ภิกฺขเว กาโย ตุมฺหากํ, นปิ อญฺเญสํ.\csromanspacer{} ปุราณมิทํ ภิกฺขเว กมฺมํ อภิสงฺขตํ อภิสญฺเจตยิตํ เวทนียํ ทฏฺฐพฺพํ.\csromanspacer{} ตตฺร โข ภิกฺขเว สุตวา อริยสาวโก ปฏิจฺจสมุปฺปาทํเยว สาธุกํ โยนิโส มนสิกโรติ “อิติ อิมสฺมิํ สติ อิทํ โหติ, อิมสฺสุปฺปาทา อิทํ อุปฺปชฺชติ, อิมสฺมิํ อสติ อิทํ น โหติ, อิมสฺส นิโรธา อิทํ นิรุชฺฌติ, ยทิทํ อวิชฺชาปจฺจยา สงฺขารา, สงฺขารปจฺจยา วิญฺญาณํ ฯเปฯ เอวเมตสฺส เกวลสฺส ทุกฺขกฺขนฺธสฺส สมุทโย โหติ.\csromanspacer{} อวิชฺชายเตฺวว อเสสวิราคนิโรธา สงฺขารนิโรโธ ฯเปฯ เอวเมตสฺส เกวลสฺส ทุกฺขกฺขนฺธสฺส นิโรโธ โหตี”ติ.\csromanspacer{} เอวมฺปิ ยสฺส นตฺถิ “อิทํ เม”ติ, ปเรสํ วาปิ กิญฺจนํ.}

\csromanheader{2}{วุตฺตมฺปิ เหตํ ภควตา\csromandash{}}
\csromangathameasure{สุญฺญโต โลกํ อเวกฺขสฺสุ, โมฆราช สทา สโต. \\ อตฺตานุทิฏฺฐิํ อูหจฺจ, เอวํ มจฺจุตโร สิยา. \\ เอวํ โลกํ อเวกฺขนฺตํ, มจฺจุราชา น ปสฺสตีติ.}
\csromangathagroup{\csromangathastanza{สุญฺญโต โลกํ อเวกฺขสฺสุ, โมฆราช สทา สโต. \\ อตฺตานุทิฏฺฐิํ อูหจฺจ\footnote{อุหจฺจ (ก) ขุ 1. 448; ขุ 10. 7 ปิฏฺเฐสุ.}, เอวํ มจฺจุตโร สิยา. \\ เอวํ โลกํ อเวกฺขนฺตํ, มจฺจุราชา น ปสฺสตีติ.}}

\prose{เอวมฺปิ ยสฺส นตฺถิ อิทํ เมติ, ปเรสํ วาปิ กิญฺจนํ.}

\csromanheader{2}{15. อตฺตทณฺฑสุตฺตนิทฺเทส}
\prose{\csromanfolio{347}วุตฺตมฺปิ เหตํ ภควตา\csromandash{}ยํ ภิกฺขเว น ตุมฺหากํ, ตํ ปชหถ, ตํ โว ปหีนํ ทีฆรตฺตํ หิตาย สุขาย ภวิสฺสติ.\csromanspacer{} กิญฺจ ภิกฺขเว น ตุมฺหากํ, รูปํ ภิกฺขเว น ตุมฺหากํ, ตํ ปชหถ, ตํ โว ปหีนํ ทีฆรตฺตํ หิตาย สุขาย ภวิสฺสติ.\csromanspacer{} เวทนา.\csromanspacer{} สญฺญา.\csromanspacer{} สงฺขารา.\csromanspacer{} วิญฺญาณํ น ตุมฺหากํ, ตํ ปชหถ, ตํ โว ปหีนํ ทีฆรตฺตํ หิตาย สุขาย ภวิสฺสติ.\csromanspacer{} ตํ กิํ มญฺญถ ภิกฺขเว, ยํ อิมสฺมิํ เชตวเน ติณกฏฺฐสาขาปลาสํ, ตํ ชโน หเรยฺย วา ฑเหยฺย\footnote{ทเหยฺย (สี, ก) สํ 2. 28 ปิฏฺเฐปิ.} วา ยถาปจฺจยํ วา กเรยฺย, อปิ นุ ตุมฺหากํ เอวมสฺส “อเมฺห ชโน หรติ วา ฑหติ วา ยถาปจฺจยํ วา กโรตี’ติ.\csromanspacer{} โน เหตํ ภนฺเต.\csromanspacer{} ตํ กิสฺส เหตุ, โน หิ โน เอตํ ภนฺเต อตฺตา วา อตฺตนิยํ วาติ.\csromanspacer{} เอวเมวํ โข ภิกฺขเว ยํ น ตุมฺหากํ ตํ ปชหถ, ตํ โว ปหีนํ ทีฆรตฺตํ หิตาย สุขาย ภวิสฺสติ.\csromanspacer{} กิญฺจ ภิกฺขเว น ตุมฺหากํ, รูปํ ภิกฺขเว น ตุมฺหากํ ตํ ปชหถ, ตํ โว ปหีนํ ทีฆรตฺตํ หิตาย สุขาย ภวิสฺสติ.\csromanspacer{} เวทนา.\csromanspacer{} สญฺญา.\csromanspacer{} สงฺขารา.\csromanspacer{} วิญฺญาณํ น ตุมฺหากํ ตํ ปชหถ, ตํ โว ปหีนํ ทีฆรตฺตํ หิตาย สุขาย ภวิสฺสตีติ.\csromanspacer{} เอวมฺปิ ยสฺส นตฺถิ “อิทํ เม”ติ, ปเรสํ วาปิ กิญฺจนํ.\csromanspacer{} ภาสิตมฺปิ เหตํ\csromandash{}}

\csromangathameasure{*สุทฺธธมฺมสมุปฺปาทํ, สุทฺธสงฺขารสนฺตติํ. \\ ปสฺสนฺตสฺส ยถาภูตํ, น ภยํ โหติ คามณิ. \\ ติณกฏฺฐสมํ โลกํ, ยทา ปญฺญาย ปสฺสติ. \\ นาญฺญํ ปตฺถยเต กิญฺจิ, อญฺญตฺร อปฺปฏิสนฺธิยาติ.}
\csromangathagroup{\csromangathastanza{\csromansymbolfootnote{*}{ขุ 8. 185 ปิฏฺเฐ.}สุทฺธธมฺมสมุปฺปาทํ, สุทฺธสงฺขารสนฺตติํ. \\ ปสฺสนฺตสฺส ยถาภูตํ, น ภยํ โหติ คามณิ.}\csromangathastanza{ติณกฏฺฐสมํ โลกํ, ยทา ปญฺญาย ปสฺสติ. \\ นาญฺญํ ปตฺถยเต กิญฺจิ, อญฺญตฺร อปฺปฏิสนฺธิยาติ\footnote{ปฏิสนฺธิยาติ (สี)}.}}

\prose{เอวมฺปิ ยสฺส นตฺถิ “อิทํ เม”ติ, ปเรสํ วาปิ กิญฺจนํ.\csromanspacer{}\csromanspacer{}วชิรา ภิกฺขุนี มารํ ปาปิมนฺตํ เอตทโวจ\csromandash{}}

\csromangathameasure{“กํ นุ สตฺโตติ ปจฺเจสิ, มาร ทิฏฺฐิคตํ นุ เต. \\ สุทฺธสงฺขารปุญฺโชยํ, นยิธ สตฺตุปลพฺภติ. \\ ยถา หิ องฺคสมฺภารา, โหติ สทฺโท รโถ อิติ. \\ เอวํ ขนฺเธสุ สนฺเตสุ, โหติ สตฺโตติ สมฺมุติ. \\ ทุกฺขเมว หิ สมฺโภติ, ทุกฺขํ ติฏฺฐติ เวติ จ. \\ นาญฺญตฺร ทุกฺขา สมฺโภติ, นาญฺญํ ทุกฺขา นิรุชฺฌตี”ติ.}
\csromangathagroup{\csromangathastanza{“กํ นุ สตฺโตติ ปจฺเจสิ, มาร ทิฏฺฐิคตํ นุ เต. \\ สุทฺธสงฺขารปุญฺโชยํ, นยิธ สตฺตุปลพฺภติ.}\csromangathastanza{ยถา หิ\footnote{ยถาปิ (พหูสุ) สํ 1. 137 ปิฏฺเฐ.} องฺคสมฺภารา, โหติ สทฺโท รโถ อิติ. \\ เอวํ ขนฺเธสุ สนฺเตสุ, โหติ สตฺโตติ สมฺมุติ\footnote{ยถาปิ (พหูสุ) สํ 1. 137 ปิฏฺเฐ.}.}\csromangathastanza{ทุกฺขเมว หิ สมฺโภติ, ทุกฺขํ ติฏฺฐติ เวติ จ. \\ นาญฺญตฺร ทุกฺขา สมฺโภติ, นาญฺญํ ทุกฺขา นิรุชฺฌตี”ติ.}}

\prose{\csromanfolio{348}\csromansymbolfootnote{*}{สํ 2. 280; ขุ 8. 105ปิฏฺเฐสุ.}เอวมฺปิ ยสฺส นตฺถิ อิทํ เมติ, ปเรสํ วาปิ กิญฺจนํ.}

\prose{วุตฺตเญฺหตํ ภควตา + เอวเมวํ โข ภิกฺขเว ภิกฺขุ รูปํ สมนฺเนสติ ยาวตา รูปสฺส คติ.\csromanspacer{} เวทนํ.\csromanspacer{} สญฺญํ.\csromanspacer{} สงฺขาเร.\csromanspacer{} วิญฺญาณํ สมนฺเนสติ ยาวตา วิญฺญาณสฺส คติ.\csromanspacer{} ตสฺส รูปํ สมนฺเนสโต ยาวตา รูปสฺส คติ.\csromanspacer{} เวทนํ.\csromanspacer{} สญฺญํ.\csromanspacer{} สงฺขาเร.\csromanspacer{} วิญฺญาณํ สมนฺเนสโต ยาวตา วิญฺญาณสฺส คติ.\csromanspacer{} ยมฺปิสฺส ตํ โหติ “อหนฺ”ติ วา “มมนฺ”ติ วา อสฺมี”ติ วา, ตมฺปิ ตสฺส นโหตีติ.\csromanspacer{} เอวมฺปิ ยสฺส นตฺถิ “อิทํ เม”ติ, ปเรสํ วาปิ กิญฺจนํ.}

\prose{อายสฺมา อานนฺโท ภควนฺตํ เอตทโวจ\csromandash{} * “สุญฺโญ โลโก สุญฺโญ โลโก”ติ ภนฺเต วุจฺจติ, กิตฺตาวตา นุ โข ภนฺเต “สุญฺโญ โลโก”ติ วุจฺจตีติ.\csromanspacer{} ยสฺมา จ โข อานนฺท สุญฺญํ อตฺเตน วา อตฺตนิเยน วา, ตสฺมา “สุญฺโญ โลโก”ติ วุจฺจติ.\csromanspacer{} กิญฺจานนฺท สุญฺญํ อตฺเตน วา อตฺตนิเยน วา, จกฺขุ โข อานนฺท สุญฺญํ อตฺเตน วา อตฺตนิเยน วา, รูปา สุญฺญา, จกฺขุวิญฺญาณํ สุญฺญํ, จกฺขุสมฺผสฺโส สุญฺโญ, ยทิทํ จกฺขุสมฺผสฺสปจฺจยา อุปฺปชฺชติ เวทยิตํ สุขํ วา ทุกฺขํ วา อทุกฺขมสุขํ วา, ตมฺปิ สุญฺญํ.\csromanspacer{} โสตํ สุญฺญํ.\csromanspacer{} สทฺทา สุญฺญา.\csromanspacer{} ฆานํ สุญฺญํ.\csromanspacer{} คนฺธา สุญฺญา.\csromanspacer{} ชิวฺหา สุญฺญา.\csromanspacer{} รสา สุญฺญา.\csromanspacer{} กาโย สุญฺโญ.\csromanspacer{} โผฏฺฐพฺพา สุญฺญา.\csromanspacer{} มโน สุญฺโญ, ธมฺมา สุญฺญา, มโนวิญฺญาณํ สุญฺญํ มโนสมฺผสฺโส สุญฺโญ, ยทิทํ มโนสมฺผสฺสปจฺจยา อุปฺปชฺชติ เวทยิตํ สุขํ วา ทุกฺขํ วา อทุกฺขมสุขํ วา, ตมฺปิ สุญฺญํ อตฺเตน วา อตฺตนิเยน วา.\csromanspacer{} ยสฺมา โข อานนฺท สุญฺญํ อตฺเตน วา อตฺตนิเยน วา, ตสฺมา “สุญฺโญ โลโก”ติ วุจฺจตีติ.\csromanspacer{} เอวมฺปิ ยสฺส นตฺถิ “อิทํ เม”ติ, ปเรสํ วาปิ กิญฺจนํ.}

\prose{\textbf{มมตฺตํ โส อสํวินฺทนฺ}ติ \textbf{มมตฺตา}ติ เทฺว \textbf{มมตฺตา} ตณฺหามมตฺตญฺจ ทิฏฺฐิมมตฺตญฺจ ฯเปฯ อิทํ ตณฺหามมตฺตํ ฯเปฯ อิทํ ทิฏฺฐิมมตฺตํ.\csromanspacer{} ตณฺหามมตฺตํ ปหาย ทิฏฺฐิมมตฺตํ ปฏินิสฺสชฺชิตฺวา มมตฺตํ อวินฺทนฺโต อสํวินฺทนฺโต อนธิคจฺฉนฺโต อปฺปฏิลภนฺโตติ มมตฺตํ โส อสํวินฺทํ.}

\prose{\textbf{“นตฺถิ เม”ติ น โสจตี}ติ วิปริณตํ วา วตฺถุํ น โสจติ, วิปริณตสฺมิํ วา วตฺถุสฺมิํ น โสจติ. “จกฺขุ เม วิปริณตนฺ”ติ น โสจติ.\csromanspacer{} โสตํ เม ฯเปฯ “สาโลหิตา เม วิปริณตา”ติ น โสจติ น กิลมติ น ปริเทวติ น อุรตฺตาฬิํ กนฺทติ น สมฺโมหํ อาปชฺชตีติ “นตฺถิ เม”ติ น โสจติ.\csromanspacer{} เตนาห ภควา\csromandash{}}

\csromanheader{2}{15. อตฺตทณฺฑสุตฺตนิทฺเทส}
\csromangathameasure{ยสฺส นตฺถิ “อิทํ เม”ติ, ปเรสํ วาปิ กิญฺจนํ. \\ มมตฺตํ โส อสํวินฺทํ, “นตฺถิ เม”ติ น โสจตีติ.}
\csromangathagroup{\csromangathastanza{\csromanfolio{349}ยสฺส นตฺถิ “อิทํ เม”ติ, ปเรสํ วาปิ กิญฺจนํ. \\ มมตฺตํ โส อสํวินฺทํ, “นตฺถิ เม”ติ น โสจตีติ.}}

\proseitem{187}{\csromansymbolfootnote{*}{ขุ 1. 425 ปิฏฺเฐปิ.}\textbf{อนิฏฺฐุรี อนนุคิทฺโธ, อเนโช สพฺพธี สโม}.}

\csromangathameasure{ตมานิสํสํ ปพฺรูมิ, ปุจฺฉิโต อวิกมฺปินํ.}
\csromangathagroup{\csromangathastanza{ตมานิสํสํ ปพฺรูมิ, ปุจฺฉิโต อวิกมฺปินํ.}}

\prose{\textbf{อนิฏฺฐุรี อนนุคิทฺโธ, อเนโช สพฺพธี สโม}ติ กตมํ \textbf{นิฏฺฐุริยํ.}\csromanspacer{} อิเธกจฺโจ นิฏฺฐุริโย โหติ, ปรลาภสกฺการครุการมานนวนฺทนปูชนาสุ อิสฺสติ อุสูยติ\footnote{อุสฺสุยฺยติ (สฺยา), อุสฺสูยติ (ก)} อิสฺสํ พนฺธติ.\csromanspacer{} ยํ เอวรูปํ นิฏฺฐุริยํ นิฏฺฐุริยกมฺมํ อิสฺสา อิสฺสายนา อิสฺสายิตตฺตํ อุสูยา อุสูยนา อุสูยิตตฺตํ.\csromanspacer{} อิทํ วุจฺจติ \textbf{นิฏฺฐุริยํ.}\csromanspacer{} ยสฺเสตํ นิฏฺฐุริยํ ปหีนํ สมุจฺฉินฺนํ ฯเปฯ ญาณคฺคินา ทฑฺฒํ.\csromanspacer{} โส วุจฺจติ อนิฏฺฐุรีติ.\csromanspacer{}\csromanspacer{}\textbf{อนนุคิทฺโธ}ติ \textbf{เคโธ} วุจฺจติ ตณฺหา, โย ราโค สาราโค ฯเปฯ อภิชฺฌา โลโภ อกุสลมูลํ.\csromanspacer{} ยสฺเสโส \textbf{เคโธ} ปหีโน สมุจฺฉินฺโน ฯเปฯ ญาณคฺคินา ทฑฺโฒ.\csromanspacer{} โส วุจฺจติ อนนุคิทฺโธ.\csromanspacer{} โส รูเป อคิทฺโธ, สทฺเท ฯเปฯ ทิฏฺฐสุตมุตวิญฺญาตพฺเพสุ ธมฺเมสุ อวิทฺโธ อคธิโต อมุจฺฉิโต อนชฺโฌสนฺโน วิต\textbf{เคโธ} วิคต\textbf{เคโธ} จตฺต\textbf{เคโธ} วนฺต\textbf{เคโธ} มุตฺต\textbf{เคโธ} ปหีน\textbf{เคโธ} ปฏินิสฺสฏฺฐ\textbf{เคโธ} วีตราโค วิคตราโค จตฺตราโค วนฺตราโค มุตฺตราโค ปหีนราโค ปฏินิสฺสฏฺฐราโค นิจฺฉาโต นิพฺพุโต สีติภูโต สุขปฺปฏิสํเวที พฺรหฺมภูเตน อตฺตนา วิหรตีติ อนิฏฺฐุรี อนนุคิทฺโธ.}

\prose{\textbf{อเนโช สพฺพธี สโม}ติ \textbf{เอชา} วุจฺจติ ตณฺหา, โย ราโค สาราโค ฯเปฯ อภิชฺฌา โลโภ อกุสลมูลํ.\csromanspacer{} ยสฺเสสา \textbf{เอชา} ตณฺหา ปหีนา สมุจฺฉินฺนา ฯเปฯ ญาณคฺคินา ทฑฺฒา, โส วุจฺจติ อเนโช, \textbf{เอชา}ย ปหีนตฺตา อเนโช.\csromanspacer{} โส ลาเภปิ น อิญฺชติ, อลาเภปิ น อิญฺชติ, ยเสปิ น อิญฺชติ, อยเสปิ น อิญฺจติ, ปสํสายปิ น อิญฺชติ, นินฺทายปิ น อิญฺชติ, สุเขปิ น อิญฺชติ, ทุกฺเขปิ น อิญฺชติ, น จลติ น เวธติ นปฺปเวธติ น สมฺปเวธตีติ อเนโช.\csromanspacer{}\csromanspacer{}\textbf{สพฺพธี สโม}ติ \textbf{สพฺพํ} วุจฺจติ ทฺวาทสายตนานิ.\csromanspacer{} จกฺขุ เจว รูปา จ ฯเปฯ.\csromanspacer{} มโน เจว ธมฺมา จ, ยโต อชฺฌตฺติกพาหิเรสุ อายตเนสุ ฉนฺทราโค ปหีโน โหติ อุจฺฉินฺนมูโล ตาลาวตฺถุกโต อนภาวํกโต อายติํ อนุปฺปาทธมฺโม, โส วุจฺจติ สพฺพธิ สโม.\csromanspacer{} โส สพฺพตฺถ ตาทิ สพฺพตฺถ มชฺฌตฺโต สพฺพตฺถ อุเปกฺขโกติ อเนโช สพฺพธี สโม.}

\prose{\csromanfolio{350}\textbf{ตมานิสํสํ ปพฺรูมิ, ปุจฺฉิโต อวิกมฺปินนฺ}ติ อวิกมฺปินํ ปุคฺคลํ ปุฏฺโฐ ปุจฺฉิโต ยาจิโต อชฺเฌสิโต ปสาทิโต อิเม จตฺตาโร อานิสํเส ปพฺรูมิ.\csromanspacer{} โย โส อนิฏฺฐุรี อนนุคิทฺโธ อเนโช สพฺพธิ สโมติ พฺรูมิ อาจิกฺขามิ ฯเปฯ ปกาเสมีติ ตมานิสํสํ ปพฺรูมิ, ปุจฺฉิโต อวิกมฺปินํ, เตนาห ภควา\csromandash{}}

\csromangathameasure{อนิฏฺฐุรี อนนุคิทฺโธ, อเนโช สพฺพธี สโม.}
\csromangathagroup{\csromangathastanza{อนิฏฺฐุรี อนนุคิทฺโธ, อเนโช สพฺพธี สโม.}}

\prose{\textbf{ตมานิสํสํ ปพฺรูมิ, ปุจฺฉิโต อวิกมฺปินนฺ}ติ.}

\csromanhangingbegin
\hangingheaditem{188}{\textbf{อเนชสฺส วิชานโต}, นตฺถิ กาจิ นิสงฺขติ\footnote{นิสงฺขิติ (พหูสุ) ขุ 1. 425 ปิฏฺเฐปิ.}.}
\hangingbody{\textbf{วิรโต โส วิยารพฺภา}2\textbf{, เขมํ ปสฺสติ สพฺพธิ.}}
\csromanhangingend

\csromanmatikaanchor{mtk.16}
\csromantocmarkref{subsection}{15. อตฺตทณฺฑสุตฺตนิทฺเทส}{mtk.16}
\bookmark[dest={mtk.16},level=2]{15. อตฺตทณฺฑสุตฺตนิทฺเทส}
\prose{\textbf{อเนชสฺส วิชานโต}ติ \textbf{เอชา} วุจฺจติ ตณฺหา, โย ราโค สาราโค ฯเปฯ อภิชฺฌา โลโภ อกุสลมูลํ, ยสฺเสสา \textbf{เอชา} ตณฺหา ปหีนา สมุจฺฉินฺนา ฯเปฯ ญาณคฺคินา ทฑฺฒา.\csromanspacer{} โส วุจฺจติ อเนโช, \textbf{เอชา}ย ปหีนตฺตา อเนโช.\csromanspacer{} โส ลาเภปิ น อิญฺชติ, อลาเภปิ น อิญฺชติ, ยเสปิ น อิญฺชติ, อยเสปิ น อิญฺชติ, ปสํสายปิ น อิญฺชติ, นินฺทายปิ น อิญฺชติ, สุเขปิ น อิญฺชติ, ทุกฺเขปิ น อิญฺชติ น จลติ น เวธติ นปฺปเวธติ น สมฺปเวธตีติ อเนชสฺส.\csromanspacer{}\csromanspacer{}\textbf{วิชานโต}ติ ชานโต อาชานโต วิชานโต ปฏิวิชานโต ปฏิวิชฺฌโต “สพฺเพ สงฺขารา อนิจฺจา”ติ ชานโต อาชานโต วิชานโต ปฏิวิชานโต ปฏิวิชฺฌโต. “สพฺเพ สงฺขารา ทุกฺขา”ติ ฯเปฯ “ยํ กิญฺจิ สมุทยธมฺมํ, สพฺพํ ตํ นิโรธธมฺมนฺ”ติ ชานโต อาชานโต วิชานโต ปฏิวิชานโต ปฏิวิชฺฌโตติ อเนชสฺส วิชานโต.}

\prose{\textbf{นตฺถิ กาจิ นิสงฺขตี}ติ \textbf{นิสงฺขติโย} วุจฺจนฺติ ปุญฺญาภิสงฺขาโร อปุญฺญาภิสงฺขาโร อาเนญฺชาภิสงฺขาโร.\csromanspacer{} ยโต ปุญฺญาภิสงฺขาโร จ อปุญฺญาภิสงฺขาโร จ อาเนญฺชาภิสงฺขาโร จ ปหีนา โหนฺติ อุจฺฉินฺนมูลา ตาลาวตฺถุกตา อนภาวํกตา อายติํ อนุปฺปาทธมฺมา.\csromanspacer{} เอตฺตาวตา \textbf{นิสงฺขติโย} นตฺถิ น สนฺติ น สํวิชฺชนฺติ นุปลพฺภนฺติ.\csromanspacer{} ปหีนา สมุจฺฉินฺนา วูปสนฺตา ปฏิปสฺสทฺธา อภพฺพุปฺปตฺติกา ญาณคฺคินา ทฑฺฒาติ นตฺถิ กาจิ นิสงฺขติ.}

\prose{\textbf{วิรโต โส วิยารพฺภา}ติ วิยารพฺโภ วุจฺจติ ปุญฺญาภิสงฺขาโร อปุญฺญาภิสงฺขาโร อาเนญฺชาภิสงฺขาโร, ยโต ปุญฺญาภิสงฺขาโร จ}

\vspace{\tipitakaparskip}
\csromancenter{อตฺตทณฺฑสุตฺตนิทฺเทส อปุญฺญาภิสงฺขาโร จ อาเนญฺชาภิสงฺขาโร จ ปหีนา โหนฺติ อุจฺฉินฺนมูลา ตาลาวตฺถุกตา อนภาวํกตา อายติํ อนุปฺปาทธมฺมา.\csromanspacer{} เอตฺตาวตา อารพฺภา วิยารพฺภา อารโต อสฺส วิรโต ปฏิวิรโต นิกฺขนฺโต นิสฺสโฏ วิปฺปมุตฺโต วิสญฺญุตฺโต วิมริยาทิกเตน เจตสา วิหรตีติ วิรโต โส วิยารพฺภา.}
\prose{\csromanfolio{351}\textbf{เขมํ ปสฺสติ สพฺพธี}ติ ภยกโร ราโค ภยกโร โทโส ภยกโร โมโห ฯเปฯ ภยกรา กิเลสา.\csromanspacer{} ภยกรสฺส ราคสฺส ปหีนตฺตา ฯเปฯ ภยกรานํ กิเลสานํ ปหีนตฺตา สพฺพตฺถ เขมํ ปสฺสติ, สพฺพตฺถ อภยํ ปสฺสติ, สพฺพตฺถ อนีติกํ ปสฺสติ, สพฺพตฺถ อนุปทฺทวํ ปสฺสติ, สพฺพตฺถ อนุปสคฺคํ ปสฺสติ, สพฺพตฺถ อนุปสฏฺฐตฺตํ\footnote{ปสฺสทฺธํ (สฺยา)} ปสฺสตีติ เขมํ ปสฺสติ สพฺพธิ.\csromanspacer{} เตนาห ภควา\csromandash{}}

\csromangathameasure{อเนชสฺส วิชานโต, นตฺถิ กาจิ นิสงฺขติ. \\ วิรโต โส วิยารพฺภา, เขมํ ปสฺสติ สพฺพธีติ.}
\csromangathagroup{\csromangathastanza{อเนชสฺส วิชานโต, นตฺถิ กาจิ นิสงฺขติ. \\ วิรโต โส วิยารพฺภา, เขมํ ปสฺสติ สพฺพธีติ.}}

\csromanhangingbegin
\hangingheaditem{189}{น สเมสุ น โอเมสุ, น อุสฺเสสุ วทเต มุนิ.}
\hangingbody{\textbf{สนฺโต โส วีตมจฺฉโร}2\textbf{, นาเทติ น นิรสฺสติ.} (อิติ ภควา)}
\csromanhangingend

\prose{\textbf{น สเมสุ น โอเมสุ, น อุสฺเสสุ วทเต มุนี}ติ \textbf{มุนี}ติ โมนํ วุจฺจติ ญาณํ ฯเปฯ สงฺคชาลมติจฺจ โส มุนิ. “เสยฺโยหมสฺมี”ติ วา “สทิโสหมสฺมี”ติ วา “หีโนหมสฺมี”ติ วา น วทติ น กเถติ น ภณติ น ทีปยติ น โวหรตีติ น สเมสุ น โอเมสุ, น อุสฺเสสุ วทเต มุนิ.}

\prose{\textbf{สนฺโต โส วีตมจฺฉโร}ติ \textbf{สนฺโต}ติ ราคสฺส สนฺตตฺตา สมิตตฺตา \textbf{สนฺโต}.\csromanspacer{} โทสสฺส.\csromanspacer{} โมหสฺส ฯเปฯ สพฺพากุสลาภิสงฺขารานํ สนฺตตฺตา สมิตตฺตา วูปสมิตตฺตา วิชฺฌาตตฺตา นิพฺพุตตฺตา วิคตตฺตา ปฏิปสฺสทฺธตฺตา \textbf{สนฺโต} อุป\textbf{สนฺโต} วูป\textbf{สนฺโต} นิพฺพุโต ปฏิปสฺสทฺโธติ \textbf{สนฺโต} โส.\csromanspacer{}\csromanspacer{}\textbf{วิตมจฺฉโร}ติ ปญฺจ มจฺฉริยานิ อาวาสมจฺฉริยํ ฯเปฯ คาโห วุจฺจติ มจฺฉริยํ.\csromanspacer{} ยสฺเสตํ มจฺฉริยํ ปหีนํ สมุจฺฉินฺนํ ฯเปฯ ญาณคฺคินา ทฑฺฒํ, โส วุจฺจติ.\csromanspacer{} วีตมจฺฉโร วิคตมจฺฉโร จตฺตมจฺฉโร วนฺตมจฺฉโร มุตฺตมจฺฉโร ปหีนมจฺฉโร ปฏินิสฺสฏฺฐมจฺฉโรติ \textbf{สนฺโต} โส วีตมจฺฉโร.}

\prose{\csromanfolio{352}\textbf{นาเทติ น นิรสฺสติ, อิติ ภควา}ติ \textbf{นาเทตี}ติ รูปํ นาทิยติ น อุปาทิยติ น คณฺหาติ น ปรามสติ นาภินิวิสติ.\csromanspacer{} เวทนํ.\csromanspacer{} สญฺญํ.\csromanspacer{} สงฺขาเร.\csromanspacer{} วิญฺญาณํ.\csromanspacer{} คติํ.\csromanspacer{} อุปปตฺติํ.\csromanspacer{} ปฏิสนฺธิํ.\csromanspacer{} ภวํ.\csromanspacer{} สํสารํ.\csromanspacer{} วฏฺฏํ นาทิยติ น อุปาทิยติ น คนฺหาติ น ปรามสติ นาภินิวิสตีติ นาเทติ.\csromanspacer{} \textbf{น} \textbf{นิรสฺสตี}ติ รูปํ น ปชหติ น วิโนเทติ น พฺยนฺติํ กโรติ น อนภาวํ คเมติ.\csromanspacer{} เวทนํ.\csromanspacer{} สญฺญํ.\csromanspacer{} สงฺขาเร.\csromanspacer{} วิญฺญาณํ.\csromanspacer{} คติํ.\csromanspacer{} อุปปตฺติํ.\csromanspacer{} ปฏิสนฺธิํ.\csromanspacer{} ภวํ.\csromanspacer{} สํสารํ.\csromanspacer{} วฏฺฏํ น ปชหติ น วิโนเทติ น พฺยนฺติํ กโรติ น อนภาวํ คเมติ.\csromanspacer{} \textbf{ภควา}ติ คารวาธิวจนํ ฯเปฯ สจฺฉิกา ปญฺญตฺติ ยทิทํ \textbf{ภควา}ติ.\csromanspacer{} เตนาห \textbf{ภควา}\csromandash{}}

\csromangathameasure{น สเมสุ น โอเมสุ, น อุสฺเสสุ วทเต มุนิ. \\ สนฺโต โส วีตมจฺฉโร, นาเทติ น นิรสฺสติ. (อิติ ภควา.) อตฺตทณฺฑสุตฺตนิทฺเทโส ปนฺนรสโม.}
\csromangathagroup{\csromangathastanza{น สเมสุ น โอเมสุ, น อุสฺเสสุ วทเต มุนิ.}\csromangathastanza{สนฺโต โส วีตมจฺฉโร, นาเทติ น นิรสฺสติ. (อิติ ภควา.) อตฺตทณฺฑสุตฺตนิทฺเทโส ปนฺนรสโม.}}

\csromanheader{2}{16. สาริปุตฺตสุตฺตนิทฺเทส}
\csromanheader{2}{อถ สาริปุตฺตสุตฺตนิทฺเทสํ วกฺขติ\csromandash{}}
\proseitem{190}{\csromansymbolfootnote{*}{ขุ 1. 426 ปิฏฺเฐ.}\textbf{น เม ทิฏฺโฐ อิโต ปุพฺเพ}, (\textbf{อิจฺจา}ยสฺมา สาริปุตฺโต,)}

\prose{\textbf{น สุโต อุท กสฺสจิ.}}

\csromangathameasure{\textbf{เอวํ วคฺคุวโท สตฺถา,} \\ \textbf{ตุสิตา คณิมาคโต.}}
\csromangathagroup{\csromangathastanza{\textbf{เอวํ วคฺคุวโท สตฺถา,} \\ \textbf{ตุสิตา คณิมาคโต.}}}

\prose{\textbf{น เม ทิฏฺโฐ อิโต ปุพฺเพ}ติ อิโต ปุพฺเพ เม มยา น ทิฏฺฐิปุพฺโพ โส \textbf{ภควา} อิมินา จกฺขุนา อิมินา อตฺตภาเวน, ยทา \textbf{ภควา} ตาวติํสภวเน ปาริจฺฉตฺตกมูเล ปณฺฑุกมฺพลสิลายํ วสฺสํวุฏฺโฐ\footnote{วสฺสํวุตฺโต (สี, ก)} เทวคณปริวุโต มชฺเฌ มณิมเยน โสปาเนน สงฺกสฺสนครํ โอติณฺโณ อิมํ ทสฺสนํ ปุพฺเพ น ทิฏฺโฐติ น เม ทิฏฺโฐ อิโต ปุพฺเพ.}

\prose{\textbf{อิจฺจายสฺมา สาริปุตฺโต}ติ \textbf{อิจฺจา}ติ ปทสนฺธิ ปทสํสคฺโค ปทปาริปูรี อกฺขรสมวาโย พฺยญฺชนสิลิฏฺฐตา ปทานุปุพฺพตาเปตํ \textbf{อิจฺจา}ติ.\csromanspacer{} \textbf{อายสฺมา}ติ ปิยวจนํ ครุวจนํ สคารวสปฺปติสฺสวจนเมตํ อายสฺมาติ.}

\vspace{\tipitakaparskip}
\csromancenter{สาริปุตฺตสุตฺตนิทฺเทส \textbf{สาริปุตฺโต}ติ ตสฺส เถรสฺส \textbf{นา}มํ สงฺขา สมญฺญา ปญฺญตฺติ โวหาโร \textbf{นา}มํ \textbf{นา}มกมฺมํ \textbf{นา}มเธยฺยํ นิรุตฺติ พฺยญฺชนํ อภิลาโปติ อิจฺจายสฺมา \textbf{สาริปุตฺโต}.}
\prose{\csromanfolio{353}\textbf{น สุโต อุท กสฺสจี}ติ \textbf{นา}ติ ปฏิกฺเขโป.\csromanspacer{} \textbf{อุทา}ติ ปทสนฺธิ ปทสํสคฺโค ปทปาริปูรี อกฺขรสมวาโย พฺยญฺชนสิลิฏฺฐตา ปทานุปุพฺพตาเปตํ อุทาติ.\csromanspacer{} \textbf{กสฺสจี}ติ ขตฺติยสฺส วา พฺราหฺมณสฺส วา เวสฺสสฺส วา สุทฺทสฺส วา คหฏฺฐสฺส วา ปพฺพชิตสฺส วา เทวสฺส วา มนุสฺสสฺส วาติ น สุโต อุท กสฺสจิ.}

\prose{\textbf{เอวํ วคฺคุวโท สตฺถา}ติ เอวํ วคฺคุวโท มธุรวโท เปมนียวโท หทยงฺคมวโท กรวีกรุตมญฺชุโฆโส\footnote{กรวิกรุทมญฺชุสฺสโร (สฺยา)}.\csromanspacer{} อฏฺฐงฺคสมนฺ\textbf{นา}คโต โข ปน ตสฺส ภควโต มุขโต โฆโส นิจฺฉรติ วิสฏฺโฐ จ วิญฺเญยฺโย จ มญฺชุ จ สวนีโย จ พินฺทุ จ อวิสารี จ คมฺภีโร จ นินฺ\textbf{นา}ที จ.\csromanspacer{} ยถา ปริสํ โข ปน โส ภควา สเรน วิญฺญาเปติ, น อสฺส พหิทฺธา ปริสาย โฆโส นิจฺฉรติ, พฺรหฺมสฺสโร โข ปน โส ภควา กรวีกภาณีติ เอวํ วคฺคุวโท.}

\prose{\textbf{สตฺถา}ติ \textbf{สตฺถา} ภควา สตฺถวาโห, ยถา สตฺถวาโห สตฺเต กนฺตารํ ตาเรติ, โจรกนฺตารํ ตาเรติ, วาฬกนฺตารํ ตาเรติ, ทุพฺภิกฺขกนฺตารํ ตาเรติ, นิรุทกกนฺตารํ ตาเรติ อุตฺตาเรติ นิตฺตาเรติ ปตาเรติ เขมนฺตภูมิํ สมฺปาเปติ, เอวเมวํ ภควา สตฺถวาโห สตฺเต กนฺตารํ ตาเรติ, ชาติกนฺตารํ ตาเรติ, ชรากนฺตารํ ตาเรติ, พฺยาธิกนฺตารํ ฯเปฯ นรณกนฺตารํ.\csromanspacer{} โสกปริเทว ทุกฺขโทมนสฺสุปายาสกนฺตารํ ตาเรติ.\csromanspacer{} ราคกนฺตารํ ตาเรติ.\csromanspacer{} โทสกนฺตารํ.\csromanspacer{} โมหกนฺตารํ.\csromanspacer{} มานกนฺตารํ.\csromanspacer{} ทิฏฺฐิกนฺตารํ.\csromanspacer{} กิเลสกนฺตารํ.\csromanspacer{} ทุจฺจริตกนฺตารํ ตาเรติ.\csromanspacer{} ราคคหนํ ตาเรติ.\csromanspacer{} โทสคหนํ.\csromanspacer{} โมหคหนํ.\csromanspacer{} มานคหนํ.\csromanspacer{} ทิฏฺฐิคหนํ.\csromanspacer{} กิเลสคหนํ.\csromanspacer{} ทุจฺจริตคหนํ ตาเรติ อุตฺตาเรติ นิตฺตาเรติ ปตาเรติ เขมนฺตํ อมตํ นิพฺพานํ สมฺปาเปตีติ เอวมฺปิ ภควา สตฺถวาโห.}

\prose{อถ วา ภควา เนตา วิเนตา อนุเนตา ปญฺญาเปตา นิชฺฌาเปตา เปกฺเขตา ปสาเทตาติ เอวมฺปิ ภควา สตฺถวาโห.}

\prose{\csromanfolio{354}อถ วา ภควา อนุปฺปนฺนสฺส มคฺคสฺส อุปฺปาเทตา, อสญฺชาตสฺส มคฺคสฺส สญฺชเนตา, อนกฺขาตสฺส มคฺคสฺส อกฺขาตา, มคฺคญฺญู มคฺควิทู มคฺคโกวิโท มคฺคานุคา จ ปน เอตรหิ สาวกา วิหรนฺติ ปจฺฉา สมนฺนาคตาติ เอวมฺปิ ภควา สตฺถวาโหติ เอวํ วคฺคุวโท สตฺถา.}

\prose{\textbf{ตุสิตา คณิมาคโต}ติ ภควา ตุสิตกายา จวิตฺวา สโต สมฺปชาโน มาตุกุจฺฉิํ โอกฺกนฺโตติ เอวมฺปิ \textbf{ตุสิตา คณิมาคโต}.}

\prose{อถ วา เทวา วุจฺจนฺติ ตุสิตา, เต ตุฏฺฐา สนฺตุฏฺฐา อตฺตมนา ปมุทิตา ปีติโสมนสฺสชาตา เทวโลกโต คณิํ อาคโตติ เอวมฺปิ \textbf{ตุสิตา คณิมาคโต}.\csromanspacer{}\csromanspacer{}อถ วา อรหนฺโต วุจฺจนฺติ ตุสิตา, เต ตุฏฺฐา สนฺตุฏฺฐา อตฺตมนา ปริปุณฺณสงฺกปฺปา อรหนฺตานํ คณิํ อาคโตติ เอวมฺปิ \textbf{ตุสิตา คณิมาคโต}.\csromanspacer{}\csromanspacer{}\textbf{คณี}ติ คณี ภควา.\csromanspacer{} คณาจริโยติ คณี, คณสฺส สตฺถาติ คณี, คณํ ปริหรตีติ คณี, คณํ โอวทตีติ คณี, คณมนุสาสตีติ คณี, วิสารโท คณํ อุปสงฺกมตีติ คณี, คณ’สฺส\footnote{คโณสฺส (สี)} สุสฺสูสติ โสตํ โอทหติ อญฺญา จิตฺตํ อุปฏฺฐเปตีติ คณิ, คณํ อกุสลา วุฏฺฐาเปตฺวา กุสเล ปติฏฺฐาเปตีติ คณี, ภิกฺขุคณสฺส คณี, ภิกฺขุนิคณสฺส คณี, อุปาสกคณสฺส คณี, อุปาสิกาคณสฺส คณี, ราชคณสฺส คณี, ขตฺติยคณสฺส.\csromanspacer{} พฺราหฺมณคณสฺส.\csromanspacer{} เวสฺสคณสฺส.\csromanspacer{} สุทฺทคณสฺส.\csromanspacer{} เทวคณสฺส.\csromanspacer{} พฺรหฺมคณสฺส คณี, สงฺฆี คณี คณาจริโย.\csromanspacer{} \textbf{อาคโต}ติ อุปคโต สมุปคโต สมุปปนฺโน\footnote{สมฺปตฺโต (สฺยา, ก)} สงฺกสฺสนครนฺติ \textbf{ตุสิตา คณิมาคโต}.\csromanspacer{} เตนาห เถโร สาริปุตฺโต\csromandash{}}

\prose{น เม ทิฏฺโฐ อิโต ปุพฺเพ, (อิจฺจายสฺมา สาริปุตฺโต,)}

\prose{น สุโต อุท กสฺสจิ.}

\csromangathameasure{เอวํ วคฺคุวโท สตฺถา, \\ \textbf{ตุสิตา คณิมาคโต}ติ.}
\csromangathagroup{\csromangathastanza{เอวํ วคฺคุวโท สตฺถา, \\ \textbf{ตุสิตา คณิมาคโต}ติ.}}

\proseitem{191}{\csromansymbolfootnote{*}{ขุ 1. 426 ปิฏฺเฐปิ.}สเทวกสฺส โลกสฺส, ยถา ทิสฺสติ จกฺขุมา.}

\csromangathameasure{สพฺพํ ตมํ วิโนเทตฺวา, เอโกว รติมชฺฌคา.}
\csromangathagroup{\csromangathastanza{สพฺพํ ตมํ วิโนเทตฺวา, เอโกว รติมชฺฌคา.}}

\csromanheader{2}{16. สาริปุตฺตสุตฺตนิทฺเทส}
\prose{\csromanfolio{355}\textbf{สเทวกสฺส โลกสฺสา}ติ สเทวกสฺส โลกสฺส สมารกสฺส สพฺรหฺมกสฺส สสฺสมณพฺราหฺมณิยา ปชาย สเทวมนุสฺสายาติ สเทวกสฺส โลกสฺส.}

\prose{\textbf{ยถา ทิสฺสติ จกฺขุมา}ติ ยถา ภควนฺตํ ตาวติํสภวเน ปาริจฺฉตฺตกมูเล ปณฺฑุกมฺพลสิลายํ นิสินฺนํ ธมฺมํ เทเสนฺตํ เทวตา ปสนฺติ, ตถา มนุสฺสา ปสฺสนฺติ.\csromanspacer{} ยถา มนุสฺสา ปสฺสนฺติ, ตถา เทวตา ปสฺสนฺติ.\csromanspacer{} ยถา เทวานํ ทิสฺสติ, ตถา มนุสฺสานํ ทิสฺสติ.\csromanspacer{} ยถา มนุสฺสานํ ทิสฺสติ, ตถา เทวานํ ทิสฺสตีติ เอวมฺปิ ยถา ทิสฺสติ จกฺขุมา.\csromanspacer{} ยถา วา ปเนเก โภนฺโต สมณพฺราหฺมณา อทนฺตา ทนฺตวณฺเณน ทิสฺสนฺติ, อสนฺตา สนฺตวณฺเณน ทิสฺสนฺติ, อนุปสนฺตา อุปสนฺตวณฺเณน ทิสฺสนฺติ, อนิพฺพุตา นิพฺพุตวณฺเณน ทิสฺสนฺติ.}

\csromangathameasure{ปติรูปโก มตฺติกากุณฺฑโลว, \\ โลหฑฺฒมาโสว สุวณฺณฉนฺโน. \\ จรนฺติ โลเก ปริวารฉนฺนา, \\ อนฺโต อสุทฺธา พหิ โสภมานาติ.}
\csromangathagroup{\csromangathastanza{ปติรูปโก มตฺติกากุณฺฑโลว, \\ โลหฑฺฒมาโสว\footnote{โลหมาโสว (สฺยา)} สุวณฺณฉนฺโน. \\ จรนฺติ โลเก ปริวารฉนฺนา, \\ อนฺโต อสุทฺธา พหิ โสภมานาติ\footnote{โลหมาโสว (สฺยา)}.}}

\prose{น ภควา เอวํ ทิสฺสติ.\csromanspacer{} ภควา ภูเตน ตจฺเฉน ตเถน ยาถาเวน อวิปรีเตน สภาเวน ทนฺโต ทนฺตวณฺเณน ทิสฺสติ, สนฺโต สนฺตวณฺเณน ทิสฺสติ, อุปสนฺโต อุปสนฺตวณฺเณน ทิสฺสติ, นิพฺพุโต นิพฺพุตวณฺเณน ทิสฺสติ, อกปฺปิต-อิริยาปถา จ พุทฺธา ภควนฺโต ปณิธิสมฺปนฺนาติ เอวมฺปิ ยถา ทิสฺสติ จกฺขุมา.}

\prose{อถ วา ภควา วิสุทฺธสทฺโท คตกิตฺติสทฺทสิโลโก\footnote{ภฏกิตฺติสทฺทสิโลโก (สฺยา)} นาคภวเน จ สุปณฺณภวเน จ ยกฺขภวเน จ อสุรภวเน จ คนฺธพฺพภวเน จ มหาราชภวเน จ อินฺทภวเน จ พฺรหฺมภวเน จ เทวภวเน จ, เอทิโส จ ตาทิโส จ ตโต จ ภิยฺโยติ เอวมฺปิ ยถา ทิสฺสติ จกฺขุมา.}

\prose{อถ วา ภควา ทสหิ พเลหิ สมนฺนาคโต จตูหิ เวสารชฺเชหิ จตูหิ ปฏิสมฺภิทาหิ ฉหิ อภิญฺญาหิ ฉหิ พุทฺธธมฺเมหิ, เตเชน จ พเลน จ คุเณน จ วีริเยน จ ปญฺญาย จ ทิสฺสติ ญายติ ปญฺญายติ.}

\csromangathameasure{ทูเร สนฺโต ปกาเสนฺติ, หิมวนฺโตว ปพฺพโต. \\ อสนฺเตตฺถ น ทิสฺสนฺติ, รตฺติํ ขิตฺตา ยถา สราติ.}
\csromangathagroup{\csromangathastanza{\csromanfolio{356}ทูเร สนฺโต ปกาเสนฺติ, หิมวนฺโตว ปพฺพโต. \\ อสนฺเตตฺถ น ทิสฺสนฺติ, รตฺติํ ขิตฺตา\footnote{รตฺติขิตฺตา (สี) ขุ 1. 56 ปิฏฺเฐ.} ยถา สราติ.}}

\prose{เอวมฺปิ ยถา ทิสฺสติ จกฺขุมา.}

\prose{\textbf{จกฺขุมา}ติ * ภควา ปญฺจหิ จกฺขูหิ จกฺขุมา, มํสจกฺขุนาปิ จกฺขุมา, ทิพฺพจกฺขุนาปิ จกฺขุมา, ปญฺญาจกฺขุนาปิ จกฺขุมา, พุทฺธจกฺขุนาปิ จกฺขุมา, สมนฺตจกฺขุนาปิ จกฺขุมา.}

\prose{กถํ ภควา มํสจกฺขุนาปิ จกฺขุมา, มํสจกฺขุมฺหิ ภควโต ปญฺจ วณฺณา สํวิชฺชนฺติ นีโล จ วณฺโณ ปีตโก จ วณฺโณ โลหิตโก จ วณฺโณ กโณฺห จ วณฺโณ โอทาโต จ วณฺโณ.\csromanspacer{} อกฺขิโลมานิ จ ภควโต.\csromanspacer{} ยตฺถ จ อกฺขิโลมานิ ปติฏฺฐิตานิ, ตํ นีลํ โหติ สุนีลํ ปาสาทิกํ ทสฺสเนยฺยํ อุมาปุปฺผสมานํ.\csromanspacer{} ตสฺส ปรโต ปีตกํ โหติ สุปีตกํ สุวณฺณวณฺณํ ปาสาทิกํ ทสฺสเนยฺยํ กณิการปุปฺผสมานํ.\csromanspacer{} อุภโต จ อกฺขิกูฏานิ ภควโต โลหิตกานิ โหนฺติ สุโลหิตกานิ ปาสาทิกานิ ทสฺสเนยฺยานิ อินฺทโคปกสมานานิ.\csromanspacer{} มชฺเฌ กณฺหํ โหติ สุกณฺหํ อลูขํ สุทฺธํ ปาสาทิกํ ทสฺสเนยฺยํ อทฺทาริฏฺฐกสมานํ.\csromanspacer{} ตสฺส ปรโต โอทาตํ โหติ สุ-โอทาตํ เสตํ ปณฺฑรํ ปาสาทิกํ ทสฺสเนยฺยํ โอสธิตารกสมานํ.\csromanspacer{} เตน ภควา ปากติเกน มํสจกฺขุนา อตฺตภาวปริยาปนฺเนน ปุริมสุจริตกมฺมาภินิพฺพตฺเตน สมนฺตา โยชนํ ปสฺสติ ทิวา เจว รตฺติญฺจ.\csromanspacer{} ยทาปิ จตุรงฺคสมนฺนาคโต อนฺธกาโร โหติ, สูริโย จ อตฺถงฺคโต โหติ, กาฬปกฺโข จ อุโปสโถ โหติ, ติพฺโพ จ วนสณฺโฑ โหติ, มหา จ กาฬเมโฆ อพฺภุฏฺฐิโต โหติ.\csromanspacer{} เอวรูเปปิ จตุรงฺคสมนฺนาคเต อนฺธกาเร สมนฺตา โยชนํ ปสฺสติ, นตฺถิ โส กุฏฺโฏ วา กวาฏํ วา ปากาโร วา ปพฺพโต วา คจฺโฉ วา ลตา วา อาวรณํ รูปานํ ทสฺสนาย.\csromanspacer{} เอกญฺเจ ติลผลํ นิมิตฺตํ กตฺวา ติลวาเห ปกฺขิเปยฺย.\csromanspacer{} ตญฺเญว ติลผลํ อุทฺธเรยฺย.\csromanspacer{} เอวํ ปริสุทฺธํ ภควโต ปากติกมํสจกฺขุ.\csromanspacer{} เอวํ ภควา มํสจกฺขุนาปิ จกฺขุมา.}

\prose{กถํ ภควา ทิพฺเพน จกฺขุนาปิ จกฺขุมา, ภควา ทิพฺเพน จกฺขุนา วิสุทฺเธน อติกฺกนฺตมานุสเกน สตฺเต ปสฺสติ จวมาเน อุปปชฺชมาเน}

\vspace{\tipitakaparskip}
\csromancenter{สาริปุตฺตสุตฺตนิทฺเทส หีเน ปณีเต สุวณฺเณ ทุพฺพณฺเณ สุคเต ทุคฺคเต, ยถากมฺมูปเค สตฺเต ปชานาติ “อิเม วต โภนฺโต สตฺตา กายทุจฺจริเตน สมนฺนาคตา วจีทุจฺจริเตน สมนฺนาคตา มโนทุจฺจริเตน สมนฺนาคตา อริยานํ อุปวาทกา มิจฺฉาทิฏฺฐิกา มิจฺฉาทิฏฺฐิกมฺมสมาทานา, เต กายสฺส เภทา ปรํ มรณา อปายํ ทุคฺคติํ วินิปาตํ นิรยํ อุปปนฺนา.\csromanspacer{} อิเม วา ปน โภนฺโต สตฺตา กายสุจริเตน สมนฺนาคตา วจีสุจริเตน สมนฺนาคตา มโนสุจริเตน สมนฺนาคตา อริยานํ อนุปวาทกา สมฺมาทิฏฺฐิกา สมฺมาทิฏฺฐิกมฺมสมาทานา, เต กายสฺส เภทา ปรํ มรณา สุคติํ สคฺคํ โลกํ อุปปนฺนา”ติ.\csromanspacer{} อิติ ทิพฺเพน จกฺขุนา วิสุทฺเธน อติกฺกนฺตมานุสเกน สตฺเต ปสฺสติ จวมาเน อุปปชฺชมาเน หีเน ปณีเต สุวณฺเณ ทุพฺพณฺเณ สุคเต ทุคฺคเต, ยถากมฺมูปเค สตฺเต ปชานาติ.\csromanspacer{} อากงฺขมาโน จ ภควา เอกมฺปิ โลกธาตุํ ปสฺเสยฺย, เทฺวปิ โลกธาตุโย ปสฺเสยฺย, ติสฺโสปิ โลกธาตุโย ปสฺเสยฺย, จตสฺโสปิ โลกธาตุโย ปสฺเสยฺย, ปญฺจปิ โลกธาตุโย ปสฺเสยฺย, ทสปิ โลกธาตุโย ปสฺเสยฺย, วีสมฺปิ\footnote{วีสติมฺปิ (สี) เหฏฺฐา 277; ขุ 8.174 ปิฏฺเฐสุปิ.} โลกธาตุโย ปสฺเสยฺย, ติํสมฺปิ โลกธาตุโย ปสฺเสยฺย, จตฺตาลีสมฺปิ โลกธาตุโย ปสฺเสยฺย, ปญฺญาสมฺปิ โลกธาตุโย ปสฺเสยฺย, สตมฺปิ โลกธาตุํ ปสฺเสยฺย, สหสฺสิมฺปิ จูฬนิกํ โลกธาตุํ ปสฺเสยฺย, ทฺวิสหสฺสิมฺปิ มชฺฌิมิกํ โลกธาตุํ ปสฺเสยฺย, ติสหสฺสิมฺปิ โลกธาตุํ ปสฺเสยฺย, มหาสหสฺสิมฺปิ โลกธาตุํ ปสฺเสยฺย, ยาวตกํ ปน อากงฺเขยฺย, ตาวตกํ ปสฺเสยฺย.\csromanspacer{} เอวํ ปริสุทฺธํ ภควโต ทิพฺพจกฺขุ.\csromanspacer{} เอวํ ภควา ทิพฺเพน จกฺขุนาปิ จกฺขุมา.}
\prose{\csromanfolio{357}กถํ ภควา ปญฺญาจกฺขุนาปิ จกฺขุมา, ภควา มหาปญฺโญ ปุถุปญฺโญ หาสปญฺโญ ชวนปญฺโญ ติกฺขปญฺโญ นิพฺเพธิกปญฺโญ ปญฺญาปเภทกุสโล ปภินฺนญาโณ อธิคตปฏิสมฺภิโท จตุเวสารชฺชปฺปตฺโต ทสพลธารี ปุริสาสโภ ปุริสสีโห ปุริสนาโค ปุริสาชญฺโญ ปุริสโธรโยฺห อนนฺตญาโณ อนนฺตเตโช อนนฺตยโส อฑฺโฒ มหทฺธโน ธนวา เนตา วิเนตา อนุเนตา ปญฺญาเปตา นิชฺฌาเปตา เปกฺเขตา ปสาเทตา, โส หิ ภควา อนุปฺปนฺนสฺส มคฺคสฺส อุปฺปาเทตา, อสญฺชาตสฺส มคฺคสฺส สญฺชเนตา, อนกฺขาตสฺส มคฺคสฺส อกฺขาตา, \csromanfolio{358}มคฺคญฺญู มคฺควิทู มคฺคโกวิโท, มคฺคานุคา จ ปน เอตรหิ สาวกา วิหรนฺติ ปจฺฉา สมนฺนาคตา.}

\prose{\csromansymbolfootnote{*}{ขุ 8. 175; ขุ 9. 375 ปิฏฺเฐสุปิ.}โส หิ ภควา ชานํ ชานาติ, ปสฺสํ ปสฺสติ, จกฺขุภูโต ญาณภูโต ธมฺมภูโต พฺรหฺมภูโต วตฺตา ปวตฺตา อตฺถส นินฺเนตา อมตสฺส ทาตา ธมฺมสฺสามี ตถาคโต, นตฺถิ ตสฺส ภควโต อญฺญาตํ อทิฏฺฐํ อวิทิตํ อสจฺฉิกตํ อผสฺสิตํ ปญฺญาย.\csromanspacer{} อตีตํ อนาคตํ ปจฺจุปฺปนฺนํ อุปาทาย สพฺเพ ธมฺมา สพฺพากาเรน พุทฺธสฺส ภควโต ญาณมุเข อาปาถํ อาคจฺฉนฺติ.\csromanspacer{} ยํ กิญฺจิ เนยฺยํ นาม อตฺถิ ชานิตพฺพํ\footnote{อตฺถิ ธมฺมํ ชานิตพฺพํ (ก) 277 ปิฏฺเฐ.} อตฺตตฺโถ วา ปรตฺโถ วา อุภยตฺโถ วา, ทิฏฺฐธมฺมิโก วา อตฺโถ สมฺปรายิโก วา อตฺโถ, อุตฺตาโน วา อตฺโถ คมฺภีโร วา อตฺโถ, คูโฬฺห วา อตฺโถ ปฏิจฺฉนฺโน วา อตฺโถ, เนยฺโย วา อตฺโถ นีโต วา อตฺโถ, อนวชฺโช วา อตฺโถ นิกฺกิเลโส วา อตฺโถ, โวทาโน วา อตฺโถ ปรมตฺโถ วา อตฺโถ, สพฺพํ ตํ อนฺโตพุทฺธญาเณ ปริวตฺตติ.}

\prose{อตีเต พุทฺธสฺส ภควโต อปฺปฏิหตํ ญาณํ, อนาคเต อปฺปฏิหตํ ญาณํ, ปจฺจุปฺปนฺเน อปฺปฏิหตํ ญาณํ, สพฺพํ กายกมฺมํ พุทฺธสฺส ภควโต ญาณานุปริวตฺติ.\csromanspacer{} สพฺพํ วจีกมฺมํ.\csromanspacer{} สพฺพํ มโนกมฺมํ พุทฺธสฺส ภควโต ญาณานุปริวตฺติ.\csromanspacer{} ยาวตกํ เนยฺยํ, ตาวตกํ ญาณํ.\csromanspacer{} ยาวตกํ ญาณํ, ตาวตกํ เนยฺยํ.\csromanspacer{} เนยฺยปริยนฺติกํ ญาณํ, ญาณปริยนฺติกํ เนยฺยํ.\csromanspacer{} เนยฺยํ อติกฺกมิตฺวา ญาณํ นปฺปวตฺตติ, ญาณํ อติกฺกมิตฺวา เนยฺยปโถ นตฺถิ, อญฺญมญฺญปริยนฺตฏฺฐายิโน เต ธมฺมา.\csromanspacer{} ยถา ทฺวินฺนํ สมุคฺคปฏลานํ สมฺมาผุสิตานํ เหฏฺฐิมํ สมุคฺคปฏลํ อุปริมํ นาติวตฺตติ, อุปริมํ สมุคฺคปฏลํ เหฏฺฐิมํ นาติวตฺตติ, อญฺญมญฺญปริยนฺตฏฺฐายิโน.\csromanspacer{} เอวเมวํ พุทฺธสฺส ภควโต เนยฺยญฺจ ญาณญฺจ อญฺญมญฺญปริยนฺตฏฺฐายิโน.\csromanspacer{} ยาวตกํ เนยฺยํ, ตาวตกํ ญาณํ.\csromanspacer{} ยาวตกํ ญาณํ, ตาวตกํ เนยฺยํ.\csromanspacer{} เนยฺยปริยนฺติกํ ญาณํ, ญาณปริยนฺติกํ เนยฺยํ.\csromanspacer{} เนยฺยํ อติกฺกมิตฺวา ญาณํ น ปวตฺตติ, ญาณํ อติกฺกมิตฺวา เนยฺยปโถ นตฺถิ, อญฺญมญฺญปริยนฺตฏฺฐายิโน เต ธมฺมา.}

\prose{สพฺพธมฺเมสุ พุทฺธสฺส ภควโต ญาณํ ปวตฺตติ สพฺเพ ธมฺมา พุทฺธสฺส ภควโต อาวชฺชนปฏิพทฺธา อากงฺขปฏิพทฺธา มนสิการปฏิพทฺธา จิตฺตุปฺปาทปฏิพทฺธา.\csromanspacer{} สพฺพสตฺเตสุ พุทฺธสฺส ภควโต ญาณํ ปวตฺตติ, สพฺเพสํ}

\vspace{\tipitakaparskip}
\csromancenter{สาริปุตฺตสุตฺตนิทฺเทส สตฺตานํ ภควา อาสยํ ชานาติ, อนุสยํ ชานาติ, จริตํ ชานาติ, อธิมุตฺติํ ชานาติ, อปฺปรชกฺเข มหารชกฺเข ติกฺขินฺทฺริเย มุทินฺทฺริเย สฺวากาเร ทฺวากาเร สุวิญฺญาปเย ทุวิญฺญาปเย ภพฺพาภพฺเพ สตฺเต ปชานาติ.\csromanspacer{} สเทวโก โลโก สมารโก สพฺรหฺมโก สสฺสมณพฺราหฺมณี ปชา สเทวมนุสฺสา อนฺโตพุธธญาเณ ปริวตฺตติ.}
\prose{\csromanfolio{359}ยถา เย เกจิ มจฺฉกจฺฉปา อนฺตมโส ติมิติมิงฺคลํ อุปาทาย อนฺโตมหาสมุทฺเท ปริวตฺตนฺติ.\csromanspacer{} เอวเมวํ สเทวโก โลโก สมารโก สพฺรหฺมโก สสฺสมณพฺราหฺมณี ปชา สเทวมนุสฺสา อนฺโตพุทฺธญาเณ ปริวตฺตติ.\csromanspacer{}\csromanspacer{}ยถา เย เกจิ ปกฺขี อนฺตมโส ครุฬฺหํ เวนเตยฺยํ อุปาทาย อากาสสฺส ปเทเส ปริวตฺตนฺติ.\csromanspacer{} เอวเมวํ เยปิ เต สาริปุตฺตสมา ปญฺญาย, เตปิ พุทฺธญาณสฺส ปเทเส ปริวตฺตนฺติ, พุทฺธญาณํ เทวมนุสฺสานํ ปญฺญํ ผริตฺวา อภิภวิตฺวา ติฏฺฐติ.\csromanspacer{} เยปิ เต ขตฺติยปณฺฑิตา พฺราหฺมณปณฺฑิตา คหปติปณฺฑิตา สมณปณฺฑิตานิปุณา กตปรปฺปวาทา วาลเวธิรูปา โวภินฺทนฺตา มญฺเญ จรนฺติ ปญฺญาคเตน ทิฏฺฐิคตานิ.\csromanspacer{} เต ปเญฺห อภิสงฺขริตฺวา อภิสงฺขริตฺวา ตถาคตํ อุปสงฺกมิตฺวา ปุจฺฉนฺติ คูฬฺหานิ จ ปฏิจฺฉนฺนานิ จ, กถิตา วิสชฺชิตาว เต ปญฺหา ภควตา โหนฺติ นิทฺทิฏฺฐการณา.\csromanspacer{} อุปกฺขิตฺตกา จ เต ภควโต สมฺปชฺชนฺติ, อถ โข ภควาว ตตฺถ อติโรจติ ยทิทํ ปญฺญายาติ.\csromanspacer{} เอวํ ภควา ปญฺญาจกฺขุนาปิ จกฺขุมา.}

\prose{กถํ ภควา พุทฺธจกฺขุนาปิ จกฺขุมา, ภควา พุทฺธจกฺขุนา โลกํ โวโลเกนฺโต\footnote{โอโลเกนฺโต (สี)} อทฺทส สตฺเต อปฺปรชกฺเข มหารชกฺเข ติกฺขินฺทฺริเย มุทินฺทฺริเย สฺวากาเร ทฺวากาเร สุวิญฺญาปเย ทุวิญฺญาปเย, อปฺเปกจฺเจ ปรโลกวชฺชภยทสฺสาวิโน วิหรนฺเต อปฺเปกจฺเจ นปรโลกวชฺชภยทสฺสาวิโน วิหรนฺเต, เสยฺยถาปิ นาม อุปฺปลินิยํ วา ปทุมินิยํ วา ปุณฺฑรีกินิยํ วา อปฺเปกจฺจานิ อุปฺปลานิ วา ปทุมานิ วา ปุณฺณรีกานิ วา อุทเก ชาตานิ อุทเก สํวฑฺฒานิ อุทกานุคฺคตานิ อนฺโตนิมุคฺคโปสีนิ, อปฺเปกจฺจานิ อุปฺปลานิ วา ปทุมานิ วา ปุณฺฑรีกานิ วา อุทเก ชาตานิ อุทเก สํวฑฺฒานิ สโมทกํ ฐิตานิ, อปฺเปกจฺจานิ อุปฺปลานิ วา ปทุมานิ วา ปุณฺฑรีกานิ วา อุทเก ชาตานิ อุทเก สํวฑฺฒานิ อุทกํ อจฺจุคฺคมฺม ติฏฺฐนฺติ \csromanfolio{360}อนุปลิตฺตานิ อุทเกน.\csromanspacer{} เอวเมวํ ภควา พุทฺธจกฺขุนา โลกํ โวโลเกนฺโต อทฺทส สตฺเต อปฺปรชกฺเข มหารชกฺเข ติกฺขินฺทฺริเย มุทินฺทฺริเย สฺวากาเร ทฺวากาเร สุวิญฺญาปเย ทุวิญฺญาปเย, อปฺเปกจฺเจ ปรโลกวชฺชภยทสฺสาวิโน วิหรนฺเต, อปฺเปกจฺเจ นปรโลกวชฺชภยทสฺสาวิโน วิหรนฺเต, ชานาติ ภควา “อยํ ปุคฺคโล ราคจริโต, อยํ โทสจริโต, อยํ โมหจริโต, อยํ วิตกฺกจริโต, อยํ สทฺธาจริโต, อยํ ญาณจริโต”ติ.\csromanspacer{} ราคจริตสฺส ภควา ปุคฺคลสฺส อสุภกถํ กเถติ, โทสจริตสฺส ภควา ปุคฺคลสฺส เมตฺตาภาวนํ อาจิกฺขติ, โมหจริตสฺส ภควา ปุคฺคลสฺส อุทฺเทเส ปริปุจฺฉาย กาเลน ธมฺมสฺสวเน กาเลน ธมฺมสากจฺฉาย ครุสํวาเส นิเวเสติ, วิตกฺกจริตสฺส ภควา ปุคฺคลสฺส อานาปานสฺสติํ อาจิกฺขติ, สทฺธาจริตสฺส ภควา ปุคฺคลสฺส ปสาทนียํ นิมิตฺตํ อาจิกฺขติ พุทฺธสุโพธิํ ธมฺมสุธมฺมตํ สํฆสุปฺปฏิปตฺติํ สีลานิ จ อตฺตโน, ญาณจริตสฺส ภควา ปุคฺคลสฺส วิปสฺสนานิมิตฺตํ อาจิกฺขติ อนิจฺจาการํ ทุกฺขาการํ อนตฺตาการํ.}

\csromangathameasure{*“เสเล ยถา ปพฺพตมุทฺธนิฏฺฐิโต, \\ ยถาปิ ปสฺเส ชนตํ สมนฺตโต. \\ ตถูปมํ ธมฺมมยํ สุเมธ, \\ ปาสาทมารุยฺห สมนฺตจกฺขุ. \\ โสกาวติณฺณํ ชนตมเปตโสโก, \\ อเวกฺขสฺสุ ชาติชราภิภูตนฺ”ติ.}
\csromangathagroup{\csromangathastanza{\csromansymbolfootnote{*}{ที 2. 33; ม 1. 225; ม 2. 292; สํ 1. 139; ขุ 8. 177ปิฏฺเฐสุปิ.}“เสเล ยถา ปพฺพตมุทฺธนิฏฺฐิโต, \\ ยถาปิ ปสฺเส ชนตํ สมนฺตโต. \\ ตถูปมํ ธมฺมมยํ สุเมธ, \\ ปาสาทมารุยฺห สมนฺตจกฺขุ. \\ โสกาวติณฺณํ ชนตมเปตโสโก, \\ อเวกฺขสฺสุ ชาติชราภิภูตนฺ”ติ.}}

\prose{เอวํ ภควา พุทฺธจกฺขุนาปิ จกฺขุมา.}

\prose{กถํ ภควา สมนฺตจกฺขุนาปิ จกฺขุมา, สมนฺตจกฺขุ วุจฺจติ สพฺพญฺญุตญาณํ.\csromanspacer{} ภควา สพฺพญฺญูตญาเณน อุเปโต สมุเปโต อุปคโต สมุปคโต อุปปนฺโน สมุปปนฺโน สมนฺนาคโต.}

\csromangathameasure{+“น ตสฺส อทิฏฺฐมิธตฺถิ กิญฺจิ, \\ อโถ อวิญฺญาตมชานิตพฺพํ. \\ สพฺพํ อภิญฺญาสิ ยทตฺถิ เนยฺยํ, \\ ตถาคโต เตน สมนฺตจกฺขู”ติ.}
\csromangathagroup{\csromangathastanza{\csromansymbolfootnote{+}{ขุ 8. 93, 178; ขุ 9. 127, 227 ปิฏฺเฐสุปิ.}“น ตสฺส อทิฏฺฐมิธตฺถิ กิญฺจิ, \\ อโถ อวิญฺญาตมชานิตพฺพํ. \\ สพฺพํ อภิญฺญาสิ ยทตฺถิ เนยฺยํ, \\ ตถาคโต เตน สมนฺตจกฺขู”ติ.}}

\prose{เอวํ ภควา สมนฺตจกฺขุนาปิ จกฺขุมาติ ยถา ทิสฺสติ จกฺขุมา.}

\csromanheader{2}{16. สาริปุตฺตสุตฺตนิทฺเทส}
\prose{\csromanfolio{361}\textbf{สพฺพํ ตมํ วิโนเทตฺวา}ติ สพฺพํ ราคตมํ โทสตมํ โมหตมํ มานตมํ ทิฏฺฐิตมํ กิเลสตมํ ทุจฺจริตตมํ อนฺธกรณํ อจกฺขุกรณํ อญฺญาณกรณํ ปญฺญานิโรธิกํ วิฆาตปกฺขิกํ อนิพฺพานสํวตฺตนิกํ นุทิตฺวา ปนุทิตฺวา ชหิตฺวา ปชหิตฺวา วิโนเทตฺวา พฺยนฺติํ กริตฺวา อนภาวํ คเมตฺวาติ สพฺพํ ตมํ วิโนเทตฺวา.}

\prose{\textbf{เอโกว รติมชฺฌคา}ติ \textbf{เอโก}ติ ภควา ปพฺพชฺชาสงฺขาเตน \textbf{เอโก}, อทุติยฏฺเฐน \textbf{เอโก}, ตณฺหาย ปหานฏฺเฐน \textbf{เอโก}, เอกนฺตวีตราโคติ \textbf{เอโก}, เอกนฺตวีตโทโสติ \textbf{เอโก}, เอกนฺตวีตโมโหติ \textbf{เอโก}, เอกนฺตนิกฺกิเลโสติ \textbf{เอโก}, เอกายนมคฺคํ คโตติ \textbf{เอโก}, อนุตฺตรํ สมฺมาสมฺโพธิํ อภิสมฺพุทฺโธติ \textbf{เอโก}.}

\prose{กถํ ภควา ปพฺพชฺชาสงฺขาเตน \textbf{เอโก}, ภควา ทหโรว สมาโน สุสุ กาฬเกโส ภเทฺรน\footnote{ภทฺเทน (สฺยา)} โยพฺพเนน สมนฺนาคโต ปฐเมน วยสา อกามกานํ มาตาปิตูนํ อสฺสุมุขานํ รุทนฺตานํ วิลปนฺตานํ ญาติสํฆํ ปหาย สพฺพํ ฆราวาสปลิโพธํ ฉินฺทิตฺวา ปุตฺตทารปลิโพธํ ฉินฺทิตฺวา ญาติปลิโพธํ ฉินฺทิตฺวา มิตฺตามจฺจปลิโพธํ ฉินฺทิตฺวา สนฺนิธิปลิโพธํ ฉินฺทิตฺวา เกสมสฺสุํ โอหาเรตฺวา กาสายานิ วตฺถานิ อจฺฉาเทตฺวา อคารสฺมา อนคาริยํ ปพฺพชิตฺวา อกิญฺจนภาวํ อุปคนฺตฺวา \textbf{เอโก} จรติ วิหรติ อิริยติ วตฺตติ ปาเลติ ยเปติ ยาเปติ.\csromanspacer{} เอวํ ภควา ปพฺพชฺชาสงฺขาเตน \textbf{เอโก}.}

\prose{กถํ ภควา อทุติยฏฺเฐน \textbf{เอโก}, โส เอวํ ปพฺพชิโต สมาโน \textbf{เอโก} อรญฺญวนปตฺถานิ ปนฺตานิ เสนาสนานิ ปฏิเสวติ อปฺปสทฺทานิ อปฺปนิคฺโฆสานิ วิชนวาตานิ มนุสฺสราหสฺเสยฺยกานิ ปฏิสลฺลานสารุปฺปานิ.\csromanspacer{} โส \textbf{เอโก} จรติ, \textbf{เอโก} คจฺฉติ, \textbf{เอโก} ติฏฺฐติ, \textbf{เอโก} นิสีทติ, \textbf{เอโก} เสยฺยํ กปฺเปติ, \textbf{เอโก} คามํ ปิณฺฑาย ปวิสติ, \textbf{เอโก} ปฏิกฺกมติ, \textbf{เอโก} รโห นิสีทติ, \textbf{เอโก} จงฺกมํ อธิฏฺฐาติ, \textbf{เอโก} จรติ วิหรติ อิริยติ วตฺตติ ปาเลติ ยเปติ ยาเปติ.\csromanspacer{} เอวํ ภควา อทุติยฏฺเฐน \textbf{เอโก}.}

\prose{กถํ ภควา ตณฺหาย ปหานฏฺเฐน \textbf{เอโก}, โส เอวํ \textbf{เอโก} อทุติโย อปฺปมตฺโต อาตาปี ปหิตตฺโต วิหรนฺโต นชฺชา}

\prose{\csromanfolio{362}เนรญฺชราย ตีเร โพธิรุกฺขมูเล มหาปธานํ ปทหนฺโต มารํ สเสนํ กณฺหํ นมุจิํ ปมตฺตพนฺธุํ วิธมิตฺวา ตณฺหาชาลินิํ วิสริตํ วิสตฺติกํ ปชหสิ วิโนเทสิ พฺยนฺติํ อกาสิ อนภาวํ คเมสิ.}

\csromangathameasure{“ตณฺหาทุติโย ปุริโส, ทีฆมทฺธานสํสรํ. \\ อิตฺถภาวญฺญถาภาวํ, สํสารํ นาติวตฺตติ. \\ เอตมาทีนวํ ญตฺวา, ตณฺหํ ทุกฺขสฺส สมฺภวํ. \\ วีตตโณฺห อนาทาโน, สโต ภิกฺขุ ปริพฺพเช”ติ.}
\csromangathagroup{\csromangathastanza{“ตณฺหาทุติโย ปุริโส, ทีฆมทฺธานสํสรํ. \\ อิตฺถภาวญฺญถาภาวํ\footnote{อิตฺถํ ภาวญฺญาถาภาวํ (ก) ขุ 1. 201, 268; อํ 1. 316; ขุ 8. 211 ปิฏฺฐาทีสุ. 2. อตริํสุ (สฺยา)}, สํสารํ นาติวตฺตติ.}\csromangathastanza{เอตมาทีนวํ ญตฺวา, ตณฺหํ ทุกฺขสฺส สมฺภวํ. \\ วีตตโณฺห อนาทาโน, สโต ภิกฺขุ ปริพฺพเช”ติ.}}

\prose{เอวํ ภควา ตณฺหาย ปหานฏฺเฐน เอโก.}

\prose{กถํ ภควา เอกนฺตวีตราโคติ เอโก, ราคสฺส ปหีนตฺตา เอกนฺตวีตราโคติ เอโก, โทสสฺส ปหีนตฺตา เอกนฺตวีตโทโสติ เอโก, โมหสฺส ปหีนตฺตา เอกนฺตวีตโมโหติ เอโก, กิเลสานํ ปหีนตฺตา เอกนฺตนิกฺกิเลโสติ เอโก.}

\prose{กถํ ภควา เอกายนมคฺคํ คโตติ เอโก, เอกายนมคฺโค วุจฺจติ จตฺตาโร สติปฏฺฐานา จตฺตาโร สมฺมปฺปธานา จตฺตาโร อิทฺธิปาทา ปญฺจินฺทฺริยานิ ปญฺจ พลานิ สตฺต โพชฺฌงฺคา อริโย อฏฺฐงฺคิโก มคฺโค.}

\csromangathameasure{*“เอกายนํ ชาติขยนฺตทสฺสี, \\ มคฺคํ ปชานาติ หิตานุกมฺปี. \\ เอเตน มคฺเคน ตริํสุ2 ปุพฺเพ, \\ ตริสฺสนฺติ เย จ ตรนฺติ โอฆนฺ”ติ.}
\csromangathagroup{\csromangathastanza{\csromansymbolfootnote{*}{สํ 3. 162; ขุ 8. 212, 231 ปิฏฺเฐสุ.}“เอกายนํ ชาติขยนฺตทสฺสี, \\ มคฺคํ ปชานาติ หิตานุกมฺปี.}\csromangathastanza{เอเตน มคฺเคน ตริํสุ2 ปุพฺเพ, \\ ตริสฺสนฺติ เย จ ตรนฺติ โอฆนฺ”ติ.}}

\prose{เอวํ ภควา เอกายนมคฺคํ คโตติ เอโก.}

\prose{กถํ ภควา เอโก อนุตฺตรํ สมฺมาสมฺ\textbf{โพธิ}ํ อภิสมฺพุทฺโธติ เอโก, \textbf{โพธิ} วุจฺจติ จตูสุ มคฺเคสุ ญาณํ, ปญฺญินฺทฺริยํ ปญฺญาพลํ ธมฺมวิจยสมฺโพชฺฌงฺโค วีมํสา วิปสฺสนา สมฺมาทิฏฺฐิ.\csromanspacer{} ภควา เตน \textbf{โพธิ}ญาเณน “สพฺเพ สงฺขารา อนิจฺจา”ติ พุชฺฌิ, “สพฺเพ สงฺขารา ทุกฺขา”ติ พุชฺฌิ, “สพฺเพ ธมฺมา อนตฺตา”ติ พุชฺฌิ, “อวิชฺชาปจฺจยา สงฺขารา”ติ พุชฺฌิ ฯเปฯ “ชาติปจฺจยา ชรามรณนฺ”ติ พุชฺฌิ. “อวิชฺชานิโรธา สงฺขารนิโรโธ”ติ พุชฺฌิ ฯเปฯ “ชาตินิโรธา ชรามรณนิโรโธ”ติ พุชฺฌิ, “อิทํ ทุกฺขนฺ”ติ พุชฺฌิ, “อยํ ทุกฺขสมุทโย”ติ พุชฺฌิ, “อยํ ทุกฺขนิโรโธ”ติ พุชฺฌิ, “อยํ}

\vspace{\tipitakaparskip}
\csromancenter{สาริปุตฺตสุตฺตนิทฺเทส ทุกฺขนิโรธคามินี ปฏิปทา”ติ พุชฺฌิ, “อิเม อาสวา”ติ พุชฺฌิ ฯเปฯ “อยํ อาสวนิโรธคามินี ปฏิปทา”ติ พุชฺฌิ, “อิเม ธมฺมา ปริญฺเญยฺยา”ติ พุชฺฌิ, ปหาตพฺพาติ.\csromanspacer{} ภาเวตพฺพาติ.\csromanspacer{} สจฺฉิกาตพฺพาติ พุชฺฌิ.\csromanspacer{} ฉนฺนํ ผสฺสายตนานํ สมุทยญฺจ อตฺถงฺคมญฺจ อสฺสาทญฺจ อาทีนวญฺจ นิสฺสรณญฺจ พุชฺฌิ, ปญฺจนฺนํ อุปาทานกฺขนฺธานํ สมุทยญฺจ อตฺถงฺคมญฺจ อสฺสาทญฺจ อาทีนวญฺจ นิสฺสรณญฺจ พุชฺฌิ, จตุนฺนํ มหาภูตานํ สมุทยญฺจ อตฺถงฺคมญฺจ อสฺสาทญฺจ อาทีนวญฺจ นิสฺสรณญฺจ พุชฺฌิ, “ยํ กิญฺจิ สมุทยธมฺมํ, สพฺพํ ตํ นิโรธธมฺมนฺ”ติ พุชฺฌิ.}
\prose{\csromanfolio{363}อถ วา ยํ กิญฺจิ พุชฺฌิตพฺพํ อนุพุชฺฌิตพฺพํ ปฏิพุชฺฌิตพฺพํ สมฺพุชฺฌิตพฺพํ อธิคนฺตพฺพํ ผสฺสิตพฺพํ สจฺฉิกาตพฺพํ, สพฺพํ ตํ เตน โพธิญาเณน พุชฺฌิ อนุพุชฺฌิ ปฏิพุชฺฌิ สมฺพุชฺฌิ สมฺมาพุชฺฌิ อธิคจฺฉิ ผสฺเสสิ สจฺฉากาสิ.\csromanspacer{} เอวํ ภควา เอโก อนุตฺตรํ สมฺมาสมฺโพธิํ อภิสมฺพุทฺโธติ เอโก.}

\prose{\textbf{รติมชฺฌคา}ติ \textbf{รตินฺ}ติ เนกฺขมฺมรติํ วิเวกรติํ อุปสมรติํ สมฺโพธิรติํ อชฺฌคา สมชฺฌคา อธิคจฺฉิ ผสฺเสสิ สจฺฉากาสีติ เอโกว รติมชฺฌคา.\csromanspacer{} เตนาห เถโร สาริปุตฺโต\csromandash{}}

\csromangathameasure{สเทวกสฺส โลกสฺส, ยถา ทิสฺสติ จกฺขุมา. \\ สพฺพํ ตมํ วิโนเทตฺวา, เอโกว รติมชฺฌคาติ.}
\csromangathagroup{\csromangathastanza{สเทวกสฺส โลกสฺส, ยถา ทิสฺสติ จกฺขุมา. \\ สพฺพํ ตมํ วิโนเทตฺวา, เอโกว รติมชฺฌคาติ.}}

\proseitem{192}{\csromansymbolfootnote{*}{ขุ 1. 426 ปิฏฺเฐปิ.}\textbf{ตํ พุทฺธํ อสิตํ ตาทิ}ํ, อกุหํ คณิมาคตํ.}

\csromangathameasure{พหูนมิธ พทฺธานํ, อตฺถิ ปเญฺหน อาคมํ.}
\csromangathagroup{\csromangathastanza{พหูนมิธ พทฺธานํ, อตฺถิ ปเญฺหน อาคมํ.}}

\prose{\textbf{ตํ พุทฺธํ อสิตํ ตาทิ}นฺติ \textbf{พุทฺโธ}ติ โย โส ภควา สยมฺภู อนาจริยโก ปุพฺเพ อนนุสฺสุเตสุ ธมฺเมสุ สามํ สจฺจานิ อภิสมฺพุชฺฌิ, ตตฺถ จ สพฺพญฺญุตํ ปาปุณิ\footnote{ปตฺโฏ (สฺยา)} พเลสุ จ วสีภาวํ, \textbf{พุทฺโธ}ติ เกนฏฺเฐน \textbf{พุทฺโธ}, พุชฺฌิตา สจฺจานีติ \textbf{พุทฺโธ}, โพเธตา ปชายาติ \textbf{พุทฺโธ}, สพฺพญฺญุตาย \textbf{พุทฺโธ}, สพฺพทสฺสาวิตาย \textbf{พุทฺโธ}, อนญฺญเนยฺยตาย \textbf{พุทฺโธ}, วิสวิตาย \textbf{พุทฺโธ}, ขีณาสวสงฺขาเตน \textbf{พุทฺโธ}, นิรุปเลปสงฺขาเตน\footnote{นิรุปกฺกิเลสสงฺขาเตน (พหูสุ) ขุ 8. 197; ขุ 9. 172 ปิฏฺเฐสุปิ.} \textbf{พุทฺโธ}, เอกนฺตวีตราโคติ \textbf{พุทฺโธ}, เอกนฺตวีตโทโสติ \textbf{พุทฺโธ}, เอกนฺตวีตโมโหติ \textbf{พุทฺโธ}, เอกนฺตนิกฺกิเลโสติ \textbf{พุทฺโธ}, เอกายนมคฺคํ คโตติ \textbf{พุทฺโธ}, เอโก อนุตฺตรํ สมฺมาสมฺโพธิํ อภิสมฺพุทฺโธติ \textbf{พุทฺโธ}, อพุทฺธิวิหตตฺตา พุทฺธิปฏิลาภตฺตา \textbf{พุทฺโธ}.\csromanspacer{} \textbf{พุทฺโธ}ติ เนตํ นามํ มาตรา กตํ, น ปิตรา กตํ, น ภาตรา กตํ, น ภคินิยา กตํ, น มิตฺตามจฺเจหิ กตํ, น \csromanfolio{364}ญาติสาโลหิเตหิ กตํ, น สมณพฺราหฺมเณหิ กตํ, น เทวตาหิ กตํ, วิโมกฺขนฺติกเมตํ พุทฺธานํ ภควนฺตานํ โพธิยา มูเล สห สพฺพญฺญุตญาณสฺส ปฏิลาภา สจฺฉิกา ปญฺญตฺติ ยทิทํ พุทฺโธติ ตํ พุทฺธํ.\csromanspacer{}\csromanspacer{}\textbf{อสิตนฺ}ติ เทฺว นิสฺสยา ตณฺหานิสฺสโย จ ทิฏฺฐินิสฺสโย จ.\csromanspacer{} กตโม ตณฺหานิสฺสโย, ยาว ตณฺหาสงฺขาเตน สีมกตํ โอธิกตํ\footnote{มริยาทิกตํ (สฺยา)} ปริยนฺตกตํ ปริคฺคหิตํ มมายิตํ, อิทํ มม, เอตํ มม, เอตฺตกํ มม, เอตฺตาวตา มม, มม รูปา สทฺทา คนฺธา รสา โผฏฺฐพฺพา, อตฺถรณา ปาวุรณา ทาสิทาสา อเชฬกา กุกฺกุฏสูกรา หตฺถิควาสฺสวฬวา เขตฺตํ วตฺถุํ หิรญฺญํ สุวณฺณํ คามนิคมราชธานิโย รฏฺฐญฺจ ชนปโท จ โกโส จ โกฏฺฐาคารญฺจ, เกวลมฺปิ มหาปถวิํ ตณฺหาวเสน มมายติ, ยาวตา อฏฺฐสตํ ตณฺหาวิจริตํ, อยํ ตณฺหานิสฺสโย.}

\prose{กตโม ทิฏฺฐินิสฺสโย, วีสติวตฺถุกา สกฺกายทิฏฺฐิ, ทสวตฺถุกา มิจฺฉาทิฏฺฐิ, ทสวตฺถุกา อนฺตคฺคาหิกา ทิฏฺฐิ.\csromanspacer{} ยา เอวรูปา ทิฏฺฐิ ทิฏฺฐิคตํ ทิฏฺฐิคหนํ ทิฏฺฐิกนฺตาโร ทิฏฺฐิวิสุกายิกํ ทิฏฺฐิวิปฺผนฺทิตํ ทิฏฺฐิสญฺโญชนํ คาโห ปฏิคฺคาโห อภินิเวโส ปรามาโส กุมฺมคฺโค มิจฺฉาปโถ มิจฺฉตฺตํ ติตฺถายตนํ วิปริเยสคฺคาโห วิปรีตคฺคาโห วิปลฺลาสคฺคาโห มิจฺฉาคาโห อยาถาวกสฺมิํ ยาถาวกนฺติ คาโห ยาวตา ทฺวาสฏฺฐิ ทิฏฺฐิคตานิ, อยํ ทิฏฺฐินิสฺสโย.}

\prose{พุทฺธสฺส ภควโต ตณฺหานิสฺสโย ปหีโน, ทิฏฺฐินิสฺสโย ปฏินิสฺสฏฺโฐ, ตณฺหานิสฺสยสฺส ปหีนตฺตา, ทิฏฺฐินิสฺสยสฺส ปฏินิสฺสฏฺฐตฺตา ภควา จกฺขุํ อสิโต, โสตํ.\csromanspacer{} ฆานํ.\csromanspacer{} ชิวฺหํ.\csromanspacer{} กายํ.\csromanspacer{} มนํ อสิโต.\csromanspacer{} รูเป.\csromanspacer{} สทฺเท.\csromanspacer{} คนฺเธ.\csromanspacer{} รเส.\csromanspacer{} โผฏฺฐพฺเพ.\csromanspacer{} กุลํ.\csromanspacer{} คณํ.\csromanspacer{} อาวาสํ.\csromanspacer{} ลาภํ.\csromanspacer{} ยสํ.\csromanspacer{} ปสํสํ.\csromanspacer{} สุขํ.\csromanspacer{} จีวรํ.\csromanspacer{} ปิณฺฑปาตํ.\csromanspacer{} เสนาสนํ.\csromanspacer{} คิลานปจฺจยเภสชฺช- ปริกฺขารํ.\csromanspacer{} กามธาตุํ.\csromanspacer{} รูปธาตุํ.\csromanspacer{} อรูปธาตุํ.\csromanspacer{} กามภวํ.\csromanspacer{} รูปภวํ.\csromanspacer{} อรูปภวํ.\csromanspacer{} สญฺญาภวํ.\csromanspacer{} อสญฺญาภวํ.\csromanspacer{} เนวสญฺญานาสญฺญาภวํ.\csromanspacer{} เอกโวการภวํ.\csromanspacer{} จตุโวการภวํ.\csromanspacer{} ปญฺจโวการภวํ.\csromanspacer{} อตีตํ.\csromanspacer{} อนาคตํ.\csromanspacer{} ปจฺจุปฺปนฺนํ.\csromanspacer{} ทิฏฺฐสุตมุตวิญฺญาตพฺเพ ธมฺเม อสิโต อนิสฺสิโต อนลฺลีโน อนุปคโต อนชฺโฌสิโต อนธิมุตฺโต นิกฺขนฺโต นิสฺสโฏ วิปฺปมุตฺโต วิสญฺญุตฺโต วิมริยาทิกเตน เจตสา วิหรตีติ ตํ พุทฺธํ อสิตํ.}

\csromanheader{2}{16. สาริปุตฺตสุตฺตนิทฺเทส}
\prose{\csromanfolio{365}\textbf{ตาทินฺ}ติ ภควา ปญฺจหากาเรหิ ตาที, อิฏฺฐานิฏฺเฐ ตาที, จตฺตาวีติ ตาที, ติณฺณาวีติ ตาที, มุตฺตาวีติ ตาที, ตํ นิทฺเทสา ตาที.}

\prose{กถํ ภควา อิฏฺฐานิฏฺเฐ ตาที, ภควา ลาเภปิ ตาที, อลาเภปิ ตาที, ยเสปิ ตาที, อยเสปิ ตาที, ปสํสายปิ ตาที, นินฺทายปิ ตาที, สุเขปิ ตาที, ทุกฺเขปิ ตาที, เอกจฺเจ พาหํ คนฺเธน ลิมฺเปยฺยุํ, เอกจฺเจ พาหํ วาสิยา ตจฺเฉยฺยุํ, อมุกสฺมิํ นตฺถิ ราโค, อมุกสฺมิํ นตฺถิ ปฏิฆํ, อนุนยปฏิฆวิปฺปหีโน อุคฺฆาตินิคฺฆาติวีติวตฺโต อนุโรธวิโรธํ สมติกนฺโต, เอวํ ภควา อิฏฺฐานิฏฺเฐ ตาที.}

\prose{กถํ ภควา จตฺตาวีติ ตาที, ภควโต\footnote{ภควตา (สี, ก)} ราโค จตฺโต วนฺโต มุตฺโต ปหีโน ปฏินิสฺสฏฺโฐ.\csromanspacer{} โทโส.\csromanspacer{} โมโห.\csromanspacer{} โกโธ.\csromanspacer{} อุปนาโห.\csromanspacer{} มกฺโข.\csromanspacer{} ปฬาโส.\csromanspacer{} อิสฺสา.\csromanspacer{} มจฺฉริยํ.\csromanspacer{} มายา.\csromanspacer{} สาเฐยฺยํ.\csromanspacer{} ถมฺโภ.\csromanspacer{} สารมฺโภ.\csromanspacer{} มาโน.\csromanspacer{} อติมาโน.\csromanspacer{} มโท.\csromanspacer{} ปมาโท.\csromanspacer{} สพฺเพ กิเลสา.\csromanspacer{} สพฺเพ ทุจฺจริตา.\csromanspacer{} สพฺเพ ทรถา.\csromanspacer{} สพฺเพ ปริฬาหา.\csromanspacer{} สพฺเพ สนฺตาปา.\csromanspacer{} สพฺพากุสลาภิสงฺขารา จตฺตา วนฺตา มุตฺตา ปหีนา ปฏินิสฺสฏฺฐา.\csromanspacer{} เอวํ ภควา จตฺตาวีติ ตาที.}

\prose{กถํ ภควา ติณฺณาวีติ ตาที, ภควา กาโมฆํ ติณฺโณ, ภโวฆํ ติณฺโณ, ทิฏฺโฐฆํ ติณฺโณ, อวิชฺโชฆํ ติณฺโณ, สพฺพํ สํสารปถํ ติณฺโณ อุตฺติณฺโณ นิตฺติณฺโณ อติกฺกนฺโต สมติกฺกนฺโต วีติวตฺโต, โส วุฏฺฐวาโส จิณฺณจรโณ คตทฺโธ คตทิโส คตโกฏิโย ปาลิตพฺรหฺมจริโย อุตฺตมทิฏฺฐิปฺปตฺโต ภาวิตมคฺโค ปหีนกิเลโส ปฏิวิทฺธากุปฺโป สจฺฉิกตนิโรโธ ทุกฺขํ ตสฺส ปริญฺญาตํ, สมุทโย ปหีโน, มคฺโค ภาวิโต, นิโรโธ สจฺฉิกโต, อภิญฺเญยฺยํ อภิญฺญาตํ, ปริญฺเญยฺยํ ปริญฺญาตํ, ปหาตพฺพํ ปหีนํ, ภาเวตพฺพํ ภาวิตํ, สจฺฉิกาตพฺพํ สจฺฉิกตํ, โส อุกฺขิตฺตปลิโฆ\footnote{อุกฺขิตฺตปฬิโฆ (สี) 16 ปิฏฺเฐปิ.} สํกิณฺณปริกฺโข อพฺพูเฬฺหสิโก นิรคฺคโฬ อริโย ปนฺนทฺธโช ปนฺนภาโร วิสญฺญุตฺโต ปญฺจงฺควิปฺปหีโน ฉฬงฺคสมนฺนาคโต เอการกฺโข จตุราปสฺเสโน ปณุนฺน ปจฺเจกสจฺโจ สมวยสฏฺเฐสโน อนาวิลสงฺกปฺโป ปสฺสทฺธกายสงฺขาโร สุวิมุตฺตจิตฺโต สุวิมุตฺตปญฺโญ เกวลี วุสิตวา อุตฺตมปุริโส ปรมปุริโส ปรมปตฺติปฺปตฺโต, โส เนวาจินติ, นาปจินติ, อปจินิตฺวา}

\prose{\csromanfolio{366}ฐิโต, เนว ปชหติ, น อุปาทิยติ, ปชหิตฺวา ฐิโต, เนว สํสิพฺพติ, น อุสิเนติ, วิสิเนตฺวา ฐิโต, เนว วิธูเปติ น สนฺธูเปติ วิธูเปตฺวา ฐิโต, อเสเขน สีลกฺขนฺเธน สมนฺนาคตตฺตา ฐิโต, อเสเขน สมาธิกฺขนฺเธน สมนฺนาคตตฺตา ฐิโต, อเสเขน ปญฺญากฺขนฺเธน สมนฺนาคตตฺตา ฐิโต, อเสเขน วิมุตฺติกฺขนฺเธน สมนฺนาคตตฺตา ฐิโต, อเสเขน วิมุตฺติญาณทสฺสนกฺขนฺเธน สมนฺนาคตตฺตา ฐิโต, สจฺจํ สมฺปฏิปาทิยิตฺวา ฐิโต, เอชํ สมติกฺกมิตฺวา ฐิโต, กิเลสคฺคิํ ปริยาทิยิตฺวา ฐิโต, อปริคมนตาย ฐิโต, กฏํ สมาทาย ฐิโต, มุตฺติปฏิเสวนตาย ฐิโต, เมตฺตาย ปาริสุทฺธิยา ฐิโต, กรุณาย ปาริสุทฺธิยา ฐิโต, มุทิตาย ปาริสุทฺธิยา ฐิโต, อุเปกฺขาย ปาริสุทฺธิยา ฐิโต, อจฺจนฺตปาริสุทฺธิยา ฐิโต, อตมฺมยตาย ปาริสุทฺธิยา ฐิโต, วิมุตฺตตฺตา ฐิโต, สนฺตุสิตตฺตา ฐิโต, ขนฺธปริยนฺเต ฐิโต, ธาตุปริยนฺเต ฐิโต, อายตนปริยนฺเต ฐิโต, คติปริยนฺเต ฐิโต, อุปปตฺติปริยนฺเต ฐิโต, ปฏิสนฺธิปริยนฺเต ฐิโต, ภวปริยนฺเต ฐิโต, สํสารปริยนฺเต ฐิโต, วฏฺฏปริยนฺเต ฐิโต, อนฺติเม ภเว ฐิโต, อนฺติเม สมุสฺสเย ฐิโต, อนฺติมเทหธโร ภควา\csromandash{}}

\csromangathameasure{ตสฺสายํ ปจฺฉิมโก ภโว, จริโมยํ สมุสฺสโย. \\ ชาติมรณสํสาโร, นตฺถิ ตสฺส ปุนพฺภโวติ.}
\csromangathagroup{\csromangathastanza{ตสฺสายํ ปจฺฉิมโก ภโว, จริโมยํ สมุสฺสโย. \\ ชาติมรณสํสาโร, นตฺถิ ตสฺส ปุนพฺภโวติ.}}

\prose{เอวํ ภควา ติณฺณาวีติ ตาที.}

\prose{กถํ ภควา มุตฺตาวีติ ตาที, ภควโต ราคา จิตฺตํ มุตฺตํ วิมุตฺตํ สุวิมุตฺตํ, โทสา จิตฺตํ.\csromanspacer{} โมหา จิตฺตํ.\csromanspacer{} โกธา.\csromanspacer{} อุปนาหา.\csromanspacer{} มกฺขา.\csromanspacer{} ปฬาสา.\csromanspacer{} อิสฺสาย.\csromanspacer{} มจฺฉริยา.\csromanspacer{} มายาย.\csromanspacer{} สาเฐยฺยา.\csromanspacer{} ถมฺภา.\csromanspacer{} สารมฺภา.\csromanspacer{} มานา.\csromanspacer{} อติมานา.\csromanspacer{} มทา.\csromanspacer{} ปมาทา.\csromanspacer{} สพฺพกิเลเสหิ.\csromanspacer{} สพฺพทุจฺจริเตหิ.\csromanspacer{} สพฺพทรเถหิ.\csromanspacer{} สพฺพปริฬาเหหิ.\csromanspacer{} สพฺพสนฺตาเปหิ.\csromanspacer{} สพฺพากุสลาภิสงฺขาเรหิ จิตฺตํ มุตฺถํ วิมุตฺตํ สุวิมุตฺตํ, เอวํ ภควา มุตฺตาวีติ ตาที.}

\prose{กถํ ภควา ตํนิทฺเทสา ตาที, ภควา สีเล สติ สีลวาติ ตํนิทฺเทสา ตาที, สทฺธาย สติ สทฺโธติ ตํนิทฺเทสา ตาที, วีริเย สติ วีริยวาติ ตํนิทฺเทสา ตาที, สติยา สติ สติมาติ ตํนิทฺเทสฺสา ตาที, สมาธิสฺมิํ สติ สมาหิโตติ ตํนิทฺเทสา ตาทิ, ปญฺญาย สติ ปญฺญวาติ ตํนิทฺเทสา ตาที, วิชฺชาย สติ เตวิชฺโชติ ตํนิทฺเทสา ตาที,}

\vspace{\tipitakaparskip}
\csromancenter{สาริปุตฺตสุตฺตนิทฺเทส อภิญฺญาย สติ ฉฬภิญฺโญติ ตํนิทฺเทสา ตาที, ทสพเล สติ ทสพโลติ ตํนิทฺเทสา ตาที, เอวํ ภควา ตํนิทฺเทสา ตาทีติ ตํ พุทฺธํ อสิตํ ตาทิํ.}
\prose{\csromanfolio{367}\textbf{อกุหํ คณิมาคตนฺ}ติ \textbf{อกุโห}ติ ตีณิ กุหนวตฺถูนิ ปจฺจยปฏิเสวนสงฺขาตํ กุหนวตฺถุ, อิริยาปถสงฺขาตํ กุหนวตฺถุ, สามนฺตชปฺปนสงฺขาตํ กุหนวตฺถุ.}

\prose{กตมํ ปจฺจยปฏิเสวนสงฺขาตํ กุหนวตฺถุ, อิธ คหปติกา ภิกฺขุํ นิมนฺเตนฺติ จีวรปิณฺฑปาตเสนาสนคิลานปจฺจยเภสชฺช- ปริกฺขาเรหิ, โส ปาปิจฺโฉ อิจฺฉาปกโต อตฺถิโก\footnote{อติตฺติโก (สี), อิตฺถิโก (สฺยา), เหฏฺฐา 172 ปิฏฺเฐ ปาฐนานตฺตํ นตฺถิ.} จีวรปิณฺฑปาตเสนาสน- คิลานปจฺจยเภสชฺชปริกฺขารานํ ภิยฺโยกมฺยตํ อุปาทาย จีวรํ ปจฺจกฺขาติ, ปิณฺฑปาตํ ปจฺจกฺขาติ, เสนาสนํ ปจฺจกฺขาติ, คิลานปจฺจยเภสชฺชปริกฺขารํ ปจฺจกฺขาติ.\csromanspacer{} โส เอวมาห “กิํ สมณสฺส มหคฺเฆน จีวเรน, เอตํ สารุปฺปํ, ยํ สมโณ สุสานา วา สงฺการกูฏา วา ปาปณิกา วา นนฺตกานิ อุจฺจินิตฺวา สงฺฆาฏิํ กริตฺวา ธาเรยฺย.\csromanspacer{} กิํ สมณสฺส มหคฺเฆน ปิณฺฑปาเตน, เอตํ สารุปฺปํ, ยํ สมโณ อุญฺฉาจริยาย ปิณฺฑิยาโลเปน ชิวิตํ กปฺเปยฺย.\csromanspacer{} กิํ สมณสฺส มหคฺเฆน เสนาสเนน, เอตํ สารุปฺปํ, ยํ สมโณ รุกฺขมูลิโก วา อสฺส โสสานิโก วา อพฺโภกาสิโก วา.\csromanspacer{} กิํ สมณสฺส มหคฺเฆน คิลานปจฺจยเภสชฺชปริกฺขาเรน, เอตํ สารุปฺปํ, ยํ สมโณ ปูติมุตฺเตน หริตกีขณฺเฑน โอสธํ กเรยฺยา”ติ, ตทุปาทาย ลูขํ จีวรํ ธาเรติ, ลูขํ ปิณฺฑปาตํ ปริภุญฺชติ, ลูขํ เสนาสนํ ปฏิเสวติ, ลูขํ คิลานปจฺจยเภสชฺชปริกฺขารํ ปฏิเสวติ.\csromanspacer{} ตเมนํ คหปติกา เอวํ ชานนฺติ “อยํ สมโณ อปฺปิจฺโฉ สนฺตุฏฺโฐ ปวิวิตฺโต อสํสฏฺโฐ อารทฺธวีริโย ธุตวาโท”ติ ภิยฺโย ภิยฺโย นิมนฺเตนฺติ จีวรปิณฺฑปาตเสนาสนคิลานปจฺจยเภสชฺชปริกฺขาเรหิ, โส เอวมาห “ติณฺณํ สมฺมุขีภาวา สทฺโธ กุลปุตฺโต พหุํ ปุญฺญํ ปสวติ, สทฺธาย สมฺมุขีภาวา สทฺโธ กุลปุตฺโต พหุํ ปุญฺญํ ปสวติ, เทยฺยธมฺมสฺส สมฺมุขีภาวา สทฺโธ กุลปุตฺโต พหุํ ปุญฺญํ ปสวติ, ทกฺขิเณยฺยานํ สมฺมุขีภาวา สทฺโธ กุลปุตฺโต พหุํ ปุญฺญํ ปสวติ.\csromanspacer{} ตุมฺหากญฺเจวายํ สทฺธา อตฺถิ, เทยฺยธมฺโม จ สํวิชฺชติ, อหญฺจ ปฏิคฺคาหโก.\csromanspacer{} สเจ อหํ น \csromanfolio{368}ปฏิคฺคเหสฺสามิ, เอวํ ตุเมฺห ปุญฺเญน ปริพาหิรา\footnote{ปฏิพาหิรา (ก)} ภวิสฺสถ, น มยฺหํ อิมินา อตฺโถ, อปิ จ ตุมฺหากํเยว อนุกมฺปาย ปฏิคฺคณฺหามี”ติ, ตทุปาทาย พหุมฺปิ จีวรํ ปฏิคฺคณฺหาติ\footnote{ปฏิคณฺหติ (สี)}, พหุมฺปิ ปิณฺฑปาตํ ปฏิคฺคณฺหาติ, พหุมฺปิ เสนาสนํ ปฏิคฺคณฺหาติ, พหุมฺปิ คิลานปจฺจยเภสชฺชปริกฺขารํ ปฏิคฺคณฺหาติ.\csromanspacer{} ยา เอวรูปา ภากุฏิกา ภากุฏิยํ กุหนา กุหายนา กุหิตตฺตํ, อิทํ ปจฺจยปฏิเสวนสงฺขาตํ กุหนวตฺถุ.}

\prose{กตมํ อิริยาปถสงฺขาตํ กุหนวตฺถุ, อิเธกจฺโจ ปาปิจฺโฉ อิจฺฉาปกโต สมฺภาวนาธิปฺปาโย “เอวํ มํ ชโน สมฺภาเวสฺสตี”ติ คมนํ สณฺฐเปติ, ฐานํ สณฺฐเปติ, นิสชฺชํ สณฺฐเปติ, สยนํ สณฺฐเปติ, ปณิธาย คจฺฉติ, ปณิธาย ติฏฺฐติ, ปณิธาย นิสีทติ, ปณิธาย เสยฺยํ กปฺเปติ, สมาหิโต วิย คจฺฉติ, สมาหิโต วิย ติฏฺฐติ, สมาหิโต วิย นิสีทติ, สมาหิโต วิย เสยฺยํ กปฺเปติ, อาปาถกชฺฌายีว โหติ.\csromanspacer{} ยา เอวรูปา อิริยาปถสฺส อาฐปนา ฐปนา สณฺฐปนา ภากุฏิกา ภากุฏิยํ กุหนา กุหายนา กุหิตตฺตํ, อิทํ อิริยาปถสงฺขาตํ กุหนวตฺถุ.}

\prose{กตมํ สามนฺตชปฺปนสงฺขาตํ กุหนวตฺถุ, อิเธกจฺโจ ปาปิจฺโฉ อิจฺฉาปกโต สมฺภาวนาธิปฺปาโย “เอวํ มํ ชโน สมฺภาเวสฺสตี”ติ อริยธมฺมสนฺนิสฺสิตํ วาจํ ภาสติ “โย เอวรูปํ จีวรํ ธาเรติ, โส สมโณ มเหสกฺโข”ติ ภณติ, “โย เอวรูปํ ปตฺตํ ธาเรติ.\csromanspacer{} โลหถาลกํ ธาเรติ.\csromanspacer{} ธมฺมกรณํ\footnote{ธมฺมกรกํ (สี, สฺยา) อภิธานปฺปทีปิกาภินวนิสฺสยํ โอโลเกตพฺพํ.} ธาเรติ.\csromanspacer{} ปริสฺสาวนํ ธาเรติ.\csromanspacer{} กุญฺจิกํ ธาเรติ.\csromanspacer{} อุปาหนํ ธาเรติ.\csromanspacer{} กายพนฺธนํ ธาเรติ.\csromanspacer{} อาโยคํ ธาเรติ, โส สมโณ มเหสกฺโข”ติ ภณติ, “ยสฺส เอวรูโป อุปชฺฌาโย, โส สมโณ มเหสกฺโข”ติ ภณติ, “ยสฺส เอวรูโป อาจริโย.\csromanspacer{} เอวรูปา สมานุปชฺฌายกา.\csromanspacer{} สมานาจริยกา.\csromanspacer{} มิตฺตา.\csromanspacer{} สนฺทิฏฺฐา.\csromanspacer{} สมฺภตฺตา.\csromanspacer{} สหายา, โส สมโณ มเหสกฺโข”ติ ภณติ, “โย เอวรูเป วิหาเร วสติ, โส สมโณ มเหสกฺโข”ติ ภณติ, โย เอวรูเป อฑฺฒโยเค วสติ.\csromanspacer{} ปาสาเท วสติ.\csromanspacer{} หมฺมิเย วสติ.\csromanspacer{} คุหายํ วสติ.\csromanspacer{} เลเณ วสติ.\csromanspacer{} กุฏิยา วสติ.\csromanspacer{} กูฏาคาเร วสติ.\csromanspacer{} อฏฺเฏ วสติ.\csromanspacer{} มาเฬ วสติ.\csromanspacer{} อุทฺทณฺเฑ วสติ.\csromanspacer{} อุปฏฺฐานสาลายํ วสติ.}

\vspace{\tipitakaparskip}
\csromancenter{สาริปุตฺตสุตฺตนิทฺเทส มณฺฑเป วสติ.\csromanspacer{} รุกฺขมูเล วสติ, โส สมโณ มเหสกฺโข”ติ ภณติ.}
\prose{\csromanfolio{369}อถ วา โกรชิกโกรชิโก ภากุฏิกภากุฏิโก กุหกกุหโก ลปกลปโก มุขสมฺภาวิโต อยํ สมโณ อิมาสํ เอวรูปานํ วิหารสมาปตฺตีนํ ลาภีติ ตาทิสํ คมฺภีรํ คูฬฺหํ นิปุณํ ปฏิจฺฉนฺนํ โลกุตฺตรํ สุญฺญตาปฏิสญฺญุตฺตํ กถํ กเถติ.\csromanspacer{} ยา เอวรูปา ภากุฏิกา ภากุฏิยํ กุหนา กุหายนา กุหิตตฺตํ, อิทํ สามนฺตชปฺปนสงฺขาตํ กุหนวตฺถุ.\csromanspacer{} พุทฺธสฺส ภควโต อิมานิ ตีณิ กุหนวตฺถูนิ ปหีนานิ สมุจฺฉินฺนานิ วูปสนฺตานิ ปฏิปสฺสทฺธานิ อภพฺพุปฺปตฺติกานิ ญาณคฺคินา ทฑฺฒานิ.\csromanspacer{} ตสฺมา พุทฺโธ อกุโหติ อกุหํ.}

\prose{\textbf{คณิมาคตนฺ}ติ \textbf{คณี}ติ \textbf{คณี} ภควา, คณาจริโยติ \textbf{คณี}, คณสฺส สตฺถาติ \textbf{คณี}, คณํ ปริหรตีติ \textbf{คณี}, คณํ โอวทตีติ \textbf{คณี}, คณํ อนุสาสตีติ \textbf{คณี}, วิสารโท คณํ อุปสงฺกมตีติ \textbf{คณี}, คณ’สฺส สุสฺสูสติ โสตํ โอทหติ อญฺญา จิตฺตํ อุปฏฺฐเปตีติ \textbf{คณี}, คณํ อกุสลา วุฏฺฐาเปตฺวา กุสเล ปติฏฺฐาเปตีติ \textbf{คณี}, ภิกฺขุคณสฺส \textbf{คณี}, ภิกฺขุนีคณสฺส \textbf{คณี}, อุปาสกคณสฺส \textbf{คณี}, อุปาสิกาคณสฺส \textbf{คณี}, ราชคณสฺส \textbf{คณี}, ขตฺติยคณสฺส \textbf{คณี}, พฺราหฺมณคณสฺส \textbf{คณี}, เวสฺสคณสฺส \textbf{คณี}, สุทฺทคณสฺส \textbf{คณี}, พฺรหฺมคณสฺส \textbf{คณี}, เทวคณสฺส \textbf{คณี}, สงฺฆิํ\footnote{สํฆคณสฺส คณี (สี)} คณิํ คณาจริยํ.\csromanspacer{} \textbf{อาคตนฺ}ติ อุปคตํ สมุปคตํ สมุปปนฺนํ\footnote{สมฺปตฺตํ (พหูสุ)} สงฺกสฺสนครนฺติ อกุหํ คณิมาคตํ.}

\prose{\textbf{พหูนมิธ พทฺธานนฺ}ติ พหูนํ ขตฺตินํ พฺราหฺมณานํ เวสฺสานํ สุทฺทานํ คหฏฺฐานํ ปพฺพชิตานํ เทวานํ มนุสฺสานํ.\csromanspacer{} \textbf{พทฺธานนฺ}ติ พทฺธานํ พทฺธจรานํ ปริจารกานํ สิสฺสานนฺติ พหูนมิธ พทฺธานํ.}

\prose{\textbf{อตฺถิ ปเญฺหน อาคมนฺ}ติ * ปเญฺหน อตฺถิโก อาคโตมฺหิ,\footnote{อตฺถิกามฺห อาคตา (สี, ก) เอวมีทิเสสุ ทฺวีสุ ปเทสุปิ พหุวจเนน.} ปญฺหํ ปุจฺฉิตุกาโม อาคโตมฺหิ, ปญฺหํ โสตุกาโม อาคโตมฺหีติ เอวมฺปิ อตฺถิ ปเญฺหน อาคมํ.\csromanspacer{}\csromanspacer{}อถ วา ปญฺหตฺถิกานํ ปญฺหํ ปุจฺฉิตุกามานํ ปญฺหํ โสตุกามานํ อาคมนํ อภิกฺกมนํ อุปสงฺกมนํ ปยิรุปาสนํ\footnote{ปยิรุปาสนา (สฺยา, ก)} สิยา อตฺถีติ เอวมฺปิ อตฺถิ ปเญฺหน อาคมํ.\csromanspacer{}\csromanspacer{}อถ วา ปญฺหาคโม ตุยฺหํ อตฺถิ, \csromanfolio{370}ตฺวมฺปิปหุ\footnote{อิมานิ ตีณิ ปทานิ นตฺถิ สฺยา-โปตฺถเก.}, ตฺวมสิ อลมตฺโถ มยา ปุจฺฉิตํ กเถตุํ วิสชฺเชตุํ.\csromanspacer{} วหสฺเสตํ ภารนฺติ เอวมฺปิ อตฺถิ ปเญฺหน อาคมํ.\csromanspacer{} เตนาห เถโร สาริปุตฺโต\csromandash{}}

\csromangathameasure{ตํ พุทฺธํ อสิตํ ตาทิํ, อกุหํ คณิมาคตํ. \\ พหูนมิธ พทฺธานํ, อตฺถิ ปเญฺหน อาคมนฺติ.}
\csromangathagroup{\csromangathastanza{ตํ พุทฺธํ อสิตํ ตาทิํ, อกุหํ คณิมาคตํ. \\ พหูนมิธ พทฺธานํ, อตฺถิ ปเญฺหน อาคมนฺติ.}}

\proseitem{193}{\csromansymbolfootnote{*}{ขุ 1. 426 ปิฏฺเฐปิ.}\textbf{ภิกฺขุโน วิชิคุจฺฉโต}, ภชโต ริตฺตมาสนํ.}

\csromangathameasure{รุกฺขมูลํ สุสานํ วา, ปพฺพตานํ คุหาสุ วา.}
\csromangathagroup{\csromangathastanza{รุกฺขมูลํ สุสานํ วา, ปพฺพตานํ คุหาสุ วา.}}

\prose{\textbf{ภิกฺขุโน วิชิคุจฺฉโต}ติ \textbf{ภิกฺขุโน}ติ ปุถุชฺชนกลฺยาณสฺส วา \textbf{ภิกฺขุโน} เสกฺขสฺส วา \textbf{ภิกฺขุโน}.\csromanspacer{} \textbf{วิชิคุจฺฉโต}ติ ชาติยา วิชิคุจฺฉโต, ชราย.\csromanspacer{} พฺยาธินา.\csromanspacer{} มรเณน.\csromanspacer{} โสเกหิ.\csromanspacer{} ปริเทเวหิ.\csromanspacer{} ทุกฺเขหิ.\csromanspacer{} โทมนสฺเสหิ.\csromanspacer{} อุปายาเสหิ วิชิคุจฺฉโต, เนรยิเกน ทุกฺเขน.\csromanspacer{} ติรจฺฉานโยนิเกน ทุกฺเขน.\csromanspacer{} เปตฺติวิสยิเกน ทุกฺเขน.\csromanspacer{} มานุสิเกน\footnote{มานุสเกน (สี, สฺยา) 319 ปิฏฺเฐ.} ทุกฺเขน.\csromanspacer{} คพฺโภกฺกนฺติมูลเกน ทุกฺเขน.\csromanspacer{} คพฺภฏฺฐิติมูลเกน ทุกฺเขน.\csromanspacer{} คพฺภวุฏฺฐานมูลเกน ทุกฺเขน.\csromanspacer{} ชาตสฺสูปนิพนฺธเกน ทุกฺเขน.\csromanspacer{} ชาตสฺส ปราเธยฺยเกน ทุกฺเขน.\csromanspacer{} อตฺตูปกฺกเมน ทุกฺเขน.\csromanspacer{} ปรูปกฺกเมน ทุกฺเขน.\csromanspacer{} ทุกฺขทุกฺเขน.\csromanspacer{} สงฺขารทุกฺเขน.\csromanspacer{} วิปริณามทุกฺเขน.\csromanspacer{} จกฺขุโรเคน ทุกฺเขน.\csromanspacer{} โสตโรเคน ทุกฺเขน.\csromanspacer{} ฆานโรเคน ทุกฺเขน.\csromanspacer{} ชิวฺหาโรเคน ทุกฺเขน.\csromanspacer{} กายโรเคน ทุกฺเขน.\csromanspacer{} สีลโรเคน ทุกฺเขน.\csromanspacer{} กณฺณโรเคน ทุกฺเขน.\csromanspacer{} มุขโรเคน ทุกฺเขน.\csromanspacer{} ทนฺตโรเคน ทุกฺเขน.\csromanspacer{} กาเสน.\csromanspacer{} สาเสน.\csromanspacer{} ปินาเสน.\csromanspacer{} qอาเหน.\csromanspacer{} ชเรน.\csromanspacer{} กุจฺฉิโรเคน.\csromanspacer{} มุจฺฉาย.\csromanspacer{} ปกฺขนฺทิกาย.\csromanspacer{} สูลาย.\csromanspacer{} วิสูจิกาย.\csromanspacer{} กุฏฺเฐน.\csromanspacer{} คณฺเฑน.\csromanspacer{} กิลาเสน.\csromanspacer{} โสเสน.\csromanspacer{} อปมาเรน.\csromanspacer{} ททฺทุยา กณฺฑุยา.\csromanspacer{} กจฺฉุยา.\csromanspacer{} รขสาย.\csromanspacer{} วิตจฺฉิกาย.\csromanspacer{} โลหิเตน.\csromanspacer{} ปิตฺเตน.\csromanspacer{} มธุเมเหน.\csromanspacer{} อํสาย.\csromanspacer{} ปิฬกาย.\csromanspacer{} ภคนฺทลาย.\csromanspacer{} ปิตฺตสมุฏฺฐาเนน อาพาเธน.\csromanspacer{} เสมฺหสมุฏฺฐาเนน อาพาเธน.\csromanspacer{} วาตสมุฏฺฐาเนน อาพาเธน.\csromanspacer{} สนฺนิปาติเกน อาพาเธน.\csromanspacer{} อุตุปริณามเชน อาพาเธน.\csromanspacer{} วิสมปริหารเชน อาพาเธน.\csromanspacer{} โอปกฺกมิเกน อาพาเธน.\csromanspacer{} กมฺมวิปากเชน อาพาเธน.\csromanspacer{} สีเตน.\csromanspacer{} อุเณฺหน.\csromanspacer{} ชิฆจฺฉาย.\csromanspacer{} ปิปาสาย.\csromanspacer{} อุจฺจาเรน.\csromanspacer{} ปสฺสาเวน.\csromanspacer{} qอํสมกสวาตาตปสรีสปสมฺผสฺเสน ทุกฺเขน.\csromanspacer{} มาตุมรเณน ทุกฺเขน.\csromanspacer{} ปิตุมรเณน ทุกฺเขน.\csromanspacer{} ภาตุมรเณน.\csromanspacer{} ภคินิมรเณน.}

\vspace{\tipitakaparskip}
\csromancenter{สาริปุตฺตสุตฺตนิทฺเทส ปุตฺตมรเณน.\csromanspacer{} ธีตุมรเณน.\csromanspacer{} ญาติพฺยสเนน.\csromanspacer{} โภคพฺยสเนน.\csromanspacer{} โรคพฺยสเนน.\csromanspacer{} สีลพฺยสเนน.\csromanspacer{} ทิฏฺฐิพฺยสเนน ทุกฺเขน วิชิคุจฺฉโต อฏฺฏียโต หรายโต ชิคุจฺฉโตติ ภิกฺขุโน วิชิคุจฺฉโต.}
\prose{\csromanfolio{371}\textbf{ภชโต ริตฺตมาสนนฺ}ติ \textbf{อาสนํ} วุจฺจติ ยตฺถ นิสีทติ.\csromanspacer{} มญฺโจ ปีฐํ ภิสิ ตฏฺฏิกา จมฺมขณฺโฑ ติณสนฺถาโร ปณฺณสนฺถาโร ปลาลสนฺถาโร\footnote{ปลาสสนฺถาโร (สี, สฺยา)}, ตํ \textbf{อาสนํ} อสปฺปายรูปทสฺสเนน ริตฺตํ วิวิตฺตํ ปวิวิตฺตํ, อสปฺปายสทฺทสฺสวเนน ริตฺตํ วิวิตฺตํ ปวิวิตฺตํ, อสปฺปาเยหิ ปญฺจหิ กามคุเณหิ ริตฺตํ วิวิตฺตํ ปวิวิตฺตํ, ตํ ปวิวิตฺตํ \textbf{อาสนํ} ภชโต สมฺภชโต เสวโต นิเสวโต สํเสวโต ปฏิเสวโตติ ภชโต ริตฺตมฺ\textbf{อาสนํ}.}

\prose{\textbf{รุกฺขมูลํ สุสานํ วา}ติ รุกฺขมูลํเยว รุกฺขมูลํ, สุสานํเยว สุสานนฺติ รุกฺขมูลํ สุสานํ วา.\csromanspacer{} \textbf{ปพฺพตานํ คุหาสุ วา}ติ ปพฺพตาเยว ปพฺพตา, กนฺทราเยว กนฺทรา, คิริคุหาเยว คิริคุหา, ปพฺพตนฺตริกาโย วุจฺจนฺติ ปพฺพตปพฺภาราติ ปพฺพตานํ คุหาสุ วา.\csromanspacer{} เตนาห เถโร สาริปุตฺโต\csromandash{}}

\csromangathameasure{ภิกฺขุโน วิชิคุจฺฉโต, ภชโต ริตฺตมาสนํ. \\ รุกฺขมูลํ สุสานํ วา, ปพฺพตานํ คุหาสุ วาติ.}
\csromangathagroup{\csromangathastanza{ภิกฺขุโน วิชิคุจฺฉโต, ภชโต ริตฺตมาสนํ. \\ รุกฺขมูลํ สุสานํ วา, ปพฺพตานํ คุหาสุ วาติ.}}

\proseitem{194}{\csromansymbolfootnote{*}{ขุ 1. 426 ปิฏฺเฐปิ.}อุจฺจาวเจสุ สยเนสุ, \textbf{กิวนฺโต}\footnote{คีวนฺโต (สฺยา) โมคฺคลฺลานพฺยากรณํ โอโลเกตพฺพํ. 3. กุวนฺโต (สี), คีวนฺโต (สฺยา)} \textbf{ตตฺถ เภรวา.}}

\csromangathameasure{เย หิ ภิกฺขุ น เวเธยฺย, นิคฺโฆเส สยนาสเน.}
\csromangathagroup{\csromangathastanza{เย หิ ภิกฺขุ น เวเธยฺย, นิคฺโฆเส สยนาสเน.}}

\prose{\textbf{อุจฺจาวเจสุ สยเนสู}ติ \textbf{อุจฺจาวเจสู}ติ อุจฺจาวเจสุ หีนปฺปณีเตสุ เฉกปาปเกสุ.\csromanspacer{} \textbf{สยนํ} วุจฺจติ เสนฺ\textbf{อาสนํ} วิหาโร อฑฺฒโยโค ปาสาโท หมฺมิยํ คุหาติ อุจฺจาวเจสุ สยเนสุ.\csromanspacer{} \textbf{กิวนฺโต} \textbf{ตตฺถเภรวา}ติ \textbf{กิวนฺโต}ติ \textbf{กิวนฺโต}3 กูชนฺโต นทนฺโต สทฺทํ กโรนฺโต.\csromanspacer{} อถ วา \textbf{กิวนฺโต}ติ กติ กิตฺตกา กีวตกา กีวพหุกา เต.\csromanspacer{} \textbf{เภรวา}ติ สีหา พฺยคฺฆา ทีปี อจฺฉา ตรจฺฉา โกกา มหิํสา หตฺถี อหี วิจฺฉิกา สตปที โจรา วา อสฺสุ มานวา วา กตกมฺมา วา อกตกมฺมา วาติ \textbf{กิวนฺโต} \textbf{ตตฺถ เภรวา.}}

\prose{\textbf{เย หิ ภิกฺขุ น เวเธยฺยา}ติ \textbf{เย หี}ติ เภรเว ปสฺสิตฺวา วา สุณิตฺวา วา น เวเธยฺย นปฺปเวเธยฺย น สมฺปเวเธยฺย น ตเสยฺย น \csromanfolio{372}อุตฺตเสยฺย น ปริตฺตเสยฺย น ภาเยยฺย น สนฺตาสํ อาปชฺเชยฺย, อภีรู อสฺส อจฺฉมฺภี อนุตฺราสี อปลายี, ปหีนภยเภรโว วิคตโลมหํโส วิหเรยฺยาติ เย หิ ภิกฺขุ น เวเธยฺย.}

\prose{\textbf{นิคฺโฆเส สยนาสเน}ติ อปฺปสทฺเท อปฺปนิคฺโฆเส วิชนวาเต มนุสฺสราหสฺเสยฺยเก ปฏิสลฺลานสารุปฺเป เสนาสเนติ นิคฺโฆเส สยนาสเน.\csromanspacer{} เตนาห เถโร สาริปุตฺโต\csromandash{}}

\csromangathameasure{อุจฺจาวเจสุ สยเนสุ, กิวนฺโต ตตฺถ เภรวา. \\ เย หิ ภิกฺขุ น เวเธยฺย, นิคฺโฆเส สยนาสเนติ.}
\csromangathagroup{\csromangathastanza{อุจฺจาวเจสุ สยเนสุ, กิวนฺโต ตตฺถ เภรวา. \\ เย หิ ภิกฺขุ น เวเธยฺย, นิคฺโฆเส สยนาสเนติ.}}

\csromanhangingbegin
\hangingheaditem{195}{กติ ปริสฺสายา โลเก, คจฺฉโต อคตํ ทิสํ.}
\hangingbody{\textbf{เย ภิกฺขุ อภิสมฺภเว, ปนฺตมฺหิ สยนาสเน.}}
\csromanhangingend

\prose{\textbf{กติ ปริสฺสยา โลเก}ติ \textbf{กตี}ติ กติ กิตฺตกา กีวตกา กีวพหุกา.\csromanspacer{} \textbf{ปริสฺสยา}ติ เทฺว ปริสฺสยา ปากฏปริสฺสยา จ ปฏิจฺฉนฺนปริสฺสยา จ.\csromanspacer{} กตเม ปากฏปริสฺสยา, สีหา พฺยคฺฆา ทีปี อจฺฉา ตรจฺฉา โกกา มหิํสา หตฺถี อหี วิจฺฉิกา สตปที โจรา วา อสฺสุ มานวา วา กตกมฺมา วา อกตกมฺมา วา, จกฺขุโรโค โสตโรโค ฆานโรโค ชิวฺหาโรโค กายโรโค สีสโรโค กณฺณโรโค มุขโรโค ทนฺตโรโค กาโส สาโส ปินาโส ฑาโห ชโร กุจฺฉิโรโค มุจฺฉา ปกฺขนฺทิกา สูลา วิสูจิกา กุฏฺฐํ คณฺโฑ กิลาโส โสโส อปมาโร ททฺทุ กณฺฑุ กจฺฉุ รขสา วิตจฺฉิกา โลหิตํ ปิตฺตํ มธุเมโห อํสา ปิฬกา ภคนฺทลา ปิตฺตสมุฏฺฐานา อาพาธา ฯเปฯ สีตํ อุณฺหํ ชิฆจฺฉา ปิปาสา อุจฺจาโร ปสฺสาโว ฑํสมกสวาตาตปสรีสปสมฺผสฺสา, อิเม วุจฺจนฺติ ปากฏปริสฺสยา.}

\prose{กตเม ปฏิจฺฉนฺนปริสฺสยา, กายทุจฺจริตํ วจีทุจฺจริตํ มโนทุจฺจริตํ กามจฺฉนฺทนีวรณํ พฺยาปาทนีวรณํ ถินมิทฺธนีวรณํ อุทฺธจฺจกุกฺกุจฺจนีวรณํ วิจิกิจฺฉานีวรณํ ราโค โทโส โมโห โกโธ อุปนาโห มกฺโข ปฬาโส อิสฺสา มจฺฉริยํ มายา สาเฐยฺยํ ถมฺโภ สารมฺโภ มาโน อติมาโน มโท ปมาโท สพฺเพ กิเลสา สพฺเพ ทุจฺจริตา สพฺเพ ทรถา สพฺเพ ปริฬาหา สพฺเพ สนฺตาปา สพฺพากุสลาภิสงฺขารา, อิเม วุจฺจนฺติ ปฏิจฺฉนฺนปริสฺสยา.}

\csromanheader{2}{16. สาริปุตฺตสุตฺตนิทฺเทส}
\prose{\csromanfolio{373}\textbf{ปริสฺสยา}ติ เกนฏฺเฐน ปริสฺสยา, ปริสหนฺตีติ ปริสฺสยา, ปริหานาย สํวตฺตนฺตีติ ปริสฺสยา, ตตฺราสยาติ ปริสฺสยา.}

\prose{กถํ ปริสหนฺตีติ ปริสฺสยา, เต ปริสฺสยา ตํ ปุคฺคลํ สหนฺติ ปริสหนฺติ อภิภวนฺติ อชฺโฌตฺถรนฺติ ปริยาทิยนฺติ มทฺทนฺติ.\csromanspacer{} เอวํ ปริสหนฺตีติ ปริสฺสยา.}

\prose{กถํ ปริหานาย สํวตฺตนฺตีติ ปริสฺสยา, เต ปริสฺสยา กุสลานํ ธมฺมานํ อนฺตรายาย ปริหานาย สํวตฺตนฺติ.\csromanspacer{} กตเมสํ กุสลานํ ธมฺมานํ, สมฺมาปฏิปทาย อนุโลมปฏิปทาย อปจฺจนีกปฏิปทาย อวิรุทฺธปฏิปทาย อนฺวตฺถปฏิปทาย ธมฺมานุธมฺมปฏิปทาย สีเลสุ ปริปูรการิตาย อินฺทฺริเยสุ คุตฺตทฺวารตาย โภชเน มตฺตญฺญุตาย ชาคริยานุโยคสฺส สติสมฺปชญฺญสฺส จตุนฺนํ สติปฏฺฐานานํ ภาวนานุโยคสฺส จตุนฺนํ สมฺมปฺปธานานํ จตุนฺนํ อิทฺธิปาทานํ ปญฺจนฺนํ อินฺทฺริยานํ ปญฺจนฺนํ พลานํ สตฺตนฺนํ โพชฺฌงฺคานํ อริยสฺส อฏฺฐงฺคิกสฺส มคฺคสฺส ภาวนานุโยคสฺส, อิเมสํ กุสลานํ ธมฺมานํ อนฺตรายาย ปริหานาย สํวตฺตนฺติ.\csromanspacer{} เอวํ ปริหานาย สํวตฺตนฺตีติ ปริสฺสยา.}

\prose{กถํ ตตฺราสยาติ ปริสฺสยา, ตตฺเถเต ปาปกา อกุสลา ธมฺมา อุปฺปชฺชนฺติ อตฺตภาวสนฺนิสฺสยา, ยถา พิเล พิลาสยา ปาณา สยนฺติ, ทเก ทกาสยา ปาณา สยนฺติ, วเน วนาสยา ปาณา สยนฺติ, รุกฺเข รุกฺขาสยา ปาณา สยนฺติ.\csromanspacer{} เอวเมวํ ตตฺเถเต ปาปกา อกุสลา ธมฺมา อุปฺปชฺชนฺติ อตฺตภาวสนฺนิสฺสยาติ เอวมฺปิ ตตฺราสยาติ ปริสฺสยา.}

\prose{วุตฺตํ เหตํ ภควตา\csromandash{} * สานฺเตวาสิโก ภิกฺขเว ภิกฺขุ สาจริยโก ทุกฺขํ น ผาสุ วิหรติ.\csromanspacer{} กถญฺจ ภิกฺขเว ภิกฺขุ สานฺเตวาสิโก สาจริยโก ทุกฺขํ น ผาสุ วิหรติ, อิธ ภิกฺขเว ภิกฺขุโน จกฺขุนา รูปํ ทิสฺวา อุปฺปชฺชนฺติ ปาปกา อกุสลา ธมฺมา สรสงฺกปฺปา สญฺโญชนิยา, ตฺยสฺส อนฺโต วสนฺติ อนฺวาสฺสวนฺติ ปาปกา อกุสลา ธมฺมาติ ตสฺมา “สานฺเตวาสิโก”ติ วุจฺจติ.\csromanspacer{} เต นํ สมุทาจรนฺติ, สมุทาจรนฺติ นํ ปาปกา อกุสลา ธมฺมาติ, ตสฺมา “สาจริยโก”ติ วุจฺจติ.}

\prose{ปุน จปรํ ภิกฺขเว ภิกฺขุโน โสเตน สทฺทํ สุตฺวา ฯเปฯ ฆาเนน คนฺธํ ฆายิตฺวา.\csromanspacer{} ชิวฺหาย รสํ สายิตฺวา.\csromanspacer{} กาเยน โผฏฺฐพฺพํ ผุสิตฺวา.\csromanspacer{} มนสา ธมฺมํ วิญฺญาย อุปฺปชฺชนฺติ ปาปกา อกุสลา ธมฺมา สรสงฺกปฺปา สญฺโญชนิยา,}

\prose{\csromanfolio{374}ตฺยสฺส อนฺโต วสนฺติ อนฺวาสฺสวนฺติ ปาปกา อกุสลา ธมฺมาติ ตสฺมา “สานฺเตวาสิโก”ติ วุจฺจติ.\csromanspacer{} เต นํ สมุทาจรนฺติ, สมุทาจรนฺติ นํ ปาปกา อกุสลา ธมฺมาติ, ตสฺมา “สาจริยโก”ติ วุจฺจติ.\csromanspacer{} เอวํ โข ภิกฺขเว ภิกฺขุ สานฺเตวาสิโก สาจริยโก ทุกฺขํ น ผาสุ วิหรตีติ เอวมฺปิ ตตฺราสยาติ ปริสฺสยา.}

\prose{วุตฺตํ เหตํ ภควตา\csromandash{} * ตโยเม ภิกฺขเว อนฺตรามลา อนฺตรา-อมิตฺตา อนฺตราสปตฺตา อนฺตราวธกา อนฺตราปจฺจตฺถิกา.\csromanspacer{} กตเม ตโย, โลโภ ภิกฺขเว อนฺตรามโล อนฺตรา-อมิตฺโต อนฺตราสปตฺโต อนฺตราวธโก อนฺตราปจฺจตฺถิโก, โทโส ภิกฺขเว ฯเปฯ โมโห ภิกฺขเว อนฺตรามโล อนฺตรา-อมิตฺโต อนฺตราสปตฺโต อนฺตราวธโก อนฺตราปจฺจตฺถิโก.\csromanspacer{} อิเม โข ภิกฺขเว ตโย อนฺตรามลา อนฺตรา-อมิตฺตา อนฺตราสปตฺตา อนฺตราวธกา อนฺตราปจฺจตฺถิกา.}

\csromangathameasure{*อนตฺถชนโน โลโภ, โลโภ จิตฺตปฺปโกปโน. \\ ภยมนฺตรโต ชาตํ, ตํ ชโน นาวพุชฺฌติ. \\ ลุทฺโธ อตฺถํ น ชานาติ, ลุทฺโธ ธมฺมํ น ปสฺสติ. \\ อนฺธนฺตมํ ตทา โหติ, ยํ โลโภ สหเต นรํ. \\ อนตฺถชนโน โทโส, โทโส จิตฺตปฺปโกปโน. \\ ภยมนฺตรโต ชาตํ, ตํ ชโน นาวพุชฺฌติ. \\ กุทฺโธ อตฺถํ น ชานาติ, กุทฺโธ ธมฺมํ น ปสฺสติ. \\ อนฺธนฺตมํ ตทา โหติ, ยํ โทโส สหเต นรํ. \\ อนตฺถชนโน โมโห, โมโห จิตฺตปฺปโกปโน. \\ ภยมนฺตรโต ชาตํ, ตํ ชโน นาวพุชฺฌติ. \\ มูโฬฺห อตฺถํ น ชานาติ, มูโฬฺห ธมฺมํ น ปสฺสติ. \\ อนฺธนฺตมํ ตทา โหติ, ยํ โมโห สหเต นรนฺติ.}
\csromangathagroup{\csromangathastanza{\csromansymbolfootnote{*}{ขุ 1. 252; ขุ 8. 251 ปิฏฺเฐสุปิ.}อนตฺถชนโน โลโภ, โลโภ จิตฺตปฺปโกปโน. \\ ภยมนฺตรโต ชาตํ, ตํ ชโน นาวพุชฺฌติ.}\csromangathastanza{ลุทฺโธ อตฺถํ น ชานาติ, ลุทฺโธ ธมฺมํ น ปสฺสติ. \\ อนฺธนฺตมํ\footnote{อนฺธตมํ (สฺยา, ก) 283 ปิฏฺเฐ; ขุ 1. 252; อํ 2. 471 ปิฏฺเฐสุปิ.} ตทา โหติ, ยํ โลโภ สหเต นรํ.}\csromangathastanza{อนตฺถชนโน โทโส, โทโส จิตฺตปฺปโกปโน. \\ ภยมนฺตรโต ชาตํ, ตํ ชโน นาวพุชฺฌติ.}\csromangathastanza{กุทฺโธ อตฺถํ น ชานาติ, กุทฺโธ ธมฺมํ น ปสฺสติ. \\ อนฺธนฺตมํ ตทา โหติ, ยํ โทโส สหเต นรํ.}\csromangathastanza{อนตฺถชนโน โมโห, โมโห จิตฺตปฺปโกปโน. \\ ภยมนฺตรโต ชาตํ, ตํ ชโน นาวพุชฺฌติ.}\csromangathastanza{มูโฬฺห อตฺถํ น ชานาติ, มูโฬฺห ธมฺมํ น ปสฺสติ. \\ อนฺธนฺตมํ ตทา โหติ, ยํ โมโห สหเต นรนฺติ.}}

\prose{เอวมฺปิ ตตฺราสยาติ ปริสฺสยา.}

\prose{วุตฺตํ เหตํ ภควตา\csromandash{} + ตโย โข มหาราช ปุริสสฺส อชฺฌตฺตํ อุปฺปชฺชมานา อุปฺปชฺชนฺติ อหิตาย ทุกฺขาย อผาสุวิหาราย.\csromanspacer{} กตเม}

\vspace{\tipitakaparskip}
\csromancenter{สาริปุตฺตสุตฺตนิทฺเทส ตโย, โลโภ โข มหาราช ปุริสสฺส ธมฺโม อชฺฌตฺตํ อุปฺปชฺชมาโน อุปฺปชฺชติ อหิตาย ทุกฺขาย อผาสุวิหาราย, โทโส โข มหาราช ฯเปฯ โมโห โข มหาราช ปุริสสฺส ธมฺโม อชฺฌตฺตํ อุปฺปชฺชมาโน อุปฺปชฺชติ อหิตาย ทุกฺขาย อผาสุวิหาราย.\csromanspacer{} อิเม โข มหาราช ตโย ปุริสสฺส ธมฺมา อชฺฌตฺตํ อุปฺปชฺชมานา อุปฺปชฺชนฺติ อหิตาย ทุกฺขาย อผาสุวิหารายาติ.}
\vspace{0.5\baselineskip}
\csromangathameasure{*โลโภ โทโส จ โมโห จ, ปุริสํ ปาปเจตสํ. \\ หิํสนฺติ อตฺตสมฺภูตา, ตจสารํว สมฺผลนฺติ.}
\csromangathagroup{\csromangathastanza{\csromanfolio{375}\csromansymbolfootnote{*}{ขุ 1. 226; ขุ 8. 252 ปิฏฺเฐสุปิ.}โลโภ โทโส จ โมโห จ, ปุริสํ ปาปเจตสํ. \\ หิํสนฺติ อตฺตสมฺภูตา, ตจสารํว สมฺผลนฺติ.}}

\prose{เอวมฺปิ ตตฺราสยาติ ปริสฺสยา.}

\csromanheader{2}{วุตฺตํ เหตํ ภควตา\csromandash{}}
\csromangathameasure{+“ราโค จ โทโส จ อิโตนิทานา, \\ อรตี รตี โลมหํโส อิโตชา. \\ อิโต สมุฏฺฐาย มโนวิตกฺกา, \\ กุมารกา ธงฺกมิโวสฺสชนฺตี”ติ.}
\csromangathagroup{\csromangathastanza{\csromansymbolfootnote{+}{สํ 1. 209; ขุ 1. 320; ขุ 8. 252; ขุ 10. 126 ปิฏฺฐาทีสุปิ.}“ราโค จ โทโส จ อิโตนิทานา, \\ อรตี รตี โลมหํโส อิโตชา. \\ อิโต สมุฏฺฐาย มโนวิตกฺกา, \\ กุมารกา ธงฺกมิโวสฺสชนฺตี”ติ.}}

\prose{เอวมฺปิ ตตฺราสยาติ ปริสฺสยา.\csromanspacer{}\csromanspacer{}\textbf{โลเก}ติ มนุสฺสโลเกติ กติ ปริสฺสยา โลเก.}

\prose{\textbf{คจฺฉโต อคตํ ทิสนฺ}ติ อคตา ทิสา วุจฺจติ อมตํ นิพฺพานํ.\csromanspacer{} โย โส สพฺพสงฺขารสมโถ สพฺพูปธิปฏินิสฺสคฺโค ตณฺหกฺขโย วิราโค นิโรโธ นิพฺพานํ, อคตปุพฺพา สา ทิสา, น สา ทิสา คตปุพฺพา อิมินา ทีเฆน อทฺธุนา.}

\csromangathameasure{สมติตฺติกํ อนวเสสํ, เตลปตฺตํ ยถา ปริหเรยฺย. \\ เอวํ สจิตฺตมนุรกฺเข, ปตฺถยาโน ทิสํ อคตปุพฺพํ.}
\csromangathagroup{\csromangathastanza{สมติตฺติกํ อนวเสสํ, เตลปตฺตํ ยถา ปริหเรยฺย. \\ เอวํ สจิตฺตมนุรกฺเข, ปตฺถยาโน\footnote{ปตฺถยมาโน (สฺยา)} ทิสํ อคตปุพฺพํ.}}

\prose{อคตปุพฺพํ ทิสํ วชโต คจฺฉโต อภิกฺกมโตติ คจฺฉโต อคตํ ทิสํ.}

\prose{\textbf{เย ภิกฺขุ อภิสมฺภเว}ติ \textbf{เย}ติ \textbf{เย} ปริสฺส\textbf{เย} อภิสมฺภเวยฺย อภิภเวยฺย อชฺโฌตฺถเรยฺย ปริยาทิ\textbf{เย}ยฺย มทฺเทยฺยาติ \textbf{เย} ภิกฺขุ อภิสมฺภเว.}

\prose{\csromanfolio{376}\textbf{ปนฺตมฺหิ สยนาสเน}ติ อนฺเต ปนฺเต ปริยนฺเต เสลนฺเต วา วนนฺเต วา นทนฺเต วา อุทกนฺเต วา ยตฺถ น กสียติ น วปียติ, ชนนฺตํ อติกฺกมิตฺวา มนุสฺสานํ อนุปจาเร เสนาสเนติ ปนฺตมฺหิ สยนาสเน.\csromanspacer{} เตนาห เถโร สาริปุตฺโต\csromandash{}}

\csromangathameasure{กติ ปริสฺสยา โลเก, คจฺฉโต อคตํ ทิสํ. \\ เย ภิกฺขุ อภิสมฺภเว, ปนฺตมฺหิ สยนาสเนติ.}
\csromangathagroup{\csromangathastanza{กติ ปริสฺสยา โลเก, คจฺฉโต อคตํ ทิสํ. \\ เย ภิกฺขุ อภิสมฺภเว, ปนฺตมฺหิ สยนาสเนติ.}}

\proseitem{196}{\csromansymbolfootnote{*}{ขุ 1. 426 ปิฏฺเฐปิ.}กฺยาสฺส พฺยปฺปถโย อสฺสุ, กฺยาสฺสสฺสุ อิธ โคจรา.}

\csromangathameasure{กานิ สีลพฺพตานาสฺสุ, ปหิตตฺตสฺส ภิกฺขุโน.}
\csromangathagroup{\csromangathastanza{กานิ สีลพฺพตานาสฺสุ, ปหิตตฺตสฺส ภิกฺขุโน.}}

\prose{\textbf{กฺยาสฺส พฺยปฺปถโย อสฺสู}ติ กีทิเสน พฺยปฺปเถน สมนฺนาคโต อสฺส กิํสณฺฐิเตน กิํปกาเรน กิํปฏิภาเคนาติ วจีปาริสุทฺธิํ ปุจฺฉติ.\csromanspacer{} กตมา วจีปาริสุทฺธิ, + อิธ ภิกฺขุ มุสาวาทํ ปหาย มุสาวาทา ปฏิวิรโต โหติ สจฺจวาที สจฺจสนฺโธ เถโต ปจฺจยิโก อวิสํวาทโก โลกสฺส.\csromanspacer{} ปิสุณํ วาจํ ปหาย ปิสุณาย วาจาย ปฏิวิรโต โหติ, อิโต สุตฺวา น อมุตฺร อกฺขาตา อิเมสํ เภทาย, อมุตฺร วา สุตฺวา น อิเมสํ อกฺขาตา อมูสํ เภทาย, อิติ ภินฺนานํ วา สนฺธาตา, สหิตานํ วา อนุปฺปทาตา, สมคฺคาราโม สมคฺครโต สมคฺคนนฺที สมคฺคกรณิํ วาจํ ภาสิตา โหติ.\csromanspacer{} ผรุสํ วาจํ ปหาย ผรุสาย วาจาย ปฏิวิรโต โหติ.\csromanspacer{} ยา สา วาจา เนลา กณฺณสุขา เปมนียา หทยงฺคมา โปรี พหุชนกนฺตา พหุชนมนาปา, ตถารูปิํ วาจํ ภาสิตา โหติ.\csromanspacer{} สมฺผปฺปลาปํ ปหาย สมฺผปฺปลาปา ปฏิวิรโต โหติ กาลวาที ภูตวาที อตฺถวาที ธมฺมวาที วินยวาที, นิธานวติํ วาจํ ภาสิตา โหติ กาเลน สาปเทสํ ปริยนฺตวติํ อตฺถสํหิตํ.\csromanspacer{} จตูหิ วจีสุจริเตหิ สมนฺนาคโต, จตุโทสาปคตํ วาจํ ภาสติ, พาตฺติํสาย ติรจฺฉานกถาย อารโต อสฺส วิรโต ปฏิวิรโต นิกฺขนฺโต นิสฺสโฏ วิปฺปมุตฺโต วิสญฺญุตฺโต วิมริยาทิกเตน เจตสา วิหรติ.\csromanspacer{} ทส กถาวตฺถูนิ กเถติ, เสยฺยถิทํ, อปฺปิจฺฉกถํ สนฺตุฏฺฐิกถํ ปวิเวกกถํ อสํสคฺคกถํ วีริยารมฺภกถํ สีลกถํ สมาธิกถํ ปญฺญากถํ วิมุตฺติกถํ วิมุตฺติญาณทสฺสนกถํ สติปฏฺฐานกถํ สมฺมปฺปธานกถํ อิทฺธิปาทกถํ อินฺทฺริยกถํ พลกถํ โพชฺฌงฺคกถํ มคฺคกถํ ผลกถํ นิพฺพานกถํ}

\vspace{\tipitakaparskip}
\csromancenter{สาริปุตฺตสุตฺตนิทฺเทส กเถติ.\csromanspacer{} วาจาย ยโต ยตฺโต ปฏิยตฺโต คุตฺโต โคปิโต รกฺขิโต สํวุโต, อยํ วจีปาริสุทฺธิ.\csromanspacer{} เอทิสาย วจีปาริสุทฺธิยา สมนฺนาคโต อสฺสาติ กฺยาสฺส พฺยปฺปถโย อสฺสุ.}
\prose{\csromanfolio{377}\textbf{กฺยาสฺสสฺสุ อิธ โคจรา}ติ กีทิเสน โคจเรน สมนฺนาคโต อสฺส กิํสณฺฐิเตน กิํปกาเรน กิํปฏิภาเคนาติ โคจรํ ปุจฺฉติ.\csromanspacer{} อตฺถิ โคจโร, อตฺถิ อโคจโร.}

\prose{กตโม อโคจโร, อิเธกจฺโจ เวสิยาโคจโร วา โหติ, วิธวาโคจโร วา โหติ, ถุลฺลกุมารีโคจโร\footnote{ถูลกุมารีโคจโร (สฺยา, ก) อภิ 2. 256 ปิฏฺเฐปิ.} วา โหติ, ปณฺฑกโคจโร วา โหติ, ภิกฺขุนีโคจโร วา โหติ, ปานาคารโคจโร วา โหติ, สํสฏฺโฐ วิหรติ ราชูหิ ราชมหามตฺเตหิ ติตฺถิเยหิ ติตฺถิยสาวเกหิ อนนุโลมิเกน สํสคฺเคน, ยานิ วา ปน ตานิ กุลานิ อสฺสทฺธานิ อปฺปสนฺนานิ อโนปานภูตานิ อกฺโกสกปริภาสกานิ อนตฺถกามานิ อหิตกามานิ อผาสุกามานิ อโยคกฺเขมกามานิ ภิกฺขูนํ ภิกฺขุนีนํ อุปาสกานํ อุปาสิกานํ, ตถารูปานิ กุลานิ เสวติ ภชติ ปยิรุปาสติ, อยํ วุจฺจติ อโคจโร.}

\prose{อถ วา อนฺตรฆรํ ปวิฏฺโฐ วีถิํ ปฏิปนฺโน อสํวุโต คจฺฉติ หตฺถิํ โอโลเกนฺโต, อสฺสํ โอโลเกนฺโต, รถํ โอโลเกนฺโต, ปตฺติํ โอโลเกนฺโต, อิตฺถิโย โอโลเกนฺโต, ปุริเส โอโลเกนฺโต, กุมาริกาโย โอโลเกนฺโต, กุมารเก โอโลเกนฺโต, อนฺตราปณํ โอโลเกนฺโต, ฆรมุขานิ โอโลเกนฺโต, อุทฺธํ โอโลเกนฺโต\footnote{อุลฺโลเกนฺโต (สี, ก) 285 ปิฏฺเฐ นตฺถิ ปาฐนานตฺตํ.}, อโธ โอโลเกนฺโต, ทิสาวิทิสํ วิเปกฺขมาโน คจฺฉติ, อยมฺปิ วุจฺจติ อโคจโร.}

\prose{อถ วา จกฺขุนา รูปํ ทิสฺวา นิมิตฺตคฺคาหี โหติ อนุพฺยญฺชนคฺคาหี, ยตฺวาธิกรณเมนํ ฯเปฯ มนินฺทฺริยํ อสํวุตํ วิหรนฺตํ อภิชฺฌา โทมนสฺสา ปาปกา อกุสลา ธมฺมา อนฺวาสฺสเวยฺยุํ, ตสฺส น สํวราย ปฏิปชฺชติ, น รกฺขติ มนินฺทฺริยํ, มนินฺทฺริเย น สํวรํ อาปชฺชติ, อยมฺปิ วุจฺจติ อโคจโร.}

\prose{\csromanfolio{378}\csromansymbolfootnote{*}{ที 1. 6 ปิฏฺเฐปิ.}ยถา วา ปเนเก โภนฺโต สมณพฺราหฺมณา สทฺธาเทยฺยานิ โภชนานิ ภุญฺชิตฺวา เต เอวรูปํ วิสูกทสฺสนํ อนุยุตฺโต วิหรนฺติ.\csromanspacer{} เสยฺยถิทํ, นจฺจํ คีตํ วาทิตํ เปกฺขํ อกฺขานํ ปาณิสฺสรํ เวตาฬํ กุมฺภถูณํ โสภนกํ จณฺฑาลํ วํสํ โธวนํ หตฺถิยุทฺธํ อสฺสยุทฺธํ มหิํสยุทฺธํ อุสภยุทฺธํ อชยุทฺธํ เมณฺฑยุทฺธํ กุกฺกุฏยุทฺธํ วฏฺฏกยุทฺธํ ทณฺฑยุทฺธํ มุฏฺฐิยุทฺธํ นิพฺพุทฺธํ อุยฺโยธิกํ พลคฺคํ เสนาพฺยูหํ อนีกทสฺสนํ อิติ วา, อิติ เอวรูปํ วิสูกทสฺสนํ อนุยุตฺโต โหติ, อยมฺปิ วุจฺจติ อโคจโร.}

\prose{ปญฺจปิ กามคุณา อโคจรา.\csromanspacer{} วุตฺตํ เหตํ ภควตา\csromandash{} + มา ภิกฺขเว อโคจเร จรตฺถ ปรวิสเย, อโคจเร ภิกฺขเว จรตํ ปรวิสเย ลจฺฉติ มาโร โอตารํ, ลจฺฉติ มาโร อารมฺมณํ.\csromanspacer{} โก จ ภิกฺขเว ภิกฺขุโน อโคจโร ปรวิสโย, ยทิทํ ปญฺจ กามคุณา.\csromanspacer{} กตเม ปญฺจ, จกฺขุวิญฺเญยฺยา รูปา อิฏฺฐา กนฺตา มนาปา ปิยรูปา กามูปสํหิตา รชนียา.\csromanspacer{} โสตวิญฺเญยฺยา สทฺทา.\csromanspacer{} ฆานวิญฺเญยฺยา คนฺธา.\csromanspacer{} ชิวฺหาวิญฺเญยฺยา รสา.\csromanspacer{} กายวิญฺเญยฺยา โผฏฺฐพฺพา อิฏฺฐา กนฺตา มนาปา ปิยรูปา กามูปสํหิตา รชนียา.\csromanspacer{} อยํ วุจฺจติ ภิกฺขเว ภิกฺขุโน อโคจโร ปรวิสโย, อยมฺปิ วุจฺจติ อโคจโร.}

\prose{กตโม โคจโร, อิธ ภิกฺขุ น เวสิยาโคจโร โหติ, น วิธวาโคจโร โหติ, น ถุลฺลกุมารีโคจโร โหติ, น ปณฺฑกโคจโร โหติ, น ภิกฺขุนีโคจโร โหติ, น ปานาคารโคจโร โหติ, อสํสฏฺโฐ วิหรติ ราชูหิ ราชมหามตฺเตหิ ติตฺถิเยหิ ติตฺถิยสาวเกหิ อนนุโลมิเกน สํสคฺเคน, ยานิ วา ปน ตานิ กุลานิ สทฺธานิ ปสนฺนานิ โอปานภูตานิ กาสาวปชฺโชตานิ อิสิวาตปฏิวาตานิ อตฺถกามานิ หิตกามานิ ผาสุกามานิ โยคกฺเขมกามานิ ภิกฺขูนํ ภิกฺขุนีนํ อุปาสกานํ อุปาสิกานํ, ตถารูปานิ กุลานิ เสวติ ภชติ ปยิรุปาสติ, อยมฺปิ วุจฺจติ โคจโร.}

\prose{อถ วา ภิกฺขุ อนฺตรฆรํ ปวิฏฺโฐ วีถิํ ปฏิปนฺโน สํวุโต คจฺฉติ, น หตฺถิํ โอโลเกนฺโต, น อสฺสํ โอโลเกนฺโต, น รถํ โอโลเกนฺโต, น ปตฺติํ โอโลเกนฺโต ฯเปฯ น ทิสาวิทิสํ วิเปกฺขมาโน คจฺฉติ, อยมฺปิ วุจฺจติ โคจโร.}

\csromanheader{2}{16. สาริปุตฺตสุตฺตนิทฺเทส}
\prose{\csromanfolio{379}อถ วา ภิกฺขุ จกฺขุนา รูปํ ทิสฺวา น นิมิตฺตคฺคาหี โหติ ฯเปฯ มนินฺทฺริเย สํวรํ อาปชฺชติ, อยมฺปิ วุจฺจติ โคจโร.}

\prose{ยถา วา ปเนเก โภนฺโต สมณพฺราหฺมณา สทฺธาเทยฺยานิ โภชนานิ ภุญฺชิตฺวา เต เอวรูปํ วิสูกทสฺสนํ อนนุยุตฺตา วิหรนฺติ.\csromanspacer{} เสยฺยถิทํ, นจฺจํ คีตํ วาทิตํ ฯเปฯ อนีกทสฺสนํ อิติ วา, อิติ เอวรูปา วิสูกทสฺสนา ปฏิวิรโต โหติ.\csromanspacer{} อยมฺปิ วุจฺจติ โคจโร.}

\prose{จตฺตาโรปิ สติปฏฺฐานา โคจโร.\csromanspacer{} วุตฺตํ เหตํ ภควตา\csromandash{} * โคจเร ภิกฺขเว จรถ สเก เปตฺติเก วิสเย, โคจเร ภิกฺขเว จรตํ สเก เปตฺติเก วิสเย น ลจฺฉติ มาโร โอตารํ, น ลจฺฉติ มาโร อารมฺมณํ.\csromanspacer{} โก จ ภิกฺขเว ภิกฺขุโน โคจโร สโก เปตฺติโก วิสโย, ยทิทํ จตฺตาโร สติปฏฺฐานา.\csromanspacer{} กตเม จตฺตาโร, อิธ ภิกฺขเว ภิกฺขุ กาเย กายานุปสฺสี วิหรติ, เวทนาสุ ฯเปฯ จิตฺเต ฯเปฯ ธมฺเมสุ ธมฺมานุปสฺสี วิหรติ อาตาปี สมฺปชาโน สติมา วิเนยฺย โลเก อภิชฺฌา โทมนสฺสํ, อยํ วุจฺจติ ภิกฺขเว ภิกฺขุโน โคจโร สโก เปตฺติโก วิสโย, อยมฺปิ วุจฺจติ โคจโร.\csromanspacer{} เอทิเสน โคจเรน สมนฺนาคโต อสฺสาติ กฺยาสฺสสฺสุ อิธ โคจรา.}

\prose{\textbf{กานิ สีลพฺพตานาสฺสู}ติ กีทิเสน สีลพฺพเตน สมนฺนาคโต อสฺส กิํสณฺฐิเตน กิํปกาเรน กิํปฏิภาเคนาติ สีลพฺพตปาริสุทฺธิํ ปุจฺฉติ.\csromanspacer{} กตมา สีลพฺพตปาริสุทฺธิ, อตฺถิ สีลญฺเจว วตญฺจ, อตฺถิ วตํ น สีลํ.\csromanspacer{} กตมํ สีลญฺเจว วตญฺจ, อิธ ภิกฺขุ สีลวา โหติ ปาติโมกฺขสํวรสํวุโต วิหรติ อาจารโคจรสมฺปนฺโน อณุมตฺเตสุ วชฺเชสุ ภยทสฺสาวี, สมาทาย สิกฺขติ สิกฺขาปเทสุ, โย ตตฺถ สญฺญโม สํวโร อวีติกฺกโม, อิทํ สีลํ.\csromanspacer{} ยํ สมาทานํ, ตํ วตํ.\csromanspacer{} สํวรฏฺเฐน สีลํ, สมาทานฏฺเฐน วตํ.\csromanspacer{} อิทํ วุจฺจติ สีลญฺเจว วตญฺจ.}

\prose{กตมํ วตํ, น สีลํ, อฏฺฐ ธุตงฺคานิ อารญฺญิกงฺคํ ปิณฺฑปาติกงฺคํ ปํสุกูลิกงฺคํ เตจีวริกงฺคํ สปทานจาริกงฺคํ ขลุปจฺฉาภตฺติกงฺคํ เนสชฺชิกงฺคํ ยถาสนฺถติกงฺคํ, อิทํ วุจฺจติ วตํ, น สีลํ.\csromanspacer{} วีริยสมาทานมฺปิ วุจฺจติ วตํ, น สีลํ. + กามํ ตโจ จ นฺหารุ จ อฏฺฐิ จ อวสิสฺสตุ สรีเร อุปสุสฺสตุ มํสโลหิตํ, ยํ ตํ ปุริสถาเมน ปุริสพเลน ปุริสวีริเยน ปุริสปรกฺกเมน ปตฺตพฺพํ, น ตํ อปาปุณิตฺวา วีริยสฺส สณฺฐานํ}

\prose{\csromanfolio{380}ภวิสฺสตีติ จิตฺตํ ปคฺคณฺหาติ ปทหติ, เอวรูปมฺปิ วีริยสมาทานํ วุจฺจติ วตํ, น สีลํ.}

\csromangathameasure{“นาสิสฺสํ น ปิวิสฺสามิ, วิหารโต น นิกฺขเม. \\ นปิ ปสฺสํ นิปาเตสฺสํ, ตณฺหาสลฺเล อนุหเต”ติ.}
\csromangathagroup{\csromangathastanza{“นาสิสฺสํ น ปิวิสฺสามิ, วิหารโต น นิกฺขเม\footnote{น นิกฺขมิํ (สฺยา) 51 ปิฏฺเฐปิ.}. \\ นปิ ปสฺสํ นิปาเตสฺสํ, ตณฺหาสลฺเล อนุหเต”ติ.}}

\prose{จิตฺตํ ปคฺคณฺหาติ ปทหติ, เอวรูปมฺปิ วีริยสมาทานํ วุจฺจติ วตํ, น สีลํ.\csromanspacer{}\csromanspacer{}น ตาวาหํ อิมํ ปลฺลงฺกํ ภินฺทิสฺสามิ, ยาว เม น อนุปาทาย อาสเวหิ จิตฺตํ วิมุจฺจิสฺสตีติ จิตฺตํ ปคฺคณฺหาติ ปทหติ, เอวรูปมฺปิ วีริยสมาทานํ วุจฺจติ วตํ, น สีลํ.\csromanspacer{}\csromanspacer{}น ตาวาหํ อิมมฺหา อาสนา วุฏฺฐหิสฺสามิ.\csromanspacer{} น ตาวาหํ อิมมฺหา จงฺกมา โอโรหิสฺสามิ.\csromanspacer{} วิหารา นิกฺขมิสฺสามิ.\csromanspacer{} อฑฺฒโยคา นิกฺขมิสฺสามิ.\csromanspacer{} ปาสาทา นิกฺขมิสฺสามิ.\csromanspacer{} หมฺมิยา.\csromanspacer{} คุหาย.\csromanspacer{} เลณา.\csromanspacer{} กุฏิยา.\csromanspacer{} กูฏาคารา.\csromanspacer{} อฏฺฏา.\csromanspacer{} มาฬา.\csromanspacer{} อุทฺทณฺฑา.\csromanspacer{} อุปฏฺฐานสาลาย.\csromanspacer{} มณฺฑปา.\csromanspacer{} รุกฺขมูลา นิกฺขมิสฺสามิ, ยาว เม น อนุปาทาย อาสเวหิ จิตฺตํ วิมุจฺจิสฺสตีติ จิตฺตํ ปคฺคณฺหาติ ปทหติ, เอวรูปมฺปิ วีริยสมาทานํ วุจฺจติ วตํ, น สีลํ.\csromanspacer{} อิมสฺมิญฺเญว ปุพฺพณฺหสมยํ อริยธมฺมํ อาหริสฺสามิ สมาหริสฺสามิ อธิคจฺฉิสฺสามิ ผสฺสยิสฺสามิ สจฺฉิกริสฺสามีติ จิตฺตํ ปคฺคณฺหาติ ปทหติ, เอวรูปมฺปิ วีริยสมาทานํ วุจฺจติ วตํ, น สีลํ.\csromanspacer{} อิมสฺมิญฺเญว มชฺฌนฺหิกสมยํ.\csromanspacer{} สายนฺหสมยํ.\csromanspacer{} ปุเรภตฺตํ.\csromanspacer{} ปจฺฉาภตฺตํ.\csromanspacer{} ปุริมยามํ.\csromanspacer{} มชฺฌิมยามํ.\csromanspacer{} ปจฺฉิมยามํ.\csromanspacer{} กาเฬ.\csromanspacer{} ชุเณฺห.\csromanspacer{} วสฺเส.\csromanspacer{} เหมนฺเต.\csromanspacer{} คิเมฺห.\csromanspacer{} ปุริเม วโยขนฺเธ.\csromanspacer{} มชฺฌิเม วโยขนฺเธ.\csromanspacer{} ปจฺฉิเม วโยขนฺเธ อริยธมฺมํ อาหริสฺสามิ สมาหริสฺสามิ อธิคจฺฉิสฺสามิ ผสฺสยิสฺสามิ สจฺฉิกริสฺสามีติ จิตฺตํ ปคฺคณฺหาติ ปทหติ, เอวรูปมฺปิ วีริยสมาทานํ วุจฺจติ วตํ, น สีลํ.\csromanspacer{} อยํ สีลพฺพตปาริสุทฺธิ.\csromanspacer{} เอทิสาย\footnote{กีทิสาย (สี), อีทิสาย (สฺยา)} สีลพฺพตปาริสุทฺธิยา สมนฺนาคโต อสฺสาติ กานิ สีลพฺพตานาสฺสุ.}

\prose{\textbf{ปหิตตฺตสฺส ภิกฺขุโน}ติ \textbf{ปหิตตฺตสฺสา}ติ อารทฺธวีริยสฺส ถามคตสฺส ทฬฺหปรกฺกมสฺส อนิกฺขิตฺตจฺฉนฺทสฺส อนิกฺขิตฺตธุรสฺส กุสเลสุ ธมฺเมสุ.\csromanspacer{} อถ วา เปสิตตฺตสฺส ยสฺสตฺถาย เปสิโต อตฺตตฺเถ จ ญาเย จ ลกฺขเณ จ การเณ จ ฐานาฐาเน จ. “สพฺเพ สงฺขารา อนิจฺจา”ติ เปสิตตฺตสฺส, “สพฺเพ สงฺขารา ทุกฺขา”ติ เปสิตตฺตสฺส, “สพฺเพ ธมฺมา อนตฺตา”ติ เปสิตตฺตสฺส. “อวิชฺชาปจฺจยา สงฺขารา”ติ}

\vspace{\tipitakaparskip}
\csromancenter{สาริปุตฺตสุตฺตนิทฺเทส เปสิตตฺตสฺส ฯเปฯ “ชาติปจฺจยา ชรามรณนฺ”ติ เปสิตตฺถสฺส. “อวิชฺชานิโรธา สงฺขารนิโรโธ”ติ เปสิตตฺตสฺส ฯเปฯ “ชาตินิโรธา ชรามรณนิโรโธ”ติ เปสิตตฺตสฺส. “อิทํ ทุกฺขนฺ”ติ เปสิตตฺตสฺส ฯเปฯ “อยํ ทุกฺขนิโรธคามินี ปฏิปทา”ติ เปสิตตฺตสฺส, “อิเม อาสคา”ติ เปสิตตฺตสฺส ฯเปฯ “อยํ อาสวนิโรธคามินี ปฏิปทา”ติ เปสิตตฺตสฺส, “อิเม ธมฺมา อภิญฺเญยฺยา”ติ เปสิตตฺตสฺส ฯเปฯ “อิเม ธมฺมา สจฺฉิกาตพฺพา”ติ เปสิตตฺตสฺส, ฉนฺนํ ผสฺสายตนานํ สมุทยญฺจ อตฺถงฺคมญฺจ อสฺสาทญฺจ อาทีนวญฺจ นิสฺสรณญฺจ เปสิตตฺตสฺส, ปญฺจนฺนํ อุปาทานกฺขนฺธานํ.\csromanspacer{} จตุนฺนํ มหาภูตานํ สมุทยญฺจ อตฺถงฺคมญฺจ อสฺสาทญฺจ อาทีนวญฺจ นิสฺสรณญฺจ เปสิตตฺตสฺส, “ยํ กิญฺจิ สมุทยธมฺมํ, สพฺพํ ตํ นิโรธธมฺมนฺ”ติ เปสิตตฺตสฺส.\csromanspacer{}\csromanspacer{}\textbf{ภิกฺขุโน}ติ ปุถุชฺชนกลฺยาณกสฺส วา ภิกฺขุโน, เสกฺขสฺส วา ภิกฺขุโนติ ปหิตตฺตสฺส ภิกฺขุโน.\csromanspacer{} เตนาห เถโร สาริปุตฺโต\csromandash{}}
\vspace{0.5\baselineskip}
\csromangathameasure{กฺยาสฺส พฺยปฺปถโย อสฺสุ, กฺยาสฺสสฺสุ อิธ โคจรา. \\ กานิ สีลพฺพตานาสฺสุ, ปหิตตฺตสฺส ภิกฺขุโนติ.}
\csromangathagroup{\csromangathastanza{\csromanfolio{381}กฺยาสฺส พฺยปฺปถโย อสฺสุ, กฺยาสฺสสฺสุ อิธ โคจรา. \\ กานิ สีลพฺพตานาสฺสุ, ปหิตตฺตสฺส ภิกฺขุโนติ.}}

\proseitem{197}{\csromansymbolfootnote{*}{ขุ 1. 426 ปิฏฺเฐปิ.}กํ โส สิกฺขํ สมาทาย, เอโกทิ นิปโก สโต.}

\csromangathameasure{กมฺมาโร รชตสฺเสว, นิทฺธเม มลมตฺตโน.}
\csromangathagroup{\csromangathastanza{กมฺมาโร รชตสฺเสว, นิทฺธเม มลมตฺตโน.}}

\prose{\textbf{กํ โส สิกฺขํ สมาทายา}ติ กํ โส สิกฺขํ อาทาย สมาทาย อาทิยิตฺวา สมาทิยิตฺวา คณฺหิตฺวา ปรามสิตฺวา อภินิวิสิตฺวาติ กํ โส สิกฺขํ สมาทาย.}

\prose{\textbf{เอโกทิ นิปโก สโต}ติ \textbf{เอโกที}ติ เอกคฺคจิตฺโต อวิกฺขิตฺตจิตฺโต อวิสาหฏมานโส สมโถ สมาธินฺทฺริยํ สมาธิพลํ ฯเปฯ สมฺมาสมาธิ.\csromanspacer{} \textbf{นิปโก}ติ นิปโก ปณฺฑิโต ปญฺญวา พุทฺธิมา ญาณี วิภาวี เมธาวี.\csromanspacer{} \textbf{สโต}ติ จตูหิ การเณหิ สโต, กาเย กายานุปสฺสนาสติปฏฺฐานํ ภาเวนฺโต สโต, เวทนาสุ ฯเปฯ จิตฺเต ฯเปฯ ธมฺเมสุ ธมฺมานุปสฺสนาสติปฏฺฐานํ ภาเวนฺโต สโต, โส วุจฺจติ สโตติ สโต.\csromanspacer{}\csromanspacer{}\textbf{กํ โส สิกฺขํ สมาทายา}ติ อธิสีลสิกฺขํ ปุจฺฉติ.\csromanspacer{} เอโกทีติ อธิจิตฺตสิกฺขํ ปุจฺฉติ.\csromanspacer{} \textbf{นิปโก}ติ อธิปญฺญาสิกฺขํ ปุจฺฉติ.\csromanspacer{} \textbf{สโต}ติ ปาริสุทฺธิํ ปุจฺฉตีติ กํ โส สิกฺขํ สมาทาย, เอโกทิ นิปโก สโต.}

\prose{\csromanfolio{382}\textbf{กมฺมาโร รชตสฺเสว, นิทฺธเม มลมตฺตโน}ติ กมฺมาโร วุจฺจติ สุวณฺณกาโร, รชตํ วุจฺจติ ชาตรูปํ.\csromanspacer{} ยถา สุวณฺณกาโร ชาตรูปสฺส โอฬาริกมฺปิ มลํ ธมติ สนฺธมติ นิทฺธมติ.\csromanspacer{} มชฺฌิมกมฺปิ มลํ ธมติ สนฺธมติ นิทฺธมติ, สุขุมกมฺปิ มลํ ธมติ สนฺธมติ นิทฺธมติ.\csromanspacer{} เอวเมว ภิกฺขุ อตฺตโน โอฬาริเกปิ กิเลเส ธมติ สนฺธมติ นิทฺธมติ ปชหติ วิโนเทติ พฺยนฺติํ กโรติ อนภาวํ คเมติ, มชฺฌิมเกปิ กิเลเส.\csromanspacer{} สุขุมเกปิ กิเลเส ธมติ สนฺธมติ นิทฺธมติ ปชหติ วิโนเทติ พฺยนฺติํ กโรติ อนภาวํ คเมติ.}

\prose{อถ วา ภิกฺขุ อตฺตโน ราคมลํ โทสมลํ โมหมลํ มานมลํ ทิฏฺฐิมลํ กิเลสมลํ ทิจฺจริตมลํ อนฺธกรณํ อจกฺขุกรณํ อญฺญาณกรณํ ปญฺญานิโรธิกํ วิฆาตปกฺขิกํ อนิพฺพานสํวตฺตนิกํ ธมติ สนฺธมติ นิทฺธมติ ปชหติ วิโนเทติ พฺยนฺติํ กโรติ อนภาวํ คเมติ.}

\prose{อถ วา สมฺมาทิฏฺฐิยา มิจฺฉาทิฏฺฐิํ ธมติ สนฺธมติ นิทฺทมติ ปชหติ วิโนเทติ พฺยนฺติํ กโรติ อนภาวํ คเมติ.\csromanspacer{} สมฺมาสงฺกปฺเปน มิจฺฉาสงฺกปฺปํ.\csromanspacer{} สมฺมาวาจาย มิจฺฉาวาจํ.\csromanspacer{} สมฺมากมฺมนฺเตน มิจฺฉากมฺมนฺตํ.\csromanspacer{} สมฺมา อาชีเวน มิจฺฉา-อาชีวํ.\csromanspacer{} สมฺมาวายาเมน มิจฺฉาวายามํ.\csromanspacer{} สมฺมาสติยา มิจฺฉาสติํ.\csromanspacer{} สมฺมาสมาธินา มิจฺฉาสมาธิํ.\csromanspacer{} สมฺมาญาเณน มิจฺฉาญาณํ.\csromanspacer{} สมฺมาวิมุตฺติยา มิจฺฉาวิมุตฺติํ ธมติ สนฺธมติ นิทฺธมติ ปชหติ วิโนเทติ พฺยนฺติํ กโรติ อนภาวํ คเมติ.}

\prose{อถ วา อริเยน อฏฺฐงฺคิเกน มคฺเคน สพฺเพ กิเลเส สพฺเพ ทุจฺจริเต สพฺเพ ทรเถ สพฺเพ ปริฬาเห สพฺเพ สนฺตาเป สพฺพากุสลาภิสงฺขาเร ธมติ สนฺธมติ นิทฺธมติ ปชหติ วิโนเทติ พฺยนฺติํ กโรติ อนภาวํ คเมตีติ กมฺมาโร รชตสฺเสว, นิทฺธเม มลมตฺตโน.\csromanspacer{} เตนาห เถโร สาริปุตฺโต\csromandash{}}

\csromangathameasure{กํ โส สิกฺขํ สมาทาย, เอโกทิ นิปโก สโต.}
\csromangathagroup{\csromangathastanza{กํ โส สิกฺขํ สมาทาย, เอโกทิ นิปโก สโต.}}

\prose{\textbf{กมฺมาโร รชตสฺเสว, นิทฺธเม มลมตฺตโน}ติ.}

\proseitem{198}{\csromansymbolfootnote{*}{ขุ 1. 426 ปิฏฺเฐ.}วิชิคุจฺฉมานสฺส ยทิทํ ผาสุ, (สาริปุตฺตาติ ภควา,)}

\prose{\textbf{ริตฺตาสนํ สยนํ เสวโต เว.}}

\csromangathameasure{\textbf{สมฺโพธิกามสฺส ยถานุธมฺมํ,} \\ \textbf{ตํ เต ปวกฺขามิ ยถา ปชานํ.}}
\csromangathagroup{\csromangathastanza{\textbf{สมฺโพธิกามสฺส ยถานุธมฺมํ,} \\ \textbf{ตํ เต ปวกฺขามิ ยถา ปชานํ.}}}

\csromanheader{2}{16. สาริปุตฺตสุตฺตนิทฺเทส}
\prose{\csromanfolio{383}\textbf{วิชิคุจฺฉมานสฺส ยทิทํ ผาสู}ติ \textbf{วิชิคุจฺฉมานสฺสา}ติ ชาติยา วิชิคุจฺฉมานสฺส, ชราย.\csromanspacer{} พฺยาธินา.\csromanspacer{} มรเณน.\csromanspacer{} โสเกหิ.\csromanspacer{} ปริเทเวหิ.\csromanspacer{} ทุกฺเขหิ.\csromanspacer{} โทมนสฺเสหิ.\csromanspacer{} อุปายาเสหิ ฯเปฯ.\csromanspacer{} ทิฏฺฐิพฺยสเนน ทุกฺเขน วิชิคุจฺฉมานสฺส อฏฺฏียมานสฺส\footnote{อฏฺฏิยมานสฺส (สฺยา, ก)} หรายมานสฺสาติ วิชิคุจฺฉมานสฺส.\csromanspacer{} \textbf{ยทิทํ ผาสู}ติ ยํ ผาสุวิหารํ, ตํ กถยิสฺสามิ.\csromanspacer{} กตโม ผาสุวิหาโร, สมฺมาปฏิปทา อนุโลมปฏิปทา อปจฺจนีกปฏิปทา อวิรุทฺธปฏิปทา อนฺวตฺถปฏิปทา ธมฺมานุธมฺมปฏิปทา สีเลสุ ปริปูรการิตา อินฺทฺริเยสุ คุตฺตทฺวารตา โภชเน มตฺตญฺญุตา ชาคริยานุโยโค สติสมฺปชญฺญํ จตฺตาโร สติปฏฺฐานา จตฺตาโร สมฺมปฺปธานา จตฺตาโร อิทฺธิปาทา ปญฺจินฺทฺริยานิ ปญฺจ พลานิ สตฺต โพชฺฌงฺคา อริโย อฏฺฐงฺคิโก มคฺโค นิพฺพานญฺจ นิพฺพานคามินี จ ปฏิปทา, อยํ ผาสุวิหาโรติ วิชิคุจฺฉมานสฺส ยทิทํ ผาสุ.}

\prose{\textbf{สาริปุตฺตาติ ภควา}ติ ตํ เถรํ นาเมนาลปติ.\csromanspacer{} \textbf{ภควา}ติ คารวาธิวจนํ.\csromanspacer{} อปิ จ ภคฺคราโคติ \textbf{ภควา}, ภคฺคโทโสติ \textbf{ภควา}, ภคฺคโมโหติ \textbf{ภควา}, ภคฺคมาโนติ \textbf{ภควา}, ภคฺคทิฏฺฐีติ \textbf{ภควา}, ภคฺคกณฺฑโกติ \textbf{ภควา}, ภคฺคกิเลโสติ \textbf{ภควา}, ภชิ วิภชิ ปวิภชิ\footnote{ปฏิภชิ (สี, สฺยา) 110 ปิฏฺเฐ นตฺถิ ปาฐนานตฺตํ.} ธมฺมรตนนฺติ \textbf{ภควา}, ภวานํ อนฺตกโรติ \textbf{ภควา}, ภาวิตกาโย, ภาวิตสีโล, ภาวิตจิตฺโต, ภาวิตปญฺโญติ \textbf{ภควา}, ภชิ วา \textbf{ภควา} อรญฺญวนปตฺถานิ ปนฺตานิ เสนาสนานิ อปฺปสทฺทานิ อปฺปนิคฺโฆสานิ วิชนวาตานิ มนุสฺสราหสฺเสยฺยกานิ ปฏิสลฺลานสารุปฺปานีติ \textbf{ภควา}, ภาคี วา \textbf{ภควา} จีวรปิณฺฑปาตเสนาสนคิลานปจฺจยเภสชฺชปริกฺขารานนฺติ \textbf{ภควา}, ภาคี วา \textbf{ภควา} อตฺถรสสฺส ธมฺมรสสฺส วิมุตฺติรสสฺส อธิสีลสฺส อธิจิตฺตสฺส อธิปญฺญายาติ \textbf{ภควา}, ภาคี วา \textbf{ภควา} จตุนฺนํ ฌานานํ จตุนฺนํ อปฺปมญฺญานํ จตุนฺนํ อารุปฺปสมาปตฺตีนนฺติ \textbf{ภควา}, ภาคี วา \textbf{ภควา} อฏฺฐนฺนํ วิโมกฺขานํ อฏฺฐนฺนํ อภิญฺญายตนานํ นวนฺนํ อนุปุพฺพวิหารสมาปตฺตีนนฺติ \textbf{ภควา}, ภาคี วา \textbf{ภควา} ทสนฺนํ สญฺญาภาวนานํ ทสนฺนํ กสิณสมาปตฺตีนํ อานาปานสฺสติสมาธิสฺส อสุภสมาปตฺติยาติ \textbf{ภควา}, ภาคี วา \textbf{ภควา} จตุนฺนํ สติปฏฺฐานานํ จตุนฺนํ สมฺมปฺปธานานํ จตุนฺนํ อิทฺธิปาทานํ ปญฺจนฺนํ อินฺทฺริยานํ ปญฺจนฺนํ พลานํ สตฺตนฺนํ โพชฺฌงฺคานํ อริยสฺส อฏฺฐงฺคิกสฺส มคฺคสฺสาติ \textbf{ภควา}, ภาคี วา \textbf{ภควา} ทสนฺนํ ตถาคตพลานํ จตุนฺนํ เวสารชฺชานํ จตุนฺนํ ปฏิสมฺภิทานํ ฉนฺนํ อภิญฺญานํ ฉนฺนํ พุทฺธธมฺมานนฺติ}

\prose{\csromanfolio{384}ภควา, ภควาติ เนตํ นามํ มาตรา กตํ, น ปิตรา กตํ, น ภาถรา กตํ, น ภคินิยา กตํ, น มิตฺตามจฺเจหิ กตํ, น ญาติสาโลหิเตหิ กตํ, น สมณพฺราหฺมเณหิ กตํ, น เทวตาหิ กตํ, วิโมกฺขนฺติกเมตํ พุทฺธานํ ภควนฺตานํ โพธิยา มูเล สห สพฺพญฺญุตญาณสฺส ปฏิลาภา สจฺฉิกา ปญฺญตฺติ ยทิทํ ภควาติ สาริปุตฺตาติ ภควา.}

\prose{\textbf{ริตฺตาสนํ สยนํ เสวโต เว}ติ อาสนํ วุจฺจติ ยตฺถ นิสีทติ มญฺโจ ปีฐํ ภิสิ ตฏฺฏิกา จมฺมขณฺโฑ ติณสนฺถาโร ปณฺณสนฺถาโร ปลาลสนฺถาโร.\csromanspacer{} สยนํ วุจฺจติ เสนาสนํ วิหาโร อฑฺฒโยโค ปาสาโท หมฺมิยํ คุหา.\csromanspacer{} ตํ สยนาสนํ อสปฺปายรูปทสฺสเนน ริตฺตํ วิวิตฺตํ ปวิวิตฺตํ, อสปฺปายสทฺทสฺสวเนน ฯเปฯ อสปฺปาเยหิ ปญฺจหิ กามคุเณหิ ริตฺตํ วิวิตฺตํ ปวิวิตฺตํ, ตํ สยนาสนํ เสวโต นิเสวโต สํเสวโต ปฏิเสวโตติ ริตฺตาสนํ สยนํ เสวโต เว.}

\prose{\textbf{สมฺโพธิกามสฺส ยถานุธมฺมนฺ}ติ สมฺโพธิ วุจฺจติ จตูสุ มคฺเคสุ ญาณํ ปญฺญา ปญฺญินฺทฺริยํ ปญฺญาพลํ ฯเปฯ ธมฺมวิจย- สมฺโพชฺฌงฺโค วีมํสา วิปสฺสนา สมฺมาทิฏฺฐิ, ตํ สมฺโพธิํ พุชฺฌิตุกามสฺส อนุพุชฺฌิตุกามสฺส ปฏิพุชฺฌิตุกามสฺส สมฺพุชฺฌิตุกามสฺส อธิคนฺตุกามสฺส ผสฺสิตุกามสฺส สจฺฉิกา ตุกามสฺสาติ สมฺโพธิกามสฺส.}

\prose{\textbf{ยถานุธมฺมนฺ}ติ กตเม โพธิยา อนุธมฺมา, สมฺมาปฏิปทา อนุโลมปฏิปทา อปจฺจนีกปฏิปทา อวิรุทฺธปฏิปทา อนฺวตฺถปฏิปทา ธมฺมานุธมฺมปฏิปทา สีเลสุ ปริปูรการิตา อินฺทฺริเยสุ คุตฺตทฺวารตา โภชเน มตฺตญฺญุตา ชาคริยานุโยโค สติสมฺปชญฺญํ, อิเม วุจฺจนฺติ โพธิยา อนุธมฺมา.\csromanspacer{} อถ วา จตุนฺนํ มคฺคานํ ปุพฺพภาเค วิปสฺสนา, อิเม วุจฺจนฺติ โพธิยา อนุธมฺมาติ สมฺโพธิกามสฺส ยถานุธมฺมํ.}

\prose{\textbf{ตํ เต ปวกฺขามิ ยถา ปชานนฺ}ติ \textbf{ตนฺ}ติ โพธิยา อนุธมฺมํ.\csromanspacer{} \textbf{ปวกฺขามี}ติ ปวกฺขามิ อาจิกฺขิสฺสามิ เทเสสฺสามิ ปญฺญเปสฺสามิ ปฐเปสฺสามิ วิวริสฺสามิ วิภชิสฺสามิ อุตฺตานีกริสฺสามิ ปกาสิสฺสามิ.\csromanspacer{} \textbf{ยถา ปชานนฺ}ติ ยถา ปชานํ ยถา ปชานนฺโต อาชานนฺโต วิชานนฺโต ปฏิวิชานนฺโต ปฏิวิชฺฌนฺโต น อิติหิติหํ น อิติกิราย น ปรํปราย น ปิฏกสมฺปทาย น ตกฺกเหตุ น นยเหตุ น อาการปริวิตกฺเกน น ทิฏฺฐินิชฺฌานกฺขนฺติยา สามํ สยมภิญฺญาตํ อตฺตปจฺจกฺขํ ธมฺมํ, ตํ กถยิสฺสามีติ ตํ เต ปวกฺขามิ ยถา ปชานํ.\csromanspacer{} เตนาห ภควา.}

\csromanheader{2}{16. สาริปุตฺตสุตฺตนิทฺเทส}
\prose{\csromanfolio{385}วิชิคุจฺฉมานสฺส ยทิทํ ผาสุ\textbf{,} (สาริปุตฺตาติ ภควา\textbf{,})}

\csromangathameasure{ริตฺตาสนํ สยนํ เสวโต เว\textbf{,} \\ สมฺโพธิกามสฺส ยถานุธมฺมํ. \\ ตํ เต ปวกฺขามิ ยถา ปชานนฺติ.}
\csromangathagroup{\csromangathastanza{ริตฺตาสนํ สยนํ เสวโต เว\textbf{,} \\ สมฺโพธิกามสฺส ยถานุธมฺมํ. \\ ตํ เต ปวกฺขามิ ยถา ปชานนฺติ.}}

\proseitem{199}{\csromansymbolfootnote{*}{ขุ 1. 426 ปิฏฺเฐปิ.}\textbf{ปญฺจนฺนํ ธีโร ภยานํ น ภาเย,}}

\prose{\textbf{ภิกฺขุ สโต สปริยนฺตจารี.}}

\csromangathameasure{q\textbf{อํสาธิปาตานํ สรีสปานํ}\textbf{,} \\ \textbf{มนุสฺสผสฺสานํ จตุปฺปทานํ.}}
\csromangathagroup{\csromangathastanza{q\textbf{อํสาธิปาตานํ สรีสปานํ}\footnote{สิริํสปานํ (สี, ก)}\textbf{,} \\ \textbf{มนุสฺสผสฺสานํ จตุปฺปทานํ.}}}

\prose{\textbf{ปญฺจนฺนํ ธีโร ภยานํ น ภาเย}ติ \textbf{ธีโร}ติ \textbf{ธีโร} ปณฺฑิโต ปญฺญวา พุทฺธิมา ญาณี วิภาวี เมธาวี \textbf{ธีโร} ปญฺจนฺนํ ภยานํ น ภาเยยฺย น ตเสยฺย น สนฺตเสยฺย น อุตฺตเสยฺย น ปริตฺตเสยฺย น สนฺตาสํ อาปชฺเชยฺย\textbf{,} อภีรู อสฺส อจฺฉมฺภี อนุตฺราสี อปลายี ปหีนภยเภรโว วิคตโลมหํโส วิหเรยฺยาติ ปญฺจนฺนํ \textbf{ธีโร} ภยานํ น ภาเย.}

\prose{\textbf{ภิกฺขุ สโต สปริยนฺตจารี}ติ \textbf{ภิกฺขู}ติ ปุถุชฺชนกลฺยาณโก วา ภิกฺขุ เสกฺโข วา ภิกฺขุ.\csromanspacer{} \textbf{สโต}ติ จตูหิ การเณหิ สโต\textbf{,} กาเย กายานุปสฺสนาสติปฏฺฐานํ ภาเวนฺโต สโต\textbf{,} เวทนาสุ ฯเปฯ จิตฺเต ฯเปฯ ธมฺเมสุ ธมฺมานุปสฺสนาสติปฏฺฐานํ ภาเวนฺโต สโต\textbf{,} โส วุจฺจติ สโต.\csromanspacer{} \textbf{สปริยนฺตจารี}ติ จตฺตาโร ปริยนฺตา สีลสํวรปริยนฺโต อินฺทฺริยสํวรปริยนฺโต โภชเน มตฺตญฺญุตาปริยนฺโต ชาคริยานุโยคปริยนฺโต.}

\prose{กตโม สีลสํวรปริยนฺโต\textbf{,} อิธ ภิกฺขุ สีลวา โหติ\textbf{,} ปาติโมกฺขสํวรสํวุโต วิหรติ อาจารโคจรสมฺปนฺโน อณุมตฺเตสุ วชฺเชสุ ภยทสฺสาวี\textbf{,} สมาทาย สิกฺขติ สิกฺขาปเทสุ.\csromanspacer{} อนฺโตปูติภาวํ ปจฺจเวกฺขมาโน อนฺโต สีลสํวรปริยนฺเต จรติ มริยาทํ น ภินฺทติ\textbf{,} อยํ สีลสํวรปริยนฺโต.}

\prose{กตโม อินฺทฺริยสํวรปริยนฺโต\textbf{,} อิธ ภิกฺขุ จกฺขุนา รูปํ ทิสฺวา น นิมิตฺตคฺคาหี โหติ นานุพฺยญฺชนคฺคาหี\textbf{,} ยตฺวาธิกรณเมนํ ฯเปฯ จกฺขุนฺทฺริเย \csromanfolio{386}สํวรํ อาปชฺชติ.\csromanspacer{} โสเตน สทฺทํ สุตฺวา.\csromanspacer{} ฆาเนน คนฺธํ ฆายิตฺวา.\csromanspacer{} ชิวฺหาย รสํ สายิตฺวา.\csromanspacer{} กาเยน โผฏฺฐพฺพํ ผุสิตฺวา.\csromanspacer{} มนสา ธมฺมํ วิญฺญาย น นิมิตฺตคฺคาหี โหติ นานุพฺยญฺชนคฺคาหี, ยตฺวาธิกรณเมนํ มนินฺทฺริยํ สุสํวุตํ วิหรนฺตํ อภิชฺฌาโทมนสฺสา ปาปกา อกุสลา ธมฺมา นานฺวาสฺสเวยฺยุํ, ตสฺส สํวราย ปฏิปชฺชติ, รกฺขติ มนินฺทฺริยํ, มนินฺทฺริเย สํวรํ อาปชฺชติ, อาทิตฺตปริยายํ ปจฺจเวกฺขมาโน อนฺโต อินฺทฺริยสํวรปริยนฺเต จรติ, มริยาทํ น ภินฺทติ.\csromanspacer{} อยํ อินฺทฺริยสํวรปริยนฺโต.}

\prose{กตโม โภชเน มตฺตญฺญุตาปริยนฺโต, อิธ ภิกฺขุ ปฏิสงฺขา โยนิโส อาหารํ อาหาเรติ, เนว ทวาย น มทาย น มณฺฑนาย น วิภูสนาย ยาวเทว อิมสฺส กายสฺส ฐิติยา ยาปนาย วิหิํสูปรติยา พฺรหฺมจริยานุคฺคหาย, อิติ ปุราณญฺจ เวทนํ ปฏิหงฺขามิ, นวญฺจ เวทนํ น อุปฺปาเทสฺสามิ, ยาตฺรา จ เม ภวิสฺสติ, อนวชฺชตา จ ผาสุวิหาโร จาติ.\csromanspacer{} อกฺขพฺภญฺชนวณปฏิจฺฉาทนปุตฺตมํสูปมํ ปจฺจเวกฺขมาโน อนฺโต โภชเน มตฺตญฺญุตาปริยนฺเต จรติ, มริยาทํ น ภินฺทติ.\csromanspacer{} อยํ โภชเน มตฺตญฺญุตาปริยนฺโต.}

\prose{กตโม ชาคริยานุโยคปริยนฺโต, อิธ ภิกฺขุ ทิวสํ จงฺกเมน นิสชฺชาย อาวรณีเยหิ ธมฺเมหิ จิตฺตํ ปริโสเธติ, รตฺติยา ปฐมํ ยามํ จงฺกเมน นิสชฺชาย อาวรณีเยหิ ธมฺเมหิ จิตฺตํ ปริโสเธติ, รตฺติยา มชฺฌิมํ ยามํ ทกฺขิเณน ปสฺเสน สีหเสยฺยํ กปฺเปติ ปาเท ปาทํ อจฺจาธาย สโต สมฺปชาโน อุฏฺฐานสญฺญํ มนสิ กริตฺวา.\csromanspacer{} รตฺติยา ปจฺฉิมํ ยามํ ปจฺจุฏฺฐาย จงฺกเมน นิสชฺชาย อาวรณีเยหิ ธมฺเมหิ จิตฺตํ ปริโสเธติ.\csromanspacer{} ภทฺเทกรตฺตวิหารํ\footnote{ภทฺเทกรตฺติวิหารํ (สี, ก)} ปจฺจเวกฺขมาโน อนฺโต ชาคริยานุโยคปริยนฺเต จรติ, มริยาทํ น ภินฺทติ.\csromanspacer{} อยํ ชาคริยานุโยคปริยนฺโตติ ภิกฺขุ สโต สปริยนฺตจารี.}

\prose{q\textbf{อํสาธิปาตานํ สรีสปานนฺ}ติ \textbf{ฑํสา} วุจฺจนฺติ ปิงฺคลมกฺขิกาโย.\csromanspacer{} \textbf{อธิปาตกา} วุจฺจนฺติ สพฺพาปิ มกฺขิกาโย.\csromanspacer{} กิํการณา \textbf{อธิปาตกา} วุจฺจนฺติ สพฺพาปิ มกฺขิกาโย.\csromanspacer{} ตา อุปฺปติตฺวา อุปฺปติตฺวา ขาทนฺติ, ตํการณา \textbf{อธิปาตกา} วุจฺจนฺติ สพฺพาปิ มกฺขิกาโย.\csromanspacer{} \textbf{สรีสปา} วุจฺจนฺติ อหีติ \textbf{ฑํสา}ธิปาตานํ สรีสปานํ.}

\csromanheader{2}{16. สาริปุตฺตสุตฺตนิทฺเทส}
\prose{\csromanfolio{387}\textbf{มนุสฺสผสฺสานํ จตุปฺปทานนฺ}ติ มนุสฺสผสฺสา วุจฺจนฺติ โจรา วา อสฺสุ มานวา วา กตกมฺมา วา อกตกมฺมา วา เต ภิกฺขุํ ปญฺหํ วา ปุจฺเฉยฺยุํ, วาทํ วา อาโรเปยฺยุํ อกฺโกเสยฺยุํ ปริภาเสยฺยุํ โรเสยฺยุํ วิโรเสยฺยุํ หิํเสยฺยุํ วิหิํเสยฺยุํ เหเฐยฺยุํ วิเหเฐยฺยุํ ฆาเตยฺยุํ อุปฆาเตยฺยุํ อุปฆาตํ วา กเรยฺยุํ.\csromanspacer{} โย โกจิ มนุสฺสโต อุปฆาโต มนุสฺสผสฺโส.\csromanspacer{} \textbf{จตุปฺปทานนฺ}ติ สีหา พฺยคฺฆา ทีปี อจฺฉา ตรจฺฉา โกกา มหิํสา หตฺถี, เต ภิกฺขุํ มทฺเทยฺยุํ ขาเทยฺยุํ หิํเสยฺยุํ วิหิํเสยฺยุํ เหเฐยฺยุํ วิเหเฐยฺยุํ ฆาเตยฺยุํ อุปฆาเตยฺยุํ อุปฆาตํ วา กเรยฺยุํ.\csromanspacer{} จตุปฺปทโต อุปฆาโต ยํ กิญฺจิ จตุปฺปทภยนฺติ มนุสฺสผสฺสานํ จตุปฺปทานํ.\csromanspacer{} เตนาห ภควา\csromandash{}}

\csromangathameasure{ปญฺจนฺนํ ธีโร ภยานํ น ภาเย, \\ ภิกฺขุ สโต สปริยนฺตจารี. \\ qอํสาธิปาตานํ สรีสปานํ,}
\csromangathagroup{\csromangathastanza{ปญฺจนฺนํ ธีโร ภยานํ น ภาเย, \\ ภิกฺขุ สโต สปริยนฺตจารี. \\ qอํสาธิปาตานํ สรีสปานํ,}}

\prose{\textbf{มนุสฺสผสฺสานํ จตุปฺปทานนฺ}ติ.}

\proseitem{200}{\csromansymbolfootnote{*}{ขุ 1. 427 ปิฏฺเฐ.}ปรธมฺมิกานมฺปิ น สนฺตเสยฺย,}

\prose{\textbf{ทิสฺวาปิ เตสํ พหุเภรวานิ.}}

\csromangathameasure{\textbf{อถาปรานิ อภิสมฺภเวยฺย,} \\ \textbf{ปริสฺสยานิ กุสลานุ-เอสี.}}
\csromangathagroup{\csromangathastanza{\textbf{อถาปรานิ อภิสมฺภเวยฺย,} \\ \textbf{ปริสฺสยานิ กุสลานุ-เอสี.}}}

\prose{\textbf{ปรธมฺมิกานมฺปิ น สนฺตเสยฺย, ทิสฺวาปิ เตสํ พหุเภรวานี}ติ ปรธมฺมิกา วุจฺจนฺติ สตฺต สหธมฺมิเก ฐเปตฺวา เย เกจิ พุทฺเธ ธมฺเม สํเฆ อปฺปสนฺนา, เต ภิกฺขุํ ปญฺหํ วา ปุจฺเฉยฺยุํ, วาทํ วา อาโรเปยฺยุํ อกฺโกเสยฺยุํ ปริภาเสยฺยุํ โรเสยฺยุํ วิโรเสยฺยุํ หิํเสยฺยุํ วิหิํเสยฺยุํ เหเฐยฺยุํ วิเหเฐยฺยุํ ฆาเตยฺยุํ อุปฆาเตยฺยุํ อุปฆาตํ วา กเรยฺยุํ.\csromanspacer{} เตสํ พหุเภรเว ปสฺสิตฺวา วา สุณิตฺวา วา น เวเธยฺย น ปเวเธยฺย น สมฺปเวเธยฺย น ตเสยฺย น สนฺตเสยฺย น อุตฺตเสยฺย น ปริตฺตเสยฺย น ภาเยยฺย น สนฺตาสํ อาปชฺเชยฺย, อภีรู อสฺส อจฺฉมฺภี อนุตฺราสี อปลายี ปหีนภยเภรโว วิคตโลมหํโส วิหเรยฺยาติ ปรธมฺมิกานมฺปิ น สนฺตเสยฺย, ทิสฺวาปิ เตสํ พหุเภรวานิ.}

\prose{\textbf{อถาปรานิ อภิสมฺภเวยฺย, ปริสฺสยานิ กุสลานุ-เอสี}ติ อถาปรานิปิ อตฺถิ อภิสมฺโภตพฺพานิ อภิภวิตพฺพานิ อชฺโฌตฺถริตพฺพานิ \csromanfolio{388}ปริยาทิยิตพฺพานิ มทฺทิตพฺพานิ.\csromanspacer{} \textbf{ปริสฺสยา}ติ เทฺว ปริสฺสยา ปากฏปริสฺสยา จ ปฏิจฺฉนฺนปริสฺสยา จ ฯเปฯ.\csromanspacer{} เอวมฺปิ ตตฺราสยาติ ปริสฺสยา.\csromanspacer{}\csromanspacer{}\textbf{กุสลานุ- เอสี}ติ สมฺมาปฏิปทํ อนุโลมปฏิปทํ อปจฺจนีกปฏิปทํ อวิรุทฺธปฏิปทํ อนฺวตฺถปฏิปทํ ฯเปฯ อริยํ อฏฺฐงฺคิกํ มคฺคํ นิพฺพานญฺจ นิพฺพานคามินิญฺจ ปฏิปทํ เอสนฺเตน คเวสนฺเตน ปริเยสนฺเตน ปริสฺสยา อภิสมฺโภตพฺพา อภิภวิตพฺพา อชฺโฌตฺถริตพฺพา ปริยาทิยิตพฺพา มทฺทิตพฺพาติ อถาปรานิ อภิสมฺภเวยฺย, ปริสฺสยานิ กุสลานุ-เอสี.\csromanspacer{} เตนาห ภควา\csromandash{}}

\csromangathameasure{ปรธมฺมิกานมฺปิ น สนฺตเสยฺย, \\ ทิสฺวาปิ เตสํ พหุเภรวานิ. \\ อถาปรานิ อภิสมฺภเวยฺย, \\ \textbf{ปริสฺสยา}นิ กุสลานุ-เอสีติ.}
\csromangathagroup{\csromangathastanza{ปรธมฺมิกานมฺปิ น สนฺตเสยฺย, \\ ทิสฺวาปิ เตสํ พหุเภรวานิ. \\ อถาปรานิ อภิสมฺภเวยฺย, \\ \textbf{ปริสฺสยา}นิ กุสลานุ-เอสีติ.}}

\proseitem{201}{\csromansymbolfootnote{*}{ขุ 1. 427 ปิฏฺเฐปิ.}\textbf{อาตงฺกผสฺเสน ขุทาย ผุฏฺโฐ},}

\prose{\textbf{สีตํ อถุณฺหํ}\footnote{อตุณฺหํ (สี, ก)} \textbf{อธิวาสเยยฺย.}}

\csromangathameasure{\textbf{โส เตหิ ผุฏฺโฐ พหุธา อโนโก,} \\ \textbf{วีริยปรกฺกมํ ทฬฺหํ กเรยฺย.}}
\csromangathagroup{\csromangathastanza{\textbf{โส เตหิ ผุฏฺโฐ พหุธา อโนโก,} \\ \textbf{วีริยปรกฺกมํ ทฬฺหํ กเรยฺย.}}}

\prose{\textbf{อาตงฺกผสฺเสน ขุทาย ผุฏฺโฐ}ติ อาตงฺกผสฺโส วุจฺจติ โรคผสฺโส, โรคผสฺเสน ผุฏฺโฐ ปเรโต สโมหิโต สมนฺนาคโต อสฺส.\csromanspacer{} จกฺขุโรเคน ผุฏฺโฐ ปเรโต สโมหิโต สมนฺนาคโต อสฺส.\csromanspacer{} โสตโรเคน.\csromanspacer{} ฆานโรเคน.\csromanspacer{} ชิวฺหาโรเคน.\csromanspacer{} กายโรเคน -ปฺ- ฑํสมกสวาตาตปสรีสปสมฺผสฺเสน ผุฏฺโฐ ปเรโต สโมหิโต สมนฺนาคโต อสฺส.\csromanspacer{} ขุทา วุจฺจติ ฉาตโก, ฉาตเกน ผุฏฺโฐ ปเรโต สโมหิโต สมนฺนาคโต อสฺสาติ อาตงฺกผสฺเสน ขุทาย ผุฏฺโฐ.}

\prose{\textbf{สีตํ อถุณฺหํ อธิวาสเยยฺยา}ติ \textbf{สีตนฺ}ติ ทฺวีหิ การเณหิ สีตํ โหติ อพฺภนฺตรธาตุปฺปโกปวเสน\footnote{อพฺภนฺตรธาตุสงฺโกปวเสน (สฺยา)} วา สีตํ โหติ, พหิทฺธา อุตุวเสน วา สีตํ โหติ.\csromanspacer{} \textbf{อุณฺห}นฺติ ทฺวีหิ การเณหิ อุณฺหํ โหติ.\csromanspacer{} อพฺภนฺตรธาตุปฺปโกปวเสน วา อุณฺหํ โหติ, พหิทฺธา อุตุวเสน วา อุณฺหํ โหตีติ สีตํ อถุณฺหํ.\csromanspacer{}\csromanspacer{}\textbf{อธิวาสเยยฺยา}ติ ขโม อสฺส สีตสฺส อุณฺหสฺส ชิฆจฺฉาย ปิปาสาย ฑํสมกสวาตาตปสรีสปสมฺผสฺสานํ}

\vspace{\tipitakaparskip}
\csromancenter{สาริปุตฺตสุตฺตนิทฺเทส ทุรุตฺตานํ ทุราคตานํ วจนปถานํ อุปฺปนฺนานํ สารีริกานํ เวทนานํ ทุกฺขานํ ติพฺพานํ\footnote{ติปฺปานํ (สี, สฺยา)} ขรานํ กฏุกานํ อสาตานํ อมนาปานํ ปาณหรานํ อธิวาสกชาติโก อสฺสาติ สีตํ อถุณฺหํ อธิวาสเยยฺย.}
\prose{\csromanfolio{389}\textbf{โส เตหิ ผุฏฺโฐ พหุธา อโนโก}ติ \textbf{โส เตหี}ติ อาตงฺกผสฺเสน จ ขุทาย จ สีเตน จ อุเณฺหน จ ผุฏฺโฐ ปเรโต สโมหิโต สมนฺนาคโต อสฺสาติ โส เตหิ ผุฏฺโฐ.\csromanspacer{}\csromanspacer{}\textbf{พหุธา}ติ อเนกวิเธหิ อากาเรหิ ผุฏฺโฐ ปเรโต สโมหิโต สมนฺนาคโต อสฺสาติ โส เตหิ ผุฏฺโฐ พหุธา.\csromanspacer{} \textbf{อโนโก}ติ อภิสงฺขารสหคตวิญฺญาณสฺส โอกาสํ น กโรตีติปิ อโนโก.\csromanspacer{} อถ วา กายทุจฺจริตสฺส วจีทุจฺจริตสฺส มโนทุจฺจริตสฺส โอกาสํ น กโรตีติปิ อโนโกติ โส เตหิ ผุฏฺโฐ พหุธา อโนโก.}

\prose{\textbf{วีริยปรกฺกมํ ทฬฺหํ กเรยฺยา}ติ วีริยปรกฺกโม วุจฺจติ โย เจตสิโก วีริยารมฺโภ นิกฺกโม ปรกฺกโม อุยฺยาโม วายาโม อุสฺสาโห อุสฺโสฬฺหี อปฺปฏิวานี ถาโม ธิติ อสิถิลปรกฺกมตา อนิกฺขิตฺตจฺฉนฺทตา อนิกฺขิตฺตธุรตา ธุรสมฺปคฺคาโห วีริยํ วีริยินฺทฺริยํ วีริยพลํ สมฺมาวายาโม.\csromanspacer{} \textbf{วีริยปรกฺกมํ ทฬฺหํ กเรยฺยา}ติ วีริยํ ปรกฺกมํ ทฬฺหํ กเรยฺย ถิรํ กเรยฺย, ทฬฺหสมาทาโน อสฺส อวตฺถิตสมาทาโนติ วีริยํ ปรกฺกมํ ทฬฺหํ กเรยฺย.\csromanspacer{} เตนาห ภควา\csromandash{}}

\csromangathameasure{อาตงฺกผสฺเสน ขุทาย ผุฏฺโฐ, \\ สีตํ อถุณฺหํ อธิวาสเยยฺย.}
\csromangathagroup{\csromangathastanza{อาตงฺกผสฺเสน ขุทาย ผุฏฺโฐ, \\ สีตํ อถุณฺหํ อธิวาสเยยฺย.}}

\prose{\textbf{โส เตหิ ผุฏฺโฐ พหุธา อโนโก},}

\prose{\textbf{วีริยปรกฺกมํ ทฬฺหํ กเรยฺยา}ติ.}

\proseitem{202}{\csromansymbolfootnote{*}{ขุ 1. 427 ปิฏฺเฐปิ.}เถยฺยํ น กาเร น มุสา ภเณยฺย,}

\prose{\textbf{เมตฺตาย ผสฺเส ตสถาวรานิ.}}

\csromangathameasure{\textbf{ยทาวิลตฺตํ มนโส วิชญฺญา,} \\ \textbf{กณฺหสฺส ปกฺโขติ วิโนทเยยฺย.}}
\csromangathagroup{\csromangathastanza{\textbf{ยทาวิลตฺตํ มนโส วิชญฺญา,} \\ \textbf{กณฺหสฺส ปกฺโขติ วิโนทเยยฺย.}}}

\csromanmatikaanchor{mtk.17}
\csromantocmarkref{subsection}{16. สาริปุตฺตสุตฺตนิทฺเทส}{mtk.17}
\bookmark[dest={mtk.17},level=2]{16. สาริปุตฺตสุตฺตนิทฺเทส}
\prose{\textbf{เถยฺยํ น กาเร น มุสา ภเณยฺยา}ติ \textbf{เถยฺยํ น กาเร}ติ อิธ ภิกฺขุ อทินฺนาทานํ ปหาย อทินฺนาทานา ปฏิวิรโต อสฺส ทินฺนาทายี ทินฺนปาฏิกงฺขี \csromanfolio{390}อเถเนน สุจิภูเตน อตฺตนา วิหเรยฺยาติ เถยฺยํ น กาเร.\csromanspacer{} \textbf{น มุสา ภเณยฺยา}ติ อิธ ภิกฺขุ มุสาวาทํ ปหาย มุสาวาทา ปฏิวิรโต อสฺส สจฺจวาที สจฺจสนฺโธ เถโต ปจฺจยิโก อวิสํวาทโก โลกสฺสาติ เถยฺยํ น กาเร น มุสา ภเณยฺย.}

\prose{\textbf{เมตฺตาย ผสฺเส ตสถาวรานี}ติ \textbf{เมตฺตา}ติ ยา สตฺเตสุ เมตฺติ \textbf{เมตฺตา}ยนา \textbf{เมตฺตา}ยิตตฺตํ อนุทา อนุทายนา อนุทายิตตฺตํ\footnote{อนุทฺทยา (สี) เอวมีทิเสสุ ทฺวีสุ ปเทสุปิ ทฺวิภาววเสน. อนุทยา อนุทยนา อนุทยิตตฺตํ (ก) อภิ 1. 215 ปิฏฺเฐ.} หิเตสิตา อนุกมฺปา อพฺยาปาโท อพฺยาปชฺโช\footnote{อพฺยาปชฺโฌ (สฺยา, ก)} อโทโส กุสลมูลํ.\csromanspacer{} \textbf{ตสา}ติ เยสํ ตสิตา ตณฺหา อปฺปหีนา, เยสญฺจ ภยเภรวา อปฺปหีนา.\csromanspacer{} กิํการณา วุจฺจนฺติ ตสา, เต ตสนฺติ อุตฺตสนฺติ ปริตสนฺติ ภายนฺติ สนฺตาสํ อาปชฺชนฺติ, ตํการณา วุจฺจนฺติ ตสา.\csromanspacer{} \textbf{ถาวรา}ติ เยสํ ตสิตา ตณฺหา ปหีนา, เยสญฺจ ภยเภรวา ปหีนา.\csromanspacer{} กิํการณา วุจฺจนฺติ ถาวรา, เต น ตสนฺติ น อุตฺตสนฺติ น ปริตสนฺติ น ภายนฺติ สนฺตาสํ น อาปชฺชนฺติ.\csromanspacer{} ตํการณา วุจฺจนฺติ ถาวรา.\csromanspacer{}\csromanspacer{}\textbf{เมตฺตาย ผสฺเส ตสถาวรานี}ติ ตเส จ ถาวเร จ \textbf{เมตฺตา}ย ผสฺเสยฺย ผเรยฺย \textbf{เมตฺตา}สหคเตน เจตสา วิปุเลน มหคฺคเตน อปฺปมาเณน อเวเรน อพฺยาปชฺเชน ผริตฺวา วิหเรยฺยาติ \textbf{เมตฺตา}ย ผสฺเส ตสถาวรานิ.}

\prose{\textbf{ยทาวิลตฺตํ มนโส วิชญฺญา}ติ \textbf{ยทา}ติ \textbf{ยทา}.\csromanspacer{} \textbf{มนโส}ติ ยํ จิตฺตํ มโน มานสํ หทยํ ปณฺฑรํ มโน มนายตนํ มนินฺทฺริยํ วิญฺญาณํ วิญฺญาณกฺขนฺโธ ตชฺชา มโนวิญฺญาณธาตุ, กายทุจฺจริเตน จิตฺตํ อาวิลํ โหติ ลุฬิตํ เอริตํ ฆฏฺฏิตํ จลิตํ ภนฺตํ อวูปสนฺตํ.\csromanspacer{} วจีทุจฺจริเตน.\csromanspacer{} มโนทุจฺจริเตน.\csromanspacer{} ราเคน.\csromanspacer{} โทเสน.\csromanspacer{} โมเหน.\csromanspacer{} โกเธน.\csromanspacer{} อุปนาเหน.\csromanspacer{} มกฺเขน.\csromanspacer{} ปฬาเสน.\csromanspacer{} อิสฺสาย.\csromanspacer{} มจฺฉริเยน.\csromanspacer{} มายาย.\csromanspacer{} สาเฐยฺเยน.\csromanspacer{} ถมฺเภน.\csromanspacer{} สารมฺเภน.\csromanspacer{} มาเนน.\csromanspacer{} อติมาเนน.\csromanspacer{} มเทน.\csromanspacer{} ปมาเทน.\csromanspacer{} สพฺพกิเลเสหิ.\csromanspacer{} สพฺพทุจฺจริเตหิ.\csromanspacer{} สพฺพทรเถหิ.\csromanspacer{} สพฺพปริฬาเหหิ.\csromanspacer{} สพฺพสนฺตาเปหิ.\csromanspacer{} สพฺพากุสลาภิสงฺขาเรหิ จิตฺตํ อาวิลํ โหติ ลุฬิตํ เอริตํ ฆฏฺฏิตํ จลิตํ ภนฺตํ อวูปสนฺตํ.\csromanspacer{} \textbf{ยทาวิลตฺตํ มนโส วิชญฺญา}ติ}

\vspace{\tipitakaparskip}
\csromancenter{สาริปุตฺตสุตฺตนิทฺเทส จิตฺตสฺส อาวิลภาวํ ชาเนยฺย อาชาเนยฺย วิชาเนยฺย ปฏิวิชาเนยฺย ปฏิวิชฺเฌยฺยาติ ยทาวิลตฺตํ มนโส วิชญฺญา.}
\prose{\csromanfolio{391}\textbf{กณฺหสฺส ปกฺโขติ วิโนทเยยฺยา}ติ \textbf{กโณฺห}ติ โย โส มาโร \textbf{กโณฺห} อธิปติ อนฺตคุ นมุจิ ปมตฺตพนฺธุ, กณฺหสฺส ปกฺโข มารปกฺโข มารปาโส มารพฬิสํ มารามิสํ มารวิสโย มารนิวาโส มารโคจโร มารพนฺธนนฺติ ปชเหยฺย วิโนเทยฺย พฺยนฺติํ กเรยฺย อนภาวํ คเมยฺยาติ เอวมฺปิ กณฺหสฺส ปกฺโขติ วิโนทเยยฺย.\csromanspacer{}\csromanspacer{}อถ วา กณฺหสฺส ปกฺโข มารปกฺโข อกุสลปกฺโข ทุกฺขุทฺทโย ทุกฺขวิปาโก นิรยสํวตฺตนิโก ติรจฺฉานโยนิสํวตฺตนิโก เปตฺติวิสยสํวตฺตนิโกติ ปชเหยฺย วิโนเทยฺย พฺยนฺติํ กเรยฺย อนภาวํ คเมยฺยาติ เอวมฺปิ กณฺหสฺส ปกฺโขติ วิโนทเยยฺย.\csromanspacer{} เตนาห ภควา\csromandash{}}

\csromangathameasure{เถยฺยํ น กาเร น มุสา ภเณยฺย, \\ เมตฺตาย ผสฺเส ตสถาวรานิ. \\ ยทาวิลตฺตํ มนโส วิชญฺญา,}
\csromangathagroup{\csromangathastanza{เถยฺยํ น กาเร น มุสา ภเณยฺย, \\ เมตฺตาย ผสฺเส ตสถาวรานิ. \\ ยทาวิลตฺตํ มนโส วิชญฺญา,}}

\prose{\textbf{กณฺหสฺส ปกฺโขติ วิโนทเยยฺยา}ติ.}

\proseitem{203}{\csromansymbolfootnote{*}{ขุ 1. 427 ปิฏฺเฐปิ.}\textbf{โกธาติมานสฺส วสํ น คจฺเฉ},}

\prose{\textbf{มูลมฺปิ เตสํ ปลิขญฺญ ติฏฺเฐ.}}

\csromangathameasure{\textbf{อถปฺปิยํ วา ปน อปฺปิยํ วา,} \\ \textbf{อทฺธา ภวนฺโต อภิสมฺภเวยฺย.}}
\csromangathagroup{\csromangathastanza{\textbf{อถปฺปิยํ วา ปน อปฺปิยํ วา,} \\ \textbf{อทฺธา ภวนฺโต อภิสมฺภเวยฺย.}}}

\prose{\textbf{โกธาติมานสฺส วสํ น คจฺเฉ}ติ \textbf{โกโธ}ติ โย จิตฺตสฺส อาฆาโต ปฏิฆาโต ฯเปฯ จณฺฑิกฺกํ อสุโรโป อนตฺตมนตา จิตฺตสฺส.\csromanspacer{} \textbf{อติมาโน}ติ อิเธกจฺโจ ปรํ อติมญฺญติ ชาติยา วา โคตฺเตน วา ฯเปฯ อญฺญตรญฺญตเรน วา วตฺถุนา.\csromanspacer{} \textbf{โกธาติมานสฺส วสํ น คจฺเฉ}ติ โกธสฺส จ อติมานสฺส จ วสํ น คจฺเฉยฺย, โกธญฺจ อติมานญฺจ ปชเหยฺย วิโนเทยฺย พฺยนฺติํ กเรยฺย อนภาวํ คเมยฺยาติ โกธาติมานสฺส วสํ น คจฺเฉ.}

\prose{\textbf{มูลมฺปิ เตสํ ปลิขญฺญ ติฏฺเฐ}ติ กตมํ โกธสฺส มูลํ, อวิชฺชา มูลํ อโยนิโส มนสิกาโร มูลํ อสฺมิมาโน มูลํ อหิริกํ มูลํ อโนตฺตปฺปํ มูลํ อุทฺธจฺจํ มูลํ, อิทํ โกธสฺส มูลํ.\csromanspacer{} กตมํ อติมานสฺส มูลํ, อวิชฺชา มูลํ อโยนิโส มนสิกาโร มูลํ อสฺมิมาโน มูลํ อหิริกํ \csromanfolio{392}มูลํ อโนตฺตปฺปํ มูลํ อุทฺธจฺจํ มูลํ, อิทํ อติมานสฺส มูลํ.\csromanspacer{} \textbf{มูลมฺปิ เตสํ ปลิขญฺญ ติฏฺเฐ}ติ โกธสฺส จ อติมานสฺส จ มูลํ ปลิขณิตฺวา อุทฺธริตฺวา สมุทฺธริตฺวา อุปฺปาฏยิตฺวา สมุปฺปาฏยิตฺวา ปชหิตฺวา วิโนเทตฺวา พฺยนฺติํ กริตฺวา อนภาวํ คเมตฺวา ติฏฺเฐยฺย สนฺติฏฺเฐยฺยาติ มูลมฺปิ เตสํ ปลิขญฺญ ติฏฺเฐ.}

\prose{\textbf{อถปฺปิยํ วา ปน อปฺปิยํ วา, อทฺธา ภวนฺโต อภิสมฺภเวยฺยา}ติ \textbf{อถา}ติ ปทสนฺธิ ปทสํสคฺโค ปทปาริปูรี อกฺขรสมวาโย พฺยญฺชนสิลิฏฺฐตา ปทานุปุพฺพตาเปตํ \textbf{อถา}ติ.\csromanspacer{} \textbf{ปิยา}ติ เทฺว ปิยา สตฺตา วา สงฺขารา วา.\csromanspacer{} กตเม สตฺตา ปิยา, อิธ ยา’สฺส\footnote{ยสฺส (สฺยา) สุตฺตมาลา โอโลเกตพฺพา.} เต โหนฺติ อตฺถกามา หิตกามา ผาสุกามา โยคกฺเขมกามา มาตา วา ปิตา วา ภาตา วา ภคินี วา ปุตฺตา วา ธีตา วา มิตฺตา วา อมจฺจา วา ญาตี วา สาโลหิตา วา, อิเม สตฺตา ปิยา.\csromanspacer{} กตเม สงฺขารา ปิยา, มนาปิกา รูปา มนาปิกา สทฺทา คนฺธา รสา โผฏฺฐพฺพา, อิเม สงฺขารา ปิยา.\csromanspacer{} \textbf{อปฺปิยา}ติ เทฺว อปฺปิยา สตฺตา วา สงฺขารา วา.\csromanspacer{} กตเม สตฺตา อปฺปิยา, อิธ ยา’สฺส เต โหนฺติ อนตฺถกามา อหิตกามา อผาสุกามา อโยคกฺเขมกามา ชีวิตา โวโรเปตุกามา, อิเม สตฺตา อปฺปิยา.\csromanspacer{} กตเม สงฺขารา อปฺปิยา, อมนาปิกา รูปา อมนาปิกา สทฺทา คนฺธา รสา โผฏฺฐพฺพา, อิเม สงฺขารา อปฺปิยา.\csromanspacer{} \textbf{อทฺธา}ติ เอกํสวจนํ นิสฺสํสยวจนํ นิกฺกงฺขวจนํ อเทฺวชฺฌวจนํ อเทฺวฬฺหกวจนํ นิยฺยานิกวจนํ อปณฺณกวจนํ อวตฺถาปนวจนเมตํ อทฺธาติ.\csromanspacer{}\csromanspacer{}\textbf{อถปฺปิยํ วา ปน อปฺปิยํ วา, อทฺธา ภวนฺโต อภิสมฺภเวยฺยา}ติ ปิยาปฺปิยํ สาตาสาตํ สุขทุกฺขํ โสมนสฺสโทมนสฺสํ อิฏฺฐานิฏฺฐํ อภิสมฺภวนฺโต วา อภิภเวยฺย อธิภวนฺโต วา อภิสมฺภเวยฺยาติ อถปฺปิยํ วา ปน อปฺปิยํ วา, อทฺธา ภวนฺโต อภิสมฺภเวยฺย.\csromanspacer{} เตนาห ภควา\csromandash{}}

\csromangathameasure{โกธาติมานสฺส วสํ น คจฺเฉ,}
\csromangathagroup{\csromangathastanza{โกธาติมานสฺส วสํ น คจฺเฉ,}}

\prosecont{\textbf{มูลมฺปิ เตสํ ปลิขญฺญ ติฏฺเฐ}.}

\prose{อถปฺปิยํ วา ปน อปฺปิยํ วา,}

\vspace{\tipitakaparskip}
\csromancenter{\textbf{อทฺธา} ภวนฺโต อภิสมฺภเวยฺยาติ.}
\csromanheader{2}{16. สาริปุตฺตสุตฺตนิทฺเทส}
\proseitem{204}{\csromanfolio{393}\csromansymbolfootnote{*}{ขุ 1. 427 ปิฏฺเฐปิ.}\textbf{ปญฺญํ ปุรกฺขตฺวา} กลฺยาณปีติ,}

\prose{\textbf{วิกฺขมฺภเย ตานิ ปริสฺสยานิ.}}

\csromangathameasure{\textbf{อรติํ สเหถ สยนมฺหิ ปนฺเต,} \\ \textbf{จตุโร สเหถ ปริเทวธมฺเม.}}
\csromangathagroup{\csromangathastanza{\textbf{อรติํ สเหถ สยนมฺหิ ปนฺเต,} \\ \textbf{จตุโร สเหถ ปริเทวธมฺเม.}}}

\prose{\textbf{ปญฺญํ ปุรกฺขตฺวา กลฺยาณปีตี}ติ \textbf{ปญฺญา}ติ ยา \textbf{ปญฺญา} ปชานนา วิจโย ปวิจโย ธมฺมวิจโย ฯเปฯ อโมโห สมฺมาทิฏฺฐิ.\csromanspacer{} \textbf{ปญฺญํ ปุรกฺขตฺวา}ติ อิเธกจฺโจ ปญฺญํ ปุรโต กตฺวา จรติ \textbf{ปญฺญา}ธโช \textbf{ปญฺญา}เกตุ \textbf{ปญฺญา}ธิปเตยฺโย วิจยพหุโล ปวิจยพหุโล เปกฺขายนพหุโล\footnote{โอกฺขายนพหุโล (ภูสุ)} สมฺเปกฺขายนพหุโล\footnote{สโมกฺขายนพหุโล (สี, สฺยา)} วิภูตวิหารี ตจฺจริโก ตพฺพหุโล ตคฺครุโก ตนฺนินฺโน ตปฺโปโณ ตปฺปพฺภาโร ตทธิมุตฺโต ตทธิปเตยฺโยติ เอวมฺปิ ปญฺญํ ปุรกฺขตฺวา.}

\prose{อถ วา คจฺฉนฺโต วา “คจฺฉามี”ติ ปชานาติ, ฐิโต วา “ฐิโตมฺหี”ติ ปชานาติ, นิสินฺโน วา “นิสินฺโนมฺหี”ติ ปชานาติ, สยาโน วา “สยาโนมฺหี”ติ ปชานาติ, ยถา ยถา วา ปนสฺส กาโย ปณิหิโต โหติ, ตถา ตถา นํ ปชานาตีติ เอวมฺปิ ปญฺญํ ปุรกฺขตฺวา.}

\prose{อถ วา อภิกฺกนฺเต ปฏิกฺกนฺเต สมฺปชานการี โหติ, อาโลกิเต วิโลกิเต สมฺปชานการี โหติ, สมิญฺชิเต ปสาริเต สมฺปชานการี โหติ, สงฺฆาฏิปตฺตจีวรธารเณ สมฺปชานการี โหติ, อสิเต ปีเต ขายิเต สายิเต สมฺปชานการี โหติ, อุจฺจารปสฺสาวกมฺเม สมฺปชานการี โหติ, คเต ฐิเต นิสินฺเน สุตฺเต ชาคริเต ภาสิเต ตุณฺหิภาเว สมฺปชานการี โหตีติ เอวมฺปิ ปญฺญํ ปุรกฺขตฺวา.}

\prose{\textbf{กลฺยาณปีตี}ติ พุทฺธานุสฺสติวเสน อุปฺปชฺชติ ปีติ ปาโมชฺชํ\footnote{ปามุชฺชํ (สฺยา)} กลฺยาณปีตีติ ธมฺมานุสฺสติ สํฆานุสฺสติ สีลานุสฺสติ จาคานุสฺสติ เทวตานุสฺสติ อานาปานสฺสติ มรณสฺสติ กายคตาสติวเสน อุปสมานุสฺสติวเสน อุปฺปชฺชติ ปีติ ปาโมชฺชํ กลฺยาณปีตีติ ปญฺญํ ปุรกฺขตฺวา กลฺยาณปีติ.}

\prose{\textbf{วิกฺขมฺภเย ตานิ ปริสฺสยานี}ติ \textbf{ปริสฺสยา}ติ เทฺว \textbf{ปริสฺสยา} ปากฏ\textbf{ปริสฺสยา} จ ปฏิจฺฉนฺน\textbf{ปริสฺสยา} จ ฯเปฯ อิเม วุจฺจนฺติ ปากฏ\textbf{ปริสฺสยา} ฯเปฯ \csromanfolio{394}อิเม วุจฺจนฺติ ปฏิจฺฉนฺนปริสฺสยา ฯเปฯ เอวมฺปิ ตตฺราสยาติ ปริสฺสยา\textbf{.}\csromanspacer{}\csromanspacer{}\textbf{วิกฺขมฺภเย ตานิ ปริสฺสยานี}ติ ตานิ ปริสฺสยานิ วิกฺขมฺเภยฺย อภิภเวยฺย อชฺโฌตฺถเรยฺย ปริยาทิเยยฺย มทฺเทยฺยาติ วิกฺขมฺภเย ตานิ ปริสฺสยานิ\textbf{.}}

\prose{\textbf{อรติํ สเหถ สยนมฺหิ ปนฺเต}ติ \textbf{อรตี}ติ ยา อรติ อรติตา อนภิรติ อนภิรมณา\footnote{อนภิรมนา (พหูสุ) อภิ 2. 365 ปิฏฺเฐ.} อุกฺกณฺฐิตา ปริตสฺสิตา\textbf{.}\csromanspacer{} \textbf{สยนมฺหิ ปนฺเต}ติ ปนฺเตสุ วา เสนาสเนสุ อญฺญตรญฺญตเรสุ วา อธิกุสเลสุ ธมฺเมสุ อรติํ สเหยฺย อภิภเวยฺย อชฺโฌตฺถเรยฺย ปริยาทิเยยฺย มทฺเทยฺยาติ อรติํ สเหถ สยนมฺหิ ปนฺเต\textbf{.}}

\prose{\textbf{จตุโร สเหถ ปริเทวธมฺเม}ติ จตฺตาโร ปริเทวนีเย ธมฺเม สเหยฺย ปริสเหยฺย อภิภเวยฺย อชฺโฌตฺถเรยฺย ปริยาทิเยยฺย มทฺเทยฺยาติ จตุโร สเหถ ปริเทวธมฺเม\textbf{.}\csromanspacer{} เตนาห ภควา\csromandash{}}

\csromangathameasure{ปญฺญํ ปุรกฺขตฺวา กลฺยาณปีติ,}
\csromangathagroup{\csromangathastanza{ปญฺญํ ปุรกฺขตฺวา กลฺยาณปีติ,}}

\prosecont{วิกฺขมฺภเย ตานิ ปริสฺสยานิ\textbf{.}}

\prose{\textbf{อรติํ สเหถ สยนมฺหิ ปนฺเต},}

\prose{\textbf{จตุโร สเหถ ปริเทวธมฺเม}ติ\textbf{.}}

\proseitem{205}{\csromansymbolfootnote{*}{ขุ 1. 427 ปิฏฺเฐปิ.}กิํสู อสิสฺสํ กุว วา\footnote{กุหิํ วา (สี), กุวํ วา (สฺยา), กุถ วา (ก)} อสิสฺสํ,}

\prose{\textbf{ทุกฺขํ วต เสตฺถ กุวชฺช เสสฺสํ.}}

\csromangathameasure{\textbf{เอเต วิตกฺเก ปริเทวเนยฺเย,} \\ \textbf{วินเยถ เสโข อนิเกตจารี} 3\textbf{.}}
\csromangathagroup{\csromangathastanza{\textbf{เอเต วิตกฺเก ปริเทวเนยฺเย,} \\ \textbf{วินเยถ เสโข อนิเกตจารี} 3\textbf{.}}}

\prose{\textbf{กิํสู อสิสฺสํ กุว วา อสิสฺสนฺ}ติ กิํสุ อสิสฺสามีติ กิํ ภุญฺชิสฺสามิ โอทนํ วา กุมฺมาสํ วา สตฺตุํ วา มจฺฉํ วา มํสํ วาติ กิํสุ อสิสฺสํ\textbf{.}\csromanspacer{} \textbf{กุว วา อสิสฺสนฺ}ติ กตฺถ ภุญฺชิสฺสามิ ขตฺติยกุเล วา พฺราหฺมณกุเล วา เวสฺสกุเล วา สุทฺทกุเล วาติ กิํสู อสิสฺสํ กุว วา อสิสฺสํ\textbf{.}}

\prose{\textbf{ทุกฺขํ วต เสตฺถ กุวชฺช เสสฺสนฺ}ติ อิมํ รตฺติํ ทุกฺขํ สยิตฺถ ผลเก วา ตฏฺฏิกาย วา จมฺมขณฺเฑ วา ติณสนฺถาเร วา ปณฺณสนฺถาเร วา ปลาลสนฺถาเร วา อาคามิรตฺติํ\footnote{อาคมนรตฺติํ (สฺยา)} กตฺถ สุขํ สยิสฺสามิ มญฺเจ วา ปีเฐ วา}

\vspace{\tipitakaparskip}
\csromancenter{สาริปุตฺตสุตฺตนิทฺเทส ภิสิยา วา พิมฺโพหเน วา วิหาเร วา อฑฺฒโยเค วา ปาสาเท วา หมฺมิเย วา คุหาย วาติ ทุกฺขํ วต เสตฺถ กุวชฺช เสสฺสํ.}
\prose{\csromanfolio{395}\textbf{เอเต วิตกฺเก ปริเทวเนยฺเย}ติ \textbf{เอเต วิตกฺเก}ติ เทฺว ปิณฺฑปาตปฏิสญฺญุตฺเต วิตกฺเก, เทฺว เสนาสนปฏิสญฺญุตฺเต วิตกฺเก.\csromanspacer{} \textbf{ปริเทวเนยฺเย}ติ อาเทวเนยฺเย ปริเทวเนยฺเยติ \textbf{เอเต วิตกฺเก} ปริเทวเนยฺเย.}

\prose{\textbf{วินเยถ เสโข อนิเกตจารี}ติ \textbf{เสโข}ติ กิํการณา วุจฺจติ \textbf{เสโข}, สิกฺขตีติ \textbf{เสโข}.\csromanspacer{} กิญฺจ สิกฺขติ, อธิสีลมฺปิ สิกฺขติ, อธิจิตฺตมฺปิ สิกฺขติ, อธิปญฺญมฺปิ สิกฺขติ.\csromanspacer{} กตมา อธิสีลสิกฺขา ฯเปฯ อยํ อธิปญฺญาสิกฺขา.\csromanspacer{} อิมา ติสฺโส สิกฺขาโย อาวชฺชนฺโต สิกฺขติ, ชานนฺโต ปสฺสนฺโต ปจฺจเวกฺขนฺโต จิตฺตํ อธิฏฺฐหนฺโต สิกฺขติ, สทฺธาย อธิมุจฺจนฺโต สิกฺขติ, วีริยํ ปคฺคณฺหนฺโต, สติํ อุปฏฺฐเปนฺโต, จิตฺตํ สมาทหนฺโต, ปญฺญาย ปชานนฺโต สิกฺขติ, อภิญฺเญยฺยํ อภิชานนฺโต สิกฺขติ, ปริญฺเญยฺยํ ปริชานนฺโต, ปหาตพฺพํ ปชหนฺโต, ภาเวตพฺพํ ภาเวนฺโต, สจฺฉิกาตพฺพํ สจฺฉิกโรนฺโต สิกฺขติ อาจรติ สมาจรติ สมาทาย สิกฺขติ.\csromanspacer{} ตํการณา วุจฺจติ \textbf{เสโข}.\csromanspacer{} เสโข วินยาย ปฏิวินยาย ปหานาย วูปสมาย ปฏินิสฺสคฺคาย ปฏิปสฺสทฺธิยา อธิสีลมฺปิ สิกฺเขยฺย, อธิจิตฺตมฺปิ สิกฺเขยฺย, อธิปญฺญมฺปิ สิกฺเขยฺย, อิมา ติสฺโส สิกฺขาโย อาวชฺชนฺโต สิกฺเขยฺย, ชานนฺโต ฯเปฯ สจฺฉิกาตพฺพํ สจฺฉิกโรนฺโต สิกฺเขยฺย อาจเรยฺย สมาจเรยฺย สมาทาย วตฺเตยฺยาติ วินเยถ \textbf{เสโข}.}

\prose{\textbf{อนิเกตจารี}ติ กถํ นิเกตจารี โหติ, อิเธกจฺโจ กุลปลิโพเธน สมนฺนาคโต โหติ, คณปลิโพเธน.\csromanspacer{} อาวาสปลิโพเธน.\csromanspacer{} จีวรปลิโพเธน.\csromanspacer{} ปิณฺฑปาตปลิโพเธน.\csromanspacer{} เสนาสนปลิโพเธน.\csromanspacer{} คิลาน ปจฺจย เภสชฺช ปริกฺขาร ปลิโพเธน สมนฺนาคโต โหติ.\csromanspacer{} เอวํ นิเกตจารี โหติ.\csromanspacer{}\csromanspacer{}กถํ อนิเกตจารี โหติ, อิธ ภิกฺขุ น กุลปลิโพเธน สมนฺนาคโต โหติ, น คณปลิโพเธน.\csromanspacer{} น อาวาสปลิโพเธน.\csromanspacer{} น จีวรปลิโพเธน.\csromanspacer{} น ปิณฺฑปาตปลิโพเธน.\csromanspacer{} น เสนาสนปลิโพเธน.\csromanspacer{} น คิลานปจฺจยเภสชฺชปริกฺขารปลิโพเธน สมนฺนาคโต โหติ.\csromanspacer{} เอวํ อนิเกตจารี โหติ.}

\csromangathameasure{“มคธํ คตา โกสลํ คตา, \\ เอกจฺจิยา ปน วชฺชิภูมิยา. \\ มิคา วิย อสงฺฆจาริโน, \\ อนิเกตา วิหรนฺติ ภิกฺขโว. \\ สาธุ จริตกํ สาธุ สุจริตํ, \\ สาธุ สทา อนิเกตวิหาโร. \\ อตฺถปุจฺฉนํ ปทกฺขิณํ กมฺมํ, \\ เอตํ สามญฺญํ อกิญฺจนสฺสา”ติ.}
\csromangathagroup{\csromangathastanza{\csromanfolio{396}“มคธํ คตา โกสลํ คตา, \\ เอกจฺจิยา ปน วชฺชิภูมิยา. \\ มิคา วิย อสงฺฆจาริโน\footnote{มาคธา วิสงฺฆจาริโน (สฺยา)}, \\ อนิเกตา วิหรนฺติ ภิกฺขโว.}\csromangathastanza{สาธุ จริตกํ สาธุ สุจริตํ, \\ สาธุ สทา อนิเกตวิหาโร. \\ อตฺถปุจฺฉนํ ปทกฺขิณํ กมฺมํ, \\ เอตํ สามญฺญํ อกิญฺจนสฺสา”ติ.}}

\csromanheader{2}{วินเยถ เสโข อนิเกตจารี. เตนาห ภควา\csromandash{}}
\csromangathameasure{กิํสู อสิสฺสํ กุว วา อสิสฺสํ, \\ ทุกฺขํ วต เสตฺถ กุวชฺช เสสฺสํ. \\ เอเต วิตกฺเก ปริเทวเนยฺเย, \\ วินเยถ เสโข อนิเกตจารีติ.}
\csromangathagroup{\csromangathastanza{กิํสู อสิสฺสํ กุว วา อสิสฺสํ, \\ ทุกฺขํ วต เสตฺถ กุวชฺช เสสฺสํ. \\ เอเต วิตกฺเก ปริเทวเนยฺเย, \\ วินเยถ เสโข อนิเกตจารีติ.}}

\proseitem{206}{\csromansymbolfootnote{*}{ขุ 1. 427 ปิฏฺเฐปิ.}\textbf{อนฺนญฺจ ลทฺธา วสนญฺจ กาเล},}

\prose{\textbf{มตฺตํ ส ชญฺญา}\footnote{มตฺตํ โส ชญฺญา (สฺยา)} \textbf{อิธ โตสนตฺถํ.}}

\csromangathameasure{\textbf{โส เตสุ คุตฺโต ยตจาริ คาเม,} \\ \textbf{รุสิโตปิ} \textbf{วาจํ ผรุสํ น วชฺชา.}}
\csromangathagroup{\csromangathastanza{\textbf{โส เตสุ คุตฺโต ยตจาริ คาเม,} \\ \textbf{รุสิโตปิ}\footnote{ทูสิโตปิ (ก)} \textbf{วาจํ ผรุสํ น วชฺชา.}}}

\prose{\textbf{อนฺนญฺจ ลทฺธา วสนญฺจ กาเล}ติ \textbf{อนฺนนฺ}ติ โอทโน กุมฺมาโส สตฺตุ มจฺโฉ มํสํ.\csromanspacer{} \textbf{วสนนฺ}ติ ฉ จีวรานิ โขมํ กปฺปาสิกํ โกเสยฺยํ กมฺพลํ สาณํ ภงฺคํ.\csromanspacer{} \textbf{อนฺนญฺจ ลทฺธา วสนญฺจ กาเล}ติ จีวรํ ลภิตฺวา ปิณฺฑปาตํ ลภิตฺวา น กุหนาย น ลปนาย, น เนมิตฺติกตาย, น นิปฺเปสิกตาย, น ลาเภน ลาภํ นิชิคีสนตาย, น ทารุทาเนน, น เวฬุทาเนน, น ปตฺตทาเนน, น ปุปฺผทาเนน, น ผลทาเนน, น สินานทาเนน, น จุณฺณทาเนน, น มตฺติกาทาเนน, น ทนฺตกฏฺฐทาเนน, น มุโขทกทาเนน, น จาฏุกมฺยตาย, น มุคฺคสูปฺยตาย, น ปาริภฏฺยตาย, น ปีฐมทฺทิกตาย, น วตฺถุวิชฺชาย, น ติรจฺฉานวิชฺชาย, น องฺควิชฺชาย, น นกฺขตฺตวิชฺชาย, น ทูตคมเนน, น ปหิณคมเนน, น ชงฺฆเปสนิเยน, น เวชฺชกมฺเมน, น ปิณฺฑปฏิปิณฺฑเกน, น ทานานุปฺปทาเนน, ธมฺเมน สเมน ลทฺธา ลภิตฺวา อธิคนฺตฺวา วินฺทิตฺวา ปฏิลภิตฺวาติ อนฺนญฺจ ลทฺธา วสนญฺจ กาเล.}

\csromanheader{2}{16. สาริปุตฺตสุตฺตนิทฺเทส}
\prose{\csromanfolio{397}\textbf{มตฺตํ ส ชญฺญา อิธ โตสนตฺถ}นฺติ \textbf{มตฺตํ ส ชญฺญา}ติ ทฺวีหิ การเณหิ มตฺตํ ชาเนยฺย ปฏิคฺคหณโต วา ปริโภคโต วา.\csromanspacer{} กถํ ปฏิคฺคหณโต มตฺตํ ชานาติ, โถเกปิ ทิยฺยมาเน กุลานุทยาย\footnote{กุลานุทฺทยาย (สี, ก)} กุลานุรกฺขาย กุลานุกมฺปาย ปฏิคฺคณฺหาติ, พหุเกปิ ทิยฺยมาเน กายปริหาริกํ จีวรํ ปฏิคฺคณฺหาติ, กุจฺฉิปริหาริกํ ปิณฺฑปาตํ ปฏิคฺคณฺหาติ, เอวํ ปฏิคฺคหณโต มตฺตํ ชานาติ.\csromanspacer{}\csromanspacer{}กถํ ปริโภคโต มตฺตํ ชานาติ\csromandash{} * ปฏิสงฺขา โยนิโส จีวรํ ปฏิเสวติ ยาวเทว สีตสฺส ปฏิฆาตาย, อุณฺหสฺส ปฏิฆาตาย, ฑํสมกสวาตาตปสรีสปสมฺผสฺสานํ ปฏิฆาตาย, ยาวเทว หิริโกปีนปฏิจฺฉาทนตฺถํ.\csromanspacer{} ปฏิสงฺขา โยนิโส ปิณฺฑปาตํ ปฏิเสวติ เนว ทวาย, น มทาย, น มณฺฑนาย, น วิภูสนาย, ยาวเทว อิมสฺส กายสฺส ฐิติยา, ยาปนาย, วิหิํสูปรติยา, พฺรหฺมจริยานุคฺคหาย, อิติ ปุราณญฺจ เวทนํ ปฏิหงฺขามิ, นวญฺจ เวทนํ น อุปฺปาเทสฺสามิ, ยาตฺรา จ เม ภวิสฺสติ อนวชฺชตา จ ผาสุวิหาโร จ.\csromanspacer{} ปฏิสงฺขา โยนิโส เสนาสนํ ปฏิเสวติยาวเทว สีตสฺส ปฏิฆาตาย, อุณฺหสฺส ปฏิฆาตาย ฑํสมกสวาตาตปสรี สปสมฺผสฺสานํ ปฏิฆาตาย, ยาวเทว อุตุปริสฺสยวิโนทนปฏิสลฺลานารามตฺถํ.\csromanspacer{} ปฏิสงฺขา โยนิโส คิลานปจฺจยเภสชฺชปริกฺขารํ ปฏิเสวติ ยาวเทว อุปฺปนฺนานํ เวยฺยาพาธิกานํ\footnote{เวยฺยาพฺยาธิกานํ (ก)} เวทนานํ ปฏิฆาตาย, อพฺยาพชฺฌปรมตาย.}

\prose{เอวํ ปริโภคโต มตฺตํ ชานาติ.\csromanspacer{} มตฺตํ ส ชญฺญาติ อิเมหิ ทฺวีหิการเณหิ มตฺตํ ชาเนยฺย อาชาเนยฺย ปฏิวิชาเนยฺย ปฏิวิชฺเฌยฺยาติ \textbf{มตฺตํ ส ชญฺญา}.}

\prose{\textbf{อิธ โตสนตฺถนฺ}ติ + อิธ ภิกฺขุ สนฺตุฏฺโฐ โหติ อิตรีตเรน จีวเรน, อิตรีตรจีวรสนฺตุฏฺฐิยา จ วณฺณวาที, น จ จีวรเหตุ อเนสนํ อปฺปติรูปํ อาปชฺชติ, อลทฺธา จ จีวรํ น ปริตสฺสติ, ลทฺธา จ จีวรํ อคธิโต อมุจฺฉิโต อนชฺโฌสนฺโน อาทีนวทสฺสาวี นิสฺสรณปญฺโญ ปริภุญฺชติ,}

\prose{\csromanfolio{398}ตาย จ ปน อิตรีตรจีวรสนฺตุฏฺฐิยา เนวตฺตานุกฺกํเสติ, น ปรํ วมฺเภติ.\csromanspacer{} โย หิ ตตฺถ ทกฺโข อนลโส สมฺปชาโน ปติสฺสโต, อยํ วุจฺจติ ภิกฺขุ โปราเณ อคฺคญฺเญ อริยวํเส ฐิโต.}

\prose{ปุน จปรํ ภิกฺขุ สนฺตุฏฺโฐ โหติ อิตรีตเรน ปิณฺฑปาเตน, อิตรีตรปิณฺฑปาตสนฺตุฏฺฐิยา จ วณฺณวาที, น จ ปิณฺฑปาตเหตุ อเนสนํ อปฺปติรูปํ อาปชฺชติ, อลทฺธา จ ปิณฺฑปาตํ น ปริตสฺสติ, ลทฺธา จ ปิณฺฑปาตํ อคธิโต อมุจฺฉิโต อนชฺโฌสนฺโน อาทีนวทสฺสาวี นิสฺสรณปญฺโญ ปริภุญฺชติ, ตาย จ ปน อิตรีตรปิณฺฑปาตสนฺตุฏฺฐิยา เนวตฺตานุกฺกํเสติ, น ปรํ วมฺเภติ.\csromanspacer{} โย หิ ตตฺถ ทกฺโข อนลโส สมฺปชาโน ปติสฺสโต, อยํ วุจฺจติ ภิกฺขุ โปราเณ อคฺคญฺเญ อริยวํเส ฐิโต.}

\prose{ปุน จปรํ ภิกฺขุ สนฺตุฏฺโฐ โหติ อิตรีตเรน เสนาสเนน, อิตรีตรเสนาสนสนฺตุฏฺฐิยา จ วณฺณวาที, น จ เสนาสนเหตุ อเนสนํ อปฺปติรูวํ อาปชฺชติ, อลทฺธา จ เสนาสนํ น ปริตสฺสติ, ลทฺธา จ เสนาสนํ อคธิโต อมุจฺฉิโต อนชฺโฌสนฺโน อาทีนวทสฺสาวี นิสฺสรณปญฺโญ ปริภุญฺชติ, ตาย จ ปน อิตรีตรเสนาสนสนฺตุฏฺฐิยา เนวตฺตานุกฺกํเสติ, น ปรํ วมฺเภติ.\csromanspacer{} โย หิ ตตฺถ ทกฺโข อนลโส สมฺปชาโน ปติสฺสโต, อยํ วุจฺจติ ภิกฺขุ โปราเณ อคฺคญฺเญ อริยวํเส ฐิโต.}

\prose{ปุน จปรํ ภิกฺขุ สนฺตุฏฺโฐ โหติ อิตรีตเรน คิลานปจฺจยเภสชฺชปริกฺขาเรน, อิตรีตรคิลานปจฺจยเภสชฺช- ปริกฺขารสนฺตุฏฺฐิยา จ วณฺณวาที, น จ คิลานปจฺจยเภสชฺช- ปริกฺขารเหตุ อเนสนํ อปฺปติรูปํ อาปชฺชติ, อลทฺธา จ คิลานปจฺจยเภสชฺชปริกฺขารํ น ปริตสฺสติ, ลทฺธา จ คิลานปจฺจยเภสชฺชปริกฺขารํ อคธิโต อมุจฺฉิโต อนชฺโฌสนฺโน อาทีนวทสฺสาวี นิสฺสรณปญฺโญ ปริภุญฺชติ, ตาย จ ปน อิตรีตรคิลาน- ปจฺจยเภสชฺชปริกฺขารสนฺตุฏฺฐิยา เนวตฺตานุกฺกํเสติ, น ปรํ วมฺเภติ.\csromanspacer{} โย หิ ตตฺถ ทกฺโข อนลโส สมฺปชาโน ปติสฺสโต, อยํ วุจฺจติ ภิกฺขุ โปราเณ อคฺคญฺเญ อริยวํเส ฐิโตติ มตฺตํ ส ชญฺญา อิธ โตสนตฺถํ.}

\csromanheader{2}{16. สาริปุตฺตสุตฺตนิทฺเทส}
\prose{\csromanfolio{399}\textbf{โส เตสุ คุตฺโต ยตจาริ คาเม}ติ \textbf{โส เตสุ คุตฺโต}ติ จีวเร ปิณฺฑปาเต เสนาสเน คิลานปจยเภสชฺชปริกฺขาเร คุตฺโต โคปิโต รกฺขิโต สํวุโตติ เอวมฺปิ \textbf{โส เตสุ คุตฺโต.}\csromanspacer{}\csromanspacer{}อถ วา อายตเนสุ คุตฺโต โคปิโต รกฺขิโต สํวุโตติ เอวมฺปิ \textbf{โส เตสุ คุตฺโต.}}

\prose{\textbf{ยตจาริ คาเม}ติ คาเม ยโต ยตฺโต ปฏิยตฺโต คุตฺโต โคปิโต รกฺขิโต สํวุโตติ \textbf{โส เตสุ คุตฺโต} ยตจาริ คาเม\textbf{.}}

\prose{\textbf{รุสิโตปิ วาจํ ผรุสํ น วชฺชา}ติ ทูสิโต ขุํสิโต วมฺภิโต ฆฏฺฏิโต ครหิโต อุปวทิโต ผรุเสน กกฺขเฬน นปฺปฏิวชฺชา นปฺปฏิภเณยฺย, อกฺโกสนฺตํ น ปจฺจกฺโกเสยฺย, โรสนฺตํ นปฺปฏิโรเสยฺย, ภณฺฑนฺตํ น ปฏิภณฺเฑยฺย, น กลหํ กเรยฺย, น ภณฺฑนํ กเรยฺย, น วิคฺคหํ กเรยฺย, น วิวาทํ กเรยฺย, น เมธคํ กเรยฺย, กลหํ ภณฺฑนํ วิคฺคหํ วิวาทํ เมธคํ ปชเหยฺย วิโนเทยฺย พฺยนฺติํ กเรยฺย อนภาวํ คเมยฺย, กลหภณฺฑนวิคฺคหวิวาทเมธคา อารโต อสฺส วิรโต ปฏิวิรโต นิกฺขนฺโต นิสฺสโฏ วิปฺปมุตฺโต วิสญฺญุตฺโต วิมริยาทิกเตน เจตสา วิหเรยฺยาติ รุสิโตปิ วาจํ ผรุสํ น วชฺชา\textbf{.}\csromanspacer{} เตนาห ภควา\csromandash{}}

\csromangathameasure{อนฺนญฺจ ลทฺธา วสนญฺจ กาเล, \\ มตฺตํ ส ชญฺญา อิธ โตสนตฺถํ\textbf{.} \\ \textbf{โส เตสุ คุตฺโต ยตจาริ คาเม},}
\csromangathagroup{\csromangathastanza{อนฺนญฺจ ลทฺธา วสนญฺจ กาเล, \\ มตฺตํ ส ชญฺญา อิธ โตสนตฺถํ\textbf{.} \\ \textbf{โส เตสุ คุตฺโต ยตจาริ คาเม},}}

\prose{\textbf{รุสิโตปิ วาจํ ผรุสํ น วชฺชา}ติ\textbf{.}}

\proseitem{207}{\csromansymbolfootnote{*}{ขุ 1. 428 ปิฏฺเฐปิ.}\textbf{โอกฺขิตฺตจกฺขุ น จ ปาทโลโล},}

\prose{\textbf{ฌานานุยุตฺโต พหุชาครสฺส.}}

\csromangathameasure{\textbf{อุเปกฺขมารพฺภ สมาหิตตฺโต,} \\ \textbf{ตกฺกาสยํ กุกฺกุจฺจญฺจุปจฺฉินฺเท}\textbf{.}}
\csromangathagroup{\csromangathastanza{\textbf{อุเปกฺขมารพฺภ สมาหิตตฺโต,} \\ \textbf{ตกฺกาสยํ กุกฺกุจฺจญฺจุปจฺฉินฺเท}\footnote{กุกฺกุจฺจิยูปจฺฉินฺเท (สฺยา)}\textbf{.}}}

\prose{\textbf{โอกฺขิตฺตจกฺขุ น จ ปาทโลโล}ติ กถํ ขิตฺตจกฺขุ โหติ, อิเธกจฺโจ ภิกฺขุ จกฺขุโลโล จกฺขุโลลิเยน สมนฺนาคโต โหติ, “อทิฏฺฐํ ทกฺขิตพฺพํ, ทิฏฺฐํ สมติกฺกมิตพฺพนฺ”ติ อาราเมน อารามํ อุยฺยาเนน อุยฺยานํ คาเมน คามํ นิคเมน นิคมํ นคเรน นครํ รฏฺเฐน รฏฺฐํ ชนปเทน ชนปทํ \csromanfolio{400}ทีฆจาริกํ อนวฏฺฐิตจาริกํ อนุยุตฺโต จ โหติ รูปทสฺสนาย.\csromanspacer{} เอวมฺปิ ขิตฺตจกฺขุ โหติ.}

\prose{อถ วา ภิกฺขุ อนฺตรฆรํ ปวิฏฺโฐ วีถิํ ปฏิปนฺโน อสํวุโต คจฺฉติ หตฺถิํ โอโลเกนฺโต, อสฺสํ โอโลเกนฺโต, รถํ โอโลเกนฺโต, ปตฺติํ โอโลเกนฺโต, อิตฺถิโย โอโลเกนฺโต, ปุริเส โอโลเกนฺโต, กุมารเก โอโลเกนฺโต, กุมาริกาโย โอโลเกนฺโต, อนฺตราปณํ โอโลเกนฺโต, ฆรมุขานิ โอโลเกนฺโต, อุทฺธํ โอโลเกนฺโต, อโธ โอโลเกนฺโต, ทิสาวิทิสํ วิเปกฺขมาโน\footnote{เปกฺขมาโน (พหูสุ) เหฏฺฐา 285 ปิฏฺเฐปิ.} คจฺฉติ, เอวมฺปิ ขิตฺตจกฺขุ โหติ.}

\prose{อถ วา ภิกฺขุ จกฺขุนา รูปํ ทิสฺวา นิมิตฺตคฺคาหี โหติ อนุพฺยญฺชนคฺคาหี, ยตฺวาธิกรณเมนํ จกฺขุนฺทฺริยํ อสํวุตํ วิรหนฺตํ อภิชฺฌาโทมนสฺสา ปาปกา อกุสลา ธมฺมา อนฺวาสฺสเวยฺยุํ, ตสฺส สํวราย นปฺปฏิปชฺชติ, น รกฺขติ จกฺขุนฺทฺริยํ, จกฺขุนฺทฺริเย น สํวรํ อาปชฺชติ.\csromanspacer{} เอวมฺปิ ขิตฺตจกฺขุ โหติ.}

\prose{ยถา วา ปเนเก โภนฺโต สมณพฺราหฺมณา สทฺธาเทยฺยานิ โภชนานิ ภุญฺชิตฺวา เต เอวรูปํ วิสูกทสฺสนํ อนุยุตฺโต วิหรนฺติ, เสยฺยถิทํ, นจฺจํ คีตํ วาทิตํ เปกฺขํ อกฺขานํ ปาณิสฺสรํ เวตาฬํ กุมฺภถูณํ โสภนกํ จณฺฑาลํ วํสํ โธวนํ หตฺถิยุทฺธํ อสฺสยุทฺธํ มหิํสยุทฺธํ อุสภยุทฺธํ อชยุทฺธํ เมณฺฑยุทฺธํ กุกฺกุฏยุทฺธํ วฏฺฏกยุทฺธํ ทณฺฑยุทฺธํ มุฏฺฐิลุทฺธํ นิพฺพุทฺธํ อุยฺโยธิกํ พลคฺคํ เสนาพฺยูหํ อนีกทสฺสนํ อิติ วา.\csromanspacer{} เอวมฺปิ ขิตฺตจกฺขุ โหติ.}

\prose{กถํ น ขิตฺตจกฺขุ โหติ, อิเธกจฺโจ ภิกฺขุ น จกฺขุโลโล น จกฺขุโลลิเยน สมนฺนาคโต โหติ, “อทิฏฺฐํ ทกฺขิตพฺพํ, ทิฏฺฐํ สมติกฺกมิตพฺพนฺ”ติ น อาราเมน อารามํ น อุยฺยาเนน อุยฺยานํ น คาเมน คามํ น นิคเมน นิคมํ น นคเรน นครํ น รฏฺเฐน รฏฺฐํ น ชนปเทน ชนปทํ ทีฆจาริกํ อนวฏฺฐิตจาริกํ อนนุยุตฺโต จ โหติ รูปทสฺสนาย.\csromanspacer{} เอวมฺปิ น ขิตฺตจกฺขุ โหติ.}

\prose{อถ วา ภิกฺขุ อนฺตรฆรํ ปวิฏฺโฐ วีถิํ ปฏิปนฺโน สํวุโต คจฺฉติ น หตฺถิํ โอโลเกนฺโต ฯเปฯ น ทิสาวิทิสํ วิเปกฺขมาโน คจฺฉติ.\csromanspacer{} เอวมฺปิ น ขิตฺตจกฺขุ โหติ.}

\csromanheader{2}{16. สาริปุตฺตสุตฺตนิทฺเทส}
\prose{\csromanfolio{401}อถ วา ภิกฺขุ จกฺขุนา รูปํ ทิสฺวา น นิมิตฺตคฺคาหี โหติ ฯเปฯ จกฺขุนฺทฺริเย สํวรํ อาปชฺชติ.\csromanspacer{} เอวมฺปิ น ขิตฺตจกฺขุ โหติ.}

\prose{ยถา วา ปเนเก โภนฺโต สมณพฺราหฺมณา สทฺธาเทยฺยานิ โภชนานิ ภุญฺชิตฺวา ฯเปฯ อนีกทสฺสนํ อิติ วา.\csromanspacer{} เอวรูปา วิสูกทสฺสนานุโยคา ปฏิวิรโต โหติ.\csromanspacer{} เอวมฺปิ น ขิตฺตจกฺขุ โหตีติ โอกฺขิตฺตจกฺขุ.}

\prose{\textbf{น จ ปาทโลโล}ติ กถํ ปาทโลโล โหติ, อิเธกจฺโจ ภิกฺขุ ปาทโลโล ปาทโลลิเยน สมนฺนาคโต โหติ, อาราเมน อารามํ ฯเปฯ ทีฆจาริกํ อนวฏฺฐิตจาริกํ อนุยุตฺโต โหติ รูปทสฺสนาย.\csromanspacer{} เอวมฺปิ ปาทโลโล โหติ.}

\prose{อถ วา ภิกฺขุ อนฺโตปิ สํฆาราเม ปาทโลโล ปาทโลลิเยน สมนฺนาคโต โหติ, น อตฺถเหตุ น การณเหตุ อุทฺธโต อวูปสนฺตจิตฺโต ปริเวณโต ปริเวณํ คจฺฉติ.\csromanspacer{} วิหารโต ฯเปฯ อิติ ภวาภวกถํ กเถติ.\csromanspacer{} เอวมฺปิ ปาทโลโล โหติ.}

\prose{\textbf{น จ ปาทโลโล}ติ ปาทโลลิยํ ปชเหยฺย วิโนเทยฺย พฺยนฺติํ กเรยฺย อนภาวํ คเมยฺย, ปาทโลลิยา อารโต อสฺส วิรโต ปฏิวิรโต นิกฺขนฺโต นิสฺสโฏ วิปฺปมุตฺโต วิสญฺญุตฺโต วิมริยาทิกเตน เจตสา วิหเรยฺย.\csromanspacer{} ปฏิสลฺลานาราโม อสฺส ปฏิสลฺลานรโต อชฺฌตฺตํ เจโตสมถมนุยุตฺโต อนิรากตชฺฌาโน วิปสฺสนาย สมนฺนาคโต พฺรูเหตา สุญฺญาคารํ ฌายี ฌานรโต เอกตฺตมนุยุตฺโต สทตฺถครุโกติ โอกฺขิตฺตจกฺขุ น จ ปาทโลโล.}

\prose{\textbf{ฌานานุยุตฺโต พหุชา ครสฺสา}ติ ฌานานุยุตฺโต ทฺวีหิ การเณหิ ฌานานุยุตฺโต อนุปฺปนฺนสฺส วา ปฐมสฺส ฌานสฺส อุปฺปาทาย ยุตฺโต ปยุตฺโต อายุตฺโต สมายุตฺโต.\csromanspacer{} อนุปฺปนฺนสฺส วา ทุติยสฺส ฌานสฺส.\csromanspacer{} ตติยสฺส ฌานสฺส.\csromanspacer{} จตุตฺถสฺส ฌานสฺส อุปฺปาทาย ยุตฺโต ปยุตฺโต อายุตฺโต สมายุตฺโตติ เอวมฺปิ ฌานานุยุตฺโต.\csromanspacer{} อถ วา อุปฺปนฺนํ วา ปฐมํ ฌานํ อาเสวติ ภาเวติ พหุลีกโรติ\footnote{พหุลิํ กโรติ (ก)}.\csromanspacer{} อุปฺปนฺนํ วา ทุติยํ}

\prose{\csromanfolio{402}ฌานํ.\csromanspacer{} ตติยํ ฌานํ.\csromanspacer{} จตุตฺถํ ฌานํ อาเสวติ ภาเวติ พหุลีกโรตีติ เอวมฺปิ ฌานานุยุตฺโต.}

\prose{\textbf{พหุชาครสฺสา}ติ อิธ ภิกฺขุ ทิวสํ จงฺกเมน นิสชฺชาย อาวรณีเยหิ ธมฺเมหิ จิตฺตํ ปริโสเธติ.\csromanspacer{} รตฺติยา ปฐมํ ยามํ จงฺกเมน นิสชฺชาย อาวรณีเยหิ ธมฺเมหิ จิตฺตํ ปริโสเธติ.\csromanspacer{} รตฺติยา มชฺฌิมํ ยามํ ทกฺขิเณน ปสฺเสน สีหเสยฺยํ กปฺเปติ ปาเท ปาทํ อจฺจาธาย สโต สมฺปชาโน อุฏฺฐานสญฺญํ มนสิ กตฺวา.\csromanspacer{} รตฺติยา ปจฺฉิมํ ยามํ ปจฺจุฏฺฐาย จงฺกเมน นิสชฺชาย อาวรณีเยหิ ธมฺเมหิ จิตฺตํ ปริโสเธตีติ ฌานานุยุตฺโต พหุชาครสฺส.}

\prose{\textbf{อุเปกฺขมารพฺภ สมาหิตตฺโต}ติ \textbf{อุเปกฺขา}ติ ยา จตุตฺเถ ฌาเน \textbf{อุเปกฺขา} อุเปกฺขนา อชฺฌุเปกฺขนา จิตฺตสมตา จิตฺตปฺปสฺสทฺธตา มชฺฌตฺตตา จิตฺตสฺส.\csromanspacer{} \textbf{สมาหิตตฺโต}ติ ยา จิตฺตสฺส ฐิติ สณฺฐิติ อวฏฺฐิติ อวิสาหาโร อวิกฺเขโป อวิสาหฏมานสตา สมโถ สมาธินฺทฺริยํ สมาธิพลํ สมฺมาสมาธิ.\csromanspacer{}\csromanspacer{}\textbf{อุเปกฺขมารพฺภ สมาหิตตฺโต}ติ จตุตฺเถ ฌาเน อุเปกฺขํ อารพฺภ เอกคฺคจิตฺโต อวิกฺขิตฺตจิตฺโต อปิสาหฏมานโสติ อุเปกฺขมารพฺภ สมาหิตตฺโต.}

\prose{\textbf{ตกฺกาสยํ กุกฺกุจฺจญฺจุปจฺฉินฺเท}ติ \textbf{ตกฺกา}ติ นว วิ\textbf{ตกฺกา}, กามวิตกฺโก พฺยาปาทวิตกฺโก วิหิํสาวิตกฺโก ญาติวิตกฺโก ชนปทวิตกฺโก อมรํวิตกฺโก ปรานุทยตาปฏิสญฺญุตฺโต วิตกฺโก ลาภสกฺการสิโลกปฏิสญฺญุตฺโต วิตกฺโก อนวญฺญตฺติปฏิสญฺญุตฺโต วิตกฺโก, อิเม วุจฺจนฺติ นว วิ\textbf{ตกฺกา}.\csromanspacer{} กามวิ\textbf{ตกฺกา}นํ กามสญฺญาสโย.\csromanspacer{} พฺยาปาทวิ\textbf{ตกฺกา}นํ พฺยาปาทสญฺญาสโย.\csromanspacer{} วิหิํสาวิ\textbf{ตกฺกา}นํ วิหิํสาสญฺญาสโย.\csromanspacer{} อถ วา \textbf{ตกฺกา}นํ วิ\textbf{ตกฺกา}นํ สงฺกปฺปานํ อวิชฺชาสโย.\csromanspacer{} อโยนิโส มนสิกาโร อาสโย.\csromanspacer{} อสฺมิมาโน อาสโย.\csromanspacer{} อโนตฺตปฺปํ อาสโย.\csromanspacer{} อุทฺธจฺจํ อาสโย.}

\prose{\textbf{กุกฺกุจฺจนฺ}ติ หตฺถกุกฺกุจฺจมฺปิ กุกฺกุจฺจํ, ปาทกุกฺกุจฺจมฺปิ กุกฺกุจฺจํ, หตฺถปาทกุกฺกุจฺจมฺปิ กุกฺกุจฺจํ, อกปฺปิเย กปฺปิยสญฺญิตา, กปฺปิเย อกปฺปิยสญฺญิตา, วิกาเล กาลสญฺญิตา, กาเล วิกาลสญฺญิตา, อวชฺเช วชฺชสญฺญิตา, วชฺเช อวชฺชสญฺญิตา.\csromanspacer{} ยํ เอวรูปํ กุกฺกุจฺจํ กุกฺกุจฺจายนา กุกฺกุจฺจายิตตฺตํ เจตโส วิปฺปฏิสาโร มโนวิเลโข.\csromanspacer{} อิทํ วุจฺจติ กุกฺกุจฺจํ.}

\csromanheader{2}{16. สาริปุตฺตสุตฺตนิทฺเทส}
\prose{\csromanfolio{403}อปิ จ ทฺวีหิ การเณหิ อุปฺปชฺชติ กุกฺกุจฺจํ เจตโส วิปฺปฏิสาโร มโนวิเลโข กตตฺตา จ อกตตฺตา จ.\csromanspacer{} กถํ กตตฺตา จ อกตตฺตา จ อุปฺปชฺชติ กุกฺกุจฺจํ เจตโส วิปฺปฏิสาโร มโนวิเลโข\csromandash{}“กตํ เม กายทุจฺจริตํ, อกตํ เม กายสุจริตนฺ”ติ อุปฺปชฺชติ กุกฺกุจฺจํ เจตโส วิปฺปฏิสาโร มโนวิเลโข. “กตํ เม วจีทุจฺจริตํ.\csromanspacer{} กตํ เม มโนทุจฺจริตํ.\csromanspacer{} กโต เม ปาณาติปาโต, อกตา เม ปาณาติปาตา เวรมณี”ติ อุปฺปชฺชติ กุกฺกุจฺจํ เจตโส วิปฺปฏิสาโร มโนวิเลโข. “กตํ เม อทินฺนาทานํ.\csromanspacer{} กโต เม กาเมสุมิจฺฉาจาโร.\csromanspacer{} กโต เม มุสาวาโท.\csromanspacer{} กตา เม ปิสุณา วาจา.\csromanspacer{} กตา เม ผรุสา วาจา.\csromanspacer{} กโต เม สมฺผปฺปลาโป.\csromanspacer{} กตา เม อภิชฺฌา.\csromanspacer{} กโต เม พฺยาปาโท.\csromanspacer{} กตา เม มิจฺฉาทิฏฺฐิ, อกตา เม สมฺมาทิฏฺฐี”ติ อุปฺปชฺชติ กุกฺกุจฺจํ เจตโส วิปฺปฏิสาโร มโนวิเลโข.\csromanspacer{} เอวํ กตตฺตา จ อกตตฺตา จ อุปฺปชฺชติ กุกฺกุจฺจํ เจตโส วิปฺปฏิสาโร มโนวิเลโข.}

\prose{อถ วา “สีเลสุมฺหิ น ปริปูรการี”ติ อุปฺปชฺชติ กุกฺกุจฺจํ เจตโส วิปฺปฏิสาโร มโนวิเลโข, “อินฺทฺริเยสุมฺหิ อคุตฺตทฺวาโร”ติ. “โภชเน อมตฺตญฺญูมฺหี”ติ. “ชาคริยํ อนนุยุตฺโตมฺหี”ติ. “น สติสมฺปชญฺเญน สมนฺนาคโตมฺหี”ติ. “อภาวิตา เม จตฺตาโร สติปฏฺฐานา”ติ. “อภาวิตา เม จตฺตาโร สมฺมปฺปธานา”ติ. “อภาวิตา เม จตฺตาโร อิทฺธิปาทา”ติ. “อภาวิตานิ เม ปญฺจินฺทฺริยานี”ติ. “อภาวิตานิ เม ปญฺจ พลานี”ติ. “อภาวิตา เม สตฺต โพชฺฌงฺคา”ติ. “อภาวิโต เม อริโย อฏฺฐงฺคิโก มคฺโค”ติ. “ทุกฺขํ เม อปริญฺญาตนฺ”ติ. “ทุกฺขสมุทโย เม อปฺปหีโน”ติ. “มคฺโค เม อภาวิโต”ติ. “นิโรโธ เม อสจฺฉิกโต”ติ อุปฺปชฺชติ กุกฺกุจฺจํ เจตโส วิปฺปฏิสาโร มโน วิเลโข.\csromanspacer{}\csromanspacer{}\textbf{ตกฺกาสยํ กุกฺกุจฺจญฺจุปจฺฉินฺเท}ติ ตกฺกญฺจ ตกฺกาสยญฺจ กุกฺกุจฺจญฺจ อุปจฺฉินฺเทยฺย ฉินฺเทยฺย อุจฺฉินฺเทยฺย สมุจฺฉินฺเทยฺย ปชเหยฺย วิโนเทยฺย พฺยนฺติํ กเรยฺย อนภาวํ คเมยฺยาติ ตกฺกาสยํ กุกฺกุจฺจญฺจุปจฺฉินฺเท.\csromanspacer{} เตนาห ภควา\csromandash{}}

\csromangathameasure{โอกฺขิตฺตจกฺขุ น จ ปาทโลโล, \\ ฌานานุยุตฺโต พหุชาครสฺส. \\ อุเปกฺขมารพฺภ สมาหิตตฺโต,}
\csromangathagroup{\csromangathastanza{โอกฺขิตฺตจกฺขุ น จ ปาทโลโล, \\ ฌานานุยุตฺโต พหุชาครสฺส. \\ อุเปกฺขมารพฺภ สมาหิตตฺโต,}}

\prose{\textbf{ตกฺกาสยํ กุกฺกุจฺจญฺจุปจฺฉินฺเท}ติ.}

\proseitem{208}{\csromanfolio{404}\csromansymbolfootnote{*}{ขุ 1. 428 ปิฏฺเฐปิ.}\textbf{จุทิโต วจีภิ สติมาภินนฺเท},}

\prose{\textbf{สพฺรหฺมจารีสุ ขิลํ ปภินฺเท.}}

\csromangathameasure{\textbf{วาจํ ปมุญฺเจ กุสลํ นาติเวลํ,} \\ \textbf{ชนวาทธมฺมาย น เจตเยยฺย.}}
\csromangathagroup{\csromangathastanza{\textbf{วาจํ ปมุญฺเจ กุสลํ นาติเวลํ,} \\ \textbf{ชนวาทธมฺมาย น เจตเยยฺย.}}}

\prose{\textbf{จุทิโต วจีภิ สติมาภินนฺเท}ติ \textbf{จุทิโต}ติ อุปชฺฌายา วา อาจริยา วา สมานุปชฺฌายกา วา สมานาจริยกา วา มิตฺตา วา สนฺทิฏฺฐา วา สมฺภตฺตา วา สหายา วา โจเทนฺติ “อิทํ เต อาวุโส อยุตฺตํ, อิทํ เต อปฺปตฺตํ, อิทํ เต อสารุปฺปํ, อิทํ เต อสีลฏฺฐนฺ”ติ.\csromanspacer{} สติํ อุปฏฺฐเปตฺวา ตํ โจทนํ นนฺเทยฺย อภินนฺเทยฺย โมเทยฺย อนุโมเทยฺย อิจฺเฉยฺย สาทิเยยฺย ปตฺถเยยฺย ปิหเยยฺย อภิชปฺเปยฺย.\csromanspacer{} ยถา อิตฺถี วา ปุริโส วา ทหโร ยุวา มณฺฑนกชาติโก สีสํนฺหาโต อุปฺปลมาลํ วา วสฺสิกมาลํ วา อธิมุตฺตกมาลํ วา ลภิตฺวา อุโภหิ หตฺเถหิ ปฏิคฺคเหตฺวา อุตฺตมงฺเค สิรสฺมิํ ปติฏฺฐาเปตฺวา นนฺเทยฺย อภินนฺเทยฺย โมเทยฺย อนุโมเทยฺย อิจฺเฉยฺย สาทิเยยฺย ปตฺถเยยฺย ปิหเยยฺย อภิชปฺเปยฺย.\csromanspacer{} เอวเมวํ สติํ อุปฏฺฐเปตฺวา ตํ โจทนํ นนฺเทยฺย อภินนฺเทยฺย โมเทยฺย อนุโมเทยฺย อิจฺเฉยฺย สาทิเยยฺย ปตฺถเยยฺย ปิหเยยฺย อภิชปฺเปยฺย.}

\csromangathameasure{“นิธีนํว ปวตฺตารํ, ยํ ปสฺเส วชฺชทสฺสินํ. \\ นิคฺคยฺหวาทิํ เมธาวิํ, ตาทิสํ ปณฺฑิตํ ภเช. \\ ตาทิสํ ภชมานสฺส, เสยฺโย โหติ น ปาปิโย. \\ โอวเทยฺยานุสาเสยฺย, อสพฺภา จ นิวารเย. \\ สตํ หิ โส ปิโย โหติ, อสตํ โหติ อปฺปิโย”ติ.}
\csromangathagroup{\csromangathastanza{“นิธีนํว\footnote{นิธินํว (ก) ขุ 1. 24 ปิฏฺเฐ.} ปวตฺตารํ, ยํ ปสฺเส วชฺชทสฺสินํ. \\ นิคฺคยฺหวาทิํ เมธาวิํ, ตาทิสํ ปณฺฑิตํ ภเช.}\csromangathastanza{ตาทิสํ ภชมานสฺส, เสยฺโย โหติ น ปาปิโย. \\ โอวเทยฺยานุสาเสยฺย, อสพฺภา จ นิวารเย. \\ สตํ หิ โส ปิโย โหติ, อสตํ โหติ อปฺปิโย”ติ.}}

\prose{\textbf{จุทิโต วจีภิ สติมาภินนฺเท, สพฺรหฺมจารีสุ ขิลํ ปภินฺเท}ติ \textbf{สพฺรหฺมจารี}ติ เอกกมฺมํ เอกุทฺเทโส สมสิกฺขตา.\csromanspacer{} \textbf{สพฺรหฺมจารีสุ ขิลํ ปภินฺเท}ติ \textbf{สพฺรหฺมจารี}สุ อาหตจิตฺตตํ ขิลชาตตํ ปภินฺเทยฺย, ปญฺจปิ เจโตขิเล ภินฺเทยฺย, ตโยปิ เจโตขิเล ภินฺเทยฺย, ราคขิลํ โทสขิลํ โมหขิลํ ภินฺเทยฺย ปภินฺเทยฺย สมฺภินฺเทยฺยาติ \textbf{สพฺรหฺมจารี}สุ ขิลํ ปภินฺเท.}

\csromanheader{2}{16. สาริปุตฺตสุตฺตนิทฺเทส}
\prose{\csromanfolio{405}\textbf{วาจํ ปมุญฺเจ กุสลํ นาติเวลนฺ}ติ ญาณสมุฏฺฐิตํ วาจํ มุญฺเจยฺย อตฺถูปสํหิตํ ธมฺมูปสํหิตํ.\csromanspacer{} กาเลน สาปเทสํ ปริยนฺตวติํ วาจํ มุญฺเจยฺย ปมุญฺเจยฺยาติ วาจํ ปมุญฺเจ กุสลํ.\csromanspacer{} \textbf{นาติเวลนฺ}ติ \textbf{เวลา}ติ เทฺว \textbf{เวลา} กาล\textbf{เวลา} จ สีล\textbf{เวลา} จ.\csromanspacer{} กตมา กาล\textbf{เวลา}, กาลาติกฺกนฺตํ วาจํ น ภาเสยฺย, \textbf{เวลา}ติกฺกนฺตํ วาจํ น ภาเสยฺย, กาล\textbf{เวลา}ติกฺกนฺตํ วาจํ น ภาเสยฺย.\csromanspacer{} กาลํ อสมฺปตฺตํ วาจํ น ภาเสยฺย, เวลํ อสมฺปตฺตํ วาจํ น ภาเสยฺย, กาลเวลํ อสมฺปตฺตํ วาจํ น ภาเสยฺย.}

\csromangathameasure{“โย เว กาเล อสมฺปตฺเต, อติเวลํ จ ภาสติ. \\ เอวํ โส นิหโต เสติ, โกกิลาเยว อตฺรโช”ติ.}
\csromangathagroup{\csromangathastanza{“โย เว\footnote{จ (สฺยา)} กาเล อสมฺปตฺเต, อติเวลํ จ ภาสติ. \\ เอวํ โส นิหโต เสติ, โกกิลาเยว\footnote{จ (สฺยา)} อตฺรโช”ติ.}}

\prose{อยํ กาล\textbf{เวลา}.\csromanspacer{}\csromanspacer{}กตมา สีล\textbf{เวลา}, รตฺโต วาจํ น ภาเสยฺย, ทุฏฺโฐ วาจํ น ภาเสยฺย, มูโฬฺห วาจํ น ภาเสยฺย, มุสาวาทํ น ภาเสยฺย, ปิสุณวาจํ น ภาเสยฺย, ผรุสวาจํ น ภาเสยฺย, สมฺผปฺปลาปํ น ภาเสยฺย น กเถยฺย น ภเณยฺย น ทีปเยยฺย น โวหเรยฺย.\csromanspacer{} อยํ สีล\textbf{เวลา}ติ วาจํ ปมุญฺเจ กุสลํ นาติเวลํ.}

\prose{\textbf{ชนวาทธมฺมาย น เจตเยยฺยา}ติ \textbf{ชนา}ติ ขตฺติยา จ พฺราหฺมณา จ เวสฺสา จ สุทฺทา จ คหฏฺฐา จ ปพฺพชิตา จ เทวา จ มนุสฺสา จ.\csromanspacer{} ชนสฺส วาทาย อุปวาทาย นินฺทาย ครหาย อกิตฺติยา อวณฺณหาริกาย สีลวิปตฺติยา วา อาจารวิปตฺติยา วา ทิฏฺฐิวิปตฺติยา วา อาชีววิปตฺติยา วา น เจตเยยฺย เจตนํ น อุปฺปาเทยฺย จิตฺตํ น อุปฺปาเทยฺย สงฺกปฺปํ น อุปฺปาเทยฺย มนสิการํ น อุปฺปาเทยฺยาติ ชนวาทธมฺมาย น เจตเยยฺย.\csromanspacer{} เตนาห ภควา\csromandash{}}

\csromangathameasure{จุทิโต วจีภิ สติมาภินนฺเท, \\ สพฺรหฺมจารีสุ ขิลํ ปภินฺเท.}
\csromangathagroup{\csromangathastanza{จุทิโต วจีภิ สติมาภินนฺเท, \\ สพฺรหฺมจารีสุ ขิลํ ปภินฺเท.}}

\prose{วาจํ ปมุญฺเจ กุสลํ นาติเวลํ,}

\prose{\textbf{ชนวาทธมฺมาย น เจตเยยฺยา}ติ.}

\proseitem{209}{\csromansymbolfootnote{*}{ขุ 1. 428 ปิฏฺเฐปิ.}อถาปรํ ปญฺจ รชานิ โลเก,}

\prose{\textbf{เยสํ สตีมา วินยาย สิกฺเข.}}

\csromangathameasure{\textbf{รูเปสุ สทฺเทสุ อโถ รเสสุ,} \\ \textbf{คนฺเธสุ ผสฺเสสุ สเหถ ราคํ.}}
\csromangathagroup{\csromangathastanza{\textbf{รูเปสุ สทฺเทสุ อโถ รเสสุ,} \\ \textbf{คนฺเธสุ ผสฺเสสุ สเหถ ราคํ.}}}

\prose{\csromanfolio{406}\textbf{อถาปรํ ปญฺจ รชานิ โลเก}ติ \textbf{อถา}ติ ปทสนฺธิ ปทสํสคฺโค ปทปาริปูรี อกฺขรสมวาโย พฺยญฺชนสิลิฏฺฐตา ปทานุปุพฺพตาเปตํ \textbf{อถา}ติ.\csromanspacer{} \textbf{ปญฺจ รชานี}ติ รูปรโช สทฺทรโช คนฺธรโช รสรโช โผฏฺฐพฺพรโช.}

\csromangathameasure{*ราโค รโช น จ ปน เรณุ วุจฺจติ, \\ ราคสฺเสตํ อธิวจนํ รโชติ. \\ เอตํ รชํ วิปฺปชหิตฺวา ปณฺฑิตา, \\ วิหรนฺติ เต วิคตรชสฺส สาสเน. \\ โทโส รโช น จ ปน เรณุ วุจฺจติ ฯเปฯ. วิหรนฺติ เต วิคตรชสฺส สาสเน. \\ โมโห รโช น จ ปน เรณุ วุจฺจติ ฯเปฯ. วิหรนฺติ เต วิคตรชสฺส สาสเน.}
\csromangathagroup{\csromangathastanza{\csromansymbolfootnote{*}{ขุ 8. 154 ปิฏฺเฐปิ.}ราโค รโช น จ ปน เรณุ วุจฺจติ, \\ ราคสฺเสตํ อธิวจนํ รโชติ. \\ เอตํ รชํ วิปฺปชหิตฺวา\footnote{ปฏิวิโนทิตฺวา (ก)} ปณฺฑิตา, \\ วิหรนฺติ เต วิคตรชสฺส สาสเน. \\ โทโส รโช น จ ปน เรณุ วุจฺจติ ฯเปฯ. วิหรนฺติ เต วิคตรชสฺส สาสเน. \\ โมโห รโช น จ ปน เรณุ วุจฺจติ ฯเปฯ. วิหรนฺติ เต วิคตรชสฺส สาสเน.}}

\prose{\textbf{โลเก}ติ อปายโลเก มนุสฺสโลเก เทวโลเก ขนฺธโลเก ธาตุโลเก อายตนโลเกติ \textbf{อถา}ปรํ ปญฺจ รชานิ โลเก.}

\prose{\textbf{เยสํ สตีมา วินยาย สิกฺเข}ติ \textbf{เยส}นฺติ รูปราคสฺส สทฺทราคสฺส คนฺธราคสฺส รสราคสฺส โผฏฺฐพฺพราคสฺส.\csromanspacer{} \textbf{สตีมา}ติ ยา สติ อนุสฺสติ ปฏิสฺสติ สติ สรณตา ธารณตา อปิลาปนตา อสมฺมุสฺสนตา สตินฺทฺริยํ สติพลํ สมฺมาสติ สติสมฺโพชฺฌงฺโค เอกายนมคฺโค, อยํ วุจฺจติ สติ, อิมาย สติยา อุเปโต สมุเปโต อุปคโต สมุปคโต อุปปนฺโน สมุปปนฺโน สมนฺนาคโต.\csromanspacer{} โส วุจฺจติ สติมา.\csromanspacer{} \textbf{สิกฺเข}ติ ติสฺโส สิกฺขา อธิสีลสิกฺขา อธิจิตฺตสิกฺขา อธิปญฺญาสิกฺขา.\csromanspacer{} กตมา อธิสีลสิกฺขา ฯเปฯ อยํ อธิปญฺญาสิกฺขา.\csromanspacer{}\csromanspacer{}\textbf{เยสํ สตีมา วินยาย สิกฺเข}ติ สติมา ปุคฺคโล \textbf{เยส}ํ รูปราคสฺส สทฺทราคสฺส คนฺธราคสฺส รสราคสฺส โผฏฺฐพฺพราคสฺส วินยาย ปฏิวินยาย ปหานาย วูปสมาย ปฏินิสฺสคฺคาย ปฏิปสฺสทฺธิยา อธิสีลมฺปิ สิกฺเขยฺย อธิจิตฺตมฺปิ สิกฺเขยฺย อธิปญฺญมฺปิ สิกฺเขยฺย, อิมา ติสฺโส สิกฺขาโย อาวชฺชนฺโต สิกฺเขยฺย.\csromanspacer{} ชานนฺโต สิกฺเขยฺย ฯเปฯ สจฺฉิกาตพฺพํ สจฺฉิกโรนฺโต สิกฺเขยฺย อาจเรยฺย สมาจเรยฺย สมาทาย วตฺเตยฺยาติ \textbf{เยส}ํ สตีมา วินยาย สิกฺเข.}

\csromanheader{2}{16. สาริปุตฺตสุตฺตนิทฺเทส}
\prose{\csromanfolio{407}\textbf{รูเปสุ สทฺเทสุ อโถ รเสสุ, คนฺเธสุ ผสฺเสสุ สเหถ ราคนฺ}ติ รูเปสุ สทฺเทสุ คนฺเธสุ รเสสุ โผฏฺฐพฺเพสุ ราคํ สเหยฺย ปริสเหยฺย อภิภเวยฺย อชฺโฌตฺถเรยฺย ปริยาทิเยยฺย มทฺเทยฺยาติ รูเปสุ สทฺเทสุ อโถ รเสสุ, คนฺเธสุ ผสฺเสสุ สเหถ ราคํ.\csromanspacer{} เตนาห ภควา\csromandash{}}

\csromangathameasure{อถาปรํ ปญฺจ รชานิ โลเก, \\ เยสํ สตีมา วินยาย สิกฺเข.}
\csromangathagroup{\csromangathastanza{อถาปรํ ปญฺจ รชานิ โลเก, \\ เยสํ สตีมา วินยาย สิกฺเข.}}

\prose{รูเปสุ สทฺเทสุ อโถ รเสสุ,}

\prose{คนฺเธสุ ผสฺเสสุ สเหถ ราคนฺติ.}

\proseitem{210}{\csromansymbolfootnote{*}{ขุ 1. 428 ปิฏฺเฐปิ.}เอเตสุ ธมฺเมสุ วิเนยฺย ฉนฺทํ,}

\prose{\textbf{ภิกฺขุ สติมา สุวิมุตฺตจิตฺโต.}}

\csromangathameasure{\textbf{กาเลน โส สมฺมา ธมฺมํ ปริวีมํสมาโน,}}
\csromangathagroup{\csromangathastanza{\textbf{กาเลน โส สมฺมา ธมฺมํ ปริวีมํสมาโน,}}}

\prose{\textbf{เอโกทิภูโต วิหเน ตมํ โส, (อิติ ภควา.)}}

\prose{\textbf{เอเตสุ ธมฺเมสุ วิเนยฺย ฉนฺทนฺ}ติ \textbf{เอเตสู}ติ รูเปสุ สทฺเทสุ คนฺเธสุ รเสสุ โผฏฺฐพฺเพสุ.\csromanspacer{} \textbf{ฉนฺโท}ติ โย กาเมสุ กามจฺฉนฺโท กามราโค กามนนฺที กามตณฺหา กามสฺเนโห กามปริฬาโห กามมุจฺฉา กามชฺโฌสานํ กาโมโฆ กามโยโค กามุปาทานํ ฯเปฯ กามจฺฉนฺทนีวรณํ.\csromanspacer{} \textbf{เอเตสุ ธมฺเมสุ วิเนยฺย ฉนฺทนฺ}ติ เอเตสุ ธมฺเมสุ ฉนฺทํ วิเนยฺย ปฏิวิเนยฺย ปชเหยฺย วิโนเทยฺย พฺยนฺติํ กเรยฺย อนภาวํ คเมยฺยาติ เอเตสุ ธมฺเมสุ วิเนยฺย ฉนฺทํ.}

\prose{\textbf{ภิกฺขุ สติมา สุวิมุตฺตจิตฺโต}ติ \textbf{ภิกฺขู}ติ ปุถุชฺชนกลฺยาณโก วา ภิกฺขุ, เสโข วา ภิกฺขุ.\csromanspacer{} \textbf{สติมา}ติ ยา สติ อนุสฺสติ ฯเปฯ สมฺมาสติ สติสมฺโพชฺฌงฺโค เอกายนมคฺโค.\csromanspacer{} อยํ วุจฺจติ สติ, อิมาย สติยา อุเปโต สมุเปโต ฯเปฯ.\csromanspacer{} โส วุจฺจติ สติมา.}

\prose{\textbf{ภิกฺขุ สติมา สุวิมุตฺตจิตฺโต}ติ ปฐมํ ฌานํ สมาปนฺนสฺส นีวรเณหิ จิตฺตํ มุตฺตํ วิมุตฺตํ สุวิมุตฺตํ.\csromanspacer{} ทุติยํ ฌานํ สมาปนฺนสฺส วิตกฺกวิจาเรหิ จิตฺตํ มุตฺตํ วิมุตฺตํ สุวิมุตฺตํ.\csromanspacer{} ตติยํ ฌานํ สมาปนฺนสฺส ปีติยา จ จิตฺตํ มุตฺตํ วิมุตฺตํ สุวิมุตฺตํ.\csromanspacer{} จตุตฺถํ ฌานํ สมาปนฺนสฺส สุขทุกฺเขหิ จิตฺตํ มุตฺตํ วิมุตฺตํ สุวิมุตฺตํ.\csromanspacer{} อากาสานญฺจายตนํ สมาปนฺนสฺส รูปสญฺญาย ปฏิฆสญฺญาย นานตฺตสญฺญาย จิตฺตํ มุตฺตํ วิมุตฺตํ สุวิมุตฺตํ.\csromanspacer{} วิญฺญาณญฺจายตนํ สมาปนฺนสฺส \csromanfolio{408}อากาสานญฺจายตนสญฺญาย จิตฺตํ.\csromanspacer{} อากิญฺจญฺญายตนํ สมาปนฺนสฺส วิญฺญาณญฺจายตนสญฺญาย จิตฺตํ.\csromanspacer{} เนวสญฺญานาสญฺญายตนํ สมาปนฺนสฺส อากิญฺจญฺญายตนสญฺญาย จิตฺตํ มุตฺตํ วิมุตฺตํ สุวิมุตฺตํ.\csromanspacer{} โสตาปนฺนสฺส สกฺกายทิฏฺฐิยา วิจิกิจฺฉาย สีลพฺพตปรามาสา ทิฏฺฐานุสยา วิจิกิจฺฉานุสยา ตเทกฏฺเฐหิ จ กิเลเสหิ จิตฺตํ มุตฺตํ วิมุตฺตํ สุวิมุตฺตํ.\csromanspacer{} สกทาคามิสฺส โอฬาริกา กามราคานุสยา ปฏิฆานุสยา ตเทกฏฺเฐหิ จ กิเลเสหิ จิตฺตํ มุตฺตํ วิมุตฺตํ สุวิมุตฺตํ.\csromanspacer{} อนาคามิสฺส อณุสหคตา กามราคสญฺโญชนา ปฏิฆสญฺโญชนา อณุสหคตา กามราคานุสยา ปฏิฆานุสยา ตเทกฏฺเฐหิ จ กิเลเสหิ จิตฺตํ มุตฺตํ วิมุตฺตํ สุวิมุตฺตํ.\csromanspacer{} อรหโต รูปราคา อรูปราคา มานา อุทฺธจฺจา อวิชฺชาย มานานุสยา ภวราคานุสยา อวิชฺชานุสยา ตเทกฏฺเฐหิ จ กิเลเสหิ พหิทฺธา จ สพฺพนิมิตฺเตหิ จิตฺตํ มุตฺตํ วิมุตฺตํ สุวิมุตฺตนฺติ ภิกฺขุ สติมา สุวิมุตฺตจิตฺโต.}

\prose{\textbf{กาเลน โส สมฺมา ธมฺมํ ปริวีมํสมาโน}ติ \textbf{กาเลนา}ติ อุทฺธเต จิตฺเต สมถสฺส\footnote{สมาธิสฺส (สี)} กาโล.\csromanspacer{} สมาหิเต จิตฺเต วิปสฺสนาย กาโล.}

\csromangathameasure{“กาเล ปคฺคณฺหติ จิตฺตํ, นิคฺคณฺหติ ปุนาปเร. \\ สมฺปหํสติ กาเลน, กาเล จิตฺตํ สมาทเห. \\ อชฺฌุเปกฺขติ กาเลน, โส โยคี กาลโกวิโท. \\ กิมฺหิ กาลมฺหิ ปคฺคาโห, กิมฺหิ กาเล วินิคฺคโห. \\ กิมฺหิ ปหํสนากาโล, สมถกาโล จ กีทิโส. \\ อุเปกฺขากาลํ จิตฺตสฺส, กถํ ทสฺเสติ โยคิโน. \\ ลีเน จิตฺตมฺหิ ปคฺคาโห, อุทฺธตสฺมิํ วินิคฺคโห. \\ นิรสฺสาทคตํ จิตฺตํ, สมฺปหํเสยฺย ตาวเท. \\ สมฺปหฏฺฐํ ยทา จิตฺตํ, อลีนํ ภวตินุทฺธตํ. \\ สมถสฺส จ โส กาโล, อชฺฌตฺตํ รมเย มโน. \\ เอเตน เมวุปาเยน, ยทา โหติ สมาหิตํ. \\ สมาหิตจิตฺตมญฺญาย, อชฺฌุเปกฺเขยฺย ตาวเท.}
\csromangathagroup{\csromangathastanza{“กาเล ปคฺคณฺหติ จิตฺตํ, นิคฺคณฺหติ ปุนาปเร\footnote{อถาปเร (สฺยา)}. \\ สมฺปหํสติ กาเลน, กาเล จิตฺตํ สมาทเห.}\csromangathastanza{อชฺฌุเปกฺขติ กาเลน, โส โยคี กาลโกวิโท. \\ กิมฺหิ กาลมฺหิ ปคฺคาโห, กิมฺหิ กาเล วินิคฺคโห.}\csromangathastanza{กิมฺหิ ปหํสนากาโล, สมถกาโล จ กีทิโส. \\ อุเปกฺขากาลํ จิตฺตสฺส, กถํ ทสฺเสติ โยคิโน.}\csromangathastanza{ลีเน จิตฺตมฺหิ ปคฺคาโห, อุทฺธตสฺมิํ วินิคฺคโห. \\ นิรสฺสาทคตํ จิตฺตํ, สมฺปหํเสยฺย ตาวเท.}\csromangathastanza{สมฺปหฏฺฐํ ยทา จิตฺตํ, อลีนํ ภวตินุทฺธตํ. \\ สมถสฺส จ โส\footnote{สมถนิมิตฺตสฺส โส (สี, ก)} กาโล, อชฺฌตฺตํ รมเย มโน.}\csromangathastanza{เอเตน เมวุปาเยน, ยทา โหติ สมาหิตํ. \\ สมาหิตจิตฺตมญฺญาย, อชฺฌุเปกฺเขยฺย ตาวเท.}}

\csromanheader{2}{16. สาริปุตฺตสุตฺตนิทฺเทส}
\csromangathameasure{เอวํ กาลวิทู ธีโร, กาลญฺญู กาลโกวิโท. \\ กาเลน กาลํ จิตฺตสฺส, นิมิตฺตมุปลกฺขเย”ติ.}
\csromangathagroup{\csromangathastanza{\csromanfolio{409}เอวํ กาลวิทู ธีโร, กาลญฺญู กาลโกวิโท. \\ กาเลน กาลํ จิตฺตสฺส, นิมิตฺตมุปลกฺขเย”ติ.}}

\prose{\textbf{กาเลน โส สมฺมา ธมฺมํ ปริวีมํสมาโน}ติ “สพฺเพ สงฺขารา อนิจฺจา”ติ สมฺมา ธมฺมํ ปริวีมํสมาโน. “สพฺเพ สงฺขารา ทุกฺขา”ติ สมฺมา ธมฺมํ ปริวีมํสมาโน. “สพฺเพ ธมฺมา อนตฺตา”ติ สมฺมา ธมฺมํ ปริวีมํสมาโน ฯเปฯ “ยํ กิญฺจิ สมุทยธมฺมํ, สพฺพํ ตํ นิโรธธมฺมนฺ”ติ สมฺมา ธมฺมํ ปริวีมํสมาโน.}

\prose{\textbf{เอโกทิภูโต วิหเน ตมํ โส, อิติ ภควา}ติ \textbf{เอโกที}ติ เอกคฺคจิตฺโต อวิกฺขิตฺตจิตฺโต อวิสาหฏมานโส สมโถ สมาธินฺทฺริยํ สมาธิพลํ สมฺมาสมาธีติ เอโกทิภูโต.\csromanspacer{} \textbf{วิหเน ตมํ โส}ติ ราคตมํ โทสตมํ โมหตมํ ทิฏฺฐิตมํ มานตมํ กิเลสตมํ ทุจฺจริตตมํ อนฺธกรณํ อจกฺขุกรณํ อญฺญาณกรณํ ปญฺญานิโรธิกํ วิฆาตปกฺขิกํ อนิพฺพานสํวตฺตนิกํ หเนยฺย วิหเนยฺย ปชเหยฺย วิโนเทยฺย พฺยนฺติํ กเรยฺย อนภาวํ คเมยฺย.}

\prose{\textbf{ภควา}ติ คารวาธิวจนํ.\csromanspacer{} อปิ จ ภคฺคราโคติ \textbf{ภควา}.\csromanspacer{} ภคฺคโทโสติ \textbf{ภควา}, ภคฺคโมโหติ \textbf{ภควา}, ภคฺคมาโนติ \textbf{ภควา}, ภคฺคทิฏฺฐีติ \textbf{ภควา}, ภคฺคกณฺฑโกติ \textbf{ภควา}, ภคฺคกิเลโสติ \textbf{ภควา}, ภชิ วิภชิ ปวิภชิ ธมฺมรตนนฺติ \textbf{ภควา}, ภวานํ อนฺตกโรติ \textbf{ภควา}, ภาวิตกาโย, ภาวิตสีโล, ภาวิตจิตฺโต, ภาวิตปญฺโญติ \textbf{ภควา}, ภชิ วา \textbf{ภควา} อรญฺญวนปตฺถานิ ปนฺตานิ เสนาสนานิ อปฺปสทฺทานิ อปฺปนิคฺโฆสานิ วิชนวาตานิ มนุสฺสราหสฺเสยฺยกานิ ปฏิสลฺลานสารุปฺปานีติ \textbf{ภควา}, ภาคี วา \textbf{ภควา} จีวรปิณฺฑปาตเสนาสนคิลานปจฺจยเภสชฺช- ปริกฺขารานนฺติ \textbf{ภควา}, ภาคี วา \textbf{ภควา} อตฺถรสสฺส ธมฺมรสสฺส วิมุตฺติรสสฺส อธิสีลสฺส อธิจิตฺตสฺส อธิปญฺญายาติ \textbf{ภควา}, ภาคี วา \textbf{ภควา} จตุนฺนํ ฌานานํ จตุนฺนํ อปฺปมญฺญานํ จตุนฺนํ อรูปสมาปตฺตีนนฺติ \textbf{ภควา}, ภาคี วา \textbf{ภควา} อฏฺฐนฺนํ วิโมกฺขานํ อฏฺฐนฺนํ อภิภายตนานํ นวนฺนํ อนุปุพฺพวิหารสมาปตฺตีนนฺติ \textbf{ภควา}, ภาคี วา \textbf{ภควา} ทสนฺนํ สญฺญาภาวนานํ ทสนฺนํ กสิณสมาปตฺตีนํ อานาปานสฺสติสมาธิสฺส อสุภสมาปตฺติยาติ \textbf{ภควา}, ภาคี วา \textbf{ภควา} จตุนฺนํ สติปฏฺฐานานํ จตุนฺนํ สมฺมปฺปธานานํ จตุนฺนํ อิทฺธิปาทานํ ปญฺจนฺนํ อินฺทฺริยานํ ปญฺจนฺนํ พลานํ สตฺตนฺนํ โพชฺฌงฺคานํ อริยสฺส อฏฺฐงฺคิกสฺส มคฺคสฺสาติ \textbf{ภควา}, ภาคี วา \textbf{ภควา} ทสนฺนํ ตถาคตพลานํ จตุนฺนํ เวสารชฺชานํ จตุนฺนํ ปฏิสมฺภิทานํ ฉนฺนํ อภิญฺญานํ ฉนฺนํ พุทฺธธมฺมานนฺติ}

\prose{\csromanfolio{410}ภควา, ภควาติ เนตํ นามํ มาตรา กตํ, น ปิตรา กตํ, น ภาตรา กตํ, น ภคินิยา กตํ, น มิตฺตามจฺเจหิ กตํ, น ญาติสาโลหิเตหิ กตํ, น สมณาพฺราหฺมเณหิ กตํ, น เทวตาหิ กตํ, วิโมกฺขนฺติกเมตํ พุทฺธานํ ภควนฺตานํ โพธิยา มูเล สห สพฺพญฺญุตญาณสฺส ปฏิลาภา สจฺฉิกา ปญฺญตฺติ ยทิทํ ภควาติ เอโกทิภูโต วิหเน ตมํ โส, อิติ ภควา.\csromanspacer{} เตนาห ภควา\csromandash{}}

\csromangathameasure{เอเตสุ ธมฺเมสุ วิเนยฺย ฉนฺทํ, \\ ภิกฺขุ สติมา สุวิมุตฺตจิตฺโต. \\ กาเลน โส สมฺมา ธมฺมํ ปริวีมํสมาโน,}
\csromangathagroup{\csromangathastanza{เอเตสุ ธมฺเมสุ วิเนยฺย ฉนฺทํ, \\ ภิกฺขุ สติมา สุวิมุตฺตจิตฺโต. \\ กาเลน โส สมฺมา ธมฺมํ ปริวีมํสมาโน,}}

\vspace{\tipitakaparskip}
\csromancenter{เอโกทิภูโต วิหเน ตมํ โส, (อิติ ภควาติ.)}
\csromancenter{สาริปุตฺตสุตฺตนิทฺเทโส โสฬสโม.}
\csromancenter{อฏฺฐกวคฺคมฺหิ โสฬส สุตฺตนิทฺเทสา สมตฺตา.}
\nitthitam{\textbf{มหานิทฺเทสปาฬิ นิฏฺฐิตา.}}
\orphannote{ม 2. 27 ปิฏฺฐาทีสุปิ.}

\orphannote{อภิ 1. 215 ปิฏฺเฐปิ.}

\orphannote{ขุ 8. 38 ปิฏฺเฐปิ.}

\orphannote{ขุ 8. 249 ปิฏฺเฐปิ.}

\orphannote{สํ 2. 349; ขุ 8. 251.}

\orphannote{สํ 1. 70 ปิฏฺเฐปิ.}

\orphannote{ม 2. 27 ปิฏฺฐาทีสุปิ.}

\orphannote{สํ 2. 383 ปิฏฺเฐ.}

\orphannote{ที 3. 224, 258 ปิฏฺเฐสุปิ.}

\orphannote{สํ 2. 311 ปิฏฺเฐ.}

\orphannote{สํ 2. 43, 45 ปิฏฺเฐสุ.}

\orphannote{สํ 1. 325; อภิ 4. 114 ปิฏฺเฐสุ.}

\orphannote{ขุ 8. 36 ปิฏฺเฐปิ.}

\orphannote{อภิ 1. 215 ปิฏฺเฐปิ.}

\orphannote{ม 1. 54 ปิฏฺเฐ.}

\orphannote{ม 3. 57; ขุ 8. 95 ปิฏฺเฐสุ.}

\orphannote{อภิ 1. 20, 168 ปิฏฺเฐสุ.}

\orphannote{อภิ 1. 222 ปิฏฺเฐ.}

\orphannote{สํ 2. 23 ปิฏฺเฐ โถกํ วิสทิสํ.}

\orphannote{สํ 2. 23 ปิฏฺเฐ.}

\orphannote{ขุ 8. 80; อุปริ 293 ปิฏฺเฐสุปิ.}

\orphannote{ม \textbf{2.} 146; สํ 1\textbf{.} \textbf{2}67; อํ 1\textbf{.} 5\textbf{2}; ขุ 8\textbf{.} \textbf{2}97 ปิฏฺเฐสุปิ\textbf{.}}

\orphannote{ขุ 8. 59 ปิฏฺเฐปิ.}

\orphannote{ขุ 8. 39 ปิฏฺเฐปิ.}

\orphannote{อภิ 2. 372 ปิฏฺเฐปิ.}

\orphannote{อภิ 2. 359 ปิฏฺฐาทีสุปิ.}

\orphannote{อภิ 1. 225 ปิฏฺเฐปิ.}

\orphannote{อภิ 1. 229, 281 ปิฏฺเฐสุปิ.}

\orphannote{อภิ 1. 93 ปิฏฺเฐ.}

\orphannote{ขุ 8. 83 ปิฏฺเฐปิ.}

\orphannote{อภิ \textbf{1.} 20, \textbf{1}68 ปิฏฺฐาทีสุปิ\textbf{.}}

\orphannote{ม 1. 190 ปิฏฺเฐ.}

\orphannote{ที 2. 244; ม 1. 82; อภิ 2. 104 ปิฏฺเฐสุปิ.}

\orphannote{ม 3. 289; สํ 2. 405 ปิฏฺเฐสุ.}

\orphannote{วิ 1. 24 ปิฏฺเฐ.}

\orphannote{เหฏฺฐา 20 ปิฏฺเฐปิ.}

\orphannote{วิ 1. 24 ปิฏฺเฐ.}

\orphannote{วิ 1. 119; วิ 2. 2 ปิฏฺฐาทีสุ.}

\orphannote{ม 1. 122; ม 3. 202; อํ 1. 50; ขุ 8. 190 ปิฏฺฐาทีสุปิ. 1. ขาราปฏิจฺฉกมฺปิ (ก) อุปริ 316 ปิฏฺเฐปิ.}

\orphannote{อภิ 1. 18 ปิฏฺเฐ.}

\orphannote{ม 3. 227 ปิฏฺเฐปิ.}

\orphannote{ขุ 8. 87 ปิฏฺเฐปิ.}

\orphannote{ม 3. 226 ปิฏฺเฐ.}

\orphannote{ม 3. 227 ปิฏฺเฐ.}

\orphannote{ม 3. 237 ปิฏฺเฐปิ.}

\orphannote{อภิ 1. 224; อภิ 2. 405; อํ 3. 375 ปิฏฺเฐสุ.}

\orphannote{ที 1. 4 ปิฏฺเฐ.}

\orphannote{ม 3. 227 ปิฏฺฐาทีสุ.}

\orphannote{อํ 1. 350 ปิฏฺเฐ. ++ ขุ 8. 272 ปิฏฺเฐปิ.}

\orphannote{สํ 3. 194 ปิฏฺเฐ.}

\orphannote{อภิ 1. 225 ปิฏฺเฐ.}

\orphannote{ที 3. 207; อํ 2. 247; ขุ 8. 64; ขุ 9. 390 ปิฏฺเฐสุปิ. 1. นาภิหสติ (สี, สฺยา)}

\orphannote{อํ 1. 280; อภิ 4. 407 ปิฏฺเฐสุ.}

\orphannote{ม 1. 158; สํ 2. 261 ปิฏฺฐาทีสุ.}

\orphannote{สํ 2. 23 ปิฏฺเฐ.}

\orphannote{ม 2. 168 ปิฏฺเฐ.}

\orphannote{ม 2\textbf{.} 168 ปิฏฺเฐ\textbf{.}}

\orphannote{ที 3. 207; อํ 2. 247; ขุ 8. 64; ขุ 9. 390 ปิฏฺเฐสุปิ.}

\orphannote{ขุ 8\textbf{.} 50\textbf{,} 62 ปิฏฺฐาทีสุ\textbf{.}}

\orphannote{ม 2. 99; อํ 3. 221; ขุ 8. 142 ปิฏฺเฐสุปิ. 1. อพฺพูหนํ (สี, สฺยา), อพฺภุหนํ (ก) ขุ 8. 77ปิฏฺเฐปิ.}

\orphannote{อุปริ 356; ขุ 8. 173 ปิฏฺเฐสุปิ.}

\orphannote{สํ 2. 349; ขุ 8. 251; เหฏฺฐา 11 ปิฏฺเฐสุปิ.}

\orphannote{ม 1. 25 ปิฏฺเฐ.}

\orphannote{วิ 3. 344 ปิฏฺเฐปิ.}

\orphannote{อํ 1. 263 ปิฏฺเฐ.}

\orphannote{ขุ 1. 270 ปิฏฺเฐ.}

\orphannote{เหฏฺฐา 167 ปิฏฺเฐปิ.}

\orphannote{อํ 2. 463 ปิฏฺเฐ.}

\orphannote{ม 1. 122; ม 3. 202; อํ 1. 50; ขุ 8. 236 ปิฏฺฐาทีสุปิ.}

\orphannote{ม 2. 311 ปิฏฺเฐ.}

\orphannote{สํ 2. 401; ขุ 8. 186 ปิฏฺเฐสุ.}

\orphannote{เหฏฺฐา 275; ขุ 8. 173 ปิฏฺเฐสุปิ.}

\orphannote{ขุ 8. 48 ปิฏฺเฐปิ.}

\orphannote{สํ 2. 349; ขุ 8. 251; เหฏฺฐา 11. 282 ปิฏฺเฐสุปิ.}

\orphannote{สํ 1. 70 ปิฏฺเฐ.}

\orphannote{ที 1. 4 ปิฏฺเฐ.}

\orphannote{สํ 3. 128 ปิฏฺเฐ.}

\orphannote{สํ 3. 128 ปิฏฺเฐ.}

\orphannote{ม 2. 146; สํ 1. 267; อํ 1. 52; ขุ 8. 297 ปิฏฺเฐสุปิ.}

\orphannote{อนิเกตสารี (สฺยา)}

\orphannote{ม 1. 12; อํ 2. 341 ปิฏฺเฐสุปิ.}

\orphannote{ที 3. 188; อํ 1. 336 ปิฏฺเฐสุปิ.}

